%% Generated by Sphinx.
\def\sphinxdocclass{report}
\documentclass[letterpaper,10pt,english]{sphinxmanual}
\ifdefined\pdfpxdimen
   \let\sphinxpxdimen\pdfpxdimen\else\newdimen\sphinxpxdimen
\fi \sphinxpxdimen=.75bp\relax

\PassOptionsToPackage{warn}{textcomp}
\usepackage[utf8]{inputenc}
\ifdefined\DeclareUnicodeCharacter
% support both utf8 and utf8x syntaxes
  \ifdefined\DeclareUnicodeCharacterAsOptional
    \def\sphinxDUC#1{\DeclareUnicodeCharacter{"#1}}
  \else
    \let\sphinxDUC\DeclareUnicodeCharacter
  \fi
  \sphinxDUC{00A0}{\nobreakspace}
  \sphinxDUC{2500}{\sphinxunichar{2500}}
  \sphinxDUC{2502}{\sphinxunichar{2502}}
  \sphinxDUC{2514}{\sphinxunichar{2514}}
  \sphinxDUC{251C}{\sphinxunichar{251C}}
  \sphinxDUC{2572}{\textbackslash}
\fi
\usepackage{cmap}
\usepackage[T1]{fontenc}
\usepackage{amsmath,amssymb,amstext}
\usepackage{babel}



\usepackage{times}
\expandafter\ifx\csname T@LGR\endcsname\relax
\else
% LGR was declared as font encoding
  \substitutefont{LGR}{\rmdefault}{cmr}
  \substitutefont{LGR}{\sfdefault}{cmss}
  \substitutefont{LGR}{\ttdefault}{cmtt}
\fi
\expandafter\ifx\csname T@X2\endcsname\relax
  \expandafter\ifx\csname T@T2A\endcsname\relax
  \else
  % T2A was declared as font encoding
    \substitutefont{T2A}{\rmdefault}{cmr}
    \substitutefont{T2A}{\sfdefault}{cmss}
    \substitutefont{T2A}{\ttdefault}{cmtt}
  \fi
\else
% X2 was declared as font encoding
  \substitutefont{X2}{\rmdefault}{cmr}
  \substitutefont{X2}{\sfdefault}{cmss}
  \substitutefont{X2}{\ttdefault}{cmtt}
\fi


\usepackage[Bjarne]{fncychap}
\usepackage{sphinx}

\fvset{fontsize=\small}
\usepackage{geometry}


% Include hyperref last.
\usepackage{hyperref}
% Fix anchor placement for figures with captions.
\usepackage{hypcap}% it must be loaded after hyperref.
% Set up styles of URL: it should be placed after hyperref.
\urlstyle{same}


\usepackage{sphinxmessages}


\newcommand{\B}[1]{{\bf #1}} 
\newcommand{\R}[1]{{\rm #1}} 
\renewcommand{\thesection}{{\hspace{-1em}}}
\renewcommand{\thesubsection}{{\hspace{-1em}}}
\renewcommand{\thesubsubsection}{{\hspace{-1em}}}


\title{cppad\_py}
\date{Sep 13, 2020}
\release{}
\author{Brad Bell}
\newcommand{\sphinxlogo}{\vbox{}}
\renewcommand{\releasename}{}
\makeindex
\begin{document}

\pagestyle{empty}
\sphinxmaketitle
\pagestyle{plain}
\sphinxtableofcontents
\pagestyle{normal}
\phantomsection\label{\detokenize{index::doc}}



\chapter{Index}
\label{\detokenize{xsrst/xsrst_automatic:index}}\begin{itemize}
\item {} 
\DUrole{xref,std,std-ref}{Link To Index}

\end{itemize}


\chapter{Table of Contents}
\label{\detokenize{xsrst/xsrst_automatic:table-of-contents}}
{\hyperref[\detokenize{xsrst/cppad_py:id1}]{\sphinxcrossref{\DUrole{std,std-ref}{A C++ Object Library and Python Interface to CppAD}}}}

\begin{DUlineblock}{0em}
\item[] 
\item[]
\begin{DUlineblock}{\DUlineblockindent}
\item[] {\hyperref[\detokenize{xsrst/setup_py:id1}]{\sphinxcrossref{\DUrole{std,std-ref}{1 Configure and Build the cppad\_py Python Module}}}}
\item[] 
\item[]    {\hyperref[\detokenize{xsrst/install_error:id1}]{\sphinxcrossref{\DUrole{std,std-ref}{1.1 Error Messages During Installation}}}}
\item[]    {\hyperref[\detokenize{xsrst/get_cppad_sh:id1}]{\sphinxcrossref{\DUrole{std,std-ref}{1.2 Get Cppad}}}}
\item[] 
\item[] {\hyperref[\detokenize{xsrst/numeric_xam:id1}]{\sphinxcrossref{\DUrole{std,std-ref}{2 Numerical Examples}}}}
\item[] 
\item[]    {\hyperref[\detokenize{xsrst/numeric_simple_inv:id1}]{\sphinxcrossref{\DUrole{std,std-ref}{2.1 An AD Compatible Matrix Inverse Routine}}}}
\item[]       {\hyperref[\detokenize{xsrst/numeric_simple_inv_xam_py:id1}]{\sphinxcrossref{\DUrole{std,std-ref}{2.1.1 Example Computing Derivatives of Matrix Inversion}}}}
\item[]    {\hyperref[\detokenize{xsrst/numeric_runge4_step:id1}]{\sphinxcrossref{\DUrole{std,std-ref}{2.2 One Fourth Order Runge\sphinxhyphen{}Kutta ODE Step}}}}
\item[]       {\hyperref[\detokenize{xsrst/numeric_runge4_step_xam_py:id1}]{\sphinxcrossref{\DUrole{std,std-ref}{2.2.1 Example Computing Derivative A Runge\sphinxhyphen{}Kutta Ode Solution}}}}
\item[]    {\hyperref[\detokenize{xsrst/numeric_rosen3_step:id1}]{\sphinxcrossref{\DUrole{std,std-ref}{2.3 One Third Order Rosenbrock ODE Step}}}}
\item[]       {\hyperref[\detokenize{xsrst/numeric_rosen3_step_xam_py:id1}]{\sphinxcrossref{\DUrole{std,std-ref}{2.3.1 Example Computing Derivative A Rosenbrock Ode Solution}}}}
\item[]    {\hyperref[\detokenize{xsrst/numeric_ode_multi_step:id1}]{\sphinxcrossref{\DUrole{std,std-ref}{2.4 Multiple Ode Steps}}}}
\item[]       {\hyperref[\detokenize{xsrst/numeric_ode_multi_step_xam_py:id1}]{\sphinxcrossref{\DUrole{std,std-ref}{2.4.1 Example Computing Derivative A Runge\sphinxhyphen{}Kutta Ode Solution}}}}
\item[]    {\hyperref[\detokenize{xsrst/numeric_optimize_fun_class:id1}]{\sphinxcrossref{\DUrole{std,std-ref}{2.5 A Helper Class That Defines Functions Needed for Optimization}}}}
\item[]       {\hyperref[\detokenize{xsrst/numeric_optimize_fun_xam_py:id1}]{\sphinxcrossref{\DUrole{std,std-ref}{2.5.1 Example Using optimize\_fun\_class with Scipy Optimization}}}}
\item[]    {\hyperref[\detokenize{xsrst/numeric_seirwd_model:id1}]{\sphinxcrossref{\DUrole{std,std-ref}{2.6 A Susceptible Exposed Infectious Recovered and Death Model}}}}
\item[]       {\hyperref[\detokenize{xsrst/numeric_seirwd_model_xam_py:id1}]{\sphinxcrossref{\DUrole{std,std-ref}{2.6.1 Example Using seris\_model}}}}
\item[]    {\hyperref[\detokenize{xsrst/numeric_covid_19_xam_py:id1}]{\sphinxcrossref{\DUrole{std,std-ref}{2.7 Example Fitting an SEIRWD Model for Covid\sphinxhyphen{}19}}}}
\item[] 
\item[] {\hyperref[\detokenize{xsrst/library:id1}]{\sphinxcrossref{\DUrole{std,std-ref}{3 The Cppad Py Libraries}}}}
\item[] 
\item[]    {\hyperref[\detokenize{xsrst/py_lib:id1}]{\sphinxcrossref{\DUrole{std,std-ref}{3.1 The Python Library}}}}
\item[] 
\item[]       {\hyperref[\detokenize{xsrst/py_fun:id1}]{\sphinxcrossref{\DUrole{std,std-ref}{3.1.1 Cppad Py AD Functions}}}}
\item[] 
\item[]          {\hyperref[\detokenize{xsrst/py_independent:id1}]{\sphinxcrossref{\DUrole{std,std-ref}{3.1.1.1 Declare Independent Variables and Start Recording}}}}
\item[]             {\hyperref[\detokenize{xsrst/fun_dynamic_xam_py:id1}]{\sphinxcrossref{\DUrole{std,std-ref}{3.1.1.1.1 Python: Using Dynamic Parameters: Example and Test}}}}
\item[]          {\hyperref[\detokenize{xsrst/py_abort_recording:id1}]{\sphinxcrossref{\DUrole{std,std-ref}{3.1.1.2 Abort Recording}}}}
\item[]             {\hyperref[\detokenize{xsrst/fun_abort_xam_py:id1}]{\sphinxcrossref{\DUrole{std,std-ref}{3.1.1.2.1 Python: Abort Recording a\_double Operations: Example and Test}}}}
\item[]          {\hyperref[\detokenize{xsrst/py_fun_ctor:id1}]{\sphinxcrossref{\DUrole{std,std-ref}{3.1.1.3 Stop Current Recording and Store Function Object}}}}
\item[]             {\hyperref[\detokenize{xsrst/a_fun_xam_py:id1}]{\sphinxcrossref{\DUrole{std,std-ref}{3.1.1.3.1 Python: Purpose of a\_fun Objects: Example and Test}}}}
\item[]          {\hyperref[\detokenize{xsrst/py_fun_property:id1}]{\sphinxcrossref{\DUrole{std,std-ref}{3.1.1.4 Properties of a Function Object}}}}
\item[]             {\hyperref[\detokenize{xsrst/fun_property_xam_py:id1}]{\sphinxcrossref{\DUrole{std,std-ref}{3.1.1.4.1 Python: d\_fun Properties: Example and Test}}}}
\item[]          {\hyperref[\detokenize{xsrst/py_fun_new_dynamic:id1}]{\sphinxcrossref{\DUrole{std,std-ref}{3.1.1.5 New Dynamic Parameters}}}}
\item[]          {\hyperref[\detokenize{xsrst/py_fun_jacobian:id1}]{\sphinxcrossref{\DUrole{std,std-ref}{3.1.1.6 Jacobian of an AD Function}}}}
\item[]             {\hyperref[\detokenize{xsrst/fun_jacobian_xam_py:id1}]{\sphinxcrossref{\DUrole{std,std-ref}{3.1.1.6.1 Python: Dense Jacobian Using AD: Example and Test}}}}
\item[]          {\hyperref[\detokenize{xsrst/py_fun_hessian:id1}]{\sphinxcrossref{\DUrole{std,std-ref}{3.1.1.7 Hessian of an AD Function}}}}
\item[]             {\hyperref[\detokenize{xsrst/fun_hessian_xam_py:id1}]{\sphinxcrossref{\DUrole{std,std-ref}{3.1.1.7.1 Python: Dense Hessian Using AD: Example and Test}}}}
\item[]          {\hyperref[\detokenize{xsrst/py_fun_forward:id1}]{\sphinxcrossref{\DUrole{std,std-ref}{3.1.1.8 Forward Mode AD}}}}
\item[]             {\hyperref[\detokenize{xsrst/fun_forward_xam_py:id1}]{\sphinxcrossref{\DUrole{std,std-ref}{3.1.1.8.1 Python: Forward Mode AD: Example and Test}}}}
\item[]          {\hyperref[\detokenize{xsrst/py_fun_reverse:id1}]{\sphinxcrossref{\DUrole{std,std-ref}{3.1.1.9 Reverse Mode AD}}}}
\item[]             {\hyperref[\detokenize{xsrst/fun_reverse_xam_py:id1}]{\sphinxcrossref{\DUrole{std,std-ref}{3.1.1.9.1 Python: Reverse Mode AD: Example and Test}}}}
\item[]          {\hyperref[\detokenize{xsrst/py_fun_optimize:id1}]{\sphinxcrossref{\DUrole{std,std-ref}{3.1.1.10 Optimize an AD Function}}}}
\item[]             {\hyperref[\detokenize{xsrst/fun_optimize_xam_py:id1}]{\sphinxcrossref{\DUrole{std,std-ref}{3.1.1.10.1 Python: Optimize an d\_fun: Example and Test}}}}
\item[]          {\hyperref[\detokenize{xsrst/py_fun_json:id1}]{\sphinxcrossref{\DUrole{std,std-ref}{3.1.1.11 Json Representation of AD Computation Graph}}}}
\item[] 
\item[]             {\hyperref[\detokenize{xsrst/fun_to_json_xam_py:id1}]{\sphinxcrossref{\DUrole{std,std-ref}{3.1.1.11.1 Python to\_json: Example and Test}}}}
\item[]             {\hyperref[\detokenize{xsrst/fun_from_json_xam_py:id1}]{\sphinxcrossref{\DUrole{std,std-ref}{3.1.1.11.2 Python from\_json: Example and Test}}}}
\item[] 
\item[]       {\hyperref[\detokenize{xsrst/py_sparse:id1}]{\sphinxcrossref{\DUrole{std,std-ref}{3.1.2 Python Sparsity Routines}}}}
\item[] 
\item[]          {\hyperref[\detokenize{xsrst/py_sparse_rc:id1}]{\sphinxcrossref{\DUrole{std,std-ref}{3.1.2.1 Sparsity Patterns}}}}
\item[]             {\hyperref[\detokenize{xsrst/sparse_rc_xam_py:id1}]{\sphinxcrossref{\DUrole{std,std-ref}{3.1.2.1.1 Python: Sparsity Patterns: Example and Test}}}}
\item[]          {\hyperref[\detokenize{xsrst/py_sparse_rcv:id1}]{\sphinxcrossref{\DUrole{std,std-ref}{3.1.2.2 Sparse Matrices}}}}
\item[]             {\hyperref[\detokenize{xsrst/sparse_rcv_xam_py:id1}]{\sphinxcrossref{\DUrole{std,std-ref}{3.1.2.2.1 Python: Sparsity Patterns: Example and Test}}}}
\item[]          {\hyperref[\detokenize{xsrst/py_jac_sparsity:id1}]{\sphinxcrossref{\DUrole{std,std-ref}{3.1.2.3 Jacobian Sparsity Patterns}}}}
\item[]             {\hyperref[\detokenize{xsrst/sparse_jac_pattern_xam_py:id1}]{\sphinxcrossref{\DUrole{std,std-ref}{3.1.2.3.1 Python: Jacobian Sparsity Patterns: Example and Test}}}}
\item[]          {\hyperref[\detokenize{xsrst/py_hes_sparsity:id1}]{\sphinxcrossref{\DUrole{std,std-ref}{3.1.2.4 Hessian Sparsity Patterns}}}}
\item[]             {\hyperref[\detokenize{xsrst/sparse_hes_pattern_xam_py:id1}]{\sphinxcrossref{\DUrole{std,std-ref}{3.1.2.4.1 Python: Hessian Sparsity Patterns: Example and Test}}}}
\item[]          {\hyperref[\detokenize{xsrst/py_sparse_jac:id1}]{\sphinxcrossref{\DUrole{std,std-ref}{3.1.2.5 Computing Sparse Jacobians}}}}
\item[]             {\hyperref[\detokenize{xsrst/sparse_jac_xam_py:id1}]{\sphinxcrossref{\DUrole{std,std-ref}{3.1.2.5.1 Python: Computing Sparse Jacobians: Example and Test}}}}
\item[]          {\hyperref[\detokenize{xsrst/py_sparse_hes:id1}]{\sphinxcrossref{\DUrole{std,std-ref}{3.1.2.6 Computing Sparse Hessians}}}}
\item[]             {\hyperref[\detokenize{xsrst/sparse_hes_xam_py:id1}]{\sphinxcrossref{\DUrole{std,std-ref}{3.1.2.6.1 Python: Hessian Sparsity Patterns: Example and Test}}}}
\item[] 
\item[]       {\hyperref[\detokenize{xsrst/py_utility:id1}]{\sphinxcrossref{\DUrole{std,std-ref}{3.1.3 Python Utilities}}}}
\item[] 
\item[]          {\hyperref[\detokenize{xsrst/numpy2vec:id1}]{\sphinxcrossref{\DUrole{std,std-ref}{3.1.3.1 Convert a Numpy Array to a cppad\_py Vector}}}}
\item[]          {\hyperref[\detokenize{xsrst/vec2numpy:id1}]{\sphinxcrossref{\DUrole{std,std-ref}{3.1.3.2 Convert a cppad\_py Vector to a Numpy Array}}}}
\item[] 
\item[]       {\hyperref[\detokenize{xsrst/more_py:id1}]{\sphinxcrossref{\DUrole{std,std-ref}{3.1.4 Steps To Add More Python Functions}}}}
\item[] 
\item[]    {\hyperref[\detokenize{xsrst/cpp_lib:id1}]{\sphinxcrossref{\DUrole{std,std-ref}{3.2 The C++ Library}}}}
\item[] 
\item[]       {\hyperref[\detokenize{xsrst/a_double:id1}]{\sphinxcrossref{\DUrole{std,std-ref}{3.2.1 Cppad Py AD Scalars}}}}
\item[] 
\item[]          {\hyperref[\detokenize{xsrst/a_double_ctor:id1}]{\sphinxcrossref{\DUrole{std,std-ref}{3.2.1.1 The a\_double Constructor}}}}
\item[]          {\hyperref[\detokenize{xsrst/a_double_unary_op:id1}]{\sphinxcrossref{\DUrole{std,std-ref}{3.2.1.2 a\_double Unary Plus and Minus}}}}
\item[] 
\item[]             {\hyperref[\detokenize{xsrst/a_double_unary_op_xam_cpp:id1}]{\sphinxcrossref{\DUrole{std,std-ref}{3.2.1.2.1 C++: a\_double Unary Plus and Minus: Example and Test}}}}
\item[]             {\hyperref[\detokenize{xsrst/a_double_unary_op_xam_py:id1}]{\sphinxcrossref{\DUrole{std,std-ref}{3.2.1.2.2 Python: a\_double Unary Plus and Minus: Example and Test}}}}
\item[] 
\item[]          {\hyperref[\detokenize{xsrst/a_double_property:id1}]{\sphinxcrossref{\DUrole{std,std-ref}{3.2.1.3 Properties of an a\_double Object}}}}
\item[] 
\item[]             {\hyperref[\detokenize{xsrst/a_double_property_xam_cpp:id1}]{\sphinxcrossref{\DUrole{std,std-ref}{3.2.1.3.1 C++: a\_double Properties: Example and Test}}}}
\item[]             {\hyperref[\detokenize{xsrst/a_double_property_xam_py:id1}]{\sphinxcrossref{\DUrole{std,std-ref}{3.2.1.3.2 Python: a\_double Properties: Example and Test}}}}
\item[] 
\item[]          {\hyperref[\detokenize{xsrst/a_double_binary:id1}]{\sphinxcrossref{\DUrole{std,std-ref}{3.2.1.4 a\_double Binary Operators with an AD Result}}}}
\item[] 
\item[]             {\hyperref[\detokenize{xsrst/a_double_binary_xam_cpp:id1}]{\sphinxcrossref{\DUrole{std,std-ref}{3.2.1.4.1 C++: a\_double Binary Operators With AD Result: Example and Test}}}}
\item[]             {\hyperref[\detokenize{xsrst/a_double_binary_xam_py:id1}]{\sphinxcrossref{\DUrole{std,std-ref}{3.2.1.4.2 Python: a\_double Binary Operators With AD Result: Example and Test}}}}
\item[] 
\item[]          {\hyperref[\detokenize{xsrst/a_double_compare:id1}]{\sphinxcrossref{\DUrole{std,std-ref}{3.2.1.5 a\_double Comparison Operators}}}}
\item[] 
\item[]             {\hyperref[\detokenize{xsrst/a_double_compare_xam_cpp:id1}]{\sphinxcrossref{\DUrole{std,std-ref}{3.2.1.5.1 C++: a\_double Comparison Operators: Example and Test}}}}
\item[]             {\hyperref[\detokenize{xsrst/a_double_compare_xam_py:id1}]{\sphinxcrossref{\DUrole{std,std-ref}{3.2.1.5.2 Python: a\_double Comparison Operators: Example and Test}}}}
\item[] 
\item[]          {\hyperref[\detokenize{xsrst/a_double_assign:id1}]{\sphinxcrossref{\DUrole{std,std-ref}{3.2.1.6 a\_double Assignment Operators}}}}
\item[] 
\item[]             {\hyperref[\detokenize{xsrst/a_double_assign_xam_cpp:id1}]{\sphinxcrossref{\DUrole{std,std-ref}{3.2.1.6.1 C++: a\_double Assignment Operators: Example and Test}}}}
\item[]             {\hyperref[\detokenize{xsrst/a_double_assign_xam_py:id1}]{\sphinxcrossref{\DUrole{std,std-ref}{3.2.1.6.2 Python: a\_double Assignment Operators: Example and Test}}}}
\item[] 
\item[]          {\hyperref[\detokenize{xsrst/a_double_unary_fun:id1}]{\sphinxcrossref{\DUrole{std,std-ref}{3.2.1.7 Unary Functions with AD Result}}}}
\item[] 
\item[]             {\hyperref[\detokenize{xsrst/a_double_unary_fun_xam_cpp:id1}]{\sphinxcrossref{\DUrole{std,std-ref}{3.2.1.7.1 C++: a\_double Unary Functions with AD Result: Example and Test}}}}
\item[]             {\hyperref[\detokenize{xsrst/a_double_unary_fun_xam_py:id1}]{\sphinxcrossref{\DUrole{std,std-ref}{3.2.1.7.2 Python: a\_double Unary Functions with AD Result: Example and Test}}}}
\item[] 
\item[]          {\hyperref[\detokenize{xsrst/a_double_cond_assign:id1}]{\sphinxcrossref{\DUrole{std,std-ref}{3.2.1.8 AD Conditional Assignment}}}}
\item[] 
\item[]             {\hyperref[\detokenize{xsrst/a_double_cond_assign_xam_cpp:id1}]{\sphinxcrossref{\DUrole{std,std-ref}{3.2.1.8.1 C++: a\_double Conditional Assignment: Example and Test}}}}
\item[]             {\hyperref[\detokenize{xsrst/a_double_cond_assign_xam_py:id1}]{\sphinxcrossref{\DUrole{std,std-ref}{3.2.1.8.2 Python: a\_double Conditional Assignment: Example and Test}}}}
\item[] 
\item[]       {\hyperref[\detokenize{xsrst/vector:id1}]{\sphinxcrossref{\DUrole{std,std-ref}{3.2.2 Cppad Py Vectors}}}}
\item[] 
\item[]          {\hyperref[\detokenize{xsrst/vector_ctor:id1}]{\sphinxcrossref{\DUrole{std,std-ref}{3.2.2.1 Cppad Py Vector Constructors}}}}
\item[]          {\hyperref[\detokenize{xsrst/vector_size:id1}]{\sphinxcrossref{\DUrole{std,std-ref}{3.2.2.2 Size of a Vector}}}}
\item[] 
\item[]             {\hyperref[\detokenize{xsrst/vector_size_xam_cpp:id1}]{\sphinxcrossref{\DUrole{std,std-ref}{3.2.2.2.1 C++: Size of Vectors: Example and Test}}}}
\item[]             {\hyperref[\detokenize{xsrst/vector_size_xam_py:id1}]{\sphinxcrossref{\DUrole{std,std-ref}{3.2.2.2.2 Python: Size of Vectors: Example and Test}}}}
\item[] 
\item[]          {\hyperref[\detokenize{xsrst/vector_set_get:id1}]{\sphinxcrossref{\DUrole{std,std-ref}{3.2.2.3 Setting and Getting Vector Elements}}}}
\item[] 
\item[]             {\hyperref[\detokenize{xsrst/vector_set_get_xam_cpp:id1}]{\sphinxcrossref{\DUrole{std,std-ref}{3.2.2.3.1 C++: Setting and Getting Vector Elements: Example and Test}}}}
\item[]             {\hyperref[\detokenize{xsrst/vector_set_get_xam_py:id1}]{\sphinxcrossref{\DUrole{std,std-ref}{3.2.2.3.2 Python: Setting and Getting Vector Elements: Example and Test}}}}
\item[] 
\item[]       {\hyperref[\detokenize{xsrst/cpp_fun:id1}]{\sphinxcrossref{\DUrole{std,std-ref}{3.2.3 Cppad Py AD Functions}}}}
\item[] 
\item[]          {\hyperref[\detokenize{xsrst/cpp_independent:id1}]{\sphinxcrossref{\DUrole{std,std-ref}{3.2.3.1 Declare Independent Variables and Start Recording}}}}
\item[]             {\hyperref[\detokenize{xsrst/fun_dynamic_xam_cpp:id1}]{\sphinxcrossref{\DUrole{std,std-ref}{3.2.3.1.1 C++: Using Dynamic Parameters: Example and Test}}}}
\item[]          {\hyperref[\detokenize{xsrst/cpp_abort_recording:id1}]{\sphinxcrossref{\DUrole{std,std-ref}{3.2.3.2 Abort Recording}}}}
\item[]             {\hyperref[\detokenize{xsrst/fun_abort_xam_cpp:id1}]{\sphinxcrossref{\DUrole{std,std-ref}{3.2.3.2.1 C++: Abort Recording a\_double Operations: Example and Test}}}}
\item[]          {\hyperref[\detokenize{xsrst/cpp_fun_ctor:id1}]{\sphinxcrossref{\DUrole{std,std-ref}{3.2.3.3 Stop Current Recording and Store Function Object}}}}
\item[]          {\hyperref[\detokenize{xsrst/cpp_fun_property:id1}]{\sphinxcrossref{\DUrole{std,std-ref}{3.2.3.4 Properties of a Function Object}}}}
\item[]             {\hyperref[\detokenize{xsrst/fun_property_xam_cpp:id1}]{\sphinxcrossref{\DUrole{std,std-ref}{3.2.3.4.1 C++: function Properties: Example and Test}}}}
\item[]          {\hyperref[\detokenize{xsrst/cpp_fun_new_dynamic:id1}]{\sphinxcrossref{\DUrole{std,std-ref}{3.2.3.5 Change The Dynamic Parameters}}}}
\item[]          {\hyperref[\detokenize{xsrst/cpp_fun_jacobian:id1}]{\sphinxcrossref{\DUrole{std,std-ref}{3.2.3.6 Jacobian of an AD Function}}}}
\item[]             {\hyperref[\detokenize{xsrst/fun_jacobian_xam_cpp:id1}]{\sphinxcrossref{\DUrole{std,std-ref}{3.2.3.6.1 C++: Dense Jacobian Using AD: Example and Test}}}}
\item[]          {\hyperref[\detokenize{xsrst/cpp_fun_hessian:id1}]{\sphinxcrossref{\DUrole{std,std-ref}{3.2.3.7 Hessian of an AD Function}}}}
\item[]             {\hyperref[\detokenize{xsrst/fun_hessian_xam_cpp:id1}]{\sphinxcrossref{\DUrole{std,std-ref}{3.2.3.7.1 C++: Dense Hessian Using AD: Example and Test}}}}
\item[]          {\hyperref[\detokenize{xsrst/cpp_fun_forward:id1}]{\sphinxcrossref{\DUrole{std,std-ref}{3.2.3.8 Forward Mode AD}}}}
\item[]             {\hyperref[\detokenize{xsrst/fun_forward_xam_cpp:id1}]{\sphinxcrossref{\DUrole{std,std-ref}{3.2.3.8.1 C++: Forward Mode AD: Example and Test}}}}
\item[]          {\hyperref[\detokenize{xsrst/cpp_fun_reverse:id1}]{\sphinxcrossref{\DUrole{std,std-ref}{3.2.3.9 Reverse Mode AD}}}}
\item[]             {\hyperref[\detokenize{xsrst/fun_reverse_xam_cpp:id1}]{\sphinxcrossref{\DUrole{std,std-ref}{3.2.3.9.1 C++: Reverse Mode AD: Example and Test}}}}
\item[]          {\hyperref[\detokenize{xsrst/cpp_fun_optimize:id1}]{\sphinxcrossref{\DUrole{std,std-ref}{3.2.3.10 Optimize an AD Function}}}}
\item[]             {\hyperref[\detokenize{xsrst/fun_optimize_xam_cpp:id1}]{\sphinxcrossref{\DUrole{std,std-ref}{3.2.3.10.1 C++: Optimize an d\_fun: Example and Test}}}}
\item[]          {\hyperref[\detokenize{xsrst/cpp_fun_json:id1}]{\sphinxcrossref{\DUrole{std,std-ref}{3.2.3.11 Json Representation of AD Computational Graph}}}}
\item[] 
\item[]             {\hyperref[\detokenize{xsrst/fun_to_json_xam_cpp:id1}]{\sphinxcrossref{\DUrole{std,std-ref}{3.2.3.11.1 C++: to\_json: Example and Test}}}}
\item[]             {\hyperref[\detokenize{xsrst/fun_from_json_xam_cpp:id1}]{\sphinxcrossref{\DUrole{std,std-ref}{3.2.3.11.2 C++: from\_json: Example and Test}}}}
\item[] 
\item[]       {\hyperref[\detokenize{xsrst/sparse:id1}]{\sphinxcrossref{\DUrole{std,std-ref}{3.2.4 Cppad Py Sparse Calculation}}}}
\item[] 
\item[]          {\hyperref[\detokenize{xsrst/cpp_sparse_rc:id1}]{\sphinxcrossref{\DUrole{std,std-ref}{3.2.4.1 Sparsity Patterns}}}}
\item[]             {\hyperref[\detokenize{xsrst/sparse_rc_xam_cpp:id1}]{\sphinxcrossref{\DUrole{std,std-ref}{3.2.4.1.1 C++: Sparsity Patterns: Example and Test}}}}
\item[]          {\hyperref[\detokenize{xsrst/cpp_sparse_rcv:id1}]{\sphinxcrossref{\DUrole{std,std-ref}{3.2.4.2 Sparse Matrices}}}}
\item[]             {\hyperref[\detokenize{xsrst/sparse_rcv_xam_cpp:id1}]{\sphinxcrossref{\DUrole{std,std-ref}{3.2.4.2.1 C++: Sparsity Patterns: Example and Test}}}}
\item[]          {\hyperref[\detokenize{xsrst/cpp_jac_sparsity:id1}]{\sphinxcrossref{\DUrole{std,std-ref}{3.2.4.3 Jacobian Sparsity Patterns}}}}
\item[]             {\hyperref[\detokenize{xsrst/sparse_jac_pattern_xam_cpp:id1}]{\sphinxcrossref{\DUrole{std,std-ref}{3.2.4.3.1 C++: Jacobian Sparsity Patterns: Example and Test}}}}
\item[]          {\hyperref[\detokenize{xsrst/cpp_sparsity:id1}]{\sphinxcrossref{\DUrole{std,std-ref}{3.2.4.4 Hessian Sparsity Patterns}}}}
\item[]             {\hyperref[\detokenize{xsrst/sparse_hes_pattern_xam_cpp:id1}]{\sphinxcrossref{\DUrole{std,std-ref}{3.2.4.4.1 C++: Hessian Sparsity Patterns: Example and Test}}}}
\item[]          {\hyperref[\detokenize{xsrst/cpp_sparse_jac:id1}]{\sphinxcrossref{\DUrole{std,std-ref}{3.2.4.5 Computing Sparse Jacobians}}}}
\item[]             {\hyperref[\detokenize{xsrst/sparse_jac_xam_cpp:id1}]{\sphinxcrossref{\DUrole{std,std-ref}{3.2.4.5.1 C++: Computing Sparse Jacobians: Example and Test}}}}
\item[]          {\hyperref[\detokenize{xsrst/cpp_sparse_hes:id1}]{\sphinxcrossref{\DUrole{std,std-ref}{3.2.4.6 Computing Sparse Hessians}}}}
\item[]             {\hyperref[\detokenize{xsrst/sparse_hes_xam_cpp:id1}]{\sphinxcrossref{\DUrole{std,std-ref}{3.2.4.6.1 C++: Hessian Sparsity Patterns: Example and Test}}}}
\item[] 
\item[]       {\hyperref[\detokenize{xsrst/cpp_utility:id1}]{\sphinxcrossref{\DUrole{std,std-ref}{3.2.5 C++ Utilities}}}}
\item[] 
\item[]          {\hyperref[\detokenize{xsrst/vec2cppad_double:id1}]{\sphinxcrossref{\DUrole{std,std-ref}{3.2.5.1 Convert an a\_double Vector to a CppAD::AD\textless{}double\textgreater{} Vector}}}}
\item[]          {\hyperref[\detokenize{xsrst/vec2a_double:id1}]{\sphinxcrossref{\DUrole{std,std-ref}{3.2.5.2 Convert a CppAD::AD\textless{}double\textgreater{} Vector to an a\_double Vector}}}}
\item[]          {\hyperref[\detokenize{xsrst/error_message:id1}]{\sphinxcrossref{\DUrole{std,std-ref}{3.2.5.3 Exception Handling}}}}
\item[] 
\item[]             {\hyperref[\detokenize{xsrst/error_message_xam_cpp:id1}]{\sphinxcrossref{\DUrole{std,std-ref}{3.2.5.3.1 C++: Cppad Py Exception Handling: Example and Test}}}}
\item[]             {\hyperref[\detokenize{xsrst/error_message_xam_py:id1}]{\sphinxcrossref{\DUrole{std,std-ref}{3.2.5.3.2 Python: Cppad Py Exception Handling: Example and Test}}}}
\item[] 
\item[]       {\hyperref[\detokenize{xsrst/more_cpp:id1}]{\sphinxcrossref{\DUrole{std,std-ref}{3.2.6 Steps To Add More C++ Functions}}}}
\item[] 
\item[] {\hyperref[\detokenize{xsrst/xsrst_py:id1}]{\sphinxcrossref{\DUrole{std,std-ref}{4 Extract Sphinx RST}}}}
\item[] 
\item[]    {\hyperref[\detokenize{xsrst/begin_cmd:id1}]{\sphinxcrossref{\DUrole{std,std-ref}{4.1 Begin and End Commands}}}}
\item[]    {\hyperref[\detokenize{xsrst/child_cmd:id1}]{\sphinxcrossref{\DUrole{std,std-ref}{4.2 Children Commands}}}}
\item[]       {\hyperref[\detokenize{xsrst/no_parent_exam:id1}]{\sphinxcrossref{\DUrole{std,std-ref}{4.2.1 No Parent Example}}}}
\item[]          {\hyperref[\detokenize{xsrst/no_parent_res:id1}]{\sphinxcrossref{\DUrole{std,std-ref}{4.2.1.1 No Parent Result}}}}
\item[] 
\item[]             {\hyperref[\detokenize{xsrst/children_exam:id1}]{\sphinxcrossref{\DUrole{std,std-ref}{4.2.1.1.1 Indent and Children Example}}}}
\item[]             {\hyperref[\detokenize{xsrst/children_res:id1}]{\sphinxcrossref{\DUrole{std,std-ref}{4.2.1.1.2 Indent and Children Result}}}}
\item[]                {\hyperref[\detokenize{xsrst/indent_exam:id1}]{\sphinxcrossref{\DUrole{std,std-ref}{4.2.1.1.2.1 Indent Example}}}}
\item[]                   {\hyperref[\detokenize{xsrst/indent_res:id1}]{\sphinxcrossref{\DUrole{std,std-ref}{4.2.1.1.2.1.1 Indent Result}}}}
\item[] 
\item[]    {\hyperref[\detokenize{xsrst/spell_cmd:id1}]{\sphinxcrossref{\DUrole{std,std-ref}{4.3 Spell Command}}}}
\item[]       {\hyperref[\detokenize{xsrst/spell_exam:id1}]{\sphinxcrossref{\DUrole{std,std-ref}{4.3.1 Spell Example}}}}
\item[]          {\hyperref[\detokenize{xsrst/spell_res:id1}]{\sphinxcrossref{\DUrole{std,std-ref}{4.3.1.1 Spell Result}}}}
\item[]    {\hyperref[\detokenize{xsrst/suspend_cmd:id1}]{\sphinxcrossref{\DUrole{std,std-ref}{4.4 Suspend and Resume Commands}}}}
\item[]       {\hyperref[\detokenize{xsrst/suspend_exam:id1}]{\sphinxcrossref{\DUrole{std,std-ref}{4.4.1 Suspend Example}}}}
\item[]          {\hyperref[\detokenize{xsrst/suspend_res:id1}]{\sphinxcrossref{\DUrole{std,std-ref}{4.4.1.1 Suspend Result}}}}
\item[]    {\hyperref[\detokenize{xsrst/code_cmd:id1}]{\sphinxcrossref{\DUrole{std,std-ref}{4.5 Code Command}}}}
\item[]       {\hyperref[\detokenize{xsrst/code_exam:id1}]{\sphinxcrossref{\DUrole{std,std-ref}{4.5.1 Code Example}}}}
\item[]          {\hyperref[\detokenize{xsrst/code_res:id1}]{\sphinxcrossref{\DUrole{std,std-ref}{4.5.1.1 Code Result}}}}
\item[]    {\hyperref[\detokenize{xsrst/file_cmd:id1}]{\sphinxcrossref{\DUrole{std,std-ref}{4.6 File Command}}}}
\item[]       {\hyperref[\detokenize{xsrst/file_exam:id1}]{\sphinxcrossref{\DUrole{std,std-ref}{4.6.1 File Example}}}}
\item[]          {\hyperref[\detokenize{xsrst/file_res:id1}]{\sphinxcrossref{\DUrole{std,std-ref}{4.6.1.1 File Result}}}}
\item[]    {\hyperref[\detokenize{xsrst/comment_ch_cmd:id1}]{\sphinxcrossref{\DUrole{std,std-ref}{4.7 Comment Character Command}}}}
\item[]       {\hyperref[\detokenize{xsrst/comment_ch_exam:id1}]{\sphinxcrossref{\DUrole{std,std-ref}{4.7.1 Comment Character Example}}}}
\item[]          {\hyperref[\detokenize{xsrst/comment_ch_res:id1}]{\sphinxcrossref{\DUrole{std,std-ref}{4.7.1.1 Comment Character Result}}}}
\item[]    {\hyperref[\detokenize{xsrst/heading_exam:id1}]{\sphinxcrossref{\DUrole{std,std-ref}{4.8 Heading Example}}}}
\item[]       {\hyperref[\detokenize{xsrst/heading_res:id1}]{\sphinxcrossref{\DUrole{std,std-ref}{4.8.1 Heading Result}}}}
\item[]    {\hyperref[\detokenize{xsrst/configure:id1}]{\sphinxcrossref{\DUrole{std,std-ref}{4.9 Example xsrst Configuration Files}}}}
\item[] 
\item[]       {\hyperref[\detokenize{xsrst/conf_py:id1}]{\sphinxcrossref{\DUrole{std,std-ref}{4.9.1 Example conf.py}}}}
\item[]       {\hyperref[\detokenize{xsrst/index_rst:id1}]{\sphinxcrossref{\DUrole{std,std-ref}{4.9.2 Example index.rst}}}}
\item[]       {\hyperref[\detokenize{xsrst/spelling:id1}]{\sphinxcrossref{\DUrole{std,std-ref}{4.9.3 Example Spelling File}}}}
\item[]       {\hyperref[\detokenize{xsrst/keyword:id1}]{\sphinxcrossref{\DUrole{std,std-ref}{4.9.4 Example Keyword File}}}}
\item[] 
\item[] {\hyperref[\detokenize{xsrst/whats_new_2020:id1}]{\sphinxcrossref{\DUrole{std,std-ref}{5 CppAD Py Changes During 2020}}}}
\item[]    {\hyperref[\detokenize{xsrst/whats_new_2018:id1}]{\sphinxcrossref{\DUrole{std,std-ref}{5.1 Cppad Py Changes During 2018}}}}
\item[] 
\end{DUlineblock}
\end{DUlineblock}


\chapter{cppad\_py}
\label{\detokenize{xsrst/cppad_py:cppad-py}}\label{\detokenize{xsrst/cppad_py::doc}}
\index{cppad\_py@\spxentry{cppad\_py}}\index{c++@\spxentry{c++}}\index{object@\spxentry{object}}\index{library@\spxentry{library}}\index{python@\spxentry{python}}\index{interface@\spxentry{interface}}\index{to@\spxentry{to}}\index{cppad@\spxentry{cppad}}\ignorespaces 

\section{A C++ Object Library and Python Interface to CppAD}
\label{\detokenize{xsrst/cppad_py:a-c-object-library-and-python-interface-to-cppad}}\label{\detokenize{xsrst/cppad_py:id1}}\begin{itemize}
\item {} 
{\hyperref[\detokenize{xsrst/cppad_py:cppad-py-version-2020-9-13}]{\sphinxcrossref{\DUrole{std,std-ref}{Version 2020.9.13}}}}

\item {} 
{\hyperref[\detokenize{xsrst/cppad_py:cppad-py-git-repository}]{\sphinxcrossref{\DUrole{std,std-ref}{Git Repository}}}}

\item {} 
{\hyperref[\detokenize{xsrst/cppad_py:cppad-py-purpose}]{\sphinxcrossref{\DUrole{std,std-ref}{Purpose}}}}

\item {} 
{\hyperref[\detokenize{xsrst/cppad_py:cppad-py-getting-started}]{\sphinxcrossref{\DUrole{std,std-ref}{Getting Started}}}}

\item {} 
{\hyperref[\detokenize{xsrst/cppad_py:cppad-py-numerical-examples}]{\sphinxcrossref{\DUrole{std,std-ref}{Numerical Examples}}}}

\item {} 
{\hyperref[\detokenize{xsrst/cppad_py:cppad-py-c-function-speed}]{\sphinxcrossref{\DUrole{std,std-ref}{C++ Function Speed}}}}

\item {} 
{\hyperref[\detokenize{xsrst/cppad_py:cppad-py-license}]{\sphinxcrossref{\DUrole{std,std-ref}{License}}}}

\item {} 
{\hyperref[\detokenize{xsrst/cppad_py:cppad-py-children}]{\sphinxcrossref{\DUrole{std,std-ref}{Children}}}}

\end{itemize}

\index{version@\spxentry{version}}\index{2020.9.13@\spxentry{2020.9.13}}\ignorespaces 

\subsection{Version 2020.9.13}
\label{\detokenize{xsrst/cppad_py:version-2020-9-13}}\label{\detokenize{xsrst/cppad_py:cppad-py-version-2020-9-13}}\label{\detokenize{xsrst/cppad_py:index-1}}
\index{git@\spxentry{git}}\index{repository@\spxentry{repository}}\ignorespaces 

\subsection{Git Repository}
\label{\detokenize{xsrst/cppad_py:git-repository}}\label{\detokenize{xsrst/cppad_py:cppad-py-git-repository}}\label{\detokenize{xsrst/cppad_py:index-2}}
\sphinxurl{https://github.com/bradbell/cppad\_py}

\index{purpose@\spxentry{purpose}}\ignorespaces 

\subsection{Purpose}
\label{\detokenize{xsrst/cppad_py:purpose}}\label{\detokenize{xsrst/cppad_py:cppad-py-purpose}}\label{\detokenize{xsrst/cppad_py:index-3}}\begin{enumerate}
\sphinxsetlistlabels{\arabic}{enumi}{enumii}{}{.}%
\item {} 
Provide a connection from Python
to the Algorithmic Differentiation (AD) package Cppad; see {\hyperref[\detokenize{xsrst/py_lib:id1}]{\sphinxcrossref{\DUrole{std,std-ref}{py\_lib}}}}.

\item {} 
Provide an AD object library; see {\hyperref[\detokenize{xsrst/cpp_lib:id1}]{\sphinxcrossref{\DUrole{std,std-ref}{cpp\_lib}}}}.

\item {} 
Provide a concrete example of how
\sphinxhref{https://github.com/bradbell/cppad\_swig}{cppad\_swig}
can be used to connect any scripting language to CppAD.

\end{enumerate}

\index{getting@\spxentry{getting}}\index{started@\spxentry{started}}\ignorespaces 

\subsection{Getting Started}
\label{\detokenize{xsrst/cppad_py:getting-started}}\label{\detokenize{xsrst/cppad_py:cppad-py-getting-started}}\label{\detokenize{xsrst/cppad_py:index-4}}
The following is a good place to get started using cppad\_py:
{\hyperref[\detokenize{xsrst/fun_jacobian_xam_py:id1}]{\sphinxcrossref{\DUrole{std,std-ref}{python}}}}.

\index{numerical@\spxentry{numerical}}\index{examples@\spxentry{examples}}\ignorespaces 

\subsection{Numerical Examples}
\label{\detokenize{xsrst/cppad_py:numerical-examples}}\label{\detokenize{xsrst/cppad_py:cppad-py-numerical-examples}}\label{\detokenize{xsrst/cppad_py:index-5}}
The following is a link to some numerical examples: {\hyperref[\detokenize{xsrst/numeric_xam:id1}]{\sphinxcrossref{\DUrole{std,std-ref}{numeric\_xam}}}}.

\index{c++@\spxentry{c++}}\index{function@\spxentry{function}}\index{speed@\spxentry{speed}}\ignorespaces 

\subsection{C++ Function Speed}
\label{\detokenize{xsrst/cppad_py:c-function-speed}}\label{\detokenize{xsrst/cppad_py:cppad-py-c-function-speed}}\label{\detokenize{xsrst/cppad_py:index-6}}
One can use Cppad Py to get faster function evaluation in scripting Python,
when the sequence of floating point operations does not depend on the
independent variables.
Once an {\hyperref[\detokenize{xsrst/py_fun:id1}]{\sphinxcrossref{\DUrole{std,std-ref}{py\_fun}}}} is recorded, zero order
{\hyperref[\detokenize{xsrst/py_fun_forward:id1}]{\sphinxcrossref{\DUrole{std,std-ref}{forward\_mode}}}} can be used to
effectively evaluate the function in C++ instead of Python.

\index{license@\spxentry{license}}\ignorespaces 

\subsection{License}
\label{\detokenize{xsrst/cppad_py:license}}\label{\detokenize{xsrst/cppad_py:cppad-py-license}}\label{\detokenize{xsrst/cppad_py:index-7}}
This program is distributed under the terms of the
GNU General Public License version 3.0 or later see
\sphinxhref{http://www.gnu.org/licenses/gpl-3.0.txt}{gpl\sphinxhyphen{}3.0.txt}.

\index{children@\spxentry{children}}\ignorespaces 

\subsection{Children}
\label{\detokenize{xsrst/cppad_py:children}}\label{\detokenize{xsrst/cppad_py:cppad-py-children}}\label{\detokenize{xsrst/cppad_py:index-8}}

\subsubsection{setup\_py}
\label{\detokenize{xsrst/setup_py:setup-py}}\label{\detokenize{xsrst/setup_py::doc}}
\index{setup\_py@\spxentry{setup\_py}}\index{configure@\spxentry{configure}}\index{build@\spxentry{build}}\index{the@\spxentry{the}}\index{cppad\_py@\spxentry{cppad\_py}}\index{python@\spxentry{python}}\index{module@\spxentry{module}}\ignorespaces 

\paragraph{1 Configure and Build the cppad\_py Python Module}
\label{\detokenize{xsrst/setup_py:configure-and-build-the-cppad-py-python-module}}\label{\detokenize{xsrst/setup_py:id1}}\begin{itemize}
\item {} 
{\hyperref[\detokenize{xsrst/setup_py:setup-py-syntax}]{\sphinxcrossref{\DUrole{std,std-ref}{Syntax}}}}

\item {} 
{\hyperref[\detokenize{xsrst/setup_py:setup-py-external-requirements}]{\sphinxcrossref{\DUrole{std,std-ref}{External Requirements}}}}

\item {} 
{\hyperref[\detokenize{xsrst/setup_py:setup-py-install-using-pip}]{\sphinxcrossref{\DUrole{std,std-ref}{Install Using Pip}}}}

\item {} 
{\hyperref[\detokenize{xsrst/setup_py:setup-py-install-errors}]{\sphinxcrossref{\DUrole{std,std-ref}{Install Errors}}}}

\item {} \begin{description}
\item[{{\hyperref[\detokenize{xsrst/setup_py:setup-py-download}]{\sphinxcrossref{\DUrole{std,std-ref}{Download}}}}}] \leavevmode\begin{itemize}
\item {} 
{\hyperref[\detokenize{xsrst/setup_py:setup-py-download-top-source-directory}]{\sphinxcrossref{\DUrole{std,std-ref}{Top Source Directory}}}}

\end{itemize}

\end{description}

\item {} 
{\hyperref[\detokenize{xsrst/setup_py:setup-py-configure}]{\sphinxcrossref{\DUrole{std,std-ref}{Configure}}}}

\item {} 
{\hyperref[\detokenize{xsrst/setup_py:setup-py-get-cppad}]{\sphinxcrossref{\DUrole{std,std-ref}{Get cppad}}}}

\item {} \begin{description}
\item[{{\hyperref[\detokenize{xsrst/setup_py:setup-py-test}]{\sphinxcrossref{\DUrole{std,std-ref}{Test}}}}}] \leavevmode\begin{itemize}
\item {} 
{\hyperref[\detokenize{xsrst/setup_py:setup-py-test-build-cppad-py}]{\sphinxcrossref{\DUrole{std,std-ref}{Build cppad\_py}}}}

\item {} 
{\hyperref[\detokenize{xsrst/setup_py:setup-py-test-python}]{\sphinxcrossref{\DUrole{std,std-ref}{python}}}}

\item {} 
{\hyperref[\detokenize{xsrst/setup_py:setup-py-test-c}]{\sphinxcrossref{\DUrole{std,std-ref}{c++}}}}

\item {} 
{\hyperref[\detokenize{xsrst/setup_py:setup-py-test-import}]{\sphinxcrossref{\DUrole{std,std-ref}{import}}}}

\end{itemize}

\end{description}

\item {} 
{\hyperref[\detokenize{xsrst/setup_py:setup-py-install}]{\sphinxcrossref{\DUrole{std,std-ref}{Install}}}}

\item {} 
{\hyperref[\detokenize{xsrst/setup_py:setup-py-python-path}]{\sphinxcrossref{\DUrole{std,std-ref}{Python Path}}}}

\end{itemize}

\index{syntax@\spxentry{syntax}}\ignorespaces 

\subparagraph{Syntax}
\label{\detokenize{xsrst/setup_py:syntax}}\label{\detokenize{xsrst/setup_py:setup-py-syntax}}\label{\detokenize{xsrst/setup_py:index-1}}
\begin{DUlineblock}{0em}
\item[] \sphinxcode{\sphinxupquote{python setup.py build\_ext}}
\end{DUlineblock}

\index{external@\spxentry{external}}\index{requirements@\spxentry{requirements}}\ignorespaces 

\subparagraph{External Requirements}
\label{\detokenize{xsrst/setup_py:external-requirements}}\label{\detokenize{xsrst/setup_py:setup-py-external-requirements}}\label{\detokenize{xsrst/setup_py:index-2}}\begin{enumerate}
\sphinxsetlistlabels{\arabic}{enumi}{enumii}{}{.}%
\item {} 
\sphinxhref{https://www.python.org/}{python} version 3.

\item {} 
\sphinxhref{https://en.wikipedia.org/wiki/Bash\_(Unix\_shell)}{bash}

\item {} 
\sphinxhref{https://en.wikipedia.org/wiki/C++}{c++}.

\item {} 
\sphinxhref{https://cmake.org}{cmake}

\item {} 
\sphinxhref{https://git-scm.com/}{git}.

\item {} 
\sphinxhref{http://www.swig.org/}{swig}:
\sphinxhref{https://github.com/bradbell/cppad\_py/issues/3}{issue 3}.

\item {} 
\sphinxhref{http://www.numpy.org/}{numpy}.

\end{enumerate}

\index{install@\spxentry{install}}\index{using@\spxentry{using}}\index{pip@\spxentry{pip}}\ignorespaces 

\subparagraph{Install Using Pip}
\label{\detokenize{xsrst/setup_py:install-using-pip}}\label{\detokenize{xsrst/setup_py:setup-py-install-using-pip}}\label{\detokenize{xsrst/setup_py:index-3}}
There is a preliminary version of cppad\_py available using \sphinxcode{\sphinxupquote{pip}} .
To install for you entire system use:

\begin{DUlineblock}{0em}
\item[]         \sphinxcode{\sphinxupquote{pip install \sphinxhyphen{}i https://test.pypi.org/simple/ cppad\_py}}
\end{DUlineblock}


\subparagraph{install\_error}
\label{\detokenize{xsrst/install_error:install-error}}\label{\detokenize{xsrst/install_error::doc}}
\index{install\_error@\spxentry{install\_error}}\index{error@\spxentry{error}}\index{messages@\spxentry{messages}}\index{during@\spxentry{during}}\index{installation@\spxentry{installation}}\ignorespaces 

\subparagraph{1.1 Error Messages During Installation}
\label{\detokenize{xsrst/install_error:error-messages-during-installation}}\label{\detokenize{xsrst/install_error:id1}}\begin{itemize}
\item {} \begin{description}
\item[{{\hyperref[\detokenize{xsrst/install_error:install-error-solved}]{\sphinxcrossref{\DUrole{std,std-ref}{Solved}}}}}] \leavevmode\begin{itemize}
\item {} 
{\hyperref[\detokenize{xsrst/install_error:install-error-solved-swig}]{\sphinxcrossref{\DUrole{std,std-ref}{swig}}}}

\item {} 
{\hyperref[\detokenize{xsrst/install_error:install-error-solved-permissions}]{\sphinxcrossref{\DUrole{std,std-ref}{Permissions}}}}

\end{itemize}

\end{description}

\item {} 
{\hyperref[\detokenize{xsrst/install_error:install-error-cppad-library-missing}]{\sphinxcrossref{\DUrole{std,std-ref}{CppAD Library Missing}}}}

\item {} 
{\hyperref[\detokenize{xsrst/install_error:install-error-fortify-source}]{\sphinxcrossref{\DUrole{std,std-ref}{Fortify Source}}}}

\item {} \begin{description}
\item[{{\hyperref[\detokenize{xsrst/install_error:install-error-unsolved}]{\sphinxcrossref{\DUrole{std,std-ref}{Unsolved}}}}}] \leavevmode\begin{itemize}
\item {} 
{\hyperref[\detokenize{xsrst/install_error:install-error-unsolved-travis}]{\sphinxcrossref{\DUrole{std,std-ref}{Travis}}}}

\end{itemize}

\end{description}

\end{itemize}

\index{solved@\spxentry{solved}}\ignorespaces 

\subparagraph{Solved}
\label{\detokenize{xsrst/install_error:solved}}\label{\detokenize{xsrst/install_error:install-error-solved}}\label{\detokenize{xsrst/install_error:index-1}}
\index{swig@\spxentry{swig}}\ignorespaces 

\subparagraph{swig}
\label{\detokenize{xsrst/install_error:swig}}\label{\detokenize{xsrst/install_error:install-error-solved-swig}}\label{\detokenize{xsrst/install_error:index-2}}
If the error message below occurs, try installing
\sphinxhref{swig}{http://www.swig.org/} on you system:

\begin{DUlineblock}{0em}
\item[]         \sphinxcode{\sphinxupquote{FileNotFoundError: {[}Errno 2{]} No such file or directory: \textquotesingle{}swig\textquotesingle{}}}
\end{DUlineblock}

The following message as also been seen:

\begin{DUlineblock}{0em}
\item[]         \sphinxcode{\sphinxupquote{Error: Unable to find \textquotesingle{}python.swg\textquotesingle{}}}
\end{DUlineblock}

In this case try installing \sphinxcode{\sphinxupquote{swig\sphinxhyphen{}python}} .

\index{permissions@\spxentry{permissions}}\ignorespaces 

\subparagraph{Permissions}
\label{\detokenize{xsrst/install_error:permissions}}\label{\detokenize{xsrst/install_error:install-error-solved-permissions}}\label{\detokenize{xsrst/install_error:index-3}}
If you try to do a system wide install, and do not have permission,
you will get the following message:

\begin{DUlineblock}{0em}
\item[]         \sphinxcode{\sphinxupquote{ERROR: Could not install packages due to}} … \sphinxcode{\sphinxupquote{Permission denied:}}
\end{DUlineblock}

Try install in your local user space. For example, if you are using pip,

\begin{DUlineblock}{0em}
\item[]         \sphinxcode{\sphinxupquote{pip install \sphinxhyphen{}i https://test.pypi.org/simple/ cppad\_py \sphinxhyphen{}\sphinxhyphen{}user \sphinxhyphen{}\sphinxhyphen{}verbose}}
\end{DUlineblock}

\index{cppad@\spxentry{cppad}}\index{library@\spxentry{library}}\index{missing@\spxentry{missing}}\ignorespaces 

\subparagraph{CppAD Library Missing}
\label{\detokenize{xsrst/install_error:cppad-library-missing}}\label{\detokenize{xsrst/install_error:install-error-cppad-library-missing}}\label{\detokenize{xsrst/install_error:index-4}}
If the error message below occurs, it is because the CppAD library
is not in your \sphinxcode{\sphinxupquote{LD\_LIBRARY\_PATH}} :

\begin{DUlineblock}{0em}
\item[]         \sphinxcode{\sphinxupquote{ImportError: libcppad\_lib.so}} … \sphinxcode{\sphinxupquote{can not open shared object file}}
\end{DUlineblock}

If you run the \sphinxcode{\sphinxupquote{setup.py}} script directory, or
if you included the \sphinxcode{\sphinxupquote{\sphinxhyphen{}\sphinxhyphen{}verbose}} option in your \sphinxcode{\sphinxupquote{pip}} command,
there should be text at the end that tells you how to modify your
\sphinxcode{\sphinxupquote{LD\_LIBRARY\_PATH}} .
If you have a Mac, you may instead need to modify \sphinxcode{\sphinxupquote{DYLD\_LIBRARY\_PATH}} .

\index{fortify@\spxentry{fortify}}\index{source@\spxentry{source}}\ignorespaces 

\subparagraph{Fortify Source}
\label{\detokenize{xsrst/install_error:fortify-source}}\label{\detokenize{xsrst/install_error:install-error-fortify-source}}\label{\detokenize{xsrst/install_error:index-5}}
If you set {\hyperref[\detokenize{xsrst/get_cppad_sh:get-cppad-sh-settings-build-type}]{\sphinxcrossref{\DUrole{std,std-ref}{build\_type}}}} to \sphinxcode{\sphinxupquote{debug}} ,
you may get the following warning during the build:

\begin{DUlineblock}{0em}
\item[]         \sphinxcode{\sphinxupquote{\#warning \_FORTIFY\_SOURCE requires compiling with optimization}}
\end{DUlineblock}

This is a problem with the python setuptools,
one can un\sphinxhyphen{}define a macro, but it does not remove a original definition.

\index{unsolved@\spxentry{unsolved}}\ignorespaces 

\subparagraph{Unsolved}
\label{\detokenize{xsrst/install_error:unsolved}}\label{\detokenize{xsrst/install_error:install-error-unsolved}}\label{\detokenize{xsrst/install_error:index-6}}
If you know how to fix and of he errors below, please open an
\sphinxhref{https://github.com/bradbell/cppad\_py/issues}{issue}
and put your solution there.

\index{travis@\spxentry{travis}}\ignorespaces 

\subparagraph{Travis}
\label{\detokenize{xsrst/install_error:travis}}\label{\detokenize{xsrst/install_error:install-error-unsolved-travis}}\label{\detokenize{xsrst/install_error:index-7}}
The following message occurs on a Travis system of unknown configuration;
i.e., it has not been reproduced outside of Travis.

\begin{DUlineblock}{0em}
\item[]         \sphinxcode{\sphinxupquote{return alpha * p * a\_exp( \sphinxhyphen{} z * z ) / numpy.sqrt(numpy.pi){[}0m}}
\item[]                 \sphinxcode{\sphinxupquote{TypeError: unsupported operand type(s) for /: \textquotesingle{}a\_double\textquotesingle{} and \textquotesingle{}float\textquotesingle{}}}
\end{DUlineblock}

Note the same code works on other systems.
The \sphinxcode{\sphinxupquote{.travis.yml}} file is included below:

\begin{DUlineblock}{0em}
\item[]         \sphinxcode{\sphinxupquote{language: python}}
\item[]         \sphinxcode{\sphinxupquote{python:}}
\item[]                 \sphinxhyphen{}  \sphinxcode{\sphinxupquote{"3.7"}}
\item[]         \sphinxcode{\sphinxupquote{install:}}
\item[]                 \sphinxhyphen{}  \sphinxcode{\sphinxupquote{sudo apt\sphinxhyphen{}get install swig}}
\item[]                 \sphinxhyphen{}  \sphinxcode{\sphinxupquote{python setup.py install}}
\item[]                 \sphinxhyphen{}  \sphinxcode{\sphinxupquote{pip install \sphinxhyphen{}i https://test.pypi.org/simple/ cppad\_py}}
\item[]         \sphinxcode{\sphinxupquote{script:}}
\item[]                 \sphinxhyphen{}  \sphinxcode{\sphinxupquote{pytest tests}}
\item[]                 \sphinxhyphen{}  \sphinxcode{\sphinxupquote{make examples}}
\end{DUlineblock}


\bigskip\hrule\bigskip


xsrst input file: \sphinxcode{\sphinxupquote{install\_error.xsrst}}


\subparagraph{get\_cppad\_sh}
\label{\detokenize{xsrst/get_cppad_sh:get-cppad-sh}}\label{\detokenize{xsrst/get_cppad_sh::doc}}
\index{get\_cppad\_sh@\spxentry{get\_cppad\_sh}}\index{get@\spxentry{get}}\index{cppad@\spxentry{cppad}}\ignorespaces 

\subparagraph{1.2 Get Cppad}
\label{\detokenize{xsrst/get_cppad_sh:get-cppad}}\label{\detokenize{xsrst/get_cppad_sh:id1}}\begin{itemize}
\item {} 
{\hyperref[\detokenize{xsrst/get_cppad_sh:get-cppad-sh-syntax}]{\sphinxcrossref{\DUrole{std,std-ref}{Syntax}}}}

\item {} 
{\hyperref[\detokenize{xsrst/get_cppad_sh:get-cppad-sh-top-source-directory}]{\sphinxcrossref{\DUrole{std,std-ref}{Top Source Directory}}}}

\item {} \begin{description}
\item[{{\hyperref[\detokenize{xsrst/get_cppad_sh:get-cppad-sh-settings}]{\sphinxcrossref{\DUrole{std,std-ref}{Settings}}}}}] \leavevmode\begin{itemize}
\item {} 
{\hyperref[\detokenize{xsrst/get_cppad_sh:get-cppad-sh-settings-cppad-prefix}]{\sphinxcrossref{\DUrole{std,std-ref}{cppad\_prefix}}}}

\item {} 
{\hyperref[\detokenize{xsrst/get_cppad_sh:get-cppad-sh-settings-extra-cxx-flags}]{\sphinxcrossref{\DUrole{std,std-ref}{extra\_cxx\_flags}}}}

\item {} 
{\hyperref[\detokenize{xsrst/get_cppad_sh:get-cppad-sh-settings-build-type}]{\sphinxcrossref{\DUrole{std,std-ref}{build\_type}}}}

\item {} 
{\hyperref[\detokenize{xsrst/get_cppad_sh:get-cppad-sh-settings-test-cppad}]{\sphinxcrossref{\DUrole{std,std-ref}{test\_cppad}}}}

\end{itemize}

\end{description}

\item {} 
{\hyperref[\detokenize{xsrst/get_cppad_sh:get-cppad-sh-caching}]{\sphinxcrossref{\DUrole{std,std-ref}{Caching}}}}

\end{itemize}

\index{syntax@\spxentry{syntax}}\ignorespaces 

\subparagraph{Syntax}
\label{\detokenize{xsrst/get_cppad_sh:syntax}}\label{\detokenize{xsrst/get_cppad_sh:get-cppad-sh-syntax}}\label{\detokenize{xsrst/get_cppad_sh:index-1}}
\sphinxcode{\sphinxupquote{bin/get\_cppad.sh}}

\index{top@\spxentry{top}}\index{source@\spxentry{source}}\index{directory@\spxentry{directory}}\ignorespaces 

\subparagraph{Top Source Directory}
\label{\detokenize{xsrst/get_cppad_sh:top-source-directory}}\label{\detokenize{xsrst/get_cppad_sh:get-cppad-sh-top-source-directory}}\label{\detokenize{xsrst/get_cppad_sh:index-2}}
This program must be run from the
{\hyperref[\detokenize{xsrst/setup_py:setup-py-download-top-source-directory}]{\sphinxcrossref{\DUrole{std,std-ref}{top\_source\_directory}}}}.
Each time it is run it removes the old \sphinxcode{\sphinxupquote{build}} directory and
starts over.

\index{settings@\spxentry{settings}}\ignorespaces 

\subparagraph{Settings}
\label{\detokenize{xsrst/get_cppad_sh:settings}}\label{\detokenize{xsrst/get_cppad_sh:get-cppad-sh-settings}}\label{\detokenize{xsrst/get_cppad_sh:index-3}}
If you change any of these settings, you must re\sphinxhyphen{}run \sphinxcode{\sphinxupquote{get\_cppad.sh}} .

\index{cppad\_prefix@\spxentry{cppad\_prefix}}\ignorespaces 

\subparagraph{cppad\_prefix}
\label{\detokenize{xsrst/get_cppad_sh:cppad-prefix}}\label{\detokenize{xsrst/get_cppad_sh:get-cppad-sh-settings-cppad-prefix}}\label{\detokenize{xsrst/get_cppad_sh:index-4}}
This prefix is used to install CppAD may be a local director; e.g.,
\sphinxcode{\sphinxupquote{build/prefix}} or an absolute path; e.g., \sphinxcode{\sphinxupquote{/usr/local}} ,
it may include the shell variable \sphinxcode{\sphinxupquote{\$HOME}} but no other variables:

\begin{sphinxVerbatim}[commandchars=\\\{\}]
\PYG{n+nv}{cppad\PYGZus{}prefix}\PYG{o}{=}\PYG{l+s+s2}{\PYGZdq{}}\PYG{n+nv}{\PYGZdl{}HOME}\PYG{l+s+s2}{/prefix/cppad}\PYG{l+s+s2}{\PYGZdq{}}
\end{sphinxVerbatim}

\index{extra\_cxx\_flags@\spxentry{extra\_cxx\_flags}}\ignorespaces 

\subparagraph{extra\_cxx\_flags}
\label{\detokenize{xsrst/get_cppad_sh:extra-cxx-flags}}\label{\detokenize{xsrst/get_cppad_sh:get-cppad-sh-settings-extra-cxx-flags}}\label{\detokenize{xsrst/get_cppad_sh:index-5}}
Extra compiler false used when compiling c++ code not including the
debugging and optimization flags.
The ones below are example flags are used by g++:

\begin{sphinxVerbatim}[commandchars=\\\{\}]
\PYG{n+nv}{extra\PYGZus{}cxx\PYGZus{}flags}\PYG{o}{=}\PYG{l+s+s1}{\PYGZsq{}\PYGZhy{}Wall \PYGZhy{}pedantic\PYGZhy{}errors \PYGZhy{}Wno\PYGZhy{}unused\PYGZhy{}result \PYGZhy{}std=c++11\PYGZsq{}}
\end{sphinxVerbatim}

\index{build\_type@\spxentry{build\_type}}\ignorespaces 

\subparagraph{build\_type}
\label{\detokenize{xsrst/get_cppad_sh:build-type}}\label{\detokenize{xsrst/get_cppad_sh:get-cppad-sh-settings-build-type}}\label{\detokenize{xsrst/get_cppad_sh:index-6}}
This must be must \sphinxcode{\sphinxupquote{debug}} or \sphinxcode{\sphinxupquote{release}} .
The debug version has more error messaging while the release
version runs faster.

\begin{sphinxVerbatim}[commandchars=\\\{\}]
\PYG{n+nv}{build\PYGZus{}type}\PYG{o}{=}\PYG{l+s+s1}{\PYGZsq{}release\PYGZsq{}}
\end{sphinxVerbatim}

If you used the \sphinxcode{\sphinxupquote{debug}} build type you may get the following warning
from the compiler (because the optimization is totally turned off):

\begin{DUlineblock}{0em}
\item[]         \sphinxcode{\sphinxupquote{warning \_FORTIFY\_SOURCE requires compiling with optimization}}
\end{DUlineblock}

\index{test\_cppad@\spxentry{test\_cppad}}\ignorespaces 

\subparagraph{test\_cppad}
\label{\detokenize{xsrst/get_cppad_sh:test-cppad}}\label{\detokenize{xsrst/get_cppad_sh:get-cppad-sh-settings-test-cppad}}\label{\detokenize{xsrst/get_cppad_sh:index-7}}
This must be must \sphinxcode{\sphinxupquote{true}} or \sphinxcode{\sphinxupquote{false}} .
Cppad has a huge test suite and this can take a significant amount of time,
but it may be useful if you have problems.

\begin{sphinxVerbatim}[commandchars=\\\{\}]
\PYG{n+nv}{test\PYGZus{}cppad}\PYG{o}{=}\PYG{l+s+s1}{\PYGZsq{}false\PYGZsq{}}
\end{sphinxVerbatim}

\index{caching@\spxentry{caching}}\ignorespaces 

\subparagraph{Caching}
\label{\detokenize{xsrst/get_cppad_sh:caching}}\label{\detokenize{xsrst/get_cppad_sh:get-cppad-sh-caching}}\label{\detokenize{xsrst/get_cppad_sh:index-8}}
This procedure cashes previous builds so that when you re\sphinxhyphen{}run
this script it does not re\sphinxhyphen{}do all the work.
If you have trouble, try deleting the directory

\begin{DUlineblock}{0em}
\item[]         \sphinxcode{\sphinxupquote{build/cppad}} \sphinxhyphen{} \sphinxstyleemphasis{yyyymmdd} . \sphinxcode{\sphinxupquote{git}}
\end{DUlineblock}

and re\sphinxhyphen{}running this script.


\bigskip\hrule\bigskip


xsrst input file: \sphinxcode{\sphinxupquote{bin/get\_cppad.sh}}

\index{install@\spxentry{install}}\index{errors@\spxentry{errors}}\ignorespaces 

\subparagraph{Install Errors}
\label{\detokenize{xsrst/setup_py:install-errors}}\label{\detokenize{xsrst/setup_py:setup-py-install-errors}}\label{\detokenize{xsrst/setup_py:index-4}}
If you get an error message during the install procedure above,
or the one below, see {\hyperref[\detokenize{xsrst/install_error:id1}]{\sphinxcrossref{\DUrole{std,std-ref}{install\_error}}}}.
This will only install the release version.
Installing a debug version is discussed below in the instructions
for downloading and building from the source code.

\index{download@\spxentry{download}}\ignorespaces 

\subparagraph{Download}
\label{\detokenize{xsrst/setup_py:download}}\label{\detokenize{xsrst/setup_py:setup-py-download}}\label{\detokenize{xsrst/setup_py:index-5}}
Use the following command to download the current version of cppad\_py:

\begin{DUlineblock}{0em}
\item[]         \sphinxcode{\sphinxupquote{git clone https://github.com/bradbell/cppad\_py.git}} \sphinxstyleemphasis{top\_srcdir}
\end{DUlineblock}

\index{top@\spxentry{top}}\index{source@\spxentry{source}}\index{directory@\spxentry{directory}}\ignorespaces 

\subparagraph{Top Source Directory}
\label{\detokenize{xsrst/setup_py:top-source-directory}}\label{\detokenize{xsrst/setup_py:setup-py-download-top-source-directory}}\label{\detokenize{xsrst/setup_py:index-6}}
The directory you choose for \sphinxstyleemphasis{top\_srcdir} is
referred to as your top source directory.

\index{configure@\spxentry{configure}}\ignorespaces 

\subparagraph{Configure}
\label{\detokenize{xsrst/setup_py:configure}}\label{\detokenize{xsrst/setup_py:setup-py-configure}}\label{\detokenize{xsrst/setup_py:index-7}}
Before running \sphinxcode{\sphinxupquote{setup.py}} or \sphinxcode{\sphinxupquote{bin/get\_cppad.sh}} ,
you should check and possibly change the
{\hyperref[\detokenize{xsrst/get_cppad_sh:get-cppad-sh-settings}]{\sphinxcrossref{\DUrole{std,std-ref}{settings}}}} in \sphinxcode{\sphinxupquote{bin/get\_cppad.sh}} .

\index{get@\spxentry{get}}\index{cppad@\spxentry{cppad}}\ignorespaces 

\subparagraph{Get cppad}
\label{\detokenize{xsrst/setup_py:get-cppad}}\label{\detokenize{xsrst/setup_py:setup-py-get-cppad}}\label{\detokenize{xsrst/setup_py:index-8}}
The next step is to get a copy of cppad using {\hyperref[\detokenize{xsrst/get_cppad_sh:id1}]{\sphinxcrossref{\DUrole{std,std-ref}{get\_cppad\_sh}}}}.

\index{test@\spxentry{test}}\ignorespaces 

\subparagraph{Test}
\label{\detokenize{xsrst/setup_py:test}}\label{\detokenize{xsrst/setup_py:setup-py-test}}\label{\detokenize{xsrst/setup_py:index-9}}
These steps are optional if you already know that cppad\_py
works on your system.

\index{build@\spxentry{build}}\index{cppad\_py@\spxentry{cppad\_py}}\ignorespaces 

\subparagraph{Build cppad\_py}
\label{\detokenize{xsrst/setup_py:build-cppad-py}}\label{\detokenize{xsrst/setup_py:setup-py-test-build-cppad-py}}\label{\detokenize{xsrst/setup_py:index-10}}
Build the Python cppad\_py module using the command:

\begin{DUlineblock}{0em}
\item[]         \sphinxcode{\sphinxupquote{python setup.py build\_ext}}
\end{DUlineblock}

\index{python@\spxentry{python}}\ignorespaces 

\subparagraph{python}
\label{\detokenize{xsrst/setup_py:setup-py-test-python}}\label{\detokenize{xsrst/setup_py:index-11}}\label{\detokenize{xsrst/setup_py:id2}}
The next step is to test the cppad\_py on your system by executing
the following commands starting in \sphinxstyleemphasis{top\_srcdir} :

\begin{DUlineblock}{0em}
\item[]         \sphinxcode{\sphinxupquote{cd example/python}}
\item[]         \sphinxcode{\sphinxupquote{python check\_all.py}}
\end{DUlineblock}

\index{c++@\spxentry{c++}}\ignorespaces 

\subparagraph{c++}
\label{\detokenize{xsrst/setup_py:setup-py-test-c}}\label{\detokenize{xsrst/setup_py:index-12}}\label{\detokenize{xsrst/setup_py:id3}}
After \sphinxcode{\sphinxupquote{setup.py}} has run,
you can also test the cppad\_py c++ interface
{\hyperref[\detokenize{xsrst/cpp_lib:id1}]{\sphinxcrossref{\DUrole{std,std-ref}{cpp\_lib}}}} on your system by executing the following commands
starting in \sphinxstyleemphasis{top\_srcdir} :

\begin{DUlineblock}{0em}
\item[]         \sphinxcode{\sphinxupquote{cd build}}
\item[]         \sphinxcode{\sphinxupquote{make check\_lib\_cplusplus}}
\end{DUlineblock}

\index{import@\spxentry{import}}\ignorespaces 

\subparagraph{import}
\label{\detokenize{xsrst/setup_py:import}}\label{\detokenize{xsrst/setup_py:setup-py-test-import}}\label{\detokenize{xsrst/setup_py:index-13}}
If you are in the \sphinxstyleemphasis{top\_srcdir} directory,
you should be able to import cppad\_py using the following commands:

\begin{DUlineblock}{0em}
\item[]         \sphinxcode{\sphinxupquote{python}}
\item[]         \sphinxcode{\sphinxupquote{import cppad\_py}}
\item[]         \sphinxcode{\sphinxupquote{quit}} ()
\end{DUlineblock}

We need to install cppad\_py so you can import it from any directory.

\index{install@\spxentry{install}}\ignorespaces 

\subparagraph{Install}
\label{\detokenize{xsrst/setup_py:install}}\label{\detokenize{xsrst/setup_py:setup-py-install}}\label{\detokenize{xsrst/setup_py:index-14}}
Use the following command to build and install cppad\_py:

\begin{DUlineblock}{0em}
\item[]         \sphinxcode{\sphinxupquote{python setup.py install \sphinxhyphen{}\sphinxhyphen{}prefix}} = \sphinxstyleemphasis{prefix}
\end{DUlineblock}

This will install \sphinxcode{\sphinxupquote{cppad\_py}} in the directory

\begin{DUlineblock}{0em}
\item[]         \sphinxstyleemphasis{prefix} \sphinxcode{\sphinxupquote{/}} \sphinxstyleemphasis{lib} \sphinxcode{\sphinxupquote{/python}} \sphinxstyleemphasis{major} . \sphinxstyleemphasis{minor} \sphinxcode{\sphinxupquote{/site\_packages/cppad\_py}}
\end{DUlineblock}

where \sphinxstyleemphasis{lib} is \sphinxcode{\sphinxupquote{lib}} or \sphinxcode{\sphinxupquote{lib64}} ,
\sphinxstyleemphasis{major} ( \sphinxstyleemphasis{minor} ) is the major (minor)
version of \sphinxstyleemphasis{python} .

\index{python@\spxentry{python}}\index{path@\spxentry{path}}\ignorespaces 

\subparagraph{Python Path}
\label{\detokenize{xsrst/setup_py:python-path}}\label{\detokenize{xsrst/setup_py:setup-py-python-path}}\label{\detokenize{xsrst/setup_py:index-15}}
Check that the directory

\begin{DUlineblock}{0em}
\item[]         \sphinxstyleemphasis{prefix} \sphinxcode{\sphinxupquote{/}} \sphinxstyleemphasis{lib} \sphinxcode{\sphinxupquote{/python}} \sphinxstyleemphasis{major} . \sphinxstyleemphasis{minor} \sphinxcode{\sphinxupquote{/site\_packages}}
\end{DUlineblock}

is in your python path.
Once it is, you should be able to execute the following commands:

\begin{DUlineblock}{0em}
\item[]         \sphinxcode{\sphinxupquote{python}}
\item[]         \sphinxcode{\sphinxupquote{import sys}}
\item[]         \sphinxcode{\sphinxupquote{print(sys.path}} )
\item[]         \sphinxcode{\sphinxupquote{quit}} ()
\end{DUlineblock}


\bigskip\hrule\bigskip


xsrst input file: \sphinxcode{\sphinxupquote{setup.py}}


\subsubsection{numeric\_xam}
\label{\detokenize{xsrst/numeric_xam:numeric-xam}}\label{\detokenize{xsrst/numeric_xam::doc}}
\index{numeric\_xam@\spxentry{numeric\_xam}}\index{numerical@\spxentry{numerical}}\index{examples@\spxentry{examples}}\ignorespaces 

\paragraph{2 Numerical Examples}
\label{\detokenize{xsrst/numeric_xam:numerical-examples}}\label{\detokenize{xsrst/numeric_xam:id1}}\begin{itemize}
\item {} 
{\hyperref[\detokenize{xsrst/numeric_xam:numeric-xam-children}]{\sphinxcrossref{\DUrole{std,std-ref}{Children}}}}

\end{itemize}

\index{children@\spxentry{children}}\ignorespaces 

\subparagraph{Children}
\label{\detokenize{xsrst/numeric_xam:children}}\label{\detokenize{xsrst/numeric_xam:numeric-xam-children}}\label{\detokenize{xsrst/numeric_xam:index-1}}

\subparagraph{numeric\_simple\_inv}
\label{\detokenize{xsrst/numeric_simple_inv:numeric-simple-inv}}\label{\detokenize{xsrst/numeric_simple_inv::doc}}
\index{numeric\_simple\_inv@\spxentry{numeric\_simple\_inv}}\index{an@\spxentry{an}}\index{ad@\spxentry{ad}}\index{compatible@\spxentry{compatible}}\index{matrix@\spxentry{matrix}}\index{inverse@\spxentry{inverse}}\index{routine@\spxentry{routine}}\ignorespaces 

\subparagraph{2.1 An AD Compatible Matrix Inverse Routine}
\label{\detokenize{xsrst/numeric_simple_inv:an-ad-compatible-matrix-inverse-routine}}\label{\detokenize{xsrst/numeric_simple_inv:id1}}\begin{itemize}
\item {} 
{\hyperref[\detokenize{xsrst/numeric_simple_inv:numeric-simple-inv-syntax}]{\sphinxcrossref{\DUrole{std,std-ref}{Syntax}}}}

\item {} 
{\hyperref[\detokenize{xsrst/numeric_simple_inv:numeric-simple-inv-purpose}]{\sphinxcrossref{\DUrole{std,std-ref}{Purpose}}}}

\item {} 
{\hyperref[\detokenize{xsrst/numeric_simple_inv:numeric-simple-inv-a}]{\sphinxcrossref{\DUrole{std,std-ref}{A}}}}

\item {} 
{\hyperref[\detokenize{xsrst/numeric_simple_inv:numeric-simple-inv-ainv}]{\sphinxcrossref{\DUrole{std,std-ref}{Ainv}}}}

\item {} 
{\hyperref[\detokenize{xsrst/numeric_simple_inv:numeric-simple-inv-example}]{\sphinxcrossref{\DUrole{std,std-ref}{Example}}}}

\item {} 
{\hyperref[\detokenize{xsrst/numeric_simple_inv:numeric-simple-inv-source-code}]{\sphinxcrossref{\DUrole{std,std-ref}{Source Code}}}}

\end{itemize}

\index{syntax@\spxentry{syntax}}\ignorespaces 

\subparagraph{Syntax}
\label{\detokenize{xsrst/numeric_simple_inv:syntax}}\label{\detokenize{xsrst/numeric_simple_inv:numeric-simple-inv-syntax}}\label{\detokenize{xsrst/numeric_simple_inv:index-1}}
\sphinxstyleemphasis{Ainv} =  \sphinxcode{\sphinxupquote{simple\_inv}} ( \sphinxstyleemphasis{A} )

\index{purpose@\spxentry{purpose}}\ignorespaces 

\subparagraph{Purpose}
\label{\detokenize{xsrst/numeric_simple_inv:purpose}}\label{\detokenize{xsrst/numeric_simple_inv:numeric-simple-inv-purpose}}\label{\detokenize{xsrst/numeric_simple_inv:index-2}}
This routine can be used with \sphinxcode{\sphinxupquote{ad\_double}}


\subparagraph{A}
\label{\detokenize{xsrst/numeric_simple_inv:a}}\label{\detokenize{xsrst/numeric_simple_inv:numeric-simple-inv-a}}
This must be an invertible square matrix (no singular detection is done).
The type of its elements can be \sphinxcode{\sphinxupquote{float}} or \sphinxcode{\sphinxupquote{a\_double}} .

\index{ainv@\spxentry{ainv}}\ignorespaces 

\subparagraph{Ainv}
\label{\detokenize{xsrst/numeric_simple_inv:ainv}}\label{\detokenize{xsrst/numeric_simple_inv:numeric-simple-inv-ainv}}\label{\detokenize{xsrst/numeric_simple_inv:index-3}}
This is the matrix inverse of \sphinxstyleemphasis{A} .


\subparagraph{numeric\_simple\_inv\_xam\_py}
\label{\detokenize{xsrst/numeric_simple_inv_xam_py:numeric-simple-inv-xam-py}}\label{\detokenize{xsrst/numeric_simple_inv_xam_py::doc}}
\index{numeric\_simple\_inv\_xam\_py@\spxentry{numeric\_simple\_inv\_xam\_py}}\index{example@\spxentry{example}}\index{computing@\spxentry{computing}}\index{derivatives@\spxentry{derivatives}}\index{matrix@\spxentry{matrix}}\index{inversion@\spxentry{inversion}}\ignorespaces 

\subparagraph{2.1.1 Example Computing Derivatives of Matrix Inversion}
\label{\detokenize{xsrst/numeric_simple_inv_xam_py:example-computing-derivatives-of-matrix-inversion}}\label{\detokenize{xsrst/numeric_simple_inv_xam_py:id1}}\begin{itemize}
\item {} 
{\hyperref[\detokenize{xsrst/numeric_simple_inv_xam_py:numeric-simple-inv-xam-py-problem}]{\sphinxcrossref{\DUrole{std,std-ref}{Problem}}}}

\item {} 
{\hyperref[\detokenize{xsrst/numeric_simple_inv_xam_py:numeric-simple-inv-xam-py-source-code}]{\sphinxcrossref{\DUrole{std,std-ref}{Source Code}}}}

\end{itemize}

\index{problem@\spxentry{problem}}\ignorespaces 

\subparagraph{Problem}
\label{\detokenize{xsrst/numeric_simple_inv_xam_py:problem}}\label{\detokenize{xsrst/numeric_simple_inv_xam_py:numeric-simple-inv-xam-py-problem}}\label{\detokenize{xsrst/numeric_simple_inv_xam_py:index-1}}
We define
\begin{equation*}
\begin{split}A(x) = \left( \begin{array}{cc}
x_0 & x_1 \\
x_2 & x_3
\end{array} \right)\end{split}
\end{equation*}
It follows that
\begin{equation*}
\begin{split}A^{-1}(x) =
\frac{1}{x_3 * x_0 - x_1 * x_2}
\left( \begin{array}{cc}
x_3 & - x_1 \\
- x_2 & x_0
\end{array} \right)\end{split}
\end{equation*}
We define
\begin{equation*}
\begin{split}f(x) = (x_3 * x_0 - x_1 * x_2) [
A_{0,0}^{-1} (x),
A_{0,1}^{-1} (x),
A_{1,0}^{-1} (x),
A_{1,1}^{-1} (x)
]
= [ x_3, -x_1, -x_2, x_0]\end{split}
\end{equation*}
The following example below check the derivative of \(f(x)\)

\index{source@\spxentry{source}}\index{code@\spxentry{code}}\ignorespaces 

\subparagraph{Source Code}
\label{\detokenize{xsrst/numeric_simple_inv_xam_py:source-code}}\label{\detokenize{xsrst/numeric_simple_inv_xam_py:numeric-simple-inv-xam-py-source-code}}\label{\detokenize{xsrst/numeric_simple_inv_xam_py:index-2}}
\begin{sphinxVerbatim}[commandchars=\\\{\}]
\PYG{k+kn}{import} \PYG{n+nn}{numpy}
\PYG{k+kn}{import} \PYG{n+nn}{cppad\PYGZus{}py}
\PYG{k+kn}{from} \PYG{n+nn}{simple\PYGZus{}inv} \PYG{k+kn}{import} \PYG{n}{simple\PYGZus{}inv}
\PYG{k}{def} \PYG{n+nf}{simple\PYGZus{}inv\PYGZus{}xam}\PYG{p}{(}\PYG{p}{)} \PYG{p}{:}
    \PYG{n}{ok} \PYG{o}{=} \PYG{k+kc}{True}
    \PYG{n}{x}  \PYG{o}{=} \PYG{n}{numpy}\PYG{o}{.}\PYG{n}{array}\PYG{p}{(} \PYG{p}{[}\PYG{l+m+mf}{1.0}\PYG{p}{,} \PYG{l+m+mf}{2.0}\PYG{p}{,} \PYG{l+m+mf}{3.0}\PYG{p}{,} \PYG{l+m+mf}{4.0}\PYG{p}{]} \PYG{p}{)}
    \PYG{n}{ax} \PYG{o}{=} \PYG{n}{cppad\PYGZus{}py}\PYG{o}{.}\PYG{n}{independent}\PYG{p}{(}\PYG{n}{x}\PYG{p}{)}
    \PYG{c+c1}{\PYGZsh{}}
    \PYG{c+c1}{\PYGZsh{} create the matrix [ [x\PYGZus{}0, x\PYGZus{}1], [x\PYGZus{}2, x\PYGZus{}3] ]}
    \PYG{n}{aA} \PYG{o}{=} \PYG{n}{numpy}\PYG{o}{.}\PYG{n}{reshape}\PYG{p}{(}\PYG{n}{ax}\PYG{p}{,} \PYG{p}{(}\PYG{l+m+mi}{2}\PYG{p}{,}\PYG{l+m+mi}{2}\PYG{p}{)}\PYG{p}{,} \PYG{n}{order}\PYG{o}{=}\PYG{l+s+s1}{\PYGZsq{}}\PYG{l+s+s1}{C}\PYG{l+s+s1}{\PYGZsq{}} \PYG{p}{)}
    \PYG{c+c1}{\PYGZsh{}}
    \PYG{c+c1}{\PYGZsh{} compute inverse of A}
    \PYG{n}{aAinv} \PYG{o}{=} \PYG{n}{simple\PYGZus{}inv}\PYG{p}{(}\PYG{n}{aA}\PYG{p}{)}
    \PYG{c+c1}{\PYGZsh{}}
    \PYG{c+c1}{\PYGZsh{} create f(x)}
    \PYG{n}{ay} \PYG{o}{=} \PYG{n}{numpy}\PYG{o}{.}\PYG{n}{reshape}\PYG{p}{(}\PYG{n}{aAinv}\PYG{p}{,} \PYG{l+m+mi}{4}\PYG{p}{,} \PYG{n}{order}\PYG{o}{=}\PYG{l+s+s1}{\PYGZsq{}}\PYG{l+s+s1}{C}\PYG{l+s+s1}{\PYGZsq{}} \PYG{p}{)}
    \PYG{n}{ay} \PYG{o}{=} \PYG{p}{(}\PYG{n}{ax}\PYG{p}{[}\PYG{l+m+mi}{3}\PYG{p}{]} \PYG{o}{*} \PYG{n}{ax}\PYG{p}{[}\PYG{l+m+mi}{0}\PYG{p}{]} \PYG{o}{\PYGZhy{}} \PYG{n}{ax}\PYG{p}{[}\PYG{l+m+mi}{1}\PYG{p}{]} \PYG{o}{*} \PYG{n}{ax}\PYG{p}{[}\PYG{l+m+mi}{2}\PYG{p}{]}\PYG{p}{)} \PYG{o}{*} \PYG{n}{ay}
    \PYG{n}{f}  \PYG{o}{=} \PYG{n}{cppad\PYGZus{}py}\PYG{o}{.}\PYG{n}{d\PYGZus{}fun}\PYG{p}{(}\PYG{n}{ax}\PYG{p}{,} \PYG{n}{ay}\PYG{p}{)}
    \PYG{c+c1}{\PYGZsh{}}
    \PYG{c+c1}{\PYGZsh{} evaluate the derivative f\PYGZsq{}(x)}
    \PYG{n}{J}      \PYG{o}{=} \PYG{n}{f}\PYG{o}{.}\PYG{n}{jacobian}\PYG{p}{(}\PYG{n}{x}\PYG{p}{)}
    \PYG{c+c1}{\PYGZsh{}}
    \PYG{c+c1}{\PYGZsh{} Check the derivative}
    \PYG{n}{check}      \PYG{o}{=} \PYG{n}{numpy}\PYG{o}{.}\PYG{n}{zeros}\PYG{p}{(} \PYG{p}{(}\PYG{l+m+mi}{4}\PYG{p}{,} \PYG{l+m+mi}{4}\PYG{p}{)} \PYG{p}{)}
    \PYG{n}{check}\PYG{p}{[}\PYG{l+m+mi}{0}\PYG{p}{,}\PYG{l+m+mi}{3}\PYG{p}{]} \PYG{o}{=}  \PYG{l+m+mf}{1.0}
    \PYG{n}{check}\PYG{p}{[}\PYG{l+m+mi}{1}\PYG{p}{,}\PYG{l+m+mi}{1}\PYG{p}{]} \PYG{o}{=} \PYG{o}{\PYGZhy{}}\PYG{l+m+mf}{1.0}
    \PYG{n}{check}\PYG{p}{[}\PYG{l+m+mi}{2}\PYG{p}{,}\PYG{l+m+mi}{2}\PYG{p}{]} \PYG{o}{=} \PYG{o}{\PYGZhy{}}\PYG{l+m+mf}{1.0}
    \PYG{n}{check}\PYG{p}{[}\PYG{l+m+mi}{3}\PYG{p}{,}\PYG{l+m+mi}{0}\PYG{p}{]} \PYG{o}{=}  \PYG{l+m+mf}{1.0}
    \PYG{c+c1}{\PYGZsh{}}
    \PYG{n}{eps99} \PYG{o}{=} \PYG{l+m+mf}{99.0} \PYG{o}{*} \PYG{n}{numpy}\PYG{o}{.}\PYG{n}{finfo}\PYG{p}{(}\PYG{n+nb}{float}\PYG{p}{)}\PYG{o}{.}\PYG{n}{eps}
    \PYG{n}{ok}   \PYG{o}{\PYGZam{}}\PYG{o}{=} \PYG{n}{numpy}\PYG{o}{.}\PYG{n}{all}\PYG{p}{(} \PYG{n}{numpy}\PYG{o}{.}\PYG{n}{fabs}\PYG{p}{(}\PYG{n}{J} \PYG{o}{\PYGZhy{}} \PYG{n}{check}\PYG{p}{)} \PYG{o}{\PYGZlt{}} \PYG{n}{eps99} \PYG{p}{)}
    \PYG{k}{return} \PYG{n}{ok}

\end{sphinxVerbatim}


\bigskip\hrule\bigskip


xsrst input file: \sphinxcode{\sphinxupquote{example/python/numeric/simple\_inv\_xam.py}}

\index{example@\spxentry{example}}\ignorespaces 

\subparagraph{Example}
\label{\detokenize{xsrst/numeric_simple_inv:example}}\label{\detokenize{xsrst/numeric_simple_inv:numeric-simple-inv-example}}\label{\detokenize{xsrst/numeric_simple_inv:index-4}}
{\hyperref[\detokenize{xsrst/numeric_simple_inv_xam_py:id1}]{\sphinxcrossref{\DUrole{std,std-ref}{numeric\_simple\_inv\_xam\_py}}}}

\index{source@\spxentry{source}}\index{code@\spxentry{code}}\ignorespaces 

\subparagraph{Source Code}
\label{\detokenize{xsrst/numeric_simple_inv:source-code}}\label{\detokenize{xsrst/numeric_simple_inv:numeric-simple-inv-source-code}}\label{\detokenize{xsrst/numeric_simple_inv:index-5}}
When viewing the source code below it is important to know that
optimizes out multiplication by the constant one while recording a function.
It also optimizes out both addition and multiplication by the constant zero.

\begin{sphinxVerbatim}[commandchars=\\\{\}]
\PYG{k+kn}{import} \PYG{n+nn}{numpy}
\PYG{k+kn}{import} \PYG{n+nn}{copy}
\PYG{k+kn}{from} \PYG{n+nn}{cppad\PYGZus{}py} \PYG{k+kn}{import} \PYG{n}{a\PYGZus{}double}
\PYG{k}{def} \PYG{n+nf}{simple\PYGZus{}inv}\PYG{p}{(}\PYG{n}{A}\PYG{p}{)} \PYG{p}{:}
    \PYG{n}{nr}\PYG{p}{,} \PYG{n}{nc} \PYG{o}{=} \PYG{n}{A}\PYG{o}{.}\PYG{n}{shape}
    \PYG{n}{has\PYGZus{}a\PYGZus{}double} \PYG{o}{=} \PYG{k+kc}{False}
    \PYG{k}{for} \PYG{n}{e} \PYG{o+ow}{in} \PYG{n}{A}\PYG{o}{.}\PYG{n}{flatten}\PYG{p}{(}\PYG{p}{)} \PYG{p}{:}
        \PYG{n}{has\PYGZus{}a\PYGZus{}double} \PYG{o}{=} \PYG{n}{has\PYGZus{}a\PYGZus{}double} \PYG{o+ow}{or} \PYG{n+nb}{type}\PYG{p}{(}\PYG{n}{e}\PYG{p}{)} \PYG{o}{==} \PYG{n}{a\PYGZus{}double}
    \PYG{k}{assert} \PYG{n}{nr} \PYG{o}{==} \PYG{n}{nc}
    \PYG{c+c1}{\PYGZsh{}}
    \PYG{k}{if} \PYG{n}{has\PYGZus{}a\PYGZus{}double} \PYG{p}{:}
        \PYG{n}{e\PYGZus{}type} \PYG{o}{=} \PYG{n}{a\PYGZus{}double}
    \PYG{k}{else} \PYG{p}{:}
        \PYG{n}{e\PYGZus{}type} \PYG{o}{=} \PYG{n+nb}{float}
    \PYG{c+c1}{\PYGZsh{}}
    \PYG{c+c1}{\PYGZsh{} initialize Tablue = [A | I]}
    \PYG{n}{Tablue} \PYG{o}{=} \PYG{n}{numpy}\PYG{o}{.}\PYG{n}{empty}\PYG{p}{(} \PYG{p}{(}\PYG{n}{nr}\PYG{p}{,} \PYG{l+m+mi}{2} \PYG{o}{*} \PYG{n}{nr}\PYG{p}{)}\PYG{p}{,} \PYG{n}{dtype}\PYG{o}{=}\PYG{n}{e\PYGZus{}type} \PYG{p}{)}
    \PYG{n}{Tablue}\PYG{p}{[}\PYG{p}{:}\PYG{p}{,}\PYG{l+m+mi}{0}\PYG{p}{:}\PYG{n}{nr}\PYG{p}{]} \PYG{o}{=} \PYG{n}{A}
    \PYG{k}{for} \PYG{n}{i} \PYG{o+ow}{in} \PYG{n+nb}{range}\PYG{p}{(}\PYG{n}{nr}\PYG{p}{)} \PYG{p}{:}
        \PYG{k}{for} \PYG{n}{j} \PYG{o+ow}{in} \PYG{n+nb}{range}\PYG{p}{(}\PYG{n}{nr}\PYG{p}{)} \PYG{p}{:}
            \PYG{n}{Tablue}\PYG{p}{[}\PYG{n}{i}\PYG{p}{,} \PYG{n}{j}\PYG{p}{]}   \PYG{o}{=} \PYG{n}{e\PYGZus{}type}\PYG{p}{(} \PYG{n}{A}\PYG{p}{[}\PYG{n}{i}\PYG{p}{,} \PYG{n}{j}\PYG{p}{]} \PYG{p}{)}
            \PYG{n}{Tablue}\PYG{p}{[}\PYG{n}{i}\PYG{p}{,}\PYG{n}{nr}\PYG{o}{+}\PYG{n}{j}\PYG{p}{]} \PYG{o}{=} \PYG{n+nb}{float}\PYG{p}{(}\PYG{n}{i} \PYG{o}{==} \PYG{n}{j}\PYG{p}{)}
    \PYG{c+c1}{\PYGZsh{} \PYGZhy{}\PYGZhy{}\PYGZhy{}\PYGZhy{}\PYGZhy{}\PYGZhy{}\PYGZhy{}\PYGZhy{}\PYGZhy{}\PYGZhy{}\PYGZhy{}\PYGZhy{}\PYGZhy{}\PYGZhy{}\PYGZhy{}\PYGZhy{}\PYGZhy{}\PYGZhy{}\PYGZhy{}\PYGZhy{}\PYGZhy{}\PYGZhy{}\PYGZhy{}\PYGZhy{}\PYGZhy{}\PYGZhy{}\PYGZhy{}\PYGZhy{}\PYGZhy{}\PYGZhy{}\PYGZhy{}\PYGZhy{}\PYGZhy{}\PYGZhy{}\PYGZhy{}\PYGZhy{}\PYGZhy{}\PYGZhy{}\PYGZhy{}\PYGZhy{}\PYGZhy{}\PYGZhy{}\PYGZhy{}\PYGZhy{}\PYGZhy{}\PYGZhy{}\PYGZhy{}\PYGZhy{}\PYGZhy{}\PYGZhy{}\PYGZhy{}\PYGZhy{}\PYGZhy{}\PYGZhy{}\PYGZhy{}\PYGZhy{}\PYGZhy{}\PYGZhy{}\PYGZhy{}\PYGZhy{}\PYGZhy{}\PYGZhy{}\PYGZhy{}\PYGZhy{}\PYGZhy{}\PYGZhy{}\PYGZhy{}\PYGZhy{}\PYGZhy{}\PYGZhy{}\PYGZhy{}\PYGZhy{}}
    \PYG{c+c1}{\PYGZsh{} Use row reduction to get convert A to an upper traingular matrix}
    \PYG{c+c1}{\PYGZsh{} with ones along the diagonal}
    \PYG{c+c1}{\PYGZsh{} \PYGZhy{}\PYGZhy{}\PYGZhy{}\PYGZhy{}\PYGZhy{}\PYGZhy{}\PYGZhy{}\PYGZhy{}\PYGZhy{}\PYGZhy{}\PYGZhy{}\PYGZhy{}\PYGZhy{}\PYGZhy{}\PYGZhy{}\PYGZhy{}\PYGZhy{}\PYGZhy{}\PYGZhy{}\PYGZhy{}\PYGZhy{}\PYGZhy{}\PYGZhy{}\PYGZhy{}\PYGZhy{}\PYGZhy{}\PYGZhy{}\PYGZhy{}\PYGZhy{}\PYGZhy{}\PYGZhy{}\PYGZhy{}\PYGZhy{}\PYGZhy{}\PYGZhy{}\PYGZhy{}\PYGZhy{}\PYGZhy{}\PYGZhy{}\PYGZhy{}\PYGZhy{}\PYGZhy{}\PYGZhy{}\PYGZhy{}\PYGZhy{}\PYGZhy{}\PYGZhy{}\PYGZhy{}\PYGZhy{}\PYGZhy{}\PYGZhy{}\PYGZhy{}\PYGZhy{}\PYGZhy{}\PYGZhy{}\PYGZhy{}\PYGZhy{}\PYGZhy{}\PYGZhy{}\PYGZhy{}\PYGZhy{}\PYGZhy{}\PYGZhy{}\PYGZhy{}\PYGZhy{}\PYGZhy{}\PYGZhy{}\PYGZhy{}\PYGZhy{}\PYGZhy{}\PYGZhy{}\PYGZhy{}}
    \PYG{c+c1}{\PYGZsh{} for each column (except the last) of the matrix}
    \PYG{k}{for} \PYG{n}{j} \PYG{o+ow}{in} \PYG{n+nb}{range}\PYG{p}{(}\PYG{n}{nr} \PYG{o}{\PYGZhy{}} \PYG{l+m+mi}{1}\PYG{p}{)} \PYG{p}{:}
        \PYG{c+c1}{\PYGZsh{} next row is one with maximum absolute element in column j}
        \PYG{k}{if} \PYG{n}{has\PYGZus{}a\PYGZus{}double} \PYG{p}{:}
            \PYG{n}{i\PYGZus{}max} \PYG{o}{=} \PYG{n}{numpy}\PYG{o}{.}\PYG{n}{argmax}\PYG{p}{(} \PYG{n}{numpy}\PYG{o}{.}\PYG{n}{fabs}\PYG{p}{(} \PYG{n}{Tablue}\PYG{p}{[}\PYG{p}{:}\PYG{p}{,}\PYG{n}{j}\PYG{p}{]} \PYG{p}{)} \PYG{p}{)}
        \PYG{k}{else} \PYG{p}{:}
            \PYG{n}{i\PYGZus{}max} \PYG{o}{=} \PYG{n}{numpy}\PYG{o}{.}\PYG{n}{argmax}\PYG{p}{(} \PYG{n}{numpy}\PYG{o}{.}\PYG{n}{abs}\PYG{p}{(} \PYG{n}{Tablue}\PYG{p}{[}\PYG{p}{:}\PYG{p}{,}\PYG{n}{j}\PYG{p}{]} \PYG{p}{)} \PYG{p}{)}
        \PYG{c+c1}{\PYGZsh{}}
        \PYG{c+c1}{\PYGZsh{} swap row j and row i\PYGZus{}max}
        \PYG{k}{if} \PYG{n}{i\PYGZus{}max} \PYG{o}{!=} \PYG{n}{j} \PYG{p}{:}
            \PYG{n}{tmp}              \PYG{o}{=} \PYG{n}{copy}\PYG{o}{.}\PYG{n}{copy}\PYG{p}{(} \PYG{n}{Tablue}\PYG{p}{[}\PYG{n}{j}\PYG{p}{,}\PYG{p}{:}\PYG{p}{]} \PYG{p}{)}
            \PYG{n}{Tablue}\PYG{p}{[}\PYG{n}{j}\PYG{p}{,}\PYG{p}{:}\PYG{p}{]}      \PYG{o}{=} \PYG{n}{Tablue}\PYG{p}{[}\PYG{n}{i\PYGZus{}max}\PYG{p}{,}\PYG{p}{:}\PYG{p}{]}
            \PYG{n}{Tablue}\PYG{p}{[}\PYG{n}{i\PYGZus{}max}\PYG{p}{,}\PYG{p}{:}\PYG{p}{]}  \PYG{o}{=} \PYG{n}{tmp}
        \PYG{c+c1}{\PYGZsh{}}
        \PYG{c+c1}{\PYGZsh{} divide row j by Tablue[j,j]}
        \PYG{n}{Tablue}\PYG{p}{[}\PYG{n}{j}\PYG{p}{,}\PYG{p}{:}\PYG{p}{]} \PYG{o}{=} \PYG{n}{Tablue}\PYG{p}{[}\PYG{n}{j}\PYG{p}{,}\PYG{p}{:}\PYG{p}{]} \PYG{o}{/} \PYG{n}{Tablue}\PYG{p}{[}\PYG{n}{j}\PYG{p}{,}\PYG{n}{j}\PYG{p}{]}
        \PYG{n}{Tablue}\PYG{p}{[}\PYG{n}{j}\PYG{p}{,}\PYG{n}{j}\PYG{p}{]} \PYG{o}{=} \PYG{n}{e\PYGZus{}type}\PYG{p}{(}\PYG{l+m+mf}{1.0}\PYG{p}{)}
        \PYG{c+c1}{\PYGZsh{}}
        \PYG{c+c1}{\PYGZsh{} zero out column j in rows after j}
        \PYG{k}{for} \PYG{n}{k} \PYG{o+ow}{in} \PYG{n+nb}{range}\PYG{p}{(}\PYG{n}{nr} \PYG{o}{\PYGZhy{}} \PYG{n}{j} \PYG{o}{\PYGZhy{}} \PYG{l+m+mi}{1}\PYG{p}{)} \PYG{p}{:}
            \PYG{n}{i} \PYG{o}{=} \PYG{n}{j} \PYG{o}{+} \PYG{n}{k} \PYG{o}{+} \PYG{l+m+mi}{1}
            \PYG{n}{Tablue}\PYG{p}{[}\PYG{n}{i}\PYG{p}{,}\PYG{p}{:}\PYG{p}{]} \PYG{o}{=} \PYG{n}{Tablue}\PYG{p}{[}\PYG{n}{i}\PYG{p}{,}\PYG{p}{:}\PYG{p}{]} \PYG{o}{\PYGZhy{}} \PYG{n}{Tablue}\PYG{p}{[}\PYG{n}{i}\PYG{p}{,}\PYG{n}{j}\PYG{p}{]} \PYG{o}{*} \PYG{n}{Tablue}\PYG{p}{[}\PYG{n}{j}\PYG{p}{,}\PYG{p}{:}\PYG{p}{]}
            \PYG{n}{Tablue}\PYG{p}{[}\PYG{n}{i}\PYG{p}{,}\PYG{n}{j}\PYG{p}{]}    \PYG{o}{=} \PYG{n}{e\PYGZus{}type}\PYG{p}{(} \PYG{l+m+mf}{0.0} \PYG{p}{)}
    \PYG{c+c1}{\PYGZsh{} divide the last row of Tablue by the last pivot element}
    \PYG{n}{Tablue}\PYG{p}{[}\PYG{n}{nr}\PYG{o}{\PYGZhy{}}\PYG{l+m+mi}{1}\PYG{p}{,}\PYG{p}{:}\PYG{p}{]}     \PYG{o}{=} \PYG{n}{Tablue}\PYG{p}{[}\PYG{n}{nr}\PYG{o}{\PYGZhy{}}\PYG{l+m+mi}{1}\PYG{p}{,}\PYG{p}{:}\PYG{p}{]} \PYG{o}{/} \PYG{n}{Tablue}\PYG{p}{[}\PYG{n}{nr}\PYG{o}{\PYGZhy{}}\PYG{l+m+mi}{1}\PYG{p}{,}\PYG{n}{nr}\PYG{o}{\PYGZhy{}}\PYG{l+m+mi}{1}\PYG{p}{]}
    \PYG{n}{Tablue}\PYG{p}{[}\PYG{n}{nr}\PYG{o}{\PYGZhy{}}\PYG{l+m+mi}{1}\PYG{p}{,} \PYG{n}{nr}\PYG{o}{\PYGZhy{}}\PYG{l+m+mi}{1}\PYG{p}{]} \PYG{o}{=} \PYG{n}{e\PYGZus{}type}\PYG{p}{(} \PYG{l+m+mf}{1.0} \PYG{p}{)}
    \PYG{c+c1}{\PYGZsh{} \PYGZhy{}\PYGZhy{}\PYGZhy{}\PYGZhy{}\PYGZhy{}\PYGZhy{}\PYGZhy{}\PYGZhy{}\PYGZhy{}\PYGZhy{}\PYGZhy{}\PYGZhy{}\PYGZhy{}\PYGZhy{}\PYGZhy{}\PYGZhy{}\PYGZhy{}\PYGZhy{}\PYGZhy{}\PYGZhy{}\PYGZhy{}\PYGZhy{}\PYGZhy{}\PYGZhy{}\PYGZhy{}\PYGZhy{}\PYGZhy{}\PYGZhy{}\PYGZhy{}\PYGZhy{}\PYGZhy{}\PYGZhy{}\PYGZhy{}\PYGZhy{}\PYGZhy{}\PYGZhy{}\PYGZhy{}\PYGZhy{}\PYGZhy{}\PYGZhy{}\PYGZhy{}\PYGZhy{}\PYGZhy{}\PYGZhy{}\PYGZhy{}\PYGZhy{}\PYGZhy{}\PYGZhy{}\PYGZhy{}\PYGZhy{}\PYGZhy{}\PYGZhy{}\PYGZhy{}\PYGZhy{}\PYGZhy{}\PYGZhy{}\PYGZhy{}\PYGZhy{}\PYGZhy{}\PYGZhy{}\PYGZhy{}\PYGZhy{}\PYGZhy{}\PYGZhy{}\PYGZhy{}\PYGZhy{}\PYGZhy{}\PYGZhy{}\PYGZhy{}\PYGZhy{}\PYGZhy{}\PYGZhy{}}
    \PYG{c+c1}{\PYGZsh{} Use row reduction to get convert upper traingular matrix to the identity}
    \PYG{c+c1}{\PYGZsh{} \PYGZhy{}\PYGZhy{}\PYGZhy{}\PYGZhy{}\PYGZhy{}\PYGZhy{}\PYGZhy{}\PYGZhy{}\PYGZhy{}\PYGZhy{}\PYGZhy{}\PYGZhy{}\PYGZhy{}\PYGZhy{}\PYGZhy{}\PYGZhy{}\PYGZhy{}\PYGZhy{}\PYGZhy{}\PYGZhy{}\PYGZhy{}\PYGZhy{}\PYGZhy{}\PYGZhy{}\PYGZhy{}\PYGZhy{}\PYGZhy{}\PYGZhy{}\PYGZhy{}\PYGZhy{}\PYGZhy{}\PYGZhy{}\PYGZhy{}\PYGZhy{}\PYGZhy{}\PYGZhy{}\PYGZhy{}\PYGZhy{}\PYGZhy{}\PYGZhy{}\PYGZhy{}\PYGZhy{}\PYGZhy{}\PYGZhy{}\PYGZhy{}\PYGZhy{}\PYGZhy{}\PYGZhy{}\PYGZhy{}\PYGZhy{}\PYGZhy{}\PYGZhy{}\PYGZhy{}\PYGZhy{}\PYGZhy{}\PYGZhy{}\PYGZhy{}\PYGZhy{}\PYGZhy{}\PYGZhy{}\PYGZhy{}\PYGZhy{}\PYGZhy{}\PYGZhy{}\PYGZhy{}\PYGZhy{}\PYGZhy{}\PYGZhy{}\PYGZhy{}\PYGZhy{}\PYGZhy{}\PYGZhy{}}
    \PYG{k}{for} \PYG{n}{jm} \PYG{o+ow}{in} \PYG{n+nb}{reversed}\PYG{p}{(}\PYG{n+nb}{range}\PYG{p}{(}\PYG{n}{nr} \PYG{o}{\PYGZhy{}} \PYG{l+m+mi}{1}\PYG{p}{)} \PYG{p}{)} \PYG{p}{:}
        \PYG{n}{j} \PYG{o}{=} \PYG{n}{jm} \PYG{o}{+} \PYG{l+m+mi}{1}
        \PYG{c+c1}{\PYGZsh{} zero out entries in column j and row less than j}
        \PYG{k}{for} \PYG{n}{i} \PYG{o+ow}{in} \PYG{n+nb}{range}\PYG{p}{(}\PYG{n}{j}\PYG{p}{)} \PYG{p}{:}
            \PYG{n}{Tablue}\PYG{p}{[}\PYG{n}{i}\PYG{p}{,}\PYG{p}{:}\PYG{p}{]} \PYG{o}{=} \PYG{n}{Tablue}\PYG{p}{[}\PYG{n}{i}\PYG{p}{,}\PYG{p}{:}\PYG{p}{]} \PYG{o}{\PYGZhy{}} \PYG{n}{Tablue}\PYG{p}{[}\PYG{n}{i}\PYG{p}{,}\PYG{n}{j}\PYG{p}{]} \PYG{o}{*} \PYG{n}{Tablue}\PYG{p}{[}\PYG{n}{j}\PYG{p}{,}\PYG{p}{:}\PYG{p}{]}
            \PYG{n}{Tablue}\PYG{p}{[}\PYG{n}{i}\PYG{p}{,}\PYG{n}{j}\PYG{p}{]} \PYG{o}{=} \PYG{n}{e\PYGZus{}type}\PYG{p}{(} \PYG{l+m+mf}{0.0} \PYG{p}{)}
    \PYG{c+c1}{\PYGZsh{} \PYGZhy{}\PYGZhy{}\PYGZhy{}\PYGZhy{}\PYGZhy{}\PYGZhy{}\PYGZhy{}\PYGZhy{}\PYGZhy{}\PYGZhy{}\PYGZhy{}\PYGZhy{}\PYGZhy{}\PYGZhy{}\PYGZhy{}\PYGZhy{}\PYGZhy{}\PYGZhy{}\PYGZhy{}\PYGZhy{}\PYGZhy{}\PYGZhy{}\PYGZhy{}\PYGZhy{}\PYGZhy{}\PYGZhy{}\PYGZhy{}\PYGZhy{}\PYGZhy{}\PYGZhy{}\PYGZhy{}\PYGZhy{}\PYGZhy{}\PYGZhy{}\PYGZhy{}\PYGZhy{}\PYGZhy{}\PYGZhy{}\PYGZhy{}\PYGZhy{}\PYGZhy{}\PYGZhy{}\PYGZhy{}\PYGZhy{}\PYGZhy{}\PYGZhy{}\PYGZhy{}\PYGZhy{}\PYGZhy{}\PYGZhy{}\PYGZhy{}\PYGZhy{}\PYGZhy{}\PYGZhy{}\PYGZhy{}\PYGZhy{}\PYGZhy{}\PYGZhy{}\PYGZhy{}\PYGZhy{}\PYGZhy{}\PYGZhy{}\PYGZhy{}\PYGZhy{}\PYGZhy{}\PYGZhy{}\PYGZhy{}\PYGZhy{}\PYGZhy{}\PYGZhy{}\PYGZhy{}\PYGZhy{}\PYGZhy{}}
    \PYG{n}{Ainv} \PYG{o}{=} \PYG{n}{Tablue}\PYG{p}{[}\PYG{p}{:}\PYG{p}{,} \PYG{n}{nr}\PYG{p}{:}\PYG{p}{]}
    \PYG{k}{return} \PYG{n}{Ainv}
\end{sphinxVerbatim}


\bigskip\hrule\bigskip


xsrst input file: \sphinxcode{\sphinxupquote{example/python/numeric/simple\_inv.py}}


\subparagraph{numeric\_runge4\_step}
\label{\detokenize{xsrst/numeric_runge4_step:numeric-runge4-step}}\label{\detokenize{xsrst/numeric_runge4_step::doc}}
\index{numeric\_runge4\_step@\spxentry{numeric\_runge4\_step}}\index{one@\spxentry{one}}\index{fourth@\spxentry{fourth}}\index{order@\spxentry{order}}\index{runge\sphinxhyphen{}kutta@\spxentry{runge\sphinxhyphen{}kutta}}\index{ode@\spxentry{ode}}\index{step@\spxentry{step}}\ignorespaces 

\subparagraph{2.2 One Fourth Order Runge\sphinxhyphen{}Kutta ODE Step}
\label{\detokenize{xsrst/numeric_runge4_step:one-fourth-order-runge-kutta-ode-step}}\label{\detokenize{xsrst/numeric_runge4_step:id1}}\begin{itemize}
\item {} 
{\hyperref[\detokenize{xsrst/numeric_runge4_step:numeric-runge4-step-syntax}]{\sphinxcrossref{\DUrole{std,std-ref}{Syntax}}}}

\item {} 
{\hyperref[\detokenize{xsrst/numeric_runge4_step:numeric-runge4-step-purpose}]{\sphinxcrossref{\DUrole{std,std-ref}{Purpose}}}}

\item {} 
{\hyperref[\detokenize{xsrst/numeric_runge4_step:numeric-runge4-step-fun}]{\sphinxcrossref{\DUrole{std,std-ref}{fun}}}}

\item {} 
{\hyperref[\detokenize{xsrst/numeric_runge4_step:numeric-runge4-step-ti}]{\sphinxcrossref{\DUrole{std,std-ref}{ti}}}}

\item {} 
{\hyperref[\detokenize{xsrst/numeric_runge4_step:numeric-runge4-step-yi}]{\sphinxcrossref{\DUrole{std,std-ref}{yi}}}}

\item {} 
{\hyperref[\detokenize{xsrst/numeric_runge4_step:numeric-runge4-step-h}]{\sphinxcrossref{\DUrole{std,std-ref}{h}}}}

\item {} 
{\hyperref[\detokenize{xsrst/numeric_runge4_step:numeric-runge4-step-yf}]{\sphinxcrossref{\DUrole{std,std-ref}{yf}}}}

\item {} 
{\hyperref[\detokenize{xsrst/numeric_runge4_step:numeric-runge4-step-example}]{\sphinxcrossref{\DUrole{std,std-ref}{Example}}}}

\item {} 
{\hyperref[\detokenize{xsrst/numeric_runge4_step:numeric-runge4-step-source-code}]{\sphinxcrossref{\DUrole{std,std-ref}{Source Code}}}}

\end{itemize}

\index{syntax@\spxentry{syntax}}\ignorespaces 

\subparagraph{Syntax}
\label{\detokenize{xsrst/numeric_runge4_step:syntax}}\label{\detokenize{xsrst/numeric_runge4_step:numeric-runge4-step-syntax}}\label{\detokenize{xsrst/numeric_runge4_step:index-1}}
\sphinxstyleemphasis{yf} =  \sphinxcode{\sphinxupquote{runge4\_step}} ( \sphinxstyleemphasis{fun} , \sphinxstyleemphasis{ti} , \sphinxstyleemphasis{yi} , \sphinxstyleemphasis{h} )

\index{purpose@\spxentry{purpose}}\ignorespaces 

\subparagraph{Purpose}
\label{\detokenize{xsrst/numeric_runge4_step:purpose}}\label{\detokenize{xsrst/numeric_runge4_step:numeric-runge4-step-purpose}}\label{\detokenize{xsrst/numeric_runge4_step:index-2}}
The routine can be used with \sphinxcode{\sphinxupquote{ad\_double}}
to solve the initial value ODE
\begin{equation*}
\begin{split}y^{(1)} (t)  = f( t , y )\end{split}
\end{equation*}
\index{fun@\spxentry{fun}}\ignorespaces 

\subparagraph{fun}
\label{\detokenize{xsrst/numeric_runge4_step:fun}}\label{\detokenize{xsrst/numeric_runge4_step:numeric-runge4-step-fun}}\label{\detokenize{xsrst/numeric_runge4_step:index-3}}
This is a function that evaluates the ordinary differential equation
using the syntax

\begin{DUlineblock}{0em}
\item[]         \sphinxstyleemphasis{yp} = \sphinxstyleemphasis{fun} . \sphinxcode{\sphinxupquote{f}} \sphinxstyleemphasis{t} , \sphinxstyleemphasis{y} )
\end{DUlineblock}

where \sphinxstyleemphasis{t} \# is the current time,
\sphinxstyleemphasis{y} is the current value of \(y(t)\), and
\sphinxstyleemphasis{yp} is the current derivative \(y^{(1)} (t)\).
The type of the elements of \sphinxstyleemphasis{t} and \sphinxstyleemphasis{y}
can be \sphinxcode{\sphinxupquote{float}} or \sphinxcode{\sphinxupquote{ad\_double}} .

\index{ti@\spxentry{ti}}\ignorespaces 

\subparagraph{ti}
\label{\detokenize{xsrst/numeric_runge4_step:ti}}\label{\detokenize{xsrst/numeric_runge4_step:numeric-runge4-step-ti}}\label{\detokenize{xsrst/numeric_runge4_step:index-4}}
This is the initial time for the Runge\sphinxhyphen{}Kutta step.
It can have type \sphinxcode{\sphinxupquote{float}} or \sphinxcode{\sphinxupquote{a\_double}} .

\index{yi@\spxentry{yi}}\ignorespaces 

\subparagraph{yi}
\label{\detokenize{xsrst/numeric_runge4_step:yi}}\label{\detokenize{xsrst/numeric_runge4_step:numeric-runge4-step-yi}}\label{\detokenize{xsrst/numeric_runge4_step:index-5}}
This is the numpy vector containing the
value of \(y(t)\) at the initial time.
The type of its elements can be \sphinxcode{\sphinxupquote{float}} or \sphinxcode{\sphinxupquote{a\_double}} .

\index{h@\spxentry{h}}\ignorespaces 

\subparagraph{h}
\label{\detokenize{xsrst/numeric_runge4_step:h}}\label{\detokenize{xsrst/numeric_runge4_step:numeric-runge4-step-h}}\label{\detokenize{xsrst/numeric_runge4_step:index-6}}
This is the step size in time; i.e., the time at the end of the step
minus the initial time.
It can have type \sphinxcode{\sphinxupquote{float}} or \sphinxcode{\sphinxupquote{a\_double}} .

\index{yf@\spxentry{yf}}\ignorespaces 

\subparagraph{yf}
\label{\detokenize{xsrst/numeric_runge4_step:yf}}\label{\detokenize{xsrst/numeric_runge4_step:numeric-runge4-step-yf}}\label{\detokenize{xsrst/numeric_runge4_step:index-7}}
This is the approximate solution for \(y(t)\) at the final time
as a numpy vector.
This solution is 4\sphinxhyphen{}th order accurate in time \(t\); e.g., if
\(y(t)\) is a polynomial in \(t\) of order four or lower,
the solution has no truncation error, only round off error.


\subparagraph{numeric\_runge4\_step\_xam\_py}
\label{\detokenize{xsrst/numeric_runge4_step_xam_py:numeric-runge4-step-xam-py}}\label{\detokenize{xsrst/numeric_runge4_step_xam_py::doc}}
\index{numeric\_runge4\_step\_xam\_py@\spxentry{numeric\_runge4\_step\_xam\_py}}\index{example@\spxentry{example}}\index{computing@\spxentry{computing}}\index{derivative@\spxentry{derivative}}\index{runge\sphinxhyphen{}kutta@\spxentry{runge\sphinxhyphen{}kutta}}\index{ode@\spxentry{ode}}\index{solution@\spxentry{solution}}\ignorespaces 

\subparagraph{2.2.1 Example Computing Derivative A Runge\sphinxhyphen{}Kutta Ode Solution}
\label{\detokenize{xsrst/numeric_runge4_step_xam_py:example-computing-derivative-a-runge-kutta-ode-solution}}\label{\detokenize{xsrst/numeric_runge4_step_xam_py:id1}}\begin{itemize}
\item {} 
{\hyperref[\detokenize{xsrst/numeric_runge4_step_xam_py:numeric-runge4-step-xam-py-ode}]{\sphinxcrossref{\DUrole{std,std-ref}{ODE}}}}

\item {} 
{\hyperref[\detokenize{xsrst/numeric_runge4_step_xam_py:numeric-runge4-step-xam-py-solution}]{\sphinxcrossref{\DUrole{std,std-ref}{Solution}}}}

\item {} 
{\hyperref[\detokenize{xsrst/numeric_runge4_step_xam_py:numeric-runge4-step-xam-py-derivative-of-solution}]{\sphinxcrossref{\DUrole{std,std-ref}{Derivative of Solution}}}}

\item {} 
{\hyperref[\detokenize{xsrst/numeric_runge4_step_xam_py:numeric-runge4-step-xam-py-source-code}]{\sphinxcrossref{\DUrole{std,std-ref}{Source Code}}}}

\end{itemize}

\index{ode@\spxentry{ode}}\ignorespaces 

\subparagraph{ODE}
\label{\detokenize{xsrst/numeric_runge4_step_xam_py:ode}}\label{\detokenize{xsrst/numeric_runge4_step_xam_py:numeric-runge4-step-xam-py-ode}}\label{\detokenize{xsrst/numeric_runge4_step_xam_py:index-1}}\begin{equation*}
\begin{split}\partial_t y_i (t, x) =  f(t, y, x) \left\{ \begin{array}{rl}
     x_0               & {\rm if} \; i = 0 \\
     x_i y_{i-1} (t)   & {\rm otherwise}
\end{array} \right.\end{split}
\end{equation*}
with the initial condition \(y(0) = 0\)

\index{solution@\spxentry{solution}}\ignorespaces 

\subparagraph{Solution}
\label{\detokenize{xsrst/numeric_runge4_step_xam_py:solution}}\label{\detokenize{xsrst/numeric_runge4_step_xam_py:numeric-runge4-step-xam-py-solution}}\label{\detokenize{xsrst/numeric_runge4_step_xam_py:index-2}}
This is a special case for which we know the solution
\begin{equation*}
\begin{split}y_i (t, x) = \left\{ \begin{array}{rl}
     t  x_0                            & {\rm if} \; i = 0 \\
     ( t^i / (i+1) ! ) \prod_{j=0}^i x_j   & {\rm otherwise}
\end{array} \right.\end{split}
\end{equation*}
\index{derivative@\spxentry{derivative}}\index{solution@\spxentry{solution}}\ignorespaces 

\subparagraph{Derivative of Solution}
\label{\detokenize{xsrst/numeric_runge4_step_xam_py:derivative-of-solution}}\label{\detokenize{xsrst/numeric_runge4_step_xam_py:numeric-runge4-step-xam-py-derivative-of-solution}}\label{\detokenize{xsrst/numeric_runge4_step_xam_py:index-3}}
For this special case, the partial derivative of the solution with respect
to the j\sphinxhyphen{}th component of the vector \(x\) is
\begin{equation*}
\begin{split}\partial_{x(j)} y_i (t, x) =  \left\{ \begin{array}{rl}
     y_i (t, x) / x_j      & {\rm if} \; j \leq i \\
     0                     & {\rm otherwise}
\end{array} \right.\end{split}
\end{equation*}
\index{source@\spxentry{source}}\index{code@\spxentry{code}}\ignorespaces 

\subparagraph{Source Code}
\label{\detokenize{xsrst/numeric_runge4_step_xam_py:source-code}}\label{\detokenize{xsrst/numeric_runge4_step_xam_py:numeric-runge4-step-xam-py-source-code}}\label{\detokenize{xsrst/numeric_runge4_step_xam_py:index-4}}
\begin{sphinxVerbatim}[commandchars=\\\{\}]
\PYG{k+kn}{import} \PYG{n+nn}{numpy}
\PYG{k+kn}{from}  \PYG{n+nn}{scipy}\PYG{n+nn}{.}\PYG{n+nn}{special} \PYG{k+kn}{import} \PYG{n}{factorial}
\PYG{k+kn}{import} \PYG{n+nn}{cppad\PYGZus{}py}
\PYG{k+kn}{from} \PYG{n+nn}{runge4\PYGZus{}step} \PYG{k+kn}{import} \PYG{n}{runge4\PYGZus{}step}
\PYG{c+c1}{\PYGZsh{}}
\PYG{k}{class} \PYG{n+nc}{fun\PYGZus{}class} \PYG{p}{:}
    \PYG{c+c1}{\PYGZsh{}}
    \PYG{k}{def} \PYG{n+nf+fm}{\PYGZus{}\PYGZus{}init\PYGZus{}\PYGZus{}}\PYG{p}{(}\PYG{n+nb+bp}{self}\PYG{p}{,} \PYG{n}{x}\PYG{p}{)} \PYG{p}{:}
        \PYG{n+nb+bp}{self}\PYG{o}{.}\PYG{n}{x} \PYG{o}{=} \PYG{n}{x}
    \PYG{c+c1}{\PYGZsh{}}
    \PYG{k}{def} \PYG{n+nf}{f}\PYG{p}{(}\PYG{n+nb+bp}{self}\PYG{p}{,} \PYG{n}{t}\PYG{p}{,} \PYG{n}{y}\PYG{p}{)} \PYG{p}{:}
        \PYG{n}{y\PYGZus{}shift} \PYG{o}{=} \PYG{n}{numpy}\PYG{o}{.}\PYG{n}{concatenate}\PYG{p}{(} \PYG{p}{(} \PYG{p}{[}\PYG{l+m+mf}{1.0}\PYG{p}{]} \PYG{p}{,} \PYG{n}{y}\PYG{p}{[}\PYG{l+m+mi}{0}\PYG{p}{:}\PYG{o}{\PYGZhy{}}\PYG{l+m+mi}{1}\PYG{p}{]} \PYG{p}{)} \PYG{p}{)}
        \PYG{k}{return} \PYG{n+nb+bp}{self}\PYG{o}{.}\PYG{n}{x} \PYG{o}{*} \PYG{n}{y\PYGZus{}shift}
\PYG{c+c1}{\PYGZsh{}}
\PYG{k}{def} \PYG{n+nf}{runge4\PYGZus{}step\PYGZus{}xam}\PYG{p}{(}\PYG{p}{)} \PYG{p}{:}
    \PYG{n}{ok}    \PYG{o}{=} \PYG{k+kc}{True}
    \PYG{n}{nx}    \PYG{o}{=} \PYG{l+m+mi}{4}
    \PYG{n}{eps99} \PYG{o}{=} \PYG{l+m+mf}{99.0} \PYG{o}{*} \PYG{n}{numpy}\PYG{o}{.}\PYG{n}{finfo}\PYG{p}{(}\PYG{n+nb}{float}\PYG{p}{)}\PYG{o}{.}\PYG{n}{eps}
    \PYG{c+c1}{\PYGZsh{}}
    \PYG{c+c1}{\PYGZsh{} independent variables for this g(x) = y(1, x)}
    \PYG{n}{x}  \PYG{o}{=} \PYG{n}{numpy}\PYG{o}{.}\PYG{n}{array}\PYG{p}{(} \PYG{n}{nx} \PYG{o}{*} \PYG{p}{[} \PYG{l+m+mf}{1.0} \PYG{p}{]} \PYG{p}{)}
    \PYG{n}{ax} \PYG{o}{=} \PYG{n}{cppad\PYGZus{}py}\PYG{o}{.}\PYG{n}{independent}\PYG{p}{(}\PYG{n}{x}\PYG{p}{)}
    \PYG{c+c1}{\PYGZsh{}}
    \PYG{c+c1}{\PYGZsh{} function to pass to runge4\PYGZus{}step}
    \PYG{n}{fun} \PYG{o}{=} \PYG{n}{fun\PYGZus{}class}\PYG{p}{(}\PYG{n}{ax}\PYG{p}{)}
    \PYG{c+c1}{\PYGZsh{}}
    \PYG{c+c1}{\PYGZsh{} initiali value for the ODE}
    \PYG{n}{ay\PYGZus{}start} \PYG{o}{=}  \PYG{n}{numpy}\PYG{o}{.}\PYG{n}{array}\PYG{p}{(} \PYG{n}{nx} \PYG{o}{*} \PYG{p}{[} \PYG{n}{cppad\PYGZus{}py}\PYG{o}{.}\PYG{n}{a\PYGZus{}double}\PYG{p}{(}\PYG{l+m+mf}{0.0}\PYG{p}{)} \PYG{p}{]} \PYG{p}{)}
    \PYG{n}{t\PYGZus{}start}  \PYG{o}{=} \PYG{l+m+mf}{0.0}
    \PYG{n}{t\PYGZus{}step}   \PYG{o}{=} \PYG{l+m+mf}{0.75}
    \PYG{c+c1}{\PYGZsh{}}
    \PYG{c+c1}{\PYGZsh{} take one step}
    \PYG{n}{ay} \PYG{o}{=} \PYG{n}{runge4\PYGZus{}step}\PYG{p}{(}\PYG{n}{fun}\PYG{p}{,} \PYG{n}{t\PYGZus{}start}\PYG{p}{,} \PYG{n}{ay\PYGZus{}start}\PYG{p}{,} \PYG{n}{t\PYGZus{}step}\PYG{p}{)}
    \PYG{c+c1}{\PYGZsh{}}
    \PYG{c+c1}{\PYGZsh{} g(x) = y(t\PYGZus{}step, x)}
    \PYG{n}{g} \PYG{o}{=} \PYG{n}{cppad\PYGZus{}py}\PYG{o}{.}\PYG{n}{d\PYGZus{}fun}\PYG{p}{(}\PYG{n}{ax}\PYG{p}{,} \PYG{n}{ay}\PYG{p}{)}
    \PYG{c+c1}{\PYGZsh{}}
    \PYG{c+c1}{\PYGZsh{} Check g\PYGZus{}i (x) = prod\PYGZus{}\PYGZob{}j=0\PYGZcb{}\PYGZca{}i x[j] t\PYGZca{}(i+1) / (i+1) !}
    \PYG{c+c1}{\PYGZsh{} 4th order method should be very accutate for functions}
    \PYG{c+c1}{\PYGZsh{} or order 4 or less in t.}
    \PYG{n}{x}  \PYG{o}{=} \PYG{n}{numpy}\PYG{o}{.}\PYG{n}{arange}\PYG{p}{(}\PYG{l+m+mi}{0}\PYG{p}{,} \PYG{n}{nx}\PYG{p}{)} \PYG{o}{+} \PYG{l+m+mf}{1.0}
    \PYG{n}{gx} \PYG{o}{=} \PYG{n}{g}\PYG{o}{.}\PYG{n}{forward}\PYG{p}{(}\PYG{l+m+mi}{0}\PYG{p}{,} \PYG{n}{x}\PYG{p}{)}
    \PYG{n}{prod} \PYG{o}{=} \PYG{l+m+mf}{1.0}
    \PYG{k}{for} \PYG{n}{i} \PYG{o+ow}{in} \PYG{n+nb}{range}\PYG{p}{(}\PYG{n}{nx}\PYG{p}{)} \PYG{p}{:}
        \PYG{n}{prod}      \PYG{o}{=} \PYG{n}{prod} \PYG{o}{*} \PYG{n}{x}\PYG{p}{[}\PYG{n}{i}\PYG{p}{]}
        \PYG{n}{check}     \PYG{o}{=} \PYG{n}{prod} \PYG{o}{*} \PYG{n}{numpy}\PYG{o}{.}\PYG{n}{power}\PYG{p}{(}\PYG{n}{t\PYGZus{}step}\PYG{p}{,} \PYG{n}{i}\PYG{o}{+}\PYG{l+m+mi}{1}\PYG{p}{)} \PYG{o}{/} \PYG{n}{factorial}\PYG{p}{(}\PYG{n}{i}\PYG{o}{+}\PYG{l+m+mi}{1}\PYG{p}{)}
        \PYG{n}{rel\PYGZus{}error} \PYG{o}{=} \PYG{n}{gx}\PYG{p}{[}\PYG{n}{i}\PYG{p}{]} \PYG{o}{/} \PYG{n}{check} \PYG{o}{\PYGZhy{}} \PYG{l+m+mf}{1.0}
        \PYG{n}{ok}       \PYG{o}{\PYGZam{}}\PYG{o}{=} \PYG{n+nb}{abs}\PYG{p}{(}\PYG{n}{rel\PYGZus{}error}\PYG{p}{)} \PYG{o}{\PYGZlt{}} \PYG{n}{eps99}
    \PYG{c+c1}{\PYGZsh{}}
    \PYG{c+c1}{\PYGZsh{} Check derivative of g\PYGZus{}i (x) w.r.t x\PYGZus{}j}
    \PYG{c+c1}{\PYGZsh{} is zero for i \PYGZlt{} j and  g\PYGZus{}i(x) / x\PYGZus{}j otherwise}
    \PYG{n}{J} \PYG{o}{=} \PYG{n}{g}\PYG{o}{.}\PYG{n}{jacobian}\PYG{p}{(}\PYG{n}{x}\PYG{p}{)}
    \PYG{k}{for} \PYG{n}{j} \PYG{o+ow}{in} \PYG{n+nb}{range}\PYG{p}{(}\PYG{n}{nx}\PYG{p}{)} \PYG{p}{:}
        \PYG{k}{for} \PYG{n}{i} \PYG{o+ow}{in} \PYG{n+nb}{range}\PYG{p}{(}\PYG{n}{nx}\PYG{p}{)} \PYG{p}{:}
            \PYG{k}{if} \PYG{n}{i} \PYG{o}{\PYGZlt{}} \PYG{n}{j} \PYG{p}{:}
                \PYG{n}{ok} \PYG{o}{\PYGZam{}}\PYG{o}{=} \PYG{n}{J}\PYG{p}{[}\PYG{n}{i}\PYG{p}{,}\PYG{n}{j}\PYG{p}{]} \PYG{o}{==} \PYG{l+m+mf}{0.0}
            \PYG{k}{else} \PYG{p}{:}
                \PYG{n}{check}     \PYG{o}{=} \PYG{n}{gx}\PYG{p}{[}\PYG{n}{i}\PYG{p}{]} \PYG{o}{/} \PYG{n}{x}\PYG{p}{[}\PYG{n}{j}\PYG{p}{]}
                \PYG{n}{rel\PYGZus{}error} \PYG{o}{=} \PYG{n}{J}\PYG{p}{[}\PYG{n}{i}\PYG{p}{,}\PYG{n}{j}\PYG{p}{]} \PYG{o}{/} \PYG{n}{check} \PYG{o}{\PYGZhy{}} \PYG{l+m+mf}{1.0}
                \PYG{n}{ok}       \PYG{o}{\PYGZam{}}\PYG{o}{=} \PYG{n+nb}{abs}\PYG{p}{(}\PYG{n}{rel\PYGZus{}error}\PYG{p}{)} \PYG{o}{\PYGZlt{}} \PYG{n}{eps99}

    \PYG{k}{return} \PYG{n}{ok}
\end{sphinxVerbatim}


\bigskip\hrule\bigskip


xsrst input file: \sphinxcode{\sphinxupquote{example/python/numeric/runge4\_step\_xam.py}}

\index{example@\spxentry{example}}\ignorespaces 

\subparagraph{Example}
\label{\detokenize{xsrst/numeric_runge4_step:example}}\label{\detokenize{xsrst/numeric_runge4_step:numeric-runge4-step-example}}\label{\detokenize{xsrst/numeric_runge4_step:index-8}}
{\hyperref[\detokenize{xsrst/numeric_runge4_step_xam_py:id1}]{\sphinxcrossref{\DUrole{std,std-ref}{numeric\_runge4\_step\_xam\_py}}}}

\index{source@\spxentry{source}}\index{code@\spxentry{code}}\ignorespaces 

\subparagraph{Source Code}
\label{\detokenize{xsrst/numeric_runge4_step:source-code}}\label{\detokenize{xsrst/numeric_runge4_step:numeric-runge4-step-source-code}}\label{\detokenize{xsrst/numeric_runge4_step:index-9}}
\begin{sphinxVerbatim}[commandchars=\\\{\}]
\PYG{k}{def} \PYG{n+nf}{runge4\PYGZus{}step}\PYG{p}{(}\PYG{n}{fun}\PYG{p}{,} \PYG{n}{ti}\PYG{p}{,} \PYG{n}{yi}\PYG{p}{,} \PYG{n}{h}\PYG{p}{)} \PYG{p}{:}
    \PYG{n}{k1}     \PYG{o}{=} \PYG{n}{h} \PYG{o}{*} \PYG{n}{fun}\PYG{o}{.}\PYG{n}{f}\PYG{p}{(}\PYG{n}{ti}\PYG{p}{,}           \PYG{n}{yi}\PYG{p}{)}
    \PYG{n}{k2}     \PYG{o}{=} \PYG{n}{h} \PYG{o}{*} \PYG{n}{fun}\PYG{o}{.}\PYG{n}{f}\PYG{p}{(}\PYG{n}{ti} \PYG{o}{+} \PYG{n}{h} \PYG{o}{/} \PYG{l+m+mf}{2.0}\PYG{p}{,} \PYG{n}{yi} \PYG{o}{+} \PYG{n}{k1} \PYG{o}{/} \PYG{l+m+mf}{2.0}\PYG{p}{)}
    \PYG{n}{k3}     \PYG{o}{=} \PYG{n}{h} \PYG{o}{*} \PYG{n}{fun}\PYG{o}{.}\PYG{n}{f}\PYG{p}{(}\PYG{n}{ti} \PYG{o}{+} \PYG{n}{h} \PYG{o}{/} \PYG{l+m+mf}{2.0}\PYG{p}{,} \PYG{n}{yi} \PYG{o}{+} \PYG{n}{k2} \PYG{o}{/} \PYG{l+m+mf}{2.0}\PYG{p}{)}
    \PYG{n}{k4}     \PYG{o}{=} \PYG{n}{h} \PYG{o}{*} \PYG{n}{fun}\PYG{o}{.}\PYG{n}{f}\PYG{p}{(}\PYG{n}{ti} \PYG{o}{+} \PYG{n}{h}\PYG{p}{,}       \PYG{n}{yi} \PYG{o}{+} \PYG{n}{k3} \PYG{p}{)}
    \PYG{n}{yf}     \PYG{o}{=} \PYG{n}{yi} \PYG{o}{+} \PYG{p}{(}\PYG{n}{k1} \PYG{o}{+} \PYG{l+m+mf}{2.0} \PYG{o}{*} \PYG{n}{k2} \PYG{o}{+} \PYG{l+m+mf}{2.0} \PYG{o}{*} \PYG{n}{k3} \PYG{o}{+} \PYG{n}{k4}\PYG{p}{)} \PYG{o}{/} \PYG{l+m+mf}{6.0}
    \PYG{k}{return} \PYG{n}{yf}
\end{sphinxVerbatim}


\bigskip\hrule\bigskip


xsrst input file: \sphinxcode{\sphinxupquote{example/python/numeric/runge4\_step.py}}


\subparagraph{numeric\_rosen3\_step}
\label{\detokenize{xsrst/numeric_rosen3_step:numeric-rosen3-step}}\label{\detokenize{xsrst/numeric_rosen3_step::doc}}
\index{numeric\_rosen3\_step@\spxentry{numeric\_rosen3\_step}}\index{one@\spxentry{one}}\index{third@\spxentry{third}}\index{order@\spxentry{order}}\index{rosenbrock@\spxentry{rosenbrock}}\index{ode@\spxentry{ode}}\index{step@\spxentry{step}}\ignorespaces 

\subparagraph{2.3 One Third Order Rosenbrock ODE Step}
\label{\detokenize{xsrst/numeric_rosen3_step:one-third-order-rosenbrock-ode-step}}\label{\detokenize{xsrst/numeric_rosen3_step:id1}}\begin{itemize}
\item {} 
{\hyperref[\detokenize{xsrst/numeric_rosen3_step:numeric-rosen3-step-syntax}]{\sphinxcrossref{\DUrole{std,std-ref}{Syntax}}}}

\item {} 
{\hyperref[\detokenize{xsrst/numeric_rosen3_step:numeric-rosen3-step-purpose}]{\sphinxcrossref{\DUrole{std,std-ref}{Purpose}}}}

\item {} 
{\hyperref[\detokenize{xsrst/numeric_rosen3_step:numeric-rosen3-step-reference}]{\sphinxcrossref{\DUrole{std,std-ref}{Reference}}}}

\item {} \begin{description}
\item[{{\hyperref[\detokenize{xsrst/numeric_rosen3_step:numeric-rosen3-step-fun}]{\sphinxcrossref{\DUrole{std,std-ref}{fun}}}}}] \leavevmode\begin{itemize}
\item {} 
{\hyperref[\detokenize{xsrst/numeric_rosen3_step:numeric-rosen3-step-fun-t}]{\sphinxcrossref{\DUrole{std,std-ref}{t}}}}

\item {} 
{\hyperref[\detokenize{xsrst/numeric_rosen3_step:numeric-rosen3-step-fun-y}]{\sphinxcrossref{\DUrole{std,std-ref}{y}}}}

\item {} 
{\hyperref[\detokenize{xsrst/numeric_rosen3_step:numeric-rosen3-step-fun-fun-f}]{\sphinxcrossref{\DUrole{std,std-ref}{fun.f}}}}

\item {} 
{\hyperref[\detokenize{xsrst/numeric_rosen3_step:numeric-rosen3-step-fun-fun-f-t}]{\sphinxcrossref{\DUrole{std,std-ref}{fun.f\_t}}}}

\item {} 
{\hyperref[\detokenize{xsrst/numeric_rosen3_step:numeric-rosen3-step-fun-fun-f-y}]{\sphinxcrossref{\DUrole{std,std-ref}{fun.f\_y}}}}

\end{itemize}

\end{description}

\item {} 
{\hyperref[\detokenize{xsrst/numeric_rosen3_step:numeric-rosen3-step-ti}]{\sphinxcrossref{\DUrole{std,std-ref}{ti}}}}

\item {} 
{\hyperref[\detokenize{xsrst/numeric_rosen3_step:numeric-rosen3-step-yi}]{\sphinxcrossref{\DUrole{std,std-ref}{yi}}}}

\item {} 
{\hyperref[\detokenize{xsrst/numeric_rosen3_step:numeric-rosen3-step-h}]{\sphinxcrossref{\DUrole{std,std-ref}{h}}}}

\item {} 
{\hyperref[\detokenize{xsrst/numeric_rosen3_step:numeric-rosen3-step-yf}]{\sphinxcrossref{\DUrole{std,std-ref}{yf}}}}

\item {} 
{\hyperref[\detokenize{xsrst/numeric_rosen3_step:numeric-rosen3-step-ok}]{\sphinxcrossref{\DUrole{std,std-ref}{ok}}}}

\item {} 
{\hyperref[\detokenize{xsrst/numeric_rosen3_step:numeric-rosen3-step-example}]{\sphinxcrossref{\DUrole{std,std-ref}{Example}}}}

\item {} \begin{description}
\item[{{\hyperref[\detokenize{xsrst/numeric_rosen3_step:numeric-rosen3-step-source-code}]{\sphinxcrossref{\DUrole{std,std-ref}{Source Code}}}}}] \leavevmode\begin{itemize}
\item {} 
{\hyperref[\detokenize{xsrst/numeric_rosen3_step:numeric-rosen3-step-source-code-rosen3-step}]{\sphinxcrossref{\DUrole{std,std-ref}{rosen3\_step}}}}

\item {} 
{\hyperref[\detokenize{xsrst/numeric_rosen3_step:numeric-rosen3-step-source-code-check-rosen3-step}]{\sphinxcrossref{\DUrole{std,std-ref}{check\_rosen3\_step}}}}

\end{itemize}

\end{description}

\end{itemize}

\index{syntax@\spxentry{syntax}}\ignorespaces 

\subparagraph{Syntax}
\label{\detokenize{xsrst/numeric_rosen3_step:syntax}}\label{\detokenize{xsrst/numeric_rosen3_step:numeric-rosen3-step-syntax}}\label{\detokenize{xsrst/numeric_rosen3_step:index-1}}
\begin{DUlineblock}{0em}
\item[] \sphinxstyleemphasis{yf} =  \sphinxcode{\sphinxupquote{rosen3\_step}} ( \sphinxstyleemphasis{fun} , \sphinxstyleemphasis{ti} , \sphinxstyleemphasis{yi} , \sphinxstyleemphasis{h} )
\item[] \sphinxstyleemphasis{ok} =  \sphinxcode{\sphinxupquote{check\_rosen3\_step}} ( \sphinxstyleemphasis{fun} , \sphinxstyleemphasis{ti} , \sphinxstyleemphasis{yi} , \sphinxstyleemphasis{h} )
\end{DUlineblock}

\index{purpose@\spxentry{purpose}}\ignorespaces 

\subparagraph{Purpose}
\label{\detokenize{xsrst/numeric_rosen3_step:purpose}}\label{\detokenize{xsrst/numeric_rosen3_step:numeric-rosen3-step-purpose}}\label{\detokenize{xsrst/numeric_rosen3_step:index-2}}
The routine \sphinxcode{\sphinxupquote{rosen3\_step}} can be used with
\sphinxcode{\sphinxupquote{ad\_double}} to solve an initial value ODE
\begin{equation*}
\begin{split}y^{(1)} (t)  = f( t , y )\end{split}
\end{equation*}
\index{reference@\spxentry{reference}}\ignorespaces 

\subparagraph{Reference}
\label{\detokenize{xsrst/numeric_rosen3_step:reference}}\label{\detokenize{xsrst/numeric_rosen3_step:numeric-rosen3-step-reference}}\label{\detokenize{xsrst/numeric_rosen3_step:index-3}}
The formulas in this method are taken from page 100 of the following
reference (except that 98/108 was correct to 97/108):
Shampine, L.F.,
\sphinxstyleemphasis{Implementation of Rosenbrock Methods} ,
ACM Transactions on Mathematical Software, Vol. 8, No. 2, June 1982.

\index{fun@\spxentry{fun}}\ignorespaces 

\subparagraph{fun}
\label{\detokenize{xsrst/numeric_rosen3_step:fun}}\label{\detokenize{xsrst/numeric_rosen3_step:numeric-rosen3-step-fun}}\label{\detokenize{xsrst/numeric_rosen3_step:index-4}}
This is a function that evaluates the ordinary differential equation,
and its partial derivatives,

\index{t@\spxentry{t}}\ignorespaces 

\subparagraph{t}
\label{\detokenize{xsrst/numeric_rosen3_step:t}}\label{\detokenize{xsrst/numeric_rosen3_step:numeric-rosen3-step-fun-t}}\label{\detokenize{xsrst/numeric_rosen3_step:index-5}}
The argument \sphinxstyleemphasis{t} below is the current time.
It can be a \sphinxcode{\sphinxupquote{float}} or \sphinxcode{\sphinxupquote{a\_double}} .

\index{y@\spxentry{y}}\ignorespaces 

\subparagraph{y}
\label{\detokenize{xsrst/numeric_rosen3_step:y}}\label{\detokenize{xsrst/numeric_rosen3_step:numeric-rosen3-step-fun-y}}\label{\detokenize{xsrst/numeric_rosen3_step:index-6}}
The argument \sphinxstyleemphasis{y} below is the current value of \(y(t)\).
The type of the elements of \sphinxstyleemphasis{y}
can be \sphinxcode{\sphinxupquote{float}} or \sphinxcode{\sphinxupquote{ad\_double}} .

\index{fun.f@\spxentry{fun.f}}\ignorespaces 

\subparagraph{fun.f}
\label{\detokenize{xsrst/numeric_rosen3_step:fun-f}}\label{\detokenize{xsrst/numeric_rosen3_step:numeric-rosen3-step-fun-fun-f}}\label{\detokenize{xsrst/numeric_rosen3_step:index-7}}
The syntax \sphinxcode{\sphinxupquote{yp}} = \sphinxcode{\sphinxupquote{fun}} . \sphinxstyleemphasis{f} \sphinxcode{\sphinxupquote{t}} , \sphinxcode{\sphinxupquote{y}} ) returns \sphinxstyleemphasis{yp} ,
the current function value \(f(t, y)\).

\index{fun.f\_t@\spxentry{fun.f\_t}}\ignorespaces 

\subparagraph{fun.f\_t}
\label{\detokenize{xsrst/numeric_rosen3_step:fun-f-t}}\label{\detokenize{xsrst/numeric_rosen3_step:numeric-rosen3-step-fun-fun-f-t}}\label{\detokenize{xsrst/numeric_rosen3_step:index-8}}
The syntax \sphinxcode{\sphinxupquote{yp\_t}} = \sphinxcode{\sphinxupquote{fun}} . \sphinxstyleemphasis{f\_t} ( \sphinxcode{\sphinxupquote{t}} , \sphinxcode{\sphinxupquote{y}} ) returns \sphinxstyleemphasis{yp\_t} ,
the current partial value \(\partial_t f(t, y)\).

\index{fun.f\_y@\spxentry{fun.f\_y}}\ignorespaces 

\subparagraph{fun.f\_y}
\label{\detokenize{xsrst/numeric_rosen3_step:fun-f-y}}\label{\detokenize{xsrst/numeric_rosen3_step:numeric-rosen3-step-fun-fun-f-y}}\label{\detokenize{xsrst/numeric_rosen3_step:index-9}}
The syntax \sphinxcode{\sphinxupquote{yp\_y}} = \sphinxcode{\sphinxupquote{fun}} . \sphinxstyleemphasis{f\_y} ( \sphinxcode{\sphinxupquote{t}} , \sphinxcode{\sphinxupquote{y}} ) returns \sphinxstyleemphasis{yp\_y} ,
the current partial value \(\partial_y f(t, y)\).

\index{ti@\spxentry{ti}}\ignorespaces 

\subparagraph{ti}
\label{\detokenize{xsrst/numeric_rosen3_step:ti}}\label{\detokenize{xsrst/numeric_rosen3_step:numeric-rosen3-step-ti}}\label{\detokenize{xsrst/numeric_rosen3_step:index-10}}
This is the initial time for the Rosenbrock step.
It can have type \sphinxcode{\sphinxupquote{float}} or \sphinxcode{\sphinxupquote{a\_double}} .
(For \sphinxcode{\sphinxupquote{check\_rosen3\_step}} only \sphinxcode{\sphinxupquote{float}} is allowed.)

\index{yi@\spxentry{yi}}\ignorespaces 

\subparagraph{yi}
\label{\detokenize{xsrst/numeric_rosen3_step:yi}}\label{\detokenize{xsrst/numeric_rosen3_step:numeric-rosen3-step-yi}}\label{\detokenize{xsrst/numeric_rosen3_step:index-11}}
This is the numpy vector containing the
value of \(y(t)\) at the initial time.
The type of its elements can be \sphinxcode{\sphinxupquote{float}} or \sphinxcode{\sphinxupquote{a\_double}} .
(For \sphinxcode{\sphinxupquote{check\_rosen3\_step}} only \sphinxcode{\sphinxupquote{float}} is allowed.)

\index{h@\spxentry{h}}\ignorespaces 

\subparagraph{h}
\label{\detokenize{xsrst/numeric_rosen3_step:h}}\label{\detokenize{xsrst/numeric_rosen3_step:numeric-rosen3-step-h}}\label{\detokenize{xsrst/numeric_rosen3_step:index-12}}
This is the step size in time; i.e., the time at the end of the step
minus the initial time.
It can have type \sphinxcode{\sphinxupquote{float}} or \sphinxcode{\sphinxupquote{a\_double}} .
(This is not used by \sphinxcode{\sphinxupquote{check\_rosen3\_step}} .)

\index{yf@\spxentry{yf}}\ignorespaces 

\subparagraph{yf}
\label{\detokenize{xsrst/numeric_rosen3_step:yf}}\label{\detokenize{xsrst/numeric_rosen3_step:numeric-rosen3-step-yf}}\label{\detokenize{xsrst/numeric_rosen3_step:index-13}}
This is the approximate solution for \(y(t)\) at the final time
as a numpy vector.
This solution is 3\sphinxhyphen{}th order accurate in time \(t\); e.g., if
\(y(t)\) is a polynomial in \(t\) of order three or lower,
the solution has no truncation error, only round off error.

\index{ok@\spxentry{ok}}\ignorespaces 

\subparagraph{ok}
\label{\detokenize{xsrst/numeric_rosen3_step:ok}}\label{\detokenize{xsrst/numeric_rosen3_step:numeric-rosen3-step-ok}}\label{\detokenize{xsrst/numeric_rosen3_step:index-14}}
This is true if the function \sphinxstyleemphasis{fun} . \sphinxcode{\sphinxupquote{f}} \sphinxstyleemphasis{t} , \sphinxstyleemphasis{y} )
and the partials \sphinxstyleemphasis{fun} . \sphinxcode{\sphinxupquote{f\_t}} ( \sphinxstyleemphasis{t} , \sphinxstyleemphasis{y} ) ,
\sphinxstyleemphasis{fun} . \sphinxcode{\sphinxupquote{f\_y}} ( \sphinxstyleemphasis{t} , \sphinxstyleemphasis{y} ) agree.
Otherwise AD has detected an error in these functions.


\subparagraph{numeric\_rosen3\_step\_xam\_py}
\label{\detokenize{xsrst/numeric_rosen3_step_xam_py:numeric-rosen3-step-xam-py}}\label{\detokenize{xsrst/numeric_rosen3_step_xam_py::doc}}
\index{numeric\_rosen3\_step\_xam\_py@\spxentry{numeric\_rosen3\_step\_xam\_py}}\index{example@\spxentry{example}}\index{computing@\spxentry{computing}}\index{derivative@\spxentry{derivative}}\index{rosenbrock@\spxentry{rosenbrock}}\index{ode@\spxentry{ode}}\index{solution@\spxentry{solution}}\ignorespaces 

\subparagraph{2.3.1 Example Computing Derivative A Rosenbrock Ode Solution}
\label{\detokenize{xsrst/numeric_rosen3_step_xam_py:example-computing-derivative-a-rosenbrock-ode-solution}}\label{\detokenize{xsrst/numeric_rosen3_step_xam_py:id1}}\begin{itemize}
\item {} 
{\hyperref[\detokenize{xsrst/numeric_rosen3_step_xam_py:numeric-rosen3-step-xam-py-y-t-x}]{\sphinxcrossref{\DUrole{std,std-ref}{y(t, x)}}}}

\item {} 
{\hyperref[\detokenize{xsrst/numeric_rosen3_step_xam_py:numeric-rosen3-step-xam-py-first-ode}]{\sphinxcrossref{\DUrole{std,std-ref}{First ODE}}}}

\item {} 
{\hyperref[\detokenize{xsrst/numeric_rosen3_step_xam_py:numeric-rosen3-step-xam-py-second-ode}]{\sphinxcrossref{\DUrole{std,std-ref}{Second ODE}}}}

\item {} 
{\hyperref[\detokenize{xsrst/numeric_rosen3_step_xam_py:numeric-rosen3-step-xam-py-derivative-of-solution}]{\sphinxcrossref{\DUrole{std,std-ref}{Derivative of Solution}}}}

\item {} 
{\hyperref[\detokenize{xsrst/numeric_rosen3_step_xam_py:numeric-rosen3-step-xam-py-source-code}]{\sphinxcrossref{\DUrole{std,std-ref}{Source Code}}}}

\end{itemize}

\index{y(t@\spxentry{y(t}}\index{x)@\spxentry{x)}}\ignorespaces 

\subparagraph{y(t, x)}
\label{\detokenize{xsrst/numeric_rosen3_step_xam_py:y-t-x}}\label{\detokenize{xsrst/numeric_rosen3_step_xam_py:numeric-rosen3-step-xam-py-y-t-x}}\label{\detokenize{xsrst/numeric_rosen3_step_xam_py:index-1}}
The two initial value problems below have the following solution:
\begin{equation*}
\begin{split}y_i (t, x) = ( t^{i+1} / (i+1) ! ) \prod_{j=0}^i x_j\end{split}
\end{equation*}
\index{first@\spxentry{first}}\index{ode@\spxentry{ode}}\ignorespaces 

\subparagraph{First ODE}
\label{\detokenize{xsrst/numeric_rosen3_step_xam_py:first-ode}}\label{\detokenize{xsrst/numeric_rosen3_step_xam_py:numeric-rosen3-step-xam-py-first-ode}}\label{\detokenize{xsrst/numeric_rosen3_step_xam_py:index-2}}\begin{equation*}
\begin{split}f(t, y, x)  =
\left\{ \begin{array}{rl}
x_0               & {\rm if} \; i = 0 \\
x_i y_{i-1} (t)   & {\rm otherwise}
\end{array} \right.\end{split}
\end{equation*}
with the initial condition \(y(0) = 0\)

\index{second@\spxentry{second}}\index{ode@\spxentry{ode}}\ignorespaces 

\subparagraph{Second ODE}
\label{\detokenize{xsrst/numeric_rosen3_step_xam_py:second-ode}}\label{\detokenize{xsrst/numeric_rosen3_step_xam_py:numeric-rosen3-step-xam-py-second-ode}}\label{\detokenize{xsrst/numeric_rosen3_step_xam_py:index-3}}\begin{equation*}
\begin{split}f(t, y, x)  = ( t^i / i ! ) \prod_{j=0}^i x_j\end{split}
\end{equation*}
with the initial condition \(y(0) = 0\)

\index{derivative@\spxentry{derivative}}\index{solution@\spxentry{solution}}\ignorespaces 

\subparagraph{Derivative of Solution}
\label{\detokenize{xsrst/numeric_rosen3_step_xam_py:derivative-of-solution}}\label{\detokenize{xsrst/numeric_rosen3_step_xam_py:numeric-rosen3-step-xam-py-derivative-of-solution}}\label{\detokenize{xsrst/numeric_rosen3_step_xam_py:index-4}}
For this special case, the partial derivative of the solution with respect
to the j\sphinxhyphen{}th component of the vector \(x\) is
\begin{equation*}
\begin{split}\partial_{x(j)} y_i (t, x) =  \left\{ \begin{array}{rl}
     y_i (t, x) / x_j      & {\rm if} \; j \leq i \\
     0                     & {\rm otherwise}
\end{array} \right.\end{split}
\end{equation*}
\index{source@\spxentry{source}}\index{code@\spxentry{code}}\ignorespaces 

\subparagraph{Source Code}
\label{\detokenize{xsrst/numeric_rosen3_step_xam_py:source-code}}\label{\detokenize{xsrst/numeric_rosen3_step_xam_py:numeric-rosen3-step-xam-py-source-code}}\label{\detokenize{xsrst/numeric_rosen3_step_xam_py:index-5}}
\begin{sphinxVerbatim}[commandchars=\\\{\}]
\PYG{k+kn}{import} \PYG{n+nn}{numpy}
\PYG{k+kn}{from}  \PYG{n+nn}{scipy}\PYG{n+nn}{.}\PYG{n+nn}{special} \PYG{k+kn}{import} \PYG{n}{factorial}
\PYG{k+kn}{import} \PYG{n+nn}{cppad\PYGZus{}py}
\PYG{k+kn}{from} \PYG{n+nn}{rosen3\PYGZus{}step} \PYG{k+kn}{import} \PYG{n}{rosen3\PYGZus{}step}
\PYG{k+kn}{from} \PYG{n+nn}{rosen3\PYGZus{}step} \PYG{k+kn}{import} \PYG{n}{check\PYGZus{}rosen3\PYGZus{}step}
\PYG{c+c1}{\PYGZsh{} \PYGZhy{}\PYGZhy{}\PYGZhy{}\PYGZhy{}\PYGZhy{}\PYGZhy{}\PYGZhy{}\PYGZhy{}\PYGZhy{}\PYGZhy{}\PYGZhy{}\PYGZhy{}\PYGZhy{}\PYGZhy{}\PYGZhy{}\PYGZhy{}\PYGZhy{}\PYGZhy{}\PYGZhy{}\PYGZhy{}\PYGZhy{}\PYGZhy{}\PYGZhy{}\PYGZhy{}\PYGZhy{}\PYGZhy{}\PYGZhy{}\PYGZhy{}\PYGZhy{}\PYGZhy{}\PYGZhy{}\PYGZhy{}\PYGZhy{}\PYGZhy{}\PYGZhy{}\PYGZhy{}\PYGZhy{}\PYGZhy{}\PYGZhy{}\PYGZhy{}\PYGZhy{}\PYGZhy{}\PYGZhy{}\PYGZhy{}\PYGZhy{}\PYGZhy{}\PYGZhy{}\PYGZhy{}\PYGZhy{}\PYGZhy{}\PYGZhy{}\PYGZhy{}\PYGZhy{}\PYGZhy{}\PYGZhy{}\PYGZhy{}\PYGZhy{}\PYGZhy{}\PYGZhy{}\PYGZhy{}\PYGZhy{}\PYGZhy{}\PYGZhy{}\PYGZhy{}\PYGZhy{}\PYGZhy{}\PYGZhy{}\PYGZhy{}\PYGZhy{}\PYGZhy{}\PYGZhy{}\PYGZhy{}\PYGZhy{}\PYGZhy{}\PYGZhy{}}
\PYG{k}{class} \PYG{n+nc}{first\PYGZus{}fun\PYGZus{}class} \PYG{p}{:}
    \PYG{k}{def} \PYG{n+nf+fm}{\PYGZus{}\PYGZus{}init\PYGZus{}\PYGZus{}}\PYG{p}{(}\PYG{n+nb+bp}{self}\PYG{p}{,} \PYG{n}{x}\PYG{p}{)} \PYG{p}{:}
        \PYG{n+nb+bp}{self}\PYG{o}{.}\PYG{n}{x} \PYG{o}{=} \PYG{n}{x}

    \PYG{c+c1}{\PYGZsh{} fun.f}
    \PYG{k}{def} \PYG{n+nf}{f}\PYG{p}{(}\PYG{n+nb+bp}{self}\PYG{p}{,} \PYG{n}{t}\PYG{p}{,} \PYG{n}{y}\PYG{p}{)} \PYG{p}{:}
        \PYG{n}{y\PYGZus{}shift} \PYG{o}{=} \PYG{n}{numpy}\PYG{o}{.}\PYG{n}{concatenate}\PYG{p}{(} \PYG{p}{(} \PYG{p}{[}\PYG{l+m+mf}{1.0}\PYG{p}{]} \PYG{p}{,} \PYG{n}{y}\PYG{p}{[}\PYG{l+m+mi}{0}\PYG{p}{:}\PYG{o}{\PYGZhy{}}\PYG{l+m+mi}{1}\PYG{p}{]} \PYG{p}{)} \PYG{p}{)}
        \PYG{k}{return} \PYG{n+nb+bp}{self}\PYG{o}{.}\PYG{n}{x} \PYG{o}{*} \PYG{n}{y\PYGZus{}shift}
    \PYG{c+c1}{\PYGZsh{}}
    \PYG{c+c1}{\PYGZsh{} fun.f\PYGZus{}t}
    \PYG{k}{def} \PYG{n+nf}{f\PYGZus{}t}\PYG{p}{(}\PYG{n+nb+bp}{self}\PYG{p}{,} \PYG{n}{t}\PYG{p}{,} \PYG{n}{y}\PYG{p}{)} \PYG{p}{:}
        \PYG{n}{ny} \PYG{o}{=} \PYG{n}{y}\PYG{o}{.}\PYG{n}{size}
        \PYG{k}{return} \PYG{n}{numpy}\PYG{o}{.}\PYG{n}{zeros}\PYG{p}{(} \PYG{n}{ny}\PYG{p}{,} \PYG{n}{dtype}\PYG{o}{=}\PYG{n+nb}{type}\PYG{p}{(}\PYG{n}{y}\PYG{p}{[}\PYG{l+m+mi}{0}\PYG{p}{]}\PYG{p}{)} \PYG{p}{)}
    \PYG{c+c1}{\PYGZsh{}}
    \PYG{c+c1}{\PYGZsh{} fun.f\PYGZus{}y}
    \PYG{k}{def} \PYG{n+nf}{f\PYGZus{}y}\PYG{p}{(}\PYG{n+nb+bp}{self}\PYG{p}{,} \PYG{n}{t}\PYG{p}{,} \PYG{n}{y}\PYG{p}{)} \PYG{p}{:}
        \PYG{n}{ny} \PYG{o}{=} \PYG{n}{y}\PYG{o}{.}\PYG{n}{size}
        \PYG{n}{J}  \PYG{o}{=} \PYG{n}{numpy}\PYG{o}{.}\PYG{n}{zeros}\PYG{p}{(} \PYG{p}{(}\PYG{n}{ny}\PYG{p}{,} \PYG{n}{ny}\PYG{p}{)}\PYG{p}{,} \PYG{n}{dtype}\PYG{o}{=}\PYG{n+nb}{type}\PYG{p}{(}\PYG{n}{y}\PYG{p}{[}\PYG{l+m+mi}{0}\PYG{p}{]}\PYG{p}{)} \PYG{p}{)}
        \PYG{k}{for} \PYG{n}{i} \PYG{o+ow}{in} \PYG{n+nb}{range}\PYG{p}{(}\PYG{n}{ny} \PYG{o}{\PYGZhy{}} \PYG{l+m+mi}{1}\PYG{p}{)} \PYG{p}{:}
            \PYG{n}{J}\PYG{p}{[}\PYG{n}{i}\PYG{o}{+}\PYG{l+m+mi}{1}\PYG{p}{,} \PYG{n}{i}\PYG{p}{]} \PYG{o}{=} \PYG{n+nb+bp}{self}\PYG{o}{.}\PYG{n}{x}\PYG{p}{[}\PYG{n}{i}\PYG{o}{+}\PYG{l+m+mi}{1}\PYG{p}{]}
        \PYG{k}{return} \PYG{n}{J}
\PYG{c+c1}{\PYGZsh{} \PYGZhy{}\PYGZhy{}\PYGZhy{}\PYGZhy{}\PYGZhy{}\PYGZhy{}\PYGZhy{}\PYGZhy{}\PYGZhy{}\PYGZhy{}\PYGZhy{}\PYGZhy{}\PYGZhy{}\PYGZhy{}\PYGZhy{}\PYGZhy{}\PYGZhy{}\PYGZhy{}\PYGZhy{}\PYGZhy{}\PYGZhy{}\PYGZhy{}\PYGZhy{}\PYGZhy{}\PYGZhy{}\PYGZhy{}\PYGZhy{}\PYGZhy{}\PYGZhy{}\PYGZhy{}\PYGZhy{}\PYGZhy{}\PYGZhy{}\PYGZhy{}\PYGZhy{}\PYGZhy{}\PYGZhy{}\PYGZhy{}\PYGZhy{}\PYGZhy{}\PYGZhy{}\PYGZhy{}\PYGZhy{}\PYGZhy{}\PYGZhy{}\PYGZhy{}\PYGZhy{}\PYGZhy{}\PYGZhy{}\PYGZhy{}\PYGZhy{}\PYGZhy{}\PYGZhy{}\PYGZhy{}\PYGZhy{}\PYGZhy{}\PYGZhy{}\PYGZhy{}\PYGZhy{}\PYGZhy{}\PYGZhy{}\PYGZhy{}\PYGZhy{}\PYGZhy{}\PYGZhy{}\PYGZhy{}\PYGZhy{}\PYGZhy{}\PYGZhy{}\PYGZhy{}\PYGZhy{}\PYGZhy{}\PYGZhy{}\PYGZhy{}\PYGZhy{}}
\PYG{k}{class} \PYG{n+nc}{second\PYGZus{}fun\PYGZus{}class} \PYG{p}{:}
    \PYG{k}{def} \PYG{n+nf+fm}{\PYGZus{}\PYGZus{}init\PYGZus{}\PYGZus{}}\PYG{p}{(}\PYG{n+nb+bp}{self}\PYG{p}{,} \PYG{n}{x}\PYG{p}{)} \PYG{p}{:}
        \PYG{n+nb+bp}{self}\PYG{o}{.}\PYG{n}{x}  \PYG{o}{=} \PYG{n}{x}
    \PYG{c+c1}{\PYGZsh{}}
    \PYG{c+c1}{\PYGZsh{} fun.f}
    \PYG{k}{def} \PYG{n+nf}{f}\PYG{p}{(}\PYG{n+nb+bp}{self}\PYG{p}{,} \PYG{n}{t}\PYG{p}{,} \PYG{n}{y}\PYG{p}{)} \PYG{p}{:}
        \PYG{n}{ny}        \PYG{o}{=} \PYG{n}{y}\PYG{o}{.}\PYG{n}{size}
        \PYG{n}{type\PYGZus{}y}    \PYG{o}{=} \PYG{n+nb}{type}\PYG{p}{(}\PYG{n}{y}\PYG{p}{[}\PYG{l+m+mi}{0}\PYG{p}{]}\PYG{p}{)}
        \PYG{n}{result}    \PYG{o}{=} \PYG{n}{numpy}\PYG{o}{.}\PYG{n}{empty}\PYG{p}{(}\PYG{n}{ny}\PYG{p}{,} \PYG{n}{dtype}\PYG{o}{=}\PYG{n+nb}{type}\PYG{p}{(}\PYG{n}{y}\PYG{p}{[}\PYG{l+m+mi}{0}\PYG{p}{]}\PYG{p}{)} \PYG{p}{)}
        \PYG{n}{result}\PYG{p}{[}\PYG{l+m+mi}{0}\PYG{p}{]} \PYG{o}{=} \PYG{n}{type\PYGZus{}y}\PYG{p}{(} \PYG{n+nb+bp}{self}\PYG{o}{.}\PYG{n}{x}\PYG{p}{[}\PYG{l+m+mi}{0}\PYG{p}{]} \PYG{p}{)}
        \PYG{k}{for} \PYG{n}{i} \PYG{o+ow}{in} \PYG{n+nb}{range}\PYG{p}{(}\PYG{n}{ny} \PYG{o}{\PYGZhy{}} \PYG{l+m+mi}{1}\PYG{p}{)} \PYG{p}{:}
            \PYG{n}{result}\PYG{p}{[}\PYG{n}{i}\PYG{o}{+}\PYG{l+m+mi}{1}\PYG{p}{]} \PYG{o}{=} \PYG{n}{result}\PYG{p}{[}\PYG{n}{i}\PYG{p}{]} \PYG{o}{*} \PYG{n+nb+bp}{self}\PYG{o}{.}\PYG{n}{x}\PYG{p}{[}\PYG{n}{i} \PYG{o}{+} \PYG{l+m+mi}{1}\PYG{p}{]} \PYG{o}{*} \PYG{n}{t} \PYG{o}{/} \PYG{n+nb}{float}\PYG{p}{(}\PYG{n}{i}\PYG{o}{+}\PYG{l+m+mi}{1}\PYG{p}{)}
        \PYG{k}{return} \PYG{n}{result}
    \PYG{c+c1}{\PYGZsh{}}
    \PYG{c+c1}{\PYGZsh{} fun.f\PYGZus{}t}
    \PYG{k}{def} \PYG{n+nf}{f\PYGZus{}t}\PYG{p}{(}\PYG{n+nb+bp}{self}\PYG{p}{,} \PYG{n}{t}\PYG{p}{,} \PYG{n}{y}\PYG{p}{)} \PYG{p}{:}
        \PYG{n}{ny} \PYG{o}{=} \PYG{n}{y}\PYG{o}{.}\PYG{n}{size}
        \PYG{n}{type\PYGZus{}y}    \PYG{o}{=} \PYG{n+nb}{type}\PYG{p}{(}\PYG{n}{y}\PYG{p}{[}\PYG{l+m+mi}{0}\PYG{p}{]}\PYG{p}{)}
        \PYG{n}{result}    \PYG{o}{=} \PYG{n}{numpy}\PYG{o}{.}\PYG{n}{empty}\PYG{p}{(}\PYG{n}{ny}\PYG{p}{,} \PYG{n}{dtype}\PYG{o}{=}\PYG{n+nb}{type}\PYG{p}{(}\PYG{n}{y}\PYG{p}{[}\PYG{l+m+mi}{0}\PYG{p}{]}\PYG{p}{)} \PYG{p}{)}
        \PYG{n}{result}\PYG{p}{[}\PYG{l+m+mi}{0}\PYG{p}{]} \PYG{o}{=} \PYG{n}{type\PYGZus{}y}\PYG{p}{(} \PYG{l+m+mf}{0.0} \PYG{p}{)}
        \PYG{n}{prod}      \PYG{o}{=} \PYG{n}{type\PYGZus{}y}\PYG{p}{(} \PYG{n+nb+bp}{self}\PYG{o}{.}\PYG{n}{x}\PYG{p}{[}\PYG{l+m+mi}{0}\PYG{p}{]} \PYG{p}{)}
        \PYG{k}{for} \PYG{n}{i} \PYG{o+ow}{in} \PYG{n+nb}{range}\PYG{p}{(}\PYG{n}{ny} \PYG{o}{\PYGZhy{}} \PYG{l+m+mi}{1}\PYG{p}{)} \PYG{p}{:}
            \PYG{n}{prod}        \PYG{o}{=} \PYG{n}{prod} \PYG{o}{*} \PYG{n+nb+bp}{self}\PYG{o}{.}\PYG{n}{x}\PYG{p}{[}\PYG{n}{i}\PYG{o}{+}\PYG{l+m+mi}{1}\PYG{p}{]}
            \PYG{n}{result}\PYG{p}{[}\PYG{n}{i}\PYG{o}{+}\PYG{l+m+mi}{1}\PYG{p}{]} \PYG{o}{=} \PYG{n}{prod}
            \PYG{n}{prod}        \PYG{o}{=} \PYG{n}{prod} \PYG{o}{*} \PYG{n}{t} \PYG{o}{/} \PYG{p}{(}\PYG{n}{i} \PYG{o}{+} \PYG{l+m+mi}{1}\PYG{p}{)}
        \PYG{k}{return} \PYG{n}{result}
    \PYG{c+c1}{\PYGZsh{}}
    \PYG{c+c1}{\PYGZsh{} fun.f\PYGZus{}y}
    \PYG{k}{def} \PYG{n+nf}{f\PYGZus{}y}\PYG{p}{(}\PYG{n+nb+bp}{self}\PYG{p}{,} \PYG{n}{t}\PYG{p}{,} \PYG{n}{y}\PYG{p}{)} \PYG{p}{:}
        \PYG{n}{ny} \PYG{o}{=} \PYG{n}{y}\PYG{o}{.}\PYG{n}{size}
        \PYG{k}{return} \PYG{n}{numpy}\PYG{o}{.}\PYG{n}{zeros}\PYG{p}{(} \PYG{n}{ny}\PYG{p}{,} \PYG{n}{dtype}\PYG{o}{=}\PYG{n+nb}{type}\PYG{p}{(}\PYG{n}{y}\PYG{p}{[}\PYG{l+m+mi}{0}\PYG{p}{]}\PYG{p}{)} \PYG{p}{)}
\PYG{c+c1}{\PYGZsh{} \PYGZhy{}\PYGZhy{}\PYGZhy{}\PYGZhy{}\PYGZhy{}\PYGZhy{}\PYGZhy{}\PYGZhy{}\PYGZhy{}\PYGZhy{}\PYGZhy{}\PYGZhy{}\PYGZhy{}\PYGZhy{}\PYGZhy{}\PYGZhy{}\PYGZhy{}\PYGZhy{}\PYGZhy{}\PYGZhy{}\PYGZhy{}\PYGZhy{}\PYGZhy{}\PYGZhy{}\PYGZhy{}\PYGZhy{}\PYGZhy{}\PYGZhy{}\PYGZhy{}\PYGZhy{}\PYGZhy{}\PYGZhy{}\PYGZhy{}\PYGZhy{}\PYGZhy{}\PYGZhy{}\PYGZhy{}\PYGZhy{}\PYGZhy{}\PYGZhy{}\PYGZhy{}\PYGZhy{}\PYGZhy{}\PYGZhy{}\PYGZhy{}\PYGZhy{}\PYGZhy{}\PYGZhy{}\PYGZhy{}\PYGZhy{}\PYGZhy{}\PYGZhy{}\PYGZhy{}\PYGZhy{}\PYGZhy{}\PYGZhy{}\PYGZhy{}\PYGZhy{}\PYGZhy{}\PYGZhy{}\PYGZhy{}\PYGZhy{}\PYGZhy{}\PYGZhy{}\PYGZhy{}\PYGZhy{}\PYGZhy{}\PYGZhy{}\PYGZhy{}\PYGZhy{}\PYGZhy{}\PYGZhy{}\PYGZhy{}\PYGZhy{}\PYGZhy{}}
\PYG{c+c1}{\PYGZsh{} Check derivative claculations in first\PYGZus{}fun\PYGZus{}class and second\PYGZus{}fun\PYGZus{}class}
\PYG{k}{def} \PYG{n+nf}{test\PYGZus{}derivative}\PYG{p}{(}\PYG{p}{)} \PYG{p}{:}
    \PYG{n}{ok}       \PYG{o}{=} \PYG{k+kc}{True}
    \PYG{n}{nx}       \PYG{o}{=} \PYG{l+m+mi}{3}
    \PYG{n}{eps99}    \PYG{o}{=} \PYG{l+m+mf}{99.0} \PYG{o}{*} \PYG{n}{numpy}\PYG{o}{.}\PYG{n}{finfo}\PYG{p}{(}\PYG{n+nb}{float}\PYG{p}{)}\PYG{o}{.}\PYG{n}{eps}
    \PYG{n}{t\PYGZus{}start}  \PYG{o}{=} \PYG{l+m+mf}{0.0}
    \PYG{n}{t\PYGZus{}step}   \PYG{o}{=} \PYG{l+m+mf}{0.75}
    \PYG{n}{x}        \PYG{o}{=} \PYG{n}{numpy}\PYG{o}{.}\PYG{n}{array}\PYG{p}{(} \PYG{n}{nx} \PYG{o}{*} \PYG{p}{[} \PYG{l+m+mf}{1.0} \PYG{p}{]} \PYG{p}{)}
    \PYG{n}{y\PYGZus{}start}  \PYG{o}{=} \PYG{n}{numpy}\PYG{o}{.}\PYG{n}{array}\PYG{p}{(} \PYG{n}{nx} \PYG{o}{*} \PYG{p}{[} \PYG{l+m+mf}{0.0} \PYG{p}{]} \PYG{p}{)}
    \PYG{c+c1}{\PYGZsh{}}
    \PYG{n}{fun\PYGZus{}1}    \PYG{o}{=} \PYG{n}{first\PYGZus{}fun\PYGZus{}class}\PYG{p}{(}\PYG{n}{x}\PYG{p}{)}
    \PYG{n}{ok}       \PYG{o}{=} \PYG{n}{ok} \PYG{o+ow}{and} \PYG{n}{check\PYGZus{}rosen3\PYGZus{}step}\PYG{p}{(}\PYG{n}{fun\PYGZus{}1}\PYG{p}{,} \PYG{n}{t\PYGZus{}start}\PYG{p}{,} \PYG{n}{y\PYGZus{}start}\PYG{p}{,} \PYG{n}{t\PYGZus{}step}\PYG{p}{)}
    \PYG{c+c1}{\PYGZsh{}}
    \PYG{n}{fun\PYGZus{}2}    \PYG{o}{=} \PYG{n}{second\PYGZus{}fun\PYGZus{}class}\PYG{p}{(}\PYG{n}{x}\PYG{p}{)}
    \PYG{n}{ok}       \PYG{o}{=} \PYG{n}{ok} \PYG{o+ow}{and} \PYG{n}{check\PYGZus{}rosen3\PYGZus{}step}\PYG{p}{(}\PYG{n}{fun\PYGZus{}1}\PYG{p}{,} \PYG{n}{t\PYGZus{}start}\PYG{p}{,} \PYG{n}{y\PYGZus{}start}\PYG{p}{,} \PYG{n}{t\PYGZus{}step}\PYG{p}{)}
    \PYG{k}{return} \PYG{n}{ok}

\PYG{k}{def} \PYG{n+nf}{test\PYGZus{}case}\PYG{p}{(}\PYG{n}{case}\PYG{p}{)} \PYG{p}{:}
    \PYG{n}{ok}    \PYG{o}{=} \PYG{k+kc}{True}
    \PYG{n}{nx}    \PYG{o}{=} \PYG{l+m+mi}{3}
    \PYG{n}{eps99} \PYG{o}{=} \PYG{l+m+mf}{99.0} \PYG{o}{*} \PYG{n}{numpy}\PYG{o}{.}\PYG{n}{finfo}\PYG{p}{(}\PYG{n+nb}{float}\PYG{p}{)}\PYG{o}{.}\PYG{n}{eps}
    \PYG{c+c1}{\PYGZsh{}}
    \PYG{c+c1}{\PYGZsh{} independent variables for this g(x) = y(1, x)}
    \PYG{n}{x}  \PYG{o}{=} \PYG{n}{numpy}\PYG{o}{.}\PYG{n}{array}\PYG{p}{(} \PYG{n}{nx} \PYG{o}{*} \PYG{p}{[} \PYG{l+m+mf}{1.0} \PYG{p}{]} \PYG{p}{)}
    \PYG{n}{ax} \PYG{o}{=} \PYG{n}{cppad\PYGZus{}py}\PYG{o}{.}\PYG{n}{independent}\PYG{p}{(}\PYG{n}{x}\PYG{p}{)}
    \PYG{c+c1}{\PYGZsh{}}
    \PYG{k}{if} \PYG{n}{case} \PYG{o}{==} \PYG{l+m+mi}{1} \PYG{p}{:}
        \PYG{n}{fun} \PYG{o}{=} \PYG{n}{first\PYGZus{}fun\PYGZus{}class}\PYG{p}{(}\PYG{n}{ax}\PYG{p}{)}
    \PYG{k}{elif} \PYG{n}{case} \PYG{o}{==} \PYG{l+m+mi}{2} \PYG{p}{:}
        \PYG{n}{fun} \PYG{o}{=} \PYG{n}{second\PYGZus{}fun\PYGZus{}class}\PYG{p}{(}\PYG{n}{ax}\PYG{p}{)}
    \PYG{k}{else} \PYG{p}{:}
        \PYG{k}{assert}\PYG{p}{(}\PYG{k+kc}{False}\PYG{p}{)}
    \PYG{c+c1}{\PYGZsh{}}
    \PYG{c+c1}{\PYGZsh{} initiali value for the ODE}
    \PYG{n}{ay\PYGZus{}start} \PYG{o}{=}  \PYG{n}{numpy}\PYG{o}{.}\PYG{n}{array}\PYG{p}{(} \PYG{n}{nx} \PYG{o}{*} \PYG{p}{[} \PYG{n}{cppad\PYGZus{}py}\PYG{o}{.}\PYG{n}{a\PYGZus{}double}\PYG{p}{(}\PYG{l+m+mf}{0.0}\PYG{p}{)} \PYG{p}{]} \PYG{p}{)}
    \PYG{n}{t\PYGZus{}start}  \PYG{o}{=} \PYG{l+m+mf}{0.0}
    \PYG{n}{t\PYGZus{}step}   \PYG{o}{=} \PYG{l+m+mf}{0.75}
    \PYG{c+c1}{\PYGZsh{}}
    \PYG{c+c1}{\PYGZsh{} take one step}
    \PYG{n}{ay} \PYG{o}{=} \PYG{n}{rosen3\PYGZus{}step}\PYG{p}{(}\PYG{n}{fun}\PYG{p}{,} \PYG{n}{t\PYGZus{}start}\PYG{p}{,} \PYG{n}{ay\PYGZus{}start}\PYG{p}{,} \PYG{n}{t\PYGZus{}step}\PYG{p}{)}
    \PYG{c+c1}{\PYGZsh{}}
    \PYG{c+c1}{\PYGZsh{} g(x) = y(t\PYGZus{}step, x)}
    \PYG{n}{g} \PYG{o}{=} \PYG{n}{cppad\PYGZus{}py}\PYG{o}{.}\PYG{n}{d\PYGZus{}fun}\PYG{p}{(}\PYG{n}{ax}\PYG{p}{,} \PYG{n}{ay}\PYG{p}{)}
    \PYG{c+c1}{\PYGZsh{}}
    \PYG{c+c1}{\PYGZsh{} Check g\PYGZus{}i (x) = prod\PYGZus{}\PYGZob{}j=0\PYGZcb{}\PYGZca{}i x[j] t\PYGZca{}(i+1) / (i+1) !}
    \PYG{c+c1}{\PYGZsh{} 3th order method should be very accutate for functions}
    \PYG{c+c1}{\PYGZsh{} or order 3 or less in t.}
    \PYG{n}{x}  \PYG{o}{=} \PYG{n}{numpy}\PYG{o}{.}\PYG{n}{arange}\PYG{p}{(}\PYG{l+m+mi}{0}\PYG{p}{,} \PYG{n}{nx}\PYG{p}{)} \PYG{o}{+} \PYG{l+m+mf}{1.0}
    \PYG{n}{gx} \PYG{o}{=} \PYG{n}{g}\PYG{o}{.}\PYG{n}{forward}\PYG{p}{(}\PYG{l+m+mi}{0}\PYG{p}{,} \PYG{n}{x}\PYG{p}{)}
    \PYG{n}{prod} \PYG{o}{=} \PYG{l+m+mf}{1.0}
    \PYG{k}{for} \PYG{n}{i} \PYG{o+ow}{in} \PYG{n+nb}{range}\PYG{p}{(}\PYG{n}{nx}\PYG{p}{)} \PYG{p}{:}
        \PYG{n}{prod}      \PYG{o}{=} \PYG{n}{prod} \PYG{o}{*} \PYG{n}{x}\PYG{p}{[}\PYG{n}{i}\PYG{p}{]}
        \PYG{n}{check}     \PYG{o}{=} \PYG{n}{prod} \PYG{o}{*} \PYG{n}{numpy}\PYG{o}{.}\PYG{n}{power}\PYG{p}{(}\PYG{n}{t\PYGZus{}step}\PYG{p}{,} \PYG{n}{i}\PYG{o}{+}\PYG{l+m+mi}{1}\PYG{p}{)} \PYG{o}{/} \PYG{n}{factorial}\PYG{p}{(}\PYG{n}{i}\PYG{o}{+}\PYG{l+m+mi}{1}\PYG{p}{)}
        \PYG{n}{rel\PYGZus{}error} \PYG{o}{=} \PYG{n}{gx}\PYG{p}{[}\PYG{n}{i}\PYG{p}{]} \PYG{o}{/} \PYG{n}{check} \PYG{o}{\PYGZhy{}} \PYG{l+m+mf}{1.0}
        \PYG{n}{ok}       \PYG{o}{\PYGZam{}}\PYG{o}{=} \PYG{n+nb}{abs}\PYG{p}{(}\PYG{n}{rel\PYGZus{}error}\PYG{p}{)} \PYG{o}{\PYGZlt{}} \PYG{n}{eps99}
    \PYG{c+c1}{\PYGZsh{}}
    \PYG{c+c1}{\PYGZsh{} Check derivative of g\PYGZus{}i (x) w.r.t x\PYGZus{}j}
    \PYG{c+c1}{\PYGZsh{} is zero for i \PYGZlt{} j and  g\PYGZus{}i(x) / x\PYGZus{}j otherwise}
    \PYG{n}{J} \PYG{o}{=} \PYG{n}{g}\PYG{o}{.}\PYG{n}{jacobian}\PYG{p}{(}\PYG{n}{x}\PYG{p}{)}
    \PYG{k}{for} \PYG{n}{j} \PYG{o+ow}{in} \PYG{n+nb}{range}\PYG{p}{(}\PYG{n}{nx}\PYG{p}{)} \PYG{p}{:}
        \PYG{k}{for} \PYG{n}{i} \PYG{o+ow}{in} \PYG{n+nb}{range}\PYG{p}{(}\PYG{n}{nx}\PYG{p}{)} \PYG{p}{:}
            \PYG{k}{if} \PYG{n}{i} \PYG{o}{\PYGZlt{}} \PYG{n}{j} \PYG{p}{:}
                \PYG{n}{ok} \PYG{o}{\PYGZam{}}\PYG{o}{=} \PYG{n}{J}\PYG{p}{[}\PYG{n}{i}\PYG{p}{,}\PYG{n}{j}\PYG{p}{]} \PYG{o}{==} \PYG{l+m+mf}{0.0}
            \PYG{k}{else} \PYG{p}{:}
                \PYG{n}{check}     \PYG{o}{=} \PYG{n}{gx}\PYG{p}{[}\PYG{n}{i}\PYG{p}{]} \PYG{o}{/} \PYG{n}{x}\PYG{p}{[}\PYG{n}{j}\PYG{p}{]}
                \PYG{n}{rel\PYGZus{}error} \PYG{o}{=} \PYG{n}{J}\PYG{p}{[}\PYG{n}{i}\PYG{p}{,}\PYG{n}{j}\PYG{p}{]} \PYG{o}{/} \PYG{n}{check} \PYG{o}{\PYGZhy{}} \PYG{l+m+mf}{1.0}
                \PYG{n}{ok}       \PYG{o}{\PYGZam{}}\PYG{o}{=} \PYG{n+nb}{abs}\PYG{p}{(}\PYG{n}{rel\PYGZus{}error}\PYG{p}{)} \PYG{o}{\PYGZlt{}} \PYG{n}{eps99}

    \PYG{k}{return} \PYG{n}{ok}
\PYG{c+c1}{\PYGZsh{} \PYGZhy{}\PYGZhy{}\PYGZhy{}\PYGZhy{}\PYGZhy{}\PYGZhy{}\PYGZhy{}\PYGZhy{}\PYGZhy{}\PYGZhy{}\PYGZhy{}\PYGZhy{}\PYGZhy{}\PYGZhy{}\PYGZhy{}\PYGZhy{}\PYGZhy{}\PYGZhy{}\PYGZhy{}\PYGZhy{}\PYGZhy{}\PYGZhy{}\PYGZhy{}\PYGZhy{}\PYGZhy{}\PYGZhy{}\PYGZhy{}\PYGZhy{}\PYGZhy{}\PYGZhy{}\PYGZhy{}\PYGZhy{}\PYGZhy{}\PYGZhy{}\PYGZhy{}\PYGZhy{}\PYGZhy{}\PYGZhy{}\PYGZhy{}\PYGZhy{}\PYGZhy{}\PYGZhy{}\PYGZhy{}\PYGZhy{}\PYGZhy{}\PYGZhy{}\PYGZhy{}\PYGZhy{}\PYGZhy{}\PYGZhy{}\PYGZhy{}\PYGZhy{}\PYGZhy{}\PYGZhy{}\PYGZhy{}\PYGZhy{}\PYGZhy{}\PYGZhy{}\PYGZhy{}\PYGZhy{}\PYGZhy{}\PYGZhy{}\PYGZhy{}\PYGZhy{}\PYGZhy{}\PYGZhy{}\PYGZhy{}\PYGZhy{}\PYGZhy{}\PYGZhy{}\PYGZhy{}\PYGZhy{}\PYGZhy{}\PYGZhy{}\PYGZhy{}}
\PYG{k}{def} \PYG{n+nf}{rosen3\PYGZus{}step\PYGZus{}xam}\PYG{p}{(}\PYG{p}{)} \PYG{p}{:}
    \PYG{n}{ok} \PYG{o}{=} \PYG{k+kc}{True}
    \PYG{n}{ok} \PYG{o}{=} \PYG{n}{ok} \PYG{o+ow}{and} \PYG{n}{test\PYGZus{}derivative}\PYG{p}{(}\PYG{p}{)}
    \PYG{n}{ok} \PYG{o}{=} \PYG{n}{ok} \PYG{o+ow}{and} \PYG{n}{test\PYGZus{}case}\PYG{p}{(}\PYG{l+m+mi}{1}\PYG{p}{)}
    \PYG{n}{ok} \PYG{o}{=} \PYG{n}{ok} \PYG{o+ow}{and} \PYG{n}{test\PYGZus{}case}\PYG{p}{(}\PYG{l+m+mi}{2}\PYG{p}{)}
    \PYG{k}{return} \PYG{n}{ok}

\end{sphinxVerbatim}


\bigskip\hrule\bigskip


xsrst input file: \sphinxcode{\sphinxupquote{example/python/numeric/rosen3\_step\_xam.py}}

\index{example@\spxentry{example}}\ignorespaces 

\subparagraph{Example}
\label{\detokenize{xsrst/numeric_rosen3_step:example}}\label{\detokenize{xsrst/numeric_rosen3_step:numeric-rosen3-step-example}}\label{\detokenize{xsrst/numeric_rosen3_step:index-15}}
{\hyperref[\detokenize{xsrst/numeric_rosen3_step_xam_py:id1}]{\sphinxcrossref{\DUrole{std,std-ref}{numeric\_rosen3\_step\_xam\_py}}}}

\index{source@\spxentry{source}}\index{code@\spxentry{code}}\ignorespaces 

\subparagraph{Source Code}
\label{\detokenize{xsrst/numeric_rosen3_step:source-code}}\label{\detokenize{xsrst/numeric_rosen3_step:numeric-rosen3-step-source-code}}\label{\detokenize{xsrst/numeric_rosen3_step:index-16}}
\index{rosen3\_step@\spxentry{rosen3\_step}}\ignorespaces 

\subparagraph{rosen3\_step}
\label{\detokenize{xsrst/numeric_rosen3_step:rosen3-step}}\label{\detokenize{xsrst/numeric_rosen3_step:numeric-rosen3-step-source-code-rosen3-step}}\label{\detokenize{xsrst/numeric_rosen3_step:index-17}}
\begin{sphinxVerbatim}[commandchars=\\\{\}]
\PYG{k+kn}{import} \PYG{n+nn}{numpy}
\PYG{k+kn}{from} \PYG{n+nn}{simple\PYGZus{}inv} \PYG{k+kn}{import} \PYG{n}{simple\PYGZus{}inv}
\PYG{k+kn}{import} \PYG{n+nn}{cppad\PYGZus{}py}
\PYG{k}{def} \PYG{n+nf}{rosen3\PYGZus{}step}\PYG{p}{(}\PYG{n}{fun}\PYG{p}{,} \PYG{n}{ti}\PYG{p}{,} \PYG{n}{yi}\PYG{p}{,} \PYG{n}{h}\PYG{p}{)} \PYG{p}{:}
    \PYG{c+c1}{\PYGZsh{}}
    \PYG{n}{r1\PYGZus{}2}     \PYG{o}{=} \PYG{l+m+mf}{1.0}   \PYG{o}{/} \PYG{l+m+mf}{2.0}
    \PYG{n}{r3\PYGZus{}2}     \PYG{o}{=} \PYG{l+m+mf}{3.0}   \PYG{o}{/} \PYG{l+m+mf}{2.0}
    \PYG{n}{r3\PYGZus{}5}     \PYG{o}{=} \PYG{l+m+mf}{3.0}   \PYG{o}{/} \PYG{l+m+mf}{5.0}
    \PYG{n}{r24\PYGZus{}25}   \PYG{o}{=} \PYG{l+m+mf}{24.0}  \PYG{o}{/} \PYG{l+m+mf}{25.0}
    \PYG{n}{r3\PYGZus{}25}    \PYG{o}{=} \PYG{l+m+mf}{3.0}   \PYG{o}{/} \PYG{l+m+mf}{25.0}
    \PYG{n}{r121\PYGZus{}50}  \PYG{o}{=} \PYG{l+m+mf}{121.0} \PYG{o}{/} \PYG{l+m+mf}{50.0}
    \PYG{n}{r186\PYGZus{}25}  \PYG{o}{=} \PYG{l+m+mf}{186.0} \PYG{o}{/} \PYG{l+m+mf}{25.0}
    \PYG{n}{r6\PYGZus{}5}     \PYG{o}{=} \PYG{l+m+mf}{6.0}   \PYG{o}{/} \PYG{l+m+mf}{5.0}
    \PYG{n}{r97\PYGZus{}108}  \PYG{o}{=} \PYG{l+m+mf}{97.0}  \PYG{o}{/} \PYG{l+m+mf}{108.0}
    \PYG{n}{r11\PYGZus{}72}   \PYG{o}{=} \PYG{l+m+mf}{11.0}  \PYG{o}{/} \PYG{l+m+mf}{72.0}
    \PYG{n}{r25\PYGZus{}216}  \PYG{o}{=} \PYG{l+m+mf}{25.0}  \PYG{o}{/} \PYG{l+m+mf}{216.0}
    \PYG{c+c1}{\PYGZsh{}}
    \PYG{c+c1}{\PYGZsh{} Einv}
    \PYG{n}{ny}   \PYG{o}{=} \PYG{n}{yi}\PYG{o}{.}\PYG{n}{size}
    \PYG{n}{I}    \PYG{o}{=} \PYG{n}{numpy}\PYG{o}{.}\PYG{n}{identity}\PYG{p}{(}\PYG{n}{ny}\PYG{p}{,} \PYG{n}{dtype}\PYG{o}{=}\PYG{n+nb}{float}\PYG{p}{)}
    \PYG{n}{E}    \PYG{o}{=} \PYG{n}{I} \PYG{o}{\PYGZhy{}} \PYG{n}{r1\PYGZus{}2} \PYG{o}{*} \PYG{n}{h} \PYG{o}{*} \PYG{n}{fun}\PYG{o}{.}\PYG{n}{f\PYGZus{}y}\PYG{p}{(}\PYG{n}{ti}\PYG{p}{,} \PYG{n}{yi}\PYG{p}{)}
    \PYG{n}{Einv} \PYG{o}{=} \PYG{n}{simple\PYGZus{}inv}\PYG{p}{(}\PYG{n}{E}\PYG{p}{)}
    \PYG{c+c1}{\PYGZsh{}}
    \PYG{c+c1}{\PYGZsh{} f\PYGZus{}t}
    \PYG{n}{f\PYGZus{}t}  \PYG{o}{=} \PYG{n}{fun}\PYG{o}{.}\PYG{n}{f\PYGZus{}t}\PYG{p}{(}\PYG{n}{ti}\PYG{p}{,} \PYG{n}{yi}\PYG{p}{)}
    \PYG{c+c1}{\PYGZsh{}}
    \PYG{c+c1}{\PYGZsh{} k1}
    \PYG{n}{k1}   \PYG{o}{=} \PYG{n}{fun}\PYG{o}{.}\PYG{n}{f}\PYG{p}{(}\PYG{n}{ti}\PYG{p}{,} \PYG{n}{yi}\PYG{p}{)}
    \PYG{n}{k1}  \PYG{o}{+}\PYG{o}{=} \PYG{n}{r1\PYGZus{}2} \PYG{o}{*} \PYG{n}{h} \PYG{o}{*} \PYG{n}{f\PYGZus{}t}
    \PYG{n}{k1}   \PYG{o}{=} \PYG{n}{numpy}\PYG{o}{.}\PYG{n}{matmul}\PYG{p}{(}\PYG{n}{Einv} \PYG{p}{,} \PYG{n}{k1}\PYG{p}{)}
    \PYG{c+c1}{\PYGZsh{}}
    \PYG{c+c1}{\PYGZsh{} k2}
    \PYG{n}{k2}   \PYG{o}{=} \PYG{n}{fun}\PYG{o}{.}\PYG{n}{f}\PYG{p}{(}\PYG{n}{ti} \PYG{o}{+} \PYG{n}{h}\PYG{p}{,} \PYG{n}{yi} \PYG{o}{+} \PYG{n}{h} \PYG{o}{*} \PYG{n}{k1}\PYG{p}{)}
    \PYG{n}{k2}  \PYG{o}{\PYGZhy{}}\PYG{o}{=} \PYG{n}{r3\PYGZus{}2} \PYG{o}{*} \PYG{n}{h} \PYG{o}{*} \PYG{n}{f\PYGZus{}t}
    \PYG{n}{k2}  \PYG{o}{\PYGZhy{}}\PYG{o}{=} \PYG{l+m+mf}{4.0} \PYG{o}{*} \PYG{n}{k1}
    \PYG{n}{k2}   \PYG{o}{=} \PYG{n}{numpy}\PYG{o}{.}\PYG{n}{matmul}\PYG{p}{(}\PYG{n}{Einv} \PYG{p}{,} \PYG{n}{k2}\PYG{p}{)}
    \PYG{c+c1}{\PYGZsh{}}
    \PYG{c+c1}{\PYGZsh{} k3}
    \PYG{n}{t}    \PYG{o}{=} \PYG{n}{ti} \PYG{o}{+} \PYG{n}{r3\PYGZus{}5} \PYG{o}{*} \PYG{n}{h}
    \PYG{n}{y}    \PYG{o}{=} \PYG{n}{yi} \PYG{o}{+} \PYG{n}{r24\PYGZus{}25} \PYG{o}{*} \PYG{n}{h} \PYG{o}{*} \PYG{n}{k1} \PYG{o}{+} \PYG{n}{r3\PYGZus{}25} \PYG{o}{*} \PYG{n}{h} \PYG{o}{*} \PYG{n}{k2}
    \PYG{n}{k3}   \PYG{o}{=} \PYG{n}{fun}\PYG{o}{.}\PYG{n}{f}\PYG{p}{(}\PYG{n}{t}\PYG{p}{,} \PYG{n}{y}\PYG{p}{)}
    \PYG{n}{k3}  \PYG{o}{+}\PYG{o}{=} \PYG{n}{r121\PYGZus{}50} \PYG{o}{*} \PYG{n}{h} \PYG{o}{*} \PYG{n}{f\PYGZus{}t}
    \PYG{n}{k3}  \PYG{o}{+}\PYG{o}{=} \PYG{n}{r186\PYGZus{}25} \PYG{o}{*} \PYG{n}{k1} \PYG{o}{+} \PYG{n}{r6\PYGZus{}5} \PYG{o}{*} \PYG{n}{k2}
    \PYG{n}{k3}   \PYG{o}{=} \PYG{n}{numpy}\PYG{o}{.}\PYG{n}{matmul}\PYG{p}{(}\PYG{n}{Einv} \PYG{p}{,} \PYG{n}{k3}\PYG{p}{)}
    \PYG{c+c1}{\PYGZsh{}}
    \PYG{n}{yf}   \PYG{o}{=} \PYG{n}{yi} \PYG{o}{+} \PYG{n}{h} \PYG{o}{*} \PYG{p}{(} \PYG{n}{r97\PYGZus{}108} \PYG{o}{*} \PYG{n}{k1} \PYG{o}{+} \PYG{n}{r11\PYGZus{}72} \PYG{o}{*} \PYG{n}{k2} \PYG{o}{+} \PYG{n}{r25\PYGZus{}216} \PYG{o}{*} \PYG{n}{k3}\PYG{p}{)}
    \PYG{c+c1}{\PYGZsh{}}
    \PYG{k}{return} \PYG{n}{yf}
\end{sphinxVerbatim}

\index{check\_rosen3\_step@\spxentry{check\_rosen3\_step}}\ignorespaces 

\subparagraph{check\_rosen3\_step}
\label{\detokenize{xsrst/numeric_rosen3_step:check-rosen3-step}}\label{\detokenize{xsrst/numeric_rosen3_step:numeric-rosen3-step-source-code-check-rosen3-step}}\label{\detokenize{xsrst/numeric_rosen3_step:index-18}}
\begin{sphinxVerbatim}[commandchars=\\\{\}]
\PYG{k}{def} \PYG{n+nf}{check\PYGZus{}rosen3\PYGZus{}step}\PYG{p}{(}\PYG{n}{fun}\PYG{p}{,} \PYG{n}{ti}\PYG{p}{,} \PYG{n}{yi}\PYG{p}{,} \PYG{n}{h}\PYG{p}{)} \PYG{p}{:}
    \PYG{n}{ok}    \PYG{o}{=} \PYG{k+kc}{True}
    \PYG{n}{eps99} \PYG{o}{=} \PYG{l+m+mf}{99.0} \PYG{o}{*} \PYG{n}{numpy}\PYG{o}{.}\PYG{n}{finfo}\PYG{p}{(}\PYG{n+nb}{float}\PYG{p}{)}\PYG{o}{.}\PYG{n}{eps}
    \PYG{c+c1}{\PYGZsh{}}
    \PYG{n}{ny}    \PYG{o}{=} \PYG{n}{yi}\PYG{o}{.}\PYG{n}{size}
    \PYG{n}{both}  \PYG{o}{=} \PYG{n}{numpy}\PYG{o}{.}\PYG{n}{concatenate}\PYG{p}{(} \PYG{p}{(}\PYG{n}{yi}\PYG{p}{,} \PYG{p}{[}\PYG{n}{ti}\PYG{p}{]}\PYG{p}{)} \PYG{p}{)}
    \PYG{n}{aboth} \PYG{o}{=} \PYG{n}{cppad\PYGZus{}py}\PYG{o}{.}\PYG{n}{independent}\PYG{p}{(}\PYG{n}{both}\PYG{p}{)}
    \PYG{n}{az}    \PYG{o}{=} \PYG{n}{fun}\PYG{o}{.}\PYG{n}{f}\PYG{p}{(}\PYG{n}{aboth}\PYG{p}{[}\PYG{n}{ny}\PYG{p}{]}\PYG{p}{,} \PYG{n}{aboth}\PYG{p}{[}\PYG{l+m+mi}{0}\PYG{p}{:}\PYG{n}{ny}\PYG{p}{]} \PYG{p}{)}
    \PYG{n}{fun\PYGZus{}d} \PYG{o}{=} \PYG{n}{cppad\PYGZus{}py}\PYG{o}{.}\PYG{n}{d\PYGZus{}fun}\PYG{p}{(}\PYG{n}{aboth}\PYG{p}{,} \PYG{n}{az}\PYG{p}{)}
    \PYG{c+c1}{\PYGZsh{}}
    \PYG{n}{J}     \PYG{o}{=} \PYG{n}{fun\PYGZus{}d}\PYG{o}{.}\PYG{n}{jacobian}\PYG{p}{(}\PYG{n}{both}\PYG{p}{)}
    \PYG{c+c1}{\PYGZsh{}}
    \PYG{c+c1}{\PYGZsh{} check fun.f\PYGZus{}t(t, y)}
    \PYG{n}{fun\PYGZus{}t}   \PYG{o}{=} \PYG{n}{fun}\PYG{o}{.}\PYG{n}{f\PYGZus{}t}\PYG{p}{(}\PYG{n}{ti}\PYG{p}{,} \PYG{n}{yi}\PYG{p}{)}
    \PYG{n}{printed} \PYG{o}{=} \PYG{k+kc}{False}
    \PYG{k}{for} \PYG{n}{i} \PYG{o+ow}{in} \PYG{n+nb}{range}\PYG{p}{(}\PYG{n}{ny}\PYG{p}{)} \PYG{p}{:}
        \PYG{k}{if} \PYG{n}{J}\PYG{p}{[}\PYG{n}{i}\PYG{p}{,} \PYG{n}{ny}\PYG{p}{]} \PYG{o}{==} \PYG{l+m+mf}{0.0} \PYG{p}{:}
            \PYG{n}{ok} \PYG{o}{=} \PYG{n}{ok} \PYG{o+ow}{and} \PYG{n}{fun\PYGZus{}t}\PYG{p}{[}\PYG{n}{i}\PYG{p}{]} \PYG{o}{==} \PYG{l+m+mf}{0.0}
        \PYG{k}{else} \PYG{p}{:}
            \PYG{n}{rel\PYGZus{}error} \PYG{o}{=} \PYG{n}{fun\PYGZus{}t}\PYG{p}{[}\PYG{n}{i}\PYG{p}{]} \PYG{o}{/} \PYG{n}{J}\PYG{p}{[}\PYG{n}{i}\PYG{p}{,}\PYG{n}{ny}\PYG{p}{]} \PYG{o}{\PYGZhy{}} \PYG{l+m+mf}{1.0}
            \PYG{k}{if} \PYG{n+nb}{abs}\PYG{p}{(}\PYG{n}{rel\PYGZus{}error}\PYG{p}{)} \PYG{o}{\PYGZgt{}} \PYG{n}{eps99} \PYG{p}{:}
                \PYG{n}{ok} \PYG{o}{=} \PYG{k+kc}{False}
                \PYG{k}{if} \PYG{o+ow}{not} \PYG{n}{printed} \PYG{p}{:}
                    \PYG{n+nb}{print}\PYG{p}{(}\PYG{l+s+s1}{\PYGZsq{}}\PYG{l+s+s1}{check\PYGZus{}rosen3\PYGZus{}step: fun.f\PYGZus{}t check failed}\PYG{l+s+s1}{\PYGZsq{}}\PYG{p}{)}
                    \PYG{n}{printed} \PYG{o}{=} \PYG{k+kc}{True}
                \PYG{n+nb}{print}\PYG{p}{(}\PYG{l+s+s1}{\PYGZsq{}}\PYG{l+s+s1}{fun\PYGZus{}t[}\PYG{l+s+s1}{\PYGZsq{}}\PYG{p}{,} \PYG{n}{i}\PYG{p}{,} \PYG{l+s+s1}{\PYGZsq{}}\PYG{l+s+s1}{] = }\PYG{l+s+s1}{\PYGZsq{}}\PYG{p}{,} \PYG{n}{fun\PYGZus{}t}\PYG{p}{[}\PYG{n}{i}\PYG{p}{]}\PYG{p}{,} \PYG{l+s+s1}{\PYGZsq{}}\PYG{l+s+s1}{, check = }\PYG{l+s+s1}{\PYGZsq{}}\PYG{p}{,} \PYG{n}{J}\PYG{p}{[}\PYG{n}{i}\PYG{p}{,}\PYG{n}{ny}\PYG{p}{]}\PYG{p}{)}
    \PYG{c+c1}{\PYGZsh{}}
    \PYG{c+c1}{\PYGZsh{} check fun.f\PYGZus{}y(t, y)}
    \PYG{n}{fun\PYGZus{}y} \PYG{o}{=} \PYG{n}{fun}\PYG{o}{.}\PYG{n}{f\PYGZus{}y}\PYG{p}{(}\PYG{n}{ti}\PYG{p}{,} \PYG{n}{yi}\PYG{p}{)}
    \PYG{n}{printed} \PYG{o}{=} \PYG{k+kc}{False}
    \PYG{k}{for} \PYG{n}{i} \PYG{o+ow}{in} \PYG{n+nb}{range}\PYG{p}{(}\PYG{n}{ny}\PYG{p}{)} \PYG{p}{:}
        \PYG{k}{for} \PYG{n}{j} \PYG{o+ow}{in} \PYG{n+nb}{range}\PYG{p}{(}\PYG{n}{ny}\PYG{p}{)} \PYG{p}{:}
            \PYG{k}{if} \PYG{n}{J}\PYG{p}{[}\PYG{n}{i}\PYG{p}{,} \PYG{n}{j}\PYG{p}{]} \PYG{o}{==} \PYG{l+m+mf}{0.0} \PYG{p}{:}
                \PYG{n}{ok\PYGZus{}ij} \PYG{o}{=}\PYG{n}{fun\PYGZus{}y}\PYG{p}{[}\PYG{n}{i}\PYG{p}{,} \PYG{n}{j}\PYG{p}{]} \PYG{o}{==} \PYG{l+m+mf}{0.0}
            \PYG{k}{else} \PYG{p}{:}
                \PYG{n}{rel\PYGZus{}error} \PYG{o}{=} \PYG{n}{fun\PYGZus{}y}\PYG{p}{[}\PYG{n}{i}\PYG{p}{,} \PYG{n}{j}\PYG{p}{]} \PYG{o}{/} \PYG{n}{J}\PYG{p}{[}\PYG{n}{i}\PYG{p}{,} \PYG{n}{j}\PYG{p}{]} \PYG{o}{\PYGZhy{}} \PYG{l+m+mf}{1.0}
                \PYG{n}{ok\PYGZus{}ij} \PYG{o}{=} \PYG{n+nb}{abs}\PYG{p}{(} \PYG{n}{rel\PYGZus{}error} \PYG{p}{)} \PYG{o}{\PYGZlt{}} \PYG{n}{eps99}
            \PYG{k}{if} \PYG{p}{(}\PYG{o+ow}{not} \PYG{n}{ok\PYGZus{}ij}\PYG{p}{)} \PYG{o+ow}{and} \PYG{p}{(}\PYG{o+ow}{not} \PYG{n}{printed}\PYG{p}{)} \PYG{p}{:}
                \PYG{n+nb}{print}\PYG{p}{(}\PYG{l+s+s1}{\PYGZsq{}}\PYG{l+s+s1}{check\PYGZus{}rosen3\PYGZus{}step: fun.f\PYGZus{}y check failed}\PYG{l+s+s1}{\PYGZsq{}}\PYG{p}{)}
                \PYG{n}{printed} \PYG{o}{=} \PYG{k+kc}{True}
            \PYG{k}{if} \PYG{o+ow}{not} \PYG{n}{ok\PYGZus{}ij} \PYG{p}{:}
                \PYG{n}{msg}  \PYG{o}{=} \PYG{l+s+s1}{\PYGZsq{}}\PYG{l+s+s1}{fun\PYGZus{}y[}\PYG{l+s+s1}{\PYGZsq{}} \PYG{o}{+} \PYG{n+nb}{str}\PYG{p}{(}\PYG{n}{i}\PYG{p}{)} \PYG{o}{+} \PYG{l+s+s1}{\PYGZsq{}}\PYG{l+s+s1}{,}\PYG{l+s+s1}{\PYGZsq{}} \PYG{o}{+} \PYG{n+nb}{str}\PYG{p}{(}\PYG{n}{j}\PYG{p}{)} \PYG{o}{+} \PYG{l+s+s1}{\PYGZsq{}}\PYG{l+s+s1}{] =}\PYG{l+s+s1}{\PYGZsq{}}
                \PYG{n}{msg} \PYG{o}{+}\PYG{o}{=} \PYG{n+nb}{str}\PYG{p}{(} \PYG{n}{fun\PYGZus{}y}\PYG{p}{[}\PYG{n}{i}\PYG{p}{,} \PYG{n}{j}\PYG{p}{]} \PYG{p}{)}
                \PYG{n}{msg} \PYG{o}{+}\PYG{o}{=} \PYG{l+s+s1}{\PYGZsq{}}\PYG{l+s+s1}{, check = }\PYG{l+s+s1}{\PYGZsq{}} \PYG{o}{+} \PYG{n+nb}{str}\PYG{p}{(} \PYG{n}{J}\PYG{p}{[}\PYG{n}{i}\PYG{p}{,} \PYG{n}{j}\PYG{p}{]} \PYG{p}{)}
                \PYG{n+nb}{print}\PYG{p}{(}\PYG{n}{msg}\PYG{p}{)}
    \PYG{c+c1}{\PYGZsh{}}
    \PYG{k}{return} \PYG{n}{ok}
\end{sphinxVerbatim}


\bigskip\hrule\bigskip


xsrst input file: \sphinxcode{\sphinxupquote{example/python/numeric/rosen3\_step.py}}


\subparagraph{numeric\_ode\_multi\_step}
\label{\detokenize{xsrst/numeric_ode_multi_step:numeric-ode-multi-step}}\label{\detokenize{xsrst/numeric_ode_multi_step::doc}}
\index{numeric\_ode\_multi\_step@\spxentry{numeric\_ode\_multi\_step}}\index{multiple@\spxentry{multiple}}\index{ode@\spxentry{ode}}\index{steps@\spxentry{steps}}\ignorespaces 

\subparagraph{2.4 Multiple Ode Steps}
\label{\detokenize{xsrst/numeric_ode_multi_step:multiple-ode-steps}}\label{\detokenize{xsrst/numeric_ode_multi_step:id1}}\begin{itemize}
\item {} 
{\hyperref[\detokenize{xsrst/numeric_ode_multi_step:numeric-ode-multi-step-syntax}]{\sphinxcrossref{\DUrole{std,std-ref}{Syntax}}}}

\item {} 
{\hyperref[\detokenize{xsrst/numeric_ode_multi_step:numeric-ode-multi-step-purpose}]{\sphinxcrossref{\DUrole{std,std-ref}{Purpose}}}}

\item {} \begin{description}
\item[{{\hyperref[\detokenize{xsrst/numeric_ode_multi_step:numeric-ode-multi-step-one-step}]{\sphinxcrossref{\DUrole{std,std-ref}{one\_step}}}}}] \leavevmode\begin{itemize}
\item {} 
{\hyperref[\detokenize{xsrst/numeric_ode_multi_step:numeric-ode-multi-step-one-step-t0}]{\sphinxcrossref{\DUrole{std,std-ref}{t0}}}}

\item {} 
{\hyperref[\detokenize{xsrst/numeric_ode_multi_step:numeric-ode-multi-step-one-step-y0}]{\sphinxcrossref{\DUrole{std,std-ref}{y0}}}}

\item {} 
{\hyperref[\detokenize{xsrst/numeric_ode_multi_step:numeric-ode-multi-step-one-step-t-step}]{\sphinxcrossref{\DUrole{std,std-ref}{t\_step}}}}

\item {} 
{\hyperref[\detokenize{xsrst/numeric_ode_multi_step:numeric-ode-multi-step-one-step-y1}]{\sphinxcrossref{\DUrole{std,std-ref}{y1}}}}

\end{itemize}

\end{description}

\item {} \begin{description}
\item[{{\hyperref[\detokenize{xsrst/numeric_ode_multi_step:numeric-ode-multi-step-fun}]{\sphinxcrossref{\DUrole{std,std-ref}{fun}}}}}] \leavevmode\begin{itemize}
\item {} 
{\hyperref[\detokenize{xsrst/numeric_ode_multi_step:numeric-ode-multi-step-fun-fun-set-t-all-index-index}]{\sphinxcrossref{\DUrole{std,std-ref}{fun.set\_t\_all\_index(index)}}}}

\item {} 
{\hyperref[\detokenize{xsrst/numeric_ode_multi_step:numeric-ode-multi-step-fun-fun-f-t-y}]{\sphinxcrossref{\DUrole{std,std-ref}{fun.f(t, y)}}}}

\item {} 
{\hyperref[\detokenize{xsrst/numeric_ode_multi_step:numeric-ode-multi-step-fun-one-step}]{\sphinxcrossref{\DUrole{std,std-ref}{one\_step}}}}

\end{itemize}

\end{description}

\item {} 
{\hyperref[\detokenize{xsrst/numeric_ode_multi_step:numeric-ode-multi-step-t-all}]{\sphinxcrossref{\DUrole{std,std-ref}{t\_all}}}}

\item {} 
{\hyperref[\detokenize{xsrst/numeric_ode_multi_step:numeric-ode-multi-step-y-init}]{\sphinxcrossref{\DUrole{std,std-ref}{y\_init}}}}

\item {} 
{\hyperref[\detokenize{xsrst/numeric_ode_multi_step:numeric-ode-multi-step-y-all}]{\sphinxcrossref{\DUrole{std,std-ref}{y\_all}}}}

\item {} 
{\hyperref[\detokenize{xsrst/numeric_ode_multi_step:numeric-ode-multi-step-example}]{\sphinxcrossref{\DUrole{std,std-ref}{Example}}}}

\item {} 
{\hyperref[\detokenize{xsrst/numeric_ode_multi_step:numeric-ode-multi-step-source-code}]{\sphinxcrossref{\DUrole{std,std-ref}{Source Code}}}}

\end{itemize}

\index{syntax@\spxentry{syntax}}\ignorespaces 

\subparagraph{Syntax}
\label{\detokenize{xsrst/numeric_ode_multi_step:syntax}}\label{\detokenize{xsrst/numeric_ode_multi_step:numeric-ode-multi-step-syntax}}\label{\detokenize{xsrst/numeric_ode_multi_step:index-1}}
\sphinxstyleemphasis{y\_all} =  \sphinxcode{\sphinxupquote{ode\_multi\_step}} ( \sphinxstyleemphasis{one\_step} , \sphinxstyleemphasis{f} , \sphinxstyleemphasis{t\_all} , \sphinxstyleemphasis{y\_init} )

\index{purpose@\spxentry{purpose}}\ignorespaces 

\subparagraph{Purpose}
\label{\detokenize{xsrst/numeric_ode_multi_step:purpose}}\label{\detokenize{xsrst/numeric_ode_multi_step:numeric-ode-multi-step-purpose}}\label{\detokenize{xsrst/numeric_ode_multi_step:index-2}}
The routine can be used with \sphinxcode{\sphinxupquote{ad\_double}} to solve an initial
value ODE
\begin{equation*}
\begin{split}y^{(1)} (t) = f( t , y )\end{split}
\end{equation*}
\index{one\_step@\spxentry{one\_step}}\ignorespaces 

\subparagraph{one\_step}
\label{\detokenize{xsrst/numeric_ode_multi_step:one-step}}\label{\detokenize{xsrst/numeric_ode_multi_step:numeric-ode-multi-step-one-step}}\label{\detokenize{xsrst/numeric_ode_multi_step:index-3}}
This routine executes one step of an ODE approximation method with the
following syntax:

\begin{DUlineblock}{0em}
\item[]         \sphinxstyleemphasis{y1} = \sphinxstyleemphasis{one\_step} ( \sphinxstyleemphasis{f} , \sphinxstyleemphasis{t0} , \sphinxstyleemphasis{y0} , \sphinxstyleemphasis{t\_step} )
\end{DUlineblock}

The elements of \sphinxstyleemphasis{y0} and the scalars above can be
\sphinxcode{\sphinxupquote{float}} or \sphinxcode{\sphinxupquote{a\_double}} .

\index{t0@\spxentry{t0}}\ignorespaces 

\subparagraph{t0}
\label{\detokenize{xsrst/numeric_ode_multi_step:t0}}\label{\detokenize{xsrst/numeric_ode_multi_step:numeric-ode-multi-step-one-step-t0}}\label{\detokenize{xsrst/numeric_ode_multi_step:index-4}}
is the initial time value for the step.

\index{y0@\spxentry{y0}}\ignorespaces 

\subparagraph{y0}
\label{\detokenize{xsrst/numeric_ode_multi_step:y0}}\label{\detokenize{xsrst/numeric_ode_multi_step:numeric-ode-multi-step-one-step-y0}}\label{\detokenize{xsrst/numeric_ode_multi_step:index-5}}
is the initial value of \(y(t)\) for the step.

\index{t\_step@\spxentry{t\_step}}\ignorespaces 

\subparagraph{t\_step}
\label{\detokenize{xsrst/numeric_ode_multi_step:t-step}}\label{\detokenize{xsrst/numeric_ode_multi_step:numeric-ode-multi-step-one-step-t-step}}\label{\detokenize{xsrst/numeric_ode_multi_step:index-6}}
is the is the size of the step (in time).

\index{y1@\spxentry{y1}}\ignorespaces 

\subparagraph{y1}
\label{\detokenize{xsrst/numeric_ode_multi_step:y1}}\label{\detokenize{xsrst/numeric_ode_multi_step:numeric-ode-multi-step-one-step-y1}}\label{\detokenize{xsrst/numeric_ode_multi_step:index-7}}
is the approximation for the ODE solution at time \sphinxstyleemphasis{t0} + \sphinxstyleemphasis{t\_step} .

\index{fun@\spxentry{fun}}\ignorespaces 

\subparagraph{fun}
\label{\detokenize{xsrst/numeric_ode_multi_step:fun}}\label{\detokenize{xsrst/numeric_ode_multi_step:numeric-ode-multi-step-fun}}\label{\detokenize{xsrst/numeric_ode_multi_step:index-8}}
\index{fun.set\_t\_all\_index(index)@\spxentry{fun.set\_t\_all\_index(index)}}\ignorespaces 

\subparagraph{fun.set\_t\_all\_index(index)}
\label{\detokenize{xsrst/numeric_ode_multi_step:fun-set-t-all-index-index}}\label{\detokenize{xsrst/numeric_ode_multi_step:numeric-ode-multi-step-fun-fun-set-t-all-index-index}}\label{\detokenize{xsrst/numeric_ode_multi_step:index-9}}
Often, we use interpolation with knots to define \(f(t, y)\).
It is one of the subtitle issues of AD that even though values are the
same, derivatives might not be the same; e.g.,
piecewise linear interpolation.
We must break the integration of the ODE at each of the knot
so we can use a method that assumes \(f(t, y)\) is smooth.
Also so that AD can be used to compute derivatives of our solutions.
The function \sphinxcode{\sphinxupquote{set\_t\_all\_index}}
informs \sphinxstyleemphasis{fun} that we are currently integrating the time interval

\begin{DUlineblock}{0em}
\item[]         {[} \sphinxstyleemphasis{t\_all} {[} \sphinxstyleemphasis{index} {]} , \sphinxstyleemphasis{t\_all} {[} \sphinxstyleemphasis{index} +1{]} {]}
\end{DUlineblock}

so that it know which smooth function to represent
even if \(t\) is at a knot and it matters if it is the interval
to the left or right of the knot.
The function \sphinxcode{\sphinxupquote{set\_t\_all\_index}} is called at the start
of each integration interval and before any of the other
\sphinxstyleemphasis{fun} member functions.

\index{fun.f(t@\spxentry{fun.f(t}}\index{y)@\spxentry{y)}}\ignorespaces 

\subparagraph{fun.f(t, y)}
\label{\detokenize{xsrst/numeric_ode_multi_step:fun-f-t-y}}\label{\detokenize{xsrst/numeric_ode_multi_step:numeric-ode-multi-step-fun-fun-f-t-y}}\label{\detokenize{xsrst/numeric_ode_multi_step:index-10}}
This call evaluates the function that defines the ODE
\begin{equation*}
\begin{split}y^{(1)} (t) = f [ t , y(t) ]\end{split}
\end{equation*}
\index{one\_step@\spxentry{one\_step}}\ignorespaces 

\subparagraph{one\_step}
\label{\detokenize{xsrst/numeric_ode_multi_step:numeric-ode-multi-step-fun-one-step}}\label{\detokenize{xsrst/numeric_ode_multi_step:index-11}}\label{\detokenize{xsrst/numeric_ode_multi_step:id2}}
The routine \sphinxstyleemphasis{one\_step} may put extra requirements on \sphinxstyleemphasis{fun} .

\index{t\_all@\spxentry{t\_all}}\ignorespaces 

\subparagraph{t\_all}
\label{\detokenize{xsrst/numeric_ode_multi_step:t-all}}\label{\detokenize{xsrst/numeric_ode_multi_step:numeric-ode-multi-step-t-all}}\label{\detokenize{xsrst/numeric_ode_multi_step:index-12}}
This is a numpy vector of time values at which the solution is calculated.
The type of its elements can be \sphinxcode{\sphinxupquote{float}} or \sphinxcode{\sphinxupquote{ad\_double}} .
It must be either monotone increasing or decreasing.

\index{y\_init@\spxentry{y\_init}}\ignorespaces 

\subparagraph{y\_init}
\label{\detokenize{xsrst/numeric_ode_multi_step:y-init}}\label{\detokenize{xsrst/numeric_ode_multi_step:numeric-ode-multi-step-y-init}}\label{\detokenize{xsrst/numeric_ode_multi_step:index-13}}
This is the value of \(y(t)\) at the initial time
\sphinxstyleemphasis{t\_all} {[} 0 {]} as a numpy vector.
The type of its elements can be \sphinxcode{\sphinxupquote{float}} or \sphinxcode{\sphinxupquote{ad\_double}} .

\index{y\_all@\spxentry{y\_all}}\ignorespaces 

\subparagraph{y\_all}
\label{\detokenize{xsrst/numeric_ode_multi_step:y-all}}\label{\detokenize{xsrst/numeric_ode_multi_step:numeric-ode-multi-step-y-all}}\label{\detokenize{xsrst/numeric_ode_multi_step:index-14}}
This is the approximate solution for \(y(t)\) at all of the
times specified by \sphinxstyleemphasis{t\_all} as a numpy array.
The value \sphinxstyleemphasis{y\_all} {[} \sphinxstyleemphasis{i} , \sphinxstyleemphasis{j} {]} is the value of the j\sphinxhyphen{}th
component of \(y(t)\) at time \sphinxstyleemphasis{t\_all} {[} \sphinxstyleemphasis{i} {]} .


\subparagraph{numeric\_ode\_multi\_step\_xam\_py}
\label{\detokenize{xsrst/numeric_ode_multi_step_xam_py:numeric-ode-multi-step-xam-py}}\label{\detokenize{xsrst/numeric_ode_multi_step_xam_py::doc}}
\index{numeric\_ode\_multi\_step\_xam\_py@\spxentry{numeric\_ode\_multi\_step\_xam\_py}}\index{example@\spxentry{example}}\index{computing@\spxentry{computing}}\index{derivative@\spxentry{derivative}}\index{runge\sphinxhyphen{}kutta@\spxentry{runge\sphinxhyphen{}kutta}}\index{ode@\spxentry{ode}}\index{solution@\spxentry{solution}}\ignorespaces 

\subparagraph{2.4.1 Example Computing Derivative A Runge\sphinxhyphen{}Kutta Ode Solution}
\label{\detokenize{xsrst/numeric_ode_multi_step_xam_py:example-computing-derivative-a-runge-kutta-ode-solution}}\label{\detokenize{xsrst/numeric_ode_multi_step_xam_py:id1}}\begin{itemize}
\item {} 
{\hyperref[\detokenize{xsrst/numeric_ode_multi_step_xam_py:numeric-ode-multi-step-xam-py-ode}]{\sphinxcrossref{\DUrole{std,std-ref}{ODE}}}}

\item {} 
{\hyperref[\detokenize{xsrst/numeric_ode_multi_step_xam_py:numeric-ode-multi-step-xam-py-solution}]{\sphinxcrossref{\DUrole{std,std-ref}{Solution}}}}

\item {} 
{\hyperref[\detokenize{xsrst/numeric_ode_multi_step_xam_py:numeric-ode-multi-step-xam-py-derivative-of-solution}]{\sphinxcrossref{\DUrole{std,std-ref}{Derivative of Solution}}}}

\item {} 
{\hyperref[\detokenize{xsrst/numeric_ode_multi_step_xam_py:numeric-ode-multi-step-xam-py-source-code}]{\sphinxcrossref{\DUrole{std,std-ref}{Source Code}}}}

\end{itemize}

\index{ode@\spxentry{ode}}\ignorespaces 

\subparagraph{ODE}
\label{\detokenize{xsrst/numeric_ode_multi_step_xam_py:ode}}\label{\detokenize{xsrst/numeric_ode_multi_step_xam_py:numeric-ode-multi-step-xam-py-ode}}\label{\detokenize{xsrst/numeric_ode_multi_step_xam_py:index-1}}\begin{equation*}
\begin{split}\partial_t y_i (t, x) =  f(t, y, x) \left\{ \begin{array}{rl}
     x_0               & {\rm if} \; i = 0 \\
     x_i y_{i-1} (t)   & {\rm otherwise}
\end{array} \right.\end{split}
\end{equation*}
with the initial condition \(y(0) = 0\)

\index{solution@\spxentry{solution}}\ignorespaces 

\subparagraph{Solution}
\label{\detokenize{xsrst/numeric_ode_multi_step_xam_py:solution}}\label{\detokenize{xsrst/numeric_ode_multi_step_xam_py:numeric-ode-multi-step-xam-py-solution}}\label{\detokenize{xsrst/numeric_ode_multi_step_xam_py:index-2}}
This is a special case for which we know the solution
\begin{equation*}
\begin{split}y_i (t, x) = \left\{ \begin{array}{rl}
     t  x_0                            & {\rm if} \; i = 0 \\
     ( t^i / (i+1) ! ) \prod_{j=0}^i x_j   & {\rm otherwise}
\end{array} \right.\end{split}
\end{equation*}
\index{derivative@\spxentry{derivative}}\index{solution@\spxentry{solution}}\ignorespaces 

\subparagraph{Derivative of Solution}
\label{\detokenize{xsrst/numeric_ode_multi_step_xam_py:derivative-of-solution}}\label{\detokenize{xsrst/numeric_ode_multi_step_xam_py:numeric-ode-multi-step-xam-py-derivative-of-solution}}\label{\detokenize{xsrst/numeric_ode_multi_step_xam_py:index-3}}
For this special case, the partial derivative of the solution with respect
to the j\sphinxhyphen{}th component of the vector \(x\) is
\begin{equation*}
\begin{split}\partial_{x(j)} y_i (t, x) =  \left\{ \begin{array}{rl}
     y_i (t, x) / x_j      & {\rm if} \; j \leq i \\
     0                     & {\rm otherwise}
\end{array} \right.\end{split}
\end{equation*}
\index{source@\spxentry{source}}\index{code@\spxentry{code}}\ignorespaces 

\subparagraph{Source Code}
\label{\detokenize{xsrst/numeric_ode_multi_step_xam_py:source-code}}\label{\detokenize{xsrst/numeric_ode_multi_step_xam_py:numeric-ode-multi-step-xam-py-source-code}}\label{\detokenize{xsrst/numeric_ode_multi_step_xam_py:index-4}}
\begin{sphinxVerbatim}[commandchars=\\\{\}]
\PYG{k+kn}{import} \PYG{n+nn}{numpy}
\PYG{k+kn}{import} \PYG{n+nn}{scipy}\PYG{n+nn}{.}\PYG{n+nn}{special}
\PYG{k+kn}{import} \PYG{n+nn}{cppad\PYGZus{}py}
\PYG{k+kn}{from} \PYG{n+nn}{ode\PYGZus{}multi\PYGZus{}step} \PYG{k+kn}{import} \PYG{n}{ode\PYGZus{}multi\PYGZus{}step}
\PYG{k+kn}{from} \PYG{n+nn}{runge4\PYGZus{}step} \PYG{k+kn}{import} \PYG{n}{runge4\PYGZus{}step}
\PYG{c+c1}{\PYGZsh{}}
\PYG{k}{class} \PYG{n+nc}{fun\PYGZus{}class} \PYG{p}{:}
    \PYG{c+c1}{\PYGZsh{}}
    \PYG{k}{def} \PYG{n+nf+fm}{\PYGZus{}\PYGZus{}init\PYGZus{}\PYGZus{}}\PYG{p}{(}\PYG{n+nb+bp}{self}\PYG{p}{,} \PYG{n}{x}\PYG{p}{)} \PYG{p}{:}
        \PYG{n+nb+bp}{self}\PYG{o}{.}\PYG{n}{x}     \PYG{o}{=} \PYG{n}{x}
        \PYG{n+nb+bp}{self}\PYG{o}{.}\PYG{n}{index} \PYG{o}{=} \PYG{k+kc}{None}
    \PYG{c+c1}{\PYGZsh{}}
    \PYG{k}{def} \PYG{n+nf}{set\PYGZus{}t\PYGZus{}all\PYGZus{}index}\PYG{p}{(}\PYG{n+nb+bp}{self}\PYG{p}{,} \PYG{n}{index}\PYG{p}{)} \PYG{p}{:}
        \PYG{n+nb+bp}{self}\PYG{o}{.}\PYG{n}{index} \PYG{o}{=} \PYG{n}{index}
    \PYG{c+c1}{\PYGZsh{}}
    \PYG{k}{def} \PYG{n+nf}{f}\PYG{p}{(}\PYG{n+nb+bp}{self}\PYG{p}{,} \PYG{n}{t}\PYG{p}{,} \PYG{n}{y}\PYG{p}{)} \PYG{p}{:}
        \PYG{k}{assert} \PYG{n+nb+bp}{self}\PYG{o}{.}\PYG{n}{index} \PYG{o}{!=} \PYG{k+kc}{None}
        \PYG{c+c1}{\PYGZsh{} This ode does not depend on the the time index}
        \PYG{c+c1}{\PYGZsh{}}
        \PYG{n}{y\PYGZus{}shift} \PYG{o}{=} \PYG{n}{numpy}\PYG{o}{.}\PYG{n}{concatenate}\PYG{p}{(} \PYG{p}{(} \PYG{p}{[}\PYG{l+m+mf}{1.0}\PYG{p}{]} \PYG{p}{,} \PYG{n}{y}\PYG{p}{[}\PYG{l+m+mi}{0}\PYG{p}{:}\PYG{o}{\PYGZhy{}}\PYG{l+m+mi}{1}\PYG{p}{]} \PYG{p}{)} \PYG{p}{)}
        \PYG{k}{return} \PYG{n+nb+bp}{self}\PYG{o}{.}\PYG{n}{x} \PYG{o}{*} \PYG{n}{y\PYGZus{}shift}
\PYG{c+c1}{\PYGZsh{}}
\PYG{k}{def} \PYG{n+nf}{ode\PYGZus{}multi\PYGZus{}step\PYGZus{}xam}\PYG{p}{(}\PYG{p}{)} \PYG{p}{:}
    \PYG{n}{ok}    \PYG{o}{=} \PYG{k+kc}{True}
    \PYG{n}{nx}    \PYG{o}{=} \PYG{l+m+mi}{4}
    \PYG{n}{eps99} \PYG{o}{=} \PYG{l+m+mf}{99.0} \PYG{o}{*} \PYG{n}{numpy}\PYG{o}{.}\PYG{n}{finfo}\PYG{p}{(}\PYG{n+nb}{float}\PYG{p}{)}\PYG{o}{.}\PYG{n}{eps}
    \PYG{c+c1}{\PYGZsh{}}
    \PYG{c+c1}{\PYGZsh{} independent variables for this g(x) = y(1, x)}
    \PYG{n}{x}  \PYG{o}{=} \PYG{n}{numpy}\PYG{o}{.}\PYG{n}{array}\PYG{p}{(} \PYG{n}{nx} \PYG{o}{*} \PYG{p}{[} \PYG{l+m+mf}{1.0} \PYG{p}{]} \PYG{p}{)}
    \PYG{n}{ax} \PYG{o}{=} \PYG{n}{cppad\PYGZus{}py}\PYG{o}{.}\PYG{n}{independent}\PYG{p}{(}\PYG{n}{x}\PYG{p}{)}
    \PYG{c+c1}{\PYGZsh{}}
    \PYG{c+c1}{\PYGZsh{} function to pass to runge4\PYGZus{}step}
    \PYG{n}{fun} \PYG{o}{=} \PYG{n}{fun\PYGZus{}class}\PYG{p}{(}\PYG{n}{ax}\PYG{p}{)}
    \PYG{c+c1}{\PYGZsh{}}
    \PYG{c+c1}{\PYGZsh{} routine we are usning for each step}
    \PYG{n}{one\PYGZus{}step} \PYG{o}{=} \PYG{n}{runge4\PYGZus{}step}
    \PYG{c+c1}{\PYGZsh{}}
    \PYG{c+c1}{\PYGZsh{} initiali value for the ODE}
    \PYG{n}{n\PYGZus{}step}  \PYG{o}{=} \PYG{l+m+mi}{10}
    \PYG{n}{ay\PYGZus{}init} \PYG{o}{=}  \PYG{n}{numpy}\PYG{o}{.}\PYG{n}{array}\PYG{p}{(} \PYG{n}{nx} \PYG{o}{*} \PYG{p}{[} \PYG{n}{cppad\PYGZus{}py}\PYG{o}{.}\PYG{n}{a\PYGZus{}double}\PYG{p}{(}\PYG{l+m+mf}{0.0}\PYG{p}{)} \PYG{p}{]} \PYG{p}{)}
    \PYG{n}{t\PYGZus{}final} \PYG{o}{=} \PYG{l+m+mf}{0.75}
    \PYG{n}{t\PYGZus{}all}   \PYG{o}{=}  \PYG{p}{[} \PYG{n}{t\PYGZus{}final} \PYG{o}{*} \PYG{n}{i} \PYG{o}{/} \PYG{p}{(}\PYG{n}{n\PYGZus{}step} \PYG{o}{\PYGZhy{}} \PYG{l+m+mi}{1}\PYG{p}{)} \PYG{k}{for} \PYG{n}{i} \PYG{o+ow}{in} \PYG{n+nb}{range}\PYG{p}{(}\PYG{n}{n\PYGZus{}step}\PYG{p}{)} \PYG{p}{]}
    \PYG{n}{t\PYGZus{}all}   \PYG{o}{=}  \PYG{n}{numpy}\PYG{o}{.}\PYG{n}{array}\PYG{p}{(} \PYG{n}{t\PYGZus{}all} \PYG{p}{)}
    \PYG{c+c1}{\PYGZsh{}}
    \PYG{c+c1}{\PYGZsh{} take multiple steps}
    \PYG{n}{ay} \PYG{o}{=} \PYG{n}{ode\PYGZus{}multi\PYGZus{}step}\PYG{p}{(}\PYG{n}{one\PYGZus{}step}\PYG{p}{,} \PYG{n}{fun}\PYG{p}{,} \PYG{n}{t\PYGZus{}all}\PYG{p}{,} \PYG{n}{ay\PYGZus{}init}\PYG{p}{)}
    \PYG{n}{ay\PYGZus{}final} \PYG{o}{=} \PYG{n}{ay}\PYG{p}{[}\PYG{n}{n\PYGZus{}step} \PYG{o}{\PYGZhy{}} \PYG{l+m+mi}{1}\PYG{p}{,}\PYG{p}{:}\PYG{p}{]}
    \PYG{c+c1}{\PYGZsh{}}
    \PYG{c+c1}{\PYGZsh{} g(x) = y(t\PYGZus{}final, x)}
    \PYG{n}{g} \PYG{o}{=} \PYG{n}{cppad\PYGZus{}py}\PYG{o}{.}\PYG{n}{d\PYGZus{}fun}\PYG{p}{(}\PYG{n}{ax}\PYG{p}{,} \PYG{n}{ay\PYGZus{}final}\PYG{p}{)}
    \PYG{c+c1}{\PYGZsh{}}
    \PYG{c+c1}{\PYGZsh{} Check g\PYGZus{}i (x) = prod\PYGZus{}\PYGZob{}j=0\PYGZcb{}\PYGZca{}i x[j] t\PYGZca{}(i+1) / (i+1) !}
    \PYG{c+c1}{\PYGZsh{} 4th order method should be very accutate for functions}
    \PYG{c+c1}{\PYGZsh{} or order 4 or less in t.}
    \PYG{n}{x}  \PYG{o}{=} \PYG{n}{numpy}\PYG{o}{.}\PYG{n}{arange}\PYG{p}{(}\PYG{l+m+mi}{0}\PYG{p}{,} \PYG{n}{nx}\PYG{p}{)} \PYG{o}{+} \PYG{l+m+mf}{1.0}
    \PYG{n}{gx} \PYG{o}{=} \PYG{n}{g}\PYG{o}{.}\PYG{n}{forward}\PYG{p}{(}\PYG{l+m+mi}{0}\PYG{p}{,} \PYG{n}{x}\PYG{p}{)}
    \PYG{n}{prod} \PYG{o}{=} \PYG{l+m+mf}{1.0}
    \PYG{k}{for} \PYG{n}{i} \PYG{o+ow}{in} \PYG{n+nb}{range}\PYG{p}{(}\PYG{n}{nx}\PYG{p}{)} \PYG{p}{:}
        \PYG{n}{prod}      \PYG{o}{=} \PYG{n}{prod} \PYG{o}{*} \PYG{n}{x}\PYG{p}{[}\PYG{n}{i}\PYG{p}{]}
        \PYG{n}{power}     \PYG{o}{=} \PYG{n}{numpy}\PYG{o}{.}\PYG{n}{power}\PYG{p}{(}\PYG{n}{t\PYGZus{}final}\PYG{p}{,} \PYG{n}{i}\PYG{o}{+}\PYG{l+m+mi}{1}\PYG{p}{)}
        \PYG{n}{factorial} \PYG{o}{=} \PYG{n}{scipy}\PYG{o}{.}\PYG{n}{special}\PYG{o}{.}\PYG{n}{factorial}\PYG{p}{(}\PYG{n}{i}\PYG{o}{+}\PYG{l+m+mi}{1}\PYG{p}{)}
        \PYG{n}{check}     \PYG{o}{=} \PYG{n}{prod} \PYG{o}{*} \PYG{n}{power} \PYG{o}{/} \PYG{n}{factorial}
        \PYG{n}{rel\PYGZus{}error} \PYG{o}{=} \PYG{n}{gx}\PYG{p}{[}\PYG{n}{i}\PYG{p}{]} \PYG{o}{/} \PYG{n}{check} \PYG{o}{\PYGZhy{}} \PYG{l+m+mf}{1.0}
        \PYG{n}{ok}       \PYG{o}{\PYGZam{}}\PYG{o}{=} \PYG{n+nb}{abs}\PYG{p}{(}\PYG{n}{rel\PYGZus{}error}\PYG{p}{)} \PYG{o}{\PYGZlt{}} \PYG{n}{eps99}
    \PYG{c+c1}{\PYGZsh{}}
    \PYG{c+c1}{\PYGZsh{} Check derivative of g\PYGZus{}i (x) w.r.t x\PYGZus{}j}
    \PYG{c+c1}{\PYGZsh{} is zero for i \PYGZlt{} j and  g\PYGZus{}i(x) / x\PYGZus{}j otherwise}
    \PYG{n}{J} \PYG{o}{=} \PYG{n}{g}\PYG{o}{.}\PYG{n}{jacobian}\PYG{p}{(}\PYG{n}{x}\PYG{p}{)}
    \PYG{k}{for} \PYG{n}{j} \PYG{o+ow}{in} \PYG{n+nb}{range}\PYG{p}{(}\PYG{n}{nx}\PYG{p}{)} \PYG{p}{:}
        \PYG{k}{for} \PYG{n}{i} \PYG{o+ow}{in} \PYG{n+nb}{range}\PYG{p}{(}\PYG{n}{nx}\PYG{p}{)} \PYG{p}{:}
            \PYG{k}{if} \PYG{n}{i} \PYG{o}{\PYGZlt{}} \PYG{n}{j} \PYG{p}{:}
                \PYG{n}{ok} \PYG{o}{\PYGZam{}}\PYG{o}{=} \PYG{n}{J}\PYG{p}{[}\PYG{n}{i}\PYG{p}{,}\PYG{n}{j}\PYG{p}{]} \PYG{o}{==} \PYG{l+m+mf}{0.0}
            \PYG{k}{else} \PYG{p}{:}
                \PYG{n}{check}     \PYG{o}{=} \PYG{n}{gx}\PYG{p}{[}\PYG{n}{i}\PYG{p}{]} \PYG{o}{/} \PYG{n}{x}\PYG{p}{[}\PYG{n}{j}\PYG{p}{]}
                \PYG{n}{rel\PYGZus{}error} \PYG{o}{=} \PYG{n}{J}\PYG{p}{[}\PYG{n}{i}\PYG{p}{,}\PYG{n}{j}\PYG{p}{]} \PYG{o}{/} \PYG{n}{check} \PYG{o}{\PYGZhy{}} \PYG{l+m+mf}{1.0}
                \PYG{n}{ok}       \PYG{o}{\PYGZam{}}\PYG{o}{=} \PYG{n+nb}{abs}\PYG{p}{(}\PYG{n}{rel\PYGZus{}error}\PYG{p}{)} \PYG{o}{\PYGZlt{}} \PYG{n}{eps99}

    \PYG{k}{return} \PYG{n}{ok}
\end{sphinxVerbatim}


\bigskip\hrule\bigskip


xsrst input file: \sphinxcode{\sphinxupquote{example/python/numeric/ode\_multi\_step\_xam.py}}

\index{example@\spxentry{example}}\ignorespaces 

\subparagraph{Example}
\label{\detokenize{xsrst/numeric_ode_multi_step:example}}\label{\detokenize{xsrst/numeric_ode_multi_step:numeric-ode-multi-step-example}}\label{\detokenize{xsrst/numeric_ode_multi_step:index-15}}
{\hyperref[\detokenize{xsrst/numeric_ode_multi_step_xam_py:id1}]{\sphinxcrossref{\DUrole{std,std-ref}{numeric\_ode\_multi\_step\_xam\_py}}}}

\index{source@\spxentry{source}}\index{code@\spxentry{code}}\ignorespaces 

\subparagraph{Source Code}
\label{\detokenize{xsrst/numeric_ode_multi_step:source-code}}\label{\detokenize{xsrst/numeric_ode_multi_step:numeric-ode-multi-step-source-code}}\label{\detokenize{xsrst/numeric_ode_multi_step:index-16}}
\begin{sphinxVerbatim}[commandchars=\\\{\}]
\PYG{k+kn}{import} \PYG{n+nn}{numpy}
\PYG{k+kn}{import} \PYG{n+nn}{copy}
\PYG{k}{def} \PYG{n+nf}{ode\PYGZus{}multi\PYGZus{}step}\PYG{p}{(}\PYG{n}{one\PYGZus{}step}\PYG{p}{,} \PYG{n}{fun}\PYG{p}{,} \PYG{n}{t\PYGZus{}all}\PYG{p}{,} \PYG{n}{y\PYGZus{}init} \PYG{p}{)} \PYG{p}{:}
    \PYG{n}{dtype}      \PYG{o}{=} \PYG{n+nb}{type}\PYG{p}{(}\PYG{n}{y\PYGZus{}init}\PYG{p}{[}\PYG{l+m+mi}{0}\PYG{p}{]}\PYG{p}{)}
    \PYG{n}{n\PYGZus{}var}      \PYG{o}{=} \PYG{n}{y\PYGZus{}init}\PYG{o}{.}\PYG{n}{size}
    \PYG{n}{n\PYGZus{}step}     \PYG{o}{=} \PYG{n}{t\PYGZus{}all}\PYG{o}{.}\PYG{n}{size} \PYG{o}{\PYGZhy{}} \PYG{l+m+mi}{1}
    \PYG{n}{y\PYGZus{}all}      \PYG{o}{=} \PYG{n}{numpy}\PYG{o}{.}\PYG{n}{empty}\PYG{p}{(} \PYG{p}{(}\PYG{n}{n\PYGZus{}step}\PYG{o}{+}\PYG{l+m+mi}{1}\PYG{p}{,} \PYG{n}{n\PYGZus{}var}\PYG{p}{)}\PYG{p}{,} \PYG{n}{dtype} \PYG{o}{=} \PYG{n}{dtype}  \PYG{p}{)}
    \PYG{n}{y1}         \PYG{o}{=} \PYG{n}{y\PYGZus{}init}
    \PYG{n}{t1}         \PYG{o}{=} \PYG{n}{t\PYGZus{}all}\PYG{p}{[}\PYG{l+m+mi}{0}\PYG{p}{]}
    \PYG{n}{y\PYGZus{}all}\PYG{p}{[}\PYG{l+m+mi}{0}\PYG{p}{,}\PYG{p}{:}\PYG{p}{]} \PYG{o}{=} \PYG{n}{y1}
    \PYG{k}{for} \PYG{n}{i} \PYG{o+ow}{in} \PYG{n+nb}{range}\PYG{p}{(}\PYG{n}{n\PYGZus{}step}\PYG{p}{)} \PYG{p}{:}
        \PYG{n}{fun}\PYG{o}{.}\PYG{n}{set\PYGZus{}t\PYGZus{}all\PYGZus{}index}\PYG{p}{(}\PYG{n}{i}\PYG{p}{)}
        \PYG{n}{y0}            \PYG{o}{=} \PYG{n}{y1}
        \PYG{n}{t0}            \PYG{o}{=} \PYG{n}{t1}
        \PYG{n}{t1}            \PYG{o}{=} \PYG{n}{t\PYGZus{}all}\PYG{p}{[}\PYG{n}{i}\PYG{o}{+}\PYG{l+m+mi}{1}\PYG{p}{]}
        \PYG{n}{t\PYGZus{}step}        \PYG{o}{=} \PYG{n}{t1} \PYG{o}{\PYGZhy{}} \PYG{n}{t0}
        \PYG{n}{y1}            \PYG{o}{=} \PYG{n}{one\PYGZus{}step}\PYG{p}{(}\PYG{n}{fun}\PYG{p}{,} \PYG{n}{t0}\PYG{p}{,} \PYG{n}{y0}\PYG{p}{,} \PYG{n}{t\PYGZus{}step}\PYG{p}{)}
        \PYG{n}{y\PYGZus{}all}\PYG{p}{[}\PYG{n}{i}\PYG{o}{+}\PYG{l+m+mi}{1}\PYG{p}{,}\PYG{p}{:}\PYG{p}{]}  \PYG{o}{=} \PYG{n}{copy}\PYG{o}{.}\PYG{n}{copy}\PYG{p}{(}\PYG{n}{y1}\PYG{p}{)}
    \PYG{k}{return} \PYG{n}{y\PYGZus{}all}
\end{sphinxVerbatim}


\bigskip\hrule\bigskip


xsrst input file: \sphinxcode{\sphinxupquote{example/python/numeric/ode\_multi\_step.py}}


\subparagraph{numeric\_optimize\_fun\_class}
\label{\detokenize{xsrst/numeric_optimize_fun_class:numeric-optimize-fun-class}}\label{\detokenize{xsrst/numeric_optimize_fun_class::doc}}
\index{numeric\_optimize\_fun\_class@\spxentry{numeric\_optimize\_fun\_class}}\index{helper@\spxentry{helper}}\index{class@\spxentry{class}}\index{that@\spxentry{that}}\index{defines@\spxentry{defines}}\index{functions@\spxentry{functions}}\index{needed@\spxentry{needed}}\index{for@\spxentry{for}}\index{optimization@\spxentry{optimization}}\ignorespaces 

\subparagraph{2.5 A Helper Class That Defines Functions Needed for Optimization}
\label{\detokenize{xsrst/numeric_optimize_fun_class:a-helper-class-that-defines-functions-needed-for-optimization}}\label{\detokenize{xsrst/numeric_optimize_fun_class:id1}}\begin{itemize}
\item {} 
{\hyperref[\detokenize{xsrst/numeric_optimize_fun_class:numeric-optimize-fun-class-syntax}]{\sphinxcrossref{\DUrole{std,std-ref}{Syntax}}}}

\item {} 
{\hyperref[\detokenize{xsrst/numeric_optimize_fun_class:numeric-optimize-fun-class-purpose}]{\sphinxcrossref{\DUrole{std,std-ref}{Purpose}}}}

\item {} 
{\hyperref[\detokenize{xsrst/numeric_optimize_fun_class:numeric-optimize-fun-class-objective-ad}]{\sphinxcrossref{\DUrole{std,std-ref}{objective\_ad}}}}

\item {} 
{\hyperref[\detokenize{xsrst/numeric_optimize_fun_class:numeric-optimize-fun-class-constraint-ad}]{\sphinxcrossref{\DUrole{std,std-ref}{constraint\_ad}}}}

\item {} \begin{description}
\item[{{\hyperref[\detokenize{xsrst/numeric_optimize_fun_class:numeric-optimize-fun-class-optimize-fun}]{\sphinxcrossref{\DUrole{std,std-ref}{optimize\_fun}}}}}] \leavevmode\begin{itemize}
\item {} 
{\hyperref[\detokenize{xsrst/numeric_optimize_fun_class:numeric-optimize-fun-class-optimize-fun-objective-fun}]{\sphinxcrossref{\DUrole{std,std-ref}{objective\_fun}}}}

\item {} 
{\hyperref[\detokenize{xsrst/numeric_optimize_fun_class:numeric-optimize-fun-class-optimize-fun-objective-grad}]{\sphinxcrossref{\DUrole{std,std-ref}{objective\_grad}}}}

\item {} 
{\hyperref[\detokenize{xsrst/numeric_optimize_fun_class:numeric-optimize-fun-class-optimize-fun-objective-hess}]{\sphinxcrossref{\DUrole{std,std-ref}{objective\_hess}}}}

\item {} 
{\hyperref[\detokenize{xsrst/numeric_optimize_fun_class:numeric-optimize-fun-class-optimize-fun-constraint-fun}]{\sphinxcrossref{\DUrole{std,std-ref}{constraint\_fun}}}}

\item {} 
{\hyperref[\detokenize{xsrst/numeric_optimize_fun_class:numeric-optimize-fun-class-optimize-fun-constraint-jac}]{\sphinxcrossref{\DUrole{std,std-ref}{constraint\_jac}}}}

\item {} 
{\hyperref[\detokenize{xsrst/numeric_optimize_fun_class:numeric-optimize-fun-class-optimize-fun-constraint-hess}]{\sphinxcrossref{\DUrole{std,std-ref}{constraint\_hess}}}}

\end{itemize}

\end{description}

\item {} 
{\hyperref[\detokenize{xsrst/numeric_optimize_fun_class:numeric-optimize-fun-class-example}]{\sphinxcrossref{\DUrole{std,std-ref}{Example}}}}

\item {} 
{\hyperref[\detokenize{xsrst/numeric_optimize_fun_class:numeric-optimize-fun-class-source-code}]{\sphinxcrossref{\DUrole{std,std-ref}{Source Code}}}}

\end{itemize}

\index{syntax@\spxentry{syntax}}\ignorespaces 

\subparagraph{Syntax}
\label{\detokenize{xsrst/numeric_optimize_fun_class:syntax}}\label{\detokenize{xsrst/numeric_optimize_fun_class:numeric-optimize-fun-class-syntax}}\label{\detokenize{xsrst/numeric_optimize_fun_class:index-1}}
\sphinxstyleemphasis{optimize\_fun} =  \sphinxcode{\sphinxupquote{optimize\_fun\_class}} ( \sphinxstyleemphasis{objective\_ad} , \sphinxstyleemphasis{constraint\_ad} )

\index{purpose@\spxentry{purpose}}\ignorespaces 

\subparagraph{Purpose}
\label{\detokenize{xsrst/numeric_optimize_fun_class:purpose}}\label{\detokenize{xsrst/numeric_optimize_fun_class:numeric-optimize-fun-class-purpose}}\label{\detokenize{xsrst/numeric_optimize_fun_class:index-2}}
This class is an aid solving optimization problems of the form
\begin{equation*}
\begin{split}\begin{array}{rl}
{\rm minimize}       & f(x) \; {\rm w.r.t} \; x \\
{\rm subject \; to}  & a \leq g(x) \leq b \\
\end{array}\end{split}
\end{equation*}
where \(x\) is a vector,
\(f(x)\) is a scalar, and
\(a, g(x), b\) are all vectors with the same length.
We use \(n\), \(m\) for the length of the vectors
\(x\) and \(g(x)\) respectively.

\index{objective\_ad@\spxentry{objective\_ad}}\ignorespaces 

\subparagraph{objective\_ad}
\label{\detokenize{xsrst/numeric_optimize_fun_class:objective-ad}}\label{\detokenize{xsrst/numeric_optimize_fun_class:numeric-optimize-fun-class-objective-ad}}\label{\detokenize{xsrst/numeric_optimize_fun_class:index-3}}
This is a {\hyperref[\detokenize{xsrst/py_fun_ctor:py-fun-ctor-syntax-d-fun}]{\sphinxcrossref{\DUrole{std,std-ref}{d\_fun}}}}
representation of the function \(f(x)\).

\index{constraint\_ad@\spxentry{constraint\_ad}}\ignorespaces 

\subparagraph{constraint\_ad}
\label{\detokenize{xsrst/numeric_optimize_fun_class:constraint-ad}}\label{\detokenize{xsrst/numeric_optimize_fun_class:numeric-optimize-fun-class-constraint-ad}}\label{\detokenize{xsrst/numeric_optimize_fun_class:index-4}}
This is a \sphinxcode{\sphinxupquote{d\_fun}} representation of the function \(g(x)\).

\index{optimize\_fun@\spxentry{optimize\_fun}}\ignorespaces 

\subparagraph{optimize\_fun}
\label{\detokenize{xsrst/numeric_optimize_fun_class:optimize-fun}}\label{\detokenize{xsrst/numeric_optimize_fun_class:numeric-optimize-fun-class-optimize-fun}}\label{\detokenize{xsrst/numeric_optimize_fun_class:index-5}}
This class object has the following functions defined:

\index{objective\_fun@\spxentry{objective\_fun}}\ignorespaces 

\subparagraph{objective\_fun}
\label{\detokenize{xsrst/numeric_optimize_fun_class:objective-fun}}\label{\detokenize{xsrst/numeric_optimize_fun_class:numeric-optimize-fun-class-optimize-fun-objective-fun}}\label{\detokenize{xsrst/numeric_optimize_fun_class:index-6}}
The syntax

\begin{DUlineblock}{0em}
\item[]         \sphinxstyleemphasis{y} = \sphinxstyleemphasis{optimize\_fun} . \sphinxcode{\sphinxupquote{objective\_fun}} ( \sphinxstyleemphasis{x} )
\end{DUlineblock}

sets \(y = f(x)\) where
\sphinxstyleemphasis{x} is a numpy vector with length \sphinxstyleemphasis{n}
and \sphinxstyleemphasis{y} is a scalar.

\index{objective\_grad@\spxentry{objective\_grad}}\ignorespaces 

\subparagraph{objective\_grad}
\label{\detokenize{xsrst/numeric_optimize_fun_class:objective-grad}}\label{\detokenize{xsrst/numeric_optimize_fun_class:numeric-optimize-fun-class-optimize-fun-objective-grad}}\label{\detokenize{xsrst/numeric_optimize_fun_class:index-7}}
The syntax

\begin{DUlineblock}{0em}
\item[]         \sphinxstyleemphasis{z} = \sphinxstyleemphasis{optimize\_fun} . \sphinxcode{\sphinxupquote{objective\_grad}} ( \sphinxstyleemphasis{x} )
\end{DUlineblock}

sets \(z = f^{(1)} (x)\) where
\sphinxstyleemphasis{x} and \sphinxstyleemphasis{z} are numpy vectors with length \sphinxstyleemphasis{n} .

\index{objective\_hess@\spxentry{objective\_hess}}\ignorespaces 

\subparagraph{objective\_hess}
\label{\detokenize{xsrst/numeric_optimize_fun_class:objective-hess}}\label{\detokenize{xsrst/numeric_optimize_fun_class:numeric-optimize-fun-class-optimize-fun-objective-hess}}\label{\detokenize{xsrst/numeric_optimize_fun_class:index-8}}
The syntax

\begin{DUlineblock}{0em}
\item[]         \sphinxstyleemphasis{h} = \sphinxstyleemphasis{optimize\_fun} . \sphinxcode{\sphinxupquote{objective\_hess}} ( \sphinxstyleemphasis{x} )
\end{DUlineblock}

sets \(h = f^{(2)} (x)\) where
\sphinxstyleemphasis{x} is a numpy vector with length \sphinxstyleemphasis{n}
and \sphinxstyleemphasis{h} is a numpy \sphinxstyleemphasis{n} by \sphinxstyleemphasis{n}  matrix.

\index{constraint\_fun@\spxentry{constraint\_fun}}\ignorespaces 

\subparagraph{constraint\_fun}
\label{\detokenize{xsrst/numeric_optimize_fun_class:constraint-fun}}\label{\detokenize{xsrst/numeric_optimize_fun_class:numeric-optimize-fun-class-optimize-fun-constraint-fun}}\label{\detokenize{xsrst/numeric_optimize_fun_class:index-9}}
The syntax

\begin{DUlineblock}{0em}
\item[]         \sphinxstyleemphasis{y} = \sphinxstyleemphasis{optimize\_fun} . \sphinxcode{\sphinxupquote{constraint\_fun}} ( \sphinxstyleemphasis{x} )
\end{DUlineblock}

sets \(y = g(x)\) where
\sphinxstyleemphasis{x} ( \sphinxstyleemphasis{y} ) is a numpy vector with length
\sphinxstyleemphasis{n} ( \sphinxstyleemphasis{m} ).

\index{constraint\_jac@\spxentry{constraint\_jac}}\ignorespaces 

\subparagraph{constraint\_jac}
\label{\detokenize{xsrst/numeric_optimize_fun_class:constraint-jac}}\label{\detokenize{xsrst/numeric_optimize_fun_class:numeric-optimize-fun-class-optimize-fun-constraint-jac}}\label{\detokenize{xsrst/numeric_optimize_fun_class:index-10}}
The syntax

\begin{DUlineblock}{0em}
\item[]         \sphinxstyleemphasis{J} = \sphinxstyleemphasis{optimize\_fun} . \sphinxcode{\sphinxupquote{constraint\_jac}} ( \sphinxstyleemphasis{x} )
\end{DUlineblock}

sets \(J = g^{(1)} (x)\) where
\sphinxstyleemphasis{x} is a numpy vector with length \sphinxstyleemphasis{n}
and \sphinxstyleemphasis{J} is a numpy \sphinxstyleemphasis{m} by \sphinxstyleemphasis{n}  matrix.

\index{constraint\_hess@\spxentry{constraint\_hess}}\ignorespaces 

\subparagraph{constraint\_hess}
\label{\detokenize{xsrst/numeric_optimize_fun_class:constraint-hess}}\label{\detokenize{xsrst/numeric_optimize_fun_class:numeric-optimize-fun-class-optimize-fun-constraint-hess}}\label{\detokenize{xsrst/numeric_optimize_fun_class:index-11}}
The syntax

\begin{DUlineblock}{0em}
\item[]         \sphinxstyleemphasis{H} = \sphinxstyleemphasis{optimize\_fun} . \sphinxcode{\sphinxupquote{constraint\_hess}} ( \sphinxstyleemphasis{x} , \sphinxstyleemphasis{v} )
\end{DUlineblock}

sets
\begin{equation*}
\begin{split}H = \sum_{i=0}^{m-1} v_k g_i^{(2)} (x)\end{split}
\end{equation*}
where \sphinxstyleemphasis{x} is a numpy vector with length \sphinxstyleemphasis{n}
and \sphinxstyleemphasis{H} is a numpy \sphinxstyleemphasis{n} by \sphinxstyleemphasis{n}  matrix.


\subparagraph{numeric\_optimize\_fun\_xam\_py}
\label{\detokenize{xsrst/numeric_optimize_fun_xam_py:numeric-optimize-fun-xam-py}}\label{\detokenize{xsrst/numeric_optimize_fun_xam_py::doc}}
\index{numeric\_optimize\_fun\_xam\_py@\spxentry{numeric\_optimize\_fun\_xam\_py}}\index{example@\spxentry{example}}\index{using@\spxentry{using}}\index{optimize\_fun\_class@\spxentry{optimize\_fun\_class}}\index{with@\spxentry{with}}\index{scipy@\spxentry{scipy}}\index{optimization@\spxentry{optimization}}\ignorespaces 

\subparagraph{2.5.1 Example Using optimize\_fun\_class with Scipy Optimization}
\label{\detokenize{xsrst/numeric_optimize_fun_xam_py:example-using-optimize-fun-class-with-scipy-optimization}}\label{\detokenize{xsrst/numeric_optimize_fun_xam_py:id1}}\begin{itemize}
\item {} 
{\hyperref[\detokenize{xsrst/numeric_optimize_fun_xam_py:numeric-optimize-fun-xam-py-reference}]{\sphinxcrossref{\DUrole{std,std-ref}{Reference}}}}

\item {} 
{\hyperref[\detokenize{xsrst/numeric_optimize_fun_xam_py:numeric-optimize-fun-xam-py-problem}]{\sphinxcrossref{\DUrole{std,std-ref}{Problem}}}}

\item {} 
{\hyperref[\detokenize{xsrst/numeric_optimize_fun_xam_py:numeric-optimize-fun-xam-py-trust-constr}]{\sphinxcrossref{\DUrole{std,std-ref}{trust\_constr}}}}

\item {} 
{\hyperref[\detokenize{xsrst/numeric_optimize_fun_xam_py:numeric-optimize-fun-xam-py-source-code}]{\sphinxcrossref{\DUrole{std,std-ref}{Source Code}}}}

\end{itemize}

\index{reference@\spxentry{reference}}\ignorespaces 

\subparagraph{Reference}
\label{\detokenize{xsrst/numeric_optimize_fun_xam_py:reference}}\label{\detokenize{xsrst/numeric_optimize_fun_xam_py:numeric-optimize-fun-xam-py-reference}}\label{\detokenize{xsrst/numeric_optimize_fun_xam_py:index-1}}
This problem comes form the
\sphinxhref{https://coin-or.github.io/Ipopt/INTERFACES.html}{Interfaces}
section of the Ipopt documentation.

\index{problem@\spxentry{problem}}\ignorespaces 

\subparagraph{Problem}
\label{\detokenize{xsrst/numeric_optimize_fun_xam_py:problem}}\label{\detokenize{xsrst/numeric_optimize_fun_xam_py:numeric-optimize-fun-xam-py-problem}}\label{\detokenize{xsrst/numeric_optimize_fun_xam_py:index-2}}\begin{equation*}
\begin{split}\begin{array}{cr}
{\rm minimize}      & x_0 x_3 ( x_0 + x_1 + x_2 ) + x_2   \\
{\rm subject \; to} &             x_0 x_1 x_2 x_3 \geq 25 \\
                     & x_0^2 + x_1^2 + x_2^2 + x_3^2 = 40  \\
                     &                    1 \leq x \leq 5
\end{array}\end{split}
\end{equation*}
with the starting point \(x = (1, 5, 5, 1)\).
The optimal value for \(x\) is
\begin{equation*}
\begin{split}\newcommand{\W}[1]{{\; #1 \;}}
(1.00000000 \W{,} 4.74299963 \W{,} 3.82114998 \W{,} 1.37940829)\end{split}
\end{equation*}
\index{trust\_constr@\spxentry{trust\_constr}}\ignorespaces 

\subparagraph{trust\_constr}
\label{\detokenize{xsrst/numeric_optimize_fun_xam_py:trust-constr}}\label{\detokenize{xsrst/numeric_optimize_fun_xam_py:numeric-optimize-fun-xam-py-trust-constr}}\label{\detokenize{xsrst/numeric_optimize_fun_xam_py:index-3}}
This is one of the
\sphinxhref{https://docs.scipy.org/doc/scipy/reference/optimize.html}{scipy.optimize}
methods.

\index{source@\spxentry{source}}\index{code@\spxentry{code}}\ignorespaces 

\subparagraph{Source Code}
\label{\detokenize{xsrst/numeric_optimize_fun_xam_py:source-code}}\label{\detokenize{xsrst/numeric_optimize_fun_xam_py:numeric-optimize-fun-xam-py-source-code}}\label{\detokenize{xsrst/numeric_optimize_fun_xam_py:index-4}}
\begin{sphinxVerbatim}[commandchars=\\\{\}]
\PYG{k}{def} \PYG{n+nf}{optimize\PYGZus{}fun\PYGZus{}xam}\PYG{p}{(}\PYG{p}{)} \PYG{p}{:}
    \PYG{c+c1}{\PYGZsh{}}
    \PYG{k+kn}{import} \PYG{n+nn}{numpy}
    \PYG{k+kn}{import} \PYG{n+nn}{cppad\PYGZus{}py}
    \PYG{k+kn}{import} \PYG{n+nn}{scipy}\PYG{n+nn}{.}\PYG{n+nn}{optimize}
    \PYG{k+kn}{from} \PYG{n+nn}{optimize\PYGZus{}fun\PYGZus{}class} \PYG{k+kn}{import} \PYG{n}{optimize\PYGZus{}fun\PYGZus{}class}
    \PYG{c+c1}{\PYGZsh{}}
    \PYG{n}{ok} \PYG{o}{=} \PYG{k+kc}{True}
    \PYG{c+c1}{\PYGZsh{}}
    \PYG{k}{def} \PYG{n+nf}{a\PYGZus{}objective}\PYG{p}{(}\PYG{n}{ax}\PYG{p}{)}  \PYG{p}{:}
        \PYG{k}{return} \PYG{n}{ax}\PYG{p}{[}\PYG{l+m+mi}{0}\PYG{p}{]} \PYG{o}{*} \PYG{n}{ax}\PYG{p}{[}\PYG{l+m+mi}{3}\PYG{p}{]} \PYG{o}{*} \PYG{p}{(} \PYG{n}{ax}\PYG{p}{[}\PYG{l+m+mi}{0}\PYG{p}{]} \PYG{o}{+} \PYG{n}{ax}\PYG{p}{[}\PYG{l+m+mi}{1}\PYG{p}{]} \PYG{o}{+} \PYG{n}{ax}\PYG{p}{[}\PYG{l+m+mi}{2}\PYG{p}{]} \PYG{p}{)} \PYG{o}{+} \PYG{n}{ax}\PYG{p}{[}\PYG{l+m+mi}{2}\PYG{p}{]}
    \PYG{k}{def} \PYG{n+nf}{a\PYGZus{}constraint}\PYG{p}{(}\PYG{n}{ax}\PYG{p}{)} \PYG{p}{:}
        \PYG{k}{return} \PYG{p}{[} \PYG{n}{numpy}\PYG{o}{.}\PYG{n}{prod}\PYG{p}{(}\PYG{n}{ax}\PYG{p}{)}\PYG{p}{,} \PYG{n}{numpy}\PYG{o}{.}\PYG{n}{sum}\PYG{p}{(} \PYG{n}{ax} \PYG{o}{*} \PYG{n}{ax}\PYG{p}{)} \PYG{p}{]}
    \PYG{c+c1}{\PYGZsh{}}
    \PYG{c+c1}{\PYGZsh{} objective\PYGZus{}ad}
    \PYG{n}{x} \PYG{o}{=} \PYG{n}{numpy}\PYG{o}{.}\PYG{n}{array}\PYG{p}{(} \PYG{p}{[} \PYG{l+m+mf}{1.0}\PYG{p}{,} \PYG{l+m+mf}{2.0}\PYG{p}{,} \PYG{l+m+mf}{3.0}\PYG{p}{,} \PYG{l+m+mf}{4.0} \PYG{p}{]} \PYG{p}{)}
    \PYG{n}{ax} \PYG{o}{=} \PYG{n}{cppad\PYGZus{}py}\PYG{o}{.}\PYG{n}{independent}\PYG{p}{(}\PYG{n}{x}\PYG{p}{)}
    \PYG{n}{ay} \PYG{o}{=} \PYG{n}{numpy}\PYG{o}{.}\PYG{n}{array}\PYG{p}{(} \PYG{p}{[}\PYG{n}{a\PYGZus{}objective}\PYG{p}{(}\PYG{n}{ax}\PYG{p}{)}\PYG{p}{]} \PYG{p}{)}
    \PYG{n}{objective\PYGZus{}ad} \PYG{o}{=} \PYG{n}{cppad\PYGZus{}py}\PYG{o}{.}\PYG{n}{d\PYGZus{}fun}\PYG{p}{(}\PYG{n}{ax}\PYG{p}{,} \PYG{n}{ay}\PYG{p}{)}
    \PYG{c+c1}{\PYGZsh{}}
    \PYG{c+c1}{\PYGZsh{} constraint\PYGZus{}ad}
    \PYG{n}{ax} \PYG{o}{=} \PYG{n}{cppad\PYGZus{}py}\PYG{o}{.}\PYG{n}{independent}\PYG{p}{(}\PYG{n}{x}\PYG{p}{)}
    \PYG{n}{ay} \PYG{o}{=} \PYG{n}{numpy}\PYG{o}{.}\PYG{n}{array}\PYG{p}{(} \PYG{n}{a\PYGZus{}constraint}\PYG{p}{(}\PYG{n}{ax}\PYG{p}{)} \PYG{p}{)}
    \PYG{n}{constraint\PYGZus{}ad} \PYG{o}{=} \PYG{n}{cppad\PYGZus{}py}\PYG{o}{.}\PYG{n}{d\PYGZus{}fun}\PYG{p}{(}\PYG{n}{ax}\PYG{p}{,} \PYG{n}{ay}\PYG{p}{)}
    \PYG{c+c1}{\PYGZsh{}}
    \PYG{c+c1}{\PYGZsh{} optimize\PYGZus{}fun}
    \PYG{n}{optimize\PYGZus{}fun} \PYG{o}{=} \PYG{n}{optimize\PYGZus{}fun\PYGZus{}class}\PYG{p}{(}\PYG{n}{objective\PYGZus{}ad}\PYG{p}{,} \PYG{n}{constraint\PYGZus{}ad}\PYG{p}{)}
    \PYG{c+c1}{\PYGZsh{}}
    \PYG{c+c1}{\PYGZsh{} constraints}
    \PYG{n}{lower\PYGZus{}bound} \PYG{o}{=} \PYG{p}{[}      \PYG{l+m+mf}{25.0}\PYG{p}{,} \PYG{l+m+mf}{40.0} \PYG{p}{]}
    \PYG{n}{upper\PYGZus{}bound} \PYG{o}{=} \PYG{p}{[} \PYG{n}{numpy}\PYG{o}{.}\PYG{n}{inf}\PYG{p}{,} \PYG{l+m+mf}{40.0} \PYG{p}{]}
    \PYG{n}{nonlinear\PYGZus{}constraint} \PYG{o}{=} \PYG{n}{scipy}\PYG{o}{.}\PYG{n}{optimize}\PYG{o}{.}\PYG{n}{NonlinearConstraint}\PYG{p}{(}
        \PYG{n}{optimize\PYGZus{}fun}\PYG{o}{.}\PYG{n}{constraint\PYGZus{}fun}\PYG{p}{,}
        \PYG{n}{lower\PYGZus{}bound}\PYG{p}{,}
        \PYG{n}{upper\PYGZus{}bound}\PYG{p}{,}
        \PYG{n}{jac}           \PYG{o}{=} \PYG{n}{optimize\PYGZus{}fun}\PYG{o}{.}\PYG{n}{constraint\PYGZus{}jac}\PYG{p}{,}
        \PYG{n}{hess}          \PYG{o}{=} \PYG{n}{optimize\PYGZus{}fun}\PYG{o}{.}\PYG{n}{constraint\PYGZus{}hess}\PYG{p}{,}
        \PYG{n}{keep\PYGZus{}feasible} \PYG{o}{=} \PYG{k+kc}{False}
    \PYG{p}{)}
    \PYG{n}{constraints}       \PYG{o}{=} \PYG{p}{[}\PYG{n}{nonlinear\PYGZus{}constraint}\PYG{p}{]}
    \PYG{c+c1}{\PYGZsh{}}
    \PYG{c+c1}{\PYGZsh{} bounds}
    \PYG{n}{lower\PYGZus{}bound} \PYG{o}{=} \PYG{l+m+mi}{4} \PYG{o}{*} \PYG{p}{[} \PYG{l+m+mf}{1.0} \PYG{p}{]}
    \PYG{n}{upper\PYGZus{}bound} \PYG{o}{=} \PYG{l+m+mi}{4} \PYG{o}{*} \PYG{p}{[} \PYG{l+m+mf}{5.0} \PYG{p}{]}
    \PYG{n}{bounds} \PYG{o}{=} \PYG{n}{scipy}\PYG{o}{.}\PYG{n}{optimize}\PYG{o}{.}\PYG{n}{Bounds}\PYG{p}{(}
        \PYG{n}{lower\PYGZus{}bound}\PYG{p}{,}
        \PYG{n}{upper\PYGZus{}bound}\PYG{p}{,}
        \PYG{n}{keep\PYGZus{}feasible} \PYG{o}{=} \PYG{k+kc}{False}
    \PYG{p}{)}
    \PYG{c+c1}{\PYGZsh{}}
    \PYG{c+c1}{\PYGZsh{} start\PYGZus{}point}
    \PYG{n}{start\PYGZus{}point} \PYG{o}{=} \PYG{p}{[} \PYG{l+m+mf}{1.0}\PYG{p}{,} \PYG{l+m+mf}{5.0}\PYG{p}{,} \PYG{l+m+mf}{5.0}\PYG{p}{,} \PYG{l+m+mf}{1.0} \PYG{p}{]}
    \PYG{c+c1}{\PYGZsh{}}
    \PYG{n}{options} \PYG{o}{=} \PYG{p}{\PYGZob{}}
                             \PYG{l+s+s1}{\PYGZsq{}}\PYG{l+s+s1}{gtol}\PYG{l+s+s1}{\PYGZsq{}} \PYG{p}{:} \PYG{l+m+mf}{1e\PYGZhy{}8}\PYG{p}{,}
                             \PYG{l+s+s1}{\PYGZsq{}}\PYG{l+s+s1}{xtol}\PYG{l+s+s1}{\PYGZsq{}} \PYG{p}{:} \PYG{l+m+mf}{1e\PYGZhy{}8}\PYG{p}{,}
                      \PYG{l+s+s1}{\PYGZsq{}}\PYG{l+s+s1}{barrier\PYGZus{}tol}\PYG{l+s+s1}{\PYGZsq{}} \PYG{p}{:} \PYG{l+m+mf}{1e\PYGZhy{}8}\PYG{p}{,}
                \PYG{l+s+s1}{\PYGZsq{}}\PYG{l+s+s1}{initial\PYGZus{}tr\PYGZus{}radius}\PYG{l+s+s1}{\PYGZsq{}} \PYG{p}{:} \PYG{l+m+mf}{1.0}\PYG{p}{,}
           \PYG{l+s+s1}{\PYGZsq{}}\PYG{l+s+s1}{initial\PYGZus{}constr\PYGZus{}penalty}\PYG{l+s+s1}{\PYGZsq{}} \PYG{p}{:} \PYG{l+m+mf}{1.0}\PYG{p}{,}
        \PYG{l+s+s1}{\PYGZsq{}}\PYG{l+s+s1}{initial\PYGZus{}barrier\PYGZus{}tolerance}\PYG{l+s+s1}{\PYGZsq{}} \PYG{p}{:} \PYG{l+m+mf}{0.1}\PYG{p}{,}
        \PYG{l+s+s1}{\PYGZsq{}}\PYG{l+s+s1}{initial\PYGZus{}barrier\PYGZus{}parameter}\PYG{l+s+s1}{\PYGZsq{}} \PYG{p}{:} \PYG{l+m+mf}{0.1}\PYG{p}{,}
             \PYG{l+s+s1}{\PYGZsq{}}\PYG{l+s+s1}{factorization\PYGZus{}method}\PYG{l+s+s1}{\PYGZsq{}} \PYG{p}{:} \PYG{k+kc}{None}\PYG{p}{,}
             \PYG{l+s+s1}{\PYGZsq{}}\PYG{l+s+s1}{finite\PYGZus{}diff\PYGZus{}rel\PYGZus{}step}\PYG{l+s+s1}{\PYGZsq{}} \PYG{p}{:} \PYG{k+kc}{None}\PYG{p}{,}
                          \PYG{l+s+s1}{\PYGZsq{}}\PYG{l+s+s1}{maxiter}\PYG{l+s+s1}{\PYGZsq{}} \PYG{p}{:} \PYG{l+m+mi}{50}\PYG{p}{,}
                          \PYG{l+s+s1}{\PYGZsq{}}\PYG{l+s+s1}{verbose}\PYG{l+s+s1}{\PYGZsq{}} \PYG{p}{:} \PYG{l+m+mi}{0}\PYG{p}{,}
    \PYG{p}{\PYGZcb{}}
    \PYG{c+c1}{\PYGZsh{}}
    \PYG{n}{result} \PYG{o}{=} \PYG{n}{scipy}\PYG{o}{.}\PYG{n}{optimize}\PYG{o}{.}\PYG{n}{minimize}\PYG{p}{(}
        \PYG{n}{optimize\PYGZus{}fun}\PYG{o}{.}\PYG{n}{objective\PYGZus{}fun}\PYG{p}{,}
        \PYG{n}{start\PYGZus{}point}\PYG{p}{,}
        \PYG{n}{method}      \PYG{o}{=} \PYG{l+s+s1}{\PYGZsq{}}\PYG{l+s+s1}{trust\PYGZhy{}constr}\PYG{l+s+s1}{\PYGZsq{}}\PYG{p}{,}
        \PYG{n}{jac}         \PYG{o}{=} \PYG{n}{optimize\PYGZus{}fun}\PYG{o}{.}\PYG{n}{objective\PYGZus{}grad}\PYG{p}{,}
        \PYG{n}{hess}        \PYG{o}{=} \PYG{n}{optimize\PYGZus{}fun}\PYG{o}{.}\PYG{n}{objective\PYGZus{}hess}\PYG{p}{,}
        \PYG{n}{constraints} \PYG{o}{=} \PYG{n}{constraints}\PYG{p}{,}
        \PYG{n}{options}     \PYG{o}{=} \PYG{n}{options}\PYG{p}{,}
        \PYG{n}{bounds}      \PYG{o}{=} \PYG{n}{bounds}\PYG{p}{,}
    \PYG{p}{)}
    \PYG{n}{ok} \PYG{o}{=} \PYG{n}{ok} \PYG{o+ow}{and} \PYG{n}{result}\PYG{o}{.}\PYG{n}{success}
    \PYG{c+c1}{\PYGZsh{}}
    \PYG{n}{optimal\PYGZus{}point} \PYG{o}{=} \PYG{n}{result}\PYG{o}{.}\PYG{n}{x}
    \PYG{n}{check}         \PYG{o}{=} \PYG{p}{[} \PYG{l+m+mf}{1.00000000}\PYG{p}{,} \PYG{l+m+mf}{4.74299963}\PYG{p}{,} \PYG{l+m+mf}{3.82114998}\PYG{p}{,} \PYG{l+m+mf}{1.37940829} \PYG{p}{]}
    \PYG{n}{rel\PYGZus{}error}     \PYG{o}{=} \PYG{n}{optimal\PYGZus{}point} \PYG{o}{/} \PYG{n}{check} \PYG{o}{\PYGZhy{}} \PYG{l+m+mf}{1.0}
    \PYG{n}{ok} \PYG{o}{=} \PYG{n}{ok} \PYG{o+ow}{and} \PYG{n}{numpy}\PYG{o}{.}\PYG{n}{all}\PYG{p}{(} \PYG{n+nb}{abs}\PYG{p}{(}\PYG{n}{rel\PYGZus{}error}\PYG{p}{)} \PYG{o}{\PYGZlt{}} \PYG{l+m+mf}{1e\PYGZhy{}5} \PYG{p}{)}
    \PYG{k}{return} \PYG{n}{ok}
\end{sphinxVerbatim}


\bigskip\hrule\bigskip


xsrst input file: \sphinxcode{\sphinxupquote{example/python/numeric/optimize\_fun\_xam.py}}

\index{example@\spxentry{example}}\ignorespaces 

\subparagraph{Example}
\label{\detokenize{xsrst/numeric_optimize_fun_class:example}}\label{\detokenize{xsrst/numeric_optimize_fun_class:numeric-optimize-fun-class-example}}\label{\detokenize{xsrst/numeric_optimize_fun_class:index-12}}
{\hyperref[\detokenize{xsrst/numeric_optimize_fun_xam_py:id1}]{\sphinxcrossref{\DUrole{std,std-ref}{numeric\_optimize\_fun\_xam\_py}}}}

\index{source@\spxentry{source}}\index{code@\spxentry{code}}\ignorespaces 

\subparagraph{Source Code}
\label{\detokenize{xsrst/numeric_optimize_fun_class:source-code}}\label{\detokenize{xsrst/numeric_optimize_fun_class:numeric-optimize-fun-class-source-code}}\label{\detokenize{xsrst/numeric_optimize_fun_class:index-13}}
\begin{sphinxVerbatim}[commandchars=\\\{\}]
\PYG{k+kn}{import} \PYG{n+nn}{numpy}
\PYG{k+kn}{import} \PYG{n+nn}{cppad\PYGZus{}py}

\PYG{k}{class} \PYG{n+nc}{optimize\PYGZus{}fun\PYGZus{}class} \PYG{p}{:}
    \PYG{k}{def} \PYG{n+nf+fm}{\PYGZus{}\PYGZus{}init\PYGZus{}\PYGZus{}}\PYG{p}{(}\PYG{n+nb+bp}{self}\PYG{p}{,} \PYG{n}{objective\PYGZus{}ad}\PYG{p}{,} \PYG{n}{constraint\PYGZus{}ad}\PYG{o}{=}\PYG{k+kc}{None}\PYG{p}{)} \PYG{p}{:}
        \PYG{n+nb+bp}{self}\PYG{o}{.}\PYG{n}{objective\PYGZus{}ad}  \PYG{o}{=} \PYG{n}{objective\PYGZus{}ad}
        \PYG{n+nb+bp}{self}\PYG{o}{.}\PYG{n}{constraint\PYGZus{}ad} \PYG{o}{=} \PYG{n}{constraint\PYGZus{}ad}
    \PYG{c+c1}{\PYGZsh{}}
    \PYG{k}{def} \PYG{n+nf}{objective\PYGZus{}fun}\PYG{p}{(}\PYG{n+nb+bp}{self}\PYG{p}{,} \PYG{n}{x}\PYG{p}{)} \PYG{p}{:}
        \PYG{c+c1}{\PYGZsh{} objective as a vector}
        \PYG{n}{y} \PYG{o}{=} \PYG{n+nb+bp}{self}\PYG{o}{.}\PYG{n}{objective\PYGZus{}ad}\PYG{o}{.}\PYG{n}{forward}\PYG{p}{(}\PYG{l+m+mi}{0}\PYG{p}{,} \PYG{n}{x}\PYG{p}{)}
        \PYG{k}{return} \PYG{n}{y}\PYG{p}{[}\PYG{l+m+mi}{0}\PYG{p}{]}
    \PYG{c+c1}{\PYGZsh{}}
    \PYG{k}{def} \PYG{n+nf}{objective\PYGZus{}grad}\PYG{p}{(}\PYG{n+nb+bp}{self}\PYG{p}{,} \PYG{n}{x}\PYG{p}{)} \PYG{p}{:}
        \PYG{c+c1}{\PYGZsh{} Jacobian as a matrix}
        \PYG{n}{J} \PYG{o}{=} \PYG{n+nb+bp}{self}\PYG{o}{.}\PYG{n}{objective\PYGZus{}ad}\PYG{o}{.}\PYG{n}{jacobian}\PYG{p}{(}\PYG{n}{x}\PYG{p}{)}
        \PYG{c+c1}{\PYGZsh{} change to a vector}
        \PYG{k}{return} \PYG{n}{J}\PYG{o}{.}\PYG{n}{flatten}\PYG{p}{(}\PYG{p}{)}
    \PYG{c+c1}{\PYGZsh{}}
    \PYG{k}{def} \PYG{n+nf}{objective\PYGZus{}hess}\PYG{p}{(}\PYG{n+nb+bp}{self}\PYG{p}{,} \PYG{n}{x}\PYG{p}{)} \PYG{p}{:}
        \PYG{n}{w} \PYG{o}{=} \PYG{n}{numpy}\PYG{o}{.}\PYG{n}{array}\PYG{p}{(} \PYG{p}{[} \PYG{l+m+mf}{1.0} \PYG{p}{]} \PYG{p}{)}
        \PYG{n}{H} \PYG{o}{=} \PYG{n+nb+bp}{self}\PYG{o}{.}\PYG{n}{objective\PYGZus{}ad}\PYG{o}{.}\PYG{n}{hessian}\PYG{p}{(}\PYG{n}{x}\PYG{p}{,} \PYG{n}{w}\PYG{p}{)}
        \PYG{k}{return} \PYG{n}{H}
    \PYG{c+c1}{\PYGZsh{}}
    \PYG{k}{def} \PYG{n+nf}{constraint\PYGZus{}fun}\PYG{p}{(}\PYG{n+nb+bp}{self}\PYG{p}{,} \PYG{n}{x}\PYG{p}{)} \PYG{p}{:}
        \PYG{k}{return} \PYG{n+nb+bp}{self}\PYG{o}{.}\PYG{n}{constraint\PYGZus{}ad}\PYG{o}{.}\PYG{n}{forward}\PYG{p}{(}\PYG{l+m+mi}{0}\PYG{p}{,} \PYG{n}{x}\PYG{p}{)}
    \PYG{c+c1}{\PYGZsh{}}
    \PYG{k}{def} \PYG{n+nf}{constraint\PYGZus{}jac}\PYG{p}{(}\PYG{n+nb+bp}{self}\PYG{p}{,} \PYG{n}{x}\PYG{p}{)} \PYG{p}{:}
        \PYG{c+c1}{\PYGZsh{} Jacobian as a matrix}
        \PYG{n}{J} \PYG{o}{=} \PYG{n+nb+bp}{self}\PYG{o}{.}\PYG{n}{constraint\PYGZus{}ad}\PYG{o}{.}\PYG{n}{jacobian}\PYG{p}{(}\PYG{n}{x}\PYG{p}{)}
        \PYG{k}{return} \PYG{n}{J}
    \PYG{c+c1}{\PYGZsh{}}
    \PYG{k}{def} \PYG{n+nf}{constraint\PYGZus{}hess}\PYG{p}{(}\PYG{n+nb+bp}{self}\PYG{p}{,} \PYG{n}{x}\PYG{p}{,} \PYG{n}{v}\PYG{p}{)} \PYG{p}{:}
        \PYG{n}{H} \PYG{o}{=} \PYG{n+nb+bp}{self}\PYG{o}{.}\PYG{n}{constraint\PYGZus{}ad}\PYG{o}{.}\PYG{n}{hessian}\PYG{p}{(}\PYG{n}{x}\PYG{p}{,} \PYG{n}{v}\PYG{p}{)}
        \PYG{k}{return} \PYG{n}{H}
\end{sphinxVerbatim}


\bigskip\hrule\bigskip


xsrst input file: \sphinxcode{\sphinxupquote{example/python/numeric/optimize\_fun\_class.py}}


\subparagraph{numeric\_seirwd\_model}
\label{\detokenize{xsrst/numeric_seirwd_model:numeric-seirwd-model}}\label{\detokenize{xsrst/numeric_seirwd_model::doc}}
\index{numeric\_seirwd\_model@\spxentry{numeric\_seirwd\_model}}\index{susceptible@\spxentry{susceptible}}\index{exposed@\spxentry{exposed}}\index{infectious@\spxentry{infectious}}\index{recovered@\spxentry{recovered}}\index{death@\spxentry{death}}\index{model@\spxentry{model}}\ignorespaces 

\subparagraph{2.6 A Susceptible Exposed Infectious Recovered and Death Model}
\label{\detokenize{xsrst/numeric_seirwd_model:a-susceptible-exposed-infectious-recovered-and-death-model}}\label{\detokenize{xsrst/numeric_seirwd_model:id1}}\begin{itemize}
\item {} 
{\hyperref[\detokenize{xsrst/numeric_seirwd_model:numeric-seirwd-model-purpose}]{\sphinxcrossref{\DUrole{std,std-ref}{Purpose}}}}

\item {} 
{\hyperref[\detokenize{xsrst/numeric_seirwd_model:numeric-seirwd-model-syntax}]{\sphinxcrossref{\DUrole{std,std-ref}{Syntax}}}}

\item {} 
{\hyperref[\detokenize{xsrst/numeric_seirwd_model:numeric-seirwd-model-notation}]{\sphinxcrossref{\DUrole{std,std-ref}{Notation}}}}

\item {} 
{\hyperref[\detokenize{xsrst/numeric_seirwd_model:numeric-seirwd-model-ode}]{\sphinxcrossref{\DUrole{std,std-ref}{ODE}}}}

\item {} 
{\hyperref[\detokenize{xsrst/numeric_seirwd_model:numeric-seirwd-model-method}]{\sphinxcrossref{\DUrole{std,std-ref}{method}}}}

\item {} 
{\hyperref[\detokenize{xsrst/numeric_seirwd_model:numeric-seirwd-model-t-all}]{\sphinxcrossref{\DUrole{std,std-ref}{t\_all}}}}

\item {} \begin{description}
\item[{{\hyperref[\detokenize{xsrst/numeric_seirwd_model:numeric-seirwd-model-p-all}]{\sphinxcrossref{\DUrole{std,std-ref}{p\_all}}}}}] \leavevmode\begin{itemize}
\item {} 
{\hyperref[\detokenize{xsrst/numeric_seirwd_model:numeric-seirwd-model-p-all-alpha}]{\sphinxcrossref{\DUrole{std,std-ref}{alpha}}}}

\end{itemize}

\end{description}

\item {} 
{\hyperref[\detokenize{xsrst/numeric_seirwd_model:numeric-seirwd-model-initial}]{\sphinxcrossref{\DUrole{std,std-ref}{initial}}}}

\item {} 
{\hyperref[\detokenize{xsrst/numeric_seirwd_model:numeric-seirwd-model-n-step}]{\sphinxcrossref{\DUrole{std,std-ref}{n\_step}}}}

\item {} 
{\hyperref[\detokenize{xsrst/numeric_seirwd_model:numeric-seirwd-model-seirwd-all}]{\sphinxcrossref{\DUrole{std,std-ref}{seirwd\_all}}}}

\item {} 
{\hyperref[\detokenize{xsrst/numeric_seirwd_model:numeric-seirwd-model-conservation-of-mass}]{\sphinxcrossref{\DUrole{std,std-ref}{Conservation of Mass}}}}

\item {} 
{\hyperref[\detokenize{xsrst/numeric_seirwd_model:numeric-seirwd-model-example}]{\sphinxcrossref{\DUrole{std,std-ref}{Example}}}}

\item {} 
{\hyperref[\detokenize{xsrst/numeric_seirwd_model:numeric-seirwd-model-source-code}]{\sphinxcrossref{\DUrole{std,std-ref}{Source Code}}}}

\end{itemize}

\index{purpose@\spxentry{purpose}}\ignorespaces 

\subparagraph{Purpose}
\label{\detokenize{xsrst/numeric_seirwd_model:purpose}}\label{\detokenize{xsrst/numeric_seirwd_model:numeric-seirwd-model-purpose}}\label{\detokenize{xsrst/numeric_seirwd_model:index-1}}
This routine can be used with \sphinxcode{\sphinxupquote{ad\_double}} .

\index{syntax@\spxentry{syntax}}\ignorespaces 

\subparagraph{Syntax}
\label{\detokenize{xsrst/numeric_seirwd_model:syntax}}\label{\detokenize{xsrst/numeric_seirwd_model:numeric-seirwd-model-syntax}}\label{\detokenize{xsrst/numeric_seirwd_model:index-2}}
\begin{DUlineblock}{0em}
\item[] \sphinxstyleemphasis{seirwd\_all} =  \sphinxcode{\sphinxupquote{seirwd\_model}} (
\item[]         \sphinxstyleemphasis{method} , \sphinxstyleemphasis{t\_all} , \sphinxstyleemphasis{p\_all} , \sphinxstyleemphasis{initial} , \sphinxstyleemphasis{n\_step} = 1
\item[] )
\end{DUlineblock}

\index{notation@\spxentry{notation}}\ignorespaces 

\subparagraph{Notation}
\label{\detokenize{xsrst/numeric_seirwd_model:notation}}\label{\detokenize{xsrst/numeric_seirwd_model:numeric-seirwd-model-notation}}\label{\detokenize{xsrst/numeric_seirwd_model:index-3}}

\begin{savenotes}\sphinxattablestart
\centering
\begin{tabular}[t]{|\X{4}{38}|\X{34}{38}|}
\hline

\(S(t)\)
&
size of the Susceptible group
\\
\hline
\(E(t)\)
&
size of the Exposed group
\\
\hline
\(I(t)\)
&
size of the Infectious group
\\
\hline
\(R(t)\)
&
size Recovered group
\\
\hline
\(W(t)\)
&
size of the group that will die
\\
\hline
\(D(t)\)
&
size of the group that has died
\\
\hline
\(\alpha(t)\)
&
infectious group size exponent
\\
\hline
\(\beta(t)\)
&
infectious rate
\\
\hline
\(\sigma(t)\)
&
incubation rate
\\
\hline
\(\gamma(t)\)
&
recovery rate
\\
\hline
\(\xi(t)\)
&
loss of immunity rate
\\
\hline
\(\chi(t)\)
&
excess mortality rate
\\
\hline
\(\delta(t)\)
&
delay between infectious and death
\\
\hline
\end{tabular}
\par
\sphinxattableend\end{savenotes}

\index{ode@\spxentry{ode}}\ignorespaces 

\subparagraph{ODE}
\label{\detokenize{xsrst/numeric_seirwd_model:ode}}\label{\detokenize{xsrst/numeric_seirwd_model:numeric-seirwd-model-ode}}\label{\detokenize{xsrst/numeric_seirwd_model:index-4}}
The ordinary differential equation for this model is:
\begin{equation*}
\begin{split}\begin{array}{rcll}
\dot{S} & = & - \beta  S I^\alpha  & + \xi    R             \\
\dot{E} & = & + \beta  S I^\alpha  & - \sigma E             \\
\dot{I} & = & + \sigma E           & - ( \gamma + \chi )  I \\
\dot{R} & = & + \gamma I           & - \xi    R             \\
\dot{W} & = & + \chi   I           & - \delta W             \\
\dot{D} & = & + \delta W           &
\end{array}\end{split}
\end{equation*}
where we dropped the time dependence in the equations above.
This model does not account for death by other causes.
It is similar to the standard
\sphinxhref{https://www.idmod.org/docs/hiv/model-seir.html}{SEIRS} model
with the following differences:
\begin{enumerate}
\sphinxsetlistlabels{\arabic}{enumi}{enumii}{}{.}%
\item {} 
The total population \(N\) is not included in this model,
so the units for \(\beta\) are different.

\item {} 
This model tracks death due to the condition using the compartments W and D.

\end{enumerate}

\index{method@\spxentry{method}}\ignorespaces 

\subparagraph{method}
\label{\detokenize{xsrst/numeric_seirwd_model:method}}\label{\detokenize{xsrst/numeric_seirwd_model:numeric-seirwd-model-method}}\label{\detokenize{xsrst/numeric_seirwd_model:index-5}}
This \sphinxcode{\sphinxupquote{str}} must be either \sphinxcode{\sphinxupquote{runge4}} or \sphinxcode{\sphinxupquote{rosen3}} .
It determines if {\hyperref[\detokenize{xsrst/numeric_runge4_step:id1}]{\sphinxcrossref{\DUrole{std,std-ref}{runge4\_step}}}} or
{\hyperref[\detokenize{xsrst/numeric_rosen3_step:id1}]{\sphinxcrossref{\DUrole{std,std-ref}{rosen3\_step}}}} is used to solve the ODE.

\index{t\_all@\spxentry{t\_all}}\ignorespaces 

\subparagraph{t\_all}
\label{\detokenize{xsrst/numeric_seirwd_model:t-all}}\label{\detokenize{xsrst/numeric_seirwd_model:numeric-seirwd-model-t-all}}\label{\detokenize{xsrst/numeric_seirwd_model:index-6}}
The argument \sphinxstyleemphasis{t\_all} is a vector that is monotone
increasing or decreasing.
The type of its elements can be \sphinxcode{\sphinxupquote{float}} or \sphinxcode{\sphinxupquote{a\_double}} .
The smaller the spacing between time points, the more accurate
the approximation is.
Note two points can be equal; i.e., no zero spacing.
We call \sphinxstyleemphasis{t\_all} {[}0{]} the initial time and
\sphinxstyleemphasis{t\_all} {[}\sphinxhyphen{}1{]} the final time.

\index{p\_all@\spxentry{p\_all}}\ignorespaces 

\subparagraph{p\_all}
\label{\detokenize{xsrst/numeric_seirwd_model:p-all}}\label{\detokenize{xsrst/numeric_seirwd_model:numeric-seirwd-model-p-all}}\label{\detokenize{xsrst/numeric_seirwd_model:index-7}}
This argument is a list of dictionaries.
The i\sphinxhyphen{}th element of the list has the following \sphinxstyleemphasis{key} ,
\sphinxstyleemphasis{value} pairs:


\begin{savenotes}\sphinxattablestart
\centering
\begin{tabular}[t]{|\X{7}{15}|\X{8}{15}|}
\hline

\sphinxstyleemphasis{key}
&
\sphinxstyleemphasis{value}
\\
\hline
\sphinxcode{\sphinxupquote{\textquotesingle{}alpha\textquotesingle{}}}
&
\(\alpha( t_i )\)
\\
\hline
\sphinxcode{\sphinxupquote{\textquotesingle{}beta\textquotesingle{}}}
&
\(\beta( t_i )\)
\\
\hline
\sphinxcode{\sphinxupquote{\textquotesingle{}sigma\textquotesingle{}}}
&
\(\sigma( t_i )\)
\\
\hline
\sphinxcode{\sphinxupquote{\textquotesingle{}gamma\textquotesingle{}}}
&
\(\gamma( t_i )\)
\\
\hline
\sphinxcode{\sphinxupquote{\textquotesingle{}xi\textquotesingle{}}}
&
\(\xi( t_i )\)
\\
\hline
\sphinxcode{\sphinxupquote{\textquotesingle{}chi\textquotesingle{}}}
&
\(\chi( t_i )\)
\\
\hline
\sphinxcode{\sphinxupquote{\textquotesingle{}delta\textquotesingle{}}}
&
\(\delta( t_i )\)
\\
\hline
\end{tabular}
\par
\sphinxattableend\end{savenotes}

where \sphinxstyleemphasis{t\_i} is the time \sphinxstyleemphasis{t\_all} {[} \sphinxstyleemphasis{i} {]} .
The type of \sphinxstyleemphasis{value} can be \sphinxcode{\sphinxupquote{float}} or \sphinxcode{\sphinxupquote{a\_double}} .
Each of the these parameters will be linearly interpolated
for times between the those in \sphinxstyleemphasis{t\_all} .

\index{alpha@\spxentry{alpha}}\ignorespaces 

\subparagraph{alpha}
\label{\detokenize{xsrst/numeric_seirwd_model:alpha}}\label{\detokenize{xsrst/numeric_seirwd_model:numeric-seirwd-model-p-all-alpha}}\label{\detokenize{xsrst/numeric_seirwd_model:index-8}}
There is a special restriction that \(\alpha(t)\) must be
constant; i.e. \(\alpha( t_i ) = \alpha( t_0 )\) for
all \(i\).
This is because talking the derivative of \(I^\alpha\)
respect to \(I\) has a special representation when \(\alpha = 1\).
Using a special representation for that case would not work with AD
unless \(\alpha\) is constant.

\index{initial@\spxentry{initial}}\ignorespaces 

\subparagraph{initial}
\label{\detokenize{xsrst/numeric_seirwd_model:initial}}\label{\detokenize{xsrst/numeric_seirwd_model:numeric-seirwd-model-initial}}\label{\detokenize{xsrst/numeric_seirwd_model:index-9}}
is a vector of length four containing the initial values for
S, E, I, R, W, D in that order.
The type of its elements can be \sphinxcode{\sphinxupquote{float}} or \sphinxcode{\sphinxupquote{a\_double}} .

\index{n\_step@\spxentry{n\_step}}\ignorespaces 

\subparagraph{n\_step}
\label{\detokenize{xsrst/numeric_seirwd_model:n-step}}\label{\detokenize{xsrst/numeric_seirwd_model:numeric-seirwd-model-n-step}}\label{\detokenize{xsrst/numeric_seirwd_model:index-10}}
This is the number of numerical integration steps to use for each
time interval in the \sphinxstyleemphasis{t\_all} array.
It must be an \sphinxcode{\sphinxupquote{int}} greater or equal one.
The larger \sphinxstyleemphasis{n\_step} the more computational effort and the more
accurate the solution.  The default value is for \sphinxstyleemphasis{n\_step} is one.

\index{seirwd\_all@\spxentry{seirwd\_all}}\ignorespaces 

\subparagraph{seirwd\_all}
\label{\detokenize{xsrst/numeric_seirwd_model:seirwd-all}}\label{\detokenize{xsrst/numeric_seirwd_model:numeric-seirwd-model-seirwd-all}}\label{\detokenize{xsrst/numeric_seirwd_model:index-11}}
The return value \sphinxstyleemphasis{seirwd\_all} is a numpy matrix with row dimension
equal to the number of elements in \sphinxstyleemphasis{t\_all} and column dimension
equal to six. The value \sphinxstyleemphasis{seirwd\_all} {[} \sphinxstyleemphasis{i} , \sphinxstyleemphasis{j} {]} is the
approximate solution for the j\sphinxhyphen{}th compartment at time \sphinxstyleemphasis{t\_all} {[} \sphinxstyleemphasis{i} {]} .
The compartments have the same order as in \sphinxstyleemphasis{initial} and
\sphinxcode{\sphinxupquote{seirwd}} {[}0, \sphinxstyleemphasis{:} is equal to \sphinxstyleemphasis{initial} .
The sequence of floating point operations only depends on \sphinxstyleemphasis{t\_all}
and the operations used to compute \sphinxstyleemphasis{p\_fun} .

\index{conservation@\spxentry{conservation}}\index{mass@\spxentry{mass}}\ignorespaces 

\subparagraph{Conservation of Mass}
\label{\detokenize{xsrst/numeric_seirwd_model:conservation-of-mass}}\label{\detokenize{xsrst/numeric_seirwd_model:numeric-seirwd-model-conservation-of-mass}}\label{\detokenize{xsrst/numeric_seirwd_model:index-12}}
Note that the sum of S, E, I, R, W, and D should be constant; i.e.,
up to numerical accuracy, it not depend on time.


\subparagraph{numeric\_seirwd\_model\_xam\_py}
\label{\detokenize{xsrst/numeric_seirwd_model_xam_py:numeric-seirwd-model-xam-py}}\label{\detokenize{xsrst/numeric_seirwd_model_xam_py::doc}}
\index{numeric\_seirwd\_model\_xam\_py@\spxentry{numeric\_seirwd\_model\_xam\_py}}\index{example@\spxentry{example}}\index{using@\spxentry{using}}\index{seris\_model@\spxentry{seris\_model}}\ignorespaces 

\subparagraph{2.6.1 Example Using seris\_model}
\label{\detokenize{xsrst/numeric_seirwd_model_xam_py:example-using-seris-model}}\label{\detokenize{xsrst/numeric_seirwd_model_xam_py:id1}}
\begin{sphinxVerbatim}[commandchars=\\\{\}]
\PYG{k+kn}{import} \PYG{n+nn}{numpy}
\PYG{k+kn}{from} \PYG{n+nn}{seirwd\PYGZus{}model} \PYG{k+kn}{import} \PYG{n}{seirwd\PYGZus{}model}
\PYG{k+kn}{from} \PYG{n+nn}{seirwd\PYGZus{}model} \PYG{k+kn}{import} \PYG{n}{test\PYGZus{}fun\PYGZus{}derivatives}
\PYG{c+c1}{\PYGZsh{}}
\PYG{k}{def} \PYG{n+nf}{check\PYGZus{}method}\PYG{p}{(}\PYG{n}{method}\PYG{p}{)} \PYG{p}{:}
    \PYG{n}{ok}    \PYG{o}{=} \PYG{k+kc}{True}
    \PYG{c+c1}{\PYGZsh{}}
    \PYG{n}{alpha} \PYG{o}{=} \PYG{l+m+mf}{9.5}           \PYG{c+c1}{\PYGZsh{} exponent in infectious group size}
    \PYG{n}{sigma} \PYG{o}{=} \PYG{l+m+mf}{1.0} \PYG{o}{/} \PYG{l+m+mf}{5.0}     \PYG{c+c1}{\PYGZsh{} average 5 days between exposure and infectious}
    \PYG{n}{gamma} \PYG{o}{=} \PYG{l+m+mf}{1.0} \PYG{o}{/} \PYG{l+m+mf}{20.0}    \PYG{c+c1}{\PYGZsh{} average 20 days between infectious and recovered}
    \PYG{n}{xi}    \PYG{o}{=} \PYG{l+m+mf}{1.0} \PYG{o}{/} \PYG{l+m+mf}{365.0}   \PYG{c+c1}{\PYGZsh{} average 1 year of immunity}
    \PYG{n}{chi}   \PYG{o}{=} \PYG{n}{gamma} \PYG{o}{/} \PYG{l+m+mf}{10.0}  \PYG{c+c1}{\PYGZsh{} probabiliy of death is 1/10 that of recovery}
    \PYG{n}{delta} \PYG{o}{=} \PYG{l+m+mf}{1.0} \PYG{o}{/} \PYG{l+m+mf}{10.0}    \PYG{c+c1}{\PYGZsh{} to days between no\PYGZhy{}longer infectious and death}
    \PYG{c+c1}{\PYGZsh{}}
    \PYG{c+c1}{\PYGZsh{} t\PYGZus{}all}
    \PYG{n}{t\PYGZus{}all}    \PYG{o}{=} \PYG{n}{numpy}\PYG{o}{.}\PYG{n}{array}\PYG{p}{(} \PYG{p}{[} \PYG{l+m+mf}{0.0}\PYG{p}{,} \PYG{l+m+mf}{0.1}\PYG{p}{,} \PYG{l+m+mf}{0.2} \PYG{p}{]} \PYG{p}{)}
    \PYG{n}{beta\PYGZus{}all} \PYG{o}{=}  \PYG{p}{[} \PYG{l+m+mf}{0.30}\PYG{p}{,} \PYG{l+m+mf}{0.20}\PYG{p}{,} \PYG{l+m+mf}{0.10} \PYG{p}{]}
    \PYG{c+c1}{\PYGZsh{}}
    \PYG{c+c1}{\PYGZsh{} p\PYGZus{}all}
    \PYG{n}{p\PYGZus{}all}    \PYG{o}{=} \PYG{n+nb}{list}\PYG{p}{(}\PYG{p}{)}
    \PYG{k}{for} \PYG{n}{i} \PYG{o+ow}{in} \PYG{n+nb}{range}\PYG{p}{(}\PYG{n}{t\PYGZus{}all}\PYG{o}{.}\PYG{n}{size}\PYG{p}{)} \PYG{p}{:}
        \PYG{n}{p} \PYG{o}{=} \PYG{p}{\PYGZob{}}
            \PYG{l+s+s1}{\PYGZsq{}}\PYG{l+s+s1}{alpha}\PYG{l+s+s1}{\PYGZsq{}} \PYG{p}{:} \PYG{n}{alpha}\PYG{p}{,}
            \PYG{l+s+s1}{\PYGZsq{}}\PYG{l+s+s1}{beta}\PYG{l+s+s1}{\PYGZsq{}}  \PYG{p}{:} \PYG{n}{beta\PYGZus{}all}\PYG{p}{[}\PYG{n}{i}\PYG{p}{]}\PYG{p}{,}
            \PYG{l+s+s1}{\PYGZsq{}}\PYG{l+s+s1}{sigma}\PYG{l+s+s1}{\PYGZsq{}} \PYG{p}{:} \PYG{n}{sigma} \PYG{p}{,}
            \PYG{l+s+s1}{\PYGZsq{}}\PYG{l+s+s1}{gamma}\PYG{l+s+s1}{\PYGZsq{}} \PYG{p}{:} \PYG{n}{gamma} \PYG{p}{,}
            \PYG{l+s+s1}{\PYGZsq{}}\PYG{l+s+s1}{xi}\PYG{l+s+s1}{\PYGZsq{}}    \PYG{p}{:} \PYG{n}{xi} \PYG{p}{,}
            \PYG{l+s+s1}{\PYGZsq{}}\PYG{l+s+s1}{chi}\PYG{l+s+s1}{\PYGZsq{}}   \PYG{p}{:} \PYG{n}{chi} \PYG{p}{,}
            \PYG{l+s+s1}{\PYGZsq{}}\PYG{l+s+s1}{delta}\PYG{l+s+s1}{\PYGZsq{}} \PYG{p}{:} \PYG{n}{delta} \PYG{p}{,}
        \PYG{p}{\PYGZcb{}}
        \PYG{n}{p\PYGZus{}all}\PYG{o}{.}\PYG{n}{append}\PYG{p}{(}\PYG{n}{p}\PYG{p}{)}
    \PYG{c+c1}{\PYGZsh{}}
    \PYG{c+c1}{\PYGZsh{} initial}
    \PYG{n}{I0}  \PYG{o}{=} \PYG{l+m+mf}{0.02}
    \PYG{n}{E0}  \PYG{o}{=} \PYG{l+m+mf}{0.02}
    \PYG{n}{R0}  \PYG{o}{=} \PYG{l+m+mf}{0.02}
    \PYG{n}{S0}  \PYG{o}{=} \PYG{l+m+mf}{1.0} \PYG{o}{\PYGZhy{}} \PYG{n}{I0} \PYG{o}{\PYGZhy{}} \PYG{n}{E0} \PYG{o}{\PYGZhy{}} \PYG{n}{R0}
    \PYG{n}{W0}  \PYG{o}{=} \PYG{l+m+mf}{0.0}
    \PYG{n}{D0}  \PYG{o}{=} \PYG{l+m+mf}{0.0}
    \PYG{n}{initial}  \PYG{o}{=} \PYG{n}{numpy}\PYG{o}{.}\PYG{n}{array}\PYG{p}{(} \PYG{p}{[} \PYG{n}{S0}\PYG{p}{,} \PYG{n}{E0}\PYG{p}{,} \PYG{n}{I0}\PYG{p}{,} \PYG{n}{R0}\PYG{p}{,} \PYG{n}{W0}\PYG{p}{,} \PYG{n}{D0} \PYG{p}{]} \PYG{p}{)}
    \PYG{c+c1}{\PYGZsh{}}
    \PYG{c+c1}{\PYGZsh{} check derivatives in seirwd\PYGZus{}model needed when method is rosen3}
    \PYG{c+c1}{\PYGZsh{} (This is not part of the User API for seirwd\PYGZus{}model)}
    \PYG{k}{if} \PYG{n}{method} \PYG{o}{==} \PYG{l+s+s1}{\PYGZsq{}}\PYG{l+s+s1}{rosen3}\PYG{l+s+s1}{\PYGZsq{}} \PYG{p}{:}
        \PYG{n}{ok} \PYG{o}{=} \PYG{n}{ok} \PYG{o+ow}{and} \PYG{n}{test\PYGZus{}fun\PYGZus{}derivatives}\PYG{p}{(}\PYG{n}{t\PYGZus{}all}\PYG{p}{,} \PYG{n}{p\PYGZus{}all}\PYG{p}{,} \PYG{n}{initial}\PYG{p}{)}
    \PYG{c+c1}{\PYGZsh{}}
    \PYG{c+c1}{\PYGZsh{} solve the ODE}
    \PYG{n}{seirwd\PYGZus{}all} \PYG{o}{=} \PYG{n}{seirwd\PYGZus{}model}\PYG{p}{(}\PYG{n}{method}\PYG{p}{,} \PYG{n}{t\PYGZus{}all}\PYG{p}{,} \PYG{n}{p\PYGZus{}all}\PYG{p}{,} \PYG{n}{initial}\PYG{p}{)}
    \PYG{c+c1}{\PYGZsh{}}
    \PYG{c+c1}{\PYGZsh{} check solution using midpoint values}
    \PYG{k}{for} \PYG{n}{i} \PYG{o+ow}{in} \PYG{n+nb}{range}\PYG{p}{(} \PYG{n+nb}{len}\PYG{p}{(}\PYG{n}{t\PYGZus{}all}\PYG{p}{)} \PYG{o}{\PYGZhy{}} \PYG{l+m+mi}{1} \PYG{p}{)} \PYG{p}{:}
        \PYG{c+c1}{\PYGZsh{} midpoint values}
        \PYG{n}{t\PYGZus{}mid}      \PYG{o}{=} \PYG{p}{(}\PYG{n}{t\PYGZus{}all}\PYG{p}{[}\PYG{n}{i}\PYG{o}{+}\PYG{l+m+mi}{1}\PYG{p}{]} \PYG{o}{+} \PYG{n}{t\PYGZus{}all}\PYG{p}{[}\PYG{n}{i}\PYG{p}{]}\PYG{p}{)} \PYG{o}{/} \PYG{l+m+mf}{2.0}
        \PYG{n}{seirwd\PYGZus{}mid} \PYG{o}{=} \PYG{p}{(}\PYG{n}{seirwd\PYGZus{}all}\PYG{p}{[}\PYG{n}{i}\PYG{p}{,}\PYG{p}{:}\PYG{p}{]} \PYG{o}{+} \PYG{n}{seirwd\PYGZus{}all}\PYG{p}{[}\PYG{n}{i}\PYG{o}{+}\PYG{l+m+mi}{1}\PYG{p}{,}\PYG{p}{:}\PYG{p}{]}\PYG{p}{)} \PYG{o}{/} \PYG{l+m+mf}{2.0}
        \PYG{n}{S}\PYG{p}{,} \PYG{n}{E}\PYG{p}{,} \PYG{n}{I}\PYG{p}{,} \PYG{n}{R}\PYG{p}{,} \PYG{n}{W}\PYG{p}{,} \PYG{n}{D} \PYG{o}{=} \PYG{n}{seirwd\PYGZus{}mid}
        \PYG{c+c1}{\PYGZsh{}}
        \PYG{c+c1}{\PYGZsh{} differential equation}
        \PYG{n}{p}       \PYG{o}{=} \PYG{n+nb}{dict}\PYG{p}{(}\PYG{p}{)}
        \PYG{k}{for} \PYG{n}{key} \PYG{o+ow}{in} \PYG{p}{[}\PYG{l+s+s1}{\PYGZsq{}}\PYG{l+s+s1}{alpha}\PYG{l+s+s1}{\PYGZsq{}}\PYG{p}{,} \PYG{l+s+s1}{\PYGZsq{}}\PYG{l+s+s1}{beta}\PYG{l+s+s1}{\PYGZsq{}}\PYG{p}{,} \PYG{l+s+s1}{\PYGZsq{}}\PYG{l+s+s1}{sigma}\PYG{l+s+s1}{\PYGZsq{}}\PYG{p}{,} \PYG{l+s+s1}{\PYGZsq{}}\PYG{l+s+s1}{gamma}\PYG{l+s+s1}{\PYGZsq{}}\PYG{p}{,} \PYG{l+s+s1}{\PYGZsq{}}\PYG{l+s+s1}{xi}\PYG{l+s+s1}{\PYGZsq{}}\PYG{p}{,} \PYG{l+s+s1}{\PYGZsq{}}\PYG{l+s+s1}{chi}\PYG{l+s+s1}{\PYGZsq{}}\PYG{p}{,} \PYG{l+s+s1}{\PYGZsq{}}\PYG{l+s+s1}{delta}\PYG{l+s+s1}{\PYGZsq{}}\PYG{p}{]} \PYG{p}{:}
            \PYG{n}{p}\PYG{p}{[}\PYG{n}{key}\PYG{p}{]} \PYG{o}{=} \PYG{p}{(}\PYG{n}{p\PYGZus{}all}\PYG{p}{[}\PYG{n}{i}\PYG{p}{]}\PYG{p}{[}\PYG{n}{key}\PYG{p}{]} \PYG{o}{+} \PYG{n}{p\PYGZus{}all}\PYG{p}{[}\PYG{n}{i}\PYG{o}{+}\PYG{l+m+mi}{1}\PYG{p}{]}\PYG{p}{[}\PYG{n}{key}\PYG{p}{]}\PYG{p}{)} \PYG{o}{/} \PYG{l+m+mf}{2.0}
        \PYG{c+c1}{\PYGZsh{}}
        \PYG{n}{dot}     \PYG{o}{=} \PYG{n}{numpy}\PYG{o}{.}\PYG{n}{empty}\PYG{p}{(} \PYG{l+m+mi}{6}\PYG{p}{,} \PYG{n}{dtype} \PYG{o}{=} \PYG{n+nb}{float}\PYG{p}{)}
        \PYG{n}{Ia}      \PYG{o}{=} \PYG{n+nb}{pow}\PYG{p}{(}\PYG{n}{I}\PYG{p}{,} \PYG{n}{p}\PYG{p}{[}\PYG{l+s+s1}{\PYGZsq{}}\PYG{l+s+s1}{alpha}\PYG{l+s+s1}{\PYGZsq{}}\PYG{p}{]}\PYG{p}{)}
        \PYG{n}{dot}\PYG{p}{[}\PYG{l+m+mi}{0}\PYG{p}{]}  \PYG{o}{=} \PYG{o}{\PYGZhy{}} \PYG{n}{p}\PYG{p}{[}\PYG{l+s+s1}{\PYGZsq{}}\PYG{l+s+s1}{beta}\PYG{l+s+s1}{\PYGZsq{}}\PYG{p}{]} \PYG{o}{*}  \PYG{n}{S} \PYG{o}{*} \PYG{n}{Ia}  \PYG{o}{+} \PYG{n}{p}\PYG{p}{[}\PYG{l+s+s1}{\PYGZsq{}}\PYG{l+s+s1}{xi}\PYG{l+s+s1}{\PYGZsq{}}\PYG{p}{]}   \PYG{o}{*} \PYG{n}{R}
        \PYG{n}{dot}\PYG{p}{[}\PYG{l+m+mi}{1}\PYG{p}{]}  \PYG{o}{=} \PYG{o}{+} \PYG{n}{p}\PYG{p}{[}\PYG{l+s+s1}{\PYGZsq{}}\PYG{l+s+s1}{beta}\PYG{l+s+s1}{\PYGZsq{}}\PYG{p}{]} \PYG{o}{*}  \PYG{n}{S} \PYG{o}{*} \PYG{n}{Ia} \PYG{o}{\PYGZhy{}} \PYG{n}{p}\PYG{p}{[}\PYG{l+s+s1}{\PYGZsq{}}\PYG{l+s+s1}{sigma}\PYG{l+s+s1}{\PYGZsq{}}\PYG{p}{]} \PYG{o}{*} \PYG{n}{E}
        \PYG{n}{dot}\PYG{p}{[}\PYG{l+m+mi}{2}\PYG{p}{]}  \PYG{o}{=} \PYG{o}{+} \PYG{n}{p}\PYG{p}{[}\PYG{l+s+s1}{\PYGZsq{}}\PYG{l+s+s1}{sigma}\PYG{l+s+s1}{\PYGZsq{}}\PYG{p}{]} \PYG{o}{*} \PYG{n}{E}      \PYG{o}{\PYGZhy{}} \PYG{p}{(}\PYG{n}{p}\PYG{p}{[}\PYG{l+s+s1}{\PYGZsq{}}\PYG{l+s+s1}{gamma}\PYG{l+s+s1}{\PYGZsq{}}\PYG{p}{]} \PYG{o}{+} \PYG{n}{p}\PYG{p}{[}\PYG{l+s+s1}{\PYGZsq{}}\PYG{l+s+s1}{chi}\PYG{l+s+s1}{\PYGZsq{}}\PYG{p}{]}\PYG{p}{)} \PYG{o}{*} \PYG{n}{I}
        \PYG{n}{dot}\PYG{p}{[}\PYG{l+m+mi}{3}\PYG{p}{]}  \PYG{o}{=} \PYG{o}{+} \PYG{n}{p}\PYG{p}{[}\PYG{l+s+s1}{\PYGZsq{}}\PYG{l+s+s1}{gamma}\PYG{l+s+s1}{\PYGZsq{}}\PYG{p}{]} \PYG{o}{*} \PYG{n}{I}      \PYG{o}{\PYGZhy{}} \PYG{n}{p}\PYG{p}{[}\PYG{l+s+s1}{\PYGZsq{}}\PYG{l+s+s1}{xi}\PYG{l+s+s1}{\PYGZsq{}}\PYG{p}{]}    \PYG{o}{*} \PYG{n}{R}
        \PYG{n}{dot}\PYG{p}{[}\PYG{l+m+mi}{4}\PYG{p}{]}  \PYG{o}{=} \PYG{o}{+} \PYG{n}{p}\PYG{p}{[}\PYG{l+s+s1}{\PYGZsq{}}\PYG{l+s+s1}{chi}\PYG{l+s+s1}{\PYGZsq{}}\PYG{p}{]}   \PYG{o}{*} \PYG{n}{I}      \PYG{o}{\PYGZhy{}} \PYG{n}{p}\PYG{p}{[}\PYG{l+s+s1}{\PYGZsq{}}\PYG{l+s+s1}{delta}\PYG{l+s+s1}{\PYGZsq{}}\PYG{p}{]} \PYG{o}{*} \PYG{n}{W}
        \PYG{n}{dot}\PYG{p}{[}\PYG{l+m+mi}{5}\PYG{p}{]}  \PYG{o}{=} \PYG{o}{+} \PYG{n}{p}\PYG{p}{[}\PYG{l+s+s1}{\PYGZsq{}}\PYG{l+s+s1}{delta}\PYG{l+s+s1}{\PYGZsq{}}\PYG{p}{]} \PYG{o}{*} \PYG{n}{W}
        \PYG{c+c1}{\PYGZsh{}}
        \PYG{c+c1}{\PYGZsh{} difference over interval}
        \PYG{n}{seirwd\PYGZus{}delta} \PYG{o}{=} \PYG{n}{seirwd\PYGZus{}all}\PYG{p}{[}\PYG{n}{i}\PYG{o}{+}\PYG{l+m+mi}{1}\PYG{p}{,}\PYG{p}{:}\PYG{p}{]} \PYG{o}{\PYGZhy{}} \PYG{n}{seirwd\PYGZus{}all}\PYG{p}{[}\PYG{n}{i}\PYG{p}{,}\PYG{p}{:}\PYG{p}{]}
        \PYG{n}{t\PYGZus{}delta}      \PYG{o}{=} \PYG{n}{t\PYGZus{}all}\PYG{p}{[}\PYG{n}{i}\PYG{o}{+}\PYG{l+m+mi}{1}\PYG{p}{]}      \PYG{o}{\PYGZhy{}} \PYG{n}{t\PYGZus{}all}\PYG{p}{[}\PYG{n}{i}\PYG{p}{]}
        \PYG{c+c1}{\PYGZsh{}}
        \PYG{n}{check}  \PYG{o}{=} \PYG{n}{seirwd\PYGZus{}delta} \PYG{o}{/} \PYG{n}{t\PYGZus{}delta}
        \PYG{n}{rel\PYGZus{}error} \PYG{o}{=} \PYG{n}{dot} \PYG{o}{/} \PYG{n}{check} \PYG{o}{\PYGZhy{}} \PYG{l+m+mf}{1.0}
        \PYG{c+c1}{\PYGZsh{} print(rel\PYGZus{}error)}
        \PYG{n}{ok}        \PYG{o}{=} \PYG{n}{ok} \PYG{o+ow}{and} \PYG{n+nb}{all}\PYG{p}{(} \PYG{n+nb}{abs}\PYG{p}{(}\PYG{n}{rel\PYGZus{}error}\PYG{p}{)} \PYG{o}{\PYGZlt{}} \PYG{l+m+mf}{1e\PYGZhy{}2} \PYG{p}{)}
    \PYG{c+c1}{\PYGZsh{}}
    \PYG{c+c1}{\PYGZsh{} now check solution using twice as many Runge\PYGZhy{}Kutta steps}
    \PYG{n}{n\PYGZus{}step}      \PYG{o}{=} \PYG{l+m+mi}{2}
    \PYG{n}{seirwd\PYGZus{}check} \PYG{o}{=} \PYG{n}{seirwd\PYGZus{}model}\PYG{p}{(}\PYG{n}{method}\PYG{p}{,} \PYG{n}{t\PYGZus{}all}\PYG{p}{,} \PYG{n}{p\PYGZus{}all}\PYG{p}{,} \PYG{n}{initial}\PYG{p}{,} \PYG{n}{n\PYGZus{}step}\PYG{p}{)}
    \PYG{n}{error}       \PYG{o}{=} \PYG{n}{seirwd\PYGZus{}all} \PYG{o}{\PYGZhy{}} \PYG{n}{seirwd\PYGZus{}check}
    \PYG{c+c1}{\PYGZsh{} print(error)}
    \PYG{n}{ok}          \PYG{o}{=} \PYG{n}{ok} \PYG{o+ow}{and} \PYG{n}{numpy}\PYG{o}{.}\PYG{n}{all}\PYG{p}{(} \PYG{n+nb}{abs}\PYG{p}{(}\PYG{n}{error}\PYG{p}{)} \PYG{o}{\PYGZlt{}} \PYG{l+m+mf}{1e\PYGZhy{}7} \PYG{p}{)}
    \PYG{c+c1}{\PYGZsh{}}
    \PYG{k}{return} \PYG{n}{ok}

\PYG{k}{def} \PYG{n+nf}{seirwd\PYGZus{}model\PYGZus{}xam}\PYG{p}{(}\PYG{p}{)} \PYG{p}{:}
    \PYG{n}{ok} \PYG{o}{=} \PYG{k+kc}{True}
    \PYG{n}{ok} \PYG{o}{=} \PYG{n}{ok} \PYG{o+ow}{and} \PYG{n}{check\PYGZus{}method}\PYG{p}{(}\PYG{l+s+s1}{\PYGZsq{}}\PYG{l+s+s1}{runge4}\PYG{l+s+s1}{\PYGZsq{}}\PYG{p}{)}\PYG{p}{;}
    \PYG{n}{ok} \PYG{o}{=} \PYG{n}{ok} \PYG{o+ow}{and} \PYG{n}{check\PYGZus{}method}\PYG{p}{(}\PYG{l+s+s1}{\PYGZsq{}}\PYG{l+s+s1}{rosen3}\PYG{l+s+s1}{\PYGZsq{}}\PYG{p}{)}\PYG{p}{;}
    \PYG{k}{return} \PYG{n}{ok}

\end{sphinxVerbatim}


\bigskip\hrule\bigskip


xsrst input file: \sphinxcode{\sphinxupquote{example/python/numeric/seirwd\_model\_xam.py}}

\index{example@\spxentry{example}}\ignorespaces 

\subparagraph{Example}
\label{\detokenize{xsrst/numeric_seirwd_model:example}}\label{\detokenize{xsrst/numeric_seirwd_model:numeric-seirwd-model-example}}\label{\detokenize{xsrst/numeric_seirwd_model:index-13}}
{\hyperref[\detokenize{xsrst/numeric_seirwd_model_xam_py:id1}]{\sphinxcrossref{\DUrole{std,std-ref}{numeric\_seirwd\_model\_xam\_py}}}}

\index{source@\spxentry{source}}\index{code@\spxentry{code}}\ignorespaces 

\subparagraph{Source Code}
\label{\detokenize{xsrst/numeric_seirwd_model:source-code}}\label{\detokenize{xsrst/numeric_seirwd_model:numeric-seirwd-model-source-code}}\label{\detokenize{xsrst/numeric_seirwd_model:index-14}}
\begin{sphinxVerbatim}[commandchars=\\\{\}]
\PYG{k+kn}{import} \PYG{n+nn}{numpy}
\PYG{k+kn}{from} \PYG{n+nn}{ode\PYGZus{}multi\PYGZus{}step} \PYG{k+kn}{import} \PYG{n}{ode\PYGZus{}multi\PYGZus{}step}
\PYG{k+kn}{from} \PYG{n+nn}{rosen3\PYGZus{}step} \PYG{k+kn}{import} \PYG{n}{rosen3\PYGZus{}step}
\PYG{k+kn}{from} \PYG{n+nn}{runge4\PYGZus{}step} \PYG{k+kn}{import} \PYG{n}{runge4\PYGZus{}step}
\PYG{k+kn}{from} \PYG{n+nn}{rosen3\PYGZus{}step} \PYG{k+kn}{import} \PYG{n}{check\PYGZus{}rosen3\PYGZus{}step}
\PYG{c+c1}{\PYGZsh{} \PYGZhy{}\PYGZhy{}\PYGZhy{}\PYGZhy{}\PYGZhy{}\PYGZhy{}\PYGZhy{}\PYGZhy{}\PYGZhy{}\PYGZhy{}\PYGZhy{}\PYGZhy{}\PYGZhy{}\PYGZhy{}\PYGZhy{}\PYGZhy{}\PYGZhy{}\PYGZhy{}\PYGZhy{}\PYGZhy{}\PYGZhy{}\PYGZhy{}\PYGZhy{}\PYGZhy{}\PYGZhy{}\PYGZhy{}\PYGZhy{}\PYGZhy{}\PYGZhy{}\PYGZhy{}\PYGZhy{}\PYGZhy{}\PYGZhy{}\PYGZhy{}\PYGZhy{}\PYGZhy{}\PYGZhy{}\PYGZhy{}\PYGZhy{}\PYGZhy{}\PYGZhy{}\PYGZhy{}\PYGZhy{}\PYGZhy{}\PYGZhy{}\PYGZhy{}\PYGZhy{}\PYGZhy{}\PYGZhy{}\PYGZhy{}\PYGZhy{}\PYGZhy{}\PYGZhy{}\PYGZhy{}\PYGZhy{}\PYGZhy{}\PYGZhy{}\PYGZhy{}\PYGZhy{}\PYGZhy{}\PYGZhy{}\PYGZhy{}\PYGZhy{}\PYGZhy{}\PYGZhy{}\PYGZhy{}\PYGZhy{}\PYGZhy{}\PYGZhy{}\PYGZhy{}\PYGZhy{}\PYGZhy{}\PYGZhy{}\PYGZhy{}\PYGZhy{}\PYGZhy{}\PYGZhy{}}
\PYG{k}{class} \PYG{n+nc}{fun\PYGZus{}class} \PYG{p}{:}
    \PYG{c+c1}{\PYGZsh{} \PYGZhy{}\PYGZhy{}\PYGZhy{}\PYGZhy{}\PYGZhy{}\PYGZhy{}\PYGZhy{}\PYGZhy{}\PYGZhy{}\PYGZhy{}\PYGZhy{}\PYGZhy{}\PYGZhy{}\PYGZhy{}\PYGZhy{}\PYGZhy{}\PYGZhy{}\PYGZhy{}\PYGZhy{}\PYGZhy{}\PYGZhy{}\PYGZhy{}\PYGZhy{}\PYGZhy{}\PYGZhy{}\PYGZhy{}\PYGZhy{}\PYGZhy{}\PYGZhy{}\PYGZhy{}\PYGZhy{}\PYGZhy{}\PYGZhy{}\PYGZhy{}\PYGZhy{}\PYGZhy{}\PYGZhy{}\PYGZhy{}\PYGZhy{}\PYGZhy{}\PYGZhy{}\PYGZhy{}\PYGZhy{}\PYGZhy{}\PYGZhy{}\PYGZhy{}\PYGZhy{}\PYGZhy{}\PYGZhy{}\PYGZhy{}\PYGZhy{}\PYGZhy{}\PYGZhy{}\PYGZhy{}\PYGZhy{}\PYGZhy{}\PYGZhy{}\PYGZhy{}\PYGZhy{}\PYGZhy{}\PYGZhy{}\PYGZhy{}\PYGZhy{}\PYGZhy{}\PYGZhy{}\PYGZhy{}\PYGZhy{}\PYGZhy{}\PYGZhy{}\PYGZhy{}\PYGZhy{}\PYGZhy{}}
    \PYG{c+c1}{\PYGZsh{} init}
    \PYG{k}{def} \PYG{n+nf+fm}{\PYGZus{}\PYGZus{}init\PYGZus{}\PYGZus{}}\PYG{p}{(}\PYG{n+nb+bp}{self}\PYG{p}{,} \PYG{n}{t\PYGZus{}all}\PYG{p}{,} \PYG{n}{p\PYGZus{}all}\PYG{p}{,} \PYG{n}{n\PYGZus{}step}\PYG{p}{)} \PYG{p}{:}
        \PYG{n+nb+bp}{self}\PYG{o}{.}\PYG{n}{t\PYGZus{}all}  \PYG{o}{=} \PYG{n}{t\PYGZus{}all}
        \PYG{n+nb+bp}{self}\PYG{o}{.}\PYG{n}{p\PYGZus{}all}  \PYG{o}{=} \PYG{n}{p\PYGZus{}all}
        \PYG{n+nb+bp}{self}\PYG{o}{.}\PYG{n}{n\PYGZus{}step} \PYG{o}{=} \PYG{n}{n\PYGZus{}step}
        \PYG{n+nb+bp}{self}\PYG{o}{.}\PYG{n}{index}  \PYG{o}{=} \PYG{k+kc}{None}
        \PYG{n+nb+bp}{self}\PYG{o}{.}\PYG{n}{p}      \PYG{o}{=} \PYG{k+kc}{None}
    \PYG{c+c1}{\PYGZsh{} \PYGZhy{}\PYGZhy{}\PYGZhy{}\PYGZhy{}\PYGZhy{}\PYGZhy{}\PYGZhy{}\PYGZhy{}\PYGZhy{}\PYGZhy{}\PYGZhy{}\PYGZhy{}\PYGZhy{}\PYGZhy{}\PYGZhy{}\PYGZhy{}\PYGZhy{}\PYGZhy{}\PYGZhy{}\PYGZhy{}\PYGZhy{}\PYGZhy{}\PYGZhy{}\PYGZhy{}\PYGZhy{}\PYGZhy{}\PYGZhy{}\PYGZhy{}\PYGZhy{}\PYGZhy{}\PYGZhy{}\PYGZhy{}\PYGZhy{}\PYGZhy{}\PYGZhy{}\PYGZhy{}\PYGZhy{}\PYGZhy{}\PYGZhy{}\PYGZhy{}\PYGZhy{}\PYGZhy{}\PYGZhy{}\PYGZhy{}\PYGZhy{}\PYGZhy{}\PYGZhy{}\PYGZhy{}\PYGZhy{}\PYGZhy{}\PYGZhy{}\PYGZhy{}\PYGZhy{}\PYGZhy{}\PYGZhy{}\PYGZhy{}\PYGZhy{}\PYGZhy{}\PYGZhy{}\PYGZhy{}\PYGZhy{}\PYGZhy{}\PYGZhy{}\PYGZhy{}\PYGZhy{}\PYGZhy{}\PYGZhy{}\PYGZhy{}\PYGZhy{}\PYGZhy{}\PYGZhy{}}
    \PYG{c+c1}{\PYGZsh{} set\PYGZus{}t\PYGZus{}all\PYGZus{}index}
    \PYG{k}{def} \PYG{n+nf}{set\PYGZus{}t\PYGZus{}all\PYGZus{}index}\PYG{p}{(}\PYG{n+nb+bp}{self}\PYG{p}{,} \PYG{n}{refine\PYGZus{}index}\PYG{p}{)} \PYG{p}{:}
        \PYG{n}{t\PYGZus{}all\PYGZus{}index}  \PYG{o}{=} \PYG{n+nb}{int}\PYG{p}{(} \PYG{n}{numpy}\PYG{o}{.}\PYG{n}{floor}\PYG{p}{(} \PYG{n}{refine\PYGZus{}index} \PYG{o}{/} \PYG{n+nb+bp}{self}\PYG{o}{.}\PYG{n}{n\PYGZus{}step} \PYG{p}{)} \PYG{p}{)}
        \PYG{n+nb+bp}{self}\PYG{o}{.}\PYG{n}{index}   \PYG{o}{=} \PYG{n}{t\PYGZus{}all\PYGZus{}index}
    \PYG{c+c1}{\PYGZsh{} \PYGZhy{}\PYGZhy{}\PYGZhy{}\PYGZhy{}\PYGZhy{}\PYGZhy{}\PYGZhy{}\PYGZhy{}\PYGZhy{}\PYGZhy{}\PYGZhy{}\PYGZhy{}\PYGZhy{}\PYGZhy{}\PYGZhy{}\PYGZhy{}\PYGZhy{}\PYGZhy{}\PYGZhy{}\PYGZhy{}\PYGZhy{}\PYGZhy{}\PYGZhy{}\PYGZhy{}\PYGZhy{}\PYGZhy{}\PYGZhy{}\PYGZhy{}\PYGZhy{}\PYGZhy{}\PYGZhy{}\PYGZhy{}\PYGZhy{}\PYGZhy{}\PYGZhy{}\PYGZhy{}\PYGZhy{}\PYGZhy{}\PYGZhy{}\PYGZhy{}\PYGZhy{}\PYGZhy{}\PYGZhy{}\PYGZhy{}\PYGZhy{}\PYGZhy{}\PYGZhy{}\PYGZhy{}\PYGZhy{}\PYGZhy{}\PYGZhy{}\PYGZhy{}\PYGZhy{}\PYGZhy{}\PYGZhy{}\PYGZhy{}\PYGZhy{}\PYGZhy{}\PYGZhy{}\PYGZhy{}\PYGZhy{}\PYGZhy{}\PYGZhy{}\PYGZhy{}\PYGZhy{}\PYGZhy{}\PYGZhy{}\PYGZhy{}\PYGZhy{}\PYGZhy{}\PYGZhy{}}
    \PYG{c+c1}{\PYGZsh{} get\PYGZus{}p}
    \PYG{k}{def} \PYG{n+nf}{get\PYGZus{}p}\PYG{p}{(}\PYG{n+nb+bp}{self}\PYG{p}{,} \PYG{n}{t}\PYG{p}{,} \PYG{n}{seirwd}\PYG{p}{)} \PYG{p}{:}
        \PYG{k}{assert} \PYG{n+nb+bp}{self}\PYG{o}{.}\PYG{n}{index} \PYG{o}{!=} \PYG{k+kc}{None}
        \PYG{n}{t\PYGZus{}all} \PYG{o}{=} \PYG{n+nb+bp}{self}\PYG{o}{.}\PYG{n}{t\PYGZus{}all}
        \PYG{n}{p\PYGZus{}all} \PYG{o}{=} \PYG{n+nb+bp}{self}\PYG{o}{.}\PYG{n}{p\PYGZus{}all}
        \PYG{n}{index} \PYG{o}{=} \PYG{n+nb+bp}{self}\PYG{o}{.}\PYG{n}{index}
        \PYG{c+c1}{\PYGZsh{}}
        \PYG{c+c1}{\PYGZsh{} linearly interpolate the parameter values}
        \PYG{n}{p}      \PYG{o}{=} \PYG{n+nb}{dict}\PYG{p}{(}\PYG{p}{)}
        \PYG{n}{t\PYGZus{}left} \PYG{o}{=} \PYG{n}{t\PYGZus{}all}\PYG{p}{[}\PYG{n}{index}\PYG{p}{]}
        \PYG{n}{t\PYGZus{}diff} \PYG{o}{=} \PYG{n}{t\PYGZus{}all}\PYG{p}{[}\PYG{n}{index} \PYG{o}{+} \PYG{l+m+mi}{1}\PYG{p}{]} \PYG{o}{\PYGZhy{}} \PYG{n}{t\PYGZus{}left}
        \PYG{k}{for} \PYG{n}{key} \PYG{o+ow}{in} \PYG{p}{[}\PYG{l+s+s1}{\PYGZsq{}}\PYG{l+s+s1}{alpha}\PYG{l+s+s1}{\PYGZsq{}}\PYG{p}{,} \PYG{l+s+s1}{\PYGZsq{}}\PYG{l+s+s1}{beta}\PYG{l+s+s1}{\PYGZsq{}}\PYG{p}{,} \PYG{l+s+s1}{\PYGZsq{}}\PYG{l+s+s1}{sigma}\PYG{l+s+s1}{\PYGZsq{}}\PYG{p}{,} \PYG{l+s+s1}{\PYGZsq{}}\PYG{l+s+s1}{gamma}\PYG{l+s+s1}{\PYGZsq{}}\PYG{p}{,} \PYG{l+s+s1}{\PYGZsq{}}\PYG{l+s+s1}{xi}\PYG{l+s+s1}{\PYGZsq{}}\PYG{p}{,} \PYG{l+s+s1}{\PYGZsq{}}\PYG{l+s+s1}{chi}\PYG{l+s+s1}{\PYGZsq{}}\PYG{p}{,} \PYG{l+s+s1}{\PYGZsq{}}\PYG{l+s+s1}{delta}\PYG{l+s+s1}{\PYGZsq{}}\PYG{p}{]} \PYG{p}{:}
            \PYG{n}{p\PYGZus{}left} \PYG{o}{=} \PYG{n}{p\PYGZus{}all}\PYG{p}{[}\PYG{n}{index}\PYG{p}{]}\PYG{p}{[}\PYG{n}{key}\PYG{p}{]}
            \PYG{n}{p\PYGZus{}diff} \PYG{o}{=} \PYG{n}{p\PYGZus{}all}\PYG{p}{[}\PYG{n}{index} \PYG{o}{+} \PYG{l+m+mi}{1}\PYG{p}{]}\PYG{p}{[}\PYG{n}{key}\PYG{p}{]} \PYG{o}{\PYGZhy{}} \PYG{n}{p\PYGZus{}left}
            \PYG{n}{p}\PYG{p}{[}\PYG{n}{key}\PYG{p}{]} \PYG{o}{=} \PYG{n}{p\PYGZus{}left} \PYG{o}{+} \PYG{p}{(}\PYG{n}{t} \PYG{o}{\PYGZhy{}} \PYG{n}{t\PYGZus{}left}\PYG{p}{)} \PYG{o}{*} \PYG{n}{p\PYGZus{}diff} \PYG{o}{/} \PYG{n}{t\PYGZus{}diff}
        \PYG{c+c1}{\PYGZsh{}}
        \PYG{k}{return} \PYG{n}{p}
    \PYG{c+c1}{\PYGZsh{} \PYGZhy{}\PYGZhy{}\PYGZhy{}\PYGZhy{}\PYGZhy{}\PYGZhy{}\PYGZhy{}\PYGZhy{}\PYGZhy{}\PYGZhy{}\PYGZhy{}\PYGZhy{}\PYGZhy{}\PYGZhy{}\PYGZhy{}\PYGZhy{}\PYGZhy{}\PYGZhy{}\PYGZhy{}\PYGZhy{}\PYGZhy{}\PYGZhy{}\PYGZhy{}\PYGZhy{}\PYGZhy{}\PYGZhy{}\PYGZhy{}\PYGZhy{}\PYGZhy{}\PYGZhy{}\PYGZhy{}\PYGZhy{}\PYGZhy{}\PYGZhy{}\PYGZhy{}\PYGZhy{}\PYGZhy{}\PYGZhy{}\PYGZhy{}\PYGZhy{}\PYGZhy{}\PYGZhy{}\PYGZhy{}\PYGZhy{}\PYGZhy{}\PYGZhy{}\PYGZhy{}\PYGZhy{}\PYGZhy{}\PYGZhy{}\PYGZhy{}\PYGZhy{}\PYGZhy{}\PYGZhy{}\PYGZhy{}\PYGZhy{}\PYGZhy{}\PYGZhy{}\PYGZhy{}\PYGZhy{}\PYGZhy{}\PYGZhy{}\PYGZhy{}\PYGZhy{}\PYGZhy{}\PYGZhy{}\PYGZhy{}\PYGZhy{}\PYGZhy{}\PYGZhy{}\PYGZhy{}}
    \PYG{c+c1}{\PYGZsh{} get\PYGZus{}p\PYGZus{}t}
    \PYG{k}{def} \PYG{n+nf}{get\PYGZus{}p\PYGZus{}t}\PYG{p}{(}\PYG{n+nb+bp}{self}\PYG{p}{)} \PYG{p}{:}
        \PYG{k}{assert} \PYG{n+nb+bp}{self}\PYG{o}{.}\PYG{n}{index} \PYG{o}{!=} \PYG{k+kc}{None}
        \PYG{n}{t\PYGZus{}all} \PYG{o}{=} \PYG{n+nb+bp}{self}\PYG{o}{.}\PYG{n}{t\PYGZus{}all}
        \PYG{n}{p\PYGZus{}all} \PYG{o}{=} \PYG{n+nb+bp}{self}\PYG{o}{.}\PYG{n}{p\PYGZus{}all}
        \PYG{n}{index} \PYG{o}{=} \PYG{n+nb+bp}{self}\PYG{o}{.}\PYG{n}{index}
        \PYG{c+c1}{\PYGZsh{}}
        \PYG{c+c1}{\PYGZsh{} linearly interpolate the parameter values}
        \PYG{n}{p\PYGZus{}t}    \PYG{o}{=} \PYG{n+nb}{dict}\PYG{p}{(}\PYG{p}{)}
        \PYG{n}{t\PYGZus{}diff} \PYG{o}{=} \PYG{n}{t\PYGZus{}all}\PYG{p}{[}\PYG{n}{index} \PYG{o}{+} \PYG{l+m+mi}{1}\PYG{p}{]} \PYG{o}{\PYGZhy{}} \PYG{n}{t\PYGZus{}all}\PYG{p}{[}\PYG{n}{index}\PYG{p}{]}
        \PYG{k}{for} \PYG{n}{key} \PYG{o+ow}{in} \PYG{p}{[}\PYG{l+s+s1}{\PYGZsq{}}\PYG{l+s+s1}{alpha}\PYG{l+s+s1}{\PYGZsq{}}\PYG{p}{,} \PYG{l+s+s1}{\PYGZsq{}}\PYG{l+s+s1}{beta}\PYG{l+s+s1}{\PYGZsq{}}\PYG{p}{,} \PYG{l+s+s1}{\PYGZsq{}}\PYG{l+s+s1}{sigma}\PYG{l+s+s1}{\PYGZsq{}}\PYG{p}{,} \PYG{l+s+s1}{\PYGZsq{}}\PYG{l+s+s1}{gamma}\PYG{l+s+s1}{\PYGZsq{}}\PYG{p}{,} \PYG{l+s+s1}{\PYGZsq{}}\PYG{l+s+s1}{xi}\PYG{l+s+s1}{\PYGZsq{}}\PYG{p}{,} \PYG{l+s+s1}{\PYGZsq{}}\PYG{l+s+s1}{chi}\PYG{l+s+s1}{\PYGZsq{}}\PYG{p}{,} \PYG{l+s+s1}{\PYGZsq{}}\PYG{l+s+s1}{delta}\PYG{l+s+s1}{\PYGZsq{}}\PYG{p}{]} \PYG{p}{:}
            \PYG{n}{p\PYGZus{}diff}   \PYG{o}{=} \PYG{n}{p\PYGZus{}all}\PYG{p}{[}\PYG{n}{index} \PYG{o}{+} \PYG{l+m+mi}{1}\PYG{p}{]}\PYG{p}{[}\PYG{n}{key}\PYG{p}{]} \PYG{o}{\PYGZhy{}} \PYG{n}{p\PYGZus{}all}\PYG{p}{[}\PYG{n}{index}\PYG{p}{]}\PYG{p}{[}\PYG{n}{key}\PYG{p}{]}
            \PYG{n}{p\PYGZus{}t}\PYG{p}{[}\PYG{n}{key}\PYG{p}{]} \PYG{o}{=} \PYG{n}{p\PYGZus{}diff} \PYG{o}{/} \PYG{n}{t\PYGZus{}diff}
        \PYG{c+c1}{\PYGZsh{}}
        \PYG{k}{return} \PYG{n}{p\PYGZus{}t}
    \PYG{c+c1}{\PYGZsh{} \PYGZhy{}\PYGZhy{}\PYGZhy{}\PYGZhy{}\PYGZhy{}\PYGZhy{}\PYGZhy{}\PYGZhy{}\PYGZhy{}\PYGZhy{}\PYGZhy{}\PYGZhy{}\PYGZhy{}\PYGZhy{}\PYGZhy{}\PYGZhy{}\PYGZhy{}\PYGZhy{}\PYGZhy{}\PYGZhy{}\PYGZhy{}\PYGZhy{}\PYGZhy{}\PYGZhy{}\PYGZhy{}\PYGZhy{}\PYGZhy{}\PYGZhy{}\PYGZhy{}\PYGZhy{}\PYGZhy{}\PYGZhy{}\PYGZhy{}\PYGZhy{}\PYGZhy{}\PYGZhy{}\PYGZhy{}\PYGZhy{}\PYGZhy{}\PYGZhy{}\PYGZhy{}\PYGZhy{}\PYGZhy{}\PYGZhy{}\PYGZhy{}\PYGZhy{}\PYGZhy{}\PYGZhy{}\PYGZhy{}\PYGZhy{}\PYGZhy{}\PYGZhy{}\PYGZhy{}\PYGZhy{}\PYGZhy{}\PYGZhy{}\PYGZhy{}\PYGZhy{}\PYGZhy{}\PYGZhy{}\PYGZhy{}\PYGZhy{}\PYGZhy{}\PYGZhy{}\PYGZhy{}\PYGZhy{}\PYGZhy{}\PYGZhy{}\PYGZhy{}\PYGZhy{}\PYGZhy{}}
    \PYG{c+c1}{\PYGZsh{} f}
    \PYG{k}{def} \PYG{n+nf}{f}\PYG{p}{(}\PYG{n+nb+bp}{self}\PYG{p}{,} \PYG{n}{t}\PYG{p}{,} \PYG{n}{seirwd}\PYG{p}{)} \PYG{p}{:}
        \PYG{n}{p}                \PYG{o}{=} \PYG{n+nb+bp}{self}\PYG{o}{.}\PYG{n}{get\PYGZus{}p}\PYG{p}{(}\PYG{n}{t}\PYG{p}{,} \PYG{n}{seirwd}\PYG{p}{)}
        \PYG{n}{S}\PYG{p}{,} \PYG{n}{E}\PYG{p}{,} \PYG{n}{I}\PYG{p}{,} \PYG{n}{R}\PYG{p}{,} \PYG{n}{W}\PYG{p}{,} \PYG{n}{D} \PYG{o}{=} \PYG{n}{seirwd}
        \PYG{c+c1}{\PYGZsh{}}
        \PYG{n}{Ia}     \PYG{o}{=} \PYG{n+nb}{pow}\PYG{p}{(}\PYG{n}{I}\PYG{p}{,} \PYG{n}{p}\PYG{p}{[}\PYG{l+s+s1}{\PYGZsq{}}\PYG{l+s+s1}{alpha}\PYG{l+s+s1}{\PYGZsq{}}\PYG{p}{]}\PYG{p}{)}
        \PYG{c+c1}{\PYGZsh{}}
        \PYG{c+c1}{\PYGZsh{} compute f(t, y)}
        \PYG{n}{Sdot}   \PYG{o}{=} \PYG{o}{\PYGZhy{}} \PYG{n}{p}\PYG{p}{[}\PYG{l+s+s1}{\PYGZsq{}}\PYG{l+s+s1}{beta}\PYG{l+s+s1}{\PYGZsq{}}\PYG{p}{]}  \PYG{o}{*}  \PYG{n}{S} \PYG{o}{*} \PYG{n}{Ia}  \PYG{o}{+} \PYG{n}{p}\PYG{p}{[}\PYG{l+s+s1}{\PYGZsq{}}\PYG{l+s+s1}{xi}\PYG{l+s+s1}{\PYGZsq{}}\PYG{p}{]}    \PYG{o}{*} \PYG{n}{R}
        \PYG{n}{Edot}   \PYG{o}{=} \PYG{o}{+} \PYG{n}{p}\PYG{p}{[}\PYG{l+s+s1}{\PYGZsq{}}\PYG{l+s+s1}{beta}\PYG{l+s+s1}{\PYGZsq{}}\PYG{p}{]}  \PYG{o}{*}  \PYG{n}{S} \PYG{o}{*} \PYG{n}{Ia}  \PYG{o}{\PYGZhy{}} \PYG{n}{p}\PYG{p}{[}\PYG{l+s+s1}{\PYGZsq{}}\PYG{l+s+s1}{sigma}\PYG{l+s+s1}{\PYGZsq{}}\PYG{p}{]} \PYG{o}{*} \PYG{n}{E}
        \PYG{n}{Idot}   \PYG{o}{=} \PYG{o}{+} \PYG{n}{p}\PYG{p}{[}\PYG{l+s+s1}{\PYGZsq{}}\PYG{l+s+s1}{sigma}\PYG{l+s+s1}{\PYGZsq{}}\PYG{p}{]} \PYG{o}{*} \PYG{n}{E}        \PYG{o}{\PYGZhy{}} \PYG{p}{(}\PYG{n}{p}\PYG{p}{[}\PYG{l+s+s1}{\PYGZsq{}}\PYG{l+s+s1}{gamma}\PYG{l+s+s1}{\PYGZsq{}}\PYG{p}{]} \PYG{o}{+} \PYG{n}{p}\PYG{p}{[}\PYG{l+s+s1}{\PYGZsq{}}\PYG{l+s+s1}{chi}\PYG{l+s+s1}{\PYGZsq{}}\PYG{p}{]}\PYG{p}{)} \PYG{o}{*} \PYG{n}{I}
        \PYG{n}{Rdot}   \PYG{o}{=} \PYG{o}{+} \PYG{n}{p}\PYG{p}{[}\PYG{l+s+s1}{\PYGZsq{}}\PYG{l+s+s1}{gamma}\PYG{l+s+s1}{\PYGZsq{}}\PYG{p}{]} \PYG{o}{*} \PYG{n}{I}        \PYG{o}{\PYGZhy{}} \PYG{n}{p}\PYG{p}{[}\PYG{l+s+s1}{\PYGZsq{}}\PYG{l+s+s1}{xi}\PYG{l+s+s1}{\PYGZsq{}}\PYG{p}{]}    \PYG{o}{*} \PYG{n}{R}
        \PYG{n}{Wdot}   \PYG{o}{=} \PYG{o}{+} \PYG{n}{p}\PYG{p}{[}\PYG{l+s+s1}{\PYGZsq{}}\PYG{l+s+s1}{chi}\PYG{l+s+s1}{\PYGZsq{}}\PYG{p}{]}   \PYG{o}{*} \PYG{n}{I}        \PYG{o}{\PYGZhy{}} \PYG{n}{p}\PYG{p}{[}\PYG{l+s+s1}{\PYGZsq{}}\PYG{l+s+s1}{delta}\PYG{l+s+s1}{\PYGZsq{}}\PYG{p}{]} \PYG{o}{*} \PYG{n}{W}
        \PYG{n}{Ddot}   \PYG{o}{=} \PYG{o}{+} \PYG{n}{p}\PYG{p}{[}\PYG{l+s+s1}{\PYGZsq{}}\PYG{l+s+s1}{delta}\PYG{l+s+s1}{\PYGZsq{}}\PYG{p}{]} \PYG{o}{*} \PYG{n}{W}
        \PYG{c+c1}{\PYGZsh{}}
        \PYG{k}{return} \PYG{n}{numpy}\PYG{o}{.}\PYG{n}{array}\PYG{p}{(}\PYG{p}{[} \PYG{n}{Sdot}\PYG{p}{,} \PYG{n}{Edot}\PYG{p}{,} \PYG{n}{Idot}\PYG{p}{,} \PYG{n}{Rdot}\PYG{p}{,} \PYG{n}{Wdot}\PYG{p}{,} \PYG{n}{Ddot}\PYG{p}{]}\PYG{p}{)}
    \PYG{c+c1}{\PYGZsh{} \PYGZhy{}\PYGZhy{}\PYGZhy{}\PYGZhy{}\PYGZhy{}\PYGZhy{}\PYGZhy{}\PYGZhy{}\PYGZhy{}\PYGZhy{}\PYGZhy{}\PYGZhy{}\PYGZhy{}\PYGZhy{}\PYGZhy{}\PYGZhy{}\PYGZhy{}\PYGZhy{}\PYGZhy{}\PYGZhy{}\PYGZhy{}\PYGZhy{}\PYGZhy{}\PYGZhy{}\PYGZhy{}\PYGZhy{}\PYGZhy{}\PYGZhy{}\PYGZhy{}\PYGZhy{}\PYGZhy{}\PYGZhy{}\PYGZhy{}\PYGZhy{}\PYGZhy{}\PYGZhy{}\PYGZhy{}\PYGZhy{}\PYGZhy{}\PYGZhy{}\PYGZhy{}\PYGZhy{}\PYGZhy{}\PYGZhy{}\PYGZhy{}\PYGZhy{}\PYGZhy{}\PYGZhy{}\PYGZhy{}\PYGZhy{}\PYGZhy{}\PYGZhy{}\PYGZhy{}\PYGZhy{}\PYGZhy{}\PYGZhy{}\PYGZhy{}\PYGZhy{}\PYGZhy{}\PYGZhy{}\PYGZhy{}\PYGZhy{}\PYGZhy{}\PYGZhy{}\PYGZhy{}\PYGZhy{}\PYGZhy{}\PYGZhy{}\PYGZhy{}\PYGZhy{}\PYGZhy{}}
    \PYG{c+c1}{\PYGZsh{} f\PYGZus{}t}
    \PYG{k}{def} \PYG{n+nf}{f\PYGZus{}t}\PYG{p}{(}\PYG{n+nb+bp}{self}\PYG{p}{,} \PYG{n}{t}\PYG{p}{,} \PYG{n}{seirwd}\PYG{p}{)} \PYG{p}{:}
        \PYG{n}{p}                \PYG{o}{=} \PYG{n+nb+bp}{self}\PYG{o}{.}\PYG{n}{get\PYGZus{}p}\PYG{p}{(}\PYG{n}{t}\PYG{p}{,} \PYG{n}{seirwd}\PYG{p}{)}
        \PYG{n}{p\PYGZus{}t}              \PYG{o}{=} \PYG{n+nb+bp}{self}\PYG{o}{.}\PYG{n}{get\PYGZus{}p\PYGZus{}t}\PYG{p}{(}\PYG{p}{)}
        \PYG{n}{S}\PYG{p}{,} \PYG{n}{E}\PYG{p}{,} \PYG{n}{I}\PYG{p}{,} \PYG{n}{R}\PYG{p}{,} \PYG{n}{W}\PYG{p}{,} \PYG{n}{D} \PYG{o}{=} \PYG{n}{seirwd}
        \PYG{c+c1}{\PYGZsh{}}
        \PYG{n}{Ia}     \PYG{o}{=} \PYG{n+nb}{pow}\PYG{p}{(}\PYG{n}{I}\PYG{p}{,} \PYG{n}{p}\PYG{p}{[}\PYG{l+s+s1}{\PYGZsq{}}\PYG{l+s+s1}{alpha}\PYG{l+s+s1}{\PYGZsq{}}\PYG{p}{]}\PYG{p}{)}
        \PYG{c+c1}{\PYGZsh{}}
        \PYG{c+c1}{\PYGZsh{}}
        \PYG{n}{Sdot\PYGZus{}t}   \PYG{o}{=} \PYG{o}{\PYGZhy{}} \PYG{n}{p\PYGZus{}t}\PYG{p}{[}\PYG{l+s+s1}{\PYGZsq{}}\PYG{l+s+s1}{beta}\PYG{l+s+s1}{\PYGZsq{}}\PYG{p}{]}  \PYG{o}{*}  \PYG{n}{S} \PYG{o}{*} \PYG{n}{Ia}  \PYG{o}{+} \PYG{n}{p\PYGZus{}t}\PYG{p}{[}\PYG{l+s+s1}{\PYGZsq{}}\PYG{l+s+s1}{xi}\PYG{l+s+s1}{\PYGZsq{}}\PYG{p}{]}    \PYG{o}{*} \PYG{n}{R}
        \PYG{n}{Edot\PYGZus{}t}   \PYG{o}{=} \PYG{o}{+} \PYG{n}{p\PYGZus{}t}\PYG{p}{[}\PYG{l+s+s1}{\PYGZsq{}}\PYG{l+s+s1}{beta}\PYG{l+s+s1}{\PYGZsq{}}\PYG{p}{]}  \PYG{o}{*}  \PYG{n}{S} \PYG{o}{*} \PYG{n}{Ia}  \PYG{o}{\PYGZhy{}} \PYG{n}{p\PYGZus{}t}\PYG{p}{[}\PYG{l+s+s1}{\PYGZsq{}}\PYG{l+s+s1}{sigma}\PYG{l+s+s1}{\PYGZsq{}}\PYG{p}{]} \PYG{o}{*} \PYG{n}{E}
        \PYG{n}{Idot\PYGZus{}t}   \PYG{o}{=} \PYG{o}{+} \PYG{n}{p\PYGZus{}t}\PYG{p}{[}\PYG{l+s+s1}{\PYGZsq{}}\PYG{l+s+s1}{sigma}\PYG{l+s+s1}{\PYGZsq{}}\PYG{p}{]} \PYG{o}{*} \PYG{n}{E}        \PYG{o}{\PYGZhy{}} \PYG{p}{(}\PYG{n}{p\PYGZus{}t}\PYG{p}{[}\PYG{l+s+s1}{\PYGZsq{}}\PYG{l+s+s1}{gamma}\PYG{l+s+s1}{\PYGZsq{}}\PYG{p}{]} \PYG{o}{+} \PYG{n}{p\PYGZus{}t}\PYG{p}{[}\PYG{l+s+s1}{\PYGZsq{}}\PYG{l+s+s1}{chi}\PYG{l+s+s1}{\PYGZsq{}}\PYG{p}{]}\PYG{p}{)} \PYG{o}{*} \PYG{n}{I}
        \PYG{n}{Rdot\PYGZus{}t}   \PYG{o}{=} \PYG{o}{+} \PYG{n}{p\PYGZus{}t}\PYG{p}{[}\PYG{l+s+s1}{\PYGZsq{}}\PYG{l+s+s1}{gamma}\PYG{l+s+s1}{\PYGZsq{}}\PYG{p}{]} \PYG{o}{*} \PYG{n}{I}        \PYG{o}{\PYGZhy{}} \PYG{n}{p\PYGZus{}t}\PYG{p}{[}\PYG{l+s+s1}{\PYGZsq{}}\PYG{l+s+s1}{xi}\PYG{l+s+s1}{\PYGZsq{}}\PYG{p}{]}    \PYG{o}{*} \PYG{n}{R}
        \PYG{n}{Wdot\PYGZus{}t}   \PYG{o}{=} \PYG{o}{+} \PYG{n}{p\PYGZus{}t}\PYG{p}{[}\PYG{l+s+s1}{\PYGZsq{}}\PYG{l+s+s1}{chi}\PYG{l+s+s1}{\PYGZsq{}}\PYG{p}{]}   \PYG{o}{*} \PYG{n}{I}        \PYG{o}{\PYGZhy{}} \PYG{n}{p\PYGZus{}t}\PYG{p}{[}\PYG{l+s+s1}{\PYGZsq{}}\PYG{l+s+s1}{delta}\PYG{l+s+s1}{\PYGZsq{}}\PYG{p}{]} \PYG{o}{*} \PYG{n}{W}
        \PYG{n}{Ddot\PYGZus{}t}   \PYG{o}{=} \PYG{o}{+} \PYG{n}{p\PYGZus{}t}\PYG{p}{[}\PYG{l+s+s1}{\PYGZsq{}}\PYG{l+s+s1}{delta}\PYG{l+s+s1}{\PYGZsq{}}\PYG{p}{]} \PYG{o}{*} \PYG{n}{W}
        \PYG{k}{return} \PYG{n}{numpy}\PYG{o}{.}\PYG{n}{array}\PYG{p}{(}\PYG{p}{[} \PYG{n}{Sdot\PYGZus{}t}\PYG{p}{,} \PYG{n}{Edot\PYGZus{}t}\PYG{p}{,} \PYG{n}{Idot\PYGZus{}t}\PYG{p}{,} \PYG{n}{Rdot\PYGZus{}t}\PYG{p}{,} \PYG{n}{Wdot\PYGZus{}t}\PYG{p}{,} \PYG{n}{Ddot\PYGZus{}t}\PYG{p}{]}\PYG{p}{)}
    \PYG{c+c1}{\PYGZsh{} \PYGZhy{}\PYGZhy{}\PYGZhy{}\PYGZhy{}\PYGZhy{}\PYGZhy{}\PYGZhy{}\PYGZhy{}\PYGZhy{}\PYGZhy{}\PYGZhy{}\PYGZhy{}\PYGZhy{}\PYGZhy{}\PYGZhy{}\PYGZhy{}\PYGZhy{}\PYGZhy{}\PYGZhy{}\PYGZhy{}\PYGZhy{}\PYGZhy{}\PYGZhy{}\PYGZhy{}\PYGZhy{}\PYGZhy{}\PYGZhy{}\PYGZhy{}\PYGZhy{}\PYGZhy{}\PYGZhy{}\PYGZhy{}\PYGZhy{}\PYGZhy{}\PYGZhy{}\PYGZhy{}\PYGZhy{}\PYGZhy{}\PYGZhy{}\PYGZhy{}\PYGZhy{}\PYGZhy{}\PYGZhy{}\PYGZhy{}\PYGZhy{}\PYGZhy{}\PYGZhy{}\PYGZhy{}\PYGZhy{}\PYGZhy{}\PYGZhy{}\PYGZhy{}\PYGZhy{}\PYGZhy{}\PYGZhy{}\PYGZhy{}\PYGZhy{}\PYGZhy{}\PYGZhy{}\PYGZhy{}\PYGZhy{}\PYGZhy{}\PYGZhy{}\PYGZhy{}\PYGZhy{}\PYGZhy{}\PYGZhy{}\PYGZhy{}\PYGZhy{}\PYGZhy{}\PYGZhy{}}
    \PYG{c+c1}{\PYGZsh{} f\PYGZus{}y}
    \PYG{k}{def} \PYG{n+nf}{f\PYGZus{}y}\PYG{p}{(}\PYG{n+nb+bp}{self}\PYG{p}{,} \PYG{n}{t}\PYG{p}{,} \PYG{n}{seirwd}\PYG{p}{)} \PYG{p}{:}
        \PYG{n}{p}                \PYG{o}{=} \PYG{n+nb+bp}{self}\PYG{o}{.}\PYG{n}{get\PYGZus{}p}\PYG{p}{(}\PYG{n}{t}\PYG{p}{,} \PYG{n}{seirwd}\PYG{p}{)}
        \PYG{n}{S}\PYG{p}{,} \PYG{n}{E}\PYG{p}{,} \PYG{n}{I}\PYG{p}{,} \PYG{n}{R}\PYG{p}{,} \PYG{n}{W}\PYG{p}{,} \PYG{n}{D} \PYG{o}{=} \PYG{n}{seirwd}
        \PYG{c+c1}{\PYGZsh{}}
        \PYG{n}{Ia}     \PYG{o}{=} \PYG{n+nb}{pow}\PYG{p}{(}\PYG{n}{I}\PYG{p}{,} \PYG{n}{p}\PYG{p}{[}\PYG{l+s+s1}{\PYGZsq{}}\PYG{l+s+s1}{alpha}\PYG{l+s+s1}{\PYGZsq{}}\PYG{p}{]}\PYG{p}{)}
        \PYG{k}{if} \PYG{n}{p}\PYG{p}{[}\PYG{l+s+s1}{\PYGZsq{}}\PYG{l+s+s1}{alpha}\PYG{l+s+s1}{\PYGZsq{}}\PYG{p}{]} \PYG{o}{==} \PYG{l+m+mf}{1.0} \PYG{p}{:}
            \PYG{n}{Ia\PYGZus{}I} \PYG{o}{=} \PYG{l+m+mf}{1.0}
        \PYG{k}{else} \PYG{p}{:}
            \PYG{n}{Ia\PYGZus{}I} \PYG{o}{=} \PYG{n}{p}\PYG{p}{[}\PYG{l+s+s1}{\PYGZsq{}}\PYG{l+s+s1}{alpha}\PYG{l+s+s1}{\PYGZsq{}}\PYG{p}{]} \PYG{o}{*} \PYG{n+nb}{pow}\PYG{p}{(}\PYG{n}{I}\PYG{p}{,} \PYG{n}{p}\PYG{p}{[}\PYG{l+s+s1}{\PYGZsq{}}\PYG{l+s+s1}{alpha}\PYG{l+s+s1}{\PYGZsq{}}\PYG{p}{]} \PYG{o}{\PYGZhy{}} \PYG{l+m+mf}{1.0} \PYG{p}{)}
        \PYG{c+c1}{\PYGZsh{}}
        \PYG{n}{ny}     \PYG{o}{=} \PYG{l+m+mi}{6}
        \PYG{n}{type\PYGZus{}y} \PYG{o}{=} \PYG{n+nb}{type}\PYG{p}{(} \PYG{n}{seirwd}\PYG{p}{[}\PYG{l+m+mi}{0}\PYG{p}{]} \PYG{p}{)}
        \PYG{n}{zero}   \PYG{o}{=} \PYG{n}{type\PYGZus{}y}\PYG{p}{(} \PYG{l+m+mf}{0.0} \PYG{p}{)}
        \PYG{n}{J}      \PYG{o}{=} \PYG{n}{numpy}\PYG{o}{.}\PYG{n}{empty}\PYG{p}{(} \PYG{p}{(}\PYG{n}{ny}\PYG{p}{,}\PYG{n}{ny}\PYG{p}{)}\PYG{p}{,} \PYG{n}{dtype} \PYG{o}{=} \PYG{n}{type\PYGZus{}y} \PYG{p}{)}
        \PYG{c+c1}{\PYGZsh{}}
        \PYG{c+c1}{\PYGZsh{} partials of Sdot(t, y) w.r.t. y}
        \PYG{n}{Sdot\PYGZus{}S}   \PYG{o}{=} \PYG{o}{\PYGZhy{}} \PYG{n}{p}\PYG{p}{[}\PYG{l+s+s1}{\PYGZsq{}}\PYG{l+s+s1}{beta}\PYG{l+s+s1}{\PYGZsq{}}\PYG{p}{]}  \PYG{o}{*}  \PYG{n}{Ia}
        \PYG{n}{Sdot\PYGZus{}I}   \PYG{o}{=} \PYG{o}{\PYGZhy{}} \PYG{n}{p}\PYG{p}{[}\PYG{l+s+s1}{\PYGZsq{}}\PYG{l+s+s1}{beta}\PYG{l+s+s1}{\PYGZsq{}}\PYG{p}{]}  \PYG{o}{*}  \PYG{n}{S} \PYG{o}{*} \PYG{n}{Ia\PYGZus{}I}
        \PYG{n}{Sdot\PYGZus{}R}   \PYG{o}{=} \PYG{o}{+} \PYG{n}{p}\PYG{p}{[}\PYG{l+s+s1}{\PYGZsq{}}\PYG{l+s+s1}{xi}\PYG{l+s+s1}{\PYGZsq{}}\PYG{p}{]}
        \PYG{n}{J}\PYG{p}{[}\PYG{l+m+mi}{0}\PYG{p}{,}\PYG{p}{:}\PYG{p}{]}   \PYG{o}{=} \PYG{p}{[} \PYG{n}{Sdot\PYGZus{}S}\PYG{p}{,} \PYG{n}{zero}\PYG{p}{,} \PYG{n}{Sdot\PYGZus{}I}\PYG{p}{,} \PYG{n}{Sdot\PYGZus{}R}\PYG{p}{,} \PYG{n}{zero}\PYG{p}{,} \PYG{n}{zero} \PYG{p}{]}
        \PYG{c+c1}{\PYGZsh{}}
        \PYG{c+c1}{\PYGZsh{} partial of Edot(t, y) w.r.t y}
        \PYG{n}{Edot\PYGZus{}S}   \PYG{o}{=} \PYG{o}{+} \PYG{n}{p}\PYG{p}{[}\PYG{l+s+s1}{\PYGZsq{}}\PYG{l+s+s1}{beta}\PYG{l+s+s1}{\PYGZsq{}}\PYG{p}{]}  \PYG{o}{*}  \PYG{n}{Ia}
        \PYG{n}{Edot\PYGZus{}I}   \PYG{o}{=} \PYG{o}{+} \PYG{n}{p}\PYG{p}{[}\PYG{l+s+s1}{\PYGZsq{}}\PYG{l+s+s1}{beta}\PYG{l+s+s1}{\PYGZsq{}}\PYG{p}{]}  \PYG{o}{*}  \PYG{n}{S} \PYG{o}{*} \PYG{n}{Ia\PYGZus{}I}
        \PYG{n}{Edot\PYGZus{}E}   \PYG{o}{=}  \PYG{o}{\PYGZhy{}} \PYG{n}{p}\PYG{p}{[}\PYG{l+s+s1}{\PYGZsq{}}\PYG{l+s+s1}{sigma}\PYG{l+s+s1}{\PYGZsq{}}\PYG{p}{]}
        \PYG{n}{J}\PYG{p}{[}\PYG{l+m+mi}{1}\PYG{p}{,}\PYG{p}{:}\PYG{p}{]}   \PYG{o}{=} \PYG{p}{[} \PYG{n}{Edot\PYGZus{}S}\PYG{p}{,} \PYG{n}{Edot\PYGZus{}E}\PYG{p}{,} \PYG{n}{Edot\PYGZus{}I}\PYG{p}{,} \PYG{n}{zero}\PYG{p}{,} \PYG{n}{zero}\PYG{p}{,} \PYG{n}{zero} \PYG{p}{]}
        \PYG{c+c1}{\PYGZsh{}}
        \PYG{c+c1}{\PYGZsh{} partial of Idot(t, y) w.r.t y}
        \PYG{n}{Idot\PYGZus{}E}   \PYG{o}{=} \PYG{o}{+} \PYG{n}{p}\PYG{p}{[}\PYG{l+s+s1}{\PYGZsq{}}\PYG{l+s+s1}{sigma}\PYG{l+s+s1}{\PYGZsq{}}\PYG{p}{]}
        \PYG{n}{Idot\PYGZus{}I}   \PYG{o}{=} \PYG{o}{\PYGZhy{}} \PYG{p}{(}\PYG{n}{p}\PYG{p}{[}\PYG{l+s+s1}{\PYGZsq{}}\PYG{l+s+s1}{gamma}\PYG{l+s+s1}{\PYGZsq{}}\PYG{p}{]} \PYG{o}{+} \PYG{n}{p}\PYG{p}{[}\PYG{l+s+s1}{\PYGZsq{}}\PYG{l+s+s1}{chi}\PYG{l+s+s1}{\PYGZsq{}}\PYG{p}{]}\PYG{p}{)}
        \PYG{n}{J}\PYG{p}{[}\PYG{l+m+mi}{2}\PYG{p}{,}\PYG{p}{:}\PYG{p}{]}   \PYG{o}{=} \PYG{p}{[} \PYG{n}{zero}\PYG{p}{,} \PYG{n}{Idot\PYGZus{}E}\PYG{p}{,} \PYG{n}{Idot\PYGZus{}I}\PYG{p}{,} \PYG{n}{zero}\PYG{p}{,} \PYG{n}{zero}\PYG{p}{,} \PYG{n}{zero} \PYG{p}{]}
        \PYG{c+c1}{\PYGZsh{}}
        \PYG{c+c1}{\PYGZsh{} partial Rdot(t, y) w.r.t y}
        \PYG{n}{Rdot\PYGZus{}I}   \PYG{o}{=} \PYG{o}{+} \PYG{n}{p}\PYG{p}{[}\PYG{l+s+s1}{\PYGZsq{}}\PYG{l+s+s1}{gamma}\PYG{l+s+s1}{\PYGZsq{}}\PYG{p}{]}
        \PYG{n}{Rdot\PYGZus{}R}   \PYG{o}{=} \PYG{o}{\PYGZhy{}} \PYG{n}{p}\PYG{p}{[}\PYG{l+s+s1}{\PYGZsq{}}\PYG{l+s+s1}{xi}\PYG{l+s+s1}{\PYGZsq{}}\PYG{p}{]}
        \PYG{n}{J}\PYG{p}{[}\PYG{l+m+mi}{3}\PYG{p}{,}\PYG{p}{:}\PYG{p}{]}   \PYG{o}{=} \PYG{p}{[} \PYG{n}{zero}\PYG{p}{,} \PYG{n}{zero}\PYG{p}{,} \PYG{n}{Rdot\PYGZus{}I}\PYG{p}{,} \PYG{n}{Rdot\PYGZus{}R}\PYG{p}{,} \PYG{n}{zero}\PYG{p}{,} \PYG{n}{zero} \PYG{p}{]}
        \PYG{c+c1}{\PYGZsh{}}
        \PYG{c+c1}{\PYGZsh{} partial Wdot(t, y) w.r.t. y}
        \PYG{n}{Wdot\PYGZus{}I}   \PYG{o}{=} \PYG{o}{+} \PYG{n}{p}\PYG{p}{[}\PYG{l+s+s1}{\PYGZsq{}}\PYG{l+s+s1}{chi}\PYG{l+s+s1}{\PYGZsq{}}\PYG{p}{]}
        \PYG{n}{Wdot\PYGZus{}W}   \PYG{o}{=} \PYG{o}{\PYGZhy{}} \PYG{n}{p}\PYG{p}{[}\PYG{l+s+s1}{\PYGZsq{}}\PYG{l+s+s1}{delta}\PYG{l+s+s1}{\PYGZsq{}}\PYG{p}{]}
        \PYG{n}{J}\PYG{p}{[}\PYG{l+m+mi}{4}\PYG{p}{,}\PYG{p}{:}\PYG{p}{]}   \PYG{o}{=} \PYG{p}{[} \PYG{n}{zero}\PYG{p}{,} \PYG{n}{zero}\PYG{p}{,} \PYG{n}{Wdot\PYGZus{}I}\PYG{p}{,} \PYG{n}{zero}\PYG{p}{,} \PYG{n}{Wdot\PYGZus{}W}\PYG{p}{,} \PYG{n}{zero} \PYG{p}{]}
        \PYG{c+c1}{\PYGZsh{}}
        \PYG{c+c1}{\PYGZsh{} partial Ddot(t, y) w.r.t y}
        \PYG{n}{Ddot\PYGZus{}W}   \PYG{o}{=} \PYG{o}{+} \PYG{n}{p}\PYG{p}{[}\PYG{l+s+s1}{\PYGZsq{}}\PYG{l+s+s1}{delta}\PYG{l+s+s1}{\PYGZsq{}}\PYG{p}{]}
        \PYG{n}{J}\PYG{p}{[}\PYG{l+m+mi}{5}\PYG{p}{,}\PYG{p}{:}\PYG{p}{]}   \PYG{o}{=} \PYG{p}{[} \PYG{n}{zero}\PYG{p}{,} \PYG{n}{zero}\PYG{p}{,} \PYG{n}{zero}\PYG{p}{,} \PYG{n}{zero}\PYG{p}{,} \PYG{n}{Ddot\PYGZus{}W}\PYG{p}{,} \PYG{n}{zero} \PYG{p}{]}
        \PYG{c+c1}{\PYGZsh{}}
        \PYG{k}{return} \PYG{n}{J}
\PYG{c+c1}{\PYGZsh{} \PYGZhy{}\PYGZhy{}\PYGZhy{}\PYGZhy{}\PYGZhy{}\PYGZhy{}\PYGZhy{}\PYGZhy{}\PYGZhy{}\PYGZhy{}\PYGZhy{}\PYGZhy{}\PYGZhy{}\PYGZhy{}\PYGZhy{}\PYGZhy{}\PYGZhy{}\PYGZhy{}\PYGZhy{}\PYGZhy{}\PYGZhy{}\PYGZhy{}\PYGZhy{}\PYGZhy{}\PYGZhy{}\PYGZhy{}\PYGZhy{}\PYGZhy{}\PYGZhy{}\PYGZhy{}\PYGZhy{}\PYGZhy{}\PYGZhy{}\PYGZhy{}\PYGZhy{}\PYGZhy{}\PYGZhy{}\PYGZhy{}\PYGZhy{}\PYGZhy{}\PYGZhy{}\PYGZhy{}\PYGZhy{}\PYGZhy{}\PYGZhy{}\PYGZhy{}\PYGZhy{}\PYGZhy{}\PYGZhy{}\PYGZhy{}\PYGZhy{}\PYGZhy{}\PYGZhy{}\PYGZhy{}\PYGZhy{}\PYGZhy{}\PYGZhy{}\PYGZhy{}\PYGZhy{}\PYGZhy{}\PYGZhy{}\PYGZhy{}\PYGZhy{}\PYGZhy{}\PYGZhy{}\PYGZhy{}\PYGZhy{}\PYGZhy{}\PYGZhy{}\PYGZhy{}\PYGZhy{}\PYGZhy{}\PYGZhy{}\PYGZhy{}\PYGZhy{}\PYGZhy{}\PYGZhy{}}
\PYG{c+c1}{\PYGZsh{} check the fun derivatives required by rosen3\PYGZus{}step}
\PYG{k}{def} \PYG{n+nf}{test\PYGZus{}fun\PYGZus{}derivatives}\PYG{p}{(}\PYG{n}{t\PYGZus{}all}\PYG{p}{,} \PYG{n}{p\PYGZus{}all}\PYG{p}{,} \PYG{n}{initial}\PYG{p}{)} \PYG{p}{:}
    \PYG{n}{ok}     \PYG{o}{=} \PYG{k+kc}{True}
    \PYG{n}{n\PYGZus{}step} \PYG{o}{=} \PYG{l+m+mi}{1}
    \PYG{n}{fun}  \PYG{o}{=} \PYG{n}{fun\PYGZus{}class}\PYG{p}{(}\PYG{n}{t\PYGZus{}all}\PYG{p}{,} \PYG{n}{p\PYGZus{}all}\PYG{p}{,} \PYG{n}{n\PYGZus{}step}\PYG{p}{)}
    \PYG{n}{index}  \PYG{o}{=} \PYG{l+m+mi}{0}
    \PYG{n}{fun}\PYG{o}{.}\PYG{n}{set\PYGZus{}t\PYGZus{}all\PYGZus{}index}\PYG{p}{(}\PYG{n}{index}\PYG{p}{)}
    \PYG{n}{ti}   \PYG{o}{=} \PYG{n}{t\PYGZus{}all}\PYG{p}{[}\PYG{n}{index}\PYG{p}{]}
    \PYG{n}{yi}   \PYG{o}{=} \PYG{n}{initial}
    \PYG{n}{h}    \PYG{o}{=} \PYG{n}{t\PYGZus{}all}\PYG{p}{[}\PYG{n}{index} \PYG{o}{+} \PYG{l+m+mi}{1}\PYG{p}{]} \PYG{o}{\PYGZhy{}} \PYG{n}{t\PYGZus{}all}\PYG{p}{[}\PYG{n}{index}\PYG{p}{]}
    \PYG{n}{ok}   \PYG{o}{=} \PYG{n}{ok} \PYG{o+ow}{and} \PYG{n}{check\PYGZus{}rosen3\PYGZus{}step}\PYG{p}{(}\PYG{n}{fun}\PYG{p}{,} \PYG{n}{ti}\PYG{p}{,} \PYG{n}{yi}\PYG{p}{,} \PYG{n}{h}\PYG{p}{)}
    \PYG{k}{return} \PYG{n}{ok}
\PYG{c+c1}{\PYGZsh{} \PYGZhy{}\PYGZhy{}\PYGZhy{}\PYGZhy{}\PYGZhy{}\PYGZhy{}\PYGZhy{}\PYGZhy{}\PYGZhy{}\PYGZhy{}\PYGZhy{}\PYGZhy{}\PYGZhy{}\PYGZhy{}\PYGZhy{}\PYGZhy{}\PYGZhy{}\PYGZhy{}\PYGZhy{}\PYGZhy{}\PYGZhy{}\PYGZhy{}\PYGZhy{}\PYGZhy{}\PYGZhy{}\PYGZhy{}\PYGZhy{}\PYGZhy{}\PYGZhy{}\PYGZhy{}\PYGZhy{}\PYGZhy{}\PYGZhy{}\PYGZhy{}\PYGZhy{}\PYGZhy{}\PYGZhy{}\PYGZhy{}\PYGZhy{}\PYGZhy{}\PYGZhy{}\PYGZhy{}\PYGZhy{}\PYGZhy{}\PYGZhy{}\PYGZhy{}\PYGZhy{}\PYGZhy{}\PYGZhy{}\PYGZhy{}\PYGZhy{}\PYGZhy{}\PYGZhy{}\PYGZhy{}\PYGZhy{}\PYGZhy{}\PYGZhy{}\PYGZhy{}\PYGZhy{}\PYGZhy{}\PYGZhy{}\PYGZhy{}\PYGZhy{}\PYGZhy{}\PYGZhy{}\PYGZhy{}\PYGZhy{}\PYGZhy{}\PYGZhy{}\PYGZhy{}\PYGZhy{}\PYGZhy{}\PYGZhy{}\PYGZhy{}\PYGZhy{}\PYGZhy{}\PYGZhy{}}
\PYG{c+c1}{\PYGZsh{} seriwd\PYGZus{}model}
\PYG{k}{def} \PYG{n+nf}{seirwd\PYGZus{}model}\PYG{p}{(}\PYG{n}{method}\PYG{p}{,} \PYG{n}{t\PYGZus{}all}\PYG{p}{,} \PYG{n}{p\PYGZus{}all}\PYG{p}{,} \PYG{n}{initial}\PYG{p}{,} \PYG{n}{n\PYGZus{}step} \PYG{o}{=} \PYG{l+m+mi}{1}\PYG{p}{)} \PYG{p}{:}
    \PYG{k}{for} \PYG{n}{i} \PYG{o+ow}{in} \PYG{n+nb}{range}\PYG{p}{(} \PYG{n+nb}{len}\PYG{p}{(}\PYG{n}{t\PYGZus{}all}\PYG{p}{)} \PYG{o}{\PYGZhy{}} \PYG{l+m+mi}{1} \PYG{p}{)} \PYG{p}{:}
        \PYG{k}{assert} \PYG{n}{p\PYGZus{}all}\PYG{p}{[}\PYG{n}{i}\PYG{o}{+}\PYG{l+m+mi}{1}\PYG{p}{]}\PYG{p}{[}\PYG{l+s+s1}{\PYGZsq{}}\PYG{l+s+s1}{alpha}\PYG{l+s+s1}{\PYGZsq{}}\PYG{p}{]} \PYG{o}{==} \PYG{n}{p\PYGZus{}all}\PYG{p}{[}\PYG{l+m+mi}{0}\PYG{p}{]}\PYG{p}{[}\PYG{l+s+s1}{\PYGZsq{}}\PYG{l+s+s1}{alpha}\PYG{l+s+s1}{\PYGZsq{}}\PYG{p}{]}
    \PYG{c+c1}{\PYGZsh{}}
    \PYG{n}{fun}      \PYG{o}{=} \PYG{n}{fun\PYGZus{}class}\PYG{p}{(}\PYG{n}{t\PYGZus{}all}\PYG{p}{,} \PYG{n}{p\PYGZus{}all}\PYG{p}{,} \PYG{n}{n\PYGZus{}step}\PYG{p}{)}
    \PYG{k}{if} \PYG{n}{method} \PYG{o}{==} \PYG{l+s+s1}{\PYGZsq{}}\PYG{l+s+s1}{runge4}\PYG{l+s+s1}{\PYGZsq{}} \PYG{p}{:}
        \PYG{n}{one\PYGZus{}step} \PYG{o}{=} \PYG{n}{runge4\PYGZus{}step}
    \PYG{k}{elif} \PYG{n}{method} \PYG{o}{==} \PYG{l+s+s1}{\PYGZsq{}}\PYG{l+s+s1}{rosen3}\PYG{l+s+s1}{\PYGZsq{}} \PYG{p}{:}
        \PYG{n}{one\PYGZus{}step} \PYG{o}{=} \PYG{n}{rosen3\PYGZus{}step}
    \PYG{k}{else} \PYG{p}{:}
        \PYG{k}{assert} \PYG{k+kc}{False}
    \PYG{k}{if} \PYG{n}{n\PYGZus{}step} \PYG{o}{==} \PYG{l+m+mi}{1} \PYG{p}{:}
        \PYG{n}{seirwd\PYGZus{}all}  \PYG{o}{=} \PYG{n}{ode\PYGZus{}multi\PYGZus{}step}\PYG{p}{(}\PYG{n}{one\PYGZus{}step}\PYG{p}{,} \PYG{n}{fun}\PYG{p}{,} \PYG{n}{t\PYGZus{}all}\PYG{p}{,} \PYG{n}{initial}\PYG{p}{)}
    \PYG{k}{else} \PYG{p}{:}
        \PYG{c+c1}{\PYGZsh{} t\PYGZus{}refine}
        \PYG{n}{t\PYGZus{}refine} \PYG{o}{=} \PYG{n+nb}{list}\PYG{p}{(}\PYG{p}{)}
        \PYG{k}{for} \PYG{n}{i} \PYG{o+ow}{in} \PYG{n+nb}{range}\PYG{p}{(} \PYG{n+nb}{len}\PYG{p}{(}\PYG{n}{t\PYGZus{}all}\PYG{p}{)} \PYG{o}{\PYGZhy{}} \PYG{l+m+mi}{1} \PYG{p}{)} \PYG{p}{:}
            \PYG{n}{step} \PYG{o}{=} \PYG{p}{(}\PYG{n}{t\PYGZus{}all}\PYG{p}{[}\PYG{n}{i}\PYG{o}{+}\PYG{l+m+mi}{1}\PYG{p}{]} \PYG{o}{\PYGZhy{}} \PYG{n}{t\PYGZus{}all}\PYG{p}{[}\PYG{n}{i}\PYG{p}{]}\PYG{p}{)} \PYG{o}{/} \PYG{n}{n\PYGZus{}step}
            \PYG{k}{for} \PYG{n}{j} \PYG{o+ow}{in} \PYG{n+nb}{range}\PYG{p}{(}\PYG{n}{n\PYGZus{}step}\PYG{p}{)} \PYG{p}{:}
                \PYG{n}{t\PYGZus{}refine}\PYG{o}{.}\PYG{n}{append}\PYG{p}{(} \PYG{n}{t\PYGZus{}all}\PYG{p}{[}\PYG{n}{i}\PYG{p}{]} \PYG{o}{+} \PYG{n}{j} \PYG{o}{*} \PYG{n}{step} \PYG{p}{)}
        \PYG{n}{t\PYGZus{}refine}\PYG{o}{.}\PYG{n}{append}\PYG{p}{(} \PYG{n}{t\PYGZus{}all}\PYG{p}{[}\PYG{o}{\PYGZhy{}}\PYG{l+m+mi}{1}\PYG{p}{]} \PYG{p}{)}
        \PYG{n}{t\PYGZus{}refine} \PYG{o}{=} \PYG{n}{numpy}\PYG{o}{.}\PYG{n}{array}\PYG{p}{(} \PYG{n}{t\PYGZus{}refine} \PYG{p}{)}
        \PYG{c+c1}{\PYGZsh{}}
        \PYG{n}{seirwd\PYGZus{}refine}   \PYG{o}{=} \PYG{n}{ode\PYGZus{}multi\PYGZus{}step}\PYG{p}{(}\PYG{n}{one\PYGZus{}step}\PYG{p}{,} \PYG{n}{fun}\PYG{p}{,} \PYG{n}{t\PYGZus{}refine}\PYG{p}{,} \PYG{n}{initial}\PYG{p}{)}
        \PYG{n}{index\PYGZus{}all}   \PYG{o}{=} \PYG{p}{[} \PYG{n}{i} \PYG{o}{*} \PYG{n}{n\PYGZus{}step} \PYG{k}{for} \PYG{n}{i} \PYG{o+ow}{in} \PYG{n+nb}{range}\PYG{p}{(}\PYG{n+nb}{len}\PYG{p}{(}\PYG{n}{t\PYGZus{}all}\PYG{p}{)}\PYG{p}{)} \PYG{p}{]}
        \PYG{n}{seirwd\PYGZus{}all}   \PYG{o}{=} \PYG{n}{seirwd\PYGZus{}refine}\PYG{p}{[} \PYG{n}{index\PYGZus{}all} \PYG{p}{]}
    \PYG{c+c1}{\PYGZsh{}}
    \PYG{k}{return} \PYG{n}{seirwd\PYGZus{}all}
\end{sphinxVerbatim}


\bigskip\hrule\bigskip


xsrst input file: \sphinxcode{\sphinxupquote{example/python/numeric/seirwd\_model.py}}


\subparagraph{numeric\_covid\_19\_xam\_py}
\label{\detokenize{xsrst/numeric_covid_19_xam_py:numeric-covid-19-xam-py}}\label{\detokenize{xsrst/numeric_covid_19_xam_py::doc}}
\index{numeric\_covid\_19\_xam\_py@\spxentry{numeric\_covid\_19\_xam\_py}}\index{example@\spxentry{example}}\index{fitting@\spxentry{fitting}}\index{an@\spxentry{an}}\index{seirwd@\spxentry{seirwd}}\index{model@\spxentry{model}}\index{for@\spxentry{for}}\index{covid\sphinxhyphen{}19@\spxentry{covid\sphinxhyphen{}19}}\ignorespaces 

\subparagraph{2.7 Example Fitting an SEIRWD Model for Covid\sphinxhyphen{}19}
\label{\detokenize{xsrst/numeric_covid_19_xam_py:example-fitting-an-seirwd-model-for-covid-19}}\label{\detokenize{xsrst/numeric_covid_19_xam_py:id1}}\begin{itemize}
\item {} 
{\hyperref[\detokenize{xsrst/numeric_covid_19_xam_py:numeric-covid-19-xam-py-covariates}]{\sphinxcrossref{\DUrole{std,std-ref}{Covariates}}}}

\item {} \begin{description}
\item[{{\hyperref[\detokenize{xsrst/numeric_covid_19_xam_py:numeric-covid-19-xam-py-model-ode}]{\sphinxcrossref{\DUrole{std,std-ref}{Model ODE}}}}}] \leavevmode\begin{itemize}
\item {} 
{\hyperref[\detokenize{xsrst/numeric_covid_19_xam_py:numeric-covid-19-xam-py-model-ode-beta-t}]{\sphinxcrossref{\DUrole{std,std-ref}{beta(t)}}}}

\item {} 
{\hyperref[\detokenize{xsrst/numeric_covid_19_xam_py:numeric-covid-19-xam-py-model-ode-other-rates}]{\sphinxcrossref{\DUrole{std,std-ref}{Other Rates}}}}

\item {} 
{\hyperref[\detokenize{xsrst/numeric_covid_19_xam_py:numeric-covid-19-xam-py-model-ode-initial-values}]{\sphinxcrossref{\DUrole{std,std-ref}{Initial Values}}}}

\item {} 
{\hyperref[\detokenize{xsrst/numeric_covid_19_xam_py:numeric-covid-19-xam-py-model-ode-ode-solver}]{\sphinxcrossref{\DUrole{std,std-ref}{Ode Solver}}}}

\end{itemize}

\end{description}

\item {} \begin{description}
\item[{{\hyperref[\detokenize{xsrst/numeric_covid_19_xam_py:numeric-covid-19-xam-py-unknown-parameters}]{\sphinxcrossref{\DUrole{std,std-ref}{Unknown Parameters}}}}}] \leavevmode\begin{itemize}
\item {} 
{\hyperref[\detokenize{xsrst/numeric_covid_19_xam_py:numeric-covid-19-xam-py-unknown-parameters-maximum-likelihood}]{\sphinxcrossref{\DUrole{std,std-ref}{Maximum Likelihood}}}}

\item {} 
{\hyperref[\detokenize{xsrst/numeric_covid_19_xam_py:numeric-covid-19-xam-py-unknown-parameters-model-bounds}]{\sphinxcrossref{\DUrole{std,std-ref}{Model Bounds}}}}

\item {} 
{\hyperref[\detokenize{xsrst/numeric_covid_19_xam_py:numeric-covid-19-xam-py-unknown-parameters-actual-bounds}]{\sphinxcrossref{\DUrole{std,std-ref}{Actual Bounds}}}}

\end{itemize}

\end{description}

\item {} \begin{description}
\item[{{\hyperref[\detokenize{xsrst/numeric_covid_19_xam_py:numeric-covid-19-xam-py-data}]{\sphinxcrossref{\DUrole{std,std-ref}{Data}}}}}] \leavevmode\begin{itemize}
\item {} 
{\hyperref[\detokenize{xsrst/numeric_covid_19_xam_py:numeric-covid-19-xam-py-data-sample-interval}]{\sphinxcrossref{\DUrole{std,std-ref}{sample\_interval}}}}

\item {} 
{\hyperref[\detokenize{xsrst/numeric_covid_19_xam_py:numeric-covid-19-xam-py-data-data-file}]{\sphinxcrossref{\DUrole{std,std-ref}{data\_file}}}}

\item {} 
{\hyperref[\detokenize{xsrst/numeric_covid_19_xam_py:numeric-covid-19-xam-py-data-coefficient-of-variation}]{\sphinxcrossref{\DUrole{std,std-ref}{Coefficient of Variation}}}}

\item {} 
{\hyperref[\detokenize{xsrst/numeric_covid_19_xam_py:numeric-covid-19-xam-py-data-simulation}]{\sphinxcrossref{\DUrole{std,std-ref}{Simulation}}}}

\item {} 
{\hyperref[\detokenize{xsrst/numeric_covid_19_xam_py:numeric-covid-19-xam-py-data-weighted-residuals}]{\sphinxcrossref{\DUrole{std,std-ref}{Weighted Residuals}}}}

\end{itemize}

\end{description}

\item {} 
{\hyperref[\detokenize{xsrst/numeric_covid_19_xam_py:numeric-covid-19-xam-py-random-seed}]{\sphinxcrossref{\DUrole{std,std-ref}{Random Seed}}}}

\item {} 
{\hyperref[\detokenize{xsrst/numeric_covid_19_xam_py:numeric-covid-19-xam-py-random-start}]{\sphinxcrossref{\DUrole{std,std-ref}{Random Start}}}}

\item {} \begin{description}
\item[{{\hyperref[\detokenize{xsrst/numeric_covid_19_xam_py:numeric-covid-19-xam-py-display-fit-results}]{\sphinxcrossref{\DUrole{std,std-ref}{Display Fit Results}}}}}] \leavevmode\begin{itemize}
\item {} 
{\hyperref[\detokenize{xsrst/numeric_covid_19_xam_py:numeric-covid-19-xam-py-display-fit-results-plot}]{\sphinxcrossref{\DUrole{std,std-ref}{Plot}}}}

\item {} 
{\hyperref[\detokenize{xsrst/numeric_covid_19_xam_py:numeric-covid-19-xam-py-display-fit-results-printout}]{\sphinxcrossref{\DUrole{std,std-ref}{Printout}}}}

\end{itemize}

\end{description}

\item {} 
{\hyperref[\detokenize{xsrst/numeric_covid_19_xam_py:numeric-covid-19-xam-py-debug-output}]{\sphinxcrossref{\DUrole{std,std-ref}{Debug Output}}}}

\item {} 
{\hyperref[\detokenize{xsrst/numeric_covid_19_xam_py:numeric-covid-19-xam-py-source-code}]{\sphinxcrossref{\DUrole{std,std-ref}{Source Code}}}}

\end{itemize}

\index{covariates@\spxentry{covariates}}\ignorespaces 

\subparagraph{Covariates}
\label{\detokenize{xsrst/numeric_covid_19_xam_py:covariates}}\label{\detokenize{xsrst/numeric_covid_19_xam_py:numeric-covid-19-xam-py-covariates}}\label{\detokenize{xsrst/numeric_covid_19_xam_py:index-1}}
In this example there are two covariates that
affect the infectious rate \(\beta\):
social mobility \(c_0 (t)\),
Covid\sphinxhyphen{}19 testing \(c_1 (t)\), and
scaled time \(c_2 (t)\).
The covariates are known functions of time.
The mobility covariate has been shifted and scaled
so it is in the interval {[}\sphinxhyphen{}1, 0{]}.
The testing covariate has been shifted and scaled
so it is in the interval {[}0, 1{]}.
Note that the maximum mobility and the minimum testing corresponds to the
normal (baseline) condition.
The scaled time covariate is shifted and scaled version of time so that
it is in the interval {[}0, 1{]} with the first data point corresponding to
time zero. It is assumed here that the baseline condition corresponds
to time zero.

\index{model@\spxentry{model}}\index{ode@\spxentry{ode}}\ignorespaces 

\subparagraph{Model ODE}
\label{\detokenize{xsrst/numeric_covid_19_xam_py:model-ode}}\label{\detokenize{xsrst/numeric_covid_19_xam_py:numeric-covid-19-xam-py-model-ode}}\label{\detokenize{xsrst/numeric_covid_19_xam_py:index-2}}
We use the {\hyperref[\detokenize{xsrst/numeric_seirwd_model:id1}]{\sphinxcrossref{\DUrole{std,std-ref}{seirwd}}}} model and notation.

\index{beta(t)@\spxentry{beta(t)}}\ignorespaces 

\subparagraph{beta(t)}
\label{\detokenize{xsrst/numeric_covid_19_xam_py:beta-t}}\label{\detokenize{xsrst/numeric_covid_19_xam_py:numeric-covid-19-xam-py-model-ode-beta-t}}\label{\detokenize{xsrst/numeric_covid_19_xam_py:index-3}}
Our model for the infectious rate is
\begin{equation*}
\begin{split}\beta(t) = \bar{\beta} \exp[ m_0 c_0 (t) + m_1 c_1 (t) + m_2 c_2 (t) ]\end{split}
\end{equation*}
where \(\bar{\beta}\) is the baseline value for the infectious rate,
\(m_0\) is the social covariate multiplier, and
\(m_1\) is the Covid\sphinxhyphen{}19 testing covariate multiplier.
The baseline value \(\bar{\beta}\) is the infectious rate corresponding
to all the covariates being zero.
The covariate multipliers, and the baseline infectious rate, are unknown.

\index{other@\spxentry{other}}\index{rates@\spxentry{rates}}\ignorespaces 

\subparagraph{Other Rates}
\label{\detokenize{xsrst/numeric_covid_19_xam_py:other-rates}}\label{\detokenize{xsrst/numeric_covid_19_xam_py:numeric-covid-19-xam-py-model-ode-other-rates}}\label{\detokenize{xsrst/numeric_covid_19_xam_py:index-4}}
The other rates
\(\alpha(t)\),
\(\sigma(t)\),
\(\gamma(t)\),
\(\xi(t)\),
\(\chi(t)\),
\(\delta(t)\),
constant functions with known values:

\begin{sphinxVerbatim}[commandchars=\\\{\}]
\PYG{n}{alpha\PYGZus{}known}  \PYG{o}{=} \PYG{l+m+mf}{0.95}
\PYG{n}{sigma\PYGZus{}known}  \PYG{o}{=} \PYG{l+m+mf}{0.2}
\PYG{n}{gamma\PYGZus{}known}  \PYG{o}{=} \PYG{l+m+mf}{0.1}
\PYG{n}{chi\PYGZus{}known}    \PYG{o}{=} \PYG{l+m+mf}{0.1}
\PYG{n}{xi\PYGZus{}known}     \PYG{o}{=} \PYG{l+m+mf}{0.00}
\PYG{n}{delta\PYGZus{}known}  \PYG{o}{=} \PYG{l+m+mf}{0.2}
\end{sphinxVerbatim}

All of theses rates must be non\sphinxhyphen{}negative.

\index{initial@\spxentry{initial}}\index{values@\spxentry{values}}\ignorespaces 

\subparagraph{Initial Values}
\label{\detokenize{xsrst/numeric_covid_19_xam_py:initial-values}}\label{\detokenize{xsrst/numeric_covid_19_xam_py:numeric-covid-19-xam-py-model-ode-initial-values}}\label{\detokenize{xsrst/numeric_covid_19_xam_py:index-5}}
The initial size of the Recovered group \(R(0)\)
and of the Death group \(D(0)\) is zero.
We use fraction of the total population for sizes, so the sum of the
other initial values is one.
We treat the initial
Infected group \(I(0)\), and
Will die group \(W(0)\),
as unknown parameters in the model.
We would like to also solve for the initial exposed population but
that model has identifiability problems, so
we use the following approximation for the initial exposed group
\begin{equation*}
\begin{split}E(0) = I(0) \gamma / \sigma\end{split}
\end{equation*}
The initial Susceptible group \(S(0)\) is
expressed as a function of the other initial conditions:
\begin{equation*}
\begin{split}S(0) = 1 - E(0) - I(0) - W(0)\end{split}
\end{equation*}
\index{ode@\spxentry{ode}}\index{solver@\spxentry{solver}}\ignorespaces 

\subparagraph{Ode Solver}
\label{\detokenize{xsrst/numeric_covid_19_xam_py:ode-solver}}\label{\detokenize{xsrst/numeric_covid_19_xam_py:numeric-covid-19-xam-py-model-ode-ode-solver}}\label{\detokenize{xsrst/numeric_covid_19_xam_py:index-6}}
There are two choices for \sphinxstyleemphasis{ode\_method} ,
the method used to solve the ODE:
{\hyperref[\detokenize{xsrst/numeric_runge4_step:id1}]{\sphinxcrossref{\DUrole{std,std-ref}{runge4}}}} and
{\hyperref[\detokenize{xsrst/numeric_rosen3_step:id1}]{\sphinxcrossref{\DUrole{std,std-ref}{rosen3}}}}.
In addition, we can choose \sphinxstyleemphasis{ode\_n\_step} ,
the number of step to take for each time interval in \sphinxstyleemphasis{t\_all} ,
before it is sub\sphinxhyphen{}sampled using the
{\hyperref[\detokenize{xsrst/numeric_covid_19_xam_py:numeric-covid-19-xam-py-data-sample-interval}]{\sphinxcrossref{\DUrole{std,std-ref}{sample\_interval}}}}.

\begin{sphinxVerbatim}[commandchars=\\\{\}]
\PYG{n}{ode\PYGZus{}method} \PYG{o}{=} \PYG{l+s+s1}{\PYGZsq{}}\PYG{l+s+s1}{runge4}\PYG{l+s+s1}{\PYGZsq{}}
\PYG{n}{ode\PYGZus{}n\PYGZus{}step} \PYG{o}{=} \PYG{l+m+mi}{4}
\end{sphinxVerbatim}

\index{unknown@\spxentry{unknown}}\index{parameters@\spxentry{parameters}}\ignorespaces 

\subparagraph{Unknown Parameters}
\label{\detokenize{xsrst/numeric_covid_19_xam_py:unknown-parameters}}\label{\detokenize{xsrst/numeric_covid_19_xam_py:numeric-covid-19-xam-py-unknown-parameters}}\label{\detokenize{xsrst/numeric_covid_19_xam_py:index-7}}
The unknown parameter vector in this model is
\begin{equation*}
\begin{split}x = [ m_0, m_1, m_2, I(0), W(0), \bar{\beta} ]\end{split}
\end{equation*}
\begin{sphinxVerbatim}[commandchars=\\\{\}]
\PYG{n}{x\PYGZus{}name} \PYG{o}{=} \PYG{p}{[} \PYG{l+s+s1}{\PYGZsq{}}\PYG{l+s+s1}{m\PYGZus{}mobility}\PYG{l+s+s1}{\PYGZsq{}}\PYG{p}{,} \PYG{l+s+s1}{\PYGZsq{}}\PYG{l+s+s1}{m\PYGZus{}testing}\PYG{l+s+s1}{\PYGZsq{}}\PYG{p}{,} \PYG{l+s+s1}{\PYGZsq{}}\PYG{l+s+s1}{m\PYGZus{}stime}\PYG{l+s+s1}{\PYGZsq{}}\PYG{p}{,} \PYG{l+s+s1}{\PYGZsq{}}\PYG{l+s+s1}{I(0)}\PYG{l+s+s1}{\PYGZsq{}}\PYG{p}{,} \PYG{l+s+s1}{\PYGZsq{}}\PYG{l+s+s1}{W(0)}\PYG{l+s+s1}{\PYGZsq{}}\PYG{p}{,} \PYG{l+s+s1}{\PYGZsq{}}\PYG{l+s+s1}{beta\PYGZus{}bar}\PYG{l+s+s1}{\PYGZsq{}} \PYG{p}{]}
\end{sphinxVerbatim}

\index{maximum@\spxentry{maximum}}\index{likelihood@\spxentry{likelihood}}\ignorespaces 

\subparagraph{Maximum Likelihood}
\label{\detokenize{xsrst/numeric_covid_19_xam_py:maximum-likelihood}}\label{\detokenize{xsrst/numeric_covid_19_xam_py:numeric-covid-19-xam-py-unknown-parameters-maximum-likelihood}}\label{\detokenize{xsrst/numeric_covid_19_xam_py:index-8}}
We use a Gaussian likelihood for each of the differences in the
cumulative deaths. The unknown parameters are estimated by maximizing the
product of these likelihoods; i.e., the differences are modeled as being
independent. The covariance of the estimates is approximated
by the inverse of the observed information matrix.
AD is used to compute first and second derivatives of the likelihood
w.r.t. the unknown parameters \(x\).
These derivatives are used during optimization as well as for
computing the observed information matrix.

\index{model@\spxentry{model}}\index{bounds@\spxentry{bounds}}\ignorespaces 

\subparagraph{Model Bounds}
\label{\detokenize{xsrst/numeric_covid_19_xam_py:model-bounds}}\label{\detokenize{xsrst/numeric_covid_19_xam_py:numeric-covid-19-xam-py-unknown-parameters-model-bounds}}\label{\detokenize{xsrst/numeric_covid_19_xam_py:index-9}}
The infection rate \(\beta(t)\) must be non\sphinxhyphen{}negative; i.e.,
\begin{equation*}
\begin{split}0 \leq \bar{\beta} \exp[ m_0 c_0 (t) + m_1 c_1 (t) + m_2 c_2 (t) ]\end{split}
\end{equation*}
is true for all \(t\).
In addition, the size of the groups can not be negative.
It is sufficient to enforce this constraint on the initial conditions; i.e.,
\begin{equation*}
\begin{split}\begin{array}{lcr}
0   &  \leq & \bar{\beta }   \\
0   &  \leq & I(0)            \\
0   & \leq  & W(0)
\end{array}\end{split}
\end{equation*}
\index{actual@\spxentry{actual}}\index{bounds@\spxentry{bounds}}\ignorespaces 

\subparagraph{Actual Bounds}
\label{\detokenize{xsrst/numeric_covid_19_xam_py:actual-bounds}}\label{\detokenize{xsrst/numeric_covid_19_xam_py:numeric-covid-19-xam-py-unknown-parameters-actual-bounds}}\label{\detokenize{xsrst/numeric_covid_19_xam_py:index-10}}
The following actual upper and lower bounds for the unknown parameters
are used as an as an aid to the optimizer:

\begin{sphinxVerbatim}[commandchars=\\\{\}]
    \PYG{n}{n\PYGZus{}x} \PYG{o}{=} \PYG{n+nb}{len}\PYG{p}{(}\PYG{n}{x\PYGZus{}name}\PYG{p}{)}
    \PYG{n}{x\PYGZus{}lower}                \PYG{o}{=} \PYG{n}{numpy}\PYG{o}{.}\PYG{n}{zeros}\PYG{p}{(} \PYG{n}{n\PYGZus{}x}\PYG{p}{,} \PYG{n}{dtype}\PYG{o}{=}\PYG{n+nb}{float}\PYG{p}{)}
    \PYG{n}{x\PYGZus{}upper}                \PYG{o}{=} \PYG{n}{numpy}\PYG{o}{.}\PYG{n}{zeros}\PYG{p}{(} \PYG{n}{n\PYGZus{}x}\PYG{p}{,} \PYG{n}{dtype}\PYG{o}{=}\PYG{n+nb}{float}\PYG{p}{)}
    \PYG{n}{x\PYGZus{}upper}\PYG{p}{[} \PYG{n}{x\PYGZus{}sim} \PYG{o}{\PYGZgt{}} \PYG{l+m+mf}{0.0} \PYG{p}{]} \PYG{o}{=} \PYG{n}{actual\PYGZus{}bound\PYGZus{}factor} \PYG{o}{*} \PYG{n}{x\PYGZus{}sim}\PYG{p}{[} \PYG{n}{x\PYGZus{}sim} \PYG{o}{\PYGZgt{}} \PYG{l+m+mf}{0.0} \PYG{p}{]}
    \PYG{n}{x\PYGZus{}lower}\PYG{p}{[} \PYG{n}{x\PYGZus{}sim} \PYG{o}{\PYGZlt{}} \PYG{l+m+mf}{0.0} \PYG{p}{]} \PYG{o}{=} \PYG{n}{actual\PYGZus{}bound\PYGZus{}factor} \PYG{o}{*} \PYG{n}{x\PYGZus{}sim}\PYG{p}{[} \PYG{n}{x\PYGZus{}sim} \PYG{o}{\PYGZlt{}} \PYG{l+m+mf}{0.0} \PYG{p}{]}
\end{sphinxVerbatim}

where \sphinxstyleemphasis{x\_sim} is the
{\hyperref[\detokenize{xsrst/numeric_covid_19_xam_py:numeric-covid-19-xam-py-data-simulation}]{\sphinxcrossref{\DUrole{std,std-ref}{simulation}}}} value
for the unknown parameters and \sphinxstyleemphasis{actual\_bound\_factor} is chosen below.
The problem has not really been solved if bounds,
other than the model bounds above, are active at the solution of the
optimization problem.

\begin{sphinxVerbatim}[commandchars=\\\{\}]
\PYG{n}{actual\PYGZus{}bound\PYGZus{}factor} \PYG{o}{=} \PYG{l+m+mf}{5.0}
\end{sphinxVerbatim}

\index{data@\spxentry{data}}\ignorespaces 

\subparagraph{Data}
\label{\detokenize{xsrst/numeric_covid_19_xam_py:data}}\label{\detokenize{xsrst/numeric_covid_19_xam_py:numeric-covid-19-xam-py-data}}\label{\detokenize{xsrst/numeric_covid_19_xam_py:index-11}}
The data in this model is the cumulative number of deaths,
as a fraction of the total population and as a function of time.
We assume that new deaths are recorded for time intervals
and the cumulative deaths is the sum of these recordings.
For this reason, we model the difference of the cumulative deaths
between time points as independent.

\index{sample\_interval@\spxentry{sample\_interval}}\ignorespaces 

\subparagraph{sample\_interval}
\label{\detokenize{xsrst/numeric_covid_19_xam_py:sample-interval}}\label{\detokenize{xsrst/numeric_covid_19_xam_py:numeric-covid-19-xam-py-data-sample-interval}}\label{\detokenize{xsrst/numeric_covid_19_xam_py:index-12}}
It is possible to sub\sphinxhyphen{}sample the data in order to reduce noise.
The cumulative death data is just sub\sphinxhyphen{}sampled since the reduces noise by
the summing the differences corresponding to a longer time period.
The covariate data is averaged over the sample interval.
The \sphinxstyleemphasis{sample\_interval} must be either one or a positive even integer
(even so an original data point corresponds to the center of the interval).

\begin{sphinxVerbatim}[commandchars=\\\{\}]
\PYG{n}{sample\PYGZus{}interval} \PYG{o}{=} \PYG{l+m+mi}{1}
\end{sphinxVerbatim}

\index{data\_file@\spxentry{data\_file}}\ignorespaces 

\subparagraph{data\_file}
\label{\detokenize{xsrst/numeric_covid_19_xam_py:data-file}}\label{\detokenize{xsrst/numeric_covid_19_xam_py:numeric-covid-19-xam-py-data-data-file}}\label{\detokenize{xsrst/numeric_covid_19_xam_py:index-13}}
If the data file name is the empty string, the cumulative death data,
and corresponding covariates, are simulated by the program.
Otherwise, the data file must be a CSV file with the following columns:
\sphinxstyleemphasis{day} , \sphinxstyleemphasis{death} , \sphinxstyleemphasis{mobility} , \sphinxstyleemphasis{testing} .
In this case the data file is used for the
cumulative death and corresponding covariates.

\begin{sphinxVerbatim}[commandchars=\\\{\}]
\PYG{n}{data\PYGZus{}file} \PYG{o}{=} \PYG{l+s+s1}{\PYGZsq{}}\PYG{l+s+s1}{/home/bradbell/Downloads/561.csv}\PYG{l+s+s1}{\PYGZsq{}}         \PYG{c+c1}{\PYGZsh{} Pennsylvania}
\PYG{n}{data\PYGZus{}file} \PYG{o}{=} \PYG{l+s+s1}{\PYGZsq{}}\PYG{l+s+s1}{/home/bradbell/trash/covid\PYGZus{}19/seirwd.csv}\PYG{l+s+s1}{\PYGZsq{}} \PYG{c+c1}{\PYGZsh{} New York}
\PYG{n}{data\PYGZus{}file} \PYG{o}{=} \PYG{l+s+s1}{\PYGZsq{}}\PYG{l+s+s1}{\PYGZsq{}}                                         \PYG{c+c1}{\PYGZsh{} empty string}
\end{sphinxVerbatim}

\index{coefficient@\spxentry{coefficient}}\index{variation@\spxentry{variation}}\ignorespaces 

\subparagraph{Coefficient of Variation}
\label{\detokenize{xsrst/numeric_covid_19_xam_py:coefficient-of-variation}}\label{\detokenize{xsrst/numeric_covid_19_xam_py:numeric-covid-19-xam-py-data-coefficient-of-variation}}\label{\detokenize{xsrst/numeric_covid_19_xam_py:index-14}}
This is the coefficient of variation for the differences
in the cumulative death data as a fraction, not a percent.
If this value is zero, a CV of zero is used for data simulation
and a CV of one in the definition of the likelihood.
This enables checking that the unknown parameters can be accurately
identified using perfect data.
For real data (when \sphinxstyleemphasis{data\_file} is not empty)
this value should be adjusted so that the average residual has variance one.

\begin{sphinxVerbatim}[commandchars=\\\{\}]
\PYG{n}{death\PYGZus{}data\PYGZus{}cv} \PYG{o}{=} \PYG{l+m+mf}{0.25}
\end{sphinxVerbatim}

Note this is the noise level in the original data before it is
sub\sphinxhyphen{}sampled using
{\hyperref[\detokenize{xsrst/numeric_covid_19_xam_py:numeric-covid-19-xam-py-data-sample-interval}]{\sphinxcrossref{\DUrole{std,std-ref}{sample\_interval}}}}.

\index{simulation@\spxentry{simulation}}\ignorespaces 

\subparagraph{Simulation}
\label{\detokenize{xsrst/numeric_covid_19_xam_py:simulation}}\label{\detokenize{xsrst/numeric_covid_19_xam_py:numeric-covid-19-xam-py-data-simulation}}\label{\detokenize{xsrst/numeric_covid_19_xam_py:index-15}}
If \sphinxstyleemphasis{data\_file} is the empty string, the data is simulated using
the following values for the
{\hyperref[\detokenize{xsrst/numeric_covid_19_xam_py:numeric-covid-19-xam-py-unknown-parameters}]{\sphinxcrossref{\DUrole{std,std-ref}{unknown\_parameters}}}}:

\begin{sphinxVerbatim}[commandchars=\\\{\}]
\PYG{n}{m\PYGZus{}mobility\PYGZus{}sim}    \PYG{o}{=}   \PYG{l+m+mf}{1.0}  \PYG{c+c1}{\PYGZsh{} m\PYGZus{}0}
\PYG{n}{m\PYGZus{}testing\PYGZus{}sim}     \PYG{o}{=} \PYG{o}{\PYGZhy{}} \PYG{l+m+mf}{1.0}  \PYG{c+c1}{\PYGZsh{} m\PYGZus{}1}
\PYG{n}{m\PYGZus{}stime\PYGZus{}sim}       \PYG{o}{=} \PYG{o}{\PYGZhy{}} \PYG{l+m+mf}{1.0}  \PYG{c+c1}{\PYGZsh{} m\PYGZus{}2}
\PYG{n}{I0\PYGZus{}sim}            \PYG{o}{=}  \PYG{l+m+mf}{2e\PYGZhy{}5}  \PYG{c+c1}{\PYGZsh{} I(0)}
\PYG{n}{W0\PYGZus{}sim}            \PYG{o}{=}  \PYG{l+m+mf}{2e\PYGZhy{}5}  \PYG{c+c1}{\PYGZsh{} W(0)}
\PYG{n}{beta\PYGZus{}bar\PYGZus{}sim}      \PYG{o}{=}  \PYG{l+m+mf}{2.0}   \PYG{c+c1}{\PYGZsh{} baseline value for beta}
\end{sphinxVerbatim}

\index{weighted@\spxentry{weighted}}\index{residuals@\spxentry{residuals}}\ignorespaces 

\subparagraph{Weighted Residuals}
\label{\detokenize{xsrst/numeric_covid_19_xam_py:weighted-residuals}}\label{\detokenize{xsrst/numeric_covid_19_xam_py:numeric-covid-19-xam-py-data-weighted-residuals}}\label{\detokenize{xsrst/numeric_covid_19_xam_py:index-16}}
If \sphinxstyleemphasis{death\_data\_cv} is zero, \(\lambda = 1\), otherwise
\(\lambda\) is equal to

\begin{DUlineblock}{0em}
\item[]         \sphinxstyleemphasis{death\_data\_cv} \sphinxcode{\sphinxupquote{* sqrt}} ( \sphinxstyleemphasis{sample\_interval} )
\end{DUlineblock}

.
(Note that the standard deviation of a sum of independent values is the
square root of the sum of the variance of each of the values.)
Let \(y_i\) be the i\sphinxhyphen{}th value for the cumulative death data.
The weighted residuals (some times referred to as just the residuals) are
\begin{equation*}
\begin{split}r_i = \frac{ ( y_{i+1} - y_i ) - [ D( t_{i+1} ) - D( t_i ) ] }{
\lambda ( y_{i+1} - y_i ) }\end{split}
\end{equation*}
where \(D(t)\) is the model for the cumulative data
given the fit results.
The time corresponding to \(r_i\) is \(( t_{i+1} + t_i ) / 2\).
We put the data difference in the denominator,
instead of the model difference,
because it is constant with respect to the unknown parameters.

\index{random@\spxentry{random}}\index{seed@\spxentry{seed}}\ignorespaces 

\subparagraph{Random Seed}
\label{\detokenize{xsrst/numeric_covid_19_xam_py:random-seed}}\label{\detokenize{xsrst/numeric_covid_19_xam_py:numeric-covid-19-xam-py-random-seed}}\label{\detokenize{xsrst/numeric_covid_19_xam_py:index-17}}
This is the random seed used to simulate noise in the data.
If this value is zero, the system clock is used to choose the random seed.

\begin{sphinxVerbatim}[commandchars=\\\{\}]
\PYG{n}{random\PYGZus{}seed} \PYG{o}{=} \PYG{l+m+mi}{0}
\end{sphinxVerbatim}

\index{random@\spxentry{random}}\index{start@\spxentry{start}}\ignorespaces 

\subparagraph{Random Start}
\label{\detokenize{xsrst/numeric_covid_19_xam_py:random-start}}\label{\detokenize{xsrst/numeric_covid_19_xam_py:numeric-covid-19-xam-py-random-start}}\label{\detokenize{xsrst/numeric_covid_19_xam_py:index-18}}
The optimizer needs a good starting point in order to succeed.
This is the number of random points, between the lower and upper limits,
that are checked. The point with the best objective value is chosen
as the starting point for the optimization.

\begin{sphinxVerbatim}[commandchars=\\\{\}]
\PYG{n}{n\PYGZus{}random\PYGZus{}start} \PYG{o}{=} \PYG{l+m+mi}{4000}
\end{sphinxVerbatim}

\index{display@\spxentry{display}}\index{fit@\spxentry{fit}}\index{results@\spxentry{results}}\ignorespaces 

\subparagraph{Display Fit Results}
\label{\detokenize{xsrst/numeric_covid_19_xam_py:display-fit-results}}\label{\detokenize{xsrst/numeric_covid_19_xam_py:numeric-covid-19-xam-py-display-fit-results}}\label{\detokenize{xsrst/numeric_covid_19_xam_py:index-19}}
If you set this variable to True,
a printout and a plot of the fit results is generated.

\begin{sphinxVerbatim}[commandchars=\\\{\}]
\PYG{n}{display\PYGZus{}fit} \PYG{o}{=} \PYG{k+kc}{False}
\end{sphinxVerbatim}

\index{plot@\spxentry{plot}}\ignorespaces 

\subparagraph{Plot}
\label{\detokenize{xsrst/numeric_covid_19_xam_py:plot}}\label{\detokenize{xsrst/numeric_covid_19_xam_py:numeric-covid-19-xam-py-display-fit-results-plot}}\label{\detokenize{xsrst/numeric_covid_19_xam_py:index-20}}
There are three plots all with time on the x\sphinxhyphen{}axis.
The first contains the size for all the compartments, except S,
as a fraction of the total population.
The second contains the model and data for the death difference values.
The third contains the weighted residuals corresponding to the death
difference data.

\index{printout@\spxentry{printout}}\ignorespaces 

\subparagraph{Printout}
\label{\detokenize{xsrst/numeric_covid_19_xam_py:printout}}\label{\detokenize{xsrst/numeric_covid_19_xam_py:numeric-covid-19-xam-py-display-fit-results-printout}}\label{\detokenize{xsrst/numeric_covid_19_xam_py:index-21}}\begin{enumerate}
\sphinxsetlistlabels{\arabic}{enumi}{enumii}{}{.}%
\item {} 
The following statistics for the weighted data residuals is printed:
the maximum, minimum, average, and average of square.

\item {} 
A table with the following columns is printed:


\begin{savenotes}\sphinxattablestart
\centering
\begin{tabular}[t]{|\X{9}{52}|\X{43}{52}|}
\hline

\sphinxstyleemphasis{x\_name}
&
name of the unknown parameter
\\
\hline
\sphinxstyleemphasis{x\_fit}
&
fit result for the unknown parameter
\\
\hline
\sphinxstyleemphasis{x\_lower}
&
lower bound used for the fit
\\
\hline
\sphinxstyleemphasis{x\_upper}
&
upper bound used for the fit
\\
\hline
\sphinxstyleemphasis{std\_error}
&
asymptotic standard error for the parameter
\\
\hline
\end{tabular}
\par
\sphinxattableend\end{savenotes}

\item {} 
If \sphinxstyleemphasis{data\_file} is empty,
a table is printed with the following columns is also printed:


\begin{savenotes}\sphinxattablestart
\centering
\begin{tabular}[t]{|\X{9}{62}|\X{53}{62}|}
\hline

\sphinxstyleemphasis{x\_name}
&
name of the unknown parameter
\\
\hline
\sphinxstyleemphasis{x\_sim}
&
known parameter value used during simulation
\\
\hline
\sphinxstyleemphasis{x\_fit}
&
fit result for the unknown parameter
\\
\hline
\sphinxstyleemphasis{rel\_error}
&
relative error for fit versus simulation
\\
\hline
\sphinxstyleemphasis{residual}
&
\sphinxstyleemphasis{std\_error} weighted residual for fit versus simulation
\\
\hline
\end{tabular}
\par
\sphinxattableend\end{savenotes}

\end{enumerate}

\index{debug@\spxentry{debug}}\index{output@\spxentry{output}}\ignorespaces 

\subparagraph{Debug Output}
\label{\detokenize{xsrst/numeric_covid_19_xam_py:debug-output}}\label{\detokenize{xsrst/numeric_covid_19_xam_py:numeric-covid-19-xam-py-debug-output}}\label{\detokenize{xsrst/numeric_covid_19_xam_py:index-22}}
If this flag is true a lot of debugging output is printed.

\begin{sphinxVerbatim}[commandchars=\\\{\}]
\PYG{n}{debug\PYGZus{}output} \PYG{o}{=} \PYG{k+kc}{False}
\end{sphinxVerbatim}

\index{source@\spxentry{source}}\index{code@\spxentry{code}}\ignorespaces 

\subparagraph{Source Code}
\label{\detokenize{xsrst/numeric_covid_19_xam_py:source-code}}\label{\detokenize{xsrst/numeric_covid_19_xam_py:numeric-covid-19-xam-py-source-code}}\label{\detokenize{xsrst/numeric_covid_19_xam_py:index-23}}
\begin{sphinxVerbatim}[commandchars=\\\{\}]
\PYG{k+kn}{from} \PYG{n+nn}{pdb} \PYG{k+kn}{import} \PYG{n}{set\PYGZus{}trace}
\PYG{k+kn}{from} \PYG{n+nn}{matplotlib} \PYG{k+kn}{import} \PYG{n}{pyplot}
\PYG{k+kn}{from} \PYG{n+nn}{matplotlib} \PYG{k+kn}{import} \PYG{n}{gridspec}
\PYG{k+kn}{import} \PYG{n+nn}{scipy}\PYG{n+nn}{.}\PYG{n+nn}{optimize}
\PYG{k+kn}{import} \PYG{n+nn}{numpy}
\PYG{k+kn}{import} \PYG{n+nn}{random}
\PYG{k+kn}{import} \PYG{n+nn}{time}
\PYG{k+kn}{import} \PYG{n+nn}{csv}
\PYG{k+kn}{import} \PYG{n+nn}{copy}
\PYG{c+c1}{\PYGZsh{}}
\PYG{k+kn}{import} \PYG{n+nn}{cppad\PYGZus{}py}
\PYG{k+kn}{from} \PYG{n+nn}{optimize\PYGZus{}fun\PYGZus{}class} \PYG{k+kn}{import} \PYG{n}{optimize\PYGZus{}fun\PYGZus{}class}
\PYG{k+kn}{from} \PYG{n+nn}{seirwd\PYGZus{}model} \PYG{k+kn}{import} \PYG{n}{seirwd\PYGZus{}model}
\PYG{c+c1}{\PYGZsh{}}
\PYG{k}{def} \PYG{n+nf}{subsample}\PYG{p}{(}\PYG{n}{file\PYGZus{}data}\PYG{p}{)} \PYG{p}{:}
    \PYG{k}{if} \PYG{n}{sample\PYGZus{}interval} \PYG{o}{==} \PYG{l+m+mi}{1} \PYG{p}{:}
        \PYG{k}{return} \PYG{n}{file\PYGZus{}data}
    \PYG{k}{assert} \PYG{n}{sample\PYGZus{}interval} \PYG{o}{\PYGZpc{}} \PYG{l+m+mi}{2} \PYG{o}{==} \PYG{l+m+mi}{0}
    \PYG{n}{sample\PYGZus{}interval\PYGZus{}2} \PYG{o}{=} \PYG{n+nb}{int}\PYG{p}{(} \PYG{n}{sample\PYGZus{}interval} \PYG{o}{/} \PYG{l+m+mi}{2} \PYG{p}{)}
    \PYG{c+c1}{\PYGZsh{}}
    \PYG{n}{n\PYGZus{}data}          \PYG{o}{=} \PYG{n+nb}{len}\PYG{p}{(}\PYG{n}{file\PYGZus{}data}\PYG{p}{)}
    \PYG{n}{n\PYGZus{}data\PYGZus{}interval} \PYG{o}{=} \PYG{n}{n\PYGZus{}data} \PYG{o}{\PYGZhy{}} \PYG{l+m+mi}{1}
    \PYG{n}{n\PYGZus{}sub\PYGZus{}interval}  \PYG{o}{=} \PYG{n+nb}{int}\PYG{p}{(} \PYG{n}{n\PYGZus{}data\PYGZus{}interval} \PYG{o}{/} \PYG{n}{sample\PYGZus{}interval} \PYG{p}{)}
    \PYG{n}{n\PYGZus{}sub}           \PYG{o}{=} \PYG{n}{n\PYGZus{}sub\PYGZus{}interval}
    \PYG{n}{sub\PYGZus{}data}        \PYG{o}{=} \PYG{n+nb}{list}\PYG{p}{(}\PYG{p}{)}
    \PYG{k}{for} \PYG{n}{j} \PYG{o+ow}{in} \PYG{n+nb}{range}\PYG{p}{(}\PYG{n}{n\PYGZus{}sub}\PYG{p}{)} \PYG{p}{:}
        \PYG{n}{i}                \PYG{o}{=} \PYG{n}{sample\PYGZus{}interval} \PYG{o}{*} \PYG{n}{j} \PYG{o}{+} \PYG{n}{sample\PYGZus{}interval\PYGZus{}2}
        \PYG{n}{row}              \PYG{o}{=} \PYG{n+nb}{dict}\PYG{p}{(}\PYG{p}{)}
        \PYG{n}{row}\PYG{p}{[}\PYG{l+s+s1}{\PYGZsq{}}\PYG{l+s+s1}{day}\PYG{l+s+s1}{\PYGZsq{}}\PYG{p}{]}       \PYG{o}{=} \PYG{n}{file\PYGZus{}data}\PYG{p}{[}\PYG{n}{i}\PYG{p}{]}\PYG{p}{[}\PYG{l+s+s1}{\PYGZsq{}}\PYG{l+s+s1}{day}\PYG{l+s+s1}{\PYGZsq{}}\PYG{p}{]}
        \PYG{n}{row}\PYG{p}{[}\PYG{l+s+s1}{\PYGZsq{}}\PYG{l+s+s1}{death}\PYG{l+s+s1}{\PYGZsq{}}\PYG{p}{]}     \PYG{o}{=} \PYG{n}{file\PYGZus{}data}\PYG{p}{[}\PYG{n}{i}\PYG{p}{]}\PYG{p}{[}\PYG{l+s+s1}{\PYGZsq{}}\PYG{l+s+s1}{death}\PYG{l+s+s1}{\PYGZsq{}}\PYG{p}{]}
        \PYG{k}{for} \PYG{n}{key} \PYG{o+ow}{in} \PYG{p}{[}\PYG{l+s+s1}{\PYGZsq{}}\PYG{l+s+s1}{mobility}\PYG{l+s+s1}{\PYGZsq{}}\PYG{p}{,} \PYG{l+s+s1}{\PYGZsq{}}\PYG{l+s+s1}{testing}\PYG{l+s+s1}{\PYGZsq{}}\PYG{p}{]} \PYG{p}{:}
            \PYG{n}{row}\PYG{p}{[}\PYG{n}{key}\PYG{p}{]}  \PYG{o}{=} \PYG{l+m+mf}{0.0}
            \PYG{k}{for} \PYG{n}{k} \PYG{o+ow}{in} \PYG{n+nb}{range}\PYG{p}{(}\PYG{n}{sample\PYGZus{}interval} \PYG{o}{+} \PYG{l+m+mi}{1}\PYG{p}{)} \PYG{p}{:}
                \PYG{n}{i}         \PYG{o}{=} \PYG{n}{sample\PYGZus{}interval} \PYG{o}{*} \PYG{n}{j} \PYG{o}{+} \PYG{n}{k}
                \PYG{n}{row}\PYG{p}{[}\PYG{n}{key}\PYG{p}{]} \PYG{o}{+}\PYG{o}{=} \PYG{n}{file\PYGZus{}data}\PYG{p}{[}\PYG{n}{i}\PYG{p}{]}\PYG{p}{[}\PYG{n}{key}\PYG{p}{]}
            \PYG{n}{row}\PYG{p}{[}\PYG{n}{key}\PYG{p}{]} \PYG{o}{/}\PYG{o}{=} \PYG{n+nb}{float}\PYG{p}{(}\PYG{n}{sample\PYGZus{}interval} \PYG{o}{+} \PYG{l+m+mi}{1}\PYG{p}{)}
        \PYG{n}{sub\PYGZus{}data}\PYG{o}{.}\PYG{n}{append}\PYG{p}{(}\PYG{n}{row}\PYG{p}{)}
    \PYG{k}{return} \PYG{n}{sub\PYGZus{}data}
\PYG{c+c1}{\PYGZsh{}}
\PYG{c+c1}{\PYGZsh{} t\PYGZus{}all, covariates}
\PYG{k}{if} \PYG{n}{data\PYGZus{}file} \PYG{o}{!=} \PYG{l+s+s1}{\PYGZsq{}}\PYG{l+s+s1}{\PYGZsq{}} \PYG{p}{:}
    \PYG{c+c1}{\PYGZsh{} getting cumulative death and covariates from data\PYGZus{}file}
    \PYG{n}{file\PYGZus{}in}   \PYG{o}{=} \PYG{n+nb}{open}\PYG{p}{(}\PYG{n}{data\PYGZus{}file}\PYG{p}{,} \PYG{l+s+s1}{\PYGZsq{}}\PYG{l+s+s1}{r}\PYG{l+s+s1}{\PYGZsq{}}\PYG{p}{)}
    \PYG{n}{file\PYGZus{}data} \PYG{o}{=} \PYG{n+nb}{list}\PYG{p}{(}\PYG{p}{)}
    \PYG{n}{reader}    \PYG{o}{=} \PYG{n}{csv}\PYG{o}{.}\PYG{n}{DictReader}\PYG{p}{(}\PYG{n}{file\PYGZus{}in}\PYG{p}{)}
    \PYG{k}{for} \PYG{n}{row} \PYG{o+ow}{in} \PYG{n}{reader} \PYG{p}{:}
        \PYG{k}{for} \PYG{n}{key} \PYG{o+ow}{in} \PYG{p}{[}\PYG{l+s+s1}{\PYGZsq{}}\PYG{l+s+s1}{day}\PYG{l+s+s1}{\PYGZsq{}}\PYG{p}{,} \PYG{l+s+s1}{\PYGZsq{}}\PYG{l+s+s1}{death}\PYG{l+s+s1}{\PYGZsq{}}\PYG{p}{,} \PYG{l+s+s1}{\PYGZsq{}}\PYG{l+s+s1}{mobility}\PYG{l+s+s1}{\PYGZsq{}}\PYG{p}{,} \PYG{l+s+s1}{\PYGZsq{}}\PYG{l+s+s1}{testing}\PYG{l+s+s1}{\PYGZsq{}}\PYG{p}{]}  \PYG{p}{:}
            \PYG{n}{row}\PYG{p}{[}\PYG{n}{key}\PYG{p}{]} \PYG{o}{=} \PYG{n+nb}{float}\PYG{p}{(} \PYG{n}{row}\PYG{p}{[}\PYG{n}{key}\PYG{p}{]} \PYG{p}{)}
        \PYG{n}{file\PYGZus{}data}\PYG{o}{.}\PYG{n}{append}\PYG{p}{(}\PYG{n}{row}\PYG{p}{)}
    \PYG{n}{file\PYGZus{}in}\PYG{o}{.}\PYG{n}{close}\PYG{p}{(}\PYG{p}{)}
    \PYG{n}{file\PYGZus{}data} \PYG{o}{=} \PYG{n}{subsample}\PYG{p}{(}\PYG{n}{file\PYGZus{}data}\PYG{p}{)}
    \PYG{c+c1}{\PYGZsh{}}
    \PYG{n}{n\PYGZus{}time}       \PYG{o}{=} \PYG{n+nb}{len}\PYG{p}{(}\PYG{n}{file\PYGZus{}data}\PYG{p}{)}
    \PYG{n}{t\PYGZus{}all}        \PYG{o}{=} \PYG{n}{numpy}\PYG{o}{.}\PYG{n}{empty}\PYG{p}{(}\PYG{n}{n\PYGZus{}time}\PYG{p}{,} \PYG{n}{dtype}\PYG{o}{=}\PYG{n+nb}{float}\PYG{p}{)}
    \PYG{n}{covariates}   \PYG{o}{=} \PYG{n}{numpy}\PYG{o}{.}\PYG{n}{empty}\PYG{p}{(} \PYG{p}{(}\PYG{n}{n\PYGZus{}time}\PYG{p}{,} \PYG{l+m+mi}{3}\PYG{p}{)} \PYG{p}{)}
    \PYG{k}{for} \PYG{n}{i} \PYG{o+ow}{in} \PYG{n+nb}{range}\PYG{p}{(}\PYG{n}{n\PYGZus{}time}\PYG{p}{)} \PYG{p}{:}
        \PYG{n}{t\PYGZus{}all}\PYG{p}{[}\PYG{n}{i}\PYG{p}{]}        \PYG{o}{=} \PYG{n}{file\PYGZus{}data}\PYG{p}{[}\PYG{n}{i}\PYG{p}{]}\PYG{p}{[}\PYG{l+s+s1}{\PYGZsq{}}\PYG{l+s+s1}{day}\PYG{l+s+s1}{\PYGZsq{}}\PYG{p}{]}
        \PYG{n}{covariates}\PYG{p}{[}\PYG{n}{i}\PYG{p}{,}\PYG{l+m+mi}{0}\PYG{p}{]} \PYG{o}{=} \PYG{n}{file\PYGZus{}data}\PYG{p}{[}\PYG{n}{i}\PYG{p}{]}\PYG{p}{[}\PYG{l+s+s1}{\PYGZsq{}}\PYG{l+s+s1}{mobility}\PYG{l+s+s1}{\PYGZsq{}}\PYG{p}{]}
        \PYG{n}{covariates}\PYG{p}{[}\PYG{n}{i}\PYG{p}{,}\PYG{l+m+mi}{1}\PYG{p}{]} \PYG{o}{=} \PYG{n}{file\PYGZus{}data}\PYG{p}{[}\PYG{n}{i}\PYG{p}{]}\PYG{p}{[}\PYG{l+s+s1}{\PYGZsq{}}\PYG{l+s+s1}{testing}\PYG{l+s+s1}{\PYGZsq{}}\PYG{p}{]}
    \PYG{k}{for} \PYG{n}{i} \PYG{o+ow}{in} \PYG{n+nb}{range}\PYG{p}{(}\PYG{n}{n\PYGZus{}time}\PYG{p}{)} \PYG{p}{:}
        \PYG{n}{stime}           \PYG{o}{=} \PYG{p}{(}\PYG{n}{t\PYGZus{}all}\PYG{p}{[}\PYG{n}{i}\PYG{p}{]} \PYG{o}{\PYGZhy{}} \PYG{n}{t\PYGZus{}all}\PYG{p}{[}\PYG{l+m+mi}{0}\PYG{p}{]}\PYG{p}{)} \PYG{o}{/} \PYG{p}{(}\PYG{n}{t\PYGZus{}all}\PYG{p}{[}\PYG{o}{\PYGZhy{}}\PYG{l+m+mi}{1}\PYG{p}{]} \PYG{o}{\PYGZhy{}} \PYG{n}{t\PYGZus{}all}\PYG{p}{[}\PYG{l+m+mi}{0}\PYG{p}{]}\PYG{p}{)}
        \PYG{n}{covariates}\PYG{p}{[}\PYG{n}{i}\PYG{p}{,}\PYG{l+m+mi}{2}\PYG{p}{]} \PYG{o}{=} \PYG{n}{stime}
\PYG{k}{else} \PYG{p}{:}
    \PYG{c+c1}{\PYGZsh{} simulating cumulative death and covariates}
    \PYG{n}{t\PYGZus{}start} \PYG{o}{=} \PYG{l+m+mf}{0.0}
    \PYG{n}{t\PYGZus{}stop}  \PYG{o}{=} \PYG{l+m+mf}{90.0}
    \PYG{n}{t\PYGZus{}step}  \PYG{o}{=} \PYG{n+nb}{float}\PYG{p}{(}\PYG{n}{sample\PYGZus{}interval}\PYG{p}{)}
    \PYG{n}{t\PYGZus{}all} \PYG{o}{=} \PYG{n}{numpy}\PYG{o}{.}\PYG{n}{arange}\PYG{p}{(}\PYG{n}{t\PYGZus{}start}\PYG{p}{,} \PYG{n}{t\PYGZus{}stop}\PYG{p}{,} \PYG{n}{t\PYGZus{}step}\PYG{p}{)}
    \PYG{c+c1}{\PYGZsh{}}
    \PYG{n}{n\PYGZus{}time} \PYG{o}{=} \PYG{n}{t\PYGZus{}all}\PYG{o}{.}\PYG{n}{size}
    \PYG{n}{covariates}   \PYG{o}{=} \PYG{n}{numpy}\PYG{o}{.}\PYG{n}{empty}\PYG{p}{(} \PYG{p}{(}\PYG{n}{n\PYGZus{}time}\PYG{p}{,} \PYG{l+m+mi}{3}\PYG{p}{)} \PYG{p}{)}
    \PYG{k}{for} \PYG{n}{i} \PYG{o+ow}{in} \PYG{n+nb}{range}\PYG{p}{(}\PYG{n}{n\PYGZus{}time}\PYG{p}{)} \PYG{p}{:}
        \PYG{n}{t}               \PYG{o}{=} \PYG{n}{t\PYGZus{}all}\PYG{p}{[}\PYG{n}{i}\PYG{p}{]}
        \PYG{n}{mobility}        \PYG{o}{=} \PYG{l+m+mf}{0.0} \PYG{k}{if} \PYG{n}{t} \PYG{o}{\PYGZlt{}} \PYG{l+m+mf}{10.0} \PYG{k}{else} \PYG{o}{\PYGZhy{}}\PYG{l+m+mf}{1.0}
        \PYG{n}{testing}         \PYG{o}{=} \PYG{l+m+mf}{0.0} \PYG{k}{if} \PYG{n}{t} \PYG{o}{\PYGZlt{}} \PYG{l+m+mf}{20.0} \PYG{k}{else} \PYG{l+m+mf}{1.0}
        \PYG{n}{stime}           \PYG{o}{=} \PYG{n}{t} \PYG{o}{/} \PYG{n}{t\PYGZus{}stop}
        \PYG{n}{covariates}\PYG{p}{[}\PYG{n}{i}\PYG{p}{,}\PYG{p}{:}\PYG{p}{]} \PYG{o}{=} \PYG{p}{[} \PYG{n}{mobility}\PYG{p}{,} \PYG{n}{testing}\PYG{p}{,} \PYG{n}{stime} \PYG{p}{]}
\PYG{c+c1}{\PYGZsh{}}
\PYG{c+c1}{\PYGZsh{} mobility in [\PYGZhy{}1, 0]}
\PYG{k}{assert} \PYG{n}{numpy}\PYG{o}{.}\PYG{n}{all}\PYG{p}{(} \PYG{n}{covariates}\PYG{p}{[}\PYG{p}{:}\PYG{p}{,}\PYG{l+m+mi}{0}\PYG{p}{]} \PYG{o}{\PYGZlt{}}\PYG{o}{=} \PYG{l+m+mf}{0.0} \PYG{p}{)}
\PYG{k}{assert} \PYG{n}{numpy}\PYG{o}{.}\PYG{n}{all}\PYG{p}{(} \PYG{o}{\PYGZhy{}}\PYG{l+m+mf}{1.0} \PYG{o}{\PYGZlt{}}\PYG{o}{=} \PYG{n}{covariates}\PYG{p}{[}\PYG{p}{:}\PYG{p}{,}\PYG{l+m+mi}{0}\PYG{p}{]} \PYG{p}{)}
\PYG{c+c1}{\PYGZsh{}}
\PYG{c+c1}{\PYGZsh{} testing in [0, 1]}
\PYG{k}{assert} \PYG{n}{numpy}\PYG{o}{.}\PYG{n}{all}\PYG{p}{(} \PYG{l+m+mf}{0.0} \PYG{o}{\PYGZlt{}}\PYG{o}{=} \PYG{n}{covariates}\PYG{p}{[}\PYG{p}{:}\PYG{p}{,}\PYG{l+m+mi}{1}\PYG{p}{]} \PYG{p}{)}
\PYG{k}{assert} \PYG{n}{numpy}\PYG{o}{.}\PYG{n}{all}\PYG{p}{(} \PYG{n}{covariates}\PYG{p}{[}\PYG{p}{:}\PYG{p}{,}\PYG{l+m+mi}{1}\PYG{p}{]} \PYG{o}{\PYGZlt{}}\PYG{o}{=} \PYG{l+m+mf}{1.0} \PYG{p}{)}
\PYG{c+c1}{\PYGZsh{}}
\PYG{c+c1}{\PYGZsh{} stime in [0, 1]}
\PYG{k}{assert} \PYG{n}{numpy}\PYG{o}{.}\PYG{n}{all}\PYG{p}{(} \PYG{l+m+mf}{0.0} \PYG{o}{\PYGZlt{}}\PYG{o}{=} \PYG{n}{covariates}\PYG{p}{[}\PYG{p}{:}\PYG{p}{,}\PYG{l+m+mi}{2}\PYG{p}{]} \PYG{p}{)}
\PYG{k}{assert} \PYG{n}{numpy}\PYG{o}{.}\PYG{n}{all}\PYG{p}{(} \PYG{n}{covariates}\PYG{p}{[}\PYG{p}{:}\PYG{p}{,}\PYG{l+m+mi}{2}\PYG{p}{]} \PYG{o}{\PYGZlt{}}\PYG{o}{=} \PYG{l+m+mf}{1.0} \PYG{p}{)}
\PYG{c+c1}{\PYGZsh{}}
\PYG{c+c1}{\PYGZsh{} other initial conditions used to simulate data}
\PYG{n}{E0\PYGZus{}sim}     \PYG{o}{=} \PYG{n}{I0\PYGZus{}sim} \PYG{o}{*} \PYG{n}{gamma\PYGZus{}known} \PYG{o}{/} \PYG{n}{sigma\PYGZus{}known}
\PYG{n}{S0\PYGZus{}sim}     \PYG{o}{=} \PYG{l+m+mf}{1.0} \PYG{o}{\PYGZhy{}} \PYG{n}{E0\PYGZus{}sim} \PYG{o}{\PYGZhy{}} \PYG{n}{I0\PYGZus{}sim} \PYG{o}{\PYGZhy{}} \PYG{n}{W0\PYGZus{}sim}
\PYG{n}{R0\PYGZus{}sim}     \PYG{o}{=} \PYG{l+m+mf}{0.0}
\PYG{n}{D0\PYGZus{}sim}     \PYG{o}{=} \PYG{l+m+mf}{0.0}
\PYG{c+c1}{\PYGZsh{}}
\PYG{c+c1}{\PYGZsh{} x\PYGZus{}sim}
\PYG{n}{x\PYGZus{}sim} \PYG{o}{=} \PYG{n}{numpy}\PYG{o}{.}\PYG{n}{array}\PYG{p}{(} \PYG{p}{[}
    \PYG{n}{m\PYGZus{}mobility\PYGZus{}sim}\PYG{p}{,} \PYG{n}{m\PYGZus{}testing\PYGZus{}sim}\PYG{p}{,} \PYG{n}{m\PYGZus{}stime\PYGZus{}sim}\PYG{p}{,} \PYG{n}{I0\PYGZus{}sim}\PYG{p}{,} \PYG{n}{W0\PYGZus{}sim}\PYG{p}{,} \PYG{n}{beta\PYGZus{}bar\PYGZus{}sim}
\PYG{p}{]} \PYG{p}{)}
\PYG{c+c1}{\PYGZsh{}}
\PYG{c+c1}{\PYGZsh{}}
\PYG{c+c1}{\PYGZsh{} actual\PYGZus{}seed}
\PYG{c+c1}{\PYGZsh{} Make a list so that we can set its element value inside a routine}
\PYG{n}{actual\PYGZus{}seed} \PYG{o}{=} \PYG{p}{[} \PYG{n}{random\PYGZus{}seed} \PYG{p}{]}
\PYG{c+c1}{\PYGZsh{}}
\PYG{k}{def} \PYG{n+nf}{x2seirwd\PYGZus{}all}\PYG{p}{(}\PYG{n}{x}\PYG{p}{)} \PYG{p}{:}
    \PYG{c+c1}{\PYGZsh{}}
    \PYG{c+c1}{\PYGZsh{} unpack x}
    \PYG{n}{cov\PYGZus{}mul}            \PYG{o}{=} \PYG{n}{x}\PYG{p}{[}\PYG{l+m+mi}{0} \PYG{p}{:} \PYG{l+m+mi}{3}\PYG{p}{]}
    \PYG{p}{[}\PYG{n}{I0}\PYG{p}{,} \PYG{n}{W0}\PYG{p}{,} \PYG{n}{beta\PYGZus{}bar}\PYG{p}{]} \PYG{o}{=} \PYG{n}{x}\PYG{p}{[}\PYG{l+m+mi}{3} \PYG{p}{:} \PYG{l+m+mi}{6}\PYG{p}{]}
    \PYG{c+c1}{\PYGZsh{}}
    \PYG{c+c1}{\PYGZsh{} beta\PYGZus{}all}
    \PYG{n}{beta\PYGZus{}all} \PYG{o}{=} \PYG{n}{beta\PYGZus{}bar} \PYG{o}{*} \PYG{n}{numpy}\PYG{o}{.}\PYG{n}{exp}\PYG{p}{(} \PYG{n}{numpy}\PYG{o}{.}\PYG{n}{matmul}\PYG{p}{(}\PYG{n}{covariates}\PYG{p}{,} \PYG{n}{cov\PYGZus{}mul}\PYG{p}{)} \PYG{p}{)}
    \PYG{c+c1}{\PYGZsh{}}
    \PYG{c+c1}{\PYGZsh{} p\PYGZus{}all}
    \PYG{n}{p\PYGZus{}all} \PYG{o}{=} \PYG{n+nb}{list}\PYG{p}{(}\PYG{p}{)}
    \PYG{k}{for} \PYG{n}{i} \PYG{o+ow}{in} \PYG{n+nb}{range}\PYG{p}{(} \PYG{n}{beta\PYGZus{}all}\PYG{o}{.}\PYG{n}{size} \PYG{p}{)} \PYG{p}{:}
        \PYG{n}{p} \PYG{o}{=} \PYG{p}{\PYGZob{}}
            \PYG{l+s+s1}{\PYGZsq{}}\PYG{l+s+s1}{alpha}\PYG{l+s+s1}{\PYGZsq{}} \PYG{p}{:} \PYG{n}{alpha\PYGZus{}known}\PYG{p}{,}
            \PYG{l+s+s1}{\PYGZsq{}}\PYG{l+s+s1}{sigma}\PYG{l+s+s1}{\PYGZsq{}} \PYG{p}{:} \PYG{n}{sigma\PYGZus{}known}\PYG{p}{,}
            \PYG{l+s+s1}{\PYGZsq{}}\PYG{l+s+s1}{gamma}\PYG{l+s+s1}{\PYGZsq{}} \PYG{p}{:} \PYG{n}{gamma\PYGZus{}known}\PYG{p}{,}
            \PYG{l+s+s1}{\PYGZsq{}}\PYG{l+s+s1}{chi}\PYG{l+s+s1}{\PYGZsq{}}   \PYG{p}{:} \PYG{n}{chi\PYGZus{}known}\PYG{p}{,}
            \PYG{l+s+s1}{\PYGZsq{}}\PYG{l+s+s1}{xi}\PYG{l+s+s1}{\PYGZsq{}}    \PYG{p}{:} \PYG{n}{xi\PYGZus{}known}\PYG{p}{,}
            \PYG{l+s+s1}{\PYGZsq{}}\PYG{l+s+s1}{delta}\PYG{l+s+s1}{\PYGZsq{}} \PYG{p}{:} \PYG{n}{delta\PYGZus{}known}\PYG{p}{,}
        \PYG{p}{\PYGZcb{}}
        \PYG{n}{p}\PYG{p}{[}\PYG{l+s+s1}{\PYGZsq{}}\PYG{l+s+s1}{beta}\PYG{l+s+s1}{\PYGZsq{}}\PYG{p}{]} \PYG{o}{=} \PYG{n}{beta\PYGZus{}all}\PYG{p}{[}\PYG{n}{i}\PYG{p}{]}
        \PYG{n}{p\PYGZus{}all}\PYG{o}{.}\PYG{n}{append}\PYG{p}{(} \PYG{n}{p} \PYG{p}{)}
    \PYG{c+c1}{\PYGZsh{}}
    \PYG{c+c1}{\PYGZsh{} initial}
    \PYG{n}{E0} \PYG{o}{=} \PYG{n}{I0} \PYG{o}{*} \PYG{n}{gamma\PYGZus{}known} \PYG{o}{/} \PYG{n}{sigma\PYGZus{}known}
    \PYG{n}{S0} \PYG{o}{=} \PYG{l+m+mf}{1.0} \PYG{o}{\PYGZhy{}} \PYG{n}{E0} \PYG{o}{\PYGZhy{}} \PYG{n}{I0} \PYG{o}{\PYGZhy{}} \PYG{n}{W0}
    \PYG{n}{R0} \PYG{o}{=} \PYG{l+m+mf}{0.0}
    \PYG{n}{D0} \PYG{o}{=} \PYG{l+m+mf}{0.0}
    \PYG{n}{initial}  \PYG{o}{=} \PYG{n}{numpy}\PYG{o}{.}\PYG{n}{array}\PYG{p}{(} \PYG{p}{[}\PYG{n}{S0}\PYG{p}{,} \PYG{n}{E0}\PYG{p}{,} \PYG{n}{I0}\PYG{p}{,} \PYG{n}{R0}\PYG{p}{,} \PYG{n}{W0}\PYG{p}{,} \PYG{n}{D0}\PYG{p}{]} \PYG{p}{)}
    \PYG{c+c1}{\PYGZsh{}}
    \PYG{c+c1}{\PYGZsh{} seirwd\PYGZus{}all}
    \PYG{n}{n\PYGZus{}step} \PYG{o}{=} \PYG{n}{sample\PYGZus{}interval} \PYG{o}{*} \PYG{n}{ode\PYGZus{}n\PYGZus{}step}
    \PYG{n}{seirwd\PYGZus{}all} \PYG{o}{=} \PYG{n}{seirwd\PYGZus{}model}\PYG{p}{(}\PYG{n}{ode\PYGZus{}method}\PYG{p}{,} \PYG{n}{t\PYGZus{}all}\PYG{p}{,} \PYG{n}{p\PYGZus{}all}\PYG{p}{,} \PYG{n}{initial}\PYG{p}{,} \PYG{n}{n\PYGZus{}step}\PYG{p}{)}
    \PYG{c+c1}{\PYGZsh{}}
    \PYG{k}{return} \PYG{n}{seirwd\PYGZus{}all}
\PYG{c+c1}{\PYGZsh{}}
\PYG{k}{def} \PYG{n+nf}{weighted\PYGZus{}data\PYGZus{}residual}\PYG{p}{(}\PYG{n}{D\PYGZus{}data}\PYG{p}{,} \PYG{n}{D\PYGZus{}model}\PYG{p}{)} \PYG{p}{:}
    \PYG{n}{Ddiff\PYGZus{}data}   \PYG{o}{=} \PYG{n}{numpy}\PYG{o}{.}\PYG{n}{diff}\PYG{p}{(}\PYG{n}{D\PYGZus{}data}\PYG{p}{)}
    \PYG{k}{if} \PYG{n}{numpy}\PYG{o}{.}\PYG{n}{any}\PYG{p}{(} \PYG{n}{Ddiff\PYGZus{}data} \PYG{o}{\PYGZlt{}}\PYG{o}{=} \PYG{l+m+mi}{0} \PYG{p}{)}  \PYG{p}{:}
        \PYG{n+nb}{print}\PYG{p}{(}\PYG{l+s+s1}{\PYGZsq{}}\PYG{l+s+s1}{covid\PYGZus{}19\PYGZus{}xam: cumutative death data not increasing at times: }\PYG{l+s+s1}{\PYGZsq{}}\PYG{p}{)}
        \PYG{n}{t\PYGZus{}msg} \PYG{o}{=} \PYG{n}{t\PYGZus{}all}\PYG{p}{[}\PYG{l+m+mi}{0} \PYG{p}{:} \PYG{o}{\PYGZhy{}}\PYG{l+m+mi}{1}\PYG{p}{]}
        \PYG{n+nb}{print}\PYG{p}{(} \PYG{n}{t\PYGZus{}msg}\PYG{p}{[} \PYG{n}{Ddiff\PYGZus{}data} \PYG{o}{\PYGZlt{}}\PYG{o}{=} \PYG{l+m+mi}{0} \PYG{p}{]} \PYG{p}{)}
    \PYG{n}{Ddiff\PYGZus{}model}  \PYG{o}{=} \PYG{n}{numpy}\PYG{o}{.}\PYG{n}{diff}\PYG{p}{(}\PYG{n}{D\PYGZus{}model}\PYG{p}{)}
    \PYG{n}{residual}     \PYG{o}{=} \PYG{p}{(}\PYG{n}{Ddiff\PYGZus{}data} \PYG{o}{\PYGZhy{}} \PYG{n}{Ddiff\PYGZus{}model}\PYG{p}{)} \PYG{o}{/} \PYG{n}{Ddiff\PYGZus{}data}
    \PYG{k}{if} \PYG{n}{death\PYGZus{}data\PYGZus{}cv} \PYG{o}{\PYGZgt{}} \PYG{l+m+mf}{0.0} \PYG{p}{:}
        \PYG{n}{residual} \PYG{o}{=} \PYG{n}{residual} \PYG{o}{/} \PYG{p}{(}\PYG{n}{numpy}\PYG{o}{.}\PYG{n}{sqrt}\PYG{p}{(}\PYG{n}{sample\PYGZus{}interval}\PYG{p}{)} \PYG{o}{*} \PYG{n}{death\PYGZus{}data\PYGZus{}cv}\PYG{p}{)}
    \PYG{k}{return} \PYG{n}{residual}
\PYG{c+c1}{\PYGZsh{}}
\PYG{c+c1}{\PYGZsh{} objective}
\PYG{c+c1}{\PYGZsh{} t\PYGZus{}all and D\PYGZus{}data, are constants relative to the objective function}
\PYG{k}{def} \PYG{n+nf}{objective}\PYG{p}{(}\PYG{n}{t\PYGZus{}all}\PYG{p}{,} \PYG{n}{D\PYGZus{}data}\PYG{p}{,} \PYG{n}{x}\PYG{p}{)} \PYG{p}{:}
    \PYG{c+c1}{\PYGZsh{}}
    \PYG{c+c1}{\PYGZsh{} compute model for data}
    \PYG{n}{seirwd\PYGZus{}all} \PYG{o}{=} \PYG{n}{x2seirwd\PYGZus{}all}\PYG{p}{(}\PYG{n}{x}\PYG{p}{)}
    \PYG{c+c1}{\PYGZsh{}}
    \PYG{c+c1}{\PYGZsh{} Model for the cumulative death as function of time}
    \PYG{n}{D\PYGZus{}model}   \PYG{o}{=} \PYG{n}{seirwd\PYGZus{}all}\PYG{p}{[}\PYG{p}{:}\PYG{p}{,}\PYG{l+m+mi}{5}\PYG{p}{]} \PYG{c+c1}{\PYGZsh{} column order is S, E, I, R, W, D}
    \PYG{c+c1}{\PYGZsh{}}
    \PYG{c+c1}{\PYGZsh{} compute negative log Gaussian likelihood dropping variance terms}
    \PYG{c+c1}{\PYGZsh{} because they are constaint w.r.t the unknown parameters}
    \PYG{n}{residual} \PYG{o}{=} \PYG{n}{weighted\PYGZus{}data\PYGZus{}residual}\PYG{p}{(}\PYG{n}{D\PYGZus{}data}\PYG{p}{,} \PYG{n}{D\PYGZus{}model}\PYG{p}{)}
    \PYG{n}{loss}     \PYG{o}{=} \PYG{l+m+mf}{0.5} \PYG{o}{*} \PYG{n}{numpy}\PYG{o}{.}\PYG{n}{sum}\PYG{p}{(} \PYG{n}{residual} \PYG{o}{*} \PYG{n}{residual} \PYG{p}{)}
    \PYG{k}{return} \PYG{n}{loss}
\PYG{c+c1}{\PYGZsh{}}
\PYG{c+c1}{\PYGZsh{} objective\PYGZus{}d\PYGZus{}fun}
\PYG{k}{def} \PYG{n+nf}{objective\PYGZus{}d\PYGZus{}fun}\PYG{p}{(}\PYG{n}{t\PYGZus{}all}\PYG{p}{,} \PYG{n}{D\PYGZus{}data}\PYG{p}{)} \PYG{p}{:}
    \PYG{c+c1}{\PYGZsh{}}
    \PYG{n}{n\PYGZus{}x}   \PYG{o}{=} \PYG{n+nb}{len}\PYG{p}{(}\PYG{n}{x\PYGZus{}name}\PYG{p}{)}
    \PYG{n}{x}     \PYG{o}{=} \PYG{l+m+mf}{0.1} \PYG{o}{*} \PYG{n}{numpy}\PYG{o}{.}\PYG{n}{ones}\PYG{p}{(}\PYG{n}{n\PYGZus{}x}\PYG{p}{)}
    \PYG{n}{ax}    \PYG{o}{=} \PYG{n}{cppad\PYGZus{}py}\PYG{o}{.}\PYG{n}{independent}\PYG{p}{(}\PYG{n}{x}\PYG{p}{)}
    \PYG{n}{aloss} \PYG{o}{=} \PYG{n}{objective}\PYG{p}{(}\PYG{n}{t\PYGZus{}all}\PYG{p}{,} \PYG{n}{D\PYGZus{}data}\PYG{p}{,} \PYG{n}{ax}\PYG{p}{)}
    \PYG{n}{aloss} \PYG{o}{=} \PYG{n}{numpy}\PYG{o}{.}\PYG{n}{array}\PYG{p}{(} \PYG{p}{[} \PYG{n}{aloss} \PYG{p}{]} \PYG{p}{)}
    \PYG{c+c1}{\PYGZsh{}}
    \PYG{n}{objective\PYGZus{}ad} \PYG{o}{=} \PYG{n}{cppad\PYGZus{}py}\PYG{o}{.}\PYG{n}{d\PYGZus{}fun}\PYG{p}{(}\PYG{n}{ax}\PYG{p}{,} \PYG{n}{aloss}\PYG{p}{)}
    \PYG{k}{return} \PYG{n}{objective\PYGZus{}ad}
\PYG{c+c1}{\PYGZsh{}}
\PYG{k}{def} \PYG{n+nf}{simulate\PYGZus{}data}\PYG{p}{(}\PYG{p}{)} \PYG{p}{:}
    \PYG{c+c1}{\PYGZsh{}}
    \PYG{c+c1}{\PYGZsh{} noiseless simulated data}
    \PYG{n}{seirwd\PYGZus{}all\PYGZus{}sim} \PYG{o}{=} \PYG{n}{x2seirwd\PYGZus{}all}\PYG{p}{(}\PYG{n}{x\PYGZus{}sim}\PYG{p}{)}
    \PYG{c+c1}{\PYGZsh{}}
    \PYG{c+c1}{\PYGZsh{} rng: numpy random number generator}
    \PYG{n}{rng} \PYG{o}{=} \PYG{n}{numpy}\PYG{o}{.}\PYG{n}{random}\PYG{o}{.}\PYG{n}{default\PYGZus{}rng}\PYG{p}{(}\PYG{n}{actual\PYGZus{}seed}\PYG{p}{[}\PYG{l+m+mi}{0}\PYG{p}{]}\PYG{p}{)}
    \PYG{c+c1}{\PYGZsh{}}
    \PYG{c+c1}{\PYGZsh{} D\PYGZus{}data}
    \PYG{n}{D\PYGZus{}sim}       \PYG{o}{=} \PYG{n}{seirwd\PYGZus{}all\PYGZus{}sim}\PYG{p}{[}\PYG{p}{:}\PYG{p}{,}\PYG{l+m+mi}{5}\PYG{p}{]}
    \PYG{n}{Ddiff\PYGZus{}sim}   \PYG{o}{=} \PYG{n}{numpy}\PYG{o}{.}\PYG{n}{diff}\PYG{p}{(} \PYG{n}{D\PYGZus{}sim} \PYG{p}{)}
    \PYG{n}{std}         \PYG{o}{=} \PYG{n}{numpy}\PYG{o}{.}\PYG{n}{sqrt}\PYG{p}{(}\PYG{n}{sample\PYGZus{}interval}\PYG{p}{)} \PYG{o}{*} \PYG{n}{death\PYGZus{}data\PYGZus{}cv} \PYG{o}{*} \PYG{n}{Ddiff\PYGZus{}sim}
    \PYG{n}{noise}       \PYG{o}{=} \PYG{n}{std} \PYG{o}{*} \PYG{n}{rng}\PYG{o}{.}\PYG{n}{standard\PYGZus{}normal}\PYG{p}{(}\PYG{n}{size} \PYG{o}{=} \PYG{n}{t\PYGZus{}all}\PYG{o}{.}\PYG{n}{size} \PYG{o}{\PYGZhy{}} \PYG{l+m+mi}{1}\PYG{p}{)}
    \PYG{n}{Ddiff\PYGZus{}data}  \PYG{o}{=} \PYG{n}{Ddiff\PYGZus{}sim} \PYG{o}{+} \PYG{n}{noise}
    \PYG{n}{D\PYGZus{}data}      \PYG{o}{=} \PYG{n}{numpy}\PYG{o}{.}\PYG{n}{cumsum}\PYG{p}{(}\PYG{n}{Ddiff\PYGZus{}data}\PYG{p}{)}
    \PYG{n}{D\PYGZus{}data}      \PYG{o}{=} \PYG{n}{numpy}\PYG{o}{.}\PYG{n}{concatenate}\PYG{p}{(} \PYG{p}{(}\PYG{p}{[}\PYG{l+m+mf}{0.0}\PYG{p}{]}\PYG{p}{,} \PYG{n}{D\PYGZus{}data}\PYG{p}{)} \PYG{p}{)}
    \PYG{c+c1}{\PYGZsh{}}
    \PYG{k}{return} \PYG{n}{D\PYGZus{}data}

\PYG{k}{def} \PYG{n+nf}{random\PYGZus{}start}\PYG{p}{(}\PYG{n}{n\PYGZus{}random}\PYG{p}{,} \PYG{n}{x\PYGZus{}lower}\PYG{p}{,} \PYG{n}{x\PYGZus{}upper}\PYG{p}{,} \PYG{n}{log\PYGZus{}scale}\PYG{p}{,} \PYG{n}{objective}\PYG{p}{)} \PYG{p}{:}
    \PYG{c+c1}{\PYGZsh{} check beta constraints}
    \PYG{n}{n\PYGZus{}x}      \PYG{o}{=} \PYG{n+nb}{len}\PYG{p}{(}\PYG{n}{x\PYGZus{}lower}\PYG{p}{)}
    \PYG{c+c1}{\PYGZsh{}}
    \PYG{n}{x\PYGZus{}best}   \PYG{o}{=} \PYG{p}{(}\PYG{n}{x\PYGZus{}lower} \PYG{o}{+} \PYG{n}{x\PYGZus{}upper}\PYG{p}{)} \PYG{o}{/} \PYG{l+m+mf}{2.0}
    \PYG{n}{obj\PYGZus{}best} \PYG{o}{=} \PYG{n}{objective}\PYG{p}{(}\PYG{n}{x\PYGZus{}best}\PYG{p}{)}
    \PYG{c+c1}{\PYGZsh{}}
    \PYG{n}{x\PYGZus{}low}    \PYG{o}{=} \PYG{n}{copy}\PYG{o}{.}\PYG{n}{copy}\PYG{p}{(}\PYG{n}{x\PYGZus{}lower}\PYG{p}{)}
    \PYG{n}{x\PYGZus{}up}     \PYG{o}{=} \PYG{n}{copy}\PYG{o}{.}\PYG{n}{copy}\PYG{p}{(}\PYG{n}{x\PYGZus{}upper}\PYG{p}{)}
    \PYG{k}{for} \PYG{n}{j} \PYG{o+ow}{in} \PYG{n+nb}{range}\PYG{p}{(}\PYG{n}{n\PYGZus{}x}\PYG{p}{)} \PYG{p}{:}
        \PYG{k}{if} \PYG{n}{log\PYGZus{}scale}\PYG{p}{[}\PYG{n}{j}\PYG{p}{]} \PYG{p}{:}
            \PYG{n}{x\PYGZus{}low}\PYG{p}{[}\PYG{n}{j}\PYG{p}{]} \PYG{o}{=} \PYG{n}{numpy}\PYG{o}{.}\PYG{n}{log}\PYG{p}{(}\PYG{n}{x\PYGZus{}lower}\PYG{p}{[}\PYG{n}{j}\PYG{p}{]}\PYG{p}{)}
            \PYG{n}{x\PYGZus{}up}\PYG{p}{[}\PYG{n}{j}\PYG{p}{]}  \PYG{o}{=} \PYG{n}{numpy}\PYG{o}{.}\PYG{n}{log}\PYG{p}{(}\PYG{n}{x\PYGZus{}upper}\PYG{p}{[}\PYG{n}{j}\PYG{p}{]}\PYG{p}{)}
    \PYG{c+c1}{\PYGZsh{}}
    \PYG{c+c1}{\PYGZsh{} rng: numpy random number generator}
    \PYG{n}{rng} \PYG{o}{=} \PYG{n}{numpy}\PYG{o}{.}\PYG{n}{random}\PYG{o}{.}\PYG{n}{default\PYGZus{}rng}\PYG{p}{(}\PYG{n}{actual\PYGZus{}seed}\PYG{p}{[}\PYG{l+m+mi}{0}\PYG{p}{]}\PYG{p}{)}
    \PYG{c+c1}{\PYGZsh{}}
    \PYG{k}{if} \PYG{n}{debug\PYGZus{}output} \PYG{p}{:}
        \PYG{n+nb}{print}\PYG{p}{(}\PYG{l+s+s1}{\PYGZsq{}}\PYG{l+s+s1}{random\PYGZus{}start:}\PYG{l+s+s1}{\PYGZsq{}}\PYG{p}{)}
    \PYG{k}{for} \PYG{n}{i} \PYG{o+ow}{in} \PYG{n+nb}{range}\PYG{p}{(}\PYG{n}{n\PYGZus{}random}\PYG{p}{)} \PYG{p}{:}
        \PYG{n}{x\PYGZus{}current}   \PYG{o}{=} \PYG{n}{rng}\PYG{o}{.}\PYG{n}{uniform}\PYG{p}{(}\PYG{n}{x\PYGZus{}low}\PYG{p}{,} \PYG{n}{x\PYGZus{}up}\PYG{p}{)}
        \PYG{k}{for} \PYG{n}{j} \PYG{o+ow}{in} \PYG{n+nb}{range}\PYG{p}{(}\PYG{n}{n\PYGZus{}x}\PYG{p}{)} \PYG{p}{:}
            \PYG{k}{if} \PYG{n}{log\PYGZus{}scale}\PYG{p}{[}\PYG{n}{j}\PYG{p}{]} \PYG{p}{:}
                \PYG{n}{x\PYGZus{}current}\PYG{p}{[}\PYG{n}{j}\PYG{p}{]} \PYG{o}{=} \PYG{n}{numpy}\PYG{o}{.}\PYG{n}{exp}\PYG{p}{(}\PYG{n}{x\PYGZus{}current}\PYG{p}{[}\PYG{n}{j}\PYG{p}{]}\PYG{p}{)}
        \PYG{n}{obj\PYGZus{}current} \PYG{o}{=} \PYG{n}{objective}\PYG{p}{(}\PYG{n}{x\PYGZus{}current}\PYG{p}{)}
        \PYG{k}{if} \PYG{n}{obj\PYGZus{}current} \PYG{o}{\PYGZlt{}} \PYG{n}{obj\PYGZus{}best} \PYG{p}{:}
            \PYG{k}{if} \PYG{n}{debug\PYGZus{}output} \PYG{p}{:}
                \PYG{n+nb}{print}\PYG{p}{(} \PYG{l+s+s1}{\PYGZsq{}}\PYG{l+s+s1}{sample \PYGZsh{} =}\PYG{l+s+s1}{\PYGZsq{}}\PYG{p}{,} \PYG{n}{i}\PYG{p}{,} \PYG{l+s+s1}{\PYGZsq{}}\PYG{l+s+s1}{, objective = }\PYG{l+s+s1}{\PYGZsq{}}\PYG{p}{,} \PYG{n}{obj\PYGZus{}current}\PYG{p}{)}
            \PYG{n}{x\PYGZus{}best} \PYG{o}{=} \PYG{n}{x\PYGZus{}current}
            \PYG{n}{obj\PYGZus{}best} \PYG{o}{=} \PYG{n}{obj\PYGZus{}current}
    \PYG{c+c1}{\PYGZsh{}}
    \PYG{k}{return} \PYG{n}{x\PYGZus{}best}

\PYG{k}{def} \PYG{n+nf}{display\PYGZus{}fit\PYGZus{}results}\PYG{p}{(}\PYG{n}{D\PYGZus{}data}\PYG{p}{,} \PYG{n}{x\PYGZus{}fit}\PYG{p}{,} \PYG{n}{x\PYGZus{}lower}\PYG{p}{,} \PYG{n}{x\PYGZus{}upper}\PYG{p}{,} \PYG{n}{std\PYGZus{}error}\PYG{p}{)} \PYG{p}{:}
    \PYG{n}{n\PYGZus{}x} \PYG{o}{=} \PYG{n+nb}{len}\PYG{p}{(}\PYG{n}{x\PYGZus{}fit}\PYG{p}{)}
    \PYG{c+c1}{\PYGZsh{}}
    \PYG{c+c1}{\PYGZsh{} seirwd\PYGZus{}all\PYGZus{}fit}
    \PYG{n}{seirwd\PYGZus{}all\PYGZus{}fit} \PYG{o}{=} \PYG{n}{x2seirwd\PYGZus{}all}\PYG{p}{(}\PYG{n}{x\PYGZus{}fit}\PYG{p}{)}
    \PYG{c+c1}{\PYGZsh{}}
    \PYG{c+c1}{\PYGZsh{} D\PYGZus{}fit}
    \PYG{n}{D\PYGZus{}fit} \PYG{o}{=} \PYG{n}{seirwd\PYGZus{}all\PYGZus{}fit}\PYG{p}{[}\PYG{p}{:}\PYG{p}{,}\PYG{l+m+mi}{5}\PYG{p}{]}
    \PYG{c+c1}{\PYGZsh{}}
    \PYG{c+c1}{\PYGZsh{} data\PYGZus{}residual}
    \PYG{n}{data\PYGZus{}residual} \PYG{o}{=} \PYG{n}{weighted\PYGZus{}data\PYGZus{}residual}\PYG{p}{(}\PYG{n}{D\PYGZus{}data}\PYG{p}{,} \PYG{n}{D\PYGZus{}fit}\PYG{p}{)}
    \PYG{c+c1}{\PYGZsh{}}
    \PYG{c+c1}{\PYGZsh{} \PYGZhy{}\PYGZhy{}\PYGZhy{}\PYGZhy{}\PYGZhy{}\PYGZhy{}\PYGZhy{}\PYGZhy{}\PYGZhy{}\PYGZhy{}\PYGZhy{}\PYGZhy{}\PYGZhy{}\PYGZhy{}\PYGZhy{}\PYGZhy{}\PYGZhy{}\PYGZhy{}\PYGZhy{}\PYGZhy{}\PYGZhy{}\PYGZhy{}\PYGZhy{}\PYGZhy{}\PYGZhy{}\PYGZhy{}\PYGZhy{}\PYGZhy{}\PYGZhy{}\PYGZhy{}\PYGZhy{}\PYGZhy{}\PYGZhy{}\PYGZhy{}\PYGZhy{}\PYGZhy{}\PYGZhy{}\PYGZhy{}\PYGZhy{}\PYGZhy{}\PYGZhy{}\PYGZhy{}\PYGZhy{}\PYGZhy{}\PYGZhy{}\PYGZhy{}\PYGZhy{}\PYGZhy{}\PYGZhy{}\PYGZhy{}\PYGZhy{}\PYGZhy{}\PYGZhy{}\PYGZhy{}\PYGZhy{}\PYGZhy{}\PYGZhy{}\PYGZhy{}\PYGZhy{}\PYGZhy{}\PYGZhy{}\PYGZhy{}\PYGZhy{}\PYGZhy{}\PYGZhy{}\PYGZhy{}\PYGZhy{}\PYGZhy{}\PYGZhy{}\PYGZhy{}\PYGZhy{}}
    \PYG{c+c1}{\PYGZsh{} printout}
    \PYG{n}{residual\PYGZus{}sq} \PYG{o}{=} \PYG{n}{data\PYGZus{}residual} \PYG{o}{*} \PYG{n}{data\PYGZus{}residual}
    \PYG{n+nb}{print}\PYG{p}{(}\PYG{l+s+s1}{\PYGZsq{}}\PYG{l+s+s1}{weighted data residuals}\PYG{l+s+s1}{\PYGZsq{}}\PYG{p}{)}
    \PYG{n+nb}{print}\PYG{p}{(}\PYG{l+s+s1}{\PYGZsq{}}\PYG{l+s+s1}{maximum           =}\PYG{l+s+s1}{\PYGZsq{}}\PYG{p}{,} \PYG{n}{numpy}\PYG{o}{.}\PYG{n}{max}\PYG{p}{(} \PYG{n}{numpy}\PYG{o}{.}\PYG{n}{max}\PYG{p}{(}\PYG{n}{data\PYGZus{}residual}\PYG{p}{)} \PYG{p}{)} \PYG{p}{)}
    \PYG{n+nb}{print}\PYG{p}{(}\PYG{l+s+s1}{\PYGZsq{}}\PYG{l+s+s1}{minimum           =}\PYG{l+s+s1}{\PYGZsq{}}\PYG{p}{,} \PYG{n}{numpy}\PYG{o}{.}\PYG{n}{max}\PYG{p}{(} \PYG{n}{numpy}\PYG{o}{.}\PYG{n}{min}\PYG{p}{(}\PYG{n}{data\PYGZus{}residual}\PYG{p}{)} \PYG{p}{)} \PYG{p}{)}
    \PYG{n+nb}{print}\PYG{p}{(}\PYG{l+s+s1}{\PYGZsq{}}\PYG{l+s+s1}{average           =}\PYG{l+s+s1}{\PYGZsq{}}\PYG{p}{,} \PYG{n}{numpy}\PYG{o}{.}\PYG{n}{mean}\PYG{p}{(}\PYG{n}{data\PYGZus{}residual}\PYG{p}{)} \PYG{p}{)}
    \PYG{n+nb}{print}\PYG{p}{(}\PYG{l+s+s1}{\PYGZsq{}}\PYG{l+s+s1}{average of square =}\PYG{l+s+s1}{\PYGZsq{}}\PYG{p}{,} \PYG{n}{numpy}\PYG{o}{.}\PYG{n}{mean}\PYG{p}{(}\PYG{n}{residual\PYGZus{}sq}\PYG{p}{)} \PYG{p}{)}
    \PYG{n+nb}{print}\PYG{p}{(}\PYG{l+s+s1}{\PYGZsq{}}\PYG{l+s+s1}{\PYGZsq{}}\PYG{p}{)}
    \PYG{c+c1}{\PYGZsh{}}
    \PYG{c+c1}{\PYGZsh{} table of fit values}
    \PYG{n}{fmt} \PYG{o}{=} \PYG{l+s+s1}{\PYGZsq{}}\PYG{l+s+si}{\PYGZob{}:\PYGZlt{}15s\PYGZcb{}}\PYG{l+s+si}{\PYGZob{}:\PYGZgt{}15s\PYGZcb{}}\PYG{l+s+si}{\PYGZob{}:\PYGZgt{}15s\PYGZcb{}}\PYG{l+s+si}{\PYGZob{}:\PYGZgt{}15s\PYGZcb{}}\PYG{l+s+si}{\PYGZob{}:\PYGZgt{}15s\PYGZcb{}}\PYG{l+s+s1}{\PYGZsq{}}
    \PYG{n}{line} \PYG{o}{=} \PYG{n}{fmt}\PYG{o}{.}\PYG{n}{format}\PYG{p}{(}
        \PYG{l+s+s1}{\PYGZsq{}}\PYG{l+s+s1}{x\PYGZus{}name}\PYG{l+s+s1}{\PYGZsq{}}\PYG{p}{,} \PYG{l+s+s1}{\PYGZsq{}}\PYG{l+s+s1}{x\PYGZus{}fit}\PYG{l+s+s1}{\PYGZsq{}}\PYG{p}{,}\PYG{l+s+s1}{\PYGZsq{}}\PYG{l+s+s1}{x\PYGZus{}lower}\PYG{l+s+s1}{\PYGZsq{}}\PYG{p}{,}\PYG{l+s+s1}{\PYGZsq{}}\PYG{l+s+s1}{x\PYGZus{}upper}\PYG{l+s+s1}{\PYGZsq{}}\PYG{p}{,} \PYG{l+s+s1}{\PYGZsq{}}\PYG{l+s+s1}{std\PYGZus{}error}\PYG{l+s+s1}{\PYGZsq{}}
    \PYG{p}{)}
    \PYG{n+nb}{print}\PYG{p}{(}\PYG{n}{line}\PYG{p}{)}
    \PYG{k}{for} \PYG{n}{i} \PYG{o+ow}{in} \PYG{n+nb}{range}\PYG{p}{(} \PYG{n}{n\PYGZus{}x} \PYG{p}{)} \PYG{p}{:}
        \PYG{n}{fmt} \PYG{o}{=} \PYG{l+s+s1}{\PYGZsq{}}\PYG{l+s+si}{\PYGZob{}:\PYGZlt{}15s\PYGZcb{}}\PYG{l+s+si}{\PYGZob{}:+15.9f\PYGZcb{}}\PYG{l+s+si}{\PYGZob{}:+15.9f\PYGZcb{}}\PYG{l+s+si}{\PYGZob{}:+15.9f\PYGZcb{}}\PYG{l+s+si}{\PYGZob{}:+15.9f\PYGZcb{}}\PYG{l+s+s1}{\PYGZsq{}}
        \PYG{n}{line} \PYG{o}{=} \PYG{n}{fmt}\PYG{o}{.}\PYG{n}{format}\PYG{p}{(}
            \PYG{n}{x\PYGZus{}name}\PYG{p}{[}\PYG{n}{i}\PYG{p}{]}\PYG{p}{,} \PYG{n}{x\PYGZus{}fit}\PYG{p}{[}\PYG{n}{i}\PYG{p}{]}\PYG{p}{,} \PYG{n}{x\PYGZus{}lower}\PYG{p}{[}\PYG{n}{i}\PYG{p}{]}\PYG{p}{,} \PYG{n}{x\PYGZus{}upper}\PYG{p}{[}\PYG{n}{i}\PYG{p}{]}\PYG{p}{,} \PYG{n}{std\PYGZus{}error}\PYG{p}{[}\PYG{n}{i}\PYG{p}{]}
        \PYG{p}{)}
        \PYG{n+nb}{print}\PYG{p}{(}\PYG{n}{line}\PYG{p}{)}
    \PYG{k}{if} \PYG{n}{data\PYGZus{}file} \PYG{o}{==} \PYG{l+s+s1}{\PYGZsq{}}\PYG{l+s+s1}{\PYGZsq{}} \PYG{p}{:}
        \PYG{c+c1}{\PYGZsh{} table of fit versus simulation}
        \PYG{n+nb}{print}\PYG{p}{(}\PYG{l+s+s1}{\PYGZsq{}}\PYG{l+s+s1}{\PYGZsq{}}\PYG{p}{)}
        \PYG{n}{fmt} \PYG{o}{=} \PYG{l+s+s1}{\PYGZsq{}}\PYG{l+s+si}{\PYGZob{}:\PYGZlt{}15s\PYGZcb{}}\PYG{l+s+si}{\PYGZob{}:\PYGZgt{}15s\PYGZcb{}}\PYG{l+s+si}{\PYGZob{}:\PYGZgt{}15s\PYGZcb{}}\PYG{l+s+si}{\PYGZob{}:\PYGZgt{}15s\PYGZcb{}}\PYG{l+s+si}{\PYGZob{}:\PYGZgt{}15s\PYGZcb{}}\PYG{l+s+s1}{\PYGZsq{}}
        \PYG{n}{line} \PYG{o}{=} \PYG{n}{fmt}\PYG{o}{.}\PYG{n}{format}\PYG{p}{(}
            \PYG{l+s+s1}{\PYGZsq{}}\PYG{l+s+s1}{x\PYGZus{}name}\PYG{l+s+s1}{\PYGZsq{}}\PYG{p}{,} \PYG{l+s+s1}{\PYGZsq{}}\PYG{l+s+s1}{x\PYGZus{}fit}\PYG{l+s+s1}{\PYGZsq{}}\PYG{p}{,}\PYG{l+s+s1}{\PYGZsq{}}\PYG{l+s+s1}{x\PYGZus{}sim}\PYG{l+s+s1}{\PYGZsq{}}\PYG{p}{,}\PYG{l+s+s1}{\PYGZsq{}}\PYG{l+s+s1}{rel\PYGZus{}error}\PYG{l+s+s1}{\PYGZsq{}}\PYG{p}{,}\PYG{l+s+s1}{\PYGZsq{}}\PYG{l+s+s1}{residual}\PYG{l+s+s1}{\PYGZsq{}}
        \PYG{p}{)}
        \PYG{n+nb}{print}\PYG{p}{(}\PYG{n}{line}\PYG{p}{)}
        \PYG{k}{for} \PYG{n}{i} \PYG{o+ow}{in} \PYG{n+nb}{range}\PYG{p}{(} \PYG{n}{n\PYGZus{}x} \PYG{p}{)} \PYG{p}{:}
            \PYG{k}{if} \PYG{n}{x\PYGZus{}sim}\PYG{p}{[}\PYG{n}{i}\PYG{p}{]} \PYG{o}{==} \PYG{l+m+mf}{0.0} \PYG{p}{:}
                \PYG{n}{rel\PYGZus{}error} \PYG{o}{=} \PYG{l+m+mf}{0.0}
            \PYG{k}{else} \PYG{p}{:}
                \PYG{n}{rel\PYGZus{}error} \PYG{o}{=} \PYG{n}{x\PYGZus{}fit}\PYG{p}{[}\PYG{n}{i}\PYG{p}{]} \PYG{o}{/} \PYG{n}{x\PYGZus{}sim}\PYG{p}{[}\PYG{n}{i}\PYG{p}{]} \PYG{o}{\PYGZhy{}} \PYG{l+m+mf}{1.0}
            \PYG{n}{residual}  \PYG{o}{=} \PYG{p}{(}\PYG{n}{x\PYGZus{}fit}\PYG{p}{[}\PYG{n}{i}\PYG{p}{]} \PYG{o}{\PYGZhy{}} \PYG{n}{x\PYGZus{}sim}\PYG{p}{[}\PYG{n}{i}\PYG{p}{]}\PYG{p}{)} \PYG{o}{/} \PYG{n}{std\PYGZus{}error}\PYG{p}{[}\PYG{n}{i}\PYG{p}{]}
            \PYG{n}{fmt} \PYG{o}{=} \PYG{l+s+s1}{\PYGZsq{}}\PYG{l+s+si}{\PYGZob{}:\PYGZlt{}15s\PYGZcb{}}\PYG{l+s+si}{\PYGZob{}:+15.9f\PYGZcb{}}\PYG{l+s+si}{\PYGZob{}:+15.9f\PYGZcb{}}\PYG{l+s+si}{\PYGZob{}:+15.9f\PYGZcb{}}\PYG{l+s+si}{\PYGZob{}:+15.9f\PYGZcb{}}\PYG{l+s+s1}{\PYGZsq{}}
            \PYG{n}{line} \PYG{o}{=} \PYG{n}{fmt}\PYG{o}{.}\PYG{n}{format}\PYG{p}{(}\PYG{n}{x\PYGZus{}name}\PYG{p}{[}\PYG{n}{i}\PYG{p}{]}\PYG{p}{,}
                \PYG{n}{x\PYGZus{}fit}\PYG{p}{[}\PYG{n}{i}\PYG{p}{]}\PYG{p}{,} \PYG{n}{x\PYGZus{}sim}\PYG{p}{[}\PYG{n}{i}\PYG{p}{]}\PYG{p}{,} \PYG{n}{rel\PYGZus{}error}\PYG{p}{,} \PYG{n}{residual}
            \PYG{p}{)}
            \PYG{n+nb}{print}\PYG{p}{(}\PYG{n}{line}\PYG{p}{)}
    \PYG{c+c1}{\PYGZsh{} \PYGZhy{}\PYGZhy{}\PYGZhy{}\PYGZhy{}\PYGZhy{}\PYGZhy{}\PYGZhy{}\PYGZhy{}\PYGZhy{}\PYGZhy{}\PYGZhy{}\PYGZhy{}\PYGZhy{}\PYGZhy{}\PYGZhy{}\PYGZhy{}\PYGZhy{}\PYGZhy{}\PYGZhy{}\PYGZhy{}\PYGZhy{}\PYGZhy{}\PYGZhy{}\PYGZhy{}\PYGZhy{}\PYGZhy{}\PYGZhy{}\PYGZhy{}\PYGZhy{}\PYGZhy{}\PYGZhy{}\PYGZhy{}\PYGZhy{}\PYGZhy{}\PYGZhy{}\PYGZhy{}\PYGZhy{}\PYGZhy{}\PYGZhy{}\PYGZhy{}\PYGZhy{}\PYGZhy{}\PYGZhy{}\PYGZhy{}\PYGZhy{}\PYGZhy{}\PYGZhy{}\PYGZhy{}\PYGZhy{}\PYGZhy{}\PYGZhy{}\PYGZhy{}\PYGZhy{}\PYGZhy{}\PYGZhy{}\PYGZhy{}\PYGZhy{}\PYGZhy{}\PYGZhy{}\PYGZhy{}\PYGZhy{}\PYGZhy{}\PYGZhy{}\PYGZhy{}\PYGZhy{}\PYGZhy{}\PYGZhy{}\PYGZhy{}\PYGZhy{}\PYGZhy{}\PYGZhy{}}
    \PYG{c+c1}{\PYGZsh{} plot}
    \PYG{c+c1}{\PYGZsh{} \PYGZhy{}\PYGZhy{}\PYGZhy{}\PYGZhy{}\PYGZhy{}\PYGZhy{}\PYGZhy{}\PYGZhy{}\PYGZhy{}\PYGZhy{}\PYGZhy{}\PYGZhy{}\PYGZhy{}\PYGZhy{}\PYGZhy{}\PYGZhy{}\PYGZhy{}\PYGZhy{}\PYGZhy{}\PYGZhy{}\PYGZhy{}\PYGZhy{}\PYGZhy{}\PYGZhy{}\PYGZhy{}\PYGZhy{}\PYGZhy{}\PYGZhy{}\PYGZhy{}\PYGZhy{}\PYGZhy{}\PYGZhy{}\PYGZhy{}\PYGZhy{}\PYGZhy{}\PYGZhy{}\PYGZhy{}\PYGZhy{}\PYGZhy{}\PYGZhy{}\PYGZhy{}\PYGZhy{}\PYGZhy{}\PYGZhy{}\PYGZhy{}\PYGZhy{}\PYGZhy{}\PYGZhy{}\PYGZhy{}\PYGZhy{}\PYGZhy{}\PYGZhy{}\PYGZhy{}\PYGZhy{}\PYGZhy{}\PYGZhy{}\PYGZhy{}\PYGZhy{}\PYGZhy{}\PYGZhy{}\PYGZhy{}\PYGZhy{}\PYGZhy{}\PYGZhy{}\PYGZhy{}\PYGZhy{}\PYGZhy{}\PYGZhy{}\PYGZhy{}\PYGZhy{}\PYGZhy{}}
    \PYG{c+c1}{\PYGZsh{}}
    \PYG{n}{fig}  \PYG{o}{=} \PYG{n}{pyplot}\PYG{o}{.}\PYG{n}{figure}\PYG{p}{(}\PYG{n}{tight\PYGZus{}layout} \PYG{o}{=} \PYG{k+kc}{True}\PYG{p}{)}
    \PYG{n}{gs}   \PYG{o}{=} \PYG{n}{gridspec}\PYG{o}{.}\PYG{n}{GridSpec}\PYG{p}{(}\PYG{l+m+mi}{3}\PYG{p}{,} \PYG{l+m+mi}{2}\PYG{p}{)}
    \PYG{n}{ax1}  \PYG{o}{=} \PYG{n}{fig}\PYG{o}{.}\PYG{n}{add\PYGZus{}subplot}\PYG{p}{(}\PYG{n}{gs}\PYG{p}{[}\PYG{p}{:}\PYG{p}{,}\PYG{l+m+mi}{0}\PYG{p}{]}\PYG{p}{)}
    \PYG{n}{ax2}  \PYG{o}{=} \PYG{n}{fig}\PYG{o}{.}\PYG{n}{add\PYGZus{}subplot}\PYG{p}{(}\PYG{n}{gs}\PYG{p}{[}\PYG{l+m+mi}{0}\PYG{p}{,}\PYG{l+m+mi}{1}\PYG{p}{]}\PYG{p}{)}
    \PYG{n}{ax3}  \PYG{o}{=} \PYG{n}{fig}\PYG{o}{.}\PYG{n}{add\PYGZus{}subplot}\PYG{p}{(}\PYG{n}{gs}\PYG{p}{[}\PYG{l+m+mi}{1}\PYG{p}{,}\PYG{l+m+mi}{1}\PYG{p}{]}\PYG{p}{)}
    \PYG{n}{ax4}  \PYG{o}{=} \PYG{n}{fig}\PYG{o}{.}\PYG{n}{add\PYGZus{}subplot}\PYG{p}{(}\PYG{n}{gs}\PYG{p}{[}\PYG{l+m+mi}{2}\PYG{p}{,}\PYG{l+m+mi}{1}\PYG{p}{]}\PYG{p}{)}
    \PYG{c+c1}{\PYGZsh{}}
    \PYG{n}{ax1}\PYG{o}{.}\PYG{n}{plot}\PYG{p}{(}\PYG{n}{t\PYGZus{}all}\PYG{p}{,} \PYG{n}{seirwd\PYGZus{}all\PYGZus{}fit}\PYG{p}{[}\PYG{p}{:}\PYG{p}{,}\PYG{l+m+mi}{1}\PYG{p}{]}\PYG{p}{,} \PYG{l+s+s1}{\PYGZsq{}}\PYG{l+s+s1}{g\PYGZhy{}}\PYG{l+s+s1}{\PYGZsq{}}\PYG{p}{,} \PYG{n}{label}\PYG{o}{=}\PYG{l+s+s1}{\PYGZsq{}}\PYG{l+s+s1}{E}\PYG{l+s+s1}{\PYGZsq{}}\PYG{p}{)}
    \PYG{n}{ax1}\PYG{o}{.}\PYG{n}{plot}\PYG{p}{(}\PYG{n}{t\PYGZus{}all}\PYG{p}{,} \PYG{n}{seirwd\PYGZus{}all\PYGZus{}fit}\PYG{p}{[}\PYG{p}{:}\PYG{p}{,}\PYG{l+m+mi}{2}\PYG{p}{]}\PYG{p}{,} \PYG{l+s+s1}{\PYGZsq{}}\PYG{l+s+s1}{r\PYGZhy{}}\PYG{l+s+s1}{\PYGZsq{}}\PYG{p}{,} \PYG{n}{label}\PYG{o}{=}\PYG{l+s+s1}{\PYGZsq{}}\PYG{l+s+s1}{I}\PYG{l+s+s1}{\PYGZsq{}}\PYG{p}{)}
    \PYG{n}{ax1}\PYG{o}{.}\PYG{n}{plot}\PYG{p}{(}\PYG{n}{t\PYGZus{}all}\PYG{p}{,} \PYG{n}{seirwd\PYGZus{}all\PYGZus{}fit}\PYG{p}{[}\PYG{p}{:}\PYG{p}{,}\PYG{l+m+mi}{3}\PYG{p}{]}\PYG{p}{,} \PYG{l+s+s1}{\PYGZsq{}}\PYG{l+s+s1}{k\PYGZhy{}}\PYG{l+s+s1}{\PYGZsq{}}\PYG{p}{,} \PYG{n}{label}\PYG{o}{=}\PYG{l+s+s1}{\PYGZsq{}}\PYG{l+s+s1}{R}\PYG{l+s+s1}{\PYGZsq{}}\PYG{p}{)}
    \PYG{n}{ax1}\PYG{o}{.}\PYG{n}{plot}\PYG{p}{(}\PYG{n}{t\PYGZus{}all}\PYG{p}{,} \PYG{n}{seirwd\PYGZus{}all\PYGZus{}fit}\PYG{p}{[}\PYG{p}{:}\PYG{p}{,}\PYG{l+m+mi}{4}\PYG{p}{]}\PYG{p}{,} \PYG{l+s+s1}{\PYGZsq{}}\PYG{l+s+s1}{y\PYGZhy{}}\PYG{l+s+s1}{\PYGZsq{}}\PYG{p}{,} \PYG{n}{label}\PYG{o}{=}\PYG{l+s+s1}{\PYGZsq{}}\PYG{l+s+s1}{W}\PYG{l+s+s1}{\PYGZsq{}}\PYG{p}{)}
    \PYG{n}{ax1}\PYG{o}{.}\PYG{n}{plot}\PYG{p}{(}\PYG{n}{t\PYGZus{}all}\PYG{p}{,} \PYG{n}{seirwd\PYGZus{}all\PYGZus{}fit}\PYG{p}{[}\PYG{p}{:}\PYG{p}{,}\PYG{l+m+mi}{5}\PYG{p}{]}\PYG{p}{,} \PYG{l+s+s1}{\PYGZsq{}}\PYG{l+s+s1}{b\PYGZhy{}}\PYG{l+s+s1}{\PYGZsq{}}\PYG{p}{,} \PYG{n}{label}\PYG{o}{=}\PYG{l+s+s1}{\PYGZsq{}}\PYG{l+s+s1}{D}\PYG{l+s+s1}{\PYGZsq{}}\PYG{p}{)}
    \PYG{n}{ax1}\PYG{o}{.}\PYG{n}{legend}\PYG{p}{(}\PYG{p}{)}
    \PYG{n}{ax1}\PYG{o}{.}\PYG{n}{set\PYGZus{}xlabel}\PYG{p}{(}\PYG{l+s+s1}{\PYGZsq{}}\PYG{l+s+s1}{time}\PYG{l+s+s1}{\PYGZsq{}}\PYG{p}{)}
    \PYG{n}{ax1}\PYG{o}{.}\PYG{n}{set\PYGZus{}ylabel}\PYG{p}{(}\PYG{l+s+s1}{\PYGZsq{}}\PYG{l+s+s1}{population fraction}\PYG{l+s+s1}{\PYGZsq{}}\PYG{p}{)}
    \PYG{c+c1}{\PYGZsh{}}
    \PYG{n}{Ddiff\PYGZus{}data} \PYG{o}{=} \PYG{n}{numpy}\PYG{o}{.}\PYG{n}{diff}\PYG{p}{(}\PYG{n}{D\PYGZus{}data}\PYG{p}{)}
    \PYG{n}{Ddiff\PYGZus{}fit}  \PYG{o}{=} \PYG{n}{numpy}\PYG{o}{.}\PYG{n}{diff}\PYG{p}{(}\PYG{n}{D\PYGZus{}fit}\PYG{p}{)}
    \PYG{n}{t\PYGZus{}mid}      \PYG{o}{=} \PYG{p}{(}\PYG{n}{t\PYGZus{}all}\PYG{p}{[}\PYG{l+m+mi}{0} \PYG{p}{:} \PYG{o}{\PYGZhy{}}\PYG{l+m+mi}{1}\PYG{p}{]} \PYG{o}{+} \PYG{n}{t\PYGZus{}all}\PYG{p}{[}\PYG{l+m+mi}{1} \PYG{p}{:}\PYG{p}{]}\PYG{p}{)} \PYG{o}{/} \PYG{l+m+mf}{2.0}
    \PYG{n}{ax2}\PYG{o}{.}\PYG{n}{plot}\PYG{p}{(}\PYG{n}{t\PYGZus{}mid}\PYG{p}{,} \PYG{n}{Ddiff\PYGZus{}data}\PYG{p}{,} \PYG{l+s+s1}{\PYGZsq{}}\PYG{l+s+s1}{k+}\PYG{l+s+s1}{\PYGZsq{}} \PYG{p}{,} \PYG{n}{label}\PYG{o}{=}\PYG{l+s+s1}{\PYGZsq{}}\PYG{l+s+s1}{data}\PYG{l+s+s1}{\PYGZsq{}}\PYG{p}{)}
    \PYG{n}{ax2}\PYG{o}{.}\PYG{n}{plot}\PYG{p}{(}\PYG{n}{t\PYGZus{}mid}\PYG{p}{,} \PYG{n}{Ddiff\PYGZus{}fit}\PYG{p}{,}  \PYG{l+s+s1}{\PYGZsq{}}\PYG{l+s+s1}{k\PYGZhy{}}\PYG{l+s+s1}{\PYGZsq{}} \PYG{p}{,} \PYG{n}{label}\PYG{o}{=}\PYG{l+s+s1}{\PYGZsq{}}\PYG{l+s+s1}{fit}\PYG{l+s+s1}{\PYGZsq{}}\PYG{p}{)}
    \PYG{n}{ax2}\PYG{o}{.}\PYG{n}{legend}\PYG{p}{(}\PYG{p}{)}
    \PYG{n}{ax2}\PYG{o}{.}\PYG{n}{set\PYGZus{}ylabel}\PYG{p}{(}\PYG{l+s+s1}{\PYGZsq{}}\PYG{l+s+s1}{death differences}\PYG{l+s+s1}{\PYGZsq{}}\PYG{p}{)}
    \PYG{c+c1}{\PYGZsh{}}
    \PYG{n}{ax3}\PYG{o}{.}\PYG{n}{plot}\PYG{p}{(} \PYG{n}{t\PYGZus{}mid}\PYG{p}{,}             \PYG{n}{data\PYGZus{}residual}\PYG{p}{,}  \PYG{l+s+s1}{\PYGZsq{}}\PYG{l+s+s1}{k+}\PYG{l+s+s1}{\PYGZsq{}} \PYG{p}{)}
    \PYG{n}{ax3}\PYG{o}{.}\PYG{n}{plot}\PYG{p}{(} \PYG{p}{[}\PYG{n}{t\PYGZus{}all}\PYG{p}{[}\PYG{l+m+mi}{0}\PYG{p}{]}\PYG{p}{,} \PYG{n}{t\PYGZus{}all}\PYG{p}{[}\PYG{o}{\PYGZhy{}}\PYG{l+m+mi}{1}\PYG{p}{]}\PYG{p}{]}\PYG{p}{,} \PYG{p}{[}\PYG{l+m+mf}{0.0}\PYG{p}{,} \PYG{l+m+mf}{0.0}\PYG{p}{]}\PYG{p}{,} \PYG{l+s+s1}{\PYGZsq{}}\PYG{l+s+s1}{k\PYGZhy{}}\PYG{l+s+s1}{\PYGZsq{}} \PYG{p}{)}
    \PYG{n}{ax3}\PYG{o}{.}\PYG{n}{set\PYGZus{}ylabel}\PYG{p}{(}\PYG{l+s+s1}{\PYGZsq{}}\PYG{l+s+s1}{weighted residuals}\PYG{l+s+s1}{\PYGZsq{}}\PYG{p}{)}
    \PYG{c+c1}{\PYGZsh{}}
    \PYG{n}{cov\PYGZus{}mul}  \PYG{o}{=} \PYG{n}{x\PYGZus{}fit}\PYG{p}{[}\PYG{l+m+mi}{0} \PYG{p}{:} \PYG{l+m+mi}{3}\PYG{p}{]}
    \PYG{n}{beta\PYGZus{}bar} \PYG{o}{=} \PYG{n}{x\PYGZus{}fit}\PYG{p}{[}\PYG{l+m+mi}{5}\PYG{p}{]}
    \PYG{n}{beta\PYGZus{}all} \PYG{o}{=} \PYG{n}{beta\PYGZus{}bar} \PYG{o}{*} \PYG{n}{numpy}\PYG{o}{.}\PYG{n}{exp}\PYG{p}{(} \PYG{n}{numpy}\PYG{o}{.}\PYG{n}{matmul}\PYG{p}{(}\PYG{n}{covariates}\PYG{p}{,} \PYG{n}{cov\PYGZus{}mul}\PYG{p}{)} \PYG{p}{)}
    \PYG{n}{ax4}\PYG{o}{.}\PYG{n}{plot}\PYG{p}{(}\PYG{n}{t\PYGZus{}all}\PYG{p}{,} \PYG{n}{beta\PYGZus{}all}\PYG{p}{,} \PYG{l+s+s1}{\PYGZsq{}}\PYG{l+s+s1}{k\PYGZhy{}}\PYG{l+s+s1}{\PYGZsq{}}\PYG{p}{)}
    \PYG{n}{ax4}\PYG{o}{.}\PYG{n}{set\PYGZus{}xlabel}\PYG{p}{(}\PYG{l+s+s1}{\PYGZsq{}}\PYG{l+s+s1}{time}\PYG{l+s+s1}{\PYGZsq{}}\PYG{p}{)}
    \PYG{n}{ax4}\PYG{o}{.}\PYG{n}{set\PYGZus{}ylabel}\PYG{p}{(}\PYG{l+s+s1}{\PYGZsq{}}\PYG{l+s+s1}{beta(t)}\PYG{l+s+s1}{\PYGZsq{}}\PYG{p}{)}
    \PYG{c+c1}{\PYGZsh{}}
    \PYG{n}{pyplot}\PYG{o}{.}\PYG{n}{show}\PYG{p}{(}\PYG{p}{)}

\PYG{k}{def} \PYG{n+nf}{covid\PYGZus{}19\PYGZus{}xam}\PYG{p}{(}\PYG{n}{call\PYGZus{}count} \PYG{o}{=} \PYG{l+m+mi}{0}\PYG{p}{)} \PYG{p}{:}
    \PYG{n}{ok} \PYG{o}{=} \PYG{k+kc}{True}
    \PYG{c+c1}{\PYGZsh{}}
    \PYG{c+c1}{\PYGZsh{} if random seed is zero, use the clock to see the generator.}
    \PYG{k}{if} \PYG{n}{random\PYGZus{}seed} \PYG{o}{==} \PYG{l+m+mi}{0} \PYG{p}{:}
        \PYG{n}{actual\PYGZus{}seed}\PYG{p}{[}\PYG{l+m+mi}{0}\PYG{p}{]} \PYG{o}{=} \PYG{n+nb}{int}\PYG{p}{(} \PYG{l+m+mi}{13} \PYG{o}{*} \PYG{n}{time}\PYG{o}{.}\PYG{n}{time}\PYG{p}{(}\PYG{p}{)} \PYG{p}{)}
    \PYG{k}{if} \PYG{n}{debug\PYGZus{}output} \PYG{p}{:}
        \PYG{n+nb}{print}\PYG{p}{(}\PYG{l+s+s1}{\PYGZsq{}}\PYG{l+s+s1}{actual\PYGZus{}seed = }\PYG{l+s+s1}{\PYGZsq{}}\PYG{p}{,} \PYG{n}{actual\PYGZus{}seed}\PYG{p}{[}\PYG{l+m+mi}{0}\PYG{p}{]}\PYG{p}{)}
    \PYG{c+c1}{\PYGZsh{}}
    \PYG{c+c1}{\PYGZsh{} D\PYGZus{}data}
    \PYG{k}{if} \PYG{n}{data\PYGZus{}file} \PYG{o}{==} \PYG{l+s+s1}{\PYGZsq{}}\PYG{l+s+s1}{\PYGZsq{}} \PYG{p}{:}
        \PYG{n}{D\PYGZus{}data} \PYG{o}{=} \PYG{n}{simulate\PYGZus{}data}\PYG{p}{(}\PYG{p}{)}
    \PYG{k}{else} \PYG{p}{:}
        \PYG{n}{D\PYGZus{}data} \PYG{o}{=} \PYG{n+nb}{list}\PYG{p}{(}\PYG{p}{)}
        \PYG{k}{for} \PYG{n}{row} \PYG{o+ow}{in} \PYG{n}{file\PYGZus{}data} \PYG{p}{:}
            \PYG{n}{D\PYGZus{}data}\PYG{o}{.}\PYG{n}{append}\PYG{p}{(} \PYG{n+nb}{float}\PYG{p}{(} \PYG{n}{row}\PYG{p}{[}\PYG{l+s+s1}{\PYGZsq{}}\PYG{l+s+s1}{death}\PYG{l+s+s1}{\PYGZsq{}}\PYG{p}{]} \PYG{p}{)} \PYG{p}{)}
        \PYG{n}{D\PYGZus{}data} \PYG{o}{=} \PYG{n}{numpy}\PYG{o}{.}\PYG{n}{array}\PYG{p}{(}\PYG{n}{D\PYGZus{}data}\PYG{p}{)}
    \PYG{c+c1}{\PYGZsh{}}
    \PYG{c+c1}{\PYGZsh{} objective\PYGZus{}ad}
    \PYG{n}{objective\PYGZus{}ad} \PYG{o}{=} \PYG{n}{objective\PYGZus{}d\PYGZus{}fun}\PYG{p}{(}\PYG{n}{t\PYGZus{}all}\PYG{p}{,} \PYG{n}{D\PYGZus{}data}\PYG{p}{)}
    \PYG{c+c1}{\PYGZsh{}}
    \PYG{c+c1}{\PYGZsh{} objective: fun, grad, hess}
    \PYG{n}{optimize\PYGZus{}fun} \PYG{o}{=} \PYG{n}{optimize\PYGZus{}fun\PYGZus{}class}\PYG{p}{(}\PYG{n}{objective\PYGZus{}ad}\PYG{p}{)}
    \PYG{c+c1}{\PYGZsh{}}
    \PYG{c+c1}{\PYGZsh{} options}
    \PYG{n}{options} \PYG{o}{=} \PYG{p}{\PYGZob{}}
        \PYG{l+s+s1}{\PYGZsq{}}\PYG{l+s+s1}{gtol}\PYG{l+s+s1}{\PYGZsq{}}    \PYG{p}{:} \PYG{l+m+mf}{1e\PYGZhy{}10}\PYG{p}{,}
        \PYG{l+s+s1}{\PYGZsq{}}\PYG{l+s+s1}{xtol}\PYG{l+s+s1}{\PYGZsq{}}    \PYG{p}{:} \PYG{l+m+mf}{1e\PYGZhy{}8}\PYG{p}{,}
        \PYG{l+s+s1}{\PYGZsq{}}\PYG{l+s+s1}{maxiter}\PYG{l+s+s1}{\PYGZsq{}} \PYG{p}{:} \PYG{l+m+mi}{300}\PYG{p}{,}
        \PYG{l+s+s1}{\PYGZsq{}}\PYG{l+s+s1}{verbose}\PYG{l+s+s1}{\PYGZsq{}} \PYG{p}{:} \PYG{l+m+mi}{0}\PYG{p}{,}
    \PYG{p}{\PYGZcb{}}
    \PYG{c+c1}{\PYGZsh{} \PYGZhy{}\PYGZhy{}\PYGZhy{}\PYGZhy{}\PYGZhy{}\PYGZhy{}\PYGZhy{}\PYGZhy{}\PYGZhy{}\PYGZhy{}\PYGZhy{}\PYGZhy{}\PYGZhy{}\PYGZhy{}\PYGZhy{}\PYGZhy{}\PYGZhy{}\PYGZhy{}\PYGZhy{}\PYGZhy{}\PYGZhy{}\PYGZhy{}\PYGZhy{}\PYGZhy{}\PYGZhy{}\PYGZhy{}\PYGZhy{}\PYGZhy{}\PYGZhy{}\PYGZhy{}\PYGZhy{}\PYGZhy{}\PYGZhy{}\PYGZhy{}\PYGZhy{}\PYGZhy{}\PYGZhy{}\PYGZhy{}\PYGZhy{}\PYGZhy{}\PYGZhy{}\PYGZhy{}\PYGZhy{}\PYGZhy{}\PYGZhy{}\PYGZhy{}\PYGZhy{}\PYGZhy{}\PYGZhy{}\PYGZhy{}\PYGZhy{}\PYGZhy{}\PYGZhy{}\PYGZhy{}\PYGZhy{}\PYGZhy{}\PYGZhy{}\PYGZhy{}\PYGZhy{}\PYGZhy{}\PYGZhy{}\PYGZhy{}\PYGZhy{}\PYGZhy{}\PYGZhy{}\PYGZhy{}\PYGZhy{}\PYGZhy{}\PYGZhy{}\PYGZhy{}\PYGZhy{}}
    \PYG{c+c1}{\PYGZsh{} This code is used by the \PYGZdl{}subhead Actual Bounds\PYGZdl{}\PYGZdl{} documentation}
\PYG{c+c1}{\PYGZsh{} BEGIN\PYGZus{}ACTUAL\PYGZus{}BOUNDS}
    \PYG{n}{n\PYGZus{}x} \PYG{o}{=} \PYG{n+nb}{len}\PYG{p}{(}\PYG{n}{x\PYGZus{}name}\PYG{p}{)}
    \PYG{n}{x\PYGZus{}lower}                \PYG{o}{=} \PYG{n}{numpy}\PYG{o}{.}\PYG{n}{zeros}\PYG{p}{(} \PYG{n}{n\PYGZus{}x}\PYG{p}{,} \PYG{n}{dtype}\PYG{o}{=}\PYG{n+nb}{float}\PYG{p}{)}
    \PYG{n}{x\PYGZus{}upper}                \PYG{o}{=} \PYG{n}{numpy}\PYG{o}{.}\PYG{n}{zeros}\PYG{p}{(} \PYG{n}{n\PYGZus{}x}\PYG{p}{,} \PYG{n}{dtype}\PYG{o}{=}\PYG{n+nb}{float}\PYG{p}{)}
    \PYG{n}{x\PYGZus{}upper}\PYG{p}{[} \PYG{n}{x\PYGZus{}sim} \PYG{o}{\PYGZgt{}} \PYG{l+m+mf}{0.0} \PYG{p}{]} \PYG{o}{=} \PYG{n}{actual\PYGZus{}bound\PYGZus{}factor} \PYG{o}{*} \PYG{n}{x\PYGZus{}sim}\PYG{p}{[} \PYG{n}{x\PYGZus{}sim} \PYG{o}{\PYGZgt{}} \PYG{l+m+mf}{0.0} \PYG{p}{]}
    \PYG{n}{x\PYGZus{}lower}\PYG{p}{[} \PYG{n}{x\PYGZus{}sim} \PYG{o}{\PYGZlt{}} \PYG{l+m+mf}{0.0} \PYG{p}{]} \PYG{o}{=} \PYG{n}{actual\PYGZus{}bound\PYGZus{}factor} \PYG{o}{*} \PYG{n}{x\PYGZus{}sim}\PYG{p}{[} \PYG{n}{x\PYGZus{}sim} \PYG{o}{\PYGZlt{}} \PYG{l+m+mf}{0.0} \PYG{p}{]}
\PYG{c+c1}{\PYGZsh{} END\PYGZus{}ACTUAL\PYGZus{}BOUNDS}
    \PYG{c+c1}{\PYGZsh{}}
    \PYG{c+c1}{\PYGZsh{} currently not using log\PYGZhy{}scaling}
    \PYG{n}{log\PYGZus{}scale}    \PYG{o}{=} \PYG{n}{numpy}\PYG{o}{.}\PYG{n}{array}\PYG{p}{(} \PYG{n}{n\PYGZus{}x} \PYG{o}{*} \PYG{p}{[} \PYG{k+kc}{False} \PYG{p}{]} \PYG{p}{)}
    \PYG{c+c1}{\PYGZsh{}}
    \PYG{k}{assert} \PYG{n}{numpy}\PYG{o}{.}\PYG{n}{all}\PYG{p}{(} \PYG{n}{x\PYGZus{}lower} \PYG{o}{\PYGZlt{}}\PYG{o}{=} \PYG{n}{x\PYGZus{}sim} \PYG{p}{)}
    \PYG{k}{assert} \PYG{n}{numpy}\PYG{o}{.}\PYG{n}{all}\PYG{p}{(} \PYG{n}{x\PYGZus{}sim} \PYG{o}{\PYGZlt{}}\PYG{o}{=} \PYG{n}{x\PYGZus{}upper} \PYG{p}{)}
    \PYG{c+c1}{\PYGZsh{}}
    \PYG{n}{bounds} \PYG{o}{=} \PYG{n}{scipy}\PYG{o}{.}\PYG{n}{optimize}\PYG{o}{.}\PYG{n}{Bounds}\PYG{p}{(}
        \PYG{n}{x\PYGZus{}lower}\PYG{p}{,}
        \PYG{n}{x\PYGZus{}upper}\PYG{p}{,}
        \PYG{n}{keep\PYGZus{}feasible} \PYG{o}{=} \PYG{k+kc}{True}
    \PYG{p}{)}
    \PYG{k}{if} \PYG{n}{debug\PYGZus{}output} \PYG{p}{:}
        \PYG{n+nb}{print}\PYG{p}{(}\PYG{l+s+s1}{\PYGZsq{}}\PYG{l+s+s1}{x\PYGZus{}lower =}\PYG{l+s+s1}{\PYGZsq{}}\PYG{p}{,} \PYG{n}{x\PYGZus{}lower}\PYG{p}{)}
        \PYG{n+nb}{print}\PYG{p}{(}\PYG{l+s+s1}{\PYGZsq{}}\PYG{l+s+s1}{x\PYGZus{}upper =}\PYG{l+s+s1}{\PYGZsq{}}\PYG{p}{,} \PYG{n}{x\PYGZus{}upper}\PYG{p}{)}
    \PYG{c+c1}{\PYGZsh{} \PYGZhy{}\PYGZhy{}\PYGZhy{}\PYGZhy{}\PYGZhy{}\PYGZhy{}\PYGZhy{}\PYGZhy{}\PYGZhy{}\PYGZhy{}\PYGZhy{}\PYGZhy{}\PYGZhy{}\PYGZhy{}\PYGZhy{}\PYGZhy{}\PYGZhy{}\PYGZhy{}\PYGZhy{}\PYGZhy{}\PYGZhy{}\PYGZhy{}\PYGZhy{}\PYGZhy{}\PYGZhy{}\PYGZhy{}\PYGZhy{}\PYGZhy{}\PYGZhy{}\PYGZhy{}\PYGZhy{}\PYGZhy{}\PYGZhy{}\PYGZhy{}\PYGZhy{}\PYGZhy{}\PYGZhy{}\PYGZhy{}\PYGZhy{}\PYGZhy{}\PYGZhy{}\PYGZhy{}\PYGZhy{}\PYGZhy{}\PYGZhy{}\PYGZhy{}\PYGZhy{}\PYGZhy{}\PYGZhy{}\PYGZhy{}\PYGZhy{}\PYGZhy{}\PYGZhy{}\PYGZhy{}\PYGZhy{}\PYGZhy{}\PYGZhy{}\PYGZhy{}\PYGZhy{}\PYGZhy{}\PYGZhy{}\PYGZhy{}\PYGZhy{}\PYGZhy{}\PYGZhy{}\PYGZhy{}\PYGZhy{}\PYGZhy{}\PYGZhy{}\PYGZhy{}\PYGZhy{}\PYGZhy{}}
    \PYG{c+c1}{\PYGZsh{} optimizer loop over beta constraints}
    \PYG{n}{start\PYGZus{}point} \PYG{o}{=} \PYG{n}{random\PYGZus{}start}\PYG{p}{(}
        \PYG{n}{n\PYGZus{}random\PYGZus{}start}\PYG{p}{,}
        \PYG{n}{x\PYGZus{}lower}\PYG{p}{,}
        \PYG{n}{x\PYGZus{}upper}\PYG{p}{,}
        \PYG{n}{log\PYGZus{}scale}\PYG{p}{,}
        \PYG{n}{optimize\PYGZus{}fun}\PYG{o}{.}\PYG{n}{objective\PYGZus{}fun}
    \PYG{p}{)}\PYG{p}{;}
    \PYG{k}{if} \PYG{n}{debug\PYGZus{}output} \PYG{p}{:}
        \PYG{n+nb}{print}\PYG{p}{(} \PYG{l+s+s1}{\PYGZsq{}}\PYG{l+s+s1}{start\PYGZus{}point = }\PYG{l+s+s1}{\PYGZsq{}}\PYG{p}{,} \PYG{n}{start\PYGZus{}point} \PYG{p}{)}
        \PYG{n+nb}{print}\PYG{p}{(} \PYG{l+s+s1}{\PYGZsq{}}\PYG{l+s+s1}{objective   = }\PYG{l+s+s1}{\PYGZsq{}}\PYG{p}{,} \PYG{n}{optimize\PYGZus{}fun}\PYG{o}{.}\PYG{n}{objective\PYGZus{}fun}\PYG{p}{(}\PYG{n}{start\PYGZus{}point}\PYG{p}{)} \PYG{p}{)}
        \PYG{n+nb}{print}\PYG{p}{(} \PYG{l+s+s1}{\PYGZsq{}}\PYG{l+s+s1}{check       = }\PYG{l+s+s1}{\PYGZsq{}}\PYG{p}{,} \PYG{n}{objective}\PYG{p}{(}\PYG{n}{t\PYGZus{}all}\PYG{p}{,} \PYG{n}{D\PYGZus{}data}\PYG{p}{,} \PYG{n}{start\PYGZus{}point}\PYG{p}{)} \PYG{p}{)}
    \PYG{n}{constraints}  \PYG{o}{=} \PYG{n+nb}{list}\PYG{p}{(}\PYG{p}{)}
    \PYG{n}{result} \PYG{o}{=} \PYG{n}{scipy}\PYG{o}{.}\PYG{n}{optimize}\PYG{o}{.}\PYG{n}{minimize}\PYG{p}{(}
        \PYG{n}{optimize\PYGZus{}fun}\PYG{o}{.}\PYG{n}{objective\PYGZus{}fun}\PYG{p}{,}
        \PYG{n}{start\PYGZus{}point}\PYG{p}{,}
        \PYG{n}{method}        \PYG{o}{=} \PYG{l+s+s1}{\PYGZsq{}}\PYG{l+s+s1}{trust\PYGZhy{}constr}\PYG{l+s+s1}{\PYGZsq{}}\PYG{p}{,}
        \PYG{n}{jac}           \PYG{o}{=} \PYG{n}{optimize\PYGZus{}fun}\PYG{o}{.}\PYG{n}{objective\PYGZus{}grad}\PYG{p}{,}
        \PYG{n}{hess}          \PYG{o}{=} \PYG{n}{optimize\PYGZus{}fun}\PYG{o}{.}\PYG{n}{objective\PYGZus{}hess}\PYG{p}{,}
        \PYG{n}{constraints}   \PYG{o}{=} \PYG{n}{constraints}\PYG{p}{,}
        \PYG{n}{options}       \PYG{o}{=} \PYG{n}{options}\PYG{p}{,}
        \PYG{n}{bounds}        \PYG{o}{=} \PYG{n}{bounds}\PYG{p}{,}
    \PYG{p}{)}
    \PYG{n}{x\PYGZus{}fit} \PYG{o}{=} \PYG{n}{result}\PYG{o}{.}\PYG{n}{x}
    \PYG{k}{if} \PYG{n}{debug\PYGZus{}output} \PYG{p}{:}
        \PYG{n+nb}{print}\PYG{p}{(}\PYG{l+s+s1}{\PYGZsq{}}\PYG{l+s+s1}{optimizer success = }\PYG{l+s+s1}{\PYGZsq{}}\PYG{p}{,} \PYG{n}{result}\PYG{o}{.}\PYG{n}{success} \PYG{p}{)}
        \PYG{n+nb}{print}\PYG{p}{(}\PYG{l+s+s1}{\PYGZsq{}}\PYG{l+s+s1}{optimal objective = }\PYG{l+s+s1}{\PYGZsq{}}\PYG{p}{,} \PYG{n}{result}\PYG{o}{.}\PYG{n}{fun} \PYG{p}{)}
    \PYG{c+c1}{\PYGZsh{} check optimizer status}
    \PYG{n}{ok} \PYG{o}{=} \PYG{n}{ok} \PYG{o+ow}{and} \PYG{n}{result}\PYG{o}{.}\PYG{n}{success}
    \PYG{c+c1}{\PYGZsh{}}
    \PYG{c+c1}{\PYGZsh{} H: the observed infromation matrix}
    \PYG{n}{H} \PYG{o}{=} \PYG{n}{optimize\PYGZus{}fun}\PYG{o}{.}\PYG{n}{objective\PYGZus{}hess}\PYG{p}{(}\PYG{n}{x\PYGZus{}fit}\PYG{p}{)}
    \PYG{c+c1}{\PYGZsh{}}
    \PYG{c+c1}{\PYGZsh{} Hinv: approxiamtion for covariance of the estimate x\PYGZus{}fit}
    \PYG{n}{Hinv} \PYG{o}{=} \PYG{n}{numpy}\PYG{o}{.}\PYG{n}{linalg}\PYG{o}{.}\PYG{n}{inv}\PYG{p}{(}\PYG{n}{H}\PYG{p}{)}
    \PYG{c+c1}{\PYGZsh{}}
    \PYG{c+c1}{\PYGZsh{} standard devaition for each component of x\PYGZus{}fit}
    \PYG{n}{std\PYGZus{}error} \PYG{o}{=} \PYG{n}{numpy}\PYG{o}{.}\PYG{n}{sqrt}\PYG{p}{(} \PYG{n}{numpy}\PYG{o}{.}\PYG{n}{diag}\PYG{p}{(}\PYG{n}{Hinv}\PYG{p}{)} \PYG{p}{)}
    \PYG{c+c1}{\PYGZsh{}}
    \PYG{c+c1}{\PYGZsh{} seirwd\PYGZus{}all\PYGZus{}fit}
    \PYG{n}{seirwd\PYGZus{}all\PYGZus{}fit} \PYG{o}{=} \PYG{n}{x2seirwd\PYGZus{}all}\PYG{p}{(}\PYG{n}{x\PYGZus{}fit}\PYG{p}{)}
    \PYG{c+c1}{\PYGZsh{}}
    \PYG{c+c1}{\PYGZsh{} D\PYGZus{}fit}
    \PYG{n}{D\PYGZus{}fit} \PYG{o}{=} \PYG{n}{seirwd\PYGZus{}all\PYGZus{}fit}\PYG{p}{[}\PYG{p}{:}\PYG{p}{,}\PYG{l+m+mi}{5}\PYG{p}{]}
    \PYG{c+c1}{\PYGZsh{}}
    \PYG{c+c1}{\PYGZsh{} data\PYGZus{}resdidual}
    \PYG{n}{data\PYGZus{}residual} \PYG{o}{=} \PYG{n}{weighted\PYGZus{}data\PYGZus{}residual}\PYG{p}{(}\PYG{n}{D\PYGZus{}data}\PYG{p}{,} \PYG{n}{D\PYGZus{}fit}\PYG{p}{)}
    \PYG{c+c1}{\PYGZsh{}}
    \PYG{c+c1}{\PYGZsh{} display\PYGZus{}fit\PYGZus{}results}
    \PYG{k}{if} \PYG{n}{display\PYGZus{}fit} \PYG{p}{:}
        \PYG{n}{display\PYGZus{}fit\PYGZus{}results}\PYG{p}{(}\PYG{n}{D\PYGZus{}data}\PYG{p}{,} \PYG{n}{x\PYGZus{}fit}\PYG{p}{,} \PYG{n}{x\PYGZus{}lower}\PYG{p}{,} \PYG{n}{x\PYGZus{}upper}\PYG{p}{,} \PYG{n}{std\PYGZus{}error}\PYG{p}{)}
    \PYG{c+c1}{\PYGZsh{}}
    \PYG{c+c1}{\PYGZsh{} check that all the data residuals an less than 3.0}
    \PYG{k}{if} \PYG{n}{numpy}\PYG{o}{.}\PYG{n}{any}\PYG{p}{(} \PYG{n}{numpy}\PYG{o}{.}\PYG{n}{abs}\PYG{p}{(} \PYG{n}{data\PYGZus{}residual} \PYG{p}{)} \PYG{o}{\PYGZgt{}}\PYG{o}{=} \PYG{l+m+mf}{4.0} \PYG{p}{)} \PYG{p}{:}
        \PYG{n+nb}{print}\PYG{p}{(}\PYG{l+s+s1}{\PYGZsq{}}\PYG{l+s+s1}{covid\PYGZus{}19\PYGZus{}xam: a weighted data residual \PYGZgt{}= 4.0}\PYG{l+s+s1}{\PYGZsq{}}\PYG{p}{)}
        \PYG{n}{ok} \PYG{o}{=} \PYG{k+kc}{False}
    \PYG{c+c1}{\PYGZsh{}}
    \PYG{c+c1}{\PYGZsh{} compare fit to simulation truth}
    \PYG{k}{if} \PYG{n}{data\PYGZus{}file} \PYG{o}{==} \PYG{l+s+s1}{\PYGZsq{}}\PYG{l+s+s1}{\PYGZsq{}} \PYG{p}{:}
        \PYG{n}{x\PYGZus{}residual} \PYG{o}{=} \PYG{n}{x\PYGZus{}fit} \PYG{o}{/} \PYG{n}{x\PYGZus{}sim} \PYG{o}{\PYGZhy{}} \PYG{l+m+mf}{1.0}
        \PYG{n}{x\PYGZus{}residual}\PYG{p}{[} \PYG{n}{x\PYGZus{}sim} \PYG{o}{==} \PYG{l+m+mf}{0.0} \PYG{p}{]} \PYG{o}{=} \PYG{l+m+mf}{0.0}
        \PYG{k}{if} \PYG{n}{death\PYGZus{}data\PYGZus{}cv} \PYG{o}{\PYGZgt{}} \PYG{l+m+mf}{0.0} \PYG{p}{:}
            \PYG{k}{if} \PYG{n}{numpy}\PYG{o}{.}\PYG{n}{any}\PYG{p}{(} \PYG{n+nb}{abs}\PYG{p}{(}\PYG{n}{x\PYGZus{}residual}\PYG{p}{)} \PYG{o}{\PYGZgt{}}\PYG{o}{=} \PYG{l+m+mf}{2.0}\PYG{p}{)} \PYG{p}{:}
                \PYG{n+nb}{print}\PYG{p}{(}\PYG{l+s+s1}{\PYGZsq{}}\PYG{l+s+s1}{covid\PYGZus{}19\PYGZus{}xam: a weight x residual \PYGZgt{}= 2.0}\PYG{l+s+s1}{\PYGZsq{}}\PYG{p}{)}
                \PYG{n}{ok} \PYG{o}{=} \PYG{k+kc}{False}
        \PYG{k}{else}\PYG{p}{:}
            \PYG{k}{if} \PYG{n}{numpy}\PYG{o}{.}\PYG{n}{any}\PYG{p}{(} \PYG{n+nb}{abs}\PYG{p}{(}\PYG{n}{x\PYGZus{}residual}\PYG{p}{)} \PYG{o}{\PYGZgt{}}\PYG{o}{=} \PYG{l+m+mf}{1e\PYGZhy{}5} \PYG{p}{)} \PYG{p}{:}
                \PYG{n+nb}{print}\PYG{p}{(}\PYG{l+s+s1}{\PYGZsq{}}\PYG{l+s+s1}{covid\PYGZus{}19\PYGZus{}xam: an x residual \PYGZgt{}= 1e\PYGZhy{}5}\PYG{l+s+s1}{\PYGZsq{}}\PYG{p}{)}
                \PYG{n+nb}{print}\PYG{p}{(}\PYG{l+s+s1}{\PYGZsq{}}\PYG{l+s+s1}{should have prefect fit because death\PYGZus{}data\PYGZus{}cv = 0}\PYG{l+s+s1}{\PYGZsq{}}\PYG{p}{)}
                \PYG{n}{ok} \PYG{o}{=} \PYG{k+kc}{False}
        \PYG{c+c1}{\PYGZsh{}}
        \PYG{k}{if} \PYG{o+ow}{not} \PYG{n}{ok} \PYG{p}{:}
            \PYG{n}{msg}  \PYG{o}{=} \PYG{l+s+s1}{\PYGZsq{}}\PYG{l+s+s1}{covid\PYGZus{}19\PYGZus{}xam: Correctness test failed, }\PYG{l+s+s1}{\PYGZsq{}}
            \PYG{n}{msg} \PYG{o}{+}\PYG{o}{=} \PYG{l+s+s1}{\PYGZsq{}}\PYG{l+s+s1}{actual random seed = }\PYG{l+s+s1}{\PYGZsq{}} \PYG{o}{+} \PYG{n+nb}{str}\PYG{p}{(}\PYG{n}{actual\PYGZus{}seed}\PYG{p}{[}\PYG{l+m+mi}{0}\PYG{p}{]}\PYG{p}{)}
            \PYG{n+nb}{print}\PYG{p}{(} \PYG{n}{msg} \PYG{p}{)}
            \PYG{n}{call\PYGZus{}count} \PYG{o}{+}\PYG{o}{=} \PYG{l+m+mi}{1}
            \PYG{k}{if} \PYG{n}{call\PYGZus{}count} \PYG{o}{\PYGZlt{}} \PYG{l+m+mi}{3} \PYG{o+ow}{and} \PYG{n}{random\PYGZus{}seed} \PYG{o}{==} \PYG{l+m+mi}{0} \PYG{p}{:}
                \PYG{n+nb}{print}\PYG{p}{(} \PYG{l+s+s1}{\PYGZsq{}}\PYG{l+s+s1}{re\PYGZhy{}trying with a differenent random seed}\PYG{l+s+s1}{\PYGZsq{}}\PYG{p}{)}
                \PYG{n}{ok} \PYG{o}{=} \PYG{n}{covid\PYGZus{}19\PYGZus{}xam}\PYG{p}{(}\PYG{n}{call\PYGZus{}count}\PYG{p}{)}
    \PYG{c+c1}{\PYGZsh{}}
    \PYG{c+c1}{\PYGZsh{} check conservation of mass in the compartmental model}
    \PYG{n}{sum\PYGZus{}all\PYGZus{}fit} \PYG{o}{=} \PYG{n}{numpy}\PYG{o}{.}\PYG{n}{sum}\PYG{p}{(}\PYG{n}{seirwd\PYGZus{}all\PYGZus{}fit}\PYG{p}{,} \PYG{n}{axis}\PYG{o}{=}\PYG{l+m+mi}{1}\PYG{p}{)}
    \PYG{n}{eps99}       \PYG{o}{=} \PYG{l+m+mf}{99.0} \PYG{o}{*} \PYG{n}{numpy}\PYG{o}{.}\PYG{n}{finfo}\PYG{p}{(}\PYG{n+nb}{float}\PYG{p}{)}\PYG{o}{.}\PYG{n}{eps}
    \PYG{k}{if} \PYG{n}{numpy}\PYG{o}{.}\PYG{n}{any}\PYG{p}{(} \PYG{n+nb}{abs}\PYG{p}{(}\PYG{n}{sum\PYGZus{}all\PYGZus{}fit} \PYG{o}{\PYGZhy{}} \PYG{l+m+mf}{1.0}\PYG{p}{)} \PYG{o}{\PYGZgt{}} \PYG{n}{eps99} \PYG{p}{)} \PYG{p}{:}
        \PYG{n+nb}{print}\PYG{p}{(}\PYG{l+s+s1}{\PYGZsq{}}\PYG{l+s+s1}{covid\PYGZus{}19\PYGZus{}xam: sum of all compartments not equal 1.0}\PYG{l+s+s1}{\PYGZsq{}}\PYG{p}{)}
        \PYG{n}{ok} \PYG{o}{=} \PYG{k+kc}{False}
    \PYG{c+c1}{\PYGZsh{}}
    \PYG{k}{return} \PYG{n}{ok}
\end{sphinxVerbatim}


\bigskip\hrule\bigskip


xsrst input file: \sphinxcode{\sphinxupquote{example/python/numeric/covid\_19\_xam.py}}
\begin{enumerate}
\sphinxsetlistlabels{\arabic}{enumi}{enumii}{}{.}%
\item {} 
numeric\_simple\_inv: {\hyperref[\detokenize{xsrst/numeric_simple_inv:id1}]{\sphinxcrossref{\DUrole{std,std-ref}{2.1 An AD Compatible Matrix Inverse Routine}}}}

\item {} 
numeric\_runge4\_step: {\hyperref[\detokenize{xsrst/numeric_runge4_step:id1}]{\sphinxcrossref{\DUrole{std,std-ref}{2.2 One Fourth Order Runge\sphinxhyphen{}Kutta ODE Step}}}}

\item {} 
numeric\_rosen3\_step: {\hyperref[\detokenize{xsrst/numeric_rosen3_step:id1}]{\sphinxcrossref{\DUrole{std,std-ref}{2.3 One Third Order Rosenbrock ODE Step}}}}

\item {} 
numeric\_ode\_multi\_step: {\hyperref[\detokenize{xsrst/numeric_ode_multi_step:id1}]{\sphinxcrossref{\DUrole{std,std-ref}{2.4 Multiple Ode Steps}}}}

\item {} 
numeric\_optimize\_fun\_class: {\hyperref[\detokenize{xsrst/numeric_optimize_fun_class:id1}]{\sphinxcrossref{\DUrole{std,std-ref}{2.5 A Helper Class That Defines Functions Needed for Optimization}}}}

\item {} 
numeric\_seirwd\_model: {\hyperref[\detokenize{xsrst/numeric_seirwd_model:id1}]{\sphinxcrossref{\DUrole{std,std-ref}{2.6 A Susceptible Exposed Infectious Recovered and Death Model}}}}

\item {} 
numeric\_covid\_19\_xam\_py: {\hyperref[\detokenize{xsrst/numeric_covid_19_xam_py:id1}]{\sphinxcrossref{\DUrole{std,std-ref}{2.7 Example Fitting an SEIRWD Model for Covid\sphinxhyphen{}19}}}}

\end{enumerate}


\bigskip\hrule\bigskip


xsrst input file: \sphinxcode{\sphinxupquote{example/python/numeric/numeric.xsrst}}


\subsubsection{library}
\label{\detokenize{xsrst/library:library}}\label{\detokenize{xsrst/library::doc}}
\index{library@\spxentry{library}}\index{the@\spxentry{the}}\index{cppad@\spxentry{cppad}}\index{py@\spxentry{py}}\index{libraries@\spxentry{libraries}}\ignorespaces 

\paragraph{3 The Cppad Py Libraries}
\label{\detokenize{xsrst/library:the-cppad-py-libraries}}\label{\detokenize{xsrst/library:id1}}\begin{itemize}
\item {} 
{\hyperref[\detokenize{xsrst/library:library-children}]{\sphinxcrossref{\DUrole{std,std-ref}{Children}}}}

\end{itemize}

\index{children@\spxentry{children}}\ignorespaces 

\subparagraph{Children}
\label{\detokenize{xsrst/library:children}}\label{\detokenize{xsrst/library:library-children}}\label{\detokenize{xsrst/library:index-1}}

\subparagraph{py\_lib}
\label{\detokenize{xsrst/py_lib:py-lib}}\label{\detokenize{xsrst/py_lib::doc}}
\index{py\_lib@\spxentry{py\_lib}}\index{the@\spxentry{the}}\index{python@\spxentry{python}}\index{library@\spxentry{library}}\ignorespaces 

\subparagraph{3.1 The Python Library}
\label{\detokenize{xsrst/py_lib:the-python-library}}\label{\detokenize{xsrst/py_lib:id1}}\begin{itemize}
\item {} 
{\hyperref[\detokenize{xsrst/py_lib:py-lib-children}]{\sphinxcrossref{\DUrole{std,std-ref}{Children}}}}

\end{itemize}

\index{children@\spxentry{children}}\ignorespaces 

\subparagraph{Children}
\label{\detokenize{xsrst/py_lib:children}}\label{\detokenize{xsrst/py_lib:py-lib-children}}\label{\detokenize{xsrst/py_lib:index-1}}

\subparagraph{py\_fun}
\label{\detokenize{xsrst/py_fun:py-fun}}\label{\detokenize{xsrst/py_fun::doc}}
\index{py\_fun@\spxentry{py\_fun}}\index{cppad@\spxentry{cppad}}\index{py@\spxentry{py}}\index{ad@\spxentry{ad}}\index{functions@\spxentry{functions}}\ignorespaces 

\subparagraph{3.1.1 Cppad Py AD Functions}
\label{\detokenize{xsrst/py_fun:cppad-py-ad-functions}}\label{\detokenize{xsrst/py_fun:id1}}\begin{itemize}
\item {} 
{\hyperref[\detokenize{xsrst/py_fun:py-fun-children}]{\sphinxcrossref{\DUrole{std,std-ref}{Children}}}}

\end{itemize}

\index{children@\spxentry{children}}\ignorespaces 

\subparagraph{Children}
\label{\detokenize{xsrst/py_fun:children}}\label{\detokenize{xsrst/py_fun:py-fun-children}}\label{\detokenize{xsrst/py_fun:index-1}}

\subparagraph{py\_independent}
\label{\detokenize{xsrst/py_independent:py-independent}}\label{\detokenize{xsrst/py_independent::doc}}
\index{py\_independent@\spxentry{py\_independent}}\index{declare@\spxentry{declare}}\index{independent@\spxentry{independent}}\index{variables@\spxentry{variables}}\index{start@\spxentry{start}}\index{recording@\spxentry{recording}}\ignorespaces 

\subparagraph{3.1.1.1 Declare Independent Variables and Start Recording}
\label{\detokenize{xsrst/py_independent:declare-independent-variables-and-start-recording}}\label{\detokenize{xsrst/py_independent:id1}}\begin{itemize}
\item {} 
{\hyperref[\detokenize{xsrst/py_independent:py-independent-syntax}]{\sphinxcrossref{\DUrole{std,std-ref}{Syntax}}}}

\item {} 
{\hyperref[\detokenize{xsrst/py_independent:py-independent-x}]{\sphinxcrossref{\DUrole{std,std-ref}{x}}}}

\item {} 
{\hyperref[\detokenize{xsrst/py_independent:py-independent-dynamic}]{\sphinxcrossref{\DUrole{std,std-ref}{dynamic}}}}

\item {} 
{\hyperref[\detokenize{xsrst/py_independent:py-independent-ax}]{\sphinxcrossref{\DUrole{std,std-ref}{ax}}}}

\item {} 
{\hyperref[\detokenize{xsrst/py_independent:py-independent-adynamic}]{\sphinxcrossref{\DUrole{std,std-ref}{adynamic}}}}

\item {} 
{\hyperref[\detokenize{xsrst/py_independent:py-independent-purpose}]{\sphinxcrossref{\DUrole{std,std-ref}{Purpose}}}}

\item {} 
{\hyperref[\detokenize{xsrst/py_independent:py-independent-example}]{\sphinxcrossref{\DUrole{std,std-ref}{Example}}}}

\end{itemize}

\index{syntax@\spxentry{syntax}}\ignorespaces 

\subparagraph{Syntax}
\label{\detokenize{xsrst/py_independent:syntax}}\label{\detokenize{xsrst/py_independent:py-independent-syntax}}\label{\detokenize{xsrst/py_independent:index-1}}
\begin{DUlineblock}{0em}
\item[] \sphinxstyleemphasis{ax} =  \sphinxcode{\sphinxupquote{cppad\_py.independent}} ( \sphinxstyleemphasis{x} )
\item[] ( \sphinxstyleemphasis{ax} , \sphinxstyleemphasis{adynamic} ) =  \sphinxcode{\sphinxupquote{cppad\_py.independent}} ( \sphinxstyleemphasis{x} , \sphinxstyleemphasis{dynamic} )
\end{DUlineblock}

\index{x@\spxentry{x}}\ignorespaces 

\subparagraph{x}
\label{\detokenize{xsrst/py_independent:x}}\label{\detokenize{xsrst/py_independent:py-independent-x}}\label{\detokenize{xsrst/py_independent:index-2}}
This argument is a numpy vector with \sphinxcode{\sphinxupquote{float}} elements.
It specifies the number of independent variables
and their values during the recording.
We use \sphinxstyleemphasis{nx} = \sphinxstyleemphasis{x} . \sphinxcode{\sphinxupquote{size}}
to denote the number of independent variables.

\index{dynamic@\spxentry{dynamic}}\ignorespaces 

\subparagraph{dynamic}
\label{\detokenize{xsrst/py_independent:dynamic}}\label{\detokenize{xsrst/py_independent:py-independent-dynamic}}\label{\detokenize{xsrst/py_independent:index-3}}
This argument is a numpy vector with \sphinxcode{\sphinxupquote{float}} elements.
It specifies the number of independent dynamic parameters
and their values during the recording.
We use \sphinxstyleemphasis{nd} = \sphinxstyleemphasis{dynamic} . \sphinxcode{\sphinxupquote{size}}
to denote the number of independent dynamic parameters.

\index{ax@\spxentry{ax}}\ignorespaces 

\subparagraph{ax}
\label{\detokenize{xsrst/py_independent:ax}}\label{\detokenize{xsrst/py_independent:py-independent-ax}}\label{\detokenize{xsrst/py_independent:index-4}}
This result is a numpy vector with \sphinxcode{\sphinxupquote{a\_double}} elements.
This is the vector of independent variables.
It has size \sphinxstyleemphasis{nx} and for
\sphinxstyleemphasis{i} = 0 to \sphinxstyleemphasis{n} \sphinxhyphen{}1

\begin{DUlineblock}{0em}
\item[]         \sphinxstyleemphasis{ax} {[} \sphinxstyleemphasis{i} {]}. \sphinxcode{\sphinxupquote{value}} () == \sphinxstyleemphasis{x} {[} \sphinxstyleemphasis{i} {]}
\end{DUlineblock}

\index{adynamic@\spxentry{adynamic}}\ignorespaces 

\subparagraph{adynamic}
\label{\detokenize{xsrst/py_independent:adynamic}}\label{\detokenize{xsrst/py_independent:py-independent-adynamic}}\label{\detokenize{xsrst/py_independent:index-5}}
This result is a numpy vector with \sphinxcode{\sphinxupquote{a\_double}} elements.
This is the vector of independent dynamic parameters.
It has size \sphinxstyleemphasis{nd} and for
\sphinxstyleemphasis{i} = 0 to \sphinxstyleemphasis{n} \sphinxhyphen{}1

\begin{DUlineblock}{0em}
\item[]         \sphinxstyleemphasis{adynamic} {[} \sphinxstyleemphasis{i} {]}. \sphinxcode{\sphinxupquote{value}} () == \sphinxstyleemphasis{dynamic} {[} \sphinxstyleemphasis{i} {]}
\end{DUlineblock}

\index{purpose@\spxentry{purpose}}\ignorespaces 

\subparagraph{Purpose}
\label{\detokenize{xsrst/py_independent:purpose}}\label{\detokenize{xsrst/py_independent:py-independent-purpose}}\label{\detokenize{xsrst/py_independent:index-6}}
This starts a recording of the {\hyperref[\detokenize{xsrst/a_double:id1}]{\sphinxcrossref{\DUrole{std,std-ref}{a\_double}}}} operations.
This recording is terminated, and the information is stored,
by calling the {\hyperref[\detokenize{xsrst/py_fun_ctor:id1}]{\sphinxcrossref{\DUrole{std,std-ref}{d\_fun\_constructor}}}}.
It is terminated, and the information is lost,
by calling {\hyperref[\detokenize{xsrst/py_abort_recording:id1}]{\sphinxcrossref{\DUrole{std,std-ref}{abort\_recording}}}}.


\subparagraph{fun\_dynamic\_xam\_py}
\label{\detokenize{xsrst/fun_dynamic_xam_py:fun-dynamic-xam-py}}\label{\detokenize{xsrst/fun_dynamic_xam_py::doc}}
\index{fun\_dynamic\_xam\_py@\spxentry{fun\_dynamic\_xam\_py}}\index{python:@\spxentry{python:}}\index{using@\spxentry{using}}\index{dynamic@\spxentry{dynamic}}\index{parameters:@\spxentry{parameters:}}\index{example@\spxentry{example}}\index{test@\spxentry{test}}\ignorespaces 

\subparagraph{3.1.1.1.1 Python: Using Dynamic Parameters: Example and Test}
\label{\detokenize{xsrst/fun_dynamic_xam_py:python-using-dynamic-parameters-example-and-test}}\label{\detokenize{xsrst/fun_dynamic_xam_py:id1}}
\begin{sphinxVerbatim}[commandchars=\\\{\}]
\PYG{k+kn}{import} \PYG{n+nn}{numpy}
\PYG{k+kn}{import} \PYG{n+nn}{cppad\PYGZus{}py}
\PYG{k}{def} \PYG{n+nf}{fun\PYGZus{}dynamic\PYGZus{}xam}\PYG{p}{(}\PYG{p}{)} \PYG{p}{:}
    \PYG{n}{ok}  \PYG{o}{=} \PYG{k+kc}{True}
    \PYG{n}{nx}  \PYG{o}{=} \PYG{l+m+mi}{2}
    \PYG{n}{nd}  \PYG{o}{=} \PYG{l+m+mi}{2}
    \PYG{c+c1}{\PYGZsh{}}
    \PYG{c+c1}{\PYGZsh{} value of independent variables during recording}
    \PYG{n}{x} \PYG{o}{=} \PYG{n}{numpy}\PYG{o}{.}\PYG{n}{empty}\PYG{p}{(}\PYG{n}{nx}\PYG{p}{,} \PYG{n}{dtype}\PYG{o}{=}\PYG{n+nb}{float}\PYG{p}{)}
    \PYG{n}{x}\PYG{p}{[}\PYG{l+m+mi}{0}\PYG{p}{]} \PYG{o}{=} \PYG{l+m+mf}{1.0}
    \PYG{n}{x}\PYG{p}{[}\PYG{l+m+mi}{1}\PYG{p}{]} \PYG{o}{=} \PYG{l+m+mf}{2.0}
    \PYG{c+c1}{\PYGZsh{}}
    \PYG{c+c1}{\PYGZsh{} value of independent dynamic parameters during recording}
    \PYG{n}{dynamic} \PYG{o}{=} \PYG{n}{numpy}\PYG{o}{.}\PYG{n}{empty}\PYG{p}{(}\PYG{n}{nd}\PYG{p}{,} \PYG{n}{dtype}\PYG{o}{=}\PYG{n+nb}{float}\PYG{p}{)}
    \PYG{n}{dynamic}\PYG{p}{[}\PYG{l+m+mi}{0}\PYG{p}{]} \PYG{o}{=} \PYG{l+m+mf}{3.0}
    \PYG{n}{dynamic}\PYG{p}{[}\PYG{l+m+mi}{1}\PYG{p}{]} \PYG{o}{=} \PYG{l+m+mf}{4.0}
    \PYG{c+c1}{\PYGZsh{}}
    \PYG{c+c1}{\PYGZsh{} start recording}
    \PYG{p}{(}\PYG{n}{ax}\PYG{p}{,} \PYG{n}{adynamic}\PYG{p}{)} \PYG{o}{=} \PYG{n}{cppad\PYGZus{}py}\PYG{o}{.}\PYG{n}{independent}\PYG{p}{(}\PYG{n}{x}\PYG{p}{,} \PYG{n}{dynamic}\PYG{p}{)}
    \PYG{c+c1}{\PYGZsh{}}
    \PYG{c+c1}{\PYGZsh{} create another dynamic paramerer}
    \PYG{n}{adyn}  \PYG{o}{=} \PYG{n}{adynamic}\PYG{p}{[}\PYG{l+m+mi}{0}\PYG{p}{]} \PYG{o}{+} \PYG{n}{adynamic}\PYG{p}{[}\PYG{l+m+mi}{1}\PYG{p}{]}
    \PYG{c+c1}{\PYGZsh{}}
    \PYG{c+c1}{\PYGZsh{} create another variable}
    \PYG{n}{avar}  \PYG{o}{=} \PYG{n}{ax}\PYG{p}{[}\PYG{l+m+mi}{0}\PYG{p}{]} \PYG{o}{+} \PYG{n}{ax}\PYG{p}{[}\PYG{l+m+mi}{1}\PYG{p}{]} \PYG{o}{+} \PYG{n}{adyn}
    \PYG{c+c1}{\PYGZsh{}}
    \PYG{c+c1}{\PYGZsh{} create f(x) = x[0] + x[1] + dynamic[0] + dynamic[1]}
    \PYG{n}{ay}    \PYG{o}{=} \PYG{n}{numpy}\PYG{o}{.}\PYG{n}{empty}\PYG{p}{(}\PYG{l+m+mi}{1}\PYG{p}{,} \PYG{n}{dtype}\PYG{o}{=}\PYG{n}{cppad\PYGZus{}py}\PYG{o}{.}\PYG{n}{a\PYGZus{}double}\PYG{p}{)}
    \PYG{n}{ay}\PYG{p}{[}\PYG{l+m+mi}{0}\PYG{p}{]} \PYG{o}{=} \PYG{n}{avar}
    \PYG{n}{f}     \PYG{o}{=} \PYG{n}{cppad\PYGZus{}py}\PYG{o}{.}\PYG{n}{d\PYGZus{}fun}\PYG{p}{(}\PYG{n}{ax}\PYG{p}{,} \PYG{n}{ay}\PYG{p}{)}
    \PYG{c+c1}{\PYGZsh{}}
    \PYG{c+c1}{\PYGZsh{} check some properties of f}
    \PYG{n}{ok} \PYG{o}{=} \PYG{n}{ok} \PYG{o+ow}{and} \PYG{n}{f}\PYG{o}{.}\PYG{n}{size\PYGZus{}domain}\PYG{p}{(}\PYG{p}{)} \PYG{o}{==} \PYG{n}{nx}
    \PYG{n}{ok} \PYG{o}{=} \PYG{n}{ok} \PYG{o+ow}{and} \PYG{n}{f}\PYG{o}{.}\PYG{n}{size\PYGZus{}order}\PYG{p}{(}\PYG{p}{)}  \PYG{o}{==} \PYG{l+m+mi}{0}
    \PYG{c+c1}{\PYGZsh{}}
    \PYG{c+c1}{\PYGZsh{} zero order forward mode using same values as during the recording}
    \PYG{n}{y}  \PYG{o}{=} \PYG{n}{f}\PYG{o}{.}\PYG{n}{forward}\PYG{p}{(}\PYG{l+m+mi}{0}\PYG{p}{,} \PYG{n}{x}\PYG{p}{)}
    \PYG{n}{ok} \PYG{o}{=} \PYG{n}{ok} \PYG{o+ow}{and} \PYG{n}{y}\PYG{p}{[}\PYG{l+m+mi}{0}\PYG{p}{]} \PYG{o}{==} \PYG{p}{(}\PYG{n}{x}\PYG{p}{[}\PYG{l+m+mi}{0}\PYG{p}{]} \PYG{o}{+} \PYG{n}{x}\PYG{p}{[}\PYG{l+m+mi}{1}\PYG{p}{]} \PYG{o}{+} \PYG{n}{dynamic}\PYG{p}{[}\PYG{l+m+mi}{0}\PYG{p}{]} \PYG{o}{+} \PYG{n}{dynamic}\PYG{p}{[}\PYG{l+m+mi}{1}\PYG{p}{]}\PYG{p}{)}
    \PYG{c+c1}{\PYGZsh{}}
    \PYG{c+c1}{\PYGZsh{}  zero order forward mode using different value for dynamic parameters}
    \PYG{n}{dynamic}\PYG{p}{[}\PYG{l+m+mi}{0}\PYG{p}{]} \PYG{o}{=} \PYG{n}{dynamic}\PYG{p}{[}\PYG{l+m+mi}{0}\PYG{p}{]} \PYG{o}{+} \PYG{l+m+mf}{1.0}
    \PYG{n}{dynamic}\PYG{p}{[}\PYG{l+m+mi}{1}\PYG{p}{]} \PYG{o}{=} \PYG{n}{dynamic}\PYG{p}{[}\PYG{l+m+mi}{1}\PYG{p}{]} \PYG{o}{+} \PYG{l+m+mf}{1.0}
    \PYG{n}{f}\PYG{o}{.}\PYG{n}{new\PYGZus{}dynamic}\PYG{p}{(}\PYG{n}{dynamic}\PYG{p}{)}
    \PYG{n}{y}   \PYG{o}{=} \PYG{n}{f}\PYG{o}{.}\PYG{n}{forward}\PYG{p}{(}\PYG{l+m+mi}{0}\PYG{p}{,} \PYG{n}{x}\PYG{p}{)}
    \PYG{n}{ok}  \PYG{o}{=} \PYG{n}{ok} \PYG{o+ow}{and} \PYG{n}{y}\PYG{p}{[}\PYG{l+m+mi}{0}\PYG{p}{]} \PYG{o}{==} \PYG{p}{(}\PYG{n}{x}\PYG{p}{[}\PYG{l+m+mi}{0}\PYG{p}{]} \PYG{o}{+} \PYG{n}{x}\PYG{p}{[}\PYG{l+m+mi}{1}\PYG{p}{]} \PYG{o}{+} \PYG{n}{dynamic}\PYG{p}{[}\PYG{l+m+mi}{0}\PYG{p}{]} \PYG{o}{+} \PYG{n}{dynamic}\PYG{p}{[}\PYG{l+m+mi}{1}\PYG{p}{]}\PYG{p}{)}
    \PYG{c+c1}{\PYGZsh{}}
    \PYG{k}{return} \PYG{n}{ok}
\end{sphinxVerbatim}


\bigskip\hrule\bigskip


xsrst input file: \sphinxcode{\sphinxupquote{example/python/core/fun\_dynamic\_xam.py}}

\index{example@\spxentry{example}}\ignorespaces 

\subparagraph{Example}
\label{\detokenize{xsrst/py_independent:example}}\label{\detokenize{xsrst/py_independent:py-independent-example}}\label{\detokenize{xsrst/py_independent:index-7}}
Most of the python \sphinxcode{\sphinxupquote{d\_fun}} examples use this function.
The {\hyperref[\detokenize{xsrst/fun_dynamic_xam_py:id1}]{\sphinxcrossref{\DUrole{std,std-ref}{fun\_dynamic\_xam\_py}}}} uses the syntax that includes
dynamic parameters.


\bigskip\hrule\bigskip


xsrst input file: \sphinxcode{\sphinxupquote{lib/python/cppad\_py/independent.py}}


\subparagraph{py\_abort\_recording}
\label{\detokenize{xsrst/py_abort_recording:py-abort-recording}}\label{\detokenize{xsrst/py_abort_recording::doc}}
\index{py\_abort\_recording@\spxentry{py\_abort\_recording}}\index{abort@\spxentry{abort}}\index{recording@\spxentry{recording}}\ignorespaces 

\subparagraph{3.1.1.2 Abort Recording}
\label{\detokenize{xsrst/py_abort_recording:abort-recording}}\label{\detokenize{xsrst/py_abort_recording:id1}}\begin{itemize}
\item {} 
{\hyperref[\detokenize{xsrst/py_abort_recording:py-abort-recording-syntax}]{\sphinxcrossref{\DUrole{std,std-ref}{Syntax}}}}

\item {} 
{\hyperref[\detokenize{xsrst/py_abort_recording:py-abort-recording-purpose}]{\sphinxcrossref{\DUrole{std,std-ref}{Purpose}}}}

\item {} 
{\hyperref[\detokenize{xsrst/py_abort_recording:py-abort-recording-example}]{\sphinxcrossref{\DUrole{std,std-ref}{Example}}}}

\end{itemize}

\index{syntax@\spxentry{syntax}}\ignorespaces 

\subparagraph{Syntax}
\label{\detokenize{xsrst/py_abort_recording:syntax}}\label{\detokenize{xsrst/py_abort_recording:py-abort-recording-syntax}}\label{\detokenize{xsrst/py_abort_recording:index-1}}
\sphinxcode{\sphinxupquote{cppad\_py.abort\_recording}} ()

\index{purpose@\spxentry{purpose}}\ignorespaces 

\subparagraph{Purpose}
\label{\detokenize{xsrst/py_abort_recording:purpose}}\label{\detokenize{xsrst/py_abort_recording:py-abort-recording-purpose}}\label{\detokenize{xsrst/py_abort_recording:index-2}}
This aborts the current recording (if it exists)
started by the most recent call to {\hyperref[\detokenize{xsrst/py_independent:id1}]{\sphinxcrossref{\DUrole{std,std-ref}{independent}}}}.


\subparagraph{fun\_abort\_xam\_py}
\label{\detokenize{xsrst/fun_abort_xam_py:fun-abort-xam-py}}\label{\detokenize{xsrst/fun_abort_xam_py::doc}}
\index{fun\_abort\_xam\_py@\spxentry{fun\_abort\_xam\_py}}\index{python:@\spxentry{python:}}\index{abort@\spxentry{abort}}\index{recording@\spxentry{recording}}\index{a\_double@\spxentry{a\_double}}\index{operations:@\spxentry{operations:}}\index{example@\spxentry{example}}\index{test@\spxentry{test}}\ignorespaces 

\subparagraph{3.1.1.2.1 Python: Abort Recording a\_double Operations: Example and Test}
\label{\detokenize{xsrst/fun_abort_xam_py:python-abort-recording-a-double-operations-example-and-test}}\label{\detokenize{xsrst/fun_abort_xam_py:id1}}
\begin{sphinxVerbatim}[commandchars=\\\{\}]
\PYG{k}{def} \PYG{n+nf}{fun\PYGZus{}abort\PYGZus{}xam}\PYG{p}{(}\PYG{p}{)} \PYG{p}{:}
    \PYG{c+c1}{\PYGZsh{}}
    \PYG{k+kn}{import} \PYG{n+nn}{numpy}
    \PYG{k+kn}{import} \PYG{n+nn}{cppad\PYGZus{}py}
    \PYG{c+c1}{\PYGZsh{}}
    \PYG{c+c1}{\PYGZsh{} initialize return variable}
    \PYG{n}{ok} \PYG{o}{=} \PYG{k+kc}{True}
    \PYG{c+c1}{\PYGZsh{} \PYGZhy{}\PYGZhy{}\PYGZhy{}\PYGZhy{}\PYGZhy{}\PYGZhy{}\PYGZhy{}\PYGZhy{}\PYGZhy{}\PYGZhy{}\PYGZhy{}\PYGZhy{}\PYGZhy{}\PYGZhy{}\PYGZhy{}\PYGZhy{}\PYGZhy{}\PYGZhy{}\PYGZhy{}\PYGZhy{}\PYGZhy{}\PYGZhy{}\PYGZhy{}\PYGZhy{}\PYGZhy{}\PYGZhy{}\PYGZhy{}\PYGZhy{}\PYGZhy{}\PYGZhy{}\PYGZhy{}\PYGZhy{}\PYGZhy{}\PYGZhy{}\PYGZhy{}\PYGZhy{}\PYGZhy{}\PYGZhy{}\PYGZhy{}\PYGZhy{}\PYGZhy{}\PYGZhy{}\PYGZhy{}\PYGZhy{}\PYGZhy{}\PYGZhy{}\PYGZhy{}\PYGZhy{}\PYGZhy{}\PYGZhy{}\PYGZhy{}\PYGZhy{}\PYGZhy{}\PYGZhy{}\PYGZhy{}\PYGZhy{}\PYGZhy{}\PYGZhy{}\PYGZhy{}\PYGZhy{}\PYGZhy{}\PYGZhy{}\PYGZhy{}\PYGZhy{}\PYGZhy{}\PYGZhy{}\PYGZhy{}\PYGZhy{}\PYGZhy{}}
    \PYG{n}{n\PYGZus{}ind} \PYG{o}{=} \PYG{l+m+mi}{2}
    \PYG{c+c1}{\PYGZsh{}}
    \PYG{c+c1}{\PYGZsh{} create ax}
    \PYG{n}{x} \PYG{o}{=} \PYG{n}{numpy}\PYG{o}{.}\PYG{n}{empty}\PYG{p}{(}\PYG{n}{n\PYGZus{}ind}\PYG{p}{,} \PYG{n}{dtype}\PYG{o}{=}\PYG{n+nb}{float}\PYG{p}{)}
    \PYG{k}{for} \PYG{n}{i} \PYG{o+ow}{in} \PYG{n+nb}{range}\PYG{p}{(} \PYG{n}{n\PYGZus{}ind}  \PYG{p}{)} \PYG{p}{:}
        \PYG{n}{x}\PYG{p}{[}\PYG{n}{i}\PYG{p}{]} \PYG{o}{=} \PYG{n}{i} \PYG{o}{+} \PYG{l+m+mf}{1.0}
    \PYG{c+c1}{\PYGZsh{}}
    \PYG{n}{ax} \PYG{o}{=} \PYG{n}{cppad\PYGZus{}py}\PYG{o}{.}\PYG{n}{independent}\PYG{p}{(}\PYG{n}{x}\PYG{p}{)}
    \PYG{c+c1}{\PYGZsh{}}
    \PYG{c+c1}{\PYGZsh{} preform some a\PYGZus{}double operations}
    \PYG{n}{ax0} \PYG{o}{=} \PYG{n}{ax}\PYG{p}{[}\PYG{l+m+mi}{0}\PYG{p}{]}
    \PYG{n}{ax1} \PYG{o}{=} \PYG{n}{ax}\PYG{p}{[}\PYG{l+m+mi}{1}\PYG{p}{]}
    \PYG{n}{ay} \PYG{o}{=} \PYG{n}{ax0} \PYG{o}{+} \PYG{n}{ax1}
    \PYG{c+c1}{\PYGZsh{}}
    \PYG{c+c1}{\PYGZsh{} check that ay is a variable; its value depends on the value of ax}
    \PYG{n}{ok} \PYG{o}{=} \PYG{n}{ok} \PYG{o+ow}{and} \PYG{n}{ay}\PYG{o}{.}\PYG{n}{variable}\PYG{p}{(}\PYG{p}{)}
    \PYG{c+c1}{\PYGZsh{}}
    \PYG{c+c1}{\PYGZsh{} abort this recording}
    \PYG{n}{cppad\PYGZus{}py}\PYG{o}{.}\PYG{n}{abort\PYGZus{}recording}\PYG{p}{(}\PYG{p}{)}
    \PYG{c+c1}{\PYGZsh{}}
    \PYG{c+c1}{\PYGZsh{} check that ay is now a parameter, no longer a variable.}
    \PYG{n}{ok} \PYG{o}{=} \PYG{n}{ok} \PYG{o+ow}{and} \PYG{n}{ay}\PYG{o}{.}\PYG{n}{parameter}\PYG{p}{(}\PYG{p}{)}
    \PYG{c+c1}{\PYGZsh{}}
    \PYG{c+c1}{\PYGZsh{} since it is a parameter, we can retrieve its value}
    \PYG{n}{y} \PYG{o}{=} \PYG{n}{ay}\PYG{o}{.}\PYG{n}{value}\PYG{p}{(}\PYG{p}{)}
    \PYG{c+c1}{\PYGZsh{}}
    \PYG{c+c1}{\PYGZsh{} its value should be x0 + x1}
    \PYG{n}{ok} \PYG{o}{=} \PYG{n}{ok} \PYG{o+ow}{and} \PYG{n}{y}  \PYG{o}{==} \PYG{n}{x}\PYG{p}{[}\PYG{l+m+mi}{0}\PYG{p}{]} \PYG{o}{+} \PYG{n}{x}\PYG{p}{[}\PYG{l+m+mi}{1}\PYG{p}{]}
    \PYG{c+c1}{\PYGZsh{}}
    \PYG{c+c1}{\PYGZsh{} an abort when not recording has no effect}
    \PYG{n}{cppad\PYGZus{}py}\PYG{o}{.}\PYG{n}{abort\PYGZus{}recording}\PYG{p}{(}\PYG{p}{)}
    \PYG{c+c1}{\PYGZsh{}}
    \PYG{k}{return}\PYG{p}{(} \PYG{n}{ok} \PYG{p}{)}
\PYG{c+c1}{\PYGZsh{}}
\end{sphinxVerbatim}


\bigskip\hrule\bigskip


xsrst input file: \sphinxcode{\sphinxupquote{example/python/core/fun\_abort\_xam.py}}

\index{example@\spxentry{example}}\ignorespaces 

\subparagraph{Example}
\label{\detokenize{xsrst/py_abort_recording:example}}\label{\detokenize{xsrst/py_abort_recording:py-abort-recording-example}}\label{\detokenize{xsrst/py_abort_recording:index-3}}
{\hyperref[\detokenize{xsrst/fun_abort_xam_py:id1}]{\sphinxcrossref{\DUrole{std,std-ref}{python}}}}


\bigskip\hrule\bigskip


xsrst input file: \sphinxcode{\sphinxupquote{lib/python/cppad\_py/abort\_recording.xsrst}}


\subparagraph{py\_fun\_ctor}
\label{\detokenize{xsrst/py_fun_ctor:py-fun-ctor}}\label{\detokenize{xsrst/py_fun_ctor::doc}}
\index{py\_fun\_ctor@\spxentry{py\_fun\_ctor}}\index{stop@\spxentry{stop}}\index{current@\spxentry{current}}\index{recording@\spxentry{recording}}\index{store@\spxentry{store}}\index{function@\spxentry{function}}\index{object@\spxentry{object}}\ignorespaces 

\subparagraph{3.1.1.3 Stop Current Recording and Store Function Object}
\label{\detokenize{xsrst/py_fun_ctor:stop-current-recording-and-store-function-object}}\label{\detokenize{xsrst/py_fun_ctor:id1}}\begin{itemize}
\item {} \begin{description}
\item[{{\hyperref[\detokenize{xsrst/py_fun_ctor:py-fun-ctor-syntax}]{\sphinxcrossref{\DUrole{std,std-ref}{Syntax}}}}}] \leavevmode\begin{itemize}
\item {} 
{\hyperref[\detokenize{xsrst/py_fun_ctor:py-fun-ctor-syntax-d-fun}]{\sphinxcrossref{\DUrole{std,std-ref}{d\_fun}}}}

\item {} 
{\hyperref[\detokenize{xsrst/py_fun_ctor:py-fun-ctor-syntax-a-fun}]{\sphinxcrossref{\DUrole{std,std-ref}{a\_fun}}}}

\end{itemize}

\end{description}

\item {} 
{\hyperref[\detokenize{xsrst/py_fun_ctor:py-fun-ctor-ax}]{\sphinxcrossref{\DUrole{std,std-ref}{ax}}}}

\item {} 
{\hyperref[\detokenize{xsrst/py_fun_ctor:py-fun-ctor-ay}]{\sphinxcrossref{\DUrole{std,std-ref}{ay}}}}

\item {} \begin{description}
\item[{{\hyperref[\detokenize{xsrst/py_fun_ctor:py-fun-ctor-f}]{\sphinxcrossref{\DUrole{std,std-ref}{f}}}}}] \leavevmode\begin{itemize}
\item {} 
{\hyperref[\detokenize{xsrst/py_fun_ctor:py-fun-ctor-f-empty-function}]{\sphinxcrossref{\DUrole{std,std-ref}{Empty Function}}}}

\end{itemize}

\end{description}

\item {} 
{\hyperref[\detokenize{xsrst/py_fun_ctor:py-fun-ctor-af}]{\sphinxcrossref{\DUrole{std,std-ref}{af}}}}

\item {} 
{\hyperref[\detokenize{xsrst/py_fun_ctor:py-fun-ctor-example}]{\sphinxcrossref{\DUrole{std,std-ref}{Example}}}}

\end{itemize}

\index{syntax@\spxentry{syntax}}\ignorespaces 

\subparagraph{Syntax}
\label{\detokenize{xsrst/py_fun_ctor:syntax}}\label{\detokenize{xsrst/py_fun_ctor:py-fun-ctor-syntax}}\label{\detokenize{xsrst/py_fun_ctor:index-1}}
\index{d\_fun@\spxentry{d\_fun}}\ignorespaces 

\subparagraph{d\_fun}
\label{\detokenize{xsrst/py_fun_ctor:d-fun}}\label{\detokenize{xsrst/py_fun_ctor:py-fun-ctor-syntax-d-fun}}\label{\detokenize{xsrst/py_fun_ctor:index-2}}
\begin{DUlineblock}{0em}
\item[] \sphinxstyleemphasis{f} =  \sphinxcode{\sphinxupquote{cppad\_py.d\_fun}} ( \sphinxstyleemphasis{ax} , \sphinxstyleemphasis{ay} )
\item[] \sphinxstyleemphasis{f} =  \sphinxcode{\sphinxupquote{cppad\_py.d\_fun}} ()
\end{DUlineblock}

\index{a\_fun@\spxentry{a\_fun}}\ignorespaces 

\subparagraph{a\_fun}
\label{\detokenize{xsrst/py_fun_ctor:a-fun}}\label{\detokenize{xsrst/py_fun_ctor:py-fun-ctor-syntax-a-fun}}\label{\detokenize{xsrst/py_fun_ctor:index-3}}
\begin{DUlineblock}{0em}
\item[] \sphinxstyleemphasis{af} =  \sphinxcode{\sphinxupquote{cppad\_py.a\_fun}} ( \sphinxstyleemphasis{f} )
\end{DUlineblock}

\index{ax@\spxentry{ax}}\ignorespaces 

\subparagraph{ax}
\label{\detokenize{xsrst/py_fun_ctor:ax}}\label{\detokenize{xsrst/py_fun_ctor:py-fun-ctor-ax}}\label{\detokenize{xsrst/py_fun_ctor:index-4}}
This argument must be the same as
{\hyperref[\detokenize{xsrst/py_independent:py-independent-ax}]{\sphinxcrossref{\DUrole{std,std-ref}{ax}}}}
returned by the previous call to \sphinxcode{\sphinxupquote{independent}} ; i.e.,
it must be the independent variable vector.
We use \sphinxstyleemphasis{n}
to denote the number of independent variables (the size of \sphinxstyleemphasis{ax} ).

\index{ay@\spxentry{ay}}\ignorespaces 

\subparagraph{ay}
\label{\detokenize{xsrst/py_fun_ctor:ay}}\label{\detokenize{xsrst/py_fun_ctor:py-fun-ctor-ay}}\label{\detokenize{xsrst/py_fun_ctor:index-5}}
This argument is a numpy array with \sphinxcode{\sphinxupquote{a\_double}} elements.
It specifies the dependent variables.
We use \sphinxstyleemphasis{m}
to denote the number of dependent variables (the size of \sphinxstyleemphasis{ay} ).

\index{f@\spxentry{f}}\ignorespaces 

\subparagraph{f}
\label{\detokenize{xsrst/py_fun_ctor:f}}\label{\detokenize{xsrst/py_fun_ctor:py-fun-ctor-f}}\label{\detokenize{xsrst/py_fun_ctor:index-6}}
This result is a function object that
has a representation for the floating point operations
that mapped the independent variables \sphinxstyleemphasis{ax}
to the dependent variables \sphinxstyleemphasis{ay} .
This object computes function and derivative values using
\sphinxcode{\sphinxupquote{double}}

\index{empty@\spxentry{empty}}\index{function@\spxentry{function}}\ignorespaces 

\subparagraph{Empty Function}
\label{\detokenize{xsrst/py_fun_ctor:empty-function}}\label{\detokenize{xsrst/py_fun_ctor:py-fun-ctor-f-empty-function}}\label{\detokenize{xsrst/py_fun_ctor:index-7}}
In the case where \sphinxstyleemphasis{ax} and \sphinxstyleemphasis{ay} are not present
the function is ‘empty’ and all its sizes are zero; see
{\hyperref[\detokenize{xsrst/cpp_fun_property:id1}]{\sphinxcrossref{\DUrole{std,std-ref}{cpp\_fun\_property}}}}.

\index{af@\spxentry{af}}\ignorespaces 

\subparagraph{af}
\label{\detokenize{xsrst/py_fun_ctor:af}}\label{\detokenize{xsrst/py_fun_ctor:py-fun-ctor-af}}\label{\detokenize{xsrst/py_fun_ctor:index-8}}
This result is a function object that
has a representation for the same function as \sphinxstyleemphasis{f} .
This object computes function and derivative values using
\sphinxcode{\sphinxupquote{a\_double}}
Initially, there are not Taylor coefficient stored in \sphinxstyleemphasis{af} ; i.e.,
{\hyperref[\detokenize{xsrst/py_fun_property:py-fun-property-size-order}]{\sphinxcrossref{\DUrole{std,std-ref}{af\_size\_order()}}}} is zero.


\subparagraph{a\_fun\_xam\_py}
\label{\detokenize{xsrst/a_fun_xam_py:a-fun-xam-py}}\label{\detokenize{xsrst/a_fun_xam_py::doc}}
\index{a\_fun\_xam\_py@\spxentry{a\_fun\_xam\_py}}\index{python:@\spxentry{python:}}\index{purpose@\spxentry{purpose}}\index{a\_fun@\spxentry{a\_fun}}\index{objects:@\spxentry{objects:}}\index{example@\spxentry{example}}\index{test@\spxentry{test}}\ignorespaces 

\subparagraph{3.1.1.3.1 Python: Purpose of a\_fun Objects: Example and Test}
\label{\detokenize{xsrst/a_fun_xam_py:python-purpose-of-a-fun-objects-example-and-test}}\label{\detokenize{xsrst/a_fun_xam_py:id1}}
\begin{sphinxVerbatim}[commandchars=\\\{\}]
\PYG{k}{def} \PYG{n+nf}{a\PYGZus{}fun\PYGZus{}xam}\PYG{p}{(}\PYG{p}{)} \PYG{p}{:}
    \PYG{c+c1}{\PYGZsh{}}
    \PYG{k+kn}{import} \PYG{n+nn}{numpy}
    \PYG{k+kn}{import} \PYG{n+nn}{cppad\PYGZus{}py}
    \PYG{c+c1}{\PYGZsh{}}
    \PYG{c+c1}{\PYGZsh{} initialize return variable}
    \PYG{n}{ok} \PYG{o}{=} \PYG{k+kc}{True}
    \PYG{c+c1}{\PYGZsh{} \PYGZhy{}\PYGZhy{}\PYGZhy{}\PYGZhy{}\PYGZhy{}\PYGZhy{}\PYGZhy{}\PYGZhy{}\PYGZhy{}\PYGZhy{}\PYGZhy{}\PYGZhy{}\PYGZhy{}\PYGZhy{}\PYGZhy{}\PYGZhy{}\PYGZhy{}\PYGZhy{}\PYGZhy{}\PYGZhy{}\PYGZhy{}\PYGZhy{}\PYGZhy{}\PYGZhy{}\PYGZhy{}\PYGZhy{}\PYGZhy{}\PYGZhy{}\PYGZhy{}\PYGZhy{}\PYGZhy{}\PYGZhy{}\PYGZhy{}\PYGZhy{}\PYGZhy{}\PYGZhy{}\PYGZhy{}\PYGZhy{}\PYGZhy{}\PYGZhy{}\PYGZhy{}\PYGZhy{}\PYGZhy{}\PYGZhy{}\PYGZhy{}\PYGZhy{}\PYGZhy{}\PYGZhy{}\PYGZhy{}\PYGZhy{}\PYGZhy{}\PYGZhy{}\PYGZhy{}\PYGZhy{}\PYGZhy{}\PYGZhy{}\PYGZhy{}\PYGZhy{}\PYGZhy{}\PYGZhy{}\PYGZhy{}\PYGZhy{}\PYGZhy{}\PYGZhy{}\PYGZhy{}\PYGZhy{}\PYGZhy{}\PYGZhy{}\PYGZhy{}}
    \PYG{c+c1}{\PYGZsh{} number of dependent and independent variables}
    \PYG{n}{m} \PYG{o}{=} \PYG{l+m+mi}{1}
    \PYG{n}{n} \PYG{o}{=} \PYG{l+m+mi}{3}
    \PYG{c+c1}{\PYGZsh{}}
    \PYG{c+c1}{\PYGZsh{} Record the function}
    \PYG{c+c1}{\PYGZsh{}   f(x) = 0.5 * ( x[0]\PYGZca{}2 + ... + x[n\PYGZhy{}1]\PYGZca{}2 )}
    \PYG{n}{x} \PYG{o}{=} \PYG{n}{numpy}\PYG{o}{.}\PYG{n}{empty}\PYG{p}{(}\PYG{n}{n}\PYG{p}{,} \PYG{n}{dtype}\PYG{o}{=}\PYG{n+nb}{float}\PYG{p}{)}
    \PYG{k}{for} \PYG{n}{i} \PYG{o+ow}{in} \PYG{n+nb}{range}\PYG{p}{(}\PYG{n}{n}\PYG{p}{)} \PYG{p}{:}
        \PYG{n}{x}\PYG{p}{[}\PYG{n}{i}\PYG{p}{]} \PYG{o}{=} \PYG{n}{i}
    \PYG{n}{ax}   \PYG{o}{=} \PYG{n}{cppad\PYGZus{}py}\PYG{o}{.}\PYG{n}{independent}\PYG{p}{(}\PYG{n}{x}\PYG{p}{)}
    \PYG{n}{ay}   \PYG{o}{=} \PYG{n}{numpy}\PYG{o}{.}\PYG{n}{empty}\PYG{p}{(}\PYG{n}{m}\PYG{p}{,} \PYG{n}{dtype}\PYG{o}{=}\PYG{n}{cppad\PYGZus{}py}\PYG{o}{.}\PYG{n}{a\PYGZus{}double}\PYG{p}{)}
    \PYG{n}{asum} \PYG{o}{=} \PYG{n}{cppad\PYGZus{}py}\PYG{o}{.}\PYG{n}{a\PYGZus{}double}\PYG{p}{(}\PYG{l+m+mf}{0.0}\PYG{p}{)}
    \PYG{k}{for} \PYG{n}{i} \PYG{o+ow}{in} \PYG{n+nb}{range}\PYG{p}{(}\PYG{n}{n}\PYG{p}{)} \PYG{p}{:}
        \PYG{n}{asum} \PYG{o}{=} \PYG{n}{asum} \PYG{o}{+} \PYG{n}{ax}\PYG{p}{[}\PYG{n}{i}\PYG{p}{]} \PYG{o}{*} \PYG{n}{ax}\PYG{p}{[}\PYG{n}{i}\PYG{p}{]}
    \PYG{n}{ay}\PYG{p}{[}\PYG{l+m+mi}{0}\PYG{p}{]} \PYG{o}{=} \PYG{n}{cppad\PYGZus{}py}\PYG{o}{.}\PYG{n}{a\PYGZus{}double}\PYG{p}{(}\PYG{l+m+mf}{0.5}\PYG{p}{)} \PYG{o}{*} \PYG{n}{asum}
    \PYG{n}{f}     \PYG{o}{=} \PYG{n}{cppad\PYGZus{}py}\PYG{o}{.}\PYG{n}{d\PYGZus{}fun}\PYG{p}{(}\PYG{n}{ax}\PYG{p}{,} \PYG{n}{ay}\PYG{p}{)}
    \PYG{n}{af}    \PYG{o}{=} \PYG{n}{cppad\PYGZus{}py}\PYG{o}{.}\PYG{n}{a\PYGZus{}fun}\PYG{p}{(}\PYG{n}{f}\PYG{p}{)}
    \PYG{c+c1}{\PYGZsh{}}
    \PYG{c+c1}{\PYGZsh{} Record the function}
    \PYG{c+c1}{\PYGZsh{}   g(x) = f\PYGZsq{}(x) * v}
    \PYG{n}{au}   \PYG{o}{=} \PYG{n}{numpy}\PYG{o}{.}\PYG{n}{empty}\PYG{p}{(}\PYG{n}{m}\PYG{p}{,} \PYG{n}{dtype}\PYG{o}{=}\PYG{n}{cppad\PYGZus{}py}\PYG{o}{.}\PYG{n}{a\PYGZus{}double}\PYG{p}{)}
    \PYG{n}{av}   \PYG{o}{=} \PYG{n}{numpy}\PYG{o}{.}\PYG{n}{empty}\PYG{p}{(}\PYG{n}{n}\PYG{p}{,} \PYG{n}{dtype}\PYG{o}{=}\PYG{n}{cppad\PYGZus{}py}\PYG{o}{.}\PYG{n}{a\PYGZus{}double}\PYG{p}{)}
    \PYG{k}{for} \PYG{n}{i} \PYG{o+ow}{in} \PYG{n+nb}{range}\PYG{p}{(}\PYG{n}{n}\PYG{p}{)} \PYG{p}{:}
        \PYG{n}{av}\PYG{p}{[}\PYG{n}{i}\PYG{p}{]} \PYG{o}{=} \PYG{n}{cppad\PYGZus{}py}\PYG{o}{.}\PYG{n}{a\PYGZus{}double}\PYG{p}{(}\PYG{n}{i}\PYG{p}{)}
    \PYG{n}{ax}   \PYG{o}{=} \PYG{n}{cppad\PYGZus{}py}\PYG{o}{.}\PYG{n}{independent}\PYG{p}{(}\PYG{n}{x}\PYG{p}{)}
    \PYG{n}{af}\PYG{o}{.}\PYG{n}{forward}\PYG{p}{(}\PYG{l+m+mi}{0}\PYG{p}{,} \PYG{n}{ax}\PYG{p}{)}
    \PYG{n}{au} \PYG{o}{=} \PYG{n}{af}\PYG{o}{.}\PYG{n}{forward}\PYG{p}{(}\PYG{l+m+mi}{1}\PYG{p}{,} \PYG{n}{av}\PYG{p}{)}
    \PYG{n}{g}  \PYG{o}{=} \PYG{n}{cppad\PYGZus{}py}\PYG{o}{.}\PYG{n}{d\PYGZus{}fun}\PYG{p}{(}\PYG{n}{ax}\PYG{p}{,} \PYG{n}{au}\PYG{p}{)}
    \PYG{c+c1}{\PYGZsh{}}
    \PYG{c+c1}{\PYGZsh{} compute g(x)}
    \PYG{n}{u}  \PYG{o}{=} \PYG{n}{g}\PYG{o}{.}\PYG{n}{forward}\PYG{p}{(}\PYG{l+m+mi}{0}\PYG{p}{,} \PYG{n}{x}\PYG{p}{)}
    \PYG{c+c1}{\PYGZsh{}}
    \PYG{c+c1}{\PYGZsh{} compute g\PYGZsq{}(x)}
    \PYG{n}{uq}       \PYG{o}{=} \PYG{n}{numpy}\PYG{o}{.}\PYG{n}{empty}\PYG{p}{(} \PYG{p}{(}\PYG{n}{m}\PYG{p}{,} \PYG{l+m+mi}{1}\PYG{p}{)}\PYG{p}{,} \PYG{n}{dtype}\PYG{o}{=}\PYG{n+nb}{float} \PYG{p}{)}
    \PYG{n}{uq}\PYG{p}{[}\PYG{l+m+mi}{0}\PYG{p}{,} \PYG{l+m+mi}{0}\PYG{p}{]} \PYG{o}{=} \PYG{l+m+mf}{1.0}
    \PYG{n}{xq}       \PYG{o}{=} \PYG{n}{g}\PYG{o}{.}\PYG{n}{reverse}\PYG{p}{(}\PYG{l+m+mi}{1}\PYG{p}{,} \PYG{n}{uq}\PYG{p}{)}
    \PYG{c+c1}{\PYGZsh{}}
    \PYG{c+c1}{\PYGZsh{} g\PYGZsq{}(x) = f\PYGZsq{}\PYGZsq{}(x) * v = v}
    \PYG{k}{for} \PYG{n}{i} \PYG{o+ow}{in} \PYG{n+nb}{range}\PYG{p}{(}\PYG{n}{n}\PYG{p}{)} \PYG{p}{:}
        \PYG{n}{ok} \PYG{o}{=} \PYG{n}{ok} \PYG{o+ow}{and} \PYG{n}{cppad\PYGZus{}py}\PYG{o}{.}\PYG{n}{a\PYGZus{}double}\PYG{p}{(}\PYG{n}{xq}\PYG{p}{[}\PYG{n}{i}\PYG{p}{,} \PYG{l+m+mi}{0}\PYG{p}{]}\PYG{p}{)} \PYG{o}{==} \PYG{n}{av}\PYG{p}{[}\PYG{n}{i}\PYG{p}{]}
    \PYG{k}{return} \PYG{n}{ok}
\end{sphinxVerbatim}


\bigskip\hrule\bigskip


xsrst input file: \sphinxcode{\sphinxupquote{example/python/core/a\_fun\_xam.py}}

\index{example@\spxentry{example}}\ignorespaces 

\subparagraph{Example}
\label{\detokenize{xsrst/py_fun_ctor:example}}\label{\detokenize{xsrst/py_fun_ctor:py-fun-ctor-example}}\label{\detokenize{xsrst/py_fun_ctor:index-9}}
All of the examples use the \sphinxcode{\sphinxupquote{d\_fun}} constructor.
The example {\hyperref[\detokenize{xsrst/a_fun_xam_py:id1}]{\sphinxcrossref{\DUrole{std,std-ref}{a\_fun\_xam\_py}}}} demonstrates the purpose of
\sphinxcode{\sphinxupquote{a\_fun}} objects.


\bigskip\hrule\bigskip


xsrst input file: \sphinxcode{\sphinxupquote{lib/python/cppad\_py/fun\_ctor.py}}


\subparagraph{py\_fun\_property}
\label{\detokenize{xsrst/py_fun_property:py-fun-property}}\label{\detokenize{xsrst/py_fun_property::doc}}
\index{py\_fun\_property@\spxentry{py\_fun\_property}}\index{properties@\spxentry{properties}}\index{function@\spxentry{function}}\index{object@\spxentry{object}}\ignorespaces 

\subparagraph{3.1.1.4 Properties of a Function Object}
\label{\detokenize{xsrst/py_fun_property:properties-of-a-function-object}}\label{\detokenize{xsrst/py_fun_property:id1}}\begin{itemize}
\item {} 
{\hyperref[\detokenize{xsrst/py_fun_property:py-fun-property-syntax}]{\sphinxcrossref{\DUrole{std,std-ref}{Syntax}}}}

\item {} 
{\hyperref[\detokenize{xsrst/py_fun_property:py-fun-property-f}]{\sphinxcrossref{\DUrole{std,std-ref}{f}}}}

\item {} 
{\hyperref[\detokenize{xsrst/py_fun_property:py-fun-property-size-domain}]{\sphinxcrossref{\DUrole{std,std-ref}{size\_domain}}}}

\item {} 
{\hyperref[\detokenize{xsrst/py_fun_property:py-fun-property-size-range}]{\sphinxcrossref{\DUrole{std,std-ref}{size\_range}}}}

\item {} 
{\hyperref[\detokenize{xsrst/py_fun_property:py-fun-property-size-var}]{\sphinxcrossref{\DUrole{std,std-ref}{size\_var}}}}

\item {} 
{\hyperref[\detokenize{xsrst/py_fun_property:py-fun-property-size-op}]{\sphinxcrossref{\DUrole{std,std-ref}{size\_op}}}}

\item {} 
{\hyperref[\detokenize{xsrst/py_fun_property:py-fun-property-size-order}]{\sphinxcrossref{\DUrole{std,std-ref}{size\_order}}}}

\item {} 
{\hyperref[\detokenize{xsrst/py_fun_property:py-fun-property-example}]{\sphinxcrossref{\DUrole{std,std-ref}{Example}}}}

\end{itemize}

\index{syntax@\spxentry{syntax}}\ignorespaces 

\subparagraph{Syntax}
\label{\detokenize{xsrst/py_fun_property:syntax}}\label{\detokenize{xsrst/py_fun_property:py-fun-property-syntax}}\label{\detokenize{xsrst/py_fun_property:index-1}}
\begin{DUlineblock}{0em}
\item[] \sphinxstyleemphasis{n} = \sphinxstyleemphasis{f} . \sphinxcode{\sphinxupquote{size\_domain}} ()
\item[] \sphinxstyleemphasis{m} = \sphinxstyleemphasis{f} . \sphinxcode{\sphinxupquote{size\_range}} ()
\item[] \sphinxstyleemphasis{v} = \sphinxstyleemphasis{f} . \sphinxcode{\sphinxupquote{size\_var}} ()
\item[] \sphinxstyleemphasis{p} = \sphinxstyleemphasis{f} . \sphinxcode{\sphinxupquote{size\_op}} ()
\item[] \sphinxstyleemphasis{q} = \sphinxstyleemphasis{f} . \sphinxcode{\sphinxupquote{size\_order}} ()
\end{DUlineblock}

\index{f@\spxentry{f}}\ignorespaces 

\subparagraph{f}
\label{\detokenize{xsrst/py_fun_property:f}}\label{\detokenize{xsrst/py_fun_property:py-fun-property-f}}\label{\detokenize{xsrst/py_fun_property:index-2}}
This is either a
{\hyperref[\detokenize{xsrst/py_fun_ctor:py-fun-ctor-syntax-d-fun}]{\sphinxcrossref{\DUrole{std,std-ref}{d\_fun}}}} or
{\hyperref[\detokenize{xsrst/py_fun_ctor:py-fun-ctor-syntax-a-fun}]{\sphinxcrossref{\DUrole{std,std-ref}{a\_fun}}}} function object
and is constant; i.e., not changed.

\index{size\_domain@\spxentry{size\_domain}}\ignorespaces 

\subparagraph{size\_domain}
\label{\detokenize{xsrst/py_fun_property:size-domain}}\label{\detokenize{xsrst/py_fun_property:py-fun-property-size-domain}}\label{\detokenize{xsrst/py_fun_property:index-3}}
The return value \sphinxstyleemphasis{n} is an \sphinxcode{\sphinxupquote{int}}
and is the size of the vector
{\hyperref[\detokenize{xsrst/py_fun_ctor:py-fun-ctor-ax}]{\sphinxcrossref{\DUrole{std,std-ref}{ax}}}} in the function constructor; i.e.,
the number of independent variables.

\index{size\_range@\spxentry{size\_range}}\ignorespaces 

\subparagraph{size\_range}
\label{\detokenize{xsrst/py_fun_property:size-range}}\label{\detokenize{xsrst/py_fun_property:py-fun-property-size-range}}\label{\detokenize{xsrst/py_fun_property:index-4}}
The return value \sphinxstyleemphasis{m} is an \sphinxcode{\sphinxupquote{int}}
and is the size of the vector
{\hyperref[\detokenize{xsrst/py_fun_ctor:py-fun-ctor-ay}]{\sphinxcrossref{\DUrole{std,std-ref}{ay}}}} in the function constructor; i.e.,
the number of dependent variables.

\index{size\_var@\spxentry{size\_var}}\ignorespaces 

\subparagraph{size\_var}
\label{\detokenize{xsrst/py_fun_property:size-var}}\label{\detokenize{xsrst/py_fun_property:py-fun-property-size-var}}\label{\detokenize{xsrst/py_fun_property:index-5}}
The return value \sphinxstyleemphasis{v} is an \sphinxcode{\sphinxupquote{int}}
and is the number of variables in the function.
This includes the independent variables, dependent variables,
and any variables that are used to compute the dependent variables
from the independent variables.

\index{size\_op@\spxentry{size\_op}}\ignorespaces 

\subparagraph{size\_op}
\label{\detokenize{xsrst/py_fun_property:size-op}}\label{\detokenize{xsrst/py_fun_property:py-fun-property-size-op}}\label{\detokenize{xsrst/py_fun_property:index-6}}
The return value \sphinxstyleemphasis{p} is an \sphinxcode{\sphinxupquote{int}}
and is the number of atomic operations that are used to express
the dependent variables as a function of the independent variables.

\index{size\_order@\spxentry{size\_order}}\ignorespaces 

\subparagraph{size\_order}
\label{\detokenize{xsrst/py_fun_property:size-order}}\label{\detokenize{xsrst/py_fun_property:py-fun-property-size-order}}\label{\detokenize{xsrst/py_fun_property:index-7}}
The return value \sphinxstyleemphasis{q} is an \sphinxcode{\sphinxupquote{int}}
and is the number of Taylor coefficients currently stored in \sphinxstyleemphasis{f} ,
for every variable in the operation sequence corresponding to \sphinxstyleemphasis{f} .
These coefficients are computed by {\hyperref[\detokenize{xsrst/py_fun_forward:id1}]{\sphinxcrossref{\DUrole{std,std-ref}{py\_fun\_forward}}}}.
This is different from the other function properties in that it can change
after each call to \sphinxstyleemphasis{f} . \sphinxcode{\sphinxupquote{forward}} ; see
{\hyperref[\detokenize{xsrst/py_fun_forward:py-fun-forward-p-size-order}]{\sphinxcrossref{\DUrole{std,std-ref}{size\_order}}}} in the forward mode section.


\subparagraph{fun\_property\_xam\_py}
\label{\detokenize{xsrst/fun_property_xam_py:fun-property-xam-py}}\label{\detokenize{xsrst/fun_property_xam_py::doc}}
\index{fun\_property\_xam\_py@\spxentry{fun\_property\_xam\_py}}\index{python:@\spxentry{python:}}\index{d\_fun@\spxentry{d\_fun}}\index{properties:@\spxentry{properties:}}\index{example@\spxentry{example}}\index{test@\spxentry{test}}\ignorespaces 

\subparagraph{3.1.1.4.1 Python: d\_fun Properties: Example and Test}
\label{\detokenize{xsrst/fun_property_xam_py:python-d-fun-properties-example-and-test}}\label{\detokenize{xsrst/fun_property_xam_py:id1}}
\begin{sphinxVerbatim}[commandchars=\\\{\}]
\PYG{k}{def} \PYG{n+nf}{fun\PYGZus{}property\PYGZus{}xam}\PYG{p}{(}\PYG{p}{)} \PYG{p}{:}
    \PYG{c+c1}{\PYGZsh{}}
    \PYG{k+kn}{import} \PYG{n+nn}{numpy}
    \PYG{k+kn}{import} \PYG{n+nn}{cppad\PYGZus{}py}
    \PYG{c+c1}{\PYGZsh{}}
    \PYG{c+c1}{\PYGZsh{} initialize return variable}
    \PYG{n}{ok} \PYG{o}{=} \PYG{k+kc}{True}
    \PYG{c+c1}{\PYGZsh{} \PYGZhy{}\PYGZhy{}\PYGZhy{}\PYGZhy{}\PYGZhy{}\PYGZhy{}\PYGZhy{}\PYGZhy{}\PYGZhy{}\PYGZhy{}\PYGZhy{}\PYGZhy{}\PYGZhy{}\PYGZhy{}\PYGZhy{}\PYGZhy{}\PYGZhy{}\PYGZhy{}\PYGZhy{}\PYGZhy{}\PYGZhy{}\PYGZhy{}\PYGZhy{}\PYGZhy{}\PYGZhy{}\PYGZhy{}\PYGZhy{}\PYGZhy{}\PYGZhy{}\PYGZhy{}\PYGZhy{}\PYGZhy{}\PYGZhy{}\PYGZhy{}\PYGZhy{}\PYGZhy{}\PYGZhy{}\PYGZhy{}\PYGZhy{}\PYGZhy{}\PYGZhy{}\PYGZhy{}\PYGZhy{}\PYGZhy{}\PYGZhy{}\PYGZhy{}\PYGZhy{}\PYGZhy{}\PYGZhy{}\PYGZhy{}\PYGZhy{}\PYGZhy{}\PYGZhy{}\PYGZhy{}\PYGZhy{}\PYGZhy{}\PYGZhy{}\PYGZhy{}\PYGZhy{}\PYGZhy{}\PYGZhy{}\PYGZhy{}\PYGZhy{}\PYGZhy{}\PYGZhy{}\PYGZhy{}\PYGZhy{}\PYGZhy{}\PYGZhy{}}
    \PYG{n}{n\PYGZus{}ind} \PYG{o}{=} \PYG{l+m+mi}{1} \PYG{c+c1}{\PYGZsh{} number of independent variables}
    \PYG{n}{n\PYGZus{}dep} \PYG{o}{=} \PYG{l+m+mi}{2} \PYG{c+c1}{\PYGZsh{} number of dependent variables}
    \PYG{n}{n\PYGZus{}var} \PYG{o}{=} \PYG{l+m+mi}{1} \PYG{c+c1}{\PYGZsh{} phantom variable at address 0}
    \PYG{n}{n\PYGZus{}op}  \PYG{o}{=} \PYG{l+m+mi}{1} \PYG{c+c1}{\PYGZsh{} special operator at beginning}
    \PYG{c+c1}{\PYGZsh{}}
    \PYG{c+c1}{\PYGZsh{} dimension some vectors}
    \PYG{c+c1}{\PYGZsh{} x  = numpy.empty(n\PYGZus{}ind, dtype=float)}
    \PYG{n}{x}  \PYG{o}{=} \PYG{n}{numpy}\PYG{o}{.}\PYG{n}{empty}\PYG{p}{(}\PYG{n}{n\PYGZus{}ind}\PYG{p}{,} \PYG{n}{dtype}\PYG{o}{=}\PYG{n+nb}{float}\PYG{p}{)}
    \PYG{n}{ay} \PYG{o}{=} \PYG{n}{numpy}\PYG{o}{.}\PYG{n}{empty}\PYG{p}{(}\PYG{n}{n\PYGZus{}dep}\PYG{p}{,} \PYG{n}{dtype}\PYG{o}{=}\PYG{n}{cppad\PYGZus{}py}\PYG{o}{.}\PYG{n}{a\PYGZus{}double}\PYG{p}{)}
    \PYG{c+c1}{\PYGZsh{}}
    \PYG{c+c1}{\PYGZsh{} independent variables}
    \PYG{n}{x}\PYG{p}{[}\PYG{l+m+mi}{0}\PYG{p}{]}  \PYG{o}{=} \PYG{l+m+mf}{1.0}
    \PYG{n}{ax}    \PYG{o}{=} \PYG{n}{cppad\PYGZus{}py}\PYG{o}{.}\PYG{n}{independent}\PYG{p}{(}\PYG{n}{x}\PYG{p}{)}
    \PYG{n}{n\PYGZus{}var} \PYG{o}{=} \PYG{n}{n\PYGZus{}var} \PYG{o}{+} \PYG{n}{n\PYGZus{}ind} \PYG{c+c1}{\PYGZsh{} one for each indpendent}
    \PYG{n}{n\PYGZus{}op}  \PYG{o}{=} \PYG{n}{n\PYGZus{}op} \PYG{o}{+} \PYG{n}{n\PYGZus{}ind}
    \PYG{c+c1}{\PYGZsh{}}
    \PYG{c+c1}{\PYGZsh{} first dependent variable}
    \PYG{n}{ay}\PYG{p}{[}\PYG{l+m+mi}{0}\PYG{p}{]} \PYG{o}{=} \PYG{n}{ax}\PYG{p}{[}\PYG{l+m+mi}{0}\PYG{p}{]} \PYG{o}{+} \PYG{n}{ax}\PYG{p}{[}\PYG{l+m+mi}{0}\PYG{p}{]}
    \PYG{n}{n\PYGZus{}var} \PYG{o}{=} \PYG{n}{n\PYGZus{}var} \PYG{o}{+} \PYG{l+m+mi}{1} \PYG{c+c1}{\PYGZsh{} one variable and operator}
    \PYG{n}{n\PYGZus{}op}  \PYG{o}{=} \PYG{n}{n\PYGZus{}op} \PYG{o}{+} \PYG{l+m+mi}{1}
    \PYG{c+c1}{\PYGZsh{}}
    \PYG{c+c1}{\PYGZsh{} second dependent variable}
    \PYG{n}{ax0}   \PYG{o}{=} \PYG{n}{ax}\PYG{p}{[}\PYG{l+m+mi}{0}\PYG{p}{]}
    \PYG{n}{ay}\PYG{p}{[}\PYG{l+m+mi}{1}\PYG{p}{]} \PYG{o}{=} \PYG{n}{ax0}\PYG{o}{.}\PYG{n}{sin}\PYG{p}{(}\PYG{p}{)}
    \PYG{n}{n\PYGZus{}var} \PYG{o}{=} \PYG{n}{n\PYGZus{}var} \PYG{o}{+} \PYG{l+m+mi}{2} \PYG{c+c1}{\PYGZsh{} two varialbes, one operator}
    \PYG{n}{n\PYGZus{}op}  \PYG{o}{=} \PYG{n}{n\PYGZus{}op} \PYG{o}{+} \PYG{l+m+mi}{1}
    \PYG{c+c1}{\PYGZsh{}}
    \PYG{c+c1}{\PYGZsh{} define f(x) = y}
    \PYG{n}{f}  \PYG{o}{=} \PYG{n}{cppad\PYGZus{}py}\PYG{o}{.}\PYG{n}{d\PYGZus{}fun}\PYG{p}{(}\PYG{n}{ax}\PYG{p}{,} \PYG{n}{ay}\PYG{p}{)}
    \PYG{n}{n\PYGZus{}op} \PYG{o}{=} \PYG{n}{n\PYGZus{}op} \PYG{o}{+} \PYG{l+m+mi}{1} \PYG{c+c1}{\PYGZsh{} speical operator at end}
    \PYG{c+c1}{\PYGZsh{}}
    \PYG{c+c1}{\PYGZsh{} check af properties except f.to\PYGZus{}json}
    \PYG{n}{ok} \PYG{o}{=} \PYG{n}{ok} \PYG{o+ow}{and} \PYG{n}{f}\PYG{o}{.}\PYG{n}{size\PYGZus{}domain}\PYG{p}{(}\PYG{p}{)} \PYG{o}{==} \PYG{n}{n\PYGZus{}ind}
    \PYG{n}{ok} \PYG{o}{=} \PYG{n}{ok} \PYG{o+ow}{and} \PYG{n}{f}\PYG{o}{.}\PYG{n}{size\PYGZus{}range}\PYG{p}{(}\PYG{p}{)}  \PYG{o}{==} \PYG{n}{n\PYGZus{}dep}
    \PYG{n}{ok} \PYG{o}{=} \PYG{n}{ok} \PYG{o+ow}{and} \PYG{n}{f}\PYG{o}{.}\PYG{n}{size\PYGZus{}var}\PYG{p}{(}\PYG{p}{)}    \PYG{o}{==} \PYG{n}{n\PYGZus{}var}
    \PYG{n}{ok} \PYG{o}{=} \PYG{n}{ok} \PYG{o+ow}{and} \PYG{n}{f}\PYG{o}{.}\PYG{n}{size\PYGZus{}op}\PYG{p}{(}\PYG{p}{)}     \PYG{o}{==} \PYG{n}{n\PYGZus{}op}
    \PYG{n}{ok} \PYG{o}{=} \PYG{n}{ok} \PYG{o+ow}{and} \PYG{n}{f}\PYG{o}{.}\PYG{n}{size\PYGZus{}order}\PYG{p}{(}\PYG{p}{)}  \PYG{o}{==} \PYG{l+m+mi}{0}
    \PYG{c+c1}{\PYGZsh{}}
    \PYG{c+c1}{\PYGZsh{} compute zero order Taylor coefficients}
    \PYG{n}{y}  \PYG{o}{=} \PYG{n}{f}\PYG{o}{.}\PYG{n}{forward}\PYG{p}{(}\PYG{l+m+mi}{0}\PYG{p}{,} \PYG{n}{x}\PYG{p}{)}
    \PYG{n}{ok} \PYG{o}{=} \PYG{n}{ok} \PYG{o+ow}{and} \PYG{n}{f}\PYG{o}{.}\PYG{n}{size\PYGZus{}order}\PYG{p}{(}\PYG{p}{)} \PYG{o}{==} \PYG{l+m+mi}{1}
    \PYG{c+c1}{\PYGZsh{} \PYGZhy{}\PYGZhy{}\PYGZhy{}\PYGZhy{}\PYGZhy{}\PYGZhy{}\PYGZhy{}\PYGZhy{}\PYGZhy{}\PYGZhy{}\PYGZhy{}\PYGZhy{}\PYGZhy{}\PYGZhy{}\PYGZhy{}\PYGZhy{}\PYGZhy{}\PYGZhy{}\PYGZhy{}\PYGZhy{}\PYGZhy{}\PYGZhy{}\PYGZhy{}\PYGZhy{}\PYGZhy{}\PYGZhy{}\PYGZhy{}\PYGZhy{}\PYGZhy{}\PYGZhy{}\PYGZhy{}\PYGZhy{}\PYGZhy{}\PYGZhy{}\PYGZhy{}\PYGZhy{}\PYGZhy{}\PYGZhy{}\PYGZhy{}\PYGZhy{}\PYGZhy{}\PYGZhy{}\PYGZhy{}\PYGZhy{}\PYGZhy{}\PYGZhy{}\PYGZhy{}\PYGZhy{}\PYGZhy{}\PYGZhy{}\PYGZhy{}\PYGZhy{}\PYGZhy{}\PYGZhy{}\PYGZhy{}\PYGZhy{}\PYGZhy{}\PYGZhy{}\PYGZhy{}\PYGZhy{}\PYGZhy{}\PYGZhy{}\PYGZhy{}\PYGZhy{}\PYGZhy{}\PYGZhy{}\PYGZhy{}\PYGZhy{}\PYGZhy{}}
    \PYG{n}{af} \PYG{o}{=} \PYG{n}{cppad\PYGZus{}py}\PYG{o}{.}\PYG{n}{a\PYGZus{}fun}\PYG{p}{(}\PYG{n}{f}\PYG{p}{)}
    \PYG{c+c1}{\PYGZsh{}}
    \PYG{c+c1}{\PYGZsh{} check af properties}
    \PYG{n}{ok} \PYG{o}{=} \PYG{n}{ok} \PYG{o+ow}{and} \PYG{n}{af}\PYG{o}{.}\PYG{n}{size\PYGZus{}domain}\PYG{p}{(}\PYG{p}{)} \PYG{o}{==} \PYG{n}{n\PYGZus{}ind}
    \PYG{n}{ok} \PYG{o}{=} \PYG{n}{ok} \PYG{o+ow}{and} \PYG{n}{af}\PYG{o}{.}\PYG{n}{size\PYGZus{}range}\PYG{p}{(}\PYG{p}{)}  \PYG{o}{==} \PYG{n}{n\PYGZus{}dep}
    \PYG{n}{ok} \PYG{o}{=} \PYG{n}{ok} \PYG{o+ow}{and} \PYG{n}{af}\PYG{o}{.}\PYG{n}{size\PYGZus{}var}\PYG{p}{(}\PYG{p}{)}    \PYG{o}{==} \PYG{n}{n\PYGZus{}var}
    \PYG{n}{ok} \PYG{o}{=} \PYG{n}{ok} \PYG{o+ow}{and} \PYG{n}{af}\PYG{o}{.}\PYG{n}{size\PYGZus{}op}\PYG{p}{(}\PYG{p}{)}     \PYG{o}{==} \PYG{n}{n\PYGZus{}op}
    \PYG{n}{ok} \PYG{o}{=} \PYG{n}{ok} \PYG{o+ow}{and} \PYG{n}{af}\PYG{o}{.}\PYG{n}{size\PYGZus{}order}\PYG{p}{(}\PYG{p}{)}  \PYG{o}{==} \PYG{l+m+mi}{0}
    \PYG{c+c1}{\PYGZsh{} \PYGZhy{}\PYGZhy{}\PYGZhy{}\PYGZhy{}\PYGZhy{}\PYGZhy{}\PYGZhy{}\PYGZhy{}\PYGZhy{}\PYGZhy{}\PYGZhy{}\PYGZhy{}\PYGZhy{}\PYGZhy{}\PYGZhy{}\PYGZhy{}\PYGZhy{}\PYGZhy{}\PYGZhy{}\PYGZhy{}\PYGZhy{}\PYGZhy{}\PYGZhy{}\PYGZhy{}\PYGZhy{}\PYGZhy{}\PYGZhy{}\PYGZhy{}\PYGZhy{}\PYGZhy{}\PYGZhy{}\PYGZhy{}\PYGZhy{}\PYGZhy{}\PYGZhy{}\PYGZhy{}\PYGZhy{}\PYGZhy{}\PYGZhy{}\PYGZhy{}\PYGZhy{}\PYGZhy{}\PYGZhy{}\PYGZhy{}\PYGZhy{}\PYGZhy{}\PYGZhy{}\PYGZhy{}\PYGZhy{}\PYGZhy{}\PYGZhy{}\PYGZhy{}\PYGZhy{}\PYGZhy{}\PYGZhy{}\PYGZhy{}\PYGZhy{}\PYGZhy{}\PYGZhy{}\PYGZhy{}\PYGZhy{}\PYGZhy{}\PYGZhy{}\PYGZhy{}\PYGZhy{}\PYGZhy{}\PYGZhy{}\PYGZhy{}\PYGZhy{}}
    \PYG{c+c1}{\PYGZsh{} The empty function}
    \PYG{n}{f} \PYG{o}{=} \PYG{n}{cppad\PYGZus{}py}\PYG{o}{.}\PYG{n}{d\PYGZus{}fun}\PYG{p}{(}\PYG{p}{)}
    \PYG{n}{ok} \PYG{o}{=} \PYG{n}{ok} \PYG{o+ow}{and} \PYG{n}{f}\PYG{o}{.}\PYG{n}{size\PYGZus{}domain}\PYG{p}{(}\PYG{p}{)} \PYG{o}{==} \PYG{l+m+mi}{0}
    \PYG{n}{ok} \PYG{o}{=} \PYG{n}{ok} \PYG{o+ow}{and} \PYG{n}{f}\PYG{o}{.}\PYG{n}{size\PYGZus{}range}\PYG{p}{(}\PYG{p}{)}  \PYG{o}{==} \PYG{l+m+mi}{0}
    \PYG{n}{ok} \PYG{o}{=} \PYG{n}{ok} \PYG{o+ow}{and} \PYG{n}{f}\PYG{o}{.}\PYG{n}{size\PYGZus{}var}\PYG{p}{(}\PYG{p}{)}    \PYG{o}{==} \PYG{l+m+mi}{0}
    \PYG{n}{ok} \PYG{o}{=} \PYG{n}{ok} \PYG{o+ow}{and} \PYG{n}{f}\PYG{o}{.}\PYG{n}{size\PYGZus{}op}\PYG{p}{(}\PYG{p}{)}     \PYG{o}{==} \PYG{l+m+mi}{0}
    \PYG{n}{ok} \PYG{o}{=} \PYG{n}{ok} \PYG{o+ow}{and} \PYG{n}{f}\PYG{o}{.}\PYG{n}{size\PYGZus{}order}\PYG{p}{(}\PYG{p}{)}  \PYG{o}{==} \PYG{l+m+mi}{0}
    \PYG{c+c1}{\PYGZsh{}}
    \PYG{k}{return}\PYG{p}{(} \PYG{n}{ok}  \PYG{p}{)}
\PYG{c+c1}{\PYGZsh{}}
\end{sphinxVerbatim}


\bigskip\hrule\bigskip


xsrst input file: \sphinxcode{\sphinxupquote{example/python/core/fun\_property\_xam.py}}

\index{example@\spxentry{example}}\ignorespaces 

\subparagraph{Example}
\label{\detokenize{xsrst/py_fun_property:example}}\label{\detokenize{xsrst/py_fun_property:py-fun-property-example}}\label{\detokenize{xsrst/py_fun_property:index-8}}
{\hyperref[\detokenize{xsrst/fun_property_xam_py:id1}]{\sphinxcrossref{\DUrole{std,std-ref}{fun\_property\_xam\_py}}}}


\bigskip\hrule\bigskip


xsrst input file: \sphinxcode{\sphinxupquote{lib/python/cppad\_py/fun\_property.xsrst}}


\subparagraph{py\_fun\_new\_dynamic}
\label{\detokenize{xsrst/py_fun_new_dynamic:py-fun-new-dynamic}}\label{\detokenize{xsrst/py_fun_new_dynamic::doc}}
\index{py\_fun\_new\_dynamic@\spxentry{py\_fun\_new\_dynamic}}\index{new@\spxentry{new}}\index{dynamic@\spxentry{dynamic}}\index{parameters@\spxentry{parameters}}\ignorespaces 

\subparagraph{3.1.1.5 New Dynamic Parameters}
\label{\detokenize{xsrst/py_fun_new_dynamic:new-dynamic-parameters}}\label{\detokenize{xsrst/py_fun_new_dynamic:id1}}\begin{itemize}
\item {} 
{\hyperref[\detokenize{xsrst/py_fun_new_dynamic:py-fun-new-dynamic-syntax}]{\sphinxcrossref{\DUrole{std,std-ref}{Syntax}}}}

\item {} 
{\hyperref[\detokenize{xsrst/py_fun_new_dynamic:py-fun-new-dynamic-f}]{\sphinxcrossref{\DUrole{std,std-ref}{f}}}}

\item {} \begin{description}
\item[{{\hyperref[\detokenize{xsrst/py_fun_new_dynamic:py-fun-new-dynamic-dynamic}]{\sphinxcrossref{\DUrole{std,std-ref}{dynamic}}}}}] \leavevmode\begin{itemize}
\item {} 
{\hyperref[\detokenize{xsrst/py_fun_new_dynamic:py-fun-new-dynamic-dynamic-size-order}]{\sphinxcrossref{\DUrole{std,std-ref}{size\_order}}}}

\end{itemize}

\end{description}

\item {} 
{\hyperref[\detokenize{xsrst/py_fun_new_dynamic:py-fun-new-dynamic-example}]{\sphinxcrossref{\DUrole{std,std-ref}{Example}}}}

\end{itemize}

\index{syntax@\spxentry{syntax}}\ignorespaces 

\subparagraph{Syntax}
\label{\detokenize{xsrst/py_fun_new_dynamic:syntax}}\label{\detokenize{xsrst/py_fun_new_dynamic:py-fun-new-dynamic-syntax}}\label{\detokenize{xsrst/py_fun_new_dynamic:index-1}}
\sphinxstyleemphasis{f} . \sphinxcode{\sphinxupquote{new\_dynamic}} ( \sphinxstyleemphasis{dynamic} )

\index{f@\spxentry{f}}\ignorespaces 

\subparagraph{f}
\label{\detokenize{xsrst/py_fun_new_dynamic:f}}\label{\detokenize{xsrst/py_fun_new_dynamic:py-fun-new-dynamic-f}}\label{\detokenize{xsrst/py_fun_new_dynamic:index-2}}
This is either a
{\hyperref[\detokenize{xsrst/py_fun_ctor:py-fun-ctor-syntax-d-fun}]{\sphinxcrossref{\DUrole{std,std-ref}{d\_fun}}}} or
{\hyperref[\detokenize{xsrst/py_fun_ctor:py-fun-ctor-syntax-a-fun}]{\sphinxcrossref{\DUrole{std,std-ref}{a\_fun}}}}.
The independent {\hyperref[\detokenize{xsrst/py_independent:py-independent-dynamic}]{\sphinxcrossref{\DUrole{std,std-ref}{dynamic}}}} parameters
are changed to have the specified values.
The other dynamic parameters are then computed.

\index{dynamic@\spxentry{dynamic}}\ignorespaces 

\subparagraph{dynamic}
\label{\detokenize{xsrst/py_fun_new_dynamic:dynamic}}\label{\detokenize{xsrst/py_fun_new_dynamic:py-fun-new-dynamic-dynamic}}\label{\detokenize{xsrst/py_fun_new_dynamic:index-3}}
If \sphinxstyleemphasis{f} is a \sphinxcode{\sphinxupquote{d\_fun}} ( \sphinxcode{\sphinxupquote{a\_fun}} ) object,
\sphinxstyleemphasis{dynamic} is a numpy vector with \sphinxcode{\sphinxupquote{float}} ( \sphinxcode{\sphinxupquote{a\_double}} )
elements and its size must be the same as the size of
{\hyperref[\detokenize{xsrst/py_independent:py-independent-dynamic}]{\sphinxcrossref{\DUrole{std,std-ref}{dynamic}}}} in the corresponding call to
\sphinxcode{\sphinxupquote{independent}} .
It specifies new values for the dynamic parameters in \sphinxstyleemphasis{f} .

\index{size\_order@\spxentry{size\_order}}\ignorespaces 

\subparagraph{size\_order}
\label{\detokenize{xsrst/py_fun_new_dynamic:size-order}}\label{\detokenize{xsrst/py_fun_new_dynamic:py-fun-new-dynamic-dynamic-size-order}}\label{\detokenize{xsrst/py_fun_new_dynamic:index-4}}
After this call,
{\hyperref[\detokenize{xsrst/py_fun_property:py-fun-property-size-order}]{\sphinxcrossref{\DUrole{std,std-ref}{f\_size\_order()}}}} is zero.

\index{example@\spxentry{example}}\ignorespaces 

\subparagraph{Example}
\label{\detokenize{xsrst/py_fun_new_dynamic:example}}\label{\detokenize{xsrst/py_fun_new_dynamic:py-fun-new-dynamic-example}}\label{\detokenize{xsrst/py_fun_new_dynamic:index-5}}
See {\hyperref[\detokenize{xsrst/fun_dynamic_xam_py:id1}]{\sphinxcrossref{\DUrole{std,std-ref}{fun\_dynamic\_xam\_py}}}}


\bigskip\hrule\bigskip


xsrst input file: \sphinxcode{\sphinxupquote{lib/python/cppad\_py/fun\_new\_dynamic.py}}


\subparagraph{py\_fun\_jacobian}
\label{\detokenize{xsrst/py_fun_jacobian:py-fun-jacobian}}\label{\detokenize{xsrst/py_fun_jacobian::doc}}
\index{py\_fun\_jacobian@\spxentry{py\_fun\_jacobian}}\index{jacobian@\spxentry{jacobian}}\index{an@\spxentry{an}}\index{ad@\spxentry{ad}}\index{function@\spxentry{function}}\ignorespaces 

\subparagraph{3.1.1.6 Jacobian of an AD Function}
\label{\detokenize{xsrst/py_fun_jacobian:jacobian-of-an-ad-function}}\label{\detokenize{xsrst/py_fun_jacobian:id1}}\begin{itemize}
\item {} 
{\hyperref[\detokenize{xsrst/py_fun_jacobian:py-fun-jacobian-syntax}]{\sphinxcrossref{\DUrole{std,std-ref}{Syntax}}}}

\item {} 
{\hyperref[\detokenize{xsrst/py_fun_jacobian:py-fun-jacobian-f}]{\sphinxcrossref{\DUrole{std,std-ref}{f}}}}

\item {} 
{\hyperref[\detokenize{xsrst/py_fun_jacobian:py-fun-jacobian-f-x}]{\sphinxcrossref{\DUrole{std,std-ref}{f(x)}}}}

\item {} 
{\hyperref[\detokenize{xsrst/py_fun_jacobian:py-fun-jacobian-x}]{\sphinxcrossref{\DUrole{std,std-ref}{x}}}}

\item {} 
{\hyperref[\detokenize{xsrst/py_fun_jacobian:py-fun-jacobian-j}]{\sphinxcrossref{\DUrole{std,std-ref}{J}}}}

\item {} 
{\hyperref[\detokenize{xsrst/py_fun_jacobian:py-fun-jacobian-example}]{\sphinxcrossref{\DUrole{std,std-ref}{Example}}}}

\end{itemize}

\index{syntax@\spxentry{syntax}}\ignorespaces 

\subparagraph{Syntax}
\label{\detokenize{xsrst/py_fun_jacobian:syntax}}\label{\detokenize{xsrst/py_fun_jacobian:py-fun-jacobian-syntax}}\label{\detokenize{xsrst/py_fun_jacobian:index-1}}
\sphinxstyleemphasis{J} = \sphinxstyleemphasis{f} . \sphinxcode{\sphinxupquote{jacobian}} ( \sphinxstyleemphasis{x} )

\index{f@\spxentry{f}}\ignorespaces 

\subparagraph{f}
\label{\detokenize{xsrst/py_fun_jacobian:f}}\label{\detokenize{xsrst/py_fun_jacobian:py-fun-jacobian-f}}\label{\detokenize{xsrst/py_fun_jacobian:index-2}}
This is either a
{\hyperref[\detokenize{xsrst/py_fun_ctor:py-fun-ctor-syntax-d-fun}]{\sphinxcrossref{\DUrole{std,std-ref}{d\_fun}}}} or
{\hyperref[\detokenize{xsrst/py_fun_ctor:py-fun-ctor-syntax-a-fun}]{\sphinxcrossref{\DUrole{std,std-ref}{a\_fun}}}} function object.
Upon return, the zero order
{\hyperref[\detokenize{xsrst/py_fun_forward:py-fun-forward-taylor-coefficient}]{\sphinxcrossref{\DUrole{std,std-ref}{taylor\_coefficients}}}}
in \sphinxstyleemphasis{f} correspond to the value of \sphinxstyleemphasis{x} .
The other Taylor coefficients in \sphinxstyleemphasis{f} are unspecified.

\index{f(x)@\spxentry{f(x)}}\ignorespaces 

\subparagraph{f(x)}
\label{\detokenize{xsrst/py_fun_jacobian:f-x}}\label{\detokenize{xsrst/py_fun_jacobian:py-fun-jacobian-f-x}}\label{\detokenize{xsrst/py_fun_jacobian:index-3}}
We use the notation \(f: \B{R}^n \rightarrow \B{R}^m\)
for the function corresponding to \sphinxstyleemphasis{f} .
Note that \sphinxstyleemphasis{n} is the size of {\hyperref[\detokenize{xsrst/py_fun_ctor:py-fun-ctor-ax}]{\sphinxcrossref{\DUrole{std,std-ref}{ax}}}}
and \sphinxstyleemphasis{m} is the size of {\hyperref[\detokenize{xsrst/py_fun_ctor:py-fun-ctor-ay}]{\sphinxcrossref{\DUrole{std,std-ref}{ay}}}}
in to the constructor for \sphinxstyleemphasis{f} .

\index{x@\spxentry{x}}\ignorespaces 

\subparagraph{x}
\label{\detokenize{xsrst/py_fun_jacobian:x}}\label{\detokenize{xsrst/py_fun_jacobian:py-fun-jacobian-x}}\label{\detokenize{xsrst/py_fun_jacobian:index-4}}
If \sphinxstyleemphasis{f} is a \sphinxcode{\sphinxupquote{d\_fun}} or \sphinxcode{\sphinxupquote{a\_fun}} ,
\sphinxstyleemphasis{x} is a numpy vector with \sphinxcode{\sphinxupquote{float}} ( \sphinxcode{\sphinxupquote{a\_double}} ) elements
and size \sphinxstyleemphasis{n} .
It specifies the argument value at we are computing the Jacobian
\(f'(x)\).

\index{j@\spxentry{j}}\ignorespaces 

\subparagraph{J}
\label{\detokenize{xsrst/py_fun_jacobian:j}}\label{\detokenize{xsrst/py_fun_jacobian:py-fun-jacobian-j}}\label{\detokenize{xsrst/py_fun_jacobian:index-5}}
The result is a numpy matrix with \sphinxcode{\sphinxupquote{float}} ( \sphinxcode{\sphinxupquote{a\_double}} ) elements,
\sphinxstyleemphasis{m} rows, and \sphinxcode{\sphinxupquote{n}} columns.
For \sphinxstyleemphasis{i} between zero and \sphinxstyleemphasis{m} \sphinxhyphen{}1
and \sphinxstyleemphasis{j} between zero and \sphinxstyleemphasis{n} \sphinxhyphen{}1 ,
\begin{equation*}
\begin{split}J [ i,  j ] = \frac{ \partial f_i }{ \partial x_j } (x)\end{split}
\end{equation*}

\subparagraph{fun\_jacobian\_xam\_py}
\label{\detokenize{xsrst/fun_jacobian_xam_py:fun-jacobian-xam-py}}\label{\detokenize{xsrst/fun_jacobian_xam_py::doc}}
\index{fun\_jacobian\_xam\_py@\spxentry{fun\_jacobian\_xam\_py}}\index{python:@\spxentry{python:}}\index{dense@\spxentry{dense}}\index{jacobian@\spxentry{jacobian}}\index{using@\spxentry{using}}\index{ad:@\spxentry{ad:}}\index{example@\spxentry{example}}\index{test@\spxentry{test}}\ignorespaces 

\subparagraph{3.1.1.6.1 Python: Dense Jacobian Using AD: Example and Test}
\label{\detokenize{xsrst/fun_jacobian_xam_py:python-dense-jacobian-using-ad-example-and-test}}\label{\detokenize{xsrst/fun_jacobian_xam_py:id1}}
\begin{sphinxVerbatim}[commandchars=\\\{\}]
\PYG{k}{def} \PYG{n+nf}{fun\PYGZus{}jacobian\PYGZus{}xam}\PYG{p}{(}\PYG{p}{)} \PYG{p}{:}
    \PYG{c+c1}{\PYGZsh{}}
    \PYG{k+kn}{import} \PYG{n+nn}{numpy}
    \PYG{k+kn}{import} \PYG{n+nn}{cppad\PYGZus{}py}
    \PYG{c+c1}{\PYGZsh{}}
    \PYG{c+c1}{\PYGZsh{} initialize return variable}
    \PYG{n}{ok} \PYG{o}{=} \PYG{k+kc}{True}
    \PYG{c+c1}{\PYGZsh{} \PYGZhy{}\PYGZhy{}\PYGZhy{}\PYGZhy{}\PYGZhy{}\PYGZhy{}\PYGZhy{}\PYGZhy{}\PYGZhy{}\PYGZhy{}\PYGZhy{}\PYGZhy{}\PYGZhy{}\PYGZhy{}\PYGZhy{}\PYGZhy{}\PYGZhy{}\PYGZhy{}\PYGZhy{}\PYGZhy{}\PYGZhy{}\PYGZhy{}\PYGZhy{}\PYGZhy{}\PYGZhy{}\PYGZhy{}\PYGZhy{}\PYGZhy{}\PYGZhy{}\PYGZhy{}\PYGZhy{}\PYGZhy{}\PYGZhy{}\PYGZhy{}\PYGZhy{}\PYGZhy{}\PYGZhy{}\PYGZhy{}\PYGZhy{}\PYGZhy{}\PYGZhy{}\PYGZhy{}\PYGZhy{}\PYGZhy{}\PYGZhy{}\PYGZhy{}\PYGZhy{}\PYGZhy{}\PYGZhy{}\PYGZhy{}\PYGZhy{}\PYGZhy{}\PYGZhy{}\PYGZhy{}\PYGZhy{}\PYGZhy{}\PYGZhy{}\PYGZhy{}\PYGZhy{}\PYGZhy{}\PYGZhy{}\PYGZhy{}\PYGZhy{}\PYGZhy{}\PYGZhy{}\PYGZhy{}\PYGZhy{}\PYGZhy{}\PYGZhy{}}
    \PYG{c+c1}{\PYGZsh{} number of dependent and independent variables}
    \PYG{n}{n\PYGZus{}dep} \PYG{o}{=} \PYG{l+m+mi}{1}
    \PYG{n}{n\PYGZus{}ind} \PYG{o}{=} \PYG{l+m+mi}{3}
    \PYG{c+c1}{\PYGZsh{}}
    \PYG{c+c1}{\PYGZsh{} create the independent variables ax}
    \PYG{n}{x} \PYG{o}{=} \PYG{n}{numpy}\PYG{o}{.}\PYG{n}{empty}\PYG{p}{(}\PYG{n}{n\PYGZus{}ind}\PYG{p}{,} \PYG{n}{dtype}\PYG{o}{=}\PYG{n+nb}{float}\PYG{p}{)}
    \PYG{k}{for} \PYG{n}{i} \PYG{o+ow}{in} \PYG{n+nb}{range}\PYG{p}{(} \PYG{n}{n\PYGZus{}ind}  \PYG{p}{)} \PYG{p}{:}
        \PYG{n}{x}\PYG{p}{[}\PYG{n}{i}\PYG{p}{]} \PYG{o}{=} \PYG{n}{i} \PYG{o}{+} \PYG{l+m+mf}{2.0}
    \PYG{c+c1}{\PYGZsh{}}
    \PYG{n}{ax} \PYG{o}{=} \PYG{n}{cppad\PYGZus{}py}\PYG{o}{.}\PYG{n}{independent}\PYG{p}{(}\PYG{n}{x}\PYG{p}{)}
    \PYG{c+c1}{\PYGZsh{}}
    \PYG{c+c1}{\PYGZsh{} create dependent variables ay with ay0 = ax\PYGZus{}0 * ax\PYGZus{}1 * ax\PYGZus{}2}
    \PYG{n}{ax\PYGZus{}0}  \PYG{o}{=} \PYG{n}{ax}\PYG{p}{[}\PYG{l+m+mi}{0}\PYG{p}{]}
    \PYG{n}{ax\PYGZus{}1}  \PYG{o}{=} \PYG{n}{ax}\PYG{p}{[}\PYG{l+m+mi}{1}\PYG{p}{]}
    \PYG{n}{ax\PYGZus{}2}  \PYG{o}{=} \PYG{n}{ax}\PYG{p}{[}\PYG{l+m+mi}{2}\PYG{p}{]}
    \PYG{n}{ay}    \PYG{o}{=} \PYG{n}{numpy}\PYG{o}{.}\PYG{n}{empty}\PYG{p}{(}\PYG{n}{n\PYGZus{}dep}\PYG{p}{,} \PYG{n}{dtype}\PYG{o}{=}\PYG{n}{cppad\PYGZus{}py}\PYG{o}{.}\PYG{n}{a\PYGZus{}double}\PYG{p}{)}
    \PYG{n}{ay}\PYG{p}{[}\PYG{l+m+mi}{0}\PYG{p}{]} \PYG{o}{=} \PYG{n}{ax\PYGZus{}0} \PYG{o}{*} \PYG{n}{ax\PYGZus{}1} \PYG{o}{*} \PYG{n}{ax\PYGZus{}2}
    \PYG{c+c1}{\PYGZsh{}}
    \PYG{c+c1}{\PYGZsh{} define af corresponding to f(x) = x\PYGZus{}0 * x\PYGZus{}1 * x\PYGZus{}2}
    \PYG{n}{f}  \PYG{o}{=} \PYG{n}{cppad\PYGZus{}py}\PYG{o}{.}\PYG{n}{d\PYGZus{}fun}\PYG{p}{(}\PYG{n}{ax}\PYG{p}{,} \PYG{n}{ay}\PYG{p}{)}
    \PYG{c+c1}{\PYGZsh{}}
    \PYG{c+c1}{\PYGZsh{} compute the Jacobian f\PYGZsq{}(x) = ( x\PYGZus{}1*x\PYGZus{}2, x\PYGZus{}0*x\PYGZus{}2, x\PYGZus{}0*x\PYGZus{}1 )}
    \PYG{n}{fp} \PYG{o}{=} \PYG{n}{f}\PYG{o}{.}\PYG{n}{jacobian}\PYG{p}{(}\PYG{n}{x}\PYG{p}{)}
    \PYG{c+c1}{\PYGZsh{}}
    \PYG{c+c1}{\PYGZsh{} check Jacobian}
    \PYG{n}{x\PYGZus{}0} \PYG{o}{=} \PYG{n}{x}\PYG{p}{[}\PYG{l+m+mi}{0}\PYG{p}{]}
    \PYG{n}{x\PYGZus{}1} \PYG{o}{=} \PYG{n}{x}\PYG{p}{[}\PYG{l+m+mi}{1}\PYG{p}{]}
    \PYG{n}{x\PYGZus{}2} \PYG{o}{=} \PYG{n}{x}\PYG{p}{[}\PYG{l+m+mi}{2}\PYG{p}{]}
    \PYG{n}{ok} \PYG{o}{=} \PYG{n}{ok} \PYG{o+ow}{and} \PYG{n}{fp}\PYG{p}{[}\PYG{l+m+mi}{0}\PYG{p}{,} \PYG{l+m+mi}{0}\PYG{p}{]} \PYG{o}{==} \PYG{n}{x\PYGZus{}1} \PYG{o}{*} \PYG{n}{x\PYGZus{}2}
    \PYG{n}{ok} \PYG{o}{=} \PYG{n}{ok} \PYG{o+ow}{and} \PYG{n}{fp}\PYG{p}{[}\PYG{l+m+mi}{0}\PYG{p}{,} \PYG{l+m+mi}{1}\PYG{p}{]} \PYG{o}{==} \PYG{n}{x\PYGZus{}0} \PYG{o}{*} \PYG{n}{x\PYGZus{}2}
    \PYG{n}{ok} \PYG{o}{=} \PYG{n}{ok} \PYG{o+ow}{and} \PYG{n}{fp}\PYG{p}{[}\PYG{l+m+mi}{0}\PYG{p}{,} \PYG{l+m+mi}{2}\PYG{p}{]} \PYG{o}{==} \PYG{n}{x\PYGZus{}0} \PYG{o}{*} \PYG{n}{x\PYGZus{}1}
    \PYG{c+c1}{\PYGZsh{} \PYGZhy{}\PYGZhy{}\PYGZhy{}\PYGZhy{}\PYGZhy{}\PYGZhy{}\PYGZhy{}\PYGZhy{}\PYGZhy{}\PYGZhy{}\PYGZhy{}\PYGZhy{}\PYGZhy{}\PYGZhy{}\PYGZhy{}\PYGZhy{}\PYGZhy{}\PYGZhy{}\PYGZhy{}\PYGZhy{}\PYGZhy{}\PYGZhy{}\PYGZhy{}\PYGZhy{}\PYGZhy{}\PYGZhy{}\PYGZhy{}\PYGZhy{}\PYGZhy{}\PYGZhy{}\PYGZhy{}\PYGZhy{}\PYGZhy{}\PYGZhy{}\PYGZhy{}\PYGZhy{}\PYGZhy{}\PYGZhy{}\PYGZhy{}\PYGZhy{}\PYGZhy{}\PYGZhy{}\PYGZhy{}\PYGZhy{}\PYGZhy{}\PYGZhy{}\PYGZhy{}\PYGZhy{}\PYGZhy{}\PYGZhy{}\PYGZhy{}\PYGZhy{}\PYGZhy{}\PYGZhy{}\PYGZhy{}\PYGZhy{}\PYGZhy{}\PYGZhy{}\PYGZhy{}\PYGZhy{}\PYGZhy{}\PYGZhy{}\PYGZhy{}\PYGZhy{}\PYGZhy{}\PYGZhy{}\PYGZhy{}\PYGZhy{}\PYGZhy{}}
    \PYG{n}{af} \PYG{o}{=} \PYG{n}{cppad\PYGZus{}py}\PYG{o}{.}\PYG{n}{a\PYGZus{}fun}\PYG{p}{(}\PYG{n}{f}\PYG{p}{)}
    \PYG{c+c1}{\PYGZsh{}}
    \PYG{n}{ax}   \PYG{o}{=} \PYG{n}{numpy}\PYG{o}{.}\PYG{n}{empty}\PYG{p}{(}\PYG{n}{n\PYGZus{}ind}\PYG{p}{,} \PYG{n}{dtype}\PYG{o}{=}\PYG{n}{cppad\PYGZus{}py}\PYG{o}{.}\PYG{n}{a\PYGZus{}double}\PYG{p}{)}
    \PYG{k}{for} \PYG{n}{i} \PYG{o+ow}{in} \PYG{n+nb}{range}\PYG{p}{(} \PYG{n}{n\PYGZus{}ind} \PYG{p}{)} \PYG{p}{:}
        \PYG{n}{ax}\PYG{p}{[}\PYG{n}{i}\PYG{p}{]} \PYG{o}{=} \PYG{n}{x}\PYG{p}{[}\PYG{n}{i}\PYG{p}{]}
    \PYG{c+c1}{\PYGZsh{}}
    \PYG{c+c1}{\PYGZsh{} compute the Jacobian f\PYGZsq{}(x) = ( x\PYGZus{}1*x\PYGZus{}2, x\PYGZus{}0*x\PYGZus{}2, x\PYGZus{}0*x\PYGZus{}1 )}
    \PYG{n}{afp} \PYG{o}{=} \PYG{n}{af}\PYG{o}{.}\PYG{n}{jacobian}\PYG{p}{(}\PYG{n}{ax}\PYG{p}{)}
    \PYG{c+c1}{\PYGZsh{}}
    \PYG{c+c1}{\PYGZsh{} check Jacobian}
    \PYG{n}{ok} \PYG{o}{=} \PYG{n}{ok} \PYG{o+ow}{and} \PYG{n}{afp}\PYG{p}{[}\PYG{l+m+mi}{0}\PYG{p}{,} \PYG{l+m+mi}{0}\PYG{p}{]} \PYG{o}{==} \PYG{n}{cppad\PYGZus{}py}\PYG{o}{.}\PYG{n}{a\PYGZus{}double}\PYG{p}{(}\PYG{n}{x\PYGZus{}1} \PYG{o}{*} \PYG{n}{x\PYGZus{}2}\PYG{p}{)}
    \PYG{n}{ok} \PYG{o}{=} \PYG{n}{ok} \PYG{o+ow}{and} \PYG{n}{afp}\PYG{p}{[}\PYG{l+m+mi}{0}\PYG{p}{,} \PYG{l+m+mi}{1}\PYG{p}{]} \PYG{o}{==} \PYG{n}{cppad\PYGZus{}py}\PYG{o}{.}\PYG{n}{a\PYGZus{}double}\PYG{p}{(}\PYG{n}{x\PYGZus{}0} \PYG{o}{*} \PYG{n}{x\PYGZus{}2}\PYG{p}{)}
    \PYG{n}{ok} \PYG{o}{=} \PYG{n}{ok} \PYG{o+ow}{and} \PYG{n}{afp}\PYG{p}{[}\PYG{l+m+mi}{0}\PYG{p}{,} \PYG{l+m+mi}{2}\PYG{p}{]} \PYG{o}{==} \PYG{n}{cppad\PYGZus{}py}\PYG{o}{.}\PYG{n}{a\PYGZus{}double}\PYG{p}{(}\PYG{n}{x\PYGZus{}0} \PYG{o}{*} \PYG{n}{x\PYGZus{}1}\PYG{p}{)}
    \PYG{c+c1}{\PYGZsh{}}
    \PYG{k}{return}\PYG{p}{(} \PYG{n}{ok} \PYG{p}{)}
\PYG{c+c1}{\PYGZsh{}}
\end{sphinxVerbatim}


\bigskip\hrule\bigskip


xsrst input file: \sphinxcode{\sphinxupquote{example/python/core/fun\_jacobian\_xam.py}}

\index{example@\spxentry{example}}\ignorespaces 

\subparagraph{Example}
\label{\detokenize{xsrst/py_fun_jacobian:example}}\label{\detokenize{xsrst/py_fun_jacobian:py-fun-jacobian-example}}\label{\detokenize{xsrst/py_fun_jacobian:index-6}}
{\hyperref[\detokenize{xsrst/fun_jacobian_xam_py:id1}]{\sphinxcrossref{\DUrole{std,std-ref}{fun\_jacobian\_xam\_py}}}}


\bigskip\hrule\bigskip


xsrst input file: \sphinxcode{\sphinxupquote{lib/python/cppad\_py/fun\_jacobian.py}}


\subparagraph{py\_fun\_hessian}
\label{\detokenize{xsrst/py_fun_hessian:py-fun-hessian}}\label{\detokenize{xsrst/py_fun_hessian::doc}}
\index{py\_fun\_hessian@\spxentry{py\_fun\_hessian}}\index{hessian@\spxentry{hessian}}\index{an@\spxentry{an}}\index{ad@\spxentry{ad}}\index{function@\spxentry{function}}\ignorespaces 

\subparagraph{3.1.1.7 Hessian of an AD Function}
\label{\detokenize{xsrst/py_fun_hessian:hessian-of-an-ad-function}}\label{\detokenize{xsrst/py_fun_hessian:id1}}\begin{itemize}
\item {} 
{\hyperref[\detokenize{xsrst/py_fun_hessian:py-fun-hessian-syntax}]{\sphinxcrossref{\DUrole{std,std-ref}{Syntax}}}}

\item {} 
{\hyperref[\detokenize{xsrst/py_fun_hessian:py-fun-hessian-f}]{\sphinxcrossref{\DUrole{std,std-ref}{f}}}}

\item {} 
{\hyperref[\detokenize{xsrst/py_fun_hessian:py-fun-hessian-f-x}]{\sphinxcrossref{\DUrole{std,std-ref}{f(x)}}}}

\item {} 
{\hyperref[\detokenize{xsrst/py_fun_hessian:py-fun-hessian-g-x}]{\sphinxcrossref{\DUrole{std,std-ref}{g(x)}}}}

\item {} 
{\hyperref[\detokenize{xsrst/py_fun_hessian:py-fun-hessian-x}]{\sphinxcrossref{\DUrole{std,std-ref}{x}}}}

\item {} 
{\hyperref[\detokenize{xsrst/py_fun_hessian:py-fun-hessian-w}]{\sphinxcrossref{\DUrole{std,std-ref}{w}}}}

\item {} 
{\hyperref[\detokenize{xsrst/py_fun_hessian:py-fun-hessian-h}]{\sphinxcrossref{\DUrole{std,std-ref}{H}}}}

\item {} 
{\hyperref[\detokenize{xsrst/py_fun_hessian:py-fun-hessian-example}]{\sphinxcrossref{\DUrole{std,std-ref}{Example}}}}

\end{itemize}

\index{syntax@\spxentry{syntax}}\ignorespaces 

\subparagraph{Syntax}
\label{\detokenize{xsrst/py_fun_hessian:syntax}}\label{\detokenize{xsrst/py_fun_hessian:py-fun-hessian-syntax}}\label{\detokenize{xsrst/py_fun_hessian:index-1}}
\sphinxstyleemphasis{H} = \sphinxstyleemphasis{f} . \sphinxcode{\sphinxupquote{hessian}} ( \sphinxstyleemphasis{x} , \sphinxstyleemphasis{w} )

\index{f@\spxentry{f}}\ignorespaces 

\subparagraph{f}
\label{\detokenize{xsrst/py_fun_hessian:f}}\label{\detokenize{xsrst/py_fun_hessian:py-fun-hessian-f}}\label{\detokenize{xsrst/py_fun_hessian:index-2}}
This is either a
{\hyperref[\detokenize{xsrst/py_fun_ctor:py-fun-ctor-syntax-d-fun}]{\sphinxcrossref{\DUrole{std,std-ref}{d\_fun}}}} or
{\hyperref[\detokenize{xsrst/py_fun_ctor:py-fun-ctor-syntax-a-fun}]{\sphinxcrossref{\DUrole{std,std-ref}{a\_fun}}}} function object.
Upon return, the zero order
{\hyperref[\detokenize{xsrst/py_fun_forward:py-fun-forward-taylor-coefficient}]{\sphinxcrossref{\DUrole{std,std-ref}{taylor\_coefficients}}}}
in \sphinxstyleemphasis{f} correspond to the value of \sphinxstyleemphasis{x} .
The other Taylor coefficients in \sphinxstyleemphasis{f} are unspecified.

\index{f(x)@\spxentry{f(x)}}\ignorespaces 

\subparagraph{f(x)}
\label{\detokenize{xsrst/py_fun_hessian:f-x}}\label{\detokenize{xsrst/py_fun_hessian:py-fun-hessian-f-x}}\label{\detokenize{xsrst/py_fun_hessian:index-3}}
We use the notation \(f: \B{R}^n \rightarrow \B{R}^m\)
for the function corresponding to \sphinxstyleemphasis{f} .
Note that \sphinxstyleemphasis{n} is the size of {\hyperref[\detokenize{xsrst/py_fun_ctor:py-fun-ctor-ax}]{\sphinxcrossref{\DUrole{std,std-ref}{ax}}}}
and \sphinxstyleemphasis{m} is the size of {\hyperref[\detokenize{xsrst/py_fun_ctor:py-fun-ctor-ay}]{\sphinxcrossref{\DUrole{std,std-ref}{ay}}}}
in to the constructor for \sphinxstyleemphasis{f} .

\index{g(x)@\spxentry{g(x)}}\ignorespaces 

\subparagraph{g(x)}
\label{\detokenize{xsrst/py_fun_hessian:g-x}}\label{\detokenize{xsrst/py_fun_hessian:py-fun-hessian-g-x}}\label{\detokenize{xsrst/py_fun_hessian:index-4}}
We use the notation \(g: \B{R}^n \rightarrow \B{R}\)
for the function defined by
\begin{equation*}
\begin{split}g(x) = \sum_{i=0}^{n-1} w_i f_i (x)\end{split}
\end{equation*}
\index{x@\spxentry{x}}\ignorespaces 

\subparagraph{x}
\label{\detokenize{xsrst/py_fun_hessian:x}}\label{\detokenize{xsrst/py_fun_hessian:py-fun-hessian-x}}\label{\detokenize{xsrst/py_fun_hessian:index-5}}
If \sphinxstyleemphasis{f} is a \sphinxcode{\sphinxupquote{d\_fun}} or \sphinxcode{\sphinxupquote{a\_fun}} ,
\sphinxstyleemphasis{x} is a numpy vector with \sphinxcode{\sphinxupquote{float}} ( \sphinxcode{\sphinxupquote{a\_double}} ) elements
and size \sphinxstyleemphasis{n} .
It specifies the argument value at we are computing the Hessian
\(g^{(2)}(x)\).

\index{w@\spxentry{w}}\ignorespaces 

\subparagraph{w}
\label{\detokenize{xsrst/py_fun_hessian:w}}\label{\detokenize{xsrst/py_fun_hessian:py-fun-hessian-w}}\label{\detokenize{xsrst/py_fun_hessian:index-6}}
If \sphinxstyleemphasis{f} is a \sphinxcode{\sphinxupquote{d\_fun}} or \sphinxcode{\sphinxupquote{a\_fun}} ,
\sphinxstyleemphasis{w} is a numpy vector with \sphinxcode{\sphinxupquote{float}} ( \sphinxcode{\sphinxupquote{a\_double}} ) elements
and size \sphinxstyleemphasis{m} .
It specifies the vector \sphinxstyleemphasis{w} in the definition of \(g(x)\) above.

\index{h@\spxentry{h}}\ignorespaces 

\subparagraph{H}
\label{\detokenize{xsrst/py_fun_hessian:h}}\label{\detokenize{xsrst/py_fun_hessian:py-fun-hessian-h}}\label{\detokenize{xsrst/py_fun_hessian:index-7}}
The result is a numpy matrix with \sphinxcode{\sphinxupquote{float}} ( \sphinxcode{\sphinxupquote{a\_double}} ) elements,
\sphinxstyleemphasis{n} rows and \sphinxcode{\sphinxupquote{n}} columns.
For \sphinxstyleemphasis{i} between zero and \sphinxstyleemphasis{n} \sphinxhyphen{}1
and \sphinxstyleemphasis{j} between zero and \sphinxstyleemphasis{n} \sphinxhyphen{}1 ,
\begin{equation*}
\begin{split}H [ i, j ] = \frac{ \partial^2 g }{ \partial x_i \partial x_j } (x)\end{split}
\end{equation*}

\subparagraph{fun\_hessian\_xam\_py}
\label{\detokenize{xsrst/fun_hessian_xam_py:fun-hessian-xam-py}}\label{\detokenize{xsrst/fun_hessian_xam_py::doc}}
\index{fun\_hessian\_xam\_py@\spxentry{fun\_hessian\_xam\_py}}\index{python:@\spxentry{python:}}\index{dense@\spxentry{dense}}\index{hessian@\spxentry{hessian}}\index{using@\spxentry{using}}\index{ad:@\spxentry{ad:}}\index{example@\spxentry{example}}\index{test@\spxentry{test}}\ignorespaces 

\subparagraph{3.1.1.7.1 Python: Dense Hessian Using AD: Example and Test}
\label{\detokenize{xsrst/fun_hessian_xam_py:python-dense-hessian-using-ad-example-and-test}}\label{\detokenize{xsrst/fun_hessian_xam_py:id1}}
\begin{sphinxVerbatim}[commandchars=\\\{\}]
\PYG{k}{def} \PYG{n+nf}{fun\PYGZus{}hessian\PYGZus{}xam}\PYG{p}{(}\PYG{p}{)} \PYG{p}{:}
    \PYG{c+c1}{\PYGZsh{}}
    \PYG{k+kn}{import} \PYG{n+nn}{numpy}
    \PYG{k+kn}{import} \PYG{n+nn}{cppad\PYGZus{}py}
    \PYG{c+c1}{\PYGZsh{}}
    \PYG{c+c1}{\PYGZsh{} initialize return variable}
    \PYG{n}{ok} \PYG{o}{=} \PYG{k+kc}{True}
    \PYG{c+c1}{\PYGZsh{} \PYGZhy{}\PYGZhy{}\PYGZhy{}\PYGZhy{}\PYGZhy{}\PYGZhy{}\PYGZhy{}\PYGZhy{}\PYGZhy{}\PYGZhy{}\PYGZhy{}\PYGZhy{}\PYGZhy{}\PYGZhy{}\PYGZhy{}\PYGZhy{}\PYGZhy{}\PYGZhy{}\PYGZhy{}\PYGZhy{}\PYGZhy{}\PYGZhy{}\PYGZhy{}\PYGZhy{}\PYGZhy{}\PYGZhy{}\PYGZhy{}\PYGZhy{}\PYGZhy{}\PYGZhy{}\PYGZhy{}\PYGZhy{}\PYGZhy{}\PYGZhy{}\PYGZhy{}\PYGZhy{}\PYGZhy{}\PYGZhy{}\PYGZhy{}\PYGZhy{}\PYGZhy{}\PYGZhy{}\PYGZhy{}\PYGZhy{}\PYGZhy{}\PYGZhy{}\PYGZhy{}\PYGZhy{}\PYGZhy{}\PYGZhy{}\PYGZhy{}\PYGZhy{}\PYGZhy{}\PYGZhy{}\PYGZhy{}\PYGZhy{}\PYGZhy{}\PYGZhy{}\PYGZhy{}\PYGZhy{}\PYGZhy{}\PYGZhy{}\PYGZhy{}\PYGZhy{}\PYGZhy{}\PYGZhy{}\PYGZhy{}\PYGZhy{}\PYGZhy{}}
    \PYG{c+c1}{\PYGZsh{} number of dependent and independent variables}
    \PYG{n}{n\PYGZus{}dep} \PYG{o}{=} \PYG{l+m+mi}{1}
    \PYG{n}{n\PYGZus{}ind} \PYG{o}{=} \PYG{l+m+mi}{3}
    \PYG{c+c1}{\PYGZsh{}}
    \PYG{c+c1}{\PYGZsh{} create the independent variables ax}
    \PYG{n}{x} \PYG{o}{=} \PYG{n}{numpy}\PYG{o}{.}\PYG{n}{empty}\PYG{p}{(}\PYG{n}{n\PYGZus{}ind}\PYG{p}{,} \PYG{n}{dtype}\PYG{o}{=}\PYG{n+nb}{float}\PYG{p}{)}
    \PYG{k}{for} \PYG{n}{i} \PYG{o+ow}{in} \PYG{n+nb}{range}\PYG{p}{(} \PYG{n}{n\PYGZus{}ind}  \PYG{p}{)} \PYG{p}{:}
        \PYG{n}{x}\PYG{p}{[}\PYG{n}{i}\PYG{p}{]} \PYG{o}{=} \PYG{n}{i} \PYG{o}{+} \PYG{l+m+mf}{2.0}
    \PYG{c+c1}{\PYGZsh{}}
    \PYG{n}{ax} \PYG{o}{=} \PYG{n}{cppad\PYGZus{}py}\PYG{o}{.}\PYG{n}{independent}\PYG{p}{(}\PYG{n}{x}\PYG{p}{)}
    \PYG{c+c1}{\PYGZsh{}}
    \PYG{c+c1}{\PYGZsh{} create dependent variables ay with ay0 = ax\PYGZus{}0 * ax\PYGZus{}1 * ax\PYGZus{}2}
    \PYG{n}{ax\PYGZus{}0} \PYG{o}{=} \PYG{n}{ax}\PYG{p}{[}\PYG{l+m+mi}{0}\PYG{p}{]}
    \PYG{n}{ax\PYGZus{}1} \PYG{o}{=} \PYG{n}{ax}\PYG{p}{[}\PYG{l+m+mi}{1}\PYG{p}{]}
    \PYG{n}{ax\PYGZus{}2} \PYG{o}{=} \PYG{n}{ax}\PYG{p}{[}\PYG{l+m+mi}{2}\PYG{p}{]}
    \PYG{n}{ay}   \PYG{o}{=} \PYG{n}{numpy}\PYG{o}{.}\PYG{n}{empty}\PYG{p}{(}\PYG{n}{n\PYGZus{}dep}\PYG{p}{,} \PYG{n}{dtype}\PYG{o}{=}\PYG{n}{cppad\PYGZus{}py}\PYG{o}{.}\PYG{n}{a\PYGZus{}double}\PYG{p}{)}
    \PYG{n}{ay}\PYG{p}{[}\PYG{l+m+mi}{0}\PYG{p}{]} \PYG{o}{=} \PYG{n}{ax\PYGZus{}0} \PYG{o}{*} \PYG{n}{ax\PYGZus{}1} \PYG{o}{*} \PYG{n}{ax\PYGZus{}2}
    \PYG{c+c1}{\PYGZsh{}}
    \PYG{c+c1}{\PYGZsh{} define af corresponding to f(x) = x\PYGZus{}0 * x\PYGZus{}1 * x\PYGZus{}2}
    \PYG{n}{f}  \PYG{o}{=} \PYG{n}{cppad\PYGZus{}py}\PYG{o}{.}\PYG{n}{d\PYGZus{}fun}\PYG{p}{(}\PYG{n}{ax}\PYG{p}{,} \PYG{n}{ay}\PYG{p}{)}
    \PYG{c+c1}{\PYGZsh{}}
    \PYG{c+c1}{\PYGZsh{} g(x) = w\PYGZus{}0 * f\PYGZus{}0 (x) = f(x)}
    \PYG{n}{w} \PYG{o}{=} \PYG{n}{numpy}\PYG{o}{.}\PYG{n}{empty}\PYG{p}{(}\PYG{n}{n\PYGZus{}dep}\PYG{p}{,} \PYG{n}{dtype}\PYG{o}{=}\PYG{n+nb}{float}\PYG{p}{)}
    \PYG{n}{w}\PYG{p}{[}\PYG{l+m+mi}{0}\PYG{p}{]} \PYG{o}{=} \PYG{l+m+mf}{1.0}
    \PYG{c+c1}{\PYGZsh{}}
    \PYG{c+c1}{\PYGZsh{} compute Hessian}
    \PYG{n}{fpp} \PYG{o}{=} \PYG{n}{f}\PYG{o}{.}\PYG{n}{hessian}\PYG{p}{(}\PYG{n}{x}\PYG{p}{,} \PYG{n}{w}\PYG{p}{)}
    \PYG{c+c1}{\PYGZsh{}}
    \PYG{c+c1}{\PYGZsh{}          [ 0.0 , x\PYGZus{}2 , x\PYGZus{}1 ]}
    \PYG{c+c1}{\PYGZsh{} f\PYGZsq{}\PYGZsq{}(x) = [ x\PYGZus{}2 , 0.0 , x\PYGZus{}0 ]}
    \PYG{c+c1}{\PYGZsh{}          [ x\PYGZus{}1 , x\PYGZus{}0 , 0.0 ]}
    \PYG{n}{ok} \PYG{o}{=} \PYG{n}{ok} \PYG{o+ow}{and} \PYG{n}{fpp}\PYG{p}{[}\PYG{l+m+mi}{0}\PYG{p}{,} \PYG{l+m+mi}{0}\PYG{p}{]} \PYG{o}{==} \PYG{l+m+mf}{0.0}
    \PYG{n}{ok} \PYG{o}{=} \PYG{n}{ok} \PYG{o+ow}{and} \PYG{n}{fpp}\PYG{p}{[}\PYG{l+m+mi}{0}\PYG{p}{,} \PYG{l+m+mi}{1}\PYG{p}{]} \PYG{o}{==} \PYG{n}{x}\PYG{p}{[}\PYG{l+m+mi}{2}\PYG{p}{]}
    \PYG{n}{ok} \PYG{o}{=} \PYG{n}{ok} \PYG{o+ow}{and} \PYG{n}{fpp}\PYG{p}{[}\PYG{l+m+mi}{0}\PYG{p}{,} \PYG{l+m+mi}{2}\PYG{p}{]} \PYG{o}{==} \PYG{n}{x}\PYG{p}{[}\PYG{l+m+mi}{1}\PYG{p}{]}
    \PYG{c+c1}{\PYGZsh{}}
    \PYG{n}{ok} \PYG{o}{=} \PYG{n}{ok} \PYG{o+ow}{and} \PYG{n}{fpp}\PYG{p}{[}\PYG{l+m+mi}{1}\PYG{p}{,} \PYG{l+m+mi}{0}\PYG{p}{]} \PYG{o}{==} \PYG{n}{x}\PYG{p}{[}\PYG{l+m+mi}{2}\PYG{p}{]}
    \PYG{n}{ok} \PYG{o}{=} \PYG{n}{ok} \PYG{o+ow}{and} \PYG{n}{fpp}\PYG{p}{[}\PYG{l+m+mi}{1}\PYG{p}{,} \PYG{l+m+mi}{1}\PYG{p}{]} \PYG{o}{==} \PYG{l+m+mf}{0.0}
    \PYG{n}{ok} \PYG{o}{=} \PYG{n}{ok} \PYG{o+ow}{and} \PYG{n}{fpp}\PYG{p}{[}\PYG{l+m+mi}{1}\PYG{p}{,} \PYG{l+m+mi}{2}\PYG{p}{]} \PYG{o}{==} \PYG{n}{x}\PYG{p}{[}\PYG{l+m+mi}{0}\PYG{p}{]}
    \PYG{c+c1}{\PYGZsh{}}
    \PYG{n}{ok} \PYG{o}{=} \PYG{n}{ok} \PYG{o+ow}{and} \PYG{n}{fpp}\PYG{p}{[}\PYG{l+m+mi}{2}\PYG{p}{,} \PYG{l+m+mi}{0}\PYG{p}{]} \PYG{o}{==} \PYG{n}{x}\PYG{p}{[}\PYG{l+m+mi}{1}\PYG{p}{]}
    \PYG{n}{ok} \PYG{o}{=} \PYG{n}{ok} \PYG{o+ow}{and} \PYG{n}{fpp}\PYG{p}{[}\PYG{l+m+mi}{2}\PYG{p}{,} \PYG{l+m+mi}{1}\PYG{p}{]} \PYG{o}{==} \PYG{n}{x}\PYG{p}{[}\PYG{l+m+mi}{0}\PYG{p}{]}
    \PYG{n}{ok} \PYG{o}{=} \PYG{n}{ok} \PYG{o+ow}{and} \PYG{n}{fpp}\PYG{p}{[}\PYG{l+m+mi}{2}\PYG{p}{,} \PYG{l+m+mi}{2}\PYG{p}{]} \PYG{o}{==} \PYG{l+m+mf}{0.0}
    \PYG{c+c1}{\PYGZsh{} \PYGZhy{}\PYGZhy{}\PYGZhy{}\PYGZhy{}\PYGZhy{}\PYGZhy{}\PYGZhy{}\PYGZhy{}\PYGZhy{}\PYGZhy{}\PYGZhy{}\PYGZhy{}\PYGZhy{}\PYGZhy{}\PYGZhy{}\PYGZhy{}\PYGZhy{}\PYGZhy{}\PYGZhy{}\PYGZhy{}\PYGZhy{}\PYGZhy{}\PYGZhy{}\PYGZhy{}\PYGZhy{}\PYGZhy{}\PYGZhy{}\PYGZhy{}\PYGZhy{}\PYGZhy{}\PYGZhy{}\PYGZhy{}\PYGZhy{}\PYGZhy{}\PYGZhy{}\PYGZhy{}\PYGZhy{}\PYGZhy{}\PYGZhy{}\PYGZhy{}\PYGZhy{}\PYGZhy{}\PYGZhy{}\PYGZhy{}\PYGZhy{}\PYGZhy{}\PYGZhy{}\PYGZhy{}\PYGZhy{}\PYGZhy{}\PYGZhy{}\PYGZhy{}\PYGZhy{}\PYGZhy{}\PYGZhy{}\PYGZhy{}\PYGZhy{}\PYGZhy{}\PYGZhy{}\PYGZhy{}\PYGZhy{}\PYGZhy{}\PYGZhy{}\PYGZhy{}\PYGZhy{}\PYGZhy{}\PYGZhy{}\PYGZhy{}\PYGZhy{}}
    \PYG{n}{af} \PYG{o}{=} \PYG{n}{cppad\PYGZus{}py}\PYG{o}{.}\PYG{n}{a\PYGZus{}fun}\PYG{p}{(}\PYG{n}{f}\PYG{p}{)}
    \PYG{c+c1}{\PYGZsh{}}
    \PYG{c+c1}{\PYGZsh{} compute and check Hessian}
    \PYG{n}{aw}    \PYG{o}{=} \PYG{n}{numpy}\PYG{o}{.}\PYG{n}{empty}\PYG{p}{(}\PYG{n}{n\PYGZus{}dep}\PYG{p}{,} \PYG{n}{dtype}\PYG{o}{=}\PYG{n}{cppad\PYGZus{}py}\PYG{o}{.}\PYG{n}{a\PYGZus{}double}\PYG{p}{)}
    \PYG{n}{aw}\PYG{p}{[}\PYG{l+m+mi}{0}\PYG{p}{]} \PYG{o}{=} \PYG{n}{w}\PYG{p}{[}\PYG{l+m+mi}{0}\PYG{p}{]}
    \PYG{n}{afpp}  \PYG{o}{=} \PYG{n}{af}\PYG{o}{.}\PYG{n}{hessian}\PYG{p}{(}\PYG{n}{ax}\PYG{p}{,} \PYG{n}{aw}\PYG{p}{)}
    \PYG{n}{ok}    \PYG{o}{=} \PYG{n}{ok} \PYG{o+ow}{and} \PYG{n}{afpp}\PYG{o}{.}\PYG{n}{shape} \PYG{o}{==} \PYG{n}{fpp}\PYG{o}{.}\PYG{n}{shape}
    \PYG{p}{(}\PYG{n}{nr}\PYG{p}{,} \PYG{n}{nc}\PYG{p}{)} \PYG{o}{=} \PYG{n}{fpp}\PYG{o}{.}\PYG{n}{shape}
    \PYG{k}{for} \PYG{n}{i} \PYG{o+ow}{in} \PYG{n+nb}{range}\PYG{p}{(}\PYG{n}{nr}\PYG{p}{)} \PYG{p}{:}
        \PYG{k}{for} \PYG{n}{j} \PYG{o+ow}{in} \PYG{n+nb}{range}\PYG{p}{(}\PYG{n}{nc}\PYG{p}{)} \PYG{p}{:}
            \PYG{n}{ok} \PYG{o}{=} \PYG{n}{ok} \PYG{o+ow}{and} \PYG{n}{afpp}\PYG{p}{[}\PYG{n}{i}\PYG{p}{,} \PYG{n}{j}\PYG{p}{]} \PYG{o}{==} \PYG{n}{cppad\PYGZus{}py}\PYG{o}{.}\PYG{n}{a\PYGZus{}double}\PYG{p}{(} \PYG{n}{fpp}\PYG{p}{[}\PYG{n}{i}\PYG{p}{,} \PYG{n}{j}\PYG{p}{]} \PYG{p}{)}
    \PYG{c+c1}{\PYGZsh{}}
    \PYG{k}{return}\PYG{p}{(} \PYG{n}{ok} \PYG{p}{)}
\PYG{c+c1}{\PYGZsh{}}
\end{sphinxVerbatim}


\bigskip\hrule\bigskip


xsrst input file: \sphinxcode{\sphinxupquote{example/python/core/fun\_hessian\_xam.py}}

\index{example@\spxentry{example}}\ignorespaces 

\subparagraph{Example}
\label{\detokenize{xsrst/py_fun_hessian:example}}\label{\detokenize{xsrst/py_fun_hessian:py-fun-hessian-example}}\label{\detokenize{xsrst/py_fun_hessian:index-8}}
{\hyperref[\detokenize{xsrst/fun_hessian_xam_py:id1}]{\sphinxcrossref{\DUrole{std,std-ref}{fun\_hessian\_xam\_py}}}}


\bigskip\hrule\bigskip


xsrst input file: \sphinxcode{\sphinxupquote{lib/python/cppad\_py/fun\_hessian.py}}


\subparagraph{py\_fun\_forward}
\label{\detokenize{xsrst/py_fun_forward:py-fun-forward}}\label{\detokenize{xsrst/py_fun_forward::doc}}
\index{py\_fun\_forward@\spxentry{py\_fun\_forward}}\index{forward@\spxentry{forward}}\index{mode@\spxentry{mode}}\index{ad@\spxentry{ad}}\ignorespaces 

\subparagraph{3.1.1.8 Forward Mode AD}
\label{\detokenize{xsrst/py_fun_forward:forward-mode-ad}}\label{\detokenize{xsrst/py_fun_forward:id1}}\begin{itemize}
\item {} 
{\hyperref[\detokenize{xsrst/py_fun_forward:py-fun-forward-syntax}]{\sphinxcrossref{\DUrole{std,std-ref}{Syntax}}}}

\item {} 
{\hyperref[\detokenize{xsrst/py_fun_forward:py-fun-forward-taylor-coefficient}]{\sphinxcrossref{\DUrole{std,std-ref}{Taylor Coefficient}}}}

\item {} 
{\hyperref[\detokenize{xsrst/py_fun_forward:py-fun-forward-f}]{\sphinxcrossref{\DUrole{std,std-ref}{f}}}}

\item {} 
{\hyperref[\detokenize{xsrst/py_fun_forward:py-fun-forward-f-x}]{\sphinxcrossref{\DUrole{std,std-ref}{f(x)}}}}

\item {} 
{\hyperref[\detokenize{xsrst/py_fun_forward:py-fun-forward-x-t}]{\sphinxcrossref{\DUrole{std,std-ref}{X(t)}}}}

\item {} 
{\hyperref[\detokenize{xsrst/py_fun_forward:py-fun-forward-y-t}]{\sphinxcrossref{\DUrole{std,std-ref}{Y(t)}}}}

\item {} \begin{description}
\item[{{\hyperref[\detokenize{xsrst/py_fun_forward:py-fun-forward-p}]{\sphinxcrossref{\DUrole{std,std-ref}{p}}}}}] \leavevmode\begin{itemize}
\item {} 
{\hyperref[\detokenize{xsrst/py_fun_forward:py-fun-forward-p-size-order}]{\sphinxcrossref{\DUrole{std,std-ref}{size\_order}}}}

\end{itemize}

\end{description}

\item {} 
{\hyperref[\detokenize{xsrst/py_fun_forward:py-fun-forward-xp}]{\sphinxcrossref{\DUrole{std,std-ref}{xp}}}}

\item {} 
{\hyperref[\detokenize{xsrst/py_fun_forward:py-fun-forward-yp}]{\sphinxcrossref{\DUrole{std,std-ref}{yp}}}}

\item {} 
{\hyperref[\detokenize{xsrst/py_fun_forward:py-fun-forward-example}]{\sphinxcrossref{\DUrole{std,std-ref}{Example}}}}

\end{itemize}

\index{syntax@\spxentry{syntax}}\ignorespaces 

\subparagraph{Syntax}
\label{\detokenize{xsrst/py_fun_forward:syntax}}\label{\detokenize{xsrst/py_fun_forward:py-fun-forward-syntax}}\label{\detokenize{xsrst/py_fun_forward:index-1}}
\sphinxstyleemphasis{yp} = \sphinxstyleemphasis{f} . \sphinxcode{\sphinxupquote{forward}} ( \sphinxstyleemphasis{p} , \sphinxstyleemphasis{xp} )

\index{taylor@\spxentry{taylor}}\index{coefficient@\spxentry{coefficient}}\ignorespaces 

\subparagraph{Taylor Coefficient}
\label{\detokenize{xsrst/py_fun_forward:taylor-coefficient}}\label{\detokenize{xsrst/py_fun_forward:py-fun-forward-taylor-coefficient}}\label{\detokenize{xsrst/py_fun_forward:index-2}}
For a function \(g(t)\) of a scalar argument \(t \in \B{R}\),
the \sphinxstyleemphasis{p}\sphinxhyphen{}th order Taylor coefficient is its
\sphinxcode{\sphinxupquote{p}} \sphinxhyphen{}th order derivative divided by \sphinxstyleemphasis{p} factorial
and evaluated at \(t = 0\); i.e.,
\begin{equation*}
\begin{split}g^{(p)} (0) /  p !\end{split}
\end{equation*}
\index{f@\spxentry{f}}\ignorespaces 

\subparagraph{f}
\label{\detokenize{xsrst/py_fun_forward:f}}\label{\detokenize{xsrst/py_fun_forward:py-fun-forward-f}}\label{\detokenize{xsrst/py_fun_forward:index-3}}
This is either a
{\hyperref[\detokenize{xsrst/py_fun_ctor:py-fun-ctor-syntax-d-fun}]{\sphinxcrossref{\DUrole{std,std-ref}{d\_fun}}}} or
{\hyperref[\detokenize{xsrst/py_fun_ctor:py-fun-ctor-syntax-a-fun}]{\sphinxcrossref{\DUrole{std,std-ref}{a\_fun}}}} function object.
Note that its state is changed by this operation because
all the Taylor coefficient that it calculates for every
variable in recording are stored.
See more discussion of this fact under the heading
{\hyperref[\detokenize{xsrst/py_fun_forward:py-fun-forward-p}]{\sphinxcrossref{\DUrole{std,std-ref}{p}}}} below.

\index{f(x)@\spxentry{f(x)}}\ignorespaces 

\subparagraph{f(x)}
\label{\detokenize{xsrst/py_fun_forward:f-x}}\label{\detokenize{xsrst/py_fun_forward:py-fun-forward-f-x}}\label{\detokenize{xsrst/py_fun_forward:index-4}}
We use the notation \(f: \B{R}^n \rightarrow \B{R}^m\)
for the function corresponding to \sphinxstyleemphasis{f} .
Note that \sphinxstyleemphasis{n} is the size of {\hyperref[\detokenize{xsrst/py_fun_ctor:py-fun-ctor-ax}]{\sphinxcrossref{\DUrole{std,std-ref}{ax}}}}
and \sphinxstyleemphasis{m} is the size of {\hyperref[\detokenize{xsrst/py_fun_ctor:py-fun-ctor-ay}]{\sphinxcrossref{\DUrole{std,std-ref}{ay}}}}
in to the constructor for \sphinxstyleemphasis{f} .

\index{x(t)@\spxentry{x(t)}}\ignorespaces 

\subparagraph{X(t)}
\label{\detokenize{xsrst/py_fun_forward:x-t}}\label{\detokenize{xsrst/py_fun_forward:py-fun-forward-x-t}}\label{\detokenize{xsrst/py_fun_forward:index-5}}
We use the notation \(X : \B{R} \rightarrow \B{R}^n\)
for a function that the calling routine chooses.

\index{y(t)@\spxentry{y(t)}}\ignorespaces 

\subparagraph{Y(t)}
\label{\detokenize{xsrst/py_fun_forward:y-t}}\label{\detokenize{xsrst/py_fun_forward:py-fun-forward-y-t}}\label{\detokenize{xsrst/py_fun_forward:index-6}}
We define the function \(Y : \B{R} \rightarrow \B{R}^n\)
by \(Y(t) = f(X(t))\).

\index{p@\spxentry{p}}\ignorespaces 

\subparagraph{p}
\label{\detokenize{xsrst/py_fun_forward:p}}\label{\detokenize{xsrst/py_fun_forward:py-fun-forward-p}}\label{\detokenize{xsrst/py_fun_forward:index-7}}
This argument has type \sphinxcode{\sphinxupquote{int}} and is non\sphinxhyphen{}negative.
Its value is the order of the Taylor coefficient being calculated.
If there was no call to \sphinxcode{\sphinxupquote{forward}} for this \sphinxstyleemphasis{f} ,
the value of \sphinxstyleemphasis{p} must be zero.
Otherwise, it must be between zero and one greater that its
value for the previous call using this \sphinxstyleemphasis{f} .
After this call, the Taylor coefficients for orders zero though \sphinxstyleemphasis{p} ,
and for every variable in the recording, will be stored in \sphinxstyleemphasis{f} .

\index{size\_order@\spxentry{size\_order}}\ignorespaces 

\subparagraph{size\_order}
\label{\detokenize{xsrst/py_fun_forward:size-order}}\label{\detokenize{xsrst/py_fun_forward:py-fun-forward-p-size-order}}\label{\detokenize{xsrst/py_fun_forward:index-8}}
After this call,
{\hyperref[\detokenize{xsrst/py_fun_property:py-fun-property-size-order}]{\sphinxcrossref{\DUrole{std,std-ref}{f\_size\_order()}}}} is \sphinxstyleemphasis{p} +1 .

\index{xp@\spxentry{xp}}\ignorespaces 

\subparagraph{xp}
\label{\detokenize{xsrst/py_fun_forward:xp}}\label{\detokenize{xsrst/py_fun_forward:py-fun-forward-xp}}\label{\detokenize{xsrst/py_fun_forward:index-9}}
If \sphinxstyleemphasis{f} is a \sphinxcode{\sphinxupquote{d\_fun}} ( \sphinxcode{\sphinxupquote{a\_fun}} ) object, \sphinxstyleemphasis{xp}
is a numpy vector with \sphinxcode{\sphinxupquote{float}} ( \sphinxcode{\sphinxupquote{a\_double}} ) elements
and size \sphinxstyleemphasis{n} .
It specifies the \sphinxstyleemphasis{p}\sphinxhyphen{}th order Taylor coefficients for \sphinxstyleemphasis{X(t} ) .

\index{yp@\spxentry{yp}}\ignorespaces 

\subparagraph{yp}
\label{\detokenize{xsrst/py_fun_forward:yp}}\label{\detokenize{xsrst/py_fun_forward:py-fun-forward-yp}}\label{\detokenize{xsrst/py_fun_forward:index-10}}
The result is a numpy vector with \sphinxcode{\sphinxupquote{float}} ( \sphinxcode{\sphinxupquote{a\_double}} ) elements
and size \sphinxstyleemphasis{m} .
It is the \sphinxstyleemphasis{p}\sphinxhyphen{}th order Taylor coefficients for \(Y(t)\).


\subparagraph{fun\_forward\_xam\_py}
\label{\detokenize{xsrst/fun_forward_xam_py:fun-forward-xam-py}}\label{\detokenize{xsrst/fun_forward_xam_py::doc}}
\index{fun\_forward\_xam\_py@\spxentry{fun\_forward\_xam\_py}}\index{python:@\spxentry{python:}}\index{forward@\spxentry{forward}}\index{mode@\spxentry{mode}}\index{ad:@\spxentry{ad:}}\index{example@\spxentry{example}}\index{test@\spxentry{test}}\ignorespaces 

\subparagraph{3.1.1.8.1 Python: Forward Mode AD: Example and Test}
\label{\detokenize{xsrst/fun_forward_xam_py:python-forward-mode-ad-example-and-test}}\label{\detokenize{xsrst/fun_forward_xam_py:id1}}
\begin{sphinxVerbatim}[commandchars=\\\{\}]
\PYG{k}{def} \PYG{n+nf}{fun\PYGZus{}forward\PYGZus{}xam}\PYG{p}{(}\PYG{p}{)} \PYG{p}{:}
    \PYG{c+c1}{\PYGZsh{}}
    \PYG{k+kn}{import} \PYG{n+nn}{numpy}
    \PYG{k+kn}{import} \PYG{n+nn}{cppad\PYGZus{}py}
    \PYG{c+c1}{\PYGZsh{}}
    \PYG{c+c1}{\PYGZsh{} initialize return variable}
    \PYG{n}{ok} \PYG{o}{=} \PYG{k+kc}{True}
    \PYG{c+c1}{\PYGZsh{} \PYGZhy{}\PYGZhy{}\PYGZhy{}\PYGZhy{}\PYGZhy{}\PYGZhy{}\PYGZhy{}\PYGZhy{}\PYGZhy{}\PYGZhy{}\PYGZhy{}\PYGZhy{}\PYGZhy{}\PYGZhy{}\PYGZhy{}\PYGZhy{}\PYGZhy{}\PYGZhy{}\PYGZhy{}\PYGZhy{}\PYGZhy{}\PYGZhy{}\PYGZhy{}\PYGZhy{}\PYGZhy{}\PYGZhy{}\PYGZhy{}\PYGZhy{}\PYGZhy{}\PYGZhy{}\PYGZhy{}\PYGZhy{}\PYGZhy{}\PYGZhy{}\PYGZhy{}\PYGZhy{}\PYGZhy{}\PYGZhy{}\PYGZhy{}\PYGZhy{}\PYGZhy{}\PYGZhy{}\PYGZhy{}\PYGZhy{}\PYGZhy{}\PYGZhy{}\PYGZhy{}\PYGZhy{}\PYGZhy{}\PYGZhy{}\PYGZhy{}\PYGZhy{}\PYGZhy{}\PYGZhy{}\PYGZhy{}\PYGZhy{}\PYGZhy{}\PYGZhy{}\PYGZhy{}\PYGZhy{}\PYGZhy{}\PYGZhy{}\PYGZhy{}\PYGZhy{}\PYGZhy{}\PYGZhy{}\PYGZhy{}\PYGZhy{}\PYGZhy{}}
    \PYG{c+c1}{\PYGZsh{} number of dependent and independent variables}
    \PYG{n}{n\PYGZus{}dep} \PYG{o}{=} \PYG{l+m+mi}{1}
    \PYG{n}{n\PYGZus{}ind} \PYG{o}{=} \PYG{l+m+mi}{2}
    \PYG{c+c1}{\PYGZsh{}}
    \PYG{c+c1}{\PYGZsh{} create the independent variables ax}
    \PYG{n}{xp} \PYG{o}{=} \PYG{n}{numpy}\PYG{o}{.}\PYG{n}{empty}\PYG{p}{(}\PYG{n}{n\PYGZus{}ind}\PYG{p}{,} \PYG{n}{dtype}\PYG{o}{=}\PYG{n+nb}{float}\PYG{p}{)}
    \PYG{k}{for} \PYG{n}{i} \PYG{o+ow}{in} \PYG{n+nb}{range}\PYG{p}{(} \PYG{n}{n\PYGZus{}ind}  \PYG{p}{)} \PYG{p}{:}
        \PYG{n}{xp}\PYG{p}{[}\PYG{n}{i}\PYG{p}{]} \PYG{o}{=} \PYG{n}{i} \PYG{o}{+} \PYG{l+m+mf}{1.0}
    \PYG{c+c1}{\PYGZsh{}}
    \PYG{n}{ax} \PYG{o}{=} \PYG{n}{cppad\PYGZus{}py}\PYG{o}{.}\PYG{n}{independent}\PYG{p}{(}\PYG{n}{xp}\PYG{p}{)}
    \PYG{c+c1}{\PYGZsh{}}
    \PYG{c+c1}{\PYGZsh{} create dependent varialbes ay with ay0 = ax0 * ax1}
    \PYG{n}{ax0}    \PYG{o}{=} \PYG{n}{ax}\PYG{p}{[}\PYG{l+m+mi}{0}\PYG{p}{]}
    \PYG{n}{ax1}    \PYG{o}{=} \PYG{n}{ax}\PYG{p}{[}\PYG{l+m+mi}{1}\PYG{p}{]}
    \PYG{n}{ay}    \PYG{o}{=} \PYG{n}{numpy}\PYG{o}{.}\PYG{n}{empty}\PYG{p}{(}\PYG{n}{n\PYGZus{}dep}\PYG{p}{,} \PYG{n}{dtype}\PYG{o}{=}\PYG{n}{cppad\PYGZus{}py}\PYG{o}{.}\PYG{n}{a\PYGZus{}double}\PYG{p}{)}
    \PYG{n}{ay}\PYG{p}{[}\PYG{l+m+mi}{0}\PYG{p}{]} \PYG{o}{=} \PYG{n}{ax0} \PYG{o}{*} \PYG{n}{ax1}
    \PYG{c+c1}{\PYGZsh{}}
    \PYG{c+c1}{\PYGZsh{} define af corresponding to f(x) = x0 * x1}
    \PYG{n}{f}  \PYG{o}{=} \PYG{n}{cppad\PYGZus{}py}\PYG{o}{.}\PYG{n}{d\PYGZus{}fun}\PYG{p}{(}\PYG{n}{ax}\PYG{p}{,} \PYG{n}{ay}\PYG{p}{)}
    \PYG{c+c1}{\PYGZsh{}}
    \PYG{c+c1}{\PYGZsh{} define X(t) = (3 + t, 2 + t)}
    \PYG{c+c1}{\PYGZsh{} it follows that Y(t) = f(X(t)) = (3 + t) * (2 + t)}
    \PYG{c+c1}{\PYGZsh{}}
    \PYG{c+c1}{\PYGZsh{} Y(0) = 6 and p ! = 1}
    \PYG{n}{p}     \PYG{o}{=} \PYG{l+m+mi}{0}
    \PYG{n}{xp}\PYG{p}{[}\PYG{l+m+mi}{0}\PYG{p}{]} \PYG{o}{=} \PYG{l+m+mf}{3.0}
    \PYG{n}{xp}\PYG{p}{[}\PYG{l+m+mi}{1}\PYG{p}{]} \PYG{o}{=} \PYG{l+m+mf}{2.0}
    \PYG{n}{yp} \PYG{o}{=} \PYG{n}{f}\PYG{o}{.}\PYG{n}{forward}\PYG{p}{(}\PYG{n}{p}\PYG{p}{,} \PYG{n}{xp}\PYG{p}{)}
    \PYG{n}{ok} \PYG{o}{=} \PYG{n}{ok} \PYG{o+ow}{and} \PYG{n}{yp}\PYG{p}{[}\PYG{l+m+mi}{0}\PYG{p}{]} \PYG{o}{==} \PYG{l+m+mf}{6.0}
    \PYG{c+c1}{\PYGZsh{}}
    \PYG{c+c1}{\PYGZsh{} first order Taylor coefficients for X(t)}
    \PYG{n}{p}     \PYG{o}{=} \PYG{l+m+mi}{1}
    \PYG{n}{xp}\PYG{p}{[}\PYG{l+m+mi}{0}\PYG{p}{]} \PYG{o}{=} \PYG{l+m+mf}{1.0}
    \PYG{n}{xp}\PYG{p}{[}\PYG{l+m+mi}{1}\PYG{p}{]} \PYG{o}{=} \PYG{l+m+mf}{1.0}
    \PYG{c+c1}{\PYGZsh{}}
    \PYG{c+c1}{\PYGZsh{} first order Taylor coefficient for Y(t)}
    \PYG{c+c1}{\PYGZsh{} Y\PYGZsq{}(0) = 3 + 2 = 5 and p ! = 1}
    \PYG{n}{yp} \PYG{o}{=} \PYG{n}{f}\PYG{o}{.}\PYG{n}{forward}\PYG{p}{(}\PYG{n}{p}\PYG{p}{,} \PYG{n}{xp}\PYG{p}{)}
    \PYG{n}{ok} \PYG{o}{=} \PYG{n}{ok} \PYG{o+ow}{and} \PYG{n}{yp}\PYG{p}{[}\PYG{l+m+mi}{0}\PYG{p}{]} \PYG{o}{==} \PYG{l+m+mf}{5.0}
    \PYG{c+c1}{\PYGZsh{}}
    \PYG{c+c1}{\PYGZsh{} second order Taylor coefficients for X(t)}
    \PYG{n}{p}     \PYG{o}{=} \PYG{l+m+mi}{2}
    \PYG{n}{xp}\PYG{p}{[}\PYG{l+m+mi}{0}\PYG{p}{]} \PYG{o}{=} \PYG{l+m+mf}{0.0}
    \PYG{n}{xp}\PYG{p}{[}\PYG{l+m+mi}{1}\PYG{p}{]} \PYG{o}{=} \PYG{l+m+mf}{0.0}
    \PYG{c+c1}{\PYGZsh{}}
    \PYG{c+c1}{\PYGZsh{} second order Taylor coefficient for Y(t)}
    \PYG{c+c1}{\PYGZsh{} Y\PYGZsq{}\PYGZsq{}(0) = 2.0 and p ! = 2}
    \PYG{n}{yp} \PYG{o}{=} \PYG{n}{f}\PYG{o}{.}\PYG{n}{forward}\PYG{p}{(}\PYG{n}{p}\PYG{p}{,} \PYG{n}{xp}\PYG{p}{)}
    \PYG{n}{ok} \PYG{o}{=} \PYG{n}{ok} \PYG{o+ow}{and} \PYG{n}{yp}\PYG{p}{[}\PYG{l+m+mi}{0}\PYG{p}{]} \PYG{o}{==} \PYG{l+m+mf}{1.0}
    \PYG{c+c1}{\PYGZsh{} \PYGZhy{}\PYGZhy{}\PYGZhy{}\PYGZhy{}\PYGZhy{}\PYGZhy{}\PYGZhy{}\PYGZhy{}\PYGZhy{}\PYGZhy{}\PYGZhy{}\PYGZhy{}\PYGZhy{}\PYGZhy{}\PYGZhy{}\PYGZhy{}\PYGZhy{}\PYGZhy{}\PYGZhy{}\PYGZhy{}\PYGZhy{}\PYGZhy{}\PYGZhy{}\PYGZhy{}\PYGZhy{}\PYGZhy{}\PYGZhy{}\PYGZhy{}\PYGZhy{}\PYGZhy{}\PYGZhy{}\PYGZhy{}\PYGZhy{}\PYGZhy{}\PYGZhy{}\PYGZhy{}\PYGZhy{}\PYGZhy{}\PYGZhy{}\PYGZhy{}\PYGZhy{}\PYGZhy{}\PYGZhy{}\PYGZhy{}\PYGZhy{}\PYGZhy{}\PYGZhy{}\PYGZhy{}\PYGZhy{}\PYGZhy{}\PYGZhy{}\PYGZhy{}\PYGZhy{}\PYGZhy{}\PYGZhy{}\PYGZhy{}\PYGZhy{}\PYGZhy{}\PYGZhy{}\PYGZhy{}\PYGZhy{}\PYGZhy{}\PYGZhy{}\PYGZhy{}\PYGZhy{}\PYGZhy{}\PYGZhy{}\PYGZhy{}\PYGZhy{}}
    \PYG{n}{af} \PYG{o}{=} \PYG{n}{cppad\PYGZus{}py}\PYG{o}{.}\PYG{n}{a\PYGZus{}fun}\PYG{p}{(}\PYG{n}{f}\PYG{p}{)}
    \PYG{n}{ok} \PYG{o}{=} \PYG{n}{ok} \PYG{o+ow}{and} \PYG{n}{af}\PYG{o}{.}\PYG{n}{size\PYGZus{}order}\PYG{p}{(}\PYG{p}{)} \PYG{o}{==} \PYG{l+m+mi}{0}
    \PYG{c+c1}{\PYGZsh{}}
    \PYG{c+c1}{\PYGZsh{} zero order forward}
    \PYG{n}{p}   \PYG{o}{=} \PYG{l+m+mi}{0}
    \PYG{n}{axp} \PYG{o}{=} \PYG{n}{numpy}\PYG{o}{.}\PYG{n}{empty}\PYG{p}{(}\PYG{n}{n\PYGZus{}ind}\PYG{p}{,} \PYG{n}{dtype}\PYG{o}{=}\PYG{n}{cppad\PYGZus{}py}\PYG{o}{.}\PYG{n}{a\PYGZus{}double}\PYG{p}{)}
    \PYG{n}{axp}\PYG{p}{[}\PYG{l+m+mi}{0}\PYG{p}{]} \PYG{o}{=} \PYG{l+m+mf}{3.0}
    \PYG{n}{axp}\PYG{p}{[}\PYG{l+m+mi}{1}\PYG{p}{]} \PYG{o}{=} \PYG{l+m+mf}{2.0}
    \PYG{n}{ayp}    \PYG{o}{=} \PYG{n}{af}\PYG{o}{.}\PYG{n}{forward}\PYG{p}{(}\PYG{n}{p}\PYG{p}{,} \PYG{n}{axp}\PYG{p}{)}
    \PYG{n}{ok}     \PYG{o}{=} \PYG{n}{ok} \PYG{o+ow}{and} \PYG{n}{ayp}\PYG{p}{[}\PYG{l+m+mi}{0}\PYG{p}{]} \PYG{o}{==} \PYG{n}{cppad\PYGZus{}py}\PYG{o}{.}\PYG{n}{a\PYGZus{}double}\PYG{p}{(}\PYG{l+m+mf}{6.0}\PYG{p}{)}
    \PYG{n}{ok}     \PYG{o}{=} \PYG{n}{ok} \PYG{o+ow}{and} \PYG{n}{af}\PYG{o}{.}\PYG{n}{size\PYGZus{}order}\PYG{p}{(}\PYG{p}{)} \PYG{o}{==} \PYG{l+m+mi}{1}
    \PYG{c+c1}{\PYGZsh{}}
    \PYG{k}{return}\PYG{p}{(} \PYG{n}{ok} \PYG{p}{)}
\PYG{c+c1}{\PYGZsh{}}
\end{sphinxVerbatim}


\bigskip\hrule\bigskip


xsrst input file: \sphinxcode{\sphinxupquote{example/python/core/fun\_forward\_xam.py}}

\index{example@\spxentry{example}}\ignorespaces 

\subparagraph{Example}
\label{\detokenize{xsrst/py_fun_forward:example}}\label{\detokenize{xsrst/py_fun_forward:py-fun-forward-example}}\label{\detokenize{xsrst/py_fun_forward:index-11}}
{\hyperref[\detokenize{xsrst/fun_forward_xam_py:id1}]{\sphinxcrossref{\DUrole{std,std-ref}{fun\_forward\_xam\_py}}}}


\bigskip\hrule\bigskip


xsrst input file: \sphinxcode{\sphinxupquote{lib/python/cppad\_py/fun\_forward.py}}


\subparagraph{py\_fun\_reverse}
\label{\detokenize{xsrst/py_fun_reverse:py-fun-reverse}}\label{\detokenize{xsrst/py_fun_reverse::doc}}
\index{py\_fun\_reverse@\spxentry{py\_fun\_reverse}}\index{reverse@\spxentry{reverse}}\index{mode@\spxentry{mode}}\index{ad@\spxentry{ad}}\ignorespaces 

\subparagraph{3.1.1.9 Reverse Mode AD}
\label{\detokenize{xsrst/py_fun_reverse:reverse-mode-ad}}\label{\detokenize{xsrst/py_fun_reverse:id1}}\begin{itemize}
\item {} 
{\hyperref[\detokenize{xsrst/py_fun_reverse:py-fun-reverse-syntax}]{\sphinxcrossref{\DUrole{std,std-ref}{Syntax}}}}

\item {} 
{\hyperref[\detokenize{xsrst/py_fun_reverse:py-fun-reverse-f}]{\sphinxcrossref{\DUrole{std,std-ref}{f}}}}

\item {} \begin{description}
\item[{{\hyperref[\detokenize{xsrst/py_fun_reverse:py-fun-reverse-notation}]{\sphinxcrossref{\DUrole{std,std-ref}{Notation}}}}}] \leavevmode\begin{itemize}
\item {} 
{\hyperref[\detokenize{xsrst/py_fun_reverse:py-fun-reverse-notation-f-x}]{\sphinxcrossref{\DUrole{std,std-ref}{f(x)}}}}

\item {} 
{\hyperref[\detokenize{xsrst/py_fun_reverse:py-fun-reverse-notation-x-t-s}]{\sphinxcrossref{\DUrole{std,std-ref}{X(t), S}}}}

\item {} 
{\hyperref[\detokenize{xsrst/py_fun_reverse:py-fun-reverse-notation-y-t-t}]{\sphinxcrossref{\DUrole{std,std-ref}{Y(t), T}}}}

\item {} 
{\hyperref[\detokenize{xsrst/py_fun_reverse:py-fun-reverse-notation-g-t}]{\sphinxcrossref{\DUrole{std,std-ref}{G(T)}}}}

\end{itemize}

\end{description}

\item {} 
{\hyperref[\detokenize{xsrst/py_fun_reverse:py-fun-reverse-q}]{\sphinxcrossref{\DUrole{std,std-ref}{q}}}}

\item {} 
{\hyperref[\detokenize{xsrst/py_fun_reverse:py-fun-reverse-yq}]{\sphinxcrossref{\DUrole{std,std-ref}{yq}}}}

\item {} 
{\hyperref[\detokenize{xsrst/py_fun_reverse:py-fun-reverse-xq}]{\sphinxcrossref{\DUrole{std,std-ref}{xq}}}}

\item {} 
{\hyperref[\detokenize{xsrst/py_fun_reverse:py-fun-reverse-example}]{\sphinxcrossref{\DUrole{std,std-ref}{Example}}}}

\end{itemize}

\index{syntax@\spxentry{syntax}}\ignorespaces 

\subparagraph{Syntax}
\label{\detokenize{xsrst/py_fun_reverse:syntax}}\label{\detokenize{xsrst/py_fun_reverse:py-fun-reverse-syntax}}\label{\detokenize{xsrst/py_fun_reverse:index-1}}
\sphinxstyleemphasis{xq} = \sphinxstyleemphasis{f} . \sphinxcode{\sphinxupquote{reverse}} ( \sphinxstyleemphasis{q} , \sphinxstyleemphasis{yq} )

\index{f@\spxentry{f}}\ignorespaces 

\subparagraph{f}
\label{\detokenize{xsrst/py_fun_reverse:f}}\label{\detokenize{xsrst/py_fun_reverse:py-fun-reverse-f}}\label{\detokenize{xsrst/py_fun_reverse:index-2}}
This is either a
{\hyperref[\detokenize{xsrst/py_fun_ctor:py-fun-ctor-syntax-d-fun}]{\sphinxcrossref{\DUrole{std,std-ref}{d\_fun}}}} or
{\hyperref[\detokenize{xsrst/py_fun_ctor:py-fun-ctor-syntax-a-fun}]{\sphinxcrossref{\DUrole{std,std-ref}{a\_fun}}}} function object
and is effectively constant; i.e., not changed.

\index{notation@\spxentry{notation}}\ignorespaces 

\subparagraph{Notation}
\label{\detokenize{xsrst/py_fun_reverse:notation}}\label{\detokenize{xsrst/py_fun_reverse:py-fun-reverse-notation}}\label{\detokenize{xsrst/py_fun_reverse:index-3}}
\index{f(x)@\spxentry{f(x)}}\ignorespaces 

\subparagraph{f(x)}
\label{\detokenize{xsrst/py_fun_reverse:f-x}}\label{\detokenize{xsrst/py_fun_reverse:py-fun-reverse-notation-f-x}}\label{\detokenize{xsrst/py_fun_reverse:index-4}}
We use the notation \(f: \B{R}^n \rightarrow \B{R}^m\)
for the function corresponding to \sphinxstyleemphasis{f} .
Note that \sphinxstyleemphasis{n} is the size of {\hyperref[\detokenize{xsrst/py_fun_ctor:py-fun-ctor-ax}]{\sphinxcrossref{\DUrole{std,std-ref}{ax}}}}
and \sphinxstyleemphasis{m} is the size of {\hyperref[\detokenize{xsrst/py_fun_ctor:py-fun-ctor-ay}]{\sphinxcrossref{\DUrole{std,std-ref}{ay}}}}
in to the constructor for \sphinxstyleemphasis{f} .

\index{x(t)@\spxentry{x(t)}}\index{s@\spxentry{s}}\ignorespaces 

\subparagraph{X(t), S}
\label{\detokenize{xsrst/py_fun_reverse:x-t-s}}\label{\detokenize{xsrst/py_fun_reverse:py-fun-reverse-notation-x-t-s}}\label{\detokenize{xsrst/py_fun_reverse:index-5}}
This is the same function as
{\hyperref[\detokenize{xsrst/py_fun_forward:py-fun-forward-x-t}]{\sphinxcrossref{\DUrole{std,std-ref}{x(t)}}}} in the previous call to
\sphinxstyleemphasis{f} . \sphinxcode{\sphinxupquote{forward}} .
We use \(S \in \B{R}^{n \times q}\) to denote the Taylor coefficients
of \(X(t)\).

\index{y(t)@\spxentry{y(t)}}\index{t@\spxentry{t}}\ignorespaces 

\subparagraph{Y(t), T}
\label{\detokenize{xsrst/py_fun_reverse:y-t-t}}\label{\detokenize{xsrst/py_fun_reverse:py-fun-reverse-notation-y-t-t}}\label{\detokenize{xsrst/py_fun_reverse:index-6}}
This is the same function as
{\hyperref[\detokenize{xsrst/py_fun_forward:py-fun-forward-y-t}]{\sphinxcrossref{\DUrole{std,std-ref}{y(t)}}}} in the previous call to
\sphinxstyleemphasis{f} . \sphinxcode{\sphinxupquote{forward}} .
We use \(T \in \B{R}^{m \times q}\) to denote the Taylor coefficients
of \(Y(t)\).
We also use the notation \(T(S)\) to express the fact that
the Taylor coefficients for \(Y(t)\) are a function of the
Taylor coefficients of \(X(t)\).

\index{g(t)@\spxentry{g(t)}}\ignorespaces 

\subparagraph{G(T)}
\label{\detokenize{xsrst/py_fun_reverse:g-t}}\label{\detokenize{xsrst/py_fun_reverse:py-fun-reverse-notation-g-t}}\label{\detokenize{xsrst/py_fun_reverse:index-7}}
We use the notation \(G : \B{R}^{m \times p} \rightarrow \B{R}\)
for a function that the calling routine chooses.

\index{q@\spxentry{q}}\ignorespaces 

\subparagraph{q}
\label{\detokenize{xsrst/py_fun_reverse:q}}\label{\detokenize{xsrst/py_fun_reverse:py-fun-reverse-q}}\label{\detokenize{xsrst/py_fun_reverse:index-8}}
This argument has type \sphinxcode{\sphinxupquote{int}} and is positive.
It is the number of the Taylor coefficient (for each variable)
that we are computing the derivative with respect to.
It must be greater than zero, and less than or equal
the number of Taylor coefficient stored in \sphinxstyleemphasis{f} ; i.e.,
{\hyperref[\detokenize{xsrst/py_fun_property:py-fun-property-size-order}]{\sphinxcrossref{\DUrole{std,std-ref}{f\_size\_order()}}}}.

\index{yq@\spxentry{yq}}\ignorespaces 

\subparagraph{yq}
\label{\detokenize{xsrst/py_fun_reverse:yq}}\label{\detokenize{xsrst/py_fun_reverse:py-fun-reverse-yq}}\label{\detokenize{xsrst/py_fun_reverse:index-9}}
If \sphinxstyleemphasis{f} is a \sphinxcode{\sphinxupquote{d\_fun}} ( \sphinxcode{\sphinxupquote{a\_fun}} ) object, \sphinxstyleemphasis{yq}
is a numpy vector with \sphinxcode{\sphinxupquote{float}} ( \sphinxcode{\sphinxupquote{a\_double}} ) elements,
\sphinxstyleemphasis{m} rows and \sphinxstyleemphasis{q} columns.
For 0 \textless{}= \sphinxstyleemphasis{i} \textless{} \sphinxstyleemphasis{m} and 0 \textless{}= \sphinxstyleemphasis{k} \textless{} \sphinxstyleemphasis{q} ,
\sphinxstyleemphasis{yq} {[} \sphinxstyleemphasis{i} , \sphinxstyleemphasis{k} {]} is the partial derivative of
\(G(T)\) with respect to the \sphinxstyleemphasis{k}\sphinxhyphen{}th order Taylor coefficient
for the \sphinxstyleemphasis{i}\sphinxhyphen{}th component function; i.e.,
the partial derivative of \(G(T)\) w.r.t. \(Y_i^{(k)} (t) / k !\).

\index{xq@\spxentry{xq}}\ignorespaces 

\subparagraph{xq}
\label{\detokenize{xsrst/py_fun_reverse:xq}}\label{\detokenize{xsrst/py_fun_reverse:py-fun-reverse-xq}}\label{\detokenize{xsrst/py_fun_reverse:index-10}}
The result is a numpy vector with \sphinxcode{\sphinxupquote{float}} ( \sphinxcode{\sphinxupquote{a\_double}} ) elements,
\sphinxstyleemphasis{n} rows and \sphinxstyleemphasis{q} columns.
For 0 \textless{}= \sphinxstyleemphasis{j} \textless{} \sphinxstyleemphasis{n} and 0 \textless{}= \sphinxstyleemphasis{k} \textless{} \sphinxstyleemphasis{q} ,
\sphinxstyleemphasis{xq} {[} \sphinxstyleemphasis{j} , \sphinxstyleemphasis{k} {]} is the partial derivative of
\(G(T(S))\) with respect to the \sphinxstyleemphasis{k}\sphinxhyphen{}th order Taylor coefficient
for the \sphinxstyleemphasis{j}\sphinxhyphen{}th component function; i.e.,
the partial derivative of
\(G(T(S))\) w.r.t. \(S_j^{(k)} (t) / k !\).


\subparagraph{fun\_reverse\_xam\_py}
\label{\detokenize{xsrst/fun_reverse_xam_py:fun-reverse-xam-py}}\label{\detokenize{xsrst/fun_reverse_xam_py::doc}}
\index{fun\_reverse\_xam\_py@\spxentry{fun\_reverse\_xam\_py}}\index{python:@\spxentry{python:}}\index{reverse@\spxentry{reverse}}\index{mode@\spxentry{mode}}\index{ad:@\spxentry{ad:}}\index{example@\spxentry{example}}\index{test@\spxentry{test}}\ignorespaces 

\subparagraph{3.1.1.9.1 Python: Reverse Mode AD: Example and Test}
\label{\detokenize{xsrst/fun_reverse_xam_py:python-reverse-mode-ad-example-and-test}}\label{\detokenize{xsrst/fun_reverse_xam_py:id1}}
\begin{sphinxVerbatim}[commandchars=\\\{\}]
\PYG{k}{def} \PYG{n+nf}{fun\PYGZus{}reverse\PYGZus{}xam}\PYG{p}{(}\PYG{p}{)} \PYG{p}{:}
    \PYG{c+c1}{\PYGZsh{}}
    \PYG{k+kn}{import} \PYG{n+nn}{numpy}
    \PYG{k+kn}{import} \PYG{n+nn}{cppad\PYGZus{}py}
    \PYG{c+c1}{\PYGZsh{}}
    \PYG{c+c1}{\PYGZsh{} initialize return variable}
    \PYG{n}{ok} \PYG{o}{=} \PYG{k+kc}{True}
    \PYG{c+c1}{\PYGZsh{} \PYGZhy{}\PYGZhy{}\PYGZhy{}\PYGZhy{}\PYGZhy{}\PYGZhy{}\PYGZhy{}\PYGZhy{}\PYGZhy{}\PYGZhy{}\PYGZhy{}\PYGZhy{}\PYGZhy{}\PYGZhy{}\PYGZhy{}\PYGZhy{}\PYGZhy{}\PYGZhy{}\PYGZhy{}\PYGZhy{}\PYGZhy{}\PYGZhy{}\PYGZhy{}\PYGZhy{}\PYGZhy{}\PYGZhy{}\PYGZhy{}\PYGZhy{}\PYGZhy{}\PYGZhy{}\PYGZhy{}\PYGZhy{}\PYGZhy{}\PYGZhy{}\PYGZhy{}\PYGZhy{}\PYGZhy{}\PYGZhy{}\PYGZhy{}\PYGZhy{}\PYGZhy{}\PYGZhy{}\PYGZhy{}\PYGZhy{}\PYGZhy{}\PYGZhy{}\PYGZhy{}\PYGZhy{}\PYGZhy{}\PYGZhy{}\PYGZhy{}\PYGZhy{}\PYGZhy{}\PYGZhy{}\PYGZhy{}\PYGZhy{}\PYGZhy{}\PYGZhy{}\PYGZhy{}\PYGZhy{}\PYGZhy{}\PYGZhy{}\PYGZhy{}\PYGZhy{}\PYGZhy{}\PYGZhy{}\PYGZhy{}\PYGZhy{}\PYGZhy{}}
    \PYG{c+c1}{\PYGZsh{} number of dependent and independent variables}
    \PYG{n}{n\PYGZus{}dep} \PYG{o}{=} \PYG{l+m+mi}{1}
    \PYG{n}{n\PYGZus{}ind} \PYG{o}{=} \PYG{l+m+mi}{3}
    \PYG{c+c1}{\PYGZsh{}}
    \PYG{c+c1}{\PYGZsh{} create the independent variables ax}
    \PYG{n}{xp} \PYG{o}{=} \PYG{n}{numpy}\PYG{o}{.}\PYG{n}{empty}\PYG{p}{(}\PYG{n}{n\PYGZus{}ind}\PYG{p}{,} \PYG{n}{dtype}\PYG{o}{=}\PYG{n+nb}{float}\PYG{p}{)}
    \PYG{k}{for} \PYG{n}{i} \PYG{o+ow}{in} \PYG{n+nb}{range}\PYG{p}{(} \PYG{n}{n\PYGZus{}ind}  \PYG{p}{)} \PYG{p}{:}
        \PYG{n}{xp}\PYG{p}{[}\PYG{n}{i}\PYG{p}{]} \PYG{o}{=} \PYG{n}{i}
    \PYG{c+c1}{\PYGZsh{}}
    \PYG{n}{ax} \PYG{o}{=} \PYG{n}{cppad\PYGZus{}py}\PYG{o}{.}\PYG{n}{independent}\PYG{p}{(}\PYG{n}{xp}\PYG{p}{)}
    \PYG{c+c1}{\PYGZsh{}}
    \PYG{c+c1}{\PYGZsh{} create dependent variables ay with ay0 = ax\PYGZus{}0 * ax\PYGZus{}1 * ax\PYGZus{}2}
    \PYG{n}{ax\PYGZus{}0}  \PYG{o}{=} \PYG{n}{ax}\PYG{p}{[}\PYG{l+m+mi}{0}\PYG{p}{]}
    \PYG{n}{ax\PYGZus{}1}  \PYG{o}{=} \PYG{n}{ax}\PYG{p}{[}\PYG{l+m+mi}{1}\PYG{p}{]}
    \PYG{n}{ax\PYGZus{}2}  \PYG{o}{=} \PYG{n}{ax}\PYG{p}{[}\PYG{l+m+mi}{2}\PYG{p}{]}
    \PYG{n}{ay}    \PYG{o}{=} \PYG{n}{numpy}\PYG{o}{.}\PYG{n}{empty}\PYG{p}{(}\PYG{n}{n\PYGZus{}dep}\PYG{p}{,} \PYG{n}{dtype}\PYG{o}{=}\PYG{n}{cppad\PYGZus{}py}\PYG{o}{.}\PYG{n}{a\PYGZus{}double}\PYG{p}{)}
    \PYG{n}{ay}\PYG{p}{[}\PYG{l+m+mi}{0}\PYG{p}{]} \PYG{o}{=} \PYG{n}{ax\PYGZus{}0} \PYG{o}{*} \PYG{n}{ax\PYGZus{}1} \PYG{o}{*} \PYG{n}{ax\PYGZus{}2}
    \PYG{c+c1}{\PYGZsh{}}
    \PYG{c+c1}{\PYGZsh{} define af corresponding to f(x) = x\PYGZus{}0 * x\PYGZus{}1 * x\PYGZus{}2}
    \PYG{n}{f}  \PYG{o}{=} \PYG{n}{cppad\PYGZus{}py}\PYG{o}{.}\PYG{n}{d\PYGZus{}fun}\PYG{p}{(}\PYG{n}{ax}\PYG{p}{,} \PYG{n}{ay}\PYG{p}{)}
    \PYG{c+c1}{\PYGZsh{} \PYGZhy{}\PYGZhy{}\PYGZhy{}\PYGZhy{}\PYGZhy{}\PYGZhy{}\PYGZhy{}\PYGZhy{}\PYGZhy{}\PYGZhy{}\PYGZhy{}\PYGZhy{}\PYGZhy{}\PYGZhy{}\PYGZhy{}\PYGZhy{}\PYGZhy{}\PYGZhy{}\PYGZhy{}\PYGZhy{}\PYGZhy{}\PYGZhy{}\PYGZhy{}\PYGZhy{}\PYGZhy{}\PYGZhy{}\PYGZhy{}\PYGZhy{}\PYGZhy{}\PYGZhy{}\PYGZhy{}\PYGZhy{}\PYGZhy{}\PYGZhy{}\PYGZhy{}\PYGZhy{}\PYGZhy{}\PYGZhy{}\PYGZhy{}\PYGZhy{}\PYGZhy{}\PYGZhy{}\PYGZhy{}\PYGZhy{}\PYGZhy{}\PYGZhy{}\PYGZhy{}\PYGZhy{}\PYGZhy{}\PYGZhy{}\PYGZhy{}\PYGZhy{}\PYGZhy{}\PYGZhy{}\PYGZhy{}\PYGZhy{}\PYGZhy{}\PYGZhy{}\PYGZhy{}\PYGZhy{}\PYGZhy{}\PYGZhy{}\PYGZhy{}\PYGZhy{}\PYGZhy{}\PYGZhy{}\PYGZhy{}\PYGZhy{}\PYGZhy{}\PYGZhy{}\PYGZhy{}}
    \PYG{c+c1}{\PYGZsh{} define          X(t) = (x\PYGZus{}0 + t, x\PYGZus{}1 + t, x\PYGZus{}2 + t)}
    \PYG{c+c1}{\PYGZsh{} it follows that Y(t) = f(X(t)) = (x\PYGZus{}0 + t) * (x\PYGZus{}1 + t) * (x\PYGZus{}2 + t)}
    \PYG{c+c1}{\PYGZsh{} and that       Y\PYGZsq{}(0) = x\PYGZus{}1 * x\PYGZus{}2 + x\PYGZus{}0 * x\PYGZus{}2 + x\PYGZus{}0 * x\PYGZus{}1}
    \PYG{c+c1}{\PYGZsh{} \PYGZhy{}\PYGZhy{}\PYGZhy{}\PYGZhy{}\PYGZhy{}\PYGZhy{}\PYGZhy{}\PYGZhy{}\PYGZhy{}\PYGZhy{}\PYGZhy{}\PYGZhy{}\PYGZhy{}\PYGZhy{}\PYGZhy{}\PYGZhy{}\PYGZhy{}\PYGZhy{}\PYGZhy{}\PYGZhy{}\PYGZhy{}\PYGZhy{}\PYGZhy{}\PYGZhy{}\PYGZhy{}\PYGZhy{}\PYGZhy{}\PYGZhy{}\PYGZhy{}\PYGZhy{}\PYGZhy{}\PYGZhy{}\PYGZhy{}\PYGZhy{}\PYGZhy{}\PYGZhy{}\PYGZhy{}\PYGZhy{}\PYGZhy{}\PYGZhy{}\PYGZhy{}\PYGZhy{}\PYGZhy{}\PYGZhy{}\PYGZhy{}\PYGZhy{}\PYGZhy{}\PYGZhy{}\PYGZhy{}\PYGZhy{}\PYGZhy{}\PYGZhy{}\PYGZhy{}\PYGZhy{}\PYGZhy{}\PYGZhy{}\PYGZhy{}\PYGZhy{}\PYGZhy{}\PYGZhy{}\PYGZhy{}\PYGZhy{}\PYGZhy{}\PYGZhy{}\PYGZhy{}\PYGZhy{}\PYGZhy{}\PYGZhy{}\PYGZhy{}\PYGZhy{}\PYGZhy{}}
    \PYG{c+c1}{\PYGZsh{} zero order forward mode}
    \PYG{n}{p}     \PYG{o}{=} \PYG{l+m+mi}{0}
    \PYG{n}{xp}\PYG{p}{[}\PYG{l+m+mi}{0}\PYG{p}{]} \PYG{o}{=} \PYG{l+m+mf}{2.0}
    \PYG{n}{xp}\PYG{p}{[}\PYG{l+m+mi}{1}\PYG{p}{]} \PYG{o}{=} \PYG{l+m+mf}{3.0}
    \PYG{n}{xp}\PYG{p}{[}\PYG{l+m+mi}{2}\PYG{p}{]} \PYG{o}{=} \PYG{l+m+mf}{4.0}
    \PYG{n}{yp} \PYG{o}{=} \PYG{n}{f}\PYG{o}{.}\PYG{n}{forward}\PYG{p}{(}\PYG{n}{p}\PYG{p}{,} \PYG{n}{xp}\PYG{p}{)}
    \PYG{n}{ok} \PYG{o}{=} \PYG{n}{ok} \PYG{o+ow}{and} \PYG{n}{yp}\PYG{p}{[}\PYG{l+m+mi}{0}\PYG{p}{]} \PYG{o}{==} \PYG{l+m+mf}{24.0}
    \PYG{c+c1}{\PYGZsh{} \PYGZhy{}\PYGZhy{}\PYGZhy{}\PYGZhy{}\PYGZhy{}\PYGZhy{}\PYGZhy{}\PYGZhy{}\PYGZhy{}\PYGZhy{}\PYGZhy{}\PYGZhy{}\PYGZhy{}\PYGZhy{}\PYGZhy{}\PYGZhy{}\PYGZhy{}\PYGZhy{}\PYGZhy{}\PYGZhy{}\PYGZhy{}\PYGZhy{}\PYGZhy{}\PYGZhy{}\PYGZhy{}\PYGZhy{}\PYGZhy{}\PYGZhy{}\PYGZhy{}\PYGZhy{}\PYGZhy{}\PYGZhy{}\PYGZhy{}\PYGZhy{}\PYGZhy{}\PYGZhy{}\PYGZhy{}\PYGZhy{}\PYGZhy{}\PYGZhy{}\PYGZhy{}\PYGZhy{}\PYGZhy{}\PYGZhy{}\PYGZhy{}\PYGZhy{}\PYGZhy{}\PYGZhy{}\PYGZhy{}\PYGZhy{}\PYGZhy{}\PYGZhy{}\PYGZhy{}\PYGZhy{}\PYGZhy{}\PYGZhy{}\PYGZhy{}\PYGZhy{}\PYGZhy{}\PYGZhy{}\PYGZhy{}\PYGZhy{}\PYGZhy{}\PYGZhy{}\PYGZhy{}\PYGZhy{}\PYGZhy{}\PYGZhy{}\PYGZhy{}\PYGZhy{}\PYGZhy{}}
    \PYG{c+c1}{\PYGZsh{} first order reverse (derivative of zero order forward)}
    \PYG{c+c1}{\PYGZsh{} define G( Y ) = y\PYGZus{}0 = x\PYGZus{}0 * x\PYGZus{}1 * x\PYGZus{}2}
    \PYG{n}{m}         \PYG{o}{=} \PYG{n}{f}\PYG{o}{.}\PYG{n}{size\PYGZus{}range}\PYG{p}{(}\PYG{p}{)}
    \PYG{n}{q}         \PYG{o}{=} \PYG{l+m+mi}{1}
    \PYG{n}{yq1}       \PYG{o}{=} \PYG{n}{numpy}\PYG{o}{.}\PYG{n}{empty}\PYG{p}{(} \PYG{p}{(}\PYG{n}{m}\PYG{p}{,} \PYG{n}{q}\PYG{p}{)}\PYG{p}{,} \PYG{n}{dtype}\PYG{o}{=}\PYG{n+nb}{float}\PYG{p}{)}
    \PYG{n}{yq1}\PYG{p}{[}\PYG{l+m+mi}{0}\PYG{p}{,} \PYG{l+m+mi}{0}\PYG{p}{]} \PYG{o}{=} \PYG{l+m+mf}{1.0}
    \PYG{n}{xq1}       \PYG{o}{=} \PYG{n}{f}\PYG{o}{.}\PYG{n}{reverse}\PYG{p}{(}\PYG{n}{q}\PYG{p}{,} \PYG{n}{yq1}\PYG{p}{)}
    \PYG{c+c1}{\PYGZsh{} partial G w.r.t x\PYGZus{}0}
    \PYG{n}{ok} \PYG{o}{=} \PYG{n}{ok} \PYG{o+ow}{and} \PYG{n}{xq1}\PYG{p}{[}\PYG{l+m+mi}{0}\PYG{p}{,}\PYG{l+m+mi}{0}\PYG{p}{]} \PYG{o}{==} \PYG{l+m+mf}{3.0} \PYG{o}{*} \PYG{l+m+mf}{4.0}
    \PYG{c+c1}{\PYGZsh{} partial G w.r.t x\PYGZus{}1}
    \PYG{n}{ok} \PYG{o}{=} \PYG{n}{ok} \PYG{o+ow}{and} \PYG{n}{xq1}\PYG{p}{[}\PYG{l+m+mi}{1}\PYG{p}{,}\PYG{l+m+mi}{0}\PYG{p}{]} \PYG{o}{==} \PYG{l+m+mf}{2.0} \PYG{o}{*} \PYG{l+m+mf}{4.0}
    \PYG{c+c1}{\PYGZsh{} partial G w.r.t x\PYGZus{}2}
    \PYG{n}{ok} \PYG{o}{=} \PYG{n}{ok} \PYG{o+ow}{and} \PYG{n}{xq1}\PYG{p}{[}\PYG{l+m+mi}{2}\PYG{p}{,}\PYG{l+m+mi}{0}\PYG{p}{]} \PYG{o}{==} \PYG{l+m+mf}{2.0} \PYG{o}{*} \PYG{l+m+mf}{3.0}
    \PYG{c+c1}{\PYGZsh{} \PYGZhy{}\PYGZhy{}\PYGZhy{}\PYGZhy{}\PYGZhy{}\PYGZhy{}\PYGZhy{}\PYGZhy{}\PYGZhy{}\PYGZhy{}\PYGZhy{}\PYGZhy{}\PYGZhy{}\PYGZhy{}\PYGZhy{}\PYGZhy{}\PYGZhy{}\PYGZhy{}\PYGZhy{}\PYGZhy{}\PYGZhy{}\PYGZhy{}\PYGZhy{}\PYGZhy{}\PYGZhy{}\PYGZhy{}\PYGZhy{}\PYGZhy{}\PYGZhy{}\PYGZhy{}\PYGZhy{}\PYGZhy{}\PYGZhy{}\PYGZhy{}\PYGZhy{}\PYGZhy{}\PYGZhy{}\PYGZhy{}\PYGZhy{}\PYGZhy{}\PYGZhy{}\PYGZhy{}\PYGZhy{}\PYGZhy{}\PYGZhy{}\PYGZhy{}\PYGZhy{}\PYGZhy{}\PYGZhy{}\PYGZhy{}\PYGZhy{}\PYGZhy{}\PYGZhy{}\PYGZhy{}\PYGZhy{}\PYGZhy{}\PYGZhy{}\PYGZhy{}\PYGZhy{}\PYGZhy{}\PYGZhy{}\PYGZhy{}\PYGZhy{}\PYGZhy{}\PYGZhy{}\PYGZhy{}\PYGZhy{}\PYGZhy{}\PYGZhy{}\PYGZhy{}\PYGZhy{}}
    \PYG{c+c1}{\PYGZsh{} first order forward mode}
    \PYG{n}{p}     \PYG{o}{=} \PYG{l+m+mi}{1}
    \PYG{n}{xp}\PYG{p}{[}\PYG{l+m+mi}{0}\PYG{p}{]} \PYG{o}{=} \PYG{l+m+mf}{1.0}
    \PYG{n}{xp}\PYG{p}{[}\PYG{l+m+mi}{1}\PYG{p}{]} \PYG{o}{=} \PYG{l+m+mf}{1.0}
    \PYG{n}{xp}\PYG{p}{[}\PYG{l+m+mi}{2}\PYG{p}{]} \PYG{o}{=} \PYG{l+m+mf}{1.0}
    \PYG{n}{yp}    \PYG{o}{=} \PYG{n}{f}\PYG{o}{.}\PYG{n}{forward}\PYG{p}{(}\PYG{n}{p}\PYG{p}{,} \PYG{n}{xp}\PYG{p}{)}
    \PYG{n}{ok}    \PYG{o}{=} \PYG{n}{ok} \PYG{o+ow}{and} \PYG{n}{yp}\PYG{p}{[}\PYG{l+m+mi}{0}\PYG{p}{]} \PYG{o}{==} \PYG{l+m+mf}{3.0}\PYG{o}{*}\PYG{l+m+mf}{4.0} \PYG{o}{+} \PYG{l+m+mf}{2.0}\PYG{o}{*}\PYG{l+m+mf}{4.0} \PYG{o}{+} \PYG{l+m+mf}{2.0}\PYG{o}{*}\PYG{l+m+mf}{3.0}
    \PYG{c+c1}{\PYGZsh{} \PYGZhy{}\PYGZhy{}\PYGZhy{}\PYGZhy{}\PYGZhy{}\PYGZhy{}\PYGZhy{}\PYGZhy{}\PYGZhy{}\PYGZhy{}\PYGZhy{}\PYGZhy{}\PYGZhy{}\PYGZhy{}\PYGZhy{}\PYGZhy{}\PYGZhy{}\PYGZhy{}\PYGZhy{}\PYGZhy{}\PYGZhy{}\PYGZhy{}\PYGZhy{}\PYGZhy{}\PYGZhy{}\PYGZhy{}\PYGZhy{}\PYGZhy{}\PYGZhy{}\PYGZhy{}\PYGZhy{}\PYGZhy{}\PYGZhy{}\PYGZhy{}\PYGZhy{}\PYGZhy{}\PYGZhy{}\PYGZhy{}\PYGZhy{}\PYGZhy{}\PYGZhy{}\PYGZhy{}\PYGZhy{}\PYGZhy{}\PYGZhy{}\PYGZhy{}\PYGZhy{}\PYGZhy{}\PYGZhy{}\PYGZhy{}\PYGZhy{}\PYGZhy{}\PYGZhy{}\PYGZhy{}\PYGZhy{}\PYGZhy{}\PYGZhy{}\PYGZhy{}\PYGZhy{}\PYGZhy{}\PYGZhy{}\PYGZhy{}\PYGZhy{}\PYGZhy{}\PYGZhy{}\PYGZhy{}\PYGZhy{}\PYGZhy{}\PYGZhy{}\PYGZhy{}\PYGZhy{}}
    \PYG{c+c1}{\PYGZsh{} second order reverse (derivative of first order forward)}
    \PYG{c+c1}{\PYGZsh{} define G( y\PYGZus{}0\PYGZca{}0 , y\PYGZus{}0\PYGZca{}1 ) = y\PYGZus{}0\PYGZca{}1}
    \PYG{c+c1}{\PYGZsh{} = x\PYGZus{}1\PYGZca{}0 * x\PYGZus{}2\PYGZca{}0  +  x\PYGZus{}0\PYGZca{}0 * x\PYGZus{}2\PYGZca{}0  +  x\PYGZus{}0\PYGZca{}0  *  x\PYGZus{}1\PYGZca{}0}
    \PYG{n}{q}         \PYG{o}{=} \PYG{l+m+mi}{2}
    \PYG{n}{yq2}       \PYG{o}{=} \PYG{n}{numpy}\PYG{o}{.}\PYG{n}{empty}\PYG{p}{(} \PYG{p}{(}\PYG{n}{m}\PYG{p}{,} \PYG{n}{q}\PYG{p}{)}\PYG{p}{,} \PYG{n}{dtype}\PYG{o}{=}\PYG{n+nb}{float}\PYG{p}{)}
    \PYG{n}{yq2}\PYG{p}{[}\PYG{l+m+mi}{0}\PYG{p}{,} \PYG{l+m+mi}{0}\PYG{p}{]} \PYG{o}{=} \PYG{l+m+mf}{0.0} \PYG{c+c1}{\PYGZsh{} partial of G w.r.t y\PYGZus{}0\PYGZca{}0}
    \PYG{n}{yq2}\PYG{p}{[}\PYG{l+m+mi}{0}\PYG{p}{,} \PYG{l+m+mi}{1}\PYG{p}{]} \PYG{o}{=} \PYG{l+m+mf}{1.0} \PYG{c+c1}{\PYGZsh{} partial of G w.r.t y\PYGZus{}0\PYGZca{}1}
    \PYG{n}{xq2}       \PYG{o}{=} \PYG{n}{f}\PYG{o}{.}\PYG{n}{reverse}\PYG{p}{(}\PYG{n}{q}\PYG{p}{,} \PYG{n}{yq2}\PYG{p}{)}
    \PYG{c+c1}{\PYGZsh{} partial G w.r.t x\PYGZus{}0\PYGZca{}0}
    \PYG{n}{ok} \PYG{o}{=} \PYG{n}{ok} \PYG{o+ow}{and} \PYG{n}{xq2}\PYG{p}{[}\PYG{l+m+mi}{0}\PYG{p}{,} \PYG{l+m+mi}{0}\PYG{p}{]} \PYG{o}{==} \PYG{l+m+mf}{3.0} \PYG{o}{+} \PYG{l+m+mf}{4.0}
    \PYG{c+c1}{\PYGZsh{} partial G w.r.t x\PYGZus{}1\PYGZca{}0}
    \PYG{n}{ok} \PYG{o}{=} \PYG{n}{ok} \PYG{o+ow}{and} \PYG{n}{xq2}\PYG{p}{[}\PYG{l+m+mi}{1}\PYG{p}{,} \PYG{l+m+mi}{0}\PYG{p}{]} \PYG{o}{==} \PYG{l+m+mf}{2.0} \PYG{o}{+} \PYG{l+m+mf}{4.0}
    \PYG{c+c1}{\PYGZsh{} partial G w.r.t x\PYGZus{}2\PYGZca{}0}
    \PYG{n}{ok} \PYG{o}{=} \PYG{n}{ok} \PYG{o+ow}{and} \PYG{n}{xq2}\PYG{p}{[}\PYG{l+m+mi}{2}\PYG{p}{,} \PYG{l+m+mi}{0}\PYG{p}{]} \PYG{o}{==} \PYG{l+m+mf}{2.0} \PYG{o}{+} \PYG{l+m+mf}{3.0}
    \PYG{c+c1}{\PYGZsh{} \PYGZhy{}\PYGZhy{}\PYGZhy{}\PYGZhy{}\PYGZhy{}\PYGZhy{}\PYGZhy{}\PYGZhy{}\PYGZhy{}\PYGZhy{}\PYGZhy{}\PYGZhy{}\PYGZhy{}\PYGZhy{}\PYGZhy{}\PYGZhy{}\PYGZhy{}\PYGZhy{}\PYGZhy{}\PYGZhy{}\PYGZhy{}\PYGZhy{}\PYGZhy{}\PYGZhy{}\PYGZhy{}\PYGZhy{}\PYGZhy{}\PYGZhy{}\PYGZhy{}\PYGZhy{}\PYGZhy{}\PYGZhy{}\PYGZhy{}\PYGZhy{}\PYGZhy{}\PYGZhy{}\PYGZhy{}\PYGZhy{}\PYGZhy{}\PYGZhy{}\PYGZhy{}\PYGZhy{}\PYGZhy{}\PYGZhy{}\PYGZhy{}\PYGZhy{}\PYGZhy{}\PYGZhy{}\PYGZhy{}\PYGZhy{}\PYGZhy{}\PYGZhy{}\PYGZhy{}\PYGZhy{}\PYGZhy{}\PYGZhy{}\PYGZhy{}\PYGZhy{}\PYGZhy{}\PYGZhy{}\PYGZhy{}\PYGZhy{}\PYGZhy{}\PYGZhy{}\PYGZhy{}\PYGZhy{}\PYGZhy{}\PYGZhy{}\PYGZhy{}\PYGZhy{}\PYGZhy{}}
    \PYG{n}{af} \PYG{o}{=} \PYG{n}{cppad\PYGZus{}py}\PYG{o}{.}\PYG{n}{a\PYGZus{}fun}\PYG{p}{(}\PYG{n}{f}\PYG{p}{)}
    \PYG{c+c1}{\PYGZsh{}}
    \PYG{c+c1}{\PYGZsh{} zero order forward}
    \PYG{n}{axp}   \PYG{o}{=} \PYG{n}{numpy}\PYG{o}{.}\PYG{n}{empty}\PYG{p}{(}\PYG{n}{n\PYGZus{}ind}\PYG{p}{,} \PYG{n}{dtype}\PYG{o}{=}\PYG{n}{cppad\PYGZus{}py}\PYG{o}{.}\PYG{n}{a\PYGZus{}double}\PYG{p}{)}
    \PYG{n}{p}     \PYG{o}{=} \PYG{l+m+mi}{0}
    \PYG{n}{axp}\PYG{p}{[}\PYG{l+m+mi}{0}\PYG{p}{]} \PYG{o}{=} \PYG{l+m+mf}{2.0}
    \PYG{n}{axp}\PYG{p}{[}\PYG{l+m+mi}{1}\PYG{p}{]} \PYG{o}{=} \PYG{l+m+mf}{3.0}
    \PYG{n}{axp}\PYG{p}{[}\PYG{l+m+mi}{2}\PYG{p}{]} \PYG{o}{=} \PYG{l+m+mf}{4.0}
    \PYG{n}{ayp} \PYG{o}{=} \PYG{n}{af}\PYG{o}{.}\PYG{n}{forward}\PYG{p}{(}\PYG{n}{p}\PYG{p}{,} \PYG{n}{axp}\PYG{p}{)}
    \PYG{n}{ok} \PYG{o}{=} \PYG{n}{ok} \PYG{o+ow}{and} \PYG{n}{ayp}\PYG{p}{[}\PYG{l+m+mi}{0}\PYG{p}{]} \PYG{o}{==} \PYG{n}{cppad\PYGZus{}py}\PYG{o}{.}\PYG{n}{a\PYGZus{}double}\PYG{p}{(}\PYG{l+m+mf}{24.0}\PYG{p}{)}
    \PYG{c+c1}{\PYGZsh{}}
    \PYG{c+c1}{\PYGZsh{} first order reverse}
    \PYG{n}{q}          \PYG{o}{=} \PYG{l+m+mi}{1}
    \PYG{n}{ayq1}       \PYG{o}{=} \PYG{n}{numpy}\PYG{o}{.}\PYG{n}{empty}\PYG{p}{(} \PYG{p}{(}\PYG{n}{m}\PYG{p}{,} \PYG{n}{q}\PYG{p}{)}\PYG{p}{,} \PYG{n}{dtype}\PYG{o}{=}\PYG{n}{cppad\PYGZus{}py}\PYG{o}{.}\PYG{n}{a\PYGZus{}double}\PYG{p}{)}
    \PYG{n}{ayq1}\PYG{p}{[}\PYG{l+m+mi}{0}\PYG{p}{,} \PYG{l+m+mi}{0}\PYG{p}{]} \PYG{o}{=} \PYG{l+m+mf}{1.0}
    \PYG{n}{axq1}       \PYG{o}{=} \PYG{n}{af}\PYG{o}{.}\PYG{n}{reverse}\PYG{p}{(}\PYG{n}{q}\PYG{p}{,} \PYG{n}{ayq1}\PYG{p}{)}
    \PYG{c+c1}{\PYGZsh{} partial G w.r.t x\PYGZus{}0}
    \PYG{n}{ok} \PYG{o}{=} \PYG{n}{ok} \PYG{o+ow}{and} \PYG{n}{axq1}\PYG{p}{[}\PYG{l+m+mi}{0}\PYG{p}{,}\PYG{l+m+mi}{0}\PYG{p}{]} \PYG{o}{==} \PYG{n}{cppad\PYGZus{}py}\PYG{o}{.}\PYG{n}{a\PYGZus{}double}\PYG{p}{(}\PYG{l+m+mf}{3.0} \PYG{o}{*} \PYG{l+m+mf}{4.0}\PYG{p}{)}
    \PYG{c+c1}{\PYGZsh{} partial G w.r.t x\PYGZus{}1}
    \PYG{n}{ok} \PYG{o}{=} \PYG{n}{ok} \PYG{o+ow}{and} \PYG{n}{axq1}\PYG{p}{[}\PYG{l+m+mi}{1}\PYG{p}{,}\PYG{l+m+mi}{0}\PYG{p}{]} \PYG{o}{==} \PYG{n}{cppad\PYGZus{}py}\PYG{o}{.}\PYG{n}{a\PYGZus{}double}\PYG{p}{(}\PYG{l+m+mf}{2.0} \PYG{o}{*} \PYG{l+m+mf}{4.0}\PYG{p}{)}
    \PYG{c+c1}{\PYGZsh{} partial G w.r.t x\PYGZus{}2}
    \PYG{n}{ok} \PYG{o}{=} \PYG{n}{ok} \PYG{o+ow}{and} \PYG{n}{axq1}\PYG{p}{[}\PYG{l+m+mi}{2}\PYG{p}{,}\PYG{l+m+mi}{0}\PYG{p}{]} \PYG{o}{==} \PYG{n}{cppad\PYGZus{}py}\PYG{o}{.}\PYG{n}{a\PYGZus{}double}\PYG{p}{(}\PYG{l+m+mf}{2.0} \PYG{o}{*} \PYG{l+m+mf}{3.0}\PYG{p}{)}
    \PYG{c+c1}{\PYGZsh{}}
    \PYG{k}{return}\PYG{p}{(} \PYG{n}{ok} \PYG{p}{)}
\PYG{c+c1}{\PYGZsh{}}
\end{sphinxVerbatim}


\bigskip\hrule\bigskip


xsrst input file: \sphinxcode{\sphinxupquote{example/python/core/fun\_reverse\_xam.py}}

\index{example@\spxentry{example}}\ignorespaces 

\subparagraph{Example}
\label{\detokenize{xsrst/py_fun_reverse:example}}\label{\detokenize{xsrst/py_fun_reverse:py-fun-reverse-example}}\label{\detokenize{xsrst/py_fun_reverse:index-11}}
{\hyperref[\detokenize{xsrst/fun_reverse_xam_py:id1}]{\sphinxcrossref{\DUrole{std,std-ref}{fun\_reverse\_xam\_py}}}}


\bigskip\hrule\bigskip


xsrst input file: \sphinxcode{\sphinxupquote{lib/python/cppad\_py/fun\_reverse.py}}


\subparagraph{py\_fun\_optimize}
\label{\detokenize{xsrst/py_fun_optimize:py-fun-optimize}}\label{\detokenize{xsrst/py_fun_optimize::doc}}
\index{py\_fun\_optimize@\spxentry{py\_fun\_optimize}}\index{optimize@\spxentry{optimize}}\index{an@\spxentry{an}}\index{ad@\spxentry{ad}}\index{function@\spxentry{function}}\ignorespaces 

\subparagraph{3.1.1.10 Optimize an AD Function}
\label{\detokenize{xsrst/py_fun_optimize:optimize-an-ad-function}}\label{\detokenize{xsrst/py_fun_optimize:id1}}\begin{itemize}
\item {} 
{\hyperref[\detokenize{xsrst/py_fun_optimize:py-fun-optimize-syntax}]{\sphinxcrossref{\DUrole{std,std-ref}{Syntax}}}}

\item {} 
{\hyperref[\detokenize{xsrst/py_fun_optimize:py-fun-optimize-purpose}]{\sphinxcrossref{\DUrole{std,std-ref}{Purpose}}}}

\item {} 
{\hyperref[\detokenize{xsrst/py_fun_optimize:py-fun-optimize-f}]{\sphinxcrossref{\DUrole{std,std-ref}{f}}}}

\item {} 
{\hyperref[\detokenize{xsrst/py_fun_optimize:py-fun-optimize-example}]{\sphinxcrossref{\DUrole{std,std-ref}{Example}}}}

\end{itemize}

\index{syntax@\spxentry{syntax}}\ignorespaces 

\subparagraph{Syntax}
\label{\detokenize{xsrst/py_fun_optimize:syntax}}\label{\detokenize{xsrst/py_fun_optimize:py-fun-optimize-syntax}}\label{\detokenize{xsrst/py_fun_optimize:index-1}}
\sphinxstyleemphasis{f} . \sphinxcode{\sphinxupquote{optimize}} ()

\index{purpose@\spxentry{purpose}}\ignorespaces 

\subparagraph{Purpose}
\label{\detokenize{xsrst/py_fun_optimize:purpose}}\label{\detokenize{xsrst/py_fun_optimize:py-fun-optimize-purpose}}\label{\detokenize{xsrst/py_fun_optimize:index-2}}
This reduces the number of operations
(hence to time and memory) used to compute the function
stored in \sphinxstyleemphasis{f}
On the other hand, the optimization may take a significant amount
of time and memory.

\index{f@\spxentry{f}}\ignorespaces 

\subparagraph{f}
\label{\detokenize{xsrst/py_fun_optimize:f}}\label{\detokenize{xsrst/py_fun_optimize:py-fun-optimize-f}}\label{\detokenize{xsrst/py_fun_optimize:index-3}}
This is a
{\hyperref[\detokenize{xsrst/cpp_fun_ctor:cpp-fun-ctor-syntax-d-fun}]{\sphinxcrossref{\DUrole{std,std-ref}{d\_fun}}}}.
Optimizing this \sphinxstyleemphasis{f} also optimizes the
corresponding {\hyperref[\detokenize{xsrst/py_fun_ctor:py-fun-ctor-syntax-a-fun}]{\sphinxcrossref{\DUrole{std,std-ref}{a\_fun}}}}.


\subparagraph{fun\_optimize\_xam\_py}
\label{\detokenize{xsrst/fun_optimize_xam_py:fun-optimize-xam-py}}\label{\detokenize{xsrst/fun_optimize_xam_py::doc}}
\index{fun\_optimize\_xam\_py@\spxentry{fun\_optimize\_xam\_py}}\index{python:@\spxentry{python:}}\index{optimize@\spxentry{optimize}}\index{an@\spxentry{an}}\index{d\_fun:@\spxentry{d\_fun:}}\index{example@\spxentry{example}}\index{test@\spxentry{test}}\ignorespaces 

\subparagraph{3.1.1.10.1 Python: Optimize an d\_fun: Example and Test}
\label{\detokenize{xsrst/fun_optimize_xam_py:python-optimize-an-d-fun-example-and-test}}\label{\detokenize{xsrst/fun_optimize_xam_py:id1}}
\begin{sphinxVerbatim}[commandchars=\\\{\}]
\PYG{k}{def} \PYG{n+nf}{fun\PYGZus{}optimize\PYGZus{}xam}\PYG{p}{(}\PYG{p}{)} \PYG{p}{:}
    \PYG{c+c1}{\PYGZsh{}}
    \PYG{k+kn}{import} \PYG{n+nn}{numpy}
    \PYG{k+kn}{import} \PYG{n+nn}{cppad\PYGZus{}py}
    \PYG{c+c1}{\PYGZsh{}}
    \PYG{c+c1}{\PYGZsh{} initialize return variable}
    \PYG{n}{ok} \PYG{o}{=} \PYG{k+kc}{True}
    \PYG{c+c1}{\PYGZsh{} \PYGZhy{}\PYGZhy{}\PYGZhy{}\PYGZhy{}\PYGZhy{}\PYGZhy{}\PYGZhy{}\PYGZhy{}\PYGZhy{}\PYGZhy{}\PYGZhy{}\PYGZhy{}\PYGZhy{}\PYGZhy{}\PYGZhy{}\PYGZhy{}\PYGZhy{}\PYGZhy{}\PYGZhy{}\PYGZhy{}\PYGZhy{}\PYGZhy{}\PYGZhy{}\PYGZhy{}\PYGZhy{}\PYGZhy{}\PYGZhy{}\PYGZhy{}\PYGZhy{}\PYGZhy{}\PYGZhy{}\PYGZhy{}\PYGZhy{}\PYGZhy{}\PYGZhy{}\PYGZhy{}\PYGZhy{}\PYGZhy{}\PYGZhy{}\PYGZhy{}\PYGZhy{}\PYGZhy{}\PYGZhy{}\PYGZhy{}\PYGZhy{}\PYGZhy{}\PYGZhy{}\PYGZhy{}\PYGZhy{}\PYGZhy{}\PYGZhy{}\PYGZhy{}\PYGZhy{}\PYGZhy{}\PYGZhy{}\PYGZhy{}\PYGZhy{}\PYGZhy{}\PYGZhy{}\PYGZhy{}\PYGZhy{}\PYGZhy{}\PYGZhy{}\PYGZhy{}\PYGZhy{}\PYGZhy{}\PYGZhy{}\PYGZhy{}\PYGZhy{}}
    \PYG{n}{n\PYGZus{}ind} \PYG{o}{=} \PYG{l+m+mi}{1} \PYG{c+c1}{\PYGZsh{} number of independent variables}
    \PYG{n}{n\PYGZus{}dep} \PYG{o}{=} \PYG{l+m+mi}{1} \PYG{c+c1}{\PYGZsh{} number of dependent variables}
    \PYG{n}{n\PYGZus{}var} \PYG{o}{=} \PYG{l+m+mi}{1} \PYG{c+c1}{\PYGZsh{} phantom variable at address 0}
    \PYG{n}{n\PYGZus{}op}  \PYG{o}{=} \PYG{l+m+mi}{1} \PYG{c+c1}{\PYGZsh{} special operator at beginning}
    \PYG{c+c1}{\PYGZsh{}}
    \PYG{c+c1}{\PYGZsh{} dimension some vectors}
    \PYG{n}{x}  \PYG{o}{=} \PYG{n}{numpy}\PYG{o}{.}\PYG{n}{empty}\PYG{p}{(}\PYG{n}{n\PYGZus{}ind}\PYG{p}{,} \PYG{n}{dtype}\PYG{o}{=}\PYG{n+nb}{float}\PYG{p}{)}
    \PYG{n}{ay} \PYG{o}{=} \PYG{n}{numpy}\PYG{o}{.}\PYG{n}{empty}\PYG{p}{(}\PYG{n}{n\PYGZus{}dep}\PYG{p}{,} \PYG{n}{dtype}\PYG{o}{=}\PYG{n}{cppad\PYGZus{}py}\PYG{o}{.}\PYG{n}{a\PYGZus{}double}\PYG{p}{)}
    \PYG{c+c1}{\PYGZsh{}}
    \PYG{c+c1}{\PYGZsh{} independent variables}
    \PYG{n}{x}\PYG{p}{[}\PYG{l+m+mi}{0}\PYG{p}{]}  \PYG{o}{=} \PYG{l+m+mf}{1.0}
    \PYG{n}{ax}    \PYG{o}{=} \PYG{n}{cppad\PYGZus{}py}\PYG{o}{.}\PYG{n}{independent}\PYG{p}{(}\PYG{n}{x}\PYG{p}{)}
    \PYG{n}{n\PYGZus{}var} \PYG{o}{=} \PYG{n}{n\PYGZus{}var} \PYG{o}{+} \PYG{n}{n\PYGZus{}ind} \PYG{c+c1}{\PYGZsh{} one for each indpendent}
    \PYG{n}{n\PYGZus{}op}  \PYG{o}{=} \PYG{n}{n\PYGZus{}op} \PYG{o}{+} \PYG{n}{n\PYGZus{}ind}
    \PYG{c+c1}{\PYGZsh{}}
    \PYG{c+c1}{\PYGZsh{} accumulate summation}
    \PYG{n}{ax0}   \PYG{o}{=} \PYG{n}{ax}\PYG{p}{[}\PYG{l+m+mi}{0}\PYG{p}{]}
    \PYG{n}{csum}  \PYG{o}{=} \PYG{n}{cppad\PYGZus{}py}\PYG{o}{.}\PYG{n}{a\PYGZus{}double}\PYG{p}{(}\PYG{l+m+mf}{0.0}\PYG{p}{)}
    \PYG{n}{csum}  \PYG{o}{=} \PYG{n}{ax0} \PYG{o}{+} \PYG{n}{ax0} \PYG{o}{+} \PYG{n}{ax0} \PYG{o}{+} \PYG{n}{ax0}
    \PYG{n}{n\PYGZus{}var} \PYG{o}{=} \PYG{n}{n\PYGZus{}var} \PYG{o}{+} \PYG{l+m+mi}{3} \PYG{c+c1}{\PYGZsh{} one per + operator}
    \PYG{n}{n\PYGZus{}op}  \PYG{o}{=} \PYG{n}{n\PYGZus{}op} \PYG{o}{+} \PYG{l+m+mi}{3}
    \PYG{c+c1}{\PYGZsh{}}
    \PYG{c+c1}{\PYGZsh{} define f(x) = y\PYGZus{}0 = csum}
    \PYG{n}{ay}\PYG{p}{[}\PYG{l+m+mi}{0}\PYG{p}{]} \PYG{o}{=} \PYG{n}{csum}
    \PYG{n}{f}     \PYG{o}{=} \PYG{n}{cppad\PYGZus{}py}\PYG{o}{.}\PYG{n}{d\PYGZus{}fun}\PYG{p}{(}\PYG{n}{ax}\PYG{p}{,} \PYG{n}{ay}\PYG{p}{)}
    \PYG{n}{n\PYGZus{}op}  \PYG{o}{=} \PYG{n}{n\PYGZus{}op} \PYG{o}{+} \PYG{l+m+mi}{1} \PYG{c+c1}{\PYGZsh{} speical operator at end}
    \PYG{c+c1}{\PYGZsh{}}
    \PYG{c+c1}{\PYGZsh{} check number of variables and operators}
    \PYG{n}{ok} \PYG{o}{=} \PYG{n}{ok} \PYG{o+ow}{and} \PYG{n}{f}\PYG{o}{.}\PYG{n}{size\PYGZus{}var}\PYG{p}{(}\PYG{p}{)} \PYG{o}{==} \PYG{n}{n\PYGZus{}var}
    \PYG{n}{ok} \PYG{o}{=} \PYG{n}{ok} \PYG{o+ow}{and} \PYG{n}{f}\PYG{o}{.}\PYG{n}{size\PYGZus{}op}\PYG{p}{(}\PYG{p}{)} \PYG{o}{==} \PYG{n}{n\PYGZus{}op}
    \PYG{c+c1}{\PYGZsh{}}
    \PYG{c+c1}{\PYGZsh{} optimize}
    \PYG{n}{f}\PYG{o}{.}\PYG{n}{optimize}\PYG{p}{(}\PYG{p}{)}
    \PYG{c+c1}{\PYGZsh{}}
    \PYG{c+c1}{\PYGZsh{} number of variables and operators has decreased by two}
    \PYG{n}{ok} \PYG{o}{=} \PYG{n}{ok} \PYG{o+ow}{and} \PYG{n}{f}\PYG{o}{.}\PYG{n}{size\PYGZus{}var}\PYG{p}{(}\PYG{p}{)} \PYG{o}{==} \PYG{n}{n\PYGZus{}var}\PYG{o}{\PYGZhy{}}\PYG{l+m+mi}{2}
    \PYG{n}{ok} \PYG{o}{=} \PYG{n}{ok} \PYG{o+ow}{and} \PYG{n}{f}\PYG{o}{.}\PYG{n}{size\PYGZus{}op}\PYG{p}{(}\PYG{p}{)} \PYG{o}{==} \PYG{n}{n\PYGZus{}op}\PYG{o}{\PYGZhy{}}\PYG{l+m+mi}{2}
    \PYG{c+c1}{\PYGZsh{}}
    \PYG{k}{return}\PYG{p}{(} \PYG{n}{ok}  \PYG{p}{)}
\PYG{c+c1}{\PYGZsh{}}
\end{sphinxVerbatim}


\bigskip\hrule\bigskip


xsrst input file: \sphinxcode{\sphinxupquote{example/python/core/fun\_optimize\_xam.py}}

\index{example@\spxentry{example}}\ignorespaces 

\subparagraph{Example}
\label{\detokenize{xsrst/py_fun_optimize:example}}\label{\detokenize{xsrst/py_fun_optimize:py-fun-optimize-example}}\label{\detokenize{xsrst/py_fun_optimize:index-4}}
{\hyperref[\detokenize{xsrst/fun_optimize_xam_py:id1}]{\sphinxcrossref{\DUrole{std,std-ref}{fun\_optimize\_xam\_py}}}}


\bigskip\hrule\bigskip


xsrst input file: \sphinxcode{\sphinxupquote{lib/python/cppad\_py/fun\_optimize.xsrst}}


\subparagraph{py\_fun\_json}
\label{\detokenize{xsrst/py_fun_json:py-fun-json}}\label{\detokenize{xsrst/py_fun_json::doc}}
\index{py\_fun\_json@\spxentry{py\_fun\_json}}\index{json@\spxentry{json}}\index{representation@\spxentry{representation}}\index{ad@\spxentry{ad}}\index{computation@\spxentry{computation}}\index{graph@\spxentry{graph}}\ignorespaces 

\subparagraph{3.1.1.11 Json Representation of AD Computation Graph}
\label{\detokenize{xsrst/py_fun_json:json-representation-of-ad-computation-graph}}\label{\detokenize{xsrst/py_fun_json:id1}}\begin{itemize}
\item {} 
{\hyperref[\detokenize{xsrst/py_fun_json:py-fun-json-syntax}]{\sphinxcrossref{\DUrole{std,std-ref}{Syntax}}}}

\item {} 
{\hyperref[\detokenize{xsrst/py_fun_json:py-fun-json-f}]{\sphinxcrossref{\DUrole{std,std-ref}{f}}}}

\item {} 
{\hyperref[\detokenize{xsrst/py_fun_json:py-fun-json-json}]{\sphinxcrossref{\DUrole{std,std-ref}{json}}}}

\item {} 
{\hyperref[\detokenize{xsrst/py_fun_json:py-fun-json-to-json}]{\sphinxcrossref{\DUrole{std,std-ref}{to\_json}}}}

\item {} 
{\hyperref[\detokenize{xsrst/py_fun_json:py-fun-json-from-json}]{\sphinxcrossref{\DUrole{std,std-ref}{from\_json}}}}

\item {} 
{\hyperref[\detokenize{xsrst/py_fun_json:py-fun-json-examples}]{\sphinxcrossref{\DUrole{std,std-ref}{Examples}}}}

\end{itemize}

\index{syntax@\spxentry{syntax}}\ignorespaces 

\subparagraph{Syntax}
\label{\detokenize{xsrst/py_fun_json:syntax}}\label{\detokenize{xsrst/py_fun_json:py-fun-json-syntax}}\label{\detokenize{xsrst/py_fun_json:index-1}}
\begin{DUlineblock}{0em}
\item[] \sphinxstyleemphasis{json} = \sphinxstyleemphasis{f} . \sphinxcode{\sphinxupquote{to\_json}} ()
\item[] \sphinxstyleemphasis{f} . \sphinxcode{\sphinxupquote{from\_json(json}} )
\end{DUlineblock}

\index{f@\spxentry{f}}\ignorespaces 

\subparagraph{f}
\label{\detokenize{xsrst/py_fun_json:f}}\label{\detokenize{xsrst/py_fun_json:py-fun-json-f}}\label{\detokenize{xsrst/py_fun_json:index-2}}
This is a {\hyperref[\detokenize{xsrst/py_fun_ctor:py-fun-ctor-syntax-d-fun}]{\sphinxcrossref{\DUrole{std,std-ref}{d\_fun}}}} function object.

\index{json@\spxentry{json}}\ignorespaces 

\subparagraph{json}
\label{\detokenize{xsrst/py_fun_json:json}}\label{\detokenize{xsrst/py_fun_json:py-fun-json-json}}\label{\detokenize{xsrst/py_fun_json:index-3}}
is a \sphinxcode{\sphinxupquote{str}} containing
a Json representation of the computation graph corresponding to
\sphinxstyleemphasis{f} ; see the CppAD documentation for
\sphinxhref{https://coin-or.github.io/CppAD/doc/json\_ad\_graph.htm}{json\_ad\_graph}.

\index{to\_json@\spxentry{to\_json}}\ignorespaces 

\subparagraph{to\_json}
\label{\detokenize{xsrst/py_fun_json:to-json}}\label{\detokenize{xsrst/py_fun_json:py-fun-json-to-json}}\label{\detokenize{xsrst/py_fun_json:index-4}}
In this case, the function object \sphinxstyleemphasis{f} is constant and
the return value \sphinxstyleemphasis{json} is created.

\index{from\_json@\spxentry{from\_json}}\ignorespaces 

\subparagraph{from\_json}
\label{\detokenize{xsrst/py_fun_json:from-json}}\label{\detokenize{xsrst/py_fun_json:py-fun-json-from-json}}\label{\detokenize{xsrst/py_fun_json:index-5}}
In this case, the argument \sphinxstyleemphasis{json} is constant and
the function \sphinxstyleemphasis{f} is changed so it corresponds to \sphinxstyleemphasis{json} .


\subparagraph{fun\_to\_json\_xam\_py}
\label{\detokenize{xsrst/fun_to_json_xam_py:fun-to-json-xam-py}}\label{\detokenize{xsrst/fun_to_json_xam_py::doc}}
\index{fun\_to\_json\_xam\_py@\spxentry{fun\_to\_json\_xam\_py}}\index{python@\spxentry{python}}\index{to\_json:@\spxentry{to\_json:}}\index{example@\spxentry{example}}\index{test@\spxentry{test}}\ignorespaces 

\subparagraph{3.1.1.11.1 Python to\_json: Example and Test}
\label{\detokenize{xsrst/fun_to_json_xam_py:python-to-json-example-and-test}}\label{\detokenize{xsrst/fun_to_json_xam_py:id1}}
\begin{sphinxVerbatim}[commandchars=\\\{\}]
\PYG{k}{def} \PYG{n+nf}{to\PYGZus{}json\PYGZus{}xam}\PYG{p}{(}\PYG{p}{)} \PYG{p}{:}
    \PYG{c+c1}{\PYGZsh{}}
    \PYG{k+kn}{import} \PYG{n+nn}{numpy}
    \PYG{k+kn}{import} \PYG{n+nn}{cppad\PYGZus{}py}
    \PYG{k+kn}{import} \PYG{n+nn}{re}
    \PYG{c+c1}{\PYGZsh{}}
    \PYG{c+c1}{\PYGZsh{} initialize return variable}
    \PYG{n}{ok} \PYG{o}{=} \PYG{k+kc}{True}
    \PYG{c+c1}{\PYGZsh{} \PYGZhy{}\PYGZhy{}\PYGZhy{}\PYGZhy{}\PYGZhy{}\PYGZhy{}\PYGZhy{}\PYGZhy{}\PYGZhy{}\PYGZhy{}\PYGZhy{}\PYGZhy{}\PYGZhy{}\PYGZhy{}\PYGZhy{}\PYGZhy{}\PYGZhy{}\PYGZhy{}\PYGZhy{}\PYGZhy{}\PYGZhy{}\PYGZhy{}\PYGZhy{}\PYGZhy{}\PYGZhy{}\PYGZhy{}\PYGZhy{}\PYGZhy{}\PYGZhy{}\PYGZhy{}\PYGZhy{}\PYGZhy{}\PYGZhy{}\PYGZhy{}\PYGZhy{}\PYGZhy{}\PYGZhy{}\PYGZhy{}\PYGZhy{}\PYGZhy{}\PYGZhy{}\PYGZhy{}\PYGZhy{}\PYGZhy{}\PYGZhy{}\PYGZhy{}\PYGZhy{}\PYGZhy{}\PYGZhy{}\PYGZhy{}\PYGZhy{}\PYGZhy{}\PYGZhy{}\PYGZhy{}\PYGZhy{}\PYGZhy{}\PYGZhy{}\PYGZhy{}\PYGZhy{}\PYGZhy{}\PYGZhy{}\PYGZhy{}\PYGZhy{}\PYGZhy{}\PYGZhy{}\PYGZhy{}\PYGZhy{}\PYGZhy{}\PYGZhy{}}
    \PYG{n}{n} \PYG{o}{=} \PYG{l+m+mi}{1} \PYG{c+c1}{\PYGZsh{} number of independent variables}
    \PYG{n}{m} \PYG{o}{=} \PYG{l+m+mi}{2} \PYG{c+c1}{\PYGZsh{} number of dependent variables}
    \PYG{c+c1}{\PYGZsh{}}
    \PYG{c+c1}{\PYGZsh{} independent variables}
    \PYG{n}{x}  \PYG{o}{=} \PYG{n}{numpy}\PYG{o}{.}\PYG{n}{array}\PYG{p}{(} \PYG{p}{[} \PYG{l+m+mf}{1.0} \PYG{p}{]} \PYG{p}{)}
    \PYG{n}{ax} \PYG{o}{=} \PYG{n}{cppad\PYGZus{}py}\PYG{o}{.}\PYG{n}{independent}\PYG{p}{(}\PYG{n}{x}\PYG{p}{)}
    \PYG{c+c1}{\PYGZsh{}}
    \PYG{c+c1}{\PYGZsh{} f(x) = [ x0 + x0, sin(x0) ]}
    \PYG{n}{ay} \PYG{o}{=} \PYG{n}{numpy}\PYG{o}{.}\PYG{n}{empty}\PYG{p}{(}\PYG{n}{m}\PYG{p}{,} \PYG{n}{dtype}\PYG{o}{=}\PYG{n}{cppad\PYGZus{}py}\PYG{o}{.}\PYG{n}{a\PYGZus{}double}\PYG{p}{)}
    \PYG{n}{ay}\PYG{p}{[}\PYG{l+m+mi}{0}\PYG{p}{]} \PYG{o}{=} \PYG{n}{ax}\PYG{p}{[}\PYG{l+m+mi}{0}\PYG{p}{]} \PYG{o}{+} \PYG{n}{ax}\PYG{p}{[}\PYG{l+m+mi}{0}\PYG{p}{]}
    \PYG{n}{ay}\PYG{p}{[}\PYG{l+m+mi}{1}\PYG{p}{]} \PYG{o}{=} \PYG{n}{ax}\PYG{p}{[}\PYG{l+m+mi}{0}\PYG{p}{]}\PYG{o}{.}\PYG{n}{sin}\PYG{p}{(}\PYG{p}{)}
    \PYG{n}{f}     \PYG{o}{=} \PYG{n}{cppad\PYGZus{}py}\PYG{o}{.}\PYG{n}{d\PYGZus{}fun}\PYG{p}{(}\PYG{n}{ax}\PYG{p}{,} \PYG{n}{ay}\PYG{p}{)}
    \PYG{c+c1}{\PYGZsh{}}
    \PYG{c+c1}{\PYGZsh{} check f.to\PYGZus{}json}
    \PYG{n}{json}     \PYG{o}{=} \PYG{n}{f}\PYG{o}{.}\PYG{n}{to\PYGZus{}json}\PYG{p}{(}\PYG{p}{)}
    \PYG{n}{pattern}  \PYG{o}{=} \PYG{l+s+sa}{r}\PYG{l+s+s1}{\PYGZsq{}}\PYG{l+s+s1}{\PYGZdq{}}\PYG{l+s+s1}{op\PYGZus{}code}\PYG{l+s+s1}{\PYGZdq{}}\PYG{l+s+s1}{ *: *([\PYGZca{},]*),}\PYG{l+s+s1}{\PYGZsq{}}
    \PYG{n}{match}    \PYG{o}{=} \PYG{n}{re}\PYG{o}{.}\PYG{n}{search}\PYG{p}{(}\PYG{n}{pattern}\PYG{p}{,} \PYG{n}{json}\PYG{p}{)}
    \PYG{n}{op\PYGZus{}code}  \PYG{o}{=} \PYG{n+nb}{int}\PYG{p}{(} \PYG{n}{match}\PYG{o}{.}\PYG{n}{group}\PYG{p}{(}\PYG{l+m+mi}{1}\PYG{p}{)} \PYG{p}{)}
    \PYG{n}{ok}      \PYG{o}{\PYGZam{}}\PYG{o}{=} \PYG{n}{op\PYGZus{}code} \PYG{o}{==} \PYG{l+m+mi}{1}
    \PYG{n}{pattern}  \PYG{o}{=} \PYG{l+s+sa}{r}\PYG{l+s+s1}{\PYGZsq{}}\PYG{l+s+s1}{\PYGZdq{}}\PYG{l+s+s1}{name}\PYG{l+s+s1}{\PYGZdq{}}\PYG{l+s+s1}{ *: *}\PYG{l+s+s1}{\PYGZdq{}}\PYG{l+s+s1}{([\PYGZca{}}\PYG{l+s+s1}{\PYGZdq{}}\PYG{l+s+s1}{]*)}\PYG{l+s+s1}{\PYGZdq{}}\PYG{l+s+s1}{ *,}\PYG{l+s+s1}{\PYGZsq{}}
    \PYG{n}{match}    \PYG{o}{=} \PYG{n}{re}\PYG{o}{.}\PYG{n}{search}\PYG{p}{(}\PYG{n}{pattern}\PYG{p}{,} \PYG{n}{json}\PYG{p}{)}
    \PYG{n}{name}     \PYG{o}{=} \PYG{n}{match}\PYG{o}{.}\PYG{n}{group}\PYG{p}{(}\PYG{l+m+mi}{1}\PYG{p}{)}\PYG{p}{;}
    \PYG{n}{ok}      \PYG{o}{\PYGZam{}}\PYG{o}{=} \PYG{n}{name} \PYG{o}{==} \PYG{l+s+s1}{\PYGZsq{}}\PYG{l+s+s1}{add}\PYG{l+s+s1}{\PYGZsq{}} \PYG{o+ow}{or} \PYG{n}{name} \PYG{o}{==} \PYG{l+s+s1}{\PYGZsq{}}\PYG{l+s+s1}{sub}\PYG{l+s+s1}{\PYGZsq{}}
    \PYG{c+c1}{\PYGZsh{}}
    \PYG{k}{return} \PYG{n}{ok}
\end{sphinxVerbatim}


\bigskip\hrule\bigskip


xsrst input file: \sphinxcode{\sphinxupquote{example/python/core/fun\_json\_xam.py}}


\subparagraph{fun\_from\_json\_xam\_py}
\label{\detokenize{xsrst/fun_from_json_xam_py:fun-from-json-xam-py}}\label{\detokenize{xsrst/fun_from_json_xam_py::doc}}
\index{fun\_from\_json\_xam\_py@\spxentry{fun\_from\_json\_xam\_py}}\index{python@\spxentry{python}}\index{from\_json:@\spxentry{from\_json:}}\index{example@\spxentry{example}}\index{test@\spxentry{test}}\ignorespaces 

\subparagraph{3.1.1.11.2 Python from\_json: Example and Test}
\label{\detokenize{xsrst/fun_from_json_xam_py:python-from-json-example-and-test}}\label{\detokenize{xsrst/fun_from_json_xam_py:id1}}
\begin{sphinxVerbatim}[commandchars=\\\{\}]
\PYG{k}{def} \PYG{n+nf}{from\PYGZus{}json\PYGZus{}xam}\PYG{p}{(}\PYG{p}{)} \PYG{p}{:}
    \PYG{c+c1}{\PYGZsh{}}
    \PYG{k+kn}{import} \PYG{n+nn}{numpy}
    \PYG{k+kn}{import} \PYG{n+nn}{cppad\PYGZus{}py}
    \PYG{k+kn}{import} \PYG{n+nn}{math}
    \PYG{c+c1}{\PYGZsh{}}
    \PYG{c+c1}{\PYGZsh{} initialize return variable}
    \PYG{n}{ok} \PYG{o}{=} \PYG{k+kc}{True}
    \PYG{c+c1}{\PYGZsh{} \PYGZhy{}\PYGZhy{}\PYGZhy{}\PYGZhy{}\PYGZhy{}\PYGZhy{}\PYGZhy{}\PYGZhy{}\PYGZhy{}\PYGZhy{}\PYGZhy{}\PYGZhy{}\PYGZhy{}\PYGZhy{}\PYGZhy{}\PYGZhy{}\PYGZhy{}\PYGZhy{}\PYGZhy{}\PYGZhy{}\PYGZhy{}\PYGZhy{}\PYGZhy{}\PYGZhy{}\PYGZhy{}\PYGZhy{}\PYGZhy{}\PYGZhy{}\PYGZhy{}\PYGZhy{}\PYGZhy{}\PYGZhy{}\PYGZhy{}\PYGZhy{}\PYGZhy{}\PYGZhy{}\PYGZhy{}\PYGZhy{}\PYGZhy{}\PYGZhy{}\PYGZhy{}\PYGZhy{}\PYGZhy{}\PYGZhy{}\PYGZhy{}\PYGZhy{}\PYGZhy{}\PYGZhy{}\PYGZhy{}\PYGZhy{}\PYGZhy{}\PYGZhy{}\PYGZhy{}\PYGZhy{}\PYGZhy{}\PYGZhy{}\PYGZhy{}\PYGZhy{}\PYGZhy{}\PYGZhy{}\PYGZhy{}\PYGZhy{}\PYGZhy{}\PYGZhy{}\PYGZhy{}\PYGZhy{}\PYGZhy{}\PYGZhy{}\PYGZhy{}}
    \PYG{c+c1}{\PYGZsh{} AD graph repersentation of f(x) = sin(x) / cos(x)}
    \PYG{c+c1}{\PYGZsh{}}
    \PYG{c+c1}{\PYGZsh{} node\PYGZus{}1 : x[0]}
    \PYG{c+c1}{\PYGZsh{} node\PYGZus{}2 : sin( x[0] )}
    \PYG{c+c1}{\PYGZsh{} node\PYGZus{}3 : cos( x[0] )}
    \PYG{c+c1}{\PYGZsh{} node\PYGZus{}4 : sin( x[0] ) / cos( x[0] )}
    \PYG{c+c1}{\PYGZsh{} y[0]   = sin( x[0] ) / cos( x[0] )}
    \PYG{n}{json} \PYG{o}{=} \PYG{l+s+s1}{\PYGZsq{}\PYGZsq{}\PYGZsq{}}
\PYG{l+s+s1}{        }\PYG{l+s+s1}{\PYGZob{}}
\PYG{l+s+s1}{            }\PYG{l+s+s1}{\PYGZdq{}}\PYG{l+s+s1}{function\PYGZus{}name}\PYG{l+s+s1}{\PYGZdq{}}\PYG{l+s+s1}{ : }\PYG{l+s+s1}{\PYGZdq{}}\PYG{l+s+s1}{tangent function}\PYG{l+s+s1}{\PYGZdq{}}\PYG{l+s+s1}{,}
\PYG{l+s+s1}{            }\PYG{l+s+s1}{\PYGZdq{}}\PYG{l+s+s1}{op\PYGZus{}define\PYGZus{}vec}\PYG{l+s+s1}{\PYGZdq{}}\PYG{l+s+s1}{ : [ 3, [}
\PYG{l+s+s1}{                }\PYG{l+s+s1}{\PYGZob{}}\PYG{l+s+s1}{ }\PYG{l+s+s1}{\PYGZdq{}}\PYG{l+s+s1}{op\PYGZus{}code}\PYG{l+s+s1}{\PYGZdq{}}\PYG{l+s+s1}{:1, }\PYG{l+s+s1}{\PYGZdq{}}\PYG{l+s+s1}{name}\PYG{l+s+s1}{\PYGZdq{}}\PYG{l+s+s1}{:}\PYG{l+s+s1}{\PYGZdq{}}\PYG{l+s+s1}{sin}\PYG{l+s+s1}{\PYGZdq{}}\PYG{l+s+s1}{, }\PYG{l+s+s1}{\PYGZdq{}}\PYG{l+s+s1}{n\PYGZus{}arg}\PYG{l+s+s1}{\PYGZdq{}}\PYG{l+s+s1}{:1 \PYGZcb{} ,}
\PYG{l+s+s1}{                }\PYG{l+s+s1}{\PYGZob{}}\PYG{l+s+s1}{ }\PYG{l+s+s1}{\PYGZdq{}}\PYG{l+s+s1}{op\PYGZus{}code}\PYG{l+s+s1}{\PYGZdq{}}\PYG{l+s+s1}{:2, }\PYG{l+s+s1}{\PYGZdq{}}\PYG{l+s+s1}{name}\PYG{l+s+s1}{\PYGZdq{}}\PYG{l+s+s1}{:}\PYG{l+s+s1}{\PYGZdq{}}\PYG{l+s+s1}{cos}\PYG{l+s+s1}{\PYGZdq{}}\PYG{l+s+s1}{, }\PYG{l+s+s1}{\PYGZdq{}}\PYG{l+s+s1}{n\PYGZus{}arg}\PYG{l+s+s1}{\PYGZdq{}}\PYG{l+s+s1}{:1 \PYGZcb{} ,}
\PYG{l+s+s1}{                }\PYG{l+s+s1}{\PYGZob{}}\PYG{l+s+s1}{ }\PYG{l+s+s1}{\PYGZdq{}}\PYG{l+s+s1}{op\PYGZus{}code}\PYG{l+s+s1}{\PYGZdq{}}\PYG{l+s+s1}{:3, }\PYG{l+s+s1}{\PYGZdq{}}\PYG{l+s+s1}{name}\PYG{l+s+s1}{\PYGZdq{}}\PYG{l+s+s1}{:}\PYG{l+s+s1}{\PYGZdq{}}\PYG{l+s+s1}{div}\PYG{l+s+s1}{\PYGZdq{}}\PYG{l+s+s1}{, }\PYG{l+s+s1}{\PYGZdq{}}\PYG{l+s+s1}{n\PYGZus{}arg}\PYG{l+s+s1}{\PYGZdq{}}\PYG{l+s+s1}{:2 \PYGZcb{} ]}
\PYG{l+s+s1}{            ],}
\PYG{l+s+s1}{            }\PYG{l+s+s1}{\PYGZdq{}}\PYG{l+s+s1}{n\PYGZus{}dynamic\PYGZus{}ind}\PYG{l+s+s1}{\PYGZdq{}}\PYG{l+s+s1}{  : 0,}
\PYG{l+s+s1}{            }\PYG{l+s+s1}{\PYGZdq{}}\PYG{l+s+s1}{n\PYGZus{}variable\PYGZus{}ind}\PYG{l+s+s1}{\PYGZdq{}}\PYG{l+s+s1}{ : 1,}
\PYG{l+s+s1}{            }\PYG{l+s+s1}{\PYGZdq{}}\PYG{l+s+s1}{constant\PYGZus{}vec}\PYG{l+s+s1}{\PYGZdq{}}\PYG{l+s+s1}{   : [ 0, [ ] ],}
\PYG{l+s+s1}{            }\PYG{l+s+s1}{\PYGZdq{}}\PYG{l+s+s1}{op\PYGZus{}usage\PYGZus{}vec}\PYG{l+s+s1}{\PYGZdq{}}\PYG{l+s+s1}{   : [ 3, [}
\PYG{l+s+s1}{                [ 1, 1 ]   ,}
\PYG{l+s+s1}{                [ 2, 1 ]   ,}
\PYG{l+s+s1}{                [ 3, 2, 3] ]}
\PYG{l+s+s1}{            ],}
\PYG{l+s+s1}{            }\PYG{l+s+s1}{\PYGZdq{}}\PYG{l+s+s1}{dependent\PYGZus{}vec}\PYG{l+s+s1}{\PYGZdq{}}\PYG{l+s+s1}{ : [ 1, [4] ]}
\PYG{l+s+s1}{        \PYGZcb{};}
\PYG{l+s+s1}{    }\PYG{l+s+s1}{\PYGZsq{}\PYGZsq{}\PYGZsq{}}
    \PYG{c+c1}{\PYGZsh{} convert json to a fucntion object}
    \PYG{n}{f} \PYG{o}{=} \PYG{n}{cppad\PYGZus{}py}\PYG{o}{.}\PYG{n}{d\PYGZus{}fun}\PYG{p}{(}\PYG{p}{)}\PYG{p}{;}
    \PYG{n}{f}\PYG{o}{.}\PYG{n}{from\PYGZus{}json}\PYG{p}{(}\PYG{n}{json}\PYG{p}{)}\PYG{p}{;}
    \PYG{c+c1}{\PYGZsh{}}
    \PYG{c+c1}{\PYGZsh{} compute y = f(x)}
    \PYG{n}{x}  \PYG{o}{=} \PYG{n}{numpy}\PYG{o}{.}\PYG{n}{array}\PYG{p}{(} \PYG{p}{[} \PYG{l+m+mf}{1.0} \PYG{p}{]} \PYG{p}{)}\PYG{p}{;}
    \PYG{n}{y}  \PYG{o}{=} \PYG{n}{f}\PYG{o}{.}\PYG{n}{forward}\PYG{p}{(}\PYG{l+m+mi}{0}\PYG{p}{,} \PYG{n}{x}\PYG{p}{)}\PYG{p}{;}
    \PYG{c+c1}{\PYGZsh{}}
    \PYG{c+c1}{\PYGZsh{} check the function value}
    \PYG{n}{eps99}     \PYG{o}{=} \PYG{l+m+mf}{99.0} \PYG{o}{*} \PYG{n}{numpy}\PYG{o}{.}\PYG{n}{finfo}\PYG{p}{(}\PYG{n+nb}{float}\PYG{p}{)}\PYG{o}{.}\PYG{n}{eps}
    \PYG{n}{check}     \PYG{o}{=} \PYG{n}{math}\PYG{o}{.}\PYG{n}{tan}\PYG{p}{(}\PYG{n}{x}\PYG{p}{[}\PYG{l+m+mi}{0}\PYG{p}{]}\PYG{p}{)}\PYG{p}{;}
    \PYG{n}{rel\PYGZus{}error} \PYG{o}{=} \PYG{n}{y}\PYG{p}{[}\PYG{l+m+mi}{0}\PYG{p}{]} \PYG{o}{/} \PYG{n}{check} \PYG{o}{\PYGZhy{}} \PYG{l+m+mf}{1.0}\PYG{p}{;}
    \PYG{n}{ok}       \PYG{o}{\PYGZam{}}\PYG{o}{=} \PYG{n+nb}{abs}\PYG{p}{(} \PYG{n}{rel\PYGZus{}error} \PYG{p}{)} \PYG{o}{\PYGZlt{}} \PYG{n}{eps99}\PYG{p}{;}
    \PYG{c+c1}{\PYGZsh{}}
    \PYG{k}{return} \PYG{n}{ok}
\end{sphinxVerbatim}


\bigskip\hrule\bigskip


xsrst input file: \sphinxcode{\sphinxupquote{example/python/core/fun\_json\_xam.py}}

\index{examples@\spxentry{examples}}\ignorespaces 

\subparagraph{Examples}
\label{\detokenize{xsrst/py_fun_json:examples}}\label{\detokenize{xsrst/py_fun_json:py-fun-json-examples}}\label{\detokenize{xsrst/py_fun_json:index-6}}
{\hyperref[\detokenize{xsrst/fun_to_json_xam_py:id1}]{\sphinxcrossref{\DUrole{std,std-ref}{fun\_to\_json\_xam\_py}}}},
{\hyperref[\detokenize{xsrst/fun_from_json_xam_py:id1}]{\sphinxcrossref{\DUrole{std,std-ref}{fun\_from\_json\_xam\_py}}}}.


\bigskip\hrule\bigskip


xsrst input file: \sphinxcode{\sphinxupquote{lib/python/cppad\_py/fun\_json.xsrst}}
\begin{enumerate}
\sphinxsetlistlabels{\arabic}{enumi}{enumii}{}{.}%
\item {} 
py\_independent: {\hyperref[\detokenize{xsrst/py_independent:id1}]{\sphinxcrossref{\DUrole{std,std-ref}{3.1.1.1 Declare Independent Variables and Start Recording}}}}

\item {} 
py\_abort\_recording: {\hyperref[\detokenize{xsrst/py_abort_recording:id1}]{\sphinxcrossref{\DUrole{std,std-ref}{3.1.1.2 Abort Recording}}}}

\item {} 
py\_fun\_ctor: {\hyperref[\detokenize{xsrst/py_fun_ctor:id1}]{\sphinxcrossref{\DUrole{std,std-ref}{3.1.1.3 Stop Current Recording and Store Function Object}}}}

\item {} 
py\_fun\_property: {\hyperref[\detokenize{xsrst/py_fun_property:id1}]{\sphinxcrossref{\DUrole{std,std-ref}{3.1.1.4 Properties of a Function Object}}}}

\item {} 
py\_fun\_new\_dynamic: {\hyperref[\detokenize{xsrst/py_fun_new_dynamic:id1}]{\sphinxcrossref{\DUrole{std,std-ref}{3.1.1.5 New Dynamic Parameters}}}}

\item {} 
py\_fun\_jacobian: {\hyperref[\detokenize{xsrst/py_fun_jacobian:id1}]{\sphinxcrossref{\DUrole{std,std-ref}{3.1.1.6 Jacobian of an AD Function}}}}

\item {} 
py\_fun\_hessian: {\hyperref[\detokenize{xsrst/py_fun_hessian:id1}]{\sphinxcrossref{\DUrole{std,std-ref}{3.1.1.7 Hessian of an AD Function}}}}

\item {} 
py\_fun\_forward: {\hyperref[\detokenize{xsrst/py_fun_forward:id1}]{\sphinxcrossref{\DUrole{std,std-ref}{3.1.1.8 Forward Mode AD}}}}

\item {} 
py\_fun\_reverse: {\hyperref[\detokenize{xsrst/py_fun_reverse:id1}]{\sphinxcrossref{\DUrole{std,std-ref}{3.1.1.9 Reverse Mode AD}}}}

\item {} 
py\_fun\_optimize: {\hyperref[\detokenize{xsrst/py_fun_optimize:id1}]{\sphinxcrossref{\DUrole{std,std-ref}{3.1.1.10 Optimize an AD Function}}}}

\item {} 
py\_fun\_json: {\hyperref[\detokenize{xsrst/py_fun_json:id1}]{\sphinxcrossref{\DUrole{std,std-ref}{3.1.1.11 Json Representation of AD Computation Graph}}}}

\end{enumerate}


\bigskip\hrule\bigskip


xsrst input file: \sphinxcode{\sphinxupquote{lib/python/cppad\_py/fun.py}}


\subparagraph{py\_sparse}
\label{\detokenize{xsrst/py_sparse:py-sparse}}\label{\detokenize{xsrst/py_sparse::doc}}
\index{py\_sparse@\spxentry{py\_sparse}}\index{python@\spxentry{python}}\index{sparsity@\spxentry{sparsity}}\index{routines@\spxentry{routines}}\ignorespaces 

\subparagraph{3.1.2 Python Sparsity Routines}
\label{\detokenize{xsrst/py_sparse:python-sparsity-routines}}\label{\detokenize{xsrst/py_sparse:id1}}\begin{itemize}
\item {} 
{\hyperref[\detokenize{xsrst/py_sparse:py-sparse-children}]{\sphinxcrossref{\DUrole{std,std-ref}{Children}}}}

\end{itemize}

\index{children@\spxentry{children}}\ignorespaces 

\subparagraph{Children}
\label{\detokenize{xsrst/py_sparse:children}}\label{\detokenize{xsrst/py_sparse:py-sparse-children}}\label{\detokenize{xsrst/py_sparse:index-1}}

\subparagraph{py\_sparse\_rc}
\label{\detokenize{xsrst/py_sparse_rc:py-sparse-rc}}\label{\detokenize{xsrst/py_sparse_rc::doc}}
\index{py\_sparse\_rc@\spxentry{py\_sparse\_rc}}\index{sparsity@\spxentry{sparsity}}\index{patterns@\spxentry{patterns}}\ignorespaces 

\subparagraph{3.1.2.1 Sparsity Patterns}
\label{\detokenize{xsrst/py_sparse_rc:sparsity-patterns}}\label{\detokenize{xsrst/py_sparse_rc:id1}}\begin{itemize}
\item {} 
{\hyperref[\detokenize{xsrst/py_sparse_rc:py-sparse-rc-syntax}]{\sphinxcrossref{\DUrole{std,std-ref}{Syntax}}}}

\item {} 
{\hyperref[\detokenize{xsrst/py_sparse_rc:py-sparse-rc-pattern}]{\sphinxcrossref{\DUrole{std,std-ref}{pattern}}}}

\item {} 
{\hyperref[\detokenize{xsrst/py_sparse_rc:py-sparse-rc-nr}]{\sphinxcrossref{\DUrole{std,std-ref}{nr}}}}

\item {} 
{\hyperref[\detokenize{xsrst/py_sparse_rc:py-sparse-rc-nc}]{\sphinxcrossref{\DUrole{std,std-ref}{nc}}}}

\item {} 
{\hyperref[\detokenize{xsrst/py_sparse_rc:py-sparse-rc-nnz}]{\sphinxcrossref{\DUrole{std,std-ref}{nnz}}}}

\item {} 
{\hyperref[\detokenize{xsrst/py_sparse_rc:py-sparse-rc-resize}]{\sphinxcrossref{\DUrole{std,std-ref}{resize}}}}

\item {} \begin{description}
\item[{{\hyperref[\detokenize{xsrst/py_sparse_rc:py-sparse-rc-put}]{\sphinxcrossref{\DUrole{std,std-ref}{put}}}}}] \leavevmode\begin{itemize}
\item {} 
{\hyperref[\detokenize{xsrst/py_sparse_rc:py-sparse-rc-put-k}]{\sphinxcrossref{\DUrole{std,std-ref}{k}}}}

\item {} 
{\hyperref[\detokenize{xsrst/py_sparse_rc:py-sparse-rc-put-r}]{\sphinxcrossref{\DUrole{std,std-ref}{r}}}}

\item {} 
{\hyperref[\detokenize{xsrst/py_sparse_rc:py-sparse-rc-put-c}]{\sphinxcrossref{\DUrole{std,std-ref}{c}}}}

\end{itemize}

\end{description}

\item {} 
{\hyperref[\detokenize{xsrst/py_sparse_rc:py-sparse-rc-row}]{\sphinxcrossref{\DUrole{std,std-ref}{row}}}}

\item {} 
{\hyperref[\detokenize{xsrst/py_sparse_rc:py-sparse-rc-col}]{\sphinxcrossref{\DUrole{std,std-ref}{col}}}}

\item {} 
{\hyperref[\detokenize{xsrst/py_sparse_rc:py-sparse-rc-row-major}]{\sphinxcrossref{\DUrole{std,std-ref}{row\_major}}}}

\item {} 
{\hyperref[\detokenize{xsrst/py_sparse_rc:py-sparse-rc-col-major}]{\sphinxcrossref{\DUrole{std,std-ref}{col\_major}}}}

\item {} 
{\hyperref[\detokenize{xsrst/py_sparse_rc:py-sparse-rc-example}]{\sphinxcrossref{\DUrole{std,std-ref}{Example}}}}

\end{itemize}

\index{syntax@\spxentry{syntax}}\ignorespaces 

\subparagraph{Syntax}
\label{\detokenize{xsrst/py_sparse_rc:syntax}}\label{\detokenize{xsrst/py_sparse_rc:py-sparse-rc-syntax}}\label{\detokenize{xsrst/py_sparse_rc:index-1}}
\begin{DUlineblock}{0em}
\item[] \sphinxstyleemphasis{pattern} =  \sphinxcode{\sphinxupquote{cppad\_py.sparse\_rc}} ()
\item[] \sphinxstyleemphasis{pattern} . \sphinxcode{\sphinxupquote{resize}} ( \sphinxstyleemphasis{nr} , \sphinxstyleemphasis{nc} , \sphinxstyleemphasis{nnz} )
\item[] \sphinxstyleemphasis{nr} = \sphinxstyleemphasis{pattern} . \sphinxcode{\sphinxupquote{nr}} ()
\item[] \sphinxstyleemphasis{nc} = \sphinxstyleemphasis{pattern} . \sphinxcode{\sphinxupquote{nc}} ()
\item[] \sphinxstyleemphasis{nnz} = \sphinxstyleemphasis{pattern} . \sphinxcode{\sphinxupquote{nnz}} ()
\item[] \sphinxstyleemphasis{pattern} . \sphinxcode{\sphinxupquote{put}} ( \sphinxstyleemphasis{k} , \sphinxstyleemphasis{r} , \sphinxstyleemphasis{c} )
\item[] \sphinxstyleemphasis{row} = \sphinxstyleemphasis{pattern} . \sphinxcode{\sphinxupquote{row}} ()
\item[] \sphinxstyleemphasis{col} = \sphinxstyleemphasis{pattern} . \sphinxcode{\sphinxupquote{col}} ()
\item[] \sphinxstyleemphasis{row\_major} = \sphinxstyleemphasis{pattern} . \sphinxcode{\sphinxupquote{row\_major}} ()
\item[] \sphinxstyleemphasis{col\_major} = \sphinxstyleemphasis{pattern} . \sphinxcode{\sphinxupquote{col\_major}} ()
\end{DUlineblock}

\index{pattern@\spxentry{pattern}}\ignorespaces 

\subparagraph{pattern}
\label{\detokenize{xsrst/py_sparse_rc:pattern}}\label{\detokenize{xsrst/py_sparse_rc:py-sparse-rc-pattern}}\label{\detokenize{xsrst/py_sparse_rc:index-2}}
This result is used to hold a sparsity pattern for a matrix.
It is constant
except during the \sphinxcode{\sphinxupquote{resize}} and \sphinxcode{\sphinxupquote{put}} operations.

\index{nr@\spxentry{nr}}\ignorespaces 

\subparagraph{nr}
\label{\detokenize{xsrst/py_sparse_rc:nr}}\label{\detokenize{xsrst/py_sparse_rc:py-sparse-rc-nr}}\label{\detokenize{xsrst/py_sparse_rc:index-3}}
This argument is a non\sphinxhyphen{}negative \sphinxcode{\sphinxupquote{int}}
and is the number of rows in the sparsity pattern.
The function \sphinxcode{\sphinxupquote{nr()}} returns the value of
\sphinxstyleemphasis{nr} in the previous \sphinxcode{\sphinxupquote{resize}} operation.

\index{nc@\spxentry{nc}}\ignorespaces 

\subparagraph{nc}
\label{\detokenize{xsrst/py_sparse_rc:nc}}\label{\detokenize{xsrst/py_sparse_rc:py-sparse-rc-nc}}\label{\detokenize{xsrst/py_sparse_rc:index-4}}
This argument is a non\sphinxhyphen{}negative \sphinxcode{\sphinxupquote{int}}
and is the number of columns in the sparsity pattern.
The function \sphinxcode{\sphinxupquote{nc()}} returns the value of
\sphinxstyleemphasis{nc} in the previous \sphinxcode{\sphinxupquote{resize}} operation.

\index{nnz@\spxentry{nnz}}\ignorespaces 

\subparagraph{nnz}
\label{\detokenize{xsrst/py_sparse_rc:nnz}}\label{\detokenize{xsrst/py_sparse_rc:py-sparse-rc-nnz}}\label{\detokenize{xsrst/py_sparse_rc:index-5}}
This argument is a non\sphinxhyphen{}negative \sphinxcode{\sphinxupquote{int}}
and is the number of possibly non\sphinxhyphen{}zero
index pairs in the sparsity pattern.
The function \sphinxcode{\sphinxupquote{nnz()}} returns the value of
\sphinxstyleemphasis{nnz} in the previous \sphinxcode{\sphinxupquote{resize}} operation.

\index{resize@\spxentry{resize}}\ignorespaces 

\subparagraph{resize}
\label{\detokenize{xsrst/py_sparse_rc:resize}}\label{\detokenize{xsrst/py_sparse_rc:py-sparse-rc-resize}}\label{\detokenize{xsrst/py_sparse_rc:index-6}}
The current sparsity pattern is lost and a new one is started
with the specified parameters.
After each \sphinxcode{\sphinxupquote{resize}} , the elements in the \sphinxstyleemphasis{row}
and \sphinxstyleemphasis{col} vectors should be assigned using \sphinxcode{\sphinxupquote{put}} .

\index{put@\spxentry{put}}\ignorespaces 

\subparagraph{put}
\label{\detokenize{xsrst/py_sparse_rc:put}}\label{\detokenize{xsrst/py_sparse_rc:py-sparse-rc-put}}\label{\detokenize{xsrst/py_sparse_rc:index-7}}
This function sets the values

\begin{DUlineblock}{0em}
\item[]         \sphinxstyleemphasis{row} {[} \sphinxstyleemphasis{k} {]} = \sphinxstyleemphasis{r}
\item[]         \sphinxstyleemphasis{col} {[} \sphinxstyleemphasis{k} {]} = \sphinxstyleemphasis{c}
\end{DUlineblock}

(The name \sphinxcode{\sphinxupquote{set}} is used by Cppad, but not used here,
because \sphinxcode{\sphinxupquote{set}} it is a built\sphinxhyphen{}in name in Python.)

\index{k@\spxentry{k}}\ignorespaces 

\subparagraph{k}
\label{\detokenize{xsrst/py_sparse_rc:k}}\label{\detokenize{xsrst/py_sparse_rc:py-sparse-rc-put-k}}\label{\detokenize{xsrst/py_sparse_rc:index-8}}
This argument is a non\sphinxhyphen{}negative \sphinxcode{\sphinxupquote{int}}
and must be less than \sphinxstyleemphasis{nnz} .

\index{r@\spxentry{r}}\ignorespaces 

\subparagraph{r}
\label{\detokenize{xsrst/py_sparse_rc:r}}\label{\detokenize{xsrst/py_sparse_rc:py-sparse-rc-put-r}}\label{\detokenize{xsrst/py_sparse_rc:index-9}}
This argument is a non\sphinxhyphen{}negative \sphinxcode{\sphinxupquote{int}}
and must be less than \sphinxstyleemphasis{nr} .
It specifies the value assigned to \sphinxstyleemphasis{row} {[} \sphinxstyleemphasis{k} {]} .

\index{c@\spxentry{c}}\ignorespaces 

\subparagraph{c}
\label{\detokenize{xsrst/py_sparse_rc:c}}\label{\detokenize{xsrst/py_sparse_rc:py-sparse-rc-put-c}}\label{\detokenize{xsrst/py_sparse_rc:index-10}}
This argument is a non\sphinxhyphen{}negative \sphinxcode{\sphinxupquote{int}}
and must be less than \sphinxstyleemphasis{nc} .
It specifies the value assigned to \sphinxstyleemphasis{col} {[} \sphinxstyleemphasis{k} {]} .

\index{row@\spxentry{row}}\ignorespaces 

\subparagraph{row}
\label{\detokenize{xsrst/py_sparse_rc:row}}\label{\detokenize{xsrst/py_sparse_rc:py-sparse-rc-row}}\label{\detokenize{xsrst/py_sparse_rc:index-11}}
This result is a numpy vector with \sphinxcode{\sphinxupquote{int}} elements
and its size is \sphinxstyleemphasis{nnz} .
For \sphinxstyleemphasis{k} = 0, … , \sphinxstyleemphasis{nnz} \sphinxhyphen{}1 ,
\sphinxstyleemphasis{row} {[} \sphinxstyleemphasis{k} {]} is the row index for the \sphinxstyleemphasis{k}\sphinxhyphen{}th possibly non\sphinxhyphen{}zero
entry in the matrix.

\index{col@\spxentry{col}}\ignorespaces 

\subparagraph{col}
\label{\detokenize{xsrst/py_sparse_rc:col}}\label{\detokenize{xsrst/py_sparse_rc:py-sparse-rc-col}}\label{\detokenize{xsrst/py_sparse_rc:index-12}}
This result is a numpy vector with \sphinxcode{\sphinxupquote{int}} elements
and its size is \sphinxstyleemphasis{nnz} .
For \sphinxstyleemphasis{k} = 0, … , \sphinxstyleemphasis{nnz} \sphinxhyphen{}1 ,
\sphinxstyleemphasis{col} {[} \sphinxstyleemphasis{k} {]} is the column index for the \sphinxstyleemphasis{k}\sphinxhyphen{}th possibly non\sphinxhyphen{}zero
entry in the matrix.

\index{row\_major@\spxentry{row\_major}}\ignorespaces 

\subparagraph{row\_major}
\label{\detokenize{xsrst/py_sparse_rc:row-major}}\label{\detokenize{xsrst/py_sparse_rc:py-sparse-rc-row-major}}\label{\detokenize{xsrst/py_sparse_rc:index-13}}
This result is a numpy vector with \sphinxcode{\sphinxupquote{int}} elements
and its size \sphinxstyleemphasis{nnz} .
It sorts the sparsity pattern in row\sphinxhyphen{}major order.
To be specific,

\begin{DUlineblock}{0em}
\item[]         \sphinxstyleemphasis{col} {[} \sphinxstyleemphasis{row\_major} {[} \sphinxstyleemphasis{k} {]} {]} \textless{}= \sphinxstyleemphasis{col} {[} \sphinxstyleemphasis{row\_major} {[} \sphinxstyleemphasis{k} +1{]} {]}
\end{DUlineblock}

and if \sphinxstyleemphasis{col} {[} \sphinxstyleemphasis{row\_major} {[} \sphinxstyleemphasis{k} {]} {]} == \sphinxstyleemphasis{col} {[} \sphinxstyleemphasis{row\_major} {[} \sphinxstyleemphasis{k} +1{]} {]} ,

\begin{DUlineblock}{0em}
\item[]         \sphinxstyleemphasis{row} {[} \sphinxstyleemphasis{row\_major} {[} \sphinxstyleemphasis{k} {]} {]} \textless{} \sphinxstyleemphasis{row} {[} \sphinxstyleemphasis{row\_major} {[} \sphinxstyleemphasis{k} +1{]} {]}
\end{DUlineblock}

This routine generates an assert if there are two entries with the same
row and column values (if \sphinxcode{\sphinxupquote{NDEBUG}} is not defined).

\index{col\_major@\spxentry{col\_major}}\ignorespaces 

\subparagraph{col\_major}
\label{\detokenize{xsrst/py_sparse_rc:col-major}}\label{\detokenize{xsrst/py_sparse_rc:py-sparse-rc-col-major}}\label{\detokenize{xsrst/py_sparse_rc:index-14}}
This result is a numpy vector with \sphinxcode{\sphinxupquote{int}} elements
and its size \sphinxstyleemphasis{nnz} .
It sorts the sparsity pattern in column\sphinxhyphen{}major order.
To be specific,

\begin{DUlineblock}{0em}
\item[]         \sphinxstyleemphasis{row} {[} \sphinxstyleemphasis{col\_major} {[} \sphinxstyleemphasis{k} {]} {]} \textless{}= \sphinxstyleemphasis{row} {[} \sphinxstyleemphasis{col\_major} {[} \sphinxstyleemphasis{k} +1{]} {]}
\end{DUlineblock}

and if \sphinxstyleemphasis{row} {[} \sphinxstyleemphasis{col\_major} {[} \sphinxstyleemphasis{k} {]} {]} == \sphinxstyleemphasis{row} {[} \sphinxstyleemphasis{col\_major} {[} \sphinxstyleemphasis{k} +1{]} {]} ,

\begin{DUlineblock}{0em}
\item[]         \sphinxstyleemphasis{col} {[} \sphinxstyleemphasis{col\_major} {[} \sphinxstyleemphasis{k} {]} {]} \textless{} \sphinxstyleemphasis{col} {[} \sphinxstyleemphasis{col\_major} {[} \sphinxstyleemphasis{k} +1{]} {]}
\end{DUlineblock}

This routine generates an assert if there are two entries with the same
row and column values (if \sphinxcode{\sphinxupquote{NDEBUG}} is not defined).


\subparagraph{sparse\_rc\_xam\_py}
\label{\detokenize{xsrst/sparse_rc_xam_py:sparse-rc-xam-py}}\label{\detokenize{xsrst/sparse_rc_xam_py::doc}}
\index{sparse\_rc\_xam\_py@\spxentry{sparse\_rc\_xam\_py}}\index{python:@\spxentry{python:}}\index{sparsity@\spxentry{sparsity}}\index{patterns:@\spxentry{patterns:}}\index{example@\spxentry{example}}\index{test@\spxentry{test}}\ignorespaces 

\subparagraph{3.1.2.1.1 Python: Sparsity Patterns: Example and Test}
\label{\detokenize{xsrst/sparse_rc_xam_py:python-sparsity-patterns-example-and-test}}\label{\detokenize{xsrst/sparse_rc_xam_py:id1}}
\begin{sphinxVerbatim}[commandchars=\\\{\}]
\PYG{k}{def} \PYG{n+nf}{sparse\PYGZus{}rc\PYGZus{}xam}\PYG{p}{(}\PYG{p}{)} \PYG{p}{:}
    \PYG{c+c1}{\PYGZsh{}}
    \PYG{k+kn}{import} \PYG{n+nn}{numpy}
    \PYG{k+kn}{import} \PYG{n+nn}{cppad\PYGZus{}py}
    \PYG{c+c1}{\PYGZsh{}}
    \PYG{c+c1}{\PYGZsh{} initialize return variable}
    \PYG{n}{ok} \PYG{o}{=} \PYG{k+kc}{True}
    \PYG{c+c1}{\PYGZsh{} \PYGZhy{}\PYGZhy{}\PYGZhy{}\PYGZhy{}\PYGZhy{}\PYGZhy{}\PYGZhy{}\PYGZhy{}\PYGZhy{}\PYGZhy{}\PYGZhy{}\PYGZhy{}\PYGZhy{}\PYGZhy{}\PYGZhy{}\PYGZhy{}\PYGZhy{}\PYGZhy{}\PYGZhy{}\PYGZhy{}\PYGZhy{}\PYGZhy{}\PYGZhy{}\PYGZhy{}\PYGZhy{}\PYGZhy{}\PYGZhy{}\PYGZhy{}\PYGZhy{}\PYGZhy{}\PYGZhy{}\PYGZhy{}\PYGZhy{}\PYGZhy{}\PYGZhy{}\PYGZhy{}\PYGZhy{}\PYGZhy{}\PYGZhy{}\PYGZhy{}\PYGZhy{}\PYGZhy{}\PYGZhy{}\PYGZhy{}\PYGZhy{}\PYGZhy{}\PYGZhy{}\PYGZhy{}\PYGZhy{}\PYGZhy{}\PYGZhy{}\PYGZhy{}\PYGZhy{}\PYGZhy{}\PYGZhy{}\PYGZhy{}\PYGZhy{}\PYGZhy{}\PYGZhy{}\PYGZhy{}\PYGZhy{}\PYGZhy{}\PYGZhy{}\PYGZhy{}\PYGZhy{}\PYGZhy{}\PYGZhy{}\PYGZhy{}\PYGZhy{}}
    \PYG{c+c1}{\PYGZsh{}}
    \PYG{c+c1}{\PYGZsh{} create a sparsity pattern}
    \PYG{n}{pattern} \PYG{o}{=} \PYG{n}{cppad\PYGZus{}py}\PYG{o}{.}\PYG{n}{sparse\PYGZus{}rc}\PYG{p}{(}\PYG{p}{)}
    \PYG{c+c1}{\PYGZsh{}}
    \PYG{n}{nr} \PYG{o}{=} \PYG{l+m+mi}{6}
    \PYG{n}{nc} \PYG{o}{=} \PYG{l+m+mi}{5}
    \PYG{n}{nnz} \PYG{o}{=} \PYG{l+m+mi}{4}
    \PYG{c+c1}{\PYGZsh{}}
    \PYG{c+c1}{\PYGZsh{} resize}
    \PYG{n}{pattern}\PYG{o}{.}\PYG{n}{resize}\PYG{p}{(}\PYG{n}{nr}\PYG{p}{,} \PYG{n}{nc}\PYG{p}{,} \PYG{n}{nnz}\PYG{p}{)}
    \PYG{c+c1}{\PYGZsh{}}
    \PYG{n}{ok} \PYG{o}{=} \PYG{n}{ok} \PYG{o+ow}{and} \PYG{n}{pattern}\PYG{o}{.}\PYG{n}{nr}\PYG{p}{(}\PYG{p}{)}  \PYG{o}{==} \PYG{n}{nr}
    \PYG{n}{ok} \PYG{o}{=} \PYG{n}{ok} \PYG{o+ow}{and} \PYG{n}{pattern}\PYG{o}{.}\PYG{n}{nc}\PYG{p}{(}\PYG{p}{)}  \PYG{o}{==} \PYG{n}{nc}
    \PYG{n}{ok} \PYG{o}{=} \PYG{n}{ok} \PYG{o+ow}{and} \PYG{n}{pattern}\PYG{o}{.}\PYG{n}{nnz}\PYG{p}{(}\PYG{p}{)} \PYG{o}{==} \PYG{n}{nnz}
    \PYG{c+c1}{\PYGZsh{}}
    \PYG{c+c1}{\PYGZsh{} indices corresponding to upper\PYGZhy{}diagonal}
    \PYG{k}{for} \PYG{n}{k} \PYG{o+ow}{in} \PYG{n+nb}{range}\PYG{p}{(} \PYG{n}{nnz} \PYG{p}{)} \PYG{p}{:}
        \PYG{n}{pattern}\PYG{o}{.}\PYG{n}{put}\PYG{p}{(}\PYG{n}{k}\PYG{p}{,} \PYG{n}{k}\PYG{p}{,} \PYG{n}{k}\PYG{o}{+}\PYG{l+m+mi}{1}\PYG{p}{)}
    \PYG{c+c1}{\PYGZsh{}}
    \PYG{c+c1}{\PYGZsh{}}
    \PYG{c+c1}{\PYGZsh{} row and column indices}
    \PYG{n}{row} \PYG{o}{=} \PYG{n}{pattern}\PYG{o}{.}\PYG{n}{row}\PYG{p}{(}\PYG{p}{)}
    \PYG{n}{col} \PYG{o}{=} \PYG{n}{pattern}\PYG{o}{.}\PYG{n}{col}\PYG{p}{(}\PYG{p}{)}
    \PYG{c+c1}{\PYGZsh{}}
    \PYG{c+c1}{\PYGZsh{} check row and column indices}
    \PYG{k}{for} \PYG{n}{k} \PYG{o+ow}{in} \PYG{n+nb}{range}\PYG{p}{(} \PYG{n}{nnz} \PYG{p}{)} \PYG{p}{:}
        \PYG{n}{ok} \PYG{o}{=} \PYG{n}{ok} \PYG{o+ow}{and} \PYG{n}{row}\PYG{p}{[}\PYG{n}{k}\PYG{p}{]} \PYG{o}{==} \PYG{n}{k}
        \PYG{n}{ok} \PYG{o}{=} \PYG{n}{ok} \PYG{o+ow}{and} \PYG{n}{col}\PYG{p}{[}\PYG{n}{k}\PYG{p}{]} \PYG{o}{==} \PYG{n}{k}\PYG{o}{+}\PYG{l+m+mi}{1}
    \PYG{c+c1}{\PYGZsh{}}
    \PYG{c+c1}{\PYGZsh{}}
    \PYG{c+c1}{\PYGZsh{} For this case, row\PYGZus{}major and col\PYGZus{}major order are same as original order}
    \PYG{n}{row\PYGZus{}major} \PYG{o}{=} \PYG{n}{pattern}\PYG{o}{.}\PYG{n}{row\PYGZus{}major}\PYG{p}{(}\PYG{p}{)}
    \PYG{n}{col\PYGZus{}major} \PYG{o}{=} \PYG{n}{pattern}\PYG{o}{.}\PYG{n}{col\PYGZus{}major}\PYG{p}{(}\PYG{p}{)}
    \PYG{k}{for} \PYG{n}{k} \PYG{o+ow}{in} \PYG{n+nb}{range}\PYG{p}{(} \PYG{n}{nnz} \PYG{p}{)} \PYG{p}{:}
        \PYG{n}{ok} \PYG{o}{=} \PYG{n}{ok} \PYG{o+ow}{and} \PYG{n}{row\PYGZus{}major}\PYG{p}{[}\PYG{n}{k}\PYG{p}{]} \PYG{o}{==} \PYG{n}{k}
        \PYG{n}{ok} \PYG{o}{=} \PYG{n}{ok} \PYG{o+ow}{and} \PYG{n}{col\PYGZus{}major}\PYG{p}{[}\PYG{n}{k}\PYG{p}{]} \PYG{o}{==} \PYG{n}{k}
    \PYG{c+c1}{\PYGZsh{}}
    \PYG{c+c1}{\PYGZsh{}}
    \PYG{k}{return}\PYG{p}{(} \PYG{n}{ok} \PYG{p}{)}
\PYG{c+c1}{\PYGZsh{}}
\end{sphinxVerbatim}


\bigskip\hrule\bigskip


xsrst input file: \sphinxcode{\sphinxupquote{example/python/core/sparse\_rc\_xam.py}}

\index{example@\spxentry{example}}\ignorespaces 

\subparagraph{Example}
\label{\detokenize{xsrst/py_sparse_rc:example}}\label{\detokenize{xsrst/py_sparse_rc:py-sparse-rc-example}}\label{\detokenize{xsrst/py_sparse_rc:index-15}}
{\hyperref[\detokenize{xsrst/sparse_rc_xam_py:id1}]{\sphinxcrossref{\DUrole{std,std-ref}{sparse\_rc\_xam\_py}}}}


\bigskip\hrule\bigskip


xsrst input file: \sphinxcode{\sphinxupquote{lib/python/cppad\_py/sparse\_rc.py}}


\subparagraph{py\_sparse\_rcv}
\label{\detokenize{xsrst/py_sparse_rcv:py-sparse-rcv}}\label{\detokenize{xsrst/py_sparse_rcv::doc}}
\index{py\_sparse\_rcv@\spxentry{py\_sparse\_rcv}}\index{sparse@\spxentry{sparse}}\index{matrices@\spxentry{matrices}}\ignorespaces 

\subparagraph{3.1.2.2 Sparse Matrices}
\label{\detokenize{xsrst/py_sparse_rcv:sparse-matrices}}\label{\detokenize{xsrst/py_sparse_rcv:id1}}\begin{itemize}
\item {} 
{\hyperref[\detokenize{xsrst/py_sparse_rcv:py-sparse-rcv-syntax}]{\sphinxcrossref{\DUrole{std,std-ref}{Syntax}}}}

\item {} 
{\hyperref[\detokenize{xsrst/py_sparse_rcv:py-sparse-rcv-pattern}]{\sphinxcrossref{\DUrole{std,std-ref}{pattern}}}}

\item {} 
{\hyperref[\detokenize{xsrst/py_sparse_rcv:py-sparse-rcv-matrix}]{\sphinxcrossref{\DUrole{std,std-ref}{matrix}}}}

\item {} 
{\hyperref[\detokenize{xsrst/py_sparse_rcv:py-sparse-rcv-nr}]{\sphinxcrossref{\DUrole{std,std-ref}{nr}}}}

\item {} 
{\hyperref[\detokenize{xsrst/py_sparse_rcv:py-sparse-rcv-nc}]{\sphinxcrossref{\DUrole{std,std-ref}{nc}}}}

\item {} 
{\hyperref[\detokenize{xsrst/py_sparse_rcv:py-sparse-rcv-nnz}]{\sphinxcrossref{\DUrole{std,std-ref}{nnz}}}}

\item {} \begin{description}
\item[{{\hyperref[\detokenize{xsrst/py_sparse_rcv:py-sparse-rcv-put}]{\sphinxcrossref{\DUrole{std,std-ref}{put}}}}}] \leavevmode\begin{itemize}
\item {} 
{\hyperref[\detokenize{xsrst/py_sparse_rcv:py-sparse-rcv-put-k}]{\sphinxcrossref{\DUrole{std,std-ref}{k}}}}

\item {} 
{\hyperref[\detokenize{xsrst/py_sparse_rcv:py-sparse-rcv-put-v}]{\sphinxcrossref{\DUrole{std,std-ref}{v}}}}

\end{itemize}

\end{description}

\item {} 
{\hyperref[\detokenize{xsrst/py_sparse_rcv:py-sparse-rcv-row}]{\sphinxcrossref{\DUrole{std,std-ref}{row}}}}

\item {} 
{\hyperref[\detokenize{xsrst/py_sparse_rcv:py-sparse-rcv-col}]{\sphinxcrossref{\DUrole{std,std-ref}{col}}}}

\item {} 
{\hyperref[\detokenize{xsrst/py_sparse_rcv:py-sparse-rcv-val}]{\sphinxcrossref{\DUrole{std,std-ref}{val}}}}

\item {} 
{\hyperref[\detokenize{xsrst/py_sparse_rcv:py-sparse-rcv-row-major}]{\sphinxcrossref{\DUrole{std,std-ref}{row\_major}}}}

\item {} 
{\hyperref[\detokenize{xsrst/py_sparse_rcv:py-sparse-rcv-col-major}]{\sphinxcrossref{\DUrole{std,std-ref}{col\_major}}}}

\item {} 
{\hyperref[\detokenize{xsrst/py_sparse_rcv:py-sparse-rcv-example}]{\sphinxcrossref{\DUrole{std,std-ref}{Example}}}}

\end{itemize}

\index{syntax@\spxentry{syntax}}\ignorespaces 

\subparagraph{Syntax}
\label{\detokenize{xsrst/py_sparse_rcv:syntax}}\label{\detokenize{xsrst/py_sparse_rcv:py-sparse-rcv-syntax}}\label{\detokenize{xsrst/py_sparse_rcv:index-1}}
\begin{DUlineblock}{0em}
\item[] \sphinxstyleemphasis{matrix} =  \sphinxcode{\sphinxupquote{cppad\_py::sparse\_rcv}} ( \sphinxstyleemphasis{pattern} )
\item[] \sphinxstyleemphasis{nr} = \sphinxstyleemphasis{matrix} . \sphinxcode{\sphinxupquote{nr}} ()
\item[] \sphinxstyleemphasis{nc} = \sphinxstyleemphasis{matrix} . \sphinxcode{\sphinxupquote{nc}} ()
\item[] \sphinxstyleemphasis{nnz} =  \sphinxcode{\sphinxupquote{matrix}} . \sphinxstyleemphasis{nnz} ()
\item[] \sphinxstyleemphasis{matrix} . \sphinxcode{\sphinxupquote{put}} ( \sphinxstyleemphasis{k} , \sphinxstyleemphasis{v} )
\item[] \sphinxstyleemphasis{row} = \sphinxstyleemphasis{matrix} . \sphinxcode{\sphinxupquote{row}} ()
\item[] \sphinxstyleemphasis{col} = \sphinxstyleemphasis{matrix} . \sphinxcode{\sphinxupquote{col}} ()
\item[] \sphinxstyleemphasis{val} = \sphinxstyleemphasis{matrix} . \sphinxcode{\sphinxupquote{val}} ()
\item[] \sphinxstyleemphasis{row\_major} = \sphinxstyleemphasis{matrix} . \sphinxcode{\sphinxupquote{row\_major}} ()
\item[] \sphinxstyleemphasis{col\_major} = \sphinxstyleemphasis{matrix} . \sphinxcode{\sphinxupquote{col\_major}} ()
\end{DUlineblock}

\index{pattern@\spxentry{pattern}}\ignorespaces 

\subparagraph{pattern}
\label{\detokenize{xsrst/py_sparse_rcv:pattern}}\label{\detokenize{xsrst/py_sparse_rcv:py-sparse-rcv-pattern}}\label{\detokenize{xsrst/py_sparse_rcv:index-2}}
This argument has prototype

\begin{DUlineblock}{0em}
\item[]         \sphinxcode{\sphinxupquote{const sparse\_rc\&}} \sphinxstyleemphasis{pattern}
\end{DUlineblock}

It specifies the number of rows, number of columns and
the possibly non\sphinxhyphen{}zero entries in the \sphinxstyleemphasis{matrix} .

\index{matrix@\spxentry{matrix}}\ignorespaces 

\subparagraph{matrix}
\label{\detokenize{xsrst/py_sparse_rcv:matrix}}\label{\detokenize{xsrst/py_sparse_rcv:py-sparse-rcv-matrix}}\label{\detokenize{xsrst/py_sparse_rcv:index-3}}
This is a sparse matrix object with the sparsity specified by \sphinxstyleemphasis{pattern} .
Only the \sphinxstyleemphasis{val} vector can be changed. All other values returned by
\sphinxstyleemphasis{matrix} are fixed during the constructor and constant there after.
The \sphinxstyleemphasis{val} vector is only changed by the constructor
and the \sphinxcode{\sphinxupquote{set}} function.

\index{nr@\spxentry{nr}}\ignorespaces 

\subparagraph{nr}
\label{\detokenize{xsrst/py_sparse_rcv:nr}}\label{\detokenize{xsrst/py_sparse_rcv:py-sparse-rcv-nr}}\label{\detokenize{xsrst/py_sparse_rcv:index-4}}
This return value is an \sphinxcode{\sphinxupquote{int}}
and is the number of rows in the matrix.

\index{nc@\spxentry{nc}}\ignorespaces 

\subparagraph{nc}
\label{\detokenize{xsrst/py_sparse_rcv:nc}}\label{\detokenize{xsrst/py_sparse_rcv:py-sparse-rcv-nc}}\label{\detokenize{xsrst/py_sparse_rcv:index-5}}
This return value is an \sphinxcode{\sphinxupquote{int}}
and is the number of columns in the matrix.

\index{nnz@\spxentry{nnz}}\ignorespaces 

\subparagraph{nnz}
\label{\detokenize{xsrst/py_sparse_rcv:nnz}}\label{\detokenize{xsrst/py_sparse_rcv:py-sparse-rcv-nnz}}\label{\detokenize{xsrst/py_sparse_rcv:index-6}}
This return value is an \sphinxcode{\sphinxupquote{int}}
and is the number of possibly non\sphinxhyphen{}zero values in the matrix.

\index{put@\spxentry{put}}\ignorespaces 

\subparagraph{put}
\label{\detokenize{xsrst/py_sparse_rcv:put}}\label{\detokenize{xsrst/py_sparse_rcv:py-sparse-rcv-put}}\label{\detokenize{xsrst/py_sparse_rcv:index-7}}
This function sets the value

\begin{DUlineblock}{0em}
\item[]         \sphinxstyleemphasis{val} {[} \sphinxstyleemphasis{k} {]} = \sphinxstyleemphasis{v}
\end{DUlineblock}

(The name \sphinxcode{\sphinxupquote{set}} is used by Cppad, but not used here,
because \sphinxcode{\sphinxupquote{set}} it is a built\sphinxhyphen{}in name in Python.)

\index{k@\spxentry{k}}\ignorespaces 

\subparagraph{k}
\label{\detokenize{xsrst/py_sparse_rcv:k}}\label{\detokenize{xsrst/py_sparse_rcv:py-sparse-rcv-put-k}}\label{\detokenize{xsrst/py_sparse_rcv:index-8}}
This is a non\sphinxhyphen{}negative \sphinxcode{\sphinxupquote{int}} and must be less than \sphinxstyleemphasis{nnz} .

\index{v@\spxentry{v}}\ignorespaces 

\subparagraph{v}
\label{\detokenize{xsrst/py_sparse_rcv:v}}\label{\detokenize{xsrst/py_sparse_rcv:py-sparse-rcv-put-v}}\label{\detokenize{xsrst/py_sparse_rcv:index-9}}
This argument has type \sphinxcode{\sphinxupquote{float}} and
specifies the value assigned to \sphinxstyleemphasis{val} {[} \sphinxstyleemphasis{k} {]} .

\index{row@\spxentry{row}}\ignorespaces 

\subparagraph{row}
\label{\detokenize{xsrst/py_sparse_rcv:row}}\label{\detokenize{xsrst/py_sparse_rcv:py-sparse-rcv-row}}\label{\detokenize{xsrst/py_sparse_rcv:index-10}}
This result is a numpy vector with \sphinxcode{\sphinxupquote{int}} elements
and its size is \sphinxstyleemphasis{nnz} .
For \sphinxstyleemphasis{k} = 0, … , \sphinxstyleemphasis{nnz} \sphinxhyphen{}1 ,
\sphinxstyleemphasis{row} {[} \sphinxstyleemphasis{k} {]} is the row index for the \sphinxstyleemphasis{k}\sphinxhyphen{}th possibly non\sphinxhyphen{}zero
entry in the matrix.

\index{col@\spxentry{col}}\ignorespaces 

\subparagraph{col}
\label{\detokenize{xsrst/py_sparse_rcv:col}}\label{\detokenize{xsrst/py_sparse_rcv:py-sparse-rcv-col}}\label{\detokenize{xsrst/py_sparse_rcv:index-11}}
This result is a numpy vector with \sphinxcode{\sphinxupquote{int}} elements
and its size is \sphinxstyleemphasis{nnz} .
For \sphinxstyleemphasis{k} = 0, … , \sphinxstyleemphasis{nnz} \sphinxhyphen{}1 ,
\sphinxstyleemphasis{col} {[} \sphinxstyleemphasis{k} {]} is the column index for the \sphinxstyleemphasis{k}\sphinxhyphen{}th possibly non\sphinxhyphen{}zero
entry in the matrix.

\index{val@\spxentry{val}}\ignorespaces 

\subparagraph{val}
\label{\detokenize{xsrst/py_sparse_rcv:val}}\label{\detokenize{xsrst/py_sparse_rcv:py-sparse-rcv-val}}\label{\detokenize{xsrst/py_sparse_rcv:index-12}}
This result is a numpy vector with \sphinxcode{\sphinxupquote{float}} elements
and its size is \sphinxstyleemphasis{nnz} .
For \sphinxstyleemphasis{k} = 0, … , \sphinxstyleemphasis{nnz} \sphinxhyphen{}1 ,
\sphinxstyleemphasis{val} {[} \sphinxstyleemphasis{k} {]} is the value of the \sphinxstyleemphasis{k}\sphinxhyphen{}th possibly non\sphinxhyphen{}zero
entry in the matrix (the value may be zero).

\index{row\_major@\spxentry{row\_major}}\ignorespaces 

\subparagraph{row\_major}
\label{\detokenize{xsrst/py_sparse_rcv:row-major}}\label{\detokenize{xsrst/py_sparse_rcv:py-sparse-rcv-row-major}}\label{\detokenize{xsrst/py_sparse_rcv:index-13}}
This result is a numpy vector with \sphinxcode{\sphinxupquote{int}} elements
and its size \sphinxstyleemphasis{nnz} .
It sorts the sparsity pattern in row\sphinxhyphen{}major order.
To be specific,

\begin{DUlineblock}{0em}
\item[]         \sphinxstyleemphasis{col} {[} \sphinxstyleemphasis{row\_major} {[} \sphinxstyleemphasis{k} {]} {]} \textless{}= \sphinxstyleemphasis{col} {[} \sphinxstyleemphasis{row\_major} {[} \sphinxstyleemphasis{k} +1{]} {]}
\end{DUlineblock}

and if \sphinxstyleemphasis{col} {[} \sphinxstyleemphasis{row\_major} {[} \sphinxstyleemphasis{k} {]} {]} == \sphinxstyleemphasis{col} {[} \sphinxstyleemphasis{row\_major} {[} \sphinxstyleemphasis{k} +1{]} {]} ,

\begin{DUlineblock}{0em}
\item[]         \sphinxstyleemphasis{row} {[} \sphinxstyleemphasis{row\_major} {[} \sphinxstyleemphasis{k} {]} {]} \textless{} \sphinxstyleemphasis{row} {[} \sphinxstyleemphasis{row\_major} {[} \sphinxstyleemphasis{k} +1{]} {]}
\end{DUlineblock}

This routine generates an assert if there are two entries with the same
row and column values (if \sphinxcode{\sphinxupquote{NDEBUG}} is not defined).

\index{col\_major@\spxentry{col\_major}}\ignorespaces 

\subparagraph{col\_major}
\label{\detokenize{xsrst/py_sparse_rcv:col-major}}\label{\detokenize{xsrst/py_sparse_rcv:py-sparse-rcv-col-major}}\label{\detokenize{xsrst/py_sparse_rcv:index-14}}
This result is a numpy vector with \sphinxcode{\sphinxupquote{int}} elements
and its size \sphinxstyleemphasis{nnz} .
It sorts the sparsity pattern in column\sphinxhyphen{}major order.
To be specific,

\begin{DUlineblock}{0em}
\item[]         \sphinxstyleemphasis{row} {[} \sphinxstyleemphasis{col\_major} {[} \sphinxstyleemphasis{k} {]} {]} \textless{}= \sphinxstyleemphasis{row} {[} \sphinxstyleemphasis{col\_major} {[} \sphinxstyleemphasis{k} +1{]} {]}
\end{DUlineblock}

and if \sphinxstyleemphasis{row} {[} \sphinxstyleemphasis{col\_major} {[} \sphinxstyleemphasis{k} {]} {]} == \sphinxstyleemphasis{row} {[} \sphinxstyleemphasis{col\_major} {[} \sphinxstyleemphasis{k} +1{]} {]} ,

\begin{DUlineblock}{0em}
\item[]         \sphinxstyleemphasis{col} {[} \sphinxstyleemphasis{col\_major} {[} \sphinxstyleemphasis{k} {]} {]} \textless{} \sphinxstyleemphasis{col} {[} \sphinxstyleemphasis{col\_major} {[} \sphinxstyleemphasis{k} +1{]} {]}
\end{DUlineblock}

This routine generates an assert if there are two entries with the same
row and column values (if \sphinxcode{\sphinxupquote{NDEBUG}} is not defined).


\subparagraph{sparse\_rcv\_xam\_py}
\label{\detokenize{xsrst/sparse_rcv_xam_py:sparse-rcv-xam-py}}\label{\detokenize{xsrst/sparse_rcv_xam_py::doc}}
\index{sparse\_rcv\_xam\_py@\spxentry{sparse\_rcv\_xam\_py}}\index{python:@\spxentry{python:}}\index{sparsity@\spxentry{sparsity}}\index{patterns:@\spxentry{patterns:}}\index{example@\spxentry{example}}\index{test@\spxentry{test}}\ignorespaces 

\subparagraph{3.1.2.2.1 Python: Sparsity Patterns: Example and Test}
\label{\detokenize{xsrst/sparse_rcv_xam_py:python-sparsity-patterns-example-and-test}}\label{\detokenize{xsrst/sparse_rcv_xam_py:id1}}
\begin{sphinxVerbatim}[commandchars=\\\{\}]
\PYG{k}{def} \PYG{n+nf}{sparse\PYGZus{}rcv\PYGZus{}xam}\PYG{p}{(}\PYG{p}{)} \PYG{p}{:}
    \PYG{c+c1}{\PYGZsh{}}
    \PYG{k+kn}{import} \PYG{n+nn}{numpy}
    \PYG{k+kn}{import} \PYG{n+nn}{cppad\PYGZus{}py}
    \PYG{c+c1}{\PYGZsh{}}
    \PYG{c+c1}{\PYGZsh{} initialize return variable}
    \PYG{n}{ok} \PYG{o}{=} \PYG{k+kc}{True}
    \PYG{c+c1}{\PYGZsh{} \PYGZhy{}\PYGZhy{}\PYGZhy{}\PYGZhy{}\PYGZhy{}\PYGZhy{}\PYGZhy{}\PYGZhy{}\PYGZhy{}\PYGZhy{}\PYGZhy{}\PYGZhy{}\PYGZhy{}\PYGZhy{}\PYGZhy{}\PYGZhy{}\PYGZhy{}\PYGZhy{}\PYGZhy{}\PYGZhy{}\PYGZhy{}\PYGZhy{}\PYGZhy{}\PYGZhy{}\PYGZhy{}\PYGZhy{}\PYGZhy{}\PYGZhy{}\PYGZhy{}\PYGZhy{}\PYGZhy{}\PYGZhy{}\PYGZhy{}\PYGZhy{}\PYGZhy{}\PYGZhy{}\PYGZhy{}\PYGZhy{}\PYGZhy{}\PYGZhy{}\PYGZhy{}\PYGZhy{}\PYGZhy{}\PYGZhy{}\PYGZhy{}\PYGZhy{}\PYGZhy{}\PYGZhy{}\PYGZhy{}\PYGZhy{}\PYGZhy{}\PYGZhy{}\PYGZhy{}\PYGZhy{}\PYGZhy{}\PYGZhy{}\PYGZhy{}\PYGZhy{}\PYGZhy{}\PYGZhy{}\PYGZhy{}\PYGZhy{}\PYGZhy{}\PYGZhy{}\PYGZhy{}\PYGZhy{}\PYGZhy{}\PYGZhy{}\PYGZhy{}}
    \PYG{c+c1}{\PYGZsh{}}
    \PYG{c+c1}{\PYGZsh{} create sparsity pattern for n by n identity matrix}
    \PYG{n}{pattern} \PYG{o}{=} \PYG{n}{cppad\PYGZus{}py}\PYG{o}{.}\PYG{n}{sparse\PYGZus{}rc}\PYG{p}{(}\PYG{p}{)}
    \PYG{n}{n} \PYG{o}{=} \PYG{l+m+mi}{5}
    \PYG{n}{pattern}\PYG{o}{.}\PYG{n}{resize}\PYG{p}{(}\PYG{n}{n}\PYG{p}{,} \PYG{n}{n}\PYG{p}{,} \PYG{n}{n}\PYG{p}{)}
    \PYG{k}{for} \PYG{n}{k} \PYG{o+ow}{in} \PYG{n+nb}{range}\PYG{p}{(} \PYG{n}{n} \PYG{p}{)} \PYG{p}{:}
        \PYG{n}{pattern}\PYG{o}{.}\PYG{n}{put}\PYG{p}{(}\PYG{n}{k}\PYG{p}{,} \PYG{n}{k}\PYG{p}{,} \PYG{n}{k}\PYG{p}{)}
    \PYG{c+c1}{\PYGZsh{}}
    \PYG{c+c1}{\PYGZsh{}}
    \PYG{c+c1}{\PYGZsh{} create n by n sparse representation of identity matrix}
    \PYG{c+c1}{\PYGZsh{} (temporarly use pattern.rc untile sparse\PYGZus{}rcv wrapper is built)}
    \PYG{n}{matrix} \PYG{o}{=} \PYG{n}{cppad\PYGZus{}py}\PYG{o}{.}\PYG{n}{sparse\PYGZus{}rcv}\PYG{p}{(}\PYG{n}{pattern}\PYG{p}{)}
    \PYG{k}{for} \PYG{n}{k} \PYG{o+ow}{in} \PYG{n+nb}{range}\PYG{p}{(} \PYG{n}{n} \PYG{p}{)} \PYG{p}{:}
        \PYG{n}{matrix}\PYG{o}{.}\PYG{n}{put}\PYG{p}{(}\PYG{n}{k}\PYG{p}{,} \PYG{l+m+mf}{1.0}\PYG{p}{)}
    \PYG{c+c1}{\PYGZsh{}}
    \PYG{c+c1}{\PYGZsh{}}
    \PYG{c+c1}{\PYGZsh{} row, column indices}
    \PYG{n}{row} \PYG{o}{=} \PYG{n}{matrix}\PYG{o}{.}\PYG{n}{row}\PYG{p}{(}\PYG{p}{)}
    \PYG{n}{col} \PYG{o}{=} \PYG{n}{matrix}\PYG{o}{.}\PYG{n}{col}\PYG{p}{(}\PYG{p}{)}
    \PYG{n}{val} \PYG{o}{=} \PYG{n}{matrix}\PYG{o}{.}\PYG{n}{val}\PYG{p}{(}\PYG{p}{)}
    \PYG{c+c1}{\PYGZsh{}}
    \PYG{c+c1}{\PYGZsh{} check results}
    \PYG{k}{for} \PYG{n}{k} \PYG{o+ow}{in} \PYG{n+nb}{range}\PYG{p}{(} \PYG{n}{n} \PYG{p}{)} \PYG{p}{:}
        \PYG{n}{ok} \PYG{o}{=} \PYG{n}{ok} \PYG{o+ow}{and} \PYG{n}{row}\PYG{p}{[}\PYG{n}{k}\PYG{p}{]} \PYG{o}{==} \PYG{n}{k}
        \PYG{n}{ok} \PYG{o}{=} \PYG{n}{ok} \PYG{o+ow}{and} \PYG{n}{col}\PYG{p}{[}\PYG{n}{k}\PYG{p}{]} \PYG{o}{==} \PYG{n}{k}
        \PYG{n}{ok} \PYG{o}{=} \PYG{n}{ok} \PYG{o+ow}{and} \PYG{n}{val}\PYG{p}{[}\PYG{n}{k}\PYG{p}{]} \PYG{o}{==} \PYG{l+m+mf}{1.0}
    \PYG{c+c1}{\PYGZsh{}}
    \PYG{c+c1}{\PYGZsh{}}
    \PYG{c+c1}{\PYGZsh{} For this case, row\PYGZus{}major and col\PYGZus{}major order are same as original order}
    \PYG{n}{row\PYGZus{}major} \PYG{o}{=} \PYG{n}{matrix}\PYG{o}{.}\PYG{n}{row\PYGZus{}major}\PYG{p}{(}\PYG{p}{)}
    \PYG{n}{col\PYGZus{}major} \PYG{o}{=} \PYG{n}{matrix}\PYG{o}{.}\PYG{n}{col\PYGZus{}major}\PYG{p}{(}\PYG{p}{)}
    \PYG{k}{for} \PYG{n}{k} \PYG{o+ow}{in} \PYG{n+nb}{range}\PYG{p}{(} \PYG{n}{n} \PYG{p}{)} \PYG{p}{:}
        \PYG{n}{ok} \PYG{o}{=} \PYG{n}{ok} \PYG{o+ow}{and} \PYG{n}{row\PYGZus{}major}\PYG{p}{[}\PYG{n}{k}\PYG{p}{]} \PYG{o}{==} \PYG{n}{k}
        \PYG{n}{ok} \PYG{o}{=} \PYG{n}{ok} \PYG{o+ow}{and} \PYG{n}{col\PYGZus{}major}\PYG{p}{[}\PYG{n}{k}\PYG{p}{]} \PYG{o}{==} \PYG{n}{k}
    \PYG{c+c1}{\PYGZsh{}}
    \PYG{c+c1}{\PYGZsh{}}
    \PYG{k}{return}\PYG{p}{(} \PYG{n}{ok} \PYG{p}{)}
\PYG{c+c1}{\PYGZsh{}}
\end{sphinxVerbatim}


\bigskip\hrule\bigskip


xsrst input file: \sphinxcode{\sphinxupquote{example/python/core/sparse\_rcv\_xam.py}}

\index{example@\spxentry{example}}\ignorespaces 

\subparagraph{Example}
\label{\detokenize{xsrst/py_sparse_rcv:example}}\label{\detokenize{xsrst/py_sparse_rcv:py-sparse-rcv-example}}\label{\detokenize{xsrst/py_sparse_rcv:index-15}}
{\hyperref[\detokenize{xsrst/sparse_rcv_xam_py:id1}]{\sphinxcrossref{\DUrole{std,std-ref}{sparse\_rcv\_xam\_py}}}}


\bigskip\hrule\bigskip


xsrst input file: \sphinxcode{\sphinxupquote{lib/python/cppad\_py/sparse\_rcv.py}}


\subparagraph{py\_jac\_sparsity}
\label{\detokenize{xsrst/py_jac_sparsity:py-jac-sparsity}}\label{\detokenize{xsrst/py_jac_sparsity::doc}}
\index{py\_jac\_sparsity@\spxentry{py\_jac\_sparsity}}\index{jacobian@\spxentry{jacobian}}\index{sparsity@\spxentry{sparsity}}\index{patterns@\spxentry{patterns}}\ignorespaces 

\subparagraph{3.1.2.3 Jacobian Sparsity Patterns}
\label{\detokenize{xsrst/py_jac_sparsity:jacobian-sparsity-patterns}}\label{\detokenize{xsrst/py_jac_sparsity:id1}}\begin{itemize}
\item {} 
{\hyperref[\detokenize{xsrst/py_jac_sparsity:py-jac-sparsity-syntax}]{\sphinxcrossref{\DUrole{std,std-ref}{Syntax}}}}

\item {} \begin{description}
\item[{{\hyperref[\detokenize{xsrst/py_jac_sparsity:py-jac-sparsity-purpose}]{\sphinxcrossref{\DUrole{std,std-ref}{Purpose}}}}}] \leavevmode\begin{itemize}
\item {} 
{\hyperref[\detokenize{xsrst/py_jac_sparsity:py-jac-sparsity-purpose-for-jac-sparsity}]{\sphinxcrossref{\DUrole{std,std-ref}{for\_jac\_sparsity}}}}

\item {} 
{\hyperref[\detokenize{xsrst/py_jac_sparsity:py-jac-sparsity-purpose-rev-jac-sparsity}]{\sphinxcrossref{\DUrole{std,std-ref}{rev\_jac\_sparsity}}}}

\end{itemize}

\end{description}

\item {} 
{\hyperref[\detokenize{xsrst/py_jac_sparsity:py-jac-sparsity-x}]{\sphinxcrossref{\DUrole{std,std-ref}{x}}}}

\item {} 
{\hyperref[\detokenize{xsrst/py_jac_sparsity:py-jac-sparsity-f}]{\sphinxcrossref{\DUrole{std,std-ref}{f}}}}

\item {} 
{\hyperref[\detokenize{xsrst/py_jac_sparsity:py-jac-sparsity-pattern-in}]{\sphinxcrossref{\DUrole{std,std-ref}{pattern\_in}}}}

\item {} 
{\hyperref[\detokenize{xsrst/py_jac_sparsity:py-jac-sparsity-pattern-out}]{\sphinxcrossref{\DUrole{std,std-ref}{pattern\_out}}}}

\item {} 
{\hyperref[\detokenize{xsrst/py_jac_sparsity:py-jac-sparsity-sparsity-for-entire-jacobian}]{\sphinxcrossref{\DUrole{std,std-ref}{Sparsity for Entire Jacobian}}}}

\item {} 
{\hyperref[\detokenize{xsrst/py_jac_sparsity:py-jac-sparsity-example}]{\sphinxcrossref{\DUrole{std,std-ref}{Example}}}}

\end{itemize}

\index{syntax@\spxentry{syntax}}\ignorespaces 

\subparagraph{Syntax}
\label{\detokenize{xsrst/py_jac_sparsity:syntax}}\label{\detokenize{xsrst/py_jac_sparsity:py-jac-sparsity-syntax}}\label{\detokenize{xsrst/py_jac_sparsity:index-1}}
\begin{DUlineblock}{0em}
\item[] \sphinxstyleemphasis{f} . \sphinxcode{\sphinxupquote{for\_jac\_sparsity}} ( \sphinxstyleemphasis{pattern\_in} , \sphinxstyleemphasis{pattern\_out} )
\item[] \sphinxstyleemphasis{f} . \sphinxcode{\sphinxupquote{rev\_jac\_sparsity}} ( \sphinxstyleemphasis{pattern\_in} , \sphinxstyleemphasis{pattern\_out} )
\end{DUlineblock}

\index{purpose@\spxentry{purpose}}\ignorespaces 

\subparagraph{Purpose}
\label{\detokenize{xsrst/py_jac_sparsity:purpose}}\label{\detokenize{xsrst/py_jac_sparsity:py-jac-sparsity-purpose}}\label{\detokenize{xsrst/py_jac_sparsity:index-2}}
We use \(F : \B{R}^n \rightarrow \B{R}^m\) to denote the
function corresponding to the operation sequence stored in \sphinxstyleemphasis{f} .

\index{for\_jac\_sparsity@\spxentry{for\_jac\_sparsity}}\ignorespaces 

\subparagraph{for\_jac\_sparsity}
\label{\detokenize{xsrst/py_jac_sparsity:for-jac-sparsity}}\label{\detokenize{xsrst/py_jac_sparsity:py-jac-sparsity-purpose-for-jac-sparsity}}\label{\detokenize{xsrst/py_jac_sparsity:index-3}}
Fix \(R \in \B{R}^{n \times \ell}\) and define the function
\begin{equation*}
\begin{split}J(x) = F^{(1)} ( x ) * R\end{split}
\end{equation*}
Given a sparsity pattern for \(R\),
\sphinxcode{\sphinxupquote{for\_jac\_sparsity}} computes a sparsity pattern for \(J(x)\).

\index{rev\_jac\_sparsity@\spxentry{rev\_jac\_sparsity}}\ignorespaces 

\subparagraph{rev\_jac\_sparsity}
\label{\detokenize{xsrst/py_jac_sparsity:rev-jac-sparsity}}\label{\detokenize{xsrst/py_jac_sparsity:py-jac-sparsity-purpose-rev-jac-sparsity}}\label{\detokenize{xsrst/py_jac_sparsity:index-4}}
Fix \(R \in \B{R}^{\ell \times m}\) and define the function
\begin{equation*}
\begin{split}J(x) = R * F^{(1)} ( x )\end{split}
\end{equation*}
Given a sparsity pattern for \(R\),
\sphinxcode{\sphinxupquote{rev\_jac\_sparsity}} computes a sparsity pattern for \(J(x)\).

\index{x@\spxentry{x}}\ignorespaces 

\subparagraph{x}
\label{\detokenize{xsrst/py_jac_sparsity:x}}\label{\detokenize{xsrst/py_jac_sparsity:py-jac-sparsity-x}}\label{\detokenize{xsrst/py_jac_sparsity:index-5}}
Note that a sparsity pattern for \(J(x)\) corresponds to the
operation sequence stored in \sphinxstyleemphasis{f} and does not depend on
the argument \sphinxstyleemphasis{x} .

\index{f@\spxentry{f}}\ignorespaces 

\subparagraph{f}
\label{\detokenize{xsrst/py_jac_sparsity:f}}\label{\detokenize{xsrst/py_jac_sparsity:py-jac-sparsity-f}}\label{\detokenize{xsrst/py_jac_sparsity:index-6}}
This object must have been returned by a previous call to the python
{\hyperref[\detokenize{xsrst/py_fun_ctor:id1}]{\sphinxcrossref{\DUrole{std,std-ref}{d\_fun}}}} constructor.
The object \sphinxstyleemphasis{f} is not constant when using
\sphinxcode{\sphinxupquote{for\_jac\_sparsity}} .
After a call to \sphinxcode{\sphinxupquote{for\_jac\_sparsity}} , a sparsity pattern
for each of the variables in the operation sequence
is held in \sphinxstyleemphasis{f} for possible later use during
reverse Hessian sparsity calculations.

\index{pattern\_in@\spxentry{pattern\_in}}\ignorespaces 

\subparagraph{pattern\_in}
\label{\detokenize{xsrst/py_jac_sparsity:pattern-in}}\label{\detokenize{xsrst/py_jac_sparsity:py-jac-sparsity-pattern-in}}\label{\detokenize{xsrst/py_jac_sparsity:index-7}}
This argument must have be a {\hyperref[\detokenize{xsrst/py_sparse_rc:py-sparse-rc-pattern}]{\sphinxcrossref{\DUrole{std,std-ref}{pattern}}}}
returned by the \sphinxcode{\sphinxupquote{sparse\_rc}} constructor.
It specifies the sparsity pattern for \(R\).

\index{pattern\_out@\spxentry{pattern\_out}}\ignorespaces 

\subparagraph{pattern\_out}
\label{\detokenize{xsrst/py_jac_sparsity:pattern-out}}\label{\detokenize{xsrst/py_jac_sparsity:py-jac-sparsity-pattern-out}}\label{\detokenize{xsrst/py_jac_sparsity:index-8}}
This argument must have be a {\hyperref[\detokenize{xsrst/py_sparse_rc:py-sparse-rc-pattern}]{\sphinxcrossref{\DUrole{std,std-ref}{pattern}}}}
returned by the \sphinxcode{\sphinxupquote{sparse\_rc}} constructor.
This input value of \sphinxstyleemphasis{pattern\_out} does not matter.
Upon return \sphinxstyleemphasis{pattern\_out} is a sparsity pattern for
\(J(x)\).

\index{sparsity@\spxentry{sparsity}}\index{for@\spxentry{for}}\index{entire@\spxentry{entire}}\index{jacobian@\spxentry{jacobian}}\ignorespaces 

\subparagraph{Sparsity for Entire Jacobian}
\label{\detokenize{xsrst/py_jac_sparsity:sparsity-for-entire-jacobian}}\label{\detokenize{xsrst/py_jac_sparsity:py-jac-sparsity-sparsity-for-entire-jacobian}}\label{\detokenize{xsrst/py_jac_sparsity:index-9}}
Suppose that \(R\) is the identity matrix.
In this case, \sphinxstyleemphasis{pattern\_out} is a sparsity pattern for
\(F^{(1)} ( x )\).


\subparagraph{sparse\_jac\_pattern\_xam\_py}
\label{\detokenize{xsrst/sparse_jac_pattern_xam_py:sparse-jac-pattern-xam-py}}\label{\detokenize{xsrst/sparse_jac_pattern_xam_py::doc}}
\index{sparse\_jac\_pattern\_xam\_py@\spxentry{sparse\_jac\_pattern\_xam\_py}}\index{python:@\spxentry{python:}}\index{jacobian@\spxentry{jacobian}}\index{sparsity@\spxentry{sparsity}}\index{patterns:@\spxentry{patterns:}}\index{example@\spxentry{example}}\index{test@\spxentry{test}}\ignorespaces 

\subparagraph{3.1.2.3.1 Python: Jacobian Sparsity Patterns: Example and Test}
\label{\detokenize{xsrst/sparse_jac_pattern_xam_py:python-jacobian-sparsity-patterns-example-and-test}}\label{\detokenize{xsrst/sparse_jac_pattern_xam_py:id1}}
\begin{sphinxVerbatim}[commandchars=\\\{\}]
\PYG{k}{def} \PYG{n+nf}{sparse\PYGZus{}jac\PYGZus{}pattern\PYGZus{}xam}\PYG{p}{(}\PYG{p}{)} \PYG{p}{:}
    \PYG{c+c1}{\PYGZsh{}}
    \PYG{k+kn}{import} \PYG{n+nn}{numpy}
    \PYG{k+kn}{import} \PYG{n+nn}{cppad\PYGZus{}py}
    \PYG{c+c1}{\PYGZsh{}}
    \PYG{c+c1}{\PYGZsh{} initialize return variable}
    \PYG{n}{ok} \PYG{o}{=} \PYG{k+kc}{True}
    \PYG{c+c1}{\PYGZsh{} \PYGZhy{}\PYGZhy{}\PYGZhy{}\PYGZhy{}\PYGZhy{}\PYGZhy{}\PYGZhy{}\PYGZhy{}\PYGZhy{}\PYGZhy{}\PYGZhy{}\PYGZhy{}\PYGZhy{}\PYGZhy{}\PYGZhy{}\PYGZhy{}\PYGZhy{}\PYGZhy{}\PYGZhy{}\PYGZhy{}\PYGZhy{}\PYGZhy{}\PYGZhy{}\PYGZhy{}\PYGZhy{}\PYGZhy{}\PYGZhy{}\PYGZhy{}\PYGZhy{}\PYGZhy{}\PYGZhy{}\PYGZhy{}\PYGZhy{}\PYGZhy{}\PYGZhy{}\PYGZhy{}\PYGZhy{}\PYGZhy{}\PYGZhy{}\PYGZhy{}\PYGZhy{}\PYGZhy{}\PYGZhy{}\PYGZhy{}\PYGZhy{}\PYGZhy{}\PYGZhy{}\PYGZhy{}\PYGZhy{}\PYGZhy{}\PYGZhy{}\PYGZhy{}\PYGZhy{}\PYGZhy{}\PYGZhy{}\PYGZhy{}\PYGZhy{}\PYGZhy{}\PYGZhy{}\PYGZhy{}\PYGZhy{}\PYGZhy{}\PYGZhy{}\PYGZhy{}\PYGZhy{}\PYGZhy{}\PYGZhy{}\PYGZhy{}\PYGZhy{}}
    \PYG{c+c1}{\PYGZsh{} number of dependent and independent variables}
    \PYG{n}{n} \PYG{o}{=} \PYG{l+m+mi}{3}
    \PYG{c+c1}{\PYGZsh{}}
    \PYG{c+c1}{\PYGZsh{} create the independent variables ax}
    \PYG{n}{x} \PYG{o}{=} \PYG{n}{numpy}\PYG{o}{.}\PYG{n}{empty}\PYG{p}{(}\PYG{n}{n}\PYG{p}{,} \PYG{n}{dtype}\PYG{o}{=}\PYG{n+nb}{float}\PYG{p}{)}
    \PYG{k}{for} \PYG{n}{i} \PYG{o+ow}{in} \PYG{n+nb}{range}\PYG{p}{(} \PYG{n}{n}  \PYG{p}{)} \PYG{p}{:}
        \PYG{n}{x}\PYG{p}{[}\PYG{n}{i}\PYG{p}{]} \PYG{o}{=} \PYG{n}{i} \PYG{o}{+} \PYG{l+m+mf}{2.0}
    \PYG{c+c1}{\PYGZsh{}}
    \PYG{n}{ax} \PYG{o}{=} \PYG{n}{cppad\PYGZus{}py}\PYG{o}{.}\PYG{n}{independent}\PYG{p}{(}\PYG{n}{x}\PYG{p}{)}
    \PYG{c+c1}{\PYGZsh{}}
    \PYG{c+c1}{\PYGZsh{} create dependent variables ay with ay[i] = ax[j]}
    \PYG{c+c1}{\PYGZsh{} where i = mod(j + 1, n)}
    \PYG{n}{ay} \PYG{o}{=} \PYG{n}{numpy}\PYG{o}{.}\PYG{n}{empty}\PYG{p}{(}\PYG{n}{n}\PYG{p}{,} \PYG{n}{dtype}\PYG{o}{=}\PYG{n}{cppad\PYGZus{}py}\PYG{o}{.}\PYG{n}{a\PYGZus{}double}\PYG{p}{)}
    \PYG{k}{for} \PYG{n}{j} \PYG{o+ow}{in} \PYG{n+nb}{range}\PYG{p}{(} \PYG{n}{n}  \PYG{p}{)} \PYG{p}{:}
        \PYG{n}{i} \PYG{o}{=} \PYG{n}{j}\PYG{o}{+}\PYG{l+m+mi}{1}
        \PYG{k}{if} \PYG{n}{i} \PYG{o}{\PYGZgt{}}\PYG{o}{=} \PYG{n}{n}  \PYG{p}{:}
            \PYG{n}{i} \PYG{o}{=} \PYG{n}{i} \PYG{o}{\PYGZhy{}} \PYG{n}{n}
        \PYG{c+c1}{\PYGZsh{}}
        \PYG{n}{ay\PYGZus{}i} \PYG{o}{=} \PYG{n}{ax}\PYG{p}{[}\PYG{n}{j}\PYG{p}{]}
        \PYG{n}{ay}\PYG{p}{[}\PYG{n}{i}\PYG{p}{]} \PYG{o}{=} \PYG{n}{ay\PYGZus{}i}
    \PYG{c+c1}{\PYGZsh{}}
    \PYG{c+c1}{\PYGZsh{}}
    \PYG{c+c1}{\PYGZsh{} define af corresponding to f(x)}
    \PYG{n}{f}  \PYG{o}{=} \PYG{n}{cppad\PYGZus{}py}\PYG{o}{.}\PYG{n}{d\PYGZus{}fun}\PYG{p}{(}\PYG{n}{ax}\PYG{p}{,} \PYG{n}{ay}\PYG{p}{)}
    \PYG{c+c1}{\PYGZsh{}}
    \PYG{c+c1}{\PYGZsh{} sparsity pattern for identity matrix}
    \PYG{n}{pat\PYGZus{}in} \PYG{o}{=} \PYG{n}{cppad\PYGZus{}py}\PYG{o}{.}\PYG{n}{sparse\PYGZus{}rc}\PYG{p}{(}\PYG{p}{)}
    \PYG{n}{pat\PYGZus{}in}\PYG{o}{.}\PYG{n}{resize}\PYG{p}{(}\PYG{n}{n}\PYG{p}{,} \PYG{n}{n}\PYG{p}{,} \PYG{n}{n}\PYG{p}{)}
    \PYG{k}{for} \PYG{n}{k} \PYG{o+ow}{in} \PYG{n+nb}{range}\PYG{p}{(} \PYG{n}{n} \PYG{p}{)} \PYG{p}{:}
        \PYG{n}{pat\PYGZus{}in}\PYG{o}{.}\PYG{n}{put}\PYG{p}{(}\PYG{n}{k}\PYG{p}{,} \PYG{n}{k}\PYG{p}{,} \PYG{n}{k}\PYG{p}{)}
    \PYG{c+c1}{\PYGZsh{}}
    \PYG{c+c1}{\PYGZsh{}}
    \PYG{c+c1}{\PYGZsh{} loop over forward and reverse mode}
    \PYG{k}{for} \PYG{n}{mode} \PYG{o+ow}{in} \PYG{n+nb}{range}\PYG{p}{(} \PYG{l+m+mi}{2} \PYG{p}{)} \PYG{p}{:}
        \PYG{n}{pat\PYGZus{}out} \PYG{o}{=} \PYG{n}{cppad\PYGZus{}py}\PYG{o}{.}\PYG{n}{sparse\PYGZus{}rc}\PYG{p}{(}\PYG{p}{)}
        \PYG{k}{if} \PYG{n}{mode} \PYG{o}{==} \PYG{l+m+mi}{0}  \PYG{p}{:}
            \PYG{n}{f}\PYG{o}{.}\PYG{n}{for\PYGZus{}jac\PYGZus{}sparsity}\PYG{p}{(}\PYG{n}{pat\PYGZus{}in}\PYG{p}{,} \PYG{n}{pat\PYGZus{}out}\PYG{p}{)}
        \PYG{c+c1}{\PYGZsh{}}
        \PYG{k}{if} \PYG{n}{mode} \PYG{o}{==} \PYG{l+m+mi}{1}  \PYG{p}{:}
            \PYG{n}{f}\PYG{o}{.}\PYG{n}{rev\PYGZus{}jac\PYGZus{}sparsity}\PYG{p}{(}\PYG{n}{pat\PYGZus{}in}\PYG{p}{,} \PYG{n}{pat\PYGZus{}out}\PYG{p}{)}
        \PYG{c+c1}{\PYGZsh{}}
        \PYG{c+c1}{\PYGZsh{}}
        \PYG{c+c1}{\PYGZsh{} check that result is sparsity pattern for Jacobian}
        \PYG{n}{ok} \PYG{o}{=} \PYG{n}{ok} \PYG{o+ow}{and} \PYG{n}{n} \PYG{o}{==} \PYG{n}{pat\PYGZus{}out}\PYG{o}{.}\PYG{n}{nnz}\PYG{p}{(}\PYG{p}{)}
        \PYG{n}{col\PYGZus{}major} \PYG{o}{=} \PYG{n}{pat\PYGZus{}out}\PYG{o}{.}\PYG{n}{col\PYGZus{}major}\PYG{p}{(}\PYG{p}{)}
        \PYG{n}{row} \PYG{o}{=} \PYG{n}{pat\PYGZus{}out}\PYG{o}{.}\PYG{n}{row}\PYG{p}{(}\PYG{p}{)}
        \PYG{n}{col} \PYG{o}{=} \PYG{n}{pat\PYGZus{}out}\PYG{o}{.}\PYG{n}{col}\PYG{p}{(}\PYG{p}{)}
        \PYG{k}{for} \PYG{n}{k} \PYG{o+ow}{in} \PYG{n+nb}{range}\PYG{p}{(} \PYG{n}{n} \PYG{p}{)} \PYG{p}{:}
            \PYG{n}{ell} \PYG{o}{=} \PYG{n}{col\PYGZus{}major}\PYG{p}{[}\PYG{n}{k}\PYG{p}{]}
            \PYG{n}{r} \PYG{o}{=} \PYG{n}{row}\PYG{p}{[}\PYG{n}{ell}\PYG{p}{]}
            \PYG{n}{c} \PYG{o}{=} \PYG{n}{col}\PYG{p}{[}\PYG{n}{ell}\PYG{p}{]}
            \PYG{n}{i} \PYG{o}{=} \PYG{n}{c}\PYG{o}{+}\PYG{l+m+mi}{1}
            \PYG{k}{if} \PYG{n}{i} \PYG{o}{\PYGZgt{}}\PYG{o}{=}  \PYG{n}{n}  \PYG{p}{:}
                \PYG{n}{i} \PYG{o}{=} \PYG{n}{i} \PYG{o}{\PYGZhy{}} \PYG{n}{n}
            \PYG{c+c1}{\PYGZsh{}}
            \PYG{n}{ok} \PYG{o}{=} \PYG{n}{ok} \PYG{o+ow}{and} \PYG{n}{c} \PYG{o}{==} \PYG{n}{k}
            \PYG{n}{ok} \PYG{o}{=} \PYG{n}{ok} \PYG{o+ow}{and} \PYG{n}{r} \PYG{o}{==} \PYG{n}{i}
        \PYG{c+c1}{\PYGZsh{}}
    \PYG{c+c1}{\PYGZsh{}}
    \PYG{c+c1}{\PYGZsh{}}
    \PYG{k}{return}\PYG{p}{(} \PYG{n}{ok} \PYG{p}{)}
\PYG{c+c1}{\PYGZsh{}}
\end{sphinxVerbatim}


\bigskip\hrule\bigskip


xsrst input file: \sphinxcode{\sphinxupquote{example/python/core/sparse\_jac\_pattern\_xam.py}}

\index{example@\spxentry{example}}\ignorespaces 

\subparagraph{Example}
\label{\detokenize{xsrst/py_jac_sparsity:example}}\label{\detokenize{xsrst/py_jac_sparsity:py-jac-sparsity-example}}\label{\detokenize{xsrst/py_jac_sparsity:index-10}}
{\hyperref[\detokenize{xsrst/sparse_jac_pattern_xam_py:id1}]{\sphinxcrossref{\DUrole{std,std-ref}{python}}}}


\bigskip\hrule\bigskip


xsrst input file: \sphinxcode{\sphinxupquote{lib/python/cppad\_py/jac\_sparsity.xsrst}}


\subparagraph{py\_hes\_sparsity}
\label{\detokenize{xsrst/py_hes_sparsity:py-hes-sparsity}}\label{\detokenize{xsrst/py_hes_sparsity::doc}}
\index{py\_hes\_sparsity@\spxentry{py\_hes\_sparsity}}\index{hessian@\spxentry{hessian}}\index{sparsity@\spxentry{sparsity}}\index{patterns@\spxentry{patterns}}\ignorespaces 

\subparagraph{3.1.2.4 Hessian Sparsity Patterns}
\label{\detokenize{xsrst/py_hes_sparsity:hessian-sparsity-patterns}}\label{\detokenize{xsrst/py_hes_sparsity:id1}}\begin{itemize}
\item {} 
{\hyperref[\detokenize{xsrst/py_hes_sparsity:py-hes-sparsity-syntax}]{\sphinxcrossref{\DUrole{std,std-ref}{Syntax}}}}

\item {} 
{\hyperref[\detokenize{xsrst/py_hes_sparsity:py-hes-sparsity-purpose}]{\sphinxcrossref{\DUrole{std,std-ref}{Purpose}}}}

\item {} 
{\hyperref[\detokenize{xsrst/py_hes_sparsity:py-hes-sparsity-x}]{\sphinxcrossref{\DUrole{std,std-ref}{x}}}}

\item {} 
{\hyperref[\detokenize{xsrst/py_hes_sparsity:py-hes-sparsity-f}]{\sphinxcrossref{\DUrole{std,std-ref}{f}}}}

\item {} 
{\hyperref[\detokenize{xsrst/py_hes_sparsity:py-hes-sparsity-select-domain}]{\sphinxcrossref{\DUrole{std,std-ref}{select\_domain}}}}

\item {} 
{\hyperref[\detokenize{xsrst/py_hes_sparsity:py-hes-sparsity-select-range}]{\sphinxcrossref{\DUrole{std,std-ref}{select\_range}}}}

\item {} 
{\hyperref[\detokenize{xsrst/py_hes_sparsity:py-hes-sparsity-pattern-out}]{\sphinxcrossref{\DUrole{std,std-ref}{pattern\_out}}}}

\item {} 
{\hyperref[\detokenize{xsrst/py_hes_sparsity:py-hes-sparsity-sparsity-for-component-wise-hessian}]{\sphinxcrossref{\DUrole{std,std-ref}{Sparsity for Component Wise Hessian}}}}

\item {} 
{\hyperref[\detokenize{xsrst/py_hes_sparsity:py-hes-sparsity-example}]{\sphinxcrossref{\DUrole{std,std-ref}{Example}}}}

\end{itemize}

\index{syntax@\spxentry{syntax}}\ignorespaces 

\subparagraph{Syntax}
\label{\detokenize{xsrst/py_hes_sparsity:syntax}}\label{\detokenize{xsrst/py_hes_sparsity:py-hes-sparsity-syntax}}\label{\detokenize{xsrst/py_hes_sparsity:index-1}}
\begin{DUlineblock}{0em}
\item[] \sphinxstyleemphasis{f} . \sphinxcode{\sphinxupquote{for\_hes\_sparsity}} ( \sphinxstyleemphasis{select\_domain} , \sphinxstyleemphasis{select\_range} , \sphinxstyleemphasis{pattern\_out} )
\end{DUlineblock}

\sphinxstyleemphasis{f} . \sphinxcode{\sphinxupquote{rev\_hes\_sparsity}} ( \sphinxstyleemphasis{select\_domain} , \sphinxstyleemphasis{select\_range} , \sphinxstyleemphasis{pattern\_out} )

\index{purpose@\spxentry{purpose}}\ignorespaces 

\subparagraph{Purpose}
\label{\detokenize{xsrst/py_hes_sparsity:purpose}}\label{\detokenize{xsrst/py_hes_sparsity:py-hes-sparsity-purpose}}\label{\detokenize{xsrst/py_hes_sparsity:index-2}}
We use \(F : \B{R}^n \rightarrow \B{R}^m\) to denote the
function corresponding to the operation sequence stored in \sphinxstyleemphasis{f} .
Fix a diagonal matrix \(D \in \B{R}^{n \times n}\), fix a vector
\(r \in \B{R}^m\), and define
\begin{equation*}
\begin{split}H(x) = D (r^\R{T} F)^{(2)} ( x ) D\end{split}
\end{equation*}
Given a sparsity pattern for \(D\) and \(r\),
these routines compute a sparsity pattern for \(H(x)\).

\index{x@\spxentry{x}}\ignorespaces 

\subparagraph{x}
\label{\detokenize{xsrst/py_hes_sparsity:x}}\label{\detokenize{xsrst/py_hes_sparsity:py-hes-sparsity-x}}\label{\detokenize{xsrst/py_hes_sparsity:index-3}}
Note that a sparsity pattern for \(H(x)\) corresponds to the
operation sequence stored in \sphinxstyleemphasis{f} and does not depend on
the argument \sphinxstyleemphasis{x} .

\index{f@\spxentry{f}}\ignorespaces 

\subparagraph{f}
\label{\detokenize{xsrst/py_hes_sparsity:f}}\label{\detokenize{xsrst/py_hes_sparsity:py-hes-sparsity-f}}\label{\detokenize{xsrst/py_hes_sparsity:index-4}}
This object must have been returned by a previous call to the python
{\hyperref[\detokenize{xsrst/py_fun_ctor:id1}]{\sphinxcrossref{\DUrole{std,std-ref}{d\_fun}}}} constructor.

\index{select\_domain@\spxentry{select\_domain}}\ignorespaces 

\subparagraph{select\_domain}
\label{\detokenize{xsrst/py_hes_sparsity:select-domain}}\label{\detokenize{xsrst/py_hes_sparsity:py-hes-sparsity-select-domain}}\label{\detokenize{xsrst/py_hes_sparsity:index-5}}
The argument is a numpy vector with \sphinxcode{\sphinxupquote{bool}} elements.
It has size \sphinxstyleemphasis{n} and is a sparsity pattern for the diagonal of
\(D\); i.e., \sphinxstyleemphasis{select\_domain} {[} \sphinxstyleemphasis{j} {]} is true if and only if
\(D_{j,j}\) is possibly non\sphinxhyphen{}zero.

\index{select\_range@\spxentry{select\_range}}\ignorespaces 

\subparagraph{select\_range}
\label{\detokenize{xsrst/py_hes_sparsity:select-range}}\label{\detokenize{xsrst/py_hes_sparsity:py-hes-sparsity-select-range}}\label{\detokenize{xsrst/py_hes_sparsity:index-6}}
The argument is a numpy vector with \sphinxcode{\sphinxupquote{bool}} elements.
It has size \sphinxstyleemphasis{m} and is a sparsity pattern for the vector
\(r\); i.e., \sphinxstyleemphasis{select\_range} {[} \sphinxstyleemphasis{i} {]} is true if and only if
\(r_i\) is possibly non\sphinxhyphen{}zero.

\index{pattern\_out@\spxentry{pattern\_out}}\ignorespaces 

\subparagraph{pattern\_out}
\label{\detokenize{xsrst/py_hes_sparsity:pattern-out}}\label{\detokenize{xsrst/py_hes_sparsity:py-hes-sparsity-pattern-out}}\label{\detokenize{xsrst/py_hes_sparsity:index-7}}
This argument must have be a {\hyperref[\detokenize{xsrst/py_sparse_rc:py-sparse-rc-pattern}]{\sphinxcrossref{\DUrole{std,std-ref}{pattern}}}}
returned by the \sphinxcode{\sphinxupquote{sparse\_rc}} constructor.
This input value of \sphinxstyleemphasis{pattern\_out} does not matter.
Upon return \sphinxstyleemphasis{pattern\_out} is a sparsity pattern for
\(H(x)\).

\index{sparsity@\spxentry{sparsity}}\index{for@\spxentry{for}}\index{component@\spxentry{component}}\index{wise@\spxentry{wise}}\index{hessian@\spxentry{hessian}}\ignorespaces 

\subparagraph{Sparsity for Component Wise Hessian}
\label{\detokenize{xsrst/py_hes_sparsity:sparsity-for-component-wise-hessian}}\label{\detokenize{xsrst/py_hes_sparsity:py-hes-sparsity-sparsity-for-component-wise-hessian}}\label{\detokenize{xsrst/py_hes_sparsity:index-8}}
Suppose that \(D\) is the identity matrix,
and only the \sphinxstyleemphasis{i}\sphinxhyphen{}th component of \sphinxstyleemphasis{r} is possibly non\sphinxhyphen{}zero.
In this case, \sphinxstyleemphasis{pattern\_out} is a sparsity pattern for
\(F_i^{(2)} ( x )\).


\subparagraph{sparse\_hes\_pattern\_xam\_py}
\label{\detokenize{xsrst/sparse_hes_pattern_xam_py:sparse-hes-pattern-xam-py}}\label{\detokenize{xsrst/sparse_hes_pattern_xam_py::doc}}
\index{sparse\_hes\_pattern\_xam\_py@\spxentry{sparse\_hes\_pattern\_xam\_py}}\index{python:@\spxentry{python:}}\index{hessian@\spxentry{hessian}}\index{sparsity@\spxentry{sparsity}}\index{patterns:@\spxentry{patterns:}}\index{example@\spxentry{example}}\index{test@\spxentry{test}}\ignorespaces 

\subparagraph{3.1.2.4.1 Python: Hessian Sparsity Patterns: Example and Test}
\label{\detokenize{xsrst/sparse_hes_pattern_xam_py:python-hessian-sparsity-patterns-example-and-test}}\label{\detokenize{xsrst/sparse_hes_pattern_xam_py:id1}}
\begin{sphinxVerbatim}[commandchars=\\\{\}]
\PYG{k}{def} \PYG{n+nf}{sparse\PYGZus{}hes\PYGZus{}pattern\PYGZus{}xam}\PYG{p}{(}\PYG{p}{)} \PYG{p}{:}
    \PYG{c+c1}{\PYGZsh{}}
    \PYG{k+kn}{import} \PYG{n+nn}{numpy}
    \PYG{k+kn}{import} \PYG{n+nn}{cppad\PYGZus{}py}
    \PYG{c+c1}{\PYGZsh{}}
    \PYG{c+c1}{\PYGZsh{} initialize return variable}
    \PYG{n}{ok} \PYG{o}{=} \PYG{k+kc}{True}
    \PYG{c+c1}{\PYGZsh{} \PYGZhy{}\PYGZhy{}\PYGZhy{}\PYGZhy{}\PYGZhy{}\PYGZhy{}\PYGZhy{}\PYGZhy{}\PYGZhy{}\PYGZhy{}\PYGZhy{}\PYGZhy{}\PYGZhy{}\PYGZhy{}\PYGZhy{}\PYGZhy{}\PYGZhy{}\PYGZhy{}\PYGZhy{}\PYGZhy{}\PYGZhy{}\PYGZhy{}\PYGZhy{}\PYGZhy{}\PYGZhy{}\PYGZhy{}\PYGZhy{}\PYGZhy{}\PYGZhy{}\PYGZhy{}\PYGZhy{}\PYGZhy{}\PYGZhy{}\PYGZhy{}\PYGZhy{}\PYGZhy{}\PYGZhy{}\PYGZhy{}\PYGZhy{}\PYGZhy{}\PYGZhy{}\PYGZhy{}\PYGZhy{}\PYGZhy{}\PYGZhy{}\PYGZhy{}\PYGZhy{}\PYGZhy{}\PYGZhy{}\PYGZhy{}\PYGZhy{}\PYGZhy{}\PYGZhy{}\PYGZhy{}\PYGZhy{}\PYGZhy{}\PYGZhy{}\PYGZhy{}\PYGZhy{}\PYGZhy{}\PYGZhy{}\PYGZhy{}\PYGZhy{}\PYGZhy{}\PYGZhy{}\PYGZhy{}\PYGZhy{}\PYGZhy{}\PYGZhy{}}
    \PYG{c+c1}{\PYGZsh{} number of dependent and independent variables}
    \PYG{n}{n} \PYG{o}{=} \PYG{l+m+mi}{3}
    \PYG{c+c1}{\PYGZsh{}}
    \PYG{c+c1}{\PYGZsh{} create the independent variables ax}
    \PYG{n}{x} \PYG{o}{=} \PYG{n}{numpy}\PYG{o}{.}\PYG{n}{empty}\PYG{p}{(}\PYG{n}{n}\PYG{p}{,} \PYG{n}{dtype}\PYG{o}{=}\PYG{n+nb}{float}\PYG{p}{)}
    \PYG{k}{for} \PYG{n}{i} \PYG{o+ow}{in} \PYG{n+nb}{range}\PYG{p}{(} \PYG{n}{n}  \PYG{p}{)} \PYG{p}{:}
        \PYG{n}{x}\PYG{p}{[}\PYG{n}{i}\PYG{p}{]} \PYG{o}{=} \PYG{n}{i} \PYG{o}{+} \PYG{l+m+mf}{2.0}
    \PYG{c+c1}{\PYGZsh{}}
    \PYG{n}{ax} \PYG{o}{=} \PYG{n}{cppad\PYGZus{}py}\PYG{o}{.}\PYG{n}{independent}\PYG{p}{(}\PYG{n}{x}\PYG{p}{)}
    \PYG{c+c1}{\PYGZsh{}}
    \PYG{c+c1}{\PYGZsh{} create dependent variables ay with ay[i] = ax[j] * ax[i]}
    \PYG{c+c1}{\PYGZsh{} where i = mod(j + 1, n)}
    \PYG{n}{ay} \PYG{o}{=} \PYG{n}{numpy}\PYG{o}{.}\PYG{n}{empty}\PYG{p}{(}\PYG{n}{n}\PYG{p}{,} \PYG{n}{dtype}\PYG{o}{=}\PYG{n}{cppad\PYGZus{}py}\PYG{o}{.}\PYG{n}{a\PYGZus{}double}\PYG{p}{)}
    \PYG{k}{for} \PYG{n}{j} \PYG{o+ow}{in} \PYG{n+nb}{range}\PYG{p}{(} \PYG{n}{n}  \PYG{p}{)} \PYG{p}{:}
        \PYG{n}{i} \PYG{o}{=} \PYG{n}{j}\PYG{o}{+}\PYG{l+m+mi}{1}
        \PYG{k}{if} \PYG{n}{i} \PYG{o}{\PYGZgt{}}\PYG{o}{=} \PYG{n}{n}  \PYG{p}{:}
            \PYG{n}{i} \PYG{o}{=} \PYG{n}{i} \PYG{o}{\PYGZhy{}} \PYG{n}{n}
        \PYG{c+c1}{\PYGZsh{}}
        \PYG{n}{ay\PYGZus{}i} \PYG{o}{=} \PYG{n}{ax}\PYG{p}{[}\PYG{n}{i}\PYG{p}{]} \PYG{o}{*} \PYG{n}{ax}\PYG{p}{[}\PYG{n}{j}\PYG{p}{]}
        \PYG{n}{ay}\PYG{p}{[}\PYG{n}{i}\PYG{p}{]} \PYG{o}{=} \PYG{n}{ay\PYGZus{}i}
    \PYG{c+c1}{\PYGZsh{}}
    \PYG{c+c1}{\PYGZsh{}}
    \PYG{c+c1}{\PYGZsh{} define af corresponding to f(x)}
    \PYG{n}{f}  \PYG{o}{=} \PYG{n}{cppad\PYGZus{}py}\PYG{o}{.}\PYG{n}{d\PYGZus{}fun}\PYG{p}{(}\PYG{n}{ax}\PYG{p}{,} \PYG{n}{ay}\PYG{p}{)}
    \PYG{c+c1}{\PYGZsh{}}
    \PYG{c+c1}{\PYGZsh{} Set select\PYGZus{}d (domain) to all true, initial select\PYGZus{}r (range) to all false}
    \PYG{n}{select\PYGZus{}d} \PYG{o}{=} \PYG{n}{numpy}\PYG{o}{.}\PYG{n}{empty}\PYG{p}{(}\PYG{n}{n}\PYG{p}{,} \PYG{n}{dtype}\PYG{o}{=}\PYG{n+nb}{bool}\PYG{p}{)}
    \PYG{n}{select\PYGZus{}r} \PYG{o}{=} \PYG{n}{numpy}\PYG{o}{.}\PYG{n}{empty}\PYG{p}{(}\PYG{n}{n}\PYG{p}{,} \PYG{n}{dtype}\PYG{o}{=}\PYG{n+nb}{bool}\PYG{p}{)}
    \PYG{k}{for} \PYG{n}{i} \PYG{o+ow}{in} \PYG{n+nb}{range}\PYG{p}{(} \PYG{n}{n} \PYG{p}{)} \PYG{p}{:}
        \PYG{n}{select\PYGZus{}d}\PYG{p}{[}\PYG{n}{i}\PYG{p}{]} \PYG{o}{=} \PYG{k+kc}{True}
        \PYG{n}{select\PYGZus{}r}\PYG{p}{[}\PYG{n}{i}\PYG{p}{]} \PYG{o}{=} \PYG{k+kc}{False}
    \PYG{c+c1}{\PYGZsh{}}
    \PYG{c+c1}{\PYGZsh{}}
    \PYG{c+c1}{\PYGZsh{} only select component 0 of the range function}
    \PYG{c+c1}{\PYGZsh{} f\PYGZus{}0 (x) = x\PYGZus{}0 * x\PYGZus{}\PYGZob{}n\PYGZhy{}1\PYGZcb{}}
    \PYG{n}{select\PYGZus{}r}\PYG{p}{[}\PYG{l+m+mi}{0}\PYG{p}{]} \PYG{o}{=} \PYG{k+kc}{True}
    \PYG{c+c1}{\PYGZsh{}}
    \PYG{c+c1}{\PYGZsh{} loop over forward and reverse mode}
    \PYG{k}{for} \PYG{n}{mode} \PYG{o+ow}{in} \PYG{n+nb}{range}\PYG{p}{(} \PYG{l+m+mi}{2} \PYG{p}{)} \PYG{p}{:}
        \PYG{n}{pat\PYGZus{}out} \PYG{o}{=} \PYG{n}{cppad\PYGZus{}py}\PYG{o}{.}\PYG{n}{sparse\PYGZus{}rc}\PYG{p}{(}\PYG{p}{)}
        \PYG{k}{if} \PYG{n}{mode} \PYG{o}{==} \PYG{l+m+mi}{0}  \PYG{p}{:}
            \PYG{n}{f}\PYG{o}{.}\PYG{n}{for\PYGZus{}hes\PYGZus{}sparsity}\PYG{p}{(}\PYG{n}{select\PYGZus{}d}\PYG{p}{,} \PYG{n}{select\PYGZus{}r}\PYG{p}{,} \PYG{n}{pat\PYGZus{}out}\PYG{p}{)}
        \PYG{c+c1}{\PYGZsh{}}
        \PYG{k}{if} \PYG{n}{mode} \PYG{o}{==} \PYG{l+m+mi}{1}  \PYG{p}{:}
            \PYG{n}{f}\PYG{o}{.}\PYG{n}{rev\PYGZus{}hes\PYGZus{}sparsity}\PYG{p}{(}\PYG{n}{select\PYGZus{}d}\PYG{p}{,} \PYG{n}{select\PYGZus{}r}\PYG{p}{,} \PYG{n}{pat\PYGZus{}out}\PYG{p}{)}
        \PYG{c+c1}{\PYGZsh{}}
        \PYG{c+c1}{\PYGZsh{}}
        \PYG{c+c1}{\PYGZsh{} check that result is sparsity pattern for Hessian of f\PYGZus{}0 (x)}
        \PYG{n}{ok} \PYG{o}{=} \PYG{n}{ok} \PYG{o+ow}{and} \PYG{n}{pat\PYGZus{}out}\PYG{o}{.}\PYG{n}{nnz}\PYG{p}{(}\PYG{p}{)} \PYG{o}{==} \PYG{l+m+mi}{2}
        \PYG{n}{row} \PYG{o}{=} \PYG{n}{pat\PYGZus{}out}\PYG{o}{.}\PYG{n}{row}\PYG{p}{(}\PYG{p}{)}
        \PYG{n}{col} \PYG{o}{=} \PYG{n}{pat\PYGZus{}out}\PYG{o}{.}\PYG{n}{col}\PYG{p}{(}\PYG{p}{)}
        \PYG{k}{for} \PYG{n}{k} \PYG{o+ow}{in} \PYG{n+nb}{range}\PYG{p}{(} \PYG{l+m+mi}{2} \PYG{p}{)} \PYG{p}{:}
            \PYG{n}{r} \PYG{o}{=} \PYG{n}{row}\PYG{p}{[}\PYG{n}{k}\PYG{p}{]}
            \PYG{n}{c} \PYG{o}{=} \PYG{n}{col}\PYG{p}{[}\PYG{n}{k}\PYG{p}{]}
            \PYG{k}{if} \PYG{n}{r} \PYG{o}{\PYGZlt{}}\PYG{o}{=} \PYG{n}{c}  \PYG{p}{:}
                \PYG{n}{ok} \PYG{o}{=} \PYG{n}{ok} \PYG{o+ow}{and} \PYG{n}{r} \PYG{o}{==} \PYG{l+m+mi}{0}
                \PYG{n}{ok} \PYG{o}{=} \PYG{n}{ok} \PYG{o+ow}{and} \PYG{n}{c} \PYG{o}{==} \PYG{n}{n}\PYG{o}{\PYGZhy{}}\PYG{l+m+mi}{1}
            \PYG{c+c1}{\PYGZsh{}}
            \PYG{k}{if} \PYG{n}{r} \PYG{o}{\PYGZgt{}}\PYG{o}{=} \PYG{n}{c}  \PYG{p}{:}
                \PYG{n}{ok} \PYG{o}{=} \PYG{n}{ok} \PYG{o+ow}{and} \PYG{n}{r} \PYG{o}{==} \PYG{n}{n}\PYG{o}{\PYGZhy{}}\PYG{l+m+mi}{1}
                \PYG{n}{ok} \PYG{o}{=} \PYG{n}{ok} \PYG{o+ow}{and} \PYG{n}{c} \PYG{o}{==} \PYG{l+m+mi}{0}
            \PYG{c+c1}{\PYGZsh{}}
        \PYG{c+c1}{\PYGZsh{}}
    \PYG{c+c1}{\PYGZsh{}}
    \PYG{c+c1}{\PYGZsh{}}
    \PYG{k}{return}\PYG{p}{(} \PYG{n}{ok} \PYG{p}{)}
\PYG{c+c1}{\PYGZsh{}}
\end{sphinxVerbatim}


\bigskip\hrule\bigskip


xsrst input file: \sphinxcode{\sphinxupquote{example/python/core/sparse\_hes\_pattern\_xam.py}}

\index{example@\spxentry{example}}\ignorespaces 

\subparagraph{Example}
\label{\detokenize{xsrst/py_hes_sparsity:example}}\label{\detokenize{xsrst/py_hes_sparsity:py-hes-sparsity-example}}\label{\detokenize{xsrst/py_hes_sparsity:index-9}}
{\hyperref[\detokenize{xsrst/sparse_hes_pattern_xam_py:id1}]{\sphinxcrossref{\DUrole{std,std-ref}{python}}}}


\bigskip\hrule\bigskip


xsrst input file: \sphinxcode{\sphinxupquote{lib/python/cppad\_py/hes\_sparsity.py}}


\subparagraph{py\_sparse\_jac}
\label{\detokenize{xsrst/py_sparse_jac:py-sparse-jac}}\label{\detokenize{xsrst/py_sparse_jac::doc}}
\index{py\_sparse\_jac@\spxentry{py\_sparse\_jac}}\index{computing@\spxentry{computing}}\index{sparse@\spxentry{sparse}}\index{jacobians@\spxentry{jacobians}}\ignorespaces 

\subparagraph{3.1.2.5 Computing Sparse Jacobians}
\label{\detokenize{xsrst/py_sparse_jac:computing-sparse-jacobians}}\label{\detokenize{xsrst/py_sparse_jac:id1}}\begin{itemize}
\item {} 
{\hyperref[\detokenize{xsrst/py_sparse_jac:py-sparse-jac-syntax}]{\sphinxcrossref{\DUrole{std,std-ref}{Syntax}}}}

\item {} 
{\hyperref[\detokenize{xsrst/py_sparse_jac:py-sparse-jac-purpose}]{\sphinxcrossref{\DUrole{std,std-ref}{Purpose}}}}

\item {} 
{\hyperref[\detokenize{xsrst/py_sparse_jac:py-sparse-jac-sparse-jac-for}]{\sphinxcrossref{\DUrole{std,std-ref}{sparse\_jac\_for}}}}

\item {} 
{\hyperref[\detokenize{xsrst/py_sparse_jac:py-sparse-jac-sparse-jac-rev}]{\sphinxcrossref{\DUrole{std,std-ref}{sparse\_jac\_rev}}}}

\item {} 
{\hyperref[\detokenize{xsrst/py_sparse_jac:py-sparse-jac-f}]{\sphinxcrossref{\DUrole{std,std-ref}{f}}}}

\item {} 
{\hyperref[\detokenize{xsrst/py_sparse_jac:py-sparse-jac-subset}]{\sphinxcrossref{\DUrole{std,std-ref}{subset}}}}

\item {} 
{\hyperref[\detokenize{xsrst/py_sparse_jac:py-sparse-jac-x}]{\sphinxcrossref{\DUrole{std,std-ref}{x}}}}

\item {} 
{\hyperref[\detokenize{xsrst/py_sparse_jac:py-sparse-jac-pattern}]{\sphinxcrossref{\DUrole{std,std-ref}{pattern}}}}

\item {} 
{\hyperref[\detokenize{xsrst/py_sparse_jac:py-sparse-jac-work}]{\sphinxcrossref{\DUrole{std,std-ref}{work}}}}

\item {} 
{\hyperref[\detokenize{xsrst/py_sparse_jac:py-sparse-jac-n-sweep}]{\sphinxcrossref{\DUrole{std,std-ref}{n\_sweep}}}}

\item {} 
{\hyperref[\detokenize{xsrst/py_sparse_jac:py-sparse-jac-uses-forward}]{\sphinxcrossref{\DUrole{std,std-ref}{Uses Forward}}}}

\item {} 
{\hyperref[\detokenize{xsrst/py_sparse_jac:py-sparse-jac-example}]{\sphinxcrossref{\DUrole{std,std-ref}{Example}}}}

\end{itemize}

\index{syntax@\spxentry{syntax}}\ignorespaces 

\subparagraph{Syntax}
\label{\detokenize{xsrst/py_sparse_jac:syntax}}\label{\detokenize{xsrst/py_sparse_jac:py-sparse-jac-syntax}}\label{\detokenize{xsrst/py_sparse_jac:index-1}}
\begin{DUlineblock}{0em}
\item[] \sphinxstyleemphasis{work} =  \sphinxcode{\sphinxupquote{cppad\_py.sparse\_jac\_work}} ()
\item[] \sphinxstyleemphasis{n\_sweep} = \sphinxstyleemphasis{f} . \sphinxcode{\sphinxupquote{sparse\_jac\_for}} ( \sphinxstyleemphasis{subset} , \sphinxstyleemphasis{x} , \sphinxstyleemphasis{pattern} , \sphinxstyleemphasis{work} )
\end{DUlineblock}

\sphinxstyleemphasis{n\_sweep} = \sphinxstyleemphasis{f} . \sphinxcode{\sphinxupquote{sparse\_jac\_rev}} ( \sphinxstyleemphasis{subset} , \sphinxstyleemphasis{x} , \sphinxstyleemphasis{pattern} , \sphinxstyleemphasis{work} )

\index{purpose@\spxentry{purpose}}\ignorespaces 

\subparagraph{Purpose}
\label{\detokenize{xsrst/py_sparse_jac:purpose}}\label{\detokenize{xsrst/py_sparse_jac:py-sparse-jac-purpose}}\label{\detokenize{xsrst/py_sparse_jac:index-2}}
We use \(F : \B{R}^n \rightarrow \B{R}^m\) to denote the
function corresponding to \sphinxstyleemphasis{f} .
The syntax above takes advantage of sparsity when computing the Jacobian
\begin{equation*}
\begin{split}J(x) = F^{(1)} (x)\end{split}
\end{equation*}
In the sparse case, this should be faster and take less memory than
{\hyperref[\detokenize{xsrst/py_fun_jacobian:id1}]{\sphinxcrossref{\DUrole{std,std-ref}{py\_fun\_jacobian}}}}.
We use the notation \(J_{i,j} (x)\) to denote the partial of
\(F_i (x)\) with respect to \(x_j\).

\index{sparse\_jac\_for@\spxentry{sparse\_jac\_for}}\ignorespaces 

\subparagraph{sparse\_jac\_for}
\label{\detokenize{xsrst/py_sparse_jac:sparse-jac-for}}\label{\detokenize{xsrst/py_sparse_jac:py-sparse-jac-sparse-jac-for}}\label{\detokenize{xsrst/py_sparse_jac:index-3}}
This function uses first order forward mode sweeps {\hyperref[\detokenize{xsrst/py_fun_forward:id1}]{\sphinxcrossref{\DUrole{std,std-ref}{py\_fun\_forward}}}}
to compute multiple columns of the Jacobian at the same time.

\index{sparse\_jac\_rev@\spxentry{sparse\_jac\_rev}}\ignorespaces 

\subparagraph{sparse\_jac\_rev}
\label{\detokenize{xsrst/py_sparse_jac:sparse-jac-rev}}\label{\detokenize{xsrst/py_sparse_jac:py-sparse-jac-sparse-jac-rev}}\label{\detokenize{xsrst/py_sparse_jac:index-4}}
This function uses first order reverse mode sweeps {\hyperref[\detokenize{xsrst/py_fun_reverse:id1}]{\sphinxcrossref{\DUrole{std,std-ref}{py\_fun\_reverse}}}}
to compute multiple rows of the Jacobian at the same time.

\index{f@\spxentry{f}}\ignorespaces 

\subparagraph{f}
\label{\detokenize{xsrst/py_sparse_jac:f}}\label{\detokenize{xsrst/py_sparse_jac:py-sparse-jac-f}}\label{\detokenize{xsrst/py_sparse_jac:index-5}}
This object must have been returned by a previous call to the python
{\hyperref[\detokenize{xsrst/py_fun_ctor:id1}]{\sphinxcrossref{\DUrole{std,std-ref}{d\_fun}}}} constructor.
Note that the Taylor coefficients stored in \sphinxstyleemphasis{f} are affected
by this operation; see
{\hyperref[\detokenize{xsrst/py_sparse_jac:py-sparse-jac-uses-forward}]{\sphinxcrossref{\DUrole{std,std-ref}{uses\_forward}}}} below.

\index{subset@\spxentry{subset}}\ignorespaces 

\subparagraph{subset}
\label{\detokenize{xsrst/py_sparse_jac:subset}}\label{\detokenize{xsrst/py_sparse_jac:py-sparse-jac-subset}}\label{\detokenize{xsrst/py_sparse_jac:index-6}}
This argument must have be a {\hyperref[\detokenize{xsrst/py_sparse_rcv:py-sparse-rcv-matrix}]{\sphinxcrossref{\DUrole{std,std-ref}{matrix}}}}
returned by the \sphinxcode{\sphinxupquote{sparse\_rcv}} constructor.
Its row size is \sphinxstyleemphasis{subset} . \sphinxcode{\sphinxupquote{nr}} () == \sphinxstyleemphasis{m} ,
and its column size is \sphinxstyleemphasis{subset} . \sphinxcode{\sphinxupquote{nc}} () == \sphinxstyleemphasis{n} .
It specifies which elements of the Jacobian are computed.
The input value of its value vector
\sphinxstyleemphasis{subset} . \sphinxcode{\sphinxupquote{val}} () does not matter.
Upon return it contains the value of the corresponding elements
of the Jacobian.
All of the row, column pairs in \sphinxstyleemphasis{subset} must also appear in
\sphinxstyleemphasis{pattern} ; i.e., they must be possibly non\sphinxhyphen{}zero.

\index{x@\spxentry{x}}\ignorespaces 

\subparagraph{x}
\label{\detokenize{xsrst/py_sparse_jac:x}}\label{\detokenize{xsrst/py_sparse_jac:py-sparse-jac-x}}\label{\detokenize{xsrst/py_sparse_jac:index-7}}
This argument is a numpy vector with \sphinxcode{\sphinxupquote{float}} elements
and size \sphinxstyleemphasis{n} .
It specifies the point at which to evaluate the Jacobian \(J(x)\).

\index{pattern@\spxentry{pattern}}\ignorespaces 

\subparagraph{pattern}
\label{\detokenize{xsrst/py_sparse_jac:pattern}}\label{\detokenize{xsrst/py_sparse_jac:py-sparse-jac-pattern}}\label{\detokenize{xsrst/py_sparse_jac:index-8}}
This argument must have be a {\hyperref[\detokenize{xsrst/py_sparse_rc:py-sparse-rc-pattern}]{\sphinxcrossref{\DUrole{std,std-ref}{pattern}}}}
returned by the \sphinxcode{\sphinxupquote{sparse\_rc}} constructor.
Its row size is \sphinxstyleemphasis{pattern} . \sphinxcode{\sphinxupquote{nr}} () == \sphinxstyleemphasis{m} ,
and its column size is \sphinxstyleemphasis{pattern} . \sphinxcode{\sphinxupquote{nc}} () == \sphinxstyleemphasis{n} .
It is a sparsity pattern for the Jacobian \(J(x)\).
This argument is not used (and need not satisfy any conditions),
when {\hyperref[\detokenize{xsrst/py_sparse_jac:py-sparse-jac-work}]{\sphinxcrossref{\DUrole{std,std-ref}{work}}}} is non\sphinxhyphen{}empty.

\index{work@\spxentry{work}}\ignorespaces 

\subparagraph{work}
\label{\detokenize{xsrst/py_sparse_jac:work}}\label{\detokenize{xsrst/py_sparse_jac:py-sparse-jac-work}}\label{\detokenize{xsrst/py_sparse_jac:index-9}}
This argument must have been constructed by the call

\begin{DUlineblock}{0em}
\item[]         \sphinxstyleemphasis{work} =  \sphinxcode{\sphinxupquote{cppad\_py.sparse\_jac\_work}} ()
\end{DUlineblock}

We refer to its initial value,
and its value after \sphinxstyleemphasis{work} . \sphinxcode{\sphinxupquote{clear}} () , as empty.
If it is empty, information is stored in \sphinxstyleemphasis{work} .
This can be used to reduce computation when
a future call is for the same object \sphinxstyleemphasis{f} ,
the same member function \sphinxcode{\sphinxupquote{sparse\_jac\_for}} or \sphinxcode{\sphinxupquote{sparse\_jac\_rev}} ,
and the same subset of the Jacobian.
If any of these values change, use \sphinxstyleemphasis{work} . \sphinxcode{\sphinxupquote{clear}} () to
empty this structure.

\index{n\_sweep@\spxentry{n\_sweep}}\ignorespaces 

\subparagraph{n\_sweep}
\label{\detokenize{xsrst/py_sparse_jac:n-sweep}}\label{\detokenize{xsrst/py_sparse_jac:py-sparse-jac-n-sweep}}\label{\detokenize{xsrst/py_sparse_jac:index-10}}
This return value is and \sphinxcode{\sphinxupquote{int}} .
If \sphinxcode{\sphinxupquote{sparse\_jac\_for}} ( \sphinxcode{\sphinxupquote{sparse\_jac\_rev}} ) is used,
\sphinxstyleemphasis{n\_sweep} is the number of first order forward (reverse) sweeps
used to compute the requested Jacobian values.
This is proportional to the total computational work,
not counting the zero order forward sweep,
or combining multiple columns (rows) into a single sweep.

\index{uses@\spxentry{uses}}\index{forward@\spxentry{forward}}\ignorespaces 

\subparagraph{Uses Forward}
\label{\detokenize{xsrst/py_sparse_jac:uses-forward}}\label{\detokenize{xsrst/py_sparse_jac:py-sparse-jac-uses-forward}}\label{\detokenize{xsrst/py_sparse_jac:index-11}}
After each call to {\hyperref[\detokenize{xsrst/py_fun_forward:id1}]{\sphinxcrossref{\DUrole{std,std-ref}{py\_fun\_forward}}}},
the object \sphinxstyleemphasis{f} contains the corresponding Taylor coefficients
for all the variables in the operation sequence..
After a call to \sphinxcode{\sphinxupquote{sparse\_jac\_forward}} or \sphinxcode{\sphinxupquote{sparse\_jac\_rev}} ,
the zero order coefficients correspond to

\begin{DUlineblock}{0em}
\item[]         \sphinxstyleemphasis{f} . \sphinxcode{\sphinxupquote{forward(0}} , \sphinxstyleemphasis{x} )
\end{DUlineblock}

All the other forward mode coefficients are unspecified.


\subparagraph{sparse\_jac\_xam\_py}
\label{\detokenize{xsrst/sparse_jac_xam_py:sparse-jac-xam-py}}\label{\detokenize{xsrst/sparse_jac_xam_py::doc}}
\index{sparse\_jac\_xam\_py@\spxentry{sparse\_jac\_xam\_py}}\index{python:@\spxentry{python:}}\index{computing@\spxentry{computing}}\index{sparse@\spxentry{sparse}}\index{jacobians:@\spxentry{jacobians:}}\index{example@\spxentry{example}}\index{test@\spxentry{test}}\ignorespaces 

\subparagraph{3.1.2.5.1 Python: Computing Sparse Jacobians: Example and Test}
\label{\detokenize{xsrst/sparse_jac_xam_py:python-computing-sparse-jacobians-example-and-test}}\label{\detokenize{xsrst/sparse_jac_xam_py:id1}}
\begin{sphinxVerbatim}[commandchars=\\\{\}]
\PYG{k}{def} \PYG{n+nf}{sparse\PYGZus{}jac\PYGZus{}xam}\PYG{p}{(}\PYG{p}{)} \PYG{p}{:}
    \PYG{c+c1}{\PYGZsh{}}
    \PYG{k+kn}{import} \PYG{n+nn}{numpy}
    \PYG{k+kn}{import} \PYG{n+nn}{cppad\PYGZus{}py}
    \PYG{c+c1}{\PYGZsh{}}
    \PYG{c+c1}{\PYGZsh{} initialize return variable}
    \PYG{n}{ok} \PYG{o}{=} \PYG{k+kc}{True}
    \PYG{c+c1}{\PYGZsh{} \PYGZhy{}\PYGZhy{}\PYGZhy{}\PYGZhy{}\PYGZhy{}\PYGZhy{}\PYGZhy{}\PYGZhy{}\PYGZhy{}\PYGZhy{}\PYGZhy{}\PYGZhy{}\PYGZhy{}\PYGZhy{}\PYGZhy{}\PYGZhy{}\PYGZhy{}\PYGZhy{}\PYGZhy{}\PYGZhy{}\PYGZhy{}\PYGZhy{}\PYGZhy{}\PYGZhy{}\PYGZhy{}\PYGZhy{}\PYGZhy{}\PYGZhy{}\PYGZhy{}\PYGZhy{}\PYGZhy{}\PYGZhy{}\PYGZhy{}\PYGZhy{}\PYGZhy{}\PYGZhy{}\PYGZhy{}\PYGZhy{}\PYGZhy{}\PYGZhy{}\PYGZhy{}\PYGZhy{}\PYGZhy{}\PYGZhy{}\PYGZhy{}\PYGZhy{}\PYGZhy{}\PYGZhy{}\PYGZhy{}\PYGZhy{}\PYGZhy{}\PYGZhy{}\PYGZhy{}\PYGZhy{}\PYGZhy{}\PYGZhy{}\PYGZhy{}\PYGZhy{}\PYGZhy{}\PYGZhy{}\PYGZhy{}\PYGZhy{}\PYGZhy{}\PYGZhy{}\PYGZhy{}\PYGZhy{}\PYGZhy{}\PYGZhy{}\PYGZhy{}}
    \PYG{c+c1}{\PYGZsh{} number of dependent and independent variables}
    \PYG{n}{n} \PYG{o}{=} \PYG{l+m+mi}{3}
    \PYG{c+c1}{\PYGZsh{} one}
    \PYG{n}{aone} \PYG{o}{=} \PYG{n}{cppad\PYGZus{}py}\PYG{o}{.}\PYG{n}{a\PYGZus{}double}\PYG{p}{(}\PYG{l+m+mf}{1.0}\PYG{p}{)}
    \PYG{c+c1}{\PYGZsh{}}
    \PYG{c+c1}{\PYGZsh{} create the independent variables ax}
    \PYG{n}{x} \PYG{o}{=} \PYG{n}{numpy}\PYG{o}{.}\PYG{n}{empty}\PYG{p}{(}\PYG{n}{n}\PYG{p}{,} \PYG{n}{dtype}\PYG{o}{=}\PYG{n+nb}{float}\PYG{p}{)}
    \PYG{k}{for} \PYG{n}{i} \PYG{o+ow}{in} \PYG{n+nb}{range}\PYG{p}{(} \PYG{n}{n}  \PYG{p}{)} \PYG{p}{:}
        \PYG{n}{x}\PYG{p}{[}\PYG{n}{i}\PYG{p}{]} \PYG{o}{=} \PYG{n}{i} \PYG{o}{+} \PYG{l+m+mf}{2.0}
    \PYG{c+c1}{\PYGZsh{}}
    \PYG{n}{ax} \PYG{o}{=} \PYG{n}{cppad\PYGZus{}py}\PYG{o}{.}\PYG{n}{independent}\PYG{p}{(}\PYG{n}{x}\PYG{p}{)}
    \PYG{c+c1}{\PYGZsh{}}
    \PYG{c+c1}{\PYGZsh{} create dependent variables ay with ay[i] = (j+1) * ax[j]}
    \PYG{c+c1}{\PYGZsh{} where i = mod(j + 1, n)}
    \PYG{n}{ay} \PYG{o}{=} \PYG{n}{numpy}\PYG{o}{.}\PYG{n}{empty}\PYG{p}{(}\PYG{n}{n}\PYG{p}{,} \PYG{n}{dtype}\PYG{o}{=}\PYG{n}{cppad\PYGZus{}py}\PYG{o}{.}\PYG{n}{a\PYGZus{}double}\PYG{p}{)}
    \PYG{k}{for} \PYG{n}{j} \PYG{o+ow}{in} \PYG{n+nb}{range}\PYG{p}{(} \PYG{n}{n}  \PYG{p}{)} \PYG{p}{:}
        \PYG{n}{i} \PYG{o}{=} \PYG{n}{j}\PYG{o}{+}\PYG{l+m+mi}{1}
        \PYG{k}{if} \PYG{n}{i} \PYG{o}{\PYGZgt{}}\PYG{o}{=} \PYG{n}{n}  \PYG{p}{:}
            \PYG{n}{i} \PYG{o}{=} \PYG{n}{i} \PYG{o}{\PYGZhy{}} \PYG{n}{n}
        \PYG{c+c1}{\PYGZsh{}}
        \PYG{n}{aj} \PYG{o}{=} \PYG{n}{cppad\PYGZus{}py}\PYG{o}{.}\PYG{n}{a\PYGZus{}double}\PYG{p}{(}\PYG{n}{j}\PYG{p}{)}
        \PYG{n}{ay\PYGZus{}i} \PYG{o}{=} \PYG{p}{(}\PYG{n}{aj} \PYG{o}{+} \PYG{n}{aone}\PYG{p}{)} \PYG{o}{*} \PYG{n}{ax}\PYG{p}{[}\PYG{n}{j}\PYG{p}{]}
        \PYG{n}{ay}\PYG{p}{[}\PYG{n}{i}\PYG{p}{]} \PYG{o}{=} \PYG{n}{ay\PYGZus{}i}
    \PYG{c+c1}{\PYGZsh{}}
    \PYG{c+c1}{\PYGZsh{}}
    \PYG{c+c1}{\PYGZsh{} define af corresponding to f(x)}
    \PYG{n}{f}  \PYG{o}{=} \PYG{n}{cppad\PYGZus{}py}\PYG{o}{.}\PYG{n}{d\PYGZus{}fun}\PYG{p}{(}\PYG{n}{ax}\PYG{p}{,} \PYG{n}{ay}\PYG{p}{)}
    \PYG{c+c1}{\PYGZsh{}}
    \PYG{c+c1}{\PYGZsh{} sparsity pattern for identity matrix}
    \PYG{n}{pat\PYGZus{}eye} \PYG{o}{=} \PYG{n}{cppad\PYGZus{}py}\PYG{o}{.}\PYG{n}{sparse\PYGZus{}rc}\PYG{p}{(}\PYG{p}{)}
    \PYG{n}{pat\PYGZus{}eye}\PYG{o}{.}\PYG{n}{resize}\PYG{p}{(}\PYG{n}{n}\PYG{p}{,} \PYG{n}{n}\PYG{p}{,} \PYG{n}{n}\PYG{p}{)}
    \PYG{k}{for} \PYG{n}{k} \PYG{o+ow}{in} \PYG{n+nb}{range}\PYG{p}{(} \PYG{n}{n} \PYG{p}{)} \PYG{p}{:}
        \PYG{n}{pat\PYGZus{}eye}\PYG{o}{.}\PYG{n}{put}\PYG{p}{(}\PYG{n}{k}\PYG{p}{,} \PYG{n}{k}\PYG{p}{,} \PYG{n}{k}\PYG{p}{)}
    \PYG{c+c1}{\PYGZsh{}}
    \PYG{c+c1}{\PYGZsh{}}
    \PYG{c+c1}{\PYGZsh{} sparsity pattern for the Jacobian}
    \PYG{n}{pat\PYGZus{}jac} \PYG{o}{=} \PYG{n}{cppad\PYGZus{}py}\PYG{o}{.}\PYG{n}{sparse\PYGZus{}rc}\PYG{p}{(}\PYG{p}{)}
    \PYG{n}{f}\PYG{o}{.}\PYG{n}{for\PYGZus{}jac\PYGZus{}sparsity}\PYG{p}{(}\PYG{n}{pat\PYGZus{}eye}\PYG{p}{,} \PYG{n}{pat\PYGZus{}jac}\PYG{p}{)}
    \PYG{c+c1}{\PYGZsh{}}
    \PYG{c+c1}{\PYGZsh{} loop over forward and reverse mode}
    \PYG{k}{for} \PYG{n}{mode} \PYG{o+ow}{in} \PYG{n+nb}{range}\PYG{p}{(} \PYG{l+m+mi}{2} \PYG{p}{)} \PYG{p}{:}
        \PYG{c+c1}{\PYGZsh{} compute all possibly non\PYGZhy{}zero entries in Jacobian}
        \PYG{n}{subset} \PYG{o}{=} \PYG{n}{cppad\PYGZus{}py}\PYG{o}{.}\PYG{n}{sparse\PYGZus{}rcv}\PYG{p}{(}\PYG{n}{pat\PYGZus{}jac}\PYG{p}{)}
        \PYG{c+c1}{\PYGZsh{} work space used to save time for multiple calls}
        \PYG{n}{work} \PYG{o}{=} \PYG{n}{cppad\PYGZus{}py}\PYG{o}{.}\PYG{n}{sparse\PYGZus{}jac\PYGZus{}work}\PYG{p}{(}\PYG{p}{)}
        \PYG{k}{if} \PYG{n}{mode} \PYG{o}{==} \PYG{l+m+mi}{0}  \PYG{p}{:}
            \PYG{n}{f}\PYG{o}{.}\PYG{n}{sparse\PYGZus{}jac\PYGZus{}for}\PYG{p}{(}\PYG{n}{subset}\PYG{p}{,} \PYG{n}{x}\PYG{p}{,} \PYG{n}{pat\PYGZus{}jac}\PYG{p}{,} \PYG{n}{work}\PYG{p}{)}
        \PYG{c+c1}{\PYGZsh{}}
        \PYG{k}{if} \PYG{n}{mode} \PYG{o}{==} \PYG{l+m+mi}{1}  \PYG{p}{:}
            \PYG{n}{f}\PYG{o}{.}\PYG{n}{sparse\PYGZus{}jac\PYGZus{}rev}\PYG{p}{(}\PYG{n}{subset}\PYG{p}{,} \PYG{n}{x}\PYG{p}{,} \PYG{n}{pat\PYGZus{}jac}\PYG{p}{,} \PYG{n}{work}\PYG{p}{)}
        \PYG{c+c1}{\PYGZsh{}}
        \PYG{c+c1}{\PYGZsh{}}
        \PYG{c+c1}{\PYGZsh{} check result}
        \PYG{n}{ok} \PYG{o}{=} \PYG{n}{ok} \PYG{o+ow}{and} \PYG{n}{n} \PYG{o}{==} \PYG{n}{subset}\PYG{o}{.}\PYG{n}{nnz}\PYG{p}{(}\PYG{p}{)}
        \PYG{n}{col\PYGZus{}major} \PYG{o}{=} \PYG{n}{subset}\PYG{o}{.}\PYG{n}{col\PYGZus{}major}\PYG{p}{(}\PYG{p}{)}
        \PYG{n}{row} \PYG{o}{=} \PYG{n}{subset}\PYG{o}{.}\PYG{n}{row}\PYG{p}{(}\PYG{p}{)}
        \PYG{n}{col} \PYG{o}{=} \PYG{n}{subset}\PYG{o}{.}\PYG{n}{col}\PYG{p}{(}\PYG{p}{)}
        \PYG{n}{val} \PYG{o}{=} \PYG{n}{subset}\PYG{o}{.}\PYG{n}{val}\PYG{p}{(}\PYG{p}{)}
        \PYG{k}{for} \PYG{n}{k} \PYG{o+ow}{in} \PYG{n+nb}{range}\PYG{p}{(} \PYG{n}{n} \PYG{p}{)} \PYG{p}{:}
            \PYG{n}{ell} \PYG{o}{=} \PYG{n}{col\PYGZus{}major}\PYG{p}{[}\PYG{n}{k}\PYG{p}{]}
            \PYG{n}{r} \PYG{o}{=} \PYG{n}{row}\PYG{p}{[}\PYG{n}{ell}\PYG{p}{]}
            \PYG{n}{c} \PYG{o}{=} \PYG{n}{col}\PYG{p}{[}\PYG{n}{ell}\PYG{p}{]}
            \PYG{n}{v} \PYG{o}{=} \PYG{n}{val}\PYG{p}{[}\PYG{n}{ell}\PYG{p}{]}
            \PYG{n}{i} \PYG{o}{=} \PYG{n}{c}\PYG{o}{+}\PYG{l+m+mi}{1}
            \PYG{k}{if} \PYG{n}{i} \PYG{o}{\PYGZgt{}}\PYG{o}{=}  \PYG{n}{n}  \PYG{p}{:}
                \PYG{n}{i} \PYG{o}{=} \PYG{n}{i} \PYG{o}{\PYGZhy{}} \PYG{n}{n}
            \PYG{c+c1}{\PYGZsh{}}
            \PYG{n}{ok} \PYG{o}{=} \PYG{n}{ok} \PYG{o+ow}{and} \PYG{n}{c} \PYG{o}{==} \PYG{n}{k}
            \PYG{n}{ok} \PYG{o}{=} \PYG{n}{ok} \PYG{o+ow}{and} \PYG{n}{r} \PYG{o}{==} \PYG{n}{i}
            \PYG{n}{ok} \PYG{o}{=} \PYG{n}{ok} \PYG{o+ow}{and} \PYG{n}{v} \PYG{o}{==} \PYG{n}{c} \PYG{o}{+} \PYG{l+m+mf}{1.0}
        \PYG{c+c1}{\PYGZsh{}}
    \PYG{c+c1}{\PYGZsh{}}
    \PYG{c+c1}{\PYGZsh{}}
    \PYG{k}{return}\PYG{p}{(} \PYG{n}{ok} \PYG{p}{)}
\PYG{c+c1}{\PYGZsh{}}
\end{sphinxVerbatim}


\bigskip\hrule\bigskip


xsrst input file: \sphinxcode{\sphinxupquote{example/python/core/sparse\_jac\_xam.py}}

\index{example@\spxentry{example}}\ignorespaces 

\subparagraph{Example}
\label{\detokenize{xsrst/py_sparse_jac:example}}\label{\detokenize{xsrst/py_sparse_jac:py-sparse-jac-example}}\label{\detokenize{xsrst/py_sparse_jac:index-12}}
{\hyperref[\detokenize{xsrst/sparse_jac_xam_py:id1}]{\sphinxcrossref{\DUrole{std,std-ref}{sparse\_jac\_xam\_py}}}}


\bigskip\hrule\bigskip


xsrst input file: \sphinxcode{\sphinxupquote{lib/python/cppad\_py/sparse\_jac.py}}


\subparagraph{py\_sparse\_hes}
\label{\detokenize{xsrst/py_sparse_hes:py-sparse-hes}}\label{\detokenize{xsrst/py_sparse_hes::doc}}
\index{py\_sparse\_hes@\spxentry{py\_sparse\_hes}}\index{computing@\spxentry{computing}}\index{sparse@\spxentry{sparse}}\index{hessians@\spxentry{hessians}}\ignorespaces 

\subparagraph{3.1.2.6 Computing Sparse Hessians}
\label{\detokenize{xsrst/py_sparse_hes:computing-sparse-hessians}}\label{\detokenize{xsrst/py_sparse_hes:id1}}\begin{itemize}
\item {} 
{\hyperref[\detokenize{xsrst/py_sparse_hes:py-sparse-hes-syntax}]{\sphinxcrossref{\DUrole{std,std-ref}{Syntax}}}}

\item {} 
{\hyperref[\detokenize{xsrst/py_sparse_hes:py-sparse-hes-purpose}]{\sphinxcrossref{\DUrole{std,std-ref}{Purpose}}}}

\item {} 
{\hyperref[\detokenize{xsrst/py_sparse_hes:py-sparse-hes-f}]{\sphinxcrossref{\DUrole{std,std-ref}{f}}}}

\item {} 
{\hyperref[\detokenize{xsrst/py_sparse_hes:py-sparse-hes-subset}]{\sphinxcrossref{\DUrole{std,std-ref}{subset}}}}

\item {} 
{\hyperref[\detokenize{xsrst/py_sparse_hes:py-sparse-hes-x}]{\sphinxcrossref{\DUrole{std,std-ref}{x}}}}

\item {} 
{\hyperref[\detokenize{xsrst/py_sparse_hes:py-sparse-hes-r}]{\sphinxcrossref{\DUrole{std,std-ref}{r}}}}

\item {} 
{\hyperref[\detokenize{xsrst/py_sparse_hes:py-sparse-hes-pattern}]{\sphinxcrossref{\DUrole{std,std-ref}{pattern}}}}

\item {} 
{\hyperref[\detokenize{xsrst/py_sparse_hes:py-sparse-hes-work}]{\sphinxcrossref{\DUrole{std,std-ref}{work}}}}

\item {} 
{\hyperref[\detokenize{xsrst/py_sparse_hes:py-sparse-hes-n-sweep}]{\sphinxcrossref{\DUrole{std,std-ref}{n\_sweep}}}}

\item {} 
{\hyperref[\detokenize{xsrst/py_sparse_hes:py-sparse-hes-uses-forward}]{\sphinxcrossref{\DUrole{std,std-ref}{Uses Forward}}}}

\item {} 
{\hyperref[\detokenize{xsrst/py_sparse_hes:py-sparse-hes-example}]{\sphinxcrossref{\DUrole{std,std-ref}{Example}}}}

\end{itemize}

\index{syntax@\spxentry{syntax}}\ignorespaces 

\subparagraph{Syntax}
\label{\detokenize{xsrst/py_sparse_hes:syntax}}\label{\detokenize{xsrst/py_sparse_hes:py-sparse-hes-syntax}}\label{\detokenize{xsrst/py_sparse_hes:index-1}}
\begin{DUlineblock}{0em}
\item[] \sphinxstyleemphasis{work} =  \sphinxcode{\sphinxupquote{cppad\_py.sparse\_hes\_work}} ()
\item[] \sphinxstyleemphasis{n\_sweep} = \sphinxstyleemphasis{f} . \sphinxcode{\sphinxupquote{sparse\_hes}} ( \sphinxstyleemphasis{subset} , \sphinxstyleemphasis{x} , \sphinxstyleemphasis{r} , \sphinxstyleemphasis{pattern} , \sphinxstyleemphasis{work} )
\end{DUlineblock}

\index{purpose@\spxentry{purpose}}\ignorespaces 

\subparagraph{Purpose}
\label{\detokenize{xsrst/py_sparse_hes:purpose}}\label{\detokenize{xsrst/py_sparse_hes:py-sparse-hes-purpose}}\label{\detokenize{xsrst/py_sparse_hes:index-2}}
We use \(F : \B{R}^n \rightarrow \B{R}^m\) to denote the
function corresponding to \sphinxstyleemphasis{f} .
Given a vector \(r \in \B{R}^m\), define
\begin{equation*}
\begin{split}H(x) = (r^\R{T} F)^{(2)} ( x )\end{split}
\end{equation*}
This routine takes advantage of sparsity when computing elements
of the Hessian \(H(x)\).

\index{f@\spxentry{f}}\ignorespaces 

\subparagraph{f}
\label{\detokenize{xsrst/py_sparse_hes:f}}\label{\detokenize{xsrst/py_sparse_hes:py-sparse-hes-f}}\label{\detokenize{xsrst/py_sparse_hes:index-3}}
This object must have been returned by a previous call to the python
{\hyperref[\detokenize{xsrst/py_fun_ctor:id1}]{\sphinxcrossref{\DUrole{std,std-ref}{d\_fun}}}} constructor.
Note that the Taylor coefficients stored in \sphinxstyleemphasis{f} are affected
by this operation; see
{\hyperref[\detokenize{xsrst/py_sparse_hes:py-sparse-hes-uses-forward}]{\sphinxcrossref{\DUrole{std,std-ref}{uses\_forward}}}} below.

\index{subset@\spxentry{subset}}\ignorespaces 

\subparagraph{subset}
\label{\detokenize{xsrst/py_sparse_hes:subset}}\label{\detokenize{xsrst/py_sparse_hes:py-sparse-hes-subset}}\label{\detokenize{xsrst/py_sparse_hes:index-4}}
This argument must have be a {\hyperref[\detokenize{xsrst/py_sparse_rcv:py-sparse-rcv-matrix}]{\sphinxcrossref{\DUrole{std,std-ref}{matrix}}}}
returned by the \sphinxcode{\sphinxupquote{sparse\_rcv}} constructor.
Its row size and column size is \sphinxstyleemphasis{n} ; i.e.,
\sphinxstyleemphasis{subset} . \sphinxcode{\sphinxupquote{nr}} () == \sphinxstyleemphasis{n} and \sphinxstyleemphasis{subset} . \sphinxcode{\sphinxupquote{nc}} () == \sphinxstyleemphasis{n} .
It specifies which elements of the Hessian are computed.
The input value of its value vector
\sphinxstyleemphasis{subset} . \sphinxcode{\sphinxupquote{val}} () does not matter.
Upon return it contains the value of the corresponding elements
of the Jacobian.
All of the row, column pairs in \sphinxstyleemphasis{subset} must also appear in
\sphinxstyleemphasis{pattern} ; i.e., they must be possibly non\sphinxhyphen{}zero.

\index{x@\spxentry{x}}\ignorespaces 

\subparagraph{x}
\label{\detokenize{xsrst/py_sparse_hes:x}}\label{\detokenize{xsrst/py_sparse_hes:py-sparse-hes-x}}\label{\detokenize{xsrst/py_sparse_hes:index-5}}
This argument is a numpy vector with \sphinxcode{\sphinxupquote{float}} elements
and size \sphinxstyleemphasis{n} .
It specifies the point at which to evaluate the Hessian \(H(x)\).

\index{r@\spxentry{r}}\ignorespaces 

\subparagraph{r}
\label{\detokenize{xsrst/py_sparse_hes:r}}\label{\detokenize{xsrst/py_sparse_hes:py-sparse-hes-r}}\label{\detokenize{xsrst/py_sparse_hes:index-6}}
This argument is a numpy vector with \sphinxcode{\sphinxupquote{float}} elements
and size \sphinxstyleemphasis{m} .
It specifies the multiplier for each component of \(F(x)\);
i.e., \(r_i\) is the multiplier for \(F_i (x)\).

\index{pattern@\spxentry{pattern}}\ignorespaces 

\subparagraph{pattern}
\label{\detokenize{xsrst/py_sparse_hes:pattern}}\label{\detokenize{xsrst/py_sparse_hes:py-sparse-hes-pattern}}\label{\detokenize{xsrst/py_sparse_hes:index-7}}
This argument must have be a {\hyperref[\detokenize{xsrst/py_sparse_rc:py-sparse-rc-pattern}]{\sphinxcrossref{\DUrole{std,std-ref}{pattern}}}}
returned by the \sphinxcode{\sphinxupquote{sparse\_rc}} constructor.
Its row size and column sizes are \sphinxstyleemphasis{n} ; i.e.,
\sphinxstyleemphasis{pattern} . \sphinxcode{\sphinxupquote{nr}} () == \sphinxstyleemphasis{n} and \sphinxstyleemphasis{pattern} . \sphinxcode{\sphinxupquote{nc}} () == \sphinxstyleemphasis{n} .
It is a sparsity pattern for the Hessian \(H(x)\).
This argument is not used (and need not satisfy any conditions),
when {\hyperref[\detokenize{xsrst/py_sparse_hes:py-sparse-hes-work}]{\sphinxcrossref{\DUrole{std,std-ref}{work}}}} is non\sphinxhyphen{}empty.

\index{work@\spxentry{work}}\ignorespaces 

\subparagraph{work}
\label{\detokenize{xsrst/py_sparse_hes:work}}\label{\detokenize{xsrst/py_sparse_hes:py-sparse-hes-work}}\label{\detokenize{xsrst/py_sparse_hes:index-8}}
This argument must have been constructed by the call

\begin{DUlineblock}{0em}
\item[]         \sphinxstyleemphasis{work} =  \sphinxcode{\sphinxupquote{cppad\_py.sparse\_hes\_work}} ()
\end{DUlineblock}

We refer to its initial value,
and its value after \sphinxstyleemphasis{work} . \sphinxcode{\sphinxupquote{clear}} () , as empty.
If it is empty, information is stored in \sphinxstyleemphasis{work} .
This can be used to reduce computation when
a future call is for the same object \sphinxstyleemphasis{f} ,
and the same subset of the Hessian.
If either of these values change, use \sphinxstyleemphasis{work} . \sphinxcode{\sphinxupquote{clear}} () to
empty this structure.

\index{n\_sweep@\spxentry{n\_sweep}}\ignorespaces 

\subparagraph{n\_sweep}
\label{\detokenize{xsrst/py_sparse_hes:n-sweep}}\label{\detokenize{xsrst/py_sparse_hes:py-sparse-hes-n-sweep}}\label{\detokenize{xsrst/py_sparse_hes:index-9}}
The return value \sphinxstyleemphasis{n\_sweep} has prototype

\begin{DUlineblock}{0em}
\item[]         \sphinxcode{\sphinxupquote{int}} \sphinxstyleemphasis{n\_sweep}
\end{DUlineblock}

It is the number of first order forward sweeps
used to compute the requested Hessian values.
Each first forward sweep is followed by a second order reverse sweep
so it is also the number of reverse sweeps.
This is proportional to the total computational work,
not counting the zero order forward sweep,
or combining multiple columns and rows into a single sweep.

\index{uses@\spxentry{uses}}\index{forward@\spxentry{forward}}\ignorespaces 

\subparagraph{Uses Forward}
\label{\detokenize{xsrst/py_sparse_hes:uses-forward}}\label{\detokenize{xsrst/py_sparse_hes:py-sparse-hes-uses-forward}}\label{\detokenize{xsrst/py_sparse_hes:index-10}}
After each call to {\hyperref[\detokenize{xsrst/py_fun_forward:id1}]{\sphinxcrossref{\DUrole{std,std-ref}{py\_fun\_forward}}}},
the object \sphinxstyleemphasis{f} contains the corresponding Taylor coefficients
for all the variables in the operation sequence..
After a call to \sphinxcode{\sphinxupquote{sparse\_hes}}
the zero order coefficients correspond to

\begin{DUlineblock}{0em}
\item[]         \sphinxstyleemphasis{f} . \sphinxcode{\sphinxupquote{forward(0}} , \sphinxstyleemphasis{x} )
\end{DUlineblock}

All the other forward mode coefficients are unspecified.


\subparagraph{sparse\_hes\_xam\_py}
\label{\detokenize{xsrst/sparse_hes_xam_py:sparse-hes-xam-py}}\label{\detokenize{xsrst/sparse_hes_xam_py::doc}}
\index{sparse\_hes\_xam\_py@\spxentry{sparse\_hes\_xam\_py}}\index{python:@\spxentry{python:}}\index{hessian@\spxentry{hessian}}\index{sparsity@\spxentry{sparsity}}\index{patterns:@\spxentry{patterns:}}\index{example@\spxentry{example}}\index{test@\spxentry{test}}\ignorespaces 

\subparagraph{3.1.2.6.1 Python: Hessian Sparsity Patterns: Example and Test}
\label{\detokenize{xsrst/sparse_hes_xam_py:python-hessian-sparsity-patterns-example-and-test}}\label{\detokenize{xsrst/sparse_hes_xam_py:id1}}
\begin{sphinxVerbatim}[commandchars=\\\{\}]
\PYG{k}{def} \PYG{n+nf}{sparse\PYGZus{}hes\PYGZus{}xam}\PYG{p}{(}\PYG{p}{)} \PYG{p}{:}
    \PYG{c+c1}{\PYGZsh{}}
    \PYG{k+kn}{import} \PYG{n+nn}{numpy}
    \PYG{k+kn}{import} \PYG{n+nn}{cppad\PYGZus{}py}
    \PYG{c+c1}{\PYGZsh{}}
    \PYG{c+c1}{\PYGZsh{} initialize return variable}
    \PYG{n}{ok} \PYG{o}{=} \PYG{k+kc}{True}
    \PYG{c+c1}{\PYGZsh{} \PYGZhy{}\PYGZhy{}\PYGZhy{}\PYGZhy{}\PYGZhy{}\PYGZhy{}\PYGZhy{}\PYGZhy{}\PYGZhy{}\PYGZhy{}\PYGZhy{}\PYGZhy{}\PYGZhy{}\PYGZhy{}\PYGZhy{}\PYGZhy{}\PYGZhy{}\PYGZhy{}\PYGZhy{}\PYGZhy{}\PYGZhy{}\PYGZhy{}\PYGZhy{}\PYGZhy{}\PYGZhy{}\PYGZhy{}\PYGZhy{}\PYGZhy{}\PYGZhy{}\PYGZhy{}\PYGZhy{}\PYGZhy{}\PYGZhy{}\PYGZhy{}\PYGZhy{}\PYGZhy{}\PYGZhy{}\PYGZhy{}\PYGZhy{}\PYGZhy{}\PYGZhy{}\PYGZhy{}\PYGZhy{}\PYGZhy{}\PYGZhy{}\PYGZhy{}\PYGZhy{}\PYGZhy{}\PYGZhy{}\PYGZhy{}\PYGZhy{}\PYGZhy{}\PYGZhy{}\PYGZhy{}\PYGZhy{}\PYGZhy{}\PYGZhy{}\PYGZhy{}\PYGZhy{}\PYGZhy{}\PYGZhy{}\PYGZhy{}\PYGZhy{}\PYGZhy{}\PYGZhy{}\PYGZhy{}\PYGZhy{}\PYGZhy{}\PYGZhy{}}
    \PYG{c+c1}{\PYGZsh{} number of dependent and independent variables}
    \PYG{n}{n} \PYG{o}{=} \PYG{l+m+mi}{3}
    \PYG{c+c1}{\PYGZsh{}}
    \PYG{c+c1}{\PYGZsh{} create the independent variables ax}
    \PYG{n}{x} \PYG{o}{=} \PYG{n}{numpy}\PYG{o}{.}\PYG{n}{empty}\PYG{p}{(}\PYG{n}{n}\PYG{p}{,} \PYG{n}{dtype}\PYG{o}{=}\PYG{n+nb}{float}\PYG{p}{)}
    \PYG{k}{for} \PYG{n}{i} \PYG{o+ow}{in} \PYG{n+nb}{range}\PYG{p}{(} \PYG{n}{n}  \PYG{p}{)} \PYG{p}{:}
        \PYG{n}{x}\PYG{p}{[}\PYG{n}{i}\PYG{p}{]} \PYG{o}{=} \PYG{n}{i} \PYG{o}{+} \PYG{l+m+mf}{2.0}
    \PYG{c+c1}{\PYGZsh{}}
    \PYG{n}{ax} \PYG{o}{=} \PYG{n}{cppad\PYGZus{}py}\PYG{o}{.}\PYG{n}{independent}\PYG{p}{(}\PYG{n}{x}\PYG{p}{)}
    \PYG{c+c1}{\PYGZsh{}}
    \PYG{c+c1}{\PYGZsh{} ay[i] = j * ax[j] * ax[i]}
    \PYG{c+c1}{\PYGZsh{} where i = mod(j + 1, n)}
    \PYG{n}{ay} \PYG{o}{=} \PYG{n}{numpy}\PYG{o}{.}\PYG{n}{empty}\PYG{p}{(}\PYG{n}{n}\PYG{p}{,} \PYG{n}{dtype}\PYG{o}{=}\PYG{n}{cppad\PYGZus{}py}\PYG{o}{.}\PYG{n}{a\PYGZus{}double}\PYG{p}{)}
    \PYG{k}{for} \PYG{n}{j} \PYG{o+ow}{in} \PYG{n+nb}{range}\PYG{p}{(} \PYG{n}{n}  \PYG{p}{)} \PYG{p}{:}
        \PYG{n}{i} \PYG{o}{=} \PYG{n}{j}\PYG{o}{+}\PYG{l+m+mi}{1}
        \PYG{k}{if} \PYG{n}{i} \PYG{o}{\PYGZgt{}}\PYG{o}{=} \PYG{n}{n}  \PYG{p}{:}
            \PYG{n}{i} \PYG{o}{=} \PYG{n}{i} \PYG{o}{\PYGZhy{}} \PYG{n}{n}
        \PYG{c+c1}{\PYGZsh{}}
        \PYG{n}{aj} \PYG{o}{=} \PYG{n}{cppad\PYGZus{}py}\PYG{o}{.}\PYG{n}{a\PYGZus{}double}\PYG{p}{(}\PYG{n}{j}\PYG{p}{)}
        \PYG{n}{ax\PYGZus{}j} \PYG{o}{=} \PYG{n}{ax}\PYG{p}{[}\PYG{n}{j}\PYG{p}{]}
        \PYG{n}{ax\PYGZus{}i} \PYG{o}{=} \PYG{n}{ax}\PYG{p}{[}\PYG{n}{i}\PYG{p}{]}
        \PYG{n}{ay}\PYG{p}{[}\PYG{n}{i}\PYG{p}{]} \PYG{o}{=} \PYG{n}{aj} \PYG{o}{*} \PYG{n}{ax\PYGZus{}j} \PYG{o}{*} \PYG{n}{ax\PYGZus{}i}
    \PYG{c+c1}{\PYGZsh{}}
    \PYG{c+c1}{\PYGZsh{}}
    \PYG{c+c1}{\PYGZsh{} define af corresponding to f(x)}
    \PYG{n}{f}  \PYG{o}{=} \PYG{n}{cppad\PYGZus{}py}\PYG{o}{.}\PYG{n}{d\PYGZus{}fun}\PYG{p}{(}\PYG{n}{ax}\PYG{p}{,} \PYG{n}{ay}\PYG{p}{)}
    \PYG{c+c1}{\PYGZsh{}}
    \PYG{c+c1}{\PYGZsh{} Set select\PYGZus{}d (domain) to all true,}
    \PYG{c+c1}{\PYGZsh{} initial select\PYGZus{}r (range) to all false}
    \PYG{c+c1}{\PYGZsh{} initialize r to all zeros}
    \PYG{n}{select\PYGZus{}d} \PYG{o}{=} \PYG{n}{numpy}\PYG{o}{.}\PYG{n}{empty}\PYG{p}{(}\PYG{n}{n}\PYG{p}{,} \PYG{n}{dtype}\PYG{o}{=}\PYG{n+nb}{bool}\PYG{p}{)}
    \PYG{n}{select\PYGZus{}r} \PYG{o}{=} \PYG{n}{numpy}\PYG{o}{.}\PYG{n}{empty}\PYG{p}{(}\PYG{n}{n}\PYG{p}{,} \PYG{n}{dtype}\PYG{o}{=}\PYG{n+nb}{bool}\PYG{p}{)}
    \PYG{n}{r} \PYG{o}{=} \PYG{n}{numpy}\PYG{o}{.}\PYG{n}{empty}\PYG{p}{(}\PYG{n}{n}\PYG{p}{,} \PYG{n}{dtype}\PYG{o}{=}\PYG{n+nb}{float}\PYG{p}{)}
    \PYG{k}{for} \PYG{n}{i} \PYG{o+ow}{in} \PYG{n+nb}{range}\PYG{p}{(} \PYG{n}{n} \PYG{p}{)} \PYG{p}{:}
        \PYG{n}{select\PYGZus{}d}\PYG{p}{[}\PYG{n}{i}\PYG{p}{]} \PYG{o}{=} \PYG{k+kc}{True}
        \PYG{n}{select\PYGZus{}r}\PYG{p}{[}\PYG{n}{i}\PYG{p}{]} \PYG{o}{=} \PYG{k+kc}{False}
        \PYG{n}{r}\PYG{p}{[}\PYG{n}{i}\PYG{p}{]} \PYG{o}{=} \PYG{l+m+mf}{0.0}
    \PYG{c+c1}{\PYGZsh{}}
    \PYG{c+c1}{\PYGZsh{}}
    \PYG{c+c1}{\PYGZsh{} only select component n\PYGZhy{}1 of the range function}
    \PYG{c+c1}{\PYGZsh{} f\PYGZus{}0 (x) = (n\PYGZhy{}1) * x\PYGZus{}\PYGZob{}n\PYGZhy{}1\PYGZcb{} * x\PYGZus{}0}
    \PYG{n}{select\PYGZus{}r}\PYG{p}{[}\PYG{l+m+mi}{0}\PYG{p}{]} \PYG{o}{=} \PYG{k+kc}{True}
    \PYG{n}{r}\PYG{p}{[}\PYG{l+m+mi}{0}\PYG{p}{]} \PYG{o}{=} \PYG{l+m+mf}{1.0}
    \PYG{c+c1}{\PYGZsh{}}
    \PYG{c+c1}{\PYGZsh{} sparisty pattern for Hessian}
    \PYG{n}{pattern} \PYG{o}{=} \PYG{n}{cppad\PYGZus{}py}\PYG{o}{.}\PYG{n}{sparse\PYGZus{}rc}\PYG{p}{(}\PYG{p}{)}
    \PYG{n}{f}\PYG{o}{.}\PYG{n}{for\PYGZus{}hes\PYGZus{}sparsity}\PYG{p}{(}\PYG{n}{select\PYGZus{}d}\PYG{p}{,} \PYG{n}{select\PYGZus{}r}\PYG{p}{,} \PYG{n}{pattern}\PYG{p}{)}
    \PYG{c+c1}{\PYGZsh{}}
    \PYG{c+c1}{\PYGZsh{} compute all possibly non\PYGZhy{}zero entries in Hessian}
    \PYG{c+c1}{\PYGZsh{} (should only compute lower triangle becuase matrix is symmetric)}
    \PYG{n}{subset} \PYG{o}{=} \PYG{n}{cppad\PYGZus{}py}\PYG{o}{.}\PYG{n}{sparse\PYGZus{}rcv}\PYG{p}{(}\PYG{n}{pattern}\PYG{p}{)}
    \PYG{c+c1}{\PYGZsh{}}
    \PYG{c+c1}{\PYGZsh{} work space used to save time for multiple calls}
    \PYG{n}{work} \PYG{o}{=} \PYG{n}{cppad\PYGZus{}py}\PYG{o}{.}\PYG{n}{sparse\PYGZus{}hes\PYGZus{}work}\PYG{p}{(}\PYG{p}{)}
    \PYG{c+c1}{\PYGZsh{}}
    \PYG{n}{f}\PYG{o}{.}\PYG{n}{sparse\PYGZus{}hes}\PYG{p}{(}\PYG{n}{subset}\PYG{p}{,} \PYG{n}{x}\PYG{p}{,} \PYG{n}{r}\PYG{p}{,} \PYG{n}{pattern}\PYG{p}{,} \PYG{n}{work}\PYG{p}{)}
    \PYG{c+c1}{\PYGZsh{}}
    \PYG{c+c1}{\PYGZsh{} check that result is sparsity pattern for Hessian of f\PYGZus{}0 (x)}
    \PYG{n}{ok} \PYG{o}{=} \PYG{n}{ok} \PYG{o+ow}{and} \PYG{n}{subset}\PYG{o}{.}\PYG{n}{nnz}\PYG{p}{(}\PYG{p}{)} \PYG{o}{==} \PYG{l+m+mi}{2}
    \PYG{n}{row} \PYG{o}{=} \PYG{n}{subset}\PYG{o}{.}\PYG{n}{row}\PYG{p}{(}\PYG{p}{)}
    \PYG{n}{col} \PYG{o}{=} \PYG{n}{subset}\PYG{o}{.}\PYG{n}{col}\PYG{p}{(}\PYG{p}{)}
    \PYG{n}{val} \PYG{o}{=} \PYG{n}{subset}\PYG{o}{.}\PYG{n}{val}\PYG{p}{(}\PYG{p}{)}
    \PYG{k}{for} \PYG{n}{k} \PYG{o+ow}{in} \PYG{n+nb}{range}\PYG{p}{(} \PYG{l+m+mi}{2} \PYG{p}{)} \PYG{p}{:}
        \PYG{n}{i} \PYG{o}{=} \PYG{n}{row}\PYG{p}{[}\PYG{n}{k}\PYG{p}{]}
        \PYG{n}{j} \PYG{o}{=} \PYG{n}{col}\PYG{p}{[}\PYG{n}{k}\PYG{p}{]}
        \PYG{n}{v} \PYG{o}{=} \PYG{n}{val}\PYG{p}{[}\PYG{n}{k}\PYG{p}{]}
        \PYG{k}{if} \PYG{n}{i} \PYG{o}{\PYGZlt{}}\PYG{o}{=} \PYG{n}{j}  \PYG{p}{:}
            \PYG{n}{ok} \PYG{o}{=} \PYG{n}{ok} \PYG{o+ow}{and} \PYG{n}{i} \PYG{o}{==} \PYG{l+m+mi}{0}
            \PYG{n}{ok} \PYG{o}{=} \PYG{n}{ok} \PYG{o+ow}{and} \PYG{n}{j} \PYG{o}{==} \PYG{n}{n}\PYG{o}{\PYGZhy{}}\PYG{l+m+mi}{1}
        \PYG{c+c1}{\PYGZsh{}}
        \PYG{k}{if} \PYG{n}{i} \PYG{o}{\PYGZgt{}}\PYG{o}{=} \PYG{n}{j}  \PYG{p}{:}
            \PYG{n}{ok} \PYG{o}{=} \PYG{n}{ok} \PYG{o+ow}{and} \PYG{n}{i} \PYG{o}{==} \PYG{n}{n}\PYG{o}{\PYGZhy{}}\PYG{l+m+mi}{1}
            \PYG{n}{ok} \PYG{o}{=} \PYG{n}{ok} \PYG{o+ow}{and} \PYG{n}{j} \PYG{o}{==} \PYG{l+m+mi}{0}
        \PYG{c+c1}{\PYGZsh{}}
        \PYG{n}{ok} \PYG{o}{=} \PYG{n}{ok} \PYG{o+ow}{and} \PYG{n}{v} \PYG{o}{==} \PYG{n}{n}\PYG{o}{\PYGZhy{}}\PYG{l+m+mi}{1}
    \PYG{c+c1}{\PYGZsh{}}
    \PYG{c+c1}{\PYGZsh{}}
    \PYG{k}{return}\PYG{p}{(} \PYG{n}{ok} \PYG{p}{)}
\PYG{c+c1}{\PYGZsh{}}
\end{sphinxVerbatim}


\bigskip\hrule\bigskip


xsrst input file: \sphinxcode{\sphinxupquote{example/python/core/sparse\_hes\_xam.py}}

\index{example@\spxentry{example}}\ignorespaces 

\subparagraph{Example}
\label{\detokenize{xsrst/py_sparse_hes:example}}\label{\detokenize{xsrst/py_sparse_hes:py-sparse-hes-example}}\label{\detokenize{xsrst/py_sparse_hes:index-11}}
{\hyperref[\detokenize{xsrst/sparse_hes_xam_py:id1}]{\sphinxcrossref{\DUrole{std,std-ref}{sparse\_hes\_xam\_py}}}}


\bigskip\hrule\bigskip


xsrst input file: \sphinxcode{\sphinxupquote{lib/python/cppad\_py/sparse\_hes.py}}
\begin{enumerate}
\sphinxsetlistlabels{\arabic}{enumi}{enumii}{}{.}%
\item {} 
py\_sparse\_rc: {\hyperref[\detokenize{xsrst/py_sparse_rc:id1}]{\sphinxcrossref{\DUrole{std,std-ref}{3.1.2.1 Sparsity Patterns}}}}

\item {} 
py\_sparse\_rcv: {\hyperref[\detokenize{xsrst/py_sparse_rcv:id1}]{\sphinxcrossref{\DUrole{std,std-ref}{3.1.2.2 Sparse Matrices}}}}

\item {} 
py\_jac\_sparsity: {\hyperref[\detokenize{xsrst/py_jac_sparsity:id1}]{\sphinxcrossref{\DUrole{std,std-ref}{3.1.2.3 Jacobian Sparsity Patterns}}}}

\item {} 
py\_hes\_sparsity: {\hyperref[\detokenize{xsrst/py_hes_sparsity:id1}]{\sphinxcrossref{\DUrole{std,std-ref}{3.1.2.4 Hessian Sparsity Patterns}}}}

\item {} 
py\_sparse\_jac: {\hyperref[\detokenize{xsrst/py_sparse_jac:id1}]{\sphinxcrossref{\DUrole{std,std-ref}{3.1.2.5 Computing Sparse Jacobians}}}}

\item {} 
py\_sparse\_hes: {\hyperref[\detokenize{xsrst/py_sparse_hes:id1}]{\sphinxcrossref{\DUrole{std,std-ref}{3.1.2.6 Computing Sparse Hessians}}}}

\end{enumerate}


\bigskip\hrule\bigskip


xsrst input file: \sphinxcode{\sphinxupquote{lib/python/cppad\_py/sparse.xsrst}}


\subparagraph{py\_utility}
\label{\detokenize{xsrst/py_utility:py-utility}}\label{\detokenize{xsrst/py_utility::doc}}
\index{py\_utility@\spxentry{py\_utility}}\index{python@\spxentry{python}}\index{utilities@\spxentry{utilities}}\ignorespaces 

\subparagraph{3.1.3 Python Utilities}
\label{\detokenize{xsrst/py_utility:python-utilities}}\label{\detokenize{xsrst/py_utility:id1}}\begin{itemize}
\item {} 
{\hyperref[\detokenize{xsrst/py_utility:py-utility-children}]{\sphinxcrossref{\DUrole{std,std-ref}{Children}}}}

\end{itemize}

\index{children@\spxentry{children}}\ignorespaces 

\subparagraph{Children}
\label{\detokenize{xsrst/py_utility:children}}\label{\detokenize{xsrst/py_utility:py-utility-children}}\label{\detokenize{xsrst/py_utility:index-1}}

\subparagraph{numpy2vec}
\label{\detokenize{xsrst/numpy2vec:numpy2vec}}\label{\detokenize{xsrst/numpy2vec::doc}}
\index{numpy2vec@\spxentry{numpy2vec}}\index{convert@\spxentry{convert}}\index{numpy@\spxentry{numpy}}\index{array@\spxentry{array}}\index{to@\spxentry{to}}\index{cppad\_py@\spxentry{cppad\_py}}\index{vector@\spxentry{vector}}\ignorespaces 

\subparagraph{3.1.3.1 Convert a Numpy Array to a cppad\_py Vector}
\label{\detokenize{xsrst/numpy2vec:convert-a-numpy-array-to-a-cppad-py-vector}}\label{\detokenize{xsrst/numpy2vec:id1}}\begin{itemize}
\item {} 
{\hyperref[\detokenize{xsrst/numpy2vec:numpy2vec-syntax}]{\sphinxcrossref{\DUrole{std,std-ref}{Syntax}}}}

\item {} 
{\hyperref[\detokenize{xsrst/numpy2vec:numpy2vec-array}]{\sphinxcrossref{\DUrole{std,std-ref}{array}}}}

\item {} 
{\hyperref[\detokenize{xsrst/numpy2vec:numpy2vec-dtype}]{\sphinxcrossref{\DUrole{std,std-ref}{dtype}}}}

\item {} 
{\hyperref[\detokenize{xsrst/numpy2vec:numpy2vec-shape}]{\sphinxcrossref{\DUrole{std,std-ref}{shape}}}}

\item {} 
{\hyperref[\detokenize{xsrst/numpy2vec:numpy2vec-context}]{\sphinxcrossref{\DUrole{std,std-ref}{context}}}}

\item {} 
{\hyperref[\detokenize{xsrst/numpy2vec:numpy2vec-name}]{\sphinxcrossref{\DUrole{std,std-ref}{name}}}}

\item {} 
{\hyperref[\detokenize{xsrst/numpy2vec:numpy2vec-vec}]{\sphinxcrossref{\DUrole{std,std-ref}{vec}}}}

\end{itemize}

\index{syntax@\spxentry{syntax}}\ignorespaces 

\subparagraph{Syntax}
\label{\detokenize{xsrst/numpy2vec:syntax}}\label{\detokenize{xsrst/numpy2vec:numpy2vec-syntax}}\label{\detokenize{xsrst/numpy2vec:index-1}}
\begin{DUlineblock}{0em}
\item[] \sphinxstyleemphasis{vec} =  \sphinxcode{\sphinxupquote{cppad\_py.utility.numpy2vec}} (
\item[]         \sphinxstyleemphasis{array} , \sphinxstyleemphasis{dtype} , \sphinxstyleemphasis{shape} , \sphinxstyleemphasis{context} , \sphinxstyleemphasis{name}
\item[] )
\end{DUlineblock}

\index{array@\spxentry{array}}\ignorespaces 

\subparagraph{array}
\label{\detokenize{xsrst/numpy2vec:array}}\label{\detokenize{xsrst/numpy2vec:numpy2vec-array}}\label{\detokenize{xsrst/numpy2vec:index-2}}
This is either a vector (only one index) or a matrix
(has two indices) that we are converting to a vector.
If this array does not match the conditions below,
an exception is raised with an appropriate error message.

\index{dtype@\spxentry{dtype}}\ignorespaces 

\subparagraph{dtype}
\label{\detokenize{xsrst/numpy2vec:dtype}}\label{\detokenize{xsrst/numpy2vec:numpy2vec-dtype}}\label{\detokenize{xsrst/numpy2vec:index-3}}
This is the expected data type for the elements of the array.
It must be one of the following:
\sphinxcode{\sphinxupquote{bool}} , \sphinxcode{\sphinxupquote{int}} , \sphinxcode{\sphinxupquote{float}} or \sphinxcode{\sphinxupquote{cppad\_py.a\_double}} .

\index{shape@\spxentry{shape}}\ignorespaces 

\subparagraph{shape}
\label{\detokenize{xsrst/numpy2vec:shape}}\label{\detokenize{xsrst/numpy2vec:numpy2vec-shape}}\label{\detokenize{xsrst/numpy2vec:index-4}}
This either a \sphinxcode{\sphinxupquote{int}} or a tuple of \sphinxcode{\sphinxupquote{int}} with length one or two.
If it is an \sphinxcode{\sphinxupquote{int}} , \sphinxstyleemphasis{array} is expected to be a vector.
If it is a tuple with length one, \sphinxstyleemphasis{array} is expected to be a vector.
Otherwise a matrix is expected.

\index{context@\spxentry{context}}\ignorespaces 

\subparagraph{context}
\label{\detokenize{xsrst/numpy2vec:context}}\label{\detokenize{xsrst/numpy2vec:numpy2vec-context}}\label{\detokenize{xsrst/numpy2vec:index-5}}
This is the context that \sphinxstyleemphasis{array} appears in
(often to syntax for some other operation).
It is a \sphinxcode{\sphinxupquote{str}} that is used for error reporting.

\index{name@\spxentry{name}}\ignorespaces 

\subparagraph{name}
\label{\detokenize{xsrst/numpy2vec:name}}\label{\detokenize{xsrst/numpy2vec:numpy2vec-name}}\label{\detokenize{xsrst/numpy2vec:index-6}}
This is the name used for \sphinxstyleemphasis{array} in \sphinxstyleemphasis{context} .
It is a \sphinxcode{\sphinxupquote{str}} that is used for error reporting.

\index{vec@\spxentry{vec}}\ignorespaces 

\subparagraph{vec}
\label{\detokenize{xsrst/numpy2vec:vec}}\label{\detokenize{xsrst/numpy2vec:numpy2vec-vec}}\label{\detokenize{xsrst/numpy2vec:index-7}}
This is the values of \sphinxstyleemphasis{array}
as a vector is row major order.
It has type \sphinxcode{\sphinxupquote{vec\_double}} if \sphinxstyleemphasis{dtype} is \sphinxcode{\sphinxupquote{float}} ,
and \sphinxcode{\sphinxupquote{vec\_a\_double}} if \sphinxstyleemphasis{dtype} is \sphinxcode{\sphinxupquote{a\_double}} .


\bigskip\hrule\bigskip


xsrst input file: \sphinxcode{\sphinxupquote{lib/python/cppad\_py/utility.py}}


\subparagraph{vec2numpy}
\label{\detokenize{xsrst/vec2numpy:vec2numpy}}\label{\detokenize{xsrst/vec2numpy::doc}}
\index{vec2numpy@\spxentry{vec2numpy}}\index{convert@\spxentry{convert}}\index{cppad\_py@\spxentry{cppad\_py}}\index{vector@\spxentry{vector}}\index{to@\spxentry{to}}\index{numpy@\spxentry{numpy}}\index{array@\spxentry{array}}\ignorespaces 

\subparagraph{3.1.3.2 Convert a cppad\_py Vector to a Numpy Array}
\label{\detokenize{xsrst/vec2numpy:convert-a-cppad-py-vector-to-a-numpy-array}}\label{\detokenize{xsrst/vec2numpy:id1}}\begin{itemize}
\item {} 
{\hyperref[\detokenize{xsrst/vec2numpy:vec2numpy-syntax}]{\sphinxcrossref{\DUrole{std,std-ref}{Syntax}}}}

\item {} 
{\hyperref[\detokenize{xsrst/vec2numpy:vec2numpy-vec}]{\sphinxcrossref{\DUrole{std,std-ref}{vec}}}}

\item {} 
{\hyperref[\detokenize{xsrst/vec2numpy:vec2numpy-nr}]{\sphinxcrossref{\DUrole{std,std-ref}{nr}}}}

\item {} 
{\hyperref[\detokenize{xsrst/vec2numpy:vec2numpy-nc}]{\sphinxcrossref{\DUrole{std,std-ref}{nc}}}}

\item {} 
{\hyperref[\detokenize{xsrst/vec2numpy:vec2numpy-array}]{\sphinxcrossref{\DUrole{std,std-ref}{array}}}}

\end{itemize}

\index{syntax@\spxentry{syntax}}\ignorespaces 

\subparagraph{Syntax}
\label{\detokenize{xsrst/vec2numpy:syntax}}\label{\detokenize{xsrst/vec2numpy:vec2numpy-syntax}}\label{\detokenize{xsrst/vec2numpy:index-1}}
\begin{DUlineblock}{0em}
\item[] \sphinxstyleemphasis{array} =  \sphinxcode{\sphinxupquote{cppad\_py.utility.vec2numpy}} ( \sphinxstyleemphasis{vec} , \sphinxstyleemphasis{nr} )
\item[] \sphinxstyleemphasis{array} =  \sphinxcode{\sphinxupquote{cppad\_py.utility.vec2numpy}} ( \sphinxstyleemphasis{vec} , \sphinxstyleemphasis{nr} , \sphinxstyleemphasis{nc} )
\end{DUlineblock}

\index{vec@\spxentry{vec}}\ignorespaces 

\subparagraph{vec}
\label{\detokenize{xsrst/vec2numpy:vec}}\label{\detokenize{xsrst/vec2numpy:vec2numpy-vec}}\label{\detokenize{xsrst/vec2numpy:index-2}}
This must have one of the following types:
\sphinxcode{\sphinxupquote{cppad\_py.vec\_int}} ,
\sphinxcode{\sphinxupquote{cppad\_py.vec\_double}} ,
\sphinxcode{\sphinxupquote{vec\_py.vec\_a\_double}} .
with size equal to \sphinxstyleemphasis{nr} \sphinxcode{\sphinxupquote{*}} \sphinxstyleemphasis{nc} .

\index{nr@\spxentry{nr}}\ignorespaces 

\subparagraph{nr}
\label{\detokenize{xsrst/vec2numpy:nr}}\label{\detokenize{xsrst/vec2numpy:vec2numpy-nr}}\label{\detokenize{xsrst/vec2numpy:index-3}}
This is an \sphinxcode{\sphinxupquote{int}} equal to the number of rows in the array.
If the argument \sphinxstyleemphasis{nc} is not present, the array is a vector; i.e.,
\sphinxcode{\sphinxupquote{len}} ( \sphinxstyleemphasis{array} . \sphinxcode{\sphinxupquote{shape ) == 1}} .

\index{nc@\spxentry{nc}}\ignorespaces 

\subparagraph{nc}
\label{\detokenize{xsrst/vec2numpy:nc}}\label{\detokenize{xsrst/vec2numpy:vec2numpy-nc}}\label{\detokenize{xsrst/vec2numpy:index-4}}
If this argument is present,
it is an \sphinxcode{\sphinxupquote{int}} equal to the number of columns in the array.
In this case the array is a matrix; i.e.,
\sphinxcode{\sphinxupquote{len}} ( \sphinxstyleemphasis{array} . \sphinxcode{\sphinxupquote{shape ) == 2}} .

\index{array@\spxentry{array}}\ignorespaces 

\subparagraph{array}
\label{\detokenize{xsrst/vec2numpy:array}}\label{\detokenize{xsrst/vec2numpy:vec2numpy-array}}\label{\detokenize{xsrst/vec2numpy:index-5}}
This is the array corresponding to \sphinxstyleemphasis{vec} in row major order.
Note that this array can be used after the vector \sphinxstyleemphasis{vec} drops
out of scope (is deleted).


\bigskip\hrule\bigskip


xsrst input file: \sphinxcode{\sphinxupquote{lib/python/cppad\_py/utility.py}}
\begin{enumerate}
\sphinxsetlistlabels{\arabic}{enumi}{enumii}{}{.}%
\item {} 
numpy2vec: {\hyperref[\detokenize{xsrst/numpy2vec:id1}]{\sphinxcrossref{\DUrole{std,std-ref}{3.1.3.1 Convert a Numpy Array to a cppad\_py Vector}}}}

\item {} 
vec2numpy: {\hyperref[\detokenize{xsrst/vec2numpy:id1}]{\sphinxcrossref{\DUrole{std,std-ref}{3.1.3.2 Convert a cppad\_py Vector to a Numpy Array}}}}

\end{enumerate}


\bigskip\hrule\bigskip


xsrst input file: \sphinxcode{\sphinxupquote{lib/python/cppad\_py/utility.xsrst}}


\subparagraph{more\_py}
\label{\detokenize{xsrst/more_py:more-py}}\label{\detokenize{xsrst/more_py::doc}}
\index{more\_py@\spxentry{more\_py}}\index{steps@\spxentry{steps}}\index{to@\spxentry{to}}\index{add@\spxentry{add}}\index{more@\spxentry{more}}\index{python@\spxentry{python}}\index{functions@\spxentry{functions}}\ignorespaces 

\subparagraph{3.1.4 Steps To Add More Python Functions}
\label{\detokenize{xsrst/more_py:steps-to-add-more-python-functions}}\label{\detokenize{xsrst/more_py:id1}}\begin{itemize}
\item {} 
{\hyperref[\detokenize{xsrst/more_py:more-py-purpose}]{\sphinxcrossref{\DUrole{std,std-ref}{Purpose}}}}

\item {} \begin{description}
\item[{{\hyperref[\detokenize{xsrst/more_py:more-py-documentation}]{\sphinxcrossref{\DUrole{std,std-ref}{Documentation}}}}}] \leavevmode\begin{itemize}
\item {} 
{\hyperref[\detokenize{xsrst/more_py:more-py-documentation-independent}]{\sphinxcrossref{\DUrole{std,std-ref}{independent}}}}

\item {} 
{\hyperref[\detokenize{xsrst/more_py:more-py-documentation-new-dynamic}]{\sphinxcrossref{\DUrole{std,std-ref}{new\_dynamic}}}}

\item {} 
{\hyperref[\detokenize{xsrst/more_py:more-py-documentation-example}]{\sphinxcrossref{\DUrole{std,std-ref}{Example}}}}

\end{itemize}

\end{description}

\item {} \begin{description}
\item[{{\hyperref[\detokenize{xsrst/more_py:more-py-implementation}]{\sphinxcrossref{\DUrole{std,std-ref}{Implementation}}}}}] \leavevmode\begin{itemize}
\item {} 
{\hyperref[\detokenize{xsrst/more_py:more-py-implementation-independent}]{\sphinxcrossref{\DUrole{std,std-ref}{independent}}}}

\item {} 
{\hyperref[\detokenize{xsrst/more_py:more-py-implementation-new-dynamic}]{\sphinxcrossref{\DUrole{std,std-ref}{new\_dynamic}}}}

\item {} 
{\hyperref[\detokenize{xsrst/more_py:more-py-implementation-fun-new-dynamic-py}]{\sphinxcrossref{\DUrole{std,std-ref}{fun\_new\_dynamic.py}}}}

\item {} 
{\hyperref[\detokenize{xsrst/more_py:more-py-implementation-init-py}]{\sphinxcrossref{\DUrole{std,std-ref}{\_\_init\_\_.py}}}}

\end{itemize}

\end{description}

\item {} 
{\hyperref[\detokenize{xsrst/more_py:more-py-example}]{\sphinxcrossref{\DUrole{std,std-ref}{Example}}}}

\item {} 
{\hyperref[\detokenize{xsrst/more_py:more-py-testing}]{\sphinxcrossref{\DUrole{std,std-ref}{Testing}}}}

\end{itemize}

\index{purpose@\spxentry{purpose}}\ignorespaces 

\subparagraph{Purpose}
\label{\detokenize{xsrst/more_py:purpose}}\label{\detokenize{xsrst/more_py:more-py-purpose}}\label{\detokenize{xsrst/more_py:index-1}}
This section outlines the steps for adding more python functionality
to cppad\_py.
This is done by example showing how the {\hyperref[\detokenize{xsrst/py_fun_new_dynamic:id1}]{\sphinxcrossref{\DUrole{std,std-ref}{py\_fun\_new\_dynamic}}}} was added.
This example case was chosen because it required both changing one
python function, {\hyperref[\detokenize{xsrst/py_independent:id1}]{\sphinxcrossref{\DUrole{std,std-ref}{py\_independent}}}},
and adding a new python function, {\hyperref[\detokenize{xsrst/py_fun_new_dynamic:id1}]{\sphinxcrossref{\DUrole{std,std-ref}{py\_fun\_new\_dynamic}}}}.

\index{documentation@\spxentry{documentation}}\ignorespaces 

\subparagraph{Documentation}
\label{\detokenize{xsrst/more_py:documentation}}\label{\detokenize{xsrst/more_py:more-py-documentation}}\label{\detokenize{xsrst/more_py:index-2}}
\index{independent@\spxentry{independent}}\ignorespaces 

\subparagraph{independent}
\label{\detokenize{xsrst/more_py:independent}}\label{\detokenize{xsrst/more_py:more-py-documentation-independent}}\label{\detokenize{xsrst/more_py:index-3}}
The python file \sphinxcode{\sphinxupquote{lib/python/cppad\_py/independent.py}}
was edited to add the following syntax documentation:

\begin{DUlineblock}{0em}
\item[]         ( \sphinxstyleemphasis{ax} , \sphinxstyleemphasis{adynamic} ) =  \sphinxcode{\sphinxupquote{cppad\_py::independent}} ( \sphinxstyleemphasis{x} , \sphinxstyleemphasis{dynamic} )
\end{DUlineblock}

and the extra return value was documented; see
{\hyperref[\detokenize{xsrst/py_independent:py-independent-adynamic}]{\sphinxcrossref{\DUrole{std,std-ref}{adynamic}}}}.

\index{new\_dynamic@\spxentry{new\_dynamic}}\ignorespaces 

\subparagraph{new\_dynamic}
\label{\detokenize{xsrst/more_py:new-dynamic}}\label{\detokenize{xsrst/more_py:more-py-documentation-new-dynamic}}\label{\detokenize{xsrst/more_py:index-4}}
The {\hyperref[\detokenize{xsrst/cpp_fun_new_dynamic:id1}]{\sphinxcrossref{\DUrole{std,std-ref}{cpp\_fun\_new\_dynamic}}}} documentation was added
in the file \sphinxcode{\sphinxupquote{lib/python/cppad\_py/fun\_new\_dynamic.py}} .
In addition, the OMhelp command

\begin{DUlineblock}{0em}
\item[]         \sphinxcode{\sphinxupquote{\%lib/python/cppad\_py/fun\_new\_dynamic.py}}
\end{DUlineblock}

was added to the file \sphinxcode{\sphinxupquote{lib/python/cppad\_py/fun.py}} .

\index{example@\spxentry{example}}\ignorespaces 

\subparagraph{Example}
\label{\detokenize{xsrst/more_py:example}}\label{\detokenize{xsrst/more_py:more-py-documentation-example}}\label{\detokenize{xsrst/more_py:index-5}}
An example file was added to the documentation,
below the {\hyperref[\detokenize{xsrst/py_independent:id1}]{\sphinxcrossref{\DUrole{std,std-ref}{py\_independent}}}} section, using the OMhelp commands:

\begin{DUlineblock}{0em}
\item[]         \{ \sphinxcode{\sphinxupquote{xsrst\_dollar\}children\%}}
\item[]                 \sphinxcode{\sphinxupquote{example/python/core/fun\_dynamic\_xam.py}}
\item[]         \sphinxcode{\sphinxupquote{\%\$\{xsrst\_dollar}} \}
\end{DUlineblock}

In addition, a reference to this example was added under the
{\hyperref[\detokenize{xsrst/py_independent:py-independent-example}]{\sphinxcrossref{\DUrole{std,std-ref}{example}}}} heading in the \sphinxcode{\sphinxupquote{independent}}
documentation.

\index{implementation@\spxentry{implementation}}\ignorespaces 

\subparagraph{Implementation}
\label{\detokenize{xsrst/more_py:implementation}}\label{\detokenize{xsrst/more_py:more-py-implementation}}\label{\detokenize{xsrst/more_py:index-6}}
\index{independent@\spxentry{independent}}\ignorespaces 

\subparagraph{independent}
\label{\detokenize{xsrst/more_py:more-py-implementation-independent}}\label{\detokenize{xsrst/more_py:index-7}}\label{\detokenize{xsrst/more_py:id2}}
The \sphinxcode{\sphinxupquote{independent}} function in \sphinxcode{\sphinxupquote{lib/python/cppad\_py/independent.py}}
was changed to handle dynamic parameters as follows:

\begin{sphinxVerbatim}[commandchars=\\\{\}]
\PYG{k+kn}{import} \PYG{n+nn}{cppad\PYGZus{}py}
\PYG{k+kn}{import} \PYG{n+nn}{numpy}
\PYG{k}{def} \PYG{n+nf}{independent}\PYG{p}{(}\PYG{n}{x}\PYG{p}{,} \PYG{n}{dynamic} \PYG{o}{=} \PYG{k+kc}{None}\PYG{p}{)} \PYG{p}{:}
    \PYG{l+s+sd}{\PYGZdq{}\PYGZdq{}\PYGZdq{}}
\PYG{l+s+sd}{    ax = independent(x)}
\PYG{l+s+sd}{    creates the indepedent numpy vector ax, with value equal numpy vector x,}
\PYG{l+s+sd}{    and starts recording a\PYGZus{}double operations.}
\PYG{l+s+sd}{    \PYGZdq{}\PYGZdq{}\PYGZdq{}}
    \PYG{c+c1}{\PYGZsh{} convert x \PYGZhy{}\PYGZgt{} u}
    \PYG{n}{dtype}    \PYG{o}{=} \PYG{n+nb}{float}
    \PYG{c+c1}{\PYGZsh{}}
    \PYG{n}{nx} \PYG{o}{=} \PYG{n}{x}\PYG{o}{.}\PYG{n}{size}
    \PYG{k}{if} \PYG{n}{dynamic} \PYG{o+ow}{is} \PYG{k+kc}{None} \PYG{p}{:}
        \PYG{n}{syntax}   \PYG{o}{=} \PYG{l+s+s1}{\PYGZsq{}}\PYG{l+s+s1}{independent(x)}\PYG{l+s+s1}{\PYGZsq{}}
        \PYG{n}{u} \PYG{o}{=} \PYG{n}{cppad\PYGZus{}py}\PYG{o}{.}\PYG{n}{utility}\PYG{o}{.}\PYG{n}{numpy2vec}\PYG{p}{(}\PYG{n}{x}\PYG{p}{,} \PYG{n}{dtype}\PYG{p}{,} \PYG{n}{nx}\PYG{p}{,} \PYG{n}{syntax}\PYG{p}{,} \PYG{l+s+s1}{\PYGZsq{}}\PYG{l+s+s1}{x}\PYG{l+s+s1}{\PYGZsq{}}\PYG{p}{)}
        \PYG{n}{av} \PYG{o}{=} \PYG{n}{cppad\PYGZus{}py}\PYG{o}{.}\PYG{n}{swig}\PYG{o}{.}\PYG{n}{independent}\PYG{p}{(}\PYG{n}{u}\PYG{p}{)}
        \PYG{n}{ax} \PYG{o}{=} \PYG{n}{cppad\PYGZus{}py}\PYG{o}{.}\PYG{n}{utility}\PYG{o}{.}\PYG{n}{vec2numpy}\PYG{p}{(}\PYG{n}{av}\PYG{p}{,} \PYG{n}{av}\PYG{o}{.}\PYG{n}{size}\PYG{p}{(}\PYG{p}{)}\PYG{p}{)}\PYG{p}{;}
        \PYG{k}{return} \PYG{n}{ax}
    \PYG{c+c1}{\PYGZsh{}}
    \PYG{n}{nd}       \PYG{o}{=} \PYG{n}{dynamic}\PYG{o}{.}\PYG{n}{size}
    \PYG{n}{syntax}   \PYG{o}{=} \PYG{l+s+s1}{\PYGZsq{}}\PYG{l+s+s1}{independent(x, dynamic)}\PYG{l+s+s1}{\PYGZsq{}}
    \PYG{n}{u} \PYG{o}{=} \PYG{n}{cppad\PYGZus{}py}\PYG{o}{.}\PYG{n}{utility}\PYG{o}{.}\PYG{n}{numpy2vec}\PYG{p}{(}\PYG{n}{x}\PYG{p}{,} \PYG{n}{dtype}\PYG{p}{,} \PYG{n}{nx}\PYG{p}{,} \PYG{n}{syntax}\PYG{p}{,} \PYG{l+s+s1}{\PYGZsq{}}\PYG{l+s+s1}{x}\PYG{l+s+s1}{\PYGZsq{}}\PYG{p}{)}
    \PYG{n}{v} \PYG{o}{=} \PYG{n}{cppad\PYGZus{}py}\PYG{o}{.}\PYG{n}{utility}\PYG{o}{.}\PYG{n}{numpy2vec}\PYG{p}{(}\PYG{n}{dynamic}\PYG{p}{,} \PYG{n}{dtype}\PYG{p}{,} \PYG{n}{nd}\PYG{p}{,} \PYG{n}{syntax}\PYG{p}{,} \PYG{l+s+s1}{\PYGZsq{}}\PYG{l+s+s1}{dynamic}\PYG{l+s+s1}{\PYGZsq{}}\PYG{p}{)}
    \PYG{n}{a\PYGZus{}both}   \PYG{o}{=} \PYG{n}{cppad\PYGZus{}py}\PYG{o}{.}\PYG{n}{swig}\PYG{o}{.}\PYG{n}{independent}\PYG{p}{(}\PYG{n}{u}\PYG{p}{,} \PYG{n}{v}\PYG{p}{)}
    \PYG{n}{ax}       \PYG{o}{=} \PYG{n}{numpy}\PYG{o}{.}\PYG{n}{empty}\PYG{p}{(}\PYG{n}{nx}\PYG{p}{,}       \PYG{n}{dtype}\PYG{o}{=}\PYG{n}{cppad\PYGZus{}py}\PYG{o}{.}\PYG{n}{a\PYGZus{}double}\PYG{p}{)}
    \PYG{n}{adynamic} \PYG{o}{=} \PYG{n}{numpy}\PYG{o}{.}\PYG{n}{empty}\PYG{p}{(}\PYG{n}{nd}\PYG{p}{,} \PYG{n}{dtype}\PYG{o}{=}\PYG{n}{cppad\PYGZus{}py}\PYG{o}{.}\PYG{n}{a\PYGZus{}double}\PYG{p}{)}
    \PYG{c+c1}{\PYGZsh{} use copy constructor so a separate copy is made for numpy arrays}
    \PYG{k}{for} \PYG{n}{i} \PYG{o+ow}{in} \PYG{n+nb}{range}\PYG{p}{(}\PYG{n}{nx}\PYG{p}{)} \PYG{p}{:}
        \PYG{n}{ax}\PYG{p}{[}\PYG{n}{i}\PYG{p}{]} \PYG{o}{=} \PYG{n}{cppad\PYGZus{}py}\PYG{o}{.}\PYG{n}{a\PYGZus{}double}\PYG{p}{(} \PYG{n}{a\PYGZus{}both}\PYG{p}{[}\PYG{n}{i}\PYG{p}{]} \PYG{p}{)}
    \PYG{k}{for} \PYG{n}{i} \PYG{o+ow}{in} \PYG{n+nb}{range}\PYG{p}{(}\PYG{n}{nd}\PYG{p}{)} \PYG{p}{:}
        \PYG{n}{adynamic}\PYG{p}{[}\PYG{n}{i}\PYG{p}{]} \PYG{o}{=} \PYG{n}{cppad\PYGZus{}py}\PYG{o}{.}\PYG{n}{a\PYGZus{}double}\PYG{p}{(} \PYG{n}{a\PYGZus{}both}\PYG{p}{[}\PYG{n}{nx} \PYG{o}{+} \PYG{n}{i}\PYG{p}{]} \PYG{p}{)}
    \PYG{c+c1}{\PYGZsh{}}
    \PYG{k}{return} \PYG{p}{(}\PYG{n}{ax}\PYG{p}{,} \PYG{n}{adynamic}\PYG{p}{)}
\end{sphinxVerbatim}

\index{new\_dynamic@\spxentry{new\_dynamic}}\ignorespaces 

\subparagraph{new\_dynamic}
\label{\detokenize{xsrst/more_py:more-py-implementation-new-dynamic}}\label{\detokenize{xsrst/more_py:index-8}}\label{\detokenize{xsrst/more_py:id3}}
The following function declaration was added to the
\sphinxcode{\sphinxupquote{d\_fun}} class
in the \sphinxcode{\sphinxupquote{lib/python/cppad\_py/fun.py}} file:

\begin{sphinxVerbatim}[commandchars=\\\{\}]
\PYG{k}{def} \PYG{n+nf}{new\PYGZus{}dynamic}\PYG{p}{(}\PYG{n+nb+bp}{self}\PYG{p}{,} \PYG{n}{dynamic}\PYG{p}{)} \PYG{p}{:}
    \PYG{k}{return} \PYG{n}{cppad\PYGZus{}py}\PYG{o}{.}\PYG{n}{d\PYGZus{}fun\PYGZus{}new\PYGZus{}dynamic}\PYG{p}{(}\PYG{n+nb+bp}{self}\PYG{o}{.}\PYG{n}{f}\PYG{p}{,} \PYG{n}{dynamic}\PYG{p}{)}
\end{sphinxVerbatim}

A similar declaration was added to the \sphinxcode{\sphinxupquote{a\_fun}} class.

\index{fun\_new\_dynamic.py@\spxentry{fun\_new\_dynamic.py}}\ignorespaces 

\subparagraph{fun\_new\_dynamic.py}
\label{\detokenize{xsrst/more_py:fun-new-dynamic-py}}\label{\detokenize{xsrst/more_py:more-py-implementation-fun-new-dynamic-py}}\label{\detokenize{xsrst/more_py:index-9}}
The implementation of \sphinxcode{\sphinxupquote{d\_fun\_new\_dynamic}} and
\sphinxcode{\sphinxupquote{a\_fun\_new\_dynamic}} were added to the file
\sphinxcode{\sphinxupquote{fun\_new\_dynamic.py}}

\index{\_\_init\_\_.py@\spxentry{\_\_init\_\_.py}}\ignorespaces 

\subparagraph{\_\_init\_\_.py}
\label{\detokenize{xsrst/more_py:init-py}}\label{\detokenize{xsrst/more_py:more-py-implementation-init-py}}\label{\detokenize{xsrst/more_py:index-10}}
The following code was added to the file \sphinxcode{\sphinxupquote{lib/python/cppad\_py/\_\_init\_\_.py}} :

\begin{DUlineblock}{0em}
\item[] \sphinxcode{\sphinxupquote{from cppad\_py.fun\_new\_dynamic import a\_fun\_new\_dynamic}}
\item[] \sphinxcode{\sphinxupquote{from cppad\_py.fun\_new\_dynamic import d\_fun\_new\_dynamic}}
\end{DUlineblock}

\index{example@\spxentry{example}}\ignorespaces 

\subparagraph{Example}
\label{\detokenize{xsrst/more_py:more-py-example}}\label{\detokenize{xsrst/more_py:index-11}}\label{\detokenize{xsrst/more_py:id4}}
The file \sphinxcode{\sphinxupquote{example/python/core/fun\_dynamic\_xam.py}} was added
with the following contents:
{\hyperref[\detokenize{xsrst/fun_dynamic_xam_py:id1}]{\sphinxcrossref{\DUrole{std,std-ref}{fun\_dynamic\_xam\_py}}}}.
In addition, in the file

\begin{DUlineblock}{0em}
\item[]         \sphinxcode{\sphinxupquote{example/python/check\_all.py.in}}
\end{DUlineblock}

the following text was added to the list of python example files.

\begin{DUlineblock}{0em}
\item[]         \sphinxcode{\sphinxupquote{\textquotesingle{}fun\_dynamic\_xam\textquotesingle{}}} ,
\item[] \sphinxcode{\sphinxupquote{\%}}
\end{DUlineblock}

\index{testing@\spxentry{testing}}\ignorespaces 

\subparagraph{Testing}
\label{\detokenize{xsrst/more_py:testing}}\label{\detokenize{xsrst/more_py:more-py-testing}}\label{\detokenize{xsrst/more_py:index-12}}
You must do a git add for all of the new files before running
\sphinxcode{\sphinxupquote{bin/check\_all.sh}}
After all the changes above were implemented,
\sphinxcode{\sphinxupquote{bin/check\_all.sh}} was run and the changes were made
until the warnings and errors were fixed.
The command

\begin{DUlineblock}{0em}
\item[]         \sphinxcode{\sphinxupquote{grep \textquotesingle{}fun\_dynamic\_xam\textquotesingle{} check\_all.log}}
\end{DUlineblock}

was used to make sure that the new python example / test was run.
Note that the python files in \sphinxcode{\sphinxupquote{cppad\_py}} are copies of the
python files in \sphinxcode{\sphinxupquote{lib/python}} .
So when you fix errors during testing, you need to fix the
\sphinxcode{\sphinxupquote{lib/python}} file.
Also note that if a particular step in \sphinxcode{\sphinxupquote{bin/check\_all.sh}} is failing,
you can just re\sphinxhyphen{}run that step to see if a particular fix works.
Once the tests were working, the changes where checked into using
\sphinxcode{\sphinxupquote{git}} .


\bigskip\hrule\bigskip


xsrst input file: \sphinxcode{\sphinxupquote{lib/python/cppad\_py/more\_py.xsrst}}
\begin{enumerate}
\sphinxsetlistlabels{\arabic}{enumi}{enumii}{}{.}%
\item {} 
py\_fun: {\hyperref[\detokenize{xsrst/py_fun:id1}]{\sphinxcrossref{\DUrole{std,std-ref}{3.1.1 Cppad Py AD Functions}}}}

\item {} 
py\_sparse: {\hyperref[\detokenize{xsrst/py_sparse:id1}]{\sphinxcrossref{\DUrole{std,std-ref}{3.1.2 Python Sparsity Routines}}}}

\item {} 
py\_utility: {\hyperref[\detokenize{xsrst/py_utility:id1}]{\sphinxcrossref{\DUrole{std,std-ref}{3.1.3 Python Utilities}}}}

\item {} 
more\_py: {\hyperref[\detokenize{xsrst/more_py:id1}]{\sphinxcrossref{\DUrole{std,std-ref}{3.1.4 Steps To Add More Python Functions}}}}

\end{enumerate}


\bigskip\hrule\bigskip


xsrst input file: \sphinxcode{\sphinxupquote{lib/python/cppad\_py/py\_lib.xsrst}}


\subparagraph{cpp\_lib}
\label{\detokenize{xsrst/cpp_lib:cpp-lib}}\label{\detokenize{xsrst/cpp_lib::doc}}
\index{cpp\_lib@\spxentry{cpp\_lib}}\index{the@\spxentry{the}}\index{c++@\spxentry{c++}}\index{library@\spxentry{library}}\ignorespaces 

\subparagraph{3.2 The C++ Library}
\label{\detokenize{xsrst/cpp_lib:the-c-library}}\label{\detokenize{xsrst/cpp_lib:id1}}\begin{itemize}
\item {} 
{\hyperref[\detokenize{xsrst/cpp_lib:cpp-lib-children}]{\sphinxcrossref{\DUrole{std,std-ref}{Children}}}}

\end{itemize}

\index{children@\spxentry{children}}\ignorespaces 

\subparagraph{Children}
\label{\detokenize{xsrst/cpp_lib:children}}\label{\detokenize{xsrst/cpp_lib:cpp-lib-children}}\label{\detokenize{xsrst/cpp_lib:index-1}}

\subparagraph{a\_double}
\label{\detokenize{xsrst/a_double:a-double}}\label{\detokenize{xsrst/a_double::doc}}
\index{a\_double@\spxentry{a\_double}}\index{cppad@\spxentry{cppad}}\index{py@\spxentry{py}}\index{ad@\spxentry{ad}}\index{scalars@\spxentry{scalars}}\ignorespaces 

\subparagraph{3.2.1 Cppad Py AD Scalars}
\label{\detokenize{xsrst/a_double:cppad-py-ad-scalars}}\label{\detokenize{xsrst/a_double:id1}}\begin{itemize}
\item {} 
{\hyperref[\detokenize{xsrst/a_double:a-double-children}]{\sphinxcrossref{\DUrole{std,std-ref}{Children}}}}

\end{itemize}

\index{children@\spxentry{children}}\ignorespaces 

\subparagraph{Children}
\label{\detokenize{xsrst/a_double:children}}\label{\detokenize{xsrst/a_double:a-double-children}}\label{\detokenize{xsrst/a_double:index-1}}

\subparagraph{a\_double\_ctor}
\label{\detokenize{xsrst/a_double_ctor:a-double-ctor}}\label{\detokenize{xsrst/a_double_ctor::doc}}
\index{a\_double\_ctor@\spxentry{a\_double\_ctor}}\index{the@\spxentry{the}}\index{a\_double@\spxentry{a\_double}}\index{constructor@\spxentry{constructor}}\ignorespaces 

\subparagraph{3.2.1.1 The a\_double Constructor}
\label{\detokenize{xsrst/a_double_ctor:the-a-double-constructor}}\label{\detokenize{xsrst/a_double_ctor:id1}}\begin{itemize}
\item {} 
{\hyperref[\detokenize{xsrst/a_double_ctor:a-double-ctor-syntax}]{\sphinxcrossref{\DUrole{std,std-ref}{Syntax}}}}

\item {} 
{\hyperref[\detokenize{xsrst/a_double_ctor:a-double-ctor-purpose}]{\sphinxcrossref{\DUrole{std,std-ref}{Purpose}}}}

\item {} 
{\hyperref[\detokenize{xsrst/a_double_ctor:a-double-ctor-d}]{\sphinxcrossref{\DUrole{std,std-ref}{d}}}}

\item {} 
{\hyperref[\detokenize{xsrst/a_double_ctor:a-double-ctor-a-other}]{\sphinxcrossref{\DUrole{std,std-ref}{a\_other}}}}

\item {} 
{\hyperref[\detokenize{xsrst/a_double_ctor:a-double-ctor-ad}]{\sphinxcrossref{\DUrole{std,std-ref}{ad}}}}

\item {} 
{\hyperref[\detokenize{xsrst/a_double_ctor:a-double-ctor-example}]{\sphinxcrossref{\DUrole{std,std-ref}{Example}}}}

\end{itemize}

\index{syntax@\spxentry{syntax}}\ignorespaces 

\subparagraph{Syntax}
\label{\detokenize{xsrst/a_double_ctor:syntax}}\label{\detokenize{xsrst/a_double_ctor:a-double-ctor-syntax}}\label{\detokenize{xsrst/a_double_ctor:index-1}}
\begin{DUlineblock}{0em}
\item[] \sphinxstyleemphasis{ad} =  \sphinxcode{\sphinxupquote{cppad\_py.a\_double}} ()
\item[] \sphinxstyleemphasis{ad} =  \sphinxcode{\sphinxupquote{cppad\_py.a\_double}} ( \sphinxstyleemphasis{d} )
\item[] \sphinxstyleemphasis{ad} =  \sphinxcode{\sphinxupquote{cppad\_py.a\_double}} ( \sphinxstyleemphasis{a\_other} )
\end{DUlineblock}

\index{purpose@\spxentry{purpose}}\ignorespaces 

\subparagraph{Purpose}
\label{\detokenize{xsrst/a_double_ctor:purpose}}\label{\detokenize{xsrst/a_double_ctor:a-double-ctor-purpose}}\label{\detokenize{xsrst/a_double_ctor:index-2}}
Creates a \sphinxcode{\sphinxupquote{cppad\_py::a\_double}} object that can be use
to track floating point operations and perform algorithmic differentiation.

\index{d@\spxentry{d}}\ignorespaces 

\subparagraph{d}
\label{\detokenize{xsrst/a_double_ctor:d}}\label{\detokenize{xsrst/a_double_ctor:a-double-ctor-d}}\label{\detokenize{xsrst/a_double_ctor:index-3}}
This argument has prototype

\begin{DUlineblock}{0em}
\item[]         \sphinxcode{\sphinxupquote{const double\&}} \sphinxstyleemphasis{d}
\end{DUlineblock}

The resulting \sphinxstyleemphasis{ad} variable represents
a constant function equal to \sphinxstyleemphasis{d} .

\index{a\_other@\spxentry{a\_other}}\ignorespaces 

\subparagraph{a\_other}
\label{\detokenize{xsrst/a_double_ctor:a-other}}\label{\detokenize{xsrst/a_double_ctor:a-double-ctor-a-other}}\label{\detokenize{xsrst/a_double_ctor:index-4}}
This argument has prototype

\begin{DUlineblock}{0em}
\item[]         \sphinxcode{\sphinxupquote{const a\_double\&}} \sphinxstyleemphasis{a\_other}
\end{DUlineblock}

The resulting \sphinxstyleemphasis{ad} variable is the same function
of the independent variables as \sphinxstyleemphasis{a\_other} .

\index{ad@\spxentry{ad}}\ignorespaces 

\subparagraph{ad}
\label{\detokenize{xsrst/a_double_ctor:ad}}\label{\detokenize{xsrst/a_double_ctor:a-double-ctor-ad}}\label{\detokenize{xsrst/a_double_ctor:index-5}}
is the \sphinxcode{\sphinxupquote{a\_double}} object that is constructed.

\index{example@\spxentry{example}}\ignorespaces 

\subparagraph{Example}
\label{\detokenize{xsrst/a_double_ctor:example}}\label{\detokenize{xsrst/a_double_ctor:a-double-ctor-example}}\label{\detokenize{xsrst/a_double_ctor:index-6}}
All of the other \sphinxcode{\sphinxupquote{a\_double}} examples use an \sphinxcode{\sphinxupquote{a\_double}}
constructor.


\bigskip\hrule\bigskip


xsrst input file: \sphinxcode{\sphinxupquote{lib/cplusplus/a\_double.cpp}}


\subparagraph{a\_double\_unary\_op}
\label{\detokenize{xsrst/a_double_unary_op:a-double-unary-op}}\label{\detokenize{xsrst/a_double_unary_op::doc}}
\index{a\_double\_unary\_op@\spxentry{a\_double\_unary\_op}}\index{a\_double@\spxentry{a\_double}}\index{unary@\spxentry{unary}}\index{plus@\spxentry{plus}}\index{minus@\spxentry{minus}}\ignorespaces 

\subparagraph{3.2.1.2 a\_double Unary Plus and Minus}
\label{\detokenize{xsrst/a_double_unary_op:a-double-unary-plus-and-minus}}\label{\detokenize{xsrst/a_double_unary_op:id1}}\begin{itemize}
\item {} 
{\hyperref[\detokenize{xsrst/a_double_unary_op:a-double-unary-op-syntax}]{\sphinxcrossref{\DUrole{std,std-ref}{Syntax}}}}

\item {} 
{\hyperref[\detokenize{xsrst/a_double_unary_op:a-double-unary-op-ax}]{\sphinxcrossref{\DUrole{std,std-ref}{ax}}}}

\item {} 
{\hyperref[\detokenize{xsrst/a_double_unary_op:a-double-unary-op-ay}]{\sphinxcrossref{\DUrole{std,std-ref}{ay}}}}

\item {} 
{\hyperref[\detokenize{xsrst/a_double_unary_op:a-double-unary-op-example}]{\sphinxcrossref{\DUrole{std,std-ref}{Example}}}}

\end{itemize}

\index{syntax@\spxentry{syntax}}\ignorespaces 

\subparagraph{Syntax}
\label{\detokenize{xsrst/a_double_unary_op:syntax}}\label{\detokenize{xsrst/a_double_unary_op:a-double-unary-op-syntax}}\label{\detokenize{xsrst/a_double_unary_op:index-1}}
\begin{DUlineblock}{0em}
\item[] \sphinxstyleemphasis{ay} = + \sphinxstyleemphasis{ax}
\item[] \sphinxstyleemphasis{ay} = \sphinxhyphen{} \sphinxstyleemphasis{ax}
\end{DUlineblock}

\index{ax@\spxentry{ax}}\ignorespaces 

\subparagraph{ax}
\label{\detokenize{xsrst/a_double_unary_op:ax}}\label{\detokenize{xsrst/a_double_unary_op:a-double-unary-op-ax}}\label{\detokenize{xsrst/a_double_unary_op:index-2}}
This object has prototype

\begin{DUlineblock}{0em}
\item[]         \sphinxcode{\sphinxupquote{const a\_double\&}} \sphinxstyleemphasis{ax}
\end{DUlineblock}

\index{ay@\spxentry{ay}}\ignorespaces 

\subparagraph{ay}
\label{\detokenize{xsrst/a_double_unary_op:ay}}\label{\detokenize{xsrst/a_double_unary_op:a-double-unary-op-ay}}\label{\detokenize{xsrst/a_double_unary_op:index-3}}
If the operator is \sphinxcode{\sphinxupquote{+}} , the result is equal to \sphinxstyleemphasis{ax} .
If it is \sphinxcode{\sphinxupquote{\sphinxhyphen{}}} , the result is the negative of \sphinxstyleemphasis{ax} .


\subparagraph{a\_double\_unary\_op\_xam\_cpp}
\label{\detokenize{xsrst/a_double_unary_op_xam_cpp:a-double-unary-op-xam-cpp}}\label{\detokenize{xsrst/a_double_unary_op_xam_cpp::doc}}
\index{a\_double\_unary\_op\_xam\_cpp@\spxentry{a\_double\_unary\_op\_xam\_cpp}}\index{c++:@\spxentry{c++:}}\index{a\_double@\spxentry{a\_double}}\index{unary@\spxentry{unary}}\index{plus@\spxentry{plus}}\index{minus:@\spxentry{minus:}}\index{example@\spxentry{example}}\index{test@\spxentry{test}}\ignorespaces 

\subparagraph{3.2.1.2.1 C++: a\_double Unary Plus and Minus: Example and Test}
\label{\detokenize{xsrst/a_double_unary_op_xam_cpp:c-a-double-unary-plus-and-minus-example-and-test}}\label{\detokenize{xsrst/a_double_unary_op_xam_cpp:id1}}
\begin{sphinxVerbatim}[commandchars=\\\{\}]
\PYG{c+cp}{\PYGZsh{}}\PYG{c+cp}{ include \PYGZlt{}cstdio\PYGZgt{}}
\PYG{c+cp}{\PYGZsh{}}\PYG{c+cp}{ include \PYGZlt{}cppad}\PYG{c+cp}{/}\PYG{c+cp}{py}\PYG{c+cp}{/}\PYG{c+cp}{cppad\PYGZus{}py.hpp\PYGZgt{}}

\PYG{k+kt}{bool} \PYG{n+nf}{a\PYGZus{}double\PYGZus{}unary\PYGZus{}op\PYGZus{}xam}\PYG{p}{(}\PYG{k+kt}{void}\PYG{p}{)} \PYG{p}{\PYGZob{}}
    \PYG{k}{using} \PYG{n}{cppad\PYGZus{}py}\PYG{o}{:}\PYG{o}{:}\PYG{n}{a\PYGZus{}double}\PYG{p}{;}
    \PYG{c+c1}{//}
    \PYG{c+c1}{// initialize return variable}
    \PYG{k+kt}{bool} \PYG{n}{ok} \PYG{o}{=} \PYG{n+nb}{true}\PYG{p}{;}
    \PYG{c+c1}{//\PYGZhy{}\PYGZhy{}\PYGZhy{}\PYGZhy{}\PYGZhy{}\PYGZhy{}\PYGZhy{}\PYGZhy{}\PYGZhy{}\PYGZhy{}\PYGZhy{}\PYGZhy{}\PYGZhy{}\PYGZhy{}\PYGZhy{}\PYGZhy{}\PYGZhy{}\PYGZhy{}\PYGZhy{}\PYGZhy{}\PYGZhy{}\PYGZhy{}\PYGZhy{}\PYGZhy{}\PYGZhy{}\PYGZhy{}\PYGZhy{}\PYGZhy{}\PYGZhy{}\PYGZhy{}\PYGZhy{}\PYGZhy{}\PYGZhy{}\PYGZhy{}\PYGZhy{}\PYGZhy{}\PYGZhy{}\PYGZhy{}\PYGZhy{}\PYGZhy{}\PYGZhy{}\PYGZhy{}\PYGZhy{}\PYGZhy{}\PYGZhy{}\PYGZhy{}\PYGZhy{}\PYGZhy{}\PYGZhy{}\PYGZhy{}\PYGZhy{}\PYGZhy{}\PYGZhy{}\PYGZhy{}\PYGZhy{}\PYGZhy{}\PYGZhy{}\PYGZhy{}\PYGZhy{}\PYGZhy{}\PYGZhy{}\PYGZhy{}\PYGZhy{}\PYGZhy{}\PYGZhy{}\PYGZhy{}\PYGZhy{}\PYGZhy{}\PYGZhy{}\PYGZhy{}\PYGZhy{}\PYGZhy{}}
    \PYG{n}{a\PYGZus{}double} \PYG{n}{a2}\PYG{p}{(}\PYG{l+m+mf}{2.0}\PYG{p}{)}\PYG{p}{;}
    \PYG{n}{a\PYGZus{}double} \PYG{n}{aplus2} \PYG{o}{=} \PYG{o}{+} \PYG{n}{a2} \PYG{p}{;}
    \PYG{n}{a\PYGZus{}double} \PYG{n}{aminus2} \PYG{o}{=} \PYG{o}{\PYGZhy{}} \PYG{n}{a2} \PYG{p}{;}
    \PYG{c+c1}{//}
    \PYG{n}{ok} \PYG{o}{=} \PYG{n}{ok} \PYG{o}{\PYGZam{}}\PYG{o}{\PYGZam{}} \PYG{n}{aplus2} \PYG{o}{=}\PYG{o}{=} \PYG{l+m+mf}{2.0}\PYG{p}{;}
    \PYG{n}{ok} \PYG{o}{=} \PYG{n}{ok} \PYG{o}{\PYGZam{}}\PYG{o}{\PYGZam{}} \PYG{n}{aminus2} \PYG{o}{=}\PYG{o}{=} \PYG{o}{\PYGZhy{}}\PYG{l+m+mf}{2.0}\PYG{p}{;}
    \PYG{c+c1}{//}
    \PYG{k}{return}\PYG{p}{(} \PYG{n}{ok} \PYG{p}{)}\PYG{p}{;}
\PYG{p}{\PYGZcb{}}
\end{sphinxVerbatim}


\bigskip\hrule\bigskip


xsrst input file: \sphinxcode{\sphinxupquote{example/cplusplus/a\_double\_unary\_op\_xam.cpp}}


\subparagraph{a\_double\_unary\_op\_xam\_py}
\label{\detokenize{xsrst/a_double_unary_op_xam_py:a-double-unary-op-xam-py}}\label{\detokenize{xsrst/a_double_unary_op_xam_py::doc}}
\index{a\_double\_unary\_op\_xam\_py@\spxentry{a\_double\_unary\_op\_xam\_py}}\index{python:@\spxentry{python:}}\index{a\_double@\spxentry{a\_double}}\index{unary@\spxentry{unary}}\index{plus@\spxentry{plus}}\index{minus:@\spxentry{minus:}}\index{example@\spxentry{example}}\index{test@\spxentry{test}}\ignorespaces 

\subparagraph{3.2.1.2.2 Python: a\_double Unary Plus and Minus: Example and Test}
\label{\detokenize{xsrst/a_double_unary_op_xam_py:python-a-double-unary-plus-and-minus-example-and-test}}\label{\detokenize{xsrst/a_double_unary_op_xam_py:id1}}
\begin{sphinxVerbatim}[commandchars=\\\{\}]
\PYG{k}{def} \PYG{n+nf}{a\PYGZus{}double\PYGZus{}unary\PYGZus{}op\PYGZus{}xam}\PYG{p}{(}\PYG{p}{)} \PYG{p}{:}
    \PYG{c+c1}{\PYGZsh{}}
    \PYG{k+kn}{import} \PYG{n+nn}{numpy}
    \PYG{k+kn}{import} \PYG{n+nn}{cppad\PYGZus{}py}
    \PYG{c+c1}{\PYGZsh{}}
    \PYG{c+c1}{\PYGZsh{} initialize return variable}
    \PYG{n}{ok} \PYG{o}{=} \PYG{k+kc}{True}
    \PYG{c+c1}{\PYGZsh{} \PYGZhy{}\PYGZhy{}\PYGZhy{}\PYGZhy{}\PYGZhy{}\PYGZhy{}\PYGZhy{}\PYGZhy{}\PYGZhy{}\PYGZhy{}\PYGZhy{}\PYGZhy{}\PYGZhy{}\PYGZhy{}\PYGZhy{}\PYGZhy{}\PYGZhy{}\PYGZhy{}\PYGZhy{}\PYGZhy{}\PYGZhy{}\PYGZhy{}\PYGZhy{}\PYGZhy{}\PYGZhy{}\PYGZhy{}\PYGZhy{}\PYGZhy{}\PYGZhy{}\PYGZhy{}\PYGZhy{}\PYGZhy{}\PYGZhy{}\PYGZhy{}\PYGZhy{}\PYGZhy{}\PYGZhy{}\PYGZhy{}\PYGZhy{}\PYGZhy{}\PYGZhy{}\PYGZhy{}\PYGZhy{}\PYGZhy{}\PYGZhy{}\PYGZhy{}\PYGZhy{}\PYGZhy{}\PYGZhy{}\PYGZhy{}\PYGZhy{}\PYGZhy{}\PYGZhy{}\PYGZhy{}\PYGZhy{}\PYGZhy{}\PYGZhy{}\PYGZhy{}\PYGZhy{}\PYGZhy{}\PYGZhy{}\PYGZhy{}\PYGZhy{}\PYGZhy{}\PYGZhy{}\PYGZhy{}\PYGZhy{}\PYGZhy{}\PYGZhy{}}
    \PYG{n}{a2} \PYG{o}{=} \PYG{n}{cppad\PYGZus{}py}\PYG{o}{.}\PYG{n}{a\PYGZus{}double}\PYG{p}{(}\PYG{l+m+mf}{2.0}\PYG{p}{)}
    \PYG{n}{aplus2} \PYG{o}{=} \PYG{o}{+} \PYG{n}{a2}
    \PYG{n}{a2\PYGZus{}minus} \PYG{o}{=} \PYG{o}{\PYGZhy{}} \PYG{n}{a2}
    \PYG{c+c1}{\PYGZsh{}}
    \PYG{n}{ok} \PYG{o}{=} \PYG{n}{ok} \PYG{o+ow}{and} \PYG{n}{aplus2} \PYG{o}{==} \PYG{l+m+mf}{2.0}
    \PYG{n}{ok} \PYG{o}{=} \PYG{n}{ok} \PYG{o+ow}{and} \PYG{n}{a2\PYGZus{}minus} \PYG{o}{==} \PYG{o}{\PYGZhy{}}\PYG{l+m+mf}{2.0}
    \PYG{c+c1}{\PYGZsh{}}
    \PYG{k}{return}\PYG{p}{(} \PYG{n}{ok} \PYG{p}{)}
\PYG{c+c1}{\PYGZsh{}}
\end{sphinxVerbatim}


\bigskip\hrule\bigskip


xsrst input file: \sphinxcode{\sphinxupquote{example/python/core/a\_double\_unary\_op\_xam.py}}

\index{example@\spxentry{example}}\ignorespaces 

\subparagraph{Example}
\label{\detokenize{xsrst/a_double_unary_op:example}}\label{\detokenize{xsrst/a_double_unary_op:a-double-unary-op-example}}\label{\detokenize{xsrst/a_double_unary_op:index-4}}
{\hyperref[\detokenize{xsrst/a_double_unary_op_xam_cpp:id1}]{\sphinxcrossref{\DUrole{std,std-ref}{c++}}}},
{\hyperref[\detokenize{xsrst/a_double_unary_op_xam_py:id1}]{\sphinxcrossref{\DUrole{std,std-ref}{python}}}}.


\bigskip\hrule\bigskip


xsrst input file: \sphinxcode{\sphinxupquote{lib/cplusplus/a\_double.cpp}}


\subparagraph{a\_double\_property}
\label{\detokenize{xsrst/a_double_property:a-double-property}}\label{\detokenize{xsrst/a_double_property::doc}}
\index{a\_double\_property@\spxentry{a\_double\_property}}\index{properties@\spxentry{properties}}\index{an@\spxentry{an}}\index{a\_double@\spxentry{a\_double}}\index{object@\spxentry{object}}\ignorespaces 

\subparagraph{3.2.1.3 Properties of an a\_double Object}
\label{\detokenize{xsrst/a_double_property:properties-of-an-a-double-object}}\label{\detokenize{xsrst/a_double_property:id1}}\begin{itemize}
\item {} 
{\hyperref[\detokenize{xsrst/a_double_property:a-double-property-syntax}]{\sphinxcrossref{\DUrole{std,std-ref}{Syntax}}}}

\item {} 
{\hyperref[\detokenize{xsrst/a_double_property:a-double-property-ad}]{\sphinxcrossref{\DUrole{std,std-ref}{ad}}}}

\item {} \begin{description}
\item[{{\hyperref[\detokenize{xsrst/a_double_property:a-double-property-value}]{\sphinxcrossref{\DUrole{std,std-ref}{value}}}}}] \leavevmode\begin{itemize}
\item {} 
{\hyperref[\detokenize{xsrst/a_double_property:a-double-property-value-restriction}]{\sphinxcrossref{\DUrole{std,std-ref}{Restriction}}}}

\end{itemize}

\end{description}

\item {} 
{\hyperref[\detokenize{xsrst/a_double_property:a-double-property-parameter}]{\sphinxcrossref{\DUrole{std,std-ref}{parameter}}}}

\item {} 
{\hyperref[\detokenize{xsrst/a_double_property:a-double-property-variable}]{\sphinxcrossref{\DUrole{std,std-ref}{variable}}}}

\item {} 
{\hyperref[\detokenize{xsrst/a_double_property:a-double-property-near-equal}]{\sphinxcrossref{\DUrole{std,std-ref}{near\_equal}}}}

\item {} 
{\hyperref[\detokenize{xsrst/a_double_property:a-double-property-var2par}]{\sphinxcrossref{\DUrole{std,std-ref}{var2par}}}}

\item {} 
{\hyperref[\detokenize{xsrst/a_double_property:a-double-property-example}]{\sphinxcrossref{\DUrole{std,std-ref}{Example}}}}

\end{itemize}

\index{syntax@\spxentry{syntax}}\ignorespaces 

\subparagraph{Syntax}
\label{\detokenize{xsrst/a_double_property:syntax}}\label{\detokenize{xsrst/a_double_property:a-double-property-syntax}}\label{\detokenize{xsrst/a_double_property:index-1}}
\begin{DUlineblock}{0em}
\item[] \sphinxstyleemphasis{d} = \sphinxstyleemphasis{ad} . \sphinxcode{\sphinxupquote{value}} ()
\item[] \sphinxstyleemphasis{p} = \sphinxstyleemphasis{ad} . \sphinxcode{\sphinxupquote{parameter}} ()
\item[] \sphinxstyleemphasis{v} = \sphinxstyleemphasis{ad} . \sphinxcode{\sphinxupquote{variable}} ()
\item[] \sphinxstyleemphasis{e} = \sphinxstyleemphasis{ad} . \sphinxcode{\sphinxupquote{near\_equal}} ( \sphinxstyleemphasis{aother} )
\item[] \sphinxstyleemphasis{ap} = \sphinxstyleemphasis{ad} . \sphinxcode{\sphinxupquote{var2par}} ()
\end{DUlineblock}

\index{ad@\spxentry{ad}}\ignorespaces 

\subparagraph{ad}
\label{\detokenize{xsrst/a_double_property:ad}}\label{\detokenize{xsrst/a_double_property:a-double-property-ad}}\label{\detokenize{xsrst/a_double_property:index-2}}
This object has prototype

\begin{DUlineblock}{0em}
\item[]         \sphinxcode{\sphinxupquote{const a\_double\&}} \sphinxstyleemphasis{ad}
\end{DUlineblock}

\index{value@\spxentry{value}}\ignorespaces 

\subparagraph{value}
\label{\detokenize{xsrst/a_double_property:value}}\label{\detokenize{xsrst/a_double_property:a-double-property-value}}\label{\detokenize{xsrst/a_double_property:index-3}}
The result \sphinxstyleemphasis{d} has prototype

\begin{DUlineblock}{0em}
\item[]         \sphinxcode{\sphinxupquote{double}} \sphinxstyleemphasis{d}
\end{DUlineblock}

It is the value of \sphinxstyleemphasis{ad} , as a constant function.

\index{restriction@\spxentry{restriction}}\ignorespaces 

\subparagraph{Restriction}
\label{\detokenize{xsrst/a_double_property:restriction}}\label{\detokenize{xsrst/a_double_property:a-double-property-value-restriction}}\label{\detokenize{xsrst/a_double_property:index-4}}
The object \sphinxstyleemphasis{ad} must not depend on the
{\hyperref[\detokenize{xsrst/cpp_independent:id1}]{\sphinxcrossref{\DUrole{std,std-ref}{independent}}}}
variables when \sphinxstyleemphasis{ad} . \sphinxcode{\sphinxupquote{value}} () is called.
If it does depend on the independent variables,
you will have to wait until the current recording is terminated
before you can access its value; see
{\hyperref[\detokenize{xsrst/a_double_property:a-double-property-var2par}]{\sphinxcrossref{\DUrole{std,std-ref}{var2par}}}} below.

\index{parameter@\spxentry{parameter}}\ignorespaces 

\subparagraph{parameter}
\label{\detokenize{xsrst/a_double_property:parameter}}\label{\detokenize{xsrst/a_double_property:a-double-property-parameter}}\label{\detokenize{xsrst/a_double_property:index-5}}
The result \sphinxstyleemphasis{p} has prototype

\begin{DUlineblock}{0em}
\item[]         \sphinxcode{\sphinxupquote{bool}} \sphinxstyleemphasis{p}
\end{DUlineblock}

It is true if \sphinxstyleemphasis{ad} represent a constant functions; i.e.,
\sphinxstyleemphasis{ad} not depend on the independent variables.

\index{variable@\spxentry{variable}}\ignorespaces 

\subparagraph{variable}
\label{\detokenize{xsrst/a_double_property:variable}}\label{\detokenize{xsrst/a_double_property:a-double-property-variable}}\label{\detokenize{xsrst/a_double_property:index-6}}
The result \sphinxstyleemphasis{v} has prototype

\begin{DUlineblock}{0em}
\item[]         \sphinxcode{\sphinxupquote{bool}} \sphinxstyleemphasis{v}
\end{DUlineblock}

It is true if \sphinxstyleemphasis{ad} is not a constant function; i.e.,
\sphinxstyleemphasis{ad} depends on the independent variables.

\index{near\_equal@\spxentry{near\_equal}}\ignorespaces 

\subparagraph{near\_equal}
\label{\detokenize{xsrst/a_double_property:near-equal}}\label{\detokenize{xsrst/a_double_property:a-double-property-near-equal}}\label{\detokenize{xsrst/a_double_property:index-7}}
The argument \sphinxstyleemphasis{aother} ,
and the result \sphinxstyleemphasis{e} , have prototype

\begin{DUlineblock}{0em}
\item[]         \sphinxcode{\sphinxupquote{const a\_double\&}} \sphinxstyleemphasis{aother}
\item[]         \sphinxcode{\sphinxupquote{bool}} \sphinxstyleemphasis{e}
\end{DUlineblock}

The result is true if \sphinxstyleemphasis{ad} is nearly equal to \sphinxstyleemphasis{aother} .
To be specific, the result is
\begin{equation*}
\begin{split}| d - o | \leq 100 \; \varepsilon \; ( |d| + |o| )\end{split}
\end{equation*}
where \sphinxstyleemphasis{d} and \sphinxstyleemphasis{o} are the value corresponding to
\sphinxstyleemphasis{ad} and \sphinxstyleemphasis{aother} and
\(\varepsilon\) is machine epsilon corresponding
to the type \sphinxcode{\sphinxupquote{double}} .

\index{var2par@\spxentry{var2par}}\ignorespaces 

\subparagraph{var2par}
\label{\detokenize{xsrst/a_double_property:var2par}}\label{\detokenize{xsrst/a_double_property:a-double-property-var2par}}\label{\detokenize{xsrst/a_double_property:index-8}}
The result has prototype

\begin{DUlineblock}{0em}
\item[]         \sphinxcode{\sphinxupquote{a\_double}} \sphinxstyleemphasis{ap}
\end{DUlineblock}

It has the same value as \sphinxstyleemphasis{ad} and is sure to be a parameter
( \sphinxstyleemphasis{ad} may or may not be a variable).
This can be useful when you want to access the value of \sphinxstyleemphasis{ad}
while is a variable; {\hyperref[\detokenize{xsrst/a_double_property:a-double-property-value}]{\sphinxcrossref{\DUrole{std,std-ref}{value}}}} above.


\subparagraph{a\_double\_property\_xam\_cpp}
\label{\detokenize{xsrst/a_double_property_xam_cpp:a-double-property-xam-cpp}}\label{\detokenize{xsrst/a_double_property_xam_cpp::doc}}
\index{a\_double\_property\_xam\_cpp@\spxentry{a\_double\_property\_xam\_cpp}}\index{c++:@\spxentry{c++:}}\index{a\_double@\spxentry{a\_double}}\index{properties:@\spxentry{properties:}}\index{example@\spxentry{example}}\index{test@\spxentry{test}}\ignorespaces 

\subparagraph{3.2.1.3.1 C++: a\_double Properties: Example and Test}
\label{\detokenize{xsrst/a_double_property_xam_cpp:c-a-double-properties-example-and-test}}\label{\detokenize{xsrst/a_double_property_xam_cpp:id1}}
\begin{sphinxVerbatim}[commandchars=\\\{\}]
\PYG{c+cp}{\PYGZsh{}}\PYG{c+cp}{ include \PYGZlt{}cstdio\PYGZgt{}}
\PYG{c+cp}{\PYGZsh{}}\PYG{c+cp}{ include \PYGZlt{}cppad}\PYG{c+cp}{/}\PYG{c+cp}{py}\PYG{c+cp}{/}\PYG{c+cp}{cppad\PYGZus{}py.hpp\PYGZgt{}}

\PYG{k+kt}{bool} \PYG{n+nf}{a\PYGZus{}double\PYGZus{}property\PYGZus{}xam}\PYG{p}{(}\PYG{k+kt}{void}\PYG{p}{)} \PYG{p}{\PYGZob{}}
    \PYG{k}{using} \PYG{n}{cppad\PYGZus{}py}\PYG{o}{:}\PYG{o}{:}\PYG{n}{a\PYGZus{}double}\PYG{p}{;}
    \PYG{c+c1}{//}
    \PYG{c+c1}{// initialize return variable}
    \PYG{k+kt}{bool} \PYG{n}{ok} \PYG{o}{=} \PYG{n+nb}{true}\PYG{p}{;}
    \PYG{c+c1}{//\PYGZhy{}\PYGZhy{}\PYGZhy{}\PYGZhy{}\PYGZhy{}\PYGZhy{}\PYGZhy{}\PYGZhy{}\PYGZhy{}\PYGZhy{}\PYGZhy{}\PYGZhy{}\PYGZhy{}\PYGZhy{}\PYGZhy{}\PYGZhy{}\PYGZhy{}\PYGZhy{}\PYGZhy{}\PYGZhy{}\PYGZhy{}\PYGZhy{}\PYGZhy{}\PYGZhy{}\PYGZhy{}\PYGZhy{}\PYGZhy{}\PYGZhy{}\PYGZhy{}\PYGZhy{}\PYGZhy{}\PYGZhy{}\PYGZhy{}\PYGZhy{}\PYGZhy{}\PYGZhy{}\PYGZhy{}\PYGZhy{}\PYGZhy{}\PYGZhy{}\PYGZhy{}\PYGZhy{}\PYGZhy{}\PYGZhy{}\PYGZhy{}\PYGZhy{}\PYGZhy{}\PYGZhy{}\PYGZhy{}\PYGZhy{}\PYGZhy{}\PYGZhy{}\PYGZhy{}\PYGZhy{}\PYGZhy{}\PYGZhy{}\PYGZhy{}\PYGZhy{}\PYGZhy{}\PYGZhy{}\PYGZhy{}\PYGZhy{}\PYGZhy{}\PYGZhy{}\PYGZhy{}\PYGZhy{}\PYGZhy{}\PYGZhy{}\PYGZhy{}\PYGZhy{}\PYGZhy{}\PYGZhy{}}
    \PYG{n}{a\PYGZus{}double} \PYG{n}{a3}\PYG{p}{(}\PYG{l+m+mf}{3.0}\PYG{p}{)}\PYG{p}{;}
    \PYG{c+c1}{//}
    \PYG{n}{ok} \PYG{o}{=} \PYG{n}{ok} \PYG{o}{\PYGZam{}}\PYG{o}{\PYGZam{}} \PYG{n}{a3}   \PYG{o}{=}\PYG{o}{=} \PYG{l+m+mf}{3.0}\PYG{p}{;}
    \PYG{n}{ok} \PYG{o}{=} \PYG{n}{ok} \PYG{o}{\PYGZam{}}\PYG{o}{\PYGZam{}} \PYG{n}{a3}\PYG{p}{.}\PYG{n}{parameter}\PYG{p}{(}\PYG{p}{)}\PYG{p}{;}
    \PYG{n}{ok} \PYG{o}{=} \PYG{n}{ok} \PYG{o}{\PYGZam{}}\PYG{o}{\PYGZam{}} \PYG{o}{!} \PYG{n}{a3}\PYG{p}{.}\PYG{n}{variable}\PYG{p}{(}\PYG{p}{)}\PYG{p}{;}
    \PYG{c+c1}{//}
    \PYG{c+c1}{// near\PYGZus{}equal}
    \PYG{n}{a\PYGZus{}double} \PYG{n}{r3} \PYG{o}{=} \PYG{n}{a3}\PYG{p}{.}\PYG{n}{sqrt}\PYG{p}{(}\PYG{p}{)} \PYG{p}{;}
    \PYG{n}{ok} \PYG{o}{=} \PYG{n}{ok} \PYG{o}{\PYGZam{}}\PYG{o}{\PYGZam{}} \PYG{n}{a3}\PYG{p}{.}\PYG{n}{near\PYGZus{}equal}\PYG{p}{(} \PYG{n}{r3} \PYG{o}{*} \PYG{n}{r3}\PYG{p}{)} \PYG{p}{;}\PYG{p}{;}
    \PYG{c+c1}{//}
    \PYG{c+c1}{// var2par}
    \PYG{n}{a\PYGZus{}double} \PYG{n}{p3} \PYG{o}{=} \PYG{n}{a3}\PYG{p}{.}\PYG{n}{var2par}\PYG{p}{(}\PYG{p}{)}\PYG{p}{;}
    \PYG{n}{ok} \PYG{o}{\PYGZam{}}\PYG{o}{=} \PYG{n}{p3}\PYG{p}{.}\PYG{n}{value}\PYG{p}{(}\PYG{p}{)} \PYG{o}{=}\PYG{o}{=} \PYG{l+m+mf}{3.0}\PYG{p}{;}
    \PYG{c+c1}{//}
    \PYG{k}{return}\PYG{p}{(} \PYG{n}{ok} \PYG{p}{)}\PYG{p}{;}
\PYG{p}{\PYGZcb{}}
\end{sphinxVerbatim}


\bigskip\hrule\bigskip


xsrst input file: \sphinxcode{\sphinxupquote{example/cplusplus/a\_double\_property\_xam.cpp}}


\subparagraph{a\_double\_property\_xam\_py}
\label{\detokenize{xsrst/a_double_property_xam_py:a-double-property-xam-py}}\label{\detokenize{xsrst/a_double_property_xam_py::doc}}
\index{a\_double\_property\_xam\_py@\spxentry{a\_double\_property\_xam\_py}}\index{python:@\spxentry{python:}}\index{a\_double@\spxentry{a\_double}}\index{properties:@\spxentry{properties:}}\index{example@\spxentry{example}}\index{test@\spxentry{test}}\ignorespaces 

\subparagraph{3.2.1.3.2 Python: a\_double Properties: Example and Test}
\label{\detokenize{xsrst/a_double_property_xam_py:python-a-double-properties-example-and-test}}\label{\detokenize{xsrst/a_double_property_xam_py:id1}}
\begin{sphinxVerbatim}[commandchars=\\\{\}]
\PYG{k}{def} \PYG{n+nf}{a\PYGZus{}double\PYGZus{}property\PYGZus{}xam}\PYG{p}{(}\PYG{p}{)} \PYG{p}{:}
    \PYG{c+c1}{\PYGZsh{}}
    \PYG{k+kn}{import} \PYG{n+nn}{numpy}
    \PYG{k+kn}{import} \PYG{n+nn}{cppad\PYGZus{}py}
    \PYG{c+c1}{\PYGZsh{}}
    \PYG{c+c1}{\PYGZsh{} initialize return variable}
    \PYG{n}{ok} \PYG{o}{=} \PYG{k+kc}{True}
    \PYG{c+c1}{\PYGZsh{} \PYGZhy{}\PYGZhy{}\PYGZhy{}\PYGZhy{}\PYGZhy{}\PYGZhy{}\PYGZhy{}\PYGZhy{}\PYGZhy{}\PYGZhy{}\PYGZhy{}\PYGZhy{}\PYGZhy{}\PYGZhy{}\PYGZhy{}\PYGZhy{}\PYGZhy{}\PYGZhy{}\PYGZhy{}\PYGZhy{}\PYGZhy{}\PYGZhy{}\PYGZhy{}\PYGZhy{}\PYGZhy{}\PYGZhy{}\PYGZhy{}\PYGZhy{}\PYGZhy{}\PYGZhy{}\PYGZhy{}\PYGZhy{}\PYGZhy{}\PYGZhy{}\PYGZhy{}\PYGZhy{}\PYGZhy{}\PYGZhy{}\PYGZhy{}\PYGZhy{}\PYGZhy{}\PYGZhy{}\PYGZhy{}\PYGZhy{}\PYGZhy{}\PYGZhy{}\PYGZhy{}\PYGZhy{}\PYGZhy{}\PYGZhy{}\PYGZhy{}\PYGZhy{}\PYGZhy{}\PYGZhy{}\PYGZhy{}\PYGZhy{}\PYGZhy{}\PYGZhy{}\PYGZhy{}\PYGZhy{}\PYGZhy{}\PYGZhy{}\PYGZhy{}\PYGZhy{}\PYGZhy{}\PYGZhy{}\PYGZhy{}\PYGZhy{}\PYGZhy{}}
    \PYG{n}{a3} \PYG{o}{=} \PYG{n}{cppad\PYGZus{}py}\PYG{o}{.}\PYG{n}{a\PYGZus{}double}\PYG{p}{(}\PYG{l+m+mf}{3.0}\PYG{p}{)}
    \PYG{c+c1}{\PYGZsh{}}
    \PYG{n}{ok} \PYG{o}{=} \PYG{n}{ok} \PYG{o+ow}{and} \PYG{n}{a3}   \PYG{o}{==} \PYG{l+m+mf}{3.0}
    \PYG{n}{ok} \PYG{o}{=} \PYG{n}{ok} \PYG{o+ow}{and} \PYG{n}{a3}\PYG{o}{.}\PYG{n}{parameter}\PYG{p}{(}\PYG{p}{)}
    \PYG{n}{ok} \PYG{o}{=} \PYG{n}{ok} \PYG{o+ow}{and} \PYG{o+ow}{not} \PYG{n}{a3}\PYG{o}{.}\PYG{n}{variable}\PYG{p}{(}\PYG{p}{)}
    \PYG{c+c1}{\PYGZsh{}}
    \PYG{c+c1}{\PYGZsh{} near\PYGZus{}equal}
    \PYG{n}{r3} \PYG{o}{=} \PYG{n}{a3}\PYG{o}{.}\PYG{n}{sqrt}\PYG{p}{(}\PYG{p}{)}
    \PYG{n}{ok} \PYG{o}{=} \PYG{n}{ok} \PYG{o+ow}{and} \PYG{n}{a3}\PYG{o}{.}\PYG{n}{near\PYGZus{}equal}\PYG{p}{(} \PYG{n}{r3} \PYG{o}{*} \PYG{n}{r3}\PYG{p}{)}
    \PYG{c+c1}{\PYGZsh{}}
    \PYG{c+c1}{\PYGZsh{} var2par}
    \PYG{n}{p3}  \PYG{o}{=} \PYG{n}{a3}\PYG{o}{.}\PYG{n}{var2par}\PYG{p}{(}\PYG{p}{)}
    \PYG{n}{ok}  \PYG{o}{=} \PYG{n}{ok} \PYG{o+ow}{and} \PYG{n}{p3}\PYG{o}{.}\PYG{n}{value}\PYG{p}{(}\PYG{p}{)} \PYG{o}{==} \PYG{l+m+mf}{3.0}
    \PYG{c+c1}{\PYGZsh{}}
    \PYG{k}{return}\PYG{p}{(} \PYG{n}{ok} \PYG{p}{)}
\PYG{c+c1}{\PYGZsh{}}
\end{sphinxVerbatim}


\bigskip\hrule\bigskip


xsrst input file: \sphinxcode{\sphinxupquote{example/python/core/a\_double\_property\_xam.py}}

\index{example@\spxentry{example}}\ignorespaces 

\subparagraph{Example}
\label{\detokenize{xsrst/a_double_property:example}}\label{\detokenize{xsrst/a_double_property:a-double-property-example}}\label{\detokenize{xsrst/a_double_property:index-9}}
{\hyperref[\detokenize{xsrst/a_double_property_xam_cpp:id1}]{\sphinxcrossref{\DUrole{std,std-ref}{c++}}}},
{\hyperref[\detokenize{xsrst/a_double_property_xam_py:id1}]{\sphinxcrossref{\DUrole{std,std-ref}{python}}}}.


\bigskip\hrule\bigskip


xsrst input file: \sphinxcode{\sphinxupquote{lib/cplusplus/a\_double.cpp}}


\subparagraph{a\_double\_binary}
\label{\detokenize{xsrst/a_double_binary:a-double-binary}}\label{\detokenize{xsrst/a_double_binary::doc}}
\index{a\_double\_binary@\spxentry{a\_double\_binary}}\index{a\_double@\spxentry{a\_double}}\index{binary@\spxentry{binary}}\index{operators@\spxentry{operators}}\index{with@\spxentry{with}}\index{an@\spxentry{an}}\index{ad@\spxentry{ad}}\index{result@\spxentry{result}}\ignorespaces 

\subparagraph{3.2.1.4 a\_double Binary Operators with an AD Result}
\label{\detokenize{xsrst/a_double_binary:a-double-binary-operators-with-an-ad-result}}\label{\detokenize{xsrst/a_double_binary:id1}}\begin{itemize}
\item {} 
{\hyperref[\detokenize{xsrst/a_double_binary:a-double-binary-syntax}]{\sphinxcrossref{\DUrole{std,std-ref}{Syntax}}}}

\item {} 
{\hyperref[\detokenize{xsrst/a_double_binary:a-double-binary-op}]{\sphinxcrossref{\DUrole{std,std-ref}{op}}}}

\item {} 
{\hyperref[\detokenize{xsrst/a_double_binary:a-double-binary-ax}]{\sphinxcrossref{\DUrole{std,std-ref}{ax}}}}

\item {} 
{\hyperref[\detokenize{xsrst/a_double_binary:a-double-binary-ay}]{\sphinxcrossref{\DUrole{std,std-ref}{ay}}}}

\item {} 
{\hyperref[\detokenize{xsrst/a_double_binary:a-double-binary-y}]{\sphinxcrossref{\DUrole{std,std-ref}{y}}}}

\item {} 
{\hyperref[\detokenize{xsrst/a_double_binary:a-double-binary-az}]{\sphinxcrossref{\DUrole{std,std-ref}{az}}}}

\item {} 
{\hyperref[\detokenize{xsrst/a_double_binary:a-double-binary-example}]{\sphinxcrossref{\DUrole{std,std-ref}{Example}}}}

\end{itemize}

\index{syntax@\spxentry{syntax}}\ignorespaces 

\subparagraph{Syntax}
\label{\detokenize{xsrst/a_double_binary:syntax}}\label{\detokenize{xsrst/a_double_binary:a-double-binary-syntax}}\label{\detokenize{xsrst/a_double_binary:index-1}}
\begin{DUlineblock}{0em}
\item[] \sphinxstyleemphasis{az} = \sphinxstyleemphasis{ax} \sphinxstyleemphasis{op} \sphinxstyleemphasis{ay}
\item[] \sphinxstyleemphasis{az} = \sphinxstyleemphasis{ax} \sphinxstyleemphasis{op} \sphinxstyleemphasis{y}
\item[] \sphinxstyleemphasis{az} = \sphinxstyleemphasis{y} \sphinxstyleemphasis{op} \sphinxstyleemphasis{ax}
\end{DUlineblock}

\index{op@\spxentry{op}}\ignorespaces 

\subparagraph{op}
\label{\detokenize{xsrst/a_double_binary:op}}\label{\detokenize{xsrst/a_double_binary:a-double-binary-op}}\label{\detokenize{xsrst/a_double_binary:index-2}}
The binary operator \sphinxstyleemphasis{op} is one of the following:
addition \sphinxcode{\sphinxupquote{+}} ,
subtraction \sphinxcode{\sphinxupquote{\sphinxhyphen{}}} ,
multiplication \sphinxcode{\sphinxupquote{*}} ,
division \sphinxcode{\sphinxupquote{/}} , or
exponentiation \sphinxcode{\sphinxupquote{**}} .
Note that exponentiation is a function is c++ and an operator in python.

\index{ax@\spxentry{ax}}\ignorespaces 

\subparagraph{ax}
\label{\detokenize{xsrst/a_double_binary:ax}}\label{\detokenize{xsrst/a_double_binary:a-double-binary-ax}}\label{\detokenize{xsrst/a_double_binary:index-3}}
This object has prototype

\begin{DUlineblock}{0em}
\item[]         \sphinxcode{\sphinxupquote{const a\_double\&}} \sphinxstyleemphasis{ax}
\end{DUlineblock}

\index{ay@\spxentry{ay}}\ignorespaces 

\subparagraph{ay}
\label{\detokenize{xsrst/a_double_binary:ay}}\label{\detokenize{xsrst/a_double_binary:a-double-binary-ay}}\label{\detokenize{xsrst/a_double_binary:index-4}}
This object has prototype

\begin{DUlineblock}{0em}
\item[]         \sphinxcode{\sphinxupquote{const a\_double\&}} \sphinxstyleemphasis{ay}
\end{DUlineblock}

\index{y@\spxentry{y}}\ignorespaces 

\subparagraph{y}
\label{\detokenize{xsrst/a_double_binary:y}}\label{\detokenize{xsrst/a_double_binary:a-double-binary-y}}\label{\detokenize{xsrst/a_double_binary:index-5}}
This object has prototype

\begin{DUlineblock}{0em}
\item[]         \sphinxcode{\sphinxupquote{const double\&}} \sphinxstyleemphasis{y}
\end{DUlineblock}

\index{az@\spxentry{az}}\ignorespaces 

\subparagraph{az}
\label{\detokenize{xsrst/a_double_binary:az}}\label{\detokenize{xsrst/a_double_binary:a-double-binary-az}}\label{\detokenize{xsrst/a_double_binary:index-6}}
The result has prototype

\begin{DUlineblock}{0em}
\item[]         \sphinxcode{\sphinxupquote{a\_double}} \sphinxstyleemphasis{az}
\end{DUlineblock}


\subparagraph{a\_double\_binary\_xam\_cpp}
\label{\detokenize{xsrst/a_double_binary_xam_cpp:a-double-binary-xam-cpp}}\label{\detokenize{xsrst/a_double_binary_xam_cpp::doc}}
\index{a\_double\_binary\_xam\_cpp@\spxentry{a\_double\_binary\_xam\_cpp}}\index{c++:@\spxentry{c++:}}\index{a\_double@\spxentry{a\_double}}\index{binary@\spxentry{binary}}\index{operators@\spxentry{operators}}\index{with@\spxentry{with}}\index{ad@\spxentry{ad}}\index{result:@\spxentry{result:}}\index{example@\spxentry{example}}\index{test@\spxentry{test}}\ignorespaces 

\subparagraph{3.2.1.4.1 C++: a\_double Binary Operators With AD Result: Example and Test}
\label{\detokenize{xsrst/a_double_binary_xam_cpp:c-a-double-binary-operators-with-ad-result-example-and-test}}\label{\detokenize{xsrst/a_double_binary_xam_cpp:id1}}
\begin{sphinxVerbatim}[commandchars=\\\{\}]
\PYG{c+cp}{\PYGZsh{}}\PYG{c+cp}{ include \PYGZlt{}cstdio\PYGZgt{}}
\PYG{c+cp}{\PYGZsh{}}\PYG{c+cp}{ include \PYGZlt{}cppad}\PYG{c+cp}{/}\PYG{c+cp}{py}\PYG{c+cp}{/}\PYG{c+cp}{cppad\PYGZus{}py.hpp\PYGZgt{}}

\PYG{k+kt}{bool} \PYG{n+nf}{a\PYGZus{}double\PYGZus{}binary\PYGZus{}xam}\PYG{p}{(}\PYG{k+kt}{void}\PYG{p}{)} \PYG{p}{\PYGZob{}}
    \PYG{k}{using} \PYG{n}{cppad\PYGZus{}py}\PYG{o}{:}\PYG{o}{:}\PYG{n}{a\PYGZus{}double}\PYG{p}{;}
    \PYG{k+kt}{bool} \PYG{n}{ok}      \PYG{o}{=} \PYG{n+nb}{true}\PYG{p}{;}
    \PYG{n}{a\PYGZus{}double} \PYG{n}{a2}  \PYG{o}{=} \PYG{l+m+mf}{2.0}\PYG{p}{;}
    \PYG{n}{a\PYGZus{}double} \PYG{n}{a3}\PYG{p}{(}\PYG{l+m+mf}{3.0}\PYG{p}{)}\PYG{p}{;}
    \PYG{c+c1}{// \PYGZhy{}\PYGZhy{}\PYGZhy{}\PYGZhy{}\PYGZhy{}\PYGZhy{}\PYGZhy{}\PYGZhy{}\PYGZhy{}\PYGZhy{}\PYGZhy{}\PYGZhy{}\PYGZhy{}\PYGZhy{}\PYGZhy{}\PYGZhy{}\PYGZhy{}\PYGZhy{}\PYGZhy{}\PYGZhy{}\PYGZhy{}\PYGZhy{}\PYGZhy{}\PYGZhy{}\PYGZhy{}\PYGZhy{}\PYGZhy{}\PYGZhy{}\PYGZhy{}\PYGZhy{}\PYGZhy{}\PYGZhy{}\PYGZhy{}\PYGZhy{}\PYGZhy{}\PYGZhy{}\PYGZhy{}\PYGZhy{}\PYGZhy{}\PYGZhy{}\PYGZhy{}\PYGZhy{}\PYGZhy{}\PYGZhy{}\PYGZhy{}\PYGZhy{}\PYGZhy{}\PYGZhy{}\PYGZhy{}\PYGZhy{}\PYGZhy{}\PYGZhy{}\PYGZhy{}\PYGZhy{}\PYGZhy{}\PYGZhy{}\PYGZhy{}\PYGZhy{}\PYGZhy{}\PYGZhy{}\PYGZhy{}\PYGZhy{}\PYGZhy{}\PYGZhy{}\PYGZhy{}\PYGZhy{}\PYGZhy{}\PYGZhy{}\PYGZhy{}\PYGZhy{}\PYGZhy{}}
    \PYG{c+c1}{// a\PYGZus{}double op a\PYGZus{}double}
    \PYG{n}{a\PYGZus{}double} \PYG{n}{a5}       \PYG{o}{=} \PYG{n}{a2} \PYG{o}{+} \PYG{n}{a3}\PYG{p}{;}
    \PYG{n}{a\PYGZus{}double} \PYG{n}{a6}       \PYG{o}{=} \PYG{n}{a2} \PYG{o}{*} \PYG{n}{a3}\PYG{p}{;}
    \PYG{n}{a\PYGZus{}double} \PYG{n}{a1\PYGZus{}minus} \PYG{o}{=} \PYG{n}{a2} \PYG{o}{\PYGZhy{}} \PYG{n}{a3}\PYG{p}{;}
    \PYG{n}{a\PYGZus{}double} \PYG{n}{a23}      \PYG{o}{=} \PYG{n}{a2} \PYG{o}{/} \PYG{n}{a3}\PYG{p}{;}
    \PYG{c+c1}{//}
    \PYG{n}{ok} \PYG{o}{=} \PYG{n}{ok} \PYG{o}{\PYGZam{}}\PYG{o}{\PYGZam{}} \PYG{n}{a5} \PYG{o}{=}\PYG{o}{=} \PYG{l+m+mf}{5.0}\PYG{p}{;}
    \PYG{n}{ok} \PYG{o}{=} \PYG{n}{ok} \PYG{o}{\PYGZam{}}\PYG{o}{\PYGZam{}} \PYG{n}{a6} \PYG{o}{=}\PYG{o}{=} \PYG{l+m+mf}{6.0}\PYG{p}{;}
    \PYG{n}{ok} \PYG{o}{=} \PYG{n}{ok} \PYG{o}{\PYGZam{}}\PYG{o}{\PYGZam{}} \PYG{n}{a1\PYGZus{}minus} \PYG{o}{=}\PYG{o}{=} \PYG{o}{\PYGZhy{}}\PYG{l+m+mf}{1.0}\PYG{p}{;}
    \PYG{n}{ok} \PYG{o}{=} \PYG{n}{ok} \PYG{o}{\PYGZam{}}\PYG{o}{\PYGZam{}} \PYG{n}{a23}\PYG{p}{.}\PYG{n}{near\PYGZus{}equal}\PYG{p}{(} \PYG{n}{a\PYGZus{}double}\PYG{p}{(}\PYG{l+m+mf}{2.0} \PYG{o}{/} \PYG{l+m+mf}{3.0} \PYG{p}{)} \PYG{p}{)}\PYG{p}{;}
    \PYG{c+c1}{// \PYGZhy{}\PYGZhy{}\PYGZhy{}\PYGZhy{}\PYGZhy{}\PYGZhy{}\PYGZhy{}\PYGZhy{}\PYGZhy{}\PYGZhy{}\PYGZhy{}\PYGZhy{}\PYGZhy{}\PYGZhy{}\PYGZhy{}\PYGZhy{}\PYGZhy{}\PYGZhy{}\PYGZhy{}\PYGZhy{}\PYGZhy{}\PYGZhy{}\PYGZhy{}\PYGZhy{}\PYGZhy{}\PYGZhy{}\PYGZhy{}\PYGZhy{}\PYGZhy{}\PYGZhy{}\PYGZhy{}\PYGZhy{}\PYGZhy{}\PYGZhy{}\PYGZhy{}\PYGZhy{}\PYGZhy{}\PYGZhy{}\PYGZhy{}\PYGZhy{}\PYGZhy{}\PYGZhy{}\PYGZhy{}\PYGZhy{}\PYGZhy{}\PYGZhy{}\PYGZhy{}\PYGZhy{}\PYGZhy{}\PYGZhy{}\PYGZhy{}\PYGZhy{}\PYGZhy{}\PYGZhy{}\PYGZhy{}\PYGZhy{}\PYGZhy{}\PYGZhy{}\PYGZhy{}\PYGZhy{}\PYGZhy{}\PYGZhy{}\PYGZhy{}\PYGZhy{}\PYGZhy{}\PYGZhy{}\PYGZhy{}\PYGZhy{}\PYGZhy{}\PYGZhy{}\PYGZhy{}}
    \PYG{c+c1}{// a\PYGZus{}double op double}
    \PYG{n}{a5}       \PYG{o}{=} \PYG{n}{a2} \PYG{o}{+} \PYG{l+m+mf}{3.0}\PYG{p}{;}
    \PYG{n}{a6}       \PYG{o}{=} \PYG{n}{a2} \PYG{o}{*} \PYG{l+m+mf}{3.0}\PYG{p}{;}
    \PYG{n}{a1\PYGZus{}minus} \PYG{o}{=} \PYG{n}{a2} \PYG{o}{\PYGZhy{}} \PYG{l+m+mf}{3.0}\PYG{p}{;}
    \PYG{n}{a23}      \PYG{o}{=} \PYG{n}{a2} \PYG{o}{/} \PYG{l+m+mf}{3.0}\PYG{p}{;}
    \PYG{c+c1}{//}
    \PYG{n}{ok} \PYG{o}{=} \PYG{n}{ok} \PYG{o}{\PYGZam{}}\PYG{o}{\PYGZam{}} \PYG{n}{a5} \PYG{o}{=}\PYG{o}{=} \PYG{l+m+mf}{5.0}\PYG{p}{;}
    \PYG{n}{ok} \PYG{o}{=} \PYG{n}{ok} \PYG{o}{\PYGZam{}}\PYG{o}{\PYGZam{}} \PYG{n}{a6} \PYG{o}{=}\PYG{o}{=} \PYG{l+m+mf}{6.0}\PYG{p}{;}
    \PYG{n}{ok} \PYG{o}{=} \PYG{n}{ok} \PYG{o}{\PYGZam{}}\PYG{o}{\PYGZam{}} \PYG{n}{a1\PYGZus{}minus} \PYG{o}{=}\PYG{o}{=} \PYG{o}{\PYGZhy{}}\PYG{l+m+mf}{1.0}\PYG{p}{;}
    \PYG{n}{ok} \PYG{o}{=} \PYG{n}{ok} \PYG{o}{\PYGZam{}}\PYG{o}{\PYGZam{}} \PYG{n}{a23}\PYG{p}{.}\PYG{n}{near\PYGZus{}equal}\PYG{p}{(} \PYG{n}{a\PYGZus{}double}\PYG{p}{(}\PYG{l+m+mf}{2.0} \PYG{o}{/} \PYG{l+m+mf}{3.0} \PYG{p}{)} \PYG{p}{)}\PYG{p}{;}
    \PYG{c+c1}{// \PYGZhy{}\PYGZhy{}\PYGZhy{}\PYGZhy{}\PYGZhy{}\PYGZhy{}\PYGZhy{}\PYGZhy{}\PYGZhy{}\PYGZhy{}\PYGZhy{}\PYGZhy{}\PYGZhy{}\PYGZhy{}\PYGZhy{}\PYGZhy{}\PYGZhy{}\PYGZhy{}\PYGZhy{}\PYGZhy{}\PYGZhy{}\PYGZhy{}\PYGZhy{}\PYGZhy{}\PYGZhy{}\PYGZhy{}\PYGZhy{}\PYGZhy{}\PYGZhy{}\PYGZhy{}\PYGZhy{}\PYGZhy{}\PYGZhy{}\PYGZhy{}\PYGZhy{}\PYGZhy{}\PYGZhy{}\PYGZhy{}\PYGZhy{}\PYGZhy{}\PYGZhy{}\PYGZhy{}\PYGZhy{}\PYGZhy{}\PYGZhy{}\PYGZhy{}\PYGZhy{}\PYGZhy{}\PYGZhy{}\PYGZhy{}\PYGZhy{}\PYGZhy{}\PYGZhy{}\PYGZhy{}\PYGZhy{}\PYGZhy{}\PYGZhy{}\PYGZhy{}\PYGZhy{}\PYGZhy{}\PYGZhy{}\PYGZhy{}\PYGZhy{}\PYGZhy{}\PYGZhy{}\PYGZhy{}\PYGZhy{}\PYGZhy{}\PYGZhy{}\PYGZhy{}\PYGZhy{}}
    \PYG{c+c1}{// double op a\PYGZus{}double}
    \PYG{n}{a5}           \PYG{o}{=} \PYG{n}{cppad\PYGZus{}py}\PYG{o}{:}\PYG{o}{:}\PYG{n}{radd}\PYG{p}{(}\PYG{l+m+mf}{3.0}\PYG{p}{,} \PYG{n}{a2}\PYG{p}{)}\PYG{p}{;}
    \PYG{n}{a6}           \PYG{o}{=} \PYG{n}{cppad\PYGZus{}py}\PYG{o}{:}\PYG{o}{:}\PYG{n}{rmul}\PYG{p}{(}\PYG{l+m+mf}{3.0}\PYG{p}{,} \PYG{n}{a2}\PYG{p}{)}\PYG{p}{;}
    \PYG{n}{a\PYGZus{}double} \PYG{n}{a1}  \PYG{o}{=} \PYG{n}{cppad\PYGZus{}py}\PYG{o}{:}\PYG{o}{:}\PYG{n}{rsub}\PYG{p}{(}\PYG{l+m+mf}{3.0}\PYG{p}{,}  \PYG{n}{a2}\PYG{p}{)}\PYG{p}{;}
    \PYG{n}{a\PYGZus{}double} \PYG{n}{a32} \PYG{o}{=} \PYG{n}{cppad\PYGZus{}py}\PYG{o}{:}\PYG{o}{:}\PYG{n}{rdiv}\PYG{p}{(}\PYG{l+m+mf}{3.0}\PYG{p}{,} \PYG{n}{a2}\PYG{p}{)}\PYG{p}{;}
    \PYG{c+c1}{//}
    \PYG{n}{ok} \PYG{o}{=} \PYG{n}{ok} \PYG{o}{\PYGZam{}}\PYG{o}{\PYGZam{}} \PYG{n}{a5} \PYG{o}{=}\PYG{o}{=} \PYG{l+m+mf}{5.0}\PYG{p}{;}
    \PYG{n}{ok} \PYG{o}{=} \PYG{n}{ok} \PYG{o}{\PYGZam{}}\PYG{o}{\PYGZam{}} \PYG{n}{a6} \PYG{o}{=}\PYG{o}{=} \PYG{l+m+mf}{6.0}\PYG{p}{;}
    \PYG{n}{ok} \PYG{o}{=} \PYG{n}{ok} \PYG{o}{\PYGZam{}}\PYG{o}{\PYGZam{}} \PYG{n}{a1} \PYG{o}{=}\PYG{o}{=} \PYG{l+m+mf}{1.0}\PYG{p}{;}
    \PYG{n}{ok} \PYG{o}{=} \PYG{n}{ok} \PYG{o}{\PYGZam{}}\PYG{o}{\PYGZam{}} \PYG{n}{a32}\PYG{p}{.}\PYG{n}{near\PYGZus{}equal}\PYG{p}{(} \PYG{n}{a\PYGZus{}double}\PYG{p}{(}\PYG{l+m+mf}{3.0} \PYG{o}{/} \PYG{l+m+mf}{2.0} \PYG{p}{)} \PYG{p}{)}\PYG{p}{;}
    \PYG{c+c1}{// \PYGZhy{}\PYGZhy{}\PYGZhy{}\PYGZhy{}\PYGZhy{}\PYGZhy{}\PYGZhy{}\PYGZhy{}\PYGZhy{}\PYGZhy{}\PYGZhy{}\PYGZhy{}\PYGZhy{}\PYGZhy{}\PYGZhy{}\PYGZhy{}\PYGZhy{}\PYGZhy{}\PYGZhy{}\PYGZhy{}\PYGZhy{}\PYGZhy{}\PYGZhy{}\PYGZhy{}\PYGZhy{}\PYGZhy{}\PYGZhy{}\PYGZhy{}\PYGZhy{}\PYGZhy{}\PYGZhy{}\PYGZhy{}\PYGZhy{}\PYGZhy{}\PYGZhy{}\PYGZhy{}\PYGZhy{}\PYGZhy{}\PYGZhy{}\PYGZhy{}\PYGZhy{}\PYGZhy{}\PYGZhy{}\PYGZhy{}\PYGZhy{}\PYGZhy{}\PYGZhy{}\PYGZhy{}\PYGZhy{}\PYGZhy{}\PYGZhy{}\PYGZhy{}\PYGZhy{}\PYGZhy{}\PYGZhy{}\PYGZhy{}\PYGZhy{}\PYGZhy{}\PYGZhy{}\PYGZhy{}\PYGZhy{}\PYGZhy{}\PYGZhy{}\PYGZhy{}\PYGZhy{}\PYGZhy{}\PYGZhy{}\PYGZhy{}\PYGZhy{}\PYGZhy{}\PYGZhy{}}
    \PYG{c+c1}{// pow}
    \PYG{n}{a\PYGZus{}double} \PYG{n}{a8} \PYG{o}{=} \PYG{n}{cppad\PYGZus{}py}\PYG{o}{:}\PYG{o}{:}\PYG{n}{pow}\PYG{p}{(}\PYG{n}{a2}\PYG{p}{,} \PYG{n}{a3}\PYG{p}{)}\PYG{p}{;}
    \PYG{n}{a\PYGZus{}double} \PYG{n}{a9} \PYG{o}{=} \PYG{n}{cppad\PYGZus{}py}\PYG{o}{:}\PYG{o}{:}\PYG{n}{pow}\PYG{p}{(}\PYG{n}{a3}\PYG{p}{,} \PYG{l+m+mf}{2.0}\PYG{p}{)}\PYG{p}{;}
    \PYG{n}{a\PYGZus{}double} \PYG{n}{a4} \PYG{o}{=} \PYG{n}{cppad\PYGZus{}py}\PYG{o}{:}\PYG{o}{:}\PYG{n}{pow}\PYG{p}{(}\PYG{l+m+mf}{2.0}\PYG{p}{,} \PYG{n}{a2}\PYG{p}{)}\PYG{p}{;}
    \PYG{n}{ok} \PYG{o}{=} \PYG{n}{ok} \PYG{o}{\PYGZam{}}\PYG{o}{\PYGZam{}} \PYG{n}{a8}\PYG{p}{.}\PYG{n}{near\PYGZus{}equal}\PYG{p}{(} \PYG{n}{a\PYGZus{}double}\PYG{p}{(}\PYG{l+m+mf}{8.0}\PYG{p}{)} \PYG{p}{)}\PYG{p}{;}
    \PYG{n}{ok} \PYG{o}{=} \PYG{n}{ok} \PYG{o}{\PYGZam{}}\PYG{o}{\PYGZam{}} \PYG{n}{a9}\PYG{p}{.}\PYG{n}{near\PYGZus{}equal}\PYG{p}{(} \PYG{n}{a\PYGZus{}double}\PYG{p}{(}\PYG{l+m+mf}{9.0}\PYG{p}{)} \PYG{p}{)}\PYG{p}{;}
    \PYG{n}{ok} \PYG{o}{=} \PYG{n}{ok} \PYG{o}{\PYGZam{}}\PYG{o}{\PYGZam{}} \PYG{n}{a4}\PYG{p}{.}\PYG{n}{near\PYGZus{}equal}\PYG{p}{(} \PYG{n}{a\PYGZus{}double}\PYG{p}{(}\PYG{l+m+mf}{4.0}\PYG{p}{)} \PYG{p}{)}\PYG{p}{;}
    \PYG{c+c1}{//}
    \PYG{k}{return}\PYG{p}{(} \PYG{n}{ok} \PYG{p}{)}\PYG{p}{;}
\PYG{p}{\PYGZcb{}}
\end{sphinxVerbatim}


\bigskip\hrule\bigskip


xsrst input file: \sphinxcode{\sphinxupquote{example/cplusplus/a\_double\_binary\_xam.cpp}}


\subparagraph{a\_double\_binary\_xam\_py}
\label{\detokenize{xsrst/a_double_binary_xam_py:a-double-binary-xam-py}}\label{\detokenize{xsrst/a_double_binary_xam_py::doc}}
\index{a\_double\_binary\_xam\_py@\spxentry{a\_double\_binary\_xam\_py}}\index{python:@\spxentry{python:}}\index{a\_double@\spxentry{a\_double}}\index{binary@\spxentry{binary}}\index{operators@\spxentry{operators}}\index{with@\spxentry{with}}\index{ad@\spxentry{ad}}\index{result:@\spxentry{result:}}\index{example@\spxentry{example}}\index{test@\spxentry{test}}\ignorespaces 

\subparagraph{3.2.1.4.2 Python: a\_double Binary Operators With AD Result: Example and Test}
\label{\detokenize{xsrst/a_double_binary_xam_py:python-a-double-binary-operators-with-ad-result-example-and-test}}\label{\detokenize{xsrst/a_double_binary_xam_py:id1}}
\begin{sphinxVerbatim}[commandchars=\\\{\}]
\PYG{k}{def} \PYG{n+nf}{a\PYGZus{}double\PYGZus{}binary\PYGZus{}xam}\PYG{p}{(}\PYG{p}{)} \PYG{p}{:}
    \PYG{c+c1}{\PYGZsh{}}
    \PYG{k+kn}{import} \PYG{n+nn}{numpy}
    \PYG{k+kn}{from} \PYG{n+nn}{cppad\PYGZus{}py} \PYG{k+kn}{import} \PYG{n}{a\PYGZus{}double}
    \PYG{n}{ok} \PYG{o}{=} \PYG{k+kc}{True}
    \PYG{n}{a2} \PYG{o}{=} \PYG{n}{a\PYGZus{}double}\PYG{p}{(}\PYG{l+m+mf}{2.0}\PYG{p}{)}
    \PYG{n}{a3} \PYG{o}{=} \PYG{n}{a\PYGZus{}double}\PYG{p}{(}\PYG{l+m+mf}{3.0}\PYG{p}{)}
    \PYG{c+c1}{\PYGZsh{} \PYGZhy{}\PYGZhy{}\PYGZhy{}\PYGZhy{}\PYGZhy{}\PYGZhy{}\PYGZhy{}\PYGZhy{}\PYGZhy{}\PYGZhy{}\PYGZhy{}\PYGZhy{}\PYGZhy{}\PYGZhy{}\PYGZhy{}\PYGZhy{}\PYGZhy{}\PYGZhy{}\PYGZhy{}\PYGZhy{}\PYGZhy{}\PYGZhy{}\PYGZhy{}\PYGZhy{}\PYGZhy{}\PYGZhy{}\PYGZhy{}\PYGZhy{}\PYGZhy{}\PYGZhy{}\PYGZhy{}\PYGZhy{}\PYGZhy{}\PYGZhy{}\PYGZhy{}\PYGZhy{}\PYGZhy{}\PYGZhy{}\PYGZhy{}\PYGZhy{}\PYGZhy{}\PYGZhy{}\PYGZhy{}\PYGZhy{}\PYGZhy{}\PYGZhy{}\PYGZhy{}\PYGZhy{}\PYGZhy{}\PYGZhy{}\PYGZhy{}\PYGZhy{}\PYGZhy{}\PYGZhy{}\PYGZhy{}\PYGZhy{}\PYGZhy{}\PYGZhy{}\PYGZhy{}\PYGZhy{}\PYGZhy{}\PYGZhy{}\PYGZhy{}\PYGZhy{}\PYGZhy{}\PYGZhy{}\PYGZhy{}\PYGZhy{}\PYGZhy{}}
    \PYG{c+c1}{\PYGZsh{} a\PYGZus{}double op a\PYGZus{}double}
    \PYG{n}{a5}       \PYG{o}{=} \PYG{n}{a2} \PYG{o}{+} \PYG{n}{a3}
    \PYG{n}{a6}       \PYG{o}{=} \PYG{n}{a2} \PYG{o}{*} \PYG{n}{a3}
    \PYG{n}{a1\PYGZus{}minus} \PYG{o}{=} \PYG{n}{a2} \PYG{o}{\PYGZhy{}} \PYG{n}{a3}
    \PYG{n}{a23}      \PYG{o}{=} \PYG{n}{a2} \PYG{o}{/} \PYG{n}{a3}
    \PYG{c+c1}{\PYGZsh{}}
    \PYG{n}{ok} \PYG{o}{=} \PYG{n}{ok} \PYG{o+ow}{and} \PYG{n}{a5} \PYG{o}{==} \PYG{l+m+mf}{5.0}
    \PYG{n}{ok} \PYG{o}{=} \PYG{n}{ok} \PYG{o+ow}{and} \PYG{n}{a6} \PYG{o}{==} \PYG{l+m+mf}{6.0}
    \PYG{n}{ok} \PYG{o}{=} \PYG{n}{ok} \PYG{o+ow}{and} \PYG{n}{a1\PYGZus{}minus} \PYG{o}{==} \PYG{o}{\PYGZhy{}}\PYG{l+m+mf}{1.0}
    \PYG{n}{ok} \PYG{o}{=} \PYG{n}{ok} \PYG{o+ow}{and} \PYG{n}{a23}\PYG{o}{.}\PYG{n}{near\PYGZus{}equal}\PYG{p}{(} \PYG{n}{a\PYGZus{}double}\PYG{p}{(}\PYG{l+m+mf}{2.0} \PYG{o}{/} \PYG{l+m+mf}{3.0}\PYG{p}{)} \PYG{p}{)}
    \PYG{c+c1}{\PYGZsh{} \PYGZhy{}\PYGZhy{}\PYGZhy{}\PYGZhy{}\PYGZhy{}\PYGZhy{}\PYGZhy{}\PYGZhy{}\PYGZhy{}\PYGZhy{}\PYGZhy{}\PYGZhy{}\PYGZhy{}\PYGZhy{}\PYGZhy{}\PYGZhy{}\PYGZhy{}\PYGZhy{}\PYGZhy{}\PYGZhy{}\PYGZhy{}\PYGZhy{}\PYGZhy{}\PYGZhy{}\PYGZhy{}\PYGZhy{}\PYGZhy{}\PYGZhy{}\PYGZhy{}\PYGZhy{}\PYGZhy{}\PYGZhy{}\PYGZhy{}\PYGZhy{}\PYGZhy{}\PYGZhy{}\PYGZhy{}\PYGZhy{}\PYGZhy{}\PYGZhy{}\PYGZhy{}\PYGZhy{}\PYGZhy{}\PYGZhy{}\PYGZhy{}\PYGZhy{}\PYGZhy{}\PYGZhy{}\PYGZhy{}\PYGZhy{}\PYGZhy{}\PYGZhy{}\PYGZhy{}\PYGZhy{}\PYGZhy{}\PYGZhy{}\PYGZhy{}\PYGZhy{}\PYGZhy{}\PYGZhy{}\PYGZhy{}\PYGZhy{}\PYGZhy{}\PYGZhy{}\PYGZhy{}\PYGZhy{}\PYGZhy{}\PYGZhy{}\PYGZhy{}}
    \PYG{c+c1}{\PYGZsh{} a\PYGZus{}double op double}
    \PYG{n}{a5}       \PYG{o}{=} \PYG{n}{a2} \PYG{o}{+} \PYG{l+m+mf}{3.0}
    \PYG{n}{a6}       \PYG{o}{=} \PYG{n}{a2} \PYG{o}{*} \PYG{l+m+mf}{3.0}
    \PYG{n}{a1\PYGZus{}minus} \PYG{o}{=} \PYG{n}{a2} \PYG{o}{\PYGZhy{}} \PYG{l+m+mf}{3.0}
    \PYG{n}{a23}      \PYG{o}{=} \PYG{n}{a2} \PYG{o}{/} \PYG{l+m+mf}{3.0}
    \PYG{c+c1}{\PYGZsh{}}
    \PYG{n}{ok} \PYG{o}{=} \PYG{n}{ok} \PYG{o+ow}{and} \PYG{n}{a5} \PYG{o}{==} \PYG{l+m+mf}{5.0}
    \PYG{n}{ok} \PYG{o}{=} \PYG{n}{ok} \PYG{o+ow}{and} \PYG{n}{a6} \PYG{o}{==} \PYG{l+m+mf}{6.0}
    \PYG{n}{ok} \PYG{o}{=} \PYG{n}{ok} \PYG{o+ow}{and} \PYG{n}{a1\PYGZus{}minus} \PYG{o}{==} \PYG{o}{\PYGZhy{}}\PYG{l+m+mf}{1.0}
    \PYG{n}{ok} \PYG{o}{=} \PYG{n}{ok} \PYG{o+ow}{and} \PYG{n}{a23}\PYG{o}{.}\PYG{n}{near\PYGZus{}equal}\PYG{p}{(} \PYG{n}{a\PYGZus{}double}\PYG{p}{(}\PYG{l+m+mf}{2.0} \PYG{o}{/} \PYG{l+m+mf}{3.0}\PYG{p}{)} \PYG{p}{)}
    \PYG{c+c1}{\PYGZsh{} \PYGZhy{}\PYGZhy{}\PYGZhy{}\PYGZhy{}\PYGZhy{}\PYGZhy{}\PYGZhy{}\PYGZhy{}\PYGZhy{}\PYGZhy{}\PYGZhy{}\PYGZhy{}\PYGZhy{}\PYGZhy{}\PYGZhy{}\PYGZhy{}\PYGZhy{}\PYGZhy{}\PYGZhy{}\PYGZhy{}\PYGZhy{}\PYGZhy{}\PYGZhy{}\PYGZhy{}\PYGZhy{}\PYGZhy{}\PYGZhy{}\PYGZhy{}\PYGZhy{}\PYGZhy{}\PYGZhy{}\PYGZhy{}\PYGZhy{}\PYGZhy{}\PYGZhy{}\PYGZhy{}\PYGZhy{}\PYGZhy{}\PYGZhy{}\PYGZhy{}\PYGZhy{}\PYGZhy{}\PYGZhy{}\PYGZhy{}\PYGZhy{}\PYGZhy{}\PYGZhy{}\PYGZhy{}\PYGZhy{}\PYGZhy{}\PYGZhy{}\PYGZhy{}\PYGZhy{}\PYGZhy{}\PYGZhy{}\PYGZhy{}\PYGZhy{}\PYGZhy{}\PYGZhy{}\PYGZhy{}\PYGZhy{}\PYGZhy{}\PYGZhy{}\PYGZhy{}\PYGZhy{}\PYGZhy{}\PYGZhy{}\PYGZhy{}\PYGZhy{}}
    \PYG{c+c1}{\PYGZsh{} double op a\PYGZus{}double}
    \PYG{n}{a5}       \PYG{o}{=} \PYG{l+m+mf}{3.0} \PYG{o}{+} \PYG{n}{a2}
    \PYG{n}{a6}       \PYG{o}{=} \PYG{l+m+mf}{3.0} \PYG{o}{*} \PYG{n}{a2}
    \PYG{n}{a1}       \PYG{o}{=} \PYG{l+m+mf}{3.0} \PYG{o}{\PYGZhy{}} \PYG{n}{a2}
    \PYG{n}{a32}      \PYG{o}{=} \PYG{l+m+mf}{3.0} \PYG{o}{/} \PYG{n}{a2}
    \PYG{c+c1}{\PYGZsh{}}
    \PYG{n}{ok} \PYG{o}{=} \PYG{n}{ok} \PYG{o+ow}{and} \PYG{n}{a5} \PYG{o}{==} \PYG{l+m+mf}{5.0}
    \PYG{n}{ok} \PYG{o}{=} \PYG{n}{ok} \PYG{o+ow}{and} \PYG{n}{a6} \PYG{o}{==} \PYG{l+m+mf}{6.0}
    \PYG{n}{ok} \PYG{o}{=} \PYG{n}{ok} \PYG{o+ow}{and} \PYG{n}{a1} \PYG{o}{==} \PYG{l+m+mf}{1.0}
    \PYG{n}{ok} \PYG{o}{=} \PYG{n}{ok} \PYG{o+ow}{and} \PYG{n}{a32}\PYG{o}{.}\PYG{n}{near\PYGZus{}equal}\PYG{p}{(} \PYG{n}{a\PYGZus{}double}\PYG{p}{(}\PYG{l+m+mf}{3.0} \PYG{o}{/} \PYG{l+m+mf}{2.0}\PYG{p}{)} \PYG{p}{)}
    \PYG{c+c1}{\PYGZsh{} \PYGZhy{}\PYGZhy{}\PYGZhy{}\PYGZhy{}\PYGZhy{}\PYGZhy{}\PYGZhy{}\PYGZhy{}\PYGZhy{}\PYGZhy{}\PYGZhy{}\PYGZhy{}\PYGZhy{}\PYGZhy{}\PYGZhy{}\PYGZhy{}\PYGZhy{}\PYGZhy{}\PYGZhy{}\PYGZhy{}\PYGZhy{}\PYGZhy{}\PYGZhy{}\PYGZhy{}\PYGZhy{}\PYGZhy{}\PYGZhy{}\PYGZhy{}\PYGZhy{}\PYGZhy{}\PYGZhy{}\PYGZhy{}\PYGZhy{}\PYGZhy{}\PYGZhy{}\PYGZhy{}\PYGZhy{}\PYGZhy{}\PYGZhy{}\PYGZhy{}\PYGZhy{}\PYGZhy{}\PYGZhy{}\PYGZhy{}\PYGZhy{}\PYGZhy{}\PYGZhy{}\PYGZhy{}\PYGZhy{}\PYGZhy{}\PYGZhy{}\PYGZhy{}\PYGZhy{}\PYGZhy{}\PYGZhy{}\PYGZhy{}\PYGZhy{}\PYGZhy{}\PYGZhy{}\PYGZhy{}\PYGZhy{}\PYGZhy{}\PYGZhy{}\PYGZhy{}\PYGZhy{}\PYGZhy{}\PYGZhy{}\PYGZhy{}\PYGZhy{}}
    \PYG{c+c1}{\PYGZsh{} pow}
    \PYG{n}{a8} \PYG{o}{=} \PYG{n}{a2}  \PYG{o}{*}\PYG{o}{*} \PYG{n}{a3}
    \PYG{n}{a9} \PYG{o}{=} \PYG{n}{a3}  \PYG{o}{*}\PYG{o}{*} \PYG{l+m+mf}{2.0}
    \PYG{n}{a4} \PYG{o}{=} \PYG{l+m+mf}{2.0} \PYG{o}{*}\PYG{o}{*} \PYG{n}{a2}
    \PYG{n}{ok} \PYG{o}{=} \PYG{n}{ok} \PYG{o+ow}{and} \PYG{n}{a8}\PYG{o}{.}\PYG{n}{near\PYGZus{}equal}\PYG{p}{(} \PYG{n}{a\PYGZus{}double}\PYG{p}{(}\PYG{l+m+mf}{8.0}\PYG{p}{)} \PYG{p}{)}
    \PYG{n}{ok} \PYG{o}{=} \PYG{n}{ok} \PYG{o+ow}{and} \PYG{n}{a9}\PYG{o}{.}\PYG{n}{near\PYGZus{}equal}\PYG{p}{(} \PYG{n}{a\PYGZus{}double}\PYG{p}{(}\PYG{l+m+mf}{9.0}\PYG{p}{)} \PYG{p}{)}
    \PYG{n}{ok} \PYG{o}{=} \PYG{n}{ok} \PYG{o+ow}{and} \PYG{n}{a4}\PYG{o}{.}\PYG{n}{near\PYGZus{}equal}\PYG{p}{(} \PYG{n}{a\PYGZus{}double}\PYG{p}{(}\PYG{l+m+mf}{4.0}\PYG{p}{)} \PYG{p}{)}
    \PYG{c+c1}{\PYGZsh{} \PYGZhy{}\PYGZhy{}\PYGZhy{}\PYGZhy{}\PYGZhy{}\PYGZhy{}\PYGZhy{}\PYGZhy{}\PYGZhy{}\PYGZhy{}\PYGZhy{}\PYGZhy{}\PYGZhy{}\PYGZhy{}\PYGZhy{}\PYGZhy{}\PYGZhy{}\PYGZhy{}\PYGZhy{}\PYGZhy{}\PYGZhy{}\PYGZhy{}\PYGZhy{}\PYGZhy{}\PYGZhy{}\PYGZhy{}\PYGZhy{}\PYGZhy{}\PYGZhy{}\PYGZhy{}\PYGZhy{}\PYGZhy{}\PYGZhy{}\PYGZhy{}\PYGZhy{}\PYGZhy{}\PYGZhy{}\PYGZhy{}\PYGZhy{}\PYGZhy{}\PYGZhy{}\PYGZhy{}\PYGZhy{}\PYGZhy{}\PYGZhy{}\PYGZhy{}\PYGZhy{}\PYGZhy{}\PYGZhy{}\PYGZhy{}\PYGZhy{}\PYGZhy{}\PYGZhy{}\PYGZhy{}\PYGZhy{}\PYGZhy{}\PYGZhy{}\PYGZhy{}\PYGZhy{}\PYGZhy{}\PYGZhy{}\PYGZhy{}\PYGZhy{}\PYGZhy{}\PYGZhy{}\PYGZhy{}\PYGZhy{}\PYGZhy{}\PYGZhy{}}

    \PYG{c+c1}{\PYGZsh{}}
    \PYG{k}{return}\PYG{p}{(} \PYG{n}{ok} \PYG{p}{)}
\PYG{c+c1}{\PYGZsh{}}
\end{sphinxVerbatim}


\bigskip\hrule\bigskip


xsrst input file: \sphinxcode{\sphinxupquote{example/python/core/a\_double\_binary\_xam.py}}

\index{example@\spxentry{example}}\ignorespaces 

\subparagraph{Example}
\label{\detokenize{xsrst/a_double_binary:example}}\label{\detokenize{xsrst/a_double_binary:a-double-binary-example}}\label{\detokenize{xsrst/a_double_binary:index-7}}
{\hyperref[\detokenize{xsrst/a_double_binary_xam_cpp:id1}]{\sphinxcrossref{\DUrole{std,std-ref}{c++}}}},
{\hyperref[\detokenize{xsrst/a_double_binary_xam_py:id1}]{\sphinxcrossref{\DUrole{std,std-ref}{python}}}}.


\bigskip\hrule\bigskip


xsrst input file: \sphinxcode{\sphinxupquote{lib/cplusplus/a\_double.cpp}}


\subparagraph{a\_double\_compare}
\label{\detokenize{xsrst/a_double_compare:a-double-compare}}\label{\detokenize{xsrst/a_double_compare::doc}}
\index{a\_double\_compare@\spxentry{a\_double\_compare}}\index{a\_double@\spxentry{a\_double}}\index{comparison@\spxentry{comparison}}\index{operators@\spxentry{operators}}\ignorespaces 

\subparagraph{3.2.1.5 a\_double Comparison Operators}
\label{\detokenize{xsrst/a_double_compare:a-double-comparison-operators}}\label{\detokenize{xsrst/a_double_compare:id1}}\begin{itemize}
\item {} 
{\hyperref[\detokenize{xsrst/a_double_compare:a-double-compare-syntax}]{\sphinxcrossref{\DUrole{std,std-ref}{Syntax}}}}

\item {} 
{\hyperref[\detokenize{xsrst/a_double_compare:a-double-compare-op}]{\sphinxcrossref{\DUrole{std,std-ref}{op}}}}

\item {} 
{\hyperref[\detokenize{xsrst/a_double_compare:a-double-compare-ax}]{\sphinxcrossref{\DUrole{std,std-ref}{ax}}}}

\item {} 
{\hyperref[\detokenize{xsrst/a_double_compare:a-double-compare-ay}]{\sphinxcrossref{\DUrole{std,std-ref}{ay}}}}

\item {} 
{\hyperref[\detokenize{xsrst/a_double_compare:a-double-compare-y}]{\sphinxcrossref{\DUrole{std,std-ref}{y}}}}

\item {} 
{\hyperref[\detokenize{xsrst/a_double_compare:a-double-compare-b}]{\sphinxcrossref{\DUrole{std,std-ref}{b}}}}

\item {} 
{\hyperref[\detokenize{xsrst/a_double_compare:a-double-compare-example}]{\sphinxcrossref{\DUrole{std,std-ref}{Example}}}}

\end{itemize}

\index{syntax@\spxentry{syntax}}\ignorespaces 

\subparagraph{Syntax}
\label{\detokenize{xsrst/a_double_compare:syntax}}\label{\detokenize{xsrst/a_double_compare:a-double-compare-syntax}}\label{\detokenize{xsrst/a_double_compare:index-1}}
\begin{DUlineblock}{0em}
\item[] \sphinxstyleemphasis{b} = \sphinxstyleemphasis{ax} \sphinxstyleemphasis{op} \sphinxstyleemphasis{ay}
\item[] \sphinxstyleemphasis{b} = \sphinxstyleemphasis{ax} \sphinxstyleemphasis{op} \sphinxstyleemphasis{y}
\end{DUlineblock}

\index{op@\spxentry{op}}\ignorespaces 

\subparagraph{op}
\label{\detokenize{xsrst/a_double_compare:op}}\label{\detokenize{xsrst/a_double_compare:a-double-compare-op}}\label{\detokenize{xsrst/a_double_compare:index-2}}
The binary operator \sphinxstyleemphasis{op} is one of the following:
\sphinxcode{\sphinxupquote{\textless{}}} (less than),
\sphinxcode{\sphinxupquote{\textless{}=}} (less than or equal),
\sphinxcode{\sphinxupquote{\textgreater{}}} (greater than),
\sphinxcode{\sphinxupquote{\textgreater{}=}} (greater than or equal),
\sphinxcode{\sphinxupquote{==}} (equal),
\sphinxcode{\sphinxupquote{!=}} (not equal).

\index{ax@\spxentry{ax}}\ignorespaces 

\subparagraph{ax}
\label{\detokenize{xsrst/a_double_compare:ax}}\label{\detokenize{xsrst/a_double_compare:a-double-compare-ax}}\label{\detokenize{xsrst/a_double_compare:index-3}}
This object has prototype

\begin{DUlineblock}{0em}
\item[]         \sphinxcode{\sphinxupquote{const a\_double\&}} \sphinxstyleemphasis{ax}
\end{DUlineblock}

\index{ay@\spxentry{ay}}\ignorespaces 

\subparagraph{ay}
\label{\detokenize{xsrst/a_double_compare:ay}}\label{\detokenize{xsrst/a_double_compare:a-double-compare-ay}}\label{\detokenize{xsrst/a_double_compare:index-4}}
This object has prototype

\begin{DUlineblock}{0em}
\item[]         \sphinxcode{\sphinxupquote{const a\_double\&}} \sphinxstyleemphasis{ay}
\end{DUlineblock}

\index{y@\spxentry{y}}\ignorespaces 

\subparagraph{y}
\label{\detokenize{xsrst/a_double_compare:y}}\label{\detokenize{xsrst/a_double_compare:a-double-compare-y}}\label{\detokenize{xsrst/a_double_compare:index-5}}
This object has prototype

\begin{DUlineblock}{0em}
\item[]         \sphinxcode{\sphinxupquote{const double\&}} \sphinxstyleemphasis{y}
\end{DUlineblock}

\index{b@\spxentry{b}}\ignorespaces 

\subparagraph{b}
\label{\detokenize{xsrst/a_double_compare:b}}\label{\detokenize{xsrst/a_double_compare:a-double-compare-b}}\label{\detokenize{xsrst/a_double_compare:index-6}}
The result has prototype

\begin{DUlineblock}{0em}
\item[]         \sphinxcode{\sphinxupquote{bool}} \sphinxstyleemphasis{b}
\end{DUlineblock}


\subparagraph{a\_double\_compare\_xam\_cpp}
\label{\detokenize{xsrst/a_double_compare_xam_cpp:a-double-compare-xam-cpp}}\label{\detokenize{xsrst/a_double_compare_xam_cpp::doc}}
\index{a\_double\_compare\_xam\_cpp@\spxentry{a\_double\_compare\_xam\_cpp}}\index{c++:@\spxentry{c++:}}\index{a\_double@\spxentry{a\_double}}\index{comparison@\spxentry{comparison}}\index{operators:@\spxentry{operators:}}\index{example@\spxentry{example}}\index{test@\spxentry{test}}\ignorespaces 

\subparagraph{3.2.1.5.1 C++: a\_double Comparison Operators: Example and Test}
\label{\detokenize{xsrst/a_double_compare_xam_cpp:c-a-double-comparison-operators-example-and-test}}\label{\detokenize{xsrst/a_double_compare_xam_cpp:id1}}
\begin{sphinxVerbatim}[commandchars=\\\{\}]
\PYG{c+cp}{\PYGZsh{}}\PYG{c+cp}{ include \PYGZlt{}cstdio\PYGZgt{}}
\PYG{c+cp}{\PYGZsh{}}\PYG{c+cp}{ include \PYGZlt{}cppad}\PYG{c+cp}{/}\PYG{c+cp}{py}\PYG{c+cp}{/}\PYG{c+cp}{cppad\PYGZus{}py.hpp\PYGZgt{}}

\PYG{k+kt}{bool} \PYG{n+nf}{a\PYGZus{}double\PYGZus{}compare\PYGZus{}xam}\PYG{p}{(}\PYG{k+kt}{void}\PYG{p}{)} \PYG{p}{\PYGZob{}}
    \PYG{k}{using} \PYG{n}{cppad\PYGZus{}py}\PYG{o}{:}\PYG{o}{:}\PYG{n}{a\PYGZus{}double}\PYG{p}{;}
    \PYG{k+kt}{bool} \PYG{n}{ok} \PYG{o}{=} \PYG{n+nb}{true}\PYG{p}{;}
    \PYG{n}{a\PYGZus{}double} \PYG{n}{a2}\PYG{p}{(}\PYG{l+m+mf}{2.0}\PYG{p}{)}\PYG{p}{;}
    \PYG{n}{a\PYGZus{}double} \PYG{n}{a3}\PYG{p}{(}\PYG{l+m+mf}{3.0}\PYG{p}{)}\PYG{p}{;}
    \PYG{c+c1}{//\PYGZhy{}\PYGZhy{}\PYGZhy{}\PYGZhy{}\PYGZhy{}\PYGZhy{}\PYGZhy{}\PYGZhy{}\PYGZhy{}\PYGZhy{}\PYGZhy{}\PYGZhy{}\PYGZhy{}\PYGZhy{}\PYGZhy{}\PYGZhy{}\PYGZhy{}\PYGZhy{}\PYGZhy{}\PYGZhy{}\PYGZhy{}\PYGZhy{}\PYGZhy{}\PYGZhy{}\PYGZhy{}\PYGZhy{}\PYGZhy{}\PYGZhy{}\PYGZhy{}\PYGZhy{}\PYGZhy{}\PYGZhy{}\PYGZhy{}\PYGZhy{}\PYGZhy{}\PYGZhy{}\PYGZhy{}\PYGZhy{}\PYGZhy{}\PYGZhy{}\PYGZhy{}\PYGZhy{}\PYGZhy{}\PYGZhy{}\PYGZhy{}\PYGZhy{}\PYGZhy{}\PYGZhy{}\PYGZhy{}\PYGZhy{}\PYGZhy{}\PYGZhy{}\PYGZhy{}\PYGZhy{}\PYGZhy{}\PYGZhy{}\PYGZhy{}\PYGZhy{}\PYGZhy{}\PYGZhy{}\PYGZhy{}\PYGZhy{}\PYGZhy{}\PYGZhy{}\PYGZhy{}\PYGZhy{}\PYGZhy{}\PYGZhy{}\PYGZhy{}\PYGZhy{}\PYGZhy{}\PYGZhy{}}
    \PYG{n}{ok} \PYG{o}{=} \PYG{n}{ok} \PYG{o}{\PYGZam{}}\PYG{o}{\PYGZam{}} \PYG{n}{a2}   \PYG{o}{\PYGZlt{}}  \PYG{n}{a3}\PYG{p}{;}
    \PYG{n}{ok} \PYG{o}{=} \PYG{n}{ok} \PYG{o}{\PYGZam{}}\PYG{o}{\PYGZam{}} \PYG{n}{a2}   \PYG{o}{\PYGZlt{}}\PYG{o}{=} \PYG{n}{a3}\PYG{p}{;}
    \PYG{n}{ok} \PYG{o}{=} \PYG{n}{ok} \PYG{o}{\PYGZam{}}\PYG{o}{\PYGZam{}} \PYG{n}{a3} \PYG{o}{\PYGZgt{}}  \PYG{n}{a2}\PYG{p}{;}
    \PYG{n}{ok} \PYG{o}{=} \PYG{n}{ok} \PYG{o}{\PYGZam{}}\PYG{o}{\PYGZam{}} \PYG{n}{a3} \PYG{o}{\PYGZgt{}}\PYG{o}{=} \PYG{n}{a2}\PYG{p}{;}
    \PYG{n}{ok} \PYG{o}{=} \PYG{n}{ok} \PYG{o}{\PYGZam{}}\PYG{o}{\PYGZam{}} \PYG{n}{a3} \PYG{o}{!}\PYG{o}{=} \PYG{n}{a2}\PYG{p}{;}
    \PYG{n}{ok} \PYG{o}{=} \PYG{n}{ok} \PYG{o}{\PYGZam{}}\PYG{o}{\PYGZam{}} \PYG{n}{a3} \PYG{o}{=}\PYG{o}{=} \PYG{n}{a3}\PYG{p}{;}
    \PYG{c+c1}{//}
    \PYG{n}{ok} \PYG{o}{=} \PYG{n}{ok} \PYG{o}{\PYGZam{}}\PYG{o}{\PYGZam{}} \PYG{o}{!} \PYG{p}{(}\PYG{n}{a2} \PYG{o}{\PYGZgt{}}  \PYG{n}{a3}\PYG{p}{)} \PYG{p}{;}
    \PYG{n}{ok} \PYG{o}{=} \PYG{n}{ok} \PYG{o}{\PYGZam{}}\PYG{o}{\PYGZam{}} \PYG{o}{!} \PYG{p}{(}\PYG{n}{a2} \PYG{o}{\PYGZgt{}}\PYG{o}{=} \PYG{n}{a3}\PYG{p}{)} \PYG{p}{;}
    \PYG{n}{ok} \PYG{o}{=} \PYG{n}{ok} \PYG{o}{\PYGZam{}}\PYG{o}{\PYGZam{}} \PYG{o}{!} \PYG{p}{(}\PYG{n}{a2} \PYG{o}{=}\PYG{o}{=} \PYG{n}{a3}\PYG{p}{)} \PYG{p}{;}
    \PYG{c+c1}{//\PYGZhy{}\PYGZhy{}\PYGZhy{}\PYGZhy{}\PYGZhy{}\PYGZhy{}\PYGZhy{}\PYGZhy{}\PYGZhy{}\PYGZhy{}\PYGZhy{}\PYGZhy{}\PYGZhy{}\PYGZhy{}\PYGZhy{}\PYGZhy{}\PYGZhy{}\PYGZhy{}\PYGZhy{}\PYGZhy{}\PYGZhy{}\PYGZhy{}\PYGZhy{}\PYGZhy{}\PYGZhy{}\PYGZhy{}\PYGZhy{}\PYGZhy{}\PYGZhy{}\PYGZhy{}\PYGZhy{}\PYGZhy{}\PYGZhy{}\PYGZhy{}\PYGZhy{}\PYGZhy{}\PYGZhy{}\PYGZhy{}\PYGZhy{}\PYGZhy{}\PYGZhy{}\PYGZhy{}\PYGZhy{}\PYGZhy{}\PYGZhy{}\PYGZhy{}\PYGZhy{}\PYGZhy{}\PYGZhy{}\PYGZhy{}\PYGZhy{}\PYGZhy{}\PYGZhy{}\PYGZhy{}\PYGZhy{}\PYGZhy{}\PYGZhy{}\PYGZhy{}\PYGZhy{}\PYGZhy{}\PYGZhy{}\PYGZhy{}\PYGZhy{}\PYGZhy{}\PYGZhy{}\PYGZhy{}\PYGZhy{}\PYGZhy{}\PYGZhy{}\PYGZhy{}\PYGZhy{}\PYGZhy{}}
    \PYG{n}{ok} \PYG{o}{=} \PYG{n}{ok} \PYG{o}{\PYGZam{}}\PYG{o}{\PYGZam{}} \PYG{n}{a2}   \PYG{o}{\PYGZlt{}}  \PYG{l+m+mf}{3.0}\PYG{p}{;}
    \PYG{n}{ok} \PYG{o}{=} \PYG{n}{ok} \PYG{o}{\PYGZam{}}\PYG{o}{\PYGZam{}} \PYG{n}{a2}   \PYG{o}{\PYGZlt{}}\PYG{o}{=} \PYG{l+m+mf}{3.0}\PYG{p}{;}
    \PYG{n}{ok} \PYG{o}{=} \PYG{n}{ok} \PYG{o}{\PYGZam{}}\PYG{o}{\PYGZam{}} \PYG{n}{a3} \PYG{o}{\PYGZgt{}}  \PYG{l+m+mf}{2.0}\PYG{p}{;}
    \PYG{n}{ok} \PYG{o}{=} \PYG{n}{ok} \PYG{o}{\PYGZam{}}\PYG{o}{\PYGZam{}} \PYG{n}{a3} \PYG{o}{\PYGZgt{}}\PYG{o}{=} \PYG{l+m+mf}{2.0}\PYG{p}{;}
    \PYG{n}{ok} \PYG{o}{=} \PYG{n}{ok} \PYG{o}{\PYGZam{}}\PYG{o}{\PYGZam{}} \PYG{n}{a3} \PYG{o}{!}\PYG{o}{=} \PYG{l+m+mf}{2.0}\PYG{p}{;}
    \PYG{n}{ok} \PYG{o}{=} \PYG{n}{ok} \PYG{o}{\PYGZam{}}\PYG{o}{\PYGZam{}} \PYG{n}{a3} \PYG{o}{=}\PYG{o}{=} \PYG{l+m+mf}{3.0}\PYG{p}{;}
    \PYG{c+c1}{//}
    \PYG{n}{ok} \PYG{o}{=} \PYG{n}{ok} \PYG{o}{\PYGZam{}}\PYG{o}{\PYGZam{}} \PYG{o}{!} \PYG{p}{(}\PYG{n}{a2} \PYG{o}{\PYGZgt{}}  \PYG{l+m+mf}{3.0}\PYG{p}{)} \PYG{p}{;}
    \PYG{n}{ok} \PYG{o}{=} \PYG{n}{ok} \PYG{o}{\PYGZam{}}\PYG{o}{\PYGZam{}} \PYG{o}{!} \PYG{p}{(}\PYG{n}{a2} \PYG{o}{\PYGZgt{}}\PYG{o}{=} \PYG{l+m+mf}{3.0}\PYG{p}{)} \PYG{p}{;}
    \PYG{n}{ok} \PYG{o}{=} \PYG{n}{ok} \PYG{o}{\PYGZam{}}\PYG{o}{\PYGZam{}} \PYG{o}{!} \PYG{p}{(}\PYG{n}{a2} \PYG{o}{=}\PYG{o}{=} \PYG{l+m+mf}{3.0}\PYG{p}{)} \PYG{p}{;}
    \PYG{c+c1}{//}
    \PYG{k}{return}\PYG{p}{(} \PYG{n}{ok} \PYG{p}{)}\PYG{p}{;}
\PYG{p}{\PYGZcb{}}
\end{sphinxVerbatim}


\bigskip\hrule\bigskip


xsrst input file: \sphinxcode{\sphinxupquote{example/cplusplus/a\_double\_compare\_xam.cpp}}


\subparagraph{a\_double\_compare\_xam\_py}
\label{\detokenize{xsrst/a_double_compare_xam_py:a-double-compare-xam-py}}\label{\detokenize{xsrst/a_double_compare_xam_py::doc}}
\index{a\_double\_compare\_xam\_py@\spxentry{a\_double\_compare\_xam\_py}}\index{python:@\spxentry{python:}}\index{a\_double@\spxentry{a\_double}}\index{comparison@\spxentry{comparison}}\index{operators:@\spxentry{operators:}}\index{example@\spxentry{example}}\index{test@\spxentry{test}}\ignorespaces 

\subparagraph{3.2.1.5.2 Python: a\_double Comparison Operators: Example and Test}
\label{\detokenize{xsrst/a_double_compare_xam_py:python-a-double-comparison-operators-example-and-test}}\label{\detokenize{xsrst/a_double_compare_xam_py:id1}}
\begin{sphinxVerbatim}[commandchars=\\\{\}]
\PYG{k}{def} \PYG{n+nf}{a\PYGZus{}double\PYGZus{}compare\PYGZus{}xam}\PYG{p}{(}\PYG{p}{)} \PYG{p}{:}
    \PYG{c+c1}{\PYGZsh{}}
    \PYG{k+kn}{import} \PYG{n+nn}{numpy}
    \PYG{k+kn}{import} \PYG{n+nn}{cppad\PYGZus{}py}
    \PYG{n}{ok} \PYG{o}{=} \PYG{k+kc}{True}
    \PYG{n}{a2} \PYG{o}{=} \PYG{n}{cppad\PYGZus{}py}\PYG{o}{.}\PYG{n}{a\PYGZus{}double}\PYG{p}{(}\PYG{l+m+mf}{2.0}\PYG{p}{)}
    \PYG{n}{a3} \PYG{o}{=} \PYG{n}{cppad\PYGZus{}py}\PYG{o}{.}\PYG{n}{a\PYGZus{}double}\PYG{p}{(}\PYG{l+m+mf}{3.0}\PYG{p}{)}
    \PYG{c+c1}{\PYGZsh{} \PYGZhy{}\PYGZhy{}\PYGZhy{}\PYGZhy{}\PYGZhy{}\PYGZhy{}\PYGZhy{}\PYGZhy{}\PYGZhy{}\PYGZhy{}\PYGZhy{}\PYGZhy{}\PYGZhy{}\PYGZhy{}\PYGZhy{}\PYGZhy{}\PYGZhy{}\PYGZhy{}\PYGZhy{}\PYGZhy{}\PYGZhy{}\PYGZhy{}\PYGZhy{}\PYGZhy{}\PYGZhy{}\PYGZhy{}\PYGZhy{}\PYGZhy{}\PYGZhy{}\PYGZhy{}\PYGZhy{}\PYGZhy{}\PYGZhy{}\PYGZhy{}\PYGZhy{}\PYGZhy{}\PYGZhy{}\PYGZhy{}\PYGZhy{}\PYGZhy{}\PYGZhy{}\PYGZhy{}\PYGZhy{}\PYGZhy{}\PYGZhy{}\PYGZhy{}\PYGZhy{}\PYGZhy{}\PYGZhy{}\PYGZhy{}\PYGZhy{}\PYGZhy{}\PYGZhy{}\PYGZhy{}\PYGZhy{}\PYGZhy{}\PYGZhy{}\PYGZhy{}\PYGZhy{}\PYGZhy{}\PYGZhy{}\PYGZhy{}\PYGZhy{}\PYGZhy{}\PYGZhy{}\PYGZhy{}\PYGZhy{}\PYGZhy{}\PYGZhy{}}
    \PYG{n}{ok} \PYG{o}{=} \PYG{n}{ok} \PYG{o+ow}{and} \PYG{n}{a2}   \PYG{o}{\PYGZlt{}}  \PYG{n}{a3}
    \PYG{n}{ok} \PYG{o}{=} \PYG{n}{ok} \PYG{o+ow}{and} \PYG{n}{a2}   \PYG{o}{\PYGZlt{}}\PYG{o}{=} \PYG{n}{a3}
    \PYG{n}{ok} \PYG{o}{=} \PYG{n}{ok} \PYG{o+ow}{and} \PYG{n}{a3} \PYG{o}{\PYGZgt{}}  \PYG{n}{a2}
    \PYG{n}{ok} \PYG{o}{=} \PYG{n}{ok} \PYG{o+ow}{and} \PYG{n}{a3} \PYG{o}{\PYGZgt{}}\PYG{o}{=} \PYG{n}{a2}
    \PYG{n}{ok} \PYG{o}{=} \PYG{n}{ok} \PYG{o+ow}{and} \PYG{n}{a3} \PYG{o}{!=} \PYG{n}{a2}
    \PYG{n}{ok} \PYG{o}{=} \PYG{n}{ok} \PYG{o+ow}{and} \PYG{n}{a3} \PYG{o}{==} \PYG{n}{a3}
    \PYG{c+c1}{\PYGZsh{}}
    \PYG{n}{ok} \PYG{o}{=} \PYG{n}{ok} \PYG{o+ow}{and} \PYG{o+ow}{not} \PYG{p}{(}\PYG{n}{a2} \PYG{o}{\PYGZgt{}}  \PYG{n}{a3}\PYG{p}{)}
    \PYG{n}{ok} \PYG{o}{=} \PYG{n}{ok} \PYG{o+ow}{and} \PYG{o+ow}{not} \PYG{p}{(}\PYG{n}{a2} \PYG{o}{\PYGZgt{}}\PYG{o}{=} \PYG{n}{a3}\PYG{p}{)}
    \PYG{n}{ok} \PYG{o}{=} \PYG{n}{ok} \PYG{o+ow}{and} \PYG{o+ow}{not} \PYG{p}{(}\PYG{n}{a2} \PYG{o}{==} \PYG{n}{a3}\PYG{p}{)}
    \PYG{c+c1}{\PYGZsh{} \PYGZhy{}\PYGZhy{}\PYGZhy{}\PYGZhy{}\PYGZhy{}\PYGZhy{}\PYGZhy{}\PYGZhy{}\PYGZhy{}\PYGZhy{}\PYGZhy{}\PYGZhy{}\PYGZhy{}\PYGZhy{}\PYGZhy{}\PYGZhy{}\PYGZhy{}\PYGZhy{}\PYGZhy{}\PYGZhy{}\PYGZhy{}\PYGZhy{}\PYGZhy{}\PYGZhy{}\PYGZhy{}\PYGZhy{}\PYGZhy{}\PYGZhy{}\PYGZhy{}\PYGZhy{}\PYGZhy{}\PYGZhy{}\PYGZhy{}\PYGZhy{}\PYGZhy{}\PYGZhy{}\PYGZhy{}\PYGZhy{}\PYGZhy{}\PYGZhy{}\PYGZhy{}\PYGZhy{}\PYGZhy{}\PYGZhy{}\PYGZhy{}\PYGZhy{}\PYGZhy{}\PYGZhy{}\PYGZhy{}\PYGZhy{}\PYGZhy{}\PYGZhy{}\PYGZhy{}\PYGZhy{}\PYGZhy{}\PYGZhy{}\PYGZhy{}\PYGZhy{}\PYGZhy{}\PYGZhy{}\PYGZhy{}\PYGZhy{}\PYGZhy{}\PYGZhy{}\PYGZhy{}\PYGZhy{}\PYGZhy{}\PYGZhy{}\PYGZhy{}}
    \PYG{n}{ok} \PYG{o}{=} \PYG{n}{ok} \PYG{o+ow}{and} \PYG{n}{a2}   \PYG{o}{\PYGZlt{}}  \PYG{l+m+mf}{3.0}
    \PYG{n}{ok} \PYG{o}{=} \PYG{n}{ok} \PYG{o+ow}{and} \PYG{n}{a2}   \PYG{o}{\PYGZlt{}}\PYG{o}{=} \PYG{l+m+mf}{3.0}
    \PYG{n}{ok} \PYG{o}{=} \PYG{n}{ok} \PYG{o+ow}{and} \PYG{n}{a3} \PYG{o}{\PYGZgt{}}  \PYG{l+m+mf}{2.0}
    \PYG{n}{ok} \PYG{o}{=} \PYG{n}{ok} \PYG{o+ow}{and} \PYG{n}{a3} \PYG{o}{\PYGZgt{}}\PYG{o}{=} \PYG{l+m+mf}{2.0}
    \PYG{n}{ok} \PYG{o}{=} \PYG{n}{ok} \PYG{o+ow}{and} \PYG{n}{a3} \PYG{o}{!=} \PYG{l+m+mf}{2.0}
    \PYG{n}{ok} \PYG{o}{=} \PYG{n}{ok} \PYG{o+ow}{and} \PYG{n}{a3} \PYG{o}{==} \PYG{l+m+mf}{3.0}
    \PYG{c+c1}{\PYGZsh{}}
    \PYG{n}{ok} \PYG{o}{=} \PYG{n}{ok} \PYG{o+ow}{and} \PYG{o+ow}{not} \PYG{p}{(}\PYG{n}{a2} \PYG{o}{\PYGZgt{}}  \PYG{l+m+mf}{3.0}\PYG{p}{)}
    \PYG{n}{ok} \PYG{o}{=} \PYG{n}{ok} \PYG{o+ow}{and} \PYG{o+ow}{not} \PYG{p}{(}\PYG{n}{a2} \PYG{o}{\PYGZgt{}}\PYG{o}{=} \PYG{l+m+mf}{3.0}\PYG{p}{)}
    \PYG{n}{ok} \PYG{o}{=} \PYG{n}{ok} \PYG{o+ow}{and} \PYG{o+ow}{not} \PYG{p}{(}\PYG{n}{a2} \PYG{o}{==} \PYG{l+m+mf}{3.0}\PYG{p}{)}
    \PYG{c+c1}{\PYGZsh{}}
    \PYG{k}{return}\PYG{p}{(} \PYG{n}{ok} \PYG{p}{)}
\PYG{c+c1}{\PYGZsh{}}
\end{sphinxVerbatim}


\bigskip\hrule\bigskip


xsrst input file: \sphinxcode{\sphinxupquote{example/python/core/a\_double\_compare\_xam.py}}

\index{example@\spxentry{example}}\ignorespaces 

\subparagraph{Example}
\label{\detokenize{xsrst/a_double_compare:example}}\label{\detokenize{xsrst/a_double_compare:a-double-compare-example}}\label{\detokenize{xsrst/a_double_compare:index-7}}
{\hyperref[\detokenize{xsrst/a_double_compare_xam_cpp:id1}]{\sphinxcrossref{\DUrole{std,std-ref}{c++}}}},
{\hyperref[\detokenize{xsrst/a_double_compare_xam_py:id1}]{\sphinxcrossref{\DUrole{std,std-ref}{python}}}}.


\bigskip\hrule\bigskip


xsrst input file: \sphinxcode{\sphinxupquote{lib/cplusplus/a\_double.cpp}}


\subparagraph{a\_double\_assign}
\label{\detokenize{xsrst/a_double_assign:a-double-assign}}\label{\detokenize{xsrst/a_double_assign::doc}}
\index{a\_double\_assign@\spxentry{a\_double\_assign}}\index{a\_double@\spxentry{a\_double}}\index{assignment@\spxentry{assignment}}\index{operators@\spxentry{operators}}\ignorespaces 

\subparagraph{3.2.1.6 a\_double Assignment Operators}
\label{\detokenize{xsrst/a_double_assign:a-double-assignment-operators}}\label{\detokenize{xsrst/a_double_assign:id1}}\begin{itemize}
\item {} 
{\hyperref[\detokenize{xsrst/a_double_assign:a-double-assign-syntax}]{\sphinxcrossref{\DUrole{std,std-ref}{Syntax}}}}

\item {} 
{\hyperref[\detokenize{xsrst/a_double_assign:a-double-assign-op}]{\sphinxcrossref{\DUrole{std,std-ref}{op}}}}

\item {} 
{\hyperref[\detokenize{xsrst/a_double_assign:a-double-assign-ax}]{\sphinxcrossref{\DUrole{std,std-ref}{ax}}}}

\item {} 
{\hyperref[\detokenize{xsrst/a_double_assign:a-double-assign-ay}]{\sphinxcrossref{\DUrole{std,std-ref}{ay}}}}

\item {} 
{\hyperref[\detokenize{xsrst/a_double_assign:a-double-assign-y}]{\sphinxcrossref{\DUrole{std,std-ref}{y}}}}

\item {} 
{\hyperref[\detokenize{xsrst/a_double_assign:a-double-assign-example}]{\sphinxcrossref{\DUrole{std,std-ref}{Example}}}}

\end{itemize}

\index{syntax@\spxentry{syntax}}\ignorespaces 

\subparagraph{Syntax}
\label{\detokenize{xsrst/a_double_assign:syntax}}\label{\detokenize{xsrst/a_double_assign:a-double-assign-syntax}}\label{\detokenize{xsrst/a_double_assign:index-1}}
\begin{DUlineblock}{0em}
\item[] \sphinxstyleemphasis{ax} \sphinxstyleemphasis{op} \sphinxstyleemphasis{ay}
\item[] \sphinxstyleemphasis{aw} \sphinxstyleemphasis{op} \sphinxstyleemphasis{y}
\end{DUlineblock}

\index{op@\spxentry{op}}\ignorespaces 

\subparagraph{op}
\label{\detokenize{xsrst/a_double_assign:op}}\label{\detokenize{xsrst/a_double_assign:a-double-assign-op}}\label{\detokenize{xsrst/a_double_assign:index-2}}
The assignment operator \sphinxstyleemphasis{op} is one of the following:


\begin{savenotes}\sphinxattablestart
\centering
\begin{tabular}[t]{|\X{2}{32}|\X{30}{32}|}
\hline

\sphinxstyleemphasis{op}
&
Meaning
\\
\hline
\sphinxcode{\sphinxupquote{=}}
&
simple assignment
\\
\hline
\sphinxcode{\sphinxupquote{+=}}
&
\sphinxstyleemphasis{ax} = \sphinxstyleemphasis{ax} + ( \sphinxstyleemphasis{ay} \sphinxcode{\sphinxupquote{or}} \sphinxstyleemphasis{y} )
\\
\hline
\sphinxcode{\sphinxupquote{\sphinxhyphen{}=}}
&
\sphinxstyleemphasis{ax} = \sphinxstyleemphasis{ax} \sphinxhyphen{} ( \sphinxstyleemphasis{ay} \sphinxcode{\sphinxupquote{or}} \sphinxstyleemphasis{y} )
\\
\hline
\sphinxcode{\sphinxupquote{*=}}
&
\sphinxstyleemphasis{ax} = \sphinxstyleemphasis{ax} \sphinxcode{\sphinxupquote{*}} \sphinxstyleemphasis{ay} \sphinxcode{\sphinxupquote{or}} \sphinxstyleemphasis{y} )
\\
\hline
\sphinxcode{\sphinxupquote{/=}}
&
\sphinxstyleemphasis{ax} = \sphinxstyleemphasis{ax} \sphinxcode{\sphinxupquote{/}} \sphinxstyleemphasis{ay} \sphinxcode{\sphinxupquote{or}} \sphinxstyleemphasis{y} )
\\
\hline
\end{tabular}
\par
\sphinxattableend\end{savenotes}

\index{ax@\spxentry{ax}}\ignorespaces 

\subparagraph{ax}
\label{\detokenize{xsrst/a_double_assign:ax}}\label{\detokenize{xsrst/a_double_assign:a-double-assign-ax}}\label{\detokenize{xsrst/a_double_assign:index-3}}
This object has prototype

\begin{DUlineblock}{0em}
\item[]         \sphinxcode{\sphinxupquote{const a\_double\&}} \sphinxstyleemphasis{ax}
\end{DUlineblock}

\index{ay@\spxentry{ay}}\ignorespaces 

\subparagraph{ay}
\label{\detokenize{xsrst/a_double_assign:ay}}\label{\detokenize{xsrst/a_double_assign:a-double-assign-ay}}\label{\detokenize{xsrst/a_double_assign:index-4}}
This object has prototype

\begin{DUlineblock}{0em}
\item[]         \sphinxcode{\sphinxupquote{a\_double\&}} \sphinxstyleemphasis{ay}
\end{DUlineblock}

\index{y@\spxentry{y}}\ignorespaces 

\subparagraph{y}
\label{\detokenize{xsrst/a_double_assign:y}}\label{\detokenize{xsrst/a_double_assign:a-double-assign-y}}\label{\detokenize{xsrst/a_double_assign:index-5}}
This object has prototype

\begin{DUlineblock}{0em}
\item[]         \sphinxcode{\sphinxupquote{const double\&}} \sphinxstyleemphasis{y}
\end{DUlineblock}


\subparagraph{a\_double\_assign\_xam\_cpp}
\label{\detokenize{xsrst/a_double_assign_xam_cpp:a-double-assign-xam-cpp}}\label{\detokenize{xsrst/a_double_assign_xam_cpp::doc}}
\index{a\_double\_assign\_xam\_cpp@\spxentry{a\_double\_assign\_xam\_cpp}}\index{c++:@\spxentry{c++:}}\index{a\_double@\spxentry{a\_double}}\index{assignment@\spxentry{assignment}}\index{operators:@\spxentry{operators:}}\index{example@\spxentry{example}}\index{test@\spxentry{test}}\ignorespaces 

\subparagraph{3.2.1.6.1 C++: a\_double Assignment Operators: Example and Test}
\label{\detokenize{xsrst/a_double_assign_xam_cpp:c-a-double-assignment-operators-example-and-test}}\label{\detokenize{xsrst/a_double_assign_xam_cpp:id1}}
\begin{sphinxVerbatim}[commandchars=\\\{\}]
\PYG{c+cp}{\PYGZsh{}}\PYG{c+cp}{ include \PYGZlt{}cstdio\PYGZgt{}}
\PYG{c+cp}{\PYGZsh{}}\PYG{c+cp}{ include \PYGZlt{}cppad}\PYG{c+cp}{/}\PYG{c+cp}{py}\PYG{c+cp}{/}\PYG{c+cp}{cppad\PYGZus{}py.hpp\PYGZgt{}}

\PYG{k+kt}{bool} \PYG{n+nf}{a\PYGZus{}double\PYGZus{}assign\PYGZus{}xam}\PYG{p}{(}\PYG{k+kt}{void}\PYG{p}{)} \PYG{p}{\PYGZob{}}
    \PYG{k}{using} \PYG{n}{cppad\PYGZus{}py}\PYG{o}{:}\PYG{o}{:}\PYG{n}{a\PYGZus{}double}\PYG{p}{;}
    \PYG{c+c1}{//}
    \PYG{c+c1}{// initialize return variable}
    \PYG{k+kt}{bool} \PYG{n}{ok} \PYG{o}{=} \PYG{n+nb}{true}\PYG{p}{;}
    \PYG{c+c1}{//\PYGZhy{}\PYGZhy{}\PYGZhy{}\PYGZhy{}\PYGZhy{}\PYGZhy{}\PYGZhy{}\PYGZhy{}\PYGZhy{}\PYGZhy{}\PYGZhy{}\PYGZhy{}\PYGZhy{}\PYGZhy{}\PYGZhy{}\PYGZhy{}\PYGZhy{}\PYGZhy{}\PYGZhy{}\PYGZhy{}\PYGZhy{}\PYGZhy{}\PYGZhy{}\PYGZhy{}\PYGZhy{}\PYGZhy{}\PYGZhy{}\PYGZhy{}\PYGZhy{}\PYGZhy{}\PYGZhy{}\PYGZhy{}\PYGZhy{}\PYGZhy{}\PYGZhy{}\PYGZhy{}\PYGZhy{}\PYGZhy{}\PYGZhy{}\PYGZhy{}\PYGZhy{}\PYGZhy{}\PYGZhy{}\PYGZhy{}\PYGZhy{}\PYGZhy{}\PYGZhy{}\PYGZhy{}\PYGZhy{}\PYGZhy{}\PYGZhy{}\PYGZhy{}\PYGZhy{}\PYGZhy{}\PYGZhy{}\PYGZhy{}\PYGZhy{}\PYGZhy{}\PYGZhy{}\PYGZhy{}\PYGZhy{}\PYGZhy{}\PYGZhy{}\PYGZhy{}\PYGZhy{}\PYGZhy{}\PYGZhy{}\PYGZhy{}\PYGZhy{}\PYGZhy{}\PYGZhy{}\PYGZhy{}}
    \PYG{n}{a\PYGZus{}double} \PYG{n}{ax}\PYG{p}{(}\PYG{l+m+mf}{2.0}\PYG{p}{)}\PYG{p}{;}
    \PYG{c+c1}{//}
    \PYG{n}{ax} \PYG{o}{=} \PYG{l+m+mf}{3.0}\PYG{p}{;}
    \PYG{n}{ok} \PYG{o}{=} \PYG{n}{ok} \PYG{o}{\PYGZam{}}\PYG{o}{\PYGZam{}} \PYG{n}{ax} \PYG{o}{=}\PYG{o}{=} \PYG{l+m+mf}{3.0}\PYG{p}{;}
    \PYG{c+c1}{//}
    \PYG{n}{ax} \PYG{o}{+}\PYG{o}{=} \PYG{n}{a\PYGZus{}double}\PYG{p}{(}\PYG{l+m+mf}{2.0}\PYG{p}{)}\PYG{p}{;}
    \PYG{n}{ok} \PYG{o}{=} \PYG{n}{ok} \PYG{o}{\PYGZam{}}\PYG{o}{\PYGZam{}} \PYG{n}{ax} \PYG{o}{=}\PYG{o}{=} \PYG{l+m+mf}{5.0}\PYG{p}{;}
    \PYG{c+c1}{//}
    \PYG{n}{ax} \PYG{o}{\PYGZhy{}}\PYG{o}{=} \PYG{l+m+mf}{1.0}\PYG{p}{;}
    \PYG{n}{ok} \PYG{o}{=} \PYG{n}{ok} \PYG{o}{\PYGZam{}}\PYG{o}{\PYGZam{}} \PYG{n}{ax} \PYG{o}{=}\PYG{o}{=} \PYG{l+m+mf}{4.0}\PYG{p}{;}
    \PYG{c+c1}{//}
    \PYG{n}{ax} \PYG{o}{*}\PYG{o}{=} \PYG{n}{a\PYGZus{}double}\PYG{p}{(}\PYG{l+m+mf}{3.0}\PYG{p}{)}\PYG{p}{;}
    \PYG{n}{ok} \PYG{o}{=} \PYG{n}{ok} \PYG{o}{\PYGZam{}}\PYG{o}{\PYGZam{}} \PYG{n}{ax} \PYG{o}{=}\PYG{o}{=} \PYG{l+m+mf}{12.0}\PYG{p}{;}
    \PYG{c+c1}{//}
    \PYG{n}{ax} \PYG{o}{/}\PYG{o}{=} \PYG{l+m+mf}{4.0}\PYG{p}{;}
    \PYG{n}{ok} \PYG{o}{=} \PYG{n}{ok} \PYG{o}{\PYGZam{}}\PYG{o}{\PYGZam{}} \PYG{n}{ax} \PYG{o}{=}\PYG{o}{=} \PYG{n}{a\PYGZus{}double}\PYG{p}{(}\PYG{l+m+mf}{3.0}\PYG{p}{)}\PYG{p}{;}
    \PYG{c+c1}{//}
    \PYG{k}{return}\PYG{p}{(} \PYG{n}{ok} \PYG{p}{)}\PYG{p}{;}
\PYG{p}{\PYGZcb{}}
\end{sphinxVerbatim}


\bigskip\hrule\bigskip


xsrst input file: \sphinxcode{\sphinxupquote{example/cplusplus/a\_double\_assign\_xam.cpp}}


\subparagraph{a\_double\_assign\_xam\_py}
\label{\detokenize{xsrst/a_double_assign_xam_py:a-double-assign-xam-py}}\label{\detokenize{xsrst/a_double_assign_xam_py::doc}}
\index{a\_double\_assign\_xam\_py@\spxentry{a\_double\_assign\_xam\_py}}\index{python:@\spxentry{python:}}\index{a\_double@\spxentry{a\_double}}\index{assignment@\spxentry{assignment}}\index{operators:@\spxentry{operators:}}\index{example@\spxentry{example}}\index{test@\spxentry{test}}\ignorespaces 

\subparagraph{3.2.1.6.2 Python: a\_double Assignment Operators: Example and Test}
\label{\detokenize{xsrst/a_double_assign_xam_py:python-a-double-assignment-operators-example-and-test}}\label{\detokenize{xsrst/a_double_assign_xam_py:id1}}
\begin{sphinxVerbatim}[commandchars=\\\{\}]
\PYG{k}{def} \PYG{n+nf}{a\PYGZus{}double\PYGZus{}assign\PYGZus{}xam}\PYG{p}{(}\PYG{p}{)} \PYG{p}{:}
    \PYG{c+c1}{\PYGZsh{}}
    \PYG{k+kn}{import} \PYG{n+nn}{numpy}
    \PYG{k+kn}{import} \PYG{n+nn}{cppad\PYGZus{}py}
    \PYG{c+c1}{\PYGZsh{}}
    \PYG{c+c1}{\PYGZsh{} initialize return variable}
    \PYG{n}{ok} \PYG{o}{=} \PYG{k+kc}{True}
    \PYG{c+c1}{\PYGZsh{} \PYGZhy{}\PYGZhy{}\PYGZhy{}\PYGZhy{}\PYGZhy{}\PYGZhy{}\PYGZhy{}\PYGZhy{}\PYGZhy{}\PYGZhy{}\PYGZhy{}\PYGZhy{}\PYGZhy{}\PYGZhy{}\PYGZhy{}\PYGZhy{}\PYGZhy{}\PYGZhy{}\PYGZhy{}\PYGZhy{}\PYGZhy{}\PYGZhy{}\PYGZhy{}\PYGZhy{}\PYGZhy{}\PYGZhy{}\PYGZhy{}\PYGZhy{}\PYGZhy{}\PYGZhy{}\PYGZhy{}\PYGZhy{}\PYGZhy{}\PYGZhy{}\PYGZhy{}\PYGZhy{}\PYGZhy{}\PYGZhy{}\PYGZhy{}\PYGZhy{}\PYGZhy{}\PYGZhy{}\PYGZhy{}\PYGZhy{}\PYGZhy{}\PYGZhy{}\PYGZhy{}\PYGZhy{}\PYGZhy{}\PYGZhy{}\PYGZhy{}\PYGZhy{}\PYGZhy{}\PYGZhy{}\PYGZhy{}\PYGZhy{}\PYGZhy{}\PYGZhy{}\PYGZhy{}\PYGZhy{}\PYGZhy{}\PYGZhy{}\PYGZhy{}\PYGZhy{}\PYGZhy{}\PYGZhy{}\PYGZhy{}\PYGZhy{}\PYGZhy{}}
    \PYG{c+c1}{\PYGZsh{}}
    \PYG{n}{ax} \PYG{o}{=} \PYG{n}{cppad\PYGZus{}py}\PYG{o}{.}\PYG{n}{a\PYGZus{}double}\PYG{p}{(}\PYG{l+m+mf}{3.0}\PYG{p}{)}\PYG{p}{;}
    \PYG{n}{ok} \PYG{o}{=} \PYG{n}{ok} \PYG{o+ow}{and} \PYG{n}{ax} \PYG{o}{==} \PYG{l+m+mf}{3.0}
    \PYG{c+c1}{\PYGZsh{}}
    \PYG{n}{ax} \PYG{o}{+}\PYG{o}{=} \PYG{n}{cppad\PYGZus{}py}\PYG{o}{.}\PYG{n}{a\PYGZus{}double}\PYG{p}{(}\PYG{l+m+mf}{2.0}\PYG{p}{)}\PYG{p}{;}
    \PYG{n}{ok} \PYG{o}{=} \PYG{n}{ok} \PYG{o+ow}{and} \PYG{n}{ax} \PYG{o}{==} \PYG{l+m+mf}{5.0}
    \PYG{c+c1}{\PYGZsh{}}
    \PYG{n}{ax} \PYG{o}{\PYGZhy{}}\PYG{o}{=} \PYG{l+m+mf}{1.0}\PYG{p}{;}
    \PYG{n}{ok} \PYG{o}{=} \PYG{n}{ok} \PYG{o+ow}{and} \PYG{n}{ax} \PYG{o}{==} \PYG{n}{cppad\PYGZus{}py}\PYG{o}{.}\PYG{n}{a\PYGZus{}double}\PYG{p}{(}\PYG{l+m+mf}{4.0}\PYG{p}{)}
    \PYG{c+c1}{\PYGZsh{}}
    \PYG{n}{ax} \PYG{o}{*}\PYG{o}{=} \PYG{n}{cppad\PYGZus{}py}\PYG{o}{.}\PYG{n}{a\PYGZus{}double}\PYG{p}{(}\PYG{l+m+mf}{3.0}\PYG{p}{)}\PYG{p}{;}
    \PYG{n}{ok} \PYG{o}{=} \PYG{n}{ok} \PYG{o+ow}{and} \PYG{n}{ax} \PYG{o}{==} \PYG{l+m+mf}{12.0}
    \PYG{c+c1}{\PYGZsh{}}
    \PYG{n}{ax} \PYG{o}{/}\PYG{o}{=} \PYG{l+m+mf}{4.0}\PYG{p}{;}
    \PYG{n}{ok} \PYG{o}{=} \PYG{n}{ok} \PYG{o+ow}{and} \PYG{n}{ax} \PYG{o}{==} \PYG{n}{cppad\PYGZus{}py}\PYG{o}{.}\PYG{n}{a\PYGZus{}double}\PYG{p}{(}\PYG{l+m+mf}{3.0}\PYG{p}{)}
    \PYG{c+c1}{\PYGZsh{}}
    \PYG{k}{return}\PYG{p}{(} \PYG{n}{ok} \PYG{p}{)}
\PYG{c+c1}{\PYGZsh{}}
\end{sphinxVerbatim}


\bigskip\hrule\bigskip


xsrst input file: \sphinxcode{\sphinxupquote{example/python/core/a\_double\_assign\_xam.py}}

\index{example@\spxentry{example}}\ignorespaces 

\subparagraph{Example}
\label{\detokenize{xsrst/a_double_assign:example}}\label{\detokenize{xsrst/a_double_assign:a-double-assign-example}}\label{\detokenize{xsrst/a_double_assign:index-6}}
{\hyperref[\detokenize{xsrst/a_double_assign_xam_cpp:id1}]{\sphinxcrossref{\DUrole{std,std-ref}{c++}}}},
{\hyperref[\detokenize{xsrst/a_double_assign_xam_py:id1}]{\sphinxcrossref{\DUrole{std,std-ref}{python}}}}.


\bigskip\hrule\bigskip


xsrst input file: \sphinxcode{\sphinxupquote{lib/cplusplus/a\_double.cpp}}


\subparagraph{a\_double\_unary\_fun}
\label{\detokenize{xsrst/a_double_unary_fun:a-double-unary-fun}}\label{\detokenize{xsrst/a_double_unary_fun::doc}}
\index{a\_double\_unary\_fun@\spxentry{a\_double\_unary\_fun}}\index{unary@\spxentry{unary}}\index{functions@\spxentry{functions}}\index{with@\spxentry{with}}\index{ad@\spxentry{ad}}\index{result@\spxentry{result}}\ignorespaces 

\subparagraph{3.2.1.7 Unary Functions with AD Result}
\label{\detokenize{xsrst/a_double_unary_fun:unary-functions-with-ad-result}}\label{\detokenize{xsrst/a_double_unary_fun:id1}}\begin{itemize}
\item {} 
{\hyperref[\detokenize{xsrst/a_double_unary_fun:a-double-unary-fun-syntax}]{\sphinxcrossref{\DUrole{std,std-ref}{Syntax}}}}

\item {} 
{\hyperref[\detokenize{xsrst/a_double_unary_fun:a-double-unary-fun-ax}]{\sphinxcrossref{\DUrole{std,std-ref}{ax}}}}

\item {} 
{\hyperref[\detokenize{xsrst/a_double_unary_fun:a-double-unary-fun-fun}]{\sphinxcrossref{\DUrole{std,std-ref}{fun}}}}

\item {} 
{\hyperref[\detokenize{xsrst/a_double_unary_fun:a-double-unary-fun-ay}]{\sphinxcrossref{\DUrole{std,std-ref}{ay}}}}

\item {} 
{\hyperref[\detokenize{xsrst/a_double_unary_fun:a-double-unary-fun-example}]{\sphinxcrossref{\DUrole{std,std-ref}{Example}}}}

\end{itemize}

\index{syntax@\spxentry{syntax}}\ignorespaces 

\subparagraph{Syntax}
\label{\detokenize{xsrst/a_double_unary_fun:syntax}}\label{\detokenize{xsrst/a_double_unary_fun:a-double-unary-fun-syntax}}\label{\detokenize{xsrst/a_double_unary_fun:index-1}}
\begin{DUlineblock}{0em}
\item[] \sphinxstyleemphasis{ay} = \sphinxstyleemphasis{ax} . \sphinxstyleemphasis{fun} ()
\end{DUlineblock}

\index{ax@\spxentry{ax}}\ignorespaces 

\subparagraph{ax}
\label{\detokenize{xsrst/a_double_unary_fun:ax}}\label{\detokenize{xsrst/a_double_unary_fun:a-double-unary-fun-ax}}\label{\detokenize{xsrst/a_double_unary_fun:index-2}}
This object has prototype

\begin{DUlineblock}{0em}
\item[]         \sphinxcode{\sphinxupquote{const a\_double\&}} \sphinxstyleemphasis{ax}
\end{DUlineblock}

This is the argument for the function evaluation.

\index{fun@\spxentry{fun}}\ignorespaces 

\subparagraph{fun}
\label{\detokenize{xsrst/a_double_unary_fun:fun}}\label{\detokenize{xsrst/a_double_unary_fun:a-double-unary-fun-fun}}\label{\detokenize{xsrst/a_double_unary_fun:index-3}}
This specifies which function is being evaluated and is one
of  following value:
\sphinxcode{\sphinxupquote{abs}} ,
\sphinxcode{\sphinxupquote{acos}} ,
\sphinxcode{\sphinxupquote{asin}} ,
\sphinxcode{\sphinxupquote{atan}} ,
\sphinxcode{\sphinxupquote{cos}} ,
\sphinxcode{\sphinxupquote{cosh}} ,
\sphinxcode{\sphinxupquote{erf}} ,
\sphinxcode{\sphinxupquote{exp}} ,
\sphinxcode{\sphinxupquote{fabs}} ,
\sphinxcode{\sphinxupquote{log}} ,
\sphinxcode{\sphinxupquote{sin}} ,
\sphinxcode{\sphinxupquote{sinh}} ,
\sphinxcode{\sphinxupquote{sqrt}} ,
\sphinxcode{\sphinxupquote{tan}} ,
\sphinxcode{\sphinxupquote{tanh}} .
2DO: Add the C++11 functions
asinh, acosh, atanh, expm1, and log1p to this list.

\index{ay@\spxentry{ay}}\ignorespaces 

\subparagraph{ay}
\label{\detokenize{xsrst/a_double_unary_fun:ay}}\label{\detokenize{xsrst/a_double_unary_fun:a-double-unary-fun-ay}}\label{\detokenize{xsrst/a_double_unary_fun:index-4}}
The result object has prototype

\begin{DUlineblock}{0em}
\item[]         \sphinxcode{\sphinxupquote{a\_double}} \sphinxstyleemphasis{ay}
\end{DUlineblock}

and is the specified function evaluated at the specified argument; i.e.,

\begin{DUlineblock}{0em}
\item[]         \sphinxstyleemphasis{ay} = \sphinxstyleemphasis{fun} ( \sphinxstyleemphasis{ax} )
\end{DUlineblock}


\subparagraph{a\_double\_unary\_fun\_xam\_cpp}
\label{\detokenize{xsrst/a_double_unary_fun_xam_cpp:a-double-unary-fun-xam-cpp}}\label{\detokenize{xsrst/a_double_unary_fun_xam_cpp::doc}}
\index{a\_double\_unary\_fun\_xam\_cpp@\spxentry{a\_double\_unary\_fun\_xam\_cpp}}\index{c++:@\spxentry{c++:}}\index{a\_double@\spxentry{a\_double}}\index{unary@\spxentry{unary}}\index{functions@\spxentry{functions}}\index{with@\spxentry{with}}\index{ad@\spxentry{ad}}\index{result:@\spxentry{result:}}\index{example@\spxentry{example}}\index{test@\spxentry{test}}\ignorespaces 

\subparagraph{3.2.1.7.1 C++: a\_double Unary Functions with AD Result: Example and Test}
\label{\detokenize{xsrst/a_double_unary_fun_xam_cpp:c-a-double-unary-functions-with-ad-result-example-and-test}}\label{\detokenize{xsrst/a_double_unary_fun_xam_cpp:id1}}
\begin{sphinxVerbatim}[commandchars=\\\{\}]
\PYG{c+cp}{\PYGZsh{}}\PYG{c+cp}{ include \PYGZlt{}cstdio\PYGZgt{}}
\PYG{c+cp}{\PYGZsh{}}\PYG{c+cp}{ include \PYGZlt{}cppad}\PYG{c+cp}{/}\PYG{c+cp}{py}\PYG{c+cp}{/}\PYG{c+cp}{cppad\PYGZus{}py.hpp\PYGZgt{}}

\PYG{k+kt}{bool} \PYG{n+nf}{a\PYGZus{}double\PYGZus{}unary\PYGZus{}fun\PYGZus{}xam}\PYG{p}{(}\PYG{k+kt}{void}\PYG{p}{)} \PYG{p}{\PYGZob{}}
    \PYG{k}{using} \PYG{n}{cppad\PYGZus{}py}\PYG{o}{:}\PYG{o}{:}\PYG{n}{a\PYGZus{}double}\PYG{p}{;}
    \PYG{c+c1}{//}
    \PYG{c+c1}{// initialize return variable}
    \PYG{k+kt}{bool} \PYG{n}{ok} \PYG{o}{=} \PYG{n+nb}{true}\PYG{p}{;}
    \PYG{c+c1}{//\PYGZhy{}\PYGZhy{}\PYGZhy{}\PYGZhy{}\PYGZhy{}\PYGZhy{}\PYGZhy{}\PYGZhy{}\PYGZhy{}\PYGZhy{}\PYGZhy{}\PYGZhy{}\PYGZhy{}\PYGZhy{}\PYGZhy{}\PYGZhy{}\PYGZhy{}\PYGZhy{}\PYGZhy{}\PYGZhy{}\PYGZhy{}\PYGZhy{}\PYGZhy{}\PYGZhy{}\PYGZhy{}\PYGZhy{}\PYGZhy{}\PYGZhy{}\PYGZhy{}\PYGZhy{}\PYGZhy{}\PYGZhy{}\PYGZhy{}\PYGZhy{}\PYGZhy{}\PYGZhy{}\PYGZhy{}\PYGZhy{}\PYGZhy{}\PYGZhy{}\PYGZhy{}\PYGZhy{}\PYGZhy{}\PYGZhy{}\PYGZhy{}\PYGZhy{}\PYGZhy{}\PYGZhy{}\PYGZhy{}\PYGZhy{}\PYGZhy{}\PYGZhy{}\PYGZhy{}\PYGZhy{}\PYGZhy{}\PYGZhy{}\PYGZhy{}\PYGZhy{}\PYGZhy{}\PYGZhy{}\PYGZhy{}\PYGZhy{}\PYGZhy{}\PYGZhy{}\PYGZhy{}\PYGZhy{}\PYGZhy{}\PYGZhy{}\PYGZhy{}\PYGZhy{}\PYGZhy{}\PYGZhy{}}
    \PYG{c+c1}{//}
    \PYG{n}{a\PYGZus{}double} \PYG{n}{a1}\PYG{p}{(}\PYG{o}{\PYGZhy{}}\PYG{l+m+mf}{1.0}\PYG{p}{)}\PYG{p}{;}
    \PYG{n}{a\PYGZus{}double} \PYG{n}{abs1} \PYG{o}{=} \PYG{n}{a1}\PYG{p}{.}\PYG{n}{abs}\PYG{p}{(}\PYG{p}{)}\PYG{p}{;}
    \PYG{n}{ok}  \PYG{o}{=} \PYG{n}{ok} \PYG{o}{\PYGZam{}}\PYG{o}{\PYGZam{}} \PYG{n}{abs1} \PYG{o}{=}\PYG{o}{=} \PYG{l+m+mf}{1.0}\PYG{p}{;}
    \PYG{c+c1}{//}
    \PYG{n}{a1}   \PYG{o}{=} \PYG{l+m+mf}{1.0}\PYG{p}{;}
    \PYG{n}{abs1} \PYG{o}{=} \PYG{n}{a1}\PYG{p}{.}\PYG{n}{fabs}\PYG{p}{(}\PYG{p}{)}\PYG{p}{;}
    \PYG{n}{ok}   \PYG{o}{=} \PYG{n}{ok} \PYG{o}{\PYGZam{}}\PYG{o}{\PYGZam{}} \PYG{n}{abs1} \PYG{o}{=}\PYG{o}{=} \PYG{l+m+mf}{1.0}\PYG{p}{;}
    \PYG{c+c1}{//}
    \PYG{c+c1}{// pi/4}
    \PYG{n}{a\PYGZus{}double} \PYG{n}{pi\PYGZus{}4} \PYG{o}{=} \PYG{n}{a1}\PYG{p}{.}\PYG{n}{atan}\PYG{p}{(}\PYG{p}{)}\PYG{p}{;}
    \PYG{c+c1}{//}
    \PYG{c+c1}{// sqrt(2)}
    \PYG{n}{a\PYGZus{}double} \PYG{n}{atmp}\PYG{p}{(}\PYG{l+m+mf}{2.0}\PYG{p}{)}\PYG{p}{;}
    \PYG{n}{a\PYGZus{}double} \PYG{n}{r2} \PYG{o}{=} \PYG{n}{atmp}\PYG{p}{.}\PYG{n}{sqrt}\PYG{p}{(}\PYG{p}{)}\PYG{p}{;}
    \PYG{c+c1}{//}
    \PYG{c+c1}{// sin(pi/4)  * sqrt(2) = 1.0;}
    \PYG{n}{atmp} \PYG{o}{=} \PYG{n}{r2} \PYG{o}{*} \PYG{n}{pi\PYGZus{}4}\PYG{p}{.}\PYG{n}{sin}\PYG{p}{(}\PYG{p}{)} \PYG{p}{;}
    \PYG{n}{ok} \PYG{o}{=} \PYG{n}{ok} \PYG{o}{\PYGZam{}}\PYG{o}{\PYGZam{}} \PYG{n}{atmp}\PYG{p}{.}\PYG{n}{near\PYGZus{}equal}\PYG{p}{(}\PYG{n}{a1}\PYG{p}{)} \PYG{p}{;}
    \PYG{c+c1}{//}
    \PYG{c+c1}{// cos(pi/4)  * sqrt(2) = 1.0;}
    \PYG{n}{atmp} \PYG{o}{=} \PYG{n}{r2} \PYG{o}{*} \PYG{n}{pi\PYGZus{}4}\PYG{p}{.}\PYG{n}{cos}\PYG{p}{(}\PYG{p}{)} \PYG{p}{;}
    \PYG{n}{ok} \PYG{o}{=} \PYG{n}{ok} \PYG{o}{\PYGZam{}}\PYG{o}{\PYGZam{}} \PYG{n}{atmp}\PYG{p}{.}\PYG{n}{near\PYGZus{}equal}\PYG{p}{(}\PYG{n}{a1}\PYG{p}{)} \PYG{p}{;}
    \PYG{c+c1}{//}
    \PYG{c+c1}{// tan(pi/4)  = 1.0;}
    \PYG{n}{atmp} \PYG{o}{=} \PYG{n}{pi\PYGZus{}4}\PYG{p}{.}\PYG{n}{tan}\PYG{p}{(}\PYG{p}{)} \PYG{p}{;}
    \PYG{n}{ok} \PYG{o}{=} \PYG{n}{ok} \PYG{o}{\PYGZam{}}\PYG{o}{\PYGZam{}} \PYG{n}{atmp}\PYG{p}{.}\PYG{n}{near\PYGZus{}equal}\PYG{p}{(}\PYG{n}{a1}\PYG{p}{)} \PYG{p}{;}
    \PYG{c+c1}{//}
    \PYG{c+c1}{// erf(0.5) = 0.5204998778130465}
    \PYG{n}{a\PYGZus{}double} \PYG{n}{acheck}\PYG{p}{(}\PYG{l+m+mf}{0.5204998778130465}\PYG{p}{)}\PYG{p}{;}
    \PYG{n}{atmp} \PYG{o}{=}  \PYG{n}{a\PYGZus{}double}\PYG{p}{(}\PYG{l+m+mf}{0.5}\PYG{p}{)}\PYG{p}{;}
    \PYG{n}{atmp} \PYG{o}{=} \PYG{n}{atmp}\PYG{p}{.}\PYG{n}{erf}\PYG{p}{(}\PYG{p}{)}\PYG{p}{;}
    \PYG{n}{ok} \PYG{o}{=} \PYG{n}{ok} \PYG{o}{\PYGZam{}}\PYG{o}{\PYGZam{}} \PYG{n}{atmp}\PYG{p}{.}\PYG{n}{near\PYGZus{}equal}\PYG{p}{(}\PYG{n}{acheck}\PYG{p}{)}\PYG{p}{;}
    \PYG{c+c1}{//}
    \PYG{k}{return}\PYG{p}{(} \PYG{n}{ok} \PYG{p}{)}\PYG{p}{;}
\PYG{p}{\PYGZcb{}}
\end{sphinxVerbatim}


\bigskip\hrule\bigskip


xsrst input file: \sphinxcode{\sphinxupquote{example/cplusplus/a\_double\_unary\_fun\_xam.cpp}}


\subparagraph{a\_double\_unary\_fun\_xam\_py}
\label{\detokenize{xsrst/a_double_unary_fun_xam_py:a-double-unary-fun-xam-py}}\label{\detokenize{xsrst/a_double_unary_fun_xam_py::doc}}
\index{a\_double\_unary\_fun\_xam\_py@\spxentry{a\_double\_unary\_fun\_xam\_py}}\index{python:@\spxentry{python:}}\index{a\_double@\spxentry{a\_double}}\index{unary@\spxentry{unary}}\index{functions@\spxentry{functions}}\index{with@\spxentry{with}}\index{ad@\spxentry{ad}}\index{result:@\spxentry{result:}}\index{example@\spxentry{example}}\index{test@\spxentry{test}}\ignorespaces 

\subparagraph{3.2.1.7.2 Python: a\_double Unary Functions with AD Result: Example and Test}
\label{\detokenize{xsrst/a_double_unary_fun_xam_py:python-a-double-unary-functions-with-ad-result-example-and-test}}\label{\detokenize{xsrst/a_double_unary_fun_xam_py:id1}}
\begin{sphinxVerbatim}[commandchars=\\\{\}]
\PYG{k}{def} \PYG{n+nf}{a\PYGZus{}double\PYGZus{}unary\PYGZus{}fun\PYGZus{}xam}\PYG{p}{(}\PYG{p}{)} \PYG{p}{:}
    \PYG{c+c1}{\PYGZsh{}}
    \PYG{k+kn}{import} \PYG{n+nn}{numpy}
    \PYG{k+kn}{import} \PYG{n+nn}{cppad\PYGZus{}py}
    \PYG{c+c1}{\PYGZsh{}}
    \PYG{c+c1}{\PYGZsh{} initialize return variable}
    \PYG{n}{ok} \PYG{o}{=} \PYG{k+kc}{True}
    \PYG{c+c1}{\PYGZsh{} \PYGZhy{}\PYGZhy{}\PYGZhy{}\PYGZhy{}\PYGZhy{}\PYGZhy{}\PYGZhy{}\PYGZhy{}\PYGZhy{}\PYGZhy{}\PYGZhy{}\PYGZhy{}\PYGZhy{}\PYGZhy{}\PYGZhy{}\PYGZhy{}\PYGZhy{}\PYGZhy{}\PYGZhy{}\PYGZhy{}\PYGZhy{}\PYGZhy{}\PYGZhy{}\PYGZhy{}\PYGZhy{}\PYGZhy{}\PYGZhy{}\PYGZhy{}\PYGZhy{}\PYGZhy{}\PYGZhy{}\PYGZhy{}\PYGZhy{}\PYGZhy{}\PYGZhy{}\PYGZhy{}\PYGZhy{}\PYGZhy{}\PYGZhy{}\PYGZhy{}\PYGZhy{}\PYGZhy{}\PYGZhy{}\PYGZhy{}\PYGZhy{}\PYGZhy{}\PYGZhy{}\PYGZhy{}\PYGZhy{}\PYGZhy{}\PYGZhy{}\PYGZhy{}\PYGZhy{}\PYGZhy{}\PYGZhy{}\PYGZhy{}\PYGZhy{}\PYGZhy{}\PYGZhy{}\PYGZhy{}\PYGZhy{}\PYGZhy{}\PYGZhy{}\PYGZhy{}\PYGZhy{}\PYGZhy{}\PYGZhy{}\PYGZhy{}\PYGZhy{}}
    \PYG{c+c1}{\PYGZsh{}}
    \PYG{c+c1}{\PYGZsh{} abs}
    \PYG{n}{a1}   \PYG{o}{=} \PYG{n}{cppad\PYGZus{}py}\PYG{o}{.}\PYG{n}{a\PYGZus{}double}\PYG{p}{(}\PYG{o}{\PYGZhy{}}\PYG{l+m+mf}{1.0}\PYG{p}{)}
    \PYG{n}{abs1} \PYG{o}{=} \PYG{n}{a1}\PYG{o}{.}\PYG{n}{abs}\PYG{p}{(}\PYG{p}{)}
    \PYG{n}{ok}   \PYG{o}{=} \PYG{n}{ok} \PYG{o+ow}{and} \PYG{n}{abs1} \PYG{o}{==} \PYG{l+m+mf}{1.0}
    \PYG{c+c1}{\PYGZsh{} fabs}
    \PYG{n}{a1}   \PYG{o}{=} \PYG{n}{cppad\PYGZus{}py}\PYG{o}{.}\PYG{n}{a\PYGZus{}double}\PYG{p}{(}\PYG{l+m+mf}{1.0}\PYG{p}{)}
    \PYG{n}{abs1} \PYG{o}{=} \PYG{n}{a1}\PYG{o}{.}\PYG{n}{fabs}\PYG{p}{(}\PYG{p}{)}
    \PYG{n}{ok}   \PYG{o}{=} \PYG{n}{ok} \PYG{o+ow}{and} \PYG{n}{abs1} \PYG{o}{==} \PYG{l+m+mf}{1.0}
    \PYG{c+c1}{\PYGZsh{}}
    \PYG{c+c1}{\PYGZsh{} pi/4}
    \PYG{n}{pi\PYGZus{}4} \PYG{o}{=} \PYG{n}{a1}\PYG{o}{.}\PYG{n}{atan}\PYG{p}{(}\PYG{p}{)}
    \PYG{c+c1}{\PYGZsh{}}
    \PYG{c+c1}{\PYGZsh{} sqrt(2)}
    \PYG{n}{atmp} \PYG{o}{=} \PYG{n}{cppad\PYGZus{}py}\PYG{o}{.}\PYG{n}{a\PYGZus{}double}\PYG{p}{(}\PYG{l+m+mf}{2.0}\PYG{p}{)}
    \PYG{n}{r2} \PYG{o}{=} \PYG{n}{atmp}\PYG{o}{.}\PYG{n}{sqrt}\PYG{p}{(}\PYG{p}{)}
    \PYG{c+c1}{\PYGZsh{}}
    \PYG{c+c1}{\PYGZsh{} sin(pi/4)  * sqrt(2) = 1.0}
    \PYG{n}{atmp} \PYG{o}{=} \PYG{n}{r2} \PYG{o}{*} \PYG{n}{pi\PYGZus{}4}\PYG{o}{.}\PYG{n}{sin}\PYG{p}{(}\PYG{p}{)}
    \PYG{n}{ok} \PYG{o}{=} \PYG{n}{ok} \PYG{o+ow}{and} \PYG{n}{atmp}\PYG{o}{.}\PYG{n}{near\PYGZus{}equal}\PYG{p}{(}\PYG{n}{a1}\PYG{p}{)}
    \PYG{c+c1}{\PYGZsh{}}
    \PYG{c+c1}{\PYGZsh{} cos(pi/4)  * sqrt(2) = 1.0}
    \PYG{n}{atmp} \PYG{o}{=} \PYG{n}{r2} \PYG{o}{*} \PYG{n}{pi\PYGZus{}4}\PYG{o}{.}\PYG{n}{cos}\PYG{p}{(}\PYG{p}{)}
    \PYG{n}{ok} \PYG{o}{=} \PYG{n}{ok} \PYG{o+ow}{and} \PYG{n}{atmp}\PYG{o}{.}\PYG{n}{near\PYGZus{}equal}\PYG{p}{(}\PYG{n}{a1}\PYG{p}{)}
    \PYG{c+c1}{\PYGZsh{}}
    \PYG{c+c1}{\PYGZsh{} tan(pi/4)  = 1.0}
    \PYG{n}{atmp} \PYG{o}{=} \PYG{n}{pi\PYGZus{}4}\PYG{o}{.}\PYG{n}{tan}\PYG{p}{(}\PYG{p}{)}
    \PYG{n}{ok} \PYG{o}{=} \PYG{n}{ok} \PYG{o+ow}{and} \PYG{n}{atmp}\PYG{o}{.}\PYG{n}{near\PYGZus{}equal}\PYG{p}{(}\PYG{n}{a1}\PYG{p}{)}
    \PYG{c+c1}{\PYGZsh{}}
    \PYG{c+c1}{\PYGZsh{} erf(0.5) = 0.5204998778130465}
    \PYG{n}{acheck} \PYG{o}{=} \PYG{n}{cppad\PYGZus{}py}\PYG{o}{.}\PYG{n}{a\PYGZus{}double}\PYG{p}{(}\PYG{l+m+mf}{0.5204998778130465}\PYG{p}{)}
    \PYG{n}{atmp}   \PYG{o}{=} \PYG{n}{cppad\PYGZus{}py}\PYG{o}{.}\PYG{n}{a\PYGZus{}double}\PYG{p}{(}\PYG{l+m+mf}{0.5}\PYG{p}{)}
    \PYG{n}{atmp}   \PYG{o}{=} \PYG{n}{atmp}\PYG{o}{.}\PYG{n}{erf}\PYG{p}{(}\PYG{p}{)}
    \PYG{n}{ok}     \PYG{o}{=} \PYG{n}{ok} \PYG{o+ow}{and} \PYG{n}{atmp}\PYG{o}{.}\PYG{n}{near\PYGZus{}equal}\PYG{p}{(}\PYG{n}{acheck}\PYG{p}{)}
    \PYG{c+c1}{\PYGZsh{}}
    \PYG{k}{return}\PYG{p}{(} \PYG{n}{ok} \PYG{p}{)}
\PYG{c+c1}{\PYGZsh{}}
\end{sphinxVerbatim}


\bigskip\hrule\bigskip


xsrst input file: \sphinxcode{\sphinxupquote{example/python/core/a\_double\_unary\_fun\_xam.py}}

\index{example@\spxentry{example}}\ignorespaces 

\subparagraph{Example}
\label{\detokenize{xsrst/a_double_unary_fun:example}}\label{\detokenize{xsrst/a_double_unary_fun:a-double-unary-fun-example}}\label{\detokenize{xsrst/a_double_unary_fun:index-5}}
{\hyperref[\detokenize{xsrst/a_double_unary_fun_xam_cpp:id1}]{\sphinxcrossref{\DUrole{std,std-ref}{c++}}}},
{\hyperref[\detokenize{xsrst/a_double_unary_fun_xam_py:id1}]{\sphinxcrossref{\DUrole{std,std-ref}{python}}}}.


\bigskip\hrule\bigskip


xsrst input file: \sphinxcode{\sphinxupquote{lib/cplusplus/a\_double.cpp}}


\subparagraph{a\_double\_cond\_assign}
\label{\detokenize{xsrst/a_double_cond_assign:a-double-cond-assign}}\label{\detokenize{xsrst/a_double_cond_assign::doc}}
\index{a\_double\_cond\_assign@\spxentry{a\_double\_cond\_assign}}\index{ad@\spxentry{ad}}\index{conditional@\spxentry{conditional}}\index{assignment@\spxentry{assignment}}\ignorespaces 

\subparagraph{3.2.1.8 AD Conditional Assignment}
\label{\detokenize{xsrst/a_double_cond_assign:ad-conditional-assignment}}\label{\detokenize{xsrst/a_double_cond_assign:id1}}\begin{itemize}
\item {} 
{\hyperref[\detokenize{xsrst/a_double_cond_assign:a-double-cond-assign-syntax}]{\sphinxcrossref{\DUrole{std,std-ref}{Syntax}}}}

\item {} 
{\hyperref[\detokenize{xsrst/a_double_cond_assign:a-double-cond-assign-purpose}]{\sphinxcrossref{\DUrole{std,std-ref}{Purpose}}}}

\item {} 
{\hyperref[\detokenize{xsrst/a_double_cond_assign:a-double-cond-assign-target}]{\sphinxcrossref{\DUrole{std,std-ref}{target}}}}

\item {} 
{\hyperref[\detokenize{xsrst/a_double_cond_assign:a-double-cond-assign-cop}]{\sphinxcrossref{\DUrole{std,std-ref}{cop}}}}

\item {} 
{\hyperref[\detokenize{xsrst/a_double_cond_assign:a-double-cond-assign-left}]{\sphinxcrossref{\DUrole{std,std-ref}{left}}}}

\item {} 
{\hyperref[\detokenize{xsrst/a_double_cond_assign:a-double-cond-assign-right}]{\sphinxcrossref{\DUrole{std,std-ref}{right}}}}

\item {} 
{\hyperref[\detokenize{xsrst/a_double_cond_assign:a-double-cond-assign-if-true}]{\sphinxcrossref{\DUrole{std,std-ref}{if\_true}}}}

\item {} 
{\hyperref[\detokenize{xsrst/a_double_cond_assign:a-double-cond-assign-if-false}]{\sphinxcrossref{\DUrole{std,std-ref}{if\_false}}}}

\item {} 
{\hyperref[\detokenize{xsrst/a_double_cond_assign:a-double-cond-assign-example}]{\sphinxcrossref{\DUrole{std,std-ref}{Example}}}}

\end{itemize}

\index{syntax@\spxentry{syntax}}\ignorespaces 

\subparagraph{Syntax}
\label{\detokenize{xsrst/a_double_cond_assign:syntax}}\label{\detokenize{xsrst/a_double_cond_assign:a-double-cond-assign-syntax}}\label{\detokenize{xsrst/a_double_cond_assign:index-1}}
\begin{DUlineblock}{0em}
\item[] \sphinxstyleemphasis{target} . \sphinxcode{\sphinxupquote{cond\_assign}} ( \sphinxstyleemphasis{cop} , \sphinxstyleemphasis{left} , \sphinxstyleemphasis{right} , \sphinxstyleemphasis{if\_true} , \sphinxstyleemphasis{if\_false} )
\end{DUlineblock}

\index{purpose@\spxentry{purpose}}\ignorespaces 

\subparagraph{Purpose}
\label{\detokenize{xsrst/a_double_cond_assign:purpose}}\label{\detokenize{xsrst/a_double_cond_assign:a-double-cond-assign-purpose}}\label{\detokenize{xsrst/a_double_cond_assign:index-2}}
The code

\begin{DUlineblock}{0em}
\item[]         \sphinxcode{\sphinxupquote{if}} ( \sphinxstyleemphasis{left} \sphinxstyleemphasis{cop} \sphinxstyleemphasis{right} )
\item[]                 \sphinxstyleemphasis{target} = \sphinxstyleemphasis{if\_true}
\item[]         \sphinxcode{\sphinxupquote{else}}
\item[]                 \sphinxstyleemphasis{target} = \sphinxstyleemphasis{if\_false}
\end{DUlineblock}

records either the true or false case depending on the value
of \sphinxstyleemphasis{left} and \sphinxstyleemphasis{right} ; see {\hyperref[\detokenize{xsrst/cpp_fun_ctor:id1}]{\sphinxcrossref{\DUrole{std,std-ref}{cpp\_fun\_ctor}}}}.
If \sphinxstyleemphasis{left} or \sphinxstyleemphasis{right} is a
{\hyperref[\detokenize{xsrst/a_double_property:a-double-property-variable}]{\sphinxcrossref{\DUrole{std,std-ref}{variable}}}},
it may be desirable to switch between \sphinxstyleemphasis{if\_true} and \sphinxstyleemphasis{if\_false}
depending of the value of the independent variable during
calls to order zero {\hyperref[\detokenize{xsrst/cpp_fun_forward:id1}]{\sphinxcrossref{\DUrole{std,std-ref}{cpp\_fun\_forward}}}}.
The \sphinxcode{\sphinxupquote{cond\_assign}} does this.

\index{target@\spxentry{target}}\ignorespaces 

\subparagraph{target}
\label{\detokenize{xsrst/a_double_cond_assign:target}}\label{\detokenize{xsrst/a_double_cond_assign:a-double-cond-assign-target}}\label{\detokenize{xsrst/a_double_cond_assign:index-3}}
This object has prototype

\begin{DUlineblock}{0em}
\item[]         \sphinxcode{\sphinxupquote{a\_double\&}} \sphinxstyleemphasis{target}
\end{DUlineblock}

\index{cop@\spxentry{cop}}\ignorespaces 

\subparagraph{cop}
\label{\detokenize{xsrst/a_double_cond_assign:cop}}\label{\detokenize{xsrst/a_double_cond_assign:a-double-cond-assign-cop}}\label{\detokenize{xsrst/a_double_cond_assign:index-4}}
This argument has prototype

\begin{DUlineblock}{0em}
\item[]         \sphinxcode{\sphinxupquote{const char *cop}}
\end{DUlineblock}

The comparison is

\begin{DUlineblock}{0em}
\item[]         \sphinxstyleemphasis{left} \sphinxstyleemphasis{cop} \sphinxstyleemphasis{right}
\end{DUlineblock}

where \sphinxstyleemphasis{cop} is one of the following:


\begin{savenotes}\sphinxattablestart
\centering
\begin{tabular}[t]{|\X{3}{24}|\X{21}{24}|}
\hline

\sphinxstyleemphasis{cop}
&\\
\hline
\sphinxcode{\sphinxupquote{\textless{}}}
&
less than
\\
\hline
\sphinxcode{\sphinxupquote{\textless{}=}}
&
less than or equal
\\
\hline
\sphinxcode{\sphinxupquote{==}}
&
equal
\\
\hline
\sphinxcode{\sphinxupquote{\textgreater{}=}}
&
greater than or equal
\\
\hline
\sphinxcode{\sphinxupquote{\textgreater{}}}
&
greater than
\\
\hline
\end{tabular}
\par
\sphinxattableend\end{savenotes}

\index{left@\spxentry{left}}\ignorespaces 

\subparagraph{left}
\label{\detokenize{xsrst/a_double_cond_assign:left}}\label{\detokenize{xsrst/a_double_cond_assign:a-double-cond-assign-left}}\label{\detokenize{xsrst/a_double_cond_assign:index-5}}
This argument has prototype

\begin{DUlineblock}{0em}
\item[]         \sphinxcode{\sphinxupquote{const a\_double\&}} \sphinxstyleemphasis{left}
\end{DUlineblock}

It specifies the left operand in the comparison.

\index{right@\spxentry{right}}\ignorespaces 

\subparagraph{right}
\label{\detokenize{xsrst/a_double_cond_assign:right}}\label{\detokenize{xsrst/a_double_cond_assign:a-double-cond-assign-right}}\label{\detokenize{xsrst/a_double_cond_assign:index-6}}
This argument has prototype

\begin{DUlineblock}{0em}
\item[]         \sphinxcode{\sphinxupquote{const a\_double\&}} \sphinxstyleemphasis{right}
\end{DUlineblock}

It specifies the right operand in the comparison.

\index{if\_true@\spxentry{if\_true}}\ignorespaces 

\subparagraph{if\_true}
\label{\detokenize{xsrst/a_double_cond_assign:if-true}}\label{\detokenize{xsrst/a_double_cond_assign:a-double-cond-assign-if-true}}\label{\detokenize{xsrst/a_double_cond_assign:index-7}}
This argument has prototype

\begin{DUlineblock}{0em}
\item[]         \sphinxcode{\sphinxupquote{const a\_double\&}} \sphinxstyleemphasis{if\_true}
\end{DUlineblock}

It specifies the value assigned to \sphinxstyleemphasis{ad} if the result
of the comparison is true.

\index{if\_false@\spxentry{if\_false}}\ignorespaces 

\subparagraph{if\_false}
\label{\detokenize{xsrst/a_double_cond_assign:if-false}}\label{\detokenize{xsrst/a_double_cond_assign:a-double-cond-assign-if-false}}\label{\detokenize{xsrst/a_double_cond_assign:index-8}}
This argument has prototype

\begin{DUlineblock}{0em}
\item[]         \sphinxcode{\sphinxupquote{const a\_double\&}} \sphinxstyleemphasis{if\_false}
\end{DUlineblock}

It specifies the value assigned to \sphinxstyleemphasis{ad} if the result
of the comparison is false.


\subparagraph{a\_double\_cond\_assign\_xam\_cpp}
\label{\detokenize{xsrst/a_double_cond_assign_xam_cpp:a-double-cond-assign-xam-cpp}}\label{\detokenize{xsrst/a_double_cond_assign_xam_cpp::doc}}
\index{a\_double\_cond\_assign\_xam\_cpp@\spxentry{a\_double\_cond\_assign\_xam\_cpp}}\index{c++:@\spxentry{c++:}}\index{a\_double@\spxentry{a\_double}}\index{conditional@\spxentry{conditional}}\index{assignment:@\spxentry{assignment:}}\index{example@\spxentry{example}}\index{test@\spxentry{test}}\ignorespaces 

\subparagraph{3.2.1.8.1 C++: a\_double Conditional Assignment: Example and Test}
\label{\detokenize{xsrst/a_double_cond_assign_xam_cpp:c-a-double-conditional-assignment-example-and-test}}\label{\detokenize{xsrst/a_double_cond_assign_xam_cpp:id1}}
\begin{sphinxVerbatim}[commandchars=\\\{\}]
\PYG{c+cp}{\PYGZsh{}}\PYG{c+cp}{ include \PYGZlt{}cstdio\PYGZgt{}}
\PYG{c+cp}{\PYGZsh{}}\PYG{c+cp}{ include \PYGZlt{}cppad}\PYG{c+cp}{/}\PYG{c+cp}{py}\PYG{c+cp}{/}\PYG{c+cp}{cppad\PYGZus{}py.hpp\PYGZgt{}}

\PYG{k+kt}{bool} \PYG{n+nf}{a\PYGZus{}double\PYGZus{}cond\PYGZus{}assign\PYGZus{}xam}\PYG{p}{(}\PYG{k+kt}{void}\PYG{p}{)} \PYG{p}{\PYGZob{}}
    \PYG{k}{using} \PYG{n}{cppad\PYGZus{}py}\PYG{o}{:}\PYG{o}{:}\PYG{n}{a\PYGZus{}double}\PYG{p}{;}
    \PYG{k}{using} \PYG{n}{cppad\PYGZus{}py}\PYG{o}{:}\PYG{o}{:}\PYG{n}{vec\PYGZus{}double}\PYG{p}{;}
    \PYG{k}{using} \PYG{n}{cppad\PYGZus{}py}\PYG{o}{:}\PYG{o}{:}\PYG{n}{vec\PYGZus{}a\PYGZus{}double}\PYG{p}{;}
    \PYG{k}{using} \PYG{n}{cppad\PYGZus{}py}\PYG{o}{:}\PYG{o}{:}\PYG{n}{d\PYGZus{}fun}\PYG{p}{;}
    \PYG{c+c1}{//}
    \PYG{c+c1}{// initialize return variable}
    \PYG{k+kt}{bool} \PYG{n}{ok} \PYG{o}{=} \PYG{n+nb}{true}\PYG{p}{;}
    \PYG{c+c1}{//\PYGZhy{}\PYGZhy{}\PYGZhy{}\PYGZhy{}\PYGZhy{}\PYGZhy{}\PYGZhy{}\PYGZhy{}\PYGZhy{}\PYGZhy{}\PYGZhy{}\PYGZhy{}\PYGZhy{}\PYGZhy{}\PYGZhy{}\PYGZhy{}\PYGZhy{}\PYGZhy{}\PYGZhy{}\PYGZhy{}\PYGZhy{}\PYGZhy{}\PYGZhy{}\PYGZhy{}\PYGZhy{}\PYGZhy{}\PYGZhy{}\PYGZhy{}\PYGZhy{}\PYGZhy{}\PYGZhy{}\PYGZhy{}\PYGZhy{}\PYGZhy{}\PYGZhy{}\PYGZhy{}\PYGZhy{}\PYGZhy{}\PYGZhy{}\PYGZhy{}\PYGZhy{}\PYGZhy{}\PYGZhy{}\PYGZhy{}\PYGZhy{}\PYGZhy{}\PYGZhy{}\PYGZhy{}\PYGZhy{}\PYGZhy{}\PYGZhy{}\PYGZhy{}\PYGZhy{}\PYGZhy{}\PYGZhy{}\PYGZhy{}\PYGZhy{}\PYGZhy{}\PYGZhy{}\PYGZhy{}\PYGZhy{}\PYGZhy{}\PYGZhy{}\PYGZhy{}\PYGZhy{}\PYGZhy{}\PYGZhy{}\PYGZhy{}\PYGZhy{}\PYGZhy{}\PYGZhy{}\PYGZhy{}}
    \PYG{k+kt}{int} \PYG{n}{n\PYGZus{}ind} \PYG{o}{=} \PYG{l+m+mi}{4}\PYG{p}{;}
    \PYG{k+kt}{int} \PYG{n}{n\PYGZus{}dep} \PYG{o}{=} \PYG{l+m+mi}{1}\PYG{p}{;}
    \PYG{c+c1}{//}
    \PYG{c+c1}{// create ax (value of independent variables does not matter)}
    \PYG{n}{vec\PYGZus{}double} \PYG{n}{x}\PYG{p}{(}\PYG{n}{n\PYGZus{}ind}\PYG{p}{)}\PYG{p}{;}
    \PYG{n}{x}\PYG{p}{[}\PYG{l+m+mi}{0}\PYG{p}{]} \PYG{o}{=} \PYG{l+m+mf}{0.0}\PYG{p}{;}
    \PYG{n}{x}\PYG{p}{[}\PYG{l+m+mi}{1}\PYG{p}{]} \PYG{o}{=} \PYG{l+m+mf}{1.0}\PYG{p}{;}
    \PYG{n}{x}\PYG{p}{[}\PYG{l+m+mi}{2}\PYG{p}{]} \PYG{o}{=} \PYG{l+m+mf}{2.0}\PYG{p}{;}
    \PYG{n}{x}\PYG{p}{[}\PYG{l+m+mi}{3}\PYG{p}{]} \PYG{o}{=} \PYG{l+m+mf}{3.0}\PYG{p}{;}
    \PYG{n}{vec\PYGZus{}a\PYGZus{}double} \PYG{n}{ax} \PYG{o}{=} \PYG{n}{cppad\PYGZus{}py}\PYG{o}{:}\PYG{o}{:}\PYG{n}{independent}\PYG{p}{(}\PYG{n}{x}\PYG{p}{)}\PYG{p}{;}
    \PYG{c+c1}{//}
    \PYG{c+c1}{// arguments to conditional assignment}
    \PYG{n}{a\PYGZus{}double} \PYG{n}{left} \PYG{o}{=} \PYG{n}{ax}\PYG{p}{[}\PYG{l+m+mi}{0}\PYG{p}{]}\PYG{p}{;}
    \PYG{n}{a\PYGZus{}double} \PYG{n}{right} \PYG{o}{=} \PYG{n}{ax}\PYG{p}{[}\PYG{l+m+mi}{1}\PYG{p}{]}\PYG{p}{;}
    \PYG{n}{a\PYGZus{}double} \PYG{n}{if\PYGZus{}true} \PYG{o}{=} \PYG{n}{ax}\PYG{p}{[}\PYG{l+m+mi}{2}\PYG{p}{]}\PYG{p}{;}
    \PYG{n}{a\PYGZus{}double} \PYG{n}{if\PYGZus{}false} \PYG{o}{=} \PYG{n}{ax}\PYG{p}{[}\PYG{l+m+mi}{3}\PYG{p}{]}\PYG{p}{;}
    \PYG{c+c1}{//}
    \PYG{c+c1}{// assignment}
    \PYG{n}{a\PYGZus{}double} \PYG{n}{target}\PYG{p}{;}
    \PYG{n}{target}\PYG{p}{.}\PYG{n}{cond\PYGZus{}assign}\PYG{p}{(}
        \PYG{l+s}{\PYGZdq{}}\PYG{l+s}{\PYGZlt{}}\PYG{l+s}{\PYGZdq{}}\PYG{p}{,}
        \PYG{n}{left}\PYG{p}{,}
        \PYG{n}{right}\PYG{p}{,}
        \PYG{n}{if\PYGZus{}true}\PYG{p}{,}
        \PYG{n}{if\PYGZus{}false}
    \PYG{p}{)}\PYG{p}{;}
    \PYG{c+c1}{//}
    \PYG{c+c1}{// f(x) = taget}
    \PYG{n}{vec\PYGZus{}a\PYGZus{}double} \PYG{n}{ay}\PYG{p}{(}\PYG{n}{n\PYGZus{}dep}\PYG{p}{)}\PYG{p}{;}
    \PYG{n}{ay}\PYG{p}{[}\PYG{l+m+mi}{0}\PYG{p}{]} \PYG{o}{=} \PYG{n}{target}\PYG{p}{;}
    \PYG{n}{d\PYGZus{}fun} \PYG{n}{f}\PYG{p}{(}\PYG{n}{ax}\PYG{p}{,} \PYG{n}{ay}\PYG{p}{)}\PYG{p}{;}
    \PYG{c+c1}{//}
    \PYG{c+c1}{// assignment with different independent variable values}
    \PYG{n}{x}\PYG{p}{[}\PYG{l+m+mi}{0}\PYG{p}{]} \PYG{o}{=} \PYG{l+m+mf}{9.0}\PYG{p}{;} \PYG{c+c1}{// left}
    \PYG{n}{x}\PYG{p}{[}\PYG{l+m+mi}{1}\PYG{p}{]} \PYG{o}{=} \PYG{l+m+mf}{8.0}\PYG{p}{;} \PYG{c+c1}{// right}
    \PYG{n}{x}\PYG{p}{[}\PYG{l+m+mi}{2}\PYG{p}{]} \PYG{o}{=} \PYG{l+m+mf}{7.0}\PYG{p}{;} \PYG{c+c1}{// if\PYGZus{}true}
    \PYG{n}{x}\PYG{p}{[}\PYG{l+m+mi}{3}\PYG{p}{]} \PYG{o}{=} \PYG{l+m+mf}{6.0}\PYG{p}{;} \PYG{c+c1}{// if\PYGZus{}false}
    \PYG{k+kt}{int} \PYG{n}{p} \PYG{o}{=} \PYG{l+m+mi}{0}\PYG{p}{;}
    \PYG{n}{vec\PYGZus{}double} \PYG{n}{y} \PYG{o}{=} \PYG{n}{f}\PYG{p}{.}\PYG{n}{forward}\PYG{p}{(}\PYG{n}{p}\PYG{p}{,} \PYG{n}{x}\PYG{p}{)}\PYG{p}{;}
    \PYG{n}{ok} \PYG{o}{=} \PYG{n}{ok} \PYG{o}{\PYGZam{}}\PYG{o}{\PYGZam{}} \PYG{n}{y}\PYG{p}{[}\PYG{l+m+mi}{0}\PYG{p}{]} \PYG{o}{=}\PYG{o}{=} \PYG{l+m+mf}{6.0}\PYG{p}{;}
    \PYG{c+c1}{//}
    \PYG{k}{return}\PYG{p}{(} \PYG{n}{ok}  \PYG{p}{)}\PYG{p}{;}
\PYG{p}{\PYGZcb{}}
\end{sphinxVerbatim}


\bigskip\hrule\bigskip


xsrst input file: \sphinxcode{\sphinxupquote{example/cplusplus/a\_double\_cond\_assign\_xam.cpp}}


\subparagraph{a\_double\_cond\_assign\_xam\_py}
\label{\detokenize{xsrst/a_double_cond_assign_xam_py:a-double-cond-assign-xam-py}}\label{\detokenize{xsrst/a_double_cond_assign_xam_py::doc}}
\index{a\_double\_cond\_assign\_xam\_py@\spxentry{a\_double\_cond\_assign\_xam\_py}}\index{python:@\spxentry{python:}}\index{a\_double@\spxentry{a\_double}}\index{conditional@\spxentry{conditional}}\index{assignment:@\spxentry{assignment:}}\index{example@\spxentry{example}}\index{test@\spxentry{test}}\ignorespaces 

\subparagraph{3.2.1.8.2 Python: a\_double Conditional Assignment: Example and Test}
\label{\detokenize{xsrst/a_double_cond_assign_xam_py:python-a-double-conditional-assignment-example-and-test}}\label{\detokenize{xsrst/a_double_cond_assign_xam_py:id1}}
\begin{sphinxVerbatim}[commandchars=\\\{\}]
\PYG{k}{def} \PYG{n+nf}{a\PYGZus{}double\PYGZus{}cond\PYGZus{}assign\PYGZus{}xam}\PYG{p}{(}\PYG{p}{)} \PYG{p}{:}
    \PYG{c+c1}{\PYGZsh{}}
    \PYG{k+kn}{import} \PYG{n+nn}{numpy}
    \PYG{k+kn}{import} \PYG{n+nn}{cppad\PYGZus{}py}
    \PYG{c+c1}{\PYGZsh{}}
    \PYG{c+c1}{\PYGZsh{} initialize return variable}
    \PYG{n}{ok} \PYG{o}{=} \PYG{k+kc}{True}
    \PYG{c+c1}{\PYGZsh{} \PYGZhy{}\PYGZhy{}\PYGZhy{}\PYGZhy{}\PYGZhy{}\PYGZhy{}\PYGZhy{}\PYGZhy{}\PYGZhy{}\PYGZhy{}\PYGZhy{}\PYGZhy{}\PYGZhy{}\PYGZhy{}\PYGZhy{}\PYGZhy{}\PYGZhy{}\PYGZhy{}\PYGZhy{}\PYGZhy{}\PYGZhy{}\PYGZhy{}\PYGZhy{}\PYGZhy{}\PYGZhy{}\PYGZhy{}\PYGZhy{}\PYGZhy{}\PYGZhy{}\PYGZhy{}\PYGZhy{}\PYGZhy{}\PYGZhy{}\PYGZhy{}\PYGZhy{}\PYGZhy{}\PYGZhy{}\PYGZhy{}\PYGZhy{}\PYGZhy{}\PYGZhy{}\PYGZhy{}\PYGZhy{}\PYGZhy{}\PYGZhy{}\PYGZhy{}\PYGZhy{}\PYGZhy{}\PYGZhy{}\PYGZhy{}\PYGZhy{}\PYGZhy{}\PYGZhy{}\PYGZhy{}\PYGZhy{}\PYGZhy{}\PYGZhy{}\PYGZhy{}\PYGZhy{}\PYGZhy{}\PYGZhy{}\PYGZhy{}\PYGZhy{}\PYGZhy{}\PYGZhy{}\PYGZhy{}\PYGZhy{}\PYGZhy{}\PYGZhy{}}
    \PYG{n}{n\PYGZus{}ind} \PYG{o}{=} \PYG{l+m+mi}{4}
    \PYG{n}{n\PYGZus{}dep} \PYG{o}{=} \PYG{l+m+mi}{1}
    \PYG{c+c1}{\PYGZsh{}}
    \PYG{c+c1}{\PYGZsh{} create ax (value of independent variables does not matter)}
    \PYG{n}{x} \PYG{o}{=} \PYG{n}{numpy}\PYG{o}{.}\PYG{n}{empty}\PYG{p}{(}\PYG{n}{n\PYGZus{}ind}\PYG{p}{,} \PYG{n}{dtype}\PYG{o}{=}\PYG{n+nb}{float}\PYG{p}{)}
    \PYG{n}{x}\PYG{p}{[}\PYG{l+m+mi}{0}\PYG{p}{]} \PYG{o}{=} \PYG{l+m+mf}{0.0}
    \PYG{n}{x}\PYG{p}{[}\PYG{l+m+mi}{1}\PYG{p}{]} \PYG{o}{=} \PYG{l+m+mf}{1.0}
    \PYG{n}{x}\PYG{p}{[}\PYG{l+m+mi}{2}\PYG{p}{]} \PYG{o}{=} \PYG{l+m+mf}{2.0}
    \PYG{n}{x}\PYG{p}{[}\PYG{l+m+mi}{3}\PYG{p}{]} \PYG{o}{=} \PYG{l+m+mf}{3.0}
    \PYG{n}{ax} \PYG{o}{=} \PYG{n}{cppad\PYGZus{}py}\PYG{o}{.}\PYG{n}{independent}\PYG{p}{(}\PYG{n}{x}\PYG{p}{)}
    \PYG{c+c1}{\PYGZsh{}}
    \PYG{c+c1}{\PYGZsh{} arguments to conditional assignment}
    \PYG{n}{left} \PYG{o}{=} \PYG{n}{ax}\PYG{p}{[}\PYG{l+m+mi}{0}\PYG{p}{]}
    \PYG{n}{right} \PYG{o}{=} \PYG{n}{ax}\PYG{p}{[}\PYG{l+m+mi}{1}\PYG{p}{]}
    \PYG{n}{if\PYGZus{}true} \PYG{o}{=} \PYG{n}{ax}\PYG{p}{[}\PYG{l+m+mi}{2}\PYG{p}{]}
    \PYG{n}{if\PYGZus{}false} \PYG{o}{=} \PYG{n}{ax}\PYG{p}{[}\PYG{l+m+mi}{3}\PYG{p}{]}
    \PYG{c+c1}{\PYGZsh{}}
    \PYG{c+c1}{\PYGZsh{} assignment}
    \PYG{n}{target} \PYG{o}{=} \PYG{n}{cppad\PYGZus{}py}\PYG{o}{.}\PYG{n}{a\PYGZus{}double}\PYG{p}{(}\PYG{p}{)}
    \PYG{n}{target}\PYG{o}{.}\PYG{n}{cond\PYGZus{}assign}\PYG{p}{(}
        \PYG{l+s+s2}{\PYGZdq{}}\PYG{l+s+s2}{\PYGZlt{}}\PYG{l+s+s2}{\PYGZdq{}}\PYG{p}{,}
        \PYG{n}{left}\PYG{p}{,}
        \PYG{n}{right}\PYG{p}{,}
        \PYG{n}{if\PYGZus{}true}\PYG{p}{,}
        \PYG{n}{if\PYGZus{}false}
    \PYG{p}{)}
    \PYG{c+c1}{\PYGZsh{}}
    \PYG{c+c1}{\PYGZsh{} f(x) = taget}
    \PYG{n}{ay} \PYG{o}{=} \PYG{n}{numpy}\PYG{o}{.}\PYG{n}{empty}\PYG{p}{(}\PYG{n}{n\PYGZus{}dep}\PYG{p}{,} \PYG{n}{dtype}\PYG{o}{=}\PYG{n}{cppad\PYGZus{}py}\PYG{o}{.}\PYG{n}{a\PYGZus{}double}\PYG{p}{)}
    \PYG{n}{ay}\PYG{p}{[}\PYG{l+m+mi}{0}\PYG{p}{]} \PYG{o}{=} \PYG{n}{target}
    \PYG{n}{f}  \PYG{o}{=} \PYG{n}{cppad\PYGZus{}py}\PYG{o}{.}\PYG{n}{d\PYGZus{}fun}\PYG{p}{(}\PYG{n}{ax}\PYG{p}{,} \PYG{n}{ay}\PYG{p}{)}
    \PYG{c+c1}{\PYGZsh{}}
    \PYG{c+c1}{\PYGZsh{} assignment with different independent variable values}
    \PYG{n}{x}\PYG{p}{[}\PYG{l+m+mi}{0}\PYG{p}{]} \PYG{o}{=} \PYG{l+m+mf}{9.0} \PYG{c+c1}{\PYGZsh{} left}
    \PYG{n}{x}\PYG{p}{[}\PYG{l+m+mi}{1}\PYG{p}{]} \PYG{o}{=} \PYG{l+m+mf}{8.0} \PYG{c+c1}{\PYGZsh{} right}
    \PYG{n}{x}\PYG{p}{[}\PYG{l+m+mi}{2}\PYG{p}{]} \PYG{o}{=} \PYG{l+m+mf}{7.0} \PYG{c+c1}{\PYGZsh{} if\PYGZus{}true}
    \PYG{n}{x}\PYG{p}{[}\PYG{l+m+mi}{3}\PYG{p}{]} \PYG{o}{=} \PYG{l+m+mf}{6.0} \PYG{c+c1}{\PYGZsh{} if\PYGZus{}false}
    \PYG{n}{p} \PYG{o}{=} \PYG{l+m+mi}{0}
    \PYG{n}{y} \PYG{o}{=} \PYG{n}{f}\PYG{o}{.}\PYG{n}{forward}\PYG{p}{(}\PYG{n}{p}\PYG{p}{,} \PYG{n}{x}\PYG{p}{)}
    \PYG{n}{ok} \PYG{o}{=} \PYG{n}{ok} \PYG{o+ow}{and} \PYG{n}{y}\PYG{p}{[}\PYG{l+m+mi}{0}\PYG{p}{]} \PYG{o}{==} \PYG{l+m+mf}{6.0}
    \PYG{c+c1}{\PYGZsh{}}
    \PYG{k}{return}\PYG{p}{(} \PYG{n}{ok}  \PYG{p}{)}
\PYG{c+c1}{\PYGZsh{}}
\end{sphinxVerbatim}


\bigskip\hrule\bigskip


xsrst input file: \sphinxcode{\sphinxupquote{example/python/core/a\_double\_cond\_assign\_xam.py}}

\index{example@\spxentry{example}}\ignorespaces 

\subparagraph{Example}
\label{\detokenize{xsrst/a_double_cond_assign:example}}\label{\detokenize{xsrst/a_double_cond_assign:a-double-cond-assign-example}}\label{\detokenize{xsrst/a_double_cond_assign:index-9}}
{\hyperref[\detokenize{xsrst/a_double_cond_assign_xam_cpp:id1}]{\sphinxcrossref{\DUrole{std,std-ref}{c++}}}},
{\hyperref[\detokenize{xsrst/a_double_cond_assign_xam_py:id1}]{\sphinxcrossref{\DUrole{std,std-ref}{python}}}}.


\bigskip\hrule\bigskip


xsrst input file: \sphinxcode{\sphinxupquote{lib/cplusplus/a\_double.cpp}}
\begin{enumerate}
\sphinxsetlistlabels{\arabic}{enumi}{enumii}{}{.}%
\item {} 
a\_double\_ctor: {\hyperref[\detokenize{xsrst/a_double_ctor:id1}]{\sphinxcrossref{\DUrole{std,std-ref}{3.2.1.1 The a\_double Constructor}}}}

\item {} 
a\_double\_unary\_op: {\hyperref[\detokenize{xsrst/a_double_unary_op:id1}]{\sphinxcrossref{\DUrole{std,std-ref}{3.2.1.2 a\_double Unary Plus and Minus}}}}

\item {} 
a\_double\_property: {\hyperref[\detokenize{xsrst/a_double_property:id1}]{\sphinxcrossref{\DUrole{std,std-ref}{3.2.1.3 Properties of an a\_double Object}}}}

\item {} 
a\_double\_binary: {\hyperref[\detokenize{xsrst/a_double_binary:id1}]{\sphinxcrossref{\DUrole{std,std-ref}{3.2.1.4 a\_double Binary Operators with an AD Result}}}}

\item {} 
a\_double\_compare: {\hyperref[\detokenize{xsrst/a_double_compare:id1}]{\sphinxcrossref{\DUrole{std,std-ref}{3.2.1.5 a\_double Comparison Operators}}}}

\item {} 
a\_double\_assign: {\hyperref[\detokenize{xsrst/a_double_assign:id1}]{\sphinxcrossref{\DUrole{std,std-ref}{3.2.1.6 a\_double Assignment Operators}}}}

\item {} 
a\_double\_unary\_fun: {\hyperref[\detokenize{xsrst/a_double_unary_fun:id1}]{\sphinxcrossref{\DUrole{std,std-ref}{3.2.1.7 Unary Functions with AD Result}}}}

\item {} 
a\_double\_cond\_assign: {\hyperref[\detokenize{xsrst/a_double_cond_assign:id1}]{\sphinxcrossref{\DUrole{std,std-ref}{3.2.1.8 AD Conditional Assignment}}}}

\end{enumerate}


\bigskip\hrule\bigskip


xsrst input file: \sphinxcode{\sphinxupquote{lib/cplusplus/cplusplus.xsrst}}


\subparagraph{vector}
\label{\detokenize{xsrst/vector:vector}}\label{\detokenize{xsrst/vector::doc}}
\index{vector@\spxentry{vector}}\index{cppad@\spxentry{cppad}}\index{py@\spxentry{py}}\index{vectors@\spxentry{vectors}}\ignorespaces 

\subparagraph{3.2.2 Cppad Py Vectors}
\label{\detokenize{xsrst/vector:cppad-py-vectors}}\label{\detokenize{xsrst/vector:id1}}\begin{itemize}
\item {} 
{\hyperref[\detokenize{xsrst/vector:vector-children}]{\sphinxcrossref{\DUrole{std,std-ref}{Children}}}}

\end{itemize}

\index{children@\spxentry{children}}\ignorespaces 

\subparagraph{Children}
\label{\detokenize{xsrst/vector:children}}\label{\detokenize{xsrst/vector:vector-children}}\label{\detokenize{xsrst/vector:index-1}}

\subparagraph{vector\_ctor}
\label{\detokenize{xsrst/vector_ctor:vector-ctor}}\label{\detokenize{xsrst/vector_ctor::doc}}
\index{vector\_ctor@\spxentry{vector\_ctor}}\index{cppad@\spxentry{cppad}}\index{py@\spxentry{py}}\index{vector@\spxentry{vector}}\index{constructors@\spxentry{constructors}}\ignorespaces 

\subparagraph{3.2.2.1 Cppad Py Vector Constructors}
\label{\detokenize{xsrst/vector_ctor:cppad-py-vector-constructors}}\label{\detokenize{xsrst/vector_ctor:id1}}\begin{itemize}
\item {} 
{\hyperref[\detokenize{xsrst/vector_ctor:vector-ctor-syntax}]{\sphinxcrossref{\DUrole{std,std-ref}{Syntax}}}}

\item {} 
{\hyperref[\detokenize{xsrst/vector_ctor:vector-ctor-purpose}]{\sphinxcrossref{\DUrole{std,std-ref}{Purpose}}}}

\item {} 
{\hyperref[\detokenize{xsrst/vector_ctor:vector-ctor-n}]{\sphinxcrossref{\DUrole{std,std-ref}{n}}}}

\item {} 
{\hyperref[\detokenize{xsrst/vector_ctor:vector-ctor-vec-bool}]{\sphinxcrossref{\DUrole{std,std-ref}{vec\_bool}}}}

\item {} 
{\hyperref[\detokenize{xsrst/vector_ctor:vector-ctor-vec-int}]{\sphinxcrossref{\DUrole{std,std-ref}{vec\_int}}}}

\item {} 
{\hyperref[\detokenize{xsrst/vector_ctor:vector-ctor-vec-double}]{\sphinxcrossref{\DUrole{std,std-ref}{vec\_double}}}}

\item {} 
{\hyperref[\detokenize{xsrst/vector_ctor:vector-ctor-vec-a-double}]{\sphinxcrossref{\DUrole{std,std-ref}{vec\_a\_double}}}}

\item {} 
{\hyperref[\detokenize{xsrst/vector_ctor:vector-ctor-example}]{\sphinxcrossref{\DUrole{std,std-ref}{Example}}}}

\end{itemize}

\index{syntax@\spxentry{syntax}}\ignorespaces 

\subparagraph{Syntax}
\label{\detokenize{xsrst/vector_ctor:syntax}}\label{\detokenize{xsrst/vector_ctor:vector-ctor-syntax}}\label{\detokenize{xsrst/vector_ctor:index-1}}
\begin{DUlineblock}{0em}
\item[] \sphinxstyleemphasis{bv} =  \sphinxcode{\sphinxupquote{cppad\_py.vec\_bool}} ( \sphinxstyleemphasis{n} )
\item[] \sphinxstyleemphasis{iv} =  \sphinxcode{\sphinxupquote{cppad\_py.vec\_int}} ( \sphinxstyleemphasis{n} )
\item[] \sphinxstyleemphasis{dv} =  \sphinxcode{\sphinxupquote{cppad\_py.vec\_double}} ( \sphinxstyleemphasis{n} )
\item[] \sphinxstyleemphasis{av} =  \sphinxcode{\sphinxupquote{cppad\_py.vec\_a\_double}} ( \sphinxstyleemphasis{n} )
\end{DUlineblock}

\index{purpose@\spxentry{purpose}}\ignorespaces 

\subparagraph{Purpose}
\label{\detokenize{xsrst/vector_ctor:purpose}}\label{\detokenize{xsrst/vector_ctor:vector-ctor-purpose}}\label{\detokenize{xsrst/vector_ctor:index-2}}
Creates a vector with \sphinxstyleemphasis{n} elements.

\index{n@\spxentry{n}}\ignorespaces 

\subparagraph{n}
\label{\detokenize{xsrst/vector_ctor:n}}\label{\detokenize{xsrst/vector_ctor:vector-ctor-n}}\label{\detokenize{xsrst/vector_ctor:index-3}}
The argument \sphinxstyleemphasis{n} is a non\sphinxhyphen{}negative integer with default value zero;
i.e., if it is not present, zero is used.

\index{vec\_bool@\spxentry{vec\_bool}}\ignorespaces 

\subparagraph{vec\_bool}
\label{\detokenize{xsrst/vector_ctor:vec-bool}}\label{\detokenize{xsrst/vector_ctor:vector-ctor-vec-bool}}\label{\detokenize{xsrst/vector_ctor:index-4}}
This result \sphinxstyleemphasis{bv} is a vector with elements of type \sphinxcode{\sphinxupquote{bool}}

\index{vec\_int@\spxentry{vec\_int}}\ignorespaces 

\subparagraph{vec\_int}
\label{\detokenize{xsrst/vector_ctor:vec-int}}\label{\detokenize{xsrst/vector_ctor:vector-ctor-vec-int}}\label{\detokenize{xsrst/vector_ctor:index-5}}
This result \sphinxstyleemphasis{bv} is a vector with elements of type \sphinxcode{\sphinxupquote{int}}

\index{vec\_double@\spxentry{vec\_double}}\ignorespaces 

\subparagraph{vec\_double}
\label{\detokenize{xsrst/vector_ctor:vec-double}}\label{\detokenize{xsrst/vector_ctor:vector-ctor-vec-double}}\label{\detokenize{xsrst/vector_ctor:index-6}}
This result \sphinxstyleemphasis{bv} is a vector with elements of type \sphinxcode{\sphinxupquote{double}}

\index{vec\_a\_double@\spxentry{vec\_a\_double}}\ignorespaces 

\subparagraph{vec\_a\_double}
\label{\detokenize{xsrst/vector_ctor:vec-a-double}}\label{\detokenize{xsrst/vector_ctor:vector-ctor-vec-a-double}}\label{\detokenize{xsrst/vector_ctor:index-7}}
This result \sphinxcode{\sphinxupquote{av}} is a vector with elements of type {\hyperref[\detokenize{xsrst/a_double:id1}]{\sphinxcrossref{\DUrole{std,std-ref}{a\_double}}}}.

\index{example@\spxentry{example}}\ignorespaces 

\subparagraph{Example}
\label{\detokenize{xsrst/vector_ctor:example}}\label{\detokenize{xsrst/vector_ctor:vector-ctor-example}}\label{\detokenize{xsrst/vector_ctor:index-8}}
All of the other vector examples use the vector constructors.


\bigskip\hrule\bigskip


xsrst input file: \sphinxcode{\sphinxupquote{lib/cplusplus/vector.xsrst}}


\subparagraph{vector\_size}
\label{\detokenize{xsrst/vector_size:vector-size}}\label{\detokenize{xsrst/vector_size::doc}}
\index{vector\_size@\spxentry{vector\_size}}\index{size@\spxentry{size}}\index{vector@\spxentry{vector}}\ignorespaces 

\subparagraph{3.2.2.2 Size of a Vector}
\label{\detokenize{xsrst/vector_size:size-of-a-vector}}\label{\detokenize{xsrst/vector_size:id1}}\begin{itemize}
\item {} 
{\hyperref[\detokenize{xsrst/vector_size:vector-size-syntax}]{\sphinxcrossref{\DUrole{std,std-ref}{Syntax}}}}

\item {} 
{\hyperref[\detokenize{xsrst/vector_size:vector-size-v}]{\sphinxcrossref{\DUrole{std,std-ref}{v}}}}

\item {} 
{\hyperref[\detokenize{xsrst/vector_size:vector-size-n}]{\sphinxcrossref{\DUrole{std,std-ref}{n}}}}

\item {} 
{\hyperref[\detokenize{xsrst/vector_size:vector-size-example}]{\sphinxcrossref{\DUrole{std,std-ref}{Example}}}}

\end{itemize}

\index{syntax@\spxentry{syntax}}\ignorespaces 

\subparagraph{Syntax}
\label{\detokenize{xsrst/vector_size:syntax}}\label{\detokenize{xsrst/vector_size:vector-size-syntax}}\label{\detokenize{xsrst/vector_size:index-1}}
\sphinxstyleemphasis{n} = \sphinxstyleemphasis{v} . \sphinxcode{\sphinxupquote{size}} ()

\index{v@\spxentry{v}}\ignorespaces 

\subparagraph{v}
\label{\detokenize{xsrst/vector_size:v}}\label{\detokenize{xsrst/vector_size:vector-size-v}}\label{\detokenize{xsrst/vector_size:index-2}}
The object \sphinxstyleemphasis{v} has one of the following prototypes:

\begin{DUlineblock}{0em}
\item[]         \sphinxcode{\sphinxupquote{const vec\_bool\&}} \sphinxstyleemphasis{v}
\item[]         \sphinxcode{\sphinxupquote{const vec\_int\&}} \sphinxstyleemphasis{v}
\item[]         \sphinxcode{\sphinxupquote{const vec\_double\&}} \sphinxstyleemphasis{v}
\item[]         \sphinxcode{\sphinxupquote{const vec\_a\_double\&}} \sphinxstyleemphasis{v}
\end{DUlineblock}

\index{n@\spxentry{n}}\ignorespaces 

\subparagraph{n}
\label{\detokenize{xsrst/vector_size:n}}\label{\detokenize{xsrst/vector_size:vector-size-n}}\label{\detokenize{xsrst/vector_size:index-3}}
The result has prototype

\begin{DUlineblock}{0em}
\item[]         \sphinxcode{\sphinxupquote{size\_t}} \sphinxstyleemphasis{n}
\end{DUlineblock}

i.e., it is a positive integer.
Its value is the number of elements in the vector \sphinxstyleemphasis{v} .


\subparagraph{vector\_size\_xam\_cpp}
\label{\detokenize{xsrst/vector_size_xam_cpp:vector-size-xam-cpp}}\label{\detokenize{xsrst/vector_size_xam_cpp::doc}}
\index{vector\_size\_xam\_cpp@\spxentry{vector\_size\_xam\_cpp}}\index{c++:@\spxentry{c++:}}\index{size@\spxentry{size}}\index{vectors:@\spxentry{vectors:}}\index{example@\spxentry{example}}\index{test@\spxentry{test}}\ignorespaces 

\subparagraph{3.2.2.2.1 C++: Size of Vectors: Example and Test}
\label{\detokenize{xsrst/vector_size_xam_cpp:c-size-of-vectors-example-and-test}}\label{\detokenize{xsrst/vector_size_xam_cpp:id1}}
\begin{sphinxVerbatim}[commandchars=\\\{\}]
\PYG{c+cp}{\PYGZsh{}}\PYG{c+cp}{ include \PYGZlt{}cstdio\PYGZgt{}}
\PYG{c+cp}{\PYGZsh{}}\PYG{c+cp}{ include \PYGZlt{}cppad}\PYG{c+cp}{/}\PYG{c+cp}{py}\PYG{c+cp}{/}\PYG{c+cp}{cppad\PYGZus{}py.hpp\PYGZgt{}}

\PYG{k+kt}{bool} \PYG{n+nf}{vector\PYGZus{}size\PYGZus{}xam}\PYG{p}{(}\PYG{k+kt}{void}\PYG{p}{)} \PYG{p}{\PYGZob{}}
    \PYG{k}{using} \PYG{n}{cppad\PYGZus{}py}\PYG{o}{:}\PYG{o}{:}\PYG{n}{a\PYGZus{}double}\PYG{p}{;}
    \PYG{k}{using} \PYG{n}{cppad\PYGZus{}py}\PYG{o}{:}\PYG{o}{:}\PYG{n}{vec\PYGZus{}bool}\PYG{p}{;}
    \PYG{k}{using} \PYG{n}{cppad\PYGZus{}py}\PYG{o}{:}\PYG{o}{:}\PYG{n}{vec\PYGZus{}int}\PYG{p}{;}
    \PYG{k}{using} \PYG{n}{cppad\PYGZus{}py}\PYG{o}{:}\PYG{o}{:}\PYG{n}{vec\PYGZus{}double}\PYG{p}{;}
    \PYG{k}{using} \PYG{n}{cppad\PYGZus{}py}\PYG{o}{:}\PYG{o}{:}\PYG{n}{vec\PYGZus{}a\PYGZus{}double}\PYG{p}{;}
    \PYG{c+c1}{//}
    \PYG{c+c1}{// initialize return variable}
    \PYG{k+kt}{bool} \PYG{n}{ok} \PYG{o}{=} \PYG{n+nb}{true}\PYG{p}{;}
    \PYG{c+c1}{//\PYGZhy{}\PYGZhy{}\PYGZhy{}\PYGZhy{}\PYGZhy{}\PYGZhy{}\PYGZhy{}\PYGZhy{}\PYGZhy{}\PYGZhy{}\PYGZhy{}\PYGZhy{}\PYGZhy{}\PYGZhy{}\PYGZhy{}\PYGZhy{}\PYGZhy{}\PYGZhy{}\PYGZhy{}\PYGZhy{}\PYGZhy{}\PYGZhy{}\PYGZhy{}\PYGZhy{}\PYGZhy{}\PYGZhy{}\PYGZhy{}\PYGZhy{}\PYGZhy{}\PYGZhy{}\PYGZhy{}\PYGZhy{}\PYGZhy{}\PYGZhy{}\PYGZhy{}\PYGZhy{}\PYGZhy{}\PYGZhy{}\PYGZhy{}\PYGZhy{}\PYGZhy{}\PYGZhy{}\PYGZhy{}\PYGZhy{}\PYGZhy{}\PYGZhy{}\PYGZhy{}\PYGZhy{}\PYGZhy{}\PYGZhy{}\PYGZhy{}\PYGZhy{}\PYGZhy{}\PYGZhy{}\PYGZhy{}\PYGZhy{}\PYGZhy{}\PYGZhy{}\PYGZhy{}\PYGZhy{}\PYGZhy{}\PYGZhy{}\PYGZhy{}\PYGZhy{}\PYGZhy{}\PYGZhy{}\PYGZhy{}\PYGZhy{}\PYGZhy{}\PYGZhy{}\PYGZhy{}\PYGZhy{}}
    \PYG{c+c1}{// create vectors}
    \PYG{n}{vec\PYGZus{}bool} \PYG{n}{bv} \PYG{o}{=} \PYG{n}{vec\PYGZus{}bool}\PYG{p}{(}\PYG{p}{)}\PYG{p}{;}
    \PYG{n}{vec\PYGZus{}int} \PYG{n}{iv} \PYG{o}{=} \PYG{n}{vec\PYGZus{}int}\PYG{p}{(}\PYG{l+m+mi}{1}\PYG{p}{)}\PYG{p}{;}
    \PYG{n}{vec\PYGZus{}double} \PYG{n}{dv}\PYG{p}{(}\PYG{l+m+mi}{2}\PYG{p}{)}\PYG{p}{;}
    \PYG{n}{vec\PYGZus{}a\PYGZus{}double} \PYG{n}{av}\PYG{p}{(}\PYG{l+m+mi}{3}\PYG{p}{)}\PYG{p}{;}
    \PYG{c+c1}{//}
    \PYG{c+c1}{// check size of vectors}
    \PYG{n}{ok} \PYG{o}{=} \PYG{n}{ok} \PYG{o}{\PYGZam{}}\PYG{o}{\PYGZam{}} \PYG{n}{bv}\PYG{p}{.}\PYG{n}{size}\PYG{p}{(}\PYG{p}{)} \PYG{o}{=}\PYG{o}{=} \PYG{l+m+mi}{0} \PYG{p}{;}
    \PYG{n}{ok} \PYG{o}{=} \PYG{n}{ok} \PYG{o}{\PYGZam{}}\PYG{o}{\PYGZam{}} \PYG{n}{iv}\PYG{p}{.}\PYG{n}{size}\PYG{p}{(}\PYG{p}{)} \PYG{o}{=}\PYG{o}{=} \PYG{l+m+mi}{1} \PYG{p}{;}
    \PYG{n}{ok} \PYG{o}{=} \PYG{n}{ok} \PYG{o}{\PYGZam{}}\PYG{o}{\PYGZam{}} \PYG{n}{dv}\PYG{p}{.}\PYG{n}{size}\PYG{p}{(}\PYG{p}{)} \PYG{o}{=}\PYG{o}{=} \PYG{l+m+mi}{2} \PYG{p}{;}
    \PYG{n}{ok} \PYG{o}{=} \PYG{n}{ok} \PYG{o}{\PYGZam{}}\PYG{o}{\PYGZam{}} \PYG{n}{av}\PYG{p}{.}\PYG{n}{size}\PYG{p}{(}\PYG{p}{)} \PYG{o}{=}\PYG{o}{=} \PYG{l+m+mi}{3} \PYG{p}{;}
    \PYG{c+c1}{//}
    \PYG{k}{return}\PYG{p}{(} \PYG{n}{ok} \PYG{p}{)}\PYG{p}{;}
\PYG{p}{\PYGZcb{}}
\end{sphinxVerbatim}


\bigskip\hrule\bigskip


xsrst input file: \sphinxcode{\sphinxupquote{example/cplusplus/vector\_size\_xam.cpp}}


\subparagraph{vector\_size\_xam\_py}
\label{\detokenize{xsrst/vector_size_xam_py:vector-size-xam-py}}\label{\detokenize{xsrst/vector_size_xam_py::doc}}
\index{vector\_size\_xam\_py@\spxentry{vector\_size\_xam\_py}}\index{python:@\spxentry{python:}}\index{size@\spxentry{size}}\index{vectors:@\spxentry{vectors:}}\index{example@\spxentry{example}}\index{test@\spxentry{test}}\ignorespaces 

\subparagraph{3.2.2.2.2 Python: Size of Vectors: Example and Test}
\label{\detokenize{xsrst/vector_size_xam_py:python-size-of-vectors-example-and-test}}\label{\detokenize{xsrst/vector_size_xam_py:id1}}
\begin{sphinxVerbatim}[commandchars=\\\{\}]
\PYG{k}{def} \PYG{n+nf}{vector\PYGZus{}size\PYGZus{}xam}\PYG{p}{(}\PYG{p}{)} \PYG{p}{:}
    \PYG{c+c1}{\PYGZsh{}}
    \PYG{k+kn}{import} \PYG{n+nn}{numpy}
    \PYG{k+kn}{import} \PYG{n+nn}{cppad\PYGZus{}py}
    \PYG{c+c1}{\PYGZsh{}}
    \PYG{c+c1}{\PYGZsh{} initialize return variable}
    \PYG{n}{ok} \PYG{o}{=} \PYG{k+kc}{True}
    \PYG{c+c1}{\PYGZsh{} \PYGZhy{}\PYGZhy{}\PYGZhy{}\PYGZhy{}\PYGZhy{}\PYGZhy{}\PYGZhy{}\PYGZhy{}\PYGZhy{}\PYGZhy{}\PYGZhy{}\PYGZhy{}\PYGZhy{}\PYGZhy{}\PYGZhy{}\PYGZhy{}\PYGZhy{}\PYGZhy{}\PYGZhy{}\PYGZhy{}\PYGZhy{}\PYGZhy{}\PYGZhy{}\PYGZhy{}\PYGZhy{}\PYGZhy{}\PYGZhy{}\PYGZhy{}\PYGZhy{}\PYGZhy{}\PYGZhy{}\PYGZhy{}\PYGZhy{}\PYGZhy{}\PYGZhy{}\PYGZhy{}\PYGZhy{}\PYGZhy{}\PYGZhy{}\PYGZhy{}\PYGZhy{}\PYGZhy{}\PYGZhy{}\PYGZhy{}\PYGZhy{}\PYGZhy{}\PYGZhy{}\PYGZhy{}\PYGZhy{}\PYGZhy{}\PYGZhy{}\PYGZhy{}\PYGZhy{}\PYGZhy{}\PYGZhy{}\PYGZhy{}\PYGZhy{}\PYGZhy{}\PYGZhy{}\PYGZhy{}\PYGZhy{}\PYGZhy{}\PYGZhy{}\PYGZhy{}\PYGZhy{}\PYGZhy{}\PYGZhy{}\PYGZhy{}\PYGZhy{}}
    \PYG{c+c1}{\PYGZsh{} create vectors}
    \PYG{n}{bv} \PYG{o}{=} \PYG{n}{cppad\PYGZus{}py}\PYG{o}{.}\PYG{n}{vec\PYGZus{}bool}\PYG{p}{(}\PYG{p}{)}
    \PYG{n}{iv} \PYG{o}{=} \PYG{n}{cppad\PYGZus{}py}\PYG{o}{.}\PYG{n}{vec\PYGZus{}int}\PYG{p}{(}\PYG{l+m+mi}{1}\PYG{p}{)}
    \PYG{n}{dv} \PYG{o}{=} \PYG{n}{cppad\PYGZus{}py}\PYG{o}{.}\PYG{n}{vec\PYGZus{}double}\PYG{p}{(}\PYG{l+m+mi}{2}\PYG{p}{)}
    \PYG{n}{av} \PYG{o}{=} \PYG{n}{cppad\PYGZus{}py}\PYG{o}{.}\PYG{n}{vec\PYGZus{}a\PYGZus{}double}\PYG{p}{(}\PYG{l+m+mi}{3}\PYG{p}{)}
    \PYG{c+c1}{\PYGZsh{}}
    \PYG{c+c1}{\PYGZsh{} check size of vectors}
    \PYG{n}{ok} \PYG{o}{=} \PYG{n}{ok} \PYG{o+ow}{and} \PYG{n}{bv}\PYG{o}{.}\PYG{n}{size}\PYG{p}{(}\PYG{p}{)} \PYG{o}{==} \PYG{l+m+mi}{0}
    \PYG{n}{ok} \PYG{o}{=} \PYG{n}{ok} \PYG{o+ow}{and} \PYG{n}{iv}\PYG{o}{.}\PYG{n}{size}\PYG{p}{(}\PYG{p}{)} \PYG{o}{==} \PYG{l+m+mi}{1}
    \PYG{n}{ok} \PYG{o}{=} \PYG{n}{ok} \PYG{o+ow}{and} \PYG{n}{dv}\PYG{o}{.}\PYG{n}{size}\PYG{p}{(}\PYG{p}{)} \PYG{o}{==} \PYG{l+m+mi}{2}
    \PYG{n}{ok} \PYG{o}{=} \PYG{n}{ok} \PYG{o+ow}{and} \PYG{n}{av}\PYG{o}{.}\PYG{n}{size}\PYG{p}{(}\PYG{p}{)} \PYG{o}{==} \PYG{l+m+mi}{3}
    \PYG{c+c1}{\PYGZsh{}}
    \PYG{k}{return}\PYG{p}{(} \PYG{n}{ok} \PYG{p}{)}
\PYG{c+c1}{\PYGZsh{}}
\end{sphinxVerbatim}


\bigskip\hrule\bigskip


xsrst input file: \sphinxcode{\sphinxupquote{example/python/core/vector\_size\_xam.py}}

\index{example@\spxentry{example}}\ignorespaces 

\subparagraph{Example}
\label{\detokenize{xsrst/vector_size:example}}\label{\detokenize{xsrst/vector_size:vector-size-example}}\label{\detokenize{xsrst/vector_size:index-4}}
{\hyperref[\detokenize{xsrst/vector_size_xam_cpp:id1}]{\sphinxcrossref{\DUrole{std,std-ref}{c++}}}},
{\hyperref[\detokenize{xsrst/vector_size_xam_py:id1}]{\sphinxcrossref{\DUrole{std,std-ref}{python}}}}.


\bigskip\hrule\bigskip


xsrst input file: \sphinxcode{\sphinxupquote{lib/cplusplus/vector.xsrst}}


\subparagraph{vector\_set\_get}
\label{\detokenize{xsrst/vector_set_get:vector-set-get}}\label{\detokenize{xsrst/vector_set_get::doc}}
\index{vector\_set\_get@\spxentry{vector\_set\_get}}\index{setting@\spxentry{setting}}\index{getting@\spxentry{getting}}\index{vector@\spxentry{vector}}\index{elements@\spxentry{elements}}\ignorespaces 

\subparagraph{3.2.2.3 Setting and Getting Vector Elements}
\label{\detokenize{xsrst/vector_set_get:setting-and-getting-vector-elements}}\label{\detokenize{xsrst/vector_set_get:id1}}\begin{itemize}
\item {} 
{\hyperref[\detokenize{xsrst/vector_set_get:vector-set-get-syntax}]{\sphinxcrossref{\DUrole{std,std-ref}{Syntax}}}}

\item {} 
{\hyperref[\detokenize{xsrst/vector_set_get:vector-set-get-element-type}]{\sphinxcrossref{\DUrole{std,std-ref}{element\_type}}}}

\item {} 
{\hyperref[\detokenize{xsrst/vector_set_get:vector-set-get-i}]{\sphinxcrossref{\DUrole{std,std-ref}{i}}}}

\item {} 
{\hyperref[\detokenize{xsrst/vector_set_get:vector-set-get-u}]{\sphinxcrossref{\DUrole{std,std-ref}{u}}}}

\item {} 
{\hyperref[\detokenize{xsrst/vector_set_get:vector-set-get-x}]{\sphinxcrossref{\DUrole{std,std-ref}{x}}}}

\item {} 
{\hyperref[\detokenize{xsrst/vector_set_get:vector-set-get-v}]{\sphinxcrossref{\DUrole{std,std-ref}{v}}}}

\item {} 
{\hyperref[\detokenize{xsrst/vector_set_get:vector-set-get-y}]{\sphinxcrossref{\DUrole{std,std-ref}{y}}}}

\item {} 
{\hyperref[\detokenize{xsrst/vector_set_get:vector-set-get-example}]{\sphinxcrossref{\DUrole{std,std-ref}{Example}}}}

\end{itemize}

\index{syntax@\spxentry{syntax}}\ignorespaces 

\subparagraph{Syntax}
\label{\detokenize{xsrst/vector_set_get:syntax}}\label{\detokenize{xsrst/vector_set_get:vector-set-get-syntax}}\label{\detokenize{xsrst/vector_set_get:index-1}}
\begin{DUlineblock}{0em}
\item[] \sphinxstyleemphasis{u} {[} \sphinxstyleemphasis{i} {]} = \sphinxstyleemphasis{x}
\item[] \sphinxstyleemphasis{y} = \sphinxstyleemphasis{v} {[} \sphinxstyleemphasis{i} {]}
\end{DUlineblock}

\index{element\_type@\spxentry{element\_type}}\ignorespaces 

\subparagraph{element\_type}
\label{\detokenize{xsrst/vector_set_get:element-type}}\label{\detokenize{xsrst/vector_set_get:vector-set-get-element-type}}\label{\detokenize{xsrst/vector_set_get:index-2}}
We use \sphinxstyleemphasis{element\_type} to denote the type of elements in the
vector. It must be one of the following types:
\sphinxcode{\sphinxupquote{bool}} , \sphinxcode{\sphinxupquote{int}} , \sphinxcode{\sphinxupquote{double}} , \sphinxcode{\sphinxupquote{a\_double}} .

\index{i@\spxentry{i}}\ignorespaces 

\subparagraph{i}
\label{\detokenize{xsrst/vector_set_get:i}}\label{\detokenize{xsrst/vector_set_get:vector-set-get-i}}\label{\detokenize{xsrst/vector_set_get:index-3}}
This argument has prototype

\begin{DUlineblock}{0em}
\item[]         \sphinxcode{\sphinxupquote{size\_t}} \sphinxstyleemphasis{i}
\end{DUlineblock}

It must be between zero and the size of the vector minus one.

\index{u@\spxentry{u}}\ignorespaces 

\subparagraph{u}
\label{\detokenize{xsrst/vector_set_get:u}}\label{\detokenize{xsrst/vector_set_get:vector-set-get-u}}\label{\detokenize{xsrst/vector_set_get:index-4}}
The object \sphinxstyleemphasis{u} has prototype

\begin{DUlineblock}{0em}
\item[]         \sphinxcode{\sphinxupquote{vec\_}} \sphinxstyleemphasis{element\_type} \sphinxcode{\sphinxupquote{\&}} \sphinxstyleemphasis{u}
\end{DUlineblock}

\index{x@\spxentry{x}}\ignorespaces 

\subparagraph{x}
\label{\detokenize{xsrst/vector_set_get:x}}\label{\detokenize{xsrst/vector_set_get:vector-set-get-x}}\label{\detokenize{xsrst/vector_set_get:index-5}}
The argument \sphinxstyleemphasis{x} has prototype

\begin{DUlineblock}{0em}
\item[]         \sphinxcode{\sphinxupquote{const}} \sphinxstyleemphasis{element\_type} \sphinxcode{\sphinxupquote{\&}} \sphinxstyleemphasis{x}
\end{DUlineblock}

\index{v@\spxentry{v}}\ignorespaces 

\subparagraph{v}
\label{\detokenize{xsrst/vector_set_get:v}}\label{\detokenize{xsrst/vector_set_get:vector-set-get-v}}\label{\detokenize{xsrst/vector_set_get:index-6}}
The object \sphinxstyleemphasis{v} has the following prototype

\begin{DUlineblock}{0em}
\item[]         \sphinxcode{\sphinxupquote{const vec\_}} \sphinxstyleemphasis{element\_type} \sphinxcode{\sphinxupquote{\&}} \sphinxstyleemphasis{v}
\end{DUlineblock}

\index{y@\spxentry{y}}\ignorespaces 

\subparagraph{y}
\label{\detokenize{xsrst/vector_set_get:y}}\label{\detokenize{xsrst/vector_set_get:vector-set-get-y}}\label{\detokenize{xsrst/vector_set_get:index-7}}
The result \sphinxstyleemphasis{y} has prototype

\begin{DUlineblock}{0em}
\item[]         \sphinxstyleemphasis{element\_type} \sphinxstyleemphasis{y}
\end{DUlineblock}


\subparagraph{vector\_set\_get\_xam\_cpp}
\label{\detokenize{xsrst/vector_set_get_xam_cpp:vector-set-get-xam-cpp}}\label{\detokenize{xsrst/vector_set_get_xam_cpp::doc}}
\index{vector\_set\_get\_xam\_cpp@\spxentry{vector\_set\_get\_xam\_cpp}}\index{c++:@\spxentry{c++:}}\index{setting@\spxentry{setting}}\index{getting@\spxentry{getting}}\index{vector@\spxentry{vector}}\index{elements:@\spxentry{elements:}}\index{example@\spxentry{example}}\index{test@\spxentry{test}}\ignorespaces 

\subparagraph{3.2.2.3.1 C++: Setting and Getting Vector Elements: Example and Test}
\label{\detokenize{xsrst/vector_set_get_xam_cpp:c-setting-and-getting-vector-elements-example-and-test}}\label{\detokenize{xsrst/vector_set_get_xam_cpp:id1}}
\begin{sphinxVerbatim}[commandchars=\\\{\}]
\PYG{c+cp}{\PYGZsh{}}\PYG{c+cp}{ include \PYGZlt{}cstdio\PYGZgt{}}
\PYG{c+cp}{\PYGZsh{}}\PYG{c+cp}{ include \PYGZlt{}cppad}\PYG{c+cp}{/}\PYG{c+cp}{py}\PYG{c+cp}{/}\PYG{c+cp}{cppad\PYGZus{}py.hpp\PYGZgt{}}

\PYG{k+kt}{bool} \PYG{n+nf}{vector\PYGZus{}set\PYGZus{}get\PYGZus{}xam}\PYG{p}{(}\PYG{k+kt}{void}\PYG{p}{)} \PYG{p}{\PYGZob{}}
    \PYG{k}{using} \PYG{n}{cppad\PYGZus{}py}\PYG{o}{:}\PYG{o}{:}\PYG{n}{a\PYGZus{}double}\PYG{p}{;}
    \PYG{k}{using} \PYG{n}{cppad\PYGZus{}py}\PYG{o}{:}\PYG{o}{:}\PYG{n}{vec\PYGZus{}bool}\PYG{p}{;}
    \PYG{k}{using} \PYG{n}{cppad\PYGZus{}py}\PYG{o}{:}\PYG{o}{:}\PYG{n}{vec\PYGZus{}int}\PYG{p}{;}
    \PYG{k}{using} \PYG{n}{cppad\PYGZus{}py}\PYG{o}{:}\PYG{o}{:}\PYG{n}{vec\PYGZus{}double}\PYG{p}{;}
    \PYG{k}{using} \PYG{n}{cppad\PYGZus{}py}\PYG{o}{:}\PYG{o}{:}\PYG{n}{vec\PYGZus{}a\PYGZus{}double}\PYG{p}{;}
    \PYG{c+c1}{//}
    \PYG{c+c1}{// initialize return variable}
    \PYG{k+kt}{bool} \PYG{n}{ok} \PYG{o}{=} \PYG{n+nb}{true}\PYG{p}{;}
    \PYG{c+c1}{//\PYGZhy{}\PYGZhy{}\PYGZhy{}\PYGZhy{}\PYGZhy{}\PYGZhy{}\PYGZhy{}\PYGZhy{}\PYGZhy{}\PYGZhy{}\PYGZhy{}\PYGZhy{}\PYGZhy{}\PYGZhy{}\PYGZhy{}\PYGZhy{}\PYGZhy{}\PYGZhy{}\PYGZhy{}\PYGZhy{}\PYGZhy{}\PYGZhy{}\PYGZhy{}\PYGZhy{}\PYGZhy{}\PYGZhy{}\PYGZhy{}\PYGZhy{}\PYGZhy{}\PYGZhy{}\PYGZhy{}\PYGZhy{}\PYGZhy{}\PYGZhy{}\PYGZhy{}\PYGZhy{}\PYGZhy{}\PYGZhy{}\PYGZhy{}\PYGZhy{}\PYGZhy{}\PYGZhy{}\PYGZhy{}\PYGZhy{}\PYGZhy{}\PYGZhy{}\PYGZhy{}\PYGZhy{}\PYGZhy{}\PYGZhy{}\PYGZhy{}\PYGZhy{}\PYGZhy{}\PYGZhy{}\PYGZhy{}\PYGZhy{}\PYGZhy{}\PYGZhy{}\PYGZhy{}\PYGZhy{}\PYGZhy{}\PYGZhy{}\PYGZhy{}\PYGZhy{}\PYGZhy{}\PYGZhy{}\PYGZhy{}\PYGZhy{}\PYGZhy{}\PYGZhy{}\PYGZhy{}\PYGZhy{}}
    \PYG{k+kt}{int} \PYG{n}{n} \PYG{o}{=} \PYG{l+m+mi}{4}\PYG{p}{;}
    \PYG{n}{vec\PYGZus{}bool} \PYG{n}{bv} \PYG{o}{=} \PYG{n}{vec\PYGZus{}bool}\PYG{p}{(}\PYG{n}{n}\PYG{p}{)}\PYG{p}{;}
    \PYG{n}{vec\PYGZus{}int} \PYG{n}{iv} \PYG{o}{=} \PYG{n}{vec\PYGZus{}int}\PYG{p}{(}\PYG{n}{n}\PYG{p}{)}\PYG{p}{;}
    \PYG{n}{vec\PYGZus{}double} \PYG{n}{dv}\PYG{p}{(}\PYG{n}{n}\PYG{p}{)}\PYG{p}{;}
    \PYG{n}{vec\PYGZus{}a\PYGZus{}double} \PYG{n}{av}\PYG{p}{(}\PYG{n}{n}\PYG{p}{)}\PYG{p}{;}
    \PYG{c+c1}{//}
    \PYG{c+c1}{// setting elements}
    \PYG{k}{for}\PYG{p}{(}\PYG{k+kt}{int} \PYG{n}{i} \PYG{o}{=} \PYG{l+m+mi}{0}\PYG{p}{;} \PYG{n}{i} \PYG{o}{\PYGZlt{}} \PYG{n}{n} \PYG{p}{;} \PYG{n}{i}\PYG{o}{+}\PYG{o}{+}\PYG{p}{)} \PYG{p}{\PYGZob{}}
        \PYG{n}{bv}\PYG{p}{[}\PYG{n}{i}\PYG{p}{]} \PYG{o}{=} \PYG{n}{i} \PYG{o}{\PYGZgt{}} \PYG{n}{n} \PYG{o}{/} \PYG{l+m+mi}{2}\PYG{p}{;}
        \PYG{n}{iv}\PYG{p}{[}\PYG{n}{i}\PYG{p}{]} \PYG{o}{=} \PYG{l+m+mi}{2} \PYG{o}{*} \PYG{n}{i}\PYG{p}{;}
        \PYG{n}{dv}\PYG{p}{[}\PYG{n}{i}\PYG{p}{]} \PYG{o}{=} \PYG{l+m+mf}{3.0} \PYG{o}{*} \PYG{n}{i}\PYG{p}{;}
        \PYG{n}{av}\PYG{p}{[}\PYG{n}{i}\PYG{p}{]} \PYG{o}{=} \PYG{l+m+mf}{4.0} \PYG{o}{*} \PYG{n}{i}\PYG{p}{;}
    \PYG{p}{\PYGZcb{}}
    \PYG{c+c1}{//}
    \PYG{k}{for}\PYG{p}{(}\PYG{k+kt}{int} \PYG{n}{i} \PYG{o}{=} \PYG{l+m+mi}{0}\PYG{p}{;} \PYG{n}{i} \PYG{o}{\PYGZlt{}} \PYG{n}{n} \PYG{p}{;} \PYG{n}{i}\PYG{o}{+}\PYG{o}{+}\PYG{p}{)} \PYG{p}{\PYGZob{}}
        \PYG{k+kt}{bool} \PYG{n}{be} \PYG{o}{=} \PYG{n}{bv}\PYG{p}{[}\PYG{n}{i}\PYG{p}{]}\PYG{p}{;}
        \PYG{n}{ok} \PYG{o}{=} \PYG{n}{ok} \PYG{o}{\PYGZam{}}\PYG{o}{\PYGZam{}} \PYG{n}{be} \PYG{o}{=}\PYG{o}{=} \PYG{p}{(}\PYG{n}{i} \PYG{o}{\PYGZgt{}} \PYG{n}{n} \PYG{o}{/} \PYG{l+m+mi}{2}\PYG{p}{)} \PYG{p}{;}
        \PYG{c+c1}{//}
        \PYG{k+kt}{int} \PYG{n}{ie} \PYG{o}{=} \PYG{n}{iv}\PYG{p}{[}\PYG{n}{i}\PYG{p}{]}\PYG{p}{;}
        \PYG{n}{ok} \PYG{o}{=} \PYG{n}{ok} \PYG{o}{\PYGZam{}}\PYG{o}{\PYGZam{}} \PYG{n}{ie} \PYG{o}{=}\PYG{o}{=} \PYG{l+m+mi}{2} \PYG{o}{*} \PYG{n}{i} \PYG{p}{;}
        \PYG{c+c1}{//}
        \PYG{k+kt}{double} \PYG{n}{de} \PYG{o}{=} \PYG{n}{dv}\PYG{p}{[}\PYG{n}{i}\PYG{p}{]}\PYG{p}{;}
        \PYG{n}{ok} \PYG{o}{=} \PYG{n}{ok} \PYG{o}{\PYGZam{}}\PYG{o}{\PYGZam{}} \PYG{n}{de} \PYG{o}{=}\PYG{o}{=} \PYG{l+m+mf}{3.0} \PYG{o}{*} \PYG{n}{i} \PYG{p}{;}
        \PYG{c+c1}{//}
        \PYG{n}{a\PYGZus{}double} \PYG{n}{ae} \PYG{o}{=} \PYG{n}{av}\PYG{p}{[}\PYG{n}{i}\PYG{p}{]}\PYG{p}{;}
        \PYG{n}{ok} \PYG{o}{=} \PYG{n}{ok} \PYG{o}{\PYGZam{}}\PYG{o}{\PYGZam{}} \PYG{n}{ae} \PYG{o}{=}\PYG{o}{=} \PYG{l+m+mf}{4.0} \PYG{o}{*} \PYG{n}{i} \PYG{p}{;}
    \PYG{p}{\PYGZcb{}}
    \PYG{c+c1}{//}
    \PYG{k}{return}\PYG{p}{(} \PYG{n}{ok} \PYG{p}{)}\PYG{p}{;}
\PYG{p}{\PYGZcb{}}
\end{sphinxVerbatim}


\bigskip\hrule\bigskip


xsrst input file: \sphinxcode{\sphinxupquote{example/cplusplus/vector\_set\_get\_xam.cpp}}


\subparagraph{vector\_set\_get\_xam\_py}
\label{\detokenize{xsrst/vector_set_get_xam_py:vector-set-get-xam-py}}\label{\detokenize{xsrst/vector_set_get_xam_py::doc}}
\index{vector\_set\_get\_xam\_py@\spxentry{vector\_set\_get\_xam\_py}}\index{python:@\spxentry{python:}}\index{setting@\spxentry{setting}}\index{getting@\spxentry{getting}}\index{vector@\spxentry{vector}}\index{elements:@\spxentry{elements:}}\index{example@\spxentry{example}}\index{test@\spxentry{test}}\ignorespaces 

\subparagraph{3.2.2.3.2 Python: Setting and Getting Vector Elements: Example and Test}
\label{\detokenize{xsrst/vector_set_get_xam_py:python-setting-and-getting-vector-elements-example-and-test}}\label{\detokenize{xsrst/vector_set_get_xam_py:id1}}
\begin{sphinxVerbatim}[commandchars=\\\{\}]
\PYG{k}{def} \PYG{n+nf}{vector\PYGZus{}set\PYGZus{}get\PYGZus{}xam}\PYG{p}{(}\PYG{p}{)} \PYG{p}{:}
    \PYG{c+c1}{\PYGZsh{}}
    \PYG{k+kn}{import} \PYG{n+nn}{numpy}
    \PYG{k+kn}{import} \PYG{n+nn}{cppad\PYGZus{}py}
    \PYG{c+c1}{\PYGZsh{}}
    \PYG{c+c1}{\PYGZsh{} initialize return variable}
    \PYG{n}{ok} \PYG{o}{=} \PYG{k+kc}{True}
    \PYG{c+c1}{\PYGZsh{} \PYGZhy{}\PYGZhy{}\PYGZhy{}\PYGZhy{}\PYGZhy{}\PYGZhy{}\PYGZhy{}\PYGZhy{}\PYGZhy{}\PYGZhy{}\PYGZhy{}\PYGZhy{}\PYGZhy{}\PYGZhy{}\PYGZhy{}\PYGZhy{}\PYGZhy{}\PYGZhy{}\PYGZhy{}\PYGZhy{}\PYGZhy{}\PYGZhy{}\PYGZhy{}\PYGZhy{}\PYGZhy{}\PYGZhy{}\PYGZhy{}\PYGZhy{}\PYGZhy{}\PYGZhy{}\PYGZhy{}\PYGZhy{}\PYGZhy{}\PYGZhy{}\PYGZhy{}\PYGZhy{}\PYGZhy{}\PYGZhy{}\PYGZhy{}\PYGZhy{}\PYGZhy{}\PYGZhy{}\PYGZhy{}\PYGZhy{}\PYGZhy{}\PYGZhy{}\PYGZhy{}\PYGZhy{}\PYGZhy{}\PYGZhy{}\PYGZhy{}\PYGZhy{}\PYGZhy{}\PYGZhy{}\PYGZhy{}\PYGZhy{}\PYGZhy{}\PYGZhy{}\PYGZhy{}\PYGZhy{}\PYGZhy{}\PYGZhy{}\PYGZhy{}\PYGZhy{}\PYGZhy{}\PYGZhy{}\PYGZhy{}\PYGZhy{}\PYGZhy{}}
    \PYG{n}{n} \PYG{o}{=} \PYG{l+m+mi}{4}
    \PYG{n}{bv} \PYG{o}{=} \PYG{n}{cppad\PYGZus{}py}\PYG{o}{.}\PYG{n}{vec\PYGZus{}bool}\PYG{p}{(}\PYG{n}{n}\PYG{p}{)}
    \PYG{n}{iv} \PYG{o}{=} \PYG{n}{cppad\PYGZus{}py}\PYG{o}{.}\PYG{n}{vec\PYGZus{}int}\PYG{p}{(}\PYG{n}{n}\PYG{p}{)}
    \PYG{n}{dv} \PYG{o}{=} \PYG{n}{numpy}\PYG{o}{.}\PYG{n}{empty}\PYG{p}{(}\PYG{n}{n}\PYG{p}{,} \PYG{n}{dtype}\PYG{o}{=}\PYG{n+nb}{float}\PYG{p}{)}
    \PYG{n}{av} \PYG{o}{=} \PYG{n}{numpy}\PYG{o}{.}\PYG{n}{empty}\PYG{p}{(}\PYG{n}{n}\PYG{p}{,} \PYG{n}{dtype}\PYG{o}{=}\PYG{n}{cppad\PYGZus{}py}\PYG{o}{.}\PYG{n}{a\PYGZus{}double}\PYG{p}{)}
    \PYG{c+c1}{\PYGZsh{}}
    \PYG{c+c1}{\PYGZsh{} setting elements}
    \PYG{k}{for} \PYG{n}{i} \PYG{o+ow}{in} \PYG{n+nb}{range}\PYG{p}{(} \PYG{n}{n}  \PYG{p}{)} \PYG{p}{:}
        \PYG{n}{bv}\PYG{p}{[}\PYG{n}{i}\PYG{p}{]} \PYG{o}{=} \PYG{n}{i} \PYG{o}{\PYGZgt{}} \PYG{n}{n} \PYG{o}{/} \PYG{l+m+mi}{2}
        \PYG{n}{iv}\PYG{p}{[}\PYG{n}{i}\PYG{p}{]} \PYG{o}{=} \PYG{l+m+mi}{2} \PYG{o}{*} \PYG{n}{i}
        \PYG{n}{dv}\PYG{p}{[}\PYG{n}{i}\PYG{p}{]} \PYG{o}{=} \PYG{l+m+mf}{3.0} \PYG{o}{*} \PYG{n}{i}
        \PYG{n}{av}\PYG{p}{[}\PYG{n}{i}\PYG{p}{]} \PYG{o}{=} \PYG{n}{cppad\PYGZus{}py}\PYG{o}{.}\PYG{n}{a\PYGZus{}double}\PYG{p}{(}\PYG{l+m+mf}{4.0} \PYG{o}{*} \PYG{n}{i}\PYG{p}{)}
    \PYG{c+c1}{\PYGZsh{}}
    \PYG{c+c1}{\PYGZsh{}}
    \PYG{k}{for} \PYG{n}{i} \PYG{o+ow}{in} \PYG{n+nb}{range}\PYG{p}{(} \PYG{n}{n}  \PYG{p}{)} \PYG{p}{:}
        \PYG{n}{be} \PYG{o}{=} \PYG{n}{bv}\PYG{p}{[}\PYG{n}{i}\PYG{p}{]}
        \PYG{n}{ok} \PYG{o}{=} \PYG{n}{ok} \PYG{o+ow}{and} \PYG{n}{be} \PYG{o}{==} \PYG{p}{(}\PYG{n}{i} \PYG{o}{\PYGZgt{}} \PYG{n}{n} \PYG{o}{/} \PYG{l+m+mi}{2}\PYG{p}{)}
        \PYG{c+c1}{\PYGZsh{}}
        \PYG{n}{ie} \PYG{o}{=} \PYG{n}{iv}\PYG{p}{[}\PYG{n}{i}\PYG{p}{]}
        \PYG{n}{ok} \PYG{o}{=} \PYG{n}{ok} \PYG{o+ow}{and} \PYG{n}{ie} \PYG{o}{==} \PYG{l+m+mi}{2} \PYG{o}{*} \PYG{n}{i}
        \PYG{c+c1}{\PYGZsh{}}
        \PYG{n}{de} \PYG{o}{=} \PYG{n}{dv}\PYG{p}{[}\PYG{n}{i}\PYG{p}{]}
        \PYG{n}{ok} \PYG{o}{=} \PYG{n}{ok} \PYG{o+ow}{and} \PYG{n}{de} \PYG{o}{==} \PYG{l+m+mf}{3.0} \PYG{o}{*} \PYG{n}{i}
        \PYG{c+c1}{\PYGZsh{}}
        \PYG{n}{ae} \PYG{o}{=} \PYG{n}{av}\PYG{p}{[}\PYG{n}{i}\PYG{p}{]}
        \PYG{n}{ok} \PYG{o}{=} \PYG{n}{ok} \PYG{o+ow}{and} \PYG{n}{ae} \PYG{o}{==} \PYG{l+m+mf}{4.0} \PYG{o}{*} \PYG{n}{i}
    \PYG{c+c1}{\PYGZsh{}}
    \PYG{c+c1}{\PYGZsh{}}
    \PYG{k}{return}\PYG{p}{(} \PYG{n}{ok} \PYG{p}{)}
\PYG{c+c1}{\PYGZsh{}}
\end{sphinxVerbatim}


\bigskip\hrule\bigskip


xsrst input file: \sphinxcode{\sphinxupquote{example/python/core/vector\_set\_get\_xam.py}}

\index{example@\spxentry{example}}\ignorespaces 

\subparagraph{Example}
\label{\detokenize{xsrst/vector_set_get:example}}\label{\detokenize{xsrst/vector_set_get:vector-set-get-example}}\label{\detokenize{xsrst/vector_set_get:index-8}}
{\hyperref[\detokenize{xsrst/vector_set_get_xam_cpp:id1}]{\sphinxcrossref{\DUrole{std,std-ref}{c++}}}},
{\hyperref[\detokenize{xsrst/vector_set_get_xam_py:id1}]{\sphinxcrossref{\DUrole{std,std-ref}{python}}}}.


\bigskip\hrule\bigskip


xsrst input file: \sphinxcode{\sphinxupquote{lib/cplusplus/vector.xsrst}}
\begin{enumerate}
\sphinxsetlistlabels{\arabic}{enumi}{enumii}{}{.}%
\item {} 
vector\_ctor: {\hyperref[\detokenize{xsrst/vector_ctor:id1}]{\sphinxcrossref{\DUrole{std,std-ref}{3.2.2.1 Cppad Py Vector Constructors}}}}

\item {} 
vector\_size: {\hyperref[\detokenize{xsrst/vector_size:id1}]{\sphinxcrossref{\DUrole{std,std-ref}{3.2.2.2 Size of a Vector}}}}

\item {} 
vector\_set\_get: {\hyperref[\detokenize{xsrst/vector_set_get:id1}]{\sphinxcrossref{\DUrole{std,std-ref}{3.2.2.3 Setting and Getting Vector Elements}}}}

\end{enumerate}


\bigskip\hrule\bigskip


xsrst input file: \sphinxcode{\sphinxupquote{lib/cplusplus/cplusplus.xsrst}}


\subparagraph{cpp\_fun}
\label{\detokenize{xsrst/cpp_fun:cpp-fun}}\label{\detokenize{xsrst/cpp_fun::doc}}
\index{cpp\_fun@\spxentry{cpp\_fun}}\index{cppad@\spxentry{cppad}}\index{py@\spxentry{py}}\index{ad@\spxentry{ad}}\index{functions@\spxentry{functions}}\ignorespaces 

\subparagraph{3.2.3 Cppad Py AD Functions}
\label{\detokenize{xsrst/cpp_fun:cppad-py-ad-functions}}\label{\detokenize{xsrst/cpp_fun:id1}}\begin{itemize}
\item {} 
{\hyperref[\detokenize{xsrst/cpp_fun:cpp-fun-children}]{\sphinxcrossref{\DUrole{std,std-ref}{Children}}}}

\end{itemize}

\index{children@\spxentry{children}}\ignorespaces 

\subparagraph{Children}
\label{\detokenize{xsrst/cpp_fun:children}}\label{\detokenize{xsrst/cpp_fun:cpp-fun-children}}\label{\detokenize{xsrst/cpp_fun:index-1}}

\subparagraph{cpp\_independent}
\label{\detokenize{xsrst/cpp_independent:cpp-independent}}\label{\detokenize{xsrst/cpp_independent::doc}}
\index{cpp\_independent@\spxentry{cpp\_independent}}\index{declare@\spxentry{declare}}\index{independent@\spxentry{independent}}\index{variables@\spxentry{variables}}\index{start@\spxentry{start}}\index{recording@\spxentry{recording}}\ignorespaces 

\subparagraph{3.2.3.1 Declare Independent Variables and Start Recording}
\label{\detokenize{xsrst/cpp_independent:declare-independent-variables-and-start-recording}}\label{\detokenize{xsrst/cpp_independent:id1}}\begin{itemize}
\item {} 
{\hyperref[\detokenize{xsrst/cpp_independent:cpp-independent-syntax}]{\sphinxcrossref{\DUrole{std,std-ref}{Syntax}}}}

\item {} 
{\hyperref[\detokenize{xsrst/cpp_independent:cpp-independent-purpose}]{\sphinxcrossref{\DUrole{std,std-ref}{Purpose}}}}

\item {} 
{\hyperref[\detokenize{xsrst/cpp_independent:cpp-independent-x}]{\sphinxcrossref{\DUrole{std,std-ref}{x}}}}

\item {} 
{\hyperref[\detokenize{xsrst/cpp_independent:cpp-independent-dynamic}]{\sphinxcrossref{\DUrole{std,std-ref}{dynamic}}}}

\item {} 
{\hyperref[\detokenize{xsrst/cpp_independent:cpp-independent-ax}]{\sphinxcrossref{\DUrole{std,std-ref}{ax}}}}

\item {} 
{\hyperref[\detokenize{xsrst/cpp_independent:cpp-independent-a-both}]{\sphinxcrossref{\DUrole{std,std-ref}{a\_both}}}}

\item {} 
{\hyperref[\detokenize{xsrst/cpp_independent:cpp-independent-example}]{\sphinxcrossref{\DUrole{std,std-ref}{Example}}}}

\end{itemize}

\index{syntax@\spxentry{syntax}}\ignorespaces 

\subparagraph{Syntax}
\label{\detokenize{xsrst/cpp_independent:syntax}}\label{\detokenize{xsrst/cpp_independent:cpp-independent-syntax}}\label{\detokenize{xsrst/cpp_independent:index-1}}
\begin{DUlineblock}{0em}
\item[] \sphinxstyleemphasis{ax} =  \sphinxcode{\sphinxupquote{cppad\_py::independent}} ( \sphinxstyleemphasis{x} )
\item[] \sphinxstyleemphasis{a\_both} =  \sphinxcode{\sphinxupquote{cppad\_py::independent}} ( \sphinxstyleemphasis{x} , \sphinxstyleemphasis{dynamic} )
\end{DUlineblock}

\index{purpose@\spxentry{purpose}}\ignorespaces 

\subparagraph{Purpose}
\label{\detokenize{xsrst/cpp_independent:purpose}}\label{\detokenize{xsrst/cpp_independent:cpp-independent-purpose}}\label{\detokenize{xsrst/cpp_independent:index-2}}
This starts recording {\hyperref[\detokenize{xsrst/a_double:id1}]{\sphinxcrossref{\DUrole{std,std-ref}{a\_double}}}} operations.
This recording is terminated, and the information is stored,
by calling the {\hyperref[\detokenize{xsrst/cpp_fun_ctor:id1}]{\sphinxcrossref{\DUrole{std,std-ref}{d\_fun\_constructor}}}}.
It can be terminated, and the information is lost,
by calling {\hyperref[\detokenize{xsrst/cpp_abort_recording:id1}]{\sphinxcrossref{\DUrole{std,std-ref}{abort\_recording}}}}.

\index{x@\spxentry{x}}\ignorespaces 

\subparagraph{x}
\label{\detokenize{xsrst/cpp_independent:x}}\label{\detokenize{xsrst/cpp_independent:cpp-independent-x}}\label{\detokenize{xsrst/cpp_independent:index-3}}
This argument has prototype

\begin{DUlineblock}{0em}
\item[]         \sphinxcode{\sphinxupquote{const vec\_double\&}} \sphinxstyleemphasis{x}
\end{DUlineblock}

Its specifies the number of independent variables
and their values during the recording.
We use the notation \sphinxstyleemphasis{nx} = \sphinxstyleemphasis{x} . \sphinxcode{\sphinxupquote{size}} ()
to denote the number of independent variables.

\index{dynamic@\spxentry{dynamic}}\ignorespaces 

\subparagraph{dynamic}
\label{\detokenize{xsrst/cpp_independent:dynamic}}\label{\detokenize{xsrst/cpp_independent:cpp-independent-dynamic}}\label{\detokenize{xsrst/cpp_independent:index-4}}
This argument has prototype

\begin{DUlineblock}{0em}
\item[]         \sphinxcode{\sphinxupquote{const vec\_double\&}} \sphinxstyleemphasis{dynamic}
\end{DUlineblock}

Its specifies the number of independent dynamic parameters
and their values during the recording.
We use the notation \sphinxstyleemphasis{nd} = \sphinxstyleemphasis{dynamic} . \sphinxcode{\sphinxupquote{size}} ()
to denote the number of independent variables.

\index{ax@\spxentry{ax}}\ignorespaces 

\subparagraph{ax}
\label{\detokenize{xsrst/cpp_independent:ax}}\label{\detokenize{xsrst/cpp_independent:cpp-independent-ax}}\label{\detokenize{xsrst/cpp_independent:index-5}}
This result has prototype

\begin{DUlineblock}{0em}
\item[]         \sphinxcode{\sphinxupquote{vec\_a\_double\&}} \sphinxstyleemphasis{ax}
\end{DUlineblock}

and is the vector of independent variables.
It has size \sphinxstyleemphasis{nx} and for
\sphinxstyleemphasis{i} = 0 to \sphinxstyleemphasis{n} \sphinxhyphen{}1

\begin{DUlineblock}{0em}
\item[]         \sphinxstyleemphasis{ax} {[} \sphinxstyleemphasis{i} {]} == \sphinxstyleemphasis{x} {[} \sphinxstyleemphasis{i} {]}
\end{DUlineblock}

\index{a\_both@\spxentry{a\_both}}\ignorespaces 

\subparagraph{a\_both}
\label{\detokenize{xsrst/cpp_independent:a-both}}\label{\detokenize{xsrst/cpp_independent:cpp-independent-a-both}}\label{\detokenize{xsrst/cpp_independent:index-6}}
this result has prototype

\begin{DUlineblock}{0em}
\item[]         \sphinxcode{\sphinxupquote{vec\_a\_double\&}} \sphinxstyleemphasis{a\_both}
\end{DUlineblock}

and is the vector of both the independent variables
and independent dynamic parameters.
It has size \sphinxstyleemphasis{nx} + \sphinxstyleemphasis{nd} .
For \sphinxstyleemphasis{i} = 0 to \sphinxstyleemphasis{nx} \sphinxhyphen{}1

\begin{DUlineblock}{0em}
\item[]         \sphinxstyleemphasis{a\_both} {[} \sphinxstyleemphasis{i} {]} == \sphinxstyleemphasis{x} {[} \sphinxstyleemphasis{i} {]}
\end{DUlineblock}

is the \sphinxstyleemphasis{i}\sphinxhyphen{}th independent variable.
For \sphinxstyleemphasis{i} = 0 to \sphinxstyleemphasis{nd} \sphinxhyphen{}1

\begin{DUlineblock}{0em}
\item[]         \sphinxstyleemphasis{a\_both} {[} \sphinxstyleemphasis{nx} + \sphinxstyleemphasis{i} {]} == \sphinxstyleemphasis{dynamic} {[} \sphinxstyleemphasis{i} {]}
\end{DUlineblock}

is the \sphinxstyleemphasis{i}\sphinxhyphen{}th independent dynamic parameter.


\subparagraph{fun\_dynamic\_xam\_cpp}
\label{\detokenize{xsrst/fun_dynamic_xam_cpp:fun-dynamic-xam-cpp}}\label{\detokenize{xsrst/fun_dynamic_xam_cpp::doc}}
\index{fun\_dynamic\_xam\_cpp@\spxentry{fun\_dynamic\_xam\_cpp}}\index{c++:@\spxentry{c++:}}\index{using@\spxentry{using}}\index{dynamic@\spxentry{dynamic}}\index{parameters:@\spxentry{parameters:}}\index{example@\spxentry{example}}\index{test@\spxentry{test}}\ignorespaces 

\subparagraph{3.2.3.1.1 C++: Using Dynamic Parameters: Example and Test}
\label{\detokenize{xsrst/fun_dynamic_xam_cpp:c-using-dynamic-parameters-example-and-test}}\label{\detokenize{xsrst/fun_dynamic_xam_cpp:id1}}
\begin{sphinxVerbatim}[commandchars=\\\{\}]
\PYG{c+cp}{\PYGZsh{}}\PYG{c+cp}{ include \PYGZlt{}cstdio\PYGZgt{}}
\PYG{c+cp}{\PYGZsh{}}\PYG{c+cp}{ include \PYGZlt{}cppad}\PYG{c+cp}{/}\PYG{c+cp}{py}\PYG{c+cp}{/}\PYG{c+cp}{cppad\PYGZus{}py.hpp\PYGZgt{}}

\PYG{k+kt}{bool} \PYG{n+nf}{fun\PYGZus{}dynamic\PYGZus{}xam}\PYG{p}{(}\PYG{k+kt}{void}\PYG{p}{)} \PYG{p}{\PYGZob{}}
    \PYG{k}{using} \PYG{n}{cppad\PYGZus{}py}\PYG{o}{:}\PYG{o}{:}\PYG{n}{a\PYGZus{}double}\PYG{p}{;}
    \PYG{k}{using} \PYG{n}{cppad\PYGZus{}py}\PYG{o}{:}\PYG{o}{:}\PYG{n}{vec\PYGZus{}double}\PYG{p}{;}
    \PYG{k}{using} \PYG{n}{cppad\PYGZus{}py}\PYG{o}{:}\PYG{o}{:}\PYG{n}{vec\PYGZus{}a\PYGZus{}double}\PYG{p}{;}
    \PYG{k}{using} \PYG{n}{cppad\PYGZus{}py}\PYG{o}{:}\PYG{o}{:}\PYG{n}{d\PYGZus{}fun}\PYG{p}{;}
    \PYG{c+c1}{//}
    \PYG{c+c1}{// initialize return variable}
    \PYG{k+kt}{bool} \PYG{n}{ok} \PYG{o}{=} \PYG{n+nb}{true}\PYG{p}{;}
    \PYG{c+c1}{//\PYGZhy{}\PYGZhy{}\PYGZhy{}\PYGZhy{}\PYGZhy{}\PYGZhy{}\PYGZhy{}\PYGZhy{}\PYGZhy{}\PYGZhy{}\PYGZhy{}\PYGZhy{}\PYGZhy{}\PYGZhy{}\PYGZhy{}\PYGZhy{}\PYGZhy{}\PYGZhy{}\PYGZhy{}\PYGZhy{}\PYGZhy{}\PYGZhy{}\PYGZhy{}\PYGZhy{}\PYGZhy{}\PYGZhy{}\PYGZhy{}\PYGZhy{}\PYGZhy{}\PYGZhy{}\PYGZhy{}\PYGZhy{}\PYGZhy{}\PYGZhy{}\PYGZhy{}\PYGZhy{}\PYGZhy{}\PYGZhy{}\PYGZhy{}\PYGZhy{}\PYGZhy{}\PYGZhy{}\PYGZhy{}\PYGZhy{}\PYGZhy{}\PYGZhy{}\PYGZhy{}\PYGZhy{}\PYGZhy{}\PYGZhy{}\PYGZhy{}\PYGZhy{}\PYGZhy{}\PYGZhy{}\PYGZhy{}\PYGZhy{}\PYGZhy{}\PYGZhy{}\PYGZhy{}\PYGZhy{}\PYGZhy{}\PYGZhy{}\PYGZhy{}\PYGZhy{}\PYGZhy{}\PYGZhy{}\PYGZhy{}\PYGZhy{}\PYGZhy{}\PYGZhy{}\PYGZhy{}\PYGZhy{}}
    \PYG{k+kt}{int} \PYG{n}{nx} \PYG{o}{=} \PYG{l+m+mi}{2}\PYG{p}{;}
    \PYG{k+kt}{int} \PYG{n}{nd} \PYG{o}{=} \PYG{l+m+mi}{2}\PYG{p}{;}
    \PYG{c+c1}{//}
    \PYG{c+c1}{// value of independent variables during recording}
    \PYG{n}{vec\PYGZus{}double} \PYG{n}{x}\PYG{p}{(}\PYG{n}{nx}\PYG{p}{)}\PYG{p}{;}
    \PYG{k}{for}\PYG{p}{(}\PYG{k+kt}{int} \PYG{n}{i} \PYG{o}{=} \PYG{l+m+mi}{0}\PYG{p}{;} \PYG{n}{i} \PYG{o}{\PYGZlt{}} \PYG{n}{nx} \PYG{p}{;} \PYG{n}{i}\PYG{o}{+}\PYG{o}{+}\PYG{p}{)} \PYG{p}{\PYGZob{}}
        \PYG{n}{x}\PYG{p}{[}\PYG{n}{i}\PYG{p}{]} \PYG{o}{=} \PYG{n}{i} \PYG{o}{+} \PYG{l+m+mf}{1.0}\PYG{p}{;}
    \PYG{p}{\PYGZcb{}}
    \PYG{c+c1}{//}
    \PYG{c+c1}{// value of independent dynamic parameters during recording}
    \PYG{n}{vec\PYGZus{}double} \PYG{n}{dynamic}\PYG{p}{(}\PYG{n}{nd}\PYG{p}{)}\PYG{p}{;}
    \PYG{k}{for}\PYG{p}{(}\PYG{k+kt}{int} \PYG{n}{i} \PYG{o}{=} \PYG{l+m+mi}{0}\PYG{p}{;} \PYG{n}{i} \PYG{o}{\PYGZlt{}} \PYG{n}{nd} \PYG{p}{;} \PYG{n}{i}\PYG{o}{+}\PYG{o}{+}\PYG{p}{)} \PYG{p}{\PYGZob{}}
        \PYG{n}{dynamic}\PYG{p}{[}\PYG{n}{i}\PYG{p}{]} \PYG{o}{=} \PYG{n}{i} \PYG{o}{+} \PYG{n}{x}\PYG{p}{[}\PYG{n}{nx}\PYG{o}{\PYGZhy{}}\PYG{l+m+mi}{1}\PYG{p}{]} \PYG{o}{+} \PYG{l+m+mf}{1.0}\PYG{p}{;}
    \PYG{p}{\PYGZcb{}}
    \PYG{n}{vec\PYGZus{}a\PYGZus{}double} \PYG{n}{a\PYGZus{}both} \PYG{o}{=} \PYG{n}{cppad\PYGZus{}py}\PYG{o}{:}\PYG{o}{:}\PYG{n}{independent}\PYG{p}{(}\PYG{n}{x}\PYG{p}{,} \PYG{n}{dynamic}\PYG{p}{)}\PYG{p}{;}
    \PYG{c+c1}{//}
    \PYG{c+c1}{// extract vector of independent variables and dynamic parameters}
    \PYG{n}{vec\PYGZus{}a\PYGZus{}double} \PYG{n}{ax}\PYG{p}{(}\PYG{n}{nx}\PYG{p}{)}\PYG{p}{,} \PYG{n}{adynamic}\PYG{p}{(}\PYG{n}{nd}\PYG{p}{)}\PYG{p}{;}
    \PYG{k}{for}\PYG{p}{(}\PYG{k+kt}{int} \PYG{n}{i} \PYG{o}{=} \PYG{l+m+mi}{0}\PYG{p}{;} \PYG{n}{i} \PYG{o}{\PYGZlt{}} \PYG{n}{nx} \PYG{p}{;} \PYG{n}{i}\PYG{o}{+}\PYG{o}{+}\PYG{p}{)} \PYG{p}{\PYGZob{}}
        \PYG{n}{ax}\PYG{p}{[}\PYG{n}{i}\PYG{p}{]} \PYG{o}{=} \PYG{n}{a\PYGZus{}both}\PYG{p}{[}\PYG{n}{i}\PYG{p}{]}\PYG{p}{;}
    \PYG{p}{\PYGZcb{}}
    \PYG{k}{for}\PYG{p}{(}\PYG{k+kt}{int} \PYG{n}{i} \PYG{o}{=} \PYG{l+m+mi}{0}\PYG{p}{;} \PYG{n}{i} \PYG{o}{\PYGZlt{}} \PYG{n}{nd} \PYG{p}{;} \PYG{n}{i}\PYG{o}{+}\PYG{o}{+}\PYG{p}{)} \PYG{p}{\PYGZob{}}
        \PYG{n}{adynamic}\PYG{p}{[}\PYG{n}{i}\PYG{p}{]} \PYG{o}{=} \PYG{n}{a\PYGZus{}both}\PYG{p}{[}\PYG{n}{nx} \PYG{o}{+} \PYG{n}{i}\PYG{p}{]}\PYG{p}{;}
    \PYG{p}{\PYGZcb{}}
    \PYG{c+c1}{//}
    \PYG{c+c1}{// create another dynamic paramerer}
    \PYG{n}{a\PYGZus{}double} \PYG{n}{adyn}  \PYG{o}{=} \PYG{n}{adynamic}\PYG{p}{[}\PYG{l+m+mi}{0}\PYG{p}{]} \PYG{o}{+} \PYG{n}{adynamic}\PYG{p}{[}\PYG{l+m+mi}{1}\PYG{p}{]}\PYG{p}{;}
    \PYG{c+c1}{//}
    \PYG{c+c1}{// create another variable}
    \PYG{n}{a\PYGZus{}double} \PYG{n}{avar}  \PYG{o}{=} \PYG{n}{ax}\PYG{p}{[}\PYG{l+m+mi}{0}\PYG{p}{]} \PYG{o}{+} \PYG{n}{ax}\PYG{p}{[}\PYG{l+m+mi}{1}\PYG{p}{]} \PYG{o}{+} \PYG{n}{adyn}\PYG{p}{;}
    \PYG{c+c1}{//}
    \PYG{c+c1}{// create f(x) = x[0] + x[1] + dynamic[0] + dynamic[1]}
    \PYG{n}{vec\PYGZus{}a\PYGZus{}double} \PYG{n}{ay}\PYG{p}{(}\PYG{l+m+mi}{1}\PYG{p}{)}\PYG{p}{;}
    \PYG{n}{ay}\PYG{p}{[}\PYG{l+m+mi}{0}\PYG{p}{]} \PYG{o}{=} \PYG{n}{avar}\PYG{p}{;}
    \PYG{n}{d\PYGZus{}fun} \PYG{n}{f}\PYG{p}{(}\PYG{n}{ax}\PYG{p}{,} \PYG{n}{ay}\PYG{p}{)}\PYG{p}{;}
    \PYG{c+c1}{//}
    \PYG{c+c1}{// check some properties of f}
    \PYG{n}{ok} \PYG{o}{\PYGZam{}}\PYG{o}{=} \PYG{n}{f}\PYG{p}{.}\PYG{n}{size\PYGZus{}domain}\PYG{p}{(}\PYG{p}{)} \PYG{o}{=}\PYG{o}{=} \PYG{n}{nx}\PYG{p}{;}
    \PYG{n}{ok} \PYG{o}{\PYGZam{}}\PYG{o}{=} \PYG{n}{f}\PYG{p}{.}\PYG{n}{size\PYGZus{}order}\PYG{p}{(}\PYG{p}{)}  \PYG{o}{=}\PYG{o}{=} \PYG{l+m+mi}{0}\PYG{p}{;}
    \PYG{c+c1}{//}
    \PYG{c+c1}{// zero order forward mode using same values as during the recording}
    \PYG{n}{vec\PYGZus{}double} \PYG{n}{y}\PYG{p}{(}\PYG{l+m+mi}{1}\PYG{p}{)}\PYG{p}{;}
    \PYG{n}{y}   \PYG{o}{=} \PYG{n}{f}\PYG{p}{.}\PYG{n}{forward}\PYG{p}{(}\PYG{l+m+mi}{0}\PYG{p}{,} \PYG{n}{x}\PYG{p}{)}\PYG{p}{;}
    \PYG{n}{ok} \PYG{o}{\PYGZam{}}\PYG{o}{=} \PYG{n}{y}\PYG{p}{[}\PYG{l+m+mi}{0}\PYG{p}{]} \PYG{o}{=}\PYG{o}{=} \PYG{p}{(}\PYG{n}{x}\PYG{p}{[}\PYG{l+m+mi}{0}\PYG{p}{]} \PYG{o}{+} \PYG{n}{x}\PYG{p}{[}\PYG{l+m+mi}{1}\PYG{p}{]} \PYG{o}{+} \PYG{n}{dynamic}\PYG{p}{[}\PYG{l+m+mi}{0}\PYG{p}{]} \PYG{o}{+} \PYG{n}{dynamic}\PYG{p}{[}\PYG{l+m+mi}{1}\PYG{p}{]}\PYG{p}{)}\PYG{p}{;}
    \PYG{c+c1}{//}
    \PYG{c+c1}{// zero order forward mode using different value for dynamic parameters}
    \PYG{n}{dynamic}\PYG{p}{[}\PYG{l+m+mi}{0}\PYG{p}{]} \PYG{o}{=} \PYG{n}{dynamic}\PYG{p}{[}\PYG{l+m+mi}{0}\PYG{p}{]} \PYG{o}{+} \PYG{l+m+mf}{1.0}\PYG{p}{;}
    \PYG{n}{dynamic}\PYG{p}{[}\PYG{l+m+mi}{1}\PYG{p}{]} \PYG{o}{=} \PYG{n}{dynamic}\PYG{p}{[}\PYG{l+m+mi}{1}\PYG{p}{]} \PYG{o}{+} \PYG{l+m+mf}{1.0}\PYG{p}{;}
    \PYG{n}{f}\PYG{p}{.}\PYG{n}{new\PYGZus{}dynamic}\PYG{p}{(}\PYG{n}{dynamic}\PYG{p}{)}\PYG{p}{;}
    \PYG{n}{y}   \PYG{o}{=} \PYG{n}{f}\PYG{p}{.}\PYG{n}{forward}\PYG{p}{(}\PYG{l+m+mi}{0}\PYG{p}{,} \PYG{n}{x}\PYG{p}{)}\PYG{p}{;}
    \PYG{n}{ok} \PYG{o}{\PYGZam{}}\PYG{o}{=} \PYG{n}{y}\PYG{p}{[}\PYG{l+m+mi}{0}\PYG{p}{]} \PYG{o}{=}\PYG{o}{=} \PYG{p}{(}\PYG{n}{x}\PYG{p}{[}\PYG{l+m+mi}{0}\PYG{p}{]} \PYG{o}{+} \PYG{n}{x}\PYG{p}{[}\PYG{l+m+mi}{1}\PYG{p}{]} \PYG{o}{+} \PYG{n}{dynamic}\PYG{p}{[}\PYG{l+m+mi}{0}\PYG{p}{]} \PYG{o}{+} \PYG{n}{dynamic}\PYG{p}{[}\PYG{l+m+mi}{1}\PYG{p}{]}\PYG{p}{)}\PYG{p}{;}
    \PYG{c+c1}{//}
    \PYG{k}{return}\PYG{p}{(} \PYG{n}{ok} \PYG{p}{)}\PYG{p}{;}
\PYG{p}{\PYGZcb{}}
\end{sphinxVerbatim}


\bigskip\hrule\bigskip


xsrst input file: \sphinxcode{\sphinxupquote{example/cplusplus/fun\_dynamic\_xam.cpp}}

\index{example@\spxentry{example}}\ignorespaces 

\subparagraph{Example}
\label{\detokenize{xsrst/cpp_independent:example}}\label{\detokenize{xsrst/cpp_independent:cpp-independent-example}}\label{\detokenize{xsrst/cpp_independent:index-7}}
Most of the c++ \sphinxcode{\sphinxupquote{d\_fun}} examples use the \sphinxstyleemphasis{ax}
return syntax.
The {\hyperref[\detokenize{xsrst/fun_dynamic_xam_cpp:id1}]{\sphinxcrossref{\DUrole{std,std-ref}{fun\_dynamic\_xam\_cpp}}}} example uses the \sphinxstyleemphasis{a\_both}
return syntax.


\bigskip\hrule\bigskip


xsrst input file: \sphinxcode{\sphinxupquote{lib/cplusplus/fun.cpp}}


\subparagraph{cpp\_abort\_recording}
\label{\detokenize{xsrst/cpp_abort_recording:cpp-abort-recording}}\label{\detokenize{xsrst/cpp_abort_recording::doc}}
\index{cpp\_abort\_recording@\spxentry{cpp\_abort\_recording}}\index{abort@\spxentry{abort}}\index{recording@\spxentry{recording}}\ignorespaces 

\subparagraph{3.2.3.2 Abort Recording}
\label{\detokenize{xsrst/cpp_abort_recording:abort-recording}}\label{\detokenize{xsrst/cpp_abort_recording:id1}}\begin{itemize}
\item {} 
{\hyperref[\detokenize{xsrst/cpp_abort_recording:cpp-abort-recording-syntax}]{\sphinxcrossref{\DUrole{std,std-ref}{Syntax}}}}

\item {} 
{\hyperref[\detokenize{xsrst/cpp_abort_recording:cpp-abort-recording-purpose}]{\sphinxcrossref{\DUrole{std,std-ref}{Purpose}}}}

\item {} 
{\hyperref[\detokenize{xsrst/cpp_abort_recording:cpp-abort-recording-example}]{\sphinxcrossref{\DUrole{std,std-ref}{Example}}}}

\end{itemize}

\index{syntax@\spxentry{syntax}}\ignorespaces 

\subparagraph{Syntax}
\label{\detokenize{xsrst/cpp_abort_recording:syntax}}\label{\detokenize{xsrst/cpp_abort_recording:cpp-abort-recording-syntax}}\label{\detokenize{xsrst/cpp_abort_recording:index-1}}
\sphinxcode{\sphinxupquote{cppad\_py::abort\_recording}} ()

\index{purpose@\spxentry{purpose}}\ignorespaces 

\subparagraph{Purpose}
\label{\detokenize{xsrst/cpp_abort_recording:purpose}}\label{\detokenize{xsrst/cpp_abort_recording:cpp-abort-recording-purpose}}\label{\detokenize{xsrst/cpp_abort_recording:index-2}}
This aborts the current recording (if it exists)
started by the most recent call to {\hyperref[\detokenize{xsrst/cpp_independent:id1}]{\sphinxcrossref{\DUrole{std,std-ref}{independent}}}}.


\subparagraph{fun\_abort\_xam\_cpp}
\label{\detokenize{xsrst/fun_abort_xam_cpp:fun-abort-xam-cpp}}\label{\detokenize{xsrst/fun_abort_xam_cpp::doc}}
\index{fun\_abort\_xam\_cpp@\spxentry{fun\_abort\_xam\_cpp}}\index{c++:@\spxentry{c++:}}\index{abort@\spxentry{abort}}\index{recording@\spxentry{recording}}\index{a\_double@\spxentry{a\_double}}\index{operations:@\spxentry{operations:}}\index{example@\spxentry{example}}\index{test@\spxentry{test}}\ignorespaces 

\subparagraph{3.2.3.2.1 C++: Abort Recording a\_double Operations: Example and Test}
\label{\detokenize{xsrst/fun_abort_xam_cpp:c-abort-recording-a-double-operations-example-and-test}}\label{\detokenize{xsrst/fun_abort_xam_cpp:id1}}
\begin{sphinxVerbatim}[commandchars=\\\{\}]
\PYG{c+cp}{\PYGZsh{}}\PYG{c+cp}{ include \PYGZlt{}cstdio\PYGZgt{}}
\PYG{c+cp}{\PYGZsh{}}\PYG{c+cp}{ include \PYGZlt{}cppad}\PYG{c+cp}{/}\PYG{c+cp}{py}\PYG{c+cp}{/}\PYG{c+cp}{cppad\PYGZus{}py.hpp\PYGZgt{}}

\PYG{k+kt}{bool} \PYG{n+nf}{fun\PYGZus{}abort\PYGZus{}xam}\PYG{p}{(}\PYG{k+kt}{void}\PYG{p}{)} \PYG{p}{\PYGZob{}}
    \PYG{k}{using} \PYG{n}{cppad\PYGZus{}py}\PYG{o}{:}\PYG{o}{:}\PYG{n}{a\PYGZus{}double}\PYG{p}{;}
    \PYG{k}{using} \PYG{n}{cppad\PYGZus{}py}\PYG{o}{:}\PYG{o}{:}\PYG{n}{vec\PYGZus{}double}\PYG{p}{;}
    \PYG{k}{using} \PYG{n}{cppad\PYGZus{}py}\PYG{o}{:}\PYG{o}{:}\PYG{n}{vec\PYGZus{}a\PYGZus{}double}\PYG{p}{;}
    \PYG{k}{using} \PYG{n}{cppad\PYGZus{}py}\PYG{o}{:}\PYG{o}{:}\PYG{n}{d\PYGZus{}fun}\PYG{p}{;}
    \PYG{c+c1}{//}
    \PYG{c+c1}{// initialize return variable}
    \PYG{k+kt}{bool} \PYG{n}{ok} \PYG{o}{=} \PYG{n+nb}{true}\PYG{p}{;}
    \PYG{c+c1}{//\PYGZhy{}\PYGZhy{}\PYGZhy{}\PYGZhy{}\PYGZhy{}\PYGZhy{}\PYGZhy{}\PYGZhy{}\PYGZhy{}\PYGZhy{}\PYGZhy{}\PYGZhy{}\PYGZhy{}\PYGZhy{}\PYGZhy{}\PYGZhy{}\PYGZhy{}\PYGZhy{}\PYGZhy{}\PYGZhy{}\PYGZhy{}\PYGZhy{}\PYGZhy{}\PYGZhy{}\PYGZhy{}\PYGZhy{}\PYGZhy{}\PYGZhy{}\PYGZhy{}\PYGZhy{}\PYGZhy{}\PYGZhy{}\PYGZhy{}\PYGZhy{}\PYGZhy{}\PYGZhy{}\PYGZhy{}\PYGZhy{}\PYGZhy{}\PYGZhy{}\PYGZhy{}\PYGZhy{}\PYGZhy{}\PYGZhy{}\PYGZhy{}\PYGZhy{}\PYGZhy{}\PYGZhy{}\PYGZhy{}\PYGZhy{}\PYGZhy{}\PYGZhy{}\PYGZhy{}\PYGZhy{}\PYGZhy{}\PYGZhy{}\PYGZhy{}\PYGZhy{}\PYGZhy{}\PYGZhy{}\PYGZhy{}\PYGZhy{}\PYGZhy{}\PYGZhy{}\PYGZhy{}\PYGZhy{}\PYGZhy{}\PYGZhy{}\PYGZhy{}\PYGZhy{}\PYGZhy{}\PYGZhy{}}
    \PYG{k+kt}{int} \PYG{n}{n\PYGZus{}ind} \PYG{o}{=} \PYG{l+m+mi}{2}\PYG{p}{;}
    \PYG{c+c1}{//}
    \PYG{c+c1}{// create ax}
    \PYG{n}{vec\PYGZus{}double} \PYG{n}{x}\PYG{p}{(}\PYG{n}{n\PYGZus{}ind}\PYG{p}{)}\PYG{p}{;}
    \PYG{k}{for}\PYG{p}{(}\PYG{k+kt}{int} \PYG{n}{i} \PYG{o}{=} \PYG{l+m+mi}{0}\PYG{p}{;} \PYG{n}{i} \PYG{o}{\PYGZlt{}} \PYG{n}{n\PYGZus{}ind} \PYG{p}{;} \PYG{n}{i}\PYG{o}{+}\PYG{o}{+}\PYG{p}{)} \PYG{p}{\PYGZob{}}
        \PYG{n}{x}\PYG{p}{[}\PYG{n}{i}\PYG{p}{]} \PYG{o}{=} \PYG{n}{i} \PYG{o}{+} \PYG{l+m+mf}{1.0}\PYG{p}{;}
    \PYG{p}{\PYGZcb{}}
    \PYG{n}{vec\PYGZus{}a\PYGZus{}double} \PYG{n}{ax} \PYG{o}{=} \PYG{n}{cppad\PYGZus{}py}\PYG{o}{:}\PYG{o}{:}\PYG{n}{independent}\PYG{p}{(}\PYG{n}{x}\PYG{p}{)}\PYG{p}{;}
    \PYG{c+c1}{//}
    \PYG{c+c1}{// preform some a\PYGZus{}double operations}
    \PYG{n}{a\PYGZus{}double} \PYG{n}{ax0} \PYG{o}{=} \PYG{n}{ax}\PYG{p}{[}\PYG{l+m+mi}{0}\PYG{p}{]}\PYG{p}{;}
    \PYG{n}{a\PYGZus{}double} \PYG{n}{ax1} \PYG{o}{=} \PYG{n}{ax}\PYG{p}{[}\PYG{l+m+mi}{1}\PYG{p}{]}\PYG{p}{;}
    \PYG{n}{a\PYGZus{}double} \PYG{n}{ay} \PYG{o}{=} \PYG{n}{ax0} \PYG{o}{+} \PYG{n}{ax1}\PYG{p}{;}
    \PYG{c+c1}{//}
    \PYG{c+c1}{// check that ay is a variable; its value depends on the value of ax}
    \PYG{n}{ok} \PYG{o}{=} \PYG{n}{ok} \PYG{o}{\PYGZam{}}\PYG{o}{\PYGZam{}} \PYG{n}{ay}\PYG{p}{.}\PYG{n}{variable}\PYG{p}{(}\PYG{p}{)}\PYG{p}{;}
    \PYG{c+c1}{//}
    \PYG{c+c1}{// abort this recording}
    \PYG{n}{cppad\PYGZus{}py}\PYG{o}{:}\PYG{o}{:}\PYG{n}{abort\PYGZus{}recording}\PYG{p}{(}\PYG{p}{)}\PYG{p}{;}
    \PYG{c+c1}{//}
    \PYG{c+c1}{// check that ay is now a parameter, no longer a variable.}
    \PYG{n}{ok} \PYG{o}{=} \PYG{n}{ok} \PYG{o}{\PYGZam{}}\PYG{o}{\PYGZam{}} \PYG{n}{ay}\PYG{p}{.}\PYG{n}{parameter}\PYG{p}{(}\PYG{p}{)}\PYG{p}{;}
    \PYG{c+c1}{//}
    \PYG{c+c1}{// since it is a parameter, we can retrieve its value}
    \PYG{k+kt}{double} \PYG{n}{y} \PYG{o}{=} \PYG{n}{ay}\PYG{p}{.}\PYG{n}{value}\PYG{p}{(}\PYG{p}{)}\PYG{p}{;}
    \PYG{c+c1}{//}
    \PYG{c+c1}{// its value should be x0 + x1}
    \PYG{n}{ok} \PYG{o}{=} \PYG{n}{ok} \PYG{o}{\PYGZam{}}\PYG{o}{\PYGZam{}} \PYG{n}{y}  \PYG{o}{=}\PYG{o}{=} \PYG{n}{x}\PYG{p}{[}\PYG{l+m+mi}{0}\PYG{p}{]} \PYG{o}{+} \PYG{n}{x}\PYG{p}{[}\PYG{l+m+mi}{1}\PYG{p}{]}\PYG{p}{;}
    \PYG{c+c1}{//}
    \PYG{c+c1}{// an abort when not recording has no effect}
    \PYG{n}{cppad\PYGZus{}py}\PYG{o}{:}\PYG{o}{:}\PYG{n}{abort\PYGZus{}recording}\PYG{p}{(}\PYG{p}{)}\PYG{p}{;}
    \PYG{c+c1}{//}
    \PYG{k}{return}\PYG{p}{(} \PYG{n}{ok} \PYG{p}{)}\PYG{p}{;}
\PYG{p}{\PYGZcb{}}
\end{sphinxVerbatim}


\bigskip\hrule\bigskip


xsrst input file: \sphinxcode{\sphinxupquote{example/cplusplus/fun\_abort\_xam.cpp}}

\index{example@\spxentry{example}}\ignorespaces 

\subparagraph{Example}
\label{\detokenize{xsrst/cpp_abort_recording:example}}\label{\detokenize{xsrst/cpp_abort_recording:cpp-abort-recording-example}}\label{\detokenize{xsrst/cpp_abort_recording:index-3}}
{\hyperref[\detokenize{xsrst/fun_abort_xam_cpp:id1}]{\sphinxcrossref{\DUrole{std,std-ref}{c++}}}}.


\bigskip\hrule\bigskip


xsrst input file: \sphinxcode{\sphinxupquote{lib/cplusplus/fun.cpp}}


\subparagraph{cpp\_fun\_ctor}
\label{\detokenize{xsrst/cpp_fun_ctor:cpp-fun-ctor}}\label{\detokenize{xsrst/cpp_fun_ctor::doc}}
\index{cpp\_fun\_ctor@\spxentry{cpp\_fun\_ctor}}\index{stop@\spxentry{stop}}\index{current@\spxentry{current}}\index{recording@\spxentry{recording}}\index{store@\spxentry{store}}\index{function@\spxentry{function}}\index{object@\spxentry{object}}\ignorespaces 

\subparagraph{3.2.3.3 Stop Current Recording and Store Function Object}
\label{\detokenize{xsrst/cpp_fun_ctor:stop-current-recording-and-store-function-object}}\label{\detokenize{xsrst/cpp_fun_ctor:id1}}\begin{itemize}
\item {} \begin{description}
\item[{{\hyperref[\detokenize{xsrst/cpp_fun_ctor:cpp-fun-ctor-syntax}]{\sphinxcrossref{\DUrole{std,std-ref}{Syntax}}}}}] \leavevmode\begin{itemize}
\item {} 
{\hyperref[\detokenize{xsrst/cpp_fun_ctor:cpp-fun-ctor-syntax-d-fun}]{\sphinxcrossref{\DUrole{std,std-ref}{d\_fun}}}}

\item {} 
{\hyperref[\detokenize{xsrst/cpp_fun_ctor:cpp-fun-ctor-syntax-a-fun}]{\sphinxcrossref{\DUrole{std,std-ref}{a\_fun}}}}

\end{itemize}

\end{description}

\item {} 
{\hyperref[\detokenize{xsrst/cpp_fun_ctor:cpp-fun-ctor-ax}]{\sphinxcrossref{\DUrole{std,std-ref}{ax}}}}

\item {} 
{\hyperref[\detokenize{xsrst/cpp_fun_ctor:cpp-fun-ctor-ay}]{\sphinxcrossref{\DUrole{std,std-ref}{ay}}}}

\item {} \begin{description}
\item[{{\hyperref[\detokenize{xsrst/cpp_fun_ctor:cpp-fun-ctor-f}]{\sphinxcrossref{\DUrole{std,std-ref}{f}}}}}] \leavevmode\begin{itemize}
\item {} 
{\hyperref[\detokenize{xsrst/cpp_fun_ctor:cpp-fun-ctor-f-empty-function}]{\sphinxcrossref{\DUrole{std,std-ref}{Empty Function}}}}

\end{itemize}

\end{description}

\item {} 
{\hyperref[\detokenize{xsrst/cpp_fun_ctor:cpp-fun-ctor-af}]{\sphinxcrossref{\DUrole{std,std-ref}{af}}}}

\item {} 
{\hyperref[\detokenize{xsrst/cpp_fun_ctor:cpp-fun-ctor-example}]{\sphinxcrossref{\DUrole{std,std-ref}{Example}}}}

\end{itemize}

\index{syntax@\spxentry{syntax}}\ignorespaces 

\subparagraph{Syntax}
\label{\detokenize{xsrst/cpp_fun_ctor:syntax}}\label{\detokenize{xsrst/cpp_fun_ctor:cpp-fun-ctor-syntax}}\label{\detokenize{xsrst/cpp_fun_ctor:index-1}}
\index{d\_fun@\spxentry{d\_fun}}\ignorespaces 

\subparagraph{d\_fun}
\label{\detokenize{xsrst/cpp_fun_ctor:d-fun}}\label{\detokenize{xsrst/cpp_fun_ctor:cpp-fun-ctor-syntax-d-fun}}\label{\detokenize{xsrst/cpp_fun_ctor:index-2}}
\begin{DUlineblock}{0em}
\item[] \sphinxstyleemphasis{f} =  \sphinxcode{\sphinxupquote{cppad\_py::d\_fun}} ( \sphinxstyleemphasis{ax} , \sphinxstyleemphasis{ay} )
\end{DUlineblock}

\index{a\_fun@\spxentry{a\_fun}}\ignorespaces 

\subparagraph{a\_fun}
\label{\detokenize{xsrst/cpp_fun_ctor:a-fun}}\label{\detokenize{xsrst/cpp_fun_ctor:cpp-fun-ctor-syntax-a-fun}}\label{\detokenize{xsrst/cpp_fun_ctor:index-3}}
\begin{DUlineblock}{0em}
\item[] \sphinxstyleemphasis{af} =  \sphinxcode{\sphinxupquote{cppad\_py::a\_fun}} ( \sphinxstyleemphasis{f} )
\end{DUlineblock}

\index{ax@\spxentry{ax}}\ignorespaces 

\subparagraph{ax}
\label{\detokenize{xsrst/cpp_fun_ctor:ax}}\label{\detokenize{xsrst/cpp_fun_ctor:cpp-fun-ctor-ax}}\label{\detokenize{xsrst/cpp_fun_ctor:index-4}}
This argument has prototype

\begin{DUlineblock}{0em}
\item[]         \sphinxcode{\sphinxupquote{const vec\_a\_double\&}} \sphinxstyleemphasis{ax}
\end{DUlineblock}

and must be the same as
{\hyperref[\detokenize{xsrst/cpp_independent:cpp-independent-ax}]{\sphinxcrossref{\DUrole{std,std-ref}{ax}}}}
returned by the previous call to \sphinxcode{\sphinxupquote{independent}} ; i.e.,
it must be the independent variable vector.
We use the notation \sphinxstyleemphasis{n} = \sphinxstyleemphasis{ax} . \sphinxcode{\sphinxupquote{size}} ()
to denote the number of independent variables.

\index{ay@\spxentry{ay}}\ignorespaces 

\subparagraph{ay}
\label{\detokenize{xsrst/cpp_fun_ctor:ay}}\label{\detokenize{xsrst/cpp_fun_ctor:cpp-fun-ctor-ay}}\label{\detokenize{xsrst/cpp_fun_ctor:index-5}}
This argument has prototype

\begin{DUlineblock}{0em}
\item[]         \sphinxcode{\sphinxupquote{const vec\_a\_double\&}} \sphinxstyleemphasis{ax}
\end{DUlineblock}

It specifies the dependent variables.
We use the notation \sphinxstyleemphasis{m} = \sphinxstyleemphasis{ay} . \sphinxcode{\sphinxupquote{size}} ()
to denote the number of dependent variables.

\index{f@\spxentry{f}}\ignorespaces 

\subparagraph{f}
\label{\detokenize{xsrst/cpp_fun_ctor:f}}\label{\detokenize{xsrst/cpp_fun_ctor:cpp-fun-ctor-f}}\label{\detokenize{xsrst/cpp_fun_ctor:index-6}}
This result has prototype

\begin{DUlineblock}{0em}
\item[]         \sphinxcode{\sphinxupquote{cppad\_py::d\_fun}} \sphinxstyleemphasis{f}
\end{DUlineblock}

It has a representation for the floating point operations
that mapped the independent variables \sphinxstyleemphasis{ax}
to the dependent variables \sphinxstyleemphasis{ay} .
This object computes function and derivative values using \sphinxcode{\sphinxupquote{double}} .

\index{empty@\spxentry{empty}}\index{function@\spxentry{function}}\ignorespaces 

\subparagraph{Empty Function}
\label{\detokenize{xsrst/cpp_fun_ctor:empty-function}}\label{\detokenize{xsrst/cpp_fun_ctor:cpp-fun-ctor-f-empty-function}}\label{\detokenize{xsrst/cpp_fun_ctor:index-7}}
In the case where \sphinxstyleemphasis{ax} and \sphinxstyleemphasis{ay} have size zero,
the function is ‘empty’ and all its sizes are zero; see
{\hyperref[\detokenize{xsrst/cpp_fun_property:id1}]{\sphinxcrossref{\DUrole{std,std-ref}{cpp\_fun\_property}}}}.

\index{af@\spxentry{af}}\ignorespaces 

\subparagraph{af}
\label{\detokenize{xsrst/cpp_fun_ctor:af}}\label{\detokenize{xsrst/cpp_fun_ctor:cpp-fun-ctor-af}}\label{\detokenize{xsrst/cpp_fun_ctor:index-8}}
This result has prototype

\begin{DUlineblock}{0em}
\item[]         \sphinxcode{\sphinxupquote{cppad\_py::a\_fun}} \sphinxstyleemphasis{af}
\end{DUlineblock}

It has a representation of the same function as \sphinxstyleemphasis{f} .
This object computes function and derivative values using \sphinxcode{\sphinxupquote{a\_double}} .
Initially, there are not Taylor coefficient stored in \sphinxstyleemphasis{af} ; i.e.,
{\hyperref[\detokenize{xsrst/cpp_fun_property:cpp-fun-property-size-order}]{\sphinxcrossref{\DUrole{std,std-ref}{af\_size\_order()}}}} is zero.

\index{example@\spxentry{example}}\ignorespaces 

\subparagraph{Example}
\label{\detokenize{xsrst/cpp_fun_ctor:example}}\label{\detokenize{xsrst/cpp_fun_ctor:cpp-fun-ctor-example}}\label{\detokenize{xsrst/cpp_fun_ctor:index-9}}
All of the examples use these constructors.


\bigskip\hrule\bigskip


xsrst input file: \sphinxcode{\sphinxupquote{lib/cplusplus/fun.cpp}}


\subparagraph{cpp\_fun\_property}
\label{\detokenize{xsrst/cpp_fun_property:cpp-fun-property}}\label{\detokenize{xsrst/cpp_fun_property::doc}}
\index{cpp\_fun\_property@\spxentry{cpp\_fun\_property}}\index{properties@\spxentry{properties}}\index{function@\spxentry{function}}\index{object@\spxentry{object}}\ignorespaces 

\subparagraph{3.2.3.4 Properties of a Function Object}
\label{\detokenize{xsrst/cpp_fun_property:properties-of-a-function-object}}\label{\detokenize{xsrst/cpp_fun_property:id1}}\begin{itemize}
\item {} 
{\hyperref[\detokenize{xsrst/cpp_fun_property:cpp-fun-property-syntax}]{\sphinxcrossref{\DUrole{std,std-ref}{Syntax}}}}

\item {} 
{\hyperref[\detokenize{xsrst/cpp_fun_property:cpp-fun-property-f}]{\sphinxcrossref{\DUrole{std,std-ref}{f}}}}

\item {} 
{\hyperref[\detokenize{xsrst/cpp_fun_property:cpp-fun-property-size-domain}]{\sphinxcrossref{\DUrole{std,std-ref}{size\_domain}}}}

\item {} 
{\hyperref[\detokenize{xsrst/cpp_fun_property:cpp-fun-property-size-range}]{\sphinxcrossref{\DUrole{std,std-ref}{size\_range}}}}

\item {} 
{\hyperref[\detokenize{xsrst/cpp_fun_property:cpp-fun-property-size-var}]{\sphinxcrossref{\DUrole{std,std-ref}{size\_var}}}}

\item {} 
{\hyperref[\detokenize{xsrst/cpp_fun_property:cpp-fun-property-size-op}]{\sphinxcrossref{\DUrole{std,std-ref}{size\_op}}}}

\item {} 
{\hyperref[\detokenize{xsrst/cpp_fun_property:cpp-fun-property-size-order}]{\sphinxcrossref{\DUrole{std,std-ref}{size\_order}}}}

\item {} 
{\hyperref[\detokenize{xsrst/cpp_fun_property:cpp-fun-property-example}]{\sphinxcrossref{\DUrole{std,std-ref}{Example}}}}

\end{itemize}

\index{syntax@\spxentry{syntax}}\ignorespaces 

\subparagraph{Syntax}
\label{\detokenize{xsrst/cpp_fun_property:syntax}}\label{\detokenize{xsrst/cpp_fun_property:cpp-fun-property-syntax}}\label{\detokenize{xsrst/cpp_fun_property:index-1}}
\begin{DUlineblock}{0em}
\item[] \sphinxcode{\sphinxupquote{a\_fun}} \sphinxstyleemphasis{f} ( \sphinxstyleemphasis{f} )
\item[] \sphinxstyleemphasis{n} = \sphinxstyleemphasis{f} . \sphinxcode{\sphinxupquote{size\_domain}} ()
\item[] \sphinxstyleemphasis{m} = \sphinxstyleemphasis{f} . \sphinxcode{\sphinxupquote{size\_range}} ()
\item[] \sphinxstyleemphasis{v} = \sphinxstyleemphasis{f} . \sphinxcode{\sphinxupquote{size\_var}} ()
\item[] \sphinxstyleemphasis{p} = \sphinxstyleemphasis{f} . \sphinxcode{\sphinxupquote{size\_op}} ()
\item[] \sphinxstyleemphasis{q} = \sphinxstyleemphasis{f} . \sphinxcode{\sphinxupquote{size\_order}} ()
\end{DUlineblock}

\index{f@\spxentry{f}}\ignorespaces 

\subparagraph{f}
\label{\detokenize{xsrst/cpp_fun_property:f}}\label{\detokenize{xsrst/cpp_fun_property:cpp-fun-property-f}}\label{\detokenize{xsrst/cpp_fun_property:index-2}}
This is either a
{\hyperref[\detokenize{xsrst/cpp_fun_ctor:cpp-fun-ctor-syntax-d-fun}]{\sphinxcrossref{\DUrole{std,std-ref}{d\_fun}}}} or
{\hyperref[\detokenize{xsrst/cpp_fun_ctor:cpp-fun-ctor-syntax-a-fun}]{\sphinxcrossref{\DUrole{std,std-ref}{a\_fun}}}} function object
and is \sphinxcode{\sphinxupquote{const}} .

\index{size\_domain@\spxentry{size\_domain}}\ignorespaces 

\subparagraph{size\_domain}
\label{\detokenize{xsrst/cpp_fun_property:size-domain}}\label{\detokenize{xsrst/cpp_fun_property:cpp-fun-property-size-domain}}\label{\detokenize{xsrst/cpp_fun_property:index-3}}
The return value has prototype

\begin{DUlineblock}{0em}
\item[]         \sphinxcode{\sphinxupquote{int}} \sphinxstyleemphasis{n}
\end{DUlineblock}

and is the size of the vector
{\hyperref[\detokenize{xsrst/cpp_fun_ctor:cpp-fun-ctor-ax}]{\sphinxcrossref{\DUrole{std,std-ref}{ax}}}} in the function constructor; i.e.,
the number of independent variables.

\index{size\_range@\spxentry{size\_range}}\ignorespaces 

\subparagraph{size\_range}
\label{\detokenize{xsrst/cpp_fun_property:size-range}}\label{\detokenize{xsrst/cpp_fun_property:cpp-fun-property-size-range}}\label{\detokenize{xsrst/cpp_fun_property:index-4}}
The return value has prototype

\begin{DUlineblock}{0em}
\item[]         \sphinxcode{\sphinxupquote{int}} \sphinxstyleemphasis{m}
\end{DUlineblock}

and is the size of the vector
{\hyperref[\detokenize{xsrst/cpp_fun_ctor:cpp-fun-ctor-ay}]{\sphinxcrossref{\DUrole{std,std-ref}{ay}}}} in the function constructor; i.e.,
the number of dependent variables.

\index{size\_var@\spxentry{size\_var}}\ignorespaces 

\subparagraph{size\_var}
\label{\detokenize{xsrst/cpp_fun_property:size-var}}\label{\detokenize{xsrst/cpp_fun_property:cpp-fun-property-size-var}}\label{\detokenize{xsrst/cpp_fun_property:index-5}}
The return value has prototype

\begin{DUlineblock}{0em}
\item[]         \sphinxcode{\sphinxupquote{int}} \sphinxstyleemphasis{v}
\end{DUlineblock}

and is the number of variables in the function.
This includes the independent variables, dependent variables,
and any variables that are used to compute the dependent variables
from the independent variables.

\index{size\_op@\spxentry{size\_op}}\ignorespaces 

\subparagraph{size\_op}
\label{\detokenize{xsrst/cpp_fun_property:size-op}}\label{\detokenize{xsrst/cpp_fun_property:cpp-fun-property-size-op}}\label{\detokenize{xsrst/cpp_fun_property:index-6}}
The return value has prototype

\begin{DUlineblock}{0em}
\item[]         \sphinxcode{\sphinxupquote{int}} \sphinxstyleemphasis{p}
\end{DUlineblock}

and is the number of atomic operations that are used to express
the dependent variables as a function of the independent variables.

\index{size\_order@\spxentry{size\_order}}\ignorespaces 

\subparagraph{size\_order}
\label{\detokenize{xsrst/cpp_fun_property:size-order}}\label{\detokenize{xsrst/cpp_fun_property:cpp-fun-property-size-order}}\label{\detokenize{xsrst/cpp_fun_property:index-7}}
The return value has prototype

\begin{DUlineblock}{0em}
\item[]         \sphinxcode{\sphinxupquote{int}} \sphinxstyleemphasis{q}
\end{DUlineblock}

and is the number of Taylor coefficients currently stored in \sphinxstyleemphasis{f} ,
for every variable in the operation sequence corresponding to \sphinxstyleemphasis{f} .
These coefficients are computed by {\hyperref[\detokenize{xsrst/cpp_fun_forward:id1}]{\sphinxcrossref{\DUrole{std,std-ref}{cpp\_fun\_forward}}}}.
This is different from the other function properties in that it can change
after each call to \sphinxstyleemphasis{f} . \sphinxcode{\sphinxupquote{forward}} ; see
{\hyperref[\detokenize{xsrst/cpp_fun_forward:cpp-fun-forward-p-size-order}]{\sphinxcrossref{\DUrole{std,std-ref}{size\_order}}}} in the forward mode section.
The initial value for this property, when the object \sphinxstyleemphasis{f}
or \sphinxstyleemphasis{af} is created, is zero.


\subparagraph{fun\_property\_xam\_cpp}
\label{\detokenize{xsrst/fun_property_xam_cpp:fun-property-xam-cpp}}\label{\detokenize{xsrst/fun_property_xam_cpp::doc}}
\index{fun\_property\_xam\_cpp@\spxentry{fun\_property\_xam\_cpp}}\index{c++:@\spxentry{c++:}}\index{function@\spxentry{function}}\index{properties:@\spxentry{properties:}}\index{example@\spxentry{example}}\index{test@\spxentry{test}}\ignorespaces 

\subparagraph{3.2.3.4.1 C++: function Properties: Example and Test}
\label{\detokenize{xsrst/fun_property_xam_cpp:c-function-properties-example-and-test}}\label{\detokenize{xsrst/fun_property_xam_cpp:id1}}
\begin{sphinxVerbatim}[commandchars=\\\{\}]
\PYG{c+cp}{\PYGZsh{}}\PYG{c+cp}{ include \PYGZlt{}cstdio\PYGZgt{}}
\PYG{c+cp}{\PYGZsh{}}\PYG{c+cp}{ include \PYGZlt{}cppad}\PYG{c+cp}{/}\PYG{c+cp}{py}\PYG{c+cp}{/}\PYG{c+cp}{cppad\PYGZus{}py.hpp\PYGZgt{}}

\PYG{k+kt}{bool} \PYG{n+nf}{fun\PYGZus{}property\PYGZus{}xam}\PYG{p}{(}\PYG{k+kt}{void}\PYG{p}{)} \PYG{p}{\PYGZob{}}
    \PYG{k}{using} \PYG{n}{cppad\PYGZus{}py}\PYG{o}{:}\PYG{o}{:}\PYG{n}{a\PYGZus{}double}\PYG{p}{;}
    \PYG{k}{using} \PYG{n}{cppad\PYGZus{}py}\PYG{o}{:}\PYG{o}{:}\PYG{n}{vec\PYGZus{}double}\PYG{p}{;}
    \PYG{k}{using} \PYG{n}{cppad\PYGZus{}py}\PYG{o}{:}\PYG{o}{:}\PYG{n}{vec\PYGZus{}a\PYGZus{}double}\PYG{p}{;}
    \PYG{k}{using} \PYG{n}{cppad\PYGZus{}py}\PYG{o}{:}\PYG{o}{:}\PYG{n}{d\PYGZus{}fun}\PYG{p}{;}
    \PYG{k}{using} \PYG{n}{cppad\PYGZus{}py}\PYG{o}{:}\PYG{o}{:}\PYG{n}{a\PYGZus{}fun}\PYG{p}{;}
    \PYG{c+c1}{//}
    \PYG{c+c1}{// initialize return variable}
    \PYG{k+kt}{bool} \PYG{n}{ok} \PYG{o}{=} \PYG{n+nb}{true}\PYG{p}{;}
    \PYG{c+c1}{// \PYGZhy{}\PYGZhy{}\PYGZhy{}\PYGZhy{}\PYGZhy{}\PYGZhy{}\PYGZhy{}\PYGZhy{}\PYGZhy{}\PYGZhy{}\PYGZhy{}\PYGZhy{}\PYGZhy{}\PYGZhy{}\PYGZhy{}\PYGZhy{}\PYGZhy{}\PYGZhy{}\PYGZhy{}\PYGZhy{}\PYGZhy{}\PYGZhy{}\PYGZhy{}\PYGZhy{}\PYGZhy{}\PYGZhy{}\PYGZhy{}\PYGZhy{}\PYGZhy{}\PYGZhy{}\PYGZhy{}\PYGZhy{}\PYGZhy{}\PYGZhy{}\PYGZhy{}\PYGZhy{}\PYGZhy{}\PYGZhy{}\PYGZhy{}\PYGZhy{}\PYGZhy{}\PYGZhy{}\PYGZhy{}\PYGZhy{}\PYGZhy{}\PYGZhy{}\PYGZhy{}\PYGZhy{}\PYGZhy{}\PYGZhy{}\PYGZhy{}\PYGZhy{}\PYGZhy{}\PYGZhy{}\PYGZhy{}\PYGZhy{}\PYGZhy{}\PYGZhy{}\PYGZhy{}\PYGZhy{}\PYGZhy{}\PYGZhy{}\PYGZhy{}\PYGZhy{}\PYGZhy{}\PYGZhy{}\PYGZhy{}\PYGZhy{}\PYGZhy{}\PYGZhy{}}
    \PYG{k+kt}{int} \PYG{n}{n\PYGZus{}ind} \PYG{o}{=} \PYG{l+m+mi}{1}\PYG{p}{;} \PYG{c+c1}{// number of independent variables}
    \PYG{k+kt}{int} \PYG{n}{n\PYGZus{}dep} \PYG{o}{=} \PYG{l+m+mi}{2}\PYG{p}{;} \PYG{c+c1}{// number of dependent variables}
    \PYG{k+kt}{int} \PYG{n}{n\PYGZus{}var} \PYG{o}{=} \PYG{l+m+mi}{1}\PYG{p}{;} \PYG{c+c1}{// phantom variable at address 0}
    \PYG{k+kt}{int} \PYG{n}{n\PYGZus{}op}  \PYG{o}{=} \PYG{l+m+mi}{1}\PYG{p}{;} \PYG{c+c1}{// special operator at beginning}
    \PYG{c+c1}{//}
    \PYG{c+c1}{// dimension some vectors}
    \PYG{n}{vec\PYGZus{}double} \PYG{n}{x}\PYG{p}{(}\PYG{n}{n\PYGZus{}ind}\PYG{p}{)}\PYG{p}{;}
    \PYG{n}{vec\PYGZus{}a\PYGZus{}double} \PYG{n}{ay}\PYG{p}{(}\PYG{n}{n\PYGZus{}dep}\PYG{p}{)}\PYG{p}{;}
    \PYG{c+c1}{//}
    \PYG{c+c1}{// independent variables}
    \PYG{n}{x}\PYG{p}{[}\PYG{l+m+mi}{0}\PYG{p}{]}            \PYG{o}{=} \PYG{l+m+mf}{1.0}\PYG{p}{;}
    \PYG{n}{vec\PYGZus{}a\PYGZus{}double} \PYG{n}{ax} \PYG{o}{=} \PYG{n}{cppad\PYGZus{}py}\PYG{o}{:}\PYG{o}{:}\PYG{n}{independent}\PYG{p}{(}\PYG{n}{x}\PYG{p}{)}\PYG{p}{;}
    \PYG{n}{n\PYGZus{}var}           \PYG{o}{=} \PYG{n}{n\PYGZus{}var} \PYG{o}{+} \PYG{n}{n\PYGZus{}ind}\PYG{p}{;} \PYG{c+c1}{// one for each indpendent}
    \PYG{n}{n\PYGZus{}op}            \PYG{o}{=} \PYG{n}{n\PYGZus{}op} \PYG{o}{+} \PYG{n}{n\PYGZus{}ind}\PYG{p}{;}
    \PYG{c+c1}{//}
    \PYG{c+c1}{// first dependent variable}
    \PYG{n}{ay}\PYG{p}{[}\PYG{l+m+mi}{0}\PYG{p}{]} \PYG{o}{=} \PYG{n}{ax}\PYG{p}{[}\PYG{l+m+mi}{0}\PYG{p}{]} \PYG{o}{+} \PYG{n}{ax}\PYG{p}{[}\PYG{l+m+mi}{0}\PYG{p}{]}\PYG{p}{;}
    \PYG{n}{n\PYGZus{}var} \PYG{o}{=} \PYG{n}{n\PYGZus{}var} \PYG{o}{+} \PYG{l+m+mi}{1}\PYG{p}{;} \PYG{c+c1}{// one variable and operator}
    \PYG{n}{n\PYGZus{}op} \PYG{o}{=} \PYG{n}{n\PYGZus{}op} \PYG{o}{+} \PYG{l+m+mi}{1}\PYG{p}{;}
    \PYG{c+c1}{//}
    \PYG{c+c1}{// second dependent variable}
    \PYG{n}{a\PYGZus{}double} \PYG{n}{ax0} \PYG{o}{=} \PYG{n}{ax}\PYG{p}{[}\PYG{l+m+mi}{0}\PYG{p}{]}\PYG{p}{;}
    \PYG{n}{ay}\PYG{p}{[}\PYG{l+m+mi}{1}\PYG{p}{]}        \PYG{o}{=} \PYG{n}{ax0}\PYG{p}{.}\PYG{n}{sin}\PYG{p}{(}\PYG{p}{)}\PYG{p}{;}
    \PYG{n}{n\PYGZus{}var}        \PYG{o}{=} \PYG{n}{n\PYGZus{}var} \PYG{o}{+} \PYG{l+m+mi}{2}\PYG{p}{;} \PYG{c+c1}{// two varialbes, one operator}
    \PYG{n}{n\PYGZus{}op}         \PYG{o}{=} \PYG{n}{n\PYGZus{}op} \PYG{o}{+} \PYG{l+m+mi}{1}\PYG{p}{;}
    \PYG{c+c1}{//}
    \PYG{c+c1}{// define f(x) = y}
    \PYG{n}{d\PYGZus{}fun} \PYG{n}{f}\PYG{p}{(}\PYG{n}{ax}\PYG{p}{,} \PYG{n}{ay}\PYG{p}{)}\PYG{p}{;}
    \PYG{n}{n\PYGZus{}op}     \PYG{o}{=} \PYG{n}{n\PYGZus{}op} \PYG{o}{+} \PYG{l+m+mi}{1}\PYG{p}{;} \PYG{c+c1}{// speical operator at end}
    \PYG{c+c1}{//}
    \PYG{c+c1}{// check f properties}
    \PYG{n}{ok} \PYG{o}{=} \PYG{n}{ok} \PYG{o}{\PYGZam{}}\PYG{o}{\PYGZam{}} \PYG{n}{f}\PYG{p}{.}\PYG{n}{size\PYGZus{}domain}\PYG{p}{(}\PYG{p}{)} \PYG{o}{=}\PYG{o}{=} \PYG{n}{n\PYGZus{}ind}\PYG{p}{;}
    \PYG{n}{ok} \PYG{o}{=} \PYG{n}{ok} \PYG{o}{\PYGZam{}}\PYG{o}{\PYGZam{}} \PYG{n}{f}\PYG{p}{.}\PYG{n}{size\PYGZus{}range}\PYG{p}{(}\PYG{p}{)}  \PYG{o}{=}\PYG{o}{=} \PYG{n}{n\PYGZus{}dep}\PYG{p}{;}
    \PYG{n}{ok} \PYG{o}{=} \PYG{n}{ok} \PYG{o}{\PYGZam{}}\PYG{o}{\PYGZam{}} \PYG{n}{f}\PYG{p}{.}\PYG{n}{size\PYGZus{}var}\PYG{p}{(}\PYG{p}{)}    \PYG{o}{=}\PYG{o}{=} \PYG{n}{n\PYGZus{}var}\PYG{p}{;}
    \PYG{n}{ok} \PYG{o}{=} \PYG{n}{ok} \PYG{o}{\PYGZam{}}\PYG{o}{\PYGZam{}} \PYG{n}{f}\PYG{p}{.}\PYG{n}{size\PYGZus{}op}\PYG{p}{(}\PYG{p}{)}     \PYG{o}{=}\PYG{o}{=} \PYG{n}{n\PYGZus{}op}\PYG{p}{;}
    \PYG{n}{ok} \PYG{o}{=} \PYG{n}{ok} \PYG{o}{\PYGZam{}}\PYG{o}{\PYGZam{}} \PYG{n}{f}\PYG{p}{.}\PYG{n}{size\PYGZus{}order}\PYG{p}{(}\PYG{p}{)}  \PYG{o}{=}\PYG{o}{=} \PYG{l+m+mi}{0}\PYG{p}{;}
    \PYG{c+c1}{//}
    \PYG{c+c1}{// compute zero order Taylor coefficients}
    \PYG{n}{vec\PYGZus{}double} \PYG{n}{y}  \PYG{o}{=} \PYG{n}{f}\PYG{p}{.}\PYG{n}{forward}\PYG{p}{(}\PYG{l+m+mi}{0}\PYG{p}{,} \PYG{n}{x}\PYG{p}{)}\PYG{p}{;}
    \PYG{n}{ok} \PYG{o}{=} \PYG{n}{ok} \PYG{o}{\PYGZam{}}\PYG{o}{\PYGZam{}} \PYG{n}{f}\PYG{p}{.}\PYG{n}{size\PYGZus{}order}\PYG{p}{(}\PYG{p}{)} \PYG{o}{=}\PYG{o}{=} \PYG{l+m+mi}{1}\PYG{p}{;}
    \PYG{c+c1}{//}
    \PYG{c+c1}{// create an a\PYGZus{}fun object}
    \PYG{n}{a\PYGZus{}fun} \PYG{n}{af}\PYG{p}{(}\PYG{n}{f}\PYG{p}{)}\PYG{p}{;}
    \PYG{c+c1}{//}
    \PYG{c+c1}{// \PYGZhy{}\PYGZhy{}\PYGZhy{}\PYGZhy{}\PYGZhy{}\PYGZhy{}\PYGZhy{}\PYGZhy{}\PYGZhy{}\PYGZhy{}\PYGZhy{}\PYGZhy{}\PYGZhy{}\PYGZhy{}\PYGZhy{}\PYGZhy{}\PYGZhy{}\PYGZhy{}\PYGZhy{}\PYGZhy{}\PYGZhy{}\PYGZhy{}\PYGZhy{}\PYGZhy{}\PYGZhy{}\PYGZhy{}\PYGZhy{}\PYGZhy{}\PYGZhy{}\PYGZhy{}\PYGZhy{}\PYGZhy{}\PYGZhy{}\PYGZhy{}\PYGZhy{}\PYGZhy{}\PYGZhy{}\PYGZhy{}\PYGZhy{}\PYGZhy{}\PYGZhy{}\PYGZhy{}\PYGZhy{}\PYGZhy{}\PYGZhy{}\PYGZhy{}\PYGZhy{}\PYGZhy{}\PYGZhy{}\PYGZhy{}\PYGZhy{}\PYGZhy{}\PYGZhy{}\PYGZhy{}\PYGZhy{}\PYGZhy{}\PYGZhy{}\PYGZhy{}\PYGZhy{}\PYGZhy{}\PYGZhy{}\PYGZhy{}\PYGZhy{}\PYGZhy{}\PYGZhy{}\PYGZhy{}\PYGZhy{}\PYGZhy{}\PYGZhy{}\PYGZhy{}}
    \PYG{c+c1}{// check af properties}
    \PYG{n}{ok} \PYG{o}{=} \PYG{n}{ok} \PYG{o}{\PYGZam{}}\PYG{o}{\PYGZam{}} \PYG{n}{af}\PYG{p}{.}\PYG{n}{size\PYGZus{}domain}\PYG{p}{(}\PYG{p}{)} \PYG{o}{=}\PYG{o}{=} \PYG{n}{n\PYGZus{}ind}\PYG{p}{;}
    \PYG{n}{ok} \PYG{o}{=} \PYG{n}{ok} \PYG{o}{\PYGZam{}}\PYG{o}{\PYGZam{}} \PYG{n}{af}\PYG{p}{.}\PYG{n}{size\PYGZus{}range}\PYG{p}{(}\PYG{p}{)}  \PYG{o}{=}\PYG{o}{=} \PYG{n}{n\PYGZus{}dep}\PYG{p}{;}
    \PYG{n}{ok} \PYG{o}{=} \PYG{n}{ok} \PYG{o}{\PYGZam{}}\PYG{o}{\PYGZam{}} \PYG{n}{af}\PYG{p}{.}\PYG{n}{size\PYGZus{}var}\PYG{p}{(}\PYG{p}{)}    \PYG{o}{=}\PYG{o}{=} \PYG{n}{n\PYGZus{}var}\PYG{p}{;}
    \PYG{n}{ok} \PYG{o}{=} \PYG{n}{ok} \PYG{o}{\PYGZam{}}\PYG{o}{\PYGZam{}} \PYG{n}{af}\PYG{p}{.}\PYG{n}{size\PYGZus{}op}\PYG{p}{(}\PYG{p}{)}     \PYG{o}{=}\PYG{o}{=} \PYG{n}{n\PYGZus{}op}\PYG{p}{;}
    \PYG{n}{ok} \PYG{o}{=} \PYG{n}{ok} \PYG{o}{\PYGZam{}}\PYG{o}{\PYGZam{}} \PYG{n}{af}\PYG{p}{.}\PYG{n}{size\PYGZus{}order}\PYG{p}{(}\PYG{p}{)}  \PYG{o}{=}\PYG{o}{=} \PYG{l+m+mi}{0}\PYG{p}{;}
    \PYG{c+c1}{// \PYGZhy{}\PYGZhy{}\PYGZhy{}\PYGZhy{}\PYGZhy{}\PYGZhy{}\PYGZhy{}\PYGZhy{}\PYGZhy{}\PYGZhy{}\PYGZhy{}\PYGZhy{}\PYGZhy{}\PYGZhy{}\PYGZhy{}\PYGZhy{}\PYGZhy{}\PYGZhy{}\PYGZhy{}\PYGZhy{}\PYGZhy{}\PYGZhy{}\PYGZhy{}\PYGZhy{}\PYGZhy{}\PYGZhy{}\PYGZhy{}\PYGZhy{}\PYGZhy{}\PYGZhy{}\PYGZhy{}\PYGZhy{}\PYGZhy{}\PYGZhy{}\PYGZhy{}\PYGZhy{}\PYGZhy{}\PYGZhy{}\PYGZhy{}\PYGZhy{}\PYGZhy{}\PYGZhy{}\PYGZhy{}\PYGZhy{}\PYGZhy{}\PYGZhy{}\PYGZhy{}\PYGZhy{}\PYGZhy{}\PYGZhy{}\PYGZhy{}\PYGZhy{}\PYGZhy{}\PYGZhy{}\PYGZhy{}\PYGZhy{}\PYGZhy{}\PYGZhy{}\PYGZhy{}\PYGZhy{}\PYGZhy{}\PYGZhy{}\PYGZhy{}\PYGZhy{}\PYGZhy{}\PYGZhy{}\PYGZhy{}\PYGZhy{}\PYGZhy{}\PYGZhy{}}
    \PYG{c+c1}{// The empty function}
    \PYG{n}{ax}\PYG{p}{.}\PYG{n}{resize}\PYG{p}{(}\PYG{l+m+mi}{0}\PYG{p}{)}\PYG{p}{;}
    \PYG{n}{ay}\PYG{p}{.}\PYG{n}{resize}\PYG{p}{(}\PYG{l+m+mi}{0}\PYG{p}{)}\PYG{p}{;}
    \PYG{n}{f} \PYG{o}{=} \PYG{n}{d\PYGZus{}fun}\PYG{p}{(}\PYG{n}{ax}\PYG{p}{,} \PYG{n}{ay}\PYG{p}{)}\PYG{p}{;}
    \PYG{n}{ok} \PYG{o}{=} \PYG{n}{ok} \PYG{o}{\PYGZam{}}\PYG{o}{\PYGZam{}} \PYG{n}{f}\PYG{p}{.}\PYG{n}{size\PYGZus{}domain}\PYG{p}{(}\PYG{p}{)} \PYG{o}{=}\PYG{o}{=} \PYG{l+m+mi}{0}\PYG{p}{;}
    \PYG{n}{ok} \PYG{o}{=} \PYG{n}{ok} \PYG{o}{\PYGZam{}}\PYG{o}{\PYGZam{}} \PYG{n}{f}\PYG{p}{.}\PYG{n}{size\PYGZus{}range}\PYG{p}{(}\PYG{p}{)}  \PYG{o}{=}\PYG{o}{=} \PYG{l+m+mi}{0}\PYG{p}{;}
    \PYG{n}{ok} \PYG{o}{=} \PYG{n}{ok} \PYG{o}{\PYGZam{}}\PYG{o}{\PYGZam{}} \PYG{n}{f}\PYG{p}{.}\PYG{n}{size\PYGZus{}var}\PYG{p}{(}\PYG{p}{)}    \PYG{o}{=}\PYG{o}{=} \PYG{l+m+mi}{0}\PYG{p}{;}
    \PYG{n}{ok} \PYG{o}{=} \PYG{n}{ok} \PYG{o}{\PYGZam{}}\PYG{o}{\PYGZam{}} \PYG{n}{f}\PYG{p}{.}\PYG{n}{size\PYGZus{}op}\PYG{p}{(}\PYG{p}{)}     \PYG{o}{=}\PYG{o}{=} \PYG{l+m+mi}{0}\PYG{p}{;}
    \PYG{n}{ok} \PYG{o}{=} \PYG{n}{ok} \PYG{o}{\PYGZam{}}\PYG{o}{\PYGZam{}} \PYG{n}{f}\PYG{p}{.}\PYG{n}{size\PYGZus{}order}\PYG{p}{(}\PYG{p}{)}  \PYG{o}{=}\PYG{o}{=} \PYG{l+m+mi}{0}\PYG{p}{;}
    \PYG{c+c1}{//}
    \PYG{k}{return}\PYG{p}{(} \PYG{n}{ok}  \PYG{p}{)}\PYG{p}{;}
\PYG{p}{\PYGZcb{}}
\end{sphinxVerbatim}


\bigskip\hrule\bigskip


xsrst input file: \sphinxcode{\sphinxupquote{example/cplusplus/fun\_property\_xam.cpp}}

\index{example@\spxentry{example}}\ignorespaces 

\subparagraph{Example}
\label{\detokenize{xsrst/cpp_fun_property:example}}\label{\detokenize{xsrst/cpp_fun_property:cpp-fun-property-example}}\label{\detokenize{xsrst/cpp_fun_property:index-8}}
{\hyperref[\detokenize{xsrst/fun_property_xam_cpp:id1}]{\sphinxcrossref{\DUrole{std,std-ref}{fun\_property\_xam\_cpp}}}}


\bigskip\hrule\bigskip


xsrst input file: \sphinxcode{\sphinxupquote{lib/cplusplus/fun.cpp}}


\subparagraph{cpp\_fun\_new\_dynamic}
\label{\detokenize{xsrst/cpp_fun_new_dynamic:cpp-fun-new-dynamic}}\label{\detokenize{xsrst/cpp_fun_new_dynamic::doc}}
\index{cpp\_fun\_new\_dynamic@\spxentry{cpp\_fun\_new\_dynamic}}\index{change@\spxentry{change}}\index{the@\spxentry{the}}\index{dynamic@\spxentry{dynamic}}\index{parameters@\spxentry{parameters}}\ignorespaces 

\subparagraph{3.2.3.5 Change The Dynamic Parameters}
\label{\detokenize{xsrst/cpp_fun_new_dynamic:change-the-dynamic-parameters}}\label{\detokenize{xsrst/cpp_fun_new_dynamic:id1}}\begin{itemize}
\item {} 
{\hyperref[\detokenize{xsrst/cpp_fun_new_dynamic:cpp-fun-new-dynamic-syntax}]{\sphinxcrossref{\DUrole{std,std-ref}{Syntax}}}}

\item {} 
{\hyperref[\detokenize{xsrst/cpp_fun_new_dynamic:cpp-fun-new-dynamic-f}]{\sphinxcrossref{\DUrole{std,std-ref}{f}}}}

\item {} 
{\hyperref[\detokenize{xsrst/cpp_fun_new_dynamic:cpp-fun-new-dynamic-dynamic}]{\sphinxcrossref{\DUrole{std,std-ref}{dynamic}}}}

\item {} 
{\hyperref[\detokenize{xsrst/cpp_fun_new_dynamic:cpp-fun-new-dynamic-size-order}]{\sphinxcrossref{\DUrole{std,std-ref}{size\_order}}}}

\item {} 
{\hyperref[\detokenize{xsrst/cpp_fun_new_dynamic:cpp-fun-new-dynamic-example}]{\sphinxcrossref{\DUrole{std,std-ref}{Example}}}}

\end{itemize}

\index{syntax@\spxentry{syntax}}\ignorespaces 

\subparagraph{Syntax}
\label{\detokenize{xsrst/cpp_fun_new_dynamic:syntax}}\label{\detokenize{xsrst/cpp_fun_new_dynamic:cpp-fun-new-dynamic-syntax}}\label{\detokenize{xsrst/cpp_fun_new_dynamic:index-1}}
\begin{DUlineblock}{0em}
\item[] \sphinxstyleemphasis{f} . \sphinxcode{\sphinxupquote{new\_dynamic}} ( \sphinxstyleemphasis{dynamic} )
\end{DUlineblock}

\index{f@\spxentry{f}}\ignorespaces 

\subparagraph{f}
\label{\detokenize{xsrst/cpp_fun_new_dynamic:f}}\label{\detokenize{xsrst/cpp_fun_new_dynamic:cpp-fun-new-dynamic-f}}\label{\detokenize{xsrst/cpp_fun_new_dynamic:index-2}}
This is either a
{\hyperref[\detokenize{xsrst/cpp_fun_ctor:cpp-fun-ctor-syntax-d-fun}]{\sphinxcrossref{\DUrole{std,std-ref}{d\_fun}}}} or
{\hyperref[\detokenize{xsrst/cpp_fun_ctor:cpp-fun-ctor-syntax-a-fun}]{\sphinxcrossref{\DUrole{std,std-ref}{a\_fun}}}} function object.

\index{dynamic@\spxentry{dynamic}}\ignorespaces 

\subparagraph{dynamic}
\label{\detokenize{xsrst/cpp_fun_new_dynamic:dynamic}}\label{\detokenize{xsrst/cpp_fun_new_dynamic:cpp-fun-new-dynamic-dynamic}}\label{\detokenize{xsrst/cpp_fun_new_dynamic:index-3}}
If \sphinxstyleemphasis{f} is a \sphinxcode{\sphinxupquote{d\_fun}} or \sphinxcode{\sphinxupquote{a\_fun}} ,
this argument has prototype

\begin{DUlineblock}{0em}
\item[]         \sphinxcode{\sphinxupquote{const vec\_double\&}} \sphinxstyleemphasis{dynamic}
\item[]         \sphinxcode{\sphinxupquote{const vec\_a\_double\&}} \sphinxstyleemphasis{dynamic}
\end{DUlineblock}

and its size must be the same as the size of
{\hyperref[\detokenize{xsrst/cpp_independent:cpp-independent-dynamic}]{\sphinxcrossref{\DUrole{std,std-ref}{dynamic}}}} in the corresponding call to
\sphinxcode{\sphinxupquote{independent}} .
It specifies new values for the dynamic parameters in \sphinxstyleemphasis{f} .

\index{size\_order@\spxentry{size\_order}}\ignorespaces 

\subparagraph{size\_order}
\label{\detokenize{xsrst/cpp_fun_new_dynamic:size-order}}\label{\detokenize{xsrst/cpp_fun_new_dynamic:cpp-fun-new-dynamic-size-order}}\label{\detokenize{xsrst/cpp_fun_new_dynamic:index-4}}
After this call
{\hyperref[\detokenize{xsrst/cpp_fun_property:cpp-fun-property-size-order}]{\sphinxcrossref{\DUrole{std,std-ref}{f\_size\_order()}}}} is zero.

\index{example@\spxentry{example}}\ignorespaces 

\subparagraph{Example}
\label{\detokenize{xsrst/cpp_fun_new_dynamic:example}}\label{\detokenize{xsrst/cpp_fun_new_dynamic:cpp-fun-new-dynamic-example}}\label{\detokenize{xsrst/cpp_fun_new_dynamic:index-5}}
See {\hyperref[\detokenize{xsrst/fun_dynamic_xam_cpp:id1}]{\sphinxcrossref{\DUrole{std,std-ref}{fun\_dynamic\_xam\_cpp}}}}.


\bigskip\hrule\bigskip


xsrst input file: \sphinxcode{\sphinxupquote{lib/cplusplus/fun.cpp}}


\subparagraph{cpp\_fun\_jacobian}
\label{\detokenize{xsrst/cpp_fun_jacobian:cpp-fun-jacobian}}\label{\detokenize{xsrst/cpp_fun_jacobian::doc}}
\index{cpp\_fun\_jacobian@\spxentry{cpp\_fun\_jacobian}}\index{jacobian@\spxentry{jacobian}}\index{an@\spxentry{an}}\index{ad@\spxentry{ad}}\index{function@\spxentry{function}}\ignorespaces 

\subparagraph{3.2.3.6 Jacobian of an AD Function}
\label{\detokenize{xsrst/cpp_fun_jacobian:jacobian-of-an-ad-function}}\label{\detokenize{xsrst/cpp_fun_jacobian:id1}}\begin{itemize}
\item {} 
{\hyperref[\detokenize{xsrst/cpp_fun_jacobian:cpp-fun-jacobian-syntax}]{\sphinxcrossref{\DUrole{std,std-ref}{Syntax}}}}

\item {} 
{\hyperref[\detokenize{xsrst/cpp_fun_jacobian:cpp-fun-jacobian-f}]{\sphinxcrossref{\DUrole{std,std-ref}{f}}}}

\item {} 
{\hyperref[\detokenize{xsrst/cpp_fun_jacobian:cpp-fun-jacobian-f-x}]{\sphinxcrossref{\DUrole{std,std-ref}{f(x)}}}}

\item {} 
{\hyperref[\detokenize{xsrst/cpp_fun_jacobian:cpp-fun-jacobian-x}]{\sphinxcrossref{\DUrole{std,std-ref}{x}}}}

\item {} 
{\hyperref[\detokenize{xsrst/cpp_fun_jacobian:cpp-fun-jacobian-j}]{\sphinxcrossref{\DUrole{std,std-ref}{J}}}}

\item {} 
{\hyperref[\detokenize{xsrst/cpp_fun_jacobian:cpp-fun-jacobian-example}]{\sphinxcrossref{\DUrole{std,std-ref}{Example}}}}

\end{itemize}

\index{syntax@\spxentry{syntax}}\ignorespaces 

\subparagraph{Syntax}
\label{\detokenize{xsrst/cpp_fun_jacobian:syntax}}\label{\detokenize{xsrst/cpp_fun_jacobian:cpp-fun-jacobian-syntax}}\label{\detokenize{xsrst/cpp_fun_jacobian:index-1}}
\sphinxstyleemphasis{J} = \sphinxstyleemphasis{f} . \sphinxcode{\sphinxupquote{jacobian}} ( \sphinxstyleemphasis{x} )

\index{f@\spxentry{f}}\ignorespaces 

\subparagraph{f}
\label{\detokenize{xsrst/cpp_fun_jacobian:f}}\label{\detokenize{xsrst/cpp_fun_jacobian:cpp-fun-jacobian-f}}\label{\detokenize{xsrst/cpp_fun_jacobian:index-2}}
This is either a
{\hyperref[\detokenize{xsrst/cpp_fun_ctor:cpp-fun-ctor-syntax-d-fun}]{\sphinxcrossref{\DUrole{std,std-ref}{d\_fun}}}} or
{\hyperref[\detokenize{xsrst/cpp_fun_ctor:cpp-fun-ctor-syntax-a-fun}]{\sphinxcrossref{\DUrole{std,std-ref}{a\_fun}}}} function object.
Upon return, the zero order
{\hyperref[\detokenize{xsrst/cpp_fun_forward:cpp-fun-forward-taylor-coefficient}]{\sphinxcrossref{\DUrole{std,std-ref}{taylor\_coefficients}}}} in \sphinxstyleemphasis{f}
correspond to the value of \sphinxstyleemphasis{x} .
The other Taylor coefficients in \sphinxstyleemphasis{f} are unspecified.

\index{f(x)@\spxentry{f(x)}}\ignorespaces 

\subparagraph{f(x)}
\label{\detokenize{xsrst/cpp_fun_jacobian:f-x}}\label{\detokenize{xsrst/cpp_fun_jacobian:cpp-fun-jacobian-f-x}}\label{\detokenize{xsrst/cpp_fun_jacobian:index-3}}
We use the notation \(f: \B{R}^n \rightarrow \B{R}^m\)
for the function corresponding to \sphinxstyleemphasis{f} .
Note that \sphinxstyleemphasis{n} is the size of {\hyperref[\detokenize{xsrst/cpp_fun_ctor:cpp-fun-ctor-ax}]{\sphinxcrossref{\DUrole{std,std-ref}{ax}}}}
and \sphinxstyleemphasis{m} is the size of {\hyperref[\detokenize{xsrst/cpp_fun_ctor:cpp-fun-ctor-ay}]{\sphinxcrossref{\DUrole{std,std-ref}{ay}}}}
in to the constructor for \sphinxstyleemphasis{f} .

\index{x@\spxentry{x}}\ignorespaces 

\subparagraph{x}
\label{\detokenize{xsrst/cpp_fun_jacobian:x}}\label{\detokenize{xsrst/cpp_fun_jacobian:cpp-fun-jacobian-x}}\label{\detokenize{xsrst/cpp_fun_jacobian:index-4}}
If \sphinxstyleemphasis{f} is a \sphinxcode{\sphinxupquote{d\_fun}} or \sphinxcode{\sphinxupquote{a\_fun}} ,
this argument has prototype

\begin{DUlineblock}{0em}
\item[]         \sphinxcode{\sphinxupquote{const vec\_double\&}} \sphinxstyleemphasis{x}
\item[]         \sphinxcode{\sphinxupquote{const vec\_a\_double\&}} \sphinxstyleemphasis{x}
\end{DUlineblock}

and its size must be \sphinxstyleemphasis{n} .
It specifies the argument value at we are computing the Jacobian
\(f'(x)\).

\index{j@\spxentry{j}}\ignorespaces 

\subparagraph{J}
\label{\detokenize{xsrst/cpp_fun_jacobian:j}}\label{\detokenize{xsrst/cpp_fun_jacobian:cpp-fun-jacobian-j}}\label{\detokenize{xsrst/cpp_fun_jacobian:index-5}}
If \sphinxstyleemphasis{f} is a \sphinxcode{\sphinxupquote{d\_fun}} or \sphinxcode{\sphinxupquote{a\_fun}} ,
the result has prototype

\begin{DUlineblock}{0em}
\item[]         \sphinxcode{\sphinxupquote{vec\_double}} \sphinxstyleemphasis{J}
\item[]         \sphinxcode{\sphinxupquote{vec\_a\_double}} \sphinxstyleemphasis{J}
\end{DUlineblock}

respectively and its size is \sphinxstyleemphasis{m} \sphinxcode{\sphinxupquote{*}} \sphinxstyleemphasis{n} .
For \sphinxstyleemphasis{i} between zero and \sphinxstyleemphasis{m} \sphinxhyphen{}1
and \sphinxstyleemphasis{j} between zero and \sphinxstyleemphasis{n} \sphinxhyphen{}1 ,
\begin{equation*}
\begin{split}J [ i * n + j ] = \frac{ \partial f_i }{ \partial x_j } (x)\end{split}
\end{equation*}

\subparagraph{fun\_jacobian\_xam\_cpp}
\label{\detokenize{xsrst/fun_jacobian_xam_cpp:fun-jacobian-xam-cpp}}\label{\detokenize{xsrst/fun_jacobian_xam_cpp::doc}}
\index{fun\_jacobian\_xam\_cpp@\spxentry{fun\_jacobian\_xam\_cpp}}\index{c++:@\spxentry{c++:}}\index{dense@\spxentry{dense}}\index{jacobian@\spxentry{jacobian}}\index{using@\spxentry{using}}\index{ad:@\spxentry{ad:}}\index{example@\spxentry{example}}\index{test@\spxentry{test}}\ignorespaces 

\subparagraph{3.2.3.6.1 C++: Dense Jacobian Using AD: Example and Test}
\label{\detokenize{xsrst/fun_jacobian_xam_cpp:c-dense-jacobian-using-ad-example-and-test}}\label{\detokenize{xsrst/fun_jacobian_xam_cpp:id1}}
\begin{sphinxVerbatim}[commandchars=\\\{\}]
\PYG{c+cp}{\PYGZsh{}}\PYG{c+cp}{ include \PYGZlt{}cstdio\PYGZgt{}}
\PYG{c+cp}{\PYGZsh{}}\PYG{c+cp}{ include \PYGZlt{}cppad}\PYG{c+cp}{/}\PYG{c+cp}{py}\PYG{c+cp}{/}\PYG{c+cp}{cppad\PYGZus{}py.hpp\PYGZgt{}}

\PYG{k+kt}{bool} \PYG{n+nf}{fun\PYGZus{}jacobian\PYGZus{}xam}\PYG{p}{(}\PYG{k+kt}{void}\PYG{p}{)} \PYG{p}{\PYGZob{}}
    \PYG{k}{using} \PYG{n}{cppad\PYGZus{}py}\PYG{o}{:}\PYG{o}{:}\PYG{n}{a\PYGZus{}double}\PYG{p}{;}
    \PYG{k}{using} \PYG{n}{cppad\PYGZus{}py}\PYG{o}{:}\PYG{o}{:}\PYG{n}{vec\PYGZus{}double}\PYG{p}{;}
    \PYG{k}{using} \PYG{n}{cppad\PYGZus{}py}\PYG{o}{:}\PYG{o}{:}\PYG{n}{vec\PYGZus{}a\PYGZus{}double}\PYG{p}{;}
    \PYG{k}{using} \PYG{n}{cppad\PYGZus{}py}\PYG{o}{:}\PYG{o}{:}\PYG{n}{d\PYGZus{}fun}\PYG{p}{;}
    \PYG{k}{using} \PYG{n}{cppad\PYGZus{}py}\PYG{o}{:}\PYG{o}{:}\PYG{n}{a\PYGZus{}fun}\PYG{p}{;}
    \PYG{c+c1}{//}
    \PYG{c+c1}{// initialize return variable}
    \PYG{k+kt}{bool} \PYG{n}{ok} \PYG{o}{=} \PYG{n+nb}{true}\PYG{p}{;}
    \PYG{c+c1}{//\PYGZhy{}\PYGZhy{}\PYGZhy{}\PYGZhy{}\PYGZhy{}\PYGZhy{}\PYGZhy{}\PYGZhy{}\PYGZhy{}\PYGZhy{}\PYGZhy{}\PYGZhy{}\PYGZhy{}\PYGZhy{}\PYGZhy{}\PYGZhy{}\PYGZhy{}\PYGZhy{}\PYGZhy{}\PYGZhy{}\PYGZhy{}\PYGZhy{}\PYGZhy{}\PYGZhy{}\PYGZhy{}\PYGZhy{}\PYGZhy{}\PYGZhy{}\PYGZhy{}\PYGZhy{}\PYGZhy{}\PYGZhy{}\PYGZhy{}\PYGZhy{}\PYGZhy{}\PYGZhy{}\PYGZhy{}\PYGZhy{}\PYGZhy{}\PYGZhy{}\PYGZhy{}\PYGZhy{}\PYGZhy{}\PYGZhy{}\PYGZhy{}\PYGZhy{}\PYGZhy{}\PYGZhy{}\PYGZhy{}\PYGZhy{}\PYGZhy{}\PYGZhy{}\PYGZhy{}\PYGZhy{}\PYGZhy{}\PYGZhy{}\PYGZhy{}\PYGZhy{}\PYGZhy{}\PYGZhy{}\PYGZhy{}\PYGZhy{}\PYGZhy{}\PYGZhy{}\PYGZhy{}\PYGZhy{}\PYGZhy{}\PYGZhy{}\PYGZhy{}\PYGZhy{}\PYGZhy{}\PYGZhy{}}
    \PYG{c+c1}{// number of dependent and independent variables}
    \PYG{k+kt}{int} \PYG{n}{n\PYGZus{}dep} \PYG{o}{=} \PYG{l+m+mi}{1}\PYG{p}{;}
    \PYG{k+kt}{int} \PYG{n}{n\PYGZus{}ind} \PYG{o}{=} \PYG{l+m+mi}{3}\PYG{p}{;}
    \PYG{c+c1}{//}
    \PYG{c+c1}{// create the independent variables ax}
    \PYG{n}{vec\PYGZus{}double} \PYG{n}{x}\PYG{p}{(}\PYG{n}{n\PYGZus{}ind}\PYG{p}{)}\PYG{p}{;}
    \PYG{k}{for}\PYG{p}{(}\PYG{k+kt}{int} \PYG{n}{i} \PYG{o}{=} \PYG{l+m+mi}{0}\PYG{p}{;} \PYG{n}{i} \PYG{o}{\PYGZlt{}} \PYG{n}{n\PYGZus{}ind} \PYG{p}{;} \PYG{n}{i}\PYG{o}{+}\PYG{o}{+}\PYG{p}{)} \PYG{p}{\PYGZob{}}
        \PYG{n}{x}\PYG{p}{[}\PYG{n}{i}\PYG{p}{]} \PYG{o}{=} \PYG{n}{i} \PYG{o}{+} \PYG{l+m+mf}{2.0}\PYG{p}{;}
    \PYG{p}{\PYGZcb{}}
    \PYG{n}{vec\PYGZus{}a\PYGZus{}double} \PYG{n}{ax} \PYG{o}{=} \PYG{n}{cppad\PYGZus{}py}\PYG{o}{:}\PYG{o}{:}\PYG{n}{independent}\PYG{p}{(}\PYG{n}{x}\PYG{p}{)}\PYG{p}{;}
    \PYG{c+c1}{//}
    \PYG{c+c1}{// create dependent variables ay with ay0 = ax\PYGZus{}0 * ax\PYGZus{}1 * ax\PYGZus{}2}
    \PYG{n}{a\PYGZus{}double} \PYG{n}{ax\PYGZus{}0} \PYG{o}{=} \PYG{n}{ax}\PYG{p}{[}\PYG{l+m+mi}{0}\PYG{p}{]}\PYG{p}{;}
    \PYG{n}{a\PYGZus{}double} \PYG{n}{ax\PYGZus{}1} \PYG{o}{=} \PYG{n}{ax}\PYG{p}{[}\PYG{l+m+mi}{1}\PYG{p}{]}\PYG{p}{;}
    \PYG{n}{a\PYGZus{}double} \PYG{n}{ax\PYGZus{}2} \PYG{o}{=} \PYG{n}{ax}\PYG{p}{[}\PYG{l+m+mi}{2}\PYG{p}{]}\PYG{p}{;}
    \PYG{n}{vec\PYGZus{}a\PYGZus{}double} \PYG{n}{ay}\PYG{p}{(}\PYG{n}{n\PYGZus{}dep}\PYG{p}{)}\PYG{p}{;}
    \PYG{n}{ay}\PYG{p}{[}\PYG{l+m+mi}{0}\PYG{p}{]} \PYG{o}{=} \PYG{n}{ax\PYGZus{}0} \PYG{o}{*} \PYG{n}{ax\PYGZus{}1} \PYG{o}{*} \PYG{n}{ax\PYGZus{}2}\PYG{p}{;}
    \PYG{c+c1}{//}
    \PYG{c+c1}{// define af corresponding to f(x) = x\PYGZus{}0 * x\PYGZus{}1 * x\PYGZus{}2}
    \PYG{n}{d\PYGZus{}fun} \PYG{n}{f}\PYG{p}{(}\PYG{n}{ax}\PYG{p}{,} \PYG{n}{ay}\PYG{p}{)}\PYG{p}{;}
    \PYG{c+c1}{//}
    \PYG{c+c1}{// compute the Jacobian f\PYGZsq{}(x) = ( x\PYGZus{}1*x\PYGZus{}2, x\PYGZus{}0*x\PYGZus{}2, x\PYGZus{}0*x\PYGZus{}1 )}
    \PYG{n}{vec\PYGZus{}double} \PYG{n}{fp} \PYG{o}{=} \PYG{n}{f}\PYG{p}{.}\PYG{n}{jacobian}\PYG{p}{(}\PYG{n}{x}\PYG{p}{)}\PYG{p}{;}
    \PYG{c+c1}{//}
    \PYG{c+c1}{// check Jacobian}
    \PYG{k+kt}{double} \PYG{n}{x\PYGZus{}0} \PYG{o}{=} \PYG{n}{x}\PYG{p}{[}\PYG{l+m+mi}{0}\PYG{p}{]}\PYG{p}{;}
    \PYG{k+kt}{double} \PYG{n}{x\PYGZus{}1} \PYG{o}{=} \PYG{n}{x}\PYG{p}{[}\PYG{l+m+mi}{1}\PYG{p}{]}\PYG{p}{;}
    \PYG{k+kt}{double} \PYG{n}{x\PYGZus{}2} \PYG{o}{=} \PYG{n}{x}\PYG{p}{[}\PYG{l+m+mi}{2}\PYG{p}{]}\PYG{p}{;}
    \PYG{n}{ok} \PYG{o}{=} \PYG{n}{ok} \PYG{o}{\PYGZam{}}\PYG{o}{\PYGZam{}} \PYG{n}{fp}\PYG{p}{[}\PYG{l+m+mi}{0} \PYG{o}{*} \PYG{n}{n\PYGZus{}ind} \PYG{o}{+} \PYG{l+m+mi}{0}\PYG{p}{]} \PYG{o}{=}\PYG{o}{=} \PYG{n}{x\PYGZus{}1} \PYG{o}{*} \PYG{n}{x\PYGZus{}2} \PYG{p}{;}
    \PYG{n}{ok} \PYG{o}{=} \PYG{n}{ok} \PYG{o}{\PYGZam{}}\PYG{o}{\PYGZam{}} \PYG{n}{fp}\PYG{p}{[}\PYG{l+m+mi}{0} \PYG{o}{*} \PYG{n}{n\PYGZus{}ind} \PYG{o}{+} \PYG{l+m+mi}{1}\PYG{p}{]} \PYG{o}{=}\PYG{o}{=} \PYG{n}{x\PYGZus{}0} \PYG{o}{*} \PYG{n}{x\PYGZus{}2} \PYG{p}{;}
    \PYG{n}{ok} \PYG{o}{=} \PYG{n}{ok} \PYG{o}{\PYGZam{}}\PYG{o}{\PYGZam{}} \PYG{n}{fp}\PYG{p}{[}\PYG{l+m+mi}{0} \PYG{o}{*} \PYG{n}{n\PYGZus{}ind} \PYG{o}{+} \PYG{l+m+mi}{2}\PYG{p}{]} \PYG{o}{=}\PYG{o}{=} \PYG{n}{x\PYGZus{}0} \PYG{o}{*} \PYG{n}{x\PYGZus{}1} \PYG{p}{;}
    \PYG{c+c1}{//\PYGZhy{}\PYGZhy{}\PYGZhy{}\PYGZhy{}\PYGZhy{}\PYGZhy{}\PYGZhy{}\PYGZhy{}\PYGZhy{}\PYGZhy{}\PYGZhy{}\PYGZhy{}\PYGZhy{}\PYGZhy{}\PYGZhy{}\PYGZhy{}\PYGZhy{}\PYGZhy{}\PYGZhy{}\PYGZhy{}\PYGZhy{}\PYGZhy{}\PYGZhy{}\PYGZhy{}\PYGZhy{}\PYGZhy{}\PYGZhy{}\PYGZhy{}\PYGZhy{}\PYGZhy{}\PYGZhy{}\PYGZhy{}\PYGZhy{}\PYGZhy{}\PYGZhy{}\PYGZhy{}\PYGZhy{}\PYGZhy{}\PYGZhy{}\PYGZhy{}\PYGZhy{}\PYGZhy{}\PYGZhy{}\PYGZhy{}\PYGZhy{}\PYGZhy{}\PYGZhy{}\PYGZhy{}\PYGZhy{}\PYGZhy{}\PYGZhy{}\PYGZhy{}\PYGZhy{}\PYGZhy{}\PYGZhy{}\PYGZhy{}\PYGZhy{}\PYGZhy{}\PYGZhy{}\PYGZhy{}\PYGZhy{}\PYGZhy{}\PYGZhy{}\PYGZhy{}\PYGZhy{}\PYGZhy{}\PYGZhy{}\PYGZhy{}\PYGZhy{}\PYGZhy{}\PYGZhy{}\PYGZhy{}}
    \PYG{n}{a\PYGZus{}fun} \PYG{n}{af}\PYG{p}{(}\PYG{n}{f}\PYG{p}{)}\PYG{p}{;}
    \PYG{c+c1}{//}
    \PYG{c+c1}{// compute the Jacobian f\PYGZsq{}(x) = ( x\PYGZus{}1*x\PYGZus{}2, x\PYGZus{}0*x\PYGZus{}2, x\PYGZus{}0*x\PYGZus{}1 )}
    \PYG{n}{vec\PYGZus{}a\PYGZus{}double} \PYG{n}{afp} \PYG{o}{=} \PYG{n}{af}\PYG{p}{.}\PYG{n}{jacobian}\PYG{p}{(}\PYG{n}{ax}\PYG{p}{)}\PYG{p}{;}
    \PYG{c+c1}{//}
    \PYG{c+c1}{// check Jacobian}
    \PYG{n}{ok} \PYG{o}{=} \PYG{n}{ok} \PYG{o}{\PYGZam{}}\PYG{o}{\PYGZam{}} \PYG{n}{afp}\PYG{p}{[}\PYG{l+m+mi}{0} \PYG{o}{*} \PYG{n}{n\PYGZus{}ind} \PYG{o}{+} \PYG{l+m+mi}{0}\PYG{p}{]} \PYG{o}{=}\PYG{o}{=} \PYG{n}{x\PYGZus{}1} \PYG{o}{*} \PYG{n}{x\PYGZus{}2} \PYG{p}{;}
    \PYG{n}{ok} \PYG{o}{=} \PYG{n}{ok} \PYG{o}{\PYGZam{}}\PYG{o}{\PYGZam{}} \PYG{n}{afp}\PYG{p}{[}\PYG{l+m+mi}{0} \PYG{o}{*} \PYG{n}{n\PYGZus{}ind} \PYG{o}{+} \PYG{l+m+mi}{1}\PYG{p}{]} \PYG{o}{=}\PYG{o}{=} \PYG{n}{x\PYGZus{}0} \PYG{o}{*} \PYG{n}{x\PYGZus{}2} \PYG{p}{;}
    \PYG{n}{ok} \PYG{o}{=} \PYG{n}{ok} \PYG{o}{\PYGZam{}}\PYG{o}{\PYGZam{}} \PYG{n}{afp}\PYG{p}{[}\PYG{l+m+mi}{0} \PYG{o}{*} \PYG{n}{n\PYGZus{}ind} \PYG{o}{+} \PYG{l+m+mi}{2}\PYG{p}{]} \PYG{o}{=}\PYG{o}{=} \PYG{n}{x\PYGZus{}0} \PYG{o}{*} \PYG{n}{x\PYGZus{}1} \PYG{p}{;}
    \PYG{c+c1}{//}
    \PYG{k}{return}\PYG{p}{(} \PYG{n}{ok} \PYG{p}{)}\PYG{p}{;}
\PYG{p}{\PYGZcb{}}
\end{sphinxVerbatim}


\bigskip\hrule\bigskip


xsrst input file: \sphinxcode{\sphinxupquote{example/cplusplus/fun\_jacobian\_xam.cpp}}

\index{example@\spxentry{example}}\ignorespaces 

\subparagraph{Example}
\label{\detokenize{xsrst/cpp_fun_jacobian:example}}\label{\detokenize{xsrst/cpp_fun_jacobian:cpp-fun-jacobian-example}}\label{\detokenize{xsrst/cpp_fun_jacobian:index-6}}
{\hyperref[\detokenize{xsrst/fun_jacobian_xam_cpp:id1}]{\sphinxcrossref{\DUrole{std,std-ref}{fun\_jacobian\_xam\_cpp}}}}


\bigskip\hrule\bigskip


xsrst input file: \sphinxcode{\sphinxupquote{lib/cplusplus/fun.cpp}}


\subparagraph{cpp\_fun\_hessian}
\label{\detokenize{xsrst/cpp_fun_hessian:cpp-fun-hessian}}\label{\detokenize{xsrst/cpp_fun_hessian::doc}}
\index{cpp\_fun\_hessian@\spxentry{cpp\_fun\_hessian}}\index{hessian@\spxentry{hessian}}\index{an@\spxentry{an}}\index{ad@\spxentry{ad}}\index{function@\spxentry{function}}\ignorespaces 

\subparagraph{3.2.3.7 Hessian of an AD Function}
\label{\detokenize{xsrst/cpp_fun_hessian:hessian-of-an-ad-function}}\label{\detokenize{xsrst/cpp_fun_hessian:id1}}\begin{itemize}
\item {} 
{\hyperref[\detokenize{xsrst/cpp_fun_hessian:cpp-fun-hessian-syntax}]{\sphinxcrossref{\DUrole{std,std-ref}{Syntax}}}}

\item {} 
{\hyperref[\detokenize{xsrst/cpp_fun_hessian:cpp-fun-hessian-f}]{\sphinxcrossref{\DUrole{std,std-ref}{f}}}}

\item {} 
{\hyperref[\detokenize{xsrst/cpp_fun_hessian:cpp-fun-hessian-f-x}]{\sphinxcrossref{\DUrole{std,std-ref}{f(x)}}}}

\item {} 
{\hyperref[\detokenize{xsrst/cpp_fun_hessian:cpp-fun-hessian-g-x}]{\sphinxcrossref{\DUrole{std,std-ref}{g(x)}}}}

\item {} 
{\hyperref[\detokenize{xsrst/cpp_fun_hessian:cpp-fun-hessian-x}]{\sphinxcrossref{\DUrole{std,std-ref}{x}}}}

\item {} 
{\hyperref[\detokenize{xsrst/cpp_fun_hessian:cpp-fun-hessian-w}]{\sphinxcrossref{\DUrole{std,std-ref}{w}}}}

\item {} 
{\hyperref[\detokenize{xsrst/cpp_fun_hessian:cpp-fun-hessian-h}]{\sphinxcrossref{\DUrole{std,std-ref}{H}}}}

\item {} 
{\hyperref[\detokenize{xsrst/cpp_fun_hessian:cpp-fun-hessian-example}]{\sphinxcrossref{\DUrole{std,std-ref}{Example}}}}

\end{itemize}

\index{syntax@\spxentry{syntax}}\ignorespaces 

\subparagraph{Syntax}
\label{\detokenize{xsrst/cpp_fun_hessian:syntax}}\label{\detokenize{xsrst/cpp_fun_hessian:cpp-fun-hessian-syntax}}\label{\detokenize{xsrst/cpp_fun_hessian:index-1}}
\sphinxstyleemphasis{H} = \sphinxstyleemphasis{f} . \sphinxcode{\sphinxupquote{hessian}} ( \sphinxstyleemphasis{x} , \sphinxstyleemphasis{w} )

\index{f@\spxentry{f}}\ignorespaces 

\subparagraph{f}
\label{\detokenize{xsrst/cpp_fun_hessian:f}}\label{\detokenize{xsrst/cpp_fun_hessian:cpp-fun-hessian-f}}\label{\detokenize{xsrst/cpp_fun_hessian:index-2}}
This is either a
{\hyperref[\detokenize{xsrst/cpp_fun_ctor:cpp-fun-ctor-syntax-d-fun}]{\sphinxcrossref{\DUrole{std,std-ref}{d\_fun}}}} or
{\hyperref[\detokenize{xsrst/cpp_fun_ctor:cpp-fun-ctor-syntax-a-fun}]{\sphinxcrossref{\DUrole{std,std-ref}{a\_fun}}}} function object.
Upon return, the zero order
{\hyperref[\detokenize{xsrst/cpp_fun_forward:cpp-fun-forward-taylor-coefficient}]{\sphinxcrossref{\DUrole{std,std-ref}{taylor\_coefficients}}}} in \sphinxstyleemphasis{f}
correspond to the value of \sphinxstyleemphasis{x} .
The other Taylor coefficients in \sphinxstyleemphasis{f} are unspecified.

\index{f(x)@\spxentry{f(x)}}\ignorespaces 

\subparagraph{f(x)}
\label{\detokenize{xsrst/cpp_fun_hessian:f-x}}\label{\detokenize{xsrst/cpp_fun_hessian:cpp-fun-hessian-f-x}}\label{\detokenize{xsrst/cpp_fun_hessian:index-3}}
We use the notation \(f: \B{R}^n \rightarrow \B{R}^m\)
for the function corresponding to \sphinxstyleemphasis{f} .
Note that \sphinxstyleemphasis{n} is the size of {\hyperref[\detokenize{xsrst/cpp_fun_ctor:cpp-fun-ctor-ax}]{\sphinxcrossref{\DUrole{std,std-ref}{ax}}}}
and \sphinxstyleemphasis{m} is the size of {\hyperref[\detokenize{xsrst/cpp_fun_ctor:cpp-fun-ctor-ay}]{\sphinxcrossref{\DUrole{std,std-ref}{ay}}}}
in to the constructor for \sphinxstyleemphasis{f} .

\index{g(x)@\spxentry{g(x)}}\ignorespaces 

\subparagraph{g(x)}
\label{\detokenize{xsrst/cpp_fun_hessian:g-x}}\label{\detokenize{xsrst/cpp_fun_hessian:cpp-fun-hessian-g-x}}\label{\detokenize{xsrst/cpp_fun_hessian:index-4}}
We use the notation \(g: \B{R}^n \rightarrow \B{R}\)
for the function defined by
\begin{equation*}
\begin{split}g(x) = \sum_{i=0}^{n-1} w_i f_i (x)\end{split}
\end{equation*}
\index{x@\spxentry{x}}\ignorespaces 

\subparagraph{x}
\label{\detokenize{xsrst/cpp_fun_hessian:x}}\label{\detokenize{xsrst/cpp_fun_hessian:cpp-fun-hessian-x}}\label{\detokenize{xsrst/cpp_fun_hessian:index-5}}
If \sphinxstyleemphasis{f} is a \sphinxcode{\sphinxupquote{d\_fun}} or \sphinxcode{\sphinxupquote{a\_fun}} ,
this argument has prototype

\begin{DUlineblock}{0em}
\item[]         \sphinxcode{\sphinxupquote{const vec\_double\&}} \sphinxstyleemphasis{x}
\item[]         \sphinxcode{\sphinxupquote{const vec\_a\_double\&}} \sphinxstyleemphasis{x}
\end{DUlineblock}

and its size must be \sphinxstyleemphasis{n} .
It specifies the argument value at we are computing the Hessian
\(g^{(2)}(x)\).

\index{w@\spxentry{w}}\ignorespaces 

\subparagraph{w}
\label{\detokenize{xsrst/cpp_fun_hessian:w}}\label{\detokenize{xsrst/cpp_fun_hessian:cpp-fun-hessian-w}}\label{\detokenize{xsrst/cpp_fun_hessian:index-6}}
If \sphinxstyleemphasis{f} is a \sphinxcode{\sphinxupquote{d\_fun}} or \sphinxcode{\sphinxupquote{a\_fun}} ,
this argument has prototype

\begin{DUlineblock}{0em}
\item[]         \sphinxcode{\sphinxupquote{const vec\_double\&}} \sphinxstyleemphasis{w}
\item[]         \sphinxcode{\sphinxupquote{const vec\_a\_double\&}} \sphinxstyleemphasis{w}
\end{DUlineblock}

and its size must be \sphinxstyleemphasis{m} .
It specifies the vector \sphinxstyleemphasis{w} in the definition of \(g(x)\) above.

\index{h@\spxentry{h}}\ignorespaces 

\subparagraph{H}
\label{\detokenize{xsrst/cpp_fun_hessian:h}}\label{\detokenize{xsrst/cpp_fun_hessian:cpp-fun-hessian-h}}\label{\detokenize{xsrst/cpp_fun_hessian:index-7}}
If \sphinxstyleemphasis{f} is a \sphinxcode{\sphinxupquote{d\_fun}} or \sphinxcode{\sphinxupquote{a\_fun}} ,
the result has prototype

\begin{DUlineblock}{0em}
\item[]         \sphinxcode{\sphinxupquote{vec\_double}} \sphinxstyleemphasis{H}
\item[]         \sphinxcode{\sphinxupquote{vec\_a\_double}} \sphinxstyleemphasis{H}
\end{DUlineblock}

and its size is \sphinxstyleemphasis{n} \sphinxcode{\sphinxupquote{*}} \sphinxstyleemphasis{n} .
For \sphinxstyleemphasis{i} between zero and \sphinxstyleemphasis{n} \sphinxhyphen{}1
and \sphinxstyleemphasis{j} between zero and \sphinxstyleemphasis{n} \sphinxhyphen{}1 ,
\begin{equation*}
\begin{split}H [ i * n + j ] = \frac{ \partial^2 g }{ \partial x_i \partial x_j } (x)\end{split}
\end{equation*}

\subparagraph{fun\_hessian\_xam\_cpp}
\label{\detokenize{xsrst/fun_hessian_xam_cpp:fun-hessian-xam-cpp}}\label{\detokenize{xsrst/fun_hessian_xam_cpp::doc}}
\index{fun\_hessian\_xam\_cpp@\spxentry{fun\_hessian\_xam\_cpp}}\index{c++:@\spxentry{c++:}}\index{dense@\spxentry{dense}}\index{hessian@\spxentry{hessian}}\index{using@\spxentry{using}}\index{ad:@\spxentry{ad:}}\index{example@\spxentry{example}}\index{test@\spxentry{test}}\ignorespaces 

\subparagraph{3.2.3.7.1 C++: Dense Hessian Using AD: Example and Test}
\label{\detokenize{xsrst/fun_hessian_xam_cpp:c-dense-hessian-using-ad-example-and-test}}\label{\detokenize{xsrst/fun_hessian_xam_cpp:id1}}
\begin{sphinxVerbatim}[commandchars=\\\{\}]
\PYG{c+cp}{\PYGZsh{}}\PYG{c+cp}{ include \PYGZlt{}cstdio\PYGZgt{}}
\PYG{c+cp}{\PYGZsh{}}\PYG{c+cp}{ include \PYGZlt{}cppad}\PYG{c+cp}{/}\PYG{c+cp}{py}\PYG{c+cp}{/}\PYG{c+cp}{cppad\PYGZus{}py.hpp\PYGZgt{}}

\PYG{k+kt}{bool} \PYG{n+nf}{fun\PYGZus{}hessian\PYGZus{}xam}\PYG{p}{(}\PYG{k+kt}{void}\PYG{p}{)} \PYG{p}{\PYGZob{}}
    \PYG{k}{using} \PYG{n}{cppad\PYGZus{}py}\PYG{o}{:}\PYG{o}{:}\PYG{n}{a\PYGZus{}double}\PYG{p}{;}
    \PYG{k}{using} \PYG{n}{cppad\PYGZus{}py}\PYG{o}{:}\PYG{o}{:}\PYG{n}{vec\PYGZus{}double}\PYG{p}{;}
    \PYG{k}{using} \PYG{n}{cppad\PYGZus{}py}\PYG{o}{:}\PYG{o}{:}\PYG{n}{vec\PYGZus{}a\PYGZus{}double}\PYG{p}{;}
    \PYG{k}{using} \PYG{n}{cppad\PYGZus{}py}\PYG{o}{:}\PYG{o}{:}\PYG{n}{d\PYGZus{}fun}\PYG{p}{;}
    \PYG{k}{using} \PYG{n}{cppad\PYGZus{}py}\PYG{o}{:}\PYG{o}{:}\PYG{n}{a\PYGZus{}fun}\PYG{p}{;}
    \PYG{c+c1}{//}
    \PYG{c+c1}{// initialize return variable}
    \PYG{k+kt}{bool} \PYG{n}{ok} \PYG{o}{=} \PYG{n+nb}{true}\PYG{p}{;}
    \PYG{c+c1}{// \PYGZhy{}\PYGZhy{}\PYGZhy{}\PYGZhy{}\PYGZhy{}\PYGZhy{}\PYGZhy{}\PYGZhy{}\PYGZhy{}\PYGZhy{}\PYGZhy{}\PYGZhy{}\PYGZhy{}\PYGZhy{}\PYGZhy{}\PYGZhy{}\PYGZhy{}\PYGZhy{}\PYGZhy{}\PYGZhy{}\PYGZhy{}\PYGZhy{}\PYGZhy{}\PYGZhy{}\PYGZhy{}\PYGZhy{}\PYGZhy{}\PYGZhy{}\PYGZhy{}\PYGZhy{}\PYGZhy{}\PYGZhy{}\PYGZhy{}\PYGZhy{}\PYGZhy{}\PYGZhy{}\PYGZhy{}\PYGZhy{}\PYGZhy{}\PYGZhy{}\PYGZhy{}\PYGZhy{}\PYGZhy{}\PYGZhy{}\PYGZhy{}\PYGZhy{}\PYGZhy{}\PYGZhy{}\PYGZhy{}\PYGZhy{}\PYGZhy{}\PYGZhy{}\PYGZhy{}\PYGZhy{}\PYGZhy{}\PYGZhy{}\PYGZhy{}\PYGZhy{}\PYGZhy{}\PYGZhy{}\PYGZhy{}\PYGZhy{}\PYGZhy{}\PYGZhy{}\PYGZhy{}\PYGZhy{}\PYGZhy{}\PYGZhy{}\PYGZhy{}\PYGZhy{}\PYGZhy{}}
    \PYG{c+c1}{// number of dependent and independent variables}
    \PYG{k+kt}{int} \PYG{n}{n\PYGZus{}dep} \PYG{o}{=} \PYG{l+m+mi}{1}\PYG{p}{;}
    \PYG{k+kt}{int} \PYG{n}{n\PYGZus{}ind} \PYG{o}{=} \PYG{l+m+mi}{3}\PYG{p}{;}
    \PYG{c+c1}{//}
    \PYG{c+c1}{// create the independent variables ax}
    \PYG{n}{vec\PYGZus{}double} \PYG{n}{x}\PYG{p}{(}\PYG{n}{n\PYGZus{}ind}\PYG{p}{)}\PYG{p}{;}
    \PYG{k}{for}\PYG{p}{(}\PYG{k+kt}{int} \PYG{n}{i} \PYG{o}{=} \PYG{l+m+mi}{0}\PYG{p}{;} \PYG{n}{i} \PYG{o}{\PYGZlt{}} \PYG{n}{n\PYGZus{}ind} \PYG{p}{;} \PYG{n}{i}\PYG{o}{+}\PYG{o}{+}\PYG{p}{)} \PYG{p}{\PYGZob{}}
        \PYG{n}{x}\PYG{p}{[}\PYG{n}{i}\PYG{p}{]} \PYG{o}{=} \PYG{n}{i} \PYG{o}{+} \PYG{l+m+mf}{2.0}\PYG{p}{;}
    \PYG{p}{\PYGZcb{}}
    \PYG{n}{vec\PYGZus{}a\PYGZus{}double} \PYG{n}{ax} \PYG{o}{=} \PYG{n}{cppad\PYGZus{}py}\PYG{o}{:}\PYG{o}{:}\PYG{n}{independent}\PYG{p}{(}\PYG{n}{x}\PYG{p}{)}\PYG{p}{;}
    \PYG{c+c1}{//}
    \PYG{c+c1}{// create dependent variables ay with ay0 = ax\PYGZus{}0 * ax\PYGZus{}1 * ax\PYGZus{}2}
    \PYG{n}{a\PYGZus{}double} \PYG{n}{ax\PYGZus{}0} \PYG{o}{=} \PYG{n}{ax}\PYG{p}{[}\PYG{l+m+mi}{0}\PYG{p}{]}\PYG{p}{;}
    \PYG{n}{a\PYGZus{}double} \PYG{n}{ax\PYGZus{}1} \PYG{o}{=} \PYG{n}{ax}\PYG{p}{[}\PYG{l+m+mi}{1}\PYG{p}{]}\PYG{p}{;}
    \PYG{n}{a\PYGZus{}double} \PYG{n}{ax\PYGZus{}2} \PYG{o}{=} \PYG{n}{ax}\PYG{p}{[}\PYG{l+m+mi}{2}\PYG{p}{]}\PYG{p}{;}
    \PYG{n}{vec\PYGZus{}a\PYGZus{}double} \PYG{n}{ay}\PYG{p}{(}\PYG{n}{n\PYGZus{}dep}\PYG{p}{)}\PYG{p}{;}
    \PYG{n}{ay}\PYG{p}{[}\PYG{l+m+mi}{0}\PYG{p}{]} \PYG{o}{=} \PYG{n}{ax\PYGZus{}0} \PYG{o}{*} \PYG{n}{ax\PYGZus{}1} \PYG{o}{*} \PYG{n}{ax\PYGZus{}2}\PYG{p}{;}
    \PYG{c+c1}{//}
    \PYG{c+c1}{// define af corresponding to f(x) = x\PYGZus{}0 * x\PYGZus{}1 * x\PYGZus{}2}
    \PYG{n}{d\PYGZus{}fun} \PYG{n}{f}\PYG{p}{(}\PYG{n}{ax}\PYG{p}{,} \PYG{n}{ay}\PYG{p}{)}\PYG{p}{;}
    \PYG{c+c1}{//}
    \PYG{c+c1}{// g(x) = w\PYGZus{}0 * f\PYGZus{}0 (x) = f(x)}
    \PYG{n}{vec\PYGZus{}double} \PYG{n}{w}\PYG{p}{(}\PYG{n}{n\PYGZus{}dep}\PYG{p}{)}\PYG{p}{;}
    \PYG{n}{w}\PYG{p}{[}\PYG{l+m+mi}{0}\PYG{p}{]} \PYG{o}{=} \PYG{l+m+mf}{1.0}\PYG{p}{;}
    \PYG{c+c1}{//}
    \PYG{c+c1}{// compute Hessian}
    \PYG{n}{vec\PYGZus{}double} \PYG{n}{fpp} \PYG{o}{=} \PYG{n}{f}\PYG{p}{.}\PYG{n}{hessian}\PYG{p}{(}\PYG{n}{x}\PYG{p}{,} \PYG{n}{w}\PYG{p}{)}\PYG{p}{;}
    \PYG{c+c1}{//}
    \PYG{c+c1}{//          [ 0.0 , x\PYGZus{}2 , x\PYGZus{}1 ]}
    \PYG{c+c1}{// f\PYGZsq{}\PYGZsq{}(x) = [ x\PYGZus{}2 , 0.0 , x\PYGZus{}0 ]}
    \PYG{c+c1}{//          [ x\PYGZus{}1 , x\PYGZus{}0 , 0.0 ]}
    \PYG{n}{ok} \PYG{o}{=} \PYG{n}{ok} \PYG{o}{\PYGZam{}}\PYG{o}{\PYGZam{}} \PYG{n}{fpp}\PYG{p}{[}\PYG{l+m+mi}{0} \PYG{o}{*} \PYG{n}{n\PYGZus{}ind} \PYG{o}{+} \PYG{l+m+mi}{0}\PYG{p}{]} \PYG{o}{=}\PYG{o}{=} \PYG{l+m+mf}{0.0} \PYG{p}{;}
    \PYG{n}{ok} \PYG{o}{=} \PYG{n}{ok} \PYG{o}{\PYGZam{}}\PYG{o}{\PYGZam{}} \PYG{n}{fpp}\PYG{p}{[}\PYG{l+m+mi}{0} \PYG{o}{*} \PYG{n}{n\PYGZus{}ind} \PYG{o}{+} \PYG{l+m+mi}{1}\PYG{p}{]} \PYG{o}{=}\PYG{o}{=} \PYG{n}{x}\PYG{p}{[}\PYG{l+m+mi}{2}\PYG{p}{]} \PYG{p}{;}
    \PYG{n}{ok} \PYG{o}{=} \PYG{n}{ok} \PYG{o}{\PYGZam{}}\PYG{o}{\PYGZam{}} \PYG{n}{fpp}\PYG{p}{[}\PYG{l+m+mi}{0} \PYG{o}{*} \PYG{n}{n\PYGZus{}ind} \PYG{o}{+} \PYG{l+m+mi}{2}\PYG{p}{]} \PYG{o}{=}\PYG{o}{=} \PYG{n}{x}\PYG{p}{[}\PYG{l+m+mi}{1}\PYG{p}{]} \PYG{p}{;}
    \PYG{c+c1}{//}
    \PYG{n}{ok} \PYG{o}{=} \PYG{n}{ok} \PYG{o}{\PYGZam{}}\PYG{o}{\PYGZam{}} \PYG{n}{fpp}\PYG{p}{[}\PYG{l+m+mi}{1} \PYG{o}{*} \PYG{n}{n\PYGZus{}ind} \PYG{o}{+} \PYG{l+m+mi}{0}\PYG{p}{]} \PYG{o}{=}\PYG{o}{=} \PYG{n}{x}\PYG{p}{[}\PYG{l+m+mi}{2}\PYG{p}{]} \PYG{p}{;}
    \PYG{n}{ok} \PYG{o}{=} \PYG{n}{ok} \PYG{o}{\PYGZam{}}\PYG{o}{\PYGZam{}} \PYG{n}{fpp}\PYG{p}{[}\PYG{l+m+mi}{1} \PYG{o}{*} \PYG{n}{n\PYGZus{}ind} \PYG{o}{+} \PYG{l+m+mi}{1}\PYG{p}{]} \PYG{o}{=}\PYG{o}{=} \PYG{l+m+mf}{0.0} \PYG{p}{;}
    \PYG{n}{ok} \PYG{o}{=} \PYG{n}{ok} \PYG{o}{\PYGZam{}}\PYG{o}{\PYGZam{}} \PYG{n}{fpp}\PYG{p}{[}\PYG{l+m+mi}{1} \PYG{o}{*} \PYG{n}{n\PYGZus{}ind} \PYG{o}{+} \PYG{l+m+mi}{2}\PYG{p}{]} \PYG{o}{=}\PYG{o}{=} \PYG{n}{x}\PYG{p}{[}\PYG{l+m+mi}{0}\PYG{p}{]} \PYG{p}{;}
    \PYG{c+c1}{//}
    \PYG{n}{ok} \PYG{o}{=} \PYG{n}{ok} \PYG{o}{\PYGZam{}}\PYG{o}{\PYGZam{}} \PYG{n}{fpp}\PYG{p}{[}\PYG{l+m+mi}{2} \PYG{o}{*} \PYG{n}{n\PYGZus{}ind} \PYG{o}{+} \PYG{l+m+mi}{0}\PYG{p}{]} \PYG{o}{=}\PYG{o}{=} \PYG{n}{x}\PYG{p}{[}\PYG{l+m+mi}{1}\PYG{p}{]} \PYG{p}{;}
    \PYG{n}{ok} \PYG{o}{=} \PYG{n}{ok} \PYG{o}{\PYGZam{}}\PYG{o}{\PYGZam{}} \PYG{n}{fpp}\PYG{p}{[}\PYG{l+m+mi}{2} \PYG{o}{*} \PYG{n}{n\PYGZus{}ind} \PYG{o}{+} \PYG{l+m+mi}{1}\PYG{p}{]} \PYG{o}{=}\PYG{o}{=} \PYG{n}{x}\PYG{p}{[}\PYG{l+m+mi}{0}\PYG{p}{]} \PYG{p}{;}
    \PYG{n}{ok} \PYG{o}{=} \PYG{n}{ok} \PYG{o}{\PYGZam{}}\PYG{o}{\PYGZam{}} \PYG{n}{fpp}\PYG{p}{[}\PYG{l+m+mi}{2} \PYG{o}{*} \PYG{n}{n\PYGZus{}ind} \PYG{o}{+} \PYG{l+m+mi}{2}\PYG{p}{]} \PYG{o}{=}\PYG{o}{=} \PYG{l+m+mf}{0.0} \PYG{p}{;}
    \PYG{c+c1}{// \PYGZhy{}\PYGZhy{}\PYGZhy{}\PYGZhy{}\PYGZhy{}\PYGZhy{}\PYGZhy{}\PYGZhy{}\PYGZhy{}\PYGZhy{}\PYGZhy{}\PYGZhy{}\PYGZhy{}\PYGZhy{}\PYGZhy{}\PYGZhy{}\PYGZhy{}\PYGZhy{}\PYGZhy{}\PYGZhy{}\PYGZhy{}\PYGZhy{}\PYGZhy{}\PYGZhy{}\PYGZhy{}\PYGZhy{}\PYGZhy{}\PYGZhy{}\PYGZhy{}\PYGZhy{}\PYGZhy{}\PYGZhy{}\PYGZhy{}\PYGZhy{}\PYGZhy{}\PYGZhy{}\PYGZhy{}\PYGZhy{}\PYGZhy{}\PYGZhy{}\PYGZhy{}\PYGZhy{}\PYGZhy{}\PYGZhy{}\PYGZhy{}\PYGZhy{}\PYGZhy{}\PYGZhy{}\PYGZhy{}\PYGZhy{}\PYGZhy{}\PYGZhy{}\PYGZhy{}\PYGZhy{}\PYGZhy{}\PYGZhy{}\PYGZhy{}\PYGZhy{}\PYGZhy{}\PYGZhy{}\PYGZhy{}\PYGZhy{}\PYGZhy{}\PYGZhy{}\PYGZhy{}\PYGZhy{}\PYGZhy{}\PYGZhy{}\PYGZhy{}\PYGZhy{}\PYGZhy{}}
    \PYG{n}{a\PYGZus{}fun} \PYG{n}{af}\PYG{p}{(}\PYG{n}{f}\PYG{p}{)}\PYG{p}{;}
    \PYG{c+c1}{//}
    \PYG{c+c1}{// compute and check Hessian}
    \PYG{n}{vec\PYGZus{}a\PYGZus{}double} \PYG{n}{aw}\PYG{p}{(}\PYG{n}{n\PYGZus{}dep}\PYG{p}{)}\PYG{p}{;}
    \PYG{n}{aw}\PYG{p}{[}\PYG{l+m+mi}{0}\PYG{p}{]} \PYG{o}{=} \PYG{n}{w}\PYG{p}{[}\PYG{l+m+mi}{0}\PYG{p}{]}\PYG{p}{;}
    \PYG{n}{vec\PYGZus{}a\PYGZus{}double} \PYG{n}{afpp} \PYG{o}{=} \PYG{n}{af}\PYG{p}{.}\PYG{n}{hessian}\PYG{p}{(}\PYG{n}{ax}\PYG{p}{,} \PYG{n}{aw}\PYG{p}{)}\PYG{p}{;}
    \PYG{n}{ok} \PYG{o}{=} \PYG{n}{ok} \PYG{o}{\PYGZam{}}\PYG{o}{\PYGZam{}} \PYG{n}{afpp}\PYG{p}{.}\PYG{n}{size}\PYG{p}{(}\PYG{p}{)} \PYG{o}{=}\PYG{o}{=} \PYG{n}{fpp}\PYG{p}{.}\PYG{n}{size}\PYG{p}{(}\PYG{p}{)}\PYG{p}{;}
    \PYG{k}{for}\PYG{p}{(}\PYG{k+kt}{size\PYGZus{}t} \PYG{n}{i} \PYG{o}{=} \PYG{l+m+mi}{0}\PYG{p}{;} \PYG{n}{i} \PYG{o}{\PYGZlt{}} \PYG{n}{fpp}\PYG{p}{.}\PYG{n}{size}\PYG{p}{(}\PYG{p}{)}\PYG{p}{;} \PYG{o}{+}\PYG{o}{+}\PYG{n}{i}\PYG{p}{)}
        \PYG{n}{ok} \PYG{o}{=} \PYG{n}{ok} \PYG{o}{\PYGZam{}}\PYG{o}{\PYGZam{}} \PYG{n}{afpp}\PYG{p}{[}\PYG{n}{i}\PYG{p}{]} \PYG{o}{=}\PYG{o}{=} \PYG{n}{fpp}\PYG{p}{[}\PYG{n}{i}\PYG{p}{]}\PYG{p}{;}
    \PYG{c+c1}{//}
    \PYG{k}{return}\PYG{p}{(} \PYG{n}{ok} \PYG{p}{)}\PYG{p}{;}
\PYG{p}{\PYGZcb{}}
\end{sphinxVerbatim}


\bigskip\hrule\bigskip


xsrst input file: \sphinxcode{\sphinxupquote{example/cplusplus/fun\_hessian\_xam.cpp}}

\index{example@\spxentry{example}}\ignorespaces 

\subparagraph{Example}
\label{\detokenize{xsrst/cpp_fun_hessian:example}}\label{\detokenize{xsrst/cpp_fun_hessian:cpp-fun-hessian-example}}\label{\detokenize{xsrst/cpp_fun_hessian:index-8}}
{\hyperref[\detokenize{xsrst/fun_hessian_xam_cpp:id1}]{\sphinxcrossref{\DUrole{std,std-ref}{fun\_hessian\_xam\_cpp}}}}


\bigskip\hrule\bigskip


xsrst input file: \sphinxcode{\sphinxupquote{lib/cplusplus/fun.cpp}}


\subparagraph{cpp\_fun\_forward}
\label{\detokenize{xsrst/cpp_fun_forward:cpp-fun-forward}}\label{\detokenize{xsrst/cpp_fun_forward::doc}}
\index{cpp\_fun\_forward@\spxentry{cpp\_fun\_forward}}\index{forward@\spxentry{forward}}\index{mode@\spxentry{mode}}\index{ad@\spxentry{ad}}\ignorespaces 

\subparagraph{3.2.3.8 Forward Mode AD}
\label{\detokenize{xsrst/cpp_fun_forward:forward-mode-ad}}\label{\detokenize{xsrst/cpp_fun_forward:id1}}\begin{itemize}
\item {} 
{\hyperref[\detokenize{xsrst/cpp_fun_forward:cpp-fun-forward-syntax}]{\sphinxcrossref{\DUrole{std,std-ref}{Syntax}}}}

\item {} 
{\hyperref[\detokenize{xsrst/cpp_fun_forward:cpp-fun-forward-taylor-coefficient}]{\sphinxcrossref{\DUrole{std,std-ref}{Taylor Coefficient}}}}

\item {} 
{\hyperref[\detokenize{xsrst/cpp_fun_forward:cpp-fun-forward-f}]{\sphinxcrossref{\DUrole{std,std-ref}{f}}}}

\item {} 
{\hyperref[\detokenize{xsrst/cpp_fun_forward:cpp-fun-forward-f-x}]{\sphinxcrossref{\DUrole{std,std-ref}{f(x)}}}}

\item {} 
{\hyperref[\detokenize{xsrst/cpp_fun_forward:cpp-fun-forward-x-t}]{\sphinxcrossref{\DUrole{std,std-ref}{X(t)}}}}

\item {} 
{\hyperref[\detokenize{xsrst/cpp_fun_forward:cpp-fun-forward-y-t}]{\sphinxcrossref{\DUrole{std,std-ref}{Y(t)}}}}

\item {} \begin{description}
\item[{{\hyperref[\detokenize{xsrst/cpp_fun_forward:cpp-fun-forward-p}]{\sphinxcrossref{\DUrole{std,std-ref}{p}}}}}] \leavevmode\begin{itemize}
\item {} 
{\hyperref[\detokenize{xsrst/cpp_fun_forward:cpp-fun-forward-p-size-order}]{\sphinxcrossref{\DUrole{std,std-ref}{size\_order}}}}

\end{itemize}

\end{description}

\item {} 
{\hyperref[\detokenize{xsrst/cpp_fun_forward:cpp-fun-forward-xp}]{\sphinxcrossref{\DUrole{std,std-ref}{xp}}}}

\item {} 
{\hyperref[\detokenize{xsrst/cpp_fun_forward:cpp-fun-forward-yp}]{\sphinxcrossref{\DUrole{std,std-ref}{yp}}}}

\item {} 
{\hyperref[\detokenize{xsrst/cpp_fun_forward:cpp-fun-forward-example}]{\sphinxcrossref{\DUrole{std,std-ref}{Example}}}}

\end{itemize}

\index{syntax@\spxentry{syntax}}\ignorespaces 

\subparagraph{Syntax}
\label{\detokenize{xsrst/cpp_fun_forward:syntax}}\label{\detokenize{xsrst/cpp_fun_forward:cpp-fun-forward-syntax}}\label{\detokenize{xsrst/cpp_fun_forward:index-1}}
\sphinxstyleemphasis{yp} = \sphinxstyleemphasis{f} . \sphinxcode{\sphinxupquote{forward}} ( \sphinxstyleemphasis{p} , \sphinxstyleemphasis{xp} )

\index{taylor@\spxentry{taylor}}\index{coefficient@\spxentry{coefficient}}\ignorespaces 

\subparagraph{Taylor Coefficient}
\label{\detokenize{xsrst/cpp_fun_forward:taylor-coefficient}}\label{\detokenize{xsrst/cpp_fun_forward:cpp-fun-forward-taylor-coefficient}}\label{\detokenize{xsrst/cpp_fun_forward:index-2}}
For a function \(g(t)\) of a scalar argument \(t \in \B{R}\),
the \sphinxstyleemphasis{p}\sphinxhyphen{}th order Taylor coefficient is its
\sphinxcode{\sphinxupquote{p}} \sphinxhyphen{}th order derivative divided by \sphinxstyleemphasis{p} factorial
and evaluated at \(t = 0\); i.e.,
\begin{equation*}
\begin{split}g^{(p)} (0) /  p !\end{split}
\end{equation*}
\index{f@\spxentry{f}}\ignorespaces 

\subparagraph{f}
\label{\detokenize{xsrst/cpp_fun_forward:f}}\label{\detokenize{xsrst/cpp_fun_forward:cpp-fun-forward-f}}\label{\detokenize{xsrst/cpp_fun_forward:index-3}}
This is either a
{\hyperref[\detokenize{xsrst/cpp_fun_ctor:cpp-fun-ctor-syntax-d-fun}]{\sphinxcrossref{\DUrole{std,std-ref}{d\_fun}}}} or
{\hyperref[\detokenize{xsrst/cpp_fun_ctor:cpp-fun-ctor-syntax-a-fun}]{\sphinxcrossref{\DUrole{std,std-ref}{a\_fun}}}} function object.
Note that its state is changed by this operation because
all the Taylor coefficient that it calculates for every
variable in recording are stored.
See more discussion of this fact under the heading
{\hyperref[\detokenize{xsrst/cpp_fun_forward:cpp-fun-forward-p}]{\sphinxcrossref{\DUrole{std,std-ref}{p}}}} below.

\index{f(x)@\spxentry{f(x)}}\ignorespaces 

\subparagraph{f(x)}
\label{\detokenize{xsrst/cpp_fun_forward:f-x}}\label{\detokenize{xsrst/cpp_fun_forward:cpp-fun-forward-f-x}}\label{\detokenize{xsrst/cpp_fun_forward:index-4}}
We use the notation \(f: \B{R}^n \rightarrow \B{R}^m\)
for the function corresponding to \sphinxstyleemphasis{f} .
Note that \sphinxstyleemphasis{n} is the size of {\hyperref[\detokenize{xsrst/cpp_fun_ctor:cpp-fun-ctor-ax}]{\sphinxcrossref{\DUrole{std,std-ref}{ax}}}}
and \sphinxstyleemphasis{m} is the size of {\hyperref[\detokenize{xsrst/cpp_fun_ctor:cpp-fun-ctor-ay}]{\sphinxcrossref{\DUrole{std,std-ref}{ay}}}}
in to the constructor for \sphinxstyleemphasis{f} .

\index{x(t)@\spxentry{x(t)}}\ignorespaces 

\subparagraph{X(t)}
\label{\detokenize{xsrst/cpp_fun_forward:x-t}}\label{\detokenize{xsrst/cpp_fun_forward:cpp-fun-forward-x-t}}\label{\detokenize{xsrst/cpp_fun_forward:index-5}}
We use the notation \(X : \B{R} \rightarrow \B{R}^n\)
for a function that the calling routine chooses.

\index{y(t)@\spxentry{y(t)}}\ignorespaces 

\subparagraph{Y(t)}
\label{\detokenize{xsrst/cpp_fun_forward:y-t}}\label{\detokenize{xsrst/cpp_fun_forward:cpp-fun-forward-y-t}}\label{\detokenize{xsrst/cpp_fun_forward:index-6}}
We define the function \(Y : \B{R} \rightarrow \B{R}^n\)
by \(Y(t) = f(X(t))\).

\index{p@\spxentry{p}}\ignorespaces 

\subparagraph{p}
\label{\detokenize{xsrst/cpp_fun_forward:p}}\label{\detokenize{xsrst/cpp_fun_forward:cpp-fun-forward-p}}\label{\detokenize{xsrst/cpp_fun_forward:index-7}}
This argument has prototype

\begin{DUlineblock}{0em}
\item[]         \sphinxcode{\sphinxupquote{int}} \sphinxstyleemphasis{p}
\end{DUlineblock}

and is non\sphinxhyphen{}negative.
It is the order of the Taylor coefficient being calculated.
If there was no call to \sphinxcode{\sphinxupquote{forward}} for this \sphinxstyleemphasis{f} ,
the value of \sphinxstyleemphasis{p} must be zero.
Otherwise, it must be between zero and one greater that its
value for the previous call using this \sphinxstyleemphasis{f} .
After this call, the Taylor coefficients for orders zero though \sphinxstyleemphasis{p} ,
and for every variable in the recording, will be stored in \sphinxstyleemphasis{f} .

\index{size\_order@\spxentry{size\_order}}\ignorespaces 

\subparagraph{size\_order}
\label{\detokenize{xsrst/cpp_fun_forward:size-order}}\label{\detokenize{xsrst/cpp_fun_forward:cpp-fun-forward-p-size-order}}\label{\detokenize{xsrst/cpp_fun_forward:index-8}}
After this call,
{\hyperref[\detokenize{xsrst/cpp_fun_property:cpp-fun-property-size-order}]{\sphinxcrossref{\DUrole{std,std-ref}{f\_size\_order()}}}} is \sphinxstyleemphasis{p} +1 .

\index{xp@\spxentry{xp}}\ignorespaces 

\subparagraph{xp}
\label{\detokenize{xsrst/cpp_fun_forward:xp}}\label{\detokenize{xsrst/cpp_fun_forward:cpp-fun-forward-xp}}\label{\detokenize{xsrst/cpp_fun_forward:index-9}}
If \sphinxstyleemphasis{f} is a \sphinxcode{\sphinxupquote{d\_fun}} or \sphinxcode{\sphinxupquote{a\_fun}} ,
this argument has prototype

\begin{DUlineblock}{0em}
\item[]         \sphinxcode{\sphinxupquote{const vec\_double\&}} \sphinxstyleemphasis{xp}
\item[]         \sphinxcode{\sphinxupquote{const vec\_a\_double\&}} \sphinxstyleemphasis{xp}
\end{DUlineblock}

respectively and its size must be \sphinxstyleemphasis{n} .
It specifies the \sphinxstyleemphasis{p}\sphinxhyphen{}th order Taylor coefficients for \sphinxstyleemphasis{X(t} ) .

\index{yp@\spxentry{yp}}\ignorespaces 

\subparagraph{yp}
\label{\detokenize{xsrst/cpp_fun_forward:yp}}\label{\detokenize{xsrst/cpp_fun_forward:cpp-fun-forward-yp}}\label{\detokenize{xsrst/cpp_fun_forward:index-10}}
If \sphinxstyleemphasis{f} is a \sphinxcode{\sphinxupquote{d\_fun}} or \sphinxcode{\sphinxupquote{a\_fun}} ,
the result has prototype

\begin{DUlineblock}{0em}
\item[]         \sphinxcode{\sphinxupquote{vec\_double\&}} \sphinxstyleemphasis{yp}
\item[]         \sphinxcode{\sphinxupquote{vec\_a\_double\&}} \sphinxstyleemphasis{yp}
\end{DUlineblock}

respectively and its size is \sphinxstyleemphasis{m} .
It is the \sphinxstyleemphasis{p}\sphinxhyphen{}th order Taylor coefficients for \(Y(t)\).


\subparagraph{fun\_forward\_xam\_cpp}
\label{\detokenize{xsrst/fun_forward_xam_cpp:fun-forward-xam-cpp}}\label{\detokenize{xsrst/fun_forward_xam_cpp::doc}}
\index{fun\_forward\_xam\_cpp@\spxentry{fun\_forward\_xam\_cpp}}\index{c++:@\spxentry{c++:}}\index{forward@\spxentry{forward}}\index{mode@\spxentry{mode}}\index{ad:@\spxentry{ad:}}\index{example@\spxentry{example}}\index{test@\spxentry{test}}\ignorespaces 

\subparagraph{3.2.3.8.1 C++: Forward Mode AD: Example and Test}
\label{\detokenize{xsrst/fun_forward_xam_cpp:c-forward-mode-ad-example-and-test}}\label{\detokenize{xsrst/fun_forward_xam_cpp:id1}}
\begin{sphinxVerbatim}[commandchars=\\\{\}]
\PYG{c+cp}{\PYGZsh{}}\PYG{c+cp}{ include \PYGZlt{}cstdio\PYGZgt{}}
\PYG{c+cp}{\PYGZsh{}}\PYG{c+cp}{ include \PYGZlt{}cppad}\PYG{c+cp}{/}\PYG{c+cp}{py}\PYG{c+cp}{/}\PYG{c+cp}{cppad\PYGZus{}py.hpp\PYGZgt{}}

\PYG{k+kt}{bool} \PYG{n+nf}{fun\PYGZus{}forward\PYGZus{}xam}\PYG{p}{(}\PYG{k+kt}{void}\PYG{p}{)} \PYG{p}{\PYGZob{}}
    \PYG{k}{using} \PYG{n}{cppad\PYGZus{}py}\PYG{o}{:}\PYG{o}{:}\PYG{n}{a\PYGZus{}double}\PYG{p}{;}
    \PYG{k}{using} \PYG{n}{cppad\PYGZus{}py}\PYG{o}{:}\PYG{o}{:}\PYG{n}{vec\PYGZus{}double}\PYG{p}{;}
    \PYG{k}{using} \PYG{n}{cppad\PYGZus{}py}\PYG{o}{:}\PYG{o}{:}\PYG{n}{vec\PYGZus{}a\PYGZus{}double}\PYG{p}{;}
    \PYG{k}{using} \PYG{n}{cppad\PYGZus{}py}\PYG{o}{:}\PYG{o}{:}\PYG{n}{d\PYGZus{}fun}\PYG{p}{;}
    \PYG{k}{using} \PYG{n}{cppad\PYGZus{}py}\PYG{o}{:}\PYG{o}{:}\PYG{n}{a\PYGZus{}fun}\PYG{p}{;}
    \PYG{c+c1}{//}
    \PYG{c+c1}{// initialize return variable}
    \PYG{k+kt}{bool} \PYG{n}{ok} \PYG{o}{=} \PYG{n+nb}{true}\PYG{p}{;}
    \PYG{c+c1}{// \PYGZhy{}\PYGZhy{}\PYGZhy{}\PYGZhy{}\PYGZhy{}\PYGZhy{}\PYGZhy{}\PYGZhy{}\PYGZhy{}\PYGZhy{}\PYGZhy{}\PYGZhy{}\PYGZhy{}\PYGZhy{}\PYGZhy{}\PYGZhy{}\PYGZhy{}\PYGZhy{}\PYGZhy{}\PYGZhy{}\PYGZhy{}\PYGZhy{}\PYGZhy{}\PYGZhy{}\PYGZhy{}\PYGZhy{}\PYGZhy{}\PYGZhy{}\PYGZhy{}\PYGZhy{}\PYGZhy{}\PYGZhy{}\PYGZhy{}\PYGZhy{}\PYGZhy{}\PYGZhy{}\PYGZhy{}\PYGZhy{}\PYGZhy{}\PYGZhy{}\PYGZhy{}\PYGZhy{}\PYGZhy{}\PYGZhy{}\PYGZhy{}\PYGZhy{}\PYGZhy{}\PYGZhy{}\PYGZhy{}\PYGZhy{}\PYGZhy{}\PYGZhy{}\PYGZhy{}\PYGZhy{}\PYGZhy{}\PYGZhy{}\PYGZhy{}\PYGZhy{}\PYGZhy{}\PYGZhy{}\PYGZhy{}\PYGZhy{}\PYGZhy{}\PYGZhy{}\PYGZhy{}\PYGZhy{}\PYGZhy{}\PYGZhy{}\PYGZhy{}\PYGZhy{}}
    \PYG{c+c1}{// number of dependent and independent variables}
    \PYG{k+kt}{int} \PYG{n}{n\PYGZus{}dep} \PYG{o}{=} \PYG{l+m+mi}{1}\PYG{p}{;}
    \PYG{k+kt}{int} \PYG{n}{n\PYGZus{}ind} \PYG{o}{=} \PYG{l+m+mi}{2}\PYG{p}{;}
    \PYG{c+c1}{//}
    \PYG{c+c1}{// create the independent variables ax}
    \PYG{n}{vec\PYGZus{}double} \PYG{n}{xp}\PYG{p}{(}\PYG{n}{n\PYGZus{}ind}\PYG{p}{)}\PYG{p}{;}
    \PYG{k}{for}\PYG{p}{(}\PYG{k+kt}{int} \PYG{n}{i} \PYG{o}{=} \PYG{l+m+mi}{0}\PYG{p}{;} \PYG{n}{i} \PYG{o}{\PYGZlt{}} \PYG{n}{n\PYGZus{}ind} \PYG{p}{;} \PYG{n}{i}\PYG{o}{+}\PYG{o}{+}\PYG{p}{)} \PYG{p}{\PYGZob{}}
        \PYG{n}{xp}\PYG{p}{[}\PYG{n}{i}\PYG{p}{]} \PYG{o}{=} \PYG{n}{i} \PYG{o}{+} \PYG{l+m+mf}{1.0}\PYG{p}{;}
    \PYG{p}{\PYGZcb{}}
    \PYG{n}{vec\PYGZus{}a\PYGZus{}double} \PYG{n}{ax} \PYG{o}{=} \PYG{n}{cppad\PYGZus{}py}\PYG{o}{:}\PYG{o}{:}\PYG{n}{independent}\PYG{p}{(}\PYG{n}{xp}\PYG{p}{)}\PYG{p}{;}
    \PYG{c+c1}{//}
    \PYG{c+c1}{// create dependent varialbes ay with ay0 = ax0 * ax1}
    \PYG{n}{a\PYGZus{}double} \PYG{n}{ax0} \PYG{o}{=} \PYG{n}{ax}\PYG{p}{[}\PYG{l+m+mi}{0}\PYG{p}{]}\PYG{p}{;}
    \PYG{n}{a\PYGZus{}double} \PYG{n}{ax1} \PYG{o}{=} \PYG{n}{ax}\PYG{p}{[}\PYG{l+m+mi}{1}\PYG{p}{]}\PYG{p}{;}
    \PYG{n}{vec\PYGZus{}a\PYGZus{}double} \PYG{n}{ay}\PYG{p}{(}\PYG{n}{n\PYGZus{}dep}\PYG{p}{)}\PYG{p}{;}
    \PYG{n}{ay}\PYG{p}{[}\PYG{l+m+mi}{0}\PYG{p}{]} \PYG{o}{=} \PYG{n}{ax0} \PYG{o}{*} \PYG{n}{ax1}\PYG{p}{;}
    \PYG{c+c1}{//}
    \PYG{c+c1}{// define af corresponding to f(x) = x0 * x1}
    \PYG{n}{d\PYGZus{}fun} \PYG{n}{f}\PYG{p}{(}\PYG{n}{ax}\PYG{p}{,} \PYG{n}{ay}\PYG{p}{)}\PYG{p}{;}
    \PYG{n}{ok} \PYG{o}{\PYGZam{}}\PYG{o}{=} \PYG{n}{f}\PYG{p}{.}\PYG{n}{size\PYGZus{}order}\PYG{p}{(}\PYG{p}{)} \PYG{o}{=}\PYG{o}{=} \PYG{l+m+mi}{0}\PYG{p}{;}
    \PYG{c+c1}{//}
    \PYG{c+c1}{// define X(t) = (3 + t, 2 + t)}
    \PYG{c+c1}{// it follows that Y(t) = f(X(t)) = (3 + t) * (2 + t)}
    \PYG{c+c1}{//}
    \PYG{c+c1}{// Y(0) = 6 and p ! = 1}
    \PYG{k+kt}{int} \PYG{n}{p} \PYG{o}{=} \PYG{l+m+mi}{0}\PYG{p}{;}
    \PYG{n}{xp}\PYG{p}{[}\PYG{l+m+mi}{0}\PYG{p}{]} \PYG{o}{=} \PYG{l+m+mf}{3.0}\PYG{p}{;}
    \PYG{n}{xp}\PYG{p}{[}\PYG{l+m+mi}{1}\PYG{p}{]} \PYG{o}{=} \PYG{l+m+mf}{2.0}\PYG{p}{;}
    \PYG{n}{vec\PYGZus{}double} \PYG{n}{yp} \PYG{o}{=} \PYG{n}{f}\PYG{p}{.}\PYG{n}{forward}\PYG{p}{(}\PYG{n}{p}\PYG{p}{,} \PYG{n}{xp}\PYG{p}{)}\PYG{p}{;}
    \PYG{n}{ok}  \PYG{o}{=} \PYG{n}{ok} \PYG{o}{\PYGZam{}}\PYG{o}{\PYGZam{}} \PYG{n}{yp}\PYG{p}{[}\PYG{l+m+mi}{0}\PYG{p}{]} \PYG{o}{=}\PYG{o}{=} \PYG{l+m+mf}{6.0}\PYG{p}{;}
    \PYG{n}{ok} \PYG{o}{\PYGZam{}}\PYG{o}{=} \PYG{n}{f}\PYG{p}{.}\PYG{n}{size\PYGZus{}order}\PYG{p}{(}\PYG{p}{)} \PYG{o}{=}\PYG{o}{=} \PYG{l+m+mi}{1}\PYG{p}{;}
    \PYG{c+c1}{//}
    \PYG{c+c1}{// first order Taylor coefficients for X(t)}
    \PYG{n}{p} \PYG{o}{=} \PYG{l+m+mi}{1}\PYG{p}{;}
    \PYG{n}{xp}\PYG{p}{[}\PYG{l+m+mi}{0}\PYG{p}{]} \PYG{o}{=} \PYG{l+m+mf}{1.0}\PYG{p}{;}
    \PYG{n}{xp}\PYG{p}{[}\PYG{l+m+mi}{1}\PYG{p}{]} \PYG{o}{=} \PYG{l+m+mf}{1.0}\PYG{p}{;}
    \PYG{c+c1}{//}
    \PYG{c+c1}{// first order Taylor coefficient for Y(t)}
    \PYG{c+c1}{// Y\PYGZsq{}(0) = 3 + 2 = 5 and p ! = 1}
    \PYG{n}{yp} \PYG{o}{=} \PYG{n}{f}\PYG{p}{.}\PYG{n}{forward}\PYG{p}{(}\PYG{n}{p}\PYG{p}{,} \PYG{n}{xp}\PYG{p}{)}\PYG{p}{;}
    \PYG{n}{ok} \PYG{o}{=} \PYG{n}{ok} \PYG{o}{\PYGZam{}}\PYG{o}{\PYGZam{}} \PYG{n}{yp}\PYG{p}{[}\PYG{l+m+mi}{0}\PYG{p}{]} \PYG{o}{=}\PYG{o}{=} \PYG{l+m+mf}{5.0}\PYG{p}{;}
    \PYG{n}{ok} \PYG{o}{\PYGZam{}}\PYG{o}{=} \PYG{n}{f}\PYG{p}{.}\PYG{n}{size\PYGZus{}order}\PYG{p}{(}\PYG{p}{)} \PYG{o}{=}\PYG{o}{=} \PYG{l+m+mi}{2}\PYG{p}{;}
    \PYG{c+c1}{//}
    \PYG{c+c1}{// second order Taylor coefficients for X(t)}
    \PYG{n}{p} \PYG{o}{=} \PYG{l+m+mi}{2}\PYG{p}{;}
    \PYG{n}{xp}\PYG{p}{[}\PYG{l+m+mi}{0}\PYG{p}{]} \PYG{o}{=} \PYG{l+m+mf}{0.0}\PYG{p}{;}
    \PYG{n}{xp}\PYG{p}{[}\PYG{l+m+mi}{1}\PYG{p}{]} \PYG{o}{=} \PYG{l+m+mf}{0.0}\PYG{p}{;}
    \PYG{c+c1}{//}
    \PYG{c+c1}{// second order Taylor coefficient for Y(t)}
    \PYG{c+c1}{// Y\PYGZsq{}\PYGZsq{}(0) = 2.0 and p ! = 2}
    \PYG{n}{yp} \PYG{o}{=} \PYG{n}{f}\PYG{p}{.}\PYG{n}{forward}\PYG{p}{(}\PYG{n}{p}\PYG{p}{,} \PYG{n}{xp}\PYG{p}{)}\PYG{p}{;}
    \PYG{n}{ok} \PYG{o}{=} \PYG{n}{ok} \PYG{o}{\PYGZam{}}\PYG{o}{\PYGZam{}} \PYG{n}{yp}\PYG{p}{[}\PYG{l+m+mi}{0}\PYG{p}{]} \PYG{o}{=}\PYG{o}{=} \PYG{l+m+mf}{1.0}\PYG{p}{;}
    \PYG{n}{ok} \PYG{o}{\PYGZam{}}\PYG{o}{=} \PYG{n}{f}\PYG{p}{.}\PYG{n}{size\PYGZus{}order}\PYG{p}{(}\PYG{p}{)} \PYG{o}{=}\PYG{o}{=} \PYG{l+m+mi}{3}\PYG{p}{;}
    \PYG{c+c1}{// \PYGZhy{}\PYGZhy{}\PYGZhy{}\PYGZhy{}\PYGZhy{}\PYGZhy{}\PYGZhy{}\PYGZhy{}\PYGZhy{}\PYGZhy{}\PYGZhy{}\PYGZhy{}\PYGZhy{}\PYGZhy{}\PYGZhy{}\PYGZhy{}\PYGZhy{}\PYGZhy{}\PYGZhy{}\PYGZhy{}\PYGZhy{}\PYGZhy{}\PYGZhy{}\PYGZhy{}\PYGZhy{}\PYGZhy{}\PYGZhy{}\PYGZhy{}\PYGZhy{}\PYGZhy{}\PYGZhy{}\PYGZhy{}\PYGZhy{}\PYGZhy{}\PYGZhy{}\PYGZhy{}\PYGZhy{}\PYGZhy{}\PYGZhy{}\PYGZhy{}\PYGZhy{}\PYGZhy{}\PYGZhy{}\PYGZhy{}\PYGZhy{}\PYGZhy{}\PYGZhy{}\PYGZhy{}\PYGZhy{}\PYGZhy{}\PYGZhy{}\PYGZhy{}\PYGZhy{}\PYGZhy{}\PYGZhy{}\PYGZhy{}\PYGZhy{}\PYGZhy{}\PYGZhy{}\PYGZhy{}\PYGZhy{}\PYGZhy{}\PYGZhy{}\PYGZhy{}\PYGZhy{}\PYGZhy{}\PYGZhy{}\PYGZhy{}\PYGZhy{}\PYGZhy{}}
    \PYG{n}{a\PYGZus{}fun} \PYG{n}{af}\PYG{p}{(}\PYG{n}{f}\PYG{p}{)}\PYG{p}{;}
    \PYG{n}{ok} \PYG{o}{\PYGZam{}}\PYG{o}{=} \PYG{n}{af}\PYG{p}{.}\PYG{n}{size\PYGZus{}order}\PYG{p}{(}\PYG{p}{)} \PYG{o}{=}\PYG{o}{=} \PYG{l+m+mi}{0}\PYG{p}{;}
    \PYG{c+c1}{//}
    \PYG{c+c1}{// zero order forward}
    \PYG{n}{vec\PYGZus{}a\PYGZus{}double} \PYG{n}{axp}\PYG{p}{(}\PYG{n}{n\PYGZus{}ind}\PYG{p}{)}\PYG{p}{,} \PYG{n}{ayp}\PYG{p}{(}\PYG{n}{n\PYGZus{}dep}\PYG{p}{)}\PYG{p}{;}
    \PYG{n}{p}      \PYG{o}{=} \PYG{l+m+mi}{0}\PYG{p}{;}
    \PYG{n}{axp}\PYG{p}{[}\PYG{l+m+mi}{0}\PYG{p}{]} \PYG{o}{=} \PYG{l+m+mf}{3.0}\PYG{p}{;}
    \PYG{n}{axp}\PYG{p}{[}\PYG{l+m+mi}{1}\PYG{p}{]} \PYG{o}{=} \PYG{l+m+mf}{2.0}\PYG{p}{;}
    \PYG{n}{ayp}    \PYG{o}{=} \PYG{n}{af}\PYG{p}{.}\PYG{n}{forward}\PYG{p}{(}\PYG{n}{p}\PYG{p}{,} \PYG{n}{axp}\PYG{p}{)}\PYG{p}{;}
    \PYG{n}{ok}     \PYG{o}{=} \PYG{n}{ok} \PYG{o}{\PYGZam{}}\PYG{o}{\PYGZam{}} \PYG{n}{ayp}\PYG{p}{[}\PYG{l+m+mi}{0}\PYG{p}{]} \PYG{o}{=}\PYG{o}{=} \PYG{l+m+mf}{6.0}\PYG{p}{;}
    \PYG{n}{ok}    \PYG{o}{\PYGZam{}}\PYG{o}{=} \PYG{n}{af}\PYG{p}{.}\PYG{n}{size\PYGZus{}order}\PYG{p}{(}\PYG{p}{)} \PYG{o}{=}\PYG{o}{=} \PYG{l+m+mi}{1}\PYG{p}{;}
    \PYG{c+c1}{//}
    \PYG{k}{return}\PYG{p}{(} \PYG{n}{ok} \PYG{p}{)}\PYG{p}{;}
\PYG{p}{\PYGZcb{}}
\end{sphinxVerbatim}


\bigskip\hrule\bigskip


xsrst input file: \sphinxcode{\sphinxupquote{example/cplusplus/fun\_forward\_xam.cpp}}

\index{example@\spxentry{example}}\ignorespaces 

\subparagraph{Example}
\label{\detokenize{xsrst/cpp_fun_forward:example}}\label{\detokenize{xsrst/cpp_fun_forward:cpp-fun-forward-example}}\label{\detokenize{xsrst/cpp_fun_forward:index-11}}
{\hyperref[\detokenize{xsrst/fun_forward_xam_cpp:id1}]{\sphinxcrossref{\DUrole{std,std-ref}{fun\_forward\_xam\_cpp}}}}


\bigskip\hrule\bigskip


xsrst input file: \sphinxcode{\sphinxupquote{lib/cplusplus/fun.cpp}}


\subparagraph{cpp\_fun\_reverse}
\label{\detokenize{xsrst/cpp_fun_reverse:cpp-fun-reverse}}\label{\detokenize{xsrst/cpp_fun_reverse::doc}}
\index{cpp\_fun\_reverse@\spxentry{cpp\_fun\_reverse}}\index{reverse@\spxentry{reverse}}\index{mode@\spxentry{mode}}\index{ad@\spxentry{ad}}\ignorespaces 

\subparagraph{3.2.3.9 Reverse Mode AD}
\label{\detokenize{xsrst/cpp_fun_reverse:reverse-mode-ad}}\label{\detokenize{xsrst/cpp_fun_reverse:id1}}\begin{itemize}
\item {} 
{\hyperref[\detokenize{xsrst/cpp_fun_reverse:cpp-fun-reverse-syntax}]{\sphinxcrossref{\DUrole{std,std-ref}{Syntax}}}}

\item {} 
{\hyperref[\detokenize{xsrst/cpp_fun_reverse:cpp-fun-reverse-f}]{\sphinxcrossref{\DUrole{std,std-ref}{f}}}}

\item {} \begin{description}
\item[{{\hyperref[\detokenize{xsrst/cpp_fun_reverse:cpp-fun-reverse-notation}]{\sphinxcrossref{\DUrole{std,std-ref}{Notation}}}}}] \leavevmode\begin{itemize}
\item {} 
{\hyperref[\detokenize{xsrst/cpp_fun_reverse:cpp-fun-reverse-notation-f-x}]{\sphinxcrossref{\DUrole{std,std-ref}{f(x)}}}}

\item {} 
{\hyperref[\detokenize{xsrst/cpp_fun_reverse:cpp-fun-reverse-notation-x-t-s}]{\sphinxcrossref{\DUrole{std,std-ref}{X(t), S}}}}

\item {} 
{\hyperref[\detokenize{xsrst/cpp_fun_reverse:cpp-fun-reverse-notation-y-t-t}]{\sphinxcrossref{\DUrole{std,std-ref}{Y(t), T}}}}

\item {} 
{\hyperref[\detokenize{xsrst/cpp_fun_reverse:cpp-fun-reverse-notation-g-t}]{\sphinxcrossref{\DUrole{std,std-ref}{G(T)}}}}

\end{itemize}

\end{description}

\item {} 
{\hyperref[\detokenize{xsrst/cpp_fun_reverse:cpp-fun-reverse-q}]{\sphinxcrossref{\DUrole{std,std-ref}{q}}}}

\item {} 
{\hyperref[\detokenize{xsrst/cpp_fun_reverse:cpp-fun-reverse-yq}]{\sphinxcrossref{\DUrole{std,std-ref}{yq}}}}

\item {} 
{\hyperref[\detokenize{xsrst/cpp_fun_reverse:cpp-fun-reverse-xq}]{\sphinxcrossref{\DUrole{std,std-ref}{xq}}}}

\item {} 
{\hyperref[\detokenize{xsrst/cpp_fun_reverse:cpp-fun-reverse-example}]{\sphinxcrossref{\DUrole{std,std-ref}{Example}}}}

\end{itemize}

\index{syntax@\spxentry{syntax}}\ignorespaces 

\subparagraph{Syntax}
\label{\detokenize{xsrst/cpp_fun_reverse:syntax}}\label{\detokenize{xsrst/cpp_fun_reverse:cpp-fun-reverse-syntax}}\label{\detokenize{xsrst/cpp_fun_reverse:index-1}}
\sphinxstyleemphasis{xq} = \sphinxstyleemphasis{f} . \sphinxcode{\sphinxupquote{reverse}} ( \sphinxstyleemphasis{q} , \sphinxstyleemphasis{yq} )

\index{f@\spxentry{f}}\ignorespaces 

\subparagraph{f}
\label{\detokenize{xsrst/cpp_fun_reverse:f}}\label{\detokenize{xsrst/cpp_fun_reverse:cpp-fun-reverse-f}}\label{\detokenize{xsrst/cpp_fun_reverse:index-2}}
This is either a
{\hyperref[\detokenize{xsrst/cpp_fun_ctor:cpp-fun-ctor-syntax-d-fun}]{\sphinxcrossref{\DUrole{std,std-ref}{d\_fun}}}} or
{\hyperref[\detokenize{xsrst/cpp_fun_ctor:cpp-fun-ctor-syntax-a-fun}]{\sphinxcrossref{\DUrole{std,std-ref}{a\_fun}}}} function object
and is effectively \sphinxcode{\sphinxupquote{const}} .
(Some details that are not visible to the user may change.)

\index{notation@\spxentry{notation}}\ignorespaces 

\subparagraph{Notation}
\label{\detokenize{xsrst/cpp_fun_reverse:notation}}\label{\detokenize{xsrst/cpp_fun_reverse:cpp-fun-reverse-notation}}\label{\detokenize{xsrst/cpp_fun_reverse:index-3}}
\index{f(x)@\spxentry{f(x)}}\ignorespaces 

\subparagraph{f(x)}
\label{\detokenize{xsrst/cpp_fun_reverse:f-x}}\label{\detokenize{xsrst/cpp_fun_reverse:cpp-fun-reverse-notation-f-x}}\label{\detokenize{xsrst/cpp_fun_reverse:index-4}}
We use the notation \(f: \B{R}^n \rightarrow \B{R}^m\)
for the function corresponding to \sphinxstyleemphasis{f} .
Note that \sphinxstyleemphasis{n} is the size of {\hyperref[\detokenize{xsrst/cpp_fun_ctor:cpp-fun-ctor-ax}]{\sphinxcrossref{\DUrole{std,std-ref}{ax}}}}
and \sphinxstyleemphasis{m} is the size of {\hyperref[\detokenize{xsrst/cpp_fun_ctor:cpp-fun-ctor-ay}]{\sphinxcrossref{\DUrole{std,std-ref}{ay}}}}
in to the constructor for \sphinxstyleemphasis{f} .

\index{x(t)@\spxentry{x(t)}}\index{s@\spxentry{s}}\ignorespaces 

\subparagraph{X(t), S}
\label{\detokenize{xsrst/cpp_fun_reverse:x-t-s}}\label{\detokenize{xsrst/cpp_fun_reverse:cpp-fun-reverse-notation-x-t-s}}\label{\detokenize{xsrst/cpp_fun_reverse:index-5}}
This is the same function as
{\hyperref[\detokenize{xsrst/cpp_fun_forward:cpp-fun-forward-x-t}]{\sphinxcrossref{\DUrole{std,std-ref}{x(t)}}}} in the previous call to
\sphinxstyleemphasis{f} . \sphinxcode{\sphinxupquote{forward}} .
We use \(S \in \B{R}^{n \times q}\) to denote the Taylor coefficients
of \(X(t)\).

\index{y(t)@\spxentry{y(t)}}\index{t@\spxentry{t}}\ignorespaces 

\subparagraph{Y(t), T}
\label{\detokenize{xsrst/cpp_fun_reverse:y-t-t}}\label{\detokenize{xsrst/cpp_fun_reverse:cpp-fun-reverse-notation-y-t-t}}\label{\detokenize{xsrst/cpp_fun_reverse:index-6}}
This is the same function as
{\hyperref[\detokenize{xsrst/cpp_fun_forward:cpp-fun-forward-y-t}]{\sphinxcrossref{\DUrole{std,std-ref}{y(t)}}}} in the previous call to
\sphinxstyleemphasis{f} . \sphinxcode{\sphinxupquote{forward}} .
We use \(T \in \B{R}^{m \times q}\) to denote the Taylor coefficients
of \(Y(t)\).
We also use the notation \(T(S)\) to express the fact that
the Taylor coefficients for \(Y(t)\) are a function of the
Taylor coefficients of \(X(t)\).

\index{g(t)@\spxentry{g(t)}}\ignorespaces 

\subparagraph{G(T)}
\label{\detokenize{xsrst/cpp_fun_reverse:g-t}}\label{\detokenize{xsrst/cpp_fun_reverse:cpp-fun-reverse-notation-g-t}}\label{\detokenize{xsrst/cpp_fun_reverse:index-7}}
We use the notation \(G : \B{R}^{m \times p} \rightarrow \B{R}\)
for a function that the calling routine chooses.

\index{q@\spxentry{q}}\ignorespaces 

\subparagraph{q}
\label{\detokenize{xsrst/cpp_fun_reverse:q}}\label{\detokenize{xsrst/cpp_fun_reverse:cpp-fun-reverse-q}}\label{\detokenize{xsrst/cpp_fun_reverse:index-8}}
This argument has prototype

\begin{DUlineblock}{0em}
\item[]         \sphinxcode{\sphinxupquote{int}} \sphinxstyleemphasis{q}
\end{DUlineblock}

and is positive.
It is the number of the Taylor coefficient (for each variable)
that we are computing the derivative with respect to.
It must be greater than zero, and
less than or equal
the number of Taylor coefficient stored in \sphinxstyleemphasis{f} ; i.e.,
{\hyperref[\detokenize{xsrst/cpp_fun_property:cpp-fun-property-size-order}]{\sphinxcrossref{\DUrole{std,std-ref}{f\_size\_order()}}}}.

\index{yq@\spxentry{yq}}\ignorespaces 

\subparagraph{yq}
\label{\detokenize{xsrst/cpp_fun_reverse:yq}}\label{\detokenize{xsrst/cpp_fun_reverse:cpp-fun-reverse-yq}}\label{\detokenize{xsrst/cpp_fun_reverse:index-9}}
If \sphinxstyleemphasis{f} is a \sphinxcode{\sphinxupquote{d\_fun}} or \sphinxcode{\sphinxupquote{a\_fun}} ,
this argument has prototype

\begin{DUlineblock}{0em}
\item[]         \sphinxcode{\sphinxupquote{const vec\_double\&}} \sphinxstyleemphasis{yq}
\item[]         \sphinxcode{\sphinxupquote{const vec\_a\_double\&}} \sphinxstyleemphasis{yq}
\end{DUlineblock}

and its size must be \sphinxstyleemphasis{m} \sphinxcode{\sphinxupquote{*}} \sphinxstyleemphasis{q} .
For 0 \textless{}= \sphinxstyleemphasis{i} \textless{} \sphinxstyleemphasis{m} and 0 \textless{}= \sphinxstyleemphasis{k} \textless{} \sphinxstyleemphasis{q} ,
\sphinxstyleemphasis{yq} {[} \sphinxstyleemphasis{i} \sphinxcode{\sphinxupquote{*}} \sphinxstyleemphasis{q} + \sphinxstyleemphasis{k} {]} is the partial derivative of
\(G(T)\) with respect to the \sphinxstyleemphasis{k}\sphinxhyphen{}th order Taylor coefficient
for the \sphinxstyleemphasis{i}\sphinxhyphen{}th component function; i.e.,
the partial derivative of \(G(T)\) w.r.t. \(Y_i^{(k)} (t) / k !\).

\index{xq@\spxentry{xq}}\ignorespaces 

\subparagraph{xq}
\label{\detokenize{xsrst/cpp_fun_reverse:xq}}\label{\detokenize{xsrst/cpp_fun_reverse:cpp-fun-reverse-xq}}\label{\detokenize{xsrst/cpp_fun_reverse:index-10}}
If \sphinxstyleemphasis{f} is a \sphinxcode{\sphinxupquote{d\_fun}} or \sphinxcode{\sphinxupquote{a\_fun}} ,
the result has prototype

\begin{DUlineblock}{0em}
\item[]         \sphinxcode{\sphinxupquote{const vec\_double\&}} \sphinxstyleemphasis{xq}
\item[]         \sphinxcode{\sphinxupquote{const vec\_a\_double\&}} \sphinxstyleemphasis{xq}
\end{DUlineblock}

respectively and its size is \sphinxstyleemphasis{n} \sphinxcode{\sphinxupquote{*}} \sphinxstyleemphasis{q} .
For 0 \textless{}= \sphinxstyleemphasis{j} \textless{} \sphinxstyleemphasis{n} and 0 \textless{}= \sphinxstyleemphasis{k} \textless{} \sphinxstyleemphasis{q} ,
\sphinxstyleemphasis{xq} {[} \sphinxstyleemphasis{j} \sphinxcode{\sphinxupquote{*}} \sphinxstyleemphasis{q} + \sphinxstyleemphasis{k} {]} is the partial derivative of
\(G(T(S))\) with respect to the \sphinxstyleemphasis{k}\sphinxhyphen{}th order Taylor coefficient
for the \sphinxstyleemphasis{j}\sphinxhyphen{}th component function; i.e.,
the partial derivative of
\(G(T(S))\) w.r.t. \(S_j^{(k)} (t) / k !\).


\subparagraph{fun\_reverse\_xam\_cpp}
\label{\detokenize{xsrst/fun_reverse_xam_cpp:fun-reverse-xam-cpp}}\label{\detokenize{xsrst/fun_reverse_xam_cpp::doc}}
\index{fun\_reverse\_xam\_cpp@\spxentry{fun\_reverse\_xam\_cpp}}\index{c++:@\spxentry{c++:}}\index{reverse@\spxentry{reverse}}\index{mode@\spxentry{mode}}\index{ad:@\spxentry{ad:}}\index{example@\spxentry{example}}\index{test@\spxentry{test}}\ignorespaces 

\subparagraph{3.2.3.9.1 C++: Reverse Mode AD: Example and Test}
\label{\detokenize{xsrst/fun_reverse_xam_cpp:c-reverse-mode-ad-example-and-test}}\label{\detokenize{xsrst/fun_reverse_xam_cpp:id1}}
\begin{sphinxVerbatim}[commandchars=\\\{\}]
\PYG{c+cp}{\PYGZsh{}}\PYG{c+cp}{ include \PYGZlt{}cstdio\PYGZgt{}}
\PYG{c+cp}{\PYGZsh{}}\PYG{c+cp}{ include \PYGZlt{}cppad}\PYG{c+cp}{/}\PYG{c+cp}{py}\PYG{c+cp}{/}\PYG{c+cp}{cppad\PYGZus{}py.hpp\PYGZgt{}}

\PYG{k+kt}{bool} \PYG{n+nf}{fun\PYGZus{}reverse\PYGZus{}xam}\PYG{p}{(}\PYG{k+kt}{void}\PYG{p}{)} \PYG{p}{\PYGZob{}}
    \PYG{k}{using} \PYG{n}{cppad\PYGZus{}py}\PYG{o}{:}\PYG{o}{:}\PYG{n}{a\PYGZus{}double}\PYG{p}{;}
    \PYG{k}{using} \PYG{n}{cppad\PYGZus{}py}\PYG{o}{:}\PYG{o}{:}\PYG{n}{vec\PYGZus{}double}\PYG{p}{;}
    \PYG{k}{using} \PYG{n}{cppad\PYGZus{}py}\PYG{o}{:}\PYG{o}{:}\PYG{n}{vec\PYGZus{}a\PYGZus{}double}\PYG{p}{;}
    \PYG{k}{using} \PYG{n}{cppad\PYGZus{}py}\PYG{o}{:}\PYG{o}{:}\PYG{n}{d\PYGZus{}fun}\PYG{p}{;}
    \PYG{k}{using} \PYG{n}{cppad\PYGZus{}py}\PYG{o}{:}\PYG{o}{:}\PYG{n}{a\PYGZus{}fun}\PYG{p}{;}
    \PYG{c+c1}{//}
    \PYG{c+c1}{// initialize return variable}
    \PYG{k+kt}{bool} \PYG{n}{ok} \PYG{o}{=} \PYG{n+nb}{true}\PYG{p}{;}
    \PYG{c+c1}{//\PYGZhy{}\PYGZhy{}\PYGZhy{}\PYGZhy{}\PYGZhy{}\PYGZhy{}\PYGZhy{}\PYGZhy{}\PYGZhy{}\PYGZhy{}\PYGZhy{}\PYGZhy{}\PYGZhy{}\PYGZhy{}\PYGZhy{}\PYGZhy{}\PYGZhy{}\PYGZhy{}\PYGZhy{}\PYGZhy{}\PYGZhy{}\PYGZhy{}\PYGZhy{}\PYGZhy{}\PYGZhy{}\PYGZhy{}\PYGZhy{}\PYGZhy{}\PYGZhy{}\PYGZhy{}\PYGZhy{}\PYGZhy{}\PYGZhy{}\PYGZhy{}\PYGZhy{}\PYGZhy{}\PYGZhy{}\PYGZhy{}\PYGZhy{}\PYGZhy{}\PYGZhy{}\PYGZhy{}\PYGZhy{}\PYGZhy{}\PYGZhy{}\PYGZhy{}\PYGZhy{}\PYGZhy{}\PYGZhy{}\PYGZhy{}\PYGZhy{}\PYGZhy{}\PYGZhy{}\PYGZhy{}\PYGZhy{}\PYGZhy{}\PYGZhy{}\PYGZhy{}\PYGZhy{}\PYGZhy{}\PYGZhy{}\PYGZhy{}\PYGZhy{}\PYGZhy{}\PYGZhy{}\PYGZhy{}\PYGZhy{}\PYGZhy{}\PYGZhy{}\PYGZhy{}\PYGZhy{}\PYGZhy{}}
    \PYG{c+c1}{// number of dependent and independent variables}
    \PYG{k+kt}{int} \PYG{n}{n\PYGZus{}dep} \PYG{o}{=} \PYG{l+m+mi}{1}\PYG{p}{;}
    \PYG{k+kt}{int} \PYG{n}{n\PYGZus{}ind} \PYG{o}{=} \PYG{l+m+mi}{3}\PYG{p}{;}
    \PYG{c+c1}{//}
    \PYG{c+c1}{// create the independent variables ax}
    \PYG{n}{vec\PYGZus{}double} \PYG{n}{xp}\PYG{p}{(}\PYG{n}{n\PYGZus{}ind}\PYG{p}{)}\PYG{p}{;}
    \PYG{k}{for}\PYG{p}{(}\PYG{k+kt}{int} \PYG{n}{i} \PYG{o}{=} \PYG{l+m+mi}{0}\PYG{p}{;} \PYG{n}{i} \PYG{o}{\PYGZlt{}} \PYG{n}{n\PYGZus{}ind} \PYG{p}{;} \PYG{n}{i}\PYG{o}{+}\PYG{o}{+}\PYG{p}{)} \PYG{p}{\PYGZob{}}
        \PYG{n}{xp}\PYG{p}{[}\PYG{n}{i}\PYG{p}{]} \PYG{o}{=} \PYG{n}{i}\PYG{p}{;}
    \PYG{p}{\PYGZcb{}}
    \PYG{n}{vec\PYGZus{}a\PYGZus{}double} \PYG{n}{ax} \PYG{o}{=} \PYG{n}{cppad\PYGZus{}py}\PYG{o}{:}\PYG{o}{:}\PYG{n}{independent}\PYG{p}{(}\PYG{n}{xp}\PYG{p}{)}\PYG{p}{;}
    \PYG{c+c1}{//}
    \PYG{c+c1}{// create dependent variables ay with ay0 = ax\PYGZus{}0 * ax\PYGZus{}1 * ax\PYGZus{}2}
    \PYG{n}{a\PYGZus{}double} \PYG{n}{ax\PYGZus{}0} \PYG{o}{=} \PYG{n}{ax}\PYG{p}{[}\PYG{l+m+mi}{0}\PYG{p}{]}\PYG{p}{;}
    \PYG{n}{a\PYGZus{}double} \PYG{n}{ax\PYGZus{}1} \PYG{o}{=} \PYG{n}{ax}\PYG{p}{[}\PYG{l+m+mi}{1}\PYG{p}{]}\PYG{p}{;}
    \PYG{n}{a\PYGZus{}double} \PYG{n}{ax\PYGZus{}2} \PYG{o}{=} \PYG{n}{ax}\PYG{p}{[}\PYG{l+m+mi}{2}\PYG{p}{]}\PYG{p}{;}
    \PYG{n}{vec\PYGZus{}a\PYGZus{}double} \PYG{n}{ay}\PYG{p}{(}\PYG{n}{n\PYGZus{}dep}\PYG{p}{)}\PYG{p}{;}
    \PYG{n}{ay}\PYG{p}{[}\PYG{l+m+mi}{0}\PYG{p}{]} \PYG{o}{=} \PYG{n}{ax\PYGZus{}0} \PYG{o}{*} \PYG{n}{ax\PYGZus{}1} \PYG{o}{*} \PYG{n}{ax\PYGZus{}2}\PYG{p}{;}
    \PYG{c+c1}{//}
    \PYG{c+c1}{// define af corresponding to f(x) = x\PYGZus{}0 * x\PYGZus{}1 * x\PYGZus{}2}
    \PYG{n}{d\PYGZus{}fun} \PYG{n}{f}\PYG{p}{(}\PYG{n}{ax}\PYG{p}{,} \PYG{n}{ay}\PYG{p}{)}\PYG{p}{;}
    \PYG{c+c1}{// \PYGZhy{}\PYGZhy{}\PYGZhy{}\PYGZhy{}\PYGZhy{}\PYGZhy{}\PYGZhy{}\PYGZhy{}\PYGZhy{}\PYGZhy{}\PYGZhy{}\PYGZhy{}\PYGZhy{}\PYGZhy{}\PYGZhy{}\PYGZhy{}\PYGZhy{}\PYGZhy{}\PYGZhy{}\PYGZhy{}\PYGZhy{}\PYGZhy{}\PYGZhy{}\PYGZhy{}\PYGZhy{}\PYGZhy{}\PYGZhy{}\PYGZhy{}\PYGZhy{}\PYGZhy{}\PYGZhy{}\PYGZhy{}\PYGZhy{}\PYGZhy{}\PYGZhy{}\PYGZhy{}\PYGZhy{}\PYGZhy{}\PYGZhy{}\PYGZhy{}\PYGZhy{}\PYGZhy{}\PYGZhy{}\PYGZhy{}\PYGZhy{}\PYGZhy{}\PYGZhy{}\PYGZhy{}\PYGZhy{}\PYGZhy{}\PYGZhy{}\PYGZhy{}\PYGZhy{}\PYGZhy{}\PYGZhy{}\PYGZhy{}\PYGZhy{}\PYGZhy{}\PYGZhy{}\PYGZhy{}\PYGZhy{}\PYGZhy{}\PYGZhy{}\PYGZhy{}\PYGZhy{}\PYGZhy{}\PYGZhy{}\PYGZhy{}\PYGZhy{}\PYGZhy{}\PYGZhy{}}
    \PYG{c+c1}{// define          X(t) = (x\PYGZus{}0 + t, x\PYGZus{}1 + t, x\PYGZus{}2 + t)}
    \PYG{c+c1}{// it follows that Y(t) = f(X(t)) = (x\PYGZus{}0 + t) * (x\PYGZus{}1 + t) * (x\PYGZus{}2 + t)}
    \PYG{c+c1}{// and that       Y\PYGZsq{}(0) = x\PYGZus{}1 * x\PYGZus{}2 + x\PYGZus{}0 * x\PYGZus{}2 + x\PYGZus{}0 * x\PYGZus{}1}
    \PYG{c+c1}{// \PYGZhy{}\PYGZhy{}\PYGZhy{}\PYGZhy{}\PYGZhy{}\PYGZhy{}\PYGZhy{}\PYGZhy{}\PYGZhy{}\PYGZhy{}\PYGZhy{}\PYGZhy{}\PYGZhy{}\PYGZhy{}\PYGZhy{}\PYGZhy{}\PYGZhy{}\PYGZhy{}\PYGZhy{}\PYGZhy{}\PYGZhy{}\PYGZhy{}\PYGZhy{}\PYGZhy{}\PYGZhy{}\PYGZhy{}\PYGZhy{}\PYGZhy{}\PYGZhy{}\PYGZhy{}\PYGZhy{}\PYGZhy{}\PYGZhy{}\PYGZhy{}\PYGZhy{}\PYGZhy{}\PYGZhy{}\PYGZhy{}\PYGZhy{}\PYGZhy{}\PYGZhy{}\PYGZhy{}\PYGZhy{}\PYGZhy{}\PYGZhy{}\PYGZhy{}\PYGZhy{}\PYGZhy{}\PYGZhy{}\PYGZhy{}\PYGZhy{}\PYGZhy{}\PYGZhy{}\PYGZhy{}\PYGZhy{}\PYGZhy{}\PYGZhy{}\PYGZhy{}\PYGZhy{}\PYGZhy{}\PYGZhy{}\PYGZhy{}\PYGZhy{}\PYGZhy{}\PYGZhy{}\PYGZhy{}\PYGZhy{}\PYGZhy{}\PYGZhy{}\PYGZhy{}\PYGZhy{}}
    \PYG{c+c1}{// zero order forward mode}
    \PYG{k+kt}{int} \PYG{n}{p} \PYG{o}{=} \PYG{l+m+mi}{0}\PYG{p}{;}
    \PYG{n}{xp}\PYG{p}{[}\PYG{l+m+mi}{0}\PYG{p}{]} \PYG{o}{=} \PYG{l+m+mf}{2.0}\PYG{p}{;}
    \PYG{n}{xp}\PYG{p}{[}\PYG{l+m+mi}{1}\PYG{p}{]} \PYG{o}{=} \PYG{l+m+mf}{3.0}\PYG{p}{;}
    \PYG{n}{xp}\PYG{p}{[}\PYG{l+m+mi}{2}\PYG{p}{]} \PYG{o}{=} \PYG{l+m+mf}{4.0}\PYG{p}{;}
    \PYG{n}{vec\PYGZus{}double} \PYG{n}{yp} \PYG{o}{=} \PYG{n}{f}\PYG{p}{.}\PYG{n}{forward}\PYG{p}{(}\PYG{n}{p}\PYG{p}{,} \PYG{n}{xp}\PYG{p}{)}\PYG{p}{;}
    \PYG{n}{ok} \PYG{o}{=} \PYG{n}{ok} \PYG{o}{\PYGZam{}}\PYG{o}{\PYGZam{}} \PYG{n}{yp}\PYG{p}{[}\PYG{l+m+mi}{0}\PYG{p}{]} \PYG{o}{=}\PYG{o}{=} \PYG{l+m+mf}{24.0}\PYG{p}{;}
    \PYG{c+c1}{// \PYGZhy{}\PYGZhy{}\PYGZhy{}\PYGZhy{}\PYGZhy{}\PYGZhy{}\PYGZhy{}\PYGZhy{}\PYGZhy{}\PYGZhy{}\PYGZhy{}\PYGZhy{}\PYGZhy{}\PYGZhy{}\PYGZhy{}\PYGZhy{}\PYGZhy{}\PYGZhy{}\PYGZhy{}\PYGZhy{}\PYGZhy{}\PYGZhy{}\PYGZhy{}\PYGZhy{}\PYGZhy{}\PYGZhy{}\PYGZhy{}\PYGZhy{}\PYGZhy{}\PYGZhy{}\PYGZhy{}\PYGZhy{}\PYGZhy{}\PYGZhy{}\PYGZhy{}\PYGZhy{}\PYGZhy{}\PYGZhy{}\PYGZhy{}\PYGZhy{}\PYGZhy{}\PYGZhy{}\PYGZhy{}\PYGZhy{}\PYGZhy{}\PYGZhy{}\PYGZhy{}\PYGZhy{}\PYGZhy{}\PYGZhy{}\PYGZhy{}\PYGZhy{}\PYGZhy{}\PYGZhy{}\PYGZhy{}\PYGZhy{}\PYGZhy{}\PYGZhy{}\PYGZhy{}\PYGZhy{}\PYGZhy{}\PYGZhy{}\PYGZhy{}\PYGZhy{}\PYGZhy{}\PYGZhy{}\PYGZhy{}\PYGZhy{}\PYGZhy{}\PYGZhy{}\PYGZhy{}}
    \PYG{c+c1}{// first order reverse (derivative of zero order forward)}
    \PYG{c+c1}{// define G( Y ) = y\PYGZus{}0 = x\PYGZus{}0 * x\PYGZus{}1 * x\PYGZus{}2}
    \PYG{k+kt}{int} \PYG{n}{q} \PYG{o}{=} \PYG{l+m+mi}{1}\PYG{p}{;}
    \PYG{n}{vec\PYGZus{}double} \PYG{n}{yq1} \PYG{o}{=} \PYG{n}{vec\PYGZus{}double}\PYG{p}{(}\PYG{n}{n\PYGZus{}dep}\PYG{p}{)}\PYG{p}{;}
    \PYG{n}{yq1}\PYG{p}{[}\PYG{l+m+mi}{0}\PYG{p}{]} \PYG{o}{=} \PYG{l+m+mf}{1.0}\PYG{p}{;}
    \PYG{n}{vec\PYGZus{}double} \PYG{n}{xq1} \PYG{o}{=} \PYG{n}{f}\PYG{p}{.}\PYG{n}{reverse}\PYG{p}{(}\PYG{n}{q}\PYG{p}{,} \PYG{n}{yq1}\PYG{p}{)}\PYG{p}{;}
    \PYG{c+c1}{// partial G w.r.t x\PYGZus{}0}
    \PYG{n}{ok} \PYG{o}{=} \PYG{n}{ok} \PYG{o}{\PYGZam{}}\PYG{o}{\PYGZam{}} \PYG{n}{xq1}\PYG{p}{[}\PYG{l+m+mi}{0}\PYG{p}{]} \PYG{o}{=}\PYG{o}{=} \PYG{l+m+mf}{3.0} \PYG{o}{*} \PYG{l+m+mf}{4.0} \PYG{p}{;}
    \PYG{c+c1}{// partial G w.r.t x\PYGZus{}1}
    \PYG{n}{ok} \PYG{o}{=} \PYG{n}{ok} \PYG{o}{\PYGZam{}}\PYG{o}{\PYGZam{}} \PYG{n}{xq1}\PYG{p}{[}\PYG{l+m+mi}{1}\PYG{p}{]} \PYG{o}{=}\PYG{o}{=} \PYG{l+m+mf}{2.0} \PYG{o}{*} \PYG{l+m+mf}{4.0} \PYG{p}{;}
    \PYG{c+c1}{// partial G w.r.t x\PYGZus{}2}
    \PYG{n}{ok} \PYG{o}{=} \PYG{n}{ok} \PYG{o}{\PYGZam{}}\PYG{o}{\PYGZam{}} \PYG{n}{xq1}\PYG{p}{[}\PYG{l+m+mi}{2}\PYG{p}{]} \PYG{o}{=}\PYG{o}{=} \PYG{l+m+mf}{2.0} \PYG{o}{*} \PYG{l+m+mf}{3.0} \PYG{p}{;}
    \PYG{c+c1}{// \PYGZhy{}\PYGZhy{}\PYGZhy{}\PYGZhy{}\PYGZhy{}\PYGZhy{}\PYGZhy{}\PYGZhy{}\PYGZhy{}\PYGZhy{}\PYGZhy{}\PYGZhy{}\PYGZhy{}\PYGZhy{}\PYGZhy{}\PYGZhy{}\PYGZhy{}\PYGZhy{}\PYGZhy{}\PYGZhy{}\PYGZhy{}\PYGZhy{}\PYGZhy{}\PYGZhy{}\PYGZhy{}\PYGZhy{}\PYGZhy{}\PYGZhy{}\PYGZhy{}\PYGZhy{}\PYGZhy{}\PYGZhy{}\PYGZhy{}\PYGZhy{}\PYGZhy{}\PYGZhy{}\PYGZhy{}\PYGZhy{}\PYGZhy{}\PYGZhy{}\PYGZhy{}\PYGZhy{}\PYGZhy{}\PYGZhy{}\PYGZhy{}\PYGZhy{}\PYGZhy{}\PYGZhy{}\PYGZhy{}\PYGZhy{}\PYGZhy{}\PYGZhy{}\PYGZhy{}\PYGZhy{}\PYGZhy{}\PYGZhy{}\PYGZhy{}\PYGZhy{}\PYGZhy{}\PYGZhy{}\PYGZhy{}\PYGZhy{}\PYGZhy{}\PYGZhy{}\PYGZhy{}\PYGZhy{}\PYGZhy{}\PYGZhy{}\PYGZhy{}\PYGZhy{}\PYGZhy{}}
    \PYG{c+c1}{// first order forward mode}
    \PYG{n}{p} \PYG{o}{=} \PYG{l+m+mi}{1}\PYG{p}{;}
    \PYG{n}{xp}\PYG{p}{[}\PYG{l+m+mi}{0}\PYG{p}{]} \PYG{o}{=} \PYG{l+m+mf}{1.0}\PYG{p}{;}
    \PYG{n}{xp}\PYG{p}{[}\PYG{l+m+mi}{1}\PYG{p}{]} \PYG{o}{=} \PYG{l+m+mf}{1.0}\PYG{p}{;}
    \PYG{n}{xp}\PYG{p}{[}\PYG{l+m+mi}{2}\PYG{p}{]} \PYG{o}{=} \PYG{l+m+mf}{1.0}\PYG{p}{;}
    \PYG{n}{yp} \PYG{o}{=} \PYG{n}{f}\PYG{p}{.}\PYG{n}{forward}\PYG{p}{(}\PYG{n}{p}\PYG{p}{,} \PYG{n}{xp}\PYG{p}{)}\PYG{p}{;}
    \PYG{n}{ok} \PYG{o}{=} \PYG{n}{ok} \PYG{o}{\PYGZam{}}\PYG{o}{\PYGZam{}} \PYG{n}{yp}\PYG{p}{[}\PYG{l+m+mi}{0}\PYG{p}{]} \PYG{o}{=}\PYG{o}{=} \PYG{l+m+mf}{3.0}\PYG{o}{*}\PYG{l+m+mf}{4.0} \PYG{o}{+} \PYG{l+m+mf}{2.0}\PYG{o}{*}\PYG{l+m+mf}{4.0} \PYG{o}{+} \PYG{l+m+mf}{2.0}\PYG{o}{*}\PYG{l+m+mf}{3.0}\PYG{p}{;}
    \PYG{c+c1}{// \PYGZhy{}\PYGZhy{}\PYGZhy{}\PYGZhy{}\PYGZhy{}\PYGZhy{}\PYGZhy{}\PYGZhy{}\PYGZhy{}\PYGZhy{}\PYGZhy{}\PYGZhy{}\PYGZhy{}\PYGZhy{}\PYGZhy{}\PYGZhy{}\PYGZhy{}\PYGZhy{}\PYGZhy{}\PYGZhy{}\PYGZhy{}\PYGZhy{}\PYGZhy{}\PYGZhy{}\PYGZhy{}\PYGZhy{}\PYGZhy{}\PYGZhy{}\PYGZhy{}\PYGZhy{}\PYGZhy{}\PYGZhy{}\PYGZhy{}\PYGZhy{}\PYGZhy{}\PYGZhy{}\PYGZhy{}\PYGZhy{}\PYGZhy{}\PYGZhy{}\PYGZhy{}\PYGZhy{}\PYGZhy{}\PYGZhy{}\PYGZhy{}\PYGZhy{}\PYGZhy{}\PYGZhy{}\PYGZhy{}\PYGZhy{}\PYGZhy{}\PYGZhy{}\PYGZhy{}\PYGZhy{}\PYGZhy{}\PYGZhy{}\PYGZhy{}\PYGZhy{}\PYGZhy{}\PYGZhy{}\PYGZhy{}\PYGZhy{}\PYGZhy{}\PYGZhy{}\PYGZhy{}\PYGZhy{}\PYGZhy{}\PYGZhy{}\PYGZhy{}\PYGZhy{}\PYGZhy{}}
    \PYG{c+c1}{// second order reverse (derivative of first order forward)}
    \PYG{c+c1}{// define G( y\PYGZus{}0\PYGZca{}0 , y\PYGZus{}0\PYGZca{}1 ) = y\PYGZus{}0\PYGZca{}1}
    \PYG{c+c1}{// = x\PYGZus{}1\PYGZca{}0 * x\PYGZus{}2\PYGZca{}0  +  x\PYGZus{}0\PYGZca{}0 * x\PYGZus{}2\PYGZca{}0  +  x\PYGZus{}0\PYGZca{}0  *  x\PYGZus{}1\PYGZca{}0}
    \PYG{n}{q} \PYG{o}{=} \PYG{l+m+mi}{2}\PYG{p}{;}
    \PYG{n}{vec\PYGZus{}double} \PYG{n}{yq2} \PYG{o}{=} \PYG{n}{vec\PYGZus{}double}\PYG{p}{(}\PYG{n}{n\PYGZus{}dep} \PYG{o}{*} \PYG{n}{q}\PYG{p}{)}\PYG{p}{;}
    \PYG{n}{yq2}\PYG{p}{[}\PYG{l+m+mi}{0} \PYG{o}{*} \PYG{n}{q} \PYG{o}{+} \PYG{l+m+mi}{0}\PYG{p}{]} \PYG{o}{=} \PYG{l+m+mf}{0.0}\PYG{p}{;} \PYG{c+c1}{// partial of G w.r.t y\PYGZus{}0\PYGZca{}0}
    \PYG{n}{yq2}\PYG{p}{[}\PYG{l+m+mi}{0} \PYG{o}{*} \PYG{n}{q} \PYG{o}{+} \PYG{l+m+mi}{1}\PYG{p}{]} \PYG{o}{=} \PYG{l+m+mf}{1.0}\PYG{p}{;} \PYG{c+c1}{// partial of G w.r.t y\PYGZus{}0\PYGZca{}1}
    \PYG{n}{vec\PYGZus{}double} \PYG{n}{xq2} \PYG{o}{=} \PYG{n}{f}\PYG{p}{.}\PYG{n}{reverse}\PYG{p}{(}\PYG{n}{q}\PYG{p}{,} \PYG{n}{yq2}\PYG{p}{)}\PYG{p}{;}
    \PYG{c+c1}{// partial G w.r.t x\PYGZus{}0\PYGZca{}0}
    \PYG{n}{ok} \PYG{o}{=} \PYG{n}{ok} \PYG{o}{\PYGZam{}}\PYG{o}{\PYGZam{}} \PYG{n}{xq2}\PYG{p}{[}\PYG{l+m+mi}{0} \PYG{o}{*} \PYG{n}{q} \PYG{o}{+} \PYG{l+m+mi}{0}\PYG{p}{]} \PYG{o}{=}\PYG{o}{=} \PYG{l+m+mf}{3.0} \PYG{o}{+} \PYG{l+m+mf}{4.0}\PYG{p}{;}
    \PYG{c+c1}{// partial G w.r.t x\PYGZus{}1\PYGZca{}0}
    \PYG{n}{ok} \PYG{o}{=} \PYG{n}{ok} \PYG{o}{\PYGZam{}}\PYG{o}{\PYGZam{}} \PYG{n}{xq2}\PYG{p}{[}\PYG{l+m+mi}{1} \PYG{o}{*} \PYG{n}{q} \PYG{o}{+} \PYG{l+m+mi}{0}\PYG{p}{]} \PYG{o}{=}\PYG{o}{=} \PYG{l+m+mf}{2.0} \PYG{o}{+} \PYG{l+m+mf}{4.0}\PYG{p}{;}
    \PYG{c+c1}{// partial G w.r.t x\PYGZus{}2\PYGZca{}0}
    \PYG{n}{ok} \PYG{o}{=} \PYG{n}{ok} \PYG{o}{\PYGZam{}}\PYG{o}{\PYGZam{}} \PYG{n}{xq2}\PYG{p}{[}\PYG{l+m+mi}{2} \PYG{o}{*} \PYG{n}{q} \PYG{o}{+} \PYG{l+m+mi}{0}\PYG{p}{]} \PYG{o}{=}\PYG{o}{=} \PYG{l+m+mf}{2.0} \PYG{o}{+} \PYG{l+m+mf}{3.0}\PYG{p}{;}
    \PYG{c+c1}{// \PYGZhy{}\PYGZhy{}\PYGZhy{}\PYGZhy{}\PYGZhy{}\PYGZhy{}\PYGZhy{}\PYGZhy{}\PYGZhy{}\PYGZhy{}\PYGZhy{}\PYGZhy{}\PYGZhy{}\PYGZhy{}\PYGZhy{}\PYGZhy{}\PYGZhy{}\PYGZhy{}\PYGZhy{}\PYGZhy{}\PYGZhy{}\PYGZhy{}\PYGZhy{}\PYGZhy{}\PYGZhy{}\PYGZhy{}\PYGZhy{}\PYGZhy{}\PYGZhy{}\PYGZhy{}\PYGZhy{}\PYGZhy{}\PYGZhy{}\PYGZhy{}\PYGZhy{}\PYGZhy{}\PYGZhy{}\PYGZhy{}\PYGZhy{}\PYGZhy{}\PYGZhy{}\PYGZhy{}\PYGZhy{}\PYGZhy{}\PYGZhy{}\PYGZhy{}\PYGZhy{}\PYGZhy{}\PYGZhy{}\PYGZhy{}\PYGZhy{}\PYGZhy{}\PYGZhy{}\PYGZhy{}\PYGZhy{}\PYGZhy{}\PYGZhy{}\PYGZhy{}\PYGZhy{}\PYGZhy{}\PYGZhy{}\PYGZhy{}\PYGZhy{}\PYGZhy{}\PYGZhy{}\PYGZhy{}\PYGZhy{}\PYGZhy{}\PYGZhy{}\PYGZhy{}\PYGZhy{}}
    \PYG{n}{a\PYGZus{}fun} \PYG{n}{af}\PYG{p}{(}\PYG{n}{f}\PYG{p}{)}\PYG{p}{;}
    \PYG{n}{ok} \PYG{o}{\PYGZam{}}\PYG{o}{=} \PYG{n}{af}\PYG{p}{.}\PYG{n}{size\PYGZus{}order}\PYG{p}{(}\PYG{p}{)} \PYG{o}{=}\PYG{o}{=} \PYG{l+m+mi}{0}\PYG{p}{;}
    \PYG{c+c1}{//}
    \PYG{c+c1}{// zero order forward}
    \PYG{n}{vec\PYGZus{}a\PYGZus{}double} \PYG{n}{axp}\PYG{p}{(}\PYG{n}{n\PYGZus{}ind}\PYG{p}{)}\PYG{p}{,} \PYG{n}{ayp}\PYG{p}{(}\PYG{n}{n\PYGZus{}dep}\PYG{p}{)}\PYG{p}{;}
    \PYG{n}{p}      \PYG{o}{=} \PYG{l+m+mi}{0}\PYG{p}{;}
    \PYG{n}{axp}\PYG{p}{[}\PYG{l+m+mi}{0}\PYG{p}{]} \PYG{o}{=} \PYG{l+m+mf}{2.0}\PYG{p}{;}
    \PYG{n}{axp}\PYG{p}{[}\PYG{l+m+mi}{1}\PYG{p}{]} \PYG{o}{=} \PYG{l+m+mf}{3.0}\PYG{p}{;}
    \PYG{n}{axp}\PYG{p}{[}\PYG{l+m+mi}{2}\PYG{p}{]} \PYG{o}{=} \PYG{l+m+mf}{4.0}\PYG{p}{;}
    \PYG{n}{ayp}    \PYG{o}{=} \PYG{n}{af}\PYG{p}{.}\PYG{n}{forward}\PYG{p}{(}\PYG{n}{p}\PYG{p}{,} \PYG{n}{axp}\PYG{p}{)}\PYG{p}{;}
    \PYG{n}{ok}     \PYG{o}{=} \PYG{n}{ok} \PYG{o}{\PYGZam{}}\PYG{o}{\PYGZam{}} \PYG{n}{ayp}\PYG{p}{[}\PYG{l+m+mi}{0}\PYG{p}{]} \PYG{o}{=}\PYG{o}{=} \PYG{l+m+mf}{24.0}\PYG{p}{;}
    \PYG{n}{ok}    \PYG{o}{\PYGZam{}}\PYG{o}{=} \PYG{n}{af}\PYG{p}{.}\PYG{n}{size\PYGZus{}order}\PYG{p}{(}\PYG{p}{)} \PYG{o}{=}\PYG{o}{=} \PYG{l+m+mi}{1}\PYG{p}{;}
    \PYG{c+c1}{//}
    \PYG{c+c1}{// first order reverse}
    \PYG{n}{q} \PYG{o}{=} \PYG{l+m+mi}{1}\PYG{p}{;}
    \PYG{n}{vec\PYGZus{}a\PYGZus{}double} \PYG{n}{ayq1} \PYG{o}{=} \PYG{n}{vec\PYGZus{}a\PYGZus{}double}\PYG{p}{(}\PYG{n}{n\PYGZus{}dep}\PYG{p}{)}\PYG{p}{;}
    \PYG{n}{ayq1}\PYG{p}{[}\PYG{l+m+mi}{0}\PYG{p}{]}           \PYG{o}{=} \PYG{l+m+mf}{1.0}\PYG{p}{;}
    \PYG{n}{vec\PYGZus{}a\PYGZus{}double} \PYG{n}{axq1} \PYG{o}{=} \PYG{n}{af}\PYG{p}{.}\PYG{n}{reverse}\PYG{p}{(}\PYG{n}{q}\PYG{p}{,} \PYG{n}{ayq1}\PYG{p}{)}\PYG{p}{;}
    \PYG{c+c1}{// partial G w.r.t x\PYGZus{}0}
    \PYG{n}{ok} \PYG{o}{=} \PYG{n}{ok} \PYG{o}{\PYGZam{}}\PYG{o}{\PYGZam{}} \PYG{n}{axq1}\PYG{p}{[}\PYG{l+m+mi}{0}\PYG{p}{]} \PYG{o}{=}\PYG{o}{=} \PYG{l+m+mf}{3.0} \PYG{o}{*} \PYG{l+m+mf}{4.0}\PYG{p}{;}
    \PYG{c+c1}{// partial G w.r.t x\PYGZus{}1}
    \PYG{n}{ok} \PYG{o}{=} \PYG{n}{ok} \PYG{o}{\PYGZam{}}\PYG{o}{\PYGZam{}} \PYG{n}{axq1}\PYG{p}{[}\PYG{l+m+mi}{1}\PYG{p}{]} \PYG{o}{=}\PYG{o}{=} \PYG{l+m+mf}{2.0} \PYG{o}{*} \PYG{l+m+mf}{4.0}\PYG{p}{;}
    \PYG{c+c1}{// partial G w.r.t x\PYGZus{}2}
    \PYG{n}{ok} \PYG{o}{=} \PYG{n}{ok} \PYG{o}{\PYGZam{}}\PYG{o}{\PYGZam{}} \PYG{n}{axq1}\PYG{p}{[}\PYG{l+m+mi}{2}\PYG{p}{]} \PYG{o}{=}\PYG{o}{=} \PYG{l+m+mf}{2.0} \PYG{o}{*} \PYG{l+m+mf}{3.0}\PYG{p}{;}
    \PYG{c+c1}{//}
    \PYG{k}{return}\PYG{p}{(} \PYG{n}{ok} \PYG{p}{)}\PYG{p}{;}
\PYG{p}{\PYGZcb{}}
\end{sphinxVerbatim}


\bigskip\hrule\bigskip


xsrst input file: \sphinxcode{\sphinxupquote{example/cplusplus/fun\_reverse\_xam.cpp}}

\index{example@\spxentry{example}}\ignorespaces 

\subparagraph{Example}
\label{\detokenize{xsrst/cpp_fun_reverse:example}}\label{\detokenize{xsrst/cpp_fun_reverse:cpp-fun-reverse-example}}\label{\detokenize{xsrst/cpp_fun_reverse:index-11}}
{\hyperref[\detokenize{xsrst/fun_reverse_xam_cpp:id1}]{\sphinxcrossref{\DUrole{std,std-ref}{fun\_reverse\_xam\_cpp}}}}


\bigskip\hrule\bigskip


xsrst input file: \sphinxcode{\sphinxupquote{lib/cplusplus/fun.cpp}}


\subparagraph{cpp\_fun\_optimize}
\label{\detokenize{xsrst/cpp_fun_optimize:cpp-fun-optimize}}\label{\detokenize{xsrst/cpp_fun_optimize::doc}}
\index{cpp\_fun\_optimize@\spxentry{cpp\_fun\_optimize}}\index{optimize@\spxentry{optimize}}\index{an@\spxentry{an}}\index{ad@\spxentry{ad}}\index{function@\spxentry{function}}\ignorespaces 

\subparagraph{3.2.3.10 Optimize an AD Function}
\label{\detokenize{xsrst/cpp_fun_optimize:optimize-an-ad-function}}\label{\detokenize{xsrst/cpp_fun_optimize:id1}}\begin{itemize}
\item {} 
{\hyperref[\detokenize{xsrst/cpp_fun_optimize:cpp-fun-optimize-syntax}]{\sphinxcrossref{\DUrole{std,std-ref}{Syntax}}}}

\item {} 
{\hyperref[\detokenize{xsrst/cpp_fun_optimize:cpp-fun-optimize-purpose}]{\sphinxcrossref{\DUrole{std,std-ref}{Purpose}}}}

\item {} 
{\hyperref[\detokenize{xsrst/cpp_fun_optimize:cpp-fun-optimize-f}]{\sphinxcrossref{\DUrole{std,std-ref}{f}}}}

\item {} 
{\hyperref[\detokenize{xsrst/cpp_fun_optimize:cpp-fun-optimize-example}]{\sphinxcrossref{\DUrole{std,std-ref}{Example}}}}

\end{itemize}

\index{syntax@\spxentry{syntax}}\ignorespaces 

\subparagraph{Syntax}
\label{\detokenize{xsrst/cpp_fun_optimize:syntax}}\label{\detokenize{xsrst/cpp_fun_optimize:cpp-fun-optimize-syntax}}\label{\detokenize{xsrst/cpp_fun_optimize:index-1}}
\sphinxstyleemphasis{f} . \sphinxcode{\sphinxupquote{optimize}} ()

\index{purpose@\spxentry{purpose}}\ignorespaces 

\subparagraph{Purpose}
\label{\detokenize{xsrst/cpp_fun_optimize:purpose}}\label{\detokenize{xsrst/cpp_fun_optimize:cpp-fun-optimize-purpose}}\label{\detokenize{xsrst/cpp_fun_optimize:index-2}}
This reduces the number of operations
(hence to time and memory) used to compute the function
stored in \sphinxstyleemphasis{f}
On the other hand, the optimization may take a significant amount
of time and memory.

\index{f@\spxentry{f}}\ignorespaces 

\subparagraph{f}
\label{\detokenize{xsrst/cpp_fun_optimize:f}}\label{\detokenize{xsrst/cpp_fun_optimize:cpp-fun-optimize-f}}\label{\detokenize{xsrst/cpp_fun_optimize:index-3}}
This object is a
{\hyperref[\detokenize{xsrst/cpp_fun_ctor:cpp-fun-ctor-syntax-d-fun}]{\sphinxcrossref{\DUrole{std,std-ref}{d\_fun}}}}.
Optimizing this \sphinxstyleemphasis{f} also optimizes the
corresponding {\hyperref[\detokenize{xsrst/cpp_fun_ctor:cpp-fun-ctor-syntax-a-fun}]{\sphinxcrossref{\DUrole{std,std-ref}{a\_fun}}}}.


\subparagraph{fun\_optimize\_xam\_cpp}
\label{\detokenize{xsrst/fun_optimize_xam_cpp:fun-optimize-xam-cpp}}\label{\detokenize{xsrst/fun_optimize_xam_cpp::doc}}
\index{fun\_optimize\_xam\_cpp@\spxentry{fun\_optimize\_xam\_cpp}}\index{c++:@\spxentry{c++:}}\index{optimize@\spxentry{optimize}}\index{an@\spxentry{an}}\index{d\_fun:@\spxentry{d\_fun:}}\index{example@\spxentry{example}}\index{test@\spxentry{test}}\ignorespaces 

\subparagraph{3.2.3.10.1 C++: Optimize an d\_fun: Example and Test}
\label{\detokenize{xsrst/fun_optimize_xam_cpp:c-optimize-an-d-fun-example-and-test}}\label{\detokenize{xsrst/fun_optimize_xam_cpp:id1}}
\begin{sphinxVerbatim}[commandchars=\\\{\}]
\PYG{c+cp}{\PYGZsh{}}\PYG{c+cp}{ include \PYGZlt{}cstdio\PYGZgt{}}
\PYG{c+cp}{\PYGZsh{}}\PYG{c+cp}{ include \PYGZlt{}cppad}\PYG{c+cp}{/}\PYG{c+cp}{py}\PYG{c+cp}{/}\PYG{c+cp}{cppad\PYGZus{}py.hpp\PYGZgt{}}

\PYG{k+kt}{bool} \PYG{n+nf}{fun\PYGZus{}optimize\PYGZus{}xam}\PYG{p}{(}\PYG{k+kt}{void}\PYG{p}{)} \PYG{p}{\PYGZob{}}
    \PYG{k}{using} \PYG{n}{cppad\PYGZus{}py}\PYG{o}{:}\PYG{o}{:}\PYG{n}{a\PYGZus{}double}\PYG{p}{;}
    \PYG{k}{using} \PYG{n}{cppad\PYGZus{}py}\PYG{o}{:}\PYG{o}{:}\PYG{n}{vec\PYGZus{}double}\PYG{p}{;}
    \PYG{k}{using} \PYG{n}{cppad\PYGZus{}py}\PYG{o}{:}\PYG{o}{:}\PYG{n}{vec\PYGZus{}a\PYGZus{}double}\PYG{p}{;}
    \PYG{k}{using} \PYG{n}{cppad\PYGZus{}py}\PYG{o}{:}\PYG{o}{:}\PYG{n}{d\PYGZus{}fun}\PYG{p}{;}
    \PYG{c+c1}{//}
    \PYG{c+c1}{// initialize return variable}
    \PYG{k+kt}{bool} \PYG{n}{ok} \PYG{o}{=} \PYG{n+nb}{true}\PYG{p}{;}
    \PYG{c+c1}{//\PYGZhy{}\PYGZhy{}\PYGZhy{}\PYGZhy{}\PYGZhy{}\PYGZhy{}\PYGZhy{}\PYGZhy{}\PYGZhy{}\PYGZhy{}\PYGZhy{}\PYGZhy{}\PYGZhy{}\PYGZhy{}\PYGZhy{}\PYGZhy{}\PYGZhy{}\PYGZhy{}\PYGZhy{}\PYGZhy{}\PYGZhy{}\PYGZhy{}\PYGZhy{}\PYGZhy{}\PYGZhy{}\PYGZhy{}\PYGZhy{}\PYGZhy{}\PYGZhy{}\PYGZhy{}\PYGZhy{}\PYGZhy{}\PYGZhy{}\PYGZhy{}\PYGZhy{}\PYGZhy{}\PYGZhy{}\PYGZhy{}\PYGZhy{}\PYGZhy{}\PYGZhy{}\PYGZhy{}\PYGZhy{}\PYGZhy{}\PYGZhy{}\PYGZhy{}\PYGZhy{}\PYGZhy{}\PYGZhy{}\PYGZhy{}\PYGZhy{}\PYGZhy{}\PYGZhy{}\PYGZhy{}\PYGZhy{}\PYGZhy{}\PYGZhy{}\PYGZhy{}\PYGZhy{}\PYGZhy{}\PYGZhy{}\PYGZhy{}\PYGZhy{}\PYGZhy{}\PYGZhy{}\PYGZhy{}\PYGZhy{}\PYGZhy{}\PYGZhy{}\PYGZhy{}\PYGZhy{}\PYGZhy{}}
    \PYG{k+kt}{int} \PYG{n}{n\PYGZus{}ind} \PYG{o}{=} \PYG{l+m+mi}{1}\PYG{p}{;} \PYG{c+c1}{// number of independent variables}
    \PYG{k+kt}{int} \PYG{n}{n\PYGZus{}dep} \PYG{o}{=} \PYG{l+m+mi}{1}\PYG{p}{;} \PYG{c+c1}{// number of dependent variables}
    \PYG{k+kt}{int} \PYG{n}{n\PYGZus{}var} \PYG{o}{=} \PYG{l+m+mi}{1}\PYG{p}{;} \PYG{c+c1}{// phantom variable at address 0}
    \PYG{k+kt}{int} \PYG{n}{n\PYGZus{}op} \PYG{o}{=} \PYG{l+m+mi}{1}\PYG{p}{;}  \PYG{c+c1}{// special operator at beginning}
    \PYG{c+c1}{//}
    \PYG{c+c1}{// dimension some vectors}
    \PYG{n}{vec\PYGZus{}double} \PYG{n}{x}\PYG{p}{(}\PYG{n}{n\PYGZus{}ind}\PYG{p}{)}\PYG{p}{;}
    \PYG{n}{vec\PYGZus{}a\PYGZus{}double} \PYG{n}{ay}\PYG{p}{(}\PYG{n}{n\PYGZus{}dep}\PYG{p}{)}\PYG{p}{;}
    \PYG{c+c1}{//}
    \PYG{c+c1}{// independent variables}
    \PYG{n}{x}\PYG{p}{[}\PYG{l+m+mi}{0}\PYG{p}{]} \PYG{o}{=} \PYG{l+m+mf}{1.0}\PYG{p}{;}
    \PYG{n}{vec\PYGZus{}a\PYGZus{}double} \PYG{n}{ax} \PYG{o}{=} \PYG{n}{cppad\PYGZus{}py}\PYG{o}{:}\PYG{o}{:}\PYG{n}{independent}\PYG{p}{(}\PYG{n}{x}\PYG{p}{)}\PYG{p}{;}
    \PYG{n}{n\PYGZus{}var} \PYG{o}{=} \PYG{n}{n\PYGZus{}var} \PYG{o}{+} \PYG{n}{n\PYGZus{}ind}\PYG{p}{;} \PYG{c+c1}{// one for each indpendent}
    \PYG{n}{n\PYGZus{}op} \PYG{o}{=} \PYG{n}{n\PYGZus{}op} \PYG{o}{+} \PYG{n}{n\PYGZus{}ind}\PYG{p}{;}
    \PYG{c+c1}{//}
    \PYG{c+c1}{// accumulate summation}
    \PYG{n}{a\PYGZus{}double} \PYG{n}{ax0} \PYG{o}{=} \PYG{n}{ax}\PYG{p}{[}\PYG{l+m+mi}{0}\PYG{p}{]}\PYG{p}{;}
    \PYG{n}{a\PYGZus{}double} \PYG{n}{csum}\PYG{p}{(}\PYG{l+m+mf}{0.0}\PYG{p}{)}\PYG{p}{;}
    \PYG{n}{csum} \PYG{o}{=} \PYG{n}{ax0} \PYG{o}{+} \PYG{n}{ax0} \PYG{o}{+} \PYG{n}{ax0} \PYG{o}{+} \PYG{n}{ax0}\PYG{p}{;}
    \PYG{n}{n\PYGZus{}var} \PYG{o}{=} \PYG{n}{n\PYGZus{}var} \PYG{o}{+} \PYG{l+m+mi}{3}\PYG{p}{;} \PYG{c+c1}{// one per + operator}
    \PYG{n}{n\PYGZus{}op} \PYG{o}{=} \PYG{n}{n\PYGZus{}op} \PYG{o}{+} \PYG{l+m+mi}{3}\PYG{p}{;}
    \PYG{c+c1}{//}
    \PYG{c+c1}{// define f(x) = y\PYGZus{}0 = csum}
    \PYG{n}{ay}\PYG{p}{[}\PYG{l+m+mi}{0}\PYG{p}{]} \PYG{o}{=} \PYG{n}{csum}\PYG{p}{;}
    \PYG{n}{d\PYGZus{}fun} \PYG{n}{f}\PYG{p}{(}\PYG{n}{ax}\PYG{p}{,} \PYG{n}{ay}\PYG{p}{)}\PYG{p}{;}
    \PYG{n}{n\PYGZus{}op} \PYG{o}{=} \PYG{n}{n\PYGZus{}op} \PYG{o}{+} \PYG{l+m+mi}{1}\PYG{p}{;} \PYG{c+c1}{// speical operator at end}
    \PYG{c+c1}{//}
    \PYG{c+c1}{// check number of variables and operators}
    \PYG{n}{ok} \PYG{o}{=} \PYG{n}{ok} \PYG{o}{\PYGZam{}}\PYG{o}{\PYGZam{}} \PYG{n}{f}\PYG{p}{.}\PYG{n}{size\PYGZus{}var}\PYG{p}{(}\PYG{p}{)} \PYG{o}{=}\PYG{o}{=} \PYG{n}{n\PYGZus{}var}\PYG{p}{;}
    \PYG{n}{ok} \PYG{o}{=} \PYG{n}{ok} \PYG{o}{\PYGZam{}}\PYG{o}{\PYGZam{}} \PYG{n}{f}\PYG{p}{.}\PYG{n}{size\PYGZus{}op}\PYG{p}{(}\PYG{p}{)} \PYG{o}{=}\PYG{o}{=} \PYG{n}{n\PYGZus{}op}\PYG{p}{;}
    \PYG{c+c1}{//}
    \PYG{c+c1}{// optimize}
    \PYG{n}{f}\PYG{p}{.}\PYG{n}{optimize}\PYG{p}{(}\PYG{p}{)}\PYG{p}{;}
    \PYG{c+c1}{//}
    \PYG{c+c1}{// number of variables and operators has decreased by two}
    \PYG{n}{ok} \PYG{o}{=} \PYG{n}{ok} \PYG{o}{\PYGZam{}}\PYG{o}{\PYGZam{}} \PYG{n}{f}\PYG{p}{.}\PYG{n}{size\PYGZus{}var}\PYG{p}{(}\PYG{p}{)} \PYG{o}{=}\PYG{o}{=} \PYG{n}{n\PYGZus{}var}\PYG{o}{\PYGZhy{}}\PYG{l+m+mi}{2}\PYG{p}{;}
    \PYG{n}{ok} \PYG{o}{=} \PYG{n}{ok} \PYG{o}{\PYGZam{}}\PYG{o}{\PYGZam{}} \PYG{n}{f}\PYG{p}{.}\PYG{n}{size\PYGZus{}op}\PYG{p}{(}\PYG{p}{)} \PYG{o}{=}\PYG{o}{=} \PYG{n}{n\PYGZus{}op}\PYG{o}{\PYGZhy{}}\PYG{l+m+mi}{2}\PYG{p}{;}
    \PYG{c+c1}{//}
    \PYG{k}{return}\PYG{p}{(} \PYG{n}{ok}  \PYG{p}{)}\PYG{p}{;}
\PYG{p}{\PYGZcb{}}
\end{sphinxVerbatim}


\bigskip\hrule\bigskip


xsrst input file: \sphinxcode{\sphinxupquote{example/cplusplus/fun\_optimize\_xam.cpp}}

\index{example@\spxentry{example}}\ignorespaces 

\subparagraph{Example}
\label{\detokenize{xsrst/cpp_fun_optimize:example}}\label{\detokenize{xsrst/cpp_fun_optimize:cpp-fun-optimize-example}}\label{\detokenize{xsrst/cpp_fun_optimize:index-4}}
{\hyperref[\detokenize{xsrst/fun_optimize_xam_cpp:id1}]{\sphinxcrossref{\DUrole{std,std-ref}{fun\_optimize\_xam\_cpp}}}}


\bigskip\hrule\bigskip


xsrst input file: \sphinxcode{\sphinxupquote{lib/cplusplus/fun.cpp}}


\subparagraph{cpp\_fun\_json}
\label{\detokenize{xsrst/cpp_fun_json:cpp-fun-json}}\label{\detokenize{xsrst/cpp_fun_json::doc}}
\index{cpp\_fun\_json@\spxentry{cpp\_fun\_json}}\index{json@\spxentry{json}}\index{representation@\spxentry{representation}}\index{ad@\spxentry{ad}}\index{computational@\spxentry{computational}}\index{graph@\spxentry{graph}}\ignorespaces 

\subparagraph{3.2.3.11 Json Representation of AD Computational Graph}
\label{\detokenize{xsrst/cpp_fun_json:json-representation-of-ad-computational-graph}}\label{\detokenize{xsrst/cpp_fun_json:id1}}\begin{itemize}
\item {} 
{\hyperref[\detokenize{xsrst/cpp_fun_json:cpp-fun-json-syntax}]{\sphinxcrossref{\DUrole{std,std-ref}{Syntax}}}}

\item {} 
{\hyperref[\detokenize{xsrst/cpp_fun_json:cpp-fun-json-f}]{\sphinxcrossref{\DUrole{std,std-ref}{f}}}}

\item {} 
{\hyperref[\detokenize{xsrst/cpp_fun_json:cpp-fun-json-json}]{\sphinxcrossref{\DUrole{std,std-ref}{json}}}}

\item {} 
{\hyperref[\detokenize{xsrst/cpp_fun_json:cpp-fun-json-to-json}]{\sphinxcrossref{\DUrole{std,std-ref}{to\_json}}}}

\item {} 
{\hyperref[\detokenize{xsrst/cpp_fun_json:cpp-fun-json-from-json}]{\sphinxcrossref{\DUrole{std,std-ref}{from\_json}}}}

\item {} 
{\hyperref[\detokenize{xsrst/cpp_fun_json:cpp-fun-json-examples}]{\sphinxcrossref{\DUrole{std,std-ref}{Examples}}}}

\end{itemize}

\index{syntax@\spxentry{syntax}}\ignorespaces 

\subparagraph{Syntax}
\label{\detokenize{xsrst/cpp_fun_json:syntax}}\label{\detokenize{xsrst/cpp_fun_json:cpp-fun-json-syntax}}\label{\detokenize{xsrst/cpp_fun_json:index-1}}
\begin{DUlineblock}{0em}
\item[] \sphinxstyleemphasis{json} = \sphinxstyleemphasis{f} . \sphinxcode{\sphinxupquote{to\_json}} ()
\item[] \sphinxstyleemphasis{f} . \sphinxcode{\sphinxupquote{from\_json}} ()
\end{DUlineblock}

\index{f@\spxentry{f}}\ignorespaces 

\subparagraph{f}
\label{\detokenize{xsrst/cpp_fun_json:f}}\label{\detokenize{xsrst/cpp_fun_json:cpp-fun-json-f}}\label{\detokenize{xsrst/cpp_fun_json:index-2}}
This is a {\hyperref[\detokenize{xsrst/cpp_fun_ctor:cpp-fun-ctor-syntax-d-fun}]{\sphinxcrossref{\DUrole{std,std-ref}{d\_fun}}}} object.

\index{json@\spxentry{json}}\ignorespaces 

\subparagraph{json}
\label{\detokenize{xsrst/cpp_fun_json:json}}\label{\detokenize{xsrst/cpp_fun_json:cpp-fun-json-json}}\label{\detokenize{xsrst/cpp_fun_json:index-3}}
is a Json representation of the computation graph corresponding to
\sphinxstyleemphasis{f} ; see the CppAD documentation for
\sphinxhref{https://coin-or.github.io/CppAD/doc/json\_ad\_graph.htm}{json\_ad\_graph}.

\index{to\_json@\spxentry{to\_json}}\ignorespaces 

\subparagraph{to\_json}
\label{\detokenize{xsrst/cpp_fun_json:to-json}}\label{\detokenize{xsrst/cpp_fun_json:cpp-fun-json-to-json}}\label{\detokenize{xsrst/cpp_fun_json:index-4}}
In this case, the function object \sphinxstyleemphasis{f} is \sphinxcode{\sphinxupquote{const}} and
the return value \sphinxstyleemphasis{json} has prototype

\begin{DUlineblock}{0em}
\item[]         \sphinxcode{\sphinxupquote{std::string}} \sphinxstyleemphasis{json}
\end{DUlineblock}

\index{from\_json@\spxentry{from\_json}}\ignorespaces 

\subparagraph{from\_json}
\label{\detokenize{xsrst/cpp_fun_json:from-json}}\label{\detokenize{xsrst/cpp_fun_json:cpp-fun-json-from-json}}\label{\detokenize{xsrst/cpp_fun_json:index-5}}
In this case, \sphinxstyleemphasis{json} has prototype

\begin{DUlineblock}{0em}
\item[]         \sphinxcode{\sphinxupquote{const std::string\&}} \sphinxstyleemphasis{json}
\end{DUlineblock}

and the function \sphinxstyleemphasis{f} so it corresponds to \sphinxstyleemphasis{json} .


\subparagraph{fun\_to\_json\_xam\_cpp}
\label{\detokenize{xsrst/fun_to_json_xam_cpp:fun-to-json-xam-cpp}}\label{\detokenize{xsrst/fun_to_json_xam_cpp::doc}}
\index{fun\_to\_json\_xam\_cpp@\spxentry{fun\_to\_json\_xam\_cpp}}\index{c++:@\spxentry{c++:}}\index{to\_json:@\spxentry{to\_json:}}\index{example@\spxentry{example}}\index{test@\spxentry{test}}\ignorespaces 

\subparagraph{3.2.3.11.1 C++: to\_json: Example and Test}
\label{\detokenize{xsrst/fun_to_json_xam_cpp:c-to-json-example-and-test}}\label{\detokenize{xsrst/fun_to_json_xam_cpp:id1}}
\begin{sphinxVerbatim}[commandchars=\\\{\}]
\PYG{k+kt}{bool} \PYG{n+nf}{to\PYGZus{}json\PYGZus{}xam}\PYG{p}{(}\PYG{k+kt}{void}\PYG{p}{)} \PYG{p}{\PYGZob{}}
    \PYG{k}{using} \PYG{n}{cppad\PYGZus{}py}\PYG{o}{:}\PYG{o}{:}\PYG{n}{a\PYGZus{}double}\PYG{p}{;}
    \PYG{k}{using} \PYG{n}{cppad\PYGZus{}py}\PYG{o}{:}\PYG{o}{:}\PYG{n}{vec\PYGZus{}double}\PYG{p}{;}
    \PYG{k}{using} \PYG{n}{cppad\PYGZus{}py}\PYG{o}{:}\PYG{o}{:}\PYG{n}{vec\PYGZus{}a\PYGZus{}double}\PYG{p}{;}
    \PYG{k}{using} \PYG{n}{cppad\PYGZus{}py}\PYG{o}{:}\PYG{o}{:}\PYG{n}{d\PYGZus{}fun}\PYG{p}{;}
    \PYG{k}{using} \PYG{n}{cppad\PYGZus{}py}\PYG{o}{:}\PYG{o}{:}\PYG{n}{a\PYGZus{}fun}\PYG{p}{;}
    \PYG{c+c1}{//}
    \PYG{c+c1}{// initialize return variable}
    \PYG{k+kt}{bool} \PYG{n}{ok} \PYG{o}{=} \PYG{n+nb}{true}\PYG{p}{;}
    \PYG{c+c1}{// \PYGZhy{}\PYGZhy{}\PYGZhy{}\PYGZhy{}\PYGZhy{}\PYGZhy{}\PYGZhy{}\PYGZhy{}\PYGZhy{}\PYGZhy{}\PYGZhy{}\PYGZhy{}\PYGZhy{}\PYGZhy{}\PYGZhy{}\PYGZhy{}\PYGZhy{}\PYGZhy{}\PYGZhy{}\PYGZhy{}\PYGZhy{}\PYGZhy{}\PYGZhy{}\PYGZhy{}\PYGZhy{}\PYGZhy{}\PYGZhy{}\PYGZhy{}\PYGZhy{}\PYGZhy{}\PYGZhy{}\PYGZhy{}\PYGZhy{}\PYGZhy{}\PYGZhy{}\PYGZhy{}\PYGZhy{}\PYGZhy{}\PYGZhy{}\PYGZhy{}\PYGZhy{}\PYGZhy{}\PYGZhy{}\PYGZhy{}\PYGZhy{}\PYGZhy{}\PYGZhy{}\PYGZhy{}\PYGZhy{}\PYGZhy{}\PYGZhy{}\PYGZhy{}\PYGZhy{}\PYGZhy{}\PYGZhy{}\PYGZhy{}\PYGZhy{}\PYGZhy{}\PYGZhy{}\PYGZhy{}\PYGZhy{}\PYGZhy{}\PYGZhy{}\PYGZhy{}\PYGZhy{}\PYGZhy{}\PYGZhy{}\PYGZhy{}\PYGZhy{}\PYGZhy{}}
    \PYG{k+kt}{int} \PYG{n}{n} \PYG{o}{=} \PYG{l+m+mi}{1}\PYG{p}{;} \PYG{c+c1}{// number of independent variables}
    \PYG{k+kt}{int} \PYG{n}{m} \PYG{o}{=} \PYG{l+m+mi}{2}\PYG{p}{;} \PYG{c+c1}{// number of dependent variables}
    \PYG{c+c1}{//}
    \PYG{c+c1}{// dimension some vectors}
    \PYG{n}{vec\PYGZus{}double} \PYG{n}{x}\PYG{p}{(}\PYG{n}{n}\PYG{p}{)}\PYG{p}{;}
    \PYG{n}{vec\PYGZus{}a\PYGZus{}double} \PYG{n}{ay}\PYG{p}{(}\PYG{n}{m}\PYG{p}{)}\PYG{p}{;}
    \PYG{c+c1}{//}
    \PYG{c+c1}{// independent variables}
    \PYG{n}{x}\PYG{p}{[}\PYG{l+m+mi}{0}\PYG{p}{]}            \PYG{o}{=} \PYG{l+m+mf}{1.0}\PYG{p}{;}
    \PYG{n}{vec\PYGZus{}a\PYGZus{}double} \PYG{n}{ax} \PYG{o}{=} \PYG{n}{cppad\PYGZus{}py}\PYG{o}{:}\PYG{o}{:}\PYG{n}{independent}\PYG{p}{(}\PYG{n}{x}\PYG{p}{)}\PYG{p}{;}
    \PYG{c+c1}{//}
    \PYG{c+c1}{// f(x) = [ x0 + x0, sin(x0) ]}
    \PYG{n}{ay}\PYG{p}{[}\PYG{l+m+mi}{0}\PYG{p}{]} \PYG{o}{=} \PYG{n}{ax}\PYG{p}{[}\PYG{l+m+mi}{0}\PYG{p}{]} \PYG{o}{+} \PYG{n}{ax}\PYG{p}{[}\PYG{l+m+mi}{0}\PYG{p}{]}\PYG{p}{;}
    \PYG{n}{ay}\PYG{p}{[}\PYG{l+m+mi}{1}\PYG{p}{]} \PYG{o}{=} \PYG{n}{ax}\PYG{p}{[}\PYG{l+m+mi}{0}\PYG{p}{]}\PYG{p}{.}\PYG{n}{sin}\PYG{p}{(}\PYG{p}{)}\PYG{p}{;}
    \PYG{n}{d\PYGZus{}fun} \PYG{n}{f}\PYG{p}{(}\PYG{n}{ax}\PYG{p}{,} \PYG{n}{ay}\PYG{p}{)}\PYG{p}{;}
    \PYG{c+c1}{//}
    \PYG{n}{std}\PYG{o}{:}\PYG{o}{:}\PYG{n}{string} \PYG{n}{json} \PYG{o}{=} \PYG{n}{f}\PYG{p}{.}\PYG{n}{to\PYGZus{}json}\PYG{p}{(}\PYG{p}{)}\PYG{p}{;}
    \PYG{k+kt}{size\PYGZus{}t} \PYG{n}{pos}       \PYG{o}{=} \PYG{n}{json}\PYG{p}{.}\PYG{n}{find}\PYG{p}{(}\PYG{l+s}{\PYGZdq{}}\PYG{l+s+se}{\PYGZbs{}\PYGZdq{}}\PYG{l+s}{op\PYGZus{}code}\PYG{l+s+se}{\PYGZbs{}\PYGZdq{}}\PYG{l+s}{\PYGZdq{}}\PYG{p}{)}\PYG{p}{;}
    \PYG{k+kt}{size\PYGZus{}t} \PYG{n}{start}     \PYG{o}{=} \PYG{n}{json}\PYG{p}{.}\PYG{n}{find}\PYG{p}{(}\PYG{l+s}{\PYGZdq{}}\PYG{l+s}{:}\PYG{l+s}{\PYGZdq{}}\PYG{p}{,} \PYG{n}{pos}\PYG{p}{)} \PYG{o}{+} \PYG{l+m+mi}{1}\PYG{p}{;}
    \PYG{k+kt}{size\PYGZus{}t} \PYG{n}{end}       \PYG{o}{=} \PYG{n}{json}\PYG{p}{.}\PYG{n}{find}\PYG{p}{(}\PYG{l+s}{\PYGZdq{}}\PYG{l+s}{,}\PYG{l+s}{\PYGZdq{}}\PYG{p}{,} \PYG{n}{pos}\PYG{p}{)}\PYG{p}{;}
    \PYG{k+kt}{int} \PYG{n}{op\PYGZus{}code}      \PYG{o}{=} \PYG{n}{std}\PYG{o}{:}\PYG{o}{:}\PYG{n}{atoi}\PYG{p}{(} \PYG{n}{json}\PYG{p}{.}\PYG{n}{substr}\PYG{p}{(}\PYG{n}{start}\PYG{p}{,} \PYG{n}{end} \PYG{o}{\PYGZhy{}} \PYG{n}{start}\PYG{p}{)}\PYG{p}{.}\PYG{n}{c\PYGZus{}str}\PYG{p}{(}\PYG{p}{)} \PYG{p}{)}\PYG{p}{;}
    \PYG{n}{ok}              \PYG{o}{\PYGZam{}}\PYG{o}{=} \PYG{n}{op\PYGZus{}code} \PYG{o}{=}\PYG{o}{=} \PYG{l+m+mi}{1}\PYG{p}{;}
    \PYG{n}{pos}              \PYG{o}{=} \PYG{n}{json}\PYG{p}{.}\PYG{n}{find}\PYG{p}{(}\PYG{l+s}{\PYGZdq{}}\PYG{l+s+se}{\PYGZbs{}\PYGZdq{}}\PYG{l+s}{name}\PYG{l+s+se}{\PYGZbs{}\PYGZdq{}}\PYG{l+s}{\PYGZdq{}}\PYG{p}{,} \PYG{n}{pos}\PYG{p}{)}\PYG{p}{;}
    \PYG{n}{start}            \PYG{o}{=} \PYG{n}{json}\PYG{p}{.}\PYG{n}{find}\PYG{p}{(}\PYG{l+s}{\PYGZdq{}}\PYG{l+s}{:}\PYG{l+s}{\PYGZdq{}}\PYG{p}{,} \PYG{n}{pos}\PYG{p}{)}\PYG{p}{;}
    \PYG{n}{start}            \PYG{o}{=} \PYG{n}{json}\PYG{p}{.}\PYG{n}{find}\PYG{p}{(}\PYG{l+s}{\PYGZdq{}}\PYG{l+s+se}{\PYGZbs{}\PYGZdq{}}\PYG{l+s}{\PYGZdq{}}\PYG{p}{,} \PYG{n}{start} \PYG{o}{+} \PYG{l+m+mi}{1}\PYG{p}{)} \PYG{o}{+} \PYG{l+m+mi}{1}\PYG{p}{;}
    \PYG{n}{end}              \PYG{o}{=} \PYG{n}{json}\PYG{p}{.}\PYG{n}{find}\PYG{p}{(}\PYG{l+s}{\PYGZdq{}}\PYG{l+s+se}{\PYGZbs{}\PYGZdq{}}\PYG{l+s}{\PYGZdq{}}\PYG{p}{,} \PYG{n}{start} \PYG{o}{+} \PYG{l+m+mi}{1}\PYG{p}{)}\PYG{p}{;}
    \PYG{n}{std}\PYG{o}{:}\PYG{o}{:}\PYG{n}{string} \PYG{n}{name} \PYG{o}{=} \PYG{n}{json}\PYG{p}{.}\PYG{n}{substr}\PYG{p}{(}\PYG{n}{start}\PYG{p}{,} \PYG{n}{end} \PYG{o}{\PYGZhy{}} \PYG{n}{start}\PYG{p}{)}\PYG{p}{;}
    \PYG{n}{ok}              \PYG{o}{\PYGZam{}}\PYG{o}{=} \PYG{n}{name} \PYG{o}{=}\PYG{o}{=} \PYG{l+s}{\PYGZdq{}}\PYG{l+s}{add}\PYG{l+s}{\PYGZdq{}} \PYG{o}{|}\PYG{o}{|} \PYG{n}{name} \PYG{o}{=}\PYG{o}{=} \PYG{l+s}{\PYGZdq{}}\PYG{l+s}{sin}\PYG{l+s}{\PYGZdq{}}\PYG{p}{;}
    \PYG{c+c1}{//}
    \PYG{k}{return}\PYG{p}{(} \PYG{n}{ok}  \PYG{p}{)}\PYG{p}{;}
\PYG{p}{\PYGZcb{}}
\end{sphinxVerbatim}


\bigskip\hrule\bigskip


xsrst input file: \sphinxcode{\sphinxupquote{example/cplusplus/fun\_json\_xam.cpp}}


\subparagraph{fun\_from\_json\_xam\_cpp}
\label{\detokenize{xsrst/fun_from_json_xam_cpp:fun-from-json-xam-cpp}}\label{\detokenize{xsrst/fun_from_json_xam_cpp::doc}}
\index{fun\_from\_json\_xam\_cpp@\spxentry{fun\_from\_json\_xam\_cpp}}\index{c++:@\spxentry{c++:}}\index{from\_json:@\spxentry{from\_json:}}\index{example@\spxentry{example}}\index{test@\spxentry{test}}\ignorespaces 

\subparagraph{3.2.3.11.2 C++: from\_json: Example and Test}
\label{\detokenize{xsrst/fun_from_json_xam_cpp:c-from-json-example-and-test}}\label{\detokenize{xsrst/fun_from_json_xam_cpp:id1}}
\begin{sphinxVerbatim}[commandchars=\\\{\}]
\PYG{k+kt}{bool} \PYG{n+nf}{from\PYGZus{}json\PYGZus{}xam}\PYG{p}{(}\PYG{k+kt}{void}\PYG{p}{)} \PYG{p}{\PYGZob{}}
    \PYG{k}{using} \PYG{n}{cppad\PYGZus{}py}\PYG{o}{:}\PYG{o}{:}\PYG{n}{a\PYGZus{}double}\PYG{p}{;}
    \PYG{k}{using} \PYG{n}{cppad\PYGZus{}py}\PYG{o}{:}\PYG{o}{:}\PYG{n}{vec\PYGZus{}double}\PYG{p}{;}
    \PYG{k}{using} \PYG{n}{cppad\PYGZus{}py}\PYG{o}{:}\PYG{o}{:}\PYG{n}{vec\PYGZus{}a\PYGZus{}double}\PYG{p}{;}
    \PYG{k}{using} \PYG{n}{cppad\PYGZus{}py}\PYG{o}{:}\PYG{o}{:}\PYG{n}{d\PYGZus{}fun}\PYG{p}{;}
    \PYG{k}{using} \PYG{n}{cppad\PYGZus{}py}\PYG{o}{:}\PYG{o}{:}\PYG{n}{a\PYGZus{}fun}\PYG{p}{;}
    \PYG{c+c1}{//}
    \PYG{c+c1}{// initialize return variable}
    \PYG{k+kt}{bool} \PYG{n}{ok} \PYG{o}{=} \PYG{n+nb}{true}\PYG{p}{;}
    \PYG{c+c1}{// \PYGZhy{}\PYGZhy{}\PYGZhy{}\PYGZhy{}\PYGZhy{}\PYGZhy{}\PYGZhy{}\PYGZhy{}\PYGZhy{}\PYGZhy{}\PYGZhy{}\PYGZhy{}\PYGZhy{}\PYGZhy{}\PYGZhy{}\PYGZhy{}\PYGZhy{}\PYGZhy{}\PYGZhy{}\PYGZhy{}\PYGZhy{}\PYGZhy{}\PYGZhy{}\PYGZhy{}\PYGZhy{}\PYGZhy{}\PYGZhy{}\PYGZhy{}\PYGZhy{}\PYGZhy{}\PYGZhy{}\PYGZhy{}\PYGZhy{}\PYGZhy{}\PYGZhy{}\PYGZhy{}\PYGZhy{}\PYGZhy{}\PYGZhy{}\PYGZhy{}\PYGZhy{}\PYGZhy{}\PYGZhy{}\PYGZhy{}\PYGZhy{}\PYGZhy{}\PYGZhy{}\PYGZhy{}\PYGZhy{}\PYGZhy{}\PYGZhy{}\PYGZhy{}\PYGZhy{}\PYGZhy{}\PYGZhy{}\PYGZhy{}\PYGZhy{}\PYGZhy{}\PYGZhy{}\PYGZhy{}\PYGZhy{}\PYGZhy{}\PYGZhy{}\PYGZhy{}\PYGZhy{}\PYGZhy{}\PYGZhy{}\PYGZhy{}\PYGZhy{}\PYGZhy{}}
    \PYG{c+c1}{// AD graph repersentation of f(x) = sin(x) / cos(x)}
    \PYG{c+c1}{//}
    \PYG{c+c1}{// node\PYGZus{}1 : x[0]}
    \PYG{c+c1}{// node\PYGZus{}2 : sin( x[0] )}
    \PYG{c+c1}{// node\PYGZus{}3 : cos( x[0] )}
    \PYG{c+c1}{// node\PYGZus{}4 : sin( x[0] ) / cos( x[0] )}
    \PYG{c+c1}{// y[0]   = sin( x[0] ) / cos( x[0] )}
    \PYG{c+c1}{// use single quote to avodi having to escape double quote}
    \PYG{n}{std}\PYG{o}{:}\PYG{o}{:}\PYG{n}{string} \PYG{n}{json} \PYG{o}{=}
        \PYG{l+s}{\PYGZdq{}}\PYG{l+s}{\PYGZob{}}\PYG{l+s+se}{\PYGZbs{}n}\PYG{l+s}{\PYGZdq{}}
        \PYG{l+s}{\PYGZdq{}}\PYG{l+s}{   \PYGZsq{}function\PYGZus{}name\PYGZsq{} : \PYGZsq{}tangent function\PYGZsq{},}\PYG{l+s+se}{\PYGZbs{}n}\PYG{l+s}{\PYGZdq{}}
        \PYG{l+s}{\PYGZdq{}}\PYG{l+s}{   \PYGZsq{}op\PYGZus{}define\PYGZus{}vec\PYGZsq{} : [ 3, [}\PYG{l+s+se}{\PYGZbs{}n}\PYG{l+s}{\PYGZdq{}}
        \PYG{l+s}{\PYGZdq{}}\PYG{l+s}{       \PYGZob{} \PYGZsq{}op\PYGZus{}code\PYGZsq{}:1, \PYGZsq{}name\PYGZsq{}:\PYGZsq{}sin\PYGZsq{}, \PYGZsq{}n\PYGZus{}arg\PYGZsq{}:1 \PYGZcb{} ,}\PYG{l+s+se}{\PYGZbs{}n}\PYG{l+s}{\PYGZdq{}}
        \PYG{l+s}{\PYGZdq{}}\PYG{l+s}{       \PYGZob{} \PYGZsq{}op\PYGZus{}code\PYGZsq{}:2, \PYGZsq{}name\PYGZsq{}:\PYGZsq{}cos\PYGZsq{}, \PYGZsq{}n\PYGZus{}arg\PYGZsq{}:1 \PYGZcb{} ,}\PYG{l+s+se}{\PYGZbs{}n}\PYG{l+s}{\PYGZdq{}}
        \PYG{l+s}{\PYGZdq{}}\PYG{l+s}{       \PYGZob{} \PYGZsq{}op\PYGZus{}code\PYGZsq{}:3, \PYGZsq{}name\PYGZsq{}:\PYGZsq{}div\PYGZsq{}, \PYGZsq{}n\PYGZus{}arg\PYGZsq{}:2 \PYGZcb{} ]}\PYG{l+s+se}{\PYGZbs{}n}\PYG{l+s}{\PYGZdq{}}
        \PYG{l+s}{\PYGZdq{}}\PYG{l+s}{   ],}\PYG{l+s+se}{\PYGZbs{}n}\PYG{l+s}{\PYGZdq{}}
        \PYG{l+s}{\PYGZdq{}}\PYG{l+s}{   \PYGZsq{}n\PYGZus{}dynamic\PYGZus{}ind\PYGZsq{}  : 0,}\PYG{l+s+se}{\PYGZbs{}n}\PYG{l+s}{\PYGZdq{}}
        \PYG{l+s}{\PYGZdq{}}\PYG{l+s}{   \PYGZsq{}n\PYGZus{}variable\PYGZus{}ind\PYGZsq{} : 1,}\PYG{l+s+se}{\PYGZbs{}n}\PYG{l+s}{\PYGZdq{}}
        \PYG{l+s}{\PYGZdq{}}\PYG{l+s}{   \PYGZsq{}constant\PYGZus{}vec\PYGZsq{}   : [ 0, [ ] ],}\PYG{l+s+se}{\PYGZbs{}n}\PYG{l+s}{\PYGZdq{}}
        \PYG{l+s}{\PYGZdq{}}\PYG{l+s}{   \PYGZsq{}op\PYGZus{}usage\PYGZus{}vec\PYGZsq{}   : [ 3, [}\PYG{l+s+se}{\PYGZbs{}n}\PYG{l+s}{\PYGZdq{}}
        \PYG{l+s}{\PYGZdq{}}\PYG{l+s}{       [ 1, 1 ]   ,}\PYG{l+s+se}{\PYGZbs{}n}\PYG{l+s}{\PYGZdq{}}  \PYG{c+c1}{// node\PYGZus{}2: sin(x[0])}
        \PYG{l+s}{\PYGZdq{}}\PYG{l+s}{       [ 2, 1 ]   ,}\PYG{l+s+se}{\PYGZbs{}n}\PYG{l+s}{\PYGZdq{}}  \PYG{c+c1}{// node\PYGZus{}3: cos(x[0])}
        \PYG{l+s}{\PYGZdq{}}\PYG{l+s}{       [ 3, 2, 3] ]}\PYG{l+s+se}{\PYGZbs{}n}\PYG{l+s}{\PYGZdq{}}  \PYG{c+c1}{// node\PYGZus{}4: sin(x[0]) / cos(x[0])}
        \PYG{l+s}{\PYGZdq{}}\PYG{l+s}{   ],}\PYG{l+s+se}{\PYGZbs{}n}\PYG{l+s}{\PYGZdq{}}
        \PYG{l+s}{\PYGZdq{}}\PYG{l+s}{   \PYGZsq{}dependent\PYGZus{}vec\PYGZsq{} : [ 1, [4] ]}\PYG{l+s+se}{\PYGZbs{}n}\PYG{l+s}{\PYGZdq{}}
        \PYG{l+s}{\PYGZdq{}}\PYG{l+s}{\PYGZcb{}}\PYG{l+s+se}{\PYGZbs{}n}\PYG{l+s}{\PYGZdq{}}\PYG{p}{;}
    \PYG{c+c1}{// convert the single quote to double quote}
    \PYG{k}{for}\PYG{p}{(}\PYG{k+kt}{size\PYGZus{}t} \PYG{n}{i} \PYG{o}{=} \PYG{l+m+mi}{0}\PYG{p}{;} \PYG{n}{i} \PYG{o}{\PYGZlt{}} \PYG{n}{json}\PYG{p}{.}\PYG{n}{size}\PYG{p}{(}\PYG{p}{)}\PYG{p}{;} \PYG{o}{+}\PYG{o}{+}\PYG{n}{i}\PYG{p}{)}
        \PYG{k}{if}\PYG{p}{(} \PYG{n}{json}\PYG{p}{[}\PYG{n}{i}\PYG{p}{]} \PYG{o}{=}\PYG{o}{=} \PYG{l+s+sc}{\PYGZsq{}}\PYG{l+s+sc}{\PYGZbs{}\PYGZsq{}}\PYG{l+s+sc}{\PYGZsq{}} \PYG{p}{)}
            \PYG{n}{json}\PYG{p}{[}\PYG{n}{i}\PYG{p}{]} \PYG{o}{=} \PYG{l+s+sc}{\PYGZsq{}}\PYG{l+s+sc}{\PYGZdq{}}\PYG{l+s+sc}{\PYGZsq{}}\PYG{p}{;}
    \PYG{c+c1}{//}
    \PYG{c+c1}{// convert json to a fucntion object}
    \PYG{n}{vec\PYGZus{}a\PYGZus{}double} \PYG{n}{ax}\PYG{p}{(}\PYG{l+m+mi}{0}\PYG{p}{)}\PYG{p}{,} \PYG{n}{ay}\PYG{p}{(}\PYG{l+m+mi}{0}\PYG{p}{)}\PYG{p}{;} \PYG{c+c1}{// construct an empty function}
    \PYG{n}{d\PYGZus{}fun} \PYG{n}{f}\PYG{p}{(}\PYG{n}{ax}\PYG{p}{,} \PYG{n}{ay}\PYG{p}{)}\PYG{p}{;}
    \PYG{n}{f}\PYG{p}{.}\PYG{n}{from\PYGZus{}json}\PYG{p}{(}\PYG{n}{json}\PYG{p}{)}\PYG{p}{;}
    \PYG{c+c1}{//}
    \PYG{c+c1}{// compute y = f(x)}
    \PYG{n}{vec\PYGZus{}double} \PYG{n}{x}\PYG{p}{(}\PYG{l+m+mi}{1}\PYG{p}{)}\PYG{p}{,} \PYG{n}{y}\PYG{p}{(}\PYG{l+m+mi}{1}\PYG{p}{)}\PYG{p}{;}
    \PYG{n}{x}\PYG{p}{[}\PYG{l+m+mi}{0}\PYG{p}{]}  \PYG{o}{=} \PYG{l+m+mf}{1.0}\PYG{p}{;}
    \PYG{n}{y}     \PYG{o}{=} \PYG{n}{f}\PYG{p}{.}\PYG{n}{forward}\PYG{p}{(}\PYG{l+m+mi}{0}\PYG{p}{,} \PYG{n}{x}\PYG{p}{)}\PYG{p}{;}
    \PYG{c+c1}{//}
    \PYG{c+c1}{// check the function value}
    \PYG{k+kt}{double} \PYG{n}{eps99}     \PYG{o}{=} \PYG{n}{std}\PYG{o}{:}\PYG{o}{:}\PYG{n}{numeric\PYGZus{}limits}\PYG{o}{\PYGZlt{}}\PYG{k+kt}{double}\PYG{o}{\PYGZgt{}}\PYG{o}{:}\PYG{o}{:}\PYG{n}{epsilon}\PYG{p}{(}\PYG{p}{)}\PYG{p}{;}
    \PYG{k+kt}{double} \PYG{n}{check}     \PYG{o}{=} \PYG{n}{std}\PYG{o}{:}\PYG{o}{:}\PYG{n}{tan}\PYG{p}{(}\PYG{n}{x}\PYG{p}{[}\PYG{l+m+mi}{0}\PYG{p}{]}\PYG{p}{)}\PYG{p}{;}
    \PYG{k+kt}{double} \PYG{n}{rel\PYGZus{}error} \PYG{o}{=} \PYG{n}{y}\PYG{p}{[}\PYG{l+m+mi}{0}\PYG{p}{]} \PYG{o}{/} \PYG{n}{check} \PYG{o}{\PYGZhy{}} \PYG{l+m+mf}{1.0}\PYG{p}{;}
    \PYG{n}{ok}              \PYG{o}{\PYGZam{}}\PYG{o}{=} \PYG{n}{std}\PYG{o}{:}\PYG{o}{:}\PYG{n}{fabs}\PYG{p}{(} \PYG{n}{rel\PYGZus{}error} \PYG{p}{)} \PYG{o}{\PYGZlt{}} \PYG{n}{eps99}\PYG{p}{;}
    \PYG{c+c1}{//}
    \PYG{k}{return} \PYG{n}{ok}\PYG{p}{;}
\PYG{p}{\PYGZcb{}}
\end{sphinxVerbatim}


\bigskip\hrule\bigskip


xsrst input file: \sphinxcode{\sphinxupquote{example/cplusplus/fun\_json\_xam.cpp}}

\index{examples@\spxentry{examples}}\ignorespaces 

\subparagraph{Examples}
\label{\detokenize{xsrst/cpp_fun_json:examples}}\label{\detokenize{xsrst/cpp_fun_json:cpp-fun-json-examples}}\label{\detokenize{xsrst/cpp_fun_json:index-6}}
{\hyperref[\detokenize{xsrst/fun_to_json_xam_cpp:id1}]{\sphinxcrossref{\DUrole{std,std-ref}{fun\_to\_json\_xam\_cpp}}}},
{\hyperref[\detokenize{xsrst/fun_from_json_xam_cpp:id1}]{\sphinxcrossref{\DUrole{std,std-ref}{fun\_from\_json\_xam\_cpp}}}}.


\bigskip\hrule\bigskip


xsrst input file: \sphinxcode{\sphinxupquote{lib/cplusplus/fun.cpp}}
\begin{enumerate}
\sphinxsetlistlabels{\arabic}{enumi}{enumii}{}{.}%
\item {} 
cpp\_independent: {\hyperref[\detokenize{xsrst/cpp_independent:id1}]{\sphinxcrossref{\DUrole{std,std-ref}{3.2.3.1 Declare Independent Variables and Start Recording}}}}

\item {} 
cpp\_abort\_recording: {\hyperref[\detokenize{xsrst/cpp_abort_recording:id1}]{\sphinxcrossref{\DUrole{std,std-ref}{3.2.3.2 Abort Recording}}}}

\item {} 
cpp\_fun\_ctor: {\hyperref[\detokenize{xsrst/cpp_fun_ctor:id1}]{\sphinxcrossref{\DUrole{std,std-ref}{3.2.3.3 Stop Current Recording and Store Function Object}}}}

\item {} 
cpp\_fun\_property: {\hyperref[\detokenize{xsrst/cpp_fun_property:id1}]{\sphinxcrossref{\DUrole{std,std-ref}{3.2.3.4 Properties of a Function Object}}}}

\item {} 
cpp\_fun\_new\_dynamic: {\hyperref[\detokenize{xsrst/cpp_fun_new_dynamic:id1}]{\sphinxcrossref{\DUrole{std,std-ref}{3.2.3.5 Change The Dynamic Parameters}}}}

\item {} 
cpp\_fun\_jacobian: {\hyperref[\detokenize{xsrst/cpp_fun_jacobian:id1}]{\sphinxcrossref{\DUrole{std,std-ref}{3.2.3.6 Jacobian of an AD Function}}}}

\item {} 
cpp\_fun\_hessian: {\hyperref[\detokenize{xsrst/cpp_fun_hessian:id1}]{\sphinxcrossref{\DUrole{std,std-ref}{3.2.3.7 Hessian of an AD Function}}}}

\item {} 
cpp\_fun\_forward: {\hyperref[\detokenize{xsrst/cpp_fun_forward:id1}]{\sphinxcrossref{\DUrole{std,std-ref}{3.2.3.8 Forward Mode AD}}}}

\item {} 
cpp\_fun\_reverse: {\hyperref[\detokenize{xsrst/cpp_fun_reverse:id1}]{\sphinxcrossref{\DUrole{std,std-ref}{3.2.3.9 Reverse Mode AD}}}}

\item {} 
cpp\_fun\_optimize: {\hyperref[\detokenize{xsrst/cpp_fun_optimize:id1}]{\sphinxcrossref{\DUrole{std,std-ref}{3.2.3.10 Optimize an AD Function}}}}

\item {} 
cpp\_fun\_json: {\hyperref[\detokenize{xsrst/cpp_fun_json:id1}]{\sphinxcrossref{\DUrole{std,std-ref}{3.2.3.11 Json Representation of AD Computational Graph}}}}

\end{enumerate}


\bigskip\hrule\bigskip


xsrst input file: \sphinxcode{\sphinxupquote{lib/cplusplus/cplusplus.xsrst}}


\subparagraph{sparse}
\label{\detokenize{xsrst/sparse:sparse}}\label{\detokenize{xsrst/sparse::doc}}
\index{sparse@\spxentry{sparse}}\index{cppad@\spxentry{cppad}}\index{py@\spxentry{py}}\index{sparse@\spxentry{sparse}}\index{calculation@\spxentry{calculation}}\ignorespaces 

\subparagraph{3.2.4 Cppad Py Sparse Calculation}
\label{\detokenize{xsrst/sparse:cppad-py-sparse-calculation}}\label{\detokenize{xsrst/sparse:id1}}\begin{itemize}
\item {} 
{\hyperref[\detokenize{xsrst/sparse:sparse-children}]{\sphinxcrossref{\DUrole{std,std-ref}{Children}}}}

\end{itemize}

\index{children@\spxentry{children}}\ignorespaces 

\subparagraph{Children}
\label{\detokenize{xsrst/sparse:children}}\label{\detokenize{xsrst/sparse:sparse-children}}\label{\detokenize{xsrst/sparse:index-1}}

\subparagraph{cpp\_sparse\_rc}
\label{\detokenize{xsrst/cpp_sparse_rc:cpp-sparse-rc}}\label{\detokenize{xsrst/cpp_sparse_rc::doc}}
\index{cpp\_sparse\_rc@\spxentry{cpp\_sparse\_rc}}\index{sparsity@\spxentry{sparsity}}\index{patterns@\spxentry{patterns}}\ignorespaces 

\subparagraph{3.2.4.1 Sparsity Patterns}
\label{\detokenize{xsrst/cpp_sparse_rc:sparsity-patterns}}\label{\detokenize{xsrst/cpp_sparse_rc:id1}}\begin{itemize}
\item {} 
{\hyperref[\detokenize{xsrst/cpp_sparse_rc:cpp-sparse-rc-syntax}]{\sphinxcrossref{\DUrole{std,std-ref}{Syntax}}}}

\item {} 
{\hyperref[\detokenize{xsrst/cpp_sparse_rc:cpp-sparse-rc-pattern}]{\sphinxcrossref{\DUrole{std,std-ref}{pattern}}}}

\item {} 
{\hyperref[\detokenize{xsrst/cpp_sparse_rc:cpp-sparse-rc-nr}]{\sphinxcrossref{\DUrole{std,std-ref}{nr}}}}

\item {} 
{\hyperref[\detokenize{xsrst/cpp_sparse_rc:cpp-sparse-rc-nc}]{\sphinxcrossref{\DUrole{std,std-ref}{nc}}}}

\item {} 
{\hyperref[\detokenize{xsrst/cpp_sparse_rc:cpp-sparse-rc-nnz}]{\sphinxcrossref{\DUrole{std,std-ref}{nnz}}}}

\item {} 
{\hyperref[\detokenize{xsrst/cpp_sparse_rc:cpp-sparse-rc-resize}]{\sphinxcrossref{\DUrole{std,std-ref}{resize}}}}

\item {} \begin{description}
\item[{{\hyperref[\detokenize{xsrst/cpp_sparse_rc:cpp-sparse-rc-put}]{\sphinxcrossref{\DUrole{std,std-ref}{put}}}}}] \leavevmode\begin{itemize}
\item {} 
{\hyperref[\detokenize{xsrst/cpp_sparse_rc:cpp-sparse-rc-put-k}]{\sphinxcrossref{\DUrole{std,std-ref}{k}}}}

\item {} 
{\hyperref[\detokenize{xsrst/cpp_sparse_rc:cpp-sparse-rc-put-r}]{\sphinxcrossref{\DUrole{std,std-ref}{r}}}}

\item {} 
{\hyperref[\detokenize{xsrst/cpp_sparse_rc:cpp-sparse-rc-put-c}]{\sphinxcrossref{\DUrole{std,std-ref}{c}}}}

\end{itemize}

\end{description}

\item {} 
{\hyperref[\detokenize{xsrst/cpp_sparse_rc:cpp-sparse-rc-row}]{\sphinxcrossref{\DUrole{std,std-ref}{row}}}}

\item {} 
{\hyperref[\detokenize{xsrst/cpp_sparse_rc:cpp-sparse-rc-col}]{\sphinxcrossref{\DUrole{std,std-ref}{col}}}}

\item {} 
{\hyperref[\detokenize{xsrst/cpp_sparse_rc:cpp-sparse-rc-row-major}]{\sphinxcrossref{\DUrole{std,std-ref}{row\_major}}}}

\item {} 
{\hyperref[\detokenize{xsrst/cpp_sparse_rc:cpp-sparse-rc-col-major}]{\sphinxcrossref{\DUrole{std,std-ref}{col\_major}}}}

\item {} 
{\hyperref[\detokenize{xsrst/cpp_sparse_rc:cpp-sparse-rc-example}]{\sphinxcrossref{\DUrole{std,std-ref}{Example}}}}

\end{itemize}

\index{syntax@\spxentry{syntax}}\ignorespaces 

\subparagraph{Syntax}
\label{\detokenize{xsrst/cpp_sparse_rc:syntax}}\label{\detokenize{xsrst/cpp_sparse_rc:cpp-sparse-rc-syntax}}\label{\detokenize{xsrst/cpp_sparse_rc:index-1}}
\begin{DUlineblock}{0em}
\item[] \sphinxstyleemphasis{pattern} =  \sphinxcode{\sphinxupquote{cppad\_py::sparse\_rc}} ()
\item[] \sphinxstyleemphasis{pattern} . \sphinxcode{\sphinxupquote{resize}} ( \sphinxstyleemphasis{nr} , \sphinxstyleemphasis{nc} , \sphinxstyleemphasis{nnz} )
\item[] \sphinxstyleemphasis{nr} = \sphinxstyleemphasis{pattern} . \sphinxcode{\sphinxupquote{nr}} ()
\item[] \sphinxstyleemphasis{nc} = \sphinxstyleemphasis{pattern} . \sphinxcode{\sphinxupquote{nc}} ()
\item[] \sphinxstyleemphasis{nnz} = \sphinxstyleemphasis{pattern} . \sphinxcode{\sphinxupquote{nnz}} ()
\item[] \sphinxstyleemphasis{pattern} . \sphinxcode{\sphinxupquote{put}} ( \sphinxstyleemphasis{k} , \sphinxstyleemphasis{r} , \sphinxstyleemphasis{c} )
\item[] \sphinxstyleemphasis{row} = \sphinxstyleemphasis{pattern} . \sphinxcode{\sphinxupquote{row}} ()
\item[] \sphinxstyleemphasis{col} = \sphinxstyleemphasis{pattern} . \sphinxcode{\sphinxupquote{col}} ()
\item[] \sphinxstyleemphasis{row\_major} = \sphinxstyleemphasis{pattern} . \sphinxcode{\sphinxupquote{row\_major}} ()
\item[] \sphinxstyleemphasis{col\_major} = \sphinxstyleemphasis{pattern} . \sphinxcode{\sphinxupquote{col\_major}} ()
\end{DUlineblock}

\index{pattern@\spxentry{pattern}}\ignorespaces 

\subparagraph{pattern}
\label{\detokenize{xsrst/cpp_sparse_rc:pattern}}\label{\detokenize{xsrst/cpp_sparse_rc:cpp-sparse-rc-pattern}}\label{\detokenize{xsrst/cpp_sparse_rc:index-2}}
This result has prototype

\begin{DUlineblock}{0em}
\item[]         \sphinxcode{\sphinxupquote{sparse\_rc}} \sphinxstyleemphasis{pattern}
\end{DUlineblock}

It is used to hold a sparsity pattern for a matrix.
The sparsity \sphinxstyleemphasis{pattern} is \sphinxcode{\sphinxupquote{const}}
except during the \sphinxcode{\sphinxupquote{resize}} and \sphinxcode{\sphinxupquote{put}} operations.

\index{nr@\spxentry{nr}}\ignorespaces 

\subparagraph{nr}
\label{\detokenize{xsrst/cpp_sparse_rc:nr}}\label{\detokenize{xsrst/cpp_sparse_rc:cpp-sparse-rc-nr}}\label{\detokenize{xsrst/cpp_sparse_rc:index-3}}
This argument has prototype

\begin{DUlineblock}{0em}
\item[]         \sphinxcode{\sphinxupquote{int}} \sphinxstyleemphasis{nr}
\end{DUlineblock}

is non\sphinxhyphen{}negative, and specifies
the number of rows in the sparsity pattern.
The function \sphinxcode{\sphinxupquote{nr()}} returns the value of
\sphinxstyleemphasis{nr} in the previous \sphinxcode{\sphinxupquote{resize}} operation.

\index{nc@\spxentry{nc}}\ignorespaces 

\subparagraph{nc}
\label{\detokenize{xsrst/cpp_sparse_rc:nc}}\label{\detokenize{xsrst/cpp_sparse_rc:cpp-sparse-rc-nc}}\label{\detokenize{xsrst/cpp_sparse_rc:index-4}}
This argument has prototype

\begin{DUlineblock}{0em}
\item[]         \sphinxcode{\sphinxupquote{int}} \sphinxstyleemphasis{nc}
\end{DUlineblock}

is non\sphinxhyphen{}negative, and specifies
the number of columns in the sparsity pattern.
The function \sphinxcode{\sphinxupquote{nc()}} returns the value of
\sphinxstyleemphasis{nc} in the previous \sphinxcode{\sphinxupquote{resize}} operation.

\index{nnz@\spxentry{nnz}}\ignorespaces 

\subparagraph{nnz}
\label{\detokenize{xsrst/cpp_sparse_rc:nnz}}\label{\detokenize{xsrst/cpp_sparse_rc:cpp-sparse-rc-nnz}}\label{\detokenize{xsrst/cpp_sparse_rc:index-5}}
This argument and result has prototype

\begin{DUlineblock}{0em}
\item[]         \sphinxcode{\sphinxupquote{int}} \sphinxstyleemphasis{nnz}
\end{DUlineblock}

It is the number of possibly non\sphinxhyphen{}zero
index pairs in the sparsity pattern.
The function \sphinxcode{\sphinxupquote{nnz()}} returns the value of
\sphinxstyleemphasis{nnz} in the previous \sphinxcode{\sphinxupquote{resize}} operation.

\index{resize@\spxentry{resize}}\ignorespaces 

\subparagraph{resize}
\label{\detokenize{xsrst/cpp_sparse_rc:resize}}\label{\detokenize{xsrst/cpp_sparse_rc:cpp-sparse-rc-resize}}\label{\detokenize{xsrst/cpp_sparse_rc:index-6}}
The current sparsity pattern is lost and a new one is started
with the specified parameters.
After each \sphinxcode{\sphinxupquote{resize}} , the elements in the \sphinxstyleemphasis{row}
and \sphinxstyleemphasis{col} vectors should be assigned using \sphinxcode{\sphinxupquote{put}} .

\index{put@\spxentry{put}}\ignorespaces 

\subparagraph{put}
\label{\detokenize{xsrst/cpp_sparse_rc:put}}\label{\detokenize{xsrst/cpp_sparse_rc:cpp-sparse-rc-put}}\label{\detokenize{xsrst/cpp_sparse_rc:index-7}}
This function sets the values

\begin{DUlineblock}{0em}
\item[]         \sphinxstyleemphasis{row} {[} \sphinxstyleemphasis{k} {]} = \sphinxstyleemphasis{r}
\item[]         \sphinxstyleemphasis{col} {[} \sphinxstyleemphasis{k} {]} = \sphinxstyleemphasis{c}
\end{DUlineblock}

(The name \sphinxcode{\sphinxupquote{set}} is used by Cppad, but not used here,
because \sphinxcode{\sphinxupquote{set}} it is a built\sphinxhyphen{}in name in Python.)

\index{k@\spxentry{k}}\ignorespaces 

\subparagraph{k}
\label{\detokenize{xsrst/cpp_sparse_rc:k}}\label{\detokenize{xsrst/cpp_sparse_rc:cpp-sparse-rc-put-k}}\label{\detokenize{xsrst/cpp_sparse_rc:index-8}}
This argument has type

\begin{DUlineblock}{0em}
\item[]         \sphinxcode{\sphinxupquote{int}} \sphinxstyleemphasis{k}
\end{DUlineblock}

is non\sphinxhyphen{}negative,
and must be less than \sphinxstyleemphasis{nnz} .

\index{r@\spxentry{r}}\ignorespaces 

\subparagraph{r}
\label{\detokenize{xsrst/cpp_sparse_rc:r}}\label{\detokenize{xsrst/cpp_sparse_rc:cpp-sparse-rc-put-r}}\label{\detokenize{xsrst/cpp_sparse_rc:index-9}}
This argument has type

\begin{DUlineblock}{0em}
\item[]         \sphinxcode{\sphinxupquote{int}} \sphinxstyleemphasis{r}
\end{DUlineblock}

is non\sphinxhyphen{}negative, and must be less than \sphinxstyleemphasis{nr} .
It specifies the value assigned to \sphinxstyleemphasis{row} {[} \sphinxstyleemphasis{k} {]} and must

\index{c@\spxentry{c}}\ignorespaces 

\subparagraph{c}
\label{\detokenize{xsrst/cpp_sparse_rc:c}}\label{\detokenize{xsrst/cpp_sparse_rc:cpp-sparse-rc-put-c}}\label{\detokenize{xsrst/cpp_sparse_rc:index-10}}
This argument has type

\begin{DUlineblock}{0em}
\item[]         \sphinxcode{\sphinxupquote{int}} \sphinxstyleemphasis{c}
\end{DUlineblock}

It specifies the value assigned to \sphinxstyleemphasis{col} {[} \sphinxstyleemphasis{k} {]} and must
be less than \sphinxstyleemphasis{nc} .

\index{row@\spxentry{row}}\ignorespaces 

\subparagraph{row}
\label{\detokenize{xsrst/cpp_sparse_rc:row}}\label{\detokenize{xsrst/cpp_sparse_rc:cpp-sparse-rc-row}}\label{\detokenize{xsrst/cpp_sparse_rc:index-11}}
This result has type

\begin{DUlineblock}{0em}
\item[]         \sphinxcode{\sphinxupquote{vec\_int}} \sphinxstyleemphasis{row}
\end{DUlineblock}

and its size is \sphinxstyleemphasis{nnz} .
For \sphinxstyleemphasis{k} = 0, … , \sphinxstyleemphasis{nnz} \sphinxhyphen{}1 ,
\sphinxstyleemphasis{row} {[} \sphinxstyleemphasis{k} {]} is the row index for the \sphinxstyleemphasis{k}\sphinxhyphen{}th possibly non\sphinxhyphen{}zero
entry in the matrix.

\index{col@\spxentry{col}}\ignorespaces 

\subparagraph{col}
\label{\detokenize{xsrst/cpp_sparse_rc:col}}\label{\detokenize{xsrst/cpp_sparse_rc:cpp-sparse-rc-col}}\label{\detokenize{xsrst/cpp_sparse_rc:index-12}}
This result has type

\begin{DUlineblock}{0em}
\item[]         \sphinxcode{\sphinxupquote{vec\_int}} \sphinxstyleemphasis{col}
\end{DUlineblock}

and its size is \sphinxstyleemphasis{nnz} .
For \sphinxstyleemphasis{k} = 0, … , \sphinxstyleemphasis{nnz} \sphinxhyphen{}1 ,
\sphinxstyleemphasis{col} {[} \sphinxstyleemphasis{k} {]} is the column index for the \sphinxstyleemphasis{k}\sphinxhyphen{}th possibly non\sphinxhyphen{}zero
entry in the matrix.

\index{row\_major@\spxentry{row\_major}}\ignorespaces 

\subparagraph{row\_major}
\label{\detokenize{xsrst/cpp_sparse_rc:row-major}}\label{\detokenize{xsrst/cpp_sparse_rc:cpp-sparse-rc-row-major}}\label{\detokenize{xsrst/cpp_sparse_rc:index-13}}
This vector has prototype

\begin{DUlineblock}{0em}
\item[]         \sphinxcode{\sphinxupquote{vec\_int}} \sphinxstyleemphasis{row\_major}
\end{DUlineblock}

and its size \sphinxstyleemphasis{nnz} .
It sorts the sparsity pattern in row\sphinxhyphen{}major order.
To be specific,

\begin{DUlineblock}{0em}
\item[]         \sphinxstyleemphasis{col} {[} \sphinxstyleemphasis{row\_major} {[} \sphinxstyleemphasis{k} {]} {]} \textless{}= \sphinxstyleemphasis{col} {[} \sphinxstyleemphasis{row\_major} {[} \sphinxstyleemphasis{k} +1{]} {]}
\end{DUlineblock}

and if \sphinxstyleemphasis{col} {[} \sphinxstyleemphasis{row\_major} {[} \sphinxstyleemphasis{k} {]} {]} == \sphinxstyleemphasis{col} {[} \sphinxstyleemphasis{row\_major} {[} \sphinxstyleemphasis{k} +1{]} {]} ,

\begin{DUlineblock}{0em}
\item[]         \sphinxstyleemphasis{row} {[} \sphinxstyleemphasis{row\_major} {[} \sphinxstyleemphasis{k} {]} {]} \textless{} \sphinxstyleemphasis{row} {[} \sphinxstyleemphasis{row\_major} {[} \sphinxstyleemphasis{k} +1{]} {]}
\end{DUlineblock}

This routine generates an assert if there are two entries with the same
row and column values (if \sphinxcode{\sphinxupquote{NDEBUG}} is not defined).

\index{col\_major@\spxentry{col\_major}}\ignorespaces 

\subparagraph{col\_major}
\label{\detokenize{xsrst/cpp_sparse_rc:col-major}}\label{\detokenize{xsrst/cpp_sparse_rc:cpp-sparse-rc-col-major}}\label{\detokenize{xsrst/cpp_sparse_rc:index-14}}
This vector has prototype

\begin{DUlineblock}{0em}
\item[]         \sphinxcode{\sphinxupquote{vec\_int}} \sphinxstyleemphasis{col\_major}
\end{DUlineblock}

and its size \sphinxstyleemphasis{nnz} .
It sorts the sparsity pattern in column\sphinxhyphen{}major order.
To be specific,

\begin{DUlineblock}{0em}
\item[]         \sphinxstyleemphasis{row} {[} \sphinxstyleemphasis{col\_major} {[} \sphinxstyleemphasis{k} {]} {]} \textless{}= \sphinxstyleemphasis{row} {[} \sphinxstyleemphasis{col\_major} {[} \sphinxstyleemphasis{k} +1{]} {]}
\end{DUlineblock}

and if \sphinxstyleemphasis{row} {[} \sphinxstyleemphasis{col\_major} {[} \sphinxstyleemphasis{k} {]} {]} == \sphinxstyleemphasis{row} {[} \sphinxstyleemphasis{col\_major} {[} \sphinxstyleemphasis{k} +1{]} {]} ,

\begin{DUlineblock}{0em}
\item[]         \sphinxstyleemphasis{col} {[} \sphinxstyleemphasis{col\_major} {[} \sphinxstyleemphasis{k} {]} {]} \textless{} \sphinxstyleemphasis{col} {[} \sphinxstyleemphasis{col\_major} {[} \sphinxstyleemphasis{k} +1{]} {]}
\end{DUlineblock}

This routine generates an assert if there are two entries with the same
row and column values (if \sphinxcode{\sphinxupquote{NDEBUG}} is not defined).


\subparagraph{sparse\_rc\_xam\_cpp}
\label{\detokenize{xsrst/sparse_rc_xam_cpp:sparse-rc-xam-cpp}}\label{\detokenize{xsrst/sparse_rc_xam_cpp::doc}}
\index{sparse\_rc\_xam\_cpp@\spxentry{sparse\_rc\_xam\_cpp}}\index{c++:@\spxentry{c++:}}\index{sparsity@\spxentry{sparsity}}\index{patterns:@\spxentry{patterns:}}\index{example@\spxentry{example}}\index{test@\spxentry{test}}\ignorespaces 

\subparagraph{3.2.4.1.1 C++: Sparsity Patterns: Example and Test}
\label{\detokenize{xsrst/sparse_rc_xam_cpp:c-sparsity-patterns-example-and-test}}\label{\detokenize{xsrst/sparse_rc_xam_cpp:id1}}
\begin{sphinxVerbatim}[commandchars=\\\{\}]
\PYG{c+cp}{\PYGZsh{}}\PYG{c+cp}{ include \PYGZlt{}cstdio\PYGZgt{}}
\PYG{c+cp}{\PYGZsh{}}\PYG{c+cp}{ include \PYGZlt{}cppad}\PYG{c+cp}{/}\PYG{c+cp}{py}\PYG{c+cp}{/}\PYG{c+cp}{cppad\PYGZus{}py.hpp\PYGZgt{}}

\PYG{k+kt}{bool} \PYG{n+nf}{sparse\PYGZus{}rc\PYGZus{}xam}\PYG{p}{(}\PYG{k+kt}{void}\PYG{p}{)} \PYG{p}{\PYGZob{}}
    \PYG{k}{using} \PYG{n}{cppad\PYGZus{}py}\PYG{o}{:}\PYG{o}{:}\PYG{n}{vec\PYGZus{}int}\PYG{p}{;}
    \PYG{k}{using} \PYG{n}{cppad\PYGZus{}py}\PYG{o}{:}\PYG{o}{:}\PYG{n}{sparse\PYGZus{}rc}\PYG{p}{;}
    \PYG{c+c1}{//}
    \PYG{c+c1}{// initialize return variable}
    \PYG{k+kt}{bool} \PYG{n}{ok} \PYG{o}{=} \PYG{n+nb}{true}\PYG{p}{;}
    \PYG{c+c1}{//\PYGZhy{}\PYGZhy{}\PYGZhy{}\PYGZhy{}\PYGZhy{}\PYGZhy{}\PYGZhy{}\PYGZhy{}\PYGZhy{}\PYGZhy{}\PYGZhy{}\PYGZhy{}\PYGZhy{}\PYGZhy{}\PYGZhy{}\PYGZhy{}\PYGZhy{}\PYGZhy{}\PYGZhy{}\PYGZhy{}\PYGZhy{}\PYGZhy{}\PYGZhy{}\PYGZhy{}\PYGZhy{}\PYGZhy{}\PYGZhy{}\PYGZhy{}\PYGZhy{}\PYGZhy{}\PYGZhy{}\PYGZhy{}\PYGZhy{}\PYGZhy{}\PYGZhy{}\PYGZhy{}\PYGZhy{}\PYGZhy{}\PYGZhy{}\PYGZhy{}\PYGZhy{}\PYGZhy{}\PYGZhy{}\PYGZhy{}\PYGZhy{}\PYGZhy{}\PYGZhy{}\PYGZhy{}\PYGZhy{}\PYGZhy{}\PYGZhy{}\PYGZhy{}\PYGZhy{}\PYGZhy{}\PYGZhy{}\PYGZhy{}\PYGZhy{}\PYGZhy{}\PYGZhy{}\PYGZhy{}\PYGZhy{}\PYGZhy{}\PYGZhy{}\PYGZhy{}\PYGZhy{}\PYGZhy{}\PYGZhy{}\PYGZhy{}\PYGZhy{}\PYGZhy{}\PYGZhy{}\PYGZhy{}}
    \PYG{c+c1}{//}
    \PYG{c+c1}{// create a sparsity pattern}
    \PYG{n}{sparse\PYGZus{}rc} \PYG{n}{pattern} \PYG{o}{=} \PYG{n}{sparse\PYGZus{}rc}\PYG{p}{(}\PYG{p}{)}\PYG{p}{;}
    \PYG{c+c1}{//}
    \PYG{k+kt}{int} \PYG{n}{nr} \PYG{o}{=} \PYG{l+m+mi}{6}\PYG{p}{;}
    \PYG{k+kt}{int} \PYG{n}{nc} \PYG{o}{=} \PYG{l+m+mi}{5}\PYG{p}{;}
    \PYG{k+kt}{int} \PYG{n}{nnz} \PYG{o}{=} \PYG{l+m+mi}{4}\PYG{p}{;}
    \PYG{c+c1}{//}
    \PYG{c+c1}{// resize}
    \PYG{n}{pattern}\PYG{p}{.}\PYG{n}{resize}\PYG{p}{(}\PYG{n}{nr}\PYG{p}{,} \PYG{n}{nc}\PYG{p}{,} \PYG{n}{nnz}\PYG{p}{)}\PYG{p}{;}
    \PYG{c+c1}{//}
    \PYG{n}{ok} \PYG{o}{=} \PYG{n}{ok} \PYG{o}{\PYGZam{}}\PYG{o}{\PYGZam{}} \PYG{n}{pattern}\PYG{p}{.}\PYG{n}{nr}\PYG{p}{(}\PYG{p}{)}  \PYG{o}{=}\PYG{o}{=} \PYG{n}{nr} \PYG{p}{;}
    \PYG{n}{ok} \PYG{o}{=} \PYG{n}{ok} \PYG{o}{\PYGZam{}}\PYG{o}{\PYGZam{}} \PYG{n}{pattern}\PYG{p}{.}\PYG{n}{nc}\PYG{p}{(}\PYG{p}{)}  \PYG{o}{=}\PYG{o}{=} \PYG{n}{nc} \PYG{p}{;}
    \PYG{n}{ok} \PYG{o}{=} \PYG{n}{ok} \PYG{o}{\PYGZam{}}\PYG{o}{\PYGZam{}} \PYG{n}{pattern}\PYG{p}{.}\PYG{n}{nnz}\PYG{p}{(}\PYG{p}{)} \PYG{o}{=}\PYG{o}{=} \PYG{n}{nnz} \PYG{p}{;}
    \PYG{c+c1}{//}
    \PYG{c+c1}{// indices corresponding to upper\PYGZhy{}diagonal}
    \PYG{k}{for}\PYG{p}{(}\PYG{k+kt}{int} \PYG{n}{k} \PYG{o}{=} \PYG{l+m+mi}{0}\PYG{p}{;} \PYG{n}{k} \PYG{o}{\PYGZlt{}} \PYG{n}{nnz}\PYG{p}{;} \PYG{n}{k}\PYG{o}{+}\PYG{o}{+}\PYG{p}{)} \PYG{p}{\PYGZob{}}
        \PYG{n}{pattern}\PYG{p}{.}\PYG{n}{put}\PYG{p}{(}\PYG{n}{k}\PYG{p}{,} \PYG{n}{k}\PYG{p}{,} \PYG{n}{k}\PYG{o}{+}\PYG{l+m+mi}{1}\PYG{p}{)}\PYG{p}{;}
    \PYG{p}{\PYGZcb{}}
    \PYG{c+c1}{//}
    \PYG{c+c1}{// row and column indices}
    \PYG{n}{vec\PYGZus{}int} \PYG{n}{row} \PYG{o}{=} \PYG{n}{pattern}\PYG{p}{.}\PYG{n}{row}\PYG{p}{(}\PYG{p}{)}\PYG{p}{;}
    \PYG{n}{vec\PYGZus{}int} \PYG{n}{col} \PYG{o}{=} \PYG{n}{pattern}\PYG{p}{.}\PYG{n}{col}\PYG{p}{(}\PYG{p}{)}\PYG{p}{;}
    \PYG{c+c1}{//}
    \PYG{c+c1}{// check row and column indices}
    \PYG{k}{for}\PYG{p}{(}\PYG{k+kt}{int} \PYG{n}{k} \PYG{o}{=} \PYG{l+m+mi}{0}\PYG{p}{;} \PYG{n}{k} \PYG{o}{\PYGZlt{}} \PYG{n}{nnz}\PYG{p}{;} \PYG{n}{k}\PYG{o}{+}\PYG{o}{+}\PYG{p}{)} \PYG{p}{\PYGZob{}}
        \PYG{n}{ok} \PYG{o}{=} \PYG{n}{ok} \PYG{o}{\PYGZam{}}\PYG{o}{\PYGZam{}} \PYG{n}{row}\PYG{p}{[}\PYG{n}{k}\PYG{p}{]} \PYG{o}{=}\PYG{o}{=} \PYG{n}{k}\PYG{p}{;}
        \PYG{n}{ok} \PYG{o}{=} \PYG{n}{ok} \PYG{o}{\PYGZam{}}\PYG{o}{\PYGZam{}} \PYG{n}{col}\PYG{p}{[}\PYG{n}{k}\PYG{p}{]} \PYG{o}{=}\PYG{o}{=} \PYG{n}{k}\PYG{o}{+}\PYG{l+m+mi}{1}\PYG{p}{;}
    \PYG{p}{\PYGZcb{}}
    \PYG{c+c1}{//}
    \PYG{c+c1}{// For this case, row\PYGZus{}major and col\PYGZus{}major order are same as original order}
    \PYG{n}{vec\PYGZus{}int} \PYG{n}{row\PYGZus{}major} \PYG{o}{=} \PYG{n}{pattern}\PYG{p}{.}\PYG{n}{row\PYGZus{}major}\PYG{p}{(}\PYG{p}{)}\PYG{p}{;}
    \PYG{n}{vec\PYGZus{}int} \PYG{n}{col\PYGZus{}major} \PYG{o}{=} \PYG{n}{pattern}\PYG{p}{.}\PYG{n}{col\PYGZus{}major}\PYG{p}{(}\PYG{p}{)}\PYG{p}{;}
    \PYG{k}{for}\PYG{p}{(}\PYG{k+kt}{int} \PYG{n}{k} \PYG{o}{=} \PYG{l+m+mi}{0}\PYG{p}{;} \PYG{n}{k} \PYG{o}{\PYGZlt{}} \PYG{n}{nnz}\PYG{p}{;} \PYG{n}{k}\PYG{o}{+}\PYG{o}{+}\PYG{p}{)} \PYG{p}{\PYGZob{}}
        \PYG{n}{ok} \PYG{o}{=} \PYG{n}{ok} \PYG{o}{\PYGZam{}}\PYG{o}{\PYGZam{}} \PYG{n}{row\PYGZus{}major}\PYG{p}{[}\PYG{n}{k}\PYG{p}{]} \PYG{o}{=}\PYG{o}{=} \PYG{n}{k}\PYG{p}{;}
        \PYG{n}{ok} \PYG{o}{=} \PYG{n}{ok} \PYG{o}{\PYGZam{}}\PYG{o}{\PYGZam{}} \PYG{n}{col\PYGZus{}major}\PYG{p}{[}\PYG{n}{k}\PYG{p}{]} \PYG{o}{=}\PYG{o}{=} \PYG{n}{k}\PYG{p}{;}
    \PYG{p}{\PYGZcb{}}
    \PYG{c+c1}{//}
    \PYG{k}{return}\PYG{p}{(} \PYG{n}{ok} \PYG{p}{)}\PYG{p}{;}
\PYG{p}{\PYGZcb{}}
\end{sphinxVerbatim}


\bigskip\hrule\bigskip


xsrst input file: \sphinxcode{\sphinxupquote{example/cplusplus/sparse\_rc\_xam.cpp}}

\index{example@\spxentry{example}}\ignorespaces 

\subparagraph{Example}
\label{\detokenize{xsrst/cpp_sparse_rc:example}}\label{\detokenize{xsrst/cpp_sparse_rc:cpp-sparse-rc-example}}\label{\detokenize{xsrst/cpp_sparse_rc:index-15}}
{\hyperref[\detokenize{xsrst/sparse_rc_xam_cpp:id1}]{\sphinxcrossref{\DUrole{std,std-ref}{sparse\_rc\_xam\_cpp}}}}


\bigskip\hrule\bigskip


xsrst input file: \sphinxcode{\sphinxupquote{lib/cplusplus/sparse.cpp}}


\subparagraph{cpp\_sparse\_rcv}
\label{\detokenize{xsrst/cpp_sparse_rcv:cpp-sparse-rcv}}\label{\detokenize{xsrst/cpp_sparse_rcv::doc}}
\index{cpp\_sparse\_rcv@\spxentry{cpp\_sparse\_rcv}}\index{sparse@\spxentry{sparse}}\index{matrices@\spxentry{matrices}}\ignorespaces 

\subparagraph{3.2.4.2 Sparse Matrices}
\label{\detokenize{xsrst/cpp_sparse_rcv:sparse-matrices}}\label{\detokenize{xsrst/cpp_sparse_rcv:id1}}\begin{itemize}
\item {} 
{\hyperref[\detokenize{xsrst/cpp_sparse_rcv:cpp-sparse-rcv-syntax}]{\sphinxcrossref{\DUrole{std,std-ref}{Syntax}}}}

\item {} 
{\hyperref[\detokenize{xsrst/cpp_sparse_rcv:cpp-sparse-rcv-pattern}]{\sphinxcrossref{\DUrole{std,std-ref}{pattern}}}}

\item {} 
{\hyperref[\detokenize{xsrst/cpp_sparse_rcv:cpp-sparse-rcv-matrix}]{\sphinxcrossref{\DUrole{std,std-ref}{matrix}}}}

\item {} 
{\hyperref[\detokenize{xsrst/cpp_sparse_rcv:cpp-sparse-rcv-nr}]{\sphinxcrossref{\DUrole{std,std-ref}{nr}}}}

\item {} 
{\hyperref[\detokenize{xsrst/cpp_sparse_rcv:cpp-sparse-rcv-nc}]{\sphinxcrossref{\DUrole{std,std-ref}{nc}}}}

\item {} 
{\hyperref[\detokenize{xsrst/cpp_sparse_rcv:cpp-sparse-rcv-nnz}]{\sphinxcrossref{\DUrole{std,std-ref}{nnz}}}}

\item {} \begin{description}
\item[{{\hyperref[\detokenize{xsrst/cpp_sparse_rcv:cpp-sparse-rcv-put}]{\sphinxcrossref{\DUrole{std,std-ref}{put}}}}}] \leavevmode\begin{itemize}
\item {} 
{\hyperref[\detokenize{xsrst/cpp_sparse_rcv:cpp-sparse-rcv-put-k}]{\sphinxcrossref{\DUrole{std,std-ref}{k}}}}

\item {} 
{\hyperref[\detokenize{xsrst/cpp_sparse_rcv:cpp-sparse-rcv-put-v}]{\sphinxcrossref{\DUrole{std,std-ref}{v}}}}

\end{itemize}

\end{description}

\item {} 
{\hyperref[\detokenize{xsrst/cpp_sparse_rcv:cpp-sparse-rcv-row}]{\sphinxcrossref{\DUrole{std,std-ref}{row}}}}

\item {} 
{\hyperref[\detokenize{xsrst/cpp_sparse_rcv:cpp-sparse-rcv-col}]{\sphinxcrossref{\DUrole{std,std-ref}{col}}}}

\item {} 
{\hyperref[\detokenize{xsrst/cpp_sparse_rcv:cpp-sparse-rcv-val}]{\sphinxcrossref{\DUrole{std,std-ref}{val}}}}

\item {} 
{\hyperref[\detokenize{xsrst/cpp_sparse_rcv:cpp-sparse-rcv-row-major}]{\sphinxcrossref{\DUrole{std,std-ref}{row\_major}}}}

\item {} 
{\hyperref[\detokenize{xsrst/cpp_sparse_rcv:cpp-sparse-rcv-col-major}]{\sphinxcrossref{\DUrole{std,std-ref}{col\_major}}}}

\item {} 
{\hyperref[\detokenize{xsrst/cpp_sparse_rcv:cpp-sparse-rcv-example}]{\sphinxcrossref{\DUrole{std,std-ref}{Example}}}}

\end{itemize}

\index{syntax@\spxentry{syntax}}\ignorespaces 

\subparagraph{Syntax}
\label{\detokenize{xsrst/cpp_sparse_rcv:syntax}}\label{\detokenize{xsrst/cpp_sparse_rcv:cpp-sparse-rcv-syntax}}\label{\detokenize{xsrst/cpp_sparse_rcv:index-1}}
\begin{DUlineblock}{0em}
\item[] \sphinxstyleemphasis{matrix} =  \sphinxcode{\sphinxupquote{cppad\_py::sparse\_rcv}} ( \sphinxstyleemphasis{pattern} )
\item[] \sphinxstyleemphasis{nr} = \sphinxstyleemphasis{matrix} . \sphinxcode{\sphinxupquote{nr}} ()
\item[] \sphinxstyleemphasis{nc} = \sphinxstyleemphasis{matrix} . \sphinxcode{\sphinxupquote{nc}} ()
\item[] \sphinxstyleemphasis{nnz} =  \sphinxcode{\sphinxupquote{matrix}} . \sphinxstyleemphasis{nnz} ()
\item[] \sphinxstyleemphasis{matrix} . \sphinxcode{\sphinxupquote{put}} ( \sphinxstyleemphasis{k} , \sphinxstyleemphasis{v} )
\item[] \sphinxstyleemphasis{row} = \sphinxstyleemphasis{matrix} . \sphinxcode{\sphinxupquote{row}} ()
\item[] \sphinxstyleemphasis{col} = \sphinxstyleemphasis{matrix} . \sphinxcode{\sphinxupquote{col}} ()
\item[] \sphinxstyleemphasis{val} = \sphinxstyleemphasis{matrix} . \sphinxcode{\sphinxupquote{val}} ()
\item[] \sphinxstyleemphasis{row\_major} = \sphinxstyleemphasis{matrix} . \sphinxcode{\sphinxupquote{row\_major}} ()
\item[] \sphinxstyleemphasis{col\_major} = \sphinxstyleemphasis{matrix} . \sphinxcode{\sphinxupquote{col\_major}} ()
\end{DUlineblock}

\index{pattern@\spxentry{pattern}}\ignorespaces 

\subparagraph{pattern}
\label{\detokenize{xsrst/cpp_sparse_rcv:pattern}}\label{\detokenize{xsrst/cpp_sparse_rcv:cpp-sparse-rcv-pattern}}\label{\detokenize{xsrst/cpp_sparse_rcv:index-2}}
This argument has prototype

\begin{DUlineblock}{0em}
\item[]         \sphinxcode{\sphinxupquote{const sparse\_rc\&}} \sphinxstyleemphasis{pattern}
\end{DUlineblock}

It specifies the number of rows, number of columns and
the possibly non\sphinxhyphen{}zero entries in the \sphinxstyleemphasis{matrix} .

\index{matrix@\spxentry{matrix}}\ignorespaces 

\subparagraph{matrix}
\label{\detokenize{xsrst/cpp_sparse_rcv:matrix}}\label{\detokenize{xsrst/cpp_sparse_rcv:cpp-sparse-rcv-matrix}}\label{\detokenize{xsrst/cpp_sparse_rcv:index-3}}
This is a sparse matrix object with the sparsity specified by \sphinxstyleemphasis{pattern} .
Only the \sphinxstyleemphasis{val} vector can be changed. All other values returned by
\sphinxstyleemphasis{matrix} are fixed during the constructor and constant there after.
The \sphinxstyleemphasis{val} vector is only changed by the constructor
and the \sphinxcode{\sphinxupquote{set}} function.

\index{nr@\spxentry{nr}}\ignorespaces 

\subparagraph{nr}
\label{\detokenize{xsrst/cpp_sparse_rcv:nr}}\label{\detokenize{xsrst/cpp_sparse_rcv:cpp-sparse-rcv-nr}}\label{\detokenize{xsrst/cpp_sparse_rcv:index-4}}
This return value has prototype

\begin{DUlineblock}{0em}
\item[]         \sphinxcode{\sphinxupquote{int}} \sphinxstyleemphasis{nr}
\end{DUlineblock}

and is the number of rows in the matrix.

\index{nc@\spxentry{nc}}\ignorespaces 

\subparagraph{nc}
\label{\detokenize{xsrst/cpp_sparse_rcv:nc}}\label{\detokenize{xsrst/cpp_sparse_rcv:cpp-sparse-rcv-nc}}\label{\detokenize{xsrst/cpp_sparse_rcv:index-5}}
This return value has prototype

\begin{DUlineblock}{0em}
\item[]         \sphinxcode{\sphinxupquote{int}} \sphinxstyleemphasis{nc}
\end{DUlineblock}

and is the number of columns in the matrix.

\index{nnz@\spxentry{nnz}}\ignorespaces 

\subparagraph{nnz}
\label{\detokenize{xsrst/cpp_sparse_rcv:nnz}}\label{\detokenize{xsrst/cpp_sparse_rcv:cpp-sparse-rcv-nnz}}\label{\detokenize{xsrst/cpp_sparse_rcv:index-6}}
This return value has prototype

\begin{DUlineblock}{0em}
\item[]         \sphinxcode{\sphinxupquote{int}} \sphinxstyleemphasis{nnz}
\end{DUlineblock}

and is the number of possibly non\sphinxhyphen{}zero values in the matrix.

\index{put@\spxentry{put}}\ignorespaces 

\subparagraph{put}
\label{\detokenize{xsrst/cpp_sparse_rcv:put}}\label{\detokenize{xsrst/cpp_sparse_rcv:cpp-sparse-rcv-put}}\label{\detokenize{xsrst/cpp_sparse_rcv:index-7}}
This function sets the value

\begin{DUlineblock}{0em}
\item[]         \sphinxstyleemphasis{val} {[} \sphinxstyleemphasis{k} {]} = \sphinxstyleemphasis{v}
\end{DUlineblock}

(The name \sphinxcode{\sphinxupquote{set}} is used by Cppad, but not used here,
because \sphinxcode{\sphinxupquote{set}} it is a built\sphinxhyphen{}in name in Python.)

\index{k@\spxentry{k}}\ignorespaces 

\subparagraph{k}
\label{\detokenize{xsrst/cpp_sparse_rcv:k}}\label{\detokenize{xsrst/cpp_sparse_rcv:cpp-sparse-rcv-put-k}}\label{\detokenize{xsrst/cpp_sparse_rcv:index-8}}
This argument has type

\begin{DUlineblock}{0em}
\item[]         \sphinxcode{\sphinxupquote{int}} \sphinxstyleemphasis{k}
\end{DUlineblock}

and must be non\sphinxhyphen{}negative and less than \sphinxstyleemphasis{nnz} .

\index{v@\spxentry{v}}\ignorespaces 

\subparagraph{v}
\label{\detokenize{xsrst/cpp_sparse_rcv:v}}\label{\detokenize{xsrst/cpp_sparse_rcv:cpp-sparse-rcv-put-v}}\label{\detokenize{xsrst/cpp_sparse_rcv:index-9}}
This argument has type

\begin{DUlineblock}{0em}
\item[]         \sphinxcode{\sphinxupquote{double}} \sphinxstyleemphasis{v}
\end{DUlineblock}

It specifies the value assigned to \sphinxstyleemphasis{val} {[} \sphinxstyleemphasis{k} {]} .

\index{row@\spxentry{row}}\ignorespaces 

\subparagraph{row}
\label{\detokenize{xsrst/cpp_sparse_rcv:row}}\label{\detokenize{xsrst/cpp_sparse_rcv:cpp-sparse-rcv-row}}\label{\detokenize{xsrst/cpp_sparse_rcv:index-10}}
This result has type

\begin{DUlineblock}{0em}
\item[]         \sphinxcode{\sphinxupquote{vec\_int}} \sphinxstyleemphasis{row}
\end{DUlineblock}

and its size is \sphinxstyleemphasis{nnz} .
For \sphinxstyleemphasis{k} = 0, … , \sphinxstyleemphasis{nnz} \sphinxhyphen{}1 ,
\sphinxstyleemphasis{row} {[} \sphinxstyleemphasis{k} {]} is the row index for the \sphinxstyleemphasis{k}\sphinxhyphen{}th possibly non\sphinxhyphen{}zero
entry in the matrix.

\index{col@\spxentry{col}}\ignorespaces 

\subparagraph{col}
\label{\detokenize{xsrst/cpp_sparse_rcv:col}}\label{\detokenize{xsrst/cpp_sparse_rcv:cpp-sparse-rcv-col}}\label{\detokenize{xsrst/cpp_sparse_rcv:index-11}}
This result has type

\begin{DUlineblock}{0em}
\item[]         \sphinxcode{\sphinxupquote{vec\_int}} \sphinxstyleemphasis{col}
\end{DUlineblock}

and its size is \sphinxstyleemphasis{nnz} .
For \sphinxstyleemphasis{k} = 0, … , \sphinxstyleemphasis{nnz} \sphinxhyphen{}1 ,
\sphinxstyleemphasis{col} {[} \sphinxstyleemphasis{k} {]} is the column index for the \sphinxstyleemphasis{k}\sphinxhyphen{}th possibly non\sphinxhyphen{}zero
entry in the matrix.

\index{val@\spxentry{val}}\ignorespaces 

\subparagraph{val}
\label{\detokenize{xsrst/cpp_sparse_rcv:val}}\label{\detokenize{xsrst/cpp_sparse_rcv:cpp-sparse-rcv-val}}\label{\detokenize{xsrst/cpp_sparse_rcv:index-12}}
This result has type

\begin{DUlineblock}{0em}
\item[]         \sphinxcode{\sphinxupquote{vec\_double}} \sphinxstyleemphasis{val}
\end{DUlineblock}

and its size is \sphinxstyleemphasis{nnz} .
For \sphinxstyleemphasis{k} = 0, … , \sphinxstyleemphasis{nnz} \sphinxhyphen{}1 ,
\sphinxstyleemphasis{val} {[} \sphinxstyleemphasis{k} {]} is the value of the \sphinxstyleemphasis{k}\sphinxhyphen{}th possibly non\sphinxhyphen{}zero
entry in the matrix (the value may be zero).

\index{row\_major@\spxentry{row\_major}}\ignorespaces 

\subparagraph{row\_major}
\label{\detokenize{xsrst/cpp_sparse_rcv:row-major}}\label{\detokenize{xsrst/cpp_sparse_rcv:cpp-sparse-rcv-row-major}}\label{\detokenize{xsrst/cpp_sparse_rcv:index-13}}
This vector has prototype

\begin{DUlineblock}{0em}
\item[]         \sphinxcode{\sphinxupquote{vec\_int}} \sphinxstyleemphasis{row\_major}
\end{DUlineblock}

and its size \sphinxstyleemphasis{nnz} .
It sorts the sparsity pattern in row\sphinxhyphen{}major order.
To be specific,

\begin{DUlineblock}{0em}
\item[]         \sphinxstyleemphasis{col} {[} \sphinxstyleemphasis{row\_major} {[} \sphinxstyleemphasis{k} {]} {]} \textless{}= \sphinxstyleemphasis{col} {[} \sphinxstyleemphasis{row\_major} {[} \sphinxstyleemphasis{k} +1{]} {]}
\end{DUlineblock}

and if \sphinxstyleemphasis{col} {[} \sphinxstyleemphasis{row\_major} {[} \sphinxstyleemphasis{k} {]} {]} == \sphinxstyleemphasis{col} {[} \sphinxstyleemphasis{row\_major} {[} \sphinxstyleemphasis{k} +1{]} {]} ,

\begin{DUlineblock}{0em}
\item[]         \sphinxstyleemphasis{row} {[} \sphinxstyleemphasis{row\_major} {[} \sphinxstyleemphasis{k} {]} {]} \textless{} \sphinxstyleemphasis{row} {[} \sphinxstyleemphasis{row\_major} {[} \sphinxstyleemphasis{k} +1{]} {]}
\end{DUlineblock}

This routine generates an assert if there are two entries with the same
row and column values (if \sphinxcode{\sphinxupquote{NDEBUG}} is not defined).

\index{col\_major@\spxentry{col\_major}}\ignorespaces 

\subparagraph{col\_major}
\label{\detokenize{xsrst/cpp_sparse_rcv:col-major}}\label{\detokenize{xsrst/cpp_sparse_rcv:cpp-sparse-rcv-col-major}}\label{\detokenize{xsrst/cpp_sparse_rcv:index-14}}
This vector has prototype

\begin{DUlineblock}{0em}
\item[]         \sphinxcode{\sphinxupquote{vec\_int}} \sphinxstyleemphasis{col\_major}
\end{DUlineblock}

and its size \sphinxstyleemphasis{nnz} .
It sorts the sparsity pattern in column\sphinxhyphen{}major order.
To be specific,

\begin{DUlineblock}{0em}
\item[]         \sphinxstyleemphasis{row} {[} \sphinxstyleemphasis{col\_major} {[} \sphinxstyleemphasis{k} {]} {]} \textless{}= \sphinxstyleemphasis{row} {[} \sphinxstyleemphasis{col\_major} {[} \sphinxstyleemphasis{k} +1{]} {]}
\end{DUlineblock}

and if \sphinxstyleemphasis{row} {[} \sphinxstyleemphasis{col\_major} {[} \sphinxstyleemphasis{k} {]} {]} == \sphinxstyleemphasis{row} {[} \sphinxstyleemphasis{col\_major} {[} \sphinxstyleemphasis{k} +1{]} {]} ,

\begin{DUlineblock}{0em}
\item[]         \sphinxstyleemphasis{col} {[} \sphinxstyleemphasis{col\_major} {[} \sphinxstyleemphasis{k} {]} {]} \textless{} \sphinxstyleemphasis{col} {[} \sphinxstyleemphasis{col\_major} {[} \sphinxstyleemphasis{k} +1{]} {]}
\end{DUlineblock}

This routine generates an assert if there are two entries with the same
row and column values (if \sphinxcode{\sphinxupquote{NDEBUG}} is not defined).


\subparagraph{sparse\_rcv\_xam\_cpp}
\label{\detokenize{xsrst/sparse_rcv_xam_cpp:sparse-rcv-xam-cpp}}\label{\detokenize{xsrst/sparse_rcv_xam_cpp::doc}}
\index{sparse\_rcv\_xam\_cpp@\spxentry{sparse\_rcv\_xam\_cpp}}\index{c++:@\spxentry{c++:}}\index{sparsity@\spxentry{sparsity}}\index{patterns:@\spxentry{patterns:}}\index{example@\spxentry{example}}\index{test@\spxentry{test}}\ignorespaces 

\subparagraph{3.2.4.2.1 C++: Sparsity Patterns: Example and Test}
\label{\detokenize{xsrst/sparse_rcv_xam_cpp:c-sparsity-patterns-example-and-test}}\label{\detokenize{xsrst/sparse_rcv_xam_cpp:id1}}
\begin{sphinxVerbatim}[commandchars=\\\{\}]
\PYG{c+cp}{\PYGZsh{}}\PYG{c+cp}{ include \PYGZlt{}cstdio\PYGZgt{}}
\PYG{c+cp}{\PYGZsh{}}\PYG{c+cp}{ include \PYGZlt{}cppad}\PYG{c+cp}{/}\PYG{c+cp}{py}\PYG{c+cp}{/}\PYG{c+cp}{cppad\PYGZus{}py.hpp\PYGZgt{}}

\PYG{k+kt}{bool} \PYG{n+nf}{sparse\PYGZus{}rcv\PYGZus{}xam}\PYG{p}{(}\PYG{k+kt}{void}\PYG{p}{)} \PYG{p}{\PYGZob{}}
    \PYG{k}{using} \PYG{n}{cppad\PYGZus{}py}\PYG{o}{:}\PYG{o}{:}\PYG{n}{vec\PYGZus{}int}\PYG{p}{;}
    \PYG{k}{using} \PYG{n}{cppad\PYGZus{}py}\PYG{o}{:}\PYG{o}{:}\PYG{n}{vec\PYGZus{}double}\PYG{p}{;}
    \PYG{k}{using} \PYG{n}{cppad\PYGZus{}py}\PYG{o}{:}\PYG{o}{:}\PYG{n}{sparse\PYGZus{}rc}\PYG{p}{;}
    \PYG{k}{using} \PYG{n}{cppad\PYGZus{}py}\PYG{o}{:}\PYG{o}{:}\PYG{n}{sparse\PYGZus{}rcv}\PYG{p}{;}
    \PYG{c+c1}{//}
    \PYG{c+c1}{// initialize return variable}
    \PYG{k+kt}{bool} \PYG{n}{ok} \PYG{o}{=} \PYG{n+nb}{true}\PYG{p}{;}
    \PYG{c+c1}{//\PYGZhy{}\PYGZhy{}\PYGZhy{}\PYGZhy{}\PYGZhy{}\PYGZhy{}\PYGZhy{}\PYGZhy{}\PYGZhy{}\PYGZhy{}\PYGZhy{}\PYGZhy{}\PYGZhy{}\PYGZhy{}\PYGZhy{}\PYGZhy{}\PYGZhy{}\PYGZhy{}\PYGZhy{}\PYGZhy{}\PYGZhy{}\PYGZhy{}\PYGZhy{}\PYGZhy{}\PYGZhy{}\PYGZhy{}\PYGZhy{}\PYGZhy{}\PYGZhy{}\PYGZhy{}\PYGZhy{}\PYGZhy{}\PYGZhy{}\PYGZhy{}\PYGZhy{}\PYGZhy{}\PYGZhy{}\PYGZhy{}\PYGZhy{}\PYGZhy{}\PYGZhy{}\PYGZhy{}\PYGZhy{}\PYGZhy{}\PYGZhy{}\PYGZhy{}\PYGZhy{}\PYGZhy{}\PYGZhy{}\PYGZhy{}\PYGZhy{}\PYGZhy{}\PYGZhy{}\PYGZhy{}\PYGZhy{}\PYGZhy{}\PYGZhy{}\PYGZhy{}\PYGZhy{}\PYGZhy{}\PYGZhy{}\PYGZhy{}\PYGZhy{}\PYGZhy{}\PYGZhy{}\PYGZhy{}\PYGZhy{}\PYGZhy{}\PYGZhy{}\PYGZhy{}\PYGZhy{}\PYGZhy{}}
    \PYG{c+c1}{//}
    \PYG{c+c1}{// create sparsity pattern for n by n identity matrix}
    \PYG{n}{sparse\PYGZus{}rc} \PYG{n}{pattern} \PYG{o}{=} \PYG{n}{sparse\PYGZus{}rc}\PYG{p}{(}\PYG{p}{)}\PYG{p}{;}
    \PYG{k+kt}{int} \PYG{n}{n} \PYG{o}{=} \PYG{l+m+mi}{5}\PYG{p}{;}
    \PYG{n}{pattern}\PYG{p}{.}\PYG{n}{resize}\PYG{p}{(}\PYG{n}{n}\PYG{p}{,} \PYG{n}{n}\PYG{p}{,} \PYG{n}{n}\PYG{p}{)}\PYG{p}{;}
    \PYG{k}{for}\PYG{p}{(}\PYG{k+kt}{int} \PYG{n}{k} \PYG{o}{=} \PYG{l+m+mi}{0}\PYG{p}{;} \PYG{n}{k} \PYG{o}{\PYGZlt{}} \PYG{n}{n}\PYG{p}{;} \PYG{n}{k}\PYG{o}{+}\PYG{o}{+}\PYG{p}{)} \PYG{p}{\PYGZob{}}
        \PYG{n}{pattern}\PYG{p}{.}\PYG{n}{put}\PYG{p}{(}\PYG{n}{k}\PYG{p}{,} \PYG{n}{k}\PYG{p}{,} \PYG{n}{k}\PYG{p}{)}\PYG{p}{;}
    \PYG{p}{\PYGZcb{}}
    \PYG{c+c1}{//}
    \PYG{c+c1}{// create n by n sparse representation of identity matrix}
    \PYG{n}{sparse\PYGZus{}rcv} \PYG{n}{matrix} \PYG{o}{=} \PYG{n}{sparse\PYGZus{}rcv}\PYG{p}{(}\PYG{n}{pattern}\PYG{p}{)}\PYG{p}{;}
    \PYG{k}{for}\PYG{p}{(}\PYG{k+kt}{int} \PYG{n}{k} \PYG{o}{=} \PYG{l+m+mi}{0}\PYG{p}{;} \PYG{n}{k} \PYG{o}{\PYGZlt{}} \PYG{n}{n}\PYG{p}{;} \PYG{n}{k}\PYG{o}{+}\PYG{o}{+}\PYG{p}{)} \PYG{p}{\PYGZob{}}
        \PYG{n}{matrix}\PYG{p}{.}\PYG{n}{put}\PYG{p}{(}\PYG{n}{k}\PYG{p}{,} \PYG{l+m+mf}{1.0}\PYG{p}{)}\PYG{p}{;}
    \PYG{p}{\PYGZcb{}}
    \PYG{c+c1}{//}
    \PYG{c+c1}{// row, column indices}
    \PYG{n}{vec\PYGZus{}int} \PYG{n}{row} \PYG{o}{=} \PYG{n}{matrix}\PYG{p}{.}\PYG{n}{row}\PYG{p}{(}\PYG{p}{)}\PYG{p}{;}
    \PYG{n}{vec\PYGZus{}int} \PYG{n}{col} \PYG{o}{=} \PYG{n}{matrix}\PYG{p}{.}\PYG{n}{col}\PYG{p}{(}\PYG{p}{)}\PYG{p}{;}
    \PYG{n}{vec\PYGZus{}double} \PYG{n}{val} \PYG{o}{=} \PYG{n}{matrix}\PYG{p}{.}\PYG{n}{val}\PYG{p}{(}\PYG{p}{)}\PYG{p}{;}
    \PYG{c+c1}{//}
    \PYG{c+c1}{// check results}
    \PYG{k}{for}\PYG{p}{(}\PYG{k+kt}{int} \PYG{n}{k} \PYG{o}{=} \PYG{l+m+mi}{0}\PYG{p}{;} \PYG{n}{k} \PYG{o}{\PYGZlt{}} \PYG{n}{n}\PYG{p}{;} \PYG{n}{k}\PYG{o}{+}\PYG{o}{+}\PYG{p}{)} \PYG{p}{\PYGZob{}}
        \PYG{n}{ok} \PYG{o}{=} \PYG{n}{ok} \PYG{o}{\PYGZam{}}\PYG{o}{\PYGZam{}} \PYG{n}{row}\PYG{p}{[}\PYG{n}{k}\PYG{p}{]} \PYG{o}{=}\PYG{o}{=} \PYG{n}{k}\PYG{p}{;}
        \PYG{n}{ok} \PYG{o}{=} \PYG{n}{ok} \PYG{o}{\PYGZam{}}\PYG{o}{\PYGZam{}} \PYG{n}{col}\PYG{p}{[}\PYG{n}{k}\PYG{p}{]} \PYG{o}{=}\PYG{o}{=} \PYG{n}{k}\PYG{p}{;}
        \PYG{n}{ok} \PYG{o}{=} \PYG{n}{ok} \PYG{o}{\PYGZam{}}\PYG{o}{\PYGZam{}} \PYG{n}{val}\PYG{p}{[}\PYG{n}{k}\PYG{p}{]} \PYG{o}{=}\PYG{o}{=} \PYG{l+m+mf}{1.0}\PYG{p}{;}
    \PYG{p}{\PYGZcb{}}
    \PYG{c+c1}{//}
    \PYG{c+c1}{// For this case, row\PYGZus{}major and col\PYGZus{}major order are same as original order}
    \PYG{n}{vec\PYGZus{}int} \PYG{n}{row\PYGZus{}major} \PYG{o}{=} \PYG{n}{matrix}\PYG{p}{.}\PYG{n}{row\PYGZus{}major}\PYG{p}{(}\PYG{p}{)}\PYG{p}{;}
    \PYG{n}{vec\PYGZus{}int} \PYG{n}{col\PYGZus{}major} \PYG{o}{=} \PYG{n}{matrix}\PYG{p}{.}\PYG{n}{col\PYGZus{}major}\PYG{p}{(}\PYG{p}{)}\PYG{p}{;}
    \PYG{k}{for}\PYG{p}{(}\PYG{k+kt}{int} \PYG{n}{k} \PYG{o}{=} \PYG{l+m+mi}{0}\PYG{p}{;} \PYG{n}{k} \PYG{o}{\PYGZlt{}} \PYG{n}{n}\PYG{p}{;} \PYG{n}{k}\PYG{o}{+}\PYG{o}{+}\PYG{p}{)} \PYG{p}{\PYGZob{}}
        \PYG{n}{ok} \PYG{o}{=} \PYG{n}{ok} \PYG{o}{\PYGZam{}}\PYG{o}{\PYGZam{}} \PYG{n}{row\PYGZus{}major}\PYG{p}{[}\PYG{n}{k}\PYG{p}{]} \PYG{o}{=}\PYG{o}{=} \PYG{n}{k}\PYG{p}{;}
        \PYG{n}{ok} \PYG{o}{=} \PYG{n}{ok} \PYG{o}{\PYGZam{}}\PYG{o}{\PYGZam{}} \PYG{n}{col\PYGZus{}major}\PYG{p}{[}\PYG{n}{k}\PYG{p}{]} \PYG{o}{=}\PYG{o}{=} \PYG{n}{k}\PYG{p}{;}
    \PYG{p}{\PYGZcb{}}
    \PYG{c+c1}{//}
    \PYG{k}{return}\PYG{p}{(} \PYG{n}{ok} \PYG{p}{)}\PYG{p}{;}
\PYG{p}{\PYGZcb{}}
\end{sphinxVerbatim}


\bigskip\hrule\bigskip


xsrst input file: \sphinxcode{\sphinxupquote{example/cplusplus/sparse\_rcv\_xam.cpp}}

\index{example@\spxentry{example}}\ignorespaces 

\subparagraph{Example}
\label{\detokenize{xsrst/cpp_sparse_rcv:example}}\label{\detokenize{xsrst/cpp_sparse_rcv:cpp-sparse-rcv-example}}\label{\detokenize{xsrst/cpp_sparse_rcv:index-15}}
{\hyperref[\detokenize{xsrst/sparse_rcv_xam_cpp:id1}]{\sphinxcrossref{\DUrole{std,std-ref}{sparse\_rcv\_xam\_cpp}}}}


\bigskip\hrule\bigskip


xsrst input file: \sphinxcode{\sphinxupquote{lib/cplusplus/sparse.cpp}}


\subparagraph{cpp\_jac\_sparsity}
\label{\detokenize{xsrst/cpp_jac_sparsity:cpp-jac-sparsity}}\label{\detokenize{xsrst/cpp_jac_sparsity::doc}}
\index{cpp\_jac\_sparsity@\spxentry{cpp\_jac\_sparsity}}\index{jacobian@\spxentry{jacobian}}\index{sparsity@\spxentry{sparsity}}\index{patterns@\spxentry{patterns}}\ignorespaces 

\subparagraph{3.2.4.3 Jacobian Sparsity Patterns}
\label{\detokenize{xsrst/cpp_jac_sparsity:jacobian-sparsity-patterns}}\label{\detokenize{xsrst/cpp_jac_sparsity:id1}}\begin{itemize}
\item {} 
{\hyperref[\detokenize{xsrst/cpp_jac_sparsity:cpp-jac-sparsity-syntax}]{\sphinxcrossref{\DUrole{std,std-ref}{Syntax}}}}

\item {} \begin{description}
\item[{{\hyperref[\detokenize{xsrst/cpp_jac_sparsity:cpp-jac-sparsity-purpose}]{\sphinxcrossref{\DUrole{std,std-ref}{Purpose}}}}}] \leavevmode\begin{itemize}
\item {} 
{\hyperref[\detokenize{xsrst/cpp_jac_sparsity:cpp-jac-sparsity-purpose-for-jac-sparsity}]{\sphinxcrossref{\DUrole{std,std-ref}{for\_jac\_sparsity}}}}

\item {} 
{\hyperref[\detokenize{xsrst/cpp_jac_sparsity:cpp-jac-sparsity-purpose-rev-jac-sparsity}]{\sphinxcrossref{\DUrole{std,std-ref}{rev\_jac\_sparsity}}}}

\end{itemize}

\end{description}

\item {} 
{\hyperref[\detokenize{xsrst/cpp_jac_sparsity:cpp-jac-sparsity-x}]{\sphinxcrossref{\DUrole{std,std-ref}{x}}}}

\item {} 
{\hyperref[\detokenize{xsrst/cpp_jac_sparsity:cpp-jac-sparsity-f}]{\sphinxcrossref{\DUrole{std,std-ref}{f}}}}

\item {} 
{\hyperref[\detokenize{xsrst/cpp_jac_sparsity:cpp-jac-sparsity-pattern-in}]{\sphinxcrossref{\DUrole{std,std-ref}{pattern\_in}}}}

\item {} 
{\hyperref[\detokenize{xsrst/cpp_jac_sparsity:cpp-jac-sparsity-pattern-out}]{\sphinxcrossref{\DUrole{std,std-ref}{pattern\_out}}}}

\item {} 
{\hyperref[\detokenize{xsrst/cpp_jac_sparsity:cpp-jac-sparsity-sparsity-for-entire-jacobian}]{\sphinxcrossref{\DUrole{std,std-ref}{Sparsity for Entire Jacobian}}}}

\item {} 
{\hyperref[\detokenize{xsrst/cpp_jac_sparsity:cpp-jac-sparsity-example}]{\sphinxcrossref{\DUrole{std,std-ref}{Example}}}}

\end{itemize}

\index{syntax@\spxentry{syntax}}\ignorespaces 

\subparagraph{Syntax}
\label{\detokenize{xsrst/cpp_jac_sparsity:syntax}}\label{\detokenize{xsrst/cpp_jac_sparsity:cpp-jac-sparsity-syntax}}\label{\detokenize{xsrst/cpp_jac_sparsity:index-1}}
\begin{DUlineblock}{0em}
\item[] \sphinxstyleemphasis{f} . \sphinxcode{\sphinxupquote{for\_jac\_sparsity}} ( \sphinxstyleemphasis{pattern\_in} , \sphinxstyleemphasis{pattern\_out} )
\item[] \sphinxstyleemphasis{f} . \sphinxcode{\sphinxupquote{rev\_jac\_sparsity}} ( \sphinxstyleemphasis{pattern\_in} , \sphinxstyleemphasis{pattern\_out} )
\end{DUlineblock}

\index{purpose@\spxentry{purpose}}\ignorespaces 

\subparagraph{Purpose}
\label{\detokenize{xsrst/cpp_jac_sparsity:purpose}}\label{\detokenize{xsrst/cpp_jac_sparsity:cpp-jac-sparsity-purpose}}\label{\detokenize{xsrst/cpp_jac_sparsity:index-2}}
We use \(F : \B{R}^n \rightarrow \B{R}^m\) to denote the
function corresponding to the operation sequence stored in \sphinxstyleemphasis{f} .

\index{for\_jac\_sparsity@\spxentry{for\_jac\_sparsity}}\ignorespaces 

\subparagraph{for\_jac\_sparsity}
\label{\detokenize{xsrst/cpp_jac_sparsity:for-jac-sparsity}}\label{\detokenize{xsrst/cpp_jac_sparsity:cpp-jac-sparsity-purpose-for-jac-sparsity}}\label{\detokenize{xsrst/cpp_jac_sparsity:index-3}}
Fix \(R \in \B{R}^{n \times \ell}\) and define the function
\begin{equation*}
\begin{split}J(x) = F^{(1)} ( x ) * R\end{split}
\end{equation*}
Given a sparsity pattern for \(R\),
\sphinxcode{\sphinxupquote{for\_jac\_sparsity}} computes a sparsity pattern for \(J(x)\).

\index{rev\_jac\_sparsity@\spxentry{rev\_jac\_sparsity}}\ignorespaces 

\subparagraph{rev\_jac\_sparsity}
\label{\detokenize{xsrst/cpp_jac_sparsity:rev-jac-sparsity}}\label{\detokenize{xsrst/cpp_jac_sparsity:cpp-jac-sparsity-purpose-rev-jac-sparsity}}\label{\detokenize{xsrst/cpp_jac_sparsity:index-4}}
Fix \(R \in \B{R}^{\ell \times m}\) and define the function
\begin{equation*}
\begin{split}J(x) = R * F^{(1)} ( x )\end{split}
\end{equation*}
Given a sparsity pattern for \(R\),
\sphinxcode{\sphinxupquote{rev\_jac\_sparsity}} computes a sparsity pattern for \(J(x)\).

\index{x@\spxentry{x}}\ignorespaces 

\subparagraph{x}
\label{\detokenize{xsrst/cpp_jac_sparsity:x}}\label{\detokenize{xsrst/cpp_jac_sparsity:cpp-jac-sparsity-x}}\label{\detokenize{xsrst/cpp_jac_sparsity:index-5}}
Note that a sparsity pattern for \(J(x)\) corresponds to the
operation sequence stored in \sphinxstyleemphasis{f} and does not depend on
the argument \sphinxstyleemphasis{x} .

\index{f@\spxentry{f}}\ignorespaces 

\subparagraph{f}
\label{\detokenize{xsrst/cpp_jac_sparsity:f}}\label{\detokenize{xsrst/cpp_jac_sparsity:cpp-jac-sparsity-f}}\label{\detokenize{xsrst/cpp_jac_sparsity:index-6}}
The object \sphinxstyleemphasis{f} has prototype

\begin{DUlineblock}{0em}
\item[]         \sphinxcode{\sphinxupquote{d\_fun}} \sphinxstyleemphasis{f}
\end{DUlineblock}

The object \sphinxstyleemphasis{f} is not \sphinxcode{\sphinxupquote{const}} when using
\sphinxcode{\sphinxupquote{for\_jac\_sparsity}} .
After a call to \sphinxcode{\sphinxupquote{for\_jac\_sparsity}} , a sparsity pattern
for each of the variables in the operation sequence
is held in \sphinxstyleemphasis{f} for possible later use during
reverse Hessian sparsity calculations.

\index{pattern\_in@\spxentry{pattern\_in}}\ignorespaces 

\subparagraph{pattern\_in}
\label{\detokenize{xsrst/cpp_jac_sparsity:pattern-in}}\label{\detokenize{xsrst/cpp_jac_sparsity:cpp-jac-sparsity-pattern-in}}\label{\detokenize{xsrst/cpp_jac_sparsity:index-7}}
The argument \sphinxstyleemphasis{pattern\_in} has prototype

\begin{DUlineblock}{0em}
\item[]         \sphinxcode{\sphinxupquote{const sparse\_rc\&}} \sphinxstyleemphasis{pattern\_in}
\end{DUlineblock}

see {\hyperref[\detokenize{xsrst/cpp_sparse_rc:id1}]{\sphinxcrossref{\DUrole{std,std-ref}{cpp\_sparse\_rc}}}}.
This is a sparsity pattern for \(R\).

\index{pattern\_out@\spxentry{pattern\_out}}\ignorespaces 

\subparagraph{pattern\_out}
\label{\detokenize{xsrst/cpp_jac_sparsity:pattern-out}}\label{\detokenize{xsrst/cpp_jac_sparsity:cpp-jac-sparsity-pattern-out}}\label{\detokenize{xsrst/cpp_jac_sparsity:index-8}}
This argument has prototype

\begin{DUlineblock}{0em}
\item[]         \sphinxcode{\sphinxupquote{sparse\_rc}} \textless{} \sphinxstyleemphasis{SizeVector} \textgreater{} \sphinxcode{\sphinxupquote{\&}} \sphinxstyleemphasis{pattern\_out}
\end{DUlineblock}

This input value of \sphinxstyleemphasis{pattern\_out} does not matter.
Upon return \sphinxstyleemphasis{pattern\_out} is a sparsity pattern for
\(J(x)\).

\index{sparsity@\spxentry{sparsity}}\index{for@\spxentry{for}}\index{entire@\spxentry{entire}}\index{jacobian@\spxentry{jacobian}}\ignorespaces 

\subparagraph{Sparsity for Entire Jacobian}
\label{\detokenize{xsrst/cpp_jac_sparsity:sparsity-for-entire-jacobian}}\label{\detokenize{xsrst/cpp_jac_sparsity:cpp-jac-sparsity-sparsity-for-entire-jacobian}}\label{\detokenize{xsrst/cpp_jac_sparsity:index-9}}
Suppose that \(R\) is the identity matrix.
In this case, \sphinxstyleemphasis{pattern\_out} is a sparsity pattern for
\(F^{(1)} ( x )\).


\subparagraph{sparse\_jac\_pattern\_xam\_cpp}
\label{\detokenize{xsrst/sparse_jac_pattern_xam_cpp:sparse-jac-pattern-xam-cpp}}\label{\detokenize{xsrst/sparse_jac_pattern_xam_cpp::doc}}
\index{sparse\_jac\_pattern\_xam\_cpp@\spxentry{sparse\_jac\_pattern\_xam\_cpp}}\index{c++:@\spxentry{c++:}}\index{jacobian@\spxentry{jacobian}}\index{sparsity@\spxentry{sparsity}}\index{patterns:@\spxentry{patterns:}}\index{example@\spxentry{example}}\index{test@\spxentry{test}}\ignorespaces 

\subparagraph{3.2.4.3.1 C++: Jacobian Sparsity Patterns: Example and Test}
\label{\detokenize{xsrst/sparse_jac_pattern_xam_cpp:c-jacobian-sparsity-patterns-example-and-test}}\label{\detokenize{xsrst/sparse_jac_pattern_xam_cpp:id1}}
\begin{sphinxVerbatim}[commandchars=\\\{\}]
\PYG{c+cp}{\PYGZsh{}}\PYG{c+cp}{ include \PYGZlt{}cstdio\PYGZgt{}}
\PYG{c+cp}{\PYGZsh{}}\PYG{c+cp}{ include \PYGZlt{}cppad}\PYG{c+cp}{/}\PYG{c+cp}{py}\PYG{c+cp}{/}\PYG{c+cp}{cppad\PYGZus{}py.hpp\PYGZgt{}}

\PYG{k+kt}{bool} \PYG{n+nf}{sparse\PYGZus{}jac\PYGZus{}pattern\PYGZus{}xam}\PYG{p}{(}\PYG{k+kt}{void}\PYG{p}{)} \PYG{p}{\PYGZob{}}
    \PYG{k}{using} \PYG{n}{cppad\PYGZus{}py}\PYG{o}{:}\PYG{o}{:}\PYG{n}{a\PYGZus{}double}\PYG{p}{;}
    \PYG{k}{using} \PYG{n}{cppad\PYGZus{}py}\PYG{o}{:}\PYG{o}{:}\PYG{n}{vec\PYGZus{}int}\PYG{p}{;}
    \PYG{k}{using} \PYG{n}{cppad\PYGZus{}py}\PYG{o}{:}\PYG{o}{:}\PYG{n}{vec\PYGZus{}double}\PYG{p}{;}
    \PYG{k}{using} \PYG{n}{cppad\PYGZus{}py}\PYG{o}{:}\PYG{o}{:}\PYG{n}{vec\PYGZus{}a\PYGZus{}double}\PYG{p}{;}
    \PYG{k}{using} \PYG{n}{cppad\PYGZus{}py}\PYG{o}{:}\PYG{o}{:}\PYG{n}{d\PYGZus{}fun}\PYG{p}{;}
    \PYG{k}{using} \PYG{n}{cppad\PYGZus{}py}\PYG{o}{:}\PYG{o}{:}\PYG{n}{sparse\PYGZus{}rc}\PYG{p}{;}
    \PYG{c+c1}{//}
    \PYG{c+c1}{// initialize return variable}
    \PYG{k+kt}{bool} \PYG{n}{ok} \PYG{o}{=} \PYG{n+nb}{true}\PYG{p}{;}
    \PYG{c+c1}{//\PYGZhy{}\PYGZhy{}\PYGZhy{}\PYGZhy{}\PYGZhy{}\PYGZhy{}\PYGZhy{}\PYGZhy{}\PYGZhy{}\PYGZhy{}\PYGZhy{}\PYGZhy{}\PYGZhy{}\PYGZhy{}\PYGZhy{}\PYGZhy{}\PYGZhy{}\PYGZhy{}\PYGZhy{}\PYGZhy{}\PYGZhy{}\PYGZhy{}\PYGZhy{}\PYGZhy{}\PYGZhy{}\PYGZhy{}\PYGZhy{}\PYGZhy{}\PYGZhy{}\PYGZhy{}\PYGZhy{}\PYGZhy{}\PYGZhy{}\PYGZhy{}\PYGZhy{}\PYGZhy{}\PYGZhy{}\PYGZhy{}\PYGZhy{}\PYGZhy{}\PYGZhy{}\PYGZhy{}\PYGZhy{}\PYGZhy{}\PYGZhy{}\PYGZhy{}\PYGZhy{}\PYGZhy{}\PYGZhy{}\PYGZhy{}\PYGZhy{}\PYGZhy{}\PYGZhy{}\PYGZhy{}\PYGZhy{}\PYGZhy{}\PYGZhy{}\PYGZhy{}\PYGZhy{}\PYGZhy{}\PYGZhy{}\PYGZhy{}\PYGZhy{}\PYGZhy{}\PYGZhy{}\PYGZhy{}\PYGZhy{}\PYGZhy{}\PYGZhy{}\PYGZhy{}\PYGZhy{}\PYGZhy{}}
    \PYG{c+c1}{// number of dependent and independent variables}
    \PYG{k+kt}{int} \PYG{n}{n} \PYG{o}{=} \PYG{l+m+mi}{3}\PYG{p}{;}
    \PYG{c+c1}{//}
    \PYG{c+c1}{// create the independent variables ax}
    \PYG{n}{vec\PYGZus{}double} \PYG{n}{x}\PYG{p}{(}\PYG{n}{n}\PYG{p}{)}\PYG{p}{;}
    \PYG{k}{for}\PYG{p}{(}\PYG{k+kt}{int} \PYG{n}{i} \PYG{o}{=} \PYG{l+m+mi}{0}\PYG{p}{;} \PYG{n}{i} \PYG{o}{\PYGZlt{}} \PYG{n}{n} \PYG{p}{;} \PYG{n}{i}\PYG{o}{+}\PYG{o}{+}\PYG{p}{)} \PYG{p}{\PYGZob{}}
        \PYG{n}{x}\PYG{p}{[}\PYG{n}{i}\PYG{p}{]} \PYG{o}{=} \PYG{n}{i} \PYG{o}{+} \PYG{l+m+mf}{2.0}\PYG{p}{;}
    \PYG{p}{\PYGZcb{}}
    \PYG{n}{vec\PYGZus{}a\PYGZus{}double} \PYG{n}{ax} \PYG{o}{=} \PYG{n}{cppad\PYGZus{}py}\PYG{o}{:}\PYG{o}{:}\PYG{n}{independent}\PYG{p}{(}\PYG{n}{x}\PYG{p}{)}\PYG{p}{;}
    \PYG{c+c1}{//}
    \PYG{c+c1}{// create dependent variables ay with ay[i] = ax[j]}
    \PYG{c+c1}{// where i = mod(j + 1, n)}
    \PYG{n}{vec\PYGZus{}a\PYGZus{}double} \PYG{n}{ay}\PYG{p}{(}\PYG{n}{n}\PYG{p}{)}\PYG{p}{;}
    \PYG{k}{for}\PYG{p}{(}\PYG{k+kt}{int} \PYG{n}{j} \PYG{o}{=} \PYG{l+m+mi}{0}\PYG{p}{;} \PYG{n}{j} \PYG{o}{\PYGZlt{}} \PYG{n}{n} \PYG{p}{;} \PYG{n}{j}\PYG{o}{+}\PYG{o}{+}\PYG{p}{)} \PYG{p}{\PYGZob{}}
        \PYG{k+kt}{int} \PYG{n}{i} \PYG{o}{=} \PYG{n}{j}\PYG{o}{+}\PYG{l+m+mi}{1}\PYG{p}{;}
        \PYG{k}{if}\PYG{p}{(} \PYG{n}{i} \PYG{o}{\PYGZgt{}}\PYG{o}{=} \PYG{n}{n}  \PYG{p}{)} \PYG{p}{\PYGZob{}}
            \PYG{n}{i} \PYG{o}{=} \PYG{n}{i} \PYG{o}{\PYGZhy{}} \PYG{n}{n}\PYG{p}{;}
        \PYG{p}{\PYGZcb{}}
        \PYG{n}{a\PYGZus{}double} \PYG{n}{ay\PYGZus{}i} \PYG{o}{=} \PYG{n}{ax}\PYG{p}{[}\PYG{n}{j}\PYG{p}{]}\PYG{p}{;}
        \PYG{n}{ay}\PYG{p}{[}\PYG{n}{i}\PYG{p}{]} \PYG{o}{=} \PYG{n}{ay\PYGZus{}i}\PYG{p}{;}
    \PYG{p}{\PYGZcb{}}
    \PYG{c+c1}{//}
    \PYG{c+c1}{// define af corresponding to f(x)}
    \PYG{n}{d\PYGZus{}fun} \PYG{n}{f}\PYG{p}{(}\PYG{n}{ax}\PYG{p}{,} \PYG{n}{ay}\PYG{p}{)}\PYG{p}{;}
    \PYG{c+c1}{//}
    \PYG{c+c1}{// sparsity pattern for identity matrix}
    \PYG{n}{sparse\PYGZus{}rc} \PYG{n}{pat\PYGZus{}in} \PYG{o}{=} \PYG{n}{sparse\PYGZus{}rc}\PYG{p}{(}\PYG{p}{)}\PYG{p}{;}
    \PYG{n}{pat\PYGZus{}in}\PYG{p}{.}\PYG{n}{resize}\PYG{p}{(}\PYG{n}{n}\PYG{p}{,} \PYG{n}{n}\PYG{p}{,} \PYG{n}{n}\PYG{p}{)}\PYG{p}{;}
    \PYG{k}{for}\PYG{p}{(}\PYG{k+kt}{int} \PYG{n}{k} \PYG{o}{=} \PYG{l+m+mi}{0}\PYG{p}{;} \PYG{n}{k} \PYG{o}{\PYGZlt{}} \PYG{n}{n}\PYG{p}{;} \PYG{n}{k}\PYG{o}{+}\PYG{o}{+}\PYG{p}{)} \PYG{p}{\PYGZob{}}
        \PYG{n}{pat\PYGZus{}in}\PYG{p}{.}\PYG{n}{put}\PYG{p}{(}\PYG{n}{k}\PYG{p}{,} \PYG{n}{k}\PYG{p}{,} \PYG{n}{k}\PYG{p}{)}\PYG{p}{;}
    \PYG{p}{\PYGZcb{}}
    \PYG{c+c1}{//}
    \PYG{c+c1}{// loop over forward and reverse mode}
    \PYG{k}{for}\PYG{p}{(}\PYG{k+kt}{int} \PYG{n}{mode} \PYG{o}{=} \PYG{l+m+mi}{0}\PYG{p}{;} \PYG{n}{mode} \PYG{o}{\PYGZlt{}} \PYG{l+m+mi}{2}\PYG{p}{;} \PYG{n}{mode}\PYG{o}{+}\PYG{o}{+}\PYG{p}{)} \PYG{p}{\PYGZob{}}
        \PYG{n}{sparse\PYGZus{}rc} \PYG{n}{pat\PYGZus{}out} \PYG{o}{=} \PYG{n}{sparse\PYGZus{}rc}\PYG{p}{(}\PYG{p}{)}\PYG{p}{;}
        \PYG{k}{if}\PYG{p}{(} \PYG{n}{mode} \PYG{o}{=}\PYG{o}{=} \PYG{l+m+mi}{0}  \PYG{p}{)} \PYG{p}{\PYGZob{}}
            \PYG{n}{f}\PYG{p}{.}\PYG{n}{for\PYGZus{}jac\PYGZus{}sparsity}\PYG{p}{(}\PYG{n}{pat\PYGZus{}in}\PYG{p}{,} \PYG{n}{pat\PYGZus{}out}\PYG{p}{)}\PYG{p}{;}
        \PYG{p}{\PYGZcb{}}
        \PYG{k}{if}\PYG{p}{(} \PYG{n}{mode} \PYG{o}{=}\PYG{o}{=} \PYG{l+m+mi}{1}  \PYG{p}{)} \PYG{p}{\PYGZob{}}
            \PYG{n}{f}\PYG{p}{.}\PYG{n}{rev\PYGZus{}jac\PYGZus{}sparsity}\PYG{p}{(}\PYG{n}{pat\PYGZus{}in}\PYG{p}{,} \PYG{n}{pat\PYGZus{}out}\PYG{p}{)}\PYG{p}{;}
        \PYG{p}{\PYGZcb{}}
        \PYG{c+c1}{//}
        \PYG{c+c1}{// check that result is sparsity pattern for Jacobian}
        \PYG{n}{ok} \PYG{o}{=} \PYG{n}{ok} \PYG{o}{\PYGZam{}}\PYG{o}{\PYGZam{}} \PYG{n}{n} \PYG{o}{=}\PYG{o}{=} \PYG{n}{pat\PYGZus{}out}\PYG{p}{.}\PYG{n}{nnz}\PYG{p}{(}\PYG{p}{)}\PYG{p}{;}
        \PYG{n}{vec\PYGZus{}int} \PYG{n}{col\PYGZus{}major} \PYG{o}{=} \PYG{n}{pat\PYGZus{}out}\PYG{p}{.}\PYG{n}{col\PYGZus{}major}\PYG{p}{(}\PYG{p}{)}\PYG{p}{;}
        \PYG{n}{vec\PYGZus{}int} \PYG{n}{row} \PYG{o}{=} \PYG{n}{pat\PYGZus{}out}\PYG{p}{.}\PYG{n}{row}\PYG{p}{(}\PYG{p}{)}\PYG{p}{;}
        \PYG{n}{vec\PYGZus{}int} \PYG{n}{col} \PYG{o}{=} \PYG{n}{pat\PYGZus{}out}\PYG{p}{.}\PYG{n}{col}\PYG{p}{(}\PYG{p}{)}\PYG{p}{;}
        \PYG{k}{for}\PYG{p}{(}\PYG{k+kt}{int} \PYG{n}{k} \PYG{o}{=} \PYG{l+m+mi}{0}\PYG{p}{;} \PYG{n}{k} \PYG{o}{\PYGZlt{}} \PYG{n}{n}\PYG{p}{;} \PYG{n}{k}\PYG{o}{+}\PYG{o}{+}\PYG{p}{)} \PYG{p}{\PYGZob{}}
            \PYG{k+kt}{int} \PYG{n}{ell} \PYG{o}{=} \PYG{n}{col\PYGZus{}major}\PYG{p}{[}\PYG{n}{k}\PYG{p}{]}\PYG{p}{;}
            \PYG{k+kt}{int} \PYG{n}{r} \PYG{o}{=} \PYG{n}{row}\PYG{p}{[}\PYG{n}{ell}\PYG{p}{]}\PYG{p}{;}
            \PYG{k+kt}{int} \PYG{n}{c} \PYG{o}{=} \PYG{n}{col}\PYG{p}{[}\PYG{n}{ell}\PYG{p}{]}\PYG{p}{;}
            \PYG{k+kt}{int} \PYG{n}{i} \PYG{o}{=} \PYG{n}{c}\PYG{o}{+}\PYG{l+m+mi}{1}\PYG{p}{;}
            \PYG{k}{if}\PYG{p}{(} \PYG{n}{i} \PYG{o}{\PYGZgt{}}\PYG{o}{=}  \PYG{n}{n}  \PYG{p}{)} \PYG{p}{\PYGZob{}}
                \PYG{n}{i} \PYG{o}{=} \PYG{n}{i} \PYG{o}{\PYGZhy{}} \PYG{n}{n}\PYG{p}{;}
            \PYG{p}{\PYGZcb{}}
            \PYG{n}{ok} \PYG{o}{=} \PYG{n}{ok} \PYG{o}{\PYGZam{}}\PYG{o}{\PYGZam{}} \PYG{n}{c} \PYG{o}{=}\PYG{o}{=} \PYG{n}{k}\PYG{p}{;}
            \PYG{n}{ok} \PYG{o}{=} \PYG{n}{ok} \PYG{o}{\PYGZam{}}\PYG{o}{\PYGZam{}} \PYG{n}{r} \PYG{o}{=}\PYG{o}{=} \PYG{n}{i}\PYG{p}{;}
        \PYG{p}{\PYGZcb{}}
    \PYG{p}{\PYGZcb{}}
    \PYG{c+c1}{//}
    \PYG{k}{return}\PYG{p}{(} \PYG{n}{ok} \PYG{p}{)}\PYG{p}{;}
\PYG{p}{\PYGZcb{}}
\end{sphinxVerbatim}


\bigskip\hrule\bigskip


xsrst input file: \sphinxcode{\sphinxupquote{example/cplusplus/sparse\_jac\_pattern\_xam.cpp}}

\index{example@\spxentry{example}}\ignorespaces 

\subparagraph{Example}
\label{\detokenize{xsrst/cpp_jac_sparsity:example}}\label{\detokenize{xsrst/cpp_jac_sparsity:cpp-jac-sparsity-example}}\label{\detokenize{xsrst/cpp_jac_sparsity:index-10}}
{\hyperref[\detokenize{xsrst/sparse_jac_pattern_xam_cpp:id1}]{\sphinxcrossref{\DUrole{std,std-ref}{c++}}}}


\bigskip\hrule\bigskip


xsrst input file: \sphinxcode{\sphinxupquote{lib/cplusplus/sparse.cpp}}


\subparagraph{cpp\_sparsity}
\label{\detokenize{xsrst/cpp_sparsity:cpp-sparsity}}\label{\detokenize{xsrst/cpp_sparsity::doc}}
\index{cpp\_sparsity@\spxentry{cpp\_sparsity}}\index{hessian@\spxentry{hessian}}\index{sparsity@\spxentry{sparsity}}\index{patterns@\spxentry{patterns}}\ignorespaces 

\subparagraph{3.2.4.4 Hessian Sparsity Patterns}
\label{\detokenize{xsrst/cpp_sparsity:hessian-sparsity-patterns}}\label{\detokenize{xsrst/cpp_sparsity:id1}}\begin{itemize}
\item {} 
{\hyperref[\detokenize{xsrst/cpp_sparsity:cpp-sparsity-syntax}]{\sphinxcrossref{\DUrole{std,std-ref}{Syntax}}}}

\item {} 
{\hyperref[\detokenize{xsrst/cpp_sparsity:cpp-sparsity-purpose}]{\sphinxcrossref{\DUrole{std,std-ref}{Purpose}}}}

\item {} 
{\hyperref[\detokenize{xsrst/cpp_sparsity:cpp-sparsity-x}]{\sphinxcrossref{\DUrole{std,std-ref}{x}}}}

\item {} 
{\hyperref[\detokenize{xsrst/cpp_sparsity:cpp-sparsity-f}]{\sphinxcrossref{\DUrole{std,std-ref}{f}}}}

\item {} 
{\hyperref[\detokenize{xsrst/cpp_sparsity:cpp-sparsity-select-domain}]{\sphinxcrossref{\DUrole{std,std-ref}{select\_domain}}}}

\item {} 
{\hyperref[\detokenize{xsrst/cpp_sparsity:cpp-sparsity-select-range}]{\sphinxcrossref{\DUrole{std,std-ref}{select\_range}}}}

\item {} 
{\hyperref[\detokenize{xsrst/cpp_sparsity:cpp-sparsity-pattern-out}]{\sphinxcrossref{\DUrole{std,std-ref}{pattern\_out}}}}

\item {} 
{\hyperref[\detokenize{xsrst/cpp_sparsity:cpp-sparsity-sparsity-for-component-wise-hessian}]{\sphinxcrossref{\DUrole{std,std-ref}{Sparsity for Component Wise Hessian}}}}

\item {} 
{\hyperref[\detokenize{xsrst/cpp_sparsity:cpp-sparsity-example}]{\sphinxcrossref{\DUrole{std,std-ref}{Example}}}}

\end{itemize}

\index{syntax@\spxentry{syntax}}\ignorespaces 

\subparagraph{Syntax}
\label{\detokenize{xsrst/cpp_sparsity:syntax}}\label{\detokenize{xsrst/cpp_sparsity:cpp-sparsity-syntax}}\label{\detokenize{xsrst/cpp_sparsity:index-1}}
\begin{DUlineblock}{0em}
\item[] \sphinxstyleemphasis{f} . \sphinxcode{\sphinxupquote{for\_hes\_sparsity}} ( \sphinxstyleemphasis{select\_domain} , \sphinxstyleemphasis{select\_range} , \sphinxstyleemphasis{pattern\_out} )
\item[] \sphinxstyleemphasis{f} . \sphinxcode{\sphinxupquote{rev\_hes\_sparsity}} ( \sphinxstyleemphasis{select\_domain} , \sphinxstyleemphasis{select\_range} , \sphinxstyleemphasis{pattern\_out} )
\end{DUlineblock}

\index{purpose@\spxentry{purpose}}\ignorespaces 

\subparagraph{Purpose}
\label{\detokenize{xsrst/cpp_sparsity:purpose}}\label{\detokenize{xsrst/cpp_sparsity:cpp-sparsity-purpose}}\label{\detokenize{xsrst/cpp_sparsity:index-2}}
We use \(F : \B{R}^n \rightarrow \B{R}^m\) to denote the
function corresponding to the operation sequence stored in \sphinxstyleemphasis{f} .
Fix a diagonal matrix \(D \in \B{R}^{n \times n}\), fix a vector
\(r \in \B{R}^m\), and define
\begin{equation*}
\begin{split}H(x) = D (r^\R{T} F)^{(2)} ( x ) D\end{split}
\end{equation*}
Given a sparsity pattern for \(D\) and \(r\),
these routines compute a sparsity pattern for \(H(x)\).

\index{x@\spxentry{x}}\ignorespaces 

\subparagraph{x}
\label{\detokenize{xsrst/cpp_sparsity:x}}\label{\detokenize{xsrst/cpp_sparsity:cpp-sparsity-x}}\label{\detokenize{xsrst/cpp_sparsity:index-3}}
Note that a sparsity pattern for \(H(x)\) corresponds to the
operation sequence stored in \sphinxstyleemphasis{f} and does not depend on
the argument \sphinxstyleemphasis{x} .

\index{f@\spxentry{f}}\ignorespaces 

\subparagraph{f}
\label{\detokenize{xsrst/cpp_sparsity:f}}\label{\detokenize{xsrst/cpp_sparsity:cpp-sparsity-f}}\label{\detokenize{xsrst/cpp_sparsity:index-4}}
The object \sphinxstyleemphasis{f} has prototype

\begin{DUlineblock}{0em}
\item[]         \sphinxcode{\sphinxupquote{d\_fun}} \sphinxstyleemphasis{f}
\end{DUlineblock}

\index{select\_domain@\spxentry{select\_domain}}\ignorespaces 

\subparagraph{select\_domain}
\label{\detokenize{xsrst/cpp_sparsity:select-domain}}\label{\detokenize{xsrst/cpp_sparsity:cpp-sparsity-select-domain}}\label{\detokenize{xsrst/cpp_sparsity:index-5}}
The argument \sphinxstyleemphasis{select\_domain} has prototype

\begin{DUlineblock}{0em}
\item[]         \sphinxcode{\sphinxupquote{const vec\_bool\&}} \sphinxstyleemphasis{select\_domain}
\end{DUlineblock}

It has size \sphinxstyleemphasis{n} and is a sparsity pattern for the diagonal of
\(D\); i.e., \sphinxstyleemphasis{select\_domain} {[} \sphinxstyleemphasis{j} {]} is true if and only if
\(D_{j,j}\) is possibly non\sphinxhyphen{}zero.

\index{select\_range@\spxentry{select\_range}}\ignorespaces 

\subparagraph{select\_range}
\label{\detokenize{xsrst/cpp_sparsity:select-range}}\label{\detokenize{xsrst/cpp_sparsity:cpp-sparsity-select-range}}\label{\detokenize{xsrst/cpp_sparsity:index-6}}
The argument \sphinxstyleemphasis{select\_range} has prototype

\begin{DUlineblock}{0em}
\item[]         \sphinxcode{\sphinxupquote{const vec\_bool\&}} \sphinxstyleemphasis{select\_range}
\end{DUlineblock}

It has size \sphinxstyleemphasis{m} and is a sparsity pattern for the vector
\(r\); i.e., \sphinxstyleemphasis{select\_range} {[} \sphinxstyleemphasis{i} {]} is true if and only if
\(r_i\) is possibly non\sphinxhyphen{}zero.

\index{pattern\_out@\spxentry{pattern\_out}}\ignorespaces 

\subparagraph{pattern\_out}
\label{\detokenize{xsrst/cpp_sparsity:pattern-out}}\label{\detokenize{xsrst/cpp_sparsity:cpp-sparsity-pattern-out}}\label{\detokenize{xsrst/cpp_sparsity:index-7}}
This argument has prototype

\begin{DUlineblock}{0em}
\item[]         \sphinxcode{\sphinxupquote{sparse\_rc}} \textless{} \sphinxstyleemphasis{SizeVector} \textgreater{} \sphinxcode{\sphinxupquote{\&}} \sphinxstyleemphasis{pattern\_out}
\end{DUlineblock}

This input value of \sphinxstyleemphasis{pattern\_out} does not matter.
Upon return \sphinxstyleemphasis{pattern\_out} is a sparsity pattern for
\(J(x)\).

\index{sparsity@\spxentry{sparsity}}\index{for@\spxentry{for}}\index{component@\spxentry{component}}\index{wise@\spxentry{wise}}\index{hessian@\spxentry{hessian}}\ignorespaces 

\subparagraph{Sparsity for Component Wise Hessian}
\label{\detokenize{xsrst/cpp_sparsity:sparsity-for-component-wise-hessian}}\label{\detokenize{xsrst/cpp_sparsity:cpp-sparsity-sparsity-for-component-wise-hessian}}\label{\detokenize{xsrst/cpp_sparsity:index-8}}
Suppose that \(D\) is the identity matrix,
and only the \sphinxstyleemphasis{i}\sphinxhyphen{}th component of \sphinxstyleemphasis{r} is possibly non\sphinxhyphen{}zero.
In this case, \sphinxstyleemphasis{pattern\_out} is a sparsity pattern for
\(F_i^{(2)} ( x )\).


\subparagraph{sparse\_hes\_pattern\_xam\_cpp}
\label{\detokenize{xsrst/sparse_hes_pattern_xam_cpp:sparse-hes-pattern-xam-cpp}}\label{\detokenize{xsrst/sparse_hes_pattern_xam_cpp::doc}}
\index{sparse\_hes\_pattern\_xam\_cpp@\spxentry{sparse\_hes\_pattern\_xam\_cpp}}\index{c++:@\spxentry{c++:}}\index{hessian@\spxentry{hessian}}\index{sparsity@\spxentry{sparsity}}\index{patterns:@\spxentry{patterns:}}\index{example@\spxentry{example}}\index{test@\spxentry{test}}\ignorespaces 

\subparagraph{3.2.4.4.1 C++: Hessian Sparsity Patterns: Example and Test}
\label{\detokenize{xsrst/sparse_hes_pattern_xam_cpp:c-hessian-sparsity-patterns-example-and-test}}\label{\detokenize{xsrst/sparse_hes_pattern_xam_cpp:id1}}
\begin{sphinxVerbatim}[commandchars=\\\{\}]
\PYG{c+cp}{\PYGZsh{}}\PYG{c+cp}{ include \PYGZlt{}cstdio\PYGZgt{}}
\PYG{c+cp}{\PYGZsh{}}\PYG{c+cp}{ include \PYGZlt{}cppad}\PYG{c+cp}{/}\PYG{c+cp}{py}\PYG{c+cp}{/}\PYG{c+cp}{cppad\PYGZus{}py.hpp\PYGZgt{}}

\PYG{k+kt}{bool} \PYG{n+nf}{sparse\PYGZus{}hes\PYGZus{}pattern\PYGZus{}xam}\PYG{p}{(}\PYG{k+kt}{void}\PYG{p}{)} \PYG{p}{\PYGZob{}}
    \PYG{k}{using} \PYG{n}{cppad\PYGZus{}py}\PYG{o}{:}\PYG{o}{:}\PYG{n}{a\PYGZus{}double}\PYG{p}{;}
    \PYG{k}{using} \PYG{n}{cppad\PYGZus{}py}\PYG{o}{:}\PYG{o}{:}\PYG{n}{vec\PYGZus{}bool}\PYG{p}{;}
    \PYG{k}{using} \PYG{n}{cppad\PYGZus{}py}\PYG{o}{:}\PYG{o}{:}\PYG{n}{vec\PYGZus{}int}\PYG{p}{;}
    \PYG{k}{using} \PYG{n}{cppad\PYGZus{}py}\PYG{o}{:}\PYG{o}{:}\PYG{n}{vec\PYGZus{}double}\PYG{p}{;}
    \PYG{k}{using} \PYG{n}{cppad\PYGZus{}py}\PYG{o}{:}\PYG{o}{:}\PYG{n}{vec\PYGZus{}a\PYGZus{}double}\PYG{p}{;}
    \PYG{k}{using} \PYG{n}{cppad\PYGZus{}py}\PYG{o}{:}\PYG{o}{:}\PYG{n}{d\PYGZus{}fun}\PYG{p}{;}
    \PYG{k}{using} \PYG{n}{cppad\PYGZus{}py}\PYG{o}{:}\PYG{o}{:}\PYG{n}{sparse\PYGZus{}rc}\PYG{p}{;}
    \PYG{c+c1}{//}
    \PYG{c+c1}{// initialize return variable}
    \PYG{k+kt}{bool} \PYG{n}{ok} \PYG{o}{=} \PYG{n+nb}{true}\PYG{p}{;}
    \PYG{c+c1}{//\PYGZhy{}\PYGZhy{}\PYGZhy{}\PYGZhy{}\PYGZhy{}\PYGZhy{}\PYGZhy{}\PYGZhy{}\PYGZhy{}\PYGZhy{}\PYGZhy{}\PYGZhy{}\PYGZhy{}\PYGZhy{}\PYGZhy{}\PYGZhy{}\PYGZhy{}\PYGZhy{}\PYGZhy{}\PYGZhy{}\PYGZhy{}\PYGZhy{}\PYGZhy{}\PYGZhy{}\PYGZhy{}\PYGZhy{}\PYGZhy{}\PYGZhy{}\PYGZhy{}\PYGZhy{}\PYGZhy{}\PYGZhy{}\PYGZhy{}\PYGZhy{}\PYGZhy{}\PYGZhy{}\PYGZhy{}\PYGZhy{}\PYGZhy{}\PYGZhy{}\PYGZhy{}\PYGZhy{}\PYGZhy{}\PYGZhy{}\PYGZhy{}\PYGZhy{}\PYGZhy{}\PYGZhy{}\PYGZhy{}\PYGZhy{}\PYGZhy{}\PYGZhy{}\PYGZhy{}\PYGZhy{}\PYGZhy{}\PYGZhy{}\PYGZhy{}\PYGZhy{}\PYGZhy{}\PYGZhy{}\PYGZhy{}\PYGZhy{}\PYGZhy{}\PYGZhy{}\PYGZhy{}\PYGZhy{}\PYGZhy{}\PYGZhy{}\PYGZhy{}\PYGZhy{}\PYGZhy{}\PYGZhy{}}
    \PYG{c+c1}{// number of dependent and independent variables}
    \PYG{k+kt}{int} \PYG{n}{n} \PYG{o}{=} \PYG{l+m+mi}{3}\PYG{p}{;}
    \PYG{c+c1}{//}
    \PYG{c+c1}{// create the independent variables ax}
    \PYG{n}{vec\PYGZus{}double} \PYG{n}{x}\PYG{p}{(}\PYG{n}{n}\PYG{p}{)}\PYG{p}{;}
    \PYG{k}{for}\PYG{p}{(}\PYG{k+kt}{int} \PYG{n}{i} \PYG{o}{=} \PYG{l+m+mi}{0}\PYG{p}{;} \PYG{n}{i} \PYG{o}{\PYGZlt{}} \PYG{n}{n} \PYG{p}{;} \PYG{n}{i}\PYG{o}{+}\PYG{o}{+}\PYG{p}{)} \PYG{p}{\PYGZob{}}
        \PYG{n}{x}\PYG{p}{[}\PYG{n}{i}\PYG{p}{]} \PYG{o}{=} \PYG{n}{i} \PYG{o}{+} \PYG{l+m+mf}{2.0}\PYG{p}{;}
    \PYG{p}{\PYGZcb{}}
    \PYG{n}{vec\PYGZus{}a\PYGZus{}double} \PYG{n}{ax} \PYG{o}{=} \PYG{n}{cppad\PYGZus{}py}\PYG{o}{:}\PYG{o}{:}\PYG{n}{independent}\PYG{p}{(}\PYG{n}{x}\PYG{p}{)}\PYG{p}{;}
    \PYG{c+c1}{//}
    \PYG{c+c1}{// create dependent variables ay with ay[i] = ax[j] * ax[i]}
    \PYG{c+c1}{// where i = mod(j + 1, n)}
    \PYG{n}{vec\PYGZus{}a\PYGZus{}double} \PYG{n}{ay}\PYG{p}{(}\PYG{n}{n}\PYG{p}{)}\PYG{p}{;}
    \PYG{k}{for}\PYG{p}{(}\PYG{k+kt}{int} \PYG{n}{j} \PYG{o}{=} \PYG{l+m+mi}{0}\PYG{p}{;} \PYG{n}{j} \PYG{o}{\PYGZlt{}} \PYG{n}{n} \PYG{p}{;} \PYG{n}{j}\PYG{o}{+}\PYG{o}{+}\PYG{p}{)} \PYG{p}{\PYGZob{}}
        \PYG{k+kt}{int} \PYG{n}{i} \PYG{o}{=} \PYG{n}{j}\PYG{o}{+}\PYG{l+m+mi}{1}\PYG{p}{;}
        \PYG{k}{if}\PYG{p}{(} \PYG{n}{i} \PYG{o}{\PYGZgt{}}\PYG{o}{=} \PYG{n}{n}  \PYG{p}{)} \PYG{p}{\PYGZob{}}
            \PYG{n}{i} \PYG{o}{=} \PYG{n}{i} \PYG{o}{\PYGZhy{}} \PYG{n}{n}\PYG{p}{;}
        \PYG{p}{\PYGZcb{}}
        \PYG{n}{a\PYGZus{}double} \PYG{n}{ay\PYGZus{}i} \PYG{o}{=} \PYG{n}{ax}\PYG{p}{[}\PYG{n}{i}\PYG{p}{]} \PYG{o}{*} \PYG{n}{ax}\PYG{p}{[}\PYG{n}{j}\PYG{p}{]}\PYG{p}{;}
        \PYG{n}{ay}\PYG{p}{[}\PYG{n}{i}\PYG{p}{]} \PYG{o}{=} \PYG{n}{ay\PYGZus{}i}\PYG{p}{;}
    \PYG{p}{\PYGZcb{}}
    \PYG{c+c1}{//}
    \PYG{c+c1}{// define af corresponding to f(x)}
    \PYG{n}{d\PYGZus{}fun} \PYG{n}{f}\PYG{p}{(}\PYG{n}{ax}\PYG{p}{,} \PYG{n}{ay}\PYG{p}{)}\PYG{p}{;}
    \PYG{c+c1}{//}
    \PYG{c+c1}{// Set select\PYGZus{}d (domain) to all true, initial select\PYGZus{}r (range) to all false}
    \PYG{n}{vec\PYGZus{}bool} \PYG{n}{select\PYGZus{}d} \PYG{o}{=} \PYG{n}{vec\PYGZus{}bool}\PYG{p}{(}\PYG{n}{n}\PYG{p}{)}\PYG{p}{;}
    \PYG{n}{vec\PYGZus{}bool} \PYG{n}{select\PYGZus{}r} \PYG{o}{=} \PYG{n}{vec\PYGZus{}bool}\PYG{p}{(}\PYG{n}{n}\PYG{p}{)}\PYG{p}{;}
    \PYG{k}{for}\PYG{p}{(}\PYG{k+kt}{int} \PYG{n}{i} \PYG{o}{=} \PYG{l+m+mi}{0}\PYG{p}{;} \PYG{n}{i} \PYG{o}{\PYGZlt{}} \PYG{n}{n}\PYG{p}{;} \PYG{n}{i}\PYG{o}{+}\PYG{o}{+}\PYG{p}{)} \PYG{p}{\PYGZob{}}
        \PYG{n}{select\PYGZus{}d}\PYG{p}{[}\PYG{n}{i}\PYG{p}{]} \PYG{o}{=} \PYG{n+nb}{true}\PYG{p}{;}
        \PYG{n}{select\PYGZus{}r}\PYG{p}{[}\PYG{n}{i}\PYG{p}{]} \PYG{o}{=} \PYG{n+nb}{false}\PYG{p}{;}
    \PYG{p}{\PYGZcb{}}
    \PYG{c+c1}{//}
    \PYG{c+c1}{// only select component 0 of the range function}
    \PYG{c+c1}{// f\PYGZus{}0 (x) = x\PYGZus{}0 * x\PYGZus{}\PYGZob{}n\PYGZhy{}1\PYGZcb{}}
    \PYG{n}{select\PYGZus{}r}\PYG{p}{[}\PYG{l+m+mi}{0}\PYG{p}{]} \PYG{o}{=} \PYG{n+nb}{true}\PYG{p}{;}
    \PYG{c+c1}{//}
    \PYG{c+c1}{// loop over forward and reverse mode}
    \PYG{k}{for}\PYG{p}{(}\PYG{k+kt}{int} \PYG{n}{mode} \PYG{o}{=} \PYG{l+m+mi}{0}\PYG{p}{;} \PYG{n}{mode} \PYG{o}{\PYGZlt{}} \PYG{l+m+mi}{2}\PYG{p}{;} \PYG{n}{mode}\PYG{o}{+}\PYG{o}{+}\PYG{p}{)} \PYG{p}{\PYGZob{}}
        \PYG{n}{sparse\PYGZus{}rc} \PYG{n}{pat\PYGZus{}out} \PYG{o}{=} \PYG{n}{sparse\PYGZus{}rc}\PYG{p}{(}\PYG{p}{)}\PYG{p}{;}
        \PYG{k}{if}\PYG{p}{(} \PYG{n}{mode} \PYG{o}{=}\PYG{o}{=} \PYG{l+m+mi}{0}  \PYG{p}{)} \PYG{p}{\PYGZob{}}
            \PYG{n}{f}\PYG{p}{.}\PYG{n}{for\PYGZus{}hes\PYGZus{}sparsity}\PYG{p}{(}\PYG{n}{select\PYGZus{}d}\PYG{p}{,} \PYG{n}{select\PYGZus{}r}\PYG{p}{,} \PYG{n}{pat\PYGZus{}out}\PYG{p}{)}\PYG{p}{;}
        \PYG{p}{\PYGZcb{}}
        \PYG{k}{if}\PYG{p}{(} \PYG{n}{mode} \PYG{o}{=}\PYG{o}{=} \PYG{l+m+mi}{1}  \PYG{p}{)} \PYG{p}{\PYGZob{}}
            \PYG{n}{f}\PYG{p}{.}\PYG{n}{rev\PYGZus{}hes\PYGZus{}sparsity}\PYG{p}{(}\PYG{n}{select\PYGZus{}d}\PYG{p}{,} \PYG{n}{select\PYGZus{}r}\PYG{p}{,} \PYG{n}{pat\PYGZus{}out}\PYG{p}{)}\PYG{p}{;}
        \PYG{p}{\PYGZcb{}}
        \PYG{c+c1}{//}
        \PYG{c+c1}{// check that result is sparsity pattern for Hessian of f\PYGZus{}0 (x)}
        \PYG{n}{ok} \PYG{o}{=} \PYG{n}{ok} \PYG{o}{\PYGZam{}}\PYG{o}{\PYGZam{}} \PYG{n}{pat\PYGZus{}out}\PYG{p}{.}\PYG{n}{nnz}\PYG{p}{(}\PYG{p}{)} \PYG{o}{=}\PYG{o}{=} \PYG{l+m+mi}{2} \PYG{p}{;}
        \PYG{n}{vec\PYGZus{}int} \PYG{n}{row} \PYG{o}{=} \PYG{n}{pat\PYGZus{}out}\PYG{p}{.}\PYG{n}{row}\PYG{p}{(}\PYG{p}{)}\PYG{p}{;}
        \PYG{n}{vec\PYGZus{}int} \PYG{n}{col} \PYG{o}{=} \PYG{n}{pat\PYGZus{}out}\PYG{p}{.}\PYG{n}{col}\PYG{p}{(}\PYG{p}{)}\PYG{p}{;}
        \PYG{k}{for}\PYG{p}{(}\PYG{k+kt}{int} \PYG{n}{k} \PYG{o}{=} \PYG{l+m+mi}{0}\PYG{p}{;} \PYG{n}{k} \PYG{o}{\PYGZlt{}} \PYG{l+m+mi}{2}\PYG{p}{;} \PYG{n}{k}\PYG{o}{+}\PYG{o}{+}\PYG{p}{)} \PYG{p}{\PYGZob{}}
            \PYG{k+kt}{int} \PYG{n}{r} \PYG{o}{=} \PYG{n}{row}\PYG{p}{[}\PYG{n}{k}\PYG{p}{]}\PYG{p}{;}
            \PYG{k+kt}{int} \PYG{n}{c} \PYG{o}{=} \PYG{n}{col}\PYG{p}{[}\PYG{n}{k}\PYG{p}{]}\PYG{p}{;}
            \PYG{k}{if}\PYG{p}{(} \PYG{n}{r} \PYG{o}{\PYGZlt{}}\PYG{o}{=} \PYG{n}{c}  \PYG{p}{)} \PYG{p}{\PYGZob{}}
                \PYG{n}{ok} \PYG{o}{=} \PYG{n}{ok} \PYG{o}{\PYGZam{}}\PYG{o}{\PYGZam{}} \PYG{n}{r} \PYG{o}{=}\PYG{o}{=} \PYG{l+m+mi}{0}\PYG{p}{;}
                \PYG{n}{ok} \PYG{o}{=} \PYG{n}{ok} \PYG{o}{\PYGZam{}}\PYG{o}{\PYGZam{}} \PYG{n}{c} \PYG{o}{=}\PYG{o}{=} \PYG{n}{n}\PYG{o}{\PYGZhy{}}\PYG{l+m+mi}{1}\PYG{p}{;}
            \PYG{p}{\PYGZcb{}}
            \PYG{k}{if}\PYG{p}{(} \PYG{n}{r} \PYG{o}{\PYGZgt{}}\PYG{o}{=} \PYG{n}{c}  \PYG{p}{)} \PYG{p}{\PYGZob{}}
                \PYG{n}{ok} \PYG{o}{=} \PYG{n}{ok} \PYG{o}{\PYGZam{}}\PYG{o}{\PYGZam{}} \PYG{n}{r} \PYG{o}{=}\PYG{o}{=} \PYG{n}{n}\PYG{o}{\PYGZhy{}}\PYG{l+m+mi}{1}\PYG{p}{;}
                \PYG{n}{ok} \PYG{o}{=} \PYG{n}{ok} \PYG{o}{\PYGZam{}}\PYG{o}{\PYGZam{}} \PYG{n}{c} \PYG{o}{=}\PYG{o}{=} \PYG{l+m+mi}{0}\PYG{p}{;}
            \PYG{p}{\PYGZcb{}}
        \PYG{p}{\PYGZcb{}}
    \PYG{p}{\PYGZcb{}}
    \PYG{c+c1}{//}
    \PYG{k}{return}\PYG{p}{(} \PYG{n}{ok} \PYG{p}{)}\PYG{p}{;}
\PYG{p}{\PYGZcb{}}
\end{sphinxVerbatim}


\bigskip\hrule\bigskip


xsrst input file: \sphinxcode{\sphinxupquote{example/cplusplus/sparse\_hes\_pattern\_xam.cpp}}

\index{example@\spxentry{example}}\ignorespaces 

\subparagraph{Example}
\label{\detokenize{xsrst/cpp_sparsity:example}}\label{\detokenize{xsrst/cpp_sparsity:cpp-sparsity-example}}\label{\detokenize{xsrst/cpp_sparsity:index-9}}
{\hyperref[\detokenize{xsrst/sparse_hes_pattern_xam_cpp:id1}]{\sphinxcrossref{\DUrole{std,std-ref}{c++}}}}


\bigskip\hrule\bigskip


xsrst input file: \sphinxcode{\sphinxupquote{lib/cplusplus/sparse.cpp}}


\subparagraph{cpp\_sparse\_jac}
\label{\detokenize{xsrst/cpp_sparse_jac:cpp-sparse-jac}}\label{\detokenize{xsrst/cpp_sparse_jac::doc}}
\index{cpp\_sparse\_jac@\spxentry{cpp\_sparse\_jac}}\index{computing@\spxentry{computing}}\index{sparse@\spxentry{sparse}}\index{jacobians@\spxentry{jacobians}}\ignorespaces 

\subparagraph{3.2.4.5 Computing Sparse Jacobians}
\label{\detokenize{xsrst/cpp_sparse_jac:computing-sparse-jacobians}}\label{\detokenize{xsrst/cpp_sparse_jac:id1}}\begin{itemize}
\item {} 
{\hyperref[\detokenize{xsrst/cpp_sparse_jac:cpp-sparse-jac-syntax}]{\sphinxcrossref{\DUrole{std,std-ref}{Syntax}}}}

\item {} 
{\hyperref[\detokenize{xsrst/cpp_sparse_jac:cpp-sparse-jac-purpose}]{\sphinxcrossref{\DUrole{std,std-ref}{Purpose}}}}

\item {} 
{\hyperref[\detokenize{xsrst/cpp_sparse_jac:cpp-sparse-jac-sparse-jac-for}]{\sphinxcrossref{\DUrole{std,std-ref}{sparse\_jac\_for}}}}

\item {} 
{\hyperref[\detokenize{xsrst/cpp_sparse_jac:cpp-sparse-jac-sparse-jac-rev}]{\sphinxcrossref{\DUrole{std,std-ref}{sparse\_jac\_rev}}}}

\item {} 
{\hyperref[\detokenize{xsrst/cpp_sparse_jac:cpp-sparse-jac-f}]{\sphinxcrossref{\DUrole{std,std-ref}{f}}}}

\item {} 
{\hyperref[\detokenize{xsrst/cpp_sparse_jac:cpp-sparse-jac-subset}]{\sphinxcrossref{\DUrole{std,std-ref}{subset}}}}

\item {} 
{\hyperref[\detokenize{xsrst/cpp_sparse_jac:cpp-sparse-jac-x}]{\sphinxcrossref{\DUrole{std,std-ref}{x}}}}

\item {} 
{\hyperref[\detokenize{xsrst/cpp_sparse_jac:cpp-sparse-jac-pattern}]{\sphinxcrossref{\DUrole{std,std-ref}{pattern}}}}

\item {} 
{\hyperref[\detokenize{xsrst/cpp_sparse_jac:cpp-sparse-jac-work}]{\sphinxcrossref{\DUrole{std,std-ref}{work}}}}

\item {} 
{\hyperref[\detokenize{xsrst/cpp_sparse_jac:cpp-sparse-jac-n-sweep}]{\sphinxcrossref{\DUrole{std,std-ref}{n\_sweep}}}}

\item {} 
{\hyperref[\detokenize{xsrst/cpp_sparse_jac:cpp-sparse-jac-uses-forward}]{\sphinxcrossref{\DUrole{std,std-ref}{Uses Forward}}}}

\item {} 
{\hyperref[\detokenize{xsrst/cpp_sparse_jac:cpp-sparse-jac-example}]{\sphinxcrossref{\DUrole{std,std-ref}{Example}}}}

\end{itemize}

\index{syntax@\spxentry{syntax}}\ignorespaces 

\subparagraph{Syntax}
\label{\detokenize{xsrst/cpp_sparse_jac:syntax}}\label{\detokenize{xsrst/cpp_sparse_jac:cpp-sparse-jac-syntax}}\label{\detokenize{xsrst/cpp_sparse_jac:index-1}}
\begin{DUlineblock}{0em}
\item[] \sphinxstyleemphasis{work} =  \sphinxcode{\sphinxupquote{cppad\_py::sparse\_jac\_work}} ()
\item[] \sphinxstyleemphasis{n\_sweep} = \sphinxstyleemphasis{f} . \sphinxcode{\sphinxupquote{sparse\_jac\_for}} ( \sphinxstyleemphasis{subset} , \sphinxstyleemphasis{x} , \sphinxstyleemphasis{pattern} , \sphinxstyleemphasis{work} )
\item[] \sphinxstyleemphasis{n\_sweep} = \sphinxstyleemphasis{f} . \sphinxcode{\sphinxupquote{sparse\_jac\_rev}} ( \sphinxstyleemphasis{subset} , \sphinxstyleemphasis{x} , \sphinxstyleemphasis{pattern} , \sphinxstyleemphasis{work} )
\end{DUlineblock}

\index{purpose@\spxentry{purpose}}\ignorespaces 

\subparagraph{Purpose}
\label{\detokenize{xsrst/cpp_sparse_jac:purpose}}\label{\detokenize{xsrst/cpp_sparse_jac:cpp-sparse-jac-purpose}}\label{\detokenize{xsrst/cpp_sparse_jac:index-2}}
We use \(F : \B{R}^n \rightarrow \B{R}^m\) to denote the
function corresponding to \sphinxstyleemphasis{f} .
The syntax above takes advantage of sparsity when computing the Jacobian
\begin{equation*}
\begin{split}J(x) = F^{(1)} (x)\end{split}
\end{equation*}
In the sparse case, this should be faster and take less memory than
{\hyperref[\detokenize{xsrst/cpp_fun_jacobian:id1}]{\sphinxcrossref{\DUrole{std,std-ref}{cpp\_fun\_jacobian}}}}.
We use the notation \(J_{i,j} (x)\) to denote the partial of
\(F_i (x)\) with respect to \(x_j\).

\index{sparse\_jac\_for@\spxentry{sparse\_jac\_for}}\ignorespaces 

\subparagraph{sparse\_jac\_for}
\label{\detokenize{xsrst/cpp_sparse_jac:sparse-jac-for}}\label{\detokenize{xsrst/cpp_sparse_jac:cpp-sparse-jac-sparse-jac-for}}\label{\detokenize{xsrst/cpp_sparse_jac:index-3}}
This function uses first order forward mode sweeps {\hyperref[\detokenize{xsrst/cpp_fun_forward:id1}]{\sphinxcrossref{\DUrole{std,std-ref}{cpp\_fun\_forward}}}}
to compute multiple columns of the Jacobian at the same time.

\index{sparse\_jac\_rev@\spxentry{sparse\_jac\_rev}}\ignorespaces 

\subparagraph{sparse\_jac\_rev}
\label{\detokenize{xsrst/cpp_sparse_jac:sparse-jac-rev}}\label{\detokenize{xsrst/cpp_sparse_jac:cpp-sparse-jac-sparse-jac-rev}}\label{\detokenize{xsrst/cpp_sparse_jac:index-4}}
This function uses first order reverse mode sweeps {\hyperref[\detokenize{xsrst/cpp_fun_reverse:id1}]{\sphinxcrossref{\DUrole{std,std-ref}{cpp\_fun\_reverse}}}}
to compute multiple rows of the Jacobian at the same time.

\index{f@\spxentry{f}}\ignorespaces 

\subparagraph{f}
\label{\detokenize{xsrst/cpp_sparse_jac:f}}\label{\detokenize{xsrst/cpp_sparse_jac:cpp-sparse-jac-f}}\label{\detokenize{xsrst/cpp_sparse_jac:index-5}}
This object has prototype

\begin{DUlineblock}{0em}
\item[]         \sphinxcode{\sphinxupquote{ADFun}} \textless{} \sphinxstyleemphasis{Base} \textgreater{} \sphinxstyleemphasis{f}
\end{DUlineblock}

Note that the Taylor coefficients stored in \sphinxstyleemphasis{f} are affected
by this operation; see
{\hyperref[\detokenize{xsrst/cpp_sparse_jac:cpp-sparse-jac-uses-forward}]{\sphinxcrossref{\DUrole{std,std-ref}{uses\_forward}}}} below.

\index{subset@\spxentry{subset}}\ignorespaces 

\subparagraph{subset}
\label{\detokenize{xsrst/cpp_sparse_jac:subset}}\label{\detokenize{xsrst/cpp_sparse_jac:cpp-sparse-jac-subset}}\label{\detokenize{xsrst/cpp_sparse_jac:index-6}}
This argument has prototype

\begin{DUlineblock}{0em}
\item[]         \sphinxcode{\sphinxupquote{sparse\_rcv\&}} \sphinxstyleemphasis{subset}
\end{DUlineblock}

Its row size is \sphinxstyleemphasis{subset} . \sphinxcode{\sphinxupquote{nr}} () == \sphinxstyleemphasis{m} ,
and its column size is \sphinxstyleemphasis{subset} . \sphinxcode{\sphinxupquote{nc}} () == \sphinxstyleemphasis{n} .
It specifies which elements of the Jacobian are computed.
The input value of its value vector
\sphinxstyleemphasis{subset} . \sphinxcode{\sphinxupquote{val}} () does not matter.
Upon return it contains the value of the corresponding elements
of the Jacobian.
All of the row, column pairs in \sphinxstyleemphasis{subset} must also appear in
\sphinxstyleemphasis{pattern} ; i.e., they must be possibly non\sphinxhyphen{}zero.

\index{x@\spxentry{x}}\ignorespaces 

\subparagraph{x}
\label{\detokenize{xsrst/cpp_sparse_jac:x}}\label{\detokenize{xsrst/cpp_sparse_jac:cpp-sparse-jac-x}}\label{\detokenize{xsrst/cpp_sparse_jac:index-7}}
This argument has prototype

\begin{DUlineblock}{0em}
\item[]         \sphinxcode{\sphinxupquote{const vec\_double\&}} \sphinxstyleemphasis{x}
\end{DUlineblock}

and its size is \sphinxstyleemphasis{n} .
It specifies the point at which to evaluate the Jacobian \(J(x)\).

\index{pattern@\spxentry{pattern}}\ignorespaces 

\subparagraph{pattern}
\label{\detokenize{xsrst/cpp_sparse_jac:pattern}}\label{\detokenize{xsrst/cpp_sparse_jac:cpp-sparse-jac-pattern}}\label{\detokenize{xsrst/cpp_sparse_jac:index-8}}
This argument has prototype

\begin{DUlineblock}{0em}
\item[]         \sphinxcode{\sphinxupquote{const sparse\_rc\&}} \sphinxstyleemphasis{pattern}
\end{DUlineblock}

Its row size is \sphinxstyleemphasis{pattern} . \sphinxcode{\sphinxupquote{nr}} () == \sphinxstyleemphasis{m} ,
and its column size is \sphinxstyleemphasis{pattern} . \sphinxcode{\sphinxupquote{nc}} () == \sphinxstyleemphasis{n} .
It is a sparsity pattern for the Jacobian \(J(x)\).
This argument is not used (and need not satisfy any conditions),
when {\hyperref[\detokenize{xsrst/cpp_sparse_jac:cpp-sparse-jac-work}]{\sphinxcrossref{\DUrole{std,std-ref}{work}}}} is non\sphinxhyphen{}empty.

\index{work@\spxentry{work}}\ignorespaces 

\subparagraph{work}
\label{\detokenize{xsrst/cpp_sparse_jac:work}}\label{\detokenize{xsrst/cpp_sparse_jac:cpp-sparse-jac-work}}\label{\detokenize{xsrst/cpp_sparse_jac:index-9}}
This argument has prototype

\begin{DUlineblock}{0em}
\item[]         \sphinxcode{\sphinxupquote{sparse\_jac\_work\&}} \sphinxstyleemphasis{work}
\end{DUlineblock}

We refer to its initial value,
and its value after \sphinxstyleemphasis{work} . \sphinxcode{\sphinxupquote{clear}} () , as empty.
If it is empty, information is stored in \sphinxstyleemphasis{work} .
This can be used to reduce computation when
a future call is for the same object \sphinxstyleemphasis{f} ,
the same member function \sphinxcode{\sphinxupquote{sparse\_jac\_for}} or \sphinxcode{\sphinxupquote{sparse\_jac\_rev}} ,
and the same subset of the Jacobian.
If any of these values change, use \sphinxstyleemphasis{work} . \sphinxcode{\sphinxupquote{clear}} () to
empty this structure.

\index{n\_sweep@\spxentry{n\_sweep}}\ignorespaces 

\subparagraph{n\_sweep}
\label{\detokenize{xsrst/cpp_sparse_jac:n-sweep}}\label{\detokenize{xsrst/cpp_sparse_jac:cpp-sparse-jac-n-sweep}}\label{\detokenize{xsrst/cpp_sparse_jac:index-10}}
The return value \sphinxstyleemphasis{n\_sweep} has prototype

\begin{DUlineblock}{0em}
\item[]         \sphinxcode{\sphinxupquote{int}} \sphinxstyleemphasis{n\_sweep}
\end{DUlineblock}

If \sphinxcode{\sphinxupquote{sparse\_jac\_for}} ( \sphinxcode{\sphinxupquote{sparse\_jac\_rev}} ) is used,
\sphinxstyleemphasis{n\_sweep} is the number of first order forward (reverse) sweeps
used to compute the requested Jacobian values.
This is proportional to the total computational work,
not counting the zero order forward sweep,
or combining multiple columns (rows) into a single sweep.

\index{uses@\spxentry{uses}}\index{forward@\spxentry{forward}}\ignorespaces 

\subparagraph{Uses Forward}
\label{\detokenize{xsrst/cpp_sparse_jac:uses-forward}}\label{\detokenize{xsrst/cpp_sparse_jac:cpp-sparse-jac-uses-forward}}\label{\detokenize{xsrst/cpp_sparse_jac:index-11}}
After each call to {\hyperref[\detokenize{xsrst/cpp_fun_forward:id1}]{\sphinxcrossref{\DUrole{std,std-ref}{cpp\_fun\_forward}}}},
the object \sphinxstyleemphasis{f} contains the corresponding Taylor coefficients
for all the variables in the operation sequence..
After a call to \sphinxcode{\sphinxupquote{sparse\_jac\_forward}} or \sphinxcode{\sphinxupquote{sparse\_jac\_rev}} ,
the zero order coefficients correspond to

\begin{DUlineblock}{0em}
\item[]         \sphinxstyleemphasis{f} . \sphinxcode{\sphinxupquote{forward(0}} , \sphinxstyleemphasis{x} )
\end{DUlineblock}

All the other forward mode coefficients are unspecified.


\subparagraph{sparse\_jac\_xam\_cpp}
\label{\detokenize{xsrst/sparse_jac_xam_cpp:sparse-jac-xam-cpp}}\label{\detokenize{xsrst/sparse_jac_xam_cpp::doc}}
\index{sparse\_jac\_xam\_cpp@\spxentry{sparse\_jac\_xam\_cpp}}\index{c++:@\spxentry{c++:}}\index{computing@\spxentry{computing}}\index{sparse@\spxentry{sparse}}\index{jacobians:@\spxentry{jacobians:}}\index{example@\spxentry{example}}\index{test@\spxentry{test}}\ignorespaces 

\subparagraph{3.2.4.5.1 C++: Computing Sparse Jacobians: Example and Test}
\label{\detokenize{xsrst/sparse_jac_xam_cpp:c-computing-sparse-jacobians-example-and-test}}\label{\detokenize{xsrst/sparse_jac_xam_cpp:id1}}
\begin{sphinxVerbatim}[commandchars=\\\{\}]
\PYG{c+cp}{\PYGZsh{}}\PYG{c+cp}{ include \PYGZlt{}cstdio\PYGZgt{}}
\PYG{c+cp}{\PYGZsh{}}\PYG{c+cp}{ include \PYGZlt{}cppad}\PYG{c+cp}{/}\PYG{c+cp}{py}\PYG{c+cp}{/}\PYG{c+cp}{cppad\PYGZus{}py.hpp\PYGZgt{}}

\PYG{k+kt}{bool} \PYG{n+nf}{sparse\PYGZus{}jac\PYGZus{}xam}\PYG{p}{(}\PYG{k+kt}{void}\PYG{p}{)} \PYG{p}{\PYGZob{}}
    \PYG{k}{using} \PYG{n}{cppad\PYGZus{}py}\PYG{o}{:}\PYG{o}{:}\PYG{n}{a\PYGZus{}double}\PYG{p}{;}
    \PYG{k}{using} \PYG{n}{cppad\PYGZus{}py}\PYG{o}{:}\PYG{o}{:}\PYG{n}{vec\PYGZus{}int}\PYG{p}{;}
    \PYG{k}{using} \PYG{n}{cppad\PYGZus{}py}\PYG{o}{:}\PYG{o}{:}\PYG{n}{vec\PYGZus{}double}\PYG{p}{;}
    \PYG{k}{using} \PYG{n}{cppad\PYGZus{}py}\PYG{o}{:}\PYG{o}{:}\PYG{n}{vec\PYGZus{}a\PYGZus{}double}\PYG{p}{;}
    \PYG{k}{using} \PYG{n}{cppad\PYGZus{}py}\PYG{o}{:}\PYG{o}{:}\PYG{n}{d\PYGZus{}fun}\PYG{p}{;}
    \PYG{k}{using} \PYG{n}{cppad\PYGZus{}py}\PYG{o}{:}\PYG{o}{:}\PYG{n}{sparse\PYGZus{}rc}\PYG{p}{;}
    \PYG{k}{using} \PYG{n}{cppad\PYGZus{}py}\PYG{o}{:}\PYG{o}{:}\PYG{n}{sparse\PYGZus{}rcv}\PYG{p}{;}
    \PYG{k}{using} \PYG{n}{cppad\PYGZus{}py}\PYG{o}{:}\PYG{o}{:}\PYG{n}{sparse\PYGZus{}jac\PYGZus{}work}\PYG{p}{;}
    \PYG{c+c1}{//}
    \PYG{c+c1}{// initialize return variable}
    \PYG{k+kt}{bool} \PYG{n}{ok} \PYG{o}{=} \PYG{n+nb}{true}\PYG{p}{;}
    \PYG{c+c1}{//\PYGZhy{}\PYGZhy{}\PYGZhy{}\PYGZhy{}\PYGZhy{}\PYGZhy{}\PYGZhy{}\PYGZhy{}\PYGZhy{}\PYGZhy{}\PYGZhy{}\PYGZhy{}\PYGZhy{}\PYGZhy{}\PYGZhy{}\PYGZhy{}\PYGZhy{}\PYGZhy{}\PYGZhy{}\PYGZhy{}\PYGZhy{}\PYGZhy{}\PYGZhy{}\PYGZhy{}\PYGZhy{}\PYGZhy{}\PYGZhy{}\PYGZhy{}\PYGZhy{}\PYGZhy{}\PYGZhy{}\PYGZhy{}\PYGZhy{}\PYGZhy{}\PYGZhy{}\PYGZhy{}\PYGZhy{}\PYGZhy{}\PYGZhy{}\PYGZhy{}\PYGZhy{}\PYGZhy{}\PYGZhy{}\PYGZhy{}\PYGZhy{}\PYGZhy{}\PYGZhy{}\PYGZhy{}\PYGZhy{}\PYGZhy{}\PYGZhy{}\PYGZhy{}\PYGZhy{}\PYGZhy{}\PYGZhy{}\PYGZhy{}\PYGZhy{}\PYGZhy{}\PYGZhy{}\PYGZhy{}\PYGZhy{}\PYGZhy{}\PYGZhy{}\PYGZhy{}\PYGZhy{}\PYGZhy{}\PYGZhy{}\PYGZhy{}\PYGZhy{}\PYGZhy{}\PYGZhy{}\PYGZhy{}}
    \PYG{c+c1}{// number of dependent and independent variables}
    \PYG{k+kt}{int} \PYG{n}{n} \PYG{o}{=} \PYG{l+m+mi}{3}\PYG{p}{;}
    \PYG{c+c1}{// one}
    \PYG{n}{a\PYGZus{}double} \PYG{n}{aone}\PYG{p}{(}\PYG{l+m+mf}{1.0}\PYG{p}{)}\PYG{p}{;}
    \PYG{c+c1}{//}
    \PYG{c+c1}{// create the independent variables ax}
    \PYG{n}{vec\PYGZus{}double} \PYG{n}{x}\PYG{p}{(}\PYG{n}{n}\PYG{p}{)}\PYG{p}{;}
    \PYG{k}{for}\PYG{p}{(}\PYG{k+kt}{int} \PYG{n}{i} \PYG{o}{=} \PYG{l+m+mi}{0}\PYG{p}{;} \PYG{n}{i} \PYG{o}{\PYGZlt{}} \PYG{n}{n} \PYG{p}{;} \PYG{n}{i}\PYG{o}{+}\PYG{o}{+}\PYG{p}{)} \PYG{p}{\PYGZob{}}
        \PYG{n}{x}\PYG{p}{[}\PYG{n}{i}\PYG{p}{]} \PYG{o}{=} \PYG{n}{i} \PYG{o}{+} \PYG{l+m+mf}{2.0}\PYG{p}{;}
    \PYG{p}{\PYGZcb{}}
    \PYG{n}{vec\PYGZus{}a\PYGZus{}double} \PYG{n}{ax} \PYG{o}{=} \PYG{n}{cppad\PYGZus{}py}\PYG{o}{:}\PYG{o}{:}\PYG{n}{independent}\PYG{p}{(}\PYG{n}{x}\PYG{p}{)}\PYG{p}{;}
    \PYG{c+c1}{//}
    \PYG{c+c1}{// create dependent variables ay with ay[i] = (j+1) * ax[j]}
    \PYG{c+c1}{// where i = mod(j + 1, n)}
    \PYG{n}{vec\PYGZus{}a\PYGZus{}double} \PYG{n}{ay}\PYG{p}{(}\PYG{n}{n}\PYG{p}{)}\PYG{p}{;}
    \PYG{k}{for}\PYG{p}{(}\PYG{k+kt}{int} \PYG{n}{j} \PYG{o}{=} \PYG{l+m+mi}{0}\PYG{p}{;} \PYG{n}{j} \PYG{o}{\PYGZlt{}} \PYG{n}{n} \PYG{p}{;} \PYG{n}{j}\PYG{o}{+}\PYG{o}{+}\PYG{p}{)} \PYG{p}{\PYGZob{}}
        \PYG{k+kt}{int} \PYG{n}{i} \PYG{o}{=} \PYG{n}{j}\PYG{o}{+}\PYG{l+m+mi}{1}\PYG{p}{;}
        \PYG{k}{if}\PYG{p}{(} \PYG{n}{i} \PYG{o}{\PYGZgt{}}\PYG{o}{=} \PYG{n}{n}  \PYG{p}{)} \PYG{p}{\PYGZob{}}
            \PYG{n}{i} \PYG{o}{=} \PYG{n}{i} \PYG{o}{\PYGZhy{}} \PYG{n}{n}\PYG{p}{;}
        \PYG{p}{\PYGZcb{}}
        \PYG{n}{a\PYGZus{}double} \PYG{n}{aj}\PYG{p}{(}\PYG{n}{j}\PYG{p}{)}\PYG{p}{;}
        \PYG{n}{a\PYGZus{}double} \PYG{n}{ay\PYGZus{}i} \PYG{o}{=} \PYG{p}{(}\PYG{n}{aj} \PYG{o}{+} \PYG{n}{aone}\PYG{p}{)} \PYG{o}{*} \PYG{n}{ax}\PYG{p}{[}\PYG{n}{j}\PYG{p}{]}\PYG{p}{;}
        \PYG{n}{ay}\PYG{p}{[}\PYG{n}{i}\PYG{p}{]} \PYG{o}{=} \PYG{n}{ay\PYGZus{}i}\PYG{p}{;}
    \PYG{p}{\PYGZcb{}}
    \PYG{c+c1}{//}
    \PYG{c+c1}{// define af corresponding to f(x)}
    \PYG{n}{d\PYGZus{}fun} \PYG{n}{f}\PYG{p}{(}\PYG{n}{ax}\PYG{p}{,} \PYG{n}{ay}\PYG{p}{)}\PYG{p}{;}
    \PYG{c+c1}{//}
    \PYG{c+c1}{// sparsity pattern for identity matrix}
    \PYG{n}{sparse\PYGZus{}rc} \PYG{n}{pat\PYGZus{}eye} \PYG{o}{=} \PYG{n}{sparse\PYGZus{}rc}\PYG{p}{(}\PYG{p}{)}\PYG{p}{;}
    \PYG{n}{pat\PYGZus{}eye}\PYG{p}{.}\PYG{n}{resize}\PYG{p}{(}\PYG{n}{n}\PYG{p}{,} \PYG{n}{n}\PYG{p}{,} \PYG{n}{n}\PYG{p}{)}\PYG{p}{;}
    \PYG{k}{for}\PYG{p}{(}\PYG{k+kt}{int} \PYG{n}{k} \PYG{o}{=} \PYG{l+m+mi}{0}\PYG{p}{;} \PYG{n}{k} \PYG{o}{\PYGZlt{}} \PYG{n}{n}\PYG{p}{;} \PYG{n}{k}\PYG{o}{+}\PYG{o}{+}\PYG{p}{)} \PYG{p}{\PYGZob{}}
        \PYG{n}{pat\PYGZus{}eye}\PYG{p}{.}\PYG{n}{put}\PYG{p}{(}\PYG{n}{k}\PYG{p}{,} \PYG{n}{k}\PYG{p}{,} \PYG{n}{k}\PYG{p}{)}\PYG{p}{;}
    \PYG{p}{\PYGZcb{}}
    \PYG{c+c1}{//}
    \PYG{c+c1}{// sparsity pattern for the Jacobian}
    \PYG{n}{sparse\PYGZus{}rc} \PYG{n}{pat\PYGZus{}jac} \PYG{o}{=} \PYG{n}{sparse\PYGZus{}rc}\PYG{p}{(}\PYG{p}{)}\PYG{p}{;}
    \PYG{n}{f}\PYG{p}{.}\PYG{n}{for\PYGZus{}jac\PYGZus{}sparsity}\PYG{p}{(}\PYG{n}{pat\PYGZus{}eye}\PYG{p}{,} \PYG{n}{pat\PYGZus{}jac}\PYG{p}{)}\PYG{p}{;}
    \PYG{c+c1}{//}
    \PYG{c+c1}{// loop over forward and reverse mode}
    \PYG{k}{for}\PYG{p}{(}\PYG{k+kt}{int} \PYG{n}{mode} \PYG{o}{=} \PYG{l+m+mi}{0}\PYG{p}{;} \PYG{n}{mode} \PYG{o}{\PYGZlt{}} \PYG{l+m+mi}{2}\PYG{p}{;} \PYG{n}{mode}\PYG{o}{+}\PYG{o}{+}\PYG{p}{)} \PYG{p}{\PYGZob{}}
        \PYG{c+c1}{// compute all possibly non\PYGZhy{}zero entries in Jacobian}
        \PYG{n}{sparse\PYGZus{}rcv} \PYG{n}{subset} \PYG{o}{=} \PYG{n}{sparse\PYGZus{}rcv}\PYG{p}{(}\PYG{n}{pat\PYGZus{}jac}\PYG{p}{)}\PYG{p}{;}
        \PYG{c+c1}{// work space used to save time for multiple calls}
        \PYG{n}{sparse\PYGZus{}jac\PYGZus{}work} \PYG{n}{work} \PYG{o}{=} \PYG{n}{sparse\PYGZus{}jac\PYGZus{}work}\PYG{p}{(}\PYG{p}{)}\PYG{p}{;}
        \PYG{k}{if}\PYG{p}{(} \PYG{n}{mode} \PYG{o}{=}\PYG{o}{=} \PYG{l+m+mi}{0}  \PYG{p}{)} \PYG{p}{\PYGZob{}}
            \PYG{n}{f}\PYG{p}{.}\PYG{n}{sparse\PYGZus{}jac\PYGZus{}for}\PYG{p}{(}\PYG{n}{subset}\PYG{p}{,} \PYG{n}{x}\PYG{p}{,} \PYG{n}{pat\PYGZus{}jac}\PYG{p}{,} \PYG{n}{work}\PYG{p}{)}\PYG{p}{;}
        \PYG{p}{\PYGZcb{}}
        \PYG{k}{if}\PYG{p}{(} \PYG{n}{mode} \PYG{o}{=}\PYG{o}{=} \PYG{l+m+mi}{1}  \PYG{p}{)} \PYG{p}{\PYGZob{}}
            \PYG{n}{f}\PYG{p}{.}\PYG{n}{sparse\PYGZus{}jac\PYGZus{}rev}\PYG{p}{(}\PYG{n}{subset}\PYG{p}{,} \PYG{n}{x}\PYG{p}{,} \PYG{n}{pat\PYGZus{}jac}\PYG{p}{,} \PYG{n}{work}\PYG{p}{)}\PYG{p}{;}
        \PYG{p}{\PYGZcb{}}
        \PYG{c+c1}{//}
        \PYG{c+c1}{// check result}
        \PYG{n}{ok} \PYG{o}{=} \PYG{n}{ok} \PYG{o}{\PYGZam{}}\PYG{o}{\PYGZam{}} \PYG{n}{n} \PYG{o}{=}\PYG{o}{=} \PYG{n}{subset}\PYG{p}{.}\PYG{n}{nnz}\PYG{p}{(}\PYG{p}{)}\PYG{p}{;}
        \PYG{n}{vec\PYGZus{}int} \PYG{n}{col\PYGZus{}major} \PYG{o}{=} \PYG{n}{subset}\PYG{p}{.}\PYG{n}{col\PYGZus{}major}\PYG{p}{(}\PYG{p}{)}\PYG{p}{;}
        \PYG{n}{vec\PYGZus{}int} \PYG{n}{row} \PYG{o}{=} \PYG{n}{subset}\PYG{p}{.}\PYG{n}{row}\PYG{p}{(}\PYG{p}{)}\PYG{p}{;}
        \PYG{n}{vec\PYGZus{}int} \PYG{n}{col} \PYG{o}{=} \PYG{n}{subset}\PYG{p}{.}\PYG{n}{col}\PYG{p}{(}\PYG{p}{)}\PYG{p}{;}
        \PYG{n}{vec\PYGZus{}double} \PYG{n}{val} \PYG{o}{=} \PYG{n}{subset}\PYG{p}{.}\PYG{n}{val}\PYG{p}{(}\PYG{p}{)}\PYG{p}{;}
        \PYG{k}{for}\PYG{p}{(}\PYG{k+kt}{int} \PYG{n}{k} \PYG{o}{=} \PYG{l+m+mi}{0}\PYG{p}{;} \PYG{n}{k} \PYG{o}{\PYGZlt{}} \PYG{n}{n}\PYG{p}{;} \PYG{n}{k}\PYG{o}{+}\PYG{o}{+}\PYG{p}{)} \PYG{p}{\PYGZob{}}
            \PYG{k+kt}{int} \PYG{n}{ell} \PYG{o}{=} \PYG{n}{col\PYGZus{}major}\PYG{p}{[}\PYG{n}{k}\PYG{p}{]}\PYG{p}{;}
            \PYG{k+kt}{int} \PYG{n}{r} \PYG{o}{=} \PYG{n}{row}\PYG{p}{[}\PYG{n}{ell}\PYG{p}{]}\PYG{p}{;}
            \PYG{k+kt}{int} \PYG{n}{c} \PYG{o}{=} \PYG{n}{col}\PYG{p}{[}\PYG{n}{ell}\PYG{p}{]}\PYG{p}{;}
            \PYG{k+kt}{double} \PYG{n}{v} \PYG{o}{=} \PYG{n}{val}\PYG{p}{[}\PYG{n}{ell}\PYG{p}{]}\PYG{p}{;}
            \PYG{k+kt}{int} \PYG{n}{i} \PYG{o}{=} \PYG{n}{c}\PYG{o}{+}\PYG{l+m+mi}{1}\PYG{p}{;}
            \PYG{k}{if}\PYG{p}{(} \PYG{n}{i} \PYG{o}{\PYGZgt{}}\PYG{o}{=}  \PYG{n}{n}  \PYG{p}{)} \PYG{p}{\PYGZob{}}
                \PYG{n}{i} \PYG{o}{=} \PYG{n}{i} \PYG{o}{\PYGZhy{}} \PYG{n}{n}\PYG{p}{;}
            \PYG{p}{\PYGZcb{}}
            \PYG{n}{ok} \PYG{o}{=} \PYG{n}{ok} \PYG{o}{\PYGZam{}}\PYG{o}{\PYGZam{}} \PYG{n}{c} \PYG{o}{=}\PYG{o}{=} \PYG{n}{k}\PYG{p}{;}
            \PYG{n}{ok} \PYG{o}{=} \PYG{n}{ok} \PYG{o}{\PYGZam{}}\PYG{o}{\PYGZam{}} \PYG{n}{r} \PYG{o}{=}\PYG{o}{=} \PYG{n}{i}\PYG{p}{;}
            \PYG{n}{ok} \PYG{o}{=} \PYG{n}{ok} \PYG{o}{\PYGZam{}}\PYG{o}{\PYGZam{}} \PYG{n}{v} \PYG{o}{=}\PYG{o}{=} \PYG{n}{c} \PYG{o}{+} \PYG{l+m+mf}{1.0}\PYG{p}{;}
        \PYG{p}{\PYGZcb{}}
    \PYG{p}{\PYGZcb{}}
    \PYG{c+c1}{//}
    \PYG{k}{return}\PYG{p}{(} \PYG{n}{ok} \PYG{p}{)}\PYG{p}{;}
\PYG{p}{\PYGZcb{}}
\end{sphinxVerbatim}


\bigskip\hrule\bigskip


xsrst input file: \sphinxcode{\sphinxupquote{example/cplusplus/sparse\_jac\_xam.cpp}}

\index{example@\spxentry{example}}\ignorespaces 

\subparagraph{Example}
\label{\detokenize{xsrst/cpp_sparse_jac:example}}\label{\detokenize{xsrst/cpp_sparse_jac:cpp-sparse-jac-example}}\label{\detokenize{xsrst/cpp_sparse_jac:index-12}}
{\hyperref[\detokenize{xsrst/sparse_jac_xam_cpp:id1}]{\sphinxcrossref{\DUrole{std,std-ref}{sparse\_jac\_xam\_cpp}}}}


\bigskip\hrule\bigskip


xsrst input file: \sphinxcode{\sphinxupquote{lib/cplusplus/sparse.cpp}}


\subparagraph{cpp\_sparse\_hes}
\label{\detokenize{xsrst/cpp_sparse_hes:cpp-sparse-hes}}\label{\detokenize{xsrst/cpp_sparse_hes::doc}}
\index{cpp\_sparse\_hes@\spxentry{cpp\_sparse\_hes}}\index{computing@\spxentry{computing}}\index{sparse@\spxentry{sparse}}\index{hessians@\spxentry{hessians}}\ignorespaces 

\subparagraph{3.2.4.6 Computing Sparse Hessians}
\label{\detokenize{xsrst/cpp_sparse_hes:computing-sparse-hessians}}\label{\detokenize{xsrst/cpp_sparse_hes:id1}}\begin{itemize}
\item {} 
{\hyperref[\detokenize{xsrst/cpp_sparse_hes:cpp-sparse-hes-syntax}]{\sphinxcrossref{\DUrole{std,std-ref}{Syntax}}}}

\item {} 
{\hyperref[\detokenize{xsrst/cpp_sparse_hes:cpp-sparse-hes-purpose}]{\sphinxcrossref{\DUrole{std,std-ref}{Purpose}}}}

\item {} 
{\hyperref[\detokenize{xsrst/cpp_sparse_hes:cpp-sparse-hes-f}]{\sphinxcrossref{\DUrole{std,std-ref}{f}}}}

\item {} 
{\hyperref[\detokenize{xsrst/cpp_sparse_hes:cpp-sparse-hes-subset}]{\sphinxcrossref{\DUrole{std,std-ref}{subset}}}}

\item {} 
{\hyperref[\detokenize{xsrst/cpp_sparse_hes:cpp-sparse-hes-x}]{\sphinxcrossref{\DUrole{std,std-ref}{x}}}}

\item {} 
{\hyperref[\detokenize{xsrst/cpp_sparse_hes:cpp-sparse-hes-r}]{\sphinxcrossref{\DUrole{std,std-ref}{r}}}}

\item {} 
{\hyperref[\detokenize{xsrst/cpp_sparse_hes:cpp-sparse-hes-pattern}]{\sphinxcrossref{\DUrole{std,std-ref}{pattern}}}}

\item {} 
{\hyperref[\detokenize{xsrst/cpp_sparse_hes:cpp-sparse-hes-work}]{\sphinxcrossref{\DUrole{std,std-ref}{work}}}}

\item {} 
{\hyperref[\detokenize{xsrst/cpp_sparse_hes:cpp-sparse-hes-n-sweep}]{\sphinxcrossref{\DUrole{std,std-ref}{n\_sweep}}}}

\item {} 
{\hyperref[\detokenize{xsrst/cpp_sparse_hes:cpp-sparse-hes-uses-forward}]{\sphinxcrossref{\DUrole{std,std-ref}{Uses Forward}}}}

\item {} 
{\hyperref[\detokenize{xsrst/cpp_sparse_hes:cpp-sparse-hes-example}]{\sphinxcrossref{\DUrole{std,std-ref}{Example}}}}

\end{itemize}

\index{syntax@\spxentry{syntax}}\ignorespaces 

\subparagraph{Syntax}
\label{\detokenize{xsrst/cpp_sparse_hes:syntax}}\label{\detokenize{xsrst/cpp_sparse_hes:cpp-sparse-hes-syntax}}\label{\detokenize{xsrst/cpp_sparse_hes:index-1}}
\begin{DUlineblock}{0em}
\item[] \sphinxstyleemphasis{work} =  \sphinxcode{\sphinxupquote{cppad\_py::sparse\_hes\_work}} ()
\item[] \sphinxstyleemphasis{n\_sweep} = \sphinxstyleemphasis{f} . \sphinxcode{\sphinxupquote{sparse\_hes}} ( \sphinxstyleemphasis{subset} , \sphinxstyleemphasis{x} , \sphinxstyleemphasis{r} , \sphinxstyleemphasis{pattern} , \sphinxstyleemphasis{work} )
\end{DUlineblock}

\index{purpose@\spxentry{purpose}}\ignorespaces 

\subparagraph{Purpose}
\label{\detokenize{xsrst/cpp_sparse_hes:purpose}}\label{\detokenize{xsrst/cpp_sparse_hes:cpp-sparse-hes-purpose}}\label{\detokenize{xsrst/cpp_sparse_hes:index-2}}
We use \(F : \B{R}^n \rightarrow \B{R}^m\) to denote the
function corresponding to \sphinxstyleemphasis{f} .
Given a vector \(r \in \B{R}^m\), define
\begin{equation*}
\begin{split}H(x) = (r^\R{T} F)^{(2)} ( x )\end{split}
\end{equation*}
This routine takes advantage of sparsity when computing elements
of the Hessian \(H(x)\).

\index{f@\spxentry{f}}\ignorespaces 

\subparagraph{f}
\label{\detokenize{xsrst/cpp_sparse_hes:f}}\label{\detokenize{xsrst/cpp_sparse_hes:cpp-sparse-hes-f}}\label{\detokenize{xsrst/cpp_sparse_hes:index-3}}
This object has prototype

\begin{DUlineblock}{0em}
\item[]         \sphinxcode{\sphinxupquote{ADFun}} \textless{} \sphinxstyleemphasis{Base} \textgreater{} \sphinxstyleemphasis{f}
\end{DUlineblock}

Note that the Taylor coefficients stored in \sphinxstyleemphasis{f} are affected
by this operation; see
{\hyperref[\detokenize{xsrst/cpp_sparse_hes:cpp-sparse-hes-uses-forward}]{\sphinxcrossref{\DUrole{std,std-ref}{uses\_forward}}}} below.

\index{subset@\spxentry{subset}}\ignorespaces 

\subparagraph{subset}
\label{\detokenize{xsrst/cpp_sparse_hes:subset}}\label{\detokenize{xsrst/cpp_sparse_hes:cpp-sparse-hes-subset}}\label{\detokenize{xsrst/cpp_sparse_hes:index-4}}
This argument has prototype

\begin{DUlineblock}{0em}
\item[]         \sphinxcode{\sphinxupquote{sparse\_rcv\&}} \sphinxstyleemphasis{subset}
\end{DUlineblock}

Its row size and column size is \sphinxstyleemphasis{n} ; i.e.,
\sphinxstyleemphasis{subset} . \sphinxcode{\sphinxupquote{nr}} () == \sphinxstyleemphasis{n} and \sphinxstyleemphasis{subset} . \sphinxcode{\sphinxupquote{nc}} () == \sphinxstyleemphasis{n} .
It specifies which elements of the Hessian are computed.
The input value of its value vector
\sphinxstyleemphasis{subset} . \sphinxcode{\sphinxupquote{val}} () does not matter.
Upon return it contains the value of the corresponding elements
of the Jacobian.
All of the row, column pairs in \sphinxstyleemphasis{subset} must also appear in
\sphinxstyleemphasis{pattern} ; i.e., they must be possibly non\sphinxhyphen{}zero.

\index{x@\spxentry{x}}\ignorespaces 

\subparagraph{x}
\label{\detokenize{xsrst/cpp_sparse_hes:x}}\label{\detokenize{xsrst/cpp_sparse_hes:cpp-sparse-hes-x}}\label{\detokenize{xsrst/cpp_sparse_hes:index-5}}
This argument has prototype

\begin{DUlineblock}{0em}
\item[]         \sphinxcode{\sphinxupquote{const vec\_double\&}} \sphinxstyleemphasis{x}
\end{DUlineblock}

and its size is \sphinxstyleemphasis{n} .
It specifies the point at which to evaluate the Hessian \(H(x)\).

\index{r@\spxentry{r}}\ignorespaces 

\subparagraph{r}
\label{\detokenize{xsrst/cpp_sparse_hes:r}}\label{\detokenize{xsrst/cpp_sparse_hes:cpp-sparse-hes-r}}\label{\detokenize{xsrst/cpp_sparse_hes:index-6}}
This argument has prototype

\begin{DUlineblock}{0em}
\item[]         \sphinxcode{\sphinxupquote{const vec\_double\&}} \sphinxstyleemphasis{r}
\end{DUlineblock}

and its size is \sphinxstyleemphasis{m} .
It specifies the multiplier for each component of \(F(x)\);
i.e., \(r_i\) is the multiplier for \(F_i (x)\).

\index{pattern@\spxentry{pattern}}\ignorespaces 

\subparagraph{pattern}
\label{\detokenize{xsrst/cpp_sparse_hes:pattern}}\label{\detokenize{xsrst/cpp_sparse_hes:cpp-sparse-hes-pattern}}\label{\detokenize{xsrst/cpp_sparse_hes:index-7}}
This argument has prototype

\begin{DUlineblock}{0em}
\item[]         \sphinxcode{\sphinxupquote{const sparse\_rc\&}} \sphinxstyleemphasis{pattern}
\end{DUlineblock}

Its row size and column sizes are \sphinxstyleemphasis{n} ; i.e.,
\sphinxstyleemphasis{pattern} . \sphinxcode{\sphinxupquote{nr}} () == \sphinxstyleemphasis{n} and \sphinxstyleemphasis{pattern} . \sphinxcode{\sphinxupquote{nc}} () == \sphinxstyleemphasis{n} .
It is a sparsity pattern for the Hessian \(H(x)\).
This argument is not used (and need not satisfy any conditions),
when {\hyperref[\detokenize{xsrst/cpp_sparse_hes:cpp-sparse-hes-work}]{\sphinxcrossref{\DUrole{std,std-ref}{work}}}} is non\sphinxhyphen{}empty.

\index{work@\spxentry{work}}\ignorespaces 

\subparagraph{work}
\label{\detokenize{xsrst/cpp_sparse_hes:work}}\label{\detokenize{xsrst/cpp_sparse_hes:cpp-sparse-hes-work}}\label{\detokenize{xsrst/cpp_sparse_hes:index-8}}
This argument has prototype

\begin{DUlineblock}{0em}
\item[]         \sphinxcode{\sphinxupquote{sparse\_hes\_work\&}} \sphinxstyleemphasis{work}
\end{DUlineblock}

We refer to its initial value,
and its value after \sphinxstyleemphasis{work} . \sphinxcode{\sphinxupquote{clear}} () , as empty.
If it is empty, information is stored in \sphinxstyleemphasis{work} .
This can be used to reduce computation when
a future call is for the same object \sphinxstyleemphasis{f} ,
and the same subset of the Hessian.
If either of these values change, use \sphinxstyleemphasis{work} . \sphinxcode{\sphinxupquote{clear}} () to
empty this structure.

\index{n\_sweep@\spxentry{n\_sweep}}\ignorespaces 

\subparagraph{n\_sweep}
\label{\detokenize{xsrst/cpp_sparse_hes:n-sweep}}\label{\detokenize{xsrst/cpp_sparse_hes:cpp-sparse-hes-n-sweep}}\label{\detokenize{xsrst/cpp_sparse_hes:index-9}}
The return value \sphinxstyleemphasis{n\_sweep} has prototype

\begin{DUlineblock}{0em}
\item[]         \sphinxcode{\sphinxupquote{int}} \sphinxstyleemphasis{n\_sweep}
\end{DUlineblock}

It is the number of first order forward sweeps
used to compute the requested Hessian values.
Each first forward sweep is followed by a second order reverse sweep
so it is also the number of reverse sweeps.
This is proportional to the total computational work,
not counting the zero order forward sweep,
or combining multiple columns and rows into a single sweep.

\index{uses@\spxentry{uses}}\index{forward@\spxentry{forward}}\ignorespaces 

\subparagraph{Uses Forward}
\label{\detokenize{xsrst/cpp_sparse_hes:uses-forward}}\label{\detokenize{xsrst/cpp_sparse_hes:cpp-sparse-hes-uses-forward}}\label{\detokenize{xsrst/cpp_sparse_hes:index-10}}
After each call to {\hyperref[\detokenize{xsrst/cpp_fun_forward:id1}]{\sphinxcrossref{\DUrole{std,std-ref}{cpp\_fun\_forward}}}},
the object \sphinxstyleemphasis{f} contains the corresponding Taylor coefficients
for all the variables in the operation sequence..
After a call to \sphinxcode{\sphinxupquote{sparse\_hes}}
the zero order coefficients correspond to

\begin{DUlineblock}{0em}
\item[]         \sphinxstyleemphasis{f} . \sphinxcode{\sphinxupquote{forward(0}} , \sphinxstyleemphasis{x} )
\end{DUlineblock}

All the other forward mode coefficients are unspecified.


\subparagraph{sparse\_hes\_xam\_cpp}
\label{\detokenize{xsrst/sparse_hes_xam_cpp:sparse-hes-xam-cpp}}\label{\detokenize{xsrst/sparse_hes_xam_cpp::doc}}
\index{sparse\_hes\_xam\_cpp@\spxentry{sparse\_hes\_xam\_cpp}}\index{c++:@\spxentry{c++:}}\index{hessian@\spxentry{hessian}}\index{sparsity@\spxentry{sparsity}}\index{patterns:@\spxentry{patterns:}}\index{example@\spxentry{example}}\index{test@\spxentry{test}}\ignorespaces 

\subparagraph{3.2.4.6.1 C++: Hessian Sparsity Patterns: Example and Test}
\label{\detokenize{xsrst/sparse_hes_xam_cpp:c-hessian-sparsity-patterns-example-and-test}}\label{\detokenize{xsrst/sparse_hes_xam_cpp:id1}}
\begin{sphinxVerbatim}[commandchars=\\\{\}]
\PYG{c+cp}{\PYGZsh{}}\PYG{c+cp}{ include \PYGZlt{}cstdio\PYGZgt{}}
\PYG{c+cp}{\PYGZsh{}}\PYG{c+cp}{ include \PYGZlt{}cppad}\PYG{c+cp}{/}\PYG{c+cp}{py}\PYG{c+cp}{/}\PYG{c+cp}{cppad\PYGZus{}py.hpp\PYGZgt{}}

\PYG{k+kt}{bool} \PYG{n+nf}{sparse\PYGZus{}hes\PYGZus{}xam}\PYG{p}{(}\PYG{k+kt}{void}\PYG{p}{)} \PYG{p}{\PYGZob{}}
    \PYG{k}{using} \PYG{n}{cppad\PYGZus{}py}\PYG{o}{:}\PYG{o}{:}\PYG{n}{a\PYGZus{}double}\PYG{p}{;}
    \PYG{k}{using} \PYG{n}{cppad\PYGZus{}py}\PYG{o}{:}\PYG{o}{:}\PYG{n}{vec\PYGZus{}bool}\PYG{p}{;}
    \PYG{k}{using} \PYG{n}{cppad\PYGZus{}py}\PYG{o}{:}\PYG{o}{:}\PYG{n}{vec\PYGZus{}int}\PYG{p}{;}
    \PYG{k}{using} \PYG{n}{cppad\PYGZus{}py}\PYG{o}{:}\PYG{o}{:}\PYG{n}{vec\PYGZus{}double}\PYG{p}{;}
    \PYG{k}{using} \PYG{n}{cppad\PYGZus{}py}\PYG{o}{:}\PYG{o}{:}\PYG{n}{vec\PYGZus{}a\PYGZus{}double}\PYG{p}{;}
    \PYG{k}{using} \PYG{n}{cppad\PYGZus{}py}\PYG{o}{:}\PYG{o}{:}\PYG{n}{d\PYGZus{}fun}\PYG{p}{;}
    \PYG{k}{using} \PYG{n}{cppad\PYGZus{}py}\PYG{o}{:}\PYG{o}{:}\PYG{n}{sparse\PYGZus{}rc}\PYG{p}{;}
    \PYG{k}{using} \PYG{n}{cppad\PYGZus{}py}\PYG{o}{:}\PYG{o}{:}\PYG{n}{sparse\PYGZus{}rcv}\PYG{p}{;}
    \PYG{k}{using} \PYG{n}{cppad\PYGZus{}py}\PYG{o}{:}\PYG{o}{:}\PYG{n}{sparse\PYGZus{}hes\PYGZus{}work}\PYG{p}{;}
    \PYG{c+c1}{//}
    \PYG{c+c1}{// initialize return variable}
    \PYG{k+kt}{bool} \PYG{n}{ok} \PYG{o}{=} \PYG{n+nb}{true}\PYG{p}{;}
    \PYG{c+c1}{//\PYGZhy{}\PYGZhy{}\PYGZhy{}\PYGZhy{}\PYGZhy{}\PYGZhy{}\PYGZhy{}\PYGZhy{}\PYGZhy{}\PYGZhy{}\PYGZhy{}\PYGZhy{}\PYGZhy{}\PYGZhy{}\PYGZhy{}\PYGZhy{}\PYGZhy{}\PYGZhy{}\PYGZhy{}\PYGZhy{}\PYGZhy{}\PYGZhy{}\PYGZhy{}\PYGZhy{}\PYGZhy{}\PYGZhy{}\PYGZhy{}\PYGZhy{}\PYGZhy{}\PYGZhy{}\PYGZhy{}\PYGZhy{}\PYGZhy{}\PYGZhy{}\PYGZhy{}\PYGZhy{}\PYGZhy{}\PYGZhy{}\PYGZhy{}\PYGZhy{}\PYGZhy{}\PYGZhy{}\PYGZhy{}\PYGZhy{}\PYGZhy{}\PYGZhy{}\PYGZhy{}\PYGZhy{}\PYGZhy{}\PYGZhy{}\PYGZhy{}\PYGZhy{}\PYGZhy{}\PYGZhy{}\PYGZhy{}\PYGZhy{}\PYGZhy{}\PYGZhy{}\PYGZhy{}\PYGZhy{}\PYGZhy{}\PYGZhy{}\PYGZhy{}\PYGZhy{}\PYGZhy{}\PYGZhy{}\PYGZhy{}\PYGZhy{}\PYGZhy{}\PYGZhy{}\PYGZhy{}\PYGZhy{}}
    \PYG{c+c1}{// number of dependent and independent variables}
    \PYG{k+kt}{int} \PYG{n}{n} \PYG{o}{=} \PYG{l+m+mi}{3}\PYG{p}{;}
    \PYG{c+c1}{//}
    \PYG{c+c1}{// create the independent variables ax}
    \PYG{n}{vec\PYGZus{}double} \PYG{n}{x}\PYG{p}{(}\PYG{n}{n}\PYG{p}{)}\PYG{p}{;}
    \PYG{k}{for}\PYG{p}{(}\PYG{k+kt}{int} \PYG{n}{i} \PYG{o}{=} \PYG{l+m+mi}{0}\PYG{p}{;} \PYG{n}{i} \PYG{o}{\PYGZlt{}} \PYG{n}{n} \PYG{p}{;} \PYG{n}{i}\PYG{o}{+}\PYG{o}{+}\PYG{p}{)} \PYG{p}{\PYGZob{}}
        \PYG{n}{x}\PYG{p}{[}\PYG{n}{i}\PYG{p}{]} \PYG{o}{=} \PYG{n}{i} \PYG{o}{+} \PYG{l+m+mf}{2.0}\PYG{p}{;}
    \PYG{p}{\PYGZcb{}}
    \PYG{n}{vec\PYGZus{}a\PYGZus{}double} \PYG{n}{ax} \PYG{o}{=} \PYG{n}{cppad\PYGZus{}py}\PYG{o}{:}\PYG{o}{:}\PYG{n}{independent}\PYG{p}{(}\PYG{n}{x}\PYG{p}{)}\PYG{p}{;}
    \PYG{c+c1}{//}
    \PYG{c+c1}{// ay[i] = j * ax[j] * ax[i]}
    \PYG{c+c1}{// where i = mod(j + 1, n)}
    \PYG{n}{vec\PYGZus{}a\PYGZus{}double} \PYG{n}{ay}\PYG{p}{(}\PYG{n}{n}\PYG{p}{)}\PYG{p}{;}
    \PYG{k}{for}\PYG{p}{(}\PYG{k+kt}{int} \PYG{n}{j} \PYG{o}{=} \PYG{l+m+mi}{0}\PYG{p}{;} \PYG{n}{j} \PYG{o}{\PYGZlt{}} \PYG{n}{n} \PYG{p}{;} \PYG{n}{j}\PYG{o}{+}\PYG{o}{+}\PYG{p}{)} \PYG{p}{\PYGZob{}}
        \PYG{k+kt}{int} \PYG{n}{i} \PYG{o}{=} \PYG{n}{j}\PYG{o}{+}\PYG{l+m+mi}{1}\PYG{p}{;}
        \PYG{k}{if}\PYG{p}{(} \PYG{n}{i} \PYG{o}{\PYGZgt{}}\PYG{o}{=} \PYG{n}{n}  \PYG{p}{)} \PYG{p}{\PYGZob{}}
            \PYG{n}{i} \PYG{o}{=} \PYG{n}{i} \PYG{o}{\PYGZhy{}} \PYG{n}{n}\PYG{p}{;}
        \PYG{p}{\PYGZcb{}}
        \PYG{n}{a\PYGZus{}double} \PYG{n}{aj}\PYG{p}{(}\PYG{n}{j}\PYG{p}{)}\PYG{p}{;}
        \PYG{n}{a\PYGZus{}double} \PYG{n}{ax\PYGZus{}j} \PYG{o}{=} \PYG{n}{ax}\PYG{p}{[}\PYG{n}{j}\PYG{p}{]}\PYG{p}{;}
        \PYG{n}{a\PYGZus{}double} \PYG{n}{ax\PYGZus{}i} \PYG{o}{=} \PYG{n}{ax}\PYG{p}{[}\PYG{n}{i}\PYG{p}{]}\PYG{p}{;}
        \PYG{n}{ay}\PYG{p}{[}\PYG{n}{i}\PYG{p}{]} \PYG{o}{=} \PYG{n}{aj} \PYG{o}{*} \PYG{n}{ax\PYGZus{}j} \PYG{o}{*} \PYG{n}{ax\PYGZus{}i}\PYG{p}{;}
    \PYG{p}{\PYGZcb{}}
    \PYG{c+c1}{//}
    \PYG{c+c1}{// define af corresponding to f(x)}
    \PYG{n}{d\PYGZus{}fun} \PYG{n}{f}\PYG{p}{(}\PYG{n}{ax}\PYG{p}{,} \PYG{n}{ay}\PYG{p}{)}\PYG{p}{;}
    \PYG{c+c1}{//}
    \PYG{c+c1}{// Set select\PYGZus{}d (domain) to all true,}
    \PYG{c+c1}{// initial select\PYGZus{}r (range) to all false}
    \PYG{c+c1}{// initialize r to all zeros}
    \PYG{n}{vec\PYGZus{}bool} \PYG{n}{select\PYGZus{}d} \PYG{o}{=} \PYG{n}{vec\PYGZus{}bool}\PYG{p}{(}\PYG{n}{n}\PYG{p}{)}\PYG{p}{;}
    \PYG{n}{vec\PYGZus{}bool} \PYG{n}{select\PYGZus{}r} \PYG{o}{=} \PYG{n}{vec\PYGZus{}bool}\PYG{p}{(}\PYG{n}{n}\PYG{p}{)}\PYG{p}{;}
    \PYG{n}{vec\PYGZus{}double} \PYG{n}{r}\PYG{p}{(}\PYG{n}{n}\PYG{p}{)}\PYG{p}{;}
    \PYG{k}{for}\PYG{p}{(}\PYG{k+kt}{int} \PYG{n}{i} \PYG{o}{=} \PYG{l+m+mi}{0}\PYG{p}{;} \PYG{n}{i} \PYG{o}{\PYGZlt{}} \PYG{n}{n}\PYG{p}{;} \PYG{n}{i}\PYG{o}{+}\PYG{o}{+}\PYG{p}{)} \PYG{p}{\PYGZob{}}
        \PYG{n}{select\PYGZus{}d}\PYG{p}{[}\PYG{n}{i}\PYG{p}{]} \PYG{o}{=} \PYG{n+nb}{true}\PYG{p}{;}
        \PYG{n}{select\PYGZus{}r}\PYG{p}{[}\PYG{n}{i}\PYG{p}{]} \PYG{o}{=} \PYG{n+nb}{false}\PYG{p}{;}
        \PYG{n}{r}\PYG{p}{[}\PYG{n}{i}\PYG{p}{]} \PYG{o}{=} \PYG{l+m+mf}{0.0}\PYG{p}{;}
    \PYG{p}{\PYGZcb{}}
    \PYG{c+c1}{//}
    \PYG{c+c1}{// only select component n\PYGZhy{}1 of the range function}
    \PYG{c+c1}{// f\PYGZus{}0 (x) = (n\PYGZhy{}1) * x\PYGZus{}\PYGZob{}n\PYGZhy{}1\PYGZcb{} * x\PYGZus{}0}
    \PYG{n}{select\PYGZus{}r}\PYG{p}{[}\PYG{l+m+mi}{0}\PYG{p}{]} \PYG{o}{=} \PYG{n+nb}{true}\PYG{p}{;}
    \PYG{n}{r}\PYG{p}{[}\PYG{l+m+mi}{0}\PYG{p}{]} \PYG{o}{=} \PYG{l+m+mf}{1.0}\PYG{p}{;}
    \PYG{c+c1}{//}
    \PYG{c+c1}{// sparisty pattern for Hessian}
    \PYG{n}{sparse\PYGZus{}rc} \PYG{n}{pattern} \PYG{o}{=} \PYG{n}{sparse\PYGZus{}rc}\PYG{p}{(}\PYG{p}{)}\PYG{p}{;}
    \PYG{n}{f}\PYG{p}{.}\PYG{n}{for\PYGZus{}hes\PYGZus{}sparsity}\PYG{p}{(}\PYG{n}{select\PYGZus{}d}\PYG{p}{,} \PYG{n}{select\PYGZus{}r}\PYG{p}{,} \PYG{n}{pattern}\PYG{p}{)}\PYG{p}{;}
    \PYG{c+c1}{//}
    \PYG{c+c1}{// compute all possibly non\PYGZhy{}zero entries in Hessian}
    \PYG{c+c1}{// (should only compute lower triangle becuase matrix is symmetric)}
    \PYG{n}{sparse\PYGZus{}rcv} \PYG{n}{subset} \PYG{o}{=} \PYG{n}{sparse\PYGZus{}rcv}\PYG{p}{(}\PYG{n}{pattern}\PYG{p}{)}\PYG{p}{;}
    \PYG{c+c1}{//}
    \PYG{c+c1}{// work space used to save time for multiple calls}
    \PYG{n}{sparse\PYGZus{}hes\PYGZus{}work} \PYG{n}{work} \PYG{o}{=} \PYG{n}{cppad\PYGZus{}py}\PYG{o}{:}\PYG{o}{:}\PYG{n}{sparse\PYGZus{}hes\PYGZus{}work}\PYG{p}{(}\PYG{p}{)}\PYG{p}{;}
    \PYG{c+c1}{//}
    \PYG{n}{f}\PYG{p}{.}\PYG{n}{sparse\PYGZus{}hes}\PYG{p}{(}\PYG{n}{subset}\PYG{p}{,} \PYG{n}{x}\PYG{p}{,} \PYG{n}{r}\PYG{p}{,} \PYG{n}{pattern}\PYG{p}{,} \PYG{n}{work}\PYG{p}{)}\PYG{p}{;}
    \PYG{c+c1}{//}
    \PYG{c+c1}{// check that result is sparsity pattern for Hessian of f\PYGZus{}0 (x)}
    \PYG{n}{ok} \PYG{o}{=} \PYG{n}{ok} \PYG{o}{\PYGZam{}}\PYG{o}{\PYGZam{}} \PYG{n}{subset}\PYG{p}{.}\PYG{n}{nnz}\PYG{p}{(}\PYG{p}{)} \PYG{o}{=}\PYG{o}{=} \PYG{l+m+mi}{2} \PYG{p}{;}
    \PYG{n}{vec\PYGZus{}int} \PYG{n}{row} \PYG{o}{=} \PYG{n}{subset}\PYG{p}{.}\PYG{n}{row}\PYG{p}{(}\PYG{p}{)}\PYG{p}{;}
    \PYG{n}{vec\PYGZus{}int} \PYG{n}{col} \PYG{o}{=} \PYG{n}{subset}\PYG{p}{.}\PYG{n}{col}\PYG{p}{(}\PYG{p}{)}\PYG{p}{;}
    \PYG{n}{vec\PYGZus{}double} \PYG{n}{val} \PYG{o}{=} \PYG{n}{subset}\PYG{p}{.}\PYG{n}{val}\PYG{p}{(}\PYG{p}{)}\PYG{p}{;}
    \PYG{k}{for}\PYG{p}{(}\PYG{k+kt}{int} \PYG{n}{k} \PYG{o}{=} \PYG{l+m+mi}{0}\PYG{p}{;} \PYG{n}{k} \PYG{o}{\PYGZlt{}} \PYG{l+m+mi}{2}\PYG{p}{;} \PYG{n}{k}\PYG{o}{+}\PYG{o}{+}\PYG{p}{)} \PYG{p}{\PYGZob{}}
        \PYG{k+kt}{int} \PYG{n}{i} \PYG{o}{=} \PYG{n}{row}\PYG{p}{[}\PYG{n}{k}\PYG{p}{]}\PYG{p}{;}
        \PYG{k+kt}{int} \PYG{n}{j} \PYG{o}{=} \PYG{n}{col}\PYG{p}{[}\PYG{n}{k}\PYG{p}{]}\PYG{p}{;}
        \PYG{k+kt}{double} \PYG{n}{v} \PYG{o}{=} \PYG{n}{val}\PYG{p}{[}\PYG{n}{k}\PYG{p}{]}\PYG{p}{;}
        \PYG{k}{if}\PYG{p}{(} \PYG{n}{i} \PYG{o}{\PYGZlt{}}\PYG{o}{=} \PYG{n}{j}  \PYG{p}{)} \PYG{p}{\PYGZob{}}
            \PYG{n}{ok} \PYG{o}{=} \PYG{n}{ok} \PYG{o}{\PYGZam{}}\PYG{o}{\PYGZam{}} \PYG{n}{i} \PYG{o}{=}\PYG{o}{=} \PYG{l+m+mi}{0}\PYG{p}{;}
            \PYG{n}{ok} \PYG{o}{=} \PYG{n}{ok} \PYG{o}{\PYGZam{}}\PYG{o}{\PYGZam{}} \PYG{n}{j} \PYG{o}{=}\PYG{o}{=} \PYG{n}{n}\PYG{o}{\PYGZhy{}}\PYG{l+m+mi}{1}\PYG{p}{;}
        \PYG{p}{\PYGZcb{}}
        \PYG{k}{if}\PYG{p}{(} \PYG{n}{i} \PYG{o}{\PYGZgt{}}\PYG{o}{=} \PYG{n}{j}  \PYG{p}{)} \PYG{p}{\PYGZob{}}
            \PYG{n}{ok} \PYG{o}{=} \PYG{n}{ok} \PYG{o}{\PYGZam{}}\PYG{o}{\PYGZam{}} \PYG{n}{i} \PYG{o}{=}\PYG{o}{=} \PYG{n}{n}\PYG{o}{\PYGZhy{}}\PYG{l+m+mi}{1}\PYG{p}{;}
            \PYG{n}{ok} \PYG{o}{=} \PYG{n}{ok} \PYG{o}{\PYGZam{}}\PYG{o}{\PYGZam{}} \PYG{n}{j} \PYG{o}{=}\PYG{o}{=} \PYG{l+m+mi}{0}\PYG{p}{;}
        \PYG{p}{\PYGZcb{}}
        \PYG{n}{ok} \PYG{o}{=} \PYG{n}{ok} \PYG{o}{\PYGZam{}}\PYG{o}{\PYGZam{}} \PYG{n}{v} \PYG{o}{=}\PYG{o}{=} \PYG{n}{n}\PYG{o}{\PYGZhy{}}\PYG{l+m+mi}{1}\PYG{p}{;}
    \PYG{p}{\PYGZcb{}}
    \PYG{c+c1}{//}
    \PYG{k}{return}\PYG{p}{(} \PYG{n}{ok} \PYG{p}{)}\PYG{p}{;}
\PYG{p}{\PYGZcb{}}
\end{sphinxVerbatim}


\bigskip\hrule\bigskip


xsrst input file: \sphinxcode{\sphinxupquote{example/cplusplus/sparse\_hes\_xam.cpp}}

\index{example@\spxentry{example}}\ignorespaces 

\subparagraph{Example}
\label{\detokenize{xsrst/cpp_sparse_hes:example}}\label{\detokenize{xsrst/cpp_sparse_hes:cpp-sparse-hes-example}}\label{\detokenize{xsrst/cpp_sparse_hes:index-11}}
{\hyperref[\detokenize{xsrst/sparse_hes_xam_cpp:id1}]{\sphinxcrossref{\DUrole{std,std-ref}{sparse\_hes\_xam\_cpp}}}}


\bigskip\hrule\bigskip


xsrst input file: \sphinxcode{\sphinxupquote{lib/cplusplus/sparse.cpp}}
\begin{enumerate}
\sphinxsetlistlabels{\arabic}{enumi}{enumii}{}{.}%
\item {} 
cpp\_sparse\_rc: {\hyperref[\detokenize{xsrst/cpp_sparse_rc:id1}]{\sphinxcrossref{\DUrole{std,std-ref}{3.2.4.1 Sparsity Patterns}}}}

\item {} 
cpp\_sparse\_rcv: {\hyperref[\detokenize{xsrst/cpp_sparse_rcv:id1}]{\sphinxcrossref{\DUrole{std,std-ref}{3.2.4.2 Sparse Matrices}}}}

\item {} 
cpp\_jac\_sparsity: {\hyperref[\detokenize{xsrst/cpp_jac_sparsity:id1}]{\sphinxcrossref{\DUrole{std,std-ref}{3.2.4.3 Jacobian Sparsity Patterns}}}}

\item {} 
cpp\_sparsity: {\hyperref[\detokenize{xsrst/cpp_sparsity:id1}]{\sphinxcrossref{\DUrole{std,std-ref}{3.2.4.4 Hessian Sparsity Patterns}}}}

\item {} 
cpp\_sparse\_jac: {\hyperref[\detokenize{xsrst/cpp_sparse_jac:id1}]{\sphinxcrossref{\DUrole{std,std-ref}{3.2.4.5 Computing Sparse Jacobians}}}}

\item {} 
cpp\_sparse\_hes: {\hyperref[\detokenize{xsrst/cpp_sparse_hes:id1}]{\sphinxcrossref{\DUrole{std,std-ref}{3.2.4.6 Computing Sparse Hessians}}}}

\end{enumerate}


\bigskip\hrule\bigskip


xsrst input file: \sphinxcode{\sphinxupquote{lib/cplusplus/cplusplus.xsrst}}


\subparagraph{cpp\_utility}
\label{\detokenize{xsrst/cpp_utility:cpp-utility}}\label{\detokenize{xsrst/cpp_utility::doc}}
\index{cpp\_utility@\spxentry{cpp\_utility}}\index{c++@\spxentry{c++}}\index{utilities@\spxentry{utilities}}\ignorespaces 

\subparagraph{3.2.5 C++ Utilities}
\label{\detokenize{xsrst/cpp_utility:c-utilities}}\label{\detokenize{xsrst/cpp_utility:id1}}\begin{itemize}
\item {} 
{\hyperref[\detokenize{xsrst/cpp_utility:cpp-utility-children}]{\sphinxcrossref{\DUrole{std,std-ref}{Children}}}}

\end{itemize}

\index{children@\spxentry{children}}\ignorespaces 

\subparagraph{Children}
\label{\detokenize{xsrst/cpp_utility:children}}\label{\detokenize{xsrst/cpp_utility:cpp-utility-children}}\label{\detokenize{xsrst/cpp_utility:index-1}}

\subparagraph{vec2cppad\_double}
\label{\detokenize{xsrst/vec2cppad_double:vec2cppad-double}}\label{\detokenize{xsrst/vec2cppad_double::doc}}
\index{vec2cppad\_double@\spxentry{vec2cppad\_double}}\index{convert@\spxentry{convert}}\index{an@\spxentry{an}}\index{a\_double@\spxentry{a\_double}}\index{vector@\spxentry{vector}}\index{to@\spxentry{to}}\index{cppad::ad\textless{}double\textgreater{}@\spxentry{cppad::ad\textless{}double\textgreater{}}}\index{vector@\spxentry{vector}}\ignorespaces 

\subparagraph{3.2.5.1 Convert an a\_double Vector to a CppAD::AD\textless{}double\textgreater{} Vector}
\label{\detokenize{xsrst/vec2cppad_double:convert-an-a-double-vector-to-a-cppad-ad-double-vector}}\label{\detokenize{xsrst/vec2cppad_double:id1}}\begin{itemize}
\item {} 
{\hyperref[\detokenize{xsrst/vec2cppad_double:vec2cppad-double-syntax}]{\sphinxcrossref{\DUrole{std,std-ref}{Syntax}}}}

\item {} 
{\hyperref[\detokenize{xsrst/vec2cppad_double:vec2cppad-double-prototype}]{\sphinxcrossref{\DUrole{std,std-ref}{Prototype}}}}

\end{itemize}

\index{syntax@\spxentry{syntax}}\ignorespaces 

\subparagraph{Syntax}
\label{\detokenize{xsrst/vec2cppad_double:syntax}}\label{\detokenize{xsrst/vec2cppad_double:vec2cppad-double-syntax}}\label{\detokenize{xsrst/vec2cppad_double:index-1}}
\begin{DUlineblock}{0em}
\item[] \sphinxstyleemphasis{v\_out} =  \sphinxcode{\sphinxupquote{cppad\_py::vec2cppad\_double}} ( \sphinxcode{\sphinxupquote{v\_in}} )
\end{DUlineblock}

\index{prototype@\spxentry{prototype}}\ignorespaces 

\subparagraph{Prototype}
\label{\detokenize{xsrst/vec2cppad_double:prototype}}\label{\detokenize{xsrst/vec2cppad_double:vec2cppad-double-prototype}}\label{\detokenize{xsrst/vec2cppad_double:index-2}}
\begin{sphinxVerbatim}[commandchars=\\\{\}]
\PYG{n}{std}\PYG{o}{:}\PYG{o}{:}\PYG{n}{vector}\PYG{o}{\PYGZlt{}} \PYG{n}{CppAD}\PYG{o}{:}\PYG{o}{:}\PYG{n}{AD}\PYG{o}{\PYGZlt{}}\PYG{k+kt}{double}\PYG{o}{\PYGZgt{}} \PYG{o}{\PYGZgt{}}
\PYG{n}{vec2cppad\PYGZus{}double}\PYG{p}{(}\PYG{k}{const} \PYG{n}{std}\PYG{o}{:}\PYG{o}{:}\PYG{n}{vector}\PYG{o}{\PYGZlt{}}\PYG{n}{a\PYGZus{}double}\PYG{o}{\PYGZgt{}}\PYG{o}{\PYGZam{}} \PYG{n}{v\PYGZus{}in} \PYG{p}{)}
\end{sphinxVerbatim}


\bigskip\hrule\bigskip


xsrst input file: \sphinxcode{\sphinxupquote{lib/cplusplus/cppad\_vec.cpp}}


\subparagraph{vec2a\_double}
\label{\detokenize{xsrst/vec2a_double:vec2a-double}}\label{\detokenize{xsrst/vec2a_double::doc}}
\index{vec2a\_double@\spxentry{vec2a\_double}}\index{convert@\spxentry{convert}}\index{cppad::ad\textless{}double\textgreater{}@\spxentry{cppad::ad\textless{}double\textgreater{}}}\index{vector@\spxentry{vector}}\index{to@\spxentry{to}}\index{an@\spxentry{an}}\index{a\_double@\spxentry{a\_double}}\index{vector@\spxentry{vector}}\ignorespaces 

\subparagraph{3.2.5.2 Convert a CppAD::AD\textless{}double\textgreater{} Vector to an a\_double Vector}
\label{\detokenize{xsrst/vec2a_double:convert-a-cppad-ad-double-vector-to-an-a-double-vector}}\label{\detokenize{xsrst/vec2a_double:id1}}\begin{itemize}
\item {} 
{\hyperref[\detokenize{xsrst/vec2a_double:vec2a-double-syntax}]{\sphinxcrossref{\DUrole{std,std-ref}{Syntax}}}}

\item {} 
{\hyperref[\detokenize{xsrst/vec2a_double:vec2a-double-prototype}]{\sphinxcrossref{\DUrole{std,std-ref}{Prototype}}}}

\end{itemize}

\index{syntax@\spxentry{syntax}}\ignorespaces 

\subparagraph{Syntax}
\label{\detokenize{xsrst/vec2a_double:syntax}}\label{\detokenize{xsrst/vec2a_double:vec2a-double-syntax}}\label{\detokenize{xsrst/vec2a_double:index-1}}
\begin{DUlineblock}{0em}
\item[] \sphinxstyleemphasis{v\_out} =  \sphinxcode{\sphinxupquote{cppad\_py::vec2a\_double}} ( \sphinxcode{\sphinxupquote{v\_in}} )
\end{DUlineblock}

\index{prototype@\spxentry{prototype}}\ignorespaces 

\subparagraph{Prototype}
\label{\detokenize{xsrst/vec2a_double:prototype}}\label{\detokenize{xsrst/vec2a_double:vec2a-double-prototype}}\label{\detokenize{xsrst/vec2a_double:index-2}}
\begin{sphinxVerbatim}[commandchars=\\\{\}]
\PYG{n}{std}\PYG{o}{:}\PYG{o}{:}\PYG{n}{vector}\PYG{o}{\PYGZlt{}}\PYG{n}{a\PYGZus{}double}\PYG{o}{\PYGZgt{}}
\PYG{n}{vec2a\PYGZus{}double}\PYG{p}{(}\PYG{k}{const} \PYG{n}{std}\PYG{o}{:}\PYG{o}{:}\PYG{n}{vector}\PYG{o}{\PYGZlt{}} \PYG{n}{CppAD}\PYG{o}{:}\PYG{o}{:}\PYG{n}{AD}\PYG{o}{\PYGZlt{}}\PYG{k+kt}{double}\PYG{o}{\PYGZgt{}} \PYG{o}{\PYGZgt{}}\PYG{o}{\PYGZam{}} \PYG{n}{v\PYGZus{}in} \PYG{p}{)}
\end{sphinxVerbatim}


\bigskip\hrule\bigskip


xsrst input file: \sphinxcode{\sphinxupquote{lib/cplusplus/cppad\_vec.cpp}}


\subparagraph{error\_message}
\label{\detokenize{xsrst/error_message:error-message}}\label{\detokenize{xsrst/error_message::doc}}
\index{error\_message@\spxentry{error\_message}}\index{exception@\spxentry{exception}}\index{handling@\spxentry{handling}}\ignorespaces 

\subparagraph{3.2.5.3 Exception Handling}
\label{\detokenize{xsrst/error_message:exception-handling}}\label{\detokenize{xsrst/error_message:id1}}\begin{itemize}
\item {} 
{\hyperref[\detokenize{xsrst/error_message:error-message-syntax}]{\sphinxcrossref{\DUrole{std,std-ref}{Syntax}}}}

\item {} \begin{description}
\item[{{\hyperref[\detokenize{xsrst/error_message:error-message-in-try-block}]{\sphinxcrossref{\DUrole{std,std-ref}{In Try Block}}}}}] \leavevmode\begin{itemize}
\item {} 
{\hyperref[\detokenize{xsrst/error_message:error-message-in-try-block-exception}]{\sphinxcrossref{\DUrole{std,std-ref}{exception}}}}

\end{itemize}

\end{description}

\item {} 
{\hyperref[\detokenize{xsrst/error_message:error-message-in-catch-block}]{\sphinxcrossref{\DUrole{std,std-ref}{In Catch Block}}}}

\item {} 
{\hyperref[\detokenize{xsrst/error_message:error-message-cppad-errors}]{\sphinxcrossref{\DUrole{std,std-ref}{Cppad Errors}}}}

\item {} 
{\hyperref[\detokenize{xsrst/error_message:error-message-not-thread-safe}]{\sphinxcrossref{\DUrole{std,std-ref}{Not Thread Safe}}}}

\item {} 
{\hyperref[\detokenize{xsrst/error_message:error-message-example}]{\sphinxcrossref{\DUrole{std,std-ref}{Example}}}}

\end{itemize}

\index{syntax@\spxentry{syntax}}\ignorespaces 

\subparagraph{Syntax}
\label{\detokenize{xsrst/error_message:syntax}}\label{\detokenize{xsrst/error_message:error-message-syntax}}\label{\detokenize{xsrst/error_message:index-1}}
\sphinxstyleemphasis{stored\_message} =  \sphinxcode{\sphinxupquote{error\_message}} ( \sphinxstyleemphasis{input\_message} )

\index{in@\spxentry{in}}\index{try@\spxentry{try}}\index{block@\spxentry{block}}\ignorespaces 

\subparagraph{In Try Block}
\label{\detokenize{xsrst/error_message:in-try-block}}\label{\detokenize{xsrst/error_message:error-message-in-try-block}}\label{\detokenize{xsrst/error_message:index-2}}
If \sphinxstyleemphasis{input\_message} is \sphinxstylestrong{not} the empty string,
it is stored in \sphinxcode{\sphinxupquote{error\_message}} and an exception is thrown.
This is intended to be done inside a \sphinxcode{\sphinxupquote{try}} block.

\index{exception@\spxentry{exception}}\ignorespaces 

\subparagraph{exception}
\label{\detokenize{xsrst/error_message:exception}}\label{\detokenize{xsrst/error_message:error-message-in-try-block-exception}}\label{\detokenize{xsrst/error_message:index-3}}
The type of the exception is \sphinxcode{\sphinxupquote{std::runtime\_error}}
which is derived from \sphinxcode{\sphinxupquote{std::exception}} .
If the standard exception \sphinxcode{\sphinxupquote{what()}} is called,
the return value will be the value of \sphinxstyleemphasis{input\_message}
when the exception was thrown.

\index{in@\spxentry{in}}\index{catch@\spxentry{catch}}\index{block@\spxentry{block}}\ignorespaces 

\subparagraph{In Catch Block}
\label{\detokenize{xsrst/error_message:in-catch-block}}\label{\detokenize{xsrst/error_message:error-message-in-catch-block}}\label{\detokenize{xsrst/error_message:index-4}}
If \sphinxstyleemphasis{input\_message} is the empty string,
the most recently stored message in \sphinxcode{\sphinxupquote{error\_message}} is returned.
In addition, the message is deleted from \sphinxcode{\sphinxupquote{error\_message}} .
If there are no more messages stored in \sphinxcode{\sphinxupquote{error\_message}} ,
the empty string is returned.
This is intended to be done inside a \sphinxcode{\sphinxupquote{catch}} block.

\index{cppad@\spxentry{cppad}}\index{errors@\spxentry{errors}}\ignorespaces 

\subparagraph{Cppad Errors}
\label{\detokenize{xsrst/error_message:cppad-errors}}\label{\detokenize{xsrst/error_message:error-message-cppad-errors}}\label{\detokenize{xsrst/error_message:index-5}}
Calls to the Cppad error handler get mapped to a call
to \sphinxcode{\sphinxupquote{error\_message}} .

\index{not@\spxentry{not}}\index{thread@\spxentry{thread}}\index{safe@\spxentry{safe}}\ignorespaces 

\subparagraph{Not Thread Safe}
\label{\detokenize{xsrst/error_message:not-thread-safe}}\label{\detokenize{xsrst/error_message:error-message-not-thread-safe}}\label{\detokenize{xsrst/error_message:index-6}}
The message storage is done using static information in
\sphinxcode{\sphinxupquote{error\_message}} and hence is not thread safe.


\subparagraph{error\_message\_xam\_cpp}
\label{\detokenize{xsrst/error_message_xam_cpp:error-message-xam-cpp}}\label{\detokenize{xsrst/error_message_xam_cpp::doc}}
\index{error\_message\_xam\_cpp@\spxentry{error\_message\_xam\_cpp}}\index{c++:@\spxentry{c++:}}\index{cppad@\spxentry{cppad}}\index{py@\spxentry{py}}\index{exception@\spxentry{exception}}\index{handling:@\spxentry{handling:}}\index{example@\spxentry{example}}\index{test@\spxentry{test}}\ignorespaces 

\subparagraph{3.2.5.3.1 C++: Cppad Py Exception Handling: Example and Test}
\label{\detokenize{xsrst/error_message_xam_cpp:c-cppad-py-exception-handling-example-and-test}}\label{\detokenize{xsrst/error_message_xam_cpp:id1}}
\begin{sphinxVerbatim}[commandchars=\\\{\}]
\PYG{c+cp}{\PYGZsh{}}\PYG{c+cp}{ include \PYGZlt{}cstdio\PYGZgt{}}
\PYG{c+cp}{\PYGZsh{}}\PYG{c+cp}{ include \PYGZlt{}string\PYGZgt{}}
\PYG{c+cp}{\PYGZsh{}}\PYG{c+cp}{ include \PYGZlt{}cppad}\PYG{c+cp}{/}\PYG{c+cp}{py}\PYG{c+cp}{/}\PYG{c+cp}{cppad\PYGZus{}py.hpp\PYGZgt{}}

\PYG{k+kt}{bool} \PYG{n+nf}{error\PYGZus{}message\PYGZus{}xam}\PYG{p}{(}\PYG{k+kt}{void}\PYG{p}{)} \PYG{p}{\PYGZob{}}
    \PYG{k}{using} \PYG{n}{std}\PYG{o}{:}\PYG{o}{:}\PYG{n}{string}\PYG{p}{;}
    \PYG{c+c1}{//}
    \PYG{c+c1}{// initialize return variable}
    \PYG{k+kt}{bool} \PYG{n}{ok} \PYG{o}{=} \PYG{n+nb}{true}\PYG{p}{;}
    \PYG{c+c1}{//\PYGZhy{}\PYGZhy{}\PYGZhy{}\PYGZhy{}\PYGZhy{}\PYGZhy{}\PYGZhy{}\PYGZhy{}\PYGZhy{}\PYGZhy{}\PYGZhy{}\PYGZhy{}\PYGZhy{}\PYGZhy{}\PYGZhy{}\PYGZhy{}\PYGZhy{}\PYGZhy{}\PYGZhy{}\PYGZhy{}\PYGZhy{}\PYGZhy{}\PYGZhy{}\PYGZhy{}\PYGZhy{}\PYGZhy{}\PYGZhy{}\PYGZhy{}\PYGZhy{}\PYGZhy{}\PYGZhy{}\PYGZhy{}\PYGZhy{}\PYGZhy{}\PYGZhy{}\PYGZhy{}\PYGZhy{}\PYGZhy{}\PYGZhy{}\PYGZhy{}\PYGZhy{}\PYGZhy{}\PYGZhy{}\PYGZhy{}\PYGZhy{}\PYGZhy{}\PYGZhy{}\PYGZhy{}\PYGZhy{}\PYGZhy{}\PYGZhy{}\PYGZhy{}\PYGZhy{}\PYGZhy{}\PYGZhy{}\PYGZhy{}\PYGZhy{}\PYGZhy{}\PYGZhy{}\PYGZhy{}\PYGZhy{}\PYGZhy{}\PYGZhy{}\PYGZhy{}\PYGZhy{}\PYGZhy{}\PYGZhy{}\PYGZhy{}\PYGZhy{}\PYGZhy{}\PYGZhy{}\PYGZhy{}}
    \PYG{n}{ok} \PYG{o}{=} \PYG{n+nb}{false}\PYG{p}{;}
    \PYG{k}{try} \PYG{p}{\PYGZob{}}
        \PYG{n}{cppad\PYGZus{}py}\PYG{o}{:}\PYG{o}{:}\PYG{n}{error\PYGZus{}message}\PYG{p}{(}\PYG{l+s}{\PYGZdq{}}\PYG{l+s}{test message}\PYG{l+s}{\PYGZdq{}}\PYG{p}{)}\PYG{p}{;}
    \PYG{p}{\PYGZcb{}} \PYG{k}{catch} \PYG{p}{(}\PYG{p}{.}\PYG{p}{.}\PYG{p}{.}\PYG{p}{)} \PYG{p}{\PYGZob{}}
        \PYG{n}{string} \PYG{n}{stored\PYGZus{}message} \PYG{o}{=} \PYG{n}{cppad\PYGZus{}py}\PYG{o}{:}\PYG{o}{:}\PYG{n}{error\PYGZus{}message}\PYG{p}{(}\PYG{l+s}{\PYGZdq{}}\PYG{l+s}{\PYGZdq{}}\PYG{p}{)}\PYG{p}{;}
        \PYG{n}{ok} \PYG{o}{=} \PYG{n}{stored\PYGZus{}message} \PYG{o}{=}\PYG{o}{=} \PYG{l+s}{\PYGZdq{}}\PYG{l+s}{test message}\PYG{l+s}{\PYGZdq{}}\PYG{p}{;}
    \PYG{p}{\PYGZcb{}}
    \PYG{k}{return}\PYG{p}{(} \PYG{n}{ok}  \PYG{p}{)}\PYG{p}{;}
\PYG{p}{\PYGZcb{}}
\end{sphinxVerbatim}


\bigskip\hrule\bigskip


xsrst input file: \sphinxcode{\sphinxupquote{example/cplusplus/error\_message\_xam.cpp}}


\subparagraph{error\_message\_xam\_py}
\label{\detokenize{xsrst/error_message_xam_py:error-message-xam-py}}\label{\detokenize{xsrst/error_message_xam_py::doc}}
\index{error\_message\_xam\_py@\spxentry{error\_message\_xam\_py}}\index{python:@\spxentry{python:}}\index{cppad@\spxentry{cppad}}\index{py@\spxentry{py}}\index{exception@\spxentry{exception}}\index{handling:@\spxentry{handling:}}\index{example@\spxentry{example}}\index{test@\spxentry{test}}\ignorespaces 

\subparagraph{3.2.5.3.2 Python: Cppad Py Exception Handling: Example and Test}
\label{\detokenize{xsrst/error_message_xam_py:python-cppad-py-exception-handling-example-and-test}}\label{\detokenize{xsrst/error_message_xam_py:id1}}
\begin{sphinxVerbatim}[commandchars=\\\{\}]
\PYG{k}{def} \PYG{n+nf}{error\PYGZus{}message\PYGZus{}xam}\PYG{p}{(}\PYG{p}{)} \PYG{p}{:}
    \PYG{c+c1}{\PYGZsh{}}
    \PYG{k+kn}{import} \PYG{n+nn}{numpy}
    \PYG{k+kn}{import} \PYG{n+nn}{cppad\PYGZus{}py}
    \PYG{c+c1}{\PYGZsh{}}
    \PYG{c+c1}{\PYGZsh{} initialize return variable}
    \PYG{n}{ok} \PYG{o}{=} \PYG{k+kc}{True}
    \PYG{c+c1}{\PYGZsh{} \PYGZhy{}\PYGZhy{}\PYGZhy{}\PYGZhy{}\PYGZhy{}\PYGZhy{}\PYGZhy{}\PYGZhy{}\PYGZhy{}\PYGZhy{}\PYGZhy{}\PYGZhy{}\PYGZhy{}\PYGZhy{}\PYGZhy{}\PYGZhy{}\PYGZhy{}\PYGZhy{}\PYGZhy{}\PYGZhy{}\PYGZhy{}\PYGZhy{}\PYGZhy{}\PYGZhy{}\PYGZhy{}\PYGZhy{}\PYGZhy{}\PYGZhy{}\PYGZhy{}\PYGZhy{}\PYGZhy{}\PYGZhy{}\PYGZhy{}\PYGZhy{}\PYGZhy{}\PYGZhy{}\PYGZhy{}\PYGZhy{}\PYGZhy{}\PYGZhy{}\PYGZhy{}\PYGZhy{}\PYGZhy{}\PYGZhy{}\PYGZhy{}\PYGZhy{}\PYGZhy{}\PYGZhy{}\PYGZhy{}\PYGZhy{}\PYGZhy{}\PYGZhy{}\PYGZhy{}\PYGZhy{}\PYGZhy{}\PYGZhy{}\PYGZhy{}\PYGZhy{}\PYGZhy{}\PYGZhy{}\PYGZhy{}\PYGZhy{}\PYGZhy{}\PYGZhy{}\PYGZhy{}\PYGZhy{}\PYGZhy{}\PYGZhy{}\PYGZhy{}}
    \PYG{n}{ok} \PYG{o}{=} \PYG{k+kc}{False}
    \PYG{k}{try} \PYG{p}{:}
        \PYG{n}{cppad\PYGZus{}py}\PYG{o}{.}\PYG{n}{error\PYGZus{}message}\PYG{p}{(}\PYG{l+s+s2}{\PYGZdq{}}\PYG{l+s+s2}{test message}\PYG{l+s+s2}{\PYGZdq{}}\PYG{p}{)}
    \PYG{k}{except} \PYG{p}{:} \PYG{c+c1}{\PYGZsh{} catch}
        \PYG{n}{stored\PYGZus{}message} \PYG{o}{=} \PYG{n}{cppad\PYGZus{}py}\PYG{o}{.}\PYG{n}{error\PYGZus{}message}\PYG{p}{(}\PYG{l+s+s2}{\PYGZdq{}}\PYG{l+s+s2}{\PYGZdq{}}\PYG{p}{)}
        \PYG{n}{ok} \PYG{o}{=} \PYG{p}{(}\PYG{n}{stored\PYGZus{}message} \PYG{o}{==} \PYG{l+s+s2}{\PYGZdq{}}\PYG{l+s+s2}{test message}\PYG{l+s+s2}{\PYGZdq{}}\PYG{p}{)}
    \PYG{c+c1}{\PYGZsh{}}
    \PYG{k}{return}\PYG{p}{(} \PYG{n}{ok}  \PYG{p}{)}
\PYG{c+c1}{\PYGZsh{}}
\end{sphinxVerbatim}


\bigskip\hrule\bigskip


xsrst input file: \sphinxcode{\sphinxupquote{example/python/core/error\_message\_xam.py}}

\index{example@\spxentry{example}}\ignorespaces 

\subparagraph{Example}
\label{\detokenize{xsrst/error_message:example}}\label{\detokenize{xsrst/error_message:error-message-example}}\label{\detokenize{xsrst/error_message:index-7}}
{\hyperref[\detokenize{xsrst/error_message_xam_cpp:id1}]{\sphinxcrossref{\DUrole{std,std-ref}{c++}}}},
{\hyperref[\detokenize{xsrst/error_message_xam_py:id1}]{\sphinxcrossref{\DUrole{std,std-ref}{python}}}}.


\bigskip\hrule\bigskip


xsrst input file: \sphinxcode{\sphinxupquote{lib/cplusplus/error.cpp}}
\begin{enumerate}
\sphinxsetlistlabels{\arabic}{enumi}{enumii}{}{.}%
\item {} 
vec2cppad\_double: {\hyperref[\detokenize{xsrst/vec2cppad_double:id1}]{\sphinxcrossref{\DUrole{std,std-ref}{3.2.5.1 Convert an a\_double Vector to a CppAD::AD\textless{}double\textgreater{} Vector}}}}

\item {} 
vec2a\_double: {\hyperref[\detokenize{xsrst/vec2a_double:id1}]{\sphinxcrossref{\DUrole{std,std-ref}{3.2.5.2 Convert a CppAD::AD\textless{}double\textgreater{} Vector to an a\_double Vector}}}}

\item {} 
error\_message: {\hyperref[\detokenize{xsrst/error_message:id1}]{\sphinxcrossref{\DUrole{std,std-ref}{3.2.5.3 Exception Handling}}}}

\end{enumerate}


\bigskip\hrule\bigskip


xsrst input file: \sphinxcode{\sphinxupquote{lib/cplusplus/cplusplus.xsrst}}


\subparagraph{more\_cpp}
\label{\detokenize{xsrst/more_cpp:more-cpp}}\label{\detokenize{xsrst/more_cpp::doc}}
\index{more\_cpp@\spxentry{more\_cpp}}\index{steps@\spxentry{steps}}\index{to@\spxentry{to}}\index{add@\spxentry{add}}\index{more@\spxentry{more}}\index{c++@\spxentry{c++}}\index{functions@\spxentry{functions}}\ignorespaces 

\subparagraph{3.2.6 Steps To Add More C++ Functions}
\label{\detokenize{xsrst/more_cpp:steps-to-add-more-c-functions}}\label{\detokenize{xsrst/more_cpp:id1}}\begin{itemize}
\item {} 
{\hyperref[\detokenize{xsrst/more_cpp:more-cpp-purpose}]{\sphinxcrossref{\DUrole{std,std-ref}{Purpose}}}}

\item {} \begin{description}
\item[{{\hyperref[\detokenize{xsrst/more_cpp:more-cpp-include-file}]{\sphinxcrossref{\DUrole{std,std-ref}{Include File}}}}}] \leavevmode\begin{itemize}
\item {} 
{\hyperref[\detokenize{xsrst/more_cpp:more-cpp-include-file-independent}]{\sphinxcrossref{\DUrole{std,std-ref}{independent}}}}

\item {} 
{\hyperref[\detokenize{xsrst/more_cpp:more-cpp-include-file-new-dynamic}]{\sphinxcrossref{\DUrole{std,std-ref}{new\_dynamic}}}}

\end{itemize}

\end{description}

\item {} \begin{description}
\item[{{\hyperref[\detokenize{xsrst/more_cpp:more-cpp-documentation}]{\sphinxcrossref{\DUrole{std,std-ref}{Documentation}}}}}] \leavevmode\begin{itemize}
\item {} 
{\hyperref[\detokenize{xsrst/more_cpp:more-cpp-documentation-independent}]{\sphinxcrossref{\DUrole{std,std-ref}{independent}}}}

\item {} 
{\hyperref[\detokenize{xsrst/more_cpp:more-cpp-documentation-new-dynamic}]{\sphinxcrossref{\DUrole{std,std-ref}{new\_dynamic}}}}

\item {} 
{\hyperref[\detokenize{xsrst/more_cpp:more-cpp-documentation-example}]{\sphinxcrossref{\DUrole{std,std-ref}{Example}}}}

\end{itemize}

\end{description}

\item {} \begin{description}
\item[{{\hyperref[\detokenize{xsrst/more_cpp:more-cpp-implementation}]{\sphinxcrossref{\DUrole{std,std-ref}{Implementation}}}}}] \leavevmode\begin{itemize}
\item {} 
{\hyperref[\detokenize{xsrst/more_cpp:more-cpp-implementation-independent}]{\sphinxcrossref{\DUrole{std,std-ref}{independent}}}}

\item {} 
{\hyperref[\detokenize{xsrst/more_cpp:more-cpp-implementation-new-dynamic}]{\sphinxcrossref{\DUrole{std,std-ref}{new\_dynamic}}}}

\end{itemize}

\end{description}

\item {} 
{\hyperref[\detokenize{xsrst/more_cpp:more-cpp-example}]{\sphinxcrossref{\DUrole{std,std-ref}{Example}}}}

\item {} 
{\hyperref[\detokenize{xsrst/more_cpp:more-cpp-testing}]{\sphinxcrossref{\DUrole{std,std-ref}{Testing}}}}

\end{itemize}

\index{purpose@\spxentry{purpose}}\ignorespaces 

\subparagraph{Purpose}
\label{\detokenize{xsrst/more_cpp:purpose}}\label{\detokenize{xsrst/more_cpp:more-cpp-purpose}}\label{\detokenize{xsrst/more_cpp:index-1}}
This section outlines the steps for adding more CppAD functionality
to cppad\_py.
This is done by example showing how the {\hyperref[\detokenize{xsrst/cpp_fun_new_dynamic:id1}]{\sphinxcrossref{\DUrole{std,std-ref}{cpp\_fun\_new\_dynamic}}}} was added.
This example case was chosen because it required both changing one
C++ function, {\hyperref[\detokenize{xsrst/cpp_independent:id1}]{\sphinxcrossref{\DUrole{std,std-ref}{cpp\_independent}}}},
and adding a new C++ function, {\hyperref[\detokenize{xsrst/cpp_fun_new_dynamic:id1}]{\sphinxcrossref{\DUrole{std,std-ref}{cpp\_fun\_new\_dynamic}}}}.

\index{include@\spxentry{include}}\index{file@\spxentry{file}}\ignorespaces 

\subparagraph{Include File}
\label{\detokenize{xsrst/more_cpp:include-file}}\label{\detokenize{xsrst/more_cpp:more-cpp-include-file}}\label{\detokenize{xsrst/more_cpp:index-2}}
\index{independent@\spxentry{independent}}\ignorespaces 

\subparagraph{independent}
\label{\detokenize{xsrst/more_cpp:independent}}\label{\detokenize{xsrst/more_cpp:more-cpp-include-file-independent}}\label{\detokenize{xsrst/more_cpp:index-3}}
The include file \sphinxcode{\sphinxupquote{include/cppad/py/fun.hpp}}
was edited to add the following prototype:

\begin{sphinxVerbatim}[commandchars=\\\{\}]
\PYG{n}{CPPAD\PYGZus{}PY\PYGZus{}LIB\PYGZus{}PUBLIC}
\PYG{n}{std}\PYG{o}{:}\PYG{o}{:}\PYG{n}{vector}\PYG{o}{\PYGZlt{}}\PYG{n}{a\PYGZus{}double}\PYG{o}{\PYGZgt{}} \PYG{n}{independent}\PYG{p}{(}
    \PYG{k}{const} \PYG{n}{std}\PYG{o}{:}\PYG{o}{:}\PYG{n}{vector}\PYG{o}{\PYGZlt{}}\PYG{k+kt}{double}\PYG{o}{\PYGZgt{}}\PYG{o}{\PYGZam{}} \PYG{n}{x}\PYG{p}{,} \PYG{k}{const} \PYG{n}{std}\PYG{o}{:}\PYG{o}{:}\PYG{n}{vector}\PYG{o}{\PYGZlt{}}\PYG{k+kt}{double}\PYG{o}{\PYGZgt{}}\PYG{o}{\PYGZam{}} \PYG{n}{dynamic}
\PYG{p}{)}\PYG{p}{;}
\end{sphinxVerbatim}

The \sphinxcode{\sphinxupquote{independent}} function is not part of the
\sphinxcode{\sphinxupquote{d\_fun}} or \sphinxcode{\sphinxupquote{a\_fun}} class, but
calling it is the first step in creating these objects.
This is why its prototype is in the \sphinxcode{\sphinxupquote{fun.hpp}} file.

\index{new\_dynamic@\spxentry{new\_dynamic}}\ignorespaces 

\subparagraph{new\_dynamic}
\label{\detokenize{xsrst/more_cpp:new-dynamic}}\label{\detokenize{xsrst/more_cpp:more-cpp-include-file-new-dynamic}}\label{\detokenize{xsrst/more_cpp:index-4}}
The include file \sphinxcode{\sphinxupquote{include/cppad/py/fun.hpp}}
was edited to add the following member function to \sphinxcode{\sphinxupquote{d\_fun}} :

\begin{sphinxVerbatim}[commandchars=\\\{\}]
\PYG{k+kt}{void} \PYG{n+nf}{new\PYGZus{}dynamic}\PYG{p}{(}\PYG{k}{const} \PYG{n}{std}\PYG{o}{:}\PYG{o}{:}\PYG{n}{vector}\PYG{o}{\PYGZlt{}}\PYG{k+kt}{double}\PYG{o}{\PYGZgt{}}\PYG{o}{\PYGZam{}} \PYG{n}{dynamic}\PYG{p}{)}\PYG{p}{;}
\end{sphinxVerbatim}

The following member function was added to \sphinxcode{\sphinxupquote{a\_fun}} :

\begin{sphinxVerbatim}[commandchars=\\\{\}]
\PYG{k+kt}{void} \PYG{n+nf}{new\PYGZus{}dynamic}\PYG{p}{(}\PYG{k}{const} \PYG{n}{std}\PYG{o}{:}\PYG{o}{:}\PYG{n}{vector}\PYG{o}{\PYGZlt{}}\PYG{n}{a\PYGZus{}double}\PYG{o}{\PYGZgt{}}\PYG{o}{\PYGZam{}} \PYG{n}{adynamic}\PYG{p}{)}\PYG{p}{;}
\end{sphinxVerbatim}

\index{documentation@\spxentry{documentation}}\ignorespaces 

\subparagraph{Documentation}
\label{\detokenize{xsrst/more_cpp:documentation}}\label{\detokenize{xsrst/more_cpp:more-cpp-documentation}}\label{\detokenize{xsrst/more_cpp:index-5}}
\index{independent@\spxentry{independent}}\ignorespaces 

\subparagraph{independent}
\label{\detokenize{xsrst/more_cpp:more-cpp-documentation-independent}}\label{\detokenize{xsrst/more_cpp:index-6}}\label{\detokenize{xsrst/more_cpp:id2}}
The C++ file \sphinxcode{\sphinxupquote{lib/cplusplus/fun.cpp}}
was edited to add the following syntax documentation:

\begin{DUlineblock}{0em}
\item[]         \sphinxstyleemphasis{a\_both} =  \sphinxcode{\sphinxupquote{cppad\_py::independent}} ( \sphinxstyleemphasis{x} , \sphinxstyleemphasis{dynamic} )
\end{DUlineblock}

and the corresponding return value was documented; see
{\hyperref[\detokenize{xsrst/cpp_independent:cpp-independent-a-both}]{\sphinxcrossref{\DUrole{std,std-ref}{a\_both}}}}.

\index{new\_dynamic@\spxentry{new\_dynamic}}\ignorespaces 

\subparagraph{new\_dynamic}
\label{\detokenize{xsrst/more_cpp:more-cpp-documentation-new-dynamic}}\label{\detokenize{xsrst/more_cpp:index-7}}\label{\detokenize{xsrst/more_cpp:id3}}
The {\hyperref[\detokenize{xsrst/cpp_fun_new_dynamic:id1}]{\sphinxcrossref{\DUrole{std,std-ref}{cpp\_fun\_new\_dynamic}}}} documentation was added.

\index{example@\spxentry{example}}\ignorespaces 

\subparagraph{Example}
\label{\detokenize{xsrst/more_cpp:example}}\label{\detokenize{xsrst/more_cpp:more-cpp-documentation-example}}\label{\detokenize{xsrst/more_cpp:index-8}}
The corresponding example file was added to the documentation,
below the {\hyperref[\detokenize{xsrst/cpp_independent:id1}]{\sphinxcrossref{\DUrole{std,std-ref}{cpp\_independent}}}} section, using the OMhelp commands:

\begin{DUlineblock}{0em}
\item[]         \{ \sphinxcode{\sphinxupquote{xsrst\_dollar\}children\%}}
\item[]                 \sphinxcode{\sphinxupquote{example/cplusplus/fun\_dynamic\_xam.cpp}}
\item[]         \sphinxcode{\sphinxupquote{\%\$\{xsrst\_dollar}} \}
\end{DUlineblock}

In addition, a reference to this example was added under the
{\hyperref[\detokenize{xsrst/cpp_independent:cpp-independent-example}]{\sphinxcrossref{\DUrole{std,std-ref}{example}}}} heading in the \sphinxcode{\sphinxupquote{independent}}
documentation.

\index{implementation@\spxentry{implementation}}\ignorespaces 

\subparagraph{Implementation}
\label{\detokenize{xsrst/more_cpp:implementation}}\label{\detokenize{xsrst/more_cpp:more-cpp-implementation}}\label{\detokenize{xsrst/more_cpp:index-9}}
\index{independent@\spxentry{independent}}\ignorespaces 

\subparagraph{independent}
\label{\detokenize{xsrst/more_cpp:more-cpp-implementation-independent}}\label{\detokenize{xsrst/more_cpp:index-10}}\label{\detokenize{xsrst/more_cpp:id4}}
The following function was added to the \sphinxcode{\sphinxupquote{lib/cplusplus/fun.cpp}} file:

\begin{sphinxVerbatim}[commandchars=\\\{\}]
\PYG{n}{std}\PYG{o}{:}\PYG{o}{:}\PYG{n}{vector}\PYG{o}{\PYGZlt{}}\PYG{n}{a\PYGZus{}double}\PYG{o}{\PYGZgt{}} \PYG{n}{independent}\PYG{p}{(}
    \PYG{k}{const} \PYG{n}{std}\PYG{o}{:}\PYG{o}{:}\PYG{n}{vector}\PYG{o}{\PYGZlt{}}\PYG{k+kt}{double}\PYG{o}{\PYGZgt{}}\PYG{o}{\PYGZam{}} \PYG{n}{x}       \PYG{p}{,}
    \PYG{k}{const} \PYG{n}{std}\PYG{o}{:}\PYG{o}{:}\PYG{n}{vector}\PYG{o}{\PYGZlt{}}\PYG{k+kt}{double}\PYG{o}{\PYGZgt{}}\PYG{o}{\PYGZam{}} \PYG{n}{dynamic} \PYG{p}{)}
\PYG{p}{\PYGZob{}}   \PYG{k}{using} \PYG{n}{CppAD}\PYG{o}{:}\PYG{o}{:}\PYG{n}{AD}\PYG{p}{;}
    \PYG{k+kt}{size\PYGZus{}t} \PYG{n}{nx} \PYG{o}{=} \PYG{n}{x}\PYG{p}{.}\PYG{n}{size}\PYG{p}{(}\PYG{p}{)}\PYG{p}{;}
    \PYG{k+kt}{size\PYGZus{}t} \PYG{n}{nd} \PYG{o}{=} \PYG{n}{dynamic}\PYG{p}{.}\PYG{n}{size}\PYG{p}{(}\PYG{p}{)}\PYG{p}{;}
    \PYG{n}{CppAD}\PYG{o}{:}\PYG{o}{:}\PYG{n}{vector}\PYG{o}{\PYGZlt{}} \PYG{n}{AD}\PYG{o}{\PYGZlt{}}\PYG{k+kt}{double}\PYG{o}{\PYGZgt{}} \PYG{o}{\PYGZgt{}} \PYG{n}{ax}\PYG{p}{(}\PYG{n}{nx}\PYG{p}{)}\PYG{p}{,} \PYG{n}{adynamic}\PYG{p}{(}\PYG{n}{nd}\PYG{p}{)}\PYG{p}{;}
    \PYG{k}{for}\PYG{p}{(}\PYG{k+kt}{size\PYGZus{}t} \PYG{n}{j} \PYG{o}{=} \PYG{l+m+mi}{0}\PYG{p}{;} \PYG{n}{j} \PYG{o}{\PYGZlt{}} \PYG{n}{nx}\PYG{p}{;} \PYG{n}{j}\PYG{o}{+}\PYG{o}{+}\PYG{p}{)}
        \PYG{n}{ax}\PYG{p}{[}\PYG{n}{j}\PYG{p}{]} \PYG{o}{=} \PYG{n}{x}\PYG{p}{[}\PYG{n}{j}\PYG{p}{]}\PYG{p}{;}
    \PYG{k}{for}\PYG{p}{(}\PYG{k+kt}{size\PYGZus{}t} \PYG{n}{j} \PYG{o}{=} \PYG{l+m+mi}{0}\PYG{p}{;} \PYG{n}{j} \PYG{o}{\PYGZlt{}} \PYG{n}{nd}\PYG{p}{;} \PYG{n}{j}\PYG{o}{+}\PYG{o}{+}\PYG{p}{)}
        \PYG{n}{adynamic}\PYG{p}{[}\PYG{n}{j}\PYG{p}{]} \PYG{o}{=} \PYG{n}{dynamic}\PYG{p}{[}\PYG{n}{j}\PYG{p}{]}\PYG{p}{;}
    \PYG{k+kt}{size\PYGZus{}t} \PYG{n}{abort\PYGZus{}op\PYGZus{}index} \PYG{o}{=} \PYG{l+m+mi}{0}\PYG{p}{;}
    \PYG{k+kt}{size\PYGZus{}t} \PYG{n}{record\PYGZus{}compare} \PYG{o}{=} \PYG{n+nb}{false}\PYG{p}{;}
    \PYG{n}{CppAD}\PYG{o}{:}\PYG{o}{:}\PYG{n}{Independent}\PYG{p}{(}\PYG{n}{ax}\PYG{p}{,} \PYG{n}{abort\PYGZus{}op\PYGZus{}index}\PYG{p}{,} \PYG{n}{record\PYGZus{}compare}\PYG{p}{,} \PYG{n}{adynamic}\PYG{p}{)}\PYG{p}{;}
    \PYG{n}{std}\PYG{o}{:}\PYG{o}{:}\PYG{n}{vector}\PYG{o}{\PYGZlt{}}\PYG{n}{a\PYGZus{}double}\PYG{o}{\PYGZgt{}} \PYG{n}{a\PYGZus{}both}\PYG{p}{(}\PYG{n}{nx} \PYG{o}{+} \PYG{n}{nd}\PYG{p}{)}\PYG{p}{;}
    \PYG{c+c1}{// use a\PYGZus{}double( *AD\PYGZlt{}double\PYGZgt{} ) constructor in these assignment loops}
    \PYG{k}{for}\PYG{p}{(}\PYG{k+kt}{size\PYGZus{}t} \PYG{n}{j} \PYG{o}{=} \PYG{l+m+mi}{0}\PYG{p}{;} \PYG{n}{j} \PYG{o}{\PYGZlt{}} \PYG{n}{nx}\PYG{p}{;} \PYG{n}{j}\PYG{o}{+}\PYG{o}{+}\PYG{p}{)}
        \PYG{n}{a\PYGZus{}both}\PYG{p}{[}\PYG{n}{j}\PYG{p}{]} \PYG{o}{=}  \PYG{o}{\PYGZam{}}\PYG{n}{ax}\PYG{p}{[}\PYG{n}{j}\PYG{p}{]} \PYG{p}{;}
    \PYG{k}{for}\PYG{p}{(}\PYG{k+kt}{size\PYGZus{}t} \PYG{n}{j} \PYG{o}{=} \PYG{l+m+mi}{0}\PYG{p}{;} \PYG{n}{j} \PYG{o}{\PYGZlt{}} \PYG{n}{nd}\PYG{p}{;} \PYG{n}{j}\PYG{o}{+}\PYG{o}{+}\PYG{p}{)}
        \PYG{n}{a\PYGZus{}both}\PYG{p}{[}\PYG{n}{nx} \PYG{o}{+} \PYG{n}{j}\PYG{p}{]} \PYG{o}{=}  \PYG{o}{\PYGZam{}}\PYG{n}{adynamic}\PYG{p}{[}\PYG{n}{j}\PYG{p}{]} \PYG{p}{;}
    \PYG{k}{return} \PYG{n}{a\PYGZus{}both}\PYG{p}{;}
\PYG{p}{\PYGZcb{}}
\end{sphinxVerbatim}

\index{new\_dynamic@\spxentry{new\_dynamic}}\ignorespaces 

\subparagraph{new\_dynamic}
\label{\detokenize{xsrst/more_cpp:more-cpp-implementation-new-dynamic}}\label{\detokenize{xsrst/more_cpp:index-11}}\label{\detokenize{xsrst/more_cpp:id5}}
The following function
was added to the \sphinxcode{\sphinxupquote{lib/cplusplus/fun.cpp}} file:

\begin{sphinxVerbatim}[commandchars=\\\{\}]
\PYG{k+kt}{void} \PYG{n}{d\PYGZus{}fun}\PYG{o}{:}\PYG{o}{:}\PYG{n}{new\PYGZus{}dynamic}\PYG{p}{(}\PYG{k}{const} \PYG{n}{std}\PYG{o}{:}\PYG{o}{:}\PYG{n}{vector}\PYG{o}{\PYGZlt{}}\PYG{k+kt}{double}\PYG{o}{\PYGZgt{}}\PYG{o}{\PYGZam{}} \PYG{n}{dynamic}\PYG{p}{)}
\PYG{p}{\PYGZob{}}   \PYG{k}{if}\PYG{p}{(} \PYG{n}{dynamic}\PYG{p}{.}\PYG{n}{size}\PYG{p}{(}\PYG{p}{)} \PYG{o}{!}\PYG{o}{=} \PYG{n}{ptr\PYGZus{}}\PYG{o}{\PYGZhy{}}\PYG{o}{\PYGZgt{}}\PYG{n}{size\PYGZus{}dyn\PYGZus{}ind}\PYG{p}{(}\PYG{p}{)} \PYG{p}{)}
        \PYG{n}{error\PYGZus{}message}\PYG{p}{(}\PYG{l+s}{\PYGZdq{}}\PYG{l+s}{cppad\PYGZus{}py::d\PYGZus{}fun::new\PYGZus{}dynamic dynamic.size() error}\PYG{l+s}{\PYGZdq{}}\PYG{p}{)}\PYG{p}{;}
    \PYG{n}{ptr\PYGZus{}}\PYG{o}{\PYGZhy{}}\PYG{o}{\PYGZgt{}}\PYG{n}{new\PYGZus{}dynamic}\PYG{p}{(}\PYG{n}{dynamic}\PYG{p}{)}\PYG{p}{;}
    \PYG{k}{return}\PYG{p}{;}
\PYG{p}{\PYGZcb{}}
\PYG{k+kt}{void} \PYG{n}{a\PYGZus{}fun}\PYG{o}{:}\PYG{o}{:}\PYG{n}{new\PYGZus{}dynamic}\PYG{p}{(}\PYG{k}{const} \PYG{n}{std}\PYG{o}{:}\PYG{o}{:}\PYG{n}{vector}\PYG{o}{\PYGZlt{}}\PYG{n}{a\PYGZus{}double}\PYG{o}{\PYGZgt{}}\PYG{o}{\PYGZam{}} \PYG{n}{adynamic}\PYG{p}{)}
\PYG{p}{\PYGZob{}}   \PYG{k}{if}\PYG{p}{(} \PYG{n}{adynamic}\PYG{p}{.}\PYG{n}{size}\PYG{p}{(}\PYG{p}{)} \PYG{o}{!}\PYG{o}{=} \PYG{n}{a\PYGZus{}ptr\PYGZus{}}\PYG{o}{\PYGZhy{}}\PYG{o}{\PYGZgt{}}\PYG{n}{size\PYGZus{}dyn\PYGZus{}ind}\PYG{p}{(}\PYG{p}{)} \PYG{p}{)}
        \PYG{n}{error\PYGZus{}message}\PYG{p}{(}\PYG{l+s}{\PYGZdq{}}\PYG{l+s}{cppad\PYGZus{}py::a\PYGZus{}fun::jacobian adynamic.size() error}\PYG{l+s}{\PYGZdq{}}\PYG{p}{)}\PYG{p}{;}
    \PYG{n}{std}\PYG{o}{:}\PYG{o}{:}\PYG{n}{vector}\PYG{o}{\PYGZlt{}} \PYG{n}{CppAD}\PYG{o}{:}\PYG{o}{:}\PYG{n}{AD}\PYG{o}{\PYGZlt{}}\PYG{k+kt}{double}\PYG{o}{\PYGZgt{}} \PYG{o}{\PYGZgt{}} \PYG{n}{au} \PYG{o}{=} \PYG{n}{vec2cppad\PYGZus{}double}\PYG{p}{(}\PYG{n}{adynamic}\PYG{p}{)}\PYG{p}{;}
    \PYG{n}{a\PYGZus{}ptr\PYGZus{}}\PYG{o}{\PYGZhy{}}\PYG{o}{\PYGZgt{}}\PYG{n}{new\PYGZus{}dynamic}\PYG{p}{(}\PYG{n}{au}\PYG{p}{)}\PYG{p}{;}
    \PYG{k}{return}\PYG{p}{;}
\PYG{p}{\PYGZcb{}}
\end{sphinxVerbatim}

A similar function was added to the \sphinxcode{\sphinxupquote{a\_fun}} class.

\index{example@\spxentry{example}}\ignorespaces 

\subparagraph{Example}
\label{\detokenize{xsrst/more_cpp:more-cpp-example}}\label{\detokenize{xsrst/more_cpp:index-12}}\label{\detokenize{xsrst/more_cpp:id6}}
The file \sphinxcode{\sphinxupquote{example/cplusplus/fun\_dynamic\_xam.cpp}} was added
with the following contents:
{\hyperref[\detokenize{xsrst/fun_dynamic_xam_cpp:id1}]{\sphinxcrossref{\DUrole{std,std-ref}{fun\_dynamic\_xam\_cpp}}}}.
In addition, in the file \sphinxcode{\sphinxupquote{example/cplusplus/check\_all.cpp}} :

\begin{DUlineblock}{0em}
\item[]         \sphinxcode{\sphinxupquote{extern bool fun\_optimize\_xam(void);}}
\end{DUlineblock}

was added to the external declarations and

\begin{DUlineblock}{0em}
\item[]         \sphinxcode{\sphinxupquote{ok \&= Run( fun\_optimize\_xam,          "fun\_optimize\_xam"        );}}
\end{DUlineblock}

was added to the main program.

\index{testing@\spxentry{testing}}\ignorespaces 

\subparagraph{Testing}
\label{\detokenize{xsrst/more_cpp:testing}}\label{\detokenize{xsrst/more_cpp:more-cpp-testing}}\label{\detokenize{xsrst/more_cpp:index-13}}
You must do a git add for all of the new files before running
\sphinxcode{\sphinxupquote{bin/check\_all.sh}} .
After all the C++ changes above were implemented,
\sphinxcode{\sphinxupquote{bin/check\_all.sh}} was run and the changes were made
until the warnings and errors were fixed.
The command

\begin{DUlineblock}{0em}
\item[]         \sphinxcode{\sphinxupquote{grep \textquotesingle{}fun\_dynamic\_xam\textquotesingle{} check\_all.log}}
\end{DUlineblock}

was used to make sure that the new C++ example / test was run.
Note that if a particular step in \sphinxcode{\sphinxupquote{bin/check\_all.sh}} is failing,
you can just re\sphinxhyphen{}run that step to see if a particular fix works.
Once the C++ tests were working, the changes where checked into using
\sphinxcode{\sphinxupquote{git}} .


\bigskip\hrule\bigskip


xsrst input file: \sphinxcode{\sphinxupquote{lib/cplusplus/more\_cpp.xsrst}}
\begin{enumerate}
\sphinxsetlistlabels{\arabic}{enumi}{enumii}{}{.}%
\item {} 
a\_double: {\hyperref[\detokenize{xsrst/a_double:id1}]{\sphinxcrossref{\DUrole{std,std-ref}{3.2.1 Cppad Py AD Scalars}}}}

\item {} 
vector: {\hyperref[\detokenize{xsrst/vector:id1}]{\sphinxcrossref{\DUrole{std,std-ref}{3.2.2 Cppad Py Vectors}}}}

\item {} 
cpp\_fun: {\hyperref[\detokenize{xsrst/cpp_fun:id1}]{\sphinxcrossref{\DUrole{std,std-ref}{3.2.3 Cppad Py AD Functions}}}}

\item {} 
sparse: {\hyperref[\detokenize{xsrst/sparse:id1}]{\sphinxcrossref{\DUrole{std,std-ref}{3.2.4 Cppad Py Sparse Calculation}}}}

\item {} 
cpp\_utility: {\hyperref[\detokenize{xsrst/cpp_utility:id1}]{\sphinxcrossref{\DUrole{std,std-ref}{3.2.5 C++ Utilities}}}}

\item {} 
more\_cpp: {\hyperref[\detokenize{xsrst/more_cpp:id1}]{\sphinxcrossref{\DUrole{std,std-ref}{3.2.6 Steps To Add More C++ Functions}}}}

\end{enumerate}


\bigskip\hrule\bigskip


xsrst input file: \sphinxcode{\sphinxupquote{lib/cplusplus/cpp\_lib.xsrst}}
\begin{enumerate}
\sphinxsetlistlabels{\arabic}{enumi}{enumii}{}{.}%
\item {} 
py\_lib: {\hyperref[\detokenize{xsrst/py_lib:id1}]{\sphinxcrossref{\DUrole{std,std-ref}{3.1 The Python Library}}}}

\item {} 
cpp\_lib: {\hyperref[\detokenize{xsrst/cpp_lib:id1}]{\sphinxcrossref{\DUrole{std,std-ref}{3.2 The C++ Library}}}}

\end{enumerate}


\bigskip\hrule\bigskip


xsrst input file: \sphinxcode{\sphinxupquote{lib/library.xsrst}}


\subsubsection{xsrst\_py}
\label{\detokenize{xsrst/xsrst_py:xsrst-py}}\label{\detokenize{xsrst/xsrst_py::doc}}
\index{xsrst\_py@\spxentry{xsrst\_py}}\index{extract@\spxentry{extract}}\index{sphinx@\spxentry{sphinx}}\index{rst@\spxentry{rst}}\ignorespaces 

\paragraph{4 Extract Sphinx RST}
\label{\detokenize{xsrst/xsrst_py:extract-sphinx-rst}}\label{\detokenize{xsrst/xsrst_py:id1}}\begin{itemize}
\item {} 
{\hyperref[\detokenize{xsrst/xsrst_py:xsrst-py-syntax}]{\sphinxcrossref{\DUrole{std,std-ref}{Syntax}}}}

\item {} 
{\hyperref[\detokenize{xsrst/xsrst_py:xsrst-py-purpose}]{\sphinxcrossref{\DUrole{std,std-ref}{Purpose}}}}

\item {} 
{\hyperref[\detokenize{xsrst/xsrst_py:xsrst-py-requirements}]{\sphinxcrossref{\DUrole{std,std-ref}{Requirements}}}}

\item {} \begin{description}
\item[{{\hyperref[\detokenize{xsrst/xsrst_py:xsrst-py-notation}]{\sphinxcrossref{\DUrole{std,std-ref}{Notation}}}}}] \leavevmode\begin{itemize}
\item {} 
{\hyperref[\detokenize{xsrst/xsrst_py:xsrst-py-notation-white-space}]{\sphinxcrossref{\DUrole{std,std-ref}{White Space}}}}

\item {} 
{\hyperref[\detokenize{xsrst/xsrst_py:xsrst-py-notation-beginning-of-a-line}]{\sphinxcrossref{\DUrole{std,std-ref}{Beginning of a Line}}}}

\end{itemize}

\end{description}

\item {} \begin{description}
\item[{{\hyperref[\detokenize{xsrst/xsrst_py:xsrst-py-command-line-arguments}]{\sphinxcrossref{\DUrole{std,std-ref}{Command Line Arguments}}}}}] \leavevmode\begin{itemize}
\item {} \begin{description}
\item[{{\hyperref[\detokenize{xsrst/xsrst_py:xsrst-py-command-line-arguments-root-file}]{\sphinxcrossref{\DUrole{std,std-ref}{root\_file}}}}}] \leavevmode\begin{itemize}
\item {} 
{\hyperref[\detokenize{xsrst/xsrst_py:xsrst-py-command-line-arguments-root-file-root-section}]{\sphinxcrossref{\DUrole{std,std-ref}{root\_section}}}}

\end{itemize}

\end{description}

\item {} \begin{description}
\item[{{\hyperref[\detokenize{xsrst/xsrst_py:xsrst-py-command-line-arguments-sphinx-dir}]{\sphinxcrossref{\DUrole{std,std-ref}{sphinx\_dir}}}}}] \leavevmode\begin{itemize}
\item {} 
{\hyperref[\detokenize{xsrst/xsrst_py:xsrst-py-command-line-arguments-sphinx-dir-example-configuration-files}]{\sphinxcrossref{\DUrole{std,std-ref}{Example Configuration Files}}}}

\end{itemize}

\end{description}

\item {} 
{\hyperref[\detokenize{xsrst/xsrst_py:xsrst-py-command-line-arguments-spelling}]{\sphinxcrossref{\DUrole{std,std-ref}{spelling}}}}

\item {} 
{\hyperref[\detokenize{xsrst/xsrst_py:xsrst-py-command-line-arguments-keyword}]{\sphinxcrossref{\DUrole{std,std-ref}{keyword}}}}

\item {} 
{\hyperref[\detokenize{xsrst/xsrst_py:xsrst-py-command-line-arguments-line-increment}]{\sphinxcrossref{\DUrole{std,std-ref}{line\_increment}}}}

\end{itemize}

\end{description}

\item {} \begin{description}
\item[{{\hyperref[\detokenize{xsrst/xsrst_py:xsrst-py-table-of-contents}]{\sphinxcrossref{\DUrole{std,std-ref}{Table Of Contents}}}}}] \leavevmode\begin{itemize}
\item {} 
{\hyperref[\detokenize{xsrst/xsrst_py:xsrst-py-table-of-contents-toctree}]{\sphinxcrossref{\DUrole{std,std-ref}{toctree}}}}

\item {} 
{\hyperref[\detokenize{xsrst/xsrst_py:xsrst-py-table-of-contents-parent-section}]{\sphinxcrossref{\DUrole{std,std-ref}{Parent Section}}}}

\end{itemize}

\end{description}

\item {} \begin{description}
\item[{{\hyperref[\detokenize{xsrst/xsrst_py:xsrst-py-heading-links}]{\sphinxcrossref{\DUrole{std,std-ref}{Heading Links}}}}}] \leavevmode\begin{itemize}
\item {} 
{\hyperref[\detokenize{xsrst/xsrst_py:xsrst-py-heading-links-section-level}]{\sphinxcrossref{\DUrole{std,std-ref}{Section Level}}}}

\item {} 
{\hyperref[\detokenize{xsrst/xsrst_py:xsrst-py-heading-links-other-levels}]{\sphinxcrossref{\DUrole{std,std-ref}{Other Levels}}}}

\item {} 
{\hyperref[\detokenize{xsrst/xsrst_py:xsrst-py-heading-links-children}]{\sphinxcrossref{\DUrole{std,std-ref}{Children}}}}

\item {} 
{\hyperref[\detokenize{xsrst/xsrst_py:xsrst-py-heading-links-example}]{\sphinxcrossref{\DUrole{std,std-ref}{Example}}}}

\end{itemize}

\end{description}

\item {} \begin{description}
\item[{{\hyperref[\detokenize{xsrst/xsrst_py:xsrst-py-indentation}]{\sphinxcrossref{\DUrole{std,std-ref}{Indentation}}}}}] \leavevmode\begin{itemize}
\item {} 
{\hyperref[\detokenize{xsrst/xsrst_py:xsrst-py-indentation-example}]{\sphinxcrossref{\DUrole{std,std-ref}{Example}}}}

\end{itemize}

\end{description}

\item {} \begin{description}
\item[{{\hyperref[\detokenize{xsrst/xsrst_py:xsrst-py-wish-list}]{\sphinxcrossref{\DUrole{std,std-ref}{Wish List}}}}}] \leavevmode\begin{itemize}
\item {} 
{\hyperref[\detokenize{xsrst/xsrst_py:xsrst-py-wish-list-tabs}]{\sphinxcrossref{\DUrole{std,std-ref}{Tabs}}}}

\item {} 
{\hyperref[\detokenize{xsrst/xsrst_py:xsrst-py-wish-list-module}]{\sphinxcrossref{\DUrole{std,std-ref}{Module}}}}

\end{itemize}

\end{description}

\item {} 
{\hyperref[\detokenize{xsrst/xsrst_py:xsrst-py-commands}]{\sphinxcrossref{\DUrole{std,std-ref}{Commands}}}}

\end{itemize}


\subparagraph{begin\_cmd}
\label{\detokenize{xsrst/begin_cmd:begin-cmd}}\label{\detokenize{xsrst/begin_cmd::doc}}
\index{begin\_cmd@\spxentry{begin\_cmd}}\index{begin@\spxentry{begin}}\index{end@\spxentry{end}}\index{commands@\spxentry{commands}}\ignorespaces 

\subparagraph{4.1 Begin and End Commands}
\label{\detokenize{xsrst/begin_cmd:begin-and-end-commands}}\label{\detokenize{xsrst/begin_cmd:id1}}\begin{itemize}
\item {} 
{\hyperref[\detokenize{xsrst/begin_cmd:begin-cmd-syntax}]{\sphinxcrossref{\DUrole{std,std-ref}{Syntax}}}}

\item {} 
{\hyperref[\detokenize{xsrst/begin_cmd:begin-cmd-section}]{\sphinxcrossref{\DUrole{std,std-ref}{Section}}}}

\item {} 
{\hyperref[\detokenize{xsrst/begin_cmd:begin-cmd-section-name}]{\sphinxcrossref{\DUrole{std,std-ref}{section\_name}}}}

\item {} 
{\hyperref[\detokenize{xsrst/begin_cmd:begin-cmd-output-file}]{\sphinxcrossref{\DUrole{std,std-ref}{Output File}}}}

\item {} 
{\hyperref[\detokenize{xsrst/begin_cmd:begin-cmd-parent-section}]{\sphinxcrossref{\DUrole{std,std-ref}{Parent Section}}}}

\end{itemize}

\index{syntax@\spxentry{syntax}}\ignorespaces 

\subparagraph{Syntax}
\label{\detokenize{xsrst/begin_cmd:syntax}}\label{\detokenize{xsrst/begin_cmd:begin-cmd-syntax}}\label{\detokenize{xsrst/begin_cmd:index-1}}\begin{itemize}
\item {} 
\sphinxcode{\sphinxupquote{\{xsrst\_begin}}        \sphinxstyleemphasis{section\_name} \sphinxcode{\sphinxupquote{\}}}

\item {} 
\sphinxcode{\sphinxupquote{\{xsrst\_begin\_parent}} \sphinxstyleemphasis{section\_name} \sphinxcode{\sphinxupquote{\}}}

\item {} 
\sphinxcode{\sphinxupquote{\{xsrst\_end}}          \sphinxstyleemphasis{section\_name} \sphinxcode{\sphinxupquote{\}}}

\end{itemize}

\index{section@\spxentry{section}}\ignorespaces 

\subparagraph{Section}
\label{\detokenize{xsrst/begin_cmd:section}}\label{\detokenize{xsrst/begin_cmd:begin-cmd-section}}\label{\detokenize{xsrst/begin_cmd:index-2}}
The start (end) of a section of the input file is indicated by a
begin (end) command at the
{\hyperref[\detokenize{xsrst/xsrst_py:xsrst-py-notation-beginning-of-a-line}]{\sphinxcrossref{\DUrole{std,std-ref}{beginning of a line}}}}.

\index{section\_name@\spxentry{section\_name}}\ignorespaces 

\subparagraph{section\_name}
\label{\detokenize{xsrst/begin_cmd:section-name}}\label{\detokenize{xsrst/begin_cmd:begin-cmd-section-name}}\label{\detokenize{xsrst/begin_cmd:index-3}}
The \sphinxstyleemphasis{section\_name} is a non\sphinxhyphen{}empty sequence of the following characters:
a\sphinxhyphen{}z, 0\sphinxhyphen{}9, and underbar \sphinxcode{\sphinxupquote{\_}}.
It can not begin with the characters \sphinxcode{\sphinxupquote{xsrst\_}}.
A link is included in the index under the section name
to the first heading the section.
The section name is also added to the html keyword meta data.

\index{output@\spxentry{output}}\index{file@\spxentry{file}}\ignorespaces 

\subparagraph{Output File}
\label{\detokenize{xsrst/begin_cmd:output-file}}\label{\detokenize{xsrst/begin_cmd:begin-cmd-output-file}}\label{\detokenize{xsrst/begin_cmd:index-4}}
The output file corresponding to \sphinxstyleemphasis{section\_name} is

        {\hyperref[\detokenize{xsrst/xsrst_py:xsrst-py-command-line-arguments-sphinx-dir}]{\sphinxcrossref{\DUrole{std,std-ref}{sphinx\_dir}}}}
\sphinxcode{\sphinxupquote{/xsrst/}} \sphinxstyleemphasis{section\_name} \sphinxcode{\sphinxupquote{.rst}}

\index{parent@\spxentry{parent}}\index{section@\spxentry{section}}\ignorespaces 

\subparagraph{Parent Section}
\label{\detokenize{xsrst/begin_cmd:parent-section}}\label{\detokenize{xsrst/begin_cmd:begin-cmd-parent-section}}\label{\detokenize{xsrst/begin_cmd:index-5}}
There can be at most one begin parent command in an input file.
In this case there must be other sections in the file
and they are children of the parent section.
The parent section is a child
of the section that included this file using a {\hyperref[\detokenize{xsrst/child_cmd:id1}]{\sphinxcrossref{\DUrole{std,std-ref}{child command}}}}.

If there is no parent command in an input file,
all the sections in the file are children
of the section that included this file using a {\hyperref[\detokenize{xsrst/child_cmd:id1}]{\sphinxcrossref{\DUrole{std,std-ref}{child command}}}}.


\bigskip\hrule\bigskip


xsrst input file: \sphinxcode{\sphinxupquote{bin/xsrst.py}}


\subparagraph{child\_cmd}
\label{\detokenize{xsrst/child_cmd:child-cmd}}\label{\detokenize{xsrst/child_cmd::doc}}
\index{child\_cmd@\spxentry{child\_cmd}}\index{children@\spxentry{children}}\index{commands@\spxentry{commands}}\ignorespaces 

\subparagraph{4.2 Children Commands}
\label{\detokenize{xsrst/child_cmd:children-commands}}\label{\detokenize{xsrst/child_cmd:id1}}\begin{itemize}
\item {} 
{\hyperref[\detokenize{xsrst/child_cmd:child-cmd-syntax}]{\sphinxcrossref{\DUrole{std,std-ref}{Syntax}}}}

\item {} 
{\hyperref[\detokenize{xsrst/child_cmd:child-cmd-purpose}]{\sphinxcrossref{\DUrole{std,std-ref}{Purpose}}}}

\item {} 
{\hyperref[\detokenize{xsrst/child_cmd:child-cmd-file-names}]{\sphinxcrossref{\DUrole{std,std-ref}{File Names}}}}

\item {} 
{\hyperref[\detokenize{xsrst/child_cmd:child-cmd-links}]{\sphinxcrossref{\DUrole{std,std-ref}{Links}}}}

\item {} 
{\hyperref[\detokenize{xsrst/child_cmd:child-cmd-example}]{\sphinxcrossref{\DUrole{std,std-ref}{Example}}}}

\end{itemize}

\index{syntax@\spxentry{syntax}}\ignorespaces 

\subparagraph{Syntax}
\label{\detokenize{xsrst/child_cmd:syntax}}\label{\detokenize{xsrst/child_cmd:child-cmd-syntax}}\label{\detokenize{xsrst/child_cmd:index-1}}
\begin{DUlineblock}{0em}
\item[] \sphinxcode{\sphinxupquote{\{xsrst\_children}}
\item[]
\begin{DUlineblock}{\DUlineblockindent}
\item[] \sphinxstyleemphasis{file\_1}
\item[] …
\item[] \sphinxstyleemphasis{file\_n}
\end{DUlineblock}
\item[] \sphinxcode{\sphinxupquote{\}}}
\item[] 
\item[] \sphinxcode{\sphinxupquote{\{xsrst\_child\_link}}
\item[]
\begin{DUlineblock}{\DUlineblockindent}
\item[] \sphinxstyleemphasis{file\_1}
\item[] …
\item[] \sphinxstyleemphasis{file\_n}
\end{DUlineblock}
\item[] \sphinxcode{\sphinxupquote{\}}}
\item[] 
\item[] \sphinxcode{\sphinxupquote{\{xsrst\_child\_list}}
\item[]
\begin{DUlineblock}{\DUlineblockindent}
\item[] \sphinxstyleemphasis{file\_1}
\item[] …
\item[] \sphinxstyleemphasis{file\_n}
\end{DUlineblock}
\item[] \sphinxcode{\sphinxupquote{\}}}
\end{DUlineblock}

\index{purpose@\spxentry{purpose}}\ignorespaces 

\subparagraph{Purpose}
\label{\detokenize{xsrst/child_cmd:purpose}}\label{\detokenize{xsrst/child_cmd:child-cmd-purpose}}\label{\detokenize{xsrst/child_cmd:index-2}}
A section can specify a set of files for which the
{\hyperref[\detokenize{xsrst/begin_cmd:begin-cmd-parent-section}]{\sphinxcrossref{\DUrole{std,std-ref}{parent section}}}} of each file
is a child of the current section.
(If there is not parent section in a file,
all the sections in the file are children of the current section.)
This is done using the commands above at the
{\hyperref[\detokenize{xsrst/xsrst_py:xsrst-py-notation-beginning-of-a-line}]{\sphinxcrossref{\DUrole{std,std-ref}{beginning of a line}}}}.

\index{file@\spxentry{file}}\index{names@\spxentry{names}}\ignorespaces 

\subparagraph{File Names}
\label{\detokenize{xsrst/child_cmd:file-names}}\label{\detokenize{xsrst/child_cmd:child-cmd-file-names}}\label{\detokenize{xsrst/child_cmd:index-3}}
A new line character must precede and follow each
of the file names \sphinxstyleemphasis{file\_1} … \sphinxstyleemphasis{file\_n}.
Leading and trailing white space is not included in the names
The file names are  relative to the directory where \sphinxcode{\sphinxupquote{xsrst.py}}
is executed; i.e., the top directory for this git repository.
This may seem verbose, but it makes it easier to write scripts
that move files and automatically change references to them.

\index{links@\spxentry{links}}\ignorespaces 

\subparagraph{Links}
\label{\detokenize{xsrst/child_cmd:links}}\label{\detokenize{xsrst/child_cmd:child-cmd-links}}\label{\detokenize{xsrst/child_cmd:index-4}}
The child link and list commands also place
links to all the children of the current at the location of the command.
The links are displayed using the title for each section.
The child list command includes the section name next to the title.
You can place a heading directly before the links to make them easier to find.

\index{example@\spxentry{example}}\ignorespaces 

\subparagraph{Example}
\label{\detokenize{xsrst/child_cmd:example}}\label{\detokenize{xsrst/child_cmd:child-cmd-example}}\label{\detokenize{xsrst/child_cmd:index-5}}

\subparagraph{no\_parent\_exam}
\label{\detokenize{xsrst/no_parent_exam:no-parent-exam}}\label{\detokenize{xsrst/no_parent_exam::doc}}
\index{no\_parent\_exam@\spxentry{no\_parent\_exam}}\index{no@\spxentry{no}}\index{parent@\spxentry{parent}}\index{example@\spxentry{example}}\ignorespaces 

\subparagraph{4.2.1 No Parent Example}
\label{\detokenize{xsrst/no_parent_exam:no-parent-example}}\label{\detokenize{xsrst/no_parent_exam:id1}}
All the sections in the file \sphinxcode{\sphinxupquote{children.py}}
are children of the section below
because \sphinxcode{\sphinxupquote{children.py}} does not have a
{\hyperref[\detokenize{xsrst/begin_cmd:begin-cmd-parent-section}]{\sphinxcrossref{\DUrole{std,std-ref}{parent section}}}}:

\begin{sphinxVerbatim}[commandchars=\\\{\}]
\PYGZob{}xsrst\PYGZus{}begin no\PYGZus{}parent\PYGZus{}res\PYGZcb{}

\PYG{g+gh}{No Parent Result}
\PYG{g+gh}{\PYGZsh{}\PYGZsh{}\PYGZsh{}\PYGZsh{}\PYGZsh{}\PYGZsh{}\PYGZsh{}\PYGZsh{}\PYGZsh{}\PYGZsh{}\PYGZsh{}\PYGZsh{}\PYGZsh{}\PYGZsh{}\PYGZsh{}\PYGZsh{}}

\PYG{g+gh}{children.py File}
\PYG{g+gh}{****************}
\PYGZob{}xsrst\PYGZus{}file
    \PYGZsh{} BEGIN\PYGZus{}FILE
    \PYGZsh{} END\PYGZus{}FILE
    sphinx/test\PYGZus{}in/children.py
\PYGZcb{}

\PYG{g+gh}{Child Link children.py}
\PYG{g+gh}{**********************}
\PYGZob{}xsrst\PYGZus{}child\PYGZus{}link
        sphinx/test\PYGZus{}in/children.py
\PYGZcb{}

\PYG{g+gh}{Example}
\PYG{g+gh}{*******}
\PYG{n+na}{:ref:}\PYG{n+nv}{`no\PYGZus{}parent\PYGZus{}exam`}

\PYGZob{}xsrst\PYGZus{}end no\PYGZus{}parent\PYGZus{}res\PYGZcb{}
\end{sphinxVerbatim}


\subparagraph{no\_parent\_res}
\label{\detokenize{xsrst/no_parent_res:no-parent-res}}\label{\detokenize{xsrst/no_parent_res::doc}}
\index{no\_parent\_res@\spxentry{no\_parent\_res}}\index{no@\spxentry{no}}\index{parent@\spxentry{parent}}\index{result@\spxentry{result}}\ignorespaces 

\subparagraph{4.2.1.1 No Parent Result}
\label{\detokenize{xsrst/no_parent_res:no-parent-result}}\label{\detokenize{xsrst/no_parent_res:id1}}\begin{itemize}
\item {} 
{\hyperref[\detokenize{xsrst/no_parent_res:no-parent-res-children-py-file}]{\sphinxcrossref{\DUrole{std,std-ref}{children.py File}}}}

\item {} 
{\hyperref[\detokenize{xsrst/no_parent_res:no-parent-res-child-link-children-py}]{\sphinxcrossref{\DUrole{std,std-ref}{Child Link children.py}}}}

\item {} 
{\hyperref[\detokenize{xsrst/no_parent_res:no-parent-res-example}]{\sphinxcrossref{\DUrole{std,std-ref}{Example}}}}

\end{itemize}

\index{children.py@\spxentry{children.py}}\index{file@\spxentry{file}}\ignorespaces 

\subparagraph{children.py File}
\label{\detokenize{xsrst/no_parent_res:children-py-file}}\label{\detokenize{xsrst/no_parent_res:no-parent-res-children-py-file}}\label{\detokenize{xsrst/no_parent_res:index-1}}
\begin{sphinxVerbatim}[commandchars=\\\{\}]
\PYG{l+s+sd}{\PYGZdq{}\PYGZdq{}\PYGZdq{}}
\PYG{l+s+sd}{\PYGZob{}xsrst\PYGZus{}begin children\PYGZus{}exam\PYGZcb{}}

\PYG{l+s+sd}{Indent and Children Example}
\PYG{l+s+sd}{\PYGZsh{}\PYGZsh{}\PYGZsh{}\PYGZsh{}\PYGZsh{}\PYGZsh{}\PYGZsh{}\PYGZsh{}\PYGZsh{}\PYGZsh{}\PYGZsh{}\PYGZsh{}\PYGZsh{}\PYGZsh{}\PYGZsh{}\PYGZsh{}\PYGZsh{}\PYGZsh{}\PYGZsh{}\PYGZsh{}\PYGZsh{}\PYGZsh{}\PYGZsh{}\PYGZsh{}\PYGZsh{}\PYGZsh{}\PYGZsh{}}

\PYG{l+s+sd}{\PYGZob{}xsrst\PYGZus{}file}
\PYG{l+s+sd}{    \PYGZsh{} BEGIN\PYGZus{}SRC}
\PYG{l+s+sd}{    \PYGZsh{} END\PYGZus{}SRC}
\PYG{l+s+sd}{\PYGZcb{}}

\PYG{l+s+sd}{Result}
\PYG{l+s+sd}{******}
\PYG{l+s+sd}{:ref:`children\PYGZus{}res`}

\PYG{l+s+sd}{\PYGZob{}xsrst\PYGZus{}end children\PYGZus{}exam\PYGZcb{}}
\PYG{l+s+sd}{\PYGZdq{}\PYGZdq{}\PYGZdq{}}
\PYG{c+c1}{\PYGZsh{} \PYGZhy{}\PYGZhy{}\PYGZhy{}\PYGZhy{}\PYGZhy{}\PYGZhy{}\PYGZhy{}\PYGZhy{}\PYGZhy{}\PYGZhy{}\PYGZhy{}\PYGZhy{}\PYGZhy{}\PYGZhy{}\PYGZhy{}\PYGZhy{}\PYGZhy{}\PYGZhy{}\PYGZhy{}\PYGZhy{}\PYGZhy{}\PYGZhy{}\PYGZhy{}\PYGZhy{}\PYGZhy{}\PYGZhy{}\PYGZhy{}\PYGZhy{}\PYGZhy{}\PYGZhy{}\PYGZhy{}\PYGZhy{}\PYGZhy{}\PYGZhy{}\PYGZhy{}\PYGZhy{}\PYGZhy{}\PYGZhy{}\PYGZhy{}\PYGZhy{}\PYGZhy{}\PYGZhy{}\PYGZhy{}\PYGZhy{}\PYGZhy{}\PYGZhy{}\PYGZhy{}\PYGZhy{}\PYGZhy{}\PYGZhy{}\PYGZhy{}\PYGZhy{}\PYGZhy{}\PYGZhy{}\PYGZhy{}\PYGZhy{}\PYGZhy{}\PYGZhy{}\PYGZhy{}\PYGZhy{}\PYGZhy{}\PYGZhy{}\PYGZhy{}\PYGZhy{}\PYGZhy{}\PYGZhy{}\PYGZhy{}\PYGZhy{}\PYGZhy{}\PYGZhy{}\PYGZhy{}\PYGZhy{}\PYGZhy{}\PYGZhy{}\PYGZhy{}\PYGZhy{}}
\PYG{c+c1}{\PYGZsh{} BEGIN\PYGZus{}SRC}
\PYG{l+s+sd}{\PYGZdq{}\PYGZdq{}\PYGZdq{}}
\PYG{l+s+sd}{\PYGZob{}xsrst\PYGZus{}begin children\PYGZus{}res\PYGZcb{}}

\PYG{l+s+sd}{Indent and Children Result}
\PYG{l+s+sd}{\PYGZsh{}\PYGZsh{}\PYGZsh{}\PYGZsh{}\PYGZsh{}\PYGZsh{}\PYGZsh{}\PYGZsh{}\PYGZsh{}\PYGZsh{}\PYGZsh{}\PYGZsh{}\PYGZsh{}\PYGZsh{}\PYGZsh{}\PYGZsh{}\PYGZsh{}\PYGZsh{}\PYGZsh{}\PYGZsh{}\PYGZsh{}\PYGZsh{}\PYGZsh{}\PYGZsh{}\PYGZsh{}\PYGZsh{}}

\PYG{l+s+sd}{Children}
\PYG{l+s+sd}{********}
\PYG{l+s+sd}{This example uses a children command and}
\PYG{l+s+sd}{the corresponding link is included.}
\PYG{l+s+sd}{If it had used a child link command, the same link would have been}
\PYG{l+s+sd}{created automatically.}

\PYG{l+s+sd}{\PYGZob{}xsrst\PYGZus{}children}
\PYG{l+s+sd}{    sphinx/test\PYGZus{}in/indent.py}
\PYG{l+s+sd}{\PYGZcb{}}

\PYG{l+s+sd}{\PYGZhy{} :ref:`indent\PYGZus{}exam`}

\PYG{l+s+sd}{Example}
\PYG{l+s+sd}{*******}
\PYG{l+s+sd}{:ref:`children\PYGZus{}exam`}

\PYG{l+s+sd}{\PYGZob{}xsrst\PYGZus{}end children\PYGZus{}res\PYGZcb{}}
\PYG{l+s+sd}{\PYGZdq{}\PYGZdq{}\PYGZdq{}}
\PYG{c+c1}{\PYGZsh{} END\PYGZus{}SRC}
\end{sphinxVerbatim}

\index{child@\spxentry{child}}\index{link@\spxentry{link}}\index{children.py@\spxentry{children.py}}\ignorespaces 

\subparagraph{Child Link children.py}
\label{\detokenize{xsrst/no_parent_res:child-link-children-py}}\label{\detokenize{xsrst/no_parent_res:no-parent-res-child-link-children-py}}\label{\detokenize{xsrst/no_parent_res:index-2}}

\subparagraph{children\_exam}
\label{\detokenize{xsrst/children_exam:children-exam}}\label{\detokenize{xsrst/children_exam::doc}}
\index{children\_exam@\spxentry{children\_exam}}\index{indent@\spxentry{indent}}\index{children@\spxentry{children}}\index{example@\spxentry{example}}\ignorespaces 

\subparagraph{4.2.1.1.1 Indent and Children Example}
\label{\detokenize{xsrst/children_exam:indent-and-children-example}}\label{\detokenize{xsrst/children_exam:id1}}\begin{itemize}
\item {} 
{\hyperref[\detokenize{xsrst/children_exam:children-exam-result}]{\sphinxcrossref{\DUrole{std,std-ref}{Result}}}}

\end{itemize}

\begin{sphinxVerbatim}[commandchars=\\\{\}]
\PYG{l+s+sd}{\PYGZdq{}\PYGZdq{}\PYGZdq{}}
\PYG{l+s+sd}{\PYGZob{}xsrst\PYGZus{}begin children\PYGZus{}res\PYGZcb{}}

\PYG{l+s+sd}{Indent and Children Result}
\PYG{l+s+sd}{\PYGZsh{}\PYGZsh{}\PYGZsh{}\PYGZsh{}\PYGZsh{}\PYGZsh{}\PYGZsh{}\PYGZsh{}\PYGZsh{}\PYGZsh{}\PYGZsh{}\PYGZsh{}\PYGZsh{}\PYGZsh{}\PYGZsh{}\PYGZsh{}\PYGZsh{}\PYGZsh{}\PYGZsh{}\PYGZsh{}\PYGZsh{}\PYGZsh{}\PYGZsh{}\PYGZsh{}\PYGZsh{}\PYGZsh{}}

\PYG{l+s+sd}{Children}
\PYG{l+s+sd}{********}
\PYG{l+s+sd}{This example uses a children command and}
\PYG{l+s+sd}{the corresponding link is included.}
\PYG{l+s+sd}{If it had used a child link command, the same link would have been}
\PYG{l+s+sd}{created automatically.}

\PYG{l+s+sd}{\PYGZob{}xsrst\PYGZus{}children}
\PYG{l+s+sd}{    sphinx/test\PYGZus{}in/indent.py}
\PYG{l+s+sd}{\PYGZcb{}}

\PYG{l+s+sd}{\PYGZhy{} :ref:`indent\PYGZus{}exam`}

\PYG{l+s+sd}{Example}
\PYG{l+s+sd}{*******}
\PYG{l+s+sd}{:ref:`children\PYGZus{}exam`}

\PYG{l+s+sd}{\PYGZob{}xsrst\PYGZus{}end children\PYGZus{}res\PYGZcb{}}
\PYG{l+s+sd}{\PYGZdq{}\PYGZdq{}\PYGZdq{}}
\end{sphinxVerbatim}

\index{result@\spxentry{result}}\ignorespaces 

\subparagraph{Result}
\label{\detokenize{xsrst/children_exam:result}}\label{\detokenize{xsrst/children_exam:children-exam-result}}\label{\detokenize{xsrst/children_exam:index-1}}
{\hyperref[\detokenize{xsrst/children_res:id1}]{\sphinxcrossref{\DUrole{std,std-ref}{4.2.1.1.2 Indent and Children Result}}}}


\bigskip\hrule\bigskip


xsrst input file: \sphinxcode{\sphinxupquote{sphinx/test\_in/children.py}}


\subparagraph{children\_res}
\label{\detokenize{xsrst/children_res:children-res}}\label{\detokenize{xsrst/children_res::doc}}
\index{children\_res@\spxentry{children\_res}}\index{indent@\spxentry{indent}}\index{children@\spxentry{children}}\index{result@\spxentry{result}}\ignorespaces 

\subparagraph{4.2.1.1.2 Indent and Children Result}
\label{\detokenize{xsrst/children_res:indent-and-children-result}}\label{\detokenize{xsrst/children_res:id1}}\begin{itemize}
\item {} 
{\hyperref[\detokenize{xsrst/children_res:children-res-children}]{\sphinxcrossref{\DUrole{std,std-ref}{Children}}}}

\item {} 
{\hyperref[\detokenize{xsrst/children_res:children-res-example}]{\sphinxcrossref{\DUrole{std,std-ref}{Example}}}}

\end{itemize}

\index{children@\spxentry{children}}\ignorespaces 

\subparagraph{Children}
\label{\detokenize{xsrst/children_res:children}}\label{\detokenize{xsrst/children_res:children-res-children}}\label{\detokenize{xsrst/children_res:index-1}}
This example uses a children command and
the corresponding link is included.
If it had used a child link command, the same link would have been
created automatically.


\subparagraph{indent\_exam}
\label{\detokenize{xsrst/indent_exam:indent-exam}}\label{\detokenize{xsrst/indent_exam::doc}}
\index{indent\_exam@\spxentry{indent\_exam}}\index{indent@\spxentry{indent}}\index{example@\spxentry{example}}\ignorespaces 

\subparagraph{4.2.1.1.2.1 Indent Example}
\label{\detokenize{xsrst/indent_exam:indent-example}}\label{\detokenize{xsrst/indent_exam:id1}}
\begin{sphinxVerbatim}[commandchars=\\\{\}]
\PYG{l+s+sd}{\PYGZdq{}\PYGZdq{}\PYGZdq{}}
\PYG{l+s+sd}{    \PYGZob{}xsrst\PYGZus{}begin indent\PYGZus{}res\PYGZcb{}}

\PYG{l+s+sd}{    =============}
\PYG{l+s+sd}{    Indent Result}
\PYG{l+s+sd}{    =============}
\PYG{l+s+sd}{    \PYGZob{}xsrst\PYGZus{}code py\PYGZcb{}\PYGZdq{}\PYGZdq{}\PYGZdq{}}
    \PYG{k}{def} \PYG{n+nf}{factorial}\PYG{p}{(}\PYG{n}{n}\PYG{p}{)} \PYG{p}{:}
        \PYG{k}{if} \PYG{n}{n} \PYG{o}{==} \PYG{l+m+mi}{1} \PYG{p}{:}
            \PYG{k}{return} \PYG{l+m+mi}{1}
        \PYG{k}{return} \PYG{n}{n} \PYG{o}{*} \PYG{n}{factorial}\PYG{p}{(}\PYG{n}{n}\PYG{o}{\PYGZhy{}}\PYG{l+m+mi}{1}\PYG{p}{)}
    \PYG{l+s+sd}{\PYGZdq{}\PYGZdq{}\PYGZdq{}\PYGZob{}xsrst\PYGZus{}code\PYGZcb{}}

\PYG{l+s+sd}{    :ref:`indent\PYGZus{}exam`}

\PYG{l+s+sd}{    \PYGZob{}xsrst\PYGZus{}end indent\PYGZus{}res\PYGZcb{}}
\PYG{l+s+sd}{\PYGZdq{}\PYGZdq{}\PYGZdq{}}
\end{sphinxVerbatim}


\subparagraph{indent\_res}
\label{\detokenize{xsrst/indent_res:indent-res}}\label{\detokenize{xsrst/indent_res::doc}}
\index{indent\_res@\spxentry{indent\_res}}\index{indent@\spxentry{indent}}\index{result@\spxentry{result}}\ignorespaces 

\subparagraph{4.2.1.1.2.1.1 Indent Result}
\label{\detokenize{xsrst/indent_res:indent-result}}\label{\detokenize{xsrst/indent_res:id1}}
\begin{sphinxVerbatim}[commandchars=\\\{\}]
\PYG{k}{def} \PYG{n+nf}{factorial}\PYG{p}{(}\PYG{n}{n}\PYG{p}{)} \PYG{p}{:}
    \PYG{k}{if} \PYG{n}{n} \PYG{o}{==} \PYG{l+m+mi}{1} \PYG{p}{:}
        \PYG{k}{return} \PYG{l+m+mi}{1}
    \PYG{k}{return} \PYG{n}{n} \PYG{o}{*} \PYG{n}{factorial}\PYG{p}{(}\PYG{n}{n}\PYG{o}{\PYGZhy{}}\PYG{l+m+mi}{1}\PYG{p}{)}
\end{sphinxVerbatim}

{\hyperref[\detokenize{xsrst/indent_exam:id1}]{\sphinxcrossref{\DUrole{std,std-ref}{4.2.1.1.2.1 Indent Example}}}}


\bigskip\hrule\bigskip


xsrst input file: \sphinxcode{\sphinxupquote{sphinx/test\_in/indent.py}}


\bigskip\hrule\bigskip


xsrst input file: \sphinxcode{\sphinxupquote{sphinx/test\_in/indent.py}}
\begin{itemize}
\item {} 
{\hyperref[\detokenize{xsrst/indent_exam:id1}]{\sphinxcrossref{\DUrole{std,std-ref}{4.2.1.1.2.1 Indent Example}}}}

\end{itemize}

\index{example@\spxentry{example}}\ignorespaces 

\subparagraph{Example}
\label{\detokenize{xsrst/children_res:example}}\label{\detokenize{xsrst/children_res:children-res-example}}\label{\detokenize{xsrst/children_res:index-2}}
{\hyperref[\detokenize{xsrst/children_exam:id1}]{\sphinxcrossref{\DUrole{std,std-ref}{4.2.1.1.1 Indent and Children Example}}}}


\bigskip\hrule\bigskip


xsrst input file: \sphinxcode{\sphinxupquote{sphinx/test\_in/children.py}}

\index{example@\spxentry{example}}\ignorespaces 

\subparagraph{Example}
\label{\detokenize{xsrst/no_parent_res:example}}\label{\detokenize{xsrst/no_parent_res:no-parent-res-example}}\label{\detokenize{xsrst/no_parent_res:index-3}}
{\hyperref[\detokenize{xsrst/no_parent_exam:id1}]{\sphinxcrossref{\DUrole{std,std-ref}{4.2.1 No Parent Example}}}}


\bigskip\hrule\bigskip


xsrst input file: \sphinxcode{\sphinxupquote{sphinx/test\_in/no\_parent.xsrst}}


\bigskip\hrule\bigskip


xsrst input file: \sphinxcode{\sphinxupquote{sphinx/test\_in/no\_parent.xsrst}}
\begin{enumerate}
\sphinxsetlistlabels{\arabic}{enumi}{enumii}{}{.}%
\item {} 
no\_parent\_exam: {\hyperref[\detokenize{xsrst/no_parent_exam:id1}]{\sphinxcrossref{\DUrole{std,std-ref}{4.2.1 No Parent Example}}}}

\end{enumerate}


\bigskip\hrule\bigskip


xsrst input file: \sphinxcode{\sphinxupquote{bin/xsrst.py}}


\subparagraph{spell\_cmd}
\label{\detokenize{xsrst/spell_cmd:spell-cmd}}\label{\detokenize{xsrst/spell_cmd::doc}}
\index{spell\_cmd@\spxentry{spell\_cmd}}\index{spell@\spxentry{spell}}\index{command@\spxentry{command}}\ignorespaces 

\subparagraph{4.3 Spell Command}
\label{\detokenize{xsrst/spell_cmd:spell-command}}\label{\detokenize{xsrst/spell_cmd:id1}}\begin{itemize}
\item {} 
{\hyperref[\detokenize{xsrst/spell_cmd:spell-cmd-syntax}]{\sphinxcrossref{\DUrole{std,std-ref}{Syntax}}}}

\item {} 
{\hyperref[\detokenize{xsrst/spell_cmd:spell-cmd-purpose}]{\sphinxcrossref{\DUrole{std,std-ref}{Purpose}}}}

\item {} 
{\hyperref[\detokenize{xsrst/spell_cmd:spell-cmd-spelling}]{\sphinxcrossref{\DUrole{std,std-ref}{spelling}}}}

\item {} 
{\hyperref[\detokenize{xsrst/spell_cmd:spell-cmd-capital-letters}]{\sphinxcrossref{\DUrole{std,std-ref}{Capital Letters}}}}

\item {} 
{\hyperref[\detokenize{xsrst/spell_cmd:spell-cmd-double-words}]{\sphinxcrossref{\DUrole{std,std-ref}{Double Words}}}}

\item {} 
{\hyperref[\detokenize{xsrst/spell_cmd:spell-cmd-example}]{\sphinxcrossref{\DUrole{std,std-ref}{Example}}}}

\end{itemize}

\index{syntax@\spxentry{syntax}}\ignorespaces 

\subparagraph{Syntax}
\label{\detokenize{xsrst/spell_cmd:syntax}}\label{\detokenize{xsrst/spell_cmd:spell-cmd-syntax}}\label{\detokenize{xsrst/spell_cmd:index-1}}
\sphinxcode{\sphinxupquote{\{xsrst\_spell}} \sphinxstyleemphasis{word\_1} …  \sphinxstyleemphasis{word\_n} \sphinxcode{\sphinxupquote{\}}}

Here \sphinxstyleemphasis{word\_1}, …, \sphinxstyleemphasis{word\_n} is the special list of words for this section.
In the syntax above the list of words is all in one line,
but they could be on different lines.
Each word starts with an upper case letter,
a lower case letter, or a back slash.
The back slash is included as a possible beginning of a word
so that latex commands can be included in the spelling list.
The rest of the characters in a word are lower case letters.

\index{purpose@\spxentry{purpose}}\ignorespaces 

\subparagraph{Purpose}
\label{\detokenize{xsrst/spell_cmd:purpose}}\label{\detokenize{xsrst/spell_cmd:spell-cmd-purpose}}\label{\detokenize{xsrst/spell_cmd:index-2}}
You can specify a special list of words
(not normally considered correct spelling)
for the current section using the command above at the
{\hyperref[\detokenize{xsrst/xsrst_py:xsrst-py-notation-beginning-of-a-line}]{\sphinxcrossref{\DUrole{std,std-ref}{beginning of a line}}}}.

\index{spelling@\spxentry{spelling}}\ignorespaces 

\subparagraph{spelling}
\label{\detokenize{xsrst/spell_cmd:spelling}}\label{\detokenize{xsrst/spell_cmd:spell-cmd-spelling}}\label{\detokenize{xsrst/spell_cmd:index-3}}
The list of words in
{\hyperref[\detokenize{xsrst/xsrst_py:xsrst-py-command-line-arguments-spelling}]{\sphinxcrossref{\DUrole{std,std-ref}{spelling}}}}
are considered correct spellings for all sections.
The latex commands corresponding to the letters in the greek alphabet
are automatically added to this list.

\index{capital@\spxentry{capital}}\index{letters@\spxentry{letters}}\ignorespaces 

\subparagraph{Capital Letters}
\label{\detokenize{xsrst/spell_cmd:capital-letters}}\label{\detokenize{xsrst/spell_cmd:spell-cmd-capital-letters}}\label{\detokenize{xsrst/spell_cmd:index-4}}
The case of the first letter does not matter when checking spelling;
e.g., if \sphinxcode{\sphinxupquote{abcd}} is \sphinxstyleemphasis{word\_1} then \sphinxcode{\sphinxupquote{Abcd}} will be considered a valid word.
Each capital letter starts a new word; e.g., \sphinxtitleref{CamelCase} is considered to
be the two words ‘camel’ and ‘case’.
Single letter words are always correct and not included in the
special word list; e.g., the word list entry \sphinxcode{\sphinxupquote{CppAD}} is the same as \sphinxcode{\sphinxupquote{Cpp}}.

\index{double@\spxentry{double}}\index{words@\spxentry{words}}\ignorespaces 

\subparagraph{Double Words}
\label{\detokenize{xsrst/spell_cmd:double-words}}\label{\detokenize{xsrst/spell_cmd:spell-cmd-double-words}}\label{\detokenize{xsrst/spell_cmd:index-5}}
It is considered an error to have only white space between two occurrences
of the same word. You can make an exception for this by entering
the same word twice (next to each other) in the special word list.

\index{example@\spxentry{example}}\ignorespaces 

\subparagraph{Example}
\label{\detokenize{xsrst/spell_cmd:example}}\label{\detokenize{xsrst/spell_cmd:spell-cmd-example}}\label{\detokenize{xsrst/spell_cmd:index-6}}

\subparagraph{spell\_exam}
\label{\detokenize{xsrst/spell_exam:spell-exam}}\label{\detokenize{xsrst/spell_exam::doc}}
\index{spell\_exam@\spxentry{spell\_exam}}\index{spell@\spxentry{spell}}\index{example@\spxentry{example}}\ignorespaces 

\subparagraph{4.3.1 Spell Example}
\label{\detokenize{xsrst/spell_exam:spell-example}}\label{\detokenize{xsrst/spell_exam:id1}}
\begin{sphinxVerbatim}[commandchars=\\\{\}]
\PYG{l+s+sd}{\PYGZdq{}\PYGZdq{}\PYGZdq{}}
\PYG{l+s+sd}{\PYGZob{}xsrst\PYGZus{}begin spell\PYGZus{}res\PYGZcb{}}
\PYG{l+s+sd}{\PYGZob{}xsrst\PYGZus{}spell}
\PYG{l+s+sd}{    iterable}
\PYG{l+s+sd}{    \PYGZbs{}cos}
\PYG{l+s+sd}{    \PYGZbs{}sin}
\PYG{l+s+sd}{    no no}
\PYG{l+s+sd}{\PYGZcb{}}


\PYG{l+s+sd}{Spell Result}
\PYG{l+s+sd}{\PYGZsh{}\PYGZsh{}\PYGZsh{}\PYGZsh{}\PYGZsh{}\PYGZsh{}\PYGZsh{}\PYGZsh{}\PYGZsh{}\PYGZsh{}\PYGZsh{}\PYGZsh{}}

\PYG{l+s+sd}{Text}
\PYG{l+s+sd}{****}
\PYG{l+s+sd}{The words ``iterable`` and ``xsrst`` are not the dictionary,}
\PYG{l+s+sd}{so we have included them in the spelling command for this section.}


\PYG{l+s+sd}{Math}
\PYG{l+s+sd}{****}
\PYG{l+s+sd}{The latex commands for greek letters}
\PYG{l+s+sd}{are automatically included as correct spelling.}
\PYG{l+s+sd}{The ``\PYGZbs{}rm`` command is included by the}
\PYG{l+s+sd}{:ref:`spelling\PYGZlt{}xsrst\PYGZus{}py.command\PYGZus{}line\PYGZus{}arguments.spelling\PYGZgt{}`.}
\PYG{l+s+sd}{The other latex commands in this section, ``\PYGZbs{}cos`` and ``\PYGZbs{}sin``,}
\PYG{l+s+sd}{have been included in the spelling command for this section.}
\PYG{l+s+sd}{An alternative would be to add them to the *spelling* file}
\PYG{l+s+sd}{which applies to all sections:}

\PYG{l+s+sd}{.. math::}

\PYG{l+s+sd}{    z = \PYGZbs{}cos( \PYGZbs{}theta ) + \PYGZob{}\PYGZbs{}rm i\PYGZcb{} \PYGZbs{}sin( \PYGZbs{}theta )}

\PYG{l+s+sd}{Double Words}
\PYG{l+s+sd}{************}
\PYG{l+s+sd}{It is consider an error to have only white space between}
\PYG{l+s+sd}{two occurrences of the same word; e.g.,}
\PYG{l+s+sd}{no no would be an error if there}
\PYG{l+s+sd}{were not two occurrences of :code:`no` next to each other in the}
\PYG{l+s+sd}{spelling command for this section.}

\PYG{l+s+sd}{:ref:`spell\PYGZus{}exam`}

\PYG{l+s+sd}{\PYGZob{}xsrst\PYGZus{}end spell\PYGZus{}res\PYGZcb{}}
\PYG{l+s+sd}{\PYGZdq{}\PYGZdq{}\PYGZdq{}}
\end{sphinxVerbatim}


\subparagraph{spell\_res}
\label{\detokenize{xsrst/spell_res:spell-res}}\label{\detokenize{xsrst/spell_res::doc}}
\index{spell\_res@\spxentry{spell\_res}}\index{spell@\spxentry{spell}}\index{result@\spxentry{result}}\ignorespaces 

\subparagraph{4.3.1.1 Spell Result}
\label{\detokenize{xsrst/spell_res:spell-result}}\label{\detokenize{xsrst/spell_res:id1}}\begin{itemize}
\item {} 
{\hyperref[\detokenize{xsrst/spell_res:spell-res-text}]{\sphinxcrossref{\DUrole{std,std-ref}{Text}}}}

\item {} 
{\hyperref[\detokenize{xsrst/spell_res:spell-res-math}]{\sphinxcrossref{\DUrole{std,std-ref}{Math}}}}

\item {} 
{\hyperref[\detokenize{xsrst/spell_res:spell-res-double-words}]{\sphinxcrossref{\DUrole{std,std-ref}{Double Words}}}}

\end{itemize}

\index{text@\spxentry{text}}\ignorespaces 

\subparagraph{Text}
\label{\detokenize{xsrst/spell_res:text}}\label{\detokenize{xsrst/spell_res:spell-res-text}}\label{\detokenize{xsrst/spell_res:index-1}}
The words \sphinxcode{\sphinxupquote{iterable}} and \sphinxcode{\sphinxupquote{xsrst}} are not the dictionary,
so we have included them in the spelling command for this section.

\index{math@\spxentry{math}}\ignorespaces 

\subparagraph{Math}
\label{\detokenize{xsrst/spell_res:math}}\label{\detokenize{xsrst/spell_res:spell-res-math}}\label{\detokenize{xsrst/spell_res:index-2}}
The latex commands for greek letters
are automatically included as correct spelling.
The \sphinxcode{\sphinxupquote{\textbackslash{}rm}} command is included by the
{\hyperref[\detokenize{xsrst/xsrst_py:xsrst-py-command-line-arguments-spelling}]{\sphinxcrossref{\DUrole{std,std-ref}{spelling}}}}.
The other latex commands in this section, \sphinxcode{\sphinxupquote{\textbackslash{}cos}} and \sphinxcode{\sphinxupquote{\textbackslash{}sin}},
have been included in the spelling command for this section.
An alternative would be to add them to the \sphinxstyleemphasis{spelling} file
which applies to all sections:
\begin{equation*}
\begin{split}z = \cos( \theta ) + {\rm i} \sin( \theta )\end{split}
\end{equation*}
\index{double@\spxentry{double}}\index{words@\spxentry{words}}\ignorespaces 

\subparagraph{Double Words}
\label{\detokenize{xsrst/spell_res:double-words}}\label{\detokenize{xsrst/spell_res:spell-res-double-words}}\label{\detokenize{xsrst/spell_res:index-3}}
It is consider an error to have only white space between
two occurrences of the same word; e.g.,
no no would be an error if there
were not two occurrences of \sphinxcode{\sphinxupquote{no}} next to each other in the
spelling command for this section.

{\hyperref[\detokenize{xsrst/spell_exam:id1}]{\sphinxcrossref{\DUrole{std,std-ref}{4.3.1 Spell Example}}}}


\bigskip\hrule\bigskip


xsrst input file: \sphinxcode{\sphinxupquote{sphinx/test\_in/spell.py}}


\bigskip\hrule\bigskip


xsrst input file: \sphinxcode{\sphinxupquote{sphinx/test\_in/spell.py}}


\bigskip\hrule\bigskip


xsrst input file: \sphinxcode{\sphinxupquote{bin/xsrst.py}}


\subparagraph{suspend\_cmd}
\label{\detokenize{xsrst/suspend_cmd:suspend-cmd}}\label{\detokenize{xsrst/suspend_cmd::doc}}
\index{suspend\_cmd@\spxentry{suspend\_cmd}}\index{suspend@\spxentry{suspend}}\index{resume@\spxentry{resume}}\index{commands@\spxentry{commands}}\ignorespaces 

\subparagraph{4.4 Suspend and Resume Commands}
\label{\detokenize{xsrst/suspend_cmd:suspend-and-resume-commands}}\label{\detokenize{xsrst/suspend_cmd:id1}}\begin{itemize}
\item {} 
{\hyperref[\detokenize{xsrst/suspend_cmd:suspend-cmd-syntax}]{\sphinxcrossref{\DUrole{std,std-ref}{Syntax}}}}

\item {} 
{\hyperref[\detokenize{xsrst/suspend_cmd:suspend-cmd-purpose}]{\sphinxcrossref{\DUrole{std,std-ref}{Purpose}}}}

\item {} 
{\hyperref[\detokenize{xsrst/suspend_cmd:suspend-cmd-example}]{\sphinxcrossref{\DUrole{std,std-ref}{Example}}}}

\end{itemize}

\index{syntax@\spxentry{syntax}}\ignorespaces 

\subparagraph{Syntax}
\label{\detokenize{xsrst/suspend_cmd:syntax}}\label{\detokenize{xsrst/suspend_cmd:suspend-cmd-syntax}}\label{\detokenize{xsrst/suspend_cmd:index-1}}\begin{itemize}
\item {} 
\sphinxcode{\sphinxupquote{\{xsrst\_suspend\}}}

\item {} 
\sphinxcode{\sphinxupquote{\{xsrst\_resume\}}}

\end{itemize}

\index{purpose@\spxentry{purpose}}\ignorespaces 

\subparagraph{Purpose}
\label{\detokenize{xsrst/suspend_cmd:purpose}}\label{\detokenize{xsrst/suspend_cmd:suspend-cmd-purpose}}\label{\detokenize{xsrst/suspend_cmd:index-2}}
It is possible to suspend (resume) the xsrst extraction during a section.
One begins (ends) the suspension with a suspend command (resume command)
at the
{\hyperref[\detokenize{xsrst/xsrst_py:xsrst-py-notation-beginning-of-a-line}]{\sphinxcrossref{\DUrole{std,std-ref}{beginning of a line}}}}.
Note that this will also suspend all other xsrst processing; e.g.,
spell checking.

\index{example@\spxentry{example}}\ignorespaces 

\subparagraph{Example}
\label{\detokenize{xsrst/suspend_cmd:example}}\label{\detokenize{xsrst/suspend_cmd:suspend-cmd-example}}\label{\detokenize{xsrst/suspend_cmd:index-3}}

\subparagraph{suspend\_exam}
\label{\detokenize{xsrst/suspend_exam:suspend-exam}}\label{\detokenize{xsrst/suspend_exam::doc}}
\index{suspend\_exam@\spxentry{suspend\_exam}}\index{suspend@\spxentry{suspend}}\index{example@\spxentry{example}}\ignorespaces 

\subparagraph{4.4.1 Suspend Example}
\label{\detokenize{xsrst/suspend_exam:suspend-example}}\label{\detokenize{xsrst/suspend_exam:id1}}
\begin{sphinxVerbatim}[commandchars=\\\{\}]
\PYG{l+s+sd}{\PYGZdq{}\PYGZdq{}\PYGZdq{}}
\PYG{l+s+sd}{\PYGZob{}xsrst\PYGZus{}begin suspend\PYGZus{}res\PYGZcb{}}
\PYG{l+s+sd}{\PYGZob{}xsrst\PYGZus{}spell}
\PYG{l+s+sd}{    iterable}
\PYG{l+s+sd}{\PYGZcb{}}

\PYG{l+s+sd}{Suspend Result}
\PYG{l+s+sd}{\PYGZsh{}\PYGZsh{}\PYGZsh{}\PYGZsh{}\PYGZsh{}\PYGZsh{}\PYGZsh{}\PYGZsh{}\PYGZsh{}\PYGZsh{}\PYGZsh{}\PYGZsh{}\PYGZsh{}\PYGZsh{}}

\PYG{l+s+sd}{Factorial}
\PYG{l+s+sd}{*********}
\PYG{l+s+sd}{*f* = ``factorial(`` *positive\PYGZus{}integer* ``)``}
\PYG{l+s+sd}{\PYGZob{}xsrst\PYGZus{}suspend\PYGZcb{}}
\PYG{l+s+sd}{\PYGZdq{}\PYGZdq{}\PYGZdq{}}
\PYG{k}{def} \PYG{n+nf}{factorial}\PYG{p}{(}\PYG{n}{n}\PYG{p}{)} \PYG{p}{:}
    \PYG{k}{if} \PYG{n}{n} \PYG{o}{==} \PYG{l+m+mi}{1} \PYG{p}{:}
        \PYG{k}{return} \PYG{l+m+mi}{1}
    \PYG{k}{return} \PYG{n}{n} \PYG{o}{*} \PYG{n}{factorial}\PYG{p}{(}\PYG{n}{n}\PYG{o}{\PYGZhy{}}\PYG{l+m+mi}{1}\PYG{p}{)}
\PYG{l+s+sd}{\PYGZdq{}\PYGZdq{}\PYGZdq{}}
\PYG{l+s+sd}{\PYGZob{}xsrst\PYGZus{}resume\PYGZcb{}}

\PYG{l+s+sd}{Product}
\PYG{l+s+sd}{*******}
\PYG{l+s+sd}{*p* = ``product(`` *iterable* ``)``}
\PYG{l+s+sd}{\PYGZob{}xsrst\PYGZus{}suspend\PYGZcb{}}
\PYG{l+s+sd}{\PYGZdq{}\PYGZdq{}\PYGZdq{}}
\PYG{k}{def} \PYG{n+nf}{product}\PYG{p}{(}\PYG{n}{itr}\PYG{p}{)} \PYG{p}{:}
    \PYG{n}{p} \PYG{o}{=} \PYG{l+m+mf}{1.0}
    \PYG{k}{for} \PYG{n}{v} \PYG{o+ow}{in} \PYG{n}{itr} \PYG{p}{:}
        \PYG{n}{p} \PYG{o}{*}\PYG{o}{=} \PYG{n}{v}
    \PYG{k}{return} \PYG{n}{p}
\PYG{l+s+sd}{\PYGZdq{}\PYGZdq{}\PYGZdq{}}
\PYG{l+s+sd}{\PYGZob{}xsrst\PYGZus{}resume\PYGZcb{}}

\PYG{l+s+sd}{:ref:`suspend\PYGZus{}exam`}

\PYG{l+s+sd}{\PYGZob{}xsrst\PYGZus{}end suspend\PYGZus{}res\PYGZcb{}}
\PYG{l+s+sd}{\PYGZdq{}\PYGZdq{}\PYGZdq{}}
\end{sphinxVerbatim}


\subparagraph{suspend\_res}
\label{\detokenize{xsrst/suspend_res:suspend-res}}\label{\detokenize{xsrst/suspend_res::doc}}
\index{suspend\_res@\spxentry{suspend\_res}}\index{suspend@\spxentry{suspend}}\index{result@\spxentry{result}}\ignorespaces 

\subparagraph{4.4.1.1 Suspend Result}
\label{\detokenize{xsrst/suspend_res:suspend-result}}\label{\detokenize{xsrst/suspend_res:id1}}\begin{itemize}
\item {} 
{\hyperref[\detokenize{xsrst/suspend_res:suspend-res-factorial}]{\sphinxcrossref{\DUrole{std,std-ref}{Factorial}}}}

\item {} 
{\hyperref[\detokenize{xsrst/suspend_res:suspend-res-product}]{\sphinxcrossref{\DUrole{std,std-ref}{Product}}}}

\end{itemize}

\index{factorial@\spxentry{factorial}}\ignorespaces 

\subparagraph{Factorial}
\label{\detokenize{xsrst/suspend_res:factorial}}\label{\detokenize{xsrst/suspend_res:suspend-res-factorial}}\label{\detokenize{xsrst/suspend_res:index-1}}
\sphinxstyleemphasis{f} = \sphinxcode{\sphinxupquote{factorial(}} \sphinxstyleemphasis{positive\_integer} \sphinxcode{\sphinxupquote{)}}

\index{product@\spxentry{product}}\ignorespaces 

\subparagraph{Product}
\label{\detokenize{xsrst/suspend_res:product}}\label{\detokenize{xsrst/suspend_res:suspend-res-product}}\label{\detokenize{xsrst/suspend_res:index-2}}
\sphinxstyleemphasis{p} = \sphinxcode{\sphinxupquote{product(}} \sphinxstyleemphasis{iterable} \sphinxcode{\sphinxupquote{)}}

{\hyperref[\detokenize{xsrst/suspend_exam:id1}]{\sphinxcrossref{\DUrole{std,std-ref}{4.4.1 Suspend Example}}}}


\bigskip\hrule\bigskip


xsrst input file: \sphinxcode{\sphinxupquote{sphinx/test\_in/suspend.py}}


\bigskip\hrule\bigskip


xsrst input file: \sphinxcode{\sphinxupquote{sphinx/test\_in/suspend.py}}


\bigskip\hrule\bigskip


xsrst input file: \sphinxcode{\sphinxupquote{bin/xsrst.py}}


\subparagraph{code\_cmd}
\label{\detokenize{xsrst/code_cmd:code-cmd}}\label{\detokenize{xsrst/code_cmd::doc}}
\index{code\_cmd@\spxentry{code\_cmd}}\index{code@\spxentry{code}}\index{command@\spxentry{command}}\ignorespaces 

\subparagraph{4.5 Code Command}
\label{\detokenize{xsrst/code_cmd:code-command}}\label{\detokenize{xsrst/code_cmd:id1}}\begin{itemize}
\item {} 
{\hyperref[\detokenize{xsrst/code_cmd:code-cmd-syntax}]{\sphinxcrossref{\DUrole{std,std-ref}{Syntax}}}}

\item {} 
{\hyperref[\detokenize{xsrst/code_cmd:code-cmd-purpose}]{\sphinxcrossref{\DUrole{std,std-ref}{Purpose}}}}

\item {} 
{\hyperref[\detokenize{xsrst/code_cmd:code-cmd-requirements}]{\sphinxcrossref{\DUrole{std,std-ref}{Requirements}}}}

\item {} 
{\hyperref[\detokenize{xsrst/code_cmd:code-cmd-language}]{\sphinxcrossref{\DUrole{std,std-ref}{language}}}}

\item {} 
{\hyperref[\detokenize{xsrst/code_cmd:code-cmd-rest-of-line}]{\sphinxcrossref{\DUrole{std,std-ref}{Rest of Line}}}}

\item {} 
{\hyperref[\detokenize{xsrst/code_cmd:code-cmd-spell-checking}]{\sphinxcrossref{\DUrole{std,std-ref}{Spell Checking}}}}

\item {} 
{\hyperref[\detokenize{xsrst/code_cmd:code-cmd-example}]{\sphinxcrossref{\DUrole{std,std-ref}{Example}}}}

\end{itemize}

\index{syntax@\spxentry{syntax}}\ignorespaces 

\subparagraph{Syntax}
\label{\detokenize{xsrst/code_cmd:syntax}}\label{\detokenize{xsrst/code_cmd:code-cmd-syntax}}\label{\detokenize{xsrst/code_cmd:index-1}}\begin{itemize}
\item {} 
\sphinxcode{\sphinxupquote{\{xsrst\_code}} \sphinxstyleemphasis{language} \sphinxcode{\sphinxupquote{\}}}

\item {} 
\sphinxcode{\sphinxupquote{\{xsrst\_code\}}}

\end{itemize}

\index{purpose@\spxentry{purpose}}\ignorespaces 

\subparagraph{Purpose}
\label{\detokenize{xsrst/code_cmd:purpose}}\label{\detokenize{xsrst/code_cmd:code-cmd-purpose}}\label{\detokenize{xsrst/code_cmd:index-2}}
A code block, directly below in the current input file, begins with
a line containing the first version ( \sphinxstyleemphasis{language} included version)
of the command above.

\index{requirements@\spxentry{requirements}}\ignorespaces 

\subparagraph{Requirements}
\label{\detokenize{xsrst/code_cmd:requirements}}\label{\detokenize{xsrst/code_cmd:code-cmd-requirements}}\label{\detokenize{xsrst/code_cmd:index-3}}
Each code command ends with
a line containing the second version of the command; i.e., \sphinxcode{\sphinxupquote{\{xsrst\_code\}}}.
Hence there must be an even number of code commands.
The back quote character \textasciigrave{} can’t be in the same line as the commands.

\index{language@\spxentry{language}}\ignorespaces 

\subparagraph{language}
\label{\detokenize{xsrst/code_cmd:language}}\label{\detokenize{xsrst/code_cmd:code-cmd-language}}\label{\detokenize{xsrst/code_cmd:index-4}}
A \sphinxstyleemphasis{language} is a non\sphinxhyphen{}empty sequence of non\sphinxhyphen{}space the characters.
It is used to determine the source code language
for highlighting the code block.

\index{rest@\spxentry{rest}}\index{line@\spxentry{line}}\ignorespaces 

\subparagraph{Rest of Line}
\label{\detokenize{xsrst/code_cmd:rest-of-line}}\label{\detokenize{xsrst/code_cmd:code-cmd-rest-of-line}}\label{\detokenize{xsrst/code_cmd:index-5}}
Other characters on the same line as a code command
are not included in the xsrst output.
This enables one to begin or end a comment block
without having the comment characters in the xsrst output.

\index{spell@\spxentry{spell}}\index{checking@\spxentry{checking}}\ignorespaces 

\subparagraph{Spell Checking}
\label{\detokenize{xsrst/code_cmd:spell-checking}}\label{\detokenize{xsrst/code_cmd:code-cmd-spell-checking}}\label{\detokenize{xsrst/code_cmd:index-6}}
Code blocks as usually small and
spell checking is done for these code blocks.
(Spell checking is not done for code blocks included using the
{\hyperref[\detokenize{xsrst/file_cmd:id1}]{\sphinxcrossref{\DUrole{std,std-ref}{file command}}}} .)

\index{example@\spxentry{example}}\ignorespaces 

\subparagraph{Example}
\label{\detokenize{xsrst/code_cmd:example}}\label{\detokenize{xsrst/code_cmd:code-cmd-example}}\label{\detokenize{xsrst/code_cmd:index-7}}

\subparagraph{code\_exam}
\label{\detokenize{xsrst/code_exam:code-exam}}\label{\detokenize{xsrst/code_exam::doc}}
\index{code\_exam@\spxentry{code\_exam}}\index{code@\spxentry{code}}\index{example@\spxentry{example}}\ignorespaces 

\subparagraph{4.5.1 Code Example}
\label{\detokenize{xsrst/code_exam:code-example}}\label{\detokenize{xsrst/code_exam:id1}}
\begin{sphinxVerbatim}[commandchars=\\\{\}]
\PYG{l+s+sd}{\PYGZdq{}\PYGZdq{}\PYGZdq{}}
\PYG{l+s+sd}{\PYGZob{}xsrst\PYGZus{}begin code\PYGZus{}res\PYGZcb{}}

\PYG{l+s+sd}{Code Result}
\PYG{l+s+sd}{\PYGZsh{}\PYGZsh{}\PYGZsh{}\PYGZsh{}\PYGZsh{}\PYGZsh{}\PYGZsh{}\PYGZsh{}\PYGZsh{}\PYGZsh{}\PYGZsh{}}
\PYG{l+s+sd}{\PYGZob{}xsrst\PYGZus{}code py\PYGZcb{}\PYGZdq{}\PYGZdq{}\PYGZdq{}}
\PYG{k}{def} \PYG{n+nf}{factorial}\PYG{p}{(}\PYG{n}{n}\PYG{p}{)} \PYG{p}{:}
    \PYG{k}{if} \PYG{n}{n} \PYG{o}{==} \PYG{l+m+mi}{1} \PYG{p}{:}
        \PYG{k}{return} \PYG{l+m+mi}{1}
    \PYG{k}{return} \PYG{n}{n} \PYG{o}{*} \PYG{n}{factorial}\PYG{p}{(}\PYG{n}{n}\PYG{o}{\PYGZhy{}}\PYG{l+m+mi}{1}\PYG{p}{)}
\PYG{l+s+sd}{\PYGZdq{}\PYGZdq{}\PYGZdq{}\PYGZob{}xsrst\PYGZus{}code\PYGZcb{}}

\PYG{l+s+sd}{:ref:`code\PYGZus{}exam`}

\PYG{l+s+sd}{\PYGZob{}xsrst\PYGZus{}end code\PYGZus{}res\PYGZcb{}}
\PYG{l+s+sd}{\PYGZdq{}\PYGZdq{}\PYGZdq{}}
\end{sphinxVerbatim}


\subparagraph{code\_res}
\label{\detokenize{xsrst/code_res:code-res}}\label{\detokenize{xsrst/code_res::doc}}
\index{code\_res@\spxentry{code\_res}}\index{code@\spxentry{code}}\index{result@\spxentry{result}}\ignorespaces 

\subparagraph{4.5.1.1 Code Result}
\label{\detokenize{xsrst/code_res:code-result}}\label{\detokenize{xsrst/code_res:id1}}
\begin{sphinxVerbatim}[commandchars=\\\{\}]
\PYG{k}{def} \PYG{n+nf}{factorial}\PYG{p}{(}\PYG{n}{n}\PYG{p}{)} \PYG{p}{:}
    \PYG{k}{if} \PYG{n}{n} \PYG{o}{==} \PYG{l+m+mi}{1} \PYG{p}{:}
        \PYG{k}{return} \PYG{l+m+mi}{1}
    \PYG{k}{return} \PYG{n}{n} \PYG{o}{*} \PYG{n}{factorial}\PYG{p}{(}\PYG{n}{n}\PYG{o}{\PYGZhy{}}\PYG{l+m+mi}{1}\PYG{p}{)}
\end{sphinxVerbatim}

{\hyperref[\detokenize{xsrst/code_exam:id1}]{\sphinxcrossref{\DUrole{std,std-ref}{4.5.1 Code Example}}}}


\bigskip\hrule\bigskip


xsrst input file: \sphinxcode{\sphinxupquote{sphinx/test\_in/code.py}}


\bigskip\hrule\bigskip


xsrst input file: \sphinxcode{\sphinxupquote{sphinx/test\_in/code.py}}


\bigskip\hrule\bigskip


xsrst input file: \sphinxcode{\sphinxupquote{bin/xsrst.py}}


\subparagraph{file\_cmd}
\label{\detokenize{xsrst/file_cmd:file-cmd}}\label{\detokenize{xsrst/file_cmd::doc}}
\index{file\_cmd@\spxentry{file\_cmd}}\index{file@\spxentry{file}}\index{command@\spxentry{command}}\ignorespaces 

\subparagraph{4.6 File Command}
\label{\detokenize{xsrst/file_cmd:file-command}}\label{\detokenize{xsrst/file_cmd:id1}}\begin{itemize}
\item {} 
{\hyperref[\detokenize{xsrst/file_cmd:file-cmd-syntax}]{\sphinxcrossref{\DUrole{std,std-ref}{Syntax}}}}

\item {} 
{\hyperref[\detokenize{xsrst/file_cmd:file-cmd-purpose}]{\sphinxcrossref{\DUrole{std,std-ref}{Purpose}}}}

\item {} 
{\hyperref[\detokenize{xsrst/file_cmd:file-cmd-white-space}]{\sphinxcrossref{\DUrole{std,std-ref}{White Space}}}}

\item {} 
{\hyperref[\detokenize{xsrst/file_cmd:file-cmd-file-name}]{\sphinxcrossref{\DUrole{std,std-ref}{file\_name}}}}

\item {} 
{\hyperref[\detokenize{xsrst/file_cmd:file-cmd-start}]{\sphinxcrossref{\DUrole{std,std-ref}{start}}}}

\item {} 
{\hyperref[\detokenize{xsrst/file_cmd:file-cmd-stop}]{\sphinxcrossref{\DUrole{std,std-ref}{stop}}}}

\item {} 
{\hyperref[\detokenize{xsrst/file_cmd:file-cmd-spell-checking}]{\sphinxcrossref{\DUrole{std,std-ref}{Spell Checking}}}}

\item {} 
{\hyperref[\detokenize{xsrst/file_cmd:file-cmd-example}]{\sphinxcrossref{\DUrole{std,std-ref}{Example}}}}

\end{itemize}

\index{syntax@\spxentry{syntax}}\ignorespaces 

\subparagraph{Syntax}
\label{\detokenize{xsrst/file_cmd:syntax}}\label{\detokenize{xsrst/file_cmd:file-cmd-syntax}}\label{\detokenize{xsrst/file_cmd:index-1}}
\begin{DUlineblock}{0em}
\item[] \sphinxcode{\sphinxupquote{\{xsrst\_file}}
\item[]         \sphinxstyleemphasis{start}
\item[]         \sphinxstyleemphasis{stop}
\item[] \sphinxcode{\sphinxupquote{\}}}
\item[] 
\item[] \sphinxcode{\sphinxupquote{\{xsrst\_file}}
\item[]         \sphinxstyleemphasis{start}
\item[]         \sphinxstyleemphasis{stop}
\item[]         \sphinxstyleemphasis{file\_name}
\item[] \sphinxcode{\sphinxupquote{\}}}
\end{DUlineblock}

\index{purpose@\spxentry{purpose}}\ignorespaces 

\subparagraph{Purpose}
\label{\detokenize{xsrst/file_cmd:purpose}}\label{\detokenize{xsrst/file_cmd:file-cmd-purpose}}\label{\detokenize{xsrst/file_cmd:index-2}}
A code block, from any where in any file,
is included by the command above at the
{\hyperref[\detokenize{xsrst/xsrst_py:xsrst-py-notation-beginning-of-a-line}]{\sphinxcrossref{\DUrole{std,std-ref}{beginning of a line}}}}.

\index{white@\spxentry{white}}\index{space@\spxentry{space}}\ignorespaces 

\subparagraph{White Space}
\label{\detokenize{xsrst/file_cmd:white-space}}\label{\detokenize{xsrst/file_cmd:file-cmd-white-space}}\label{\detokenize{xsrst/file_cmd:index-3}}
Leading and trailing white space is not included in
\sphinxstyleemphasis{start}, \sphinxstyleemphasis{stop} or \sphinxstyleemphasis{file\_name}.
The new line character separates these tokens.

\index{file\_name@\spxentry{file\_name}}\ignorespaces 

\subparagraph{file\_name}
\label{\detokenize{xsrst/file_cmd:file-name}}\label{\detokenize{xsrst/file_cmd:file-cmd-file-name}}\label{\detokenize{xsrst/file_cmd:index-4}}
If \sphinxstyleemphasis{file\_name} is not in the syntax,
the code block is in the current input file.
Otherwise, the code block is in \sphinxstyleemphasis{file\_name}.
This file name is relative to the directory where \sphinxcode{\sphinxupquote{xsrst.py}}
is executed; i.e., the top directory for this git repository.
This may seem verbose, but it makes it easier to write scripts
that move files and automatically change references to them.

\index{start@\spxentry{start}}\ignorespaces 

\subparagraph{start}
\label{\detokenize{xsrst/file_cmd:start}}\label{\detokenize{xsrst/file_cmd:file-cmd-start}}\label{\detokenize{xsrst/file_cmd:index-5}}
The code block starts with the line following the occurence
of the text \sphinxstyleemphasis{start} in \sphinxstyleemphasis{file\_name}.
If this is the same as the file containing the command,
the text \sphinxstyleemphasis{start} will not match any text in the command.
There must be one and only one occurence of \sphinxstyleemphasis{start} in \sphinxstyleemphasis{file\_name},
not counting the command itself when the files are the same.

\index{stop@\spxentry{stop}}\ignorespaces 

\subparagraph{stop}
\label{\detokenize{xsrst/file_cmd:stop}}\label{\detokenize{xsrst/file_cmd:file-cmd-stop}}\label{\detokenize{xsrst/file_cmd:index-6}}
The code block ends with the line before the occurence
of the text \sphinxstyleemphasis{start} in \sphinxstyleemphasis{file\_name}.
If this is the same as the file containing the command,
the text \sphinxstyleemphasis{stop} will not match any text in the command.
There must be one and only one occurence of \sphinxstyleemphasis{stop} in \sphinxstyleemphasis{file\_name},
not counting the command itself when the files are the same.

\index{spell@\spxentry{spell}}\index{checking@\spxentry{checking}}\ignorespaces 

\subparagraph{Spell Checking}
\label{\detokenize{xsrst/file_cmd:spell-checking}}\label{\detokenize{xsrst/file_cmd:file-cmd-spell-checking}}\label{\detokenize{xsrst/file_cmd:index-7}}
Spell checking is \sphinxstylestrong{not} done for these code blocks.

\index{example@\spxentry{example}}\ignorespaces 

\subparagraph{Example}
\label{\detokenize{xsrst/file_cmd:example}}\label{\detokenize{xsrst/file_cmd:file-cmd-example}}\label{\detokenize{xsrst/file_cmd:index-8}}

\subparagraph{file\_exam}
\label{\detokenize{xsrst/file_exam:file-exam}}\label{\detokenize{xsrst/file_exam::doc}}
\index{file\_exam@\spxentry{file\_exam}}\index{file@\spxentry{file}}\index{example@\spxentry{example}}\ignorespaces 

\subparagraph{4.6.1 File Example}
\label{\detokenize{xsrst/file_exam:file-example}}\label{\detokenize{xsrst/file_exam:id1}}
\begin{sphinxVerbatim}[commandchars=\\\{\}]
\PYG{c+c1}{// BEGIN\PYGZus{}FACTORIAL}
\PYG{k}{template}\PYG{o}{\PYGZlt{}}\PYG{k}{class} \PYG{n+nc}{T}\PYG{o}{\PYGZgt{}} \PYG{n}{factorial}\PYG{p}{(}\PYG{k}{const} \PYG{n}{T}\PYG{o}{\PYGZam{}} \PYG{n}{n}\PYG{p}{)}
\PYG{p}{\PYGZob{}}   \PYG{k}{if} \PYG{n}{n} \PYG{o}{=}\PYG{o}{=} \PYG{k}{static\PYGZus{}cast}\PYG{o}{\PYGZlt{}}\PYG{n}{T}\PYG{o}{\PYGZgt{}}\PYG{p}{(}\PYG{l+m+mi}{1}\PYG{p}{)}
        \PYG{k}{return} \PYG{n}{n}\PYG{p}{;}
    \PYG{k}{return} \PYG{n}{n} \PYG{o}{*} \PYG{n+nf}{factorial}\PYG{p}{(}\PYG{n}{n} \PYG{o}{\PYGZhy{}} \PYG{l+m+mi}{1}\PYG{p}{)}\PYG{p}{;}
\PYG{p}{\PYGZcb{}}
\PYG{c+c1}{// END\PYGZus{}FACTORIAL}
\PYG{c+c1}{//}
\PYG{c+c1}{// BEGIN\PYGZus{}SQUARE}
\PYG{k}{template}\PYG{o}{\PYGZlt{}}\PYG{k}{class} \PYG{n+nc}{T}\PYG{o}{\PYGZgt{}} \PYG{n}{square}\PYG{p}{(}\PYG{k}{const} \PYG{n}{T}\PYG{o}{\PYGZam{}} \PYG{n}{x}\PYG{p}{)}
\PYG{c+c1}{// END\PYGZus{}SQUARE}
\PYG{p}{\PYGZob{}}   \PYG{k}{return} \PYG{n}{x} \PYG{o}{*} \PYG{n}{x}\PYG{p}{;}
\PYG{p}{\PYGZcb{}}
\PYG{c+cm}{/*}
\PYG{c+cm}{\PYGZhy{}\PYGZhy{}\PYGZhy{}\PYGZhy{}\PYGZhy{}\PYGZhy{}\PYGZhy{}\PYGZhy{}\PYGZhy{}\PYGZhy{}\PYGZhy{}\PYGZhy{}\PYGZhy{}\PYGZhy{}\PYGZhy{}\PYGZhy{}\PYGZhy{}\PYGZhy{}\PYGZhy{}\PYGZhy{}\PYGZhy{}\PYGZhy{}\PYGZhy{}\PYGZhy{}\PYGZhy{}\PYGZhy{}\PYGZhy{}\PYGZhy{}\PYGZhy{}\PYGZhy{}\PYGZhy{}\PYGZhy{}\PYGZhy{}\PYGZhy{}\PYGZhy{}\PYGZhy{}\PYGZhy{}\PYGZhy{}\PYGZhy{}\PYGZhy{}\PYGZhy{}\PYGZhy{}\PYGZhy{}\PYGZhy{}\PYGZhy{}\PYGZhy{}\PYGZhy{}\PYGZhy{}\PYGZhy{}\PYGZhy{}\PYGZhy{}\PYGZhy{}\PYGZhy{}\PYGZhy{}\PYGZhy{}\PYGZhy{}\PYGZhy{}\PYGZhy{}\PYGZhy{}\PYGZhy{}\PYGZhy{}\PYGZhy{}\PYGZhy{}\PYGZhy{}\PYGZhy{}\PYGZhy{}\PYGZhy{}\PYGZhy{}\PYGZhy{}\PYGZhy{}\PYGZhy{}\PYGZhy{}\PYGZhy{}\PYGZhy{}\PYGZhy{}\PYGZhy{}\PYGZhy{}\PYGZhy{}}
\PYG{c+cm}{\PYGZob{}xsrst\PYGZus{}begin file\PYGZus{}res\PYGZcb{}}

\PYG{c+cm}{File Result}
\PYG{c+cm}{\PYGZsh{}\PYGZsh{}\PYGZsh{}\PYGZsh{}\PYGZsh{}\PYGZsh{}\PYGZsh{}\PYGZsh{}\PYGZsh{}\PYGZsh{}\PYGZsh{}}

\PYG{c+cm}{factorial}
\PYG{c+cm}{*********}
\PYG{c+cm}{\PYGZob{}xsrst\PYGZus{}file}
\PYG{c+cm}{    // BEGIN\PYGZus{}FACTORIAL}
\PYG{c+cm}{    // END\PYGZus{}FACTORIAL}
\PYG{c+cm}{\PYGZcb{}}

\PYG{c+cm}{square}
\PYG{c+cm}{******}
\PYG{c+cm}{\PYGZob{}xsrst\PYGZus{}file}
\PYG{c+cm}{    // BEGIN\PYGZus{}SQUARE}
\PYG{c+cm}{    // END\PYGZus{}SQUARE}
\PYG{c+cm}{\PYGZcb{}}

\PYG{c+cm}{\PYGZob{}xsrst\PYGZus{}end file\PYGZus{}res\PYGZcb{}}
\PYG{c+cm}{\PYGZhy{}\PYGZhy{}\PYGZhy{}\PYGZhy{}\PYGZhy{}\PYGZhy{}\PYGZhy{}\PYGZhy{}\PYGZhy{}\PYGZhy{}\PYGZhy{}\PYGZhy{}\PYGZhy{}\PYGZhy{}\PYGZhy{}\PYGZhy{}\PYGZhy{}\PYGZhy{}\PYGZhy{}\PYGZhy{}\PYGZhy{}\PYGZhy{}\PYGZhy{}\PYGZhy{}\PYGZhy{}\PYGZhy{}\PYGZhy{}\PYGZhy{}\PYGZhy{}\PYGZhy{}\PYGZhy{}\PYGZhy{}\PYGZhy{}\PYGZhy{}\PYGZhy{}\PYGZhy{}\PYGZhy{}\PYGZhy{}\PYGZhy{}\PYGZhy{}\PYGZhy{}\PYGZhy{}\PYGZhy{}\PYGZhy{}\PYGZhy{}\PYGZhy{}\PYGZhy{}\PYGZhy{}\PYGZhy{}\PYGZhy{}\PYGZhy{}\PYGZhy{}\PYGZhy{}\PYGZhy{}\PYGZhy{}\PYGZhy{}\PYGZhy{}\PYGZhy{}\PYGZhy{}\PYGZhy{}\PYGZhy{}\PYGZhy{}\PYGZhy{}\PYGZhy{}\PYGZhy{}\PYGZhy{}\PYGZhy{}\PYGZhy{}\PYGZhy{}\PYGZhy{}\PYGZhy{}\PYGZhy{}\PYGZhy{}\PYGZhy{}\PYGZhy{}\PYGZhy{}\PYGZhy{}\PYGZhy{}}
\PYG{c+cm}{*/}
\end{sphinxVerbatim}


\subparagraph{file\_res}
\label{\detokenize{xsrst/file_res:file-res}}\label{\detokenize{xsrst/file_res::doc}}
\index{file\_res@\spxentry{file\_res}}\index{file@\spxentry{file}}\index{result@\spxentry{result}}\ignorespaces 

\subparagraph{4.6.1.1 File Result}
\label{\detokenize{xsrst/file_res:file-result}}\label{\detokenize{xsrst/file_res:id1}}\begin{itemize}
\item {} 
{\hyperref[\detokenize{xsrst/file_res:file-res-factorial}]{\sphinxcrossref{\DUrole{std,std-ref}{factorial}}}}

\item {} 
{\hyperref[\detokenize{xsrst/file_res:file-res-square}]{\sphinxcrossref{\DUrole{std,std-ref}{square}}}}

\end{itemize}

\index{factorial@\spxentry{factorial}}\ignorespaces 

\subparagraph{factorial}
\label{\detokenize{xsrst/file_res:factorial}}\label{\detokenize{xsrst/file_res:file-res-factorial}}\label{\detokenize{xsrst/file_res:index-1}}
\begin{sphinxVerbatim}[commandchars=\\\{\}]
\PYG{k}{template}\PYG{o}{\PYGZlt{}}\PYG{k}{class} \PYG{n+nc}{T}\PYG{o}{\PYGZgt{}} \PYG{n}{factorial}\PYG{p}{(}\PYG{k}{const} \PYG{n}{T}\PYG{o}{\PYGZam{}} \PYG{n}{n}\PYG{p}{)}
\PYG{p}{\PYGZob{}}   \PYG{k}{if} \PYG{n}{n} \PYG{o}{=}\PYG{o}{=} \PYG{k}{static\PYGZus{}cast}\PYG{o}{\PYGZlt{}}\PYG{n}{T}\PYG{o}{\PYGZgt{}}\PYG{p}{(}\PYG{l+m+mi}{1}\PYG{p}{)}
        \PYG{k}{return} \PYG{n}{n}\PYG{p}{;}
    \PYG{k}{return} \PYG{n}{n} \PYG{o}{*} \PYG{n+nf}{factorial}\PYG{p}{(}\PYG{n}{n} \PYG{o}{\PYGZhy{}} \PYG{l+m+mi}{1}\PYG{p}{)}\PYG{p}{;}
\PYG{p}{\PYGZcb{}}
\end{sphinxVerbatim}

\index{square@\spxentry{square}}\ignorespaces 

\subparagraph{square}
\label{\detokenize{xsrst/file_res:square}}\label{\detokenize{xsrst/file_res:file-res-square}}\label{\detokenize{xsrst/file_res:index-2}}
\begin{sphinxVerbatim}[commandchars=\\\{\}]
\PYG{k}{template}\PYG{o}{\PYGZlt{}}\PYG{k}{class} \PYG{n+nc}{T}\PYG{o}{\PYGZgt{}} \PYG{n}{square}\PYG{p}{(}\PYG{k}{const} \PYG{n}{T}\PYG{o}{\PYGZam{}} \PYG{n}{x}\PYG{p}{)}
\end{sphinxVerbatim}


\bigskip\hrule\bigskip


xsrst input file: \sphinxcode{\sphinxupquote{sphinx/test\_in/file.cpp}}


\bigskip\hrule\bigskip


xsrst input file: \sphinxcode{\sphinxupquote{sphinx/test\_in/file.cpp}}


\bigskip\hrule\bigskip


xsrst input file: \sphinxcode{\sphinxupquote{bin/xsrst.py}}


\subparagraph{comment\_ch\_cmd}
\label{\detokenize{xsrst/comment_ch_cmd:comment-ch-cmd}}\label{\detokenize{xsrst/comment_ch_cmd::doc}}
\index{comment\_ch\_cmd@\spxentry{comment\_ch\_cmd}}\index{comment@\spxentry{comment}}\index{character@\spxentry{character}}\index{command@\spxentry{command}}\ignorespaces 

\subparagraph{4.7 Comment Character Command}
\label{\detokenize{xsrst/comment_ch_cmd:comment-character-command}}\label{\detokenize{xsrst/comment_ch_cmd:id1}}\begin{itemize}
\item {} 
{\hyperref[\detokenize{xsrst/comment_ch_cmd:comment-ch-cmd-syntax}]{\sphinxcrossref{\DUrole{std,std-ref}{Syntax}}}}

\item {} \begin{description}
\item[{{\hyperref[\detokenize{xsrst/comment_ch_cmd:comment-ch-cmd-purpose}]{\sphinxcrossref{\DUrole{std,std-ref}{Purpose}}}}}] \leavevmode\begin{itemize}
\item {} 
{\hyperref[\detokenize{xsrst/comment_ch_cmd:comment-ch-cmd-purpose-ch}]{\sphinxcrossref{\DUrole{std,std-ref}{ch}}}}

\end{itemize}

\end{description}

\item {} 
{\hyperref[\detokenize{xsrst/comment_ch_cmd:comment-ch-cmd-beginning-of-a-line}]{\sphinxcrossref{\DUrole{std,std-ref}{Beginning of a Line}}}}

\item {} 
{\hyperref[\detokenize{xsrst/comment_ch_cmd:comment-ch-cmd-input-stream}]{\sphinxcrossref{\DUrole{std,std-ref}{Input Stream}}}}

\item {} 
{\hyperref[\detokenize{xsrst/comment_ch_cmd:comment-ch-cmd-example}]{\sphinxcrossref{\DUrole{std,std-ref}{Example}}}}

\end{itemize}

\index{syntax@\spxentry{syntax}}\ignorespaces 

\subparagraph{Syntax}
\label{\detokenize{xsrst/comment_ch_cmd:syntax}}\label{\detokenize{xsrst/comment_ch_cmd:comment-ch-cmd-syntax}}\label{\detokenize{xsrst/comment_ch_cmd:index-1}}
\sphinxcode{\sphinxupquote{\{xsrst\_comment\_ch}} \sphinxstyleemphasis{ch} \sphinxcode{\sphinxupquote{\}}}

\index{purpose@\spxentry{purpose}}\ignorespaces 

\subparagraph{Purpose}
\label{\detokenize{xsrst/comment_ch_cmd:purpose}}\label{\detokenize{xsrst/comment_ch_cmd:comment-ch-cmd-purpose}}\label{\detokenize{xsrst/comment_ch_cmd:index-2}}
Some languages have a special character that
indicates the rest of the line is a comment.
If you embed sphinx documentation in this type of comment,
you need to inform xsrst of the special character so it does
not end up in your \sphinxcode{\sphinxupquote{.rst}} output file.

\index{ch@\spxentry{ch}}\ignorespaces 

\subparagraph{ch}
\label{\detokenize{xsrst/comment_ch_cmd:ch}}\label{\detokenize{xsrst/comment_ch_cmd:comment-ch-cmd-purpose-ch}}\label{\detokenize{xsrst/comment_ch_cmd:index-3}}
The value of \sphinxstyleemphasis{ch} must be one non white space character.
There must be at least one white space character
between \sphinxcode{\sphinxupquote{xsrst\_comment\_ch}} and \sphinxstyleemphasis{ch}.
Leading and trailing white space around \sphinxstyleemphasis{ch} is ignored.
There can be only one occurence of this command within a file,
it’s effect lasts for the entire file, and
it must come before the first {\hyperref[\detokenize{xsrst/begin_cmd:id1}]{\sphinxcrossref{\DUrole{std,std-ref}{4.1 Begin and End Commands}}}} in the file.

\index{beginning@\spxentry{beginning}}\index{line@\spxentry{line}}\ignorespaces 

\subparagraph{Beginning of a Line}
\label{\detokenize{xsrst/comment_ch_cmd:beginning-of-a-line}}\label{\detokenize{xsrst/comment_ch_cmd:comment-ch-cmd-beginning-of-a-line}}\label{\detokenize{xsrst/comment_ch_cmd:index-4}}
A sequence of characters \sphinxstyleemphasis{text} is at the beginning of a line if there
are only space characters
between the previous new line character and \sphinxstyleemphasis{text}.
In addition, the special character \sphinxstyleemphasis{ch} can be the first character
after the new line and before \sphinxstyleemphasis{text}.

\index{input@\spxentry{input}}\index{stream@\spxentry{stream}}\ignorespaces 

\subparagraph{Input Stream}
\label{\detokenize{xsrst/comment_ch_cmd:input-stream}}\label{\detokenize{xsrst/comment_ch_cmd:comment-ch-cmd-input-stream}}\label{\detokenize{xsrst/comment_ch_cmd:index-5}}
The special character (and one space if present directly after)
is removed from the input stream before any xsrst processing; e.g.,
calculating the amount of
{\hyperref[\detokenize{xsrst/xsrst_py:xsrst-py-indentation}]{\sphinxcrossref{\DUrole{std,std-ref}{Indentation}}}} for the current section.
For example, if \sphinxcode{\sphinxupquote{\#}} is the special character,
the following input has the heading Factorial
and the \sphinxcode{\sphinxupquote{def}} token indented the same amount:

\begin{sphinxVerbatim}[commandchars=\\\{\}]
\PYG{c+c1}{\PYGZsh{} Factorial}
\PYG{c+c1}{\PYGZsh{} \PYGZhy{}\PYGZhy{}\PYGZhy{}\PYGZhy{}\PYGZhy{}\PYGZhy{}\PYGZhy{}\PYGZhy{}\PYGZhy{}}
\PYG{k}{def} \PYG{n+nf}{factorial}\PYG{p}{(}\PYG{n}{n}\PYG{p}{)} \PYG{p}{:}
    \PYG{k}{if} \PYG{n}{n} \PYG{o}{==} \PYG{l+m+mi}{1} \PYG{p}{:}
        \PYG{k}{return} \PYG{l+m+mi}{1}
    \PYG{k}{return} \PYG{n}{n} \PYG{o}{*} \PYG{n}{factorial}\PYG{p}{(}\PYG{n}{n}\PYG{o}{\PYGZhy{}}\PYG{l+m+mi}{1}\PYG{p}{)}
\end{sphinxVerbatim}

\index{example@\spxentry{example}}\ignorespaces 

\subparagraph{Example}
\label{\detokenize{xsrst/comment_ch_cmd:example}}\label{\detokenize{xsrst/comment_ch_cmd:comment-ch-cmd-example}}\label{\detokenize{xsrst/comment_ch_cmd:index-6}}

\subparagraph{comment\_ch\_exam}
\label{\detokenize{xsrst/comment_ch_exam:comment-ch-exam}}\label{\detokenize{xsrst/comment_ch_exam::doc}}
\index{comment\_ch\_exam@\spxentry{comment\_ch\_exam}}\index{comment@\spxentry{comment}}\index{character@\spxentry{character}}\index{example@\spxentry{example}}\ignorespaces 

\subparagraph{4.7.1 Comment Character Example}
\label{\detokenize{xsrst/comment_ch_exam:comment-character-example}}\label{\detokenize{xsrst/comment_ch_exam:id1}}
\begin{sphinxVerbatim}[commandchars=\\\{\}]
\PYG{c+c1}{\PYGZsh{} \PYGZob{}xsrst\PYGZus{}begin comment\PYGZus{}ch\PYGZus{}res\PYGZcb{}}
\PYG{c+c1}{\PYGZsh{}}
\PYG{c+c1}{\PYGZsh{} ========================}
\PYG{c+c1}{\PYGZsh{} Comment Character Result}
\PYG{c+c1}{\PYGZsh{} ========================}
\PYG{c+c1}{\PYGZsh{} In the original source, the leading characters \PYGZsq{}\PYGZsh{}\PYGZsq{} and \PYGZsq{} \PYGZsq{} get removed}
\PYG{c+c1}{\PYGZsh{} and the remaining text lines up with the ``def`` below:}
\PYG{c+c1}{\PYGZsh{}}
\PYG{c+c1}{\PYGZsh{} \PYGZob{}xsrst\PYGZus{}code py\PYGZcb{}}
\PYG{k}{def} \PYG{n+nf}{factorial}\PYG{p}{(}\PYG{n}{n}\PYG{p}{)} \PYG{p}{:}
    \PYG{k}{if} \PYG{n}{n} \PYG{o}{==} \PYG{l+m+mi}{1} \PYG{p}{:}
        \PYG{k}{return} \PYG{l+m+mi}{1}
    \PYG{k}{return} \PYG{n}{n} \PYG{o}{*} \PYG{n}{factorial}\PYG{p}{(}\PYG{n}{n}\PYG{o}{\PYGZhy{}}\PYG{l+m+mi}{1}\PYG{p}{)}
\PYG{c+c1}{\PYGZsh{} \PYGZob{}xsrst\PYGZus{}code\PYGZcb{}}
\PYG{c+c1}{\PYGZsh{}}
\PYG{c+c1}{\PYGZsh{} :ref:`comment\PYGZus{}ch\PYGZus{}exam`}
\PYG{c+c1}{\PYGZsh{}}
\PYG{c+c1}{\PYGZsh{} \PYGZob{}xsrst\PYGZus{}end comment\PYGZus{}ch\PYGZus{}res\PYGZcb{}}
\end{sphinxVerbatim}


\subparagraph{comment\_ch\_res}
\label{\detokenize{xsrst/comment_ch_res:comment-ch-res}}\label{\detokenize{xsrst/comment_ch_res::doc}}
\index{comment\_ch\_res@\spxentry{comment\_ch\_res}}\index{comment@\spxentry{comment}}\index{character@\spxentry{character}}\index{result@\spxentry{result}}\ignorespaces 

\subparagraph{4.7.1.1 Comment Character Result}
\label{\detokenize{xsrst/comment_ch_res:comment-character-result}}\label{\detokenize{xsrst/comment_ch_res:id1}}
In the original source, the leading characters ‘\#’ and ‘ ‘ get removed
and the remaining text lines up with the \sphinxcode{\sphinxupquote{def}} below:

\begin{sphinxVerbatim}[commandchars=\\\{\}]
\PYG{k}{def} \PYG{n+nf}{factorial}\PYG{p}{(}\PYG{n}{n}\PYG{p}{)} \PYG{p}{:}
    \PYG{k}{if} \PYG{n}{n} \PYG{o}{==} \PYG{l+m+mi}{1} \PYG{p}{:}
        \PYG{k}{return} \PYG{l+m+mi}{1}
    \PYG{k}{return} \PYG{n}{n} \PYG{o}{*} \PYG{n}{factorial}\PYG{p}{(}\PYG{n}{n}\PYG{o}{\PYGZhy{}}\PYG{l+m+mi}{1}\PYG{p}{)}
\end{sphinxVerbatim}

{\hyperref[\detokenize{xsrst/comment_ch_exam:id1}]{\sphinxcrossref{\DUrole{std,std-ref}{4.7.1 Comment Character Example}}}}


\bigskip\hrule\bigskip


xsrst input file: \sphinxcode{\sphinxupquote{sphinx/test\_in/comment\_ch.py}}


\bigskip\hrule\bigskip


xsrst input file: \sphinxcode{\sphinxupquote{sphinx/test\_in/comment\_ch.py}}


\bigskip\hrule\bigskip


xsrst input file: \sphinxcode{\sphinxupquote{bin/xsrst.py}}


\subparagraph{heading\_exam}
\label{\detokenize{xsrst/heading_exam:heading-exam}}\label{\detokenize{xsrst/heading_exam::doc}}
\index{heading\_exam@\spxentry{heading\_exam}}\index{heading@\spxentry{heading}}\index{example@\spxentry{example}}\ignorespaces 

\subparagraph{4.8 Heading Example}
\label{\detokenize{xsrst/heading_exam:heading-example}}\label{\detokenize{xsrst/heading_exam:id1}}\begin{itemize}
\item {} 
{\hyperref[\detokenize{xsrst/heading_exam:heading-exam-child-sections}]{\sphinxcrossref{\DUrole{std,std-ref}{Child Sections}}}}

\end{itemize}

\begin{sphinxVerbatim}[commandchars=\\\{\}]
\PYG{l+s+sd}{\PYGZdq{}\PYGZdq{}\PYGZdq{}}
\PYG{l+s+sd}{\PYGZob{}xsrst\PYGZus{}begin heading\PYGZus{}res\PYGZcb{}}

\PYG{l+s+sd}{Heading Result}
\PYG{l+s+sd}{\PYGZsh{}\PYGZsh{}\PYGZsh{}\PYGZsh{}\PYGZsh{}\PYGZsh{}\PYGZsh{}\PYGZsh{}\PYGZsh{}\PYGZsh{}\PYGZsh{}\PYGZsh{}\PYGZsh{}\PYGZsh{}}
\PYG{l+s+sd}{The label for this heading is the section name ``heading\PYGZus{}res``.}

\PYG{l+s+sd}{Second Level}
\PYG{l+s+sd}{************}
\PYG{l+s+sd}{The label for this heading is ``heading\PYGZus{}res.second\PYGZus{}level``.}

\PYG{l+s+sd}{Third Level}
\PYG{l+s+sd}{===========}
\PYG{l+s+sd}{The label for this heading is ``heading\PYGZus{}res.second\PYGZus{}level.third\PYGZus{}level``.}

\PYG{l+s+sd}{Another Second Level}
\PYG{l+s+sd}{********************}
\PYG{l+s+sd}{The label for this heading is ``heading\PYGZus{}res.another\PYGZus{}second\PYGZus{}level``.}

\PYG{l+s+sd}{Third Level}
\PYG{l+s+sd}{===========}
\PYG{l+s+sd}{The label for this heading is}
\PYG{l+s+sd}{``heading\PYGZus{}res.another\PYGZus{}second\PYGZus{}level.third\PYGZus{}level``.}

\PYG{l+s+sd}{Links}
\PYG{l+s+sd}{*****}
\PYG{l+s+sd}{These links would also work from any other section because the}
\PYG{l+s+sd}{:ref:`section\PYGZus{}name\PYGZlt{}begin\PYGZus{}cmd.section\PYGZus{}name\PYGZgt{}`}
\PYG{l+s+sd}{(which is ``heading\PYGZus{}res`` in this case)}
\PYG{l+s+sd}{is included at the beginning of the target for the link:}

\PYG{l+s+sd}{1. :ref:`heading\PYGZus{}res`}
\PYG{l+s+sd}{2. :ref:`heading\PYGZus{}res.second\PYGZus{}level`}
\PYG{l+s+sd}{3. :ref:`heading\PYGZus{}res.second\PYGZus{}level.third\PYGZus{}level`}
\PYG{l+s+sd}{4. :ref:`heading\PYGZus{}res.another\PYGZus{}second\PYGZus{}level`}
\PYG{l+s+sd}{5. :ref:`heading\PYGZus{}res.another\PYGZus{}second\PYGZus{}level.third\PYGZus{}level`}

\PYG{l+s+sd}{:ref:`heading\PYGZus{}exam`}

\PYG{l+s+sd}{\PYGZob{}xsrst\PYGZus{}end heading\PYGZus{}res\PYGZcb{}}
\PYG{l+s+sd}{\PYGZdq{}\PYGZdq{}\PYGZdq{}}
\end{sphinxVerbatim}

\index{child@\spxentry{child}}\index{sections@\spxentry{sections}}\ignorespaces 

\subparagraph{Child Sections}
\label{\detokenize{xsrst/heading_exam:child-sections}}\label{\detokenize{xsrst/heading_exam:heading-exam-child-sections}}\label{\detokenize{xsrst/heading_exam:index-1}}
The heading above (Child Sections) is an example heading for the
{\hyperref[\detokenize{xsrst/xsrst_py:xsrst-py-heading-links-children}]{\sphinxcrossref{\DUrole{std,std-ref}{children}}}}
of a
{\hyperref[\detokenize{xsrst/xsrst_py:xsrst-py-table-of-contents-parent-section}]{\sphinxcrossref{\DUrole{std,std-ref}{parent section}}}}.


\subparagraph{heading\_res}
\label{\detokenize{xsrst/heading_res:heading-res}}\label{\detokenize{xsrst/heading_res::doc}}
\index{heading\_res@\spxentry{heading\_res}}\index{heading@\spxentry{heading}}\index{result@\spxentry{result}}\ignorespaces 

\subparagraph{4.8.1 Heading Result}
\label{\detokenize{xsrst/heading_res:heading-result}}\label{\detokenize{xsrst/heading_res:id1}}\begin{itemize}
\item {} \begin{description}
\item[{{\hyperref[\detokenize{xsrst/heading_res:heading-res-second-level}]{\sphinxcrossref{\DUrole{std,std-ref}{Second Level}}}}}] \leavevmode\begin{itemize}
\item {} 
{\hyperref[\detokenize{xsrst/heading_res:heading-res-second-level-third-level}]{\sphinxcrossref{\DUrole{std,std-ref}{Third Level}}}}

\end{itemize}

\end{description}

\item {} \begin{description}
\item[{{\hyperref[\detokenize{xsrst/heading_res:heading-res-another-second-level}]{\sphinxcrossref{\DUrole{std,std-ref}{Another Second Level}}}}}] \leavevmode\begin{itemize}
\item {} 
{\hyperref[\detokenize{xsrst/heading_res:heading-res-another-second-level-third-level}]{\sphinxcrossref{\DUrole{std,std-ref}{Third Level}}}}

\end{itemize}

\end{description}

\item {} 
{\hyperref[\detokenize{xsrst/heading_res:heading-res-links}]{\sphinxcrossref{\DUrole{std,std-ref}{Links}}}}

\end{itemize}

The label for this heading is the section name \sphinxcode{\sphinxupquote{heading\_res}}.

\index{second@\spxentry{second}}\index{level@\spxentry{level}}\ignorespaces 

\subparagraph{Second Level}
\label{\detokenize{xsrst/heading_res:second-level}}\label{\detokenize{xsrst/heading_res:heading-res-second-level}}\label{\detokenize{xsrst/heading_res:index-1}}
The label for this heading is \sphinxcode{\sphinxupquote{heading\_res.second\_level}}.

\index{third@\spxentry{third}}\index{level@\spxentry{level}}\ignorespaces 

\subparagraph{Third Level}
\label{\detokenize{xsrst/heading_res:third-level}}\label{\detokenize{xsrst/heading_res:heading-res-second-level-third-level}}\label{\detokenize{xsrst/heading_res:index-2}}
The label for this heading is \sphinxcode{\sphinxupquote{heading\_res.second\_level.third\_level}}.

\index{another@\spxentry{another}}\index{second@\spxentry{second}}\index{level@\spxentry{level}}\ignorespaces 

\subparagraph{Another Second Level}
\label{\detokenize{xsrst/heading_res:another-second-level}}\label{\detokenize{xsrst/heading_res:heading-res-another-second-level}}\label{\detokenize{xsrst/heading_res:index-3}}
The label for this heading is \sphinxcode{\sphinxupquote{heading\_res.another\_second\_level}}.

\index{third@\spxentry{third}}\index{level@\spxentry{level}}\ignorespaces 

\subparagraph{Third Level}
\label{\detokenize{xsrst/heading_res:heading-res-another-second-level-third-level}}\label{\detokenize{xsrst/heading_res:index-4}}\label{\detokenize{xsrst/heading_res:id2}}
The label for this heading is
\sphinxcode{\sphinxupquote{heading\_res.another\_second\_level.third\_level}}.

\index{links@\spxentry{links}}\ignorespaces 

\subparagraph{Links}
\label{\detokenize{xsrst/heading_res:links}}\label{\detokenize{xsrst/heading_res:heading-res-links}}\label{\detokenize{xsrst/heading_res:index-5}}
These links would also work from any other section because the
{\hyperref[\detokenize{xsrst/begin_cmd:begin-cmd-section-name}]{\sphinxcrossref{\DUrole{std,std-ref}{section\_name}}}}
(which is \sphinxcode{\sphinxupquote{heading\_res}} in this case)
is included at the beginning of the target for the link:
\begin{enumerate}
\sphinxsetlistlabels{\arabic}{enumi}{enumii}{}{.}%
\item {} 
{\hyperref[\detokenize{xsrst/heading_res:id1}]{\sphinxcrossref{\DUrole{std,std-ref}{4.8.1 Heading Result}}}}

\item {} 
{\hyperref[\detokenize{xsrst/heading_res:heading-res-second-level}]{\sphinxcrossref{\DUrole{std,std-ref}{Second Level}}}}

\item {} 
{\hyperref[\detokenize{xsrst/heading_res:heading-res-second-level-third-level}]{\sphinxcrossref{\DUrole{std,std-ref}{Third Level}}}}

\item {} 
{\hyperref[\detokenize{xsrst/heading_res:heading-res-another-second-level}]{\sphinxcrossref{\DUrole{std,std-ref}{Another Second Level}}}}

\item {} 
{\hyperref[\detokenize{xsrst/heading_res:heading-res-another-second-level-third-level}]{\sphinxcrossref{\DUrole{std,std-ref}{Third Level}}}}

\end{enumerate}

{\hyperref[\detokenize{xsrst/heading_exam:id1}]{\sphinxcrossref{\DUrole{std,std-ref}{4.8 Heading Example}}}}


\bigskip\hrule\bigskip


xsrst input file: \sphinxcode{\sphinxupquote{sphinx/test\_in/heading.py}}


\bigskip\hrule\bigskip


xsrst input file: \sphinxcode{\sphinxupquote{sphinx/test\_in/heading.py}}


\subparagraph{configure}
\label{\detokenize{xsrst/configure:configure}}\label{\detokenize{xsrst/configure::doc}}
\index{configure@\spxentry{configure}}\index{example@\spxentry{example}}\index{xsrst@\spxentry{xsrst}}\index{configuration@\spxentry{configuration}}\index{files@\spxentry{files}}\ignorespaces 

\subparagraph{4.9 Example xsrst Configuration Files}
\label{\detokenize{xsrst/configure:example-xsrst-configuration-files}}\label{\detokenize{xsrst/configure:id1}}

\subparagraph{conf\_py}
\label{\detokenize{xsrst/conf_py:conf-py}}\label{\detokenize{xsrst/conf_py::doc}}
\index{conf\_py@\spxentry{conf\_py}}\index{example@\spxentry{example}}\index{conf.py@\spxentry{conf.py}}\ignorespaces 

\subparagraph{4.9.1 Example conf.py}
\label{\detokenize{xsrst/conf_py:example-conf-py}}\label{\detokenize{xsrst/conf_py:id1}}
\begin{sphinxVerbatim}[commandchars=\\\{\}]
\PYG{c+c1}{\PYGZsh{} Configuration file for the Sphinx documentation builder.}
\PYG{c+c1}{\PYGZsh{}}
\PYG{c+c1}{\PYGZsh{} For conf.py documentation see}
\PYG{c+c1}{\PYGZsh{} http://www.sphinx\PYGZhy{}doc.org/en/master/config}
\PYG{c+c1}{\PYGZsh{}}
\PYG{c+c1}{\PYGZsh{} \PYGZhy{}\PYGZhy{} Path setup \PYGZhy{}\PYGZhy{}\PYGZhy{}\PYGZhy{}\PYGZhy{}\PYGZhy{}\PYGZhy{}\PYGZhy{}\PYGZhy{}\PYGZhy{}\PYGZhy{}\PYGZhy{}\PYGZhy{}\PYGZhy{}\PYGZhy{}\PYGZhy{}\PYGZhy{}\PYGZhy{}\PYGZhy{}\PYGZhy{}\PYGZhy{}\PYGZhy{}\PYGZhy{}\PYGZhy{}\PYGZhy{}\PYGZhy{}\PYGZhy{}\PYGZhy{}\PYGZhy{}\PYGZhy{}\PYGZhy{}\PYGZhy{}\PYGZhy{}\PYGZhy{}\PYGZhy{}\PYGZhy{}\PYGZhy{}\PYGZhy{}\PYGZhy{}\PYGZhy{}\PYGZhy{}\PYGZhy{}\PYGZhy{}\PYGZhy{}\PYGZhy{}\PYGZhy{}\PYGZhy{}\PYGZhy{}\PYGZhy{}\PYGZhy{}\PYGZhy{}\PYGZhy{}\PYGZhy{}\PYGZhy{}\PYGZhy{}\PYGZhy{}\PYGZhy{}\PYGZhy{}\PYGZhy{}\PYGZhy{}\PYGZhy{}\PYGZhy{}}
\PYG{k+kn}{import} \PYG{n+nn}{sphinx\PYGZus{}rtd\PYGZus{}theme}

\PYG{c+c1}{\PYGZsh{} \PYGZhy{}\PYGZhy{} Project information \PYGZhy{}\PYGZhy{}\PYGZhy{}\PYGZhy{}\PYGZhy{}\PYGZhy{}\PYGZhy{}\PYGZhy{}\PYGZhy{}\PYGZhy{}\PYGZhy{}\PYGZhy{}\PYGZhy{}\PYGZhy{}\PYGZhy{}\PYGZhy{}\PYGZhy{}\PYGZhy{}\PYGZhy{}\PYGZhy{}\PYGZhy{}\PYGZhy{}\PYGZhy{}\PYGZhy{}\PYGZhy{}\PYGZhy{}\PYGZhy{}\PYGZhy{}\PYGZhy{}\PYGZhy{}\PYGZhy{}\PYGZhy{}\PYGZhy{}\PYGZhy{}\PYGZhy{}\PYGZhy{}\PYGZhy{}\PYGZhy{}\PYGZhy{}\PYGZhy{}\PYGZhy{}\PYGZhy{}\PYGZhy{}\PYGZhy{}\PYGZhy{}\PYGZhy{}\PYGZhy{}\PYGZhy{}\PYGZhy{}\PYGZhy{}\PYGZhy{}\PYGZhy{}\PYGZhy{}}
\PYG{n}{project}   \PYG{o}{=} \PYG{l+s+s1}{\PYGZsq{}}\PYG{l+s+s1}{cppad\PYGZus{}py}\PYG{l+s+s1}{\PYGZsq{}}
\PYG{n}{copyright} \PYG{o}{=} \PYG{l+s+s1}{\PYGZsq{}}\PYG{l+s+s1}{2020}\PYG{l+s+s1}{\PYGZsq{}}
\PYG{n}{author}    \PYG{o}{=} \PYG{l+s+s1}{\PYGZsq{}}\PYG{l+s+s1}{Brad Bell}\PYG{l+s+s1}{\PYGZsq{}}

\PYG{c+c1}{\PYGZsh{} \PYGZhy{}\PYGZhy{} General configuration \PYGZhy{}\PYGZhy{}\PYGZhy{}\PYGZhy{}\PYGZhy{}\PYGZhy{}\PYGZhy{}\PYGZhy{}\PYGZhy{}\PYGZhy{}\PYGZhy{}\PYGZhy{}\PYGZhy{}\PYGZhy{}\PYGZhy{}\PYGZhy{}\PYGZhy{}\PYGZhy{}\PYGZhy{}\PYGZhy{}\PYGZhy{}\PYGZhy{}\PYGZhy{}\PYGZhy{}\PYGZhy{}\PYGZhy{}\PYGZhy{}\PYGZhy{}\PYGZhy{}\PYGZhy{}\PYGZhy{}\PYGZhy{}\PYGZhy{}\PYGZhy{}\PYGZhy{}\PYGZhy{}\PYGZhy{}\PYGZhy{}\PYGZhy{}\PYGZhy{}\PYGZhy{}\PYGZhy{}\PYGZhy{}\PYGZhy{}\PYGZhy{}\PYGZhy{}\PYGZhy{}\PYGZhy{}\PYGZhy{}\PYGZhy{}\PYGZhy{}}
\PYG{n}{extensions} \PYG{o}{=} \PYG{p}{[}
    \PYG{l+s+s1}{\PYGZsq{}}\PYG{l+s+s1}{sphinx.ext.mathjax}\PYG{l+s+s1}{\PYGZsq{}}\PYG{p}{,}
    \PYG{l+s+s1}{\PYGZsq{}}\PYG{l+s+s1}{sphinx\PYGZus{}rtd\PYGZus{}theme}\PYG{l+s+s1}{\PYGZsq{}}\PYG{p}{,}
\PYG{p}{]}
\PYG{n}{exclude\PYGZus{}patterns} \PYG{o}{=} \PYG{p}{[}
    \PYG{l+s+s1}{\PYGZsq{}}\PYG{l+s+s1}{\PYGZus{}build}\PYG{l+s+s1}{\PYGZsq{}}\PYG{p}{,} \PYG{l+s+s1}{\PYGZsq{}}\PYG{l+s+s1}{Thumbs.db}\PYG{l+s+s1}{\PYGZsq{}}\PYG{p}{,} \PYG{l+s+s1}{\PYGZsq{}}\PYG{l+s+s1}{.DS\PYGZus{}Store}\PYG{l+s+s1}{\PYGZsq{}}\PYG{p}{,} \PYG{l+s+s1}{\PYGZsq{}}\PYG{l+s+s1}{test\PYGZus{}out}\PYG{l+s+s1}{\PYGZsq{}}\PYG{p}{,} \PYG{l+s+s1}{\PYGZsq{}}\PYG{l+s+s1}{preamble.rst}\PYG{l+s+s1}{\PYGZsq{}}
\PYG{p}{]}

\PYG{c+c1}{\PYGZsh{} \PYGZhy{}\PYGZhy{} Options for HTML output \PYGZhy{}\PYGZhy{}\PYGZhy{}\PYGZhy{}\PYGZhy{}\PYGZhy{}\PYGZhy{}\PYGZhy{}\PYGZhy{}\PYGZhy{}\PYGZhy{}\PYGZhy{}\PYGZhy{}\PYGZhy{}\PYGZhy{}\PYGZhy{}\PYGZhy{}\PYGZhy{}\PYGZhy{}\PYGZhy{}\PYGZhy{}\PYGZhy{}\PYGZhy{}\PYGZhy{}\PYGZhy{}\PYGZhy{}\PYGZhy{}\PYGZhy{}\PYGZhy{}\PYGZhy{}\PYGZhy{}\PYGZhy{}\PYGZhy{}\PYGZhy{}\PYGZhy{}\PYGZhy{}\PYGZhy{}\PYGZhy{}\PYGZhy{}\PYGZhy{}\PYGZhy{}\PYGZhy{}\PYGZhy{}\PYGZhy{}\PYGZhy{}\PYGZhy{}\PYGZhy{}\PYGZhy{}\PYGZhy{}}
\PYG{n}{html\PYGZus{}theme} \PYG{o}{=} \PYG{l+s+s1}{\PYGZsq{}}\PYG{l+s+s1}{sphinx\PYGZus{}rtd\PYGZus{}theme}\PYG{l+s+s1}{\PYGZsq{}}
\PYG{n}{html\PYGZus{}theme\PYGZus{}options} \PYG{o}{=} \PYG{p}{\PYGZob{}}
    \PYG{l+s+s1}{\PYGZsq{}}\PYG{l+s+s1}{navigation\PYGZus{}depth}\PYG{l+s+s1}{\PYGZsq{}} \PYG{p}{:} \PYG{o}{\PYGZhy{}}\PYG{l+m+mi}{1}   \PYG{p}{,}
    \PYG{l+s+s1}{\PYGZsq{}}\PYG{l+s+s1}{titles\PYGZus{}only}\PYG{l+s+s1}{\PYGZsq{}}      \PYG{p}{:} \PYG{k+kc}{True} \PYG{p}{,}
\PYG{p}{\PYGZcb{}}

\PYG{c+c1}{\PYGZsh{} \PYGZhy{}\PYGZhy{} These folders are copied to the documentation\PYGZsq{}s HTML output \PYGZhy{}\PYGZhy{}\PYGZhy{}\PYGZhy{}\PYGZhy{}\PYGZhy{}\PYGZhy{}\PYGZhy{}\PYGZhy{}\PYGZhy{}\PYGZhy{}\PYGZhy{}}
\PYG{c+c1}{\PYGZsh{} html\PYGZus{}static\PYGZus{}path = [ \PYGZsq{}\PYGZus{}static\PYGZsq{} ]}

\PYG{c+c1}{\PYGZsh{} \PYGZhy{}\PYGZhy{} These paths are either relative to html\PYGZus{}static\PYGZus{}path \PYGZhy{}\PYGZhy{}\PYGZhy{}\PYGZhy{}\PYGZhy{}\PYGZhy{}\PYGZhy{}\PYGZhy{}\PYGZhy{}\PYGZhy{}\PYGZhy{}\PYGZhy{}\PYGZhy{}\PYGZhy{}\PYGZhy{}\PYGZhy{}\PYGZhy{}\PYGZhy{}\PYGZhy{}\PYGZhy{}}
\PYG{c+c1}{\PYGZsh{} or fully qualified paths (eg. https://...)}
\PYG{n}{html\PYGZus{}css\PYGZus{}files} \PYG{o}{=} \PYG{p}{[}
    \PYG{l+s+s1}{\PYGZsq{}}\PYG{l+s+s1}{css/custom.css}\PYG{l+s+s1}{\PYGZsq{}}\PYG{p}{,}
\PYG{p}{]}

\PYG{c+c1}{\PYGZsh{} \PYGZhy{}\PYGZhy{} Latex commands used by all sections \PYGZhy{}\PYGZhy{}\PYGZhy{}\PYGZhy{}\PYGZhy{}\PYGZhy{}\PYGZhy{}\PYGZhy{}\PYGZhy{}\PYGZhy{}\PYGZhy{}\PYGZhy{}\PYGZhy{}\PYGZhy{}\PYGZhy{}\PYGZhy{}\PYGZhy{}\PYGZhy{}\PYGZhy{}\PYGZhy{}\PYGZhy{}\PYGZhy{}\PYGZhy{}\PYGZhy{}\PYGZhy{}\PYGZhy{}\PYGZhy{}\PYGZhy{}\PYGZhy{}\PYGZhy{}\PYGZhy{}\PYGZhy{}\PYGZhy{}\PYGZhy{}\PYGZhy{}\PYGZhy{}\PYGZhy{}}
\PYG{n}{latex\PYGZus{}elements} \PYG{o}{=} \PYG{p}{\PYGZob{}}
    \PYG{l+s+s1}{\PYGZsq{}}\PYG{l+s+s1}{preamble}\PYG{l+s+s1}{\PYGZsq{}} \PYG{p}{:}  
        \PYG{l+s+sa}{r}\PYG{l+s+s1}{\PYGZsq{}}\PYG{l+s+s1}{\PYGZbs{}}\PYG{l+s+s1}{newcommand}\PYG{l+s+s1}{\PYGZob{}}\PYG{l+s+s1}{\PYGZbs{}}\PYG{l+s+s1}{B\PYGZcb{}[1]}\PYG{l+s+s1}{\PYGZob{}\PYGZob{}}\PYG{l+s+s1}{\PYGZbs{}}\PYG{l+s+s1}{bf \PYGZsh{}1\PYGZcb{}\PYGZcb{} }\PYG{l+s+s1}{\PYGZsq{}}  \PYG{o}{+} \PYG{l+s+s1}{\PYGZsq{}}\PYG{l+s+se}{\PYGZbs{}n}\PYG{l+s+s1}{\PYGZsq{}} \PYG{o}{+}
        \PYG{l+s+sa}{r}\PYG{l+s+s1}{\PYGZsq{}}\PYG{l+s+s1}{\PYGZbs{}}\PYG{l+s+s1}{newcommand}\PYG{l+s+s1}{\PYGZob{}}\PYG{l+s+s1}{\PYGZbs{}}\PYG{l+s+s1}{R\PYGZcb{}[1]}\PYG{l+s+s1}{\PYGZob{}\PYGZob{}}\PYG{l+s+s1}{\PYGZbs{}}\PYG{l+s+s1}{rm \PYGZsh{}1\PYGZcb{}\PYGZcb{} }\PYG{l+s+s1}{\PYGZsq{}}  \PYG{o}{+} \PYG{l+s+s1}{\PYGZsq{}}\PYG{l+s+se}{\PYGZbs{}n}\PYG{l+s+s1}{\PYGZsq{}} \PYG{o}{+}
        \PYG{l+s+sa}{r}\PYG{l+s+s1}{\PYGZsq{}}\PYG{l+s+s1}{\PYGZbs{}}\PYG{l+s+s1}{renewcommand}\PYG{l+s+s1}{\PYGZob{}}\PYG{l+s+s1}{\PYGZbs{}}\PYG{l+s+s1}{thesection\PYGZcb{}}\PYG{l+s+s1}{\PYGZob{}\PYGZob{}}\PYG{l+s+s1}{\PYGZbs{}}\PYG{l+s+s1}{hspace}\PYG{l+s+s1}{\PYGZob{}}\PYG{l+s+s1}{\PYGZhy{}1em\PYGZcb{}\PYGZcb{}\PYGZcb{}}\PYG{l+s+s1}{\PYGZsq{}}        \PYG{o}{+} \PYG{l+s+s1}{\PYGZsq{}}\PYG{l+s+se}{\PYGZbs{}n}\PYG{l+s+s1}{\PYGZsq{}} \PYG{o}{+}
        \PYG{l+s+sa}{r}\PYG{l+s+s1}{\PYGZsq{}}\PYG{l+s+s1}{\PYGZbs{}}\PYG{l+s+s1}{renewcommand}\PYG{l+s+s1}{\PYGZob{}}\PYG{l+s+s1}{\PYGZbs{}}\PYG{l+s+s1}{thesubsection\PYGZcb{}}\PYG{l+s+s1}{\PYGZob{}\PYGZob{}}\PYG{l+s+s1}{\PYGZbs{}}\PYG{l+s+s1}{hspace}\PYG{l+s+s1}{\PYGZob{}}\PYG{l+s+s1}{\PYGZhy{}1em\PYGZcb{}\PYGZcb{}\PYGZcb{}}\PYG{l+s+s1}{\PYGZsq{}}     \PYG{o}{+} \PYG{l+s+s1}{\PYGZsq{}}\PYG{l+s+se}{\PYGZbs{}n}\PYG{l+s+s1}{\PYGZsq{}} \PYG{o}{+}
        \PYG{l+s+sa}{r}\PYG{l+s+s1}{\PYGZsq{}}\PYG{l+s+s1}{\PYGZbs{}}\PYG{l+s+s1}{renewcommand}\PYG{l+s+s1}{\PYGZob{}}\PYG{l+s+s1}{\PYGZbs{}}\PYG{l+s+s1}{thesubsubsection\PYGZcb{}}\PYG{l+s+s1}{\PYGZob{}\PYGZob{}}\PYG{l+s+s1}{\PYGZbs{}}\PYG{l+s+s1}{hspace}\PYG{l+s+s1}{\PYGZob{}}\PYG{l+s+s1}{\PYGZhy{}1em\PYGZcb{}\PYGZcb{}\PYGZcb{}}\PYG{l+s+s1}{\PYGZsq{}}  \PYG{o}{+} \PYG{l+s+s1}{\PYGZsq{}}\PYG{l+s+se}{\PYGZbs{}n}\PYG{l+s+s1}{\PYGZsq{}} \PYG{o}{+}
        \PYG{l+s+s1}{\PYGZsq{}}\PYG{l+s+s1}{\PYGZsq{}}
    \PYG{p}{,}
\PYG{p}{\PYGZcb{}}
\end{sphinxVerbatim}


\bigskip\hrule\bigskip


xsrst input file: \sphinxcode{\sphinxupquote{sphinx/conf.py}}


\subparagraph{index\_rst}
\label{\detokenize{xsrst/index_rst:index-rst}}\label{\detokenize{xsrst/index_rst::doc}}
\index{index\_rst@\spxentry{index\_rst}}\index{example@\spxentry{example}}\index{index.rst@\spxentry{index.rst}}\ignorespaces 

\subparagraph{4.9.2 Example index.rst}
\label{\detokenize{xsrst/index_rst:example-index-rst}}\label{\detokenize{xsrst/index_rst:id1}}
\begin{sphinxVerbatim}[commandchars=\\\{\}]
\PYG{p}{..} \PYG{o+ow}{toctree}\PYG{p}{::}
    xsrst/xsrst\PYGZus{}automatic
    xsrst/cppad\PYGZus{}py
\end{sphinxVerbatim}


\bigskip\hrule\bigskip


xsrst input file: \sphinxcode{\sphinxupquote{sphinx/index.rst}}


\subparagraph{spelling}
\label{\detokenize{xsrst/spelling:spelling}}\label{\detokenize{xsrst/spelling::doc}}
\index{spelling@\spxentry{spelling}}\index{example@\spxentry{example}}\index{spelling@\spxentry{spelling}}\index{file@\spxentry{file}}\ignorespaces 

\subparagraph{4.9.3 Example Spelling File}
\label{\detokenize{xsrst/spelling:example-spelling-file}}\label{\detokenize{xsrst/spelling:id1}}
\begin{sphinxVerbatim}[commandchars=\\\{\}]
\PYG{c+c1}{\PYGZsh{} standard mathematical abbreviations}
\PYG{n}{grad}
\PYG{n}{hess}
\PYG{n}{sinh}
\PYG{c+c1}{\PYGZsh{}}
\PYG{c+c1}{\PYGZsh{} example words special for this projext}
\PYG{n}{ctor}
\PYG{n}{jac}
\PYG{n}{nc}
\PYG{n}{nr}
\PYG{n}{op}
\PYG{n}{jac}
\PYG{n}{rc}
\PYG{n}{var}
\PYG{n}{vec}
\PYG{n}{xam}
\PYG{c+c1}{\PYGZsh{}}
\PYG{c+c1}{\PYGZsh{} example words in xsrst commands}
\PYG{n}{xsrst}
\PYG{n}{rst}
\PYG{c+c1}{\PYGZsh{}}
\PYG{c+c1}{\PYGZsh{} example words in sphinx commands}
\PYG{n}{http}
\PYG{n}{https}
\PYG{n}{www}
\PYG{n}{html}
\PYG{n}{github}
\PYG{n}{csv}
\PYG{c+c1}{\PYGZsh{}}
\PYG{c+c1}{\PYGZsh{} example python commands}
\PYG{n+nb}{dict}
\PYG{n}{dtype}
\PYG{n+nb}{len}
\PYG{n}{numpy}
\PYG{n}{py}
\PYG{n}{scipy}
\PYG{n+nb}{str}
\PYG{n}{sys}
\PYG{c+c1}{\PYGZsh{}}
\PYG{c+c1}{\PYGZsh{} example c++ commands}
\PYG{n}{extern}
\PYG{n}{const}
\PYG{n}{cpp}
\PYG{n}{std}
\PYG{c+c1}{\PYGZsh{}}
\PYG{c+c1}{\PYGZsh{} example latex commands}
\PYGZbs{}\PYG{n}{bar}
\PYGZbs{}\PYG{n}{begin}
\PYGZbs{}\PYG{n}{bf}
\PYGZbs{}\PYG{n}{cdot}
\PYGZbs{}\PYG{n}{circ}
\PYGZbs{}\PYG{n}{dot}
\PYGZbs{}\PYG{n}{ell}
\PYGZbs{}\PYG{n}{end}
\PYGZbs{}\PYG{n}{exp}
\PYGZbs{}\PYG{n}{frac}
\PYGZbs{}\PYG{n}{geq}
\PYGZbs{}\PYG{n+nb}{int}
\PYGZbs{}\PYG{o+ow}{in}
\PYGZbs{}\PYG{n}{ldots}
\PYGZbs{}\PYG{n}{leq}
\PYGZbs{}\PYG{n}{log}
\PYGZbs{}\PYG{n}{left}
\PYGZbs{}\PYG{n}{mbox}
\PYGZbs{}\PYG{n}{newcommand}
\PYGZbs{}\PYG{n}{partial}
\PYGZbs{}\PYG{n}{prod}
\PYGZbs{}\PYG{n}{right}
\PYGZbs{}\PYG{n}{rightarrow}
\PYGZbs{}\PYG{n}{rm}
\PYGZbs{}\PYG{n}{sqrt}
\PYGZbs{}\PYG{n+nb}{sum}
\PYGZbs{}\PYG{n}{times}
\PYGZbs{}\PYG{n}{varepsilon}
\end{sphinxVerbatim}


\bigskip\hrule\bigskip


xsrst input file: \sphinxcode{\sphinxupquote{sphinx/spelling}}


\subparagraph{keyword}
\label{\detokenize{xsrst/keyword:keyword}}\label{\detokenize{xsrst/keyword::doc}}
\index{keyword@\spxentry{keyword}}\index{example@\spxentry{example}}\index{keyword@\spxentry{keyword}}\index{file@\spxentry{file}}\ignorespaces 

\subparagraph{4.9.4 Example Keyword File}
\label{\detokenize{xsrst/keyword:example-keyword-file}}\label{\detokenize{xsrst/keyword:id1}}
\begin{sphinxVerbatim}[commandchars=\\\{\}]
\PYG{c+c1}{\PYGZsh{} List of patterns to exclude from the index}
\PYG{n}{a}
\PYG{o+ow}{and}
\PYG{n}{of}
\PYG{c+c1}{\PYGZsh{} exclude whats new mm\PYGZhy{}dd headings used by whats new sections for each year}
\PYG{p}{[}\PYG{l+m+mi}{0}\PYG{o}{\PYGZhy{}}\PYG{l+m+mi}{9}\PYG{p}{]}\PYG{p}{[}\PYG{l+m+mi}{0}\PYG{o}{\PYGZhy{}}\PYG{l+m+mi}{9}\PYG{p}{]}\PYG{o}{\PYGZhy{}}\PYG{p}{[}\PYG{l+m+mi}{0}\PYG{o}{\PYGZhy{}}\PYG{l+m+mi}{9}\PYG{p}{]}\PYG{p}{[}\PYG{l+m+mi}{0}\PYG{o}{\PYGZhy{}}\PYG{l+m+mi}{9}\PYG{p}{]}
\end{sphinxVerbatim}


\bigskip\hrule\bigskip


xsrst input file: \sphinxcode{\sphinxupquote{sphinx/keyword}}


\bigskip\hrule\bigskip


xsrst input file: \sphinxcode{\sphinxupquote{sphinx/configure.xsrst}}

\index{syntax@\spxentry{syntax}}\ignorespaces 

\subparagraph{Syntax}
\label{\detokenize{xsrst/xsrst_py:syntax}}\label{\detokenize{xsrst/xsrst_py:xsrst-py-syntax}}\label{\detokenize{xsrst/xsrst_py:index-1}}\begin{itemize}
\item {} 
\sphinxcode{\sphinxupquote{xsrst.py}} \sphinxstyleemphasis{sphinx\_dir} \sphinxstyleemphasis{root\_file} \sphinxstyleemphasis{spelling} \sphinxstyleemphasis{keyword}

\item {} 
\sphinxcode{\sphinxupquote{xsrst.py}} \sphinxstyleemphasis{sphinx\_dir} \sphinxstyleemphasis{root\_file} \sphinxstyleemphasis{spelling} \sphinxstyleemphasis{keyword} \sphinxstyleemphasis{line\_increment}

\end{itemize}

\index{purpose@\spxentry{purpose}}\ignorespaces 

\subparagraph{Purpose}
\label{\detokenize{xsrst/xsrst_py:purpose}}\label{\detokenize{xsrst/xsrst_py:xsrst-py-purpose}}\label{\detokenize{xsrst/xsrst_py:index-2}}
This is a pseudo sphinx extension that provides the following features:
\begin{enumerate}
\sphinxsetlistlabels{\arabic}{enumi}{enumii}{}{.}%
\item {} 
The file name for each section is also an abbreviated title used
in the navigation bar and for linking to the section. This makes the
navigation bar more useful while also having long descriptive titles.
It also makes cross reference linking from other sections easier.

\item {} 
Enables documentation in the comments for source code
even when multiple computer languages are used for one package.
Allows the documentation for one section to span multiple locations
in the source code; see {\hyperref[\detokenize{xsrst/suspend_cmd:id1}]{\sphinxcrossref{\DUrole{std,std-ref}{suspend command}}}}.

\item {} 
Allows for multiple sections (rst output files) to be specified by one
input file. In addition, one section can be the parent for the
other sections in a file.

\item {} 
Generates the table of contents from the specification
of which files are included; see {\hyperref[\detokenize{xsrst/child_cmd:id1}]{\sphinxcrossref{\DUrole{std,std-ref}{child commands}}}}.
Generates a jump table to the headings for each section
so that the navigation bar need not include this information.

\item {} 
Includes a configurable {\hyperref[\detokenize{xsrst/spell_cmd:id1}]{\sphinxcrossref{\DUrole{std,std-ref}{spell checker}}}} and
\DUrole{xref,std,std-ref}{index}.
Words in each heading are automatically included in the index.

\item {} 
Makes it easy to include source code, that also executes, from
directly below the {\hyperref[\detokenize{xsrst/code_cmd:id1}]{\sphinxcrossref{\DUrole{std,std-ref}{code command}}}} or from
a different location in a {\hyperref[\detokenize{xsrst/file_cmd:id1}]{\sphinxcrossref{\DUrole{std,std-ref}{file}}}}.
Uses tokens in the source, not line numbers in the source,
to signify start and stop of inclusion from a file.

\end{enumerate}

\index{requirements@\spxentry{requirements}}\ignorespaces 

\subparagraph{Requirements}
\label{\detokenize{xsrst/xsrst_py:requirements}}\label{\detokenize{xsrst/xsrst_py:xsrst-py-requirements}}\label{\detokenize{xsrst/xsrst_py:index-3}}\begin{itemize}
\item {} 
\sphinxcode{\sphinxupquote{pip install \sphinxhyphen{}\sphinxhyphen{}user pyspellchecker}}

\item {} 
\sphinxcode{\sphinxupquote{pip install \sphinxhyphen{}\sphinxhyphen{}user sphinx}}

\item {} 
The directory \sphinxstyleemphasis{cppad\_prefix} \sphinxcode{\sphinxupquote{/bin}} must be in your execution path
where \sphinxstyleemphasis{cppad\_prefix} is the install prefix for cppad.

\end{itemize}

\index{notation@\spxentry{notation}}\ignorespaces 

\subparagraph{Notation}
\label{\detokenize{xsrst/xsrst_py:notation}}\label{\detokenize{xsrst/xsrst_py:xsrst-py-notation}}\label{\detokenize{xsrst/xsrst_py:index-4}}
\index{white@\spxentry{white}}\index{space@\spxentry{space}}\ignorespaces 

\subparagraph{White Space}
\label{\detokenize{xsrst/xsrst_py:white-space}}\label{\detokenize{xsrst/xsrst_py:xsrst-py-notation-white-space}}\label{\detokenize{xsrst/xsrst_py:index-5}}
We define white space to be a sequence of space characters; e.g.,
tabs are not consider white space by xsrst.

\index{beginning@\spxentry{beginning}}\index{line@\spxentry{line}}\ignorespaces 

\subparagraph{Beginning of a Line}
\label{\detokenize{xsrst/xsrst_py:beginning-of-a-line}}\label{\detokenize{xsrst/xsrst_py:xsrst-py-notation-beginning-of-a-line}}\label{\detokenize{xsrst/xsrst_py:index-6}}
We say that a string \sphinxstyleemphasis{text} is a the beginning of a line if
only white space, or nothing, comes before \sphinxstyleemphasis{text} in the line.

\index{command@\spxentry{command}}\index{line@\spxentry{line}}\index{arguments@\spxentry{arguments}}\ignorespaces 

\subparagraph{Command Line Arguments}
\label{\detokenize{xsrst/xsrst_py:command-line-arguments}}\label{\detokenize{xsrst/xsrst_py:xsrst-py-command-line-arguments}}\label{\detokenize{xsrst/xsrst_py:index-7}}
\index{root\_file@\spxentry{root\_file}}\ignorespaces 

\subparagraph{root\_file}
\label{\detokenize{xsrst/xsrst_py:root-file}}\label{\detokenize{xsrst/xsrst_py:xsrst-py-command-line-arguments-root-file}}\label{\detokenize{xsrst/xsrst_py:index-8}}
The command line argument \sphinxstyleemphasis{root\_file} is the name of a file,
relative to the top git repository directory.

\index{root\_section@\spxentry{root\_section}}\ignorespaces 

\subparagraph{root\_section}
\label{\detokenize{xsrst/xsrst_py:root-section}}\label{\detokenize{xsrst/xsrst_py:xsrst-py-command-line-arguments-root-file-root-section}}\label{\detokenize{xsrst/xsrst_py:index-9}}
If there is only one section in the \sphinxstyleemphasis{root\_file} it is called
the \sphinxstyleemphasis{root\_section}; i.e., it is the top section in that table of contents.
If there is more than one section in the \sphinxstyleemphasis{root\_file},
the file must have a
{\hyperref[\detokenize{xsrst/begin_cmd:begin-cmd-parent-section}]{\sphinxcrossref{\DUrole{std,std-ref}{Parent Section}}}} and it is the \sphinxstyleemphasis{root\_section}.
The file \sphinxstyleemphasis{sphinx\_dir} \sphinxcode{\sphinxupquote{/index.rst}} must contain the line

        \sphinxcode{\sphinxupquote{xsrst/}} \sphinxstyleemphasis{section\_name}

where \sphinxstyleemphasis{section\_name} is the name of the \sphinxstyleemphasis{root\_section}.

\index{sphinx\_dir@\spxentry{sphinx\_dir}}\ignorespaces 

\subparagraph{sphinx\_dir}
\label{\detokenize{xsrst/xsrst_py:sphinx-dir}}\label{\detokenize{xsrst/xsrst_py:xsrst-py-command-line-arguments-sphinx-dir}}\label{\detokenize{xsrst/xsrst_py:index-10}}
The command line argument \sphinxstyleemphasis{sphinx\_dir} is a sub\sphinxhyphen{}directory,
of the top git repository directory.
The  sphinx \sphinxcode{\sphinxupquote{conf.py}}, \sphinxcode{\sphinxupquote{index.rst}}, \sphinxstyleemphasis{spelling}, and \sphinxstyleemphasis{keyword}
files are located in this directory.
Any files that have names ending in \sphinxcode{\sphinxupquote{.rst}},
and that are in the directory \sphinxstyleemphasis{sphinx\_dir} \sphinxcode{\sphinxupquote{/xsrst}},
are removed at the beginning of execution of \sphinxcode{\sphinxupquote{xsrst.py}}.
All the \sphinxcode{\sphinxupquote{.rst}} files in \sphinxstyleemphasis{sphinx\_dir} \sphinxcode{\sphinxupquote{/xsrst}}
were extracted from the source code the last time that \sphinxcode{\sphinxupquote{xsrst.py}}
was executed.

\index{example@\spxentry{example}}\index{configuration@\spxentry{configuration}}\index{files@\spxentry{files}}\ignorespaces 

\subparagraph{Example Configuration Files}
\label{\detokenize{xsrst/xsrst_py:example-configuration-files}}\label{\detokenize{xsrst/xsrst_py:xsrst-py-command-line-arguments-sphinx-dir-example-configuration-files}}\label{\detokenize{xsrst/xsrst_py:index-11}}
\begin{DUlineblock}{0em}
\item[]         conf.py: {\hyperref[\detokenize{xsrst/conf_py:id1}]{\sphinxcrossref{\DUrole{std,std-ref}{4.9.1 Example conf.py}}}}
\item[]         index.rst: {\hyperref[\detokenize{xsrst/index_rst:id1}]{\sphinxcrossref{\DUrole{std,std-ref}{4.9.2 Example index.rst}}}}
\item[]         keyword: {\hyperref[\detokenize{xsrst/keyword:id1}]{\sphinxcrossref{\DUrole{std,std-ref}{4.9.4 Example Keyword File}}}}
\item[]         spelling: {\hyperref[\detokenize{xsrst/spelling:id1}]{\sphinxcrossref{\DUrole{std,std-ref}{4.9.3 Example Spelling File}}}}
\end{DUlineblock}

\index{spelling@\spxentry{spelling}}\ignorespaces 

\subparagraph{spelling}
\label{\detokenize{xsrst/xsrst_py:spelling}}\label{\detokenize{xsrst/xsrst_py:xsrst-py-command-line-arguments-spelling}}\label{\detokenize{xsrst/xsrst_py:index-12}}
The command line argument \sphinxstyleemphasis{spelling} is the name of a file,
relative to the \sphinxstyleemphasis{sphinx\_dir} directory.
This file contains a list of words
that the spell checker will consider correct for all sections
(it can be an empty file).
A line that begins with \sphinxcode{\sphinxupquote{\#}} is a comment (not included in the list).
The words are one per line and
leading and trailing white space in a word are ignored.
Special words, for a particular section, are specified using the
{\hyperref[\detokenize{xsrst/spell_cmd:id1}]{\sphinxcrossref{\DUrole{std,std-ref}{spell command}}}}.

\index{keyword@\spxentry{keyword}}\ignorespaces 

\subparagraph{keyword}
\label{\detokenize{xsrst/xsrst_py:keyword}}\label{\detokenize{xsrst/xsrst_py:xsrst-py-command-line-arguments-keyword}}\label{\detokenize{xsrst/xsrst_py:index-13}}
The command line argument \sphinxstyleemphasis{keyword} is the name of a file,
relative to the \sphinxstyleemphasis{sphinx\_dir} directory.
This file contains a list of python regular expressions for heading tokens
that should not be included in the index (it can be an empty file).
A heading token is any sequence of non space or new line characters
with upper case letters converted to lower case.
For example, a heading might contain the token \sphinxcode{\sphinxupquote{The}} but you
might not want to include \sphinxcode{\sphinxupquote{the}} as a entry in the \DUrole{xref,std,std-ref}{genindex}.
In this case you could have a line containing just \sphinxcode{\sphinxupquote{the}} in \sphinxstyleemphasis{keyword}.
For another example, you might want to exclude all tokens that are numbers.
In this case you could have a line containing just \sphinxcode{\sphinxupquote{{[}0\sphinxhyphen{}9{]}*}} in \sphinxstyleemphasis{keyword}.
The regular expressions are one per line and
leading and trailing spaces are ignored.
A line that begins with \sphinxcode{\sphinxupquote{\#}} is a comment
(not included in the list of python regular expressions).

\index{line\_increment@\spxentry{line\_increment}}\ignorespaces 

\subparagraph{line\_increment}
\label{\detokenize{xsrst/xsrst_py:line-increment}}\label{\detokenize{xsrst/xsrst_py:xsrst-py-command-line-arguments-line-increment}}\label{\detokenize{xsrst/xsrst_py:index-14}}
This optional argument helps find the source of errors reported by sphinx.
If the argument \sphinxstyleemphasis{line\_increment} is present,
a table is generated at the end of each output file.
This table maps line numbers in the output file to
line numbers in the corresponding xsrst input file.
The argument \sphinxstyleemphasis{line\_increment} is a positive integer specifying the minimum
difference between xsrst input line numbers for entries in the table.
The value \sphinxcode{\sphinxupquote{1}} will give the maximum resolution.
For example, the sphinx warning

\begin{DUlineblock}{0em}
\item[]         … \sphinxcode{\sphinxupquote{/xsrst/children\_exam.rst:30: WARNING:}} …
\end{DUlineblock}

corresponds to line number 30 in the file \sphinxcode{\sphinxupquote{children\_exam.rst}}.
The table at the bottom of that file maps line numbers in
\sphinxcode{\sphinxupquote{children\_exam.rst}} to line numbers in the corresponding xsrst input file.

\index{table@\spxentry{table}}\index{contents@\spxentry{contents}}\ignorespaces 

\subparagraph{Table Of Contents}
\label{\detokenize{xsrst/xsrst_py:table-of-contents}}\label{\detokenize{xsrst/xsrst_py:xsrst-py-table-of-contents}}\label{\detokenize{xsrst/xsrst_py:index-15}}
\index{toctree@\spxentry{toctree}}\ignorespaces 

\subparagraph{toctree}
\label{\detokenize{xsrst/xsrst_py:toctree}}\label{\detokenize{xsrst/xsrst_py:xsrst-py-table-of-contents-toctree}}\label{\detokenize{xsrst/xsrst_py:index-16}}
The sphinx \sphinxcode{\sphinxupquote{toctree}} directives are automatically generated
for sections. The only such directive you should directly edit
is in the file \sphinxstyleemphasis{sphinx\_dir} \sphinxcode{\sphinxupquote{/index.rst}}.
One entry in this file specifies the
{\hyperref[\detokenize{xsrst/xsrst_py:xsrst-py-command-line-arguments-root-file-root-section}]{\sphinxcrossref{\DUrole{std,std-ref}{root\_section}}}}.
Other entries are for \sphinxcode{\sphinxupquote{.rst}} files that are not extracted by
\sphinxcode{\sphinxupquote{xsrst.py}}.

\index{parent@\spxentry{parent}}\index{section@\spxentry{section}}\ignorespaces 

\subparagraph{Parent Section}
\label{\detokenize{xsrst/xsrst_py:parent-section}}\label{\detokenize{xsrst/xsrst_py:xsrst-py-table-of-contents-parent-section}}\label{\detokenize{xsrst/xsrst_py:index-17}}
A single input file may contain multiple
{\hyperref[\detokenize{xsrst/begin_cmd:begin-cmd-section}]{\sphinxcrossref{\DUrole{std,std-ref}{sections}}}}.
One (and at most one) of these sections may use begin with a
{\hyperref[\detokenize{xsrst/begin_cmd:begin-cmd-parent-section}]{\sphinxcrossref{\DUrole{std,std-ref}{parent begin}}}} command.
In this case, the other sections in the file are children of this section
and this section is a child of the section containing the
{\hyperref[\detokenize{xsrst/child_cmd:id1}]{\sphinxcrossref{\DUrole{std,std-ref}{child command}}}} that included this file.

If there is no parent section for a file,
all the sections in the file are children of the section containing the
child command that included the file.

\index{heading@\spxentry{heading}}\index{links@\spxentry{links}}\ignorespaces 

\subparagraph{Heading Links}
\label{\detokenize{xsrst/xsrst_py:heading-links}}\label{\detokenize{xsrst/xsrst_py:xsrst-py-heading-links}}\label{\detokenize{xsrst/xsrst_py:index-18}}\begin{itemize}
\item {} 
For each word in a heading,
a link is included in the index from the word to the heading.

\item {} 
Each word in a heading is added to the html keyword meta data.

\item {} 
A cross reference label is defined for linking
from anywhere to a heading. The details of how to use
these labels are described below.

\item {} 
Headings can also be used to help find links to children
of the current section; see the heading
{\hyperref[\detokenize{xsrst/xsrst_py:xsrst-py-heading-links-children}]{\sphinxcrossref{\DUrole{std,std-ref}{Children}}}} below.

\end{itemize}

\index{section@\spxentry{section}}\index{level@\spxentry{level}}\ignorespaces 

\subparagraph{Section Level}
\label{\detokenize{xsrst/xsrst_py:section-level}}\label{\detokenize{xsrst/xsrst_py:xsrst-py-heading-links-section-level}}\label{\detokenize{xsrst/xsrst_py:index-19}}
Each {\hyperref[\detokenize{xsrst/begin_cmd:begin-cmd-section}]{\sphinxcrossref{\DUrole{std,std-ref}{section}}}} can have only one header at
the first level which is a title for the section.
The {\hyperref[\detokenize{xsrst/begin_cmd:begin-cmd-section-name}]{\sphinxcrossref{\DUrole{std,std-ref}{section\_name}}}}
is automatically used
as a label for linking the title for a section; i.e., the
following will link to the title for \sphinxstyleemphasis{section\_name}:

        \sphinxcode{\sphinxupquote{:ref:}}  \textasciigrave{} \sphinxstyleemphasis{linking\_text} \sphinxcode{\sphinxupquote{\textless{}}} \sphinxstyleemphasis{section\_name} \sphinxcode{\sphinxupquote{\textgreater{}}} \textasciigrave{}

where \sphinxstyleemphasis{linking\_text} is the text the user sees.

\index{other@\spxentry{other}}\index{levels@\spxentry{levels}}\ignorespaces 

\subparagraph{Other Levels}
\label{\detokenize{xsrst/xsrst_py:other-levels}}\label{\detokenize{xsrst/xsrst_py:xsrst-py-heading-links-other-levels}}\label{\detokenize{xsrst/xsrst_py:index-20}}
The label for linking a heading that is not at the first level
is the label for the heading directly above it plus a dot character \sphinxcode{\sphinxupquote{.}},
plus a lower case version of the heading with spaces converted to
underbars \sphinxcode{\sphinxupquote{\_}}. For example, the label for the heading for this
paragraph is

        \sphinxcode{\sphinxupquote{xsrst\_py.headers\_and\_links.other\_levels}}.

This may seem verbose, but it helps keep the links up to date.
If a heading changes, all the links to that heading will break.
This identifies the links that should be checked
to make sure they are still valid.

\index{children@\spxentry{children}}\ignorespaces 

\subparagraph{Children}
\label{\detokenize{xsrst/xsrst_py:children}}\label{\detokenize{xsrst/xsrst_py:xsrst-py-heading-links-children}}\label{\detokenize{xsrst/xsrst_py:index-21}}
If a xsrst input file has a
{\hyperref[\detokenize{xsrst/xsrst_py:xsrst-py-table-of-contents-parent-section}]{\sphinxcrossref{\DUrole{std,std-ref}{parent section}}}}
the other sections in the file are children of the parent.
\begin{itemize}
\item {} 
If a section has a {\hyperref[\detokenize{xsrst/child_cmd:id1}]{\sphinxcrossref{\DUrole{std,std-ref}{child link or list command}}}}
links to all the children of the section are placed where the
child link command is located.

\item {} 
If a section has a {\hyperref[\detokenize{xsrst/child_cmd:id1}]{\sphinxcrossref{\DUrole{std,std-ref}{children command}}}}
no automatic links to the children of the current section are generated.

\item {} 
Otherwise, the links to the children of a section are placed
at the end of the section.

\end{itemize}

You can place a heading directly before the links to make them easier to find.

\index{example@\spxentry{example}}\ignorespaces 

\subparagraph{Example}
\label{\detokenize{xsrst/xsrst_py:example}}\label{\detokenize{xsrst/xsrst_py:xsrst-py-heading-links-example}}\label{\detokenize{xsrst/xsrst_py:index-22}}
{\hyperref[\detokenize{xsrst/heading_exam:id1}]{\sphinxcrossref{\DUrole{std,std-ref}{4.8 Heading Example}}}}

\index{indentation@\spxentry{indentation}}\ignorespaces 

\subparagraph{Indentation}
\label{\detokenize{xsrst/xsrst_py:indentation}}\label{\detokenize{xsrst/xsrst_py:xsrst-py-indentation}}\label{\detokenize{xsrst/xsrst_py:index-23}}
If there are a number of spaces before
all of the xsrst documentation for a section,
those characters are not included in the xsrst output.
This enables one to indent the
xsrst so it is grouped with the proper code block in the source.
An error message will result if
you use tabs in the indentation.

\index{example@\spxentry{example}}\ignorespaces 

\subparagraph{Example}
\label{\detokenize{xsrst/xsrst_py:xsrst-py-indentation-example}}\label{\detokenize{xsrst/xsrst_py:index-24}}\label{\detokenize{xsrst/xsrst_py:id2}}\begin{itemize}
\item {} 
{\hyperref[\detokenize{xsrst/indent_exam:id1}]{\sphinxcrossref{\DUrole{std,std-ref}{4.2.1.1.2.1 Indent Example}}}}

\end{itemize}

\index{wish@\spxentry{wish}}\index{list@\spxentry{list}}\ignorespaces 

\subparagraph{Wish List}
\label{\detokenize{xsrst/xsrst_py:wish-list}}\label{\detokenize{xsrst/xsrst_py:xsrst-py-wish-list}}\label{\detokenize{xsrst/xsrst_py:index-25}}
The following is a wish list for future improvements to \sphinxcode{\sphinxupquote{xsrst.py}}:

\index{tabs@\spxentry{tabs}}\ignorespaces 

\subparagraph{Tabs}
\label{\detokenize{xsrst/xsrst_py:tabs}}\label{\detokenize{xsrst/xsrst_py:xsrst-py-wish-list-tabs}}\label{\detokenize{xsrst/xsrst_py:index-26}}
Tabs in a code blocks get expanded to 8 spaces; see \sphinxhref{https://stackoverflow.com/questions/1686837/sphinx-documentation-tool-set-tab-width-in-output}{stackoverflow}.
It would be nice to have a way to control the size of tabs in the code blocks
displayed by {\hyperref[\detokenize{xsrst/code_cmd:id1}]{\sphinxcrossref{\DUrole{std,std-ref}{4.5 Code Command}}}} and {\hyperref[\detokenize{xsrst/file_cmd:id1}]{\sphinxcrossref{\DUrole{std,std-ref}{4.6 File Command}}}}.
Perhaps it would be good to support tabs as a method for
indenting xsrst input sections.

\index{module@\spxentry{module}}\ignorespaces 

\subparagraph{Module}
\label{\detokenize{xsrst/xsrst_py:module}}\label{\detokenize{xsrst/xsrst_py:xsrst-py-wish-list-module}}\label{\detokenize{xsrst/xsrst_py:index-27}}
Convert the program into a python module and provide a pip distribution for it.
It would at least be nice for cppad\_py to install the \sphinxcode{\sphinxupquote{xsrst.py}} program
so that users would not have to copy it to a directory in
their execution path.

\index{commands@\spxentry{commands}}\ignorespaces 

\subparagraph{Commands}
\label{\detokenize{xsrst/xsrst_py:commands}}\label{\detokenize{xsrst/xsrst_py:xsrst-py-commands}}\label{\detokenize{xsrst/xsrst_py:index-28}}\begin{itemize}
\item {} 
{\hyperref[\detokenize{xsrst/begin_cmd:id1}]{\sphinxcrossref{\DUrole{std,std-ref}{4.1 Begin and End Commands}}}}

\item {} 
{\hyperref[\detokenize{xsrst/child_cmd:id1}]{\sphinxcrossref{\DUrole{std,std-ref}{4.2 Children Commands}}}}

\item {} 
{\hyperref[\detokenize{xsrst/spell_cmd:id1}]{\sphinxcrossref{\DUrole{std,std-ref}{4.3 Spell Command}}}}

\item {} 
{\hyperref[\detokenize{xsrst/suspend_cmd:id1}]{\sphinxcrossref{\DUrole{std,std-ref}{4.4 Suspend and Resume Commands}}}}

\item {} 
{\hyperref[\detokenize{xsrst/code_cmd:id1}]{\sphinxcrossref{\DUrole{std,std-ref}{4.5 Code Command}}}}

\item {} 
{\hyperref[\detokenize{xsrst/file_cmd:id1}]{\sphinxcrossref{\DUrole{std,std-ref}{4.6 File Command}}}}

\item {} 
{\hyperref[\detokenize{xsrst/comment_ch_cmd:id1}]{\sphinxcrossref{\DUrole{std,std-ref}{4.7 Comment Character Command}}}}

\end{itemize}


\bigskip\hrule\bigskip


xsrst input file: \sphinxcode{\sphinxupquote{bin/xsrst.py}}


\subsubsection{whats\_new\_2020}
\label{\detokenize{xsrst/whats_new_2020:whats-new-2020}}\label{\detokenize{xsrst/whats_new_2020::doc}}
\index{whats\_new\_2020@\spxentry{whats\_new\_2020}}\index{cppad@\spxentry{cppad}}\index{py@\spxentry{py}}\index{changes@\spxentry{changes}}\index{during@\spxentry{during}}\index{2020@\spxentry{2020}}\ignorespaces 

\paragraph{5 CppAD Py Changes During 2020}
\label{\detokenize{xsrst/whats_new_2020:cppad-py-changes-during-2020}}\label{\detokenize{xsrst/whats_new_2020:id1}}\begin{itemize}
\item {} 
{\hyperref[\detokenize{xsrst/whats_new_2020:whats-new-2020-previous-years}]{\sphinxcrossref{\DUrole{std,std-ref}{Previous Years}}}}

\item {} 
{\hyperref[\detokenize{xsrst/whats_new_2020:whats-new-2020-09-13}]{\sphinxcrossref{\DUrole{std,std-ref}{09\sphinxhyphen{}13}}}}

\item {} 
{\hyperref[\detokenize{xsrst/whats_new_2020:whats-new-2020-09-12}]{\sphinxcrossref{\DUrole{std,std-ref}{09\sphinxhyphen{}12}}}}

\item {} 
{\hyperref[\detokenize{xsrst/whats_new_2020:whats-new-2020-07-31}]{\sphinxcrossref{\DUrole{std,std-ref}{07\sphinxhyphen{}31}}}}

\item {} 
{\hyperref[\detokenize{xsrst/whats_new_2020:whats-new-2020-07-30}]{\sphinxcrossref{\DUrole{std,std-ref}{07\sphinxhyphen{}30}}}}

\item {} 
{\hyperref[\detokenize{xsrst/whats_new_2020:whats-new-2020-07-29}]{\sphinxcrossref{\DUrole{std,std-ref}{07\sphinxhyphen{}29}}}}

\item {} 
{\hyperref[\detokenize{xsrst/whats_new_2020:whats-new-2020-07-28}]{\sphinxcrossref{\DUrole{std,std-ref}{07\sphinxhyphen{}28}}}}

\item {} 
{\hyperref[\detokenize{xsrst/whats_new_2020:whats-new-2020-07-25}]{\sphinxcrossref{\DUrole{std,std-ref}{07\sphinxhyphen{}25}}}}

\item {} 
{\hyperref[\detokenize{xsrst/whats_new_2020:whats-new-2020-07-24}]{\sphinxcrossref{\DUrole{std,std-ref}{07\sphinxhyphen{}24}}}}

\item {} 
{\hyperref[\detokenize{xsrst/whats_new_2020:whats-new-2020-07-18}]{\sphinxcrossref{\DUrole{std,std-ref}{07\sphinxhyphen{}18}}}}

\item {} 
{\hyperref[\detokenize{xsrst/whats_new_2020:whats-new-2020-07-05}]{\sphinxcrossref{\DUrole{std,std-ref}{07\sphinxhyphen{}05}}}}

\item {} 
{\hyperref[\detokenize{xsrst/whats_new_2020:whats-new-2020-05-17}]{\sphinxcrossref{\DUrole{std,std-ref}{05\sphinxhyphen{}17}}}}

\item {} 
{\hyperref[\detokenize{xsrst/whats_new_2020:whats-new-2020-05-16}]{\sphinxcrossref{\DUrole{std,std-ref}{05\sphinxhyphen{}16}}}}

\item {} 
{\hyperref[\detokenize{xsrst/whats_new_2020:whats-new-2020-05-15}]{\sphinxcrossref{\DUrole{std,std-ref}{05\sphinxhyphen{}15}}}}

\item {} 
{\hyperref[\detokenize{xsrst/whats_new_2020:whats-new-2020-05-14}]{\sphinxcrossref{\DUrole{std,std-ref}{05\sphinxhyphen{}14}}}}

\item {} 
{\hyperref[\detokenize{xsrst/whats_new_2020:whats-new-2020-05-12}]{\sphinxcrossref{\DUrole{std,std-ref}{05\sphinxhyphen{}12}}}}

\item {} 
{\hyperref[\detokenize{xsrst/whats_new_2020:whats-new-2020-05-09}]{\sphinxcrossref{\DUrole{std,std-ref}{05\sphinxhyphen{}09}}}}

\item {} 
{\hyperref[\detokenize{xsrst/whats_new_2020:whats-new-2020-05-08}]{\sphinxcrossref{\DUrole{std,std-ref}{05\sphinxhyphen{}08}}}}

\item {} 
{\hyperref[\detokenize{xsrst/whats_new_2020:whats-new-2020-05-07}]{\sphinxcrossref{\DUrole{std,std-ref}{05\sphinxhyphen{}07}}}}

\item {} 
{\hyperref[\detokenize{xsrst/whats_new_2020:whats-new-2020-05-06}]{\sphinxcrossref{\DUrole{std,std-ref}{05\sphinxhyphen{}06}}}}

\item {} 
{\hyperref[\detokenize{xsrst/whats_new_2020:whats-new-2020-05-05}]{\sphinxcrossref{\DUrole{std,std-ref}{05\sphinxhyphen{}05}}}}

\item {} 
{\hyperref[\detokenize{xsrst/whats_new_2020:whats-new-2020-05-04}]{\sphinxcrossref{\DUrole{std,std-ref}{05\sphinxhyphen{}04}}}}

\item {} 
{\hyperref[\detokenize{xsrst/whats_new_2020:whats-new-2020-04-28}]{\sphinxcrossref{\DUrole{std,std-ref}{04\sphinxhyphen{}28}}}}

\item {} 
{\hyperref[\detokenize{xsrst/whats_new_2020:whats-new-2020-04-27}]{\sphinxcrossref{\DUrole{std,std-ref}{04\sphinxhyphen{}27}}}}

\item {} 
{\hyperref[\detokenize{xsrst/whats_new_2020:whats-new-2020-04-26}]{\sphinxcrossref{\DUrole{std,std-ref}{04\sphinxhyphen{}26}}}}

\item {} 
{\hyperref[\detokenize{xsrst/whats_new_2020:whats-new-2020-04-25}]{\sphinxcrossref{\DUrole{std,std-ref}{04\sphinxhyphen{}25}}}}

\item {} 
{\hyperref[\detokenize{xsrst/whats_new_2020:whats-new-2020-04-24}]{\sphinxcrossref{\DUrole{std,std-ref}{04\sphinxhyphen{}24}}}}

\item {} 
{\hyperref[\detokenize{xsrst/whats_new_2020:whats-new-2020-04-23}]{\sphinxcrossref{\DUrole{std,std-ref}{04\sphinxhyphen{}23}}}}

\item {} 
{\hyperref[\detokenize{xsrst/whats_new_2020:whats-new-2020-04-22}]{\sphinxcrossref{\DUrole{std,std-ref}{04\sphinxhyphen{}22}}}}

\item {} 
{\hyperref[\detokenize{xsrst/whats_new_2020:whats-new-2020-04-20}]{\sphinxcrossref{\DUrole{std,std-ref}{04\sphinxhyphen{}20}}}}

\item {} 
{\hyperref[\detokenize{xsrst/whats_new_2020:whats-new-2020-04-19}]{\sphinxcrossref{\DUrole{std,std-ref}{04\sphinxhyphen{}19}}}}

\item {} 
{\hyperref[\detokenize{xsrst/whats_new_2020:whats-new-2020-04-18}]{\sphinxcrossref{\DUrole{std,std-ref}{04\sphinxhyphen{}18}}}}

\item {} 
{\hyperref[\detokenize{xsrst/whats_new_2020:whats-new-2020-04-13}]{\sphinxcrossref{\DUrole{std,std-ref}{04\sphinxhyphen{}13}}}}

\item {} 
{\hyperref[\detokenize{xsrst/whats_new_2020:whats-new-2020-04-12}]{\sphinxcrossref{\DUrole{std,std-ref}{04\sphinxhyphen{}12}}}}

\item {} 
{\hyperref[\detokenize{xsrst/whats_new_2020:whats-new-2020-04-10}]{\sphinxcrossref{\DUrole{std,std-ref}{04\sphinxhyphen{}10}}}}

\end{itemize}


\subparagraph{whats\_new\_2018}
\label{\detokenize{xsrst/whats_new_2018:whats-new-2018}}\label{\detokenize{xsrst/whats_new_2018::doc}}
\index{whats\_new\_2018@\spxentry{whats\_new\_2018}}\index{cppad@\spxentry{cppad}}\index{py@\spxentry{py}}\index{changes@\spxentry{changes}}\index{during@\spxentry{during}}\index{2018@\spxentry{2018}}\ignorespaces 

\subparagraph{5.1 Cppad Py Changes During 2018}
\label{\detokenize{xsrst/whats_new_2018:cppad-py-changes-during-2018}}\label{\detokenize{xsrst/whats_new_2018:id1}}\begin{itemize}
\item {} 
{\hyperref[\detokenize{xsrst/whats_new_2018:whats-new-2018-11-10}]{\sphinxcrossref{\DUrole{std,std-ref}{11\sphinxhyphen{}10}}}}

\item {} 
{\hyperref[\detokenize{xsrst/whats_new_2018:whats-new-2018-11-09}]{\sphinxcrossref{\DUrole{std,std-ref}{11\sphinxhyphen{}09}}}}

\item {} 
{\hyperref[\detokenize{xsrst/whats_new_2018:whats-new-2018-11-07}]{\sphinxcrossref{\DUrole{std,std-ref}{11\sphinxhyphen{}07}}}}

\item {} 
{\hyperref[\detokenize{xsrst/whats_new_2018:whats-new-2018-11-05}]{\sphinxcrossref{\DUrole{std,std-ref}{11\sphinxhyphen{}05}}}}

\item {} 
{\hyperref[\detokenize{xsrst/whats_new_2018:whats-new-2018-08-13}]{\sphinxcrossref{\DUrole{std,std-ref}{08\sphinxhyphen{}13}}}}

\item {} 
{\hyperref[\detokenize{xsrst/whats_new_2018:whats-new-2018-07-31}]{\sphinxcrossref{\DUrole{std,std-ref}{07\sphinxhyphen{}31}}}}

\item {} 
{\hyperref[\detokenize{xsrst/whats_new_2018:whats-new-2018-07-26}]{\sphinxcrossref{\DUrole{std,std-ref}{07\sphinxhyphen{}26}}}}

\item {} 
{\hyperref[\detokenize{xsrst/whats_new_2018:whats-new-2018-07-19}]{\sphinxcrossref{\DUrole{std,std-ref}{07\sphinxhyphen{}19}}}}

\item {} 
{\hyperref[\detokenize{xsrst/whats_new_2018:whats-new-2018-07-15}]{\sphinxcrossref{\DUrole{std,std-ref}{07\sphinxhyphen{}15}}}}

\item {} 
{\hyperref[\detokenize{xsrst/whats_new_2018:whats-new-2018-07-14}]{\sphinxcrossref{\DUrole{std,std-ref}{07\sphinxhyphen{}14}}}}

\item {} 
{\hyperref[\detokenize{xsrst/whats_new_2018:whats-new-2018-07-13}]{\sphinxcrossref{\DUrole{std,std-ref}{07\sphinxhyphen{}13}}}}

\item {} 
{\hyperref[\detokenize{xsrst/whats_new_2018:whats-new-2018-07-12}]{\sphinxcrossref{\DUrole{std,std-ref}{07\sphinxhyphen{}12}}}}

\item {} 
{\hyperref[\detokenize{xsrst/whats_new_2018:whats-new-2018-07-10}]{\sphinxcrossref{\DUrole{std,std-ref}{07\sphinxhyphen{}10}}}}

\item {} 
{\hyperref[\detokenize{xsrst/whats_new_2018:whats-new-2018-07-08}]{\sphinxcrossref{\DUrole{std,std-ref}{07\sphinxhyphen{}08}}}}

\item {} 
{\hyperref[\detokenize{xsrst/whats_new_2018:whats-new-2018-07-07}]{\sphinxcrossref{\DUrole{std,std-ref}{07\sphinxhyphen{}07}}}}

\item {} 
{\hyperref[\detokenize{xsrst/whats_new_2018:whats-new-2018-07-03}]{\sphinxcrossref{\DUrole{std,std-ref}{07\sphinxhyphen{}03}}}}

\end{itemize}


\subparagraph{11\sphinxhyphen{}10}
\label{\detokenize{xsrst/whats_new_2018:whats-new-2018-11-10}}\label{\detokenize{xsrst/whats_new_2018:id2}}\begin{enumerate}
\sphinxsetlistlabels{\arabic}{enumi}{enumii}{}{.}%
\item {} 
Added the {\hyperref[\detokenize{xsrst/py_independent:py-independent-dynamic}]{\sphinxcrossref{\DUrole{std,std-ref}{dynamic}}}}
argument to the \sphinxcode{\sphinxupquote{independent}} function and the
{\hyperref[\detokenize{xsrst/py_fun_new_dynamic:id1}]{\sphinxcrossref{\DUrole{std,std-ref}{f\_new\_dynamic}}}} function.

\item {} 
Added the sections {\hyperref[\detokenize{xsrst/more_cpp:id1}]{\sphinxcrossref{\DUrole{std,std-ref}{more\_cpp}}}} and {\hyperref[\detokenize{xsrst/more_py:id1}]{\sphinxcrossref{\DUrole{std,std-ref}{more\_py}}}}.

\end{enumerate}


\subparagraph{11\sphinxhyphen{}09}
\label{\detokenize{xsrst/whats_new_2018:whats-new-2018-11-09}}\label{\detokenize{xsrst/whats_new_2018:id3}}
Include binary operations where first operand is an \sphinxcode{\sphinxupquote{a\_double}}
and second operand is a \sphinxcode{\sphinxupquote{double}} ; see
{\hyperref[\detokenize{xsrst/a_double_binary:id1}]{\sphinxcrossref{\DUrole{std,std-ref}{a\_double\_binary}}}},
{\hyperref[\detokenize{xsrst/a_double_compare:id1}]{\sphinxcrossref{\DUrole{std,std-ref}{a\_double\_compare}}}}, and
{\hyperref[\detokenize{xsrst/a_double_assign:id1}]{\sphinxcrossref{\DUrole{std,std-ref}{a\_double\_assign}}}}.


\subparagraph{11\sphinxhyphen{}07}
\label{\detokenize{xsrst/whats_new_2018:whats-new-2018-11-07}}\label{\detokenize{xsrst/whats_new_2018:id4}}\begin{enumerate}
\sphinxsetlistlabels{\arabic}{enumi}{enumii}{}{.}%
\item {} 
Advance to cppad\sphinxhyphen{}20181106 (must re\sphinxhyphen{}run \sphinxcode{\sphinxupquote{bin/get\_cppad.sh}} ).

\item {} 
The {\hyperref[\detokenize{xsrst/py_fun_ctor:py-fun-ctor-syntax-a-fun}]{\sphinxcrossref{\DUrole{std,std-ref}{a\_fun}}}} objects were added.
This support all the operations that
{\hyperref[\detokenize{xsrst/py_fun_ctor:py-fun-ctor-syntax-d-fun}]{\sphinxcrossref{\DUrole{std,std-ref}{d\_fun}}}} support all the
{\hyperref[\detokenize{xsrst/py_fun:id1}]{\sphinxcrossref{\DUrole{std,std-ref}{py\_fun}}}} operations with the exception of
{\hyperref[\detokenize{xsrst/py_fun_optimize:id1}]{\sphinxcrossref{\DUrole{std,std-ref}{optimize}}}}.
The {\hyperref[\detokenize{xsrst/py_sparse:id1}]{\sphinxcrossref{\DUrole{std,std-ref}{py\_sparse}}}} operations are not yet supported.

\end{enumerate}


\subparagraph{11\sphinxhyphen{}05}
\label{\detokenize{xsrst/whats_new_2018:whats-new-2018-11-05}}\label{\detokenize{xsrst/whats_new_2018:id5}}\begin{enumerate}
\sphinxsetlistlabels{\arabic}{enumi}{enumii}{}{.}%
\item {} 
\sphinxstylestrong{API Change} :
change \sphinxcode{\sphinxupquote{cppad\_py.a\_fun}} to \sphinxcode{\sphinxupquote{cppad\_py.d\_fun}}
because it is a function that evaluates in \sphinxcode{\sphinxupquote{double}}
and plan to use \sphinxcode{\sphinxupquote{cppad\_py.a\_fun}} for a function that
evaluates in \sphinxcode{\sphinxupquote{a\_double}} .

\item {} 
Advance to cppad\sphinxhyphen{}20181105 (must re\sphinxhyphen{}run \sphinxcode{\sphinxupquote{bin/get\_cppad.sh}} ).

\end{enumerate}


\subparagraph{08\sphinxhyphen{}13}
\label{\detokenize{xsrst/whats_new_2018:whats-new-2018-08-13}}\label{\detokenize{xsrst/whats_new_2018:id6}}\begin{enumerate}
\sphinxsetlistlabels{\arabic}{enumi}{enumii}{}{.}%
\item {} 
Change \sphinxcode{\sphinxupquote{cppad\_cxx\_flags}} to
{\hyperref[\detokenize{xsrst/get_cppad_sh:get-cppad-sh-settings-extra-cxx-flags}]{\sphinxcrossref{\DUrole{std,std-ref}{extra\_cxx\_flags}}}} in the user
configuration settings.
Also remove \sphinxcode{\sphinxupquote{swig\_cxx\_flags}} from these user settings
(set automatically).

\item {} 
Improve the install instructions; see {\hyperref[\detokenize{xsrst/setup_py:id1}]{\sphinxcrossref{\DUrole{std,std-ref}{setup\_py}}}}.

\end{enumerate}


\subparagraph{07\sphinxhyphen{}31}
\label{\detokenize{xsrst/whats_new_2018:whats-new-2018-07-31}}\label{\detokenize{xsrst/whats_new_2018:id7}}\begin{enumerate}
\sphinxsetlistlabels{\arabic}{enumi}{enumii}{}{.}%
\item {} 
The general purpose Swig example \sphinxcode{\sphinxupquote{swig\_xam}} ,
which was not specific to cppad\_py,
has been removed.

\item {} 
A change to \sphinxcode{\sphinxupquote{setup.py}} was making the
{\hyperref[\detokenize{xsrst/get_cppad_sh:id1}]{\sphinxcrossref{\DUrole{std,std-ref}{get\_cppad\_sh}}}} script fail.
This has been fixed.

\item {} 
The python tests are no run in \sphinxcode{\sphinxupquote{example/python}}
(not copied to the build directory).

\end{enumerate}


\subparagraph{07\sphinxhyphen{}26}
\label{\detokenize{xsrst/whats_new_2018:whats-new-2018-07-26}}\label{\detokenize{xsrst/whats_new_2018:id8}}
Change \sphinxcode{\sphinxupquote{cppad\_py/python\_major\_version}}
to \sphinxcode{\sphinxupquote{cppad\_py/python\_version}}
because it now contains both major and minor version numbers.


\subparagraph{07\sphinxhyphen{}19}
\label{\detokenize{xsrst/whats_new_2018:whats-new-2018-07-19}}\label{\detokenize{xsrst/whats_new_2018:id9}}
Add \sphinxcode{\sphinxupquote{\sphinxhyphen{}std=c++11}} to the default compile flags in
{\hyperref[\detokenize{xsrst/setup_py:id1}]{\sphinxcrossref{\DUrole{std,std-ref}{setup\_py}}}}.
Add \sphinxcode{\sphinxupquote{\sphinxhyphen{}py3}} to the \sphinxcode{\sphinxupquote{swig}} builds
when running \sphinxcode{\sphinxupquote{setup.py}} using Python 3; see
\sphinxhref{http://www.swig.org/Doc1.3/Python.html\#Python\_python3support}{swig python 3}.


\subparagraph{07\sphinxhyphen{}15}
\label{\detokenize{xsrst/whats_new_2018:whats-new-2018-07-15}}\label{\detokenize{xsrst/whats_new_2018:id10}}\begin{enumerate}
\sphinxsetlistlabels{\arabic}{enumi}{enumii}{}{.}%
\item {} 
Better error reporting for type and size errors in
\sphinxcode{\sphinxupquote{numpy.array}} arguments to \sphinxcode{\sphinxupquote{cppad\_py}} functions using
{\hyperref[\detokenize{xsrst/numpy2vec:id1}]{\sphinxcrossref{\DUrole{std,std-ref}{numpy2vec}}}}.

\item {} 
Finish converting all vectors and matrices in
{\hyperref[\detokenize{xsrst/py_lib:id1}]{\sphinxcrossref{\DUrole{std,std-ref}{py\_lib}}}} to numpy arrays.
Note the {\hyperref[\detokenize{xsrst/py_utility:id1}]{\sphinxcrossref{\DUrole{std,std-ref}{py\_utility}}}} routines do the conversion.

\end{enumerate}


\subparagraph{07\sphinxhyphen{}14}
\label{\detokenize{xsrst/whats_new_2018:whats-new-2018-07-14}}\label{\detokenize{xsrst/whats_new_2018:id11}}\begin{enumerate}
\sphinxsetlistlabels{\arabic}{enumi}{enumii}{}{.}%
\item {} 
Continue conversion of python library to using \sphinxcode{\sphinxupquote{numpy.array}} ; see
{\hyperref[\detokenize{xsrst/py_fun_forward:id1}]{\sphinxcrossref{\DUrole{std,std-ref}{py\_fun\_forward}}}}, {\hyperref[\detokenize{xsrst/py_fun_reverse:id1}]{\sphinxcrossref{\DUrole{std,std-ref}{py\_fun\_reverse}}}}.

\item {} 
Add the {\hyperref[\detokenize{xsrst/py_fun_property:py-fun-property-size-order}]{\sphinxcrossref{\DUrole{std,std-ref}{f\_size\_order()}}}} function.

\end{enumerate}


\subparagraph{07\sphinxhyphen{}13}
\label{\detokenize{xsrst/whats_new_2018:whats-new-2018-07-13}}\label{\detokenize{xsrst/whats_new_2018:id12}}
Continue conversion of python library to using \sphinxcode{\sphinxupquote{numpy.array}} ; see
{\hyperref[\detokenize{xsrst/py_fun_jacobian:id1}]{\sphinxcrossref{\DUrole{std,std-ref}{py\_fun\_jacobian}}}}, {\hyperref[\detokenize{xsrst/py_fun_hessian:id1}]{\sphinxcrossref{\DUrole{std,std-ref}{py\_fun\_hessian}}}}.


\subparagraph{07\sphinxhyphen{}12}
\label{\detokenize{xsrst/whats_new_2018:whats-new-2018-07-12}}\label{\detokenize{xsrst/whats_new_2018:id13}}\begin{enumerate}
\sphinxsetlistlabels{\arabic}{enumi}{enumii}{}{.}%
\item {} 
Continue conversion of python library to using \sphinxcode{\sphinxupquote{numpy.array}} ; see
{\hyperref[\detokenize{xsrst/py_independent:id1}]{\sphinxcrossref{\DUrole{std,std-ref}{py\_independent}}}}, {\hyperref[\detokenize{xsrst/py_fun_ctor:id1}]{\sphinxcrossref{\DUrole{std,std-ref}{py\_fun\_ctor}}}}.

\item {} 
Add the file \sphinxcode{\sphinxupquote{cppad\_py/python\_major\_version}} which contains
the major version of python that this build of \sphinxcode{\sphinxupquote{cppad\_py}} is for.

\item {} 
Automatically use {\hyperref[\detokenize{xsrst/get_cppad_sh:get-cppad-sh-settings-build-type}]{\sphinxcrossref{\DUrole{std,std-ref}{build\_type}}}}
control how setup.py compiles C++
(one used to use flags on setup.py command line).

\item {} 
setup.py was using the git repository version of Cppad instead of the
one installed in \sphinxcode{\sphinxupquote{build/prefix}} by {\hyperref[\detokenize{xsrst/get_cppad_sh:id1}]{\sphinxcrossref{\DUrole{std,std-ref}{get\_cppad\_sh}}}}.
This has been fixed.

\item {} 
Change \sphinxcode{\sphinxupquote{size\_ind}} to {\hyperref[\detokenize{xsrst/py_fun_property:py-fun-property-size-domain}]{\sphinxcrossref{\DUrole{std,std-ref}{size\_domain}}}}
and \sphinxcode{\sphinxupquote{size\_dep}} to {\hyperref[\detokenize{xsrst/py_fun_property:py-fun-property-size-range}]{\sphinxcrossref{\DUrole{std,std-ref}{size\_range}}}}.

\end{enumerate}


\subparagraph{07\sphinxhyphen{}10}
\label{\detokenize{xsrst/whats_new_2018:whats-new-2018-07-10}}\label{\detokenize{xsrst/whats_new_2018:id14}}
Start conversion of python library from using
\sphinxcode{\sphinxupquote{vector}} to using \sphinxcode{\sphinxupquote{numpy.array}} ; see
{\hyperref[\detokenize{xsrst/py_independent:id1}]{\sphinxcrossref{\DUrole{std,std-ref}{py\_independent}}}}.


\subparagraph{07\sphinxhyphen{}08}
\label{\detokenize{xsrst/whats_new_2018:whats-new-2018-07-08}}\label{\detokenize{xsrst/whats_new_2018:id15}}\begin{enumerate}
\sphinxsetlistlabels{\arabic}{enumi}{enumii}{}{.}%
\item {} 
Advanced to \sphinxcode{\sphinxupquote{cppad\sphinxhyphen{}20180703}} .

\item {} 
Fix the following warning:
\sphinxcode{\sphinxupquote{catching polymorphic type class \textquotesingle{}std::runtime\_error\textquotesingle{} by value}} .

\item {} 
Suppress the following warning when running {\hyperref[\detokenize{xsrst/setup_py:id1}]{\sphinxcrossref{\DUrole{std,std-ref}{setup\_py}}}}:
\sphinxcode{\sphinxupquote{clearing an object of non\sphinxhyphen{}trivial type class \textquotesingle{}cppad\_py::a\_double\textquotesingle{}}} ;
see \sphinxhref{https://github.com/swig/swig/issues/1259}{swig issue 1259}.

\item {} 
Move all settings to {\hyperref[\detokenize{xsrst/setup_py:id1}]{\sphinxcrossref{\DUrole{std,std-ref}{setup\_py}}}} and run \sphinxcode{\sphinxupquote{cmake}} from there.

\item {} 
The file \sphinxcode{\sphinxupquote{cppad\_py/\_\_init\_\_.py}} was not being created.
This has been fixed.

\end{enumerate}


\subparagraph{07\sphinxhyphen{}07}
\label{\detokenize{xsrst/whats_new_2018:whats-new-2018-07-07}}\label{\detokenize{xsrst/whats_new_2018:id16}}
Change build of Python module to use Python 3
(now it can use either 2 or 3).


\subparagraph{07\sphinxhyphen{}03}
\label{\detokenize{xsrst/whats_new_2018:whats-new-2018-07-03}}\label{\detokenize{xsrst/whats_new_2018:id17}}
Birthday when Cppad Py was first extracted from Cppad Swig.


\bigskip\hrule\bigskip


xsrst input file: \sphinxcode{\sphinxupquote{whats\_new/2018.xsrst}}

\index{previous@\spxentry{previous}}\index{years@\spxentry{years}}\ignorespaces 

\subparagraph{Previous Years}
\label{\detokenize{xsrst/whats_new_2020:previous-years}}\label{\detokenize{xsrst/whats_new_2020:whats-new-2020-previous-years}}\label{\detokenize{xsrst/whats_new_2020:index-1}}
{\hyperref[\detokenize{xsrst/whats_new_2018:id1}]{\sphinxcrossref{\DUrole{std,std-ref}{whats\_new\_2018}}}}


\subparagraph{09\sphinxhyphen{}13}
\label{\detokenize{xsrst/whats_new_2020:whats-new-2020-09-13}}\label{\detokenize{xsrst/whats_new_2020:id2}}\begin{enumerate}
\sphinxsetlistlabels{\arabic}{enumi}{enumii}{}{.}%
\item {} 
The text \sphinxcode{\sphinxupquote{xsrst:}} was added to the beginning of each whats new item
(see below) that only pertained to the {\hyperref[\detokenize{xsrst/xsrst_py:id1}]{\sphinxcrossref{\DUrole{std,std-ref}{xsrst.py}}}} program.

\item {} 
xsrst: The order of the sections in the pdf file and the
table of contents was corrected.

\item {} 
xsrst: Put spaces around table of contents levels that have more that one
entry. Also remove the level of indentation in table of contents by one
because the root section stands out has having not number in front of it.

\item {} 
The latex numbering of sections is incorrect, so suppress it in the pdf
version of the output.

\end{enumerate}


\subparagraph{09\sphinxhyphen{}12}
\label{\detokenize{xsrst/whats_new_2020:whats-new-2020-09-12}}\label{\detokenize{xsrst/whats_new_2020:id3}}\begin{enumerate}
\sphinxsetlistlabels{\arabic}{enumi}{enumii}{}{.}%
\item {} 
xsrst: A table of contents that only contains section titles,
not headings within a section, was added.

\item {} 
xsrst: The following restriction was added: a
{\hyperref[\detokenize{xsrst/begin_cmd:begin-cmd-section-name}]{\sphinxcrossref{\DUrole{std,std-ref}{section\_name}}}}
can not begin with \sphinxcode{\sphinxupquote{xsrst\_}}.

\end{enumerate}


\subparagraph{07\sphinxhyphen{}31}
\label{\detokenize{xsrst/whats_new_2020:whats-new-2020-07-31}}\label{\detokenize{xsrst/whats_new_2020:id4}}\begin{enumerate}
\sphinxsetlistlabels{\arabic}{enumi}{enumii}{}{.}%
\item {} 
xsrst: An empty line at beginning of xsrst input file was causing it to crash
(this has been fixed).

\item {} 
xsrst: Putting version number in the documentation navigation side bar
changes every gh\sphinxhyphen{}pages file when version changes. Move it to
the first heading under the title {\hyperref[\detokenize{xsrst/cppad_py:id1}]{\sphinxcrossref{\DUrole{std,std-ref}{cppad\_py}}}}

\item {} 
xsrst: spell checking is not done for a \sphinxstyleemphasis{url} of the form
\textasciigrave{} \textless{} \sphinxstyleemphasis{url} \textgreater{} \textasciigrave{}\_ or \textasciigrave{} \sphinxstyleemphasis{text} \textless{} \sphinxstyleemphasis{url} \textgreater{} \textasciigrave{}\_.

\end{enumerate}


\subparagraph{07\sphinxhyphen{}30}
\label{\detokenize{xsrst/whats_new_2020:whats-new-2020-07-30}}\label{\detokenize{xsrst/whats_new_2020:id5}}
xsrst: The words ‘anl’, ‘dir’, ‘mcs’ were removed from the xsrst dictionary.
The word ‘initialization’ was added to the xsrst dictionary.


\subparagraph{07\sphinxhyphen{}29}
\label{\detokenize{xsrst/whats_new_2020:whats-new-2020-07-29}}\label{\detokenize{xsrst/whats_new_2020:id6}}\begin{enumerate}
\sphinxsetlistlabels{\arabic}{enumi}{enumii}{}{.}%
\item {} 
xsrst: Add the optional command line argument
{\hyperref[\detokenize{xsrst/xsrst_py:xsrst-py-command-line-arguments-line-increment}]{\sphinxcrossref{\DUrole{std,std-ref}{line\_increment}}}}.
This is an aid for finding the source of errors and warnings
reported by sphinx.

\item {} 
xsrst: Sometimes xsrst would not recognize a command that came directly after a
Fix {\hyperref[\detokenize{xsrst/file_cmd:id1}]{\sphinxcrossref{\DUrole{std,std-ref}{4.6 File Command}}}}. This has been fixed.

\end{enumerate}


\subparagraph{07\sphinxhyphen{}28}
\label{\detokenize{xsrst/whats_new_2020:whats-new-2020-07-28}}\label{\detokenize{xsrst/whats_new_2020:id7}}\begin{enumerate}
\sphinxsetlistlabels{\arabic}{enumi}{enumii}{}{.}%
\item {} 
xsrst: Move version number from the title in {\hyperref[\detokenize{xsrst/cppad_py:id1}]{\sphinxcrossref{\DUrole{std,std-ref}{cppad\_py}}}} to
the documentation navigation side bar.

\item {} 
xsrst: Add the {\hyperref[\detokenize{xsrst/configure:id1}]{\sphinxcrossref{\DUrole{std,std-ref}{4.9 Example xsrst Configuration Files}}}} section.

\item {} 
xsrst: The {\hyperref[\detokenize{xsrst/begin_cmd:id1}]{\sphinxcrossref{\DUrole{std,std-ref}{xsrst begin command}}}} was not recognized
at the beginning of a file
(when there is no new line before it). This has been fixed.

\item {} 
xsrst: Change the \sphinxstyleemphasis{index\_file} to
{\hyperref[\detokenize{xsrst/xsrst_py:xsrst-py-command-line-arguments-keyword}]{\sphinxcrossref{\DUrole{std,std-ref}{keyword}}}}
and \sphinxstyleemphasis{spell\_file} to
{\hyperref[\detokenize{xsrst/xsrst_py:xsrst-py-command-line-arguments-spelling}]{\sphinxcrossref{\DUrole{std,std-ref}{spelling}}}}

\end{enumerate}


\subparagraph{07\sphinxhyphen{}25}
\label{\detokenize{xsrst/whats_new_2020:whats-new-2020-07-25}}\label{\detokenize{xsrst/whats_new_2020:id8}}\begin{enumerate}
\sphinxsetlistlabels{\arabic}{enumi}{enumii}{}{.}%
\item {} 
xsrst: Install {\hyperref[\detokenize{xsrst/xsrst_py:id1}]{\sphinxcrossref{\DUrole{std,std-ref}{xsrst.py}}}} during the
{\hyperref[\detokenize{xsrst/setup_py:setup-py-install}]{\sphinxcrossref{\DUrole{std,std-ref}{install}}}} procedure
(not yet available from \sphinxcode{\sphinxupquote{pip}}).

\item {} 
xsrst: Add a {\hyperref[\detokenize{xsrst/xsrst_py:xsrst-py-purpose}]{\sphinxcrossref{\DUrole{std,std-ref}{purpose}}}} paragraph to xsrst documentation.

\item {} 
Move the \sphinxcode{\sphinxupquote{*.omh}} files to \sphinxcode{\sphinxupquote{*.xsrst}} files (they have been converted
from omhelp input files to xsrst input files).

\item {} 
xsrst: Change the {\hyperref[\detokenize{xsrst/xsrst_py:xsrst-py-command-line-arguments-keyword}]{\sphinxcrossref{\DUrole{std,std-ref}{keyword}}}}
to use python regular expressions and remove the whats new month\sphinxhyphen{}day
headings from the \DUrole{xref,std,std-ref}{index}.

\end{enumerate}


\subparagraph{07\sphinxhyphen{}24}
\label{\detokenize{xsrst/whats_new_2020:whats-new-2020-07-24}}\label{\detokenize{xsrst/whats_new_2020:id9}}\begin{enumerate}
\sphinxsetlistlabels{\arabic}{enumi}{enumii}{}{.}%
\item {} 
xsrst: Add the {\hyperref[\detokenize{xsrst/xsrst_py:id1}]{\sphinxcrossref{\DUrole{std,std-ref}{xsrst.py}}}} program
which extracts sphinx \sphinxcode{\sphinxupquote{*.rst}} files from source code.
This program is not yet installed but you can use it by placing
the script \sphinxcode{\sphinxupquote{bin/xsrst.py}} in your execution path.

\item {} 
Convert the cppad\_py documentation to xsrst input files,
which are mostly sphinx rst with a few extra xsrst commands.

\item {} 
Use xsrst and sphinx to generate this documentation for cppad\_py.

\end{enumerate}


\subparagraph{07\sphinxhyphen{}18}
\label{\detokenize{xsrst/whats_new_2020:whats-new-2020-07-18}}\label{\detokenize{xsrst/whats_new_2020:id10}}\begin{enumerate}
\sphinxsetlistlabels{\arabic}{enumi}{enumii}{}{.}%
\item {} 
The experimental \sphinxcode{\sphinxupquote{bin/xsrst.py}} program was added.
This program runs its tests and builds its documentation
(in the \sphinxcode{\sphinxupquote{sphinx}} ) directory with the command \sphinxcode{\sphinxupquote{bin/check\_xsrst.sh}}.
The intention is to convert the cppad\_py documentation from omhelp to
sphinx.

\item {} 
Convert all the cppad\_py source code to use spaces instead of tabs
(with tab stops at multiples of 4 spaces).

\end{enumerate}


\subparagraph{07\sphinxhyphen{}05}
\label{\detokenize{xsrst/whats_new_2020:whats-new-2020-07-05}}\label{\detokenize{xsrst/whats_new_2020:id11}}\begin{enumerate}
\sphinxsetlistlabels{\arabic}{enumi}{enumii}{}{.}%
\item {} 
The {\hyperref[\detokenize{xsrst/vec2a_double:vec2a-double-prototype}]{\sphinxcrossref{\DUrole{std,std-ref}{prototype}}}} for \sphinxcode{\sphinxupquote{vec2a\_double}}
was the same as for {\hyperref[\detokenize{xsrst/vec2cppad_double:id1}]{\sphinxcrossref{\DUrole{std,std-ref}{vec2cppad\_double}}}}.
This has been fixed.

\item {} 
The limits for the {\hyperref[\detokenize{xsrst/fun_from_json_xam_py:id1}]{\sphinxcrossref{\DUrole{std,std-ref}{fun\_from\_json\_xam\_py}}}} and {\hyperref[\detokenize{xsrst/fun_to_json_xam_py:id1}]{\sphinxcrossref{\DUrole{std,std-ref}{fun\_to\_json\_xam\_py}}}}
example code were incorrect (this has been fixed).

\item {} 
The in {\hyperref[\detokenize{xsrst/fun_from_json_xam_cpp:id1}]{\sphinxcrossref{\DUrole{std,std-ref}{fun\_from\_json\_xam\_cpp}}}} the title was corrected; to be specific,
to\_json was changed to from\_json

\end{enumerate}


\subparagraph{05\sphinxhyphen{}17}
\label{\detokenize{xsrst/whats_new_2020:whats-new-2020-05-17}}\label{\detokenize{xsrst/whats_new_2020:id12}}
The \sphinxcode{\sphinxupquote{abs}} {\hyperref[\detokenize{xsrst/a_double_unary_fun:id1}]{\sphinxcrossref{\DUrole{std,std-ref}{a\_double\_unary\_fun}}}} was added
(works the same as the \sphinxcode{\sphinxupquote{fabs}} function).


\subparagraph{05\sphinxhyphen{}16}
\label{\detokenize{xsrst/whats_new_2020:whats-new-2020-05-16}}\label{\detokenize{xsrst/whats_new_2020:id13}}\begin{enumerate}
\sphinxsetlistlabels{\arabic}{enumi}{enumii}{}{.}%
\item {} 
A discussion of some subtle issues,
when interpolation is used to define an ODE,
was added to the \sphinxcode{\sphinxupquote{ode\_multi\_step}} ; see
{\hyperref[\detokenize{xsrst/numeric_ode_multi_step:id1}]{\sphinxcrossref{\DUrole{std,std-ref}{set\_t\_all\_index}}}}.

\item {} 
Add the {\hyperref[\detokenize{xsrst/numeric_seirwd_model:numeric-seirwd-model-p-all-alpha}]{\sphinxcrossref{\DUrole{std,std-ref}{alpha}}}} parameter to the
SEIRWD model.

\end{enumerate}


\subparagraph{05\sphinxhyphen{}15}
\label{\detokenize{xsrst/whats_new_2020:whats-new-2020-05-15}}\label{\detokenize{xsrst/whats_new_2020:id14}}
The file \sphinxcode{\sphinxupquote{ode\_solve.py}} was moved to \sphinxcode{\sphinxupquote{ode\_multi\_step.py}}
and {\hyperref[\detokenize{xsrst/numeric_ode_multi_step:id1}]{\sphinxcrossref{\DUrole{std,std-ref}{numeric\_ode\_multi\_step}}}} was extended to allow for different ODE solvers
(steppers).


\subparagraph{05\sphinxhyphen{}14}
\label{\detokenize{xsrst/whats_new_2020:whats-new-2020-05-14}}\label{\detokenize{xsrst/whats_new_2020:id15}}\begin{enumerate}
\sphinxsetlistlabels{\arabic}{enumi}{enumii}{}{.}%
\item {} 
The example semi\sphinxhyphen{}stiff integrator {\hyperref[\detokenize{xsrst/numeric_rosen3_step:id1}]{\sphinxcrossref{\DUrole{std,std-ref}{numeric\_rosen3\_step}}}} was added.

\item {} 
The example file \sphinxcode{\sphinxupquote{runge4.py}} was moved to \sphinxcode{\sphinxupquote{ode\_solve.py}}
in preparation for other solvers.

\item {} 
Add the {\hyperref[\detokenize{xsrst/numeric_simple_inv:id1}]{\sphinxcrossref{\DUrole{std,std-ref}{numeric\_simple\_inv}}}} example routine for
AD inversion of matrices.

\item {} 
Add mention of which {\hyperref[\detokenize{xsrst/numeric_xam:id1}]{\sphinxcrossref{\DUrole{std,std-ref}{numeric\_xam}}}} routines can take \sphinxcode{\sphinxupquote{a\_double}} values.

\item {} 
Simplify the {\hyperref[\detokenize{xsrst/numeric_seirwd_model_xam_py:id1}]{\sphinxcrossref{\DUrole{std,std-ref}{numeric\_seirwd\_model\_xam\_py}}}} example using numpy vector operations.

\item {} 
More improvements to were made to
{\hyperref[\detokenize{xsrst/numeric_covid_19_xam_py:id1}]{\sphinxcrossref{\DUrole{std,std-ref}{covid\sphinxhyphen{}19}}}} example.

\end{enumerate}


\subparagraph{05\sphinxhyphen{}12}
\label{\detokenize{xsrst/whats_new_2020:whats-new-2020-05-12}}\label{\detokenize{xsrst/whats_new_2020:id16}}\begin{enumerate}
\sphinxsetlistlabels{\arabic}{enumi}{enumii}{}{.}%
\item {} 
More improvements to were made to
{\hyperref[\detokenize{xsrst/numeric_covid_19_xam_py:id1}]{\sphinxcrossref{\DUrole{std,std-ref}{covid\sphinxhyphen{}19}}}} example.

\item {} 
The {\hyperref[\detokenize{xsrst/numeric_seirwd_model:numeric-seirwd-model-n-step}]{\sphinxcrossref{\DUrole{std,std-ref}{n\_step}}}} option was added to the
SEIRWD model example.

\item {} 
A delay between when is no longer infectious and when one dies was
added to the SEIRWD model using a W compartment (Will die compartment).

\end{enumerate}


\subparagraph{05\sphinxhyphen{}09}
\label{\detokenize{xsrst/whats_new_2020:whats-new-2020-05-09}}\label{\detokenize{xsrst/whats_new_2020:id17}}
The {\hyperref[\detokenize{xsrst/numeric_covid_19_xam_py:id1}]{\sphinxcrossref{\DUrole{std,std-ref}{covid\sphinxhyphen{}19\_modeling\_example}}}} application
of AD is changing on a regular basis, so the details of the
changes will no longer be tracked in this file, just see current
example if you are interested.


\subparagraph{05\sphinxhyphen{}08}
\label{\detokenize{xsrst/whats_new_2020:whats-new-2020-05-08}}\label{\detokenize{xsrst/whats_new_2020:id18}}\begin{enumerate}
\sphinxsetlistlabels{\arabic}{enumi}{enumii}{}{.}%
\item {} 
The \sphinxcode{\sphinxupquote{numeric\_seirs\_fit\_xam.py}} example was renamed
{\hyperref[\detokenize{xsrst/numeric_covid_19_xam_py:id1}]{\sphinxcrossref{\DUrole{std,std-ref}{numeric\_covid\_19\_xam\_py}}}}.
And it was change to fit covariate multipliers that model
the infectious rate \(\beta\).

\item {} 
The \sphinxcode{\sphinxupquote{numeric\_seirs\_model}} was renamed {\hyperref[\detokenize{xsrst/numeric_seirwd_model:id1}]{\sphinxcrossref{\DUrole{std,std-ref}{numeric\_seirwd\_model}}}},
the D compartment was added to track total deaths,
the data was changed from Infectious to cumulative death, and
noise was added to the data.

\item {} 
In \sphinxcode{\sphinxupquote{numeric\_covid\_a9\_xam}} ,
reduce to just E and I as unknowns in initial conditions.
Initial R = D = 0, S = 1 \sphinxhyphen{} E \sphinxhyphen{} R.

\item {} 
Use the observed information matrix to estimate the covariance
of the optimal parameters.

\end{enumerate}


\subparagraph{05\sphinxhyphen{}07}
\label{\detokenize{xsrst/whats_new_2020:whats-new-2020-05-07}}\label{\detokenize{xsrst/whats_new_2020:id19}}\begin{enumerate}
\sphinxsetlistlabels{\arabic}{enumi}{enumii}{}{.}%
\item {} 
The function \sphinxcode{\sphinxupquote{scipy.misc.factorial}}
has been deprecated so change it to
\sphinxcode{\sphinxupquote{scipy.special.factorial}} .

\item {} 
Add the \sphinxcode{\sphinxupquote{numeric\_seirs\_model}} utility example
which allows rates in the ODE to change with time.
This enabled changing the {\hyperref[\detokenize{xsrst/numeric_covid_19_xam_py:id1}]{\sphinxcrossref{\DUrole{std,std-ref}{numeric\_covid\_19\_xam\_py}}}}
example to estimate a \(\beta\) coefficient that
changes with time.

\item {} 
The \sphinxcode{\sphinxupquote{numeric\_seirs\_class}} example was removed because
\sphinxcode{\sphinxupquote{numeric\_seirwd\_model}} replaces it with more general
capabilities.

\end{enumerate}


\subparagraph{05\sphinxhyphen{}06}
\label{\detokenize{xsrst/whats_new_2020:whats-new-2020-05-06}}\label{\detokenize{xsrst/whats_new_2020:id20}}\begin{enumerate}
\sphinxsetlistlabels{\arabic}{enumi}{enumii}{}{.}%
\item {} 
Add the following utility examples:
{\hyperref[\detokenize{xsrst/numeric_optimize_fun_class:id1}]{\sphinxcrossref{\DUrole{std,std-ref}{numeric\_optimize\_fun\_class}}}},
{\hyperref[\detokenize{xsrst/numeric_runge4_step:id1}]{\sphinxcrossref{\DUrole{std,std-ref}{numeric\_runge4\_step}}}},
{\hyperref[\detokenize{xsrst/numeric_ode_multi_step:id1}]{\sphinxcrossref{\DUrole{std,std-ref}{numeric\_ode\_multi\_step}}}},
\sphinxcode{\sphinxupquote{numeric\_seirs\_class}} .
(The \sphinxcode{\sphinxupquote{numeric\_seirs\_class}} was later replaced by
{\hyperref[\detokenize{xsrst/numeric_seirwd_model:id1}]{\sphinxcrossref{\DUrole{std,std-ref}{numeric\_seirwd\_model}}}}.)

\item {} 
Add mention of using \sphinxcode{\sphinxupquote{DYLD\_LIBRARY\_PATH}} in the
{\hyperref[\detokenize{xsrst/install_error:install-error-cppad-library-missing}]{\sphinxcrossref{\DUrole{std,std-ref}{cppad\_library\_missing}}}}
install instructions for the Mac.

\end{enumerate}


\subparagraph{05\sphinxhyphen{}05}
\label{\detokenize{xsrst/whats_new_2020:whats-new-2020-05-05}}\label{\detokenize{xsrst/whats_new_2020:id21}}
Add the {\hyperref[\detokenize{xsrst/numeric_runge4_step_xam_py:id1}]{\sphinxcrossref{\DUrole{std,std-ref}{numeric\_runge4\_step\_xam\_py}}}} example.


\subparagraph{05\sphinxhyphen{}04}
\label{\detokenize{xsrst/whats_new_2020:whats-new-2020-05-04}}\label{\detokenize{xsrst/whats_new_2020:id22}}
Add the {\hyperref[\detokenize{xsrst/numeric_optimize_fun_xam_py:id1}]{\sphinxcrossref{\DUrole{std,std-ref}{numeric\_optimize\_fun\_xam\_py}}}} example.


\subparagraph{04\sphinxhyphen{}28}
\label{\detokenize{xsrst/whats_new_2020:whats-new-2020-04-28}}\label{\detokenize{xsrst/whats_new_2020:id23}}
The {\hyperref[\detokenize{xsrst/py_fun_json:py-fun-json-from-json}]{\sphinxcrossref{\DUrole{std,std-ref}{f\_from\_json()}}}} function was added.
In addition, the
{\hyperref[\detokenize{xsrst/py_fun_ctor:py-fun-ctor-f-empty-function}]{\sphinxcrossref{\DUrole{std,std-ref}{empty\_function}}}} constructor was added.


\subparagraph{04\sphinxhyphen{}27}
\label{\detokenize{xsrst/whats_new_2020:whats-new-2020-04-27}}\label{\detokenize{xsrst/whats_new_2020:id24}}
The \sphinxcode{\sphinxupquote{setup.py}} program now installs a separate copy of CppAD
and provides instructions at the end for modifying your
\sphinxcode{\sphinxupquote{LD\_LIBRARY\_PATH}} ; see
{\hyperref[\detokenize{xsrst/install_error:install-error-cppad-library-missing}]{\sphinxcrossref{\DUrole{std,std-ref}{cppad\_library\_missing}}}}.


\subparagraph{04\sphinxhyphen{}26}
\label{\detokenize{xsrst/whats_new_2020:whats-new-2020-04-26}}\label{\detokenize{xsrst/whats_new_2020:id25}}
The \sphinxcode{\sphinxupquote{setup.py}} program was modified to try to automatically solve the
{\hyperref[\detokenize{xsrst/install_error:install-error-cppad-library-missing}]{\sphinxcrossref{\DUrole{std,std-ref}{cppad\_library\_missing}}}} problem.


\subparagraph{04\sphinxhyphen{}25}
\label{\detokenize{xsrst/whats_new_2020:whats-new-2020-04-25}}\label{\detokenize{xsrst/whats_new_2020:id26}}
The {\hyperref[\detokenize{xsrst/py_fun_json:py-fun-json-to-json}]{\sphinxcrossref{\DUrole{std,std-ref}{f\_to\_json()}}}} function was added.


\subparagraph{04\sphinxhyphen{}24}
\label{\detokenize{xsrst/whats_new_2020:whats-new-2020-04-24}}\label{\detokenize{xsrst/whats_new_2020:id27}}
The newer Mac systems seems to require that one use
\sphinxcode{\sphinxupquote{\sphinxhyphen{}stdlib=libc++}} compile and link flag.
The install has been changed to check for and adapt to this condition.
In addition, {\hyperref[\detokenize{xsrst/setup_py:id1}]{\sphinxcrossref{\DUrole{std,std-ref}{setup\_py}}}} now runs the \sphinxcode{\sphinxupquote{cmake}} command; i.e.,
the user no longer needs to run \sphinxcode{\sphinxupquote{cmake}} to test the
{\hyperref[\detokenize{xsrst/setup_py:setup-py-test-c}]{\sphinxcrossref{\DUrole{std,std-ref}{c++}}}} library.


\subparagraph{04\sphinxhyphen{}23}
\label{\detokenize{xsrst/whats_new_2020:whats-new-2020-04-23}}\label{\detokenize{xsrst/whats_new_2020:id28}}
Add an {\hyperref[\detokenize{xsrst/install_error:id1}]{\sphinxcrossref{\DUrole{std,std-ref}{install\_error}}}} section to the documentation.


\subparagraph{04\sphinxhyphen{}22}
\label{\detokenize{xsrst/whats_new_2020:whats-new-2020-04-22}}\label{\detokenize{xsrst/whats_new_2020:id29}}\begin{enumerate}
\sphinxsetlistlabels{\arabic}{enumi}{enumii}{}{.}%
\item {} 
Instructions were added for
{\hyperref[\detokenize{xsrst/setup_py:setup-py-install-using-pip}]{\sphinxcrossref{\DUrole{std,std-ref}{installing\_using\_pip}}}}.

\item {} 
The binary operators were extended to include \sphinxstyleemphasis{x} \sphinxcode{\sphinxupquote{op}} \sphinxstyleemphasis{ay} where:
\sphinxstyleemphasis{x} is a double (python \sphinxcode{\sphinxupquote{float}} ),
\sphinxstyleemphasis{ay} is an \sphinxcode{\sphinxupquote{a\_double}} ,
and \sphinxstyleemphasis{op} is
\sphinxcode{\sphinxupquote{+}} ,
\sphinxcode{\sphinxupquote{\sphinxhyphen{}}} ,
\sphinxcode{\sphinxupquote{*}} ,
\sphinxcode{\sphinxupquote{/}} , or
\sphinxcode{\sphinxupquote{**}} .
Note that this automatically transfers to numpy arrays; e.g
\sphinxstyleemphasis{x} \sphinxcode{\sphinxupquote{*}} \sphinxstyleemphasis{ay} where \sphinxstyleemphasis{x} is a double and \sphinxstyleemphasis{ay}
is a numpy array of \sphinxcode{\sphinxupquote{a\_double}} .

\end{enumerate}


\subparagraph{04\sphinxhyphen{}20}
\label{\detokenize{xsrst/whats_new_2020:whats-new-2020-04-20}}\label{\detokenize{xsrst/whats_new_2020:id30}}\begin{enumerate}
\sphinxsetlistlabels{\arabic}{enumi}{enumii}{}{.}%
\item {} 
Move configuration setting from {\hyperref[\detokenize{xsrst/setup_py:id1}]{\sphinxcrossref{\DUrole{std,std-ref}{setup\_py}}}} to
{\hyperref[\detokenize{xsrst/get_cppad_sh:get-cppad-sh-settings}]{\sphinxcrossref{\DUrole{std,std-ref}{bin/get\_cppad\_sh}}}}.

\item {} 
First version that installs using \sphinxcode{\sphinxupquote{pip}} .
Install instructions for pip will be added soon.

\end{enumerate}


\subparagraph{04\sphinxhyphen{}19}
\label{\detokenize{xsrst/whats_new_2020:whats-new-2020-04-19}}\label{\detokenize{xsrst/whats_new_2020:id31}}\begin{enumerate}
\sphinxsetlistlabels{\arabic}{enumi}{enumii}{}{.}%
\item {} 
Move the python source that gets distributed from \sphinxcode{\sphinxupquote{lib/python}}
to \sphinxcode{\sphinxupquote{lib/python/cppad\_py}} so that more like a standard python package.

\item {} 
Drop support for python2. It is not consistent with python3 in
some of the \sphinxcode{\sphinxupquote{setup.py}} actions.

\end{enumerate}


\subparagraph{04\sphinxhyphen{}18}
\label{\detokenize{xsrst/whats_new_2020:whats-new-2020-04-18}}\label{\detokenize{xsrst/whats_new_2020:id32}}\begin{enumerate}
\sphinxsetlistlabels{\arabic}{enumi}{enumii}{}{.}%
\item {} 
Change \sphinxstyleemphasis{yq} to \sphinxstyleemphasis{xq} , correct documentation,
for \sphinxstyleemphasis{xq} in the
{\hyperref[\detokenize{xsrst/cpp_fun_reverse:cpp-fun-reverse-xq}]{\sphinxcrossref{\DUrole{std,std-ref}{c++}}}} and {\hyperref[\detokenize{xsrst/py_fun_reverse:py-fun-reverse-xq}]{\sphinxcrossref{\DUrole{std,std-ref}{python}}}}
reverse mode documentation.

\item {} 
Remove the \sphinxcode{\sphinxupquote{\sphinxhyphen{}\sphinxhyphen{}inplace}} option from the
{\hyperref[\detokenize{xsrst/setup_py:setup-py-syntax}]{\sphinxcrossref{\DUrole{std,std-ref}{syntax}}}} for building the cppad\_py python module.

\end{enumerate}


\subparagraph{04\sphinxhyphen{}13}
\label{\detokenize{xsrst/whats_new_2020:whats-new-2020-04-13}}\label{\detokenize{xsrst/whats_new_2020:id33}}\begin{enumerate}
\sphinxsetlistlabels{\arabic}{enumi}{enumii}{}{.}%
\item {} 
The \sphinxstyleemphasis{z} =  \sphinxcode{\sphinxupquote{pow}} ( \sphinxstyleemphasis{x} , \sphinxstyleemphasis{y} ) functions was added; see
{\hyperref[\detokenize{xsrst/a_double_binary:id1}]{\sphinxcrossref{\DUrole{std,std-ref}{a\_double\_binary}}}}.

\item {} 
Add the {\hyperref[\detokenize{xsrst/a_double_property:a-double-property-var2par}]{\sphinxcrossref{\DUrole{std,std-ref}{var2par}}}} function
and improve the notation in the
{\hyperref[\detokenize{xsrst/a_double_property:a-double-property-near-equal}]{\sphinxcrossref{\DUrole{std,std-ref}{near\_equal}}}} notation.

\end{enumerate}


\subparagraph{04\sphinxhyphen{}12}
\label{\detokenize{xsrst/whats_new_2020:whats-new-2020-04-12}}\label{\detokenize{xsrst/whats_new_2020:id34}}\begin{enumerate}
\sphinxsetlistlabels{\arabic}{enumi}{enumii}{}{.}%
\item {} 
Add the \sphinxcode{\sphinxupquote{erf}} function was added to the
list of \sphinxcode{\sphinxupquote{a\_double}} unary {\hyperref[\detokenize{xsrst/a_double_unary_fun:a-double-unary-fun-fun}]{\sphinxcrossref{\DUrole{std,std-ref}{fun}}}}
that have been implemented.

\item {} 
The dynamic parameter argument was missing from the
{\hyperref[\detokenize{xsrst/py_independent:py-independent-syntax}]{\sphinxcrossref{\DUrole{std,std-ref}{syntax}}}} for the python version
of the \sphinxcode{\sphinxupquote{independent}} function.  This has been fixed.

\item {} 
Improve the {\hyperref[\detokenize{xsrst/setup_py:setup-py-test}]{\sphinxcrossref{\DUrole{std,std-ref}{test}}}} and {\hyperref[\detokenize{xsrst/setup_py:setup-py-install}]{\sphinxcrossref{\DUrole{std,std-ref}{install}}}}
discussion in \sphinxcode{\sphinxupquote{setup.py}} .

\end{enumerate}


\subparagraph{04\sphinxhyphen{}10}
\label{\detokenize{xsrst/whats_new_2020:whats-new-2020-04-10}}\label{\detokenize{xsrst/whats_new_2020:id35}}\begin{enumerate}
\sphinxsetlistlabels{\arabic}{enumi}{enumii}{}{.}%
\item {} 
Change the documentation display on the web using a more recent version of
the documentation program \sphinxcode{\sphinxupquote{omhelp\sphinxhyphen{}20200130}} .

\item {} 
Add {\hyperref[\detokenize{xsrst/get_cppad_sh:get-cppad-sh-caching}]{\sphinxcrossref{\DUrole{std,std-ref}{caching}}}} to the Cppad install.

\end{enumerate}


\bigskip\hrule\bigskip


xsrst input file: \sphinxcode{\sphinxupquote{whats\_new/2020.xsrst}}
\begin{enumerate}
\sphinxsetlistlabels{\arabic}{enumi}{enumii}{}{.}%
\item {} 
setup\_py: {\hyperref[\detokenize{xsrst/setup_py:id1}]{\sphinxcrossref{\DUrole{std,std-ref}{1 Configure and Build the cppad\_py Python Module}}}}

\item {} 
numeric\_xam: {\hyperref[\detokenize{xsrst/numeric_xam:id1}]{\sphinxcrossref{\DUrole{std,std-ref}{2 Numerical Examples}}}}

\item {} 
library: {\hyperref[\detokenize{xsrst/library:id1}]{\sphinxcrossref{\DUrole{std,std-ref}{3 The Cppad Py Libraries}}}}

\item {} 
xsrst\_py: {\hyperref[\detokenize{xsrst/xsrst_py:id1}]{\sphinxcrossref{\DUrole{std,std-ref}{4 Extract Sphinx RST}}}}

\item {} 
whats\_new\_2020: {\hyperref[\detokenize{xsrst/whats_new_2020:id1}]{\sphinxcrossref{\DUrole{std,std-ref}{5 CppAD Py Changes During 2020}}}}

\end{enumerate}


\bigskip\hrule\bigskip


xsrst input file: \sphinxcode{\sphinxupquote{doc.xsrst}}



\renewcommand{\indexname}{Index}
\printindex
\end{document}